\chapter{Galois representations and local Langlands}\label{chapter-galois-representations}

This chapter states the results on Galois representations that the proof uses: a local Langlands
correspondence for $\mathrm{GL}_n$ over $p$-adic fields, and $\ell$-adic Galois representations
attached to regular algebraic polarized cuspidal representations of $\mathrm{GL}_a$ over CM fields.
Both hold in weaker forms than the published theorems. The correspondence is pinned down by a list of
laws that involve no $L$- or $\varepsilon$-factors and determine it only on good representations, and
the Galois representations carry local-global compatibility only at the places where the laws
determine the correspondence, with no polarization and no purity. The last section states the
ingredients from which the Galois representations are built: cyclic base change, base change from
definite unitary groups, theta lifts to unitary groups of large rank, and seed representations from
the cohomology of PEL unitary Shimura varieties.

\section{Weil--Deligne representations}\label{section-gal-weil-deligne}

\begin{situation}\label{gal-local-situation}
Let $E$ be a number field, let $\Omega$ be an algebraic closure of $E$, and let
$G_E=\operatorname{Gal}(\Omega/E)$. Let $w$ be a finite place of $E$, with completion $E_w$,
valuation ring $\mathcal O_w$, residue field $k_w$ and $q_w=|k_w|$. Let $A\subset\Omega$ be a
valuation ring whose maximal ideal $\mathfrak m_A$ meets $\mathcal O_E$ in the prime ideal of $w$,
and let $v_A$ be its valuation. We write:
\begin{enumerate}
\item $D_A\subset G_E$ for the stabilizer of $A$;
\item $I_A\subset D_A$ for the inertia group, the $\sigma\in D_A$ with $v_A(\sigma(z)-z)>0$ for all
  $z\in A$;
\item $P_A$ for the wild inertia group, the $\sigma\in G_E$ with $v_A(\sigma(z)/z-1)>0$ for all
  $z\in\Omega^\times$;
\item $W_A\subset D_A$ for the Weil group, the $\sigma\in D_A$ for which there are integers
  $\alpha,\beta\geq0$ with $\sigma(z)^{q_w^{\beta}}\equiv z^{q_w^{\alpha}}\pmod{\mathfrak m_A}$ for
  all $z\in A$.
\end{enumerate}
For $\sigma\in W_A$ the integer $d(\sigma)=\alpha-\beta$ depends only on $\sigma$, and
$d:W_A\to\mathbf Z$ is a homomorphism with kernel $I_A$. An element of degree $1$ is an arithmetic
Frobenius element and an element of degree $-1$ a geometric Frobenius element.
\end{situation}
\leanref{Development/AbelianVarietyTheory/Definitions/AutomorphicMembership/ExistsNotCongPow.lean}{GlobalTheory.WeilDeligne.Weil}
\leanref{Development/AbelianVarietyTheory/Definitions/AutomorphicMembership/ExistsNotCongPow.lean}{GlobalTheory.WeilDeligne.deg}
\leanref{ReciprocityTheory/GaloisValuations/Global.lean}{ReciprocityTheory.Gal.Ine}

\begin{definition}\label{gal-local-artin}
In Situation~\ref{gal-local-situation}, a \emph{local Artin map} is a surjective homomorphism
$a:W_A\to E_w^\times$ with the following properties.
\begin{enumerate}
\item The restriction of $a$ to $I_A$ is continuous.
\item $v_w(a(\sigma))=-d(\sigma)$ for every $\sigma\in W_A$.
\item Every homomorphism $\xi:W_A\to\mathbf C^\times$ that is trivial on an open subgroup of $I_A$
  is $\chi\circ a$ for a smooth character $\chi$ of $E_w^\times$.
\item For every reciprocity map $\operatorname{art}:\mathbf A_E^\times\to G_E^{\mathrm{ab}}$ of
  global class field theory and every $\sigma\in W_A$, $\operatorname{art}$ sends the idele equal to
  $a(\sigma)$ at $w$ and to $1$ elsewhere to the image of $\sigma$ in $G_E^{\mathrm{ab}}$.
\end{enumerate}
\end{definition}
\leanref{Development/AbelianVarietyTheory/PublishedTheorems/AutomorphicForms₂/HasExponentZeroIff.lean}{GlobalTheory.LocLanglands.IsLocalArtin}

\begin{proposition}\label{gal-local-artin-unique}
In Situation~\ref{gal-local-situation}, there is exactly one local Artin map
(Definition~\ref{gal-local-artin}) \cite{Ser67}.
\end{proposition}
\leanref{Development/AutomorphicArithmetic/Foundations/LocalLanglands/LocalCentralUnits/ArtEx.lean}{UpperTheory.Zero.LLC.artEx}

\begin{definition}\label{gal-weil-deligne}
In Situation~\ref{gal-local-situation}, a \emph{Weil--Deligne representation} of $W_A$ is a triple
$(V,\rho,N)$ in which $V$ is a finite-dimensional complex vector space,
$\rho:W_A\to\operatorname{GL}(V)$ is a homomorphism that is trivial on an open subgroup of $I_A$ and
such that every $\rho(\sigma)$ is semisimple, and $N\in\operatorname{End}(V)$ is nilpotent with
$\rho(\sigma)N=q_w^{d(\sigma)}N\rho(\sigma)$ for all $\sigma\in W_A$. Write $\mathrm{WD}(A)$ for the
set of isomorphism classes of Weil--Deligne representations of $W_A$. Direct sums, tensor products
(with $N\otimes1+1\otimes N$) and duals (with $-N^{\vee}$) are taken as usual. For a homomorphism
$\xi:W_A\to\mathbf C^\times$ trivial on an open subgroup of $I_A$, $x\otimes\xi$ is
$(V,\xi\rho,N)$. For $x=(V,\rho,N)$ and $\sigma\in W_A$ write $\det x(\sigma)$ for the product of
the eigenvalues of $\rho(\sigma)$. For an integer $l\geq1$, $\operatorname{Sp}(l)$ is $\mathbf C^l$
with basis $e_0,\ldots,e_{l-1}$, with $\rho(\sigma)e_j=q_w^{(j-(l-1)/2)d(\sigma)}e_j$, and with
$Ne_j=e_{j+1}$ for $j<l-1$ and $Ne_{l-1}=0$.
\end{definition}
\leanref{Development/AbelianVarietyTheory/Definitions/AutomorphicMembership/MonodromyZeroOfWeilDeligneRepClassIff.lean}{GlobalTheory.WeilDeligne.WDRep}
\leanref{Development/AbelianVarietyTheory/Definitions/AutomorphicMembership/MonodromyZeroOfWeilDeligneRepClassIff.lean}{GlobalTheory.WeilDeligne.WD}
\leanref{Development/AbelianVarietyTheory/Definitions/AutomorphicMembership/MonodromyZeroOfWeilDeligneRepClassIff.lean}{GlobalTheory.WeilDeligne.WDRep.sp}

\begin{definition}\label{gal-wd-classes}
In Situation~\ref{gal-local-situation}, let $x=(V,\rho,N)$ be a Weil--Deligne representation of
$W_A$ (Definition~\ref{gal-weil-deligne}) and let $n=\dim V$.
\begin{enumerate}
\item $x$ is \emph{monodromy-free} if $N=0$.
\item $x$ is \emph{irreducible} if $\rho$ is irreducible.
\item $x$ is \emph{tame} if $\rho(\sigma)$ has characteristic polynomial $(X-1)^n$ for every
  $\sigma\in P_A$, that is, $\rho$ is trivial on $P_A$.
\item $x$ is \emph{abelian} if $\rho(\sigma\tau\sigma^{-1}\tau^{-1})$ has characteristic polynomial
  $(X-1)^n$ for all $\sigma,\tau\in W_A$. Since every $\rho(\sigma)$ is semisimple, this holds if and
  only if $\rho$ is a direct sum of characters.
\item $x$ has \emph{multiplicity-free tame part} if for every $\gamma\in I_A$ such that the closed
  subgroup generated by $\gamma$ and $P_A$ contains $I_A$, every eigenvalue of $\rho(\gamma)$ of
  multiplicity at least $2$ is a $(q_w-1)$-th root of unity.
\end{enumerate}
\end{definition}
\leanref{Development/AbelianVarietyTheory/Definitions/AutomorphicMembership/MonodromyZeroOfWeilDeligneRepClassIff.lean}{GlobalTheory.WeilDeligne.WD.mono0}
\leanref{Development/AbelianVarietyTheory/PublishedTheorems/AutomorphicForms₂/LocalLanglandsGL.lean}{GlobalTheory.LocLanglands.IsIrrWD}
\leanref{Development/AutomorphicArithmetic/Foundations/GaloisRepresentations/TypeConditions/WeightsConj.lean}{UpperTheory.GalS.TmW}
\leanref{Development/AbelianVarietyTheory/PublishedTheorems/AutomorphicForms₂/LocalLanglandsGL.lean}{GlobalTheory.LocLanglands.IsTameWD}
\leanref{Development/AutomorphicArithmetic/Foundations/GaloisRepresentations/TypeConditions/WeightsConjAlt.lean}{UpperTheory.GalStar.Tor}
\leanref{Development/AutomorphicArithmetic/Foundations/GaloisRepresentations/TypeConditions/WeightsConjAlt.lean}{UpperTheory.GalStar.Sqf}

\begin{definition}\label{gal-wd-galois}
In Situation~\ref{gal-local-situation}, let $\ell$ be a prime that does not lie under $w$, let
$\iota:\bar{\mathbf Q}_\ell\to\mathbf C$ be a field isomorphism, and let
$r:G_E\to\operatorname{GL}_n(\bar{\mathbf Q}_\ell)$ be a homomorphism. Suppose there are a
continuous surjective homomorphism $t_\ell:I_A\to\mathbf Z_\ell$, a nilpotent
$N\in M_n(\bar{\mathbf Q}_\ell)$ with $r(\sigma)=\exp(t_\ell(\sigma)N)$ for all $\sigma$ in an open
subgroup of $I_A$, and an element $\varphi\in W_A$ of degree $1$. Put
$r'(\sigma)=r(\sigma)\exp(-t_\ell(\varphi^{-d(\sigma)}\sigma)N)$ for $\sigma\in W_A$. Then
$\mathrm{WD}_\iota(r|_{W_A})\in\mathrm{WD}(A)$ is the class of the Frobenius semisimplification of
$(\mathbf C^n,\iota\circ r',\iota(N))$; it does not depend on the choices. If no such data exist,
$\mathrm{WD}_\iota(r|_{W_A})$ is the trivial representation of dimension $n$. By Grothendieck's
$\ell$-adic monodromy theorem \cite[Appendix]{ST68}, such data exist whenever $r$ is continuous.
\end{definition}
\leanref{Development/AbelianVarietyTheory/Definitions/AutomorphicMembership/WeilDeligneParamMonodromyEqZeroIff.lean}{GlobalTheory.WeilDeligne.IsWDOf}
\leanref{Development/AbelianVarietyTheory/Definitions/AutomorphicMembership/WeilDeligneParamMonodromyEqZeroIff.lean}{GlobalTheory.WeilDeligne.wdOf}

\begin{definition}\label{gal-p-adic-hodge}
In Situation~\ref{gal-local-situation}, let $\ell$ be the prime under $w$, let
$\iota:\bar{\mathbf Q}_\ell\to\mathbf C$ be a field isomorphism, and let
$r:G_E\to\operatorname{GL}_n(\bar{\mathbf Q}_\ell)$ be a continuous homomorphism. The restriction
$r|_{D_A}$ is \emph{crystalline}, \emph{de Rham} or \emph{potentially semistable} in Fontaine's
sense \cite{Fon94}. For an embedding $\tau:E\to\bar{\mathbf Q}_\ell$ that extends continuously to
$E_w$, the \emph{Hodge--Tate multiset} $\mathrm{HT}_\tau(r|_{D_A})$ contains each integer $j$ with
multiplicity the dimension of the $j$-th graded piece of the Hodge filtration on the $\tau$-part of
$D_{\mathrm{dR}}(r|_{D_A})$. If $r|_{D_A}$ is potentially semistable,
$\mathrm{WD}_\iota(r|_{D_A})\in\mathrm{WD}(A)$ is the Frobenius semisimplification of the
Weil--Deligne representation attached to $D_{\mathrm{pst}}(r|_{D_A})$ \cite{Fon94}, transported to
$\mathbf C$ by $\iota$; otherwise it is the trivial representation of dimension $n$.
\end{definition}
\leanref{Development/AbelianVarietyTheory/Definitions/AutomorphicMembership/FiniteDimensionalWdSpace.lean}{GlobalTheory.Pst.IsCrysAt}
\leanref{Development/AbelianVarietyTheory/Definitions/AutomorphicMembership/FiniteDimensionalWdSpace.lean}{GlobalTheory.Pst.IsdRAt}
\leanref{Development/AbelianVarietyTheory/Definitions/AutomorphicMembership/FiniteDimensionalWdSpace.lean}{GlobalTheory.Pst.IsPstAt}
\leanref{Development/AbelianVarietyTheory/Definitions/AutomorphicMembership/FiniteDimensionalWdSpace.lean}{GlobalTheory.Pst.IsJumpsAt}
\leanref{Development/AbelianVarietyTheory/Definitions/AutomorphicMembership/FiniteDimensionalWdSpace.lean}{GlobalTheory.Pst.wdAt}

\section{A local Langlands correspondence on good representations}\label{section-gal-local-langlands}

\begin{definition}\label{gal-irreducible-classes}
In Situation~\ref{gal-local-situation}, let $\operatorname{Irr}(E_w)$ be the set of pairs
$(n,[\pi])$ with $n\geq0$ and $[\pi]$ an isomorphism class of irreducible smooth complex
representations of $\operatorname{GL}_n(E_w)$; $n$ is the \emph{degree}. We write $\pi$ for
$(n,[\pi])$.
\begin{enumerate}
\item A smooth character $\chi$ of $E_w^\times$ is an element of degree $1$. For $\pi$ of degree $n$,
  $\pi\otimes\chi$ is $\pi\otimes(\chi\circ\det)$, and $\pi\nu^s=\pi\otimes|\det|^s$ for $s\in\mathbf R$.
  If the centre $E_w^\times$ acts on $\pi$ by a character, it is written $\omega_\pi$.
\item $\pi$ is \emph{supercuspidal} if $n>0$ and $\pi$ is cuspidal.
\item A multiset $\{\rho_1,\ldots,\rho_r\}$ of supercuspidals is the \emph{supercuspidal support} of
  $\pi$ if $\pi$ is a subquotient of the normalized parabolic induction
  $\rho_1\times\cdots\times\rho_r$ \cite{BZ77}.
\item A \emph{segment} is a pair $(\rho,l)$ with $\rho$ supercuspidal and $l\geq1$, standing for
  $\rho,\rho\nu,\ldots,\rho\nu^{l-1}$. The representation $\pi$ is the \emph{Langlands quotient} of a
  multiset of segments if it is the unique irreducible quotient of the product, in a standard order,
  of the essentially square-integrable representations attached to the segments \cite{Zel80}.
\item $\pi$ has \emph{depth zero} if it has nonzero vectors fixed by
  $K_1=\ker(\operatorname{GL}_n(\mathcal O_w)\to\operatorname{GL}_n(k_w))$.
\end{enumerate}
\end{definition}
\leanref{Development/AbelianVarietyTheory/Definitions/NonGaloisDescent/NonemptyEquivMultiProdOne.lean}{GlobalTheory.LocGL.IrrGL}
\leanref{Development/AbelianVarietyTheory/Definitions/NonGaloisDescent/NonemptyEquivMultiProdOne.lean}{GlobalTheory.LocGL.IrrGL.IsCusp}
\leanref{Development/AbelianVarietyTheory/Definitions/NonGaloisDescent/NonemptyEquivMultiProdOne.lean}{GlobalTheory.LocGL.HasSupport}
\leanref{Development/AbelianVarietyTheory/Definitions/NonGaloisDescent/NonemptyEquivMultiProdOne.lean}{GlobalTheory.LocGL.IsL}
\leanref{Development/AbelianVarietyTheory/Definitions/Foundations/MatrixMeasurableSpace.lean}{GlobalTheory.SplitWeil.KOne}

\begin{definition}\label{gal-good}
In Situation~\ref{gal-local-situation}, a supercuspidal $\rho\in\operatorname{Irr}(E_w)$
(Definition~\ref{gal-irreducible-classes}) is \emph{good} if either $\rho$ has degree $1$, or there
are a prime $g$, a continuous field isomorphism $\mathbf Q_g\to E_w$, an element
$\rho_0\in\operatorname{Irr}(E_w)$ of depth zero and a smooth character $\chi$ of $E_w^\times$ with
$\rho=\rho_0\otimes\chi$. An element $\pi\in\operatorname{Irr}(E_w)$ is \emph{good} if its
supercuspidal support consists of good supercuspidals.
\end{definition}
\leanref{Development/AutomorphicArithmetic/Foundations/LocalLanglands/IsRecTOfIsRec.lean}{UpperTheory.Zero.LLC.GoodCusp}
\leanref{Development/AutomorphicArithmetic/Foundations/LocalLanglands/IsRecTOfIsRec.lean}{UpperTheory.Zero.LLC.Good}

\begin{definition}\label{gal-green-datum}
In Situation~\ref{gal-local-situation}, suppose that there are a prime $g$ and a continuous field
isomorphism $i:\mathbf Q_g\to E_w$, and let $n\geq1$.
\begin{enumerate}
\item A \emph{Green datum of degree $n$} consists of a number field $E'$ with
  $E\subset E'\subset\Omega$, a finite place $w'$ of $E'$ whose prime ideal is
  $\mathfrak m_A\cap\mathcal O_{E'}$, with $g$ not in the square of that ideal and with residue field
  $k'$ of $g^n$ elements, and a local Artin map $a'$ for $E'$, $w'$ and $A$
  (Definition~\ref{gal-local-artin}). Since $E_w\cong\mathbf Q_g$ and $g$ is not in the square of the
  ideal of $w'$, the place $w'$ is unramified over $w$, so the inertia group of $A$ in
  $\operatorname{Gal}(\Omega/E')$ is $I_A$. For $t\in I_A$, write $\bar a'(t)\in k'^\times$ for the
  residue of the unit $a'(t)$.
\item A character $\theta:k'^\times\to\mathbf C^\times$ is \emph{regular} if
  $\theta^{g^i}\neq\theta^{g^j}$ for $0\leq i<j<n$.
\item A cuspidal irreducible complex representation $\tau$ of $\operatorname{GL}_n(\mathbf F_g)$ is
  \emph{attached} to a regular $\theta$ if
  $\operatorname{tr}\tau(y)=(-1)^{n-1}\sum_{j=0}^{n-1}\theta(z^{g^j})$ for every
  $y\in\operatorname{GL}_n(\mathbf F_g)$ with irreducible characteristic polynomial and every root
  $z\in k'^\times$ of that polynomial. By Green's theorem (Theorem~\ref{lit-green-cuspidal}) every
  cuspidal irreducible representation of $\operatorname{GL}_n(\mathbf F_g)$ is attached to some
  regular $\theta$.
\item An irreducible smooth representation $\sigma$ of $\operatorname{GL}_n(E_w)$ \emph{contains}
  $\tau$ if there is a nonzero $\operatorname{GL}_n(\mathbf Z_g)$-equivariant map from $\tau$,
  inflated along reduction modulo $g$, to $\sigma\circ i$.
\end{enumerate}
\end{definition}
\leanref{Development/AbelianVarietyTheory/PublishedTheorems/GaloisCohomology₂/GreenRegularCharacters.lean}{GlobalTheory.Published.IsRegularChar}
\leanref{Development/AbelianVarietyTheory/Proofs/InertiaEigenvalues/PrincipalCongruenceFixedVectorsLeSmoothVectors.lean}{GlobalTheory.SerSys.TameGen.Inflates}
\leanref{Development/AutomorphicArithmetic/Foundations/LocalLanglands/RecTTransport.lean}{UpperTheory.Zero.LLC.Names}

\begin{definition}\label{gal-llc-laws}
In Situation~\ref{gal-local-situation}, let $a$ be a local Artin map
(Definition~\ref{gal-local-artin}). A map
$\operatorname{rc}:\operatorname{Irr}(E_w)\to\mathrm{WD}(A)$
(Definitions~\ref{gal-weil-deligne} and~\ref{gal-irreducible-classes}) \emph{satisfies the laws}
with respect to $a$ if the following hold for all $\pi,\pi'\in\operatorname{Irr}(E_w)$.
\begin{enumerate}
\item \emph{Dimension.} $\operatorname{rc}(\pi)$ has dimension the degree of $\pi$.
\item \emph{Characters.} $\operatorname{rc}(\chi)=\chi\circ a$ for every smooth character $\chi$ of
  $E_w^\times$.
\item \emph{Classification.} If $\pi$ is the Langlands quotient of the segments
  $(\rho_1,l_1),\ldots,(\rho_r,l_r)$, then
  \[
   \operatorname{rc}(\pi)=\bigoplus_{i=1}^r\bigl(\operatorname{rc}(\rho_i)\otimes
   (|\cdot|^{(l_i-1)/2}\circ a)\bigr)\otimes\operatorname{Sp}(l_i).
  \]
\item \emph{Supercuspidals.} If $\pi$ is supercuspidal, $\operatorname{rc}(\pi)$ is irreducible.
\item \emph{Unramified twists.} $\operatorname{rc}(\pi\nu^s)=\operatorname{rc}(\pi)\otimes
  (|\cdot|^s\circ a)$ for every $s\in\mathbf R$.
\item \emph{Absolute value of the central character.} If $\pi$ has central character $\omega_\pi$,
  then $|\det\operatorname{rc}(\pi)(\sigma)|=|\omega_\pi(a(\sigma))|$ for every $\sigma\in W_A$.
\item \emph{Twists of good representations.} If $\pi$ is good (Definition~\ref{gal-good}), then
  $\operatorname{rc}(\pi\otimes\chi)=\operatorname{rc}(\pi)\otimes\operatorname{rc}(\chi)$ for every
  smooth character $\chi$.
\item \emph{Central characters of good representations.} If $\pi$ is good with central character
  $\omega_\pi$, then $\det\operatorname{rc}(\pi)(\sigma)=\omega_\pi(a(\sigma))$ for every
  $\sigma\in W_A$.
\item \emph{Contragredients of good representations.} If $\pi$ is good, then
  $\operatorname{rc}(\pi^\vee)=\operatorname{rc}(\pi)^\vee$.
\item \emph{Injectivity at good representations.} If $\pi'$ is good and
  $\operatorname{rc}(\pi)=\operatorname{rc}(\pi')$, then $\pi=\pi'$.
\item \emph{Depth.} If $\pi$ is good and has depth zero, then $\operatorname{rc}(\pi)$ is tame.
\item \emph{Types.} Suppose there are a prime $g$ and a continuous field isomorphism
  $\mathbf Q_g\to E_w$. Let $n\geq1$, take a Green datum of degree $n$, a regular character $\theta$
  of $k'^\times$ and a cuspidal representation $\tau$ attached to $\theta$
  (Definition~\ref{gal-green-datum}), and let $\sigma$ be an irreducible smooth representation of
  $\operatorname{GL}_n(E_w)$ that contains $\tau$. Then $\operatorname{rc}(\sigma)$ is
  monodromy-free, and for every $t\in I_A$ the characteristic polynomial of
  $\operatorname{rc}(\sigma)(t)$ is
  $\prod_{j=0}^{n-1}\bigl(X-\theta(\bar a'(t))^{g^j}\bigr)$.
\end{enumerate}
A map $\operatorname{rc}$ \emph{satisfies the laws} if it satisfies them with respect to some local
Artin map.
\end{definition}
\leanref{Development/AutomorphicArithmetic/Foundations/LocalLanglands/IsRecTOfIsRec.lean}{UpperTheory.Zero.LLC.LawsT}
\leanref{Development/AutomorphicArithmetic/Foundations/LocalLanglands/IsRecTOfIsRec.lean}{UpperTheory.Zero.LLC.TypeLaw}
\leanref{Development/AutomorphicArithmetic/Foundations/LocalLanglands/IsRecTOfIsRec.lean}{UpperTheory.Zero.LLC.IsRecT}
\leanref{Development/AbelianVarietyTheory/Definitions/NonGaloisDescent/Langlands.lean}{GlobalTheory.LocLanglands.segSum}

\begin{theorem}\label{gal-llc-exists}
Let $E$ be a number field, $\Omega$ an algebraic closure of $E$, $w$ a finite place of $E$ and
$A\subset\Omega$ a valuation ring over $w$, as in Situation~\ref{gal-local-situation}. Then there are
a local Artin map $a$ and a map $\operatorname{rc}:\operatorname{Irr}(E_w)\to\mathrm{WD}(A)$ that
satisfies the laws with respect to $a$ (Definition~\ref{gal-llc-laws}).
\end{theorem}
\leanref{Development/AutomorphicArithmetic/Foundations/LocalLanglands/IsRecTOfIsRec.lean}{UpperTheory.Zero.LLC.RecTExists}
\leanref{Development/AutomorphicArithmetic/Foundations/LocalLanglands/LocalCentralUnits/RecTExists.lean}{UpperTheory.Zero.LLC.recTExists}

\begin{proof}
The construction is in Section~\ref{section-gal-local-construction}. Let $a$ be the local Artin map of
Proposition~\ref{gal-local-artin-unique}, let $r$ be a map that satisfies the laws on supercuspidals
with respect to $a$ (Proposition~\ref{gal-cusp-map}), and let $\operatorname{rc}$ be its Langlands
extension (Lemma~\ref{gal-langlands-extension}). We check the laws of
Definition~\ref{gal-llc-laws} in turn. Law~(1) is Proposition~\ref{gal-extension-laws}(1). For a
character $\chi$, $\operatorname{rc}(\chi)=r(\chi)=\chi\circ a$, which is law~(2). Law~(3) holds by
the definition of $\operatorname{rc}$, since $\operatorname{rc}$ and $r$ agree on supercuspidals.
Law~(4) is law~(1) on supercuspidals. Laws~(5) and~(7) are Proposition~\ref{gal-extension-laws}(3),
with $r(\chi)=\chi\circ a$; laws~(6) and~(8) are Proposition~\ref{gal-extension-laws}(2); law~(9) is
Proposition~\ref{gal-extension-laws}(4), law~(10) is Proposition~\ref{gal-extension-laws}(5) and
law~(11) is Proposition~\ref{gal-extension-laws}(6). For law~(12), let $\sigma$ be a representation
of $\operatorname{GL}_n(E_w)$ that contains a cuspidal representation of
$\operatorname{GL}_n(\mathbf F_g)$. By Proposition~\ref{gal-depth-zero-facts}(2), $\sigma$ is
supercuspidal, so $\operatorname{rc}(\sigma)=r(\sigma)$, and law~(9) on supercuspidals gives the
conclusion of law~(12).
\end{proof}

\begin{lemma}\label{gal-depth-zero-parameter}
In Situation~\ref{gal-local-situation}, suppose that there are a prime $g$ and a continuous field
isomorphism $\mathbf Q_g\to E_w$, and let $a$ be a local Artin map. Let $\rho\in\operatorname{Irr}(E_w)$
be a supercuspidal of degree $n$ and depth zero, with central character $\omega_\rho$. Call
$V\in\mathrm{WD}(A)$ a \emph{parameter} for $\rho$ if $V$ is irreducible, tame and of dimension $n$,
if $\det V(\sigma)=\omega_\rho(a(\sigma))$ for all $\sigma\in W_A$, and if $V$ satisfies the
conclusion of law~(12) of Definition~\ref{gal-llc-laws} for every Green datum of degree $n$, every
regular $\theta$, every $\tau$ attached to $\theta$ and every representation in the class of $\rho$
that contains $\tau$. Then:
\begin{enumerate}
\item any two parameters for $\rho$ are equal;
\item if $\rho'$ is a supercuspidal of depth zero and a parameter for $\rho$ is also a parameter for
  $\rho'$, then $\rho'=\rho$.
\end{enumerate}
\end{lemma}
\leanref{Development/AutomorphicArithmetic/Foundations/LocalLanglands/RecTTransport.lean}{UpperTheory.Zero.LLC.Par}
\leanref{Development/AutomorphicArithmetic/Foundations/LocalLanglands/RecTTransport.lean}{UpperTheory.Zero.LLC.ParUniq}
\leanref{Development/AutomorphicArithmetic/Foundations/LocalLanglands/LocalCentralUnits/ParUniq.lean}{UpperTheory.Zero.LLC.parUniq}

\begin{proof}
Fix a Green datum of degree $n$; one exists because the unramified extension of degree $n$ of
$\mathbf Q_g$ is the completion of a number field at a place of residue degree $n$ over $g$. By local
class field theory, $t\mapsto\bar a'(t)$ is a surjection $I_A\to k'^\times$. A supercuspidal of depth
zero contains a cuspidal irreducible representation $\tau$ of $\operatorname{GL}_n(\mathbf F_g)$
(Proposition~\ref{gal-depth-zero-facts}(1)), and by Theorem~\ref{lit-green-cuspidal} $\tau$ is attached to a regular character
$\theta$ of $k'^\times$.

For (1), let $V$ and $V'$ be parameters for $\rho$. Since $E_w\cong\mathbf Q_g$, $q_w=g$. Let
$t_0\in I_A$ be an element whose image topologically generates the tame inertia group $I_A/P_A$, and
let $\varphi\in W_A$ have degree $-1$. Put $M=V(t_0)$ and $G=V(\varphi)$, and define $M'$ and $G'$
from $V'$ in the same way. These are invertible.
\begin{enumerate}
\item[(a)] Since $V$ is tame and $\varphi^{-1}t_0\varphi\equiv t_0^{g}$ modulo $P_A$, we have
  $MG=GM^g$. A subspace of $V$ stable under $M$ and $G$ is stable under $W_A$: the elements of $W_A$
  that preserve it form a subgroup that contains $P_A$, on which $V$ is trivial, $t_0$, $\varphi$ and
  the open subgroup of $I_A$ on which $V$ is trivial, and these generate $W_A$. As $V$ is
  irreducible, such a subspace is $0$ or $V$. The same holds for $M'$ and $G'$.
\item[(b)] If $Mv=zv$, then $M(Gv)=GM^gv=z^gGv$, so $z\mapsto z^g$ maps the finite set of
  eigenvalues of $M$ to itself. Choose an eigenvalue $z$ and the least $d\geq1$ with
  $z^{g^d}=z$. The operator $G^d$ preserves the $z$-eigenspace of $M$, since $MG^d=G^dM^{g^d}$; let
  $v\neq0$ be an eigenvector of $G^d$ in it, $G^dv=\mu v$. The vectors $v,Gv,\ldots,G^{d-1}v$ are
  eigenvectors of $M$ for the distinct eigenvalues $z,z^g,\ldots,z^{g^{d-1}}$, so they are linearly
  independent, and their span is stable under $M$ and $G$. By (a) it is $V$, so $d=n$ and
  $z^{g^j}\neq z$ for $0<j<n$.
\item[(c)] By the condition of law~(12) in the definition of a parameter, $M$ and $M'$ both have
  characteristic polynomial $\prod_{j=0}^{n-1}(X-\theta(\bar a'(t_0))^{g^j})$. So $z$ is an
  eigenvalue of $M'$ as well, and (b) for $V'$, started from $z$, gives a basis
  $v',G'v',\ldots,G'^{n-1}v'$ of $V'$ with $M'v'=zv'$ and $G'^nv'=\mu'v'$.
\item[(d)] In the basis $(G^jv)$, $G$ is the companion matrix of $X^n-\mu$, so
  $\det G=(-1)^{n-1}\mu$; likewise $\det G'=(-1)^{n-1}\mu'$. Both parameters satisfy
  $\det V(\sigma)=\omega_\rho(a(\sigma))$, so $\det G=\det G'$ and $\mu=\mu'$. Hence the linear
  isomorphism $e:V\to V'$ with $e(G^jv)=G'^jv'$ for $0\leq j<n$ satisfies $eM=M'e$ and $eG=G'e$.
\item[(e)] The monodromy operators of $V$ and $V'$ vanish, because the kernel of $N$ is a nonzero
  subrepresentation of an irreducible representation. The elements $\sigma\in W_A$ with
  $eV(\sigma)=V'(\sigma)e$ form a subgroup that contains $P_A$, $t_0$, $\varphi$ and an open subgroup
  of $I_A$ on which $V$ and $V'$ are trivial, so as in (a) it is $W_A$. Hence $V'=V$.
\end{enumerate}

For (2), let $\tau'$ be a cuspidal representation contained in $\rho'$, attached to a regular
character $\theta'$. Law~(12) for $\rho'$ shows that the eigenvalues of $V(t)$ are also the
$\theta'(\bar a'(t))^{g^j}$. Comparing the sums of the eigenvalues for all $t$, using that $\bar a'$
is surjective and that characters of $k'^\times$ are linearly independent, gives
$\theta'=\theta^{g^j}$ for some $j<n$, so $\tau'$ and $\tau$ have the same character by
Theorem~\ref{lit-green-cuspidal}. For every $\sigma\in W_A$,
$\omega_\rho(a(\sigma))=\det V(\sigma)=\omega_{\rho'}(a(\sigma))$, and $a$ is surjective, so
$\omega_\rho=\omega_{\rho'}$. Two irreducible smooth representations of $\operatorname{GL}_n(E_w)$
that contain cuspidal representations of $\operatorname{GL}_n(\mathbf F_g)$ with the same character
and have the same central character are isomorphic (Proposition~\ref{gal-depth-zero-facts}(3)).
Hence $\rho'=\rho$.
\end{proof}

\begin{proposition}\label{gal-llc-unique}
Let $E$ be a number field, $\Omega$ an algebraic closure of $E$, $w$ a finite place of $E$ and
$A\subset\Omega$ a valuation ring over $w$, as in Situation~\ref{gal-local-situation}. Let
$\operatorname{rc},\operatorname{rc}':\operatorname{Irr}(E_w)\to\mathrm{WD}(A)$ be maps that satisfy
the laws (Definition~\ref{gal-llc-laws}). Then $\operatorname{rc}(\pi)=\operatorname{rc}'(\pi)$ for
every good $\pi\in\operatorname{Irr}(E_w)$ (Definition~\ref{gal-good}).
\end{proposition}
\leanref{Development/AutomorphicArithmetic/Foundations/LocalLanglands/IsRecTOfIsRec.lean}{UpperTheory.Zero.LLC.RecTUnique}
\leanref{Development/AutomorphicArithmetic/Foundations/LocalLanglands/LocalCentralUnits/RecTUnique.lean}{UpperTheory.Zero.LLC.recTUnique}

\begin{proof}
Let $a$ and $a'$ be the local Artin maps for $\operatorname{rc}$ and $\operatorname{rc}'$. By
Proposition~\ref{gal-local-artin-unique}, $a=a'$.

First let $\rho$ be a good supercuspidal. If $\rho$ is a character $\chi$, law~(2) gives
$\operatorname{rc}(\chi)=\chi\circ a=\operatorname{rc}'(\chi)$. Otherwise $E_w\cong\mathbf Q_g$ and
$\rho=\rho_0\otimes\chi$ with $\rho_0$ of depth zero; $\rho_0$ is supercuspidal and good. By
laws~(1), (4), (8), (11) and~(12), $\operatorname{rc}(\rho_0)$ and $\operatorname{rc}'(\rho_0)$ are
both parameters for $\rho_0$ in the sense of Lemma~\ref{gal-depth-zero-parameter}, so they are equal.
Law~(7) and law~(2) then give
$\operatorname{rc}(\rho)=\operatorname{rc}(\rho_0)\otimes(\chi\circ a)=\operatorname{rc}'(\rho)$.

Now let $\pi$ be good. It is the Langlands quotient of a multiset of segments $(\rho_i,l_i)$
\cite{Zel80}, and every $\rho_i$ lies in the supercuspidal support of $\pi$, so it is good. Law~(3)
expresses $\operatorname{rc}(\pi)$ and $\operatorname{rc}'(\pi)$ by the same formula in the values at
the $\rho_i$, which agree by the first step.
\end{proof}

\begin{definition}\label{gal-llc-characterization}
In Situation~\ref{gal-local-situation}, let $a$ be a local Artin map and let $\varepsilon$ be a
system of local constants: for every finite extension of $E$ inside $\Omega$, every finite place of
it under $A$, every nontrivial continuous additive character $\psi$ of the completion and every
Weil--Deligne representation $x$ of the corresponding Weil group, a number
$\varepsilon(x,\psi)\in\mathbf C^\times$; it is required to be additive in $x$, to be Tate's local
constant on characters, and to be inductive in degree zero \cite{Del73}. A map
$\operatorname{rc}:\operatorname{Irr}(E_w)\to\mathrm{WD}(A)$ \emph{satisfies the characterization of
the local Langlands correspondence} with respect to $a$ and $\varepsilon$ if:
\begin{enumerate}
\item $\operatorname{rc}$ is bijective, and laws~(1), (2) and~(3) of Definition~\ref{gal-llc-laws}
  hold;
\item for every nontrivial continuous additive character $\psi$ of $E_w$, every supercuspidal $\pi$
  of degree $m+1$ and every $\tau\in\operatorname{Irr}(E_w)$ of degree $m$, the functional equation
  of Jacquet, Piatetski-Shapiro and Shalika \cite{JPSS83} at $s=1/2$, for $\psi$-Whittaker
  functionals on $\pi$ and $\psi^{-1}$-Whittaker functionals on $\tau$, holds with the constant
  $\varepsilon(\operatorname{rc}(\pi)\otimes\operatorname{rc}(\tau),\psi)$;
\item laws~(7), (8) and~(9) hold for all $\pi$, good or not;
\item $\pi$ is supercuspidal if and only if $\operatorname{rc}(\pi)$ is irreducible;
\item $\operatorname{rc}(\chi\circ\det_t)=\bigoplus_{j=0}^{t-1}\operatorname{rc}(\chi)\otimes
  (|\cdot|^{(t-1)/2-j}\circ a)$ for every smooth character $\chi$ and $t\geq0$;
\item if $\sigma$ and $\tau$ have degrees $r$ and $t$ and $\pi\cong\sigma\times\tau$ is
  irreducible, then $\operatorname{rc}(\pi)=\operatorname{rc}(\sigma)\oplus\operatorname{rc}(\tau)$;
\item if $\pi$ has supercuspidal support $\{\rho_1,\ldots,\rho_r\}$, the characteristic polynomial
  of $\operatorname{rc}(\pi)(\sigma)$ is the product of those of the $\operatorname{rc}(\rho_i)(\sigma)$;
\item $\pi$ is spherical if and only if $\operatorname{rc}(\pi)$ is unramified, and if $\pi$ is
  spherical with Satake parameters $\alpha_1,\ldots,\alpha_n$, then $\operatorname{rc}(\pi)$ is the
  monodromy-free representation $\sigma\mapsto\operatorname{diag}(\alpha_i^{-d(\sigma)})$;
\item $\pi$ has depth zero if and only if $\operatorname{rc}(\pi)$ is tame.
\end{enumerate}
\end{definition}
\leanref{Development/AbelianVarietyTheory/PublishedTheorems/AutomorphicForms₂/LocalLanglandsGL.lean}{GlobalTheory.LocLanglands.Laws}
\leanref{Development/AbelianVarietyTheory/PublishedTheorems/AutomorphicForms₂/LocalLanglandsGL.lean}{GlobalTheory.LocLanglands.IsRec}
\leanref{Development/AbelianVarietyTheory/PublishedTheorems/AutomorphicForms₂/HasExponentZeroIff.lean}{GlobalTheory.LocLanglands.HasPairEps}
\leanref{Development/AbelianVarietyTheory/PublishedTheorems/AutomorphicForms₂/HasExponentZeroIff.lean}{GlobalTheory.LocLanglands.IsDeligneEps}

\begin{proposition}\label{gal-llc-of-rec}
Assume the following statement $(\mathrm{T})$: for every number field $E$, every algebraic closure
$\Omega$ of $E$, every finite place $w$ of $E$ and valuation ring $A\subset\Omega$ over $w$, every map
$\operatorname{rc}:\operatorname{Irr}(E_w)\to\mathrm{WD}(A)$ that satisfies the characterization of
the local Langlands correspondence (Definition~\ref{gal-llc-characterization}) with respect to some
local Artin map and some system of local constants also satisfies law~(12) of
Definition~\ref{gal-llc-laws}. Then, in Situation~\ref{gal-local-situation}, every map that
satisfies the characterization of the local Langlands correspondence satisfies the laws.
\end{proposition}
\leanref{Development/AbelianVarietyTheory/PublishedTheorems/GaloisCohomology₂/DepthZeroTypeInertia.lean}{GlobalTheory.Published.DepthZeroTypeInertia}
\leanref{Development/AutomorphicArithmetic/Foundations/LocalLanglands/IsRecTOfIsRec.lean}{UpperTheory.Zero.LLC.lawsT_of_laws}
\leanref{Development/AutomorphicArithmetic/Foundations/LocalLanglands/IsRecTOfIsRec.lean}{UpperTheory.Zero.LLC.isRecT_of_isRec}

\begin{proof}
Let $\operatorname{rc}$ satisfy the characterization with respect to $a$ and $\varepsilon$. Laws~(1),
(2) and~(3) are part of the characterization. Law~(4) is one direction of item~(4) of
Definition~\ref{gal-llc-characterization}, and laws~(7), (8) and~(9) are the restrictions of
item~(3) to good representations. Law~(5) is law~(7) for the character $|\cdot|^s$, together with
law~(2). Law~(6) follows from law~(8) for all $\pi$ by taking absolute values. Law~(10) follows from
injectivity, and law~(11) is one direction of item~(9). Law~(12) is the hypothesis $(\mathrm{T})$.
\end{proof}

\begin{remark}\label{gal-llc-scope}
The laws of Definition~\ref{gal-llc-laws} involve no $L$-factors and no $\varepsilon$-factors. By
Proposition~\ref{gal-llc-unique} they determine $\operatorname{rc}(\pi)$ for good $\pi$; for a
supercuspidal $\rho$ that is not good they impose only laws~(1), (4), (5), (6) and~(10). The
characterization of Definition~\ref{gal-llc-characterization} lists properties of the kind that
characterize the local Langlands correspondence of Harris and Taylor \cite{HT01} and of Henniart
\cite{Hen00} (see also \cite{Sch13}). If a map with that characterization exists and the statement
$(\mathrm{T})$ holds, Propositions~\ref{gal-llc-of-rec} and~\ref{gal-llc-unique} show that every map
satisfying the laws agrees with it on good representations. Neither the existence of such a map nor
$(\mathrm{T})$ is proved in this book, and the later chapters use only maps that satisfy the laws.
\end{remark}

\section{Construction of a map satisfying the laws}\label{section-gal-local-construction}

\begin{proposition}\label{gal-depth-zero-facts}
In Situation~\ref{gal-local-situation}, let $n\geq1$.
\begin{enumerate}
\item Suppose there are a prime $g$ and a continuous field isomorphism $i:\mathbf Q_g\to E_w$, and use
  it to identify $\operatorname{GL}_n(\mathbf Z_g)$ with $\operatorname{GL}_n(\mathcal O_w)$. Every
  irreducible smooth representation $\sigma$ of $\operatorname{GL}_n(E_w)$ that is supercuspidal and
  has depth zero contains a cuspidal irreducible representation $\tau$ of
  $\operatorname{GL}_n(\mathbf F_g)$ (Definition~\ref{gal-green-datum}). If $\sigma\cong\sigma'$ and
  $\sigma$ and $\sigma'$ contain cuspidal representations $\tau$ and $\tau'$, then $\tau$ and $\tau'$
  have the same character.
\item In the setting of (1), an irreducible smooth representation of $\operatorname{GL}_n(E_w)$ that
  contains a cuspidal irreducible representation of $\operatorname{GL}_n(\mathbf F_g)$ is cuspidal and
  has depth zero.
\item In the setting of (1), let $\sigma$ and $\sigma'$ be irreducible smooth representations of
  $\operatorname{GL}_n(E_w)$ with the same central character that contain cuspidal representations
  $\tau$ and $\tau'$ of $\operatorname{GL}_n(\mathbf F_g)$ with the same character. Then
  $\sigma\cong\sigma'$.
\item Let $k'$ be a field with $g^n$ elements, let $\theta$ be a regular character of $k'^\times$ and
  let $\tau$ be a cuspidal representation of $\operatorname{GL}_n(\mathbf F_g)$ attached to $\theta$
  (Definition~\ref{gal-green-datum}). Then every scalar $c\in\mathbf F_g^\times$ acts on $\tau$ by
  $\theta(c)$.
\item Let $x\in\operatorname{Irr}(E_w)$ have positive degree and depth zero
  (Definition~\ref{gal-irreducible-classes}) and let $\chi$ be a smooth character of $E_w^\times$.
  Then $x\otimes\chi$ has depth zero if and only if $\chi$ is \emph{tame}, that is, trivial on
  $1+\mathfrak m_w$.
\item If $x\in\operatorname{Irr}(E_w)$ has depth zero, every member of its supercuspidal support has
  depth zero.
\item Let $a$ be a local Artin map (Definition~\ref{gal-local-artin}). A smooth character $\chi$ of
  $E_w^\times$ is tame if and only if $\chi\circ a$ is trivial on $P_A$.
\end{enumerate}
Items (1), (2), (3) and (6) are the theory of depth-zero types \cite{MP96}, \cite{BK93}, \cite{Car84};
item (4) follows from Green's character formula \cite{Gre55}; item (7) is local class field theory
\cite{Ser67}, since $a$ maps $P_A$ onto $1+\mathfrak m_w$.
\end{proposition}
\leanref{Development/AutomorphicArithmetic/Foundations/LocalLanglands/LocalCentralUnits/TypeEx.lean}{UpperTheory.Zero.LLC.typeEx}
\leanref{Development/AutomorphicArithmetic/Foundations/LocalLanglands/LocalCentralUnits/DetByType.lean}{UpperTheory.Zero.LLC.detByType}
\leanref{Development/AutomorphicArithmetic/Foundations/LocalLanglands/LocalCentralUnits/NameCentral.lean}{UpperTheory.Zero.LLC.nameCentral}
\leanref{Development/AutomorphicArithmetic/Foundations/LocalLanglands/LocalCentralUnits/TwTame.lean}{UpperTheory.Zero.LLC.twTame}
\leanref{Development/AutomorphicArithmetic/Foundations/LocalLanglands/LocalCentralUnits/TameArt.lean}{UpperTheory.Zero.LLC.tameArt}

\begin{lemma}\label{gal-tame-inertia-character}
In Situation~\ref{gal-local-situation}, suppose there are a prime $g$ and a continuous field
isomorphism $i:\mathbf Q_g\to E_w$. Let $n\geq1$ and let $a$ be a local Artin map for $E$, $w$ and $A$
(Definition~\ref{gal-local-artin}).
\begin{enumerate}
\item There is a Green datum of degree $n$ (Definition~\ref{gal-green-datum}).
\item Take a Green datum of degree $n$, with residue field $k'$ and local Artin map $a'$. Then
  $q_w=g$, $k'$ has degree $n$ over $\mathbf F_g$, and $t\mapsto\bar a'(t)$ is a surjective
  homomorphism $I_A\to k'^\times$ with open kernel, trivial on $P_A$, such that
  $\bar a'(\sigma t\sigma^{-1})=\bar a'(t)^g$ for every $t\in I_A$ and every $\sigma\in W_A$ of degree
  $1$.
\item In the setting of (2), for every $t\in I_A$ there is a unit $u\in\mathbf Z_g^\times$ with
  $a(t)=i(u)$ whose residue in $\mathbf F_g$ is $\prod_{j=0}^{n-1}\bar a'(t)^{g^j}$, the norm of
  $\bar a'(t)$ from $k'$ to $\mathbf F_g$.
\item In the setting of (2), let $\varpi\in\Omega$ satisfy $\varpi^{g^n-1}=g$. For every $t\in I_A$,
  $t(\varpi)/\varpi$ lies in $A$, and its residue is the image of $\bar a'(t)$ under the inclusion of
  $k'$ in the residue field of $A$. In particular, for two Green data of degree $n$, with residue fields
  $k'$ and $k''$, there is a field isomorphism $e:k'\to k''$ that carries the first map
  $t\mapsto\bar a'(t)$ to the second.
\end{enumerate}
\end{lemma}
\leanref{Development/AutomorphicArithmetic/Foundations/LocalLanglands/LocalCentralUnits/ExtEx.lean}{UpperTheory.Zero.LLC.extEx}
\leanref{Development/AutomorphicArithmetic/Foundations/LocalLanglands/RecTTransport.lean}{UpperTheory.Zero.LLC.TameSys}
\leanref{Development/AutomorphicArithmetic/Foundations/LocalLanglands/LocalCentralUnits/FinrankK.lean}{UpperTheory.Zero.LLC.resEx}
\leanref{Development/AutomorphicArithmetic/Foundations/LocalLanglands/LocalCentralUnits/ResBase.lean}{UpperTheory.Zero.LLC.resBase}
\leanref{Development/AutomorphicArithmetic/Foundations/LocalLanglands/LocalCentralUnits/TameRecAll.lean}{UpperTheory.Zero.LLC.tameRecAll}
\leanref{Development/AutomorphicArithmetic/Foundations/LocalLanglands/LocalCentralUnits/ResIndep.lean}{UpperTheory.Zero.LLC.resIndep}

\begin{proof}
For (1), there is a number field $E'$ with $E\subset E'\subset\Omega$, containing a primitive
$(g^n-1)$-th root of unity and fixed by $I_A$, such that the place $w'$ of $E'$ under $A$ is unramified
over $g$ with $g^n$ elements in its residue field: a primitive $(g^n-1)$-th root of unity generates
the unramified extension of degree $n$ of $E_w\cong\mathbf Q_g$. A local Artin map for $E'$, $w'$ and
$A$ exists by Proposition~\ref{gal-local-artin-unique}.

For (2), $q_w=g$ because $E_w\cong\mathbf Q_g$, and $k'$ has $g^n$ elements. By local class field
theory $a'$ maps the inertia group $I_A$ of $E'$ onto the units of $E'_{w'}$, and the residue map
from these units onto $k'^\times$ is surjective, so $t\mapsto\bar a'(t)$ is a surjective homomorphism.
It is trivial on the preimage under $a'$ of $1+\mathfrak m_{w'}$, which is open since $a'$ is
continuous on $I_A$, and which contains $P_A$ since $a'$ maps $P_A$ into $1+\mathfrak m_{w'}$. An
element $\sigma\in W_A$ of degree $1$ acts on $E'_{w'}$ as the arithmetic Frobenius, which raises
residues to the $g$-th power, and $a'(\sigma t\sigma^{-1})=\sigma(a'(t))$; hence
$\bar a'(\sigma t\sigma^{-1})=\bar a'(t)^g$.

For (3), $a(t)$ is the norm of $a'(t)$ from $E'_{w'}$ to $E_w$ (compatibility of local class field
theory with restriction from $E$ to $E'$, which holds here because $I_A$ fixes $E'$ and both maps are
compatible with the global reciprocity maps). The norm of a unit of the unramified extension
$E'_{w'}$ is a unit of $\mathbf Z_g$ whose residue is the norm of the residue.

For (4), the explicit tame reciprocity law for the unramified extension $E'_{w'}$ of $\mathbf Q_g$,
in which $g$ is a uniformizer, says that $t(\varpi)/\varpi$ and $a'(t)$ have the same residue
\cite{Ser67}. The residue of $t(\varpi)/\varpi$ does not involve the Green datum. Hence the two maps
$t\mapsto\bar a'(t)$ coming from two Green data agree after the identification of both residue fields
with the subfield of $g^n$ elements of the residue field of $A$.
\end{proof}

\begin{lemma}\label{gal-induced-parameter}
In Situation~\ref{gal-local-situation}, let $\varphi\in W_A$ have degree $1$, let $n\geq1$, and let
$(\eta_j)_{j\in\mathbf Z}$ be homomorphisms $I_A\to\mathbf C^\times$ such that $\eta_0$ is trivial on
an open subgroup of $I_A$, $\eta_j(\varphi t\varphi^{-1})=\eta_{j+1}(t)$ for all $j$ and $t\in I_A$,
and $\eta_j=\eta_{j'}$ if and only if $n$ divides $j-j'$. Let $c\in\mathbf C^\times$. Then there is
an irreducible $V\in\mathrm{WD}(A)$ of dimension $n$ (Definitions~\ref{gal-weil-deligne}
and~\ref{gal-wd-classes}) such that $\det V(\varphi)=c$, the characteristic polynomial of $V(t)$ is
$\prod_{j=0}^{n-1}(X-\eta_j(t))$ for every $t\in I_A$, and $V$ is tame if every $\eta_j$ is trivial on
$P_A$.
\end{lemma}
\leanref{Development/AutomorphicArithmetic/Foundations/LocalLanglands/RecTTransport.lean}{UpperTheory.Zero.LLC.IndEx}
\leanref{Development/AutomorphicArithmetic/Foundations/LocalLanglands/LocalCentralUnits/IsSemisimpleInducedParam.lean}{UpperTheory.Zero.LLC.indEx}

\begin{proof}
Put $b=(-1)^{n-1}c$ and let $V$ be the $n$-dimensional space of functions $f:\mathbf Z\to\mathbf C$
with $f(j+n)=bf(j)$ for all $j$. Every $\sigma\in W_A$ is $\varphi^{m}t$ with $m=d(\sigma)$ and
$t\in I_A$, and we let it act by
\[
 (\rho(\varphi^mt)f)(j)=\eta_{j+m}(t)\,f(j+m).
\]
The relation $\eta_j(\varphi t\varphi^{-1})=\eta_{j+1}(t)$ shows that $\rho$ is a homomorphism. The
kernel of $\rho$ on $I_A$ contains $\bigcap_{j<n}\ker\eta_j$, which is open since
$\eta_j=\eta_0\circ(\text{conjugation by }\varphi^{-j})$. Some power of each $\rho(\sigma)$ is a
scalar, so $\rho(\sigma)$ is semisimple, and $(V,\rho,0)$ is a Weil--Deligne representation.

In the basis of $V$ given by the values at $0,\ldots,n-1$, the group $I_A$ acts diagonally through
$\eta_0,\ldots,\eta_{n-1}$, which gives the characteristic polynomial of $V(t)$. These characters are
distinct, so a nonzero subrepresentation contains one of the basis vectors, and then all of them,
because $\rho(\varphi)$ shifts the basis cyclically. Hence $V$ is irreducible. The matrix of
$\rho(\varphi)$ is the companion matrix of $X^n-b$, whose determinant is $(-1)^{n-1}b=c$. If every
$\eta_j$ is trivial on $P_A$, so is $\rho$, and $V$ is tame.
\end{proof}

\begin{proposition}\label{gal-depth-zero-parameter-exists}
In Situation~\ref{gal-local-situation}, suppose there are a prime $g$ and a continuous field
isomorphism $\mathbf Q_g\to E_w$, and let $a$ be a local Artin map. Then every supercuspidal
$\rho\in\operatorname{Irr}(E_w)$ of depth zero has a parameter in the sense of
Lemma~\ref{gal-depth-zero-parameter}.
\end{proposition}
\leanref{Development/AutomorphicArithmetic/Foundations/LocalLanglands/RecTTransport.lean}{UpperTheory.Zero.LLC.ParEx}
\leanref{Development/AutomorphicArithmetic/Foundations/LocalLanglands/LocalCentralUnits/CentralUnit.lean}{UpperTheory.Zero.LLC.parEx}

\begin{proof}
Let $n$ be the degree of $\rho$ and $\omega_\rho$ its central character, and choose $\varphi\in W_A$ of
degree $1$. By Proposition~\ref{gal-depth-zero-facts}(1), $\rho$ contains a cuspidal representation
$\tau$ of $\operatorname{GL}_n(\mathbf F_g)$. Fix a Green datum of degree $n$, with residue field
$k'$ (Lemma~\ref{gal-tame-inertia-character}). By Green's theorem
(Theorem~\ref{lit-green-cuspidal}), $\tau$ is attached to a regular character $\theta$ of
$k'^\times$. Put $\eta_j(t)=\theta(\bar a'(t)^{g^{j'}})$, where $j'\in\{0,\ldots,n-1\}$ is the residue
of $j$ modulo $n$. Since every element of $k'^\times$ is fixed by $x\mapsto x^{g^n}$, the relation
$\bar a'(\varphi t\varphi^{-1})=\bar a'(t)^g$ of Lemma~\ref{gal-tame-inertia-character}(2) gives
$\eta_j(\varphi t\varphi^{-1})=\eta_{j+1}(t)$. The kernel of $\eta_0$ contains the open kernel of
$\bar a'$. If $\eta_j=\eta_{j'}$, then $\theta^{g^{j}}=\theta^{g^{j'}}$ because $\bar a'$ is surjective,
and regularity of $\theta$ gives that $n$ divides $j-j'$.

Lemma~\ref{gal-induced-parameter}, with $c=\omega_\rho(a(\varphi))$, gives an irreducible
$V\in\mathrm{WD}(A)$ of dimension $n$ with $\det V(\varphi)=\omega_\rho(a(\varphi))$ and with
characteristic polynomials $\prod_{j<n}(X-\theta(\bar a'(t))^{g^j})$ on $I_A$. It is tame since each
$\eta_j$ is trivial on $P_A$. We check that $V$ is a parameter for $\rho$.

\emph{Determinant.} The homomorphism $\sigma\mapsto\det V(\sigma)$ and $\omega_\rho\circ a$ agree at
$\varphi$. For $t\in I_A$, $\det V(t)=\theta(N(\bar a'(t)))$, where $N$ is the norm from $k'$ to
$\mathbf F_g$. By Lemma~\ref{gal-tame-inertia-character}(3), $a(t)=i(u)$ with $u\in\mathbf Z_g^\times$
reducing to $N(\bar a'(t))$. The scalar matrix $u$ lies in $\operatorname{GL}_n(\mathbf Z_g)$; it acts
on $\rho$ by $\omega_\rho(a(t))$ and on $\tau$ by $\theta(N(\bar a'(t)))$
(Proposition~\ref{gal-depth-zero-facts}(4)), and a nonzero equivariant map from $\tau$ to $\rho$
shows that these numbers are equal. Since $W_A$ is generated by $\varphi$ and $I_A$, $\det V$ equals
$\omega_\rho\circ a$.

\emph{Types.} Let a second Green datum of degree $n$, a regular character $\theta'$, a cuspidal $\tau'$
attached to $\theta'$ and a representation in the class of $\rho$ containing $\tau'$ be given. By
Proposition~\ref{gal-depth-zero-facts}(1), $\tau'$ and $\tau$ have the same character. By
Lemma~\ref{gal-tame-inertia-character}(4), an isomorphism $e$ of the two residue fields carries the
first map $\bar a'$ to the second, and $\tau$ is attached to $\theta'\circ e$. By Green's theorem
$\theta'\circ e=\theta^{g^j}$ for some $j<n$. The multiset $\{\theta(\bar a'(t))^{g^{j+k}}:k<n\}$ is
the multiset $\{\theta(\bar a'(t))^{g^{k}}:k<n\}$, since $\theta^{g^n}=\theta$. So $V(t)$ has the
characteristic polynomial required by law~(12), and $V$ is monodromy-free since it is irreducible.
\end{proof}

\begin{lemma}\label{gal-parameter-operations}
In Situation~\ref{gal-local-situation}, suppose there are a prime $g$ and a continuous field
isomorphism $\mathbf Q_g\to E_w$, and let $a$ be a local Artin map. Parameters are taken in the sense
of Lemma~\ref{gal-depth-zero-parameter}.
\begin{enumerate}
\item Let $\chi$ be a smooth character of $E_w^\times$, viewed as an element of degree $1$ of
  $\operatorname{Irr}(E_w)$. Then $\chi$ has depth zero if and only if it is tame, and if $\chi$ is
  tame then $\chi\circ a$ is a parameter for $\chi$.
\item Let $\rho$ be a supercuspidal of depth zero with a parameter $V$, and let $\chi$ be a tame
  character. Then $V\otimes(\chi\circ a)$ is a parameter for $\rho\otimes\chi$.
\item Let $\rho$ be a supercuspidal of depth zero with a parameter $V$. Then $V^\vee$ is a parameter
  for the contragredient $\rho^\vee$.
\end{enumerate}
\end{lemma}
\leanref{Development/AutomorphicArithmetic/Foundations/LocalLanglands/LocalCentralUnits/ParOne.lean}{UpperTheory.Zero.LLC.parOne}
\leanref{Development/AutomorphicArithmetic/Foundations/LocalLanglands/LocalCentralUnits/NormEqProd.lean}{UpperTheory.Zero.LLC.parTw}
\leanref{Development/AutomorphicArithmetic/Foundations/LocalLanglands/LocalCentralUnits/ParDual.lean}{UpperTheory.Zero.LLC.parDual}

\begin{proof}
For (1), a character has nonzero vectors fixed by $1+\mathfrak m_w$ if and only if it is trivial
there. If $\chi$ is tame, $\chi\circ a$ is irreducible of dimension $1$ and tame by
Proposition~\ref{gal-depth-zero-facts}(7), its determinant is $\chi\circ a$, and the type condition
for $n=1$ is the statement of Lemma~\ref{gal-tame-inertia-character}(3) for $n=1$: a character of
$\operatorname{GL}_1(\mathbf F_g)$ contained in $\chi$ is the reduction $\chi_0$ of $\chi$ on
$\mathbf Z_g^\times$, attached to $\theta=\chi_0$, and $\chi(a(t))=\chi_0(\bar a'(t))$.

For (2), $V\otimes(\chi\circ a)$ is irreducible of dimension $n$, and it is tame because
$\chi\circ a$ is trivial on $P_A$ (Proposition~\ref{gal-depth-zero-facts}(7)). The central character
of $\rho\otimes\chi$ is $\omega_\rho\chi^n$, and $\det(V\otimes(\chi\circ a))=\det V\cdot(\chi\circ
a)^n$. For the type condition, let $\chi_0$ be the character of $\mathbf F_g^\times$ through which
$\chi$ factors on $\mathbf Z_g^\times$. If a representation in the class of $\rho\otimes\chi$ contains
a cuspidal $\tau$ attached to $\theta$, then twisting by $\chi^{-1}$ shows that $\rho$ contains
$\tau\otimes(\chi_0^{-1}\circ\det)$, which is cuspidal and attached to the regular character
$\theta\cdot(\chi_0^{-1}\circ N)$. The type condition for $V$ gives the eigenvalues
$\theta(\bar a'(t))^{g^j}\chi_0(N\bar a'(t))^{-1}$ of $V(t)$, and $\chi(a(t))=\chi_0(N\bar a'(t))$ by
Lemma~\ref{gal-tame-inertia-character}(3). Multiplying, $(V\otimes(\chi\circ a))(t)$ has the
eigenvalues $\theta(\bar a'(t))^{g^j}$.

For (3), $V^\vee$ is irreducible, tame and of dimension $n$, and its determinant is
$(\omega_\rho\circ a)^{-1}=\omega_{\rho^\vee}\circ a$. The contragredient $\rho^\vee$ has depth zero,
and it contains $\tau^\vee$ when $\rho$ contains $\tau$; the contragredient of a cuspidal
representation attached to $\theta$ is attached to $\theta^{-1}$ by Green's character formula.
Hence the inertial eigenvalues for $\rho^\vee$ are the inverses of those for $\rho$, which are the
eigenvalues of $V^\vee(t)$.
\end{proof}

\begin{lemma}\label{gal-good-map}
In Situation~\ref{gal-local-situation}, let $a$ be a local Artin map. There is a map
$r_0:\operatorname{Irr}(E_w)\to\mathrm{WD}(A)$ such that:
\begin{enumerate}
\item $r_0(\chi)=\chi\circ a$ for every smooth character $\chi$ of $E_w^\times$;
\item $r_0(x\otimes\chi)=r_0(x)\otimes(\chi\circ a)$ for every good supercuspidal $x$
  (Definition~\ref{gal-good}) and every smooth character $\chi$;
\item if there are a prime $g$ and a continuous field isomorphism $\mathbf Q_g\to E_w$, then
  $r_0(\rho)$ is a parameter for every supercuspidal $\rho$ of depth zero, in the sense of
  Lemma~\ref{gal-depth-zero-parameter}.
\end{enumerate}
\end{lemma}
\leanref{Development/AutomorphicArithmetic/Foundations/LocalLanglands/RecTTransport.lean}{UpperTheory.Zero.LLC.GoodMap}
\leanref{Development/AutomorphicArithmetic/Foundations/LocalLanglands/LocalCentralUnits/GoodMap.lean}{UpperTheory.Zero.LLC.goodMap}

\begin{proof}
Say that $V\in\mathrm{WD}(A)$ \emph{fits} $x\in\operatorname{Irr}(E_w)$ if either $x=\chi$ is a
character and $V=\chi\circ a$, or $E_w\cong\mathbf Q_g$ for some prime $g$ and $x=\rho_0\otimes\chi$
and $V=V_0\otimes(\chi\circ a)$, where $\rho_0$ is a supercuspidal of depth zero, $V_0$ is a parameter
for $\rho_0$ and $\chi$ is a smooth character. Every good supercuspidal has a $V$ that fits it, by
Proposition~\ref{gal-depth-zero-parameter-exists}. We show that at most one $V$ fits a given $x$.

Two characters $\chi$, $\chi'$ that are equal as representations are equal, so the first case gives
one value. Suppose $\rho_0\otimes\chi=\rho_1\otimes\chi_1$ with parameters $V_0$, $V_1$. Then
$\rho_1=\rho_0\otimes\chi\chi_1^{-1}$, and since $\rho_0$ and $\rho_1$ have depth zero,
$\chi\chi_1^{-1}$ is tame (Proposition~\ref{gal-depth-zero-facts}(5)). By
Lemma~\ref{gal-parameter-operations}(2), $V_0\otimes(\chi\chi_1^{-1}\circ a)$ is a parameter for
$\rho_1$, so it equals $V_1$ by Lemma~\ref{gal-depth-zero-parameter}(1), and $V_0\otimes(\chi\circ
a)=V_1\otimes(\chi_1\circ a)$. If $\chi'=\rho_0\otimes\chi$ is a character, then $\rho_0$ is the
character $\chi'\chi^{-1}$ of depth zero; it is tame by Lemma~\ref{gal-parameter-operations}(1), its
parameter is $(\chi'\chi^{-1})\circ a$ by the same lemma and
Lemma~\ref{gal-depth-zero-parameter}(1), and both candidate values equal $\chi'\circ a$.

Let $r_0(x)$ be the unique $V$ that fits $x$ when there is one, and an arbitrary element otherwise.
Then (1) and (3) hold, the latter with $\chi=1$. For (2), if $V$ fits $x$ then
$V\otimes(\chi\circ a)$ fits $x\otimes\chi$.
\end{proof}

\begin{lemma}\label{gal-placeholder}
In Situation~\ref{gal-local-situation}, let $n\geq2$. There is an irreducible $J\in\mathrm{WD}(A)$ of
dimension $n$ with $|\det J(\sigma)|=1$ for every $\sigma\in W_A$, such that, if there are a prime $g$
and a continuous field isomorphism $\mathbf Q_g\to E_w$, then $J$ is not of the form $V\otimes\xi$
with $V$ tame and $\xi:W_A\to\mathbf C^\times$ a homomorphism trivial on an open subgroup of $I_A$.
\end{lemma}
\leanref{Development/AutomorphicArithmetic/Foundations/LocalLanglands/RecTTransport.lean}{UpperTheory.Zero.LLC.JunkEx}
\leanref{Development/AutomorphicArithmetic/Foundations/LocalLanglands/LocalCentralUnits/JunkEx.lean}{UpperTheory.Zero.LLC.junkEx}
\leanref{Development/AutomorphicArithmetic/WhittakerHecke/TameInertiaCharacterFrobeniusConjEqPow.lean}{UpperTheory.tameInertiaCharacter_frobenius_conj_eq_pow}

\begin{proof}
Let $\varphi\in W_A$ be a Frobenius element of degree $1$. Choose a number field $L$ with
$E\subset L\subset\Omega$ such that the place $u$ of $L$ under $A$ is unramified over $w$ of residue
degree $n$, and choose $c\in\mathcal O_E$ in the prime of $w$ and not in its square, a uniformizer of
$L_u$. Let $K$ be the residue field of $L_u$, with $q_w^n$ elements, and $f$ its automorphism
$x\mapsto x^{q_w}$, of order $n$. The map $\lambda$ sending a unit $z$ of $L_u$ to the residue of
$(z^{q_w^n-1}-1)/c$ is a homomorphism onto the additive group of $K$, trivial on $1+c^2\mathcal O_{L_u}$. Let $a_L$ be the local Artin map for $L$, $u$ and $A$, and put
$\pi(t)=\lambda(a_L(t))$ for $t\in I_A$, which is also the inertia group of $A$ over $L$. Local class
field theory shows that $\pi:I_A\to K$ is surjective with open kernel, and since $\varphi$ acts on
$L_u$ as a Frobenius that fixes $c$, $\pi(\varphi t\varphi^{-1})=f(\pi(t))$.

Choose a character $\psi$ of the additive group of $K$ with $\psi\circ f^m\neq\psi$ unless $f^m=1$,
and put $\eta_j=\psi\circ f^j\circ\pi$. These satisfy the hypotheses of
Lemma~\ref{gal-induced-parameter}, which with $c=1$ gives an irreducible $J$ of dimension $n$ with
$\det J(\varphi)=1$ and inertial eigenvalues $\eta_0(t),\ldots,\eta_{n-1}(t)$. These are $p$-th roots
of unity, $p$ the residue characteristic, so $|\det J|=1$ on $I_A$ and on $\varphi$, hence on $W_A$.

Suppose $J=V\otimes\xi$ with $V$ tame. Then on $P_A$ all eigenvalues of $J(t)$ equal $\xi(t)$, so
$\eta_0$ and $\eta_1$ agree on $P_A$. Put $\kappa=\eta_0/\eta_1$. Since $\varphi s\varphi^{-1}s^{-q_w}$
lies in $P_A$ for $s\in I_A$, $\kappa(\varphi s\varphi^{-1})=\kappa(s)^{q_w}$, and by induction
$\kappa(\varphi^ns\varphi^{-n})=\kappa(s)^{q_w^n}=1$. But $\pi(\varphi^ns\varphi^{-n})=f^n(\pi(s))=\pi(s)$
and $\kappa$ factors through $\pi$, so $\kappa(s)=1$. Hence $\eta_0=\eta_1$, contradicting
$\psi\neq\psi\circ f$.
\end{proof}

\begin{lemma}\label{gal-exponent}
In Situation~\ref{gal-local-situation}, let $a$ be a local Artin map. There is a function
$e:\operatorname{Irr}(E_w)\to\mathbf R$ such that for every $x$ of positive degree $n$ with central
character $\omega_x$:
\begin{enumerate}
\item $e(x\nu^s)=e(x)+s$ for every $s\in\mathbf R$;
\item for every $J\in\mathrm{WD}(A)$ of dimension $n$ with $|\det J(\sigma)|=1$ for all $\sigma$,
  $|\det(J\otimes(|\cdot|^{e(x)}\circ a))(\sigma)|=|\omega_x(a(\sigma))|$ for all $\sigma\in W_A$.
\end{enumerate}
\end{lemma}
\leanref{Development/AutomorphicArithmetic/Foundations/LocalLanglands/RecTTransport.lean}{UpperTheory.Zero.LLC.Expo}
\leanref{Development/AutomorphicArithmetic/Foundations/LocalLanglands/LocalCentralUnits/ExponentMulAbsPow.lean}{UpperTheory.Zero.LLC.expo}

\begin{proof}
Every smooth character $\omega$ of $E_w^\times$ satisfies $|\omega(z)|=|z|^{s(\omega)}$ for a unique
real $s(\omega)$. Put $e(x)=s(\omega_x)/n$. The central character of $x\nu^s$ is
$\omega_x|\cdot|^{ns}$, which gives (1), and
$|\det(J\otimes(|\cdot|^{e(x)}\circ a))(\sigma)|=|a(\sigma)|^{ne(x)}=|\omega_x(a(\sigma))|$, which
gives (2).
\end{proof}

\begin{definition}\label{gal-cusp-laws}
In Situation~\ref{gal-local-situation}, let $a$ be a local Artin map. A map
$r:\operatorname{Irr}(E_w)\to\mathrm{WD}(A)$ \emph{satisfies the laws on supercuspidals} with respect
to $a$ if the following hold for all supercuspidals $x,y\in\operatorname{Irr}(E_w)$:
\begin{enumerate}
\item $r(x)$ is irreducible of dimension the degree of $x$;
\item $r(\chi)=\chi\circ a$ for every smooth character $\chi$;
\item $r(x\nu^s)=r(x)\otimes(|\cdot|^s\circ a)$ for every $s\in\mathbf R$;
\item $|\det r(x)(\sigma)|=|\omega_x(a(\sigma))|$ for all $\sigma\in W_A$;
\item if $x$ is good, $r(x\otimes\chi)=r(x)\otimes(\chi\circ a)$ for every smooth $\chi$, and
  $\det r(x)(\sigma)=\omega_x(a(\sigma))$ for all $\sigma$;
\item if $x$ is good and $y\cong x^\vee$, then $r(y)=r(x)^\vee$;
\item if $y$ is good and $r(x)=r(y)$, then $x=y$;
\item if $x$ is good and has depth zero, then $r(x)$ is tame;
\item $r(x)$ satisfies the conclusion of law~(12) of Definition~\ref{gal-llc-laws} at $x$: for every
  Green datum, regular $\theta$, cuspidal $\tau$ attached to $\theta$ and representation in the class
  of $x$ that contains $\tau$, $r(x)$ is monodromy-free and $r(x)(t)$ has characteristic polynomial
  $\prod_{j<n}(X-\theta(\bar a'(t))^{g^j})$ for $t\in I_A$.
\end{enumerate}
\end{definition}
\leanref{Development/AutomorphicArithmetic/Foundations/LocalLanglands/RecTTransport.lean}{UpperTheory.Zero.LLC.CuspT}

\begin{proposition}\label{gal-cusp-map}
In Situation~\ref{gal-local-situation}, let $a$ be a local Artin map. There is a map
$r:\operatorname{Irr}(E_w)\to\mathrm{WD}(A)$ that satisfies the laws on supercuspidals with respect to
$a$ (Definition~\ref{gal-cusp-laws}). One such map is $r(x)=r_0(x)$ for good $x$ and
$r(x)=J_n\otimes(|\cdot|^{e(x)}\circ a)$ for $x$ of degree $n$ that is not good, where $r_0$ is the
map of Lemma~\ref{gal-good-map}, $J_n$ the representation of Lemma~\ref{gal-placeholder} and $e$ the
function of Lemma~\ref{gal-exponent}.
\end{proposition}
\leanref{Development/AutomorphicArithmetic/Foundations/LocalLanglands/RecTTransport.lean}{UpperTheory.Zero.LLC.CuspEx}
\leanref{Development/AutomorphicArithmetic/Foundations/LocalLanglands/LocalCentralUnits/CuspEx.lean}{UpperTheory.Zero.LLC.cuspEx}

\begin{proof}
A supercuspidal that is not good has degree at least $2$, since characters are good, so $J_n$ is
defined. Twists $x\otimes\chi$, unramified twists $x\nu^s$ and contragredients of a good supercuspidal
are good, and those of a supercuspidal that is not good are not good.

(1) For $x=\chi$ a character, $r(x)=\chi\circ a$. For $x=\rho_0\otimes\chi$ with $\rho_0$ of depth
zero, $r(x)=r_0(\rho_0)\otimes(\chi\circ a)$ is a twist of a parameter, hence irreducible of the right
dimension. For $x$ not good, $J_n$ is irreducible and so are its twists. (2) is
Lemma~\ref{gal-good-map}(1). (3) For good $x$ it is Lemma~\ref{gal-good-map}(2) with
$\chi=|\cdot|^s$; otherwise it is Lemma~\ref{gal-exponent}(1). (4) For good $x$ it follows from the
equality in (5); otherwise it is Lemma~\ref{gal-exponent}(2). (5) The first part is
Lemma~\ref{gal-good-map}(2). For the determinant: if $x=\chi$, $\det r(x)=\chi\circ a=\omega_x\circ
a$; if $x=\rho_0\otimes\chi$, $\det r(x)=(\omega_{\rho_0}\circ a)(\chi\circ a)^n=\omega_x\circ a$ by
the determinant condition on parameters.

(6) If $x=\chi$, then $y=\chi^{-1}$ and $r(y)=\chi^{-1}\circ a=r(x)^\vee$. If $x=\rho_0\otimes\chi$,
then $y\cong\rho_0^\vee\otimes\chi^{-1}$; by Lemma~\ref{gal-parameter-operations}(3) and
Lemma~\ref{gal-depth-zero-parameter}(1), $r_0(\rho_0^\vee)=r_0(\rho_0)^\vee$, and twisting gives
$r(y)=r(x)^\vee$.

(7) The dimensions show that $x$ and $y$ have the same degree. Suppose first that $x$ is good. If both
are characters $\chi,\chi'$, then $\chi\circ a=\chi'\circ a$ and $a$ is surjective, so $\chi=\chi'$.
Otherwise $x=u\otimes\chi$ and $y=v\otimes\chi'$ with $u,v$ supercuspidal of depth zero and parameters
$V_u=r_0(u)$, $V_v=r_0(v)$, and $V_u=V_v\otimes(\psi\circ a)$ with $\psi=\chi'\chi^{-1}$. Both $V_u$ and
$V_v$ are tame, so $\psi\circ a$ is trivial on $P_A$, and $\psi$ is tame by
Proposition~\ref{gal-depth-zero-facts}(7). Then $v\otimes\psi$ has depth zero and $V_u$ is a
parameter for it (Lemma~\ref{gal-parameter-operations}(2)), so $u=v\otimes\psi$ by
Lemma~\ref{gal-depth-zero-parameter}(2), and $x=y$. If $x$ is not good, then $y$ has degree at least
$2$, so $y=v\otimes\chi$ as above and $r(y)=V_v\otimes(\chi\circ a)$ with $V_v$ tame. Then
$J_n=V_v\otimes(\chi|\cdot|^{-e(x)}\circ a)$, which Lemma~\ref{gal-placeholder} excludes.

(8) A good character of depth zero is tame (Lemma~\ref{gal-parameter-operations}(1)), so $\chi\circ a$
is tame by Proposition~\ref{gal-depth-zero-facts}(7). A good supercuspidal of degree at least $2$ and
depth zero is itself a supercuspidal $\rho$ of depth zero with $E_w\cong\mathbf Q_g$, and
$r(\rho)=r_0(\rho)$ is a parameter, hence tame.

(9) If a representation in the class of $x$ contains a cuspidal representation of
$\operatorname{GL}_n(\mathbf F_g)$, then $x$ has depth zero by
Proposition~\ref{gal-depth-zero-facts}(2), so it is good, and $r(x)=r_0(x)$ is a parameter, which
satisfies the type condition. Otherwise there is nothing to prove.
\end{proof}

\begin{lemma}\label{gal-langlands-extension}
In Situation~\ref{gal-local-situation}, let $r:\operatorname{Irr}(E_w)\to\mathrm{WD}(A)$ be any map.
Every $x\in\operatorname{Irr}(E_w)$ is the Langlands quotient of a multiset of segments
$(\rho_1,l_1),\ldots,(\rho_k,l_k)$, unique up to order (Definition~\ref{gal-irreducible-classes}),
and
\[
 \operatorname{rc}(x)=\bigoplus_{i=1}^k\bigl(r(\rho_i)\otimes(|\cdot|^{(l_i-1)/2}\circ a)\bigr)\otimes
 \operatorname{Sp}(l_i)
\]
defines a map $\operatorname{rc}$ that satisfies law~(3) of Definition~\ref{gal-llc-laws} with respect
to $r$ in place of $\operatorname{rc}$ on the right, and that agrees with $r$ on supercuspidals. Call
$\operatorname{rc}$ the \emph{Langlands extension} of $r$. The existence and uniqueness of the
multiset are the Langlands classification of Zelevinsky \cite{Zel80}.
\end{lemma}
\leanref{Development/AbelianVarietyTheory/Definitions/NonGaloisDescent/CuspLaws.lean}{GlobalTheory.LocLanglands.IsLExt}
\leanref{Development/AbelianVarietyTheory/Proofs/InertiaEigenvalues/LanglandsExtensionUniqueOnCuspidal.lean}{GlobalTheory.LocLanglands.exists_isLExt}
\leanref{Development/AbelianVarietyTheory/PublishedTheorems/GaloisCohomology₂/LanglandsExhaustionGL.lean}{GlobalTheory.Published.LanglandsExhaustionGL}
\leanref{Development/AbelianVarietyTheory/PublishedTheorems/GaloisCohomology₂/LanglandsMultisegments.lean}{GlobalTheory.Published.LanglandsMultisegments}

\begin{proof}
The sum depends only on the multiset, so $\operatorname{rc}$ is well defined. A supercuspidal $x$ is the
Langlands quotient of the single segment $(x,1)$, and the summand is $r(x)\otimes\operatorname{Sp}(1)=
r(x)$.
\end{proof}

\begin{proposition}\label{gal-extension-laws}
In Situation~\ref{gal-local-situation}, let $a$ be a local Artin map, let $r$ satisfy the laws on
supercuspidals with respect to $a$ (Definition~\ref{gal-cusp-laws}), and let $\operatorname{rc}$ be its
Langlands extension (Lemma~\ref{gal-langlands-extension}). Let $x\in\operatorname{Irr}(E_w)$ have
degree $n$, central character $\omega_x$ and supercuspidal support $s$.
\begin{enumerate}
\item $\operatorname{rc}(x)$ has dimension $n$, and for every $\sigma\in W_A$ the characteristic
  polynomial of $\operatorname{rc}(x)(\sigma)$ is the product over $y\in s$ of those of $r(y)(\sigma)$.
\item $|\det\operatorname{rc}(x)(\sigma)|=|\omega_x(a(\sigma))|$, and if $x$ is good,
  $\det\operatorname{rc}(x)(\sigma)=\omega_x(a(\sigma))$.
\item $\operatorname{rc}(x\nu^s)=\operatorname{rc}(x)\otimes(|\cdot|^s\circ a)$, and if $x$ is good,
  $\operatorname{rc}(x\otimes\chi)=\operatorname{rc}(x)\otimes(\chi\circ a)$ for every smooth $\chi$.
\item If $x$ is good, $\operatorname{rc}(x^\vee)=\operatorname{rc}(x)^\vee$.
\item If $y$ is good and $\operatorname{rc}(x)=\operatorname{rc}(y)$, then $x=y$.
\item If $x$ is good and has depth zero, $\operatorname{rc}(x)$ is tame.
\end{enumerate}
\end{proposition}
\leanref{Development/AutomorphicArithmetic/Foundations/LocalLanglands/LocalCentralUnits/ExtDim.lean}{UpperTheory.Zero.LLC.extDim}
\leanref{Development/AutomorphicArithmetic/Foundations/LocalLanglands/LocalCentralUnits/RcTwOfSegments.lean}{UpperTheory.Zero.LLC.extTw}
\leanref{Development/AutomorphicArithmetic/Foundations/LocalLanglands/LocalCentralUnits/ExtDual.lean}{UpperTheory.Zero.LLC.extDual}
\leanref{Development/AutomorphicArithmetic/Foundations/LocalLanglands/LocalCentralUnits/ExtInj.lean}{UpperTheory.Zero.LLC.extInj}
\leanref{Development/AutomorphicArithmetic/Foundations/LocalLanglands/LocalCentralUnits/ExtDepth.lean}{UpperTheory.Zero.LLC.extDepth}

\begin{proof}
Write $x$ as the Langlands quotient of segments $(\rho_i,l_i)$. Its supercuspidal support is the
multiset of the $\rho_i\nu^j$ with $j<l_i$ \cite{Zel80}, and $x$ is good if and only if every $\rho_i$
is good, since twisting by $\nu^j$ preserves goodness and the supercuspidal support is unique
\cite{BZ77}.

(1) Since $v_w(a(\sigma))=-d(\sigma)$, $|a(\sigma)|=q_w^{d(\sigma)}$. On the basis vector $e_j$ of
$\operatorname{Sp}(l)$, $\sigma$ acts by $q_w^{(j-(l-1)/2)d(\sigma)}$, so the summand for a segment
$(\rho,l)$ has at $\sigma$ the eigenvalues of $r(\rho)(\sigma)$ multiplied by
$|a(\sigma)|^{(l-1)/2}q_w^{(j-(l-1)/2)d(\sigma)}=|a(\sigma)|^j$ for $j<l$. By law~(3) on
supercuspidals these are the eigenvalues of the $r(\rho\nu^j)(\sigma)$, $j<l$. Taking degrees gives
the dimension.

(2) The central character of $x$ is the product of those of the members of its support \cite{BZ77}.
Taking determinants in (1) and applying laws~(4) and~(5) on supercuspidals to each member gives both
statements.

(3) The segments of $x\nu^s$ are the $(\rho_i\nu^s,l_i)$, and those of $x\otimes\chi$ are the
$(\rho_i\otimes\chi,l_i)$; apply laws~(3) and~(5) on supercuspidals to each summand.

(4) The contragredient of the Langlands quotient of the segments $(\rho_i,l_i)$ is the Langlands
quotient of the segments $(\rho_i^\vee\nu^{1-l_i},l_i)$ \cite{Zel80}. For good $\rho_i$, law~(6)
applied to $\rho_i\nu^{l_i-1}$ and law~(3) on supercuspidals give $r(\rho_i^\vee\nu^{1-l_i})=r(\rho_i)^\vee\otimes(|\cdot|^{1-l_i}\circ a)$, and
the dual of $(V\otimes(|\cdot|^{(l-1)/2}\circ a))\otimes\operatorname{Sp}(l)$ is
$(V^\vee\otimes(|\cdot|^{-(l-1)/2}\circ a))\otimes\operatorname{Sp}(l)$.

(5) A Weil--Deligne representation that is a sum of terms $V_i\otimes\operatorname{Sp}(m_i)$ with $V_i$
irreducible determines the multiset of pairs $(V_i,m_i)$. Every $r(\rho)$ is irreducible, so
$\operatorname{rc}(x)=\operatorname{rc}(y)$ gives a bijection between the segments of $x$ and of $y$
under which $r(\rho_i)\otimes(|\cdot|^{(l_i-1)/2}\circ a)$ and the lengths agree. The segments of $y$
start at good supercuspidals, so law~(7) on supercuspidals and law~(3) give equal segments. By the
uniqueness in the Langlands classification, $x=y$.

(6) By Proposition~\ref{gal-depth-zero-facts}(6), every member $y$ of the support of $x$ has depth zero,
and it is good, so $r(y)$ is tame by law~(8) on supercuspidals; by (1), $\operatorname{rc}(x)(\sigma)$
has characteristic polynomial $(X-1)^n$ for $\sigma\in P_A$.
\end{proof}

\section{Galois representations attached to polarized cuspidal representations}\label{section-gal-galois}

\begin{definition}\label{gal-polarized}
Let $E$ be a CM field with maximal totally real subfield $F$ and complex conjugation $c$, and let
$a\geq1$.
\begin{enumerate}
\item A \emph{cuspidal automorphic representation} of $\operatorname{GL}_a(\mathbf A_E)$ is a nonzero
  space $\pi$ of smooth cusp forms on $\operatorname{GL}_a(E)\backslash\operatorname{GL}_a(\mathbf A_E)$
  that is stable and irreducible under $\operatorname{GL}_a(\mathbf A_{E,f})$, the maximal compact
  subgroup at infinity and the Lie algebra at infinity. For a finite place $w$ of $E$, a
  \emph{local component} of $\pi$ at $w$ is an irreducible smooth representation $\pi_w$ of
  $\operatorname{GL}_a(E_w)$ with a nonzero $\operatorname{GL}_a(E_w)$-equivariant map
  $\pi_w\to\pi$.
\item For an embedding $\tau:E\to\mathbf C$, the \emph{infinitesimal character} $Z_\tau(\pi)$ of
  $\pi$ at $\tau$ is the multiset of $a$ complex numbers by which the centre of the enveloping
  algebra of the $\tau$-factor acts on $\pi$. The representation $\pi$ is \emph{regular algebraic}
  if for every $\tau$ the multiset $Z_\tau(\pi)$ consists of $a$ distinct elements of
  $\frac{a-1}{2}+\mathbf Z$.
\item Let $\chi$ be a Hecke character of $F$, a continuous character of
  $F^\times\backslash\mathbf A_F^\times$. The pair $(\pi,\chi)$ is \emph{polarized} if
  $\chi_v(-1)=(-1)^a$ for every real place $v$ of $F$ and
  $\pi^c\cong\pi^\vee\otimes(\chi\circ N_{E/F}\circ\det)$, in the sense that there is a
  nondegenerate pairing between $\pi\otimes(\chi\circ N_{E/F}\circ\det)^{-1}$ and $\pi$ that is
  invariant under $g$ acting on the first factor together with $c(g)$ acting on the second.
\end{enumerate}
\end{definition}
\leanref{Development/AbelianVarietyTheory/Definitions/AutomorphicMembership/ExpoConjOfIsCsd.lean}{GlobalTheory.CuspRep.CuspSub}
\leanref{Development/AbelianVarietyTheory/Definitions/AutomorphicMembership/ExpoConjOfIsCsd.lean}{GlobalTheory.CuspRep.CuspSub.IsLocComp}
\leanref{Development/AbelianVarietyTheory/Definitions/AutomorphicMembership/ExpoConjOfIsCsd.lean}{GlobalTheory.CuspRep.CG.infChar}
\leanref{Development/AbelianVarietyTheory/Definitions/AutomorphicMembership/ExpoConjOfIsCsd.lean}{GlobalTheory.CuspRep.CG.IsRegAlg}
\leanref{Development/AbelianVarietyTheory/Definitions/GaloisGeometry/DualIsConjByOne.lean}{GlobalTheory.ContraPair.PolarizedPair}

\begin{definition}\label{gal-compatible}
In Situation~\ref{gal-local-situation}, let $x,y\in\mathrm{WD}(A)$ and let $n$ be the dimension of
$y$. Then $x$ is \emph{compatible} with $y$ if:
\begin{enumerate}
\item $x$ has dimension $n$;
\item if $y$ is abelian and monodromy-free (Definition~\ref{gal-wd-classes}), then $x=y$;
\item if there are a prime $g$ and a continuous field isomorphism $\mathbf Q_g\to E_w$ with $n=1$ or
  $g>16n^2$, and $y$ is tame, monodromy-free and has multiplicity-free tame part, then $x$ is
  monodromy-free and $x(\sigma)$ and $y(\sigma)$ have the same characteristic polynomial for every
  $\sigma\in I_A$.
\end{enumerate}
\end{definition}
\leanref{Development/AutomorphicArithmetic/Foundations/GaloisRepresentations/TypeConditions/WeightsConjBis.lean}{UpperTheory.GalStar.IsQpBig}
\leanref{Development/AutomorphicArithmetic/Foundations/GaloisRepresentations/TypeConditions/WeightsConjBis.lean}{UpperTheory.GalStar.AwayB}

\begin{definition}\label{gal-attached}
Let $\mathcal R$ assign to every number field $E$, algebraic closure $\Omega$, finite place $w$ and
valuation ring $A\subset\Omega$ over $w$ a set $\mathcal R(A)$ of maps
$\operatorname{Irr}(E_w)\to\mathrm{WD}(A)$. Let $E$ be a CM field with complex conjugation $c$, let
$\Omega$ be an algebraic closure of $E$, let $\ell$ be a prime, let
$\iota:\bar{\mathbf Q}_\ell\to\mathbf C$ be a field isomorphism, let $a\geq1$ and let $\pi$ be a
cuspidal automorphic representation of $\operatorname{GL}_a(\mathbf A_E)$
(Definition~\ref{gal-polarized}). A triple $(r,m,\mathrm{HT})$, consisting of a homomorphism
$r:G_E\to\operatorname{GL}_a(\bar{\mathbf Q}_\ell)$, an integer $m$ and a multiset of integers
$\mathrm{HT}_\tau$ for each embedding $\tau:E\to\bar{\mathbf Q}_\ell$, is \emph{attached to $\pi$
relative to $\mathcal R$} if the following hold.
\begin{enumerate}
\item $r$ is continuous and semisimple.
\item \emph{Away from $\ell$.} For every finite place $w$ of $E$ not over $\ell$, every valuation ring
  $A$ over $w$, every $\operatorname{rc}\in\mathcal R(A)$ and every local component $\pi_w$ of $\pi$,
  $\mathrm{WD}_\iota(r|_{W_A})$ (Definition~\ref{gal-wd-galois}) is compatible with
  $\operatorname{rc}(\pi_w\nu^{(1-a)/2})$ (Definition~\ref{gal-compatible}).
\item \emph{Weights.} $\mathrm{HT}_\tau=\{\frac{a-1}{2}-z:z\in Z_{\iota\circ\tau}(\pi)\}$ for every
  $\tau$.
\item \emph{Symmetry.} $\mathrm{HT}_{\tau\circ c}=\{m-h:h\in\mathrm{HT}_\tau\}$ for every $\tau$.
\item \emph{Above $\ell$.} Let $w$ be a place of $E$ over $\ell$, $A$ a valuation ring over $w$ and
  $\pi_w$ a local component of $\pi$. In the notation of Definition~\ref{gal-p-adic-hodge}:
  \begin{enumerate}
  \item if $\pi_w$ has nonzero vectors fixed by $\operatorname{GL}_a(\mathcal O_w)$, then $r|_{D_A}$ is
    crystalline and de Rham, and $\mathrm{HT}_\tau(r|_{D_A})=\mathrm{HT}_\tau$ for every $\tau$ that
    extends continuously to $E_w$;
  \item for every $\operatorname{rc}\in\mathcal R(A)$ such that $y=\operatorname{rc}(\pi_w\nu^{(1-a)/2})$
    is monodromy-free and either abelian, or tame with multiplicity-free tame part where $E_w$ is
    isomorphic to $\mathbf Q_g$ for a prime $g$ with $a=1$ or $g>16a^2$, the restriction $r|_{D_A}$
    is de Rham and potentially semistable, $\mathrm{HT}_\tau(r|_{D_A})=\mathrm{HT}_\tau$ for every
    $\tau$ that extends continuously to $E_w$, and $\mathrm{WD}_\iota(r|_{D_A})$ is compatible with
    $y$.
  \end{enumerate}
\end{enumerate}
\end{definition}
\leanref{Development/AutomorphicArithmetic/Foundations/GaloisRepresentations/TypeConditions/WeightsConjBis.lean}{UpperTheory.GalStar.AboveB}
\leanref{Development/AutomorphicArithmetic/Foundations/GaloisRepresentations/TypeConditions/WeightsConjBis.lean}{UpperTheory.GalStar.PropsB}
\leanref{Development/AutomorphicArithmetic/Foundations/GaloisRepresentations/TypeConditions/WeightsConjAlt.lean}{UpperTheory.GalStar.RecNotion}

\begin{definition}\label{gal-galois-representations}
Let $\mathcal R$ be as in Definition~\ref{gal-attached}. \emph{Galois representations exist relative
to $\mathcal R$} if for every CM field $E$ with maximal totally real subfield $F$, every algebraic
closure $\Omega$ of $E$, every prime $\ell$, every field isomorphism
$\iota:\bar{\mathbf Q}_\ell\to\mathbf C$, every $a\geq1$, every regular algebraic cuspidal
automorphic representation $\pi$ of $\operatorname{GL}_a(\mathbf A_E)$ and every Hecke character
$\chi$ of $F$ with $(\pi,\chi)$ polarized (Definition~\ref{gal-polarized}), there is a triple
$(r,m,\mathrm{HT})$ attached to $\pi$ relative to $\mathcal R$.
\end{definition}
\leanref{Development/AutomorphicArithmetic/Foundations/GaloisRepresentations/TypeConditions/WeightsConjBis.lean}{UpperTheory.GalStar.Row0b}

\begin{theorem}\label{gal-galois-representations-any-rc}
For every number field $E$, algebraic closure $\Omega$, finite place $w$ and valuation ring $A$ over
$w$, let $\mathcal R(A)$ be the set of maps $\operatorname{Irr}(E_w)\to\mathrm{WD}(A)$ that satisfy
the laws (Definition~\ref{gal-llc-laws}). Then Galois representations exist relative to
$\mathcal R$ (Definition~\ref{gal-galois-representations}).
\end{theorem}
\leanref{Development/AutomorphicArithmetic/Foundations/GaloisRepresentations/IntegralStructures/Development.lean}{UpperTheory.Gal.Int.recT}
\leanref{Development/AutomorphicArithmetic/Foundations/GaloisRepresentations/IntegralStructures/RowSrc.lean}{UpperTheory.Gal.Int.galoisLocalGlobalStrong_recParameter}
\leanref{Development/AutomorphicArithmetic/Foundations/GaloisRepresentations/IntegralStructures/Development.lean}{UpperTheory.Gal.Int.galoisLocalGlobal_of_geometry_lift}
\leanref{Development/AutomorphicArithmetic/Foundations/GaloisRepresentations/IntegralStructures/Development.lean}{UpperTheory.Gal.Int.coreAll_of}

\begin{proof}
The construction is in Section~\ref{section-gal-construction}. By Proposition~\ref{gal-core}, which
rests on the seeds of Theorem~\ref{gal-seed-representations}, the limits of
Proposition~\ref{gal-limits} and the theta lifts of Theorem~\ref{gal-theta-lift-large-rank}, existence
and the four conditions of Definition~\ref{gal-scope} hold in the self-dual scope. By
Proposition~\ref{gal-core-all} they hold in the first scope. Proposition~\ref{gal-row-from-core} then
gives, for every $(\pi,\chi)$ in scope, a triple attached to $\pi$ relative to $\mathcal R$ in which
clause~(3) of Definition~\ref{gal-compatible} and clause~(5)(b) of Definition~\ref{gal-attached} hold
whenever $E_w\cong\mathbf Q_g$, with no condition on $g$. Restricting these clauses to the case $a=1$
or $g>16a^2$ gives Definition~\ref{gal-attached} as stated.
\end{proof}

\begin{remark}\label{gal-galois-scope}
Let $(r,m,\mathrm{HT})$ be attached to $\pi$ relative to the maps satisfying the laws, let $w$ be a
place not over $\ell$, and put $y=\operatorname{rc}(\pi_w\nu^{(1-a)/2})$. By the laws, $y$ is
abelian and monodromy-free if and only if the supercuspidal support of $\pi_w$ consists of characters
and $\pi_w$ is the Langlands quotient of segments of length one. Indeed, by law~(3), $y$ is the sum
over the segments $(\rho_i,l_i)$ of $\pi_w$ of twists of $\operatorname{rc}(\rho_i)\otimes
\operatorname{Sp}(l_i)$; this is monodromy-free only if every $l_i=1$, and abelian only if every
$\operatorname{rc}(\rho_i)$, which is irreducible by law~(4), has dimension $1$, that is, every
$\rho_i$ is a character by law~(1); conversely, law~(2) makes such a $y$ a sum of characters. In this
case item~(2) of Definition~\ref{gal-attached} gives $\mathrm{WD}_\iota(r|_{W_A})=y$. When $E_w\cong
\mathbf Q_g$ with $a=1$ or $g>16a^2$ and $y$ is tame, monodromy-free and has multiplicity-free tame
part, it gives only that $\mathrm{WD}_\iota(r|_{W_A})$ is monodromy-free and that its elements of
$I_A$ have the same characteristic polynomials as those of $y$. At the other places it gives nothing
beyond the dimension; for example
when $\pi_w$ is a Steinberg representation, or when $E_w$ is not isomorphic to some $\mathbf Q_g$ and
the supercuspidal support of $\pi_w$ contains a supercuspidal of degree at least $2$. Definition~\ref{gal-attached}
asserts nothing about $r^c$ and $r^\vee$, that is, no polarization of $r$, and nothing about the
absolute values of the eigenvalues of Frobenius elements, that is, no purity. The published
theorem \cite[Theorem~2.1.1]{BLGGT14}, which rests on \cite{Shi11}, \cite{CH13} and \cite{Car12},
asserts in addition the polarization, the full local-global compatibility at every finite place,
including the places over $\ell$, with the local Langlands correspondence, and potential
semistability at every place over $\ell$.
\end{remark}

\section{Ingredients of the construction}\label{section-gal-ingredients}

\begin{theorem}\label{gal-cyclic-base-change}
Let $K$ be a totally complex number field and $M$ a totally complex Galois extension of $K$ of prime
degree $p$. Let $a\geq1$ and let $\pi$ be a cuspidal automorphic representation of
$\operatorname{GL}_a(\mathbf A_K)$, in the sense of Definition~\ref{gal-polarized}. Suppose that there
is no Hecke character $\eta$ of $K$ with the following properties: $\eta\neq1$, $\eta^p=1$; for all
but finitely many places $u$ of $M$, with $v$ the place of $K$ under $u$ and $f$ with
$N(u)=N(v)^f$, $\eta(\varpi_v)^f=1$; and for all but finitely many places $v$ of $K$, $\eta$ is
unramified at $v$ and $\pi$ has Satake parameters $B$ at $v$ with
$\{\eta(\varpi_v)b:b\in B\}=B$. Then there is a cuspidal automorphic representation $\Pi$ of
$\operatorname{GL}_a(\mathbf A_M)$ such that:
\begin{enumerate}
\item for every place $u$ of $M$, with $v$ the place of $K$ under $u$ and $N(u)=N(v)^f$, if $\pi$ has
  Satake parameters $B$ at $v$ then $\Pi$ has Satake parameters $\{b^f:b\in B\}$ at $u$;
\item $Z_\tau(\Pi)=Z_{\tau|_K}(\pi)$ for every embedding $\tau:M\to\mathbf C$;
\item for every place $v$ of $K$, place $u$ of $M$ and continuous isomorphism $i:K_v\to M_u$
  extending $K\to M$, if $\sigma$ and $\sigma'$ are local components of $\pi$ at $v$ and of $\Pi$ at
  $u$, then $\sigma'\circ i\cong\sigma$.
\end{enumerate}
This is part of the cyclic base change of Arthur and Clozel \cite{AC89}.
\end{theorem}
\leanref{Development/AbelianVarietyTheory/PublishedTheorems/GaloisGeometry₂/ArthurClozelStrongTC.lean}{GlobalTheory.Published.ArthurClozelStrongTC}
\leanref{Development/AutomorphicArithmetic/AutomorphicVarieties/PublishedTheorems/CyclicBaseChangeC.lean}{UpperTheory.V.Published.IsNormChar}
\leanref{Development/AutomorphicArithmetic/AutomorphicVarieties/PublishedTheorems/CyclicBaseChangeC.lean}{UpperTheory.V.Published.SelfTw}
\leanref{Development/AbelianVarietyTheory/Definitions/GaloisGeometry/IsWeight.lean}{GlobalTheory.GalCut.IsLift}
\leanref{Development/AutomorphicArithmetic/Foundations/GaloisRepresentations/LiftFactorization/FirstQuadraticSquareHolds.lean}{UpperTheory.Gal.AbelianGaloisNormality.strongTC_holds}

\begin{definition}\label{gal-definite-unitary}
Let $E$ be a CM field with maximal totally real subfield $F$ and complex conjugation $c$, and let
$n\geq1$.
\begin{enumerate}
\item $U_n=\{g\in\operatorname{GL}_n(\mathbf A_E):{}^t\!c(g)\,g=1\}$, with $c$ applied entrywise, is
  the group of adelic points of the definite unitary group of the standard Hermitian form; $U_n(F)$
  is $U_n\cap\operatorname{GL}_n(E)$.
\item An \emph{automorphic representation} of $U_n$ is a nonzero space $U$ of continuous functions
  on $U_n$, left invariant under $U_n(F)$, each invariant on the right under an open subgroup of the
  finite part and finite under the part at infinity, that is stable and irreducible under right
  translation by $U_n$.
\item For a finite place $w$ of $E$, let $U_{n,w}\subset U_n$ be the subgroup of elements equal to
  $1$ at infinity and at every place not over the place of $F$ under $w$, and
  $c_w:U_{n,w}\to\operatorname{GL}_n(E_w)$ the projection. An irreducible smooth representation
  $\sigma$ of $\operatorname{GL}_n(E_w)$ is a \emph{local component} of $U$ at $w$ if there is a
  nonzero map $\sigma\to U$ that intertwines $\sigma\circ c_w$ with right translation. $U$ is
  \emph{unramified} at $w$ if it has a nonzero vector fixed by $c_w^{-1}(\operatorname{GL}_n(\mathcal
  O_w))$, and $U$ has \emph{Satake parameters} $B$ at $w$ if it has a spherical local component at
  $w$ with Satake parameters $B$, normalized so that the $j$-th Hecke operator acts by
  $q_w^{j(n-j)/2}e_j(B)$.
\item For an embedding $\tau:E\to\mathbf C$, the $\tau$-factor of $U_n$ at infinity is a compact
  unitary group $U(n)$. For integers $\lambda_0\geq\cdots\geq\lambda_{n-1}$, $U$ has \emph{weight}
  $\lambda$ at $\tau$ if there is a nonzero map to $U$ from the irreducible representation
  $\nabla(\lambda)$ of $\operatorname{GL}_n(\mathbf C)$ attached to $\lambda$ that intertwines
  $\nabla(\lambda)|_{U(n)}$ with right translation by the $\tau$-factor.
\end{enumerate}
\end{definition}
\leanref{Development/AbelianVarietyTheory/PublishedTheorems/GaloisCohomology₂/BaseChangeDefiniteUnitary.lean}{GlobalTheory.UnitaryBC.UA}
\leanref{Development/AbelianVarietyTheory/PublishedTheorems/GaloisCohomology₂/BaseChangeDefiniteUnitary.lean}{GlobalTheory.UnitaryBC.IsAutRepU}
\leanref{Development/AbelianVarietyTheory/PublishedTheorems/GaloisCohomology₂/BaseChangeDefiniteUnitary.lean}{GlobalTheory.UnitaryBC.IsCompAlong}
\leanref{Development/AbelianVarietyTheory/PublishedTheorems/GaloisCohomology₂/BaseChangeDefiniteUnitary.lean}{GlobalTheory.UnitaryBC.IsUnramU}
\leanref{Development/AbelianVarietyTheory/PublishedTheorems/GaloisCohomology₂/BaseChangeDefiniteUnitary.lean}{GlobalTheory.UnitaryBC.HasWeightAt}
\leanref{Development/AbelianVarietyTheory/Definitions/GaloisGeometry/LatticeDataThicknessBound.lean}{GlobalTheory.CongL.SatU}

\begin{theorem}\label{gal-base-change-definite}
Let $E$ be a CM field with maximal totally real subfield $F$ such that $E/F$ is unramified at every
finite place. Let $n\geq1$ and let $U$ be an automorphic representation of $U_n$
(Definition~\ref{gal-definite-unitary}). Then there are positive integers $N_1,\ldots,N_c$ with
$N_1+\cdots+N_c=n$, and for each $i$ an irreducible space $V_i$ of square-integrable automorphic
forms on $\operatorname{GL}_{N_i}(E)\backslash\operatorname{GL}_{N_i}(\mathbf A_E)$, invariant under
the positive reals embedded in the centre at infinity, with local components $\sigma_{i,w}$ and
infinitesimal characters $m_{i,\tau}$, and for each finite place $w$ of $E$ an irreducible smooth
representation $\Pi_w$ of $\operatorname{GL}_n(E_w)$, such that:
\begin{enumerate}
\item for each $i$ there are $a_i,d_i\geq1$ with $N_i=a_id_i$ and a cuspidal automorphic
  representation $P_i$ of $\operatorname{GL}_{a_i}(\mathbf A_E)$ with unitary central character,
  invariant under the positive reals at infinity and with $P_i^c\cong P_i^\vee$, such that $V_i$
  embeds equivariantly into the space of forms parabolically induced from
  $P_i|\det|^{-s_0}\otimes\cdots\otimes P_i|\det|^{-s_{d_i-1}}$, $s_j=\frac{d_i-1}{2}-j$, so
  that $V_i$ is a Speh representation in the sense of \cite{MW89};
\item $\Pi_w$ is a subquotient of $\sigma_{1,w}\times\cdots\times\sigma_{c,w}$;
\item if $U$ is unramified at $w$, then $\Pi_w$ is spherical;
\item if $w$ is inert over $F$ and the local component of $U$ at $w$ is unramified with Satake
  parameters $\alpha_1,\ldots,\alpha_{\lfloor n/2\rfloor}$ for the quasi-split unitary group, then
  $\Pi_w$ has Satake parameters
  $\{\alpha_1^{\pm1},\ldots,\alpha_{\lfloor n/2\rfloor}^{\pm1}\}$, together with $1$ when $n$ is odd;
\item if $w$ is split over $F$, then $\Pi_w$ is a local component of $U$ at $w$;
\item if $U$ has weight $\lambda$ at $\tau:E\to\mathbf C$, then
  $m_{1,\tau}+\cdots+m_{c,\tau}=\{\lambda_k+\frac{n-1}{2}-k:0\leq k<n\}$ as multisets.
\end{enumerate}
This is the form of Labesse's base change \cite{Lab11} used here.
\end{theorem}
\leanref{Development/WeilBound/RepresentationTheoryGL/ArchimedeanBaseChange/ArchimedeanBaseChangeProof/BaseChangeIsobaricNOf.lean}{PositivityTheory.Rg.BcN.IsBaseChangeN}
\leanref{Development/WeilBound/RepresentationTheoryGL/ArchimedeanBaseChange/ArchimedeanBaseChangeProof/BaseChangeIsobaricNOf.lean}{PositivityTheory.Rg.BcN.BaseChangeIsobaricN}
\leanref{Development/AutomorphicArithmetic/Foundations/BaseChangeIdentities/DefiniteNHolds.lean}{UpperTheory.BaseChangeIdentities.baseChangeN}

\begin{definition}\label{gal-good-base}
Let $E$ be a CM field with maximal totally real subfield $F$ and complex conjugation $c$.
\begin{enumerate}
\item A \emph{good base} for $E$ consists of a continuous additive character $\psi$ of
  $F\backslash\mathbf A_F$ with $\psi(x)=\exp(-2\pi i\operatorname{Tr}(x))$ on $F\otimes\mathbf R$, an
  embedding $\tau_v:E\to\mathbf C$ over each infinite place $v$ of $F$, integers $m_\alpha(v)$ and
  $m_\beta(v)$ for each infinite place $v$ of $F$, a finite set $P$ of finite places of $F$ with an
  $F$-embedding $\tau_v:E\to F_v$ for each $v\in P$, and a finite place $v_1$ of $F$ with an $F$-embedding
  $\tau_1:E\to F_{v_1}$, such that: $E$ contains $\sqrt{-d}$ for some integer $d>0$; $E/F$ is unramified
  at every finite place; $\psi_v$ has conductor $\mathcal O_v$ for every $v\in P$; and every finite place
  $v$ of $F$ inert in $E$ has $v\nmid2$, $N(v)>3$ and $\psi_v$ of conductor $\mathcal O_v$. No condition is
  imposed on the integers $m_\alpha(v)$ and $m_\beta(v)$ or on $(v_1,\tau_1)$, and they are not used in this
  chapter.
\item Let $M\in M_n(\mathbf C)$. A \emph{Cartan involution} of $M$ is a matrix $\theta\in M_n(\mathbf C)$ with
  $\theta^2=1$, ${}^t\theta M\bar\theta=M$, and $\operatorname{Re}({}^t(\theta x)M\bar x)>0$ for every nonzero
  $x\in\mathbf C^n$. For a finite place $v$ of $F$, an $F$-embedding $t:E\to F_v$ and an
  $\mathcal O_E$-lattice $L\subset E^n$, let $L_t\subset F_v^n$ be the $\mathcal O_v$-module spanned by the
  vectors $(t(x_i))_i$ for $x\in L$. A matrix $\beta\in\operatorname{GL}_n(F_v)$ is a \emph{basis of $L$ at
  $t$} if, for $u\in F_v^n$, $u\in\mathcal O_v^n$ if and only if $\beta u\in L_t$.
\item Fix a good base. For $n\geq1$, a \emph{standard Hermitian space} of rank $n$ is $E^n$ with
  $h(x,y)=\sum_ix_ic(y_i)$, an $\mathcal O_E$-lattice $L$ with
  $L=\{y:h(y,x)\in\mathcal O_E\text{ for all }x\in L\}$, a Cartan involution $\theta_v$ of the identity
  matrix for each infinite place $v$ of $F$, and a matrix $\beta_v\in\operatorname{GL}_n(F_v)$ for each
  finite place $v$ of $F$ that is a basis of $L$ at $\tau_v$ for every $v\in P$.
\item Fix a good base. For $d\geq1$ and $\delta\in E^\times$ with $c(\delta)=-\delta$ and $i\tau_v(\delta)>0$
  for every infinite place $v$, a \emph{scalar skew-Hermitian space} of rank $d$ is $E^d$ with
  $\langle x,y\rangle=\delta\sum_ix_ic(y_i)$, an $\mathcal O_E$-lattice $L$ with
  $L=\{y:\operatorname{Tr}_{E/\mathbf Q}(o\langle y,x\rangle)\in\mathbf Z\text{ for all }x\in L,\
  o\in\mathcal O_E\}$, a Cartan involution $\theta_v$ of the scalar matrix $i\tau_v(\delta)$ for each infinite
  place $v$ of $F$, and a matrix $\beta_v\in\operatorname{GL}_d(F_v)$ for each finite place $v$ of $F$ that is
  a basis of $L$ at $\tau_v$ for every $v\in P$.
\end{enumerate}
The matrices $\beta_v$ for $v\notin P$ are subject to no condition.
\end{definition}
\leanref{Development/AutomorphicArithmetic/Foundations/GaloisRepresentations/ThetaParseval/HalfIndexLattices/SplitStab.lean}{UpperTheory.TP.W.GoodK}
\leanref{Development/AbelianVarietyTheory/BlockTheory/DivisorBase/WeilTensorModels.lean}{GlobalTheory.IndefiniteRowSideFullHecke.WeilTensorModels.Base}
\leanref{Development/AbelianVarietyTheory/BlockTheory/DivisorBase/WeilTensorModels.lean}{GlobalTheory.IndefiniteRowSideFullHecke.WeilTensorModels.Spc.IsHerm}
\leanref{Development/AbelianVarietyTheory/BlockTheory/DivisorBase/WeilTensorModels.lean}{GlobalTheory.IndefiniteRowSideFullHecke.WeilTensorModels.Spc.IsSkew}

\begin{theorem}\label{gal-theta-lift-large-rank}
Let $E$ be a CM field with maximal totally real subfield $F$, complex conjugation $c$, quadratic
character $\eta$ of $\mathbf A_F^\times$ attached to $E/F$, and a good base
(Definition~\ref{gal-good-base}). Let $n\geq1$ and let $X$ be a standard Hermitian space of rank $n$.
Then there is $d_0$ such that the following holds for every $d\geq d_0$. Let $W$ be a scalar
skew-Hermitian space of rank $d$. Let $\chi_X,\chi_W$ be unitary Hecke characters of $E$ with
$\chi_X|_{\mathbf A_F^\times}=\eta^n$ and $\chi_W|_{\mathbf A_F^\times}=\eta^d$, and with
$\chi_X(z)=\prod_v(\tau_v(z_v)/|\tau_v(z_v)|)^{m_X(v)}$ and
$\chi_W(z)=\prod_v(\tau_v(z_v)/|\tau_v(z_v)|)^{m_W(v)}$ on $(E\otimes\mathbf R)^\times$, for integers
$m_X(v),m_W(v)$. Let $U$ be an automorphic representation of $U_n$ with weight $\lambda(v)$ at
$\tau_v$ for every infinite place $v$ (Definition~\ref{gal-definite-unitary}), such that
$d+m_W(v)\leq2\lambda(v)_k$ for all $v$ and $k$. Then there is an automorphic representation $U'$ of
$U_d$ such that:
\begin{enumerate}
\item for every infinite place $v$, $U'$ has weight $\lambda'(v)$ at $\tau_v$ with
  $2\lambda'(v)_k=2\lambda(v)_k-d-m_W(v)+n+m_X(v)$ for $k<n$ and $2\lambda'(v)_k=n+m_X(v)$ for
  $n\leq k<d$;
\item for every finite place $v$ of $F$, every $F$-embedding $t:E\to F_v$, with $w$ the place of $E$
  it defines, and all local components $\tau$ of $U$ and $\tau'$ of $U'$ at $w$, there is a nonzero
  map $\beta$ from pairs $(\Phi,x)$, with $\Phi$ a Schwartz--Bruhat function on
  $M_{n\times d}(F_v)$ and $x\in\tau$, to $\tau'$, linear in $\Phi$ and conjugate-linear in $x$, with
  $\beta(\omega(g,h)\Phi,\tau(g)x)=\tau'(h)\beta(\Phi,x)$ for
  $g\in\operatorname{GL}_n(F_v)$ and $h\in\operatorname{GL}_d(F_v)$, where
  \[
   \omega(g,h)\Phi(Y)=\chi_{W,w}(\det g)\,\chi_{X,w}(\det h)\,|\det g|^{-d/2}|\det h|^{-n/2}\,
   \Phi(g^{-1}Y\,{}^t\!h^{-1}).
  \]
\end{enumerate}
\end{theorem}
\leanref{Development/AutomorphicArithmetic/Foundations/GaloisRepresentations/ThetaParseval/SchurStep.lean}{UpperTheory.TP.IsLiftWt}
\leanref{Development/AutomorphicArithmetic/Foundations/GaloisRepresentations/ThetaParseval/SchurStep.lean}{UpperTheory.TP.LocPair}
\leanref{Development/AutomorphicArithmetic/Foundations/GaloisRepresentations/ThetaParseval/SchurStep.lean}{UpperTheory.TP.IsLift}
\leanref{Development/AutomorphicArithmetic/Foundations/GaloisRepresentations/ThetaParseval/ThetaKernels/GalTPDefs.lean}{UpperTheory.TP.TLift_k}
\leanref{Development/AutomorphicArithmetic/Foundations/GaloisRepresentations/ThetaParseval/ThetaKernels/TPGHolds.lean}{UpperTheory.TP.tLift_holds_k}

\begin{definition}\label{gal-seed}
Let $E$ be a CM field with maximal totally real subfield $F$, let $\Omega$ be an algebraic closure of
$E$, let $\ell$ be a prime, let $\iota:\bar{\mathbf Q}_\ell\to\mathbf C$ be a field isomorphism, and
let $a,d\geq1$.
\begin{enumerate}
\item A place $w$ of $E$ is \emph{$d$-split} if it is split over $F$ and $-d$ is a square in
  $\mathbf Q_p$ for the prime $p$ under $w$.
\item A finite set $J$ of integers is \emph{lacunary} if any two disjoint nonempty subsets of $J$
  have different averages.
\item Let $U$ be an automorphic representation of $U_a$ (Definition~\ref{gal-definite-unitary}). A
  \emph{seed representation} for $U$ is a continuous semisimple homomorphism
  $\rho:G_E\to\operatorname{GL}_N(\bar{\mathbf Q}_\ell)$, for some $N$, such that:
  \begin{enumerate}
  \item for all but finitely many $d$-split places $w$ not over $\ell$, and every $B$ such that $U$
    has Satake parameters $B$ at $w$, $\rho$ is unramified at $w$ and
    $\prod_{b\in B}\bigl(\rho(\varphi)-\iota^{-1}(q_w^{(a-1)/2}b)\bigr)=0$ for every geometric
    Frobenius element $\varphi$ at $w$;
  \item for some place $w$ over $\ell$, some valuation ring $A$ over $w$ and some
    $\tau:E\to\bar{\mathbf Q}_\ell$ extending continuously to $E_w$, the multiset
    $\mathrm{HT}_\tau(\rho|_{D_A})$ (Definition~\ref{gal-p-adic-hodge}) contains a lacunary set of
    $a$ integers.
  \end{enumerate}
\end{enumerate}
\end{definition}
\leanref{Development/AbelianVarietyTheory/Definitions/GaloisGeometry/ApartAllEmpty.lean}{GlobalTheory.CongL.UsPl}
\leanref{Development/AbelianVarietyTheory/Definitions/GaloisGeometry/DivisionOfDeterminantCharacter.lean}{GlobalTheory.RoadWGeo.Lacunary}
\leanref{Development/AbelianVarietyTheory/Definitions/GaloisGeometry/DivisionOfDeterminantCharacter.lean}{GlobalTheory.RoadWGeo.ESAt}
\leanref{Development/AutomorphicArithmetic/Foundations/GaloisRepresentations/IntegralStructures/ZetaCharacterDivisibility/InsideRefl.lean}{UpperTheory.Gal.Int.Div.HTRichEx}
\leanref{Development/AutomorphicArithmetic/Foundations/GaloisRepresentations/IntegralStructures/ZetaCharacterDivisibility/InsideRefl.lean}{UpperTheory.Gal.Int.Div.VBase}

\begin{theorem}\label{gal-seed-representations}
Let $E$ be a CM field with maximal totally real subfield $F$, let $\Omega$ be an algebraic closure of
$E$, let $\ell$ be a prime and let $\iota:\bar{\mathbf Q}_\ell\to\mathbf C$ be a field isomorphism.
Let $a\geq1$ be odd and $d\geq1$, and suppose that $-d$ is a square in $E$ and in $\mathbf Q_\ell$,
that $[F:\mathbf Q]$ is even, and that every place of $E$ not split over $F$ is prime to $\ell$ and
unramified over $F$. Let $x\in E$ with $x^2=-d$. Let $T$ be a finite set of places of $E$ not split
over $F$, let $w_0$ and $v'$ be places of $E$ over $\ell$ that lie over different places of $F$ and
over the same place of $\mathbf Q(x)$, and let $S$ be a finite set of places of $E$. Let $U$ be an
automorphic representation of $U_a$ (Definition~\ref{gal-definite-unitary}) such that:
\begin{enumerate}
\item $U$ is unramified at every place not split over $F$ outside $T$;
\item there are $m\geq1$, a compact open subgroup $J\subset\operatorname{GL}_a(E_{w_0})$ without
  torsion, a surjective homomorphism $\vartheta:J\to\mathbf Z/\ell^m$ and a nontrivial character
  $\chi$ of $\mathbf Z/\ell^m$, such that every irreducible smooth representation of
  $\operatorname{GL}_a(E_{w_0})$ with a nonzero vector on which $J$ acts by $\chi\circ\vartheta$ is
  supercuspidal, and $U$ has a nonzero vector on which $c_{w_0}^{-1}(J)$ acts by
  $\chi\circ\vartheta\circ c_{w_0}$;
\item $U$ has Satake parameters at every place split over $F$ that lies over a place of $F$ different
  from the one under $w_0$ and from those under the places of $S$;
\item $U$ has Satake parameters at $v'$;
\item there are an embedding $\tau:E\to\mathbf C$ for which $\iota^{-1}\circ\tau$ induces $v'$ and a
  weight $\lambda$ of $U$ at $\tau$ for which $\{k-\lambda_k:0\leq k<a\}$ is lacunary
  (Definition~\ref{gal-seed}).
\end{enumerate}
Then:
\begin{enumerate}
\item[(a)] there is a seed representation for $U$;
\item[(b)] for every $d$-split place $w$ not over $\ell$ at which $U$ has Satake parameters, there is a
  seed representation $\rho$ for $U$ that also satisfies at $w$ the condition of item~(3)(a) of
  Definition~\ref{gal-seed};
\item[(c)] for every place $w$ over $\ell$ that lies over the same place of $\mathbf Q(x)$ as $v'$,
  over a place of $F$ different from those under $w_0$ and $v'$, and at which $U$ has Satake
  parameters, there are seed representations $\rho$ and $\rho'$ for $U$ with the following
  properties for every valuation ring $A$ over $w$: $\rho|_{D_A}$ and $\rho'|_{D_A}$ are
  crystalline; $\mathrm{HT}_\tau(\rho|_{D_A})$ contains $\{k-\lambda_k:0\leq k<a\}$ for every
  embedding $\tau:E\to\bar{\mathbf Q}_\ell$ extending continuously to $E_w$ and every weight
  $\lambda$ of $U$ at $\iota\circ\tau$; and for every $B$ such that $U$ has Satake parameters $B$ at
  $w$, there is a geometric Frobenius element $\varphi\in W_A$ for which the characteristic
  polynomial of $\mathrm{WD}_\iota(\rho'|_{D_A})(\varphi)$ divides a power of
  $\prod_{b\in B}(X-q_w^{(a-1)/2}b)$.
\end{enumerate}
\end{theorem}
\leanref{Development/AbelianVarietyTheory/Definitions/GaloisGeometry/LatticeDataThicknessBound.lean}{GlobalTheory.CongL.FieldOK}
\leanref{Development/AbelianVarietyTheory/Definitions/GaloisGeometry/LatticeDataThicknessBound.lean}{GlobalTheory.CongL.AutOffT}
\leanref{Development/AbelianVarietyTheory/Definitions/Divisors/GenOf.lean}{GlobalTheory.CongW.TyEll}
\leanref{Development/AbelianVarietyTheory/Definitions/GaloisGeometry/LatticeDataThicknessBound.lean}{GlobalTheory.CongL.SphOffU}
\leanref{Development/AbelianVarietyTheory/Definitions/GaloisGeometry/DivisionOfDeterminantCharacter.lean}{GlobalTheory.RoadWGeo.LacWtAt}
\leanref{Development/AutomorphicArithmetic/Foundations/GaloisRepresentations/IntegralStructures/ZetaCharacterDivisibility/InsideRefl.lean}{UpperTheory.Gal.Int.Div.JumpsOcc}
\leanref{Development/AutomorphicArithmetic/Foundations/GaloisRepresentations/IntegralStructures/ZetaCharacterDivisibility/InsideRefl.lean}{UpperTheory.Gal.Int.Div.PhiKilled}
\leanref{Development/AutomorphicArithmetic/Foundations/GaloisRepresentations/IntegralStructures/GeometricFlagBases/HodgeTrAt.lean}{UpperTheory.Gal.Int.Geo.GeoOne}
\leanref{Development/AutomorphicArithmetic/Foundations/GaloisRepresentations/IntegralStructures/GeometricFlagBases/WeightOneGeometricRealization.lean}{UpperTheory.Gal.Int.Geo.geoOne_holds}

\begin{proof}
The construction is in Section~\ref{section-gal-seed-construction}; the hypotheses (1)--(4) are those
of Situation~\ref{gal-seed-situation}. For (a), Proposition~\ref{gal-packs-exist}(1) gives a pack with
Hodge data at $v'$, and Lemma~\ref{gal-pack-seed}, applied at $v=v'$ with hypothesis~(5), gives a seed
representation. For (b), Proposition~\ref{gal-packs-exist}(2) gives a pack with Hodge data at $v'$ that
is congruent at $w$, and Lemma~\ref{gal-pack-seed}(1) gives the condition at $w$. For (c),
Proposition~\ref{gal-packs-exist}(3) gives a pack with Hodge data at $v'$ and $w$, whose seed $\rho$ has
the properties of Lemma~\ref{gal-pack-seed}(2) at $w$, and a pack with Hodge data at $v'$ congruent at
$w$, whose seed $\rho'$ has the properties of Lemma~\ref{gal-pack-seed}(3) at $w$.
\end{proof}

\section{Seed representations from unitary Shimura varieties}\label{section-gal-seed-construction}

\subsection{Hecke towers, packs and \'etale cohomology}\label{subsection-gal-hecke-towers}

\begin{definition}\label{gal-hecke-tower}
Let $k$ be a field with an algebraic closure $\bar k$, let $\ell$ be a prime and let $L$ be an
algebraic closure of $\mathbf Q_\ell$.
\begin{enumerate}
\item A \emph{Hecke tower} $H$ over $k$ consists of a $k$-scheme $Z$, a set $\Gamma$, and for each
  $\gamma\in\Gamma$ a scheme $Y_\gamma$ with two $k$-morphisms $s_\gamma,t_\gamma:Y_\gamma\to Z$. It is
  \emph{smooth projective} if $Z$ is smooth and projective over $k$ and every $s_\gamma$ is finite
  \'etale of constant degree. Then each $\gamma$ defines a \emph{Hecke operator}
  $T_\gamma=t_{\gamma*}s_\gamma^*$ on the total \'etale cohomology
  $H^*(Z_{\bar k},\mathbf Q_\ell)$, commuting with $\operatorname{Gal}(\bar k/k)$.
\item For $\theta:\Gamma\to L$, the \emph{$\theta$-part} is the subspace of
  $L\otimes_{\mathbf Q_\ell}H^*(Z_{\bar k},\mathbf Q_\ell)$ on which every $T_\gamma$ acts by
  $\theta(\gamma)$. A homomorphism $\rho:\operatorname{Gal}(\bar k/k)\to\operatorname{GL}_N(L)$ is a
  \emph{piece} of $H$ at $\theta$ if, in some basis of the $\theta$-part, the Galois action is block
  upper triangular and $\rho$ is its block diagonal part.
\item For $\theta:\Gamma\to k$ and $j\geq0$, $\theta$ \emph{occurs in Hodge degree $j$} if some
  nonzero class $x\in H^b(Z,\Omega^j_{Z/k})$, for some $b$, satisfies
  $s_{\gamma*}t_\gamma^*x=\theta(\gamma)x$ in Hodge cohomology for every $\gamma$.
\end{enumerate}
\end{definition}
\leanref{Development/AbelianVarietyTheory/Definitions/Divisors/IdTower.lean}{GlobalTheory.GalV.HTower}
\leanref{Development/AutomorphicArithmetic/Foundations/GaloisRepresentations/IntegralStructures/GeometricFlagBases/IsGrRefl.lean}{UpperTheory.Gal.Int.Geo.Nice}
\leanref{Development/AutomorphicArithmetic/Foundations/GaloisRepresentations/IntegralStructures/GeometricFlagBases/IsGrRefl.lean}{UpperTheory.Gal.Int.Geo.IsMat}
\leanref{Development/AutomorphicArithmetic/Foundations/GaloisRepresentations/IntegralStructures/GeometricFlagBases/HodgeTrAt.lean}{UpperTheory.Gal.Int.Geo.IsPiece}
\leanref{Development/AutomorphicArithmetic/Foundations/GaloisRepresentations/IntegralStructures/GeometricFlagBases/IsGrRefl.lean}{UpperTheory.Gal.Int.Geo.HOcc}

\begin{definition}\label{gal-tower-reduction}
Let $E$ be a number field, $H$ a smooth projective Hecke tower over $E$
(Definition~\ref{gal-hecke-tower}), $\theta:\Gamma\to L$ a function and $w$ a finite place of $E$ with
valuation ring $\mathcal O_w$ of $E$ at $w$.
\begin{enumerate}
\item A \emph{model} of $H$ at $w$ consists of a smooth projective $\mathcal O_w$-scheme $Y$ of
  constant relative dimension and a morphism $Z\to Y$ over $\mathcal O_w$ such that $Z\to Y_E$ is
  finite \'etale and split by a finite group $\Delta$ acting on $Z$ through Hecke correspondences
  $\gamma$ with $s_\gamma$ an isomorphism and $\theta(\gamma)=1$. It is \emph{surjective} if
  $Z\to Y_E$ is surjective, and \emph{good} if every correspondence $(s_\gamma,t_\gamma)$ lies over a
  correspondence of $Y$ that is finite \'etale over $Y$.
\item Let $p$ be the prime under $w$ and $q=|k_w|=p^f$. A \emph{finite correspondence} of $Y$ is a
  scheme $W$ with two morphisms $p_W,r_W:W\to Y$ over $\mathcal O_w$ such that $p_W$ is finite, \'etale
  over $Y[1/p]$, and has the same number of geometric points over every geometric point of $Y$ at
  which $p$ is invertible. It \emph{lies over} $\gamma\in\Gamma$ \emph{with multiplicity} $m>0$ if
  there is a morphism $\kappa:Y_\gamma\to W$ with $p_W\kappa=\pi s_\gamma$ and $r_W\kappa=\pi t_\gamma$,
  where $\pi:Z\to Y$ is the structure map, such that $Y_\gamma\to Z\times_YW$ is finite \'etale of
  degree $m$.
\item A model satisfies the \emph{congruence relation of degree $n$ with coefficients
  $\vartheta_0,\ldots,\vartheta_n\in L$} if there are a uniformizer $\varpi$ of $\mathcal O_w$ and
  finite correspondences $T_j=(W_j,p_j,r_j)$ of $Y$, $0\leq j\leq n$, with $p_j$ and $r_j$ over the
  same point of $\operatorname{Spec}\mathcal O_w$, lying over Hecke correspondences $\gamma_j$ with
  multiplicities $m_j$ and $\theta(\gamma_j)=m_j\vartheta_j$, such that the \emph{Eichler--Shimura
  relation} holds at the generic points of the special fibre: for every valuation ring $V$ with
  algebraically closed fraction field in which $p\neq0$, every morphism $y:\operatorname{Spec}V\to Y$
  that maps the closed point to a generic point of the special fibre $Y_{k_w}$, and the $q$-power
  Frobenius $\phi$ of the residue field $\kappa_V$ of $V$, the two multisets of $\kappa_V$-points of
  $Y$
  \[
   \sum_{j\text{ even}}q^{j(j-1)/2}\bigl\{\phi^{\,n-j}(\bar r_j(z))\bigr\}_z,\qquad
   \sum_{j\text{ odd}}q^{j(j-1)/2}\bigl\{\phi^{\,n-j}(\bar r_j(z))\bigr\}_z
  \]
  are equal, where $z$ runs over the $V$-points of $W_j$ over $y$, $\bar r_j(z)$ is the reduction of
  $r_j\circ z$, and $\phi$ acts on $\kappa_V$-points by composition. The \emph{Hecke polynomial} of the
  coefficients is $P_\vartheta(X)=\sum_{j=0}^n(-1)^jq^{j(j-1)/2}\vartheta_jX^{n-j}$.
\end{enumerate}
\end{definition}
\leanref{Development/AutomorphicArithmetic/Foundations/GaloisRepresentations/IntegralStructures/GeometricFlagBases/IsGrRefl.lean}{UpperTheory.Gal.Int.Geo.Cover}
\leanref{Development/AutomorphicArithmetic/Foundations/GaloisRepresentations/IntegralStructures/GeometricFlagBases/IsGrRefl.lean}{UpperTheory.Gal.Int.Geo.GoodAtW}
\leanref{Development/AutomorphicArithmetic/Foundations/GaloisRepresentations/IntegralStructures/GeometricFlagBases/IsGrRefl.lean}{UpperTheory.Gal.Int.Geo.CongAtW}
\leanref{Development/AutomorphicArithmetic/Foundations/GaloisRepresentations/AnnihilatorModels/GenericMulOfProd.lean}{UpperTheory.Gal.crystallineComparison.heckePol}

\begin{situation}\label{gal-seed-situation}
Let $E$ be a CM field with maximal totally real subfield $F$ and complex conjugation $c$, let $\Omega$
be an algebraic closure of $E$, let $\ell$ be a prime and let $\iota:\bar{\mathbf Q}_\ell\to\mathbf C$
be a field isomorphism. Let $a\geq1$ be odd and $d\geq1$, and suppose that $-d$ is a square in $E$
and in $\mathbf Q_\ell$, that $[F:\mathbf Q]$ is even, and that every place of $E$ not split over $F$ is
prime to $\ell$ and unramified over $F$. Let $x\in E$ with $x^2=-d$. Let $T$ be a finite set of places
of $E$ not split over $F$, let $w_0$ and $v'$ be places of $E$ over $\ell$ that lie over different
places of $F$ and over the same place of $\mathbf Q(x)$, and let $S$ be a finite set of places of $E$.
Let $U$ be an automorphic representation of $U_a$ (Definition~\ref{gal-definite-unitary}) that
satisfies hypotheses (1)--(4) of Theorem~\ref{gal-seed-representations}.
\end{situation}
\leanref{Development/AutomorphicArithmetic/Foundations/GaloisRepresentations/IntegralStructures/GeometricFlagBases/HodgeTrAt.lean}{UpperTheory.Gal.Int.Geo.OneK}
\leanref{Development/AbelianVarietyTheory/Definitions/GaloisGeometry/LatticeDataThicknessBound.lean}{GlobalTheory.CongL.FieldOK}

\begin{definition}\label{gal-pack}
In Situation~\ref{gal-seed-situation}, a \emph{pack} for $U$ consists of a smooth projective Hecke
tower $H$ over $E$, a function $\theta:\Gamma\to\bar{\mathbf Q}_\ell$, a Hecke character $\kappa$ of
$E$ that is algebraic of infinity type $m=(m_\tau)_{\tau:E\to\mathbf C}$ (so that
$\kappa(z)=\prod_\tau\tau(z)^{m_\tau}$ on totally positive units $z$), and a finite set $X$ of places,
such that for every $d$-split place $w\notin X$ not over $\ell$ (Definition~\ref{gal-seed}) and every
$B$ such that $U$ has Satake parameters $B$ at $w$, $\kappa$ is unramified at $w$ and $H$ has a model
at $w$ satisfying a congruence relation of degree $a$ whose Hecke polynomial is
$\vartheta_0\prod_{b\in B}(X-\iota^{-1}(q_w^{(a-1)/2}b\,\kappa(\varpi_w)^{-1}))$ with
$\vartheta_0\neq0$ (Definition~\ref{gal-tower-reduction}). For such a pack and a place $w$:
\begin{enumerate}
\item the pack is \emph{congruent at $w$} if the last condition holds at $w$ for every $B$ such that
  $U$ has Satake parameters $B$ at $w$;
\item the pack has \emph{Hodge data at $w$} if $\kappa$ is unramified at $w$, $H$ has a good model at
  $w$, and for every embedding $\tau:E\to\mathbf C$ such that $\iota^{-1}\circ\tau$ induces $w$ there
  is an automorphism $\sigma$ of $\mathbf C$ such that for every weight $\lambda$ of $U$ at $\tau$ and
  every $k<a$, $\iota\circ\theta$ composed with $\sigma$ occurs in Hodge degree $k-\lambda_k+m_\tau$ of
  the base change of $H$ along $\sigma\circ\tau$ (Definition~\ref{gal-hecke-tower}).
\end{enumerate}
\end{definition}
\leanref{Development/AutomorphicArithmetic/Foundations/GaloisRepresentations/IntegralStructures/GeometricFlagBases/HodgeTrAt.lean}{UpperTheory.Gal.Int.Geo.Pack}
\leanref{Development/AutomorphicArithmetic/Foundations/GaloisRepresentations/IntegralStructures/GeometricFlagBases/HodgeTrAt.lean}{UpperTheory.Gal.Int.Geo.CongFor}
\leanref{Development/AutomorphicArithmetic/Foundations/GaloisRepresentations/IntegralStructures/GeometricFlagBases/HodgeTrAt.lean}{UpperTheory.Gal.Int.Geo.Pack.HTFor}
\leanref{Development/AutomorphicArithmetic/Foundations/GaloisRepresentations/IntegralStructures/GeometricFlagBases/HodgeTrAt.lean}{UpperTheory.Gal.Int.Geo.Pack.CongAll}

\begin{theorem}\label{gal-etale-foundations}
Let $\ell$ be a prime.
\begin{enumerate}
\item Let $k$ be an algebraically closed field in which $\ell$ is invertible and $X$ a smooth
  projective $k$-scheme. Then every $H^n(X,\mathbf Q_\ell)$ is finite-dimensional, and there is $N$
  with $H^n(X,\mathbf Q_\ell)=0$ for $n<0$ and for $n>N$.
\item Let $\mathcal O$ be a discrete valuation ring with fraction field $K$ and residue field $\kappa$
  in which $\ell$ is invertible, let $Y\to\operatorname{Spec}\mathcal O$ be smooth and proper, let
  $\Omega$ be an algebraically closed field containing $K$ and $V\subset\Omega$ a valuation ring over
  $\mathcal O$ with residue field $\bar\kappa$. Then for every $n$ the specialization map
  $H^n(Y_{\bar\kappa},\mathbf Q_\ell)\to H^n(Y_\Omega,\mathbf Q_\ell)$ is an isomorphism, compatible
  with pullback along morphisms of smooth proper $\mathcal O$-schemes and equivariant for every
  $\sigma\in\operatorname{Aut}(\Omega/K)$ with $\sigma(V)=V$, which acts on $Y_{\bar\kappa}$ through
  its action on $\bar\kappa$. In particular such a $\sigma$ that acts trivially on $\bar\kappa$ acts
  trivially on $H^n(Y_\Omega,\mathbf Q_\ell)$.
\end{enumerate}
These are the finiteness theorem and the smooth and proper base change theorems of \'etale cohomology
\cite[Expos\'es XIV, XVI]{AGV73}, \cite[Th.~finitude]{Del77}.
\end{theorem}
\leanref{Development/AutomorphicArithmetic/Foundations/GaloisRepresentations/IntegralStructures/GeometricFlagBases/FlagBasis.lean}{UpperTheory.Gal.Int.Geo.DegFin}
\leanref{Development/AutomorphicArithmetic/Foundations/GaloisRepresentations/KunnethLefschetz/EtSideGood.lean}{UpperTheory.Gal.KunnethLefschetz.Core}
\leanref{Development/AutomorphicArithmetic/Foundations/GaloisRepresentations/KunnethLefschetz/CoreHolds.lean}{UpperTheory.Gal.KunnethLefschetz.core_holds}
\leanref{Development/AutomorphicArithmetic/Foundations/GaloisRepresentations/BasicOpenSections/EsGenAll.lean}{UpperTheory.Gal.Mod.specGood}

\begin{theorem}\label{gal-padic-comparison}
Let $\ell$ be a prime, $K$ a finite extension of $\mathbf Q_\ell$ inside an algebraic closure
$\bar K$, $\mathcal O_K$ its valuation ring with residue field $k$, and $Y\to\operatorname{Spec}
\mathcal O_K$ a smooth projective morphism. For $n\geq0$ let $V_n=H^n(Y_{\bar K},\mathbf Q_\ell)$ with
its action of $\operatorname{Gal}(\bar K/K)$.
\begin{enumerate}
\item $V_n$ is crystalline \cite{Fal89}.
\item \emph{Hodge--Tate decomposition.} There is a $\operatorname{Gal}(\bar K/K)$-equivariant
  isomorphism $\mathbf C_K\otimes_{\mathbf Q_\ell}V_n\cong\bigoplus_j\mathbf C_K(-j)\otimes_K
  H^{n-j}(Y_K,\Omega^j_{Y_K/K})$, natural in $Y$ for morphisms of smooth projective
  $\mathcal O_K$-schemes \cite{Fal88}.
\item The crystalline comparison identifies $D_{\mathrm{cris}}(V_n)$ with the rational crystalline
  cohomology $H^n_{\mathrm{cris}}(Y_k/W(k))[1/\ell]$, compatibly with Frobenius and with pullback along
  morphisms of smooth projective $\mathcal O_K$-schemes \cite{Fal89}.
\end{enumerate}
\end{theorem}
\leanref{Development/WeilBound/CoreStatements/PadicHodgeComparison/EtaleCohomologyComparison/SemistableModelSemistable.lean}{CrystallineSemistableComparison.SmoothModelCrystalline}
\leanref{Development/WeilBound/CoreStatements/PadicHodgeComparison/EtaleCohomologyComparison/BocksteinBaseChange/CstSmoothOf.lean}{CrystallineSemistableComparison.CstSmooth}
\leanref{Development/WeilBound/CoreStatements/PadicHodgeComparison/EtaleCohomologyComparison/HodgeTateResidues/HodgeTateSpectral/HTSpecGood.lean}{CrystallineSemistableComparison.HTSpecB.HTSpecGood}
\leanref{Development/WeilBound/CoreStatements/PadicHodgeComparison/EtaleCohomologyComparison/BocksteinBaseChange/StmtPhi.lean}{CrystallineSemistableComparison.CrysPhi}
\leanref{Development/AutomorphicArithmetic/Foundations/GaloisRepresentations/CrystallineTowers/Closed.lean}{UpperTheory.Gal.Seam.cstSmooth}

\begin{lemma}\label{gal-weighted-trace}
Let $p:Z\to S$ be a finite morphism of schemes and $q\geq1$ an integer invertible on $S$. A
\emph{weight} for $p$ is a function $\mu$ on the geometric points of $Z$ with values in $\mathbf Z$
such that for every morphism $S'\to S$, with $p':Z'=Z\times_SS'\to S'$, and every locally constant
function $\phi$ on $Z'$, the function $x\mapsto\sum_{p'(z)=x}\mu(z)\phi(z)$ on the geometric points of
$S'$ comes from a locally constant function on $S'$.
\begin{enumerate}
\item For every weight $\mu$ there are homomorphisms $\operatorname{tr}_{p,\mu}:H^n(Z,\mathbf
  Z/q(i))\to H^n(S,\mathbf Z/q(i))$, induced by a map of sheaves whose effect on the sections over
  every $S'\to S$ is the sum over the geometric points of the fibres of $p'$ weighted by $\mu$.
\item $\operatorname{tr}_{p,\mu}$ is additive in $\mu$, $\operatorname{tr}_{p,d\mu}=d\operatorname{tr}_{p,\mu}$
  for $d\in\mathbf Z$, and $\operatorname{tr}_{p,0}=0$.
\item Let $h:Z_1\to Z$ be finite and $\mu$ a weight for $p\circ h$. Then
  $h_*\mu(z)=\sum_{h(z_1)=z}\mu(z_1)$ is a weight for $p$, and
  $\operatorname{tr}_{p\circ h,\mu}\circ h^*=\operatorname{tr}_{p,h_*\mu}$.
\item A weight for $p$ that vanishes at every geometric point of $Z$ lying over a point of a dense
  subset of $S$ is zero.
\item For $q=\ell^k$ the maps of (1) are compatible with reduction from $\ell^{k+1}$ to $\ell^k$, and
  so define $\operatorname{tr}_{p,\mu}$ on $H^n(Z,\mathbf Z_\ell(i))$ and $H^n(Z,\mathbf
  Q_\ell(i))$, with the properties (2) and~(3).
\end{enumerate}
When $p$ is flat and of finite presentation and $\mu$ is the multiplicity, (1) is the trace morphism
of \cite[Expos\'e~XVII, 6.2]{AGV73}.
\end{lemma}
\leanref{Development/AutomorphicArithmetic/Foundations/GaloisRepresentations/TranscendenceTraces/IsWtZero.lean}{UpperTheory.Gal.TranscendenceTraces.IsWt}
\leanref{Development/AutomorphicArithmetic/Foundations/GaloisRepresentations/TranscendenceTraces/IsWtZero.lean}{UpperTheory.Gal.TranscendenceTraces.tr}
\leanref{Development/AutomorphicArithmetic/Foundations/GaloisRepresentations/TranscendenceTraces/WtDegConst.lean}{UpperTheory.Gal.TranscendenceTraces.TrSum}
\leanref{Development/AutomorphicArithmetic/Foundations/GaloisRepresentations/TranscendenceTraces/WtDegConst.lean}{UpperTheory.Gal.TranscendenceTraces.TrLin}
\leanref{Development/AutomorphicArithmetic/Foundations/GaloisRepresentations/TranscendenceTraces/WtDegConst.lean}{UpperTheory.Gal.TranscendenceTraces.WtPush}
\leanref{Development/AutomorphicArithmetic/Foundations/GaloisRepresentations/TranscendenceTraces/WtDegConst.lean}{UpperTheory.Gal.TranscendenceTraces.TrPush}
\leanref{Development/AutomorphicArithmetic/Foundations/GaloisRepresentations/TranscendenceTraces/IsolateSpec.lean}{UpperTheory.Gal.TranscendenceTraces.WtVanish}
\leanref{Development/AutomorphicArithmetic/Foundations/GaloisRepresentations/TranscendenceTraces/WtDegConst.lean}{UpperTheory.Gal.TranscendenceTraces.TrRed}
\leanref{Development/AutomorphicArithmetic/Foundations/GaloisRepresentations/TranscendenceTraces/WtDegConst.lean}{UpperTheory.Gal.TranscendenceTraces.TrHl}
\leanref{Development/AutomorphicArithmetic/Foundations/GaloisRepresentations/TranscendenceTraces/IsolateSpecHolds.lean}{UpperTheory.Gal.TranscendenceTraces.wtVanish_holds}

\begin{lemma}\label{gal-weighted-correspondences}
Let $\ell$ be a prime, let $S$ and $Y$ be schemes over a scheme $B$, with $Y$ separated and locally
of finite type over $B$ and $\ell$ invertible on $S$. A \emph{weighted finite correspondence} from $S$
to $Y$ is a finite morphism $p:W\to S$, a morphism $f:W\to Y$ with $W\to Y\to B$ equal to
$W\to S\to B$, and a weight $\mu$ for $p$ (Lemma~\ref{gal-weighted-trace}); it acts on cohomology by
$c^*=\operatorname{tr}_{p,\mu}\circ f^*:H^n(Y,\mathbf Q_\ell(i))\to H^n(S,\mathbf Q_\ell(i))$. Let
$c_1,\ldots,c_k$ be weighted finite correspondences from $S$ to $Y$, let $m_1,\ldots,m_k\in\mathbf Z$,
and let $D\subset S$ be dense. Suppose that for every geometric point $x$ of $S$ lying over a point of
$D$ and every geometric point $y$ of $Y$, $\sum_am_a\sum_z\mu_a(z)=0$, where $z$ runs over the
geometric points of $W_a$ over $x$ with $f_a(z)=y$. Then $\sum_am_ac_a^*=0$ on every
$H^n(S,\mathbf Q_\ell)$.
\end{lemma}
\leanref{Development/AutomorphicArithmetic/Foundations/GaloisRepresentations/TranscendenceTraces/IsWtZero.lean}{UpperTheory.Gal.TranscendenceTraces.Corr}
\leanref{Development/AutomorphicArithmetic/Foundations/GaloisRepresentations/TranscendenceTraces/IsWtZero.lean}{UpperTheory.Gal.TranscendenceTraces.Corr.opQ}
\leanref{Development/AutomorphicArithmetic/Foundations/GaloisRepresentations/TranscendenceTraces/WtDegConst.lean}{UpperTheory.Gal.TranscendenceTraces.CorrEqQ}
\leanref{Development/AutomorphicArithmetic/Foundations/GaloisRepresentations/TranscendenceTraces/AllHolds.lean}{UpperTheory.Gal.TranscendenceTraces.all_holds}

\begin{proof}
Write $c_a=(p_a:W_a\to S,f_a,\mu_a)$. We first work with coefficients $\mathbf Z/\ell^k(i)$. Let
$G\subset S\times_BY$ be the scheme-theoretic image of $\coprod_aW_a$, with projections
$p_G:G\to S$ and $f_G:G\to Y$, and let $h_a:W_a\to G$ be the induced morphisms. Since each $W_a$ is
finite over $S$ and $Y$ is separated and locally of finite type over $B$, $p_G$ and the $h_a$ are
finite. Moreover $p_a=p_Gh_a$, $f_a=f_Gh_a$, and $G\to S\times_BY$ is a monomorphism, so a geometric
point of $G$ is determined by its images in $S$ and in $Y$.

By Lemma~\ref{gal-weighted-trace}(3), $\nu_a=h_{a*}\mu_a$ is a weight for $p_G$ and
$c_a^*=\operatorname{tr}_{p_a,\mu_a}\circ h_a^*\circ f_G^*=\operatorname{tr}_{p_G,\nu_a}\circ f_G^*$.
By Lemma~\ref{gal-weighted-trace}(2), $\sum_am_ac_a^*=\operatorname{tr}_{p_G,\nu}\circ f_G^*$ with
$\nu=\sum_am_a\nu_a$, which is again a weight for $p_G$. Let $g$ be a geometric point of $G$, with
image $x$ in $S$ and $y=f_G(g)$ in $Y$. The geometric points of $W_a$ over $g$ are the points $z$ of
$W_a$ over $x$ with $f_a(z)=y$, since $G$ embeds in $S\times_BY$. Hence
$\nu(g)=\sum_am_a\sum_z\mu_a(z)$, which vanishes by hypothesis when $x$ lies over a point of $D$. By
Lemma~\ref{gal-weighted-trace}(4), $\nu=0$, and by Lemma~\ref{gal-weighted-trace}(2) the operator
$\sum_am_ac_a^*$ vanishes with coefficients $\mathbf Z/\ell^k(i)$. This holds for every $k$, so by
Lemma~\ref{gal-weighted-trace}(5) it holds with coefficients $\mathbf Z_\ell(i)$, and hence with
$\mathbf Q_\ell(i)$.
\end{proof}
\leanref{Development/AutomorphicArithmetic/Foundations/GaloisRepresentations/TranscendenceTraces/PCorrEq.lean}{UpperTheory.Gal.TranscendenceTraces.corrEq_of}
\leanref{Development/AutomorphicArithmetic/Foundations/GaloisRepresentations/TranscendenceTraces/PCorrEqQ.lean}{UpperTheory.Gal.TranscendenceTraces.corrEqQ_of}
\leanref{Development/AutomorphicArithmetic/Foundations/GaloisRepresentations/TranscendenceTraces/PTrRel.lean}{UpperTheory.Gal.TranscendenceTraces.trRel_of}
\leanref{Development/AutomorphicArithmetic/Foundations/GaloisRepresentations/TranscendenceTraces/IsolateSpec.lean}{UpperTheory.Gal.TranscendenceTraces.ImageFin}

\begin{lemma}\label{gal-eichler-shimura-cohomology}
Let $\mathcal O$ be a discrete valuation ring with uniformizer $\varpi$, fraction field $K$ and residue
field $\kappa$ with $q=p^f$ elements, let $\ell\neq p$ be a prime, and let $Y\to\operatorname{Spec}
\mathcal O$ be smooth and proper. Let $T_j=(W_j,p_j,r_j)$, $0\leq j\leq n$, be finite correspondences
of $Y$ over $\mathcal O$ that satisfy the Eichler--Shimura relation at the generic points of the special
fibre (Definition~\ref{gal-tower-reduction}). Let $\Omega$ be an algebraically closed field containing
$K$, $V\subset\Omega$ a valuation ring over $\mathcal O$, and $t_j=\operatorname{tr}_{p_j}\circ r_j^*$
the endomorphism of $H^k(Y_\Omega,\mathbf Q_\ell)$ defined by $T_j$. Let $\sigma\in\operatorname{Aut}
(\Omega/K)$ with $\sigma(V)=V$.
\begin{enumerate}
\item If $\sigma$ acts trivially on the residue field of $V$, then $\sigma$ acts trivially on
  $H^k(Y_\Omega,\mathbf Q_\ell)$.
\item If $\sigma$ acts on the residue field of $V$ by $x\mapsto x^q$, then
  $\sum_{j=0}^n(-1)^jq^{j(j-1)/2}\,t_j\circ(\sigma^{-1})^{n-j}=0$ on $H^k(Y_\Omega,\mathbf Q_\ell)$.
\end{enumerate}
\end{lemma}
\leanref{Development/AutomorphicArithmetic/Foundations/GaloisRepresentations/BasicOpenSections/BaseFacts.lean}{UpperTheory.Gal.Mod.ESGen}
\leanref{Development/AutomorphicArithmetic/Foundations/GaloisRepresentations/BasicOpenSections/ResClosed.lean}{UpperTheory.Gal.Mod.CycAll}
\leanref{Development/AutomorphicArithmetic/Foundations/GaloisRepresentations/BasicOpenSections/ResClosed.lean}{UpperTheory.Gal.Mod.ESPt}
\leanref{Development/AutomorphicArithmetic/Foundations/GaloisRepresentations/BasicOpenSections/PullHQComp.lean}{UpperTheory.Gal.Mod.esPt}
\leanref{Development/AutomorphicArithmetic/Foundations/GaloisRepresentations/BasicOpenSections/EsGenAll.lean}{UpperTheory.Gal.Mod.esGen_all}

\begin{proof}
Item (1) is Theorem~\ref{gal-etale-foundations}(2). For (2), let $\bar\kappa$ be the residue field of
$V$ and $F$ the $q$-power Frobenius endomorphism of $Y_{\bar\kappa}$ over $\bar\kappa$. By
Theorem~\ref{gal-etale-foundations}(2), the specialization isomorphism carries $\sigma^{-1}$ to $F^*$,
and it carries $t_j$ to the endomorphism $\bar t_j$ defined by the special fibre of $T_j$, since the
specialization map is compatible with pullback and with the traces of finite morphisms. So it suffices
to show that $\sum_j(-1)^jq^{j(j-1)/2}\,\bar t_j\circ(F^*)^{n-j}=0$ on $H^k(Y_{\bar\kappa},\mathbf
Q_\ell)$, and $\bar t_j\circ(F^*)^{n-j}$ is the action of the correspondence
$(W_{j,\bar\kappa},p_j,F^{n-j}\circ r_j)$.

Give the geometric points of $W_{j,\bar\kappa}$ the weight $\mu_j(z_0)$ equal to the number of points
$z$ of $W_j$ over a valuation-ring point $y$ of $Y$ whose reduction is $z_0$; at points over
$Y[1/p]$ this is $1$, and on the special fibre it is a weight in the sense of
Lemma~\ref{gal-weighted-correspondences} because $Y$ is normal. Every geometric point $x$ of
$Y_{\bar\kappa}$ over a generic point of the special fibre is the reduction of a point
$y:\operatorname{Spec}V'\to Y$ for a valuation ring $V'$ with algebraically closed fraction field,
and the Eichler--Shimura relation at $y$ says that the weighted $0$-cycles
$\sum_j(-1)^jq^{j(j-1)/2}\sum_z\mu_j(z)[F^{n-j}r_j(z)]$, $z$ over $x$, vanish. These points are dense
in $Y_{\bar\kappa}$, and Lemma~\ref{gal-weighted-correspondences} gives the identity of operators.
\end{proof}

\subsection{Crystalline test algebras}\label{subsection-gal-test-algebras}

\begin{definition}\label{gal-crystalline-test-algebras}
Let $\ell$ be a prime, $C$ the completion of an algebraic closure of $\mathbf Q_\ell$,
$\mathcal O_C^\flat=\varprojlim_{x\mapsto x^\ell}\mathcal O_C/\ell$ the tilt of its valuation ring,
$\ell^\flat=(\ell,\ell^{1/\ell},\ell^{1/\ell^2},\ldots)\in\mathcal O_C^\flat$ for a compatible system of
roots, and $A_{\mathrm{cris}}$ Fontaine's ring with its Frobenius $\phi$ \cite{Fon94}.
\begin{enumerate}
\item A \emph{test algebra} is a perfect $\mathcal O_C^\flat$-algebra $A$, complete for the
  $\ell^\flat$-adic topology, in which $\ell^\flat$ is a nonzerodivisor and every monic polynomial of
  positive degree has a root, and such that for $a,b\in A$ with $(\ell^\flat)^M\in(a,b)$ for some $M$
  there is an idempotent $e\in A$ with $ea\in bA$ and $(1-e)b\in aA$. Its \emph{untilt} is
  $A^\sharp=W(A)/\xi W(A)$, $\xi=\ell-[\ell^\flat]$, with the quotient map $\theta:W(A)\to A^\sharp$ and
  the map $\mathcal O_C=W(\mathcal O_C^\flat)/\xi\to A^\sharp$; then $A^\sharp/\ell=A/\ell^\flat$, with
  $\theta(x)\mapsto x_0$ for the first Witt component $x_0$ of $x$. Its \emph{crystalline ring}
  $A_{\mathrm{cris}}(A)$ is the $W(A)$-subalgebra of $W(A)[1/\ell]$ generated by the $\xi^i/i!$, with the
  Frobenius $\phi$ of $W(A)$; each $A_{\mathrm{cris}}(A)/\ell^\nu$ is an $A_{\mathrm{cris}}$-algebra. A
  \emph{decomposition} of $A$ is a finite family of maps of test algebras $A\to A_s$ such that
  $A\to\prod_sA_s$ is bijective.
\item For a functor $P$ from test algebras to sets, let $\mathrm{El}(P)$ be the category of pairs
  $(A,y)$ with $y\in P(A)$. The $A_{\mathrm{cris}}$-module $H^k(P,A_{\mathrm{cris}}/\ell^\nu)$ is the
  $k$-th cohomology of the complex whose $k$-cochains assign to every chain
  $(A_0,y_0)\to\cdots\to(A_k,y_k)$ in $\mathrm{El}(P)$ an element of $A_{\mathrm{cris}}(A_k)/\ell^\nu$,
  with the alternating sum of the face maps as differential. A natural transformation $u:Q\to P$
  induces $u^*:H^k(P,A_{\mathrm{cris}}/\ell^\nu)\to H^k(Q,A_{\mathrm{cris}}/\ell^\nu)$, and the Frobenius
  of the rings $A_{\mathrm{cris}}(A)$ induces a $\phi$-semilinear endomorphism $\phi$ of
  $H^k(P,A_{\mathrm{cris}}/\ell^\nu)$.
\item Let $u_i:Q\to P$, $i\in I$, be finitely many natural transformations and $a_i\in\mathbf Z$. The
  sum $\sum_ia_iu_i$ \emph{vanishes locally} if for every test algebra $A$ and $y\in Q(A)$ there is a
  decomposition $(A\to A_s)$ of $A$ such that $\sum_ia_i[u_i(y_s)]=0$ in the free abelian group on
  $P(A_s)$ for every $s$ with $A_s\neq0$, where $y_s$ is the image of $y$ in $Q(A_s)$.
\item Let $\mathcal O$ be a ring, $X$ an $\mathcal O$-scheme, and $\kappa:\mathcal O\to\mathcal O_C$ and
  $g:\mathcal O\to\mathcal O_C/\ell=\mathcal O_C^\flat/\ell^\flat$ ring maps. The functor $X^\sharp$ sends
  $A$ to the set of $\mathcal O$-morphisms $\operatorname{Spec}A^\sharp\to X$, with $A^\sharp$ an
  $\mathcal O$-algebra through $\kappa$, and $X_g$ sends $A$ to the set of $\mathcal O$-morphisms
  $\operatorname{Spec}A/\ell^\flat\to X$, with $A/\ell^\flat$ an $\mathcal O$-algebra through $g$. An
  $\mathcal O$-morphism $u:X\to X'$ gives $u_*:X^\sharp\to X'^\sharp$ by composition. If $f\geq1$ and
  $g(x)^{\ell^f}=g(x)$ for all $x$, the $\ell^f$-power maps of the rings $A/\ell^\flat$ give a natural
  transformation $F:X_g\to X_g$. If $m\geq0$ and $(\kappa(x)\bmod\ell)^{\ell^m}=g(x)$ for all
  $x\in\mathcal O$, the maps $\psi_m:A^\sharp\to A/\ell^\flat$, $\theta(x)\mapsto x_0^{\ell^m}$, give a
  natural transformation $\rho_m:X^\sharp\to X_g$, $y\mapsto y\circ\operatorname{Spec}\psi_m$.
\end{enumerate}
\end{definition}
\leanref{Development/WeilBound/CoreStatements/PadicHodgeComparison/EtaleCohomologyComparison/MapWXi.lean}{CrystallineSemistableComparison.PAlg}
\leanref{Development/WeilBound/CoreStatements/PadicHodgeComparison/EtaleCohomologyComparison/MapWXi.lean}{CrystallineSemistableComparison.PAlg.IsTest}
\leanref{Development/WeilBound/CoreStatements/PadicHodgeComparison/EtaleCohomologyComparison/MapWXi.lean}{CrystallineSemistableComparison.PAlg.Acr}
\leanref{Development/WeilBound/CoreStatements/PadicHodgeComparison/EtaleCohomologyComparison/AcrisKummerUnits/GaloisTwistIdentitySection.lean}{CrystallineSemistableComparison.MH}
\leanref{Development/WeilBound/CoreStatements/PadicHodgeComparison/EtaleCohomologyComparison/AcrisKummerUnits/GaloisTwistIdentitySection.lean}{CrystallineSemistableComparison.PX}
\leanref{Development/AutomorphicArithmetic/CommutativeAlgebra/EtaleCohomology/TestFacts.lean}{UpperTheory.Ru.crystallineComparison.IsDec}
\leanref{Development/AutomorphicArithmetic/CommutativeAlgebra/EtaleCohomology/TestFacts.lean}{UpperTheory.Ru.crystallineComparison.TZero}
\leanref{Development/AutomorphicArithmetic/CommutativeAlgebra/EtaleCohomology/TestFacts.lean}{UpperTheory.Ru.crystallineComparison.PQ}
\leanref{Development/AutomorphicArithmetic/CommutativeAlgebra/EtaleCohomology/TestFacts.lean}{UpperTheory.Ru.crystallineComparison.IsFr}
\leanref{Development/AutomorphicArithmetic/CommutativeAlgebra/EtaleCohomology/TestFacts.lean}{UpperTheory.Ru.crystallineComparison.IsRd}
\leanref{Development/AutomorphicArithmetic/CommutativeAlgebra/EtaleCohomology/TestFacts.lean}{UpperTheory.Ru.crystallineComparison.phiH}
\leanref{Development/AutomorphicArithmetic/CommutativeAlgebra/EtaleCohomology/TestFacts.lean}{UpperTheory.Ru.crystallineComparison.pullH}

\begin{definition}\label{gal-passage-ordered-rings}
Let $S$ be a ring and $\varpi\in S$.
\begin{enumerate}
\item $S$ is \emph{absolutely integrally closed} if every monic polynomial of positive degree in $S[X]$
  has a root in $S$.
\item An idempotent $e\in S$ \emph{orders} a pair $f,g\in S$ if $ef\in gS$ and $(1-e)g\in fS$. The ring
  $S$ is \emph{ordered} if every pair of elements of $S$ is ordered by an idempotent, and
  \emph{$\varpi$-ordered} if every pair $f,g\in S$ with $\varpi^N\in fS+gS$ for some $N\geq0$ is ordered
  by an idempotent.
\item The element $\varpi$ is \emph{chained} in $S$ if for every pair $f,g\in S$ with
  $\varpi^N\in fS+gS$ for some $N\geq0$ one has $f\in gS$ or $g\in fS$.
\item $S$ is \emph{componentwise local} if for every $f\in S$ there is an idempotent $e\in fS$ with
  $1-e\in(1-f)S$.
\item For a prime ideal $P$ of $S$, $I_P$ is the ideal of $S$ generated by the idempotents that lie in
  $P$.
\item The element $\varpi$ \emph{admits Newton steps} in $S$ if for every monic $Q\in S[X]$ of degree
  $d\geq1$ and every $n\geq0$ with $\varpi^n\mid Q(0)$ there is $z\in S$ with $\varpi^{n+1}\mid Q(z)$
  and $\varpi^n\mid z^d$.
\end{enumerate}
\end{definition}
\leanref{Development/WeilBound/CoreStatements/PadicHodgeComparison/EtaleCohomologyComparison/AcrisKummerUnits/Splits.lean}{CrystallineSemistableComparison.R.Aic}
\leanref{Development/WeilBound/CoreStatements/PadicHodgeComparison/EtaleCohomologyComparison/AcrisKummerUnits/Splits.lean}{CrystallineSemistableComparison.R.Val0}
\leanref{Development/WeilBound/CoreStatements/PadicHodgeComparison/EtaleCohomologyComparison/AcrisKummerUnits/Splits.lean}{CrystallineSemistableComparison.R.Val}
\leanref{Development/WeilBound/CoreStatements/PadicHodgeComparison/EtaleCohomologyComparison/AcrisKummerUnits/Splits.lean}{CrystallineSemistableComparison.R.Dich}
\leanref{Development/WeilBound/CoreStatements/PadicHodgeComparison/EtaleCohomologyComparison/AcrisKummerUnits/Splits.lean}{CrystallineSemistableComparison.R.Loc}
\leanref{Development/WeilBound/CoreStatements/PadicHodgeComparison/EtaleCohomologyComparison/AcrisKummerUnits/Splits.lean}{CrystallineSemistableComparison.R.compI}
\leanref{Development/WeilBound/CoreStatements/PadicHodgeComparison/EtaleCohomologyComparison/AcrisKummerUnits/Splits.lean}{CrystallineSemistableComparison.R.Step}

\begin{lemma}\label{gal-passage-ordered-basics}
Let $S$ be a ring and $\varpi\in S$ (Definition~\ref{gal-passage-ordered-rings}).
\begin{enumerate}
\item If $S\to T$ is a surjective ring map and $S$ is absolutely integrally closed, or ordered, then so
  is $T$.
\item If $S$ is $\varpi$-ordered, then $S/\varpi S$ is ordered.
\item For $k\geq1$, $S$ is $\varpi^k$-ordered if and only if it is $\varpi$-ordered, and if $S$ is
  $\varpi^k$-adically complete, then it is $\varpi$-adically complete.
\item If $S$ is absolutely integrally closed, then every monic polynomial in $S[X]$ is a product of
  monic polynomials of degree $1$, and $S[1/\varpi]$ is absolutely integrally closed.
\item If $S$ is $\varpi$-ordered, then $S$ and $S[1/\varpi]$ are componentwise local.
\end{enumerate}
\end{lemma}
\leanref{Development/WeilBound/CoreStatements/PadicHodgeComparison/EtaleCohomologyComparison/AcrisKummerUnits/ValuationSquareOfValuationMod.lean}{CrystallineSemistableComparison.RL.surj_pass}
\leanref{Development/AutomorphicSpectrum/Proofs/BlockDecomposition/ComparableQuotientOfValuation.lean}{GlobalTheory.CompanionTheory.ValQuot.valuation_comparable_quotient}
\leanref{Development/WeilBound/CoreStatements/PadicHodgeComparison/EtaleCohomologyComparison/AcrisKummerUnits/ValuationSquareOfValuationMod.lean}{CrystallineSemistableComparison.RL.val_pow}
\leanref{Development/AutomorphicSpectrum/Proofs/BlockDecomposition/AbsolutelyClosedPolynomialSplitsStatement.lean}{GlobalTheory.T.absolutely_closed_polynomial_splits}
\leanref{Development/WeilBound/CoreStatements/PadicHodgeComparison/EtaleCohomologyComparison/AcrisKummerUnits/ValuationSquareOfValuationMod.lean}{CrystallineSemistableComparison.RL.loc_aic}

\begin{proof}
For (1), a monic polynomial over $T$ lifts to a monic polynomial over $S$ of the same degree, and a pair
of elements of $T$ lifts to a pair of elements of $S$. The image of a root, or of an ordering
idempotent, does what is required.

For (2), let $f,g\in S$. The pairs $f,\varpi$ and $g,\varpi$ generate ideals containing $\varpi$, so
there are idempotents $e_1$ and $e_2$ with $e_1f,e_2g\in\varpi S$, $(1-e_1)\varpi\in fS$ and
$(1-e_2)\varpi\in gS$. Let $u=(1-e_1)(1-e_2)$, an idempotent with $e_1u=0$, and let $f_1=uf+1-u$ and
$g_1=ug+1-u$. The ideal $f_1S+g_1S$ contains $uf=uf_1$ and $1-u=(1-u)f_1$. It also contains
$u\varpi=u(1-e_2)(1-e_1)\varpi\in ufS$, hence $\varpi=u\varpi+(1-u)\varpi$. So an idempotent $e_3$
orders $f_1,g_1$. Then $e=e_1+ue_3$ is an idempotent. Modulo $\varpi S$ we have
$ef\equiv ue_3f=ue_3f_1\in ug_1S=ugS$. Since $1-e_1=u+e_2(1-e_1)$ and $e_2g\in\varpi S$, also
$(1-e)g\equiv u(1-e_3)g=u(1-e_3)g_1\in uf_1S=ufS$. So the image of $e$ orders the images of $f$ and $g$.

For (3), an ideal contains a power of $\varpi$ if and only if it contains a power of $\varpi^k$, and the
$\varpi^k$-adic and $\varpi$-adic topologies on $S$ agree.

For (4), a root $b$ of a monic $P$ of positive degree gives $P=(X-b)P_1$ with $P_1$ monic, and induction
on the degree applies. Let $P\in S[1/\varpi][X]$ be monic of degree $d\geq1$, and choose $N$ such that
$\varpi^N$ times each coefficient of $P$ lies in the image of $S$. Then $\varpi^{Nd}P(Y/\varpi^N)$ is a
monic polynomial in $Y$ with coefficients in the image of $S$. A root $y\in S$ of a monic lift of it
gives the root $y/\varpi^N$ of $P$.

For (5), let $f\in S$. The pair $f,1-f$ generates $S$, which contains $\varpi^0$, so an idempotent $e$
satisfies $ef\in(1-f)S$ and $(1-e)(1-f)\in fS$. Then $e'=1-e$ satisfies $e'=e'(1-f)+e'f\in fS$ and
$1-e'=ef+e(1-f)\in(1-f)S$. An element of $S[1/\varpi]$ has the form $f=a/\varpi^n$ with $a\in S$. The
pair $a,\varpi^n-a$ generates an ideal containing $\varpi^n$, and in $S[1/\varpi]$ its members are unit
multiples of $f$ and $1-f$. So the image of an idempotent ordering $a,\varpi^n-a$ orders $f,1-f$, and the
same construction applies.
\end{proof}

\begin{lemma}\label{gal-passage-integral-closure}
Let $S$ be a ring and let $\varpi\in S$ be a nonzerodivisor such that $S$ is $\varpi$-ordered
(Definition~\ref{gal-passage-ordered-rings}).
\begin{enumerate}
\item If $z\in S$, $m\geq0$ and $d\geq1$ satisfy $\varpi^{md}\mid z^d$, then $\varpi^m\mid z$.
\item $S$ is integrally closed in $S[1/\varpi]$.
\end{enumerate}
\end{lemma}
\leanref{Development/WeilBound/CoreStatements/PadicHodgeComparison/EtaleCohomologyComparison/AcrisKummerUnits/ValuationSquareOfValuationMod.lean}{CrystallineSemistableComparison.RL.pow_dvd}
\leanref{Development/WeilBound/CoreStatements/PadicHodgeComparison/EtaleCohomologyComparison/AcrisKummerUnits/ValuationSquareOfValuationMod.lean}{CrystallineSemistableComparison.RL.int_closed}

\begin{proof}
For (1), write $z^d=\varpi^{md}w$. The pair $\varpi^m,z$ generates an ideal containing $\varpi^m$, so an
idempotent $e$ orders it: $e\varpi^m=c'z$ and $(1-e)z=c\varpi^m$ for some $c,c'\in S$. Then
$e\varpi^{md}=(c'z)^d=c'^dw\varpi^{md}$, so $e=c'^dw$, since $\varpi$ is a nonzerodivisor. Hence
$ez=c'^dwz=wc'^{d-1}e\varpi^m$, and $z=ez+(1-e)z=(wc'^{d-1}e+c)\varpi^m$.

For (2), let $x\in S[1/\varpi]$ be a root of a monic $P=\sum_{i=0}^dp_iX^i\in S[X]$, and write
$x=f/\varpi^N$ with $f\in S$. An idempotent $e$ orders the pair $f,\varpi^N$. Let $u=1-e$, so that
$ef\in\varpi^NS$ and $u\varpi^N=fh$ for some $h\in S$. Then $hx=u$ in $S[1/\varpi]$. Multiplying
$P(x)=0$ by $uh^d$ and using $u(hx)^i=u$ gives $u(1+hk)=0$ with $k=\sum_{i<d}p_ih^{d-1-i}$. This
holds in $S$, since $\varpi$ is a nonzerodivisor and $S\to S[1/\varpi]$ is injective. Hence $u=-uhk$,
and $uf=-ku\,fh=-ku\varpi^N$. So $f=ef+uf\in\varpi^NS$, and $x$ lies in $S$.
\end{proof}

\begin{lemma}\label{gal-passage-components}
Let $S$ be a ring (Definition~\ref{gal-passage-ordered-rings}).
\begin{enumerate}
\item Let $P$ be a prime ideal of $S$. An element $x\in S$ lies in $I_P$ if and only if $ex=0$ for an
  idempotent $e\notin P$.
\item Let $\mathcal Q\subset S$ be a subset with $0\in\mathcal Q$, with $ee'\in\mathcal Q$ for all
  idempotents $e\in\mathcal Q$ and $e'\in S$, and with $e+e'\in\mathcal Q$ for all idempotents
  $e,e'\in\mathcal Q$ with $ee'=0$. Suppose that for every maximal ideal $\mathfrak m$ of $S$ there is
  an idempotent $e\in\mathcal Q$ with $e\notin\mathfrak m$. Then $1\in\mathcal Q$.
\item Let $\varpi\in S$ and let $P$ be a prime ideal of $S$. If $\varpi$ is a nonzerodivisor in $S$,
  then its image in $S/I_P$ is a nonzerodivisor. If $S$ is $\varpi$-ordered, then the image of
  $\varpi$ is chained in $S/I_P$.
\item Let $\varpi\in S$. If the image of $\varpi$ is chained in $S/I_{\mathfrak m}$ for every maximal
  ideal $\mathfrak m$ of $S$, then $S$ is $\varpi$-ordered.
\item Let $\varpi\in S$ and $k\geq1$, and let $S$ be $\varpi$-adically complete. Every idempotent of
  $S/\varpi^kS$ is the image of a unique idempotent of $S$.
\end{enumerate}
\end{lemma}
\leanref{Development/AutomorphicSpectrum/Proofs/BlockDecomposition/MemIdempotentComplementIffStatement.lean}{GlobalTheory.T.mem_idempotent_complement_iff}
\leanref{Development/WeilBound/CoreStatements/PadicHodgeComparison/EtaleCohomologyComparison/AcrisKummerUnits/ValuationSquareOfValuationMod.lean}{CrystallineSemistableComparison.RL.glue}
\leanref{Development/WeilBound/CoreStatements/PadicHodgeComparison/EtaleCohomologyComparison/AcrisKummerUnits/ValuationSquareOfValuationMod.lean}{CrystallineSemistableComparison.RL.dich_comp}
\leanref{Development/WeilBound/CoreStatements/PadicHodgeComparison/EtaleCohomologyComparison/AcrisKummerUnits/ValuationSquareOfValuationMod.lean}{CrystallineSemistableComparison.RL.val_of_dich}
\leanref{Development/AutomorphicSpectrum/Proofs/BlockDecomposition/LiftIdempotentAdicComplete.lean}{GlobalTheory.CompanionTheory.IdemLift.existsUnique_idempotent_lift}

\begin{proof}
For (1), let $e\notin P$ be an idempotent. Then $e(1-e)=0$ gives $1-e\in P$, so $1-e\in I_P$, and if
$ex=0$, then $x=(1-e)x\in I_P$. Conversely, the elements killed by some idempotent outside $P$ form an
ideal: if $e$ kills $x$ and $e'$ kills $x'$, then $ee'\notin P$ kills $x+x'$. This ideal contains every
idempotent $e\in P$, which $1-e\notin P$ kills.

For (2), let $I$ be the set of $x\in S$ with $ex=x$ for some idempotent $e\in\mathcal Q$. If
idempotents $e,e'\in\mathcal Q$ fix $x$ and $x'$, then $e'(1-e)\in\mathcal Q$ is orthogonal to $e$, so
$e+e'(1-e)\in\mathcal Q$, and it fixes $x$ and $x'$. So $I$ is an ideal. An idempotent $e\in\mathcal Q$
with $e\notin\mathfrak m$ lies in $I$, so $I$ lies in no maximal ideal and $1\in I$. Then $e=e\cdot1=1$
for some $e\in\mathcal Q$.

For (3), if $x\varpi\in I_P$, then by (1) $ex\varpi=0$ for an idempotent $e\notin P$, so $ex=0$ and
$x\in I_P$. Now let $S$ be $\varpi$-ordered, and let $f,g,a,b\in S$ with $\varpi^N-af-bg\in I_P$. By (1)
there is an idempotent $e\notin P$ with $e(\varpi^N-af-bg)=0$. Then
$\varpi^N=(ae+(1-e)\varpi^N)(ef+1-e)+beg$, so an idempotent $e_1$ orders the pair $ef+1-e,g$. As
$1-e\in I_P$, the image of $ef+1-e$ in $S/I_P$ is that of $f$. If $e_1\in P$, then $e_1\in I_P$, and
$(1-e_1)g\in(ef+1-e)S$ shows that $f$ divides $g$ in $S/I_P$. Otherwise $1-e_1\in P$, so $1-e_1\in I_P$,
and $e_1(ef+1-e)\in gS$ shows that $g$ divides $f$ in $S/I_P$.

For (4), let $f,g\in S$ with $\varpi^N\in fS+gS$. Let $\mathcal Q$ be the set of $e\in S$ for which there
is an idempotent $d$ with $de=d$, $df\in gS$ and $(e-d)g\in fS$. It contains $0$, with $d=0$. If
$e\in\mathcal Q$ with $d$ and $e'$ is an idempotent, then $ee'\in\mathcal Q$ with $de'$. If
$e,e'\in\mathcal Q$ are orthogonal idempotents, with $d$ and $d'$, then $e+e'\in\mathcal Q$ with $d+d'$.
Let $\mathfrak m$ be a maximal ideal. In $S/I_{\mathfrak m}$ the image of $\varpi^N$ lies in the ideal
generated by the images of $f$ and $g$, so there is $a\in S$ with $f-ga\in I_{\mathfrak m}$ or
$g-fa\in I_{\mathfrak m}$. By (1) there is an idempotent $e\notin\mathfrak m$ with $ef=ega$ in the first
case, and then $e\in\mathcal Q$ with $d=e$. In the second case $eg=efa$, and $e\in\mathcal Q$ with $d=0$.
By (2), $1\in\mathcal Q$, and the corresponding $d$ orders $f,g$.

For (5), let $a\in S$ lift an idempotent of $S/\varpi^kS$, so that $a^2-a\in\varpi S$. As $S$ is
$\varpi$-adically complete, it is henselian along $\varpi S$, and the derivative $2a-1$ of $X^2-X$ at $a$
is a unit modulo $\varpi$, since $(2a-1)^2=1+4(a^2-a)$. So $X^2-X$ has a root $y$ with
$y\equiv a$ modulo $\varpi$. Two idempotents whose difference $\delta$ is nilpotent are equal, since
$\delta^3=\delta$; so $y$ and $a$ have the same image in $S/\varpi^kS$. If $y'$ is another idempotent
lift, then $\delta=y-y'$ satisfies $\delta^3=\delta$ and lies in $\varpi S$, which is contained in the
Jacobson radical. So $1-\delta^2$ is a unit, and $\delta(1-\delta^2)=0$ gives $\delta=0$.
\end{proof}

\begin{lemma}\label{gal-passage-chained-mod-square}
Let $S$ be a ring and $\varpi\in S$ (Definition~\ref{gal-passage-ordered-rings}).
\begin{enumerate}
\item Suppose that $\varpi$ lies in the Jacobson radical of $S$, and that for all $f,g\in S$ one has
  $f\in gS+\varpi^2S$ or $g\in fS+\varpi^2S$. Then $\varpi$ is chained in $S$.
\item Suppose that $S$ is $\varpi$-adically complete and that $S/\varpi^2S$ is ordered. Let $P$ be a
  prime ideal of $S$, and $\bar\varpi$ the image of $\varpi$ in $\bar S=S/I_P$. For all
  $f,g\in\bar S$ one has $f\in g\bar S+\bar\varpi^2\bar S$ or $g\in f\bar S+\bar\varpi^2\bar S$.
\end{enumerate}
\end{lemma}
\leanref{Development/WeilBound/CoreStatements/PadicHodgeComparison/EtaleCohomologyComparison/AcrisKummerUnits/ValuationSquareOfValuationMod.lean}{CrystallineSemistableComparison.RL.dich_of_mod}
\leanref{Development/WeilBound/CoreStatements/PadicHodgeComparison/EtaleCohomologyComparison/AcrisKummerUnits/ValuationSquareOfValuationMod.lean}{CrystallineSemistableComparison.RL.valuation_square_of_valuation_mod}

\begin{proof}
For (1), first let $f\in S$, and apply the hypothesis to $f$ and $\varpi$. Either $\varpi\mid f$, or
$\varpi=cf+w\varpi^2$ for some $c,w\in S$; then $f$ divides $(1-w\varpi)\varpi$, hence $\varpi$, as
$1-w\varpi$ is a unit. By induction on $M$, every $f\in S$ is divisible by $\varpi^M$ or has the form
$\varpi^jf'$ with $f'\mid\varpi$. Now let $af+bg=\varpi^N$. If $f=\varpi^{N+1}t$, then
$bg=(1-at\varpi)\varpi^N$, so $g\mid\varpi^N$ and $g\mid f$. The same holds with $f$ and $g$ exchanged.
Otherwise $f=\varpi^if'$ and $g=\varpi^jg'$ with $f',g'\mid\varpi$. If $i<j$, then
$f\mid\varpi^{i+1}\mid g$, and symmetrically if $j<i$. If $i=j$, the hypothesis applied to $f',g'$ gives
$f'\in g'S+\varpi^2S=g'S$ or $g'\in f'S$, since $f'$ and $g'$ divide $\varpi^2$. Multiplying by
$\varpi^i$ gives $f\in gS$ or $g\in fS$.

For (2), lift $f,g$ to $S$. An idempotent of $S/\varpi^2S$ orders their images, and by
Lemma~\ref{gal-passage-components}(5) it is the image of an idempotent $e\in S$. So
$ef\in gS+\varpi^2S$ and $(1-e)g\in fS+\varpi^2S$. If $e\in P$, then $e\in I_P$ and the image of $1-e$
is $1$, so $g\in f\bar S+\bar\varpi^2\bar S$. Otherwise $1-e\in P$, so $1-e\in I_P$ and
$f\in g\bar S+\bar\varpi^2\bar S$.
\end{proof}

\begin{lemma}\label{gal-passage-newton}
Let $S$ be a ring and $\varpi\in S$ a nonzerodivisor that lies in the Jacobson radical of $S$ and is
chained in $S$, and suppose that $S/\varpi^2S$ is absolutely integrally closed
(Definition~\ref{gal-passage-ordered-rings}).
\begin{enumerate}
\item For every divisor $f$ of a power of $\varpi$ and every $i\geq1$ there are $r\in S$ and a unit $u$
  of $S$ with $f=ur^i$.
\item Let $d\geq1$ and $q_0,\ldots,q_d\in S$ with $q_d=1$ and $q_0$ dividing a power of $\varpi$. There
  are $c\in S$ and $1\leq i_0\leq d$ such that $q_0\mid q_ic^i$ for $1\leq i\leq d$ and
  $q_{i_0}c^{i_0}\mid q_0$.
\item $\varpi$ admits Newton steps in $S$.
\end{enumerate}
\end{lemma}
\leanref{Development/AutomorphicSpectrum/Proofs/BlockDecomposition/PolynomialRootsOfClosedQuotientStatement.lean}{GlobalTheory.T.polynomial_roots_of_closed_quotient}
\leanref{Development/AutomorphicSpectrum/Proofs/BlockDecomposition/NewtonPolygonSlopeExistenceStatement.lean}{GlobalTheory.T.newton_polygon_slope_existence}
\leanref{Development/AutomorphicSpectrum/Proofs/BlockDecomposition/AbsolutelyClosedUnitCoefficientRootStatement.lean}{GlobalTheory.T.absolutely_closed_unit_coefficient_root}
\leanref{Development/WeilBound/CoreStatements/PadicHodgeComparison/EtaleCohomologyComparison/AcrisKummerUnits/ValuationSquareOfValuationMod.lean}{CrystallineSemistableComparison.RL.step_comp}

\begin{proof}
We may assume $S\neq0$. Since $\varpi$ is chained, every $f\in S$ satisfies $\varpi\mid f$ or
$f\mid\varpi$, as the pair $f,\varpi$ generates an ideal containing $\varpi$. By induction on $M$, every
$f\in S$ is divisible by $\varpi^M$ or has the form $\varpi^jf'$ with $f'\mid\varpi$.

For (1), first let $f\mid\varpi$, say $\varpi=fh$. As $S/\varpi^2S$ is absolutely integrally closed,
$X^i-f$ has a root modulo $\varpi^2$: there are $r,t\in S$ with
$r^i=f+t\varpi^2=f(1+th\varpi)$, and $1+th\varpi$ is a unit. In general let $f\mid\varpi^N$. If
$\varpi^{N+1}\mid f$, then $\varpi^N=\varpi^{N+1}k$ for some $k$, so $\varpi k=1$, which is impossible
for $\varpi$ in the Jacobson radical of $S\neq0$. So $f=\varpi^jf'$ with $f'\mid\varpi$. Writing
$\varpi=u_0r_0^i$ and $f'=ur^i$ gives $f=u_0^ju(r_0^jr)^i$.

For (2), let $T$ be the set of $i\in\{1,\ldots,d\}$ with $q_i\mid q_0$; it contains $d$. Let
$q_0\mid\varpi^M$. For $i\in T$ write $q_0=q_it_i$; then $t_i\mid\varpi^M$, so by (1) $t_i=u_is_i^i$
with $u_i$ a unit. The $s_i$, $i\in T$, divide $\varpi^M$, so any two of them generate an ideal
containing $\varpi^M$ and one of them divides the other. Choose $i_0\in T$ with $s_i\mid s_{i_0}$ for
all $i\in T$, and let $c=s_{i_0}$. Then $q_{i_0}c^{i_0}=u_{i_0}^{-1}q_0$. For $i\in T$,
$q_0=u_iq_is_i^i$ divides $q_ic^i$. For $i\notin T$, the pair $q_0,q_i$ generates an ideal containing
$\varpi^M$ and $q_i\nmid q_0$, so $q_0\mid q_i$.

For (3), let $Q=\sum_{i=0}^dq_iX^i$ be monic of degree $d\geq1$ with $\varpi^n\mid q_0$. As $\varpi$ is
chained, $\varpi^{n+1}\mid q_0$, and then $z=0$ will do, or $q_0\mid\varpi^{n+1}$. In the second case
let $c$ and $i_0$ be as in (2), and write $q_ic^i=q_0s_i$ for $1\leq i\leq d$ and $s_0=1$. Since $q_0$
divides $\varpi^{n+1}$, it is a nonzerodivisor, and $q_0=q_{i_0}c^{i_0}t=q_0s_{i_0}t$ for some $t\in S$
shows that $s_{i_0}$ is a unit. Let $R=\sum_{i=0}^ds_iX^i$. In $\bar S=S/\varpi^2S$, if $a+b$ is a unit,
then so is $a$ or $b$: lifts of $a$ and $b$ to $S$ have a sum that is a unit of $S$, as $\varpi$ lies in
the Jacobson radical, so the lifts generate $S$ and one of them divides the other, hence divides their
sum and is a unit. So $\bar S$ is a local ring. The polynomial $X^dR(1/X)$ is monic over $\bar S$, so by
Lemma~\ref{gal-passage-ordered-basics}(4) it is a product $\prod_k(X-b_k)$ there. If every $b_k$ were in
the maximal ideal of $\bar S$, then so would be every coefficient of $X^dR(1/X)$ except the leading one,
including the unit $s_{i_0}$. So some $b_k$ is a unit, and its inverse is a root of $R$ in $\bar S$.
This gives $y\in S$ with $\varpi^2\mid R(y)$. Then $Q(cy)=\sum_iq_ic^iy^i=q_0R(y)$ is divisible by
$\varpi^{n+2}$, and $\varpi^n\mid q_0\mid q_dc^d=c^d$ gives $\varpi^n\mid(cy)^d$. So $z=cy$ will do.
\end{proof}

\begin{lemma}\label{gal-passage-newton-limit}
Let $S$ be a ring and $\varpi\in S$ (Definition~\ref{gal-passage-ordered-rings}).
\begin{enumerate}
\item If the image of $\varpi$ admits Newton steps in $S/I_{\mathfrak m}$ for every maximal ideal
  $\mathfrak m$ of $S$, then $\varpi$ admits Newton steps in $S$.
\item Suppose that $S$ is $\varpi$-adically complete, that $\varpi$ admits Newton steps in $S$, and that
  $\varpi^{md}\mid z^d$ implies $\varpi^m\mid z$ for all $z\in S$, $m\geq0$ and $d\geq1$. Then $S$ is
  absolutely integrally closed.
\end{enumerate}
\end{lemma}
\leanref{Development/WeilBound/CoreStatements/PadicHodgeComparison/EtaleCohomologyComparison/AcrisKummerUnits/ValuationSquareOfValuationMod.lean}{CrystallineSemistableComparison.RL.step_glob}
\leanref{Development/WeilBound/CoreStatements/PadicHodgeComparison/EtaleCohomologyComparison/AcrisKummerUnits/ValuationSquareOfValuationMod.lean}{CrystallineSemistableComparison.RL.aic_limit}

\begin{proof}
For (1), let $Q\in S[X]$ be monic of degree $d\geq1$, and $\varpi^n\mid Q(0)$. Let $\mathcal Q$ be the
set of $e\in S$ for which some $z\in S$ has $\varpi^{n+1}\mid eQ(z)$ and $\varpi^n\mid ez^d$. It contains
$0$ and is closed under multiplication by idempotents. Let $e,e'\in\mathcal Q$ be orthogonal idempotents,
with $z$ and $z'$. Then $e(ez+e'z')=ez$ and $e'(ez+e'z')=e'z'$, so $eQ(ez+e'z')=eQ(z)$,
$e'Q(ez+e'z')=e'Q(z')$, and likewise for $d$-th powers; so $e+e'\in\mathcal Q$, with $ez+e'z'$. Let
$\mathfrak m$ be a maximal ideal. A Newton step for the image of $Q$ in $S/I_{\mathfrak m}$ gives
$z,a,b\in S$ with $Q(z)-\varpi^{n+1}a$ and $z^d-\varpi^nb$ in $I_{\mathfrak m}$. By
Lemma~\ref{gal-passage-components}(1) they are killed by idempotents $e_1,e_2\notin\mathfrak m$, and
$e_1e_2\notin\mathfrak m$ lies in $\mathcal Q$. By Lemma~\ref{gal-passage-components}(2), $1\in\mathcal Q$.

For (2), let $Q\in S[X]$ be monic of degree $d\geq1$. Newton steps for the monic polynomials $Q(X+x)$ give
$x_0=0,x_1,x_2,\ldots$ in $S$ with $\varpi^n\mid Q(x_n)$ and $\varpi^n\mid(x_{n+1}-x_n)^d$ for all $n$.
By hypothesis $\varpi^{\lfloor n/d\rfloor}\mid x_{n+1}-x_n$, so $\varpi^k\mid x_{(k+1)d}-x_{kd}$ for all
$k$. The sequence $(x_{kd})_k$ converges to some $L\in S$ with $\varpi^k\mid x_{kd}-L$. Then $\varpi^k$
divides $Q(x_{kd})$ and $Q(x_{kd})-Q(L)$, hence $Q(L)$, for every $k$, and $Q(L)=0$ as $S$ is
$\varpi$-adically separated.
\end{proof}

\begin{lemma}\label{gal-passage-lift-square}
Let $S$ be a ring and $\varpi\in S$ a nonzerodivisor. Suppose that $S$ is $\varpi$-adically complete
and that $S/\varpi^2S$ is ordered and absolutely integrally closed
(Definition~\ref{gal-passage-ordered-rings}). Then $S$ is $\varpi$-ordered and absolutely integrally
closed.
\end{lemma}
\leanref{Development/AutomorphicArithmetic/CommutativeAlgebra/EtaleCohomology/TestFactsAlt.lean}{val_aic}

\begin{proof}
Since $S$ is $\varpi$-adically complete, $\varpi$ lies in the Jacobson radical of $S$, so its image lies
in the Jacobson radical of every quotient of $S$. Let $P$ be a prime ideal of $S$. By
Lemma~\ref{gal-passage-chained-mod-square}(2) and~(1), the image of $\varpi$ is chained in $S/I_P$. By
Lemma~\ref{gal-passage-components}(4), $S$ is $\varpi$-ordered.

Let $\mathfrak m$ be a maximal ideal. The image of $\varpi$ in $S/I_{\mathfrak m}$ is a nonzerodivisor by
Lemma~\ref{gal-passage-components}(3); it lies in the Jacobson radical and is chained. The quotient of
$S/I_{\mathfrak m}$ by the square of the image of $\varpi$ is a quotient of $S/\varpi^2S$, so it is
absolutely integrally closed by Lemma~\ref{gal-passage-ordered-basics}(1). By
Lemma~\ref{gal-passage-newton}(3) the image of $\varpi$ admits Newton steps in $S/I_{\mathfrak m}$, so
$\varpi$ admits Newton steps in $S$ by Lemma~\ref{gal-passage-newton-limit}(1).
Lemma~\ref{gal-passage-integral-closure}(1) gives the last hypothesis of
Lemma~\ref{gal-passage-newton-limit}(2), which shows that $S$ is absolutely integrally closed.
\end{proof}

\begin{lemma}\label{gal-passage-etale-sections}
Let $S$ be an absolutely integrally closed and componentwise local ring
(Definition~\ref{gal-passage-ordered-rings}). Every \'etale surjective morphism of schemes
$Y\to\operatorname{Spec}S$ has a section.
\end{lemma}
\leanref{Development/AutomorphicArithmetic/CommutativeAlgebra/EtaleCohomology/TestFactsAlt.lean}{splits}
\leanref{Development/WeilBound/CoreStatements/PadicHodgeComparison/EtaleCohomologyComparison/AcrisKummerUnits/ValuationSquareOfValuationMod.lean}{CrystallineSemistableComparison.RL.split_pt}
\leanref{Development/WeilBound/CoreStatements/PadicHodgeComparison/EtaleCohomologyComparison/AcrisKummerUnits/ValuationSquareOfValuationMod.lean}{CrystallineSemistableComparison.RL.sec_rules}
\leanref{Development/WeilBound/CoreStatements/PadicHodgeComparison/EtaleCohomologyComparison/AcrisKummerUnits/ValuationSquareOfValuationMod.lean}{CrystallineSemistableComparison.RL.sec_add}

\begin{proof}
Let $\pi:Y\to\operatorname{Spec}S$ be \'etale and surjective. Let $\mathcal Q$ be the set of $e\in S$
such that there is a morphism $s:\operatorname{Spec}S[1/e]\to Y$ with $\pi s$ the canonical morphism. It
contains $0$, as $S[1/0]=0$. If $e\in\mathcal Q$ and $e'\in S$, then $ee'\in\mathcal Q$, since
$\operatorname{Spec}S[1/ee']\to\operatorname{Spec}S$ factors through $\operatorname{Spec}S[1/e]$. For
orthogonal idempotents $e,e'$ one has $S[1/(e+e')]=S[1/e]\times S[1/e']$, so morphisms over the two
factors give one over $\operatorname{Spec}S[1/(e+e')]$. By Lemma~\ref{gal-passage-components}(2) it
suffices to find, for each maximal ideal $\mathfrak m$ of $S$, an idempotent $e\in\mathcal Q$ with
$e\notin\mathfrak m$: then $1\in\mathcal Q$, which gives a section of $\pi$.

Let $y\in Y$ lie over $\mathfrak m$, let $\operatorname{Spec}B\subset Y$ be an affine open neighbourhood
of $y$, so that $S\to B$ is \'etale, and let $\mathfrak q\subset B$ be the prime ideal of $y$. By
\cite[Tag 00UE]{Sta} there is $g\in B$ with $g\notin\mathfrak q$ such that $B_g\cong S[X]_h/(f)$ as
$S$-algebras, for a monic $f\in S[X]$ and an $h\in S[X]$ such that $f'$ is invertible in $S[X]_h/(f)$.
Let $x\in B_g$ be the image of $X$, and $\mathfrak q'=\mathfrak qB_g$, a prime ideal over $\mathfrak m$.
By Lemma~\ref{gal-passage-ordered-basics}(4), $f=\prod_k(X-b_k)$ with $b_k\in S$. Since
$\prod_k(x-b_k)=0$, there is a root $b$ of $f$ with $x-b\in\mathfrak q'$. Then $c=h(b)\notin\mathfrak m$:
otherwise $h(x)\in\mathfrak q'$, while $h(x)$ is a unit of $B_g$. As $\mathfrak m$ is maximal, there is
$a\in S$ with $1-ac\in\mathfrak m$. Since $S$ is componentwise local, there is an idempotent $e\in acS$
with $1-e\in(1-ac)S\subset\mathfrak m$. So $e\notin\mathfrak m$, and $c$ is a unit in $S[1/e]$. Hence
$X\mapsto b$ defines a homomorphism of $S$-algebras $S[X]_h/(f)\to S[1/e]$, and its composite with
$B\to B_g$ gives a morphism $\operatorname{Spec}S[1/e]\to\operatorname{Spec}B\subset Y$ over
$\operatorname{Spec}S$. So $e\in\mathcal Q$.
\end{proof}

\begin{lemma}\label{gal-passage-untilt}
Let $A$ be a test algebra, with untilt $A^\sharp=W(A)/\xi W(A)$ and quotient map
$\theta:W(A)\to A^\sharp$ (Definition~\ref{gal-crystalline-test-algebras}).
\begin{enumerate}
\item $\ell$ is a nonzerodivisor in $A^\sharp$, and $A^\sharp$ is $\ell$-adically complete.
\item The map $A\to A^\sharp/\ell A^\sharp$, $a\mapsto\theta([a])\bmod\ell$, is a surjective ring
  homomorphism with kernel $\ell^\flat A$.
\item Let $t\in A$ with $t^\ell=\ell^\flat$, and let $\varpi=\theta([t])$. Then $\varpi^\ell=\ell$,
  $\varpi$ is a nonzerodivisor in $A^\sharp$, $A^\sharp$ is $\varpi$-adically complete, and
  $A^\sharp/\varpi^2A^\sharp$ is ordered and absolutely integrally closed
  (Definition~\ref{gal-passage-ordered-rings}).
\end{enumerate}
\end{lemma}
\leanref{Development/WeilBound/CoreStatements/PadicHodgeComparison/EtaleCohomologyComparison/AcrisKummerUnits/FiltrationSpanDividedPowers.lean}{CrystallineSemistableComparison.RL.u_tf}
\leanref{Development/WeilBound/CoreStatements/PadicHodgeComparison/EtaleCohomologyComparison/AcrisKummerUnits/FiltrationSpanDividedPowers.lean}{CrystallineSemistableComparison.RL.u_complete}
\leanref{Development/WeilBound/CoreStatements/PadicHodgeComparison/EtaleCohomologyComparison/AcrisKummerUnits/FiltrationSpanDividedPowers.lean}{CrystallineSemistableComparison.RL.u_mod}
\leanref{Development/WeilBound/CoreStatements/PadicHodgeComparison/EtaleCohomologyComparison/AcrisKummerUnits/FiltrationSpanDividedPowers.lean}{CrystallineSemistableComparison.RL.untilt_valuation_mod_square}

\begin{proof}
For (1), since $A$ is perfect of characteristic $\ell$, $\ell$ is a nonzerodivisor in $W(A)$, $W(A)$ is
$\ell$-adically complete, and $x\mapsto x_0$ identifies $W(A)/\ell W(A)$ with $A$. It sends $\xi$ to
$-\ell^\flat$, a nonzerodivisor in $A$. So $\xi y\in\ell W(A)$ implies $y\in\ell W(A)$, and by induction
on $n$, $\xi y\in\ell^nW(A)$ implies $y\in\ell^nW(A)$. If $\ell y=\xi c$ in $W(A)$, then $c=\ell c'$ and
$y=\xi c'$; so $\ell$ is a nonzerodivisor in $A^\sharp$. If $a\in\xi W(A)+\ell^nW(A)$ for all $n$, write
$a=\xi y_n+\ell^nz_n$. Then $\xi(y_{n+1}-y_n)\in\ell^nW(A)$, so $y_{n+1}-y_n\in\ell^nW(A)$, and
$a=\xi\lim_ny_n$. So $A^\sharp$ is $\ell$-adically separated. A sequence in $A^\sharp$ whose successive
differences lie in $\ell^nA^\sharp$ lifts, difference by difference, to such a sequence in $W(A)$, whose
limit maps to a limit of the given sequence.

For (2), every $x\in W(A)$ satisfies $x-[x_0]\in\ell W(A)$, so $\theta(x)\equiv\theta([x_0])$ modulo
$\ell$. Hence the map is additive and surjective. It kills $\ell^\flat$, since
$\theta([\ell^\flat])=\theta(\ell-\xi)=\ell$. If $\theta([a])=\ell\theta(y)$, then $[a]-\ell y=\xi c$ for
some $c\in W(A)$, and comparing first components gives $a=-\ell^\flat c_0$.

For (3), $t$ exists because $A$ is perfect, and $\varpi^\ell=\theta([\ell^\flat])=\ell$. So
$\varpi^{\ell-1}\cdot\varpi$ is a nonzerodivisor, hence so is $\varpi$, and $A^\sharp$ is
$\varpi$-adically complete by (1) and Lemma~\ref{gal-passage-ordered-basics}(3). Since $\ell\geq2$, $\ell$
lies in $\varpi^2A^\sharp$. So the map of (2), followed by $A^\sharp/\ell\to A^\sharp/\varpi^2$, is a
surjective ring map $A\to A^\sharp/\varpi^2A^\sharp$ that kills $\ell^\flat$. The test algebra $A$ is
absolutely integrally closed and $\ell^\flat$-ordered, so $A/\ell^\flat A$ is ordered by
Lemma~\ref{gal-passage-ordered-basics}(2). Lemma~\ref{gal-passage-ordered-basics}(1), applied to the
surjections from $A$ and from $A/\ell^\flat A$ onto $A^\sharp/\varpi^2A^\sharp$, completes the proof.
\end{proof}

\begin{lemma}\label{gal-crystalline-test-facts}
Let $A$ be a test algebra (Definition~\ref{gal-crystalline-test-algebras}).
\begin{enumerate}
\item $\ell$ is a nonzerodivisor in $A^\sharp$, and $A^\sharp$ is integrally closed in $A^\sharp[1/\ell]$.
\item Every \'etale surjective morphism onto $\operatorname{Spec}A^\sharp$ or onto
  $\operatorname{Spec}A^\sharp[1/\ell]$ has a section.
\item Every idempotent of $A^\sharp$ is $\theta([e_0])$ for a unique idempotent $e_0$ of $A$.
\item For orthogonal idempotents $e_s$ of $A$ with sum $1$ there is a decomposition $(A\to A_s)$ with
  $e_s\mapsto1$ in $A_s$. Every open cover of $\operatorname{Spec}A^\sharp$ is refined by a
  decomposition $(A\to A_s)$: each $\operatorname{Spec}A_s^\sharp\to\operatorname{Spec}A^\sharp$ lands in
  a member of the cover.
\end{enumerate}
\end{lemma}
\leanref{Development/AutomorphicArithmetic/CommutativeAlgebra/EtaleCohomology/TestFacts.lean}{UpperTheory.Ru.crystallineComparison.TestFacts}
\leanref{Development/AutomorphicArithmetic/CommutativeAlgebra/EtaleCohomology/TestFactsAlt.lean}{UpperTheory.Ru.crystallineComparison.testFacts}
\leanref{Development/AutomorphicArithmetic/CommutativeAlgebra/EtaleCohomology/CrystallineComparison.lean}{UpperTheory.Ru.crystallineComparison.decEx}
\leanref{Development/AutomorphicArithmetic/EtaleIdempotents/ExistsOrthogonalIdempotentsSubordinateOfEtaleSec.lean}{UpperTheory.exists_orthogonalIdempotents_subordinate_of_etale_sections}
\leanref{Development/AutomorphicArithmetic/EtaleIdempotents/ExistsOrthogonalIdempotentsLiftThroughSharp.lean}{UpperTheory.exists_orthogonalIdempotents_lift_through_sharp}

\begin{proof}
Let $t\in A$ with $t^\ell=\ell^\flat$, and let $\varpi=\theta([t])$. By
Lemma~\ref{gal-passage-untilt}(3) and Lemma~\ref{gal-passage-lift-square}, $A^\sharp$ is
$\varpi$-ordered and absolutely integrally closed (Definition~\ref{gal-passage-ordered-rings}); these
are the two defining properties of $A$, with $\varpi$ in place of $\ell^\flat$. Since $\ell=\varpi^\ell$,
$A^\sharp$ is $\ell$-ordered by Lemma~\ref{gal-passage-ordered-basics}(3). Item (1) follows from
Lemma~\ref{gal-passage-untilt}(1) and Lemma~\ref{gal-passage-integral-closure}(2), applied to $\ell$. By
Lemma~\ref{gal-passage-ordered-basics}(4) and~(5), $A^\sharp$ and $A^\sharp[1/\ell]$ are absolutely
integrally closed and componentwise local, and Lemma~\ref{gal-passage-etale-sections} gives (2).

For (3), an idempotent of $A^\sharp$ reduces to an idempotent of $A^\sharp/\ell=A/\ell^\flat$, which lifts
to a unique idempotent $e_0$ of the $\ell^\flat$-adically complete ring $A$; then $\theta([e_0])$ and the
given idempotent have the same reduction, so they agree, since $A^\sharp$ is $\ell$-adically complete.
In (4), $A$ is the product of the rings $A/(1-e_s)$, which are again test algebras. An open cover of
$\operatorname{Spec}A^\sharp$ is refined by finitely many basic opens, whose disjoint union is \'etale and
surjective over $\operatorname{Spec}A^\sharp$; a section as in (2) gives orthogonal idempotents of
$A^\sharp$ subordinate to the cover, and (3) carries them to $A$.
\end{proof}

\begin{definition}\label{gal-shifting-coefficient-systems}
Let $\ell$ be a prime, with test algebras and their decompositions as in
Definition~\ref{gal-crystalline-test-algebras}.
\begin{enumerate}
\item A \emph{coefficient system} is a functor $F$ from test algebras to abelian groups; a map of
  coefficient systems is a natural transformation, and a sequence of coefficient systems is exact if it
  is exact at every test algebra. A coefficient system $F$ \emph{carries decompositions to products} if
  $F(A)\to\prod_sF(A_s)$ is bijective for every test algebra $A$ and every decomposition $(A\to A_s)$
  of $A$.
\item Let $P$ be a functor from test algebras to sets and $F$ a coefficient system. A $k$-cochain of
  $P$ with coefficients in $F$ assigns to every chain $\sigma=((A_0,y_0)\to\cdots\to(A_k,y_k))$ in
  $\mathrm{El}(P)$ an element $c(\sigma)\in F(A_k)$. Its differential is the $(k+1)$-cochain
  \[
   (dc)(\sigma)=\sum_{i=0}^{k}(-1)^ic(\partial_i\sigma)+(-1)^{k+1}F(\gamma)\,c(\partial_{k+1}\sigma),
  \]
  where $\sigma$ has length $k+1$, $\partial_i\sigma$ omits $(A_i,y_i)$, and $\gamma:A_k\to A_{k+1}$ is
  the last map of $\sigma$. $H^k(P,F)$ is the $k$-th cohomology of this complex. A natural
  transformation $u:Q\to P$ induces $u^*:H^k(P,F)\to H^k(Q,F)$ through $(u^*c)(\sigma)=c(u(\sigma))$,
  where $u(\sigma)=((A_0,u(y_0))\to\cdots\to(A_k,u(y_k)))$, and a map $\alpha:F\to G$ of coefficient
  systems induces $\alpha_*:H^k(P,F)\to H^k(P,G)$ through $\alpha_*c=\alpha\circ c$. For the coefficient
  system $A\mapsto A_{\mathrm{cris}}(A)/\ell^\nu$ these are the groups $H^k(P,A_{\mathrm{cris}}/\ell^\nu)$
  and the maps $u^*$ of Definition~\ref{gal-crystalline-test-algebras}(2).
\end{enumerate}
\end{definition}
\leanref{Development/AutomorphicArithmetic/CommutativeAlgebra/EtaleCohomology/ZeroPull.lean}{UpperTheory.Ru.crystallineComparison.Coef}
\leanref{Development/AutomorphicArithmetic/CommutativeAlgebra/EtaleCohomology/ZeroPull.lean}{UpperTheory.Ru.crystallineComparison.Fam}
\leanref{Development/AutomorphicArithmetic/CommutativeAlgebra/EtaleCohomology/ZeroPull.lean}{UpperTheory.Ru.crystallineComparison.Parts.Add}
\leanref{Development/AutomorphicArithmetic/CommutativeAlgebra/EtaleCohomology/TPullTMap.lean}{UpperTheory.Ru.crystallineComparison.Th.TH}
\leanref{Development/AutomorphicArithmetic/CommutativeAlgebra/EtaleCohomology/TPullTMap.lean}{UpperTheory.Ru.crystallineComparison.Th.tPull}
\leanref{Development/AutomorphicArithmetic/CommutativeAlgebra/EtaleCohomology/TPullTMap.lean}{UpperTheory.Ru.crystallineComparison.Th.tMap}

\begin{lemma}\label{gal-shifting-cohomology}
Let $\ell$ be a prime, $P$ and $Q$ functors from test algebras to sets, $u:Q\to P$ a natural
transformation and $F$ a coefficient system (Definition~\ref{gal-shifting-coefficient-systems}).
\begin{enumerate}
\item $H^0(P,F)$ is the group of families $c=(c(A,y))$, with $c(A,y)\in F(A)$ for $(A,y)$ in
  $\mathrm{El}(P)$, such that $F(\gamma)c(A,y)=c(B,P(\gamma)y)$ for every map $\gamma:A\to B$ of test
  algebras. On such families, $u^*c$ is $(A,y)\mapsto c(A,u(y))$, and $\alpha_*c=\alpha\circ c$ for a
  map $\alpha$ of coefficient systems.
\item An exact sequence $0\to F_1\xrightarrow{\alpha}F_2\xrightarrow{\beta}F_3\to0$ of coefficient
  systems gives maps $\delta:H^k(P,F_3)\to H^{k+1}(P,F_1)$, $k\geq0$, such that the sequence
  \[
   0\to H^0(P,F_1)\xrightarrow{\alpha_*}H^0(P,F_2)\xrightarrow{\beta_*}H^0(P,F_3)\xrightarrow{\delta}
   H^1(P,F_1)\xrightarrow{\alpha_*}H^1(P,F_2)\xrightarrow{\beta_*}\cdots
  \]
  is exact, and $u^*\circ\delta=\delta\circ u^*$.
\item If $K$ is an injective object of the category of coefficient systems, then $H^k(P,K)=0$ for every
  $k\geq1$.
\end{enumerate}
\end{lemma}
\leanref{Development/AutomorphicArithmetic/CommutativeAlgebra/EtaleCohomology/ZeroPull.lean}{UpperTheory.Ru.crystallineComparison.Sec}
\leanref{Development/AutomorphicArithmetic/CommutativeAlgebra/EtaleCohomology/ZeroPull.lean}{UpperTheory.Ru.crystallineComparison.Theory}
\leanref{Development/AutomorphicArithmetic/CommutativeAlgebra/EtaleCohomology/TheoryEx.lean}{UpperTheory.Ru.crystallineComparison.theoryEx}
\leanref{Development/AutomorphicArithmetic/EtaleIdempotents/ExistsCohomologyZeroEquivSectionsNatural.lean}{UpperTheory.exists_cohomology_zero_equiv_sections_natural}
\leanref{Development/AutomorphicArithmetic/CommutativeAlgebra/EtaleCohomology/TPullTMap.lean}{UpperTheory.Ru.crystallineComparison.Th.les_ex}
\leanref{Development/AutomorphicArithmetic/CommutativeAlgebra/EtaleCohomology/TPullTMap.lean}{UpperTheory.Ru.crystallineComparison.Th.tPull_tδ}
\leanref{Development/AutomorphicArithmetic/CommutativeAlgebra/EtaleCohomology/GFaceSnGCastSucc.lean}{UpperTheory.Ru.crystallineComparison.Th.vanish}

\begin{proof}
For (1): there are no cochains in degree $-1$, so $H^0(P,F)$ is the group of $0$-cocycles. A
$0$-cochain is a family $c(A,y)\in F(A)$, and for a map $\gamma:(A_0,y_0)\to(A_1,y_1)$ in
$\mathrm{El}(P)$, $(dc)(\gamma)=c(A_1,y_1)-F(\gamma)c(A_0,y_0)$, with $y_1=P(\gamma)y_0$. The formulas
for $u^*$ and $\alpha_*$ are those on cochains.

For (2): the group of $k$-cochains with coefficients in $F_j$ is the product, over the chains $\sigma$
of length $k$, of the groups $F_j(A_k)$. A product of exact sequences of abelian groups is exact, so
$\alpha$ and $\beta$ give a short exact sequence of cochain complexes, and its long exact cohomology
sequence is the sequence of (2) \cite[Tag 0117]{Sta}. The maps $u^*$ on cochains commute with the
differentials and with $\alpha$ and $\beta$, so they form a map of short exact sequences of complexes,
and $u^*\circ\delta=\delta\circ u^*$ because the long exact sequence is functorial in the short exact
sequence \cite[Tag 0117]{Sta}.

For (3): for a test algebra $B$ and $m\geq0$, let $L_m(B)$ be the free abelian group on the pairs
$(\sigma,\gamma)$ of a chain $\sigma=((A_0,y_0)\to\cdots\to(A_m,y_m))$ in $\mathrm{El}(P)$ and a map
$\gamma:A_m\to B$. A map $B\to B'$ acts by composition with $\gamma$, which makes $L_m$ a coefficient
system. Let $d:L_{m+1}\to L_m$ send $(\sigma,\gamma)$ to
$\sum_{i=0}^m(-1)^i(\partial_i\sigma,\gamma)+(-1)^{m+1}(\partial_{m+1}\sigma,\gamma\circ\gamma')$, where
$\gamma':A_m\to A_{m+1}$ is the last map of $\sigma$. A map $f:L_m\to K$ of coefficient systems is
determined by the elements $f(\sigma,\mathrm{id}_{A_m})\in K(A_m)$, since $(\sigma,\gamma)$ is the image
of $(\sigma,\mathrm{id}_{A_m})$ under $L_m(\gamma)$, and these elements form an arbitrary $m$-cochain with
coefficients in $K$. Under this identification, composition with $d$ is the differential of cochains.
So the complex of cochains with coefficients in $K$ is the complex of maps $L_\bullet\to K$.

The complex $L_\bullet$ is exact at $L_m$ for $m\geq1$. Let $h:L_m(B)\to L_{m+1}(B)$ send
$(\sigma,\gamma)$ to $(\sigma',\mathrm{id}_B)$, where $\sigma'$ extends $\sigma$ by the map
$\gamma:(A_m,y_m)\to(B,P(\gamma)y_m)$. Omitting the last object of $\sigma'$ gives back
$(\sigma,\gamma)$, and for $i\leq m$ omitting the $i$-th object commutes with the extension. Hence
$d(hx)=h(dx)+(-1)^{m+1}x$ for $x\in L_m(B)$ and $m\geq1$, and a cycle $x$ is the boundary of
$(-1)^{m+1}hx$.

Now let $m\geq1$ and let $f:L_m\to K$ be a cocycle, so $f\circ d=0$ on $L_{m+1}$. Then $f$ vanishes on
the image of $d:L_{m+1}\to L_m$, which is the kernel of $d:L_m\to L_{m-1}$, so $f=f'\circ d$ for a map
$f'$ from the image of $d:L_m\to L_{m-1}$ to $K$. Since $K$ is injective, $f'$ extends to a map
$g:L_{m-1}\to K$, and $f=g\circ d$ is a coboundary.
\end{proof}
\leanref{Development/AutomorphicArithmetic/CommutativeAlgebra/EtaleCohomology/VanishOfExact.lean}{UpperTheory.Ru.crystallineComparison.Th.L}
\leanref{Development/AutomorphicArithmetic/CommutativeAlgebra/EtaleCohomology/VanishOfExact.lean}{UpperTheory.Ru.crystallineComparison.Th.dL}
\leanref{Development/AutomorphicArithmetic/CommutativeAlgebra/EtaleCohomology/GFaceSnGCastSucc.lean}{UpperTheory.Ru.crystallineComparison.Th.exact_dL}
\leanref{Development/AutomorphicArithmetic/CommutativeAlgebra/EtaleCohomology/VanishOfExact.lean}{UpperTheory.Ru.crystallineComparison.Th.vanish_of_exact}

\begin{lemma}\label{gal-shifting-embedding}
Let $\ell$ be a prime (Definition~\ref{gal-shifting-coefficient-systems}).
\begin{enumerate}
\item For every $\nu\geq1$ the coefficient system $A\mapsto A_{\mathrm{cris}}(A)/\ell^\nu$ carries
  decompositions to products.
\item Let $0\to F_1\to F_2\to F_3\to0$ be an exact sequence of coefficient systems. If $F_1$ and $F_2$
  carry decompositions to products, then so does $F_3$.
\item Let $F$ be a coefficient system that carries decompositions to products. There are an injective
  object $K$ of the category of coefficient systems that carries decompositions to products and a map
  $F\to K$ that is injective at every test algebra.
\end{enumerate}
\end{lemma}
\leanref{Development/AutomorphicArithmetic/CommutativeAlgebra/EtaleCohomology/TestFacts.lean}{UpperTheory.Ru.crystallineComparison.AcrAdd}
\leanref{Development/AutomorphicArithmetic/CommutativeAlgebra/EtaleCohomology/CrystallineAlgebraAlgebraMap.lean}{UpperTheory.Ru.crystallineComparison.acrAdd}
\leanref{Development/AutomorphicArithmetic/CommutativeAlgebra/EtaleCohomology/ZeroPull.lean}{UpperTheory.Ru.crystallineComparison.QuotAdd}
\leanref{Development/AutomorphicArithmetic/CommutativeAlgebra/EtaleCohomology/EvalExact.lean}{UpperTheory.Ru.crystallineComparison.quotAdd}
\leanref{Development/AutomorphicArithmetic/CommutativeAlgebra/EtaleCohomology/ZeroPull.lean}{UpperTheory.Ru.crystallineComparison.Emb}
\leanref{Development/AutomorphicArithmetic/CommutativeAlgebra/EtaleCohomology/AddIff.lean}{UpperTheory.Ru.crystallineComparison.emb}

\begin{proof}
For (1), let $(A\to A_s)$ be a decomposition. Witt vectors are formed componentwise, so
$W(A)\to\prod_sW(A_s)$ is bijective, and so is $W(A)[1/\ell]\to\prod_sW(A_s)[1/\ell]$, since the
product is finite. This ring isomorphism sends $\xi$ to $(\xi)_s$, so it maps $A_{\mathrm{cris}}(A)$
injectively into $\prod_sA_{\mathrm{cris}}(A_s)$. The images in $A_{\mathrm{cris}}(A_s)$ of elements of
$A_{\mathrm{cris}}(A)$ form a subring that contains $W(A_s)$ and the $\xi^i/i!$, hence all of
$A_{\mathrm{cris}}(A_s)$. Let $E_s\in W(A)$ map to $1$ in $W(A_s)$ and to $0$ in $W(A_t)$ for $t\neq s$.
If $z_s\in A_{\mathrm{cris}}(A)$ maps to $x_s\in A_{\mathrm{cris}}(A_s)$, then $\sum_sE_sz_s$ maps to
$(x_s)_s$. So $A_{\mathrm{cris}}(A)\to\prod_sA_{\mathrm{cris}}(A_s)$ is a ring isomorphism, and so is its
reduction modulo $\ell^\nu$.

For (2), the maps into the products give a map from the exact sequence
$0\to F_1(A)\to F_2(A)\to F_3(A)\to0$ to the exact sequence of the groups $\prod_sF_j(A_s)$. Its first
two components are bijective, so the third is bijective by the five lemma.

For (3), take test algebras with underlying sets in a fixed universe and coefficient systems with values
in abelian groups of a larger universe, so that test algebras form a small category. Decompositions,
read as families of maps in the opposite category, are the coverings of a site \cite[Tag 00VH]{Sta}. A
single isomorphism is a decomposition, and the composites of a decomposition $(A\to A_s)$ with
decompositions of the $A_s$ form a decomposition of $A$. For base change, let $(A\to A_s)$ be a
decomposition and let $e_s\in A$ map to $1$ in $A_s$ and to $0$ in $A_t$ for $t\neq s$. The $e_s$ are
orthogonal idempotents with sum $1$, and $A\to A_s$ is surjective with kernel $(1-e_s)A$. For a map
$\gamma:A\to A'$ of test algebras, Lemma~\ref{gal-crystalline-test-facts}(4) gives a decomposition
$(A'\to A'_s)$ with $\gamma(e_s)\mapsto1$, and then $A'_s=A'/(1-\gamma(e_s))A'=A'\otimes_AA_s$. In
particular $A_s\otimes_AA_t$ is $A_s$ for $s=t$ and the zero ring for $s\neq t$. The zero ring has the
empty family as a decomposition, so both a sheaf on this site and a coefficient system that carries
decompositions to products vanish on it. Hence the sheaf condition for a decomposition $(A\to A_s)$
says that $F(A)\to\prod_sF(A_s)$ is bijective, and the abelian sheaves on this site are the coefficient
systems that carry decompositions to products.

Abelian sheaves on a site have enough injectives \cite[Tag 01DP]{Sta}, so $F$ embeds into an injective
abelian sheaf $K$. Sheafification is exact and left adjoint to the inclusion of abelian sheaves into
abelian presheaves \cite[Tag 01DL]{Sta}, and abelian presheaves on the opposite category of test
algebras are coefficient systems. Hence the inclusion carries $K$ to an injective coefficient system
\cite[Tag 015Z]{Sta}, and, as a right adjoint, it carries the monomorphism $F\to K$ to a monomorphism of
coefficient systems, which is injective at every test algebra.
\end{proof}
\leanref{Development/AutomorphicArithmetic/CommutativeAlgebra/EtaleCohomology/TestFacts.lean}{UpperTheory.Ru.crystallineComparison.PartsEx}
\leanref{Development/AutomorphicArithmetic/CommutativeAlgebra/EtaleCohomology/PullbackDecomposition.lean}{UpperTheory.Ru.crystallineComparison.partsEx}

\begin{lemma}\label{gal-crystalline-test-cohomology}
Let $P$, $P'$ and $Q$ be functors from test algebras to sets, $k\geq0$ and $\nu\geq1$
(Definition~\ref{gal-crystalline-test-algebras}).
\begin{enumerate}
\item For natural transformations $v:Q\to P$ and $u:P\to P'$, $(u\circ v)^*=v^*\circ u^*$, and $u^*$ is
  $A_{\mathrm{cris}}$-linear and commutes with $\phi$.
\item If $u_i:Q\to P$ and $a_i\in\mathbf Z$, $i\in I$, are such that $\sum_ia_iu_i$ vanishes locally,
  then $\sum_ia_iu_i^*=0$ on $H^k(P,A_{\mathrm{cris}}/\ell^\nu)$.
\item Let $p:Q\to P$ and $g_1,\ldots,g_h:Q\to Q$ be natural transformations with $p\circ g_i=p$.
  Suppose that $Q(A)\to P(A)$ is surjective for every test algebra $A$, and that for $y,y'\in Q(A)$
  with $p(y)=p(y')$ there is a decomposition $(A\to A_s)$ of $A$ such that the multisets
  $\{g_i(y_s)\}_i$ and $\{g_i(y'_s)\}_i$ agree for every $s$. If
  $x\in H^k(P,A_{\mathrm{cris}}/\ell^\nu)$ and $p^*x=0$, then $hx=0$.
\end{enumerate}
\end{lemma}
\leanref{Development/AutomorphicArithmetic/CommutativeAlgebra/EtaleCohomology/TestFacts.lean}{UpperTheory.Ru.crystallineComparison.FrEx}
\leanref{Development/AutomorphicArithmetic/CommutativeAlgebra/EtaleCohomology/NatCastQb.lean}{UpperTheory.Ru.crystallineComparison.frEx}
\leanref{Development/AutomorphicArithmetic/CommutativeAlgebra/EtaleCohomology/TestFacts.lean}{UpperTheory.Ru.crystallineComparison.SumMH}
\leanref{Development/AutomorphicArithmetic/CommutativeAlgebra/EtaleCohomology/SumTransportAlt.lean}{UpperTheory.Ru.crystallineComparison.sumMH}
\leanref{Development/AutomorphicArithmetic/CommutativeAlgebra/EtaleCohomology/TestFacts.lean}{UpperTheory.Ru.crystallineComparison.TTrans}
\leanref{Development/AutomorphicArithmetic/CommutativeAlgebra/EtaleCohomology/TestFacts.lean}{UpperTheory.Ru.crystallineComparison.AvgMH}
\leanref{Development/AutomorphicArithmetic/CommutativeAlgebra/EtaleCohomology/AveragingTransferKills.lean}{UpperTheory.Ru.crystallineComparison.avgMH}

\begin{proof}
Item (1) holds on cochains, since $u^*$ is induced by the functor $\mathrm{El}(P)\to\mathrm{El}(P')$ and
$\phi$ acts on the coefficients.

Items (2) and (3) hold with $H^k(P,F)$ in place of $H^k(P,A_{\mathrm{cris}}/\ell^\nu)$ for every
coefficient system $F$ that carries decompositions to products
(Definition~\ref{gal-shifting-coefficient-systems}), and $A\mapsto A_{\mathrm{cris}}(A)/\ell^\nu$ is one
by Lemma~\ref{gal-shifting-embedding}(1). We prove this by induction on $k$, for all such $F$ at once.
Such an $F$ vanishes on the zero ring, which has the empty family as a decomposition.

The induction step from $k$ to $k+1$ uses the following. Lemma~\ref{gal-shifting-embedding}(3) gives an
injective coefficient system $K$ that carries decompositions to products and a map $F\to K$ that is
injective at every test algebra. Let $\pi:K\to L$ be its cokernel; $L$ carries decompositions to
products by Lemma~\ref{gal-shifting-embedding}(2). Since $H^{k+1}(P,K)=0$
(Lemma~\ref{gal-shifting-cohomology}(3)), the long exact sequence of
Lemma~\ref{gal-shifting-cohomology}(2) shows that every $x\in H^{k+1}(P,F)$ is $\delta w$ for some
$w\in H^k(P,L)$.

Item (2) in degree $0$. By Lemma~\ref{gal-shifting-cohomology}(1), $x\in H^0(P,F)$ is a compatible family
$c$, and $\sum_ia_iu_i^*x$ is the family $(A,y)\mapsto\sum_ia_ic(A,u_i(y))$ on $\mathrm{El}(Q)$. Let $A$
be a test algebra, $y\in Q(A)$, and $(A\to A_s)$ a decomposition as in
Definition~\ref{gal-crystalline-test-algebras}(3). Since $c$ is compatible and the $u_i$ are natural,
the image of $\sum_ia_ic(A,u_i(y))$ in $F(A_s)$ is $\sum_ia_ic(A_s,u_i(y_s))$. This is the image of
$\sum_ia_i[u_i(y_s)]$ under the additive map from the free abelian group on $P(A_s)$ to $F(A_s)$ that
sends $[z]$ to $c(A_s,z)$, so it vanishes if $A_s\neq0$; if $A_s=0$, then $F(A_s)=0$. Since
$F(A)\to\prod_sF(A_s)$ is injective, $\sum_ia_ic(A,u_i(y))=0$.

Item (2) in degree $k+1$. Write $x=\delta w$ as above. By Lemma~\ref{gal-shifting-cohomology}(2),
$\sum_ia_iu_i^*x=\delta\bigl(\sum_ia_iu_i^*w\bigr)$, and $\sum_ia_iu_i^*w=0$ by the induction hypothesis
for $L$.

Item (3) in degree $0$. Let $x\in H^0(P,F)$ be the family $c$, with $p^*x=0$. Then $c(A,p(y))=0$ for
every test algebra $A$ and $y\in Q(A)$. Since $Q(A)\to P(A)$ is surjective, $x=0$.

Item (3) in degree $k+1$. Let $p^*x=0$ and write $x=\delta w$ as above. Then $\delta(p^*w)=p^*x=0$, so
$p^*w=\pi_*v$ for some $v\in H^k(Q,K)$ by exactness. It suffices to show that $hw$ lies in the image of
$\pi_*:H^k(P,K)\to H^k(P,L)$, since then $hx=\delta(hw)=0$ by exactness. If $k\geq1$, then
$H^k(Q,K)=0$, so $p^*w=0$, and $hw=0$ by the induction hypothesis for $L$. If $k=0$, average. Let $G$
be a coefficient system that carries decompositions to products and $c\in H^0(Q,G)$. For $y\in P(A)$
put $\operatorname{av}(c)(A,y)=\sum_ic(A,g_i(\tilde y))$, where $\tilde y\in Q(A)$ is any element with
$p(\tilde y)=y$. This does not depend on $\tilde y$: for two choices, take a decomposition $(A\to A_s)$
on whose pieces the multisets $\{g_i(\tilde y_s)\}_i$ agree; the two sums have the same image in every
$G(A_s)$, and $G(A)\to\prod_sG(A_s)$ is injective. For a map $\gamma:A\to B$, $Q(\gamma)\tilde y$ lies
over $P(\gamma)y$, so $G(\gamma)\operatorname{av}(c)(A,y)=\operatorname{av}(c)(B,P(\gamma)y)$, and
$\operatorname{av}(c)\in H^0(P,G)$. Since $p\circ g_i=p$, $\operatorname{av}(p^*z)=hz$ for
$z\in H^0(P,G)$, and $\operatorname{av}(\pi\circ c)=\pi\circ\operatorname{av}(c)$ for $c\in H^0(Q,K)$.
Hence $hw=\operatorname{av}(p^*w)=\operatorname{av}(\pi_*v)=\pi_*\operatorname{av}(v)$.
\end{proof}
\leanref{Development/AutomorphicArithmetic/CommutativeAlgebra/EtaleCohomology/ZeroPull.lean}{UpperTheory.Ru.crystallineComparison.Parts.LocZero}
\leanref{Development/AutomorphicArithmetic/CommutativeAlgebra/EtaleCohomology/ZeroPull.lean}{UpperTheory.Ru.crystallineComparison.LocSum}
\leanref{Development/AutomorphicArithmetic/CommutativeAlgebra/EtaleCohomology/SumSecEqZero.lean}{UpperTheory.Ru.crystallineComparison.locSum}
\leanref{Development/AutomorphicArithmetic/CommutativeAlgebra/EtaleCohomology/ZeroPull.lean}{UpperTheory.Ru.crystallineComparison.Parts.LocTrans}
\leanref{Development/AutomorphicArithmetic/CommutativeAlgebra/EtaleCohomology/ZeroPull.lean}{UpperTheory.Ru.crystallineComparison.LocAvg}
\leanref{Development/AutomorphicArithmetic/CommutativeAlgebra/EtaleCohomology/SumTransport.lean}{UpperTheory.Ru.crystallineComparison.locAvg}

\subsection{Comparison maps, period data and Frobenius data}\label{subsection-gal-comparison-maps}

\begin{definition}\label{gal-realization-split-diagrams}
Let $\mathcal C$ be a category.
\begin{enumerate}
\item For a functor $F$ from $\mathcal C$ to abelian groups, $H^k(\mathcal C,F)$ is the $k$-th cohomology
  of the complex whose $k$-cochains assign to every chain $c_0\to\cdots\to c_k$ in $\mathcal C$ an element
  of $F(c_k)$, with the alternating sum of the face maps as differential. A functor $m:\mathcal
  C'\to\mathcal C$ and homomorphisms $F(m(c))\to F'(c)$, natural in $c\in\mathcal C'$, induce
  $m^*:H^k(\mathcal C,F)\to H^k(\mathcal C',F')$. If $F$ takes values in commutative rings, the cup
  product of cochains, $(a\cup b)(c_0\to\cdots\to c_{i+j})=a(c_0\to\cdots\to c_i)\,b(c_i\to\cdots\to
  c_{i+j})$ with the first factor carried to $F(c_{i+j})$, makes $\bigoplus_kH^k(\mathcal C,F)$ a graded
  ring. In Definition~\ref{gal-crystalline-test-algebras}, $H^k(P,A_{\mathrm{cris}}/\ell^\nu)$ is
  $H^k(\mathrm{El}(P),F)$ for $F(A,y)=A_{\mathrm{cris}}(A)/\ell^\nu$.
\item A \emph{split diagram} over $\mathcal C$ consists of schemes $U_c$ for $c\in\mathcal C$ and morphisms
  $U_\varphi:U_{c'}\to U_c$ for $\varphi:c\to c'$, with $U_{\mathrm{id}}=\mathrm{id}$ and
  $U_{\psi\varphi}=U_\varphi U_\psi$, such that every \'etale surjective morphism onto any $U_c$ has a
  section. \emph{Points} of it in a scheme $X$ are morphisms $y_c:U_c\to X$ with $y_cU_\varphi=y_{c'}$
  for every $\varphi:c\to c'$. For $q\geq1$ its \emph{coefficient system} $L_q$ sends $c$ to the ring of
  locally constant functions $U_c\to\mathbf Z/q$, and $\varphi$ to pullback along $U_\varphi$.
\item Let $\ell$ be a prime, $C$ and $A_{\mathrm{cris}}$ as in Definition~\ref{gal-crystalline-test-algebras},
  $K$ a finite extension of $\mathbf Q_\ell$ in $C$ with valuation ring $\mathcal O_K$, $\bar K$ its
  algebraic closure in $C$ and $G_K=\operatorname{Gal}(\bar K/K)$, which acts on $C$, on
  $\mathcal O_C^\flat$ and on $A_{\mathrm{cris}}$. For $\sigma\in G_K$ and a test algebra $A$ let
  $A^\sigma$ be the test algebra with the same ring and with its map from $\mathcal O_C^\flat$ composed
  with $\sigma$. For an $\mathcal O_K$-scheme $X$, with $X^\sharp$ formed for the inclusion
  $\mathcal O_K\subset\mathcal O_C$, transport of structure along $A\mapsto A^\sigma$ gives an action of
  $G_K$ on $H^k(X^\sharp,A_{\mathrm{cris}}/\ell^\nu)$, semilinear over its action on $A_{\mathrm{cris}}$.
\end{enumerate}
\end{definition}
\leanref{Development/WeilBound/CoreStatements/PadicHodgeComparison/EtaleCohomologyComparison/CechPairings/InstAddCommGroup.lean}{CrystallineSemistableComparison.CH.H}
\leanref{Development/WeilBound/CoreStatements/PadicHodgeComparison/EtaleCohomologyComparison/CechPairings/InstAddCommGroup.lean}{CrystallineSemistableComparison.CH.Mor.pull}
\leanref{Development/WeilBound/CoreStatements/PadicHodgeComparison/EtaleCohomologyComparison/CechPairings/InstAddCommGroup.lean}{CrystallineSemistableComparison.CH.Pairing.cupC}
\leanref{Development/WeilBound/CoreStatements/PadicHodgeComparison/EtaleCohomologyComparison/KummerCupProducts/IsKumCoc.lean}{CrystallineSemistableComparison.E.Splits}
\leanref{Development/WeilBound/CoreStatements/PadicHodgeComparison/EtaleCohomologyComparison/KummerCupProducts/IsKumCoc.lean}{CrystallineSemistableComparison.E.Probe}
\leanref{Development/WeilBound/CoreStatements/PadicHodgeComparison/EtaleCohomologyComparison/KummerCupProducts/IsKumCoc.lean}{CrystallineSemistableComparison.E.Probe.Pts}
\leanref{Development/WeilBound/CoreStatements/PadicHodgeComparison/EtaleCohomologyComparison/KummerCupProducts/IsKumCoc.lean}{CrystallineSemistableComparison.E.lcF}
\leanref{Development/WeilBound/CoreStatements/PadicHodgeComparison/EtaleCohomologyComparison/AcrisKummerUnits/GaloisTwistIdentitySection.lean}{CrystallineSemistableComparison.PAlg.tw}
\leanref{Development/WeilBound/CoreStatements/PadicHodgeComparison/EtaleCohomologyComparison/AcrisKummerUnits/GalOCOkToOC.lean}{CrystallineSemistableComparison.Model.galMH}

\begin{lemma}\label{gal-realization-edge-map}
Let $\mathcal C$ be a category, $(U_c)$ a split diagram over $\mathcal C$ with points $(y_c)$ in a scheme
$X$, $q\geq1$ and $j\geq0$ (Definition~\ref{gal-realization-split-diagrams}). Let $\mathbf Z/q\to
I^\bullet$ be an injective resolution of the constant sheaf on the big \'etale site of schemes, so that
$H^j(X,\mathbf Z/q)$ is the cohomology of $\Gamma(X,I^\bullet)$. For a cochain $b$ of $\mathcal C$ with
values in $c\mapsto\Gamma(U_c,I^s)$, write $db$ for the differential of $I^\bullet$ applied to its values
and $\delta b$ for its differential as a cochain, and view $L_q(c)=\Gamma(U_c,\mathbf Z/q)$ inside
$\Gamma(U_c,I^0)$.
\begin{enumerate}
\item There is a unique homomorphism $\varepsilon_y:H^j(X,\mathbf Z/q)\to H^j(\mathcal C,L_q)$ with the
  following property. Let $a\in\Gamma(X,I^j)$ with $da=0$ represent $x$, let $b_k$, $0\le k<j$, be
  $k$-cochains with values in $c\mapsto\Gamma(U_c,I^{j-1-k})$ and $z$ a $j$-cocycle with values in $L_q$,
  such that $db_0=(y_c^*a)_c$, $db_k=(-1)^{k-1}\delta b_{k-1}$ for $1\le k<j$, and
  $z=(-1)^{j-1}\delta b_{j-1}$ (for $j=0$: $z=(y_c^*a)_c$). Then $\varepsilon_y(x)$ is the class of $z$.
  Such $b_k$ and $z$ exist for every representative $a$ of every $x$.
\item Let $(U'_c)$ be a split diagram over a category $\mathcal C'$ with points $(y'_c)$ in a scheme $X'$,
  $m:\mathcal C'\to\mathcal C$ a functor, $g_c:U'_c\to U_{m(c)}$ morphisms with
  $U_{m(\varphi)}g_{c'}=g_cU'_\varphi$ for $\varphi:c\to c'$, and $u:X'\to X$ a morphism with
  $uy'_c=y_{m(c)}g_c$ for all $c$. Let $m^*:H^j(\mathcal C,L_q)\to H^j(\mathcal C',L'_q)$ be induced by $m$
  and the pullbacks along the $g_c$. Then $m^*\circ\varepsilon_y=\varepsilon_{y'}\circ u^*$.
\item Let $q\mid Q$, and let $(U'_c)$, $\mathcal C'$, $(y'_c)$, $m$ and $g_c$ be as in (2), with $X'=X$ and $u$ the
  identity, so that $y'_c=y_{m(c)}g_c$. Let $m^*:H^j(\mathcal C,L_Q)\to H^j(\mathcal C',L'_q)$ be induced by $m$ and
  the maps $L_Q(m(c))\to L'_q(c)$, $f\mapsto g_c^*f$ reduced modulo $q$, and $\mathrm{red}:H^j(X,\mathbf Z/Q)\to
  H^j(X,\mathbf Z/q)$ the reduction. Then $m^*\circ\varepsilon_y=\varepsilon_{y'}\circ\mathrm{red}$, where
  $\varepsilon_y$ is formed for $Q$ and $\varepsilon_{y'}$ for $q$. For $\mathcal C'=\mathcal C$, $m$ and the $g_c$ the
  identities, this says that $\varepsilon_y$ commutes with the reduction maps $H^j(X,\mathbf Z/Q)\to H^j(X,\mathbf Z/q)$
  and $H^j(\mathcal C,L_Q)\to H^j(\mathcal C,L_q)$.
\end{enumerate}
\end{lemma}
\leanref{Development/WeilBound/CoreStatements/PadicHodgeComparison/EtaleCohomologyComparison/KummerCupProducts/IsKumCoc.lean}{CrystallineSemistableComparison.E.Probe.Pts.et}
\leanref{Development/WeilBound/CoreStatements/PadicHodgeComparison/EtaleCohomologyComparison/KummerCupProducts/EtaleChernClassEqKummerClass.lean}{CrystallineSemistableComparison.E.Main.resSplit}
\leanref{Development/WeilBound/CoreStatements/PadicHodgeComparison/EtaleCohomologyComparison/KummerCupProducts/EtaleChernClassEqKummerClass.lean}{CrystallineSemistableComparison.E.Main.resGood}
\leanref{Development/WeilBound/CoreStatements/PadicHodgeComparison/EtaleCohomologyComparison/KummerCupProducts/EtaleChernClassEqKummerClass.lean}{CrystallineSemistableComparison.E.Main.etSpec}
\leanref{Development/WeilBound/CoreStatements/PadicHodgeComparison/EtaleCohomologyComparison/KummerCupProducts/EtaleChernClassEqKummerClass.lean}{CrystallineSemistableComparison.E.Main.etNat}
\leanref{Development/WeilBound/CoreStatements/PadicHodgeComparison/EtaleCohomologyComparison/KummerCupProducts/EtaleChernClassEqKummerClass.lean}{CrystallineSemistableComparison.E.Main.etRed}
\leanref{Development/WeilBound/CoreStatements/PadicHodgeComparison/EtaleCohomologyComparison/KummerCupProducts/EtaleChernClassEqKummerClass.lean}{CrystallineSemistableComparison.E.Main.stEx}
\leanref{Development/WeilBound/CoreStatements/PadicHodgeComparison/EtaleCohomologyComparison/KummerCupProducts/EtaleChernClassEqKummerClass.lean}{CrystallineSemistableComparison.E.Main.stAdd}
\leanref{Development/WeilBound/CoreStatements/PadicHodgeComparison/EtaleCohomologyComparison/KummerCupProducts/EtaleChernClassEqKummerClass.lean}{CrystallineSemistableComparison.E.Main.stUniq}
\leanref{Development/WeilBound/CoreStatements/PadicHodgeComparison/EtaleCohomologyComparison/KummerCupProducts/EtaleChernClassEqKummerClass.lean}{CrystallineSemistableComparison.E.Main.edgeUniq}

\begin{proof}
For each $c$ the sequence $0\to L_q(c)\to\Gamma(U_c,I^0)\to\Gamma(U_c,I^1)\to\cdots$ is exact. Sections
of the constant sheaf $\mathbf Z/q$ over a scheme are the locally constant functions to $\mathbf Z/q$, so
it suffices to see that $\Gamma(U_c,-)$ keeps $\mathbf Z/q\to I^\bullet$ exact. A section $x$ of $I^{s+1}$
over $U_c$ with $dx=0$ is $d$ of a section on each member of some \'etale covering of $U_c$, since
$I^\bullet$ is exact. The disjoint union of the members is \'etale and surjective over $U_c$, so it has a
section. Pulling back along it writes $U_c$ as a disjoint union of open subschemes on each of which $x$
is a coboundary. These preimages glue, so $x$ is a coboundary on $U_c$. Exactness at $L_q(c)$ and at
$\Gamma(U_c,I^0)$ holds over every scheme, since $\Gamma(U_c,-)$ is left exact.

(1) Let $T^{k,s}$ be the group of $k$-cochains with values in $c\mapsto\Gamma(U_c,I^s)$, so that the
$k$-cochains with values in $L_q$ form a subgroup $C^k$ of $T^{k,0}$. Applying the exact sequences above
value by value, the sequence $0\to C^k\to T^{k,0}\xrightarrow{d}T^{k,1}\xrightarrow{d}\cdots$ is exact
for every $k$. The maps along the $U_\varphi$ commute with $d$ and with the inclusion of $L_q$, so
$d\delta=\delta d$, and $\delta$ on $T^{k,0}$ restricts to the differential of $C^\bullet$.

For $s+k=j$, $w\in T^{k,s}$ and $z\in C^j$, a \emph{staircase} from $w$ to $z$ is defined by induction on
$s$: for $s=0$ it is the equality $w=z$; for $s\geq1$ it is an element $e\in T^{k,s-1}$ with $de=w$
together with a staircase from $(-1)^k\delta e$ to $z$. The conditions on the $b_k$ and $z$ in (1) say
that $b_0,\ldots,b_{j-1}$ is a staircase from $(y_c^*a)_c\in T^{0,j}$ to $z$. Three facts are proved by
induction on $s$.
\begin{enumerate}
\item[(i)] If $w\in T^{k,s}$ satisfies $dw=0$ and $\delta w=0$, there is a staircase from $w$ to some
  $z\in C^j$ with $\delta z=0$. For $s=0$, $w$ lies in $C^k$ by exactness, and $\delta w=0$. For
  $s\geq1$, exactness gives $e\in T^{k,s-1}$ with $de=w$; then $w'=(-1)^k\delta e$ satisfies
  $dw'=(-1)^k\delta w=0$ and $\delta w'=0$, and the induction hypothesis applies to $w'$.
\item[(ii)] Staircases from $w$ to $z$ and from $w'$ to $z'$ give, by adding or subtracting their
  elements $e$ step by step, staircases from $w+w'$ to $z+z'$ and from $w-w'$ to $z-z'$.
\item[(iii)] Suppose there is a staircase from $w\in T^{k,s}$ to $z$, and that $w=0$ if $k=0$, and
  $w=\delta b$ for some $b\in T^{k-1,s}$ with $db=0$ if $k\geq1$. Then $z$ is a coboundary in $C^\bullet$.
  For $s=0$ and $k=0$, $z=w=0$. For $s=0$ and $k\geq1$, $b$ lies in $C^{k-1}$ by exactness, and
  $z=w=\delta b$. For $s\geq1$, let $e$ be the first step, so that $de=w$ and the rest is a staircase from
  $w'=(-1)^k\delta e$ to $z$. If $k=0$, then $de=0$ and $w'=\delta e$, so the hypothesis holds for $w'$
  with $b=e$. If $k\geq1$, exactness gives $b'\in T^{k-1,s-1}$ with $db'=b$, and
  $b''=(-1)^k(e-\delta b')$ satisfies $db''=(-1)^k(w-\delta b)=0$ and $\delta b''=w'$. In both cases the
  induction hypothesis shows that $z$ is a coboundary.
\end{enumerate}

Let $a$ represent $x$. The $0$-cochain $w=(y_c^*a)_c$ satisfies $dw=(y_c^*da)_c=0$, and $\delta w=0$
because $y_cU_\varphi=y_{c'}$. By (i) there is a staircase from $w$ to a cocycle $z$, which gives the
$b_k$ and $z$ of (1). Let $a'$ be another representative of $x$, with a staircase $b'_0,\ldots,b'_{j-1}$
from $(y_c^*a')_c$ to $z'$. By (ii) the $b_k-b'_k$ form a staircase from $(y_c^*(a-a'))_c$ to $z-z'$. If
$j=0$, then $a=a'$, since $H^0(X,\mathbf Z/q)$ is the kernel of $d$ on $\Gamma(X,I^0)$, and $z=z'$. If
$j\geq1$, then $a-a'=da''$ for some $a''\in\Gamma(X,I^{j-1})$, and $g=(y_c^*a'')_c\in T^{0,j-1}$ satisfies
$dg=(y_c^*(a-a'))_c$ and $\delta g=0$. The first step of the staircase is $b_0-b'_0$, and the rest is a
staircase from $\delta(b_0-b'_0)=\delta(b_0-b'_0-g)$ to $z-z'$, where $d(b_0-b'_0-g)=0$. By (iii) for
$(s,k)=(j-1,1)$, $z-z'$ is a coboundary. So the class of $z$ depends only on $x$; call it
$\varepsilon_y(x)$. If $a$ and $a'$ represent $x$ and $x'$, then $a+a'$ represents $x+x'$, and (ii) shows
that $\varepsilon_y$ is additive. It is the unique homomorphism with the property of (1).

(2) Pullback along $u$ carries a representative $a$ of $x$ to the representative $u^*a$ of $u^*x$. For
a cochain $b$ of $\mathcal C$ put $(m^*b)(c_0\to\cdots\to c_k)=g_{c_k}^*b(m(c_0)\to\cdots\to m(c_k))$;
this commutes with $d$ and with $\delta$, and $m^*(y_c^*a)_c=(y'^*_cu^*a)_c$ since $uy'_c=y_{m(c)}g_c$.
So the $m^*b_k$ and $m^*z$ satisfy the conditions of (1) for $u^*a$ and the points $y'$.

(3) The reduction $\mathbf Z/Q\to\mathbf Z/q$ extends to a map $\rho$ from the injective resolution $I^\bullet_Q$ of
$\mathbf Z/Q$ to the injective resolution $I^\bullet_q$ of $\mathbf Z/q$, and $\rho$ restricts to the reduction
$L_Q\to L_q$ on $\Gamma(U_c,\mathbf Z/Q)\subset\Gamma(U_c,I^0_Q)$. If $a$ represents $x$, then $\rho(a)$ represents
$\mathrm{red}(x)$. Let $b_k$ and $z$ be as in (1) for $a$, and put $(m^*b)(c_0\to\cdots\to c_k)=g_{c_k}^*\rho\bigl(b(m(c_0)
\to\cdots\to m(c_k))\bigr)$ for a cochain $b$ of $\mathcal C$. As in (2), $m^*$ commutes with $d$ and with $\delta$, and
$m^*(y_c^*a)_c=(y'^*_c\rho(a))_c$, as $y'_c=y_{m(c)}g_c$ and $\rho$ commutes with pullback. So the $m^*b_k$ and $m^*z$
satisfy the conditions of (1) for $\rho(a)$ and the points $y'$, and $m^*z$ represents $m^*\varepsilon_y(x)$.
% Lean: E.Main.etRed = EdgeRed.etRed: a representative a (repBasic), the staircase for Q (etSpec), its image under the
% morphism of resolved diagrams redResMor (EdgeMap), the reduction redMap Q q of the representative (repOp), and the
% naturality of redMap for the points.
\end{proof}

\begin{remark}\label{gal-realization-edge-assembly}
In the notation of the proof of Lemma~\ref{gal-realization-edge-map}, give the double complex
$T^{\bullet,\bullet}$ the total differential $\delta+(-1)^kd$ on $T^{k,s}$. The conditions on the $b_k$
and $z$ in Lemma~\ref{gal-realization-edge-map}(1) say that $(y_c^*a)_c-(-1)^jz$ is the total differential
of $(b_0,\ldots,b_{j-1})$. The sequences $0\to C^k\to T^{k,0}\to T^{k,1}\to\cdots$ are exact, so $C^\bullet$
is a subcomplex of the total complex with the same cohomology, by the acyclic assembly lemma
\cite[Lemma~2.7.3]{Wei94}. Hence $\varepsilon_y$ is $(-1)^j$ times the map $x\mapsto[(y_c^*a)_c]$ into the
cohomology of the total complex followed by the inverse of this isomorphism. This shows again that the
class of $z$ depends only on $x$ and is additive in $x$, without the facts (ii) and (iii) of the proof.
\end{remark}

\begin{lemma}\label{gal-realization-locally-constant}
Let $A$ be a test algebra, $\nu\geq1$, and $f$ a locally constant function on
$\operatorname{Spec}A^\sharp[1/\ell]$ with values in $\mathbf Z/\ell^\nu$
(Definition~\ref{gal-crystalline-test-algebras}).
\begin{enumerate}
\item There are unique orthogonal idempotents $e_a\in A$, $a\in\mathbf Z/\ell^\nu$, with sum $1$, such
  that for every $a$ the set $f^{-1}(a)$ consists of the primes of $A^\sharp[1/\ell]$ that do not
  contain $\theta([e_a])$.
\item Put $\lambda_A(f)=\sum_a\tilde a\,[e_a]\in A_{\mathrm{cris}}(A)/\ell^\nu$, with $\tilde a\in\mathbf Z$
  lifting $a$. Then $\lambda_A$ is a ring homomorphism and $\phi\circ\lambda_A=\lambda_A$. Let $B$ be a test algebra,
  $h:A\to B$ a ring map such that the elements $W(h)(\xi)^k/k!$ of $W(B)[1/\ell]$ lie in $A_{\mathrm{cris}}(B)$ for every
  $k\geq1$, so that $W(h)$ induces $h:A_{\mathrm{cris}}(A)/\ell^\nu\to A_{\mathrm{cris}}(B)/\ell^\nu$, and
  $g:A^\sharp[1/\ell]\to B^\sharp[1/\ell]$ a ring map with $g(\theta([e]))=\theta([h(e)])$ for every idempotent $e$ of $A$.
  Then $\lambda_B(f\circ\operatorname{Spec}g)=h(\lambda_A(f))$. A map $h:A\to B$ of test algebras satisfies this, with
  $W(h)(\xi)=\xi$ and $g$ induced by $h$. If $f$ takes values in $\mathbf Z/\ell^{\nu+1}$, the image of $\lambda_A(f)$ in
  $A_{\mathrm{cris}}(A)/\ell^\nu$ is $\lambda_A$ of the reduction of $f$ modulo $\ell^\nu$.
\end{enumerate}
\end{lemma}
\leanref{Development/WeilBound/CoreStatements/PadicHodgeComparison/EtaleCohomologyComparison/KummerCupProducts/MapUgCToUg.lean}{CrystallineSemistableComparison.E.lcA}
\leanref{Development/WeilBound/CoreStatements/PadicHodgeComparison/EtaleCohomologyComparison/KummerCupProducts/EtaleChernClassEqKummerClass.lean}{CrystallineSemistableComparison.E.Main.lcIdem}
\leanref{Development/WeilBound/CoreStatements/PadicHodgeComparison/EtaleCohomologyComparison/KummerCupProducts/EtaleChernClassEqKummerClass.lean}{CrystallineSemistableComparison.E.Main.lcSpec}
\leanref{Development/WeilBound/CoreStatements/PadicHodgeComparison/EtaleCohomologyComparison/KummerCupProducts/EtaleChernClassEqKummerClass.lean}{CrystallineSemistableComparison.E.Main.lcNat}
\leanref{Development/AutomorphicSpectrum/Proofs/BlockDecomposition/FrobIdem.lean}{GlobalTheory.CompanionTheory.FrobLc.phi_lcM}
\leanref{Development/WeilBound/CoreStatements/PadicHodgeComparison/EtaleCohomologyComparison/AcrisKummerUnits/ValuationOfValuationModSquare.lean}{CrystallineSemistableComparison.PAlg.IsTest.idem}

\begin{proof}
(1) The spectrum is quasi-compact, so $f$ takes finitely many values, and each $f^{-1}(a)$ is open and
closed. It is the set of primes not containing a unique idempotent $\epsilon_a$ of $A^\sharp[1/\ell]$, and
the $\epsilon_a$ are orthogonal with sum $1$. An idempotent of $A^\sharp[1/\ell]$ is integral over
$A^\sharp$, so it lies in $A^\sharp$ by Lemma~\ref{gal-crystalline-test-facts}(1), and it is
$\theta([e_a])$ for a unique idempotent $e_a$ of $A$ by Lemma~\ref{gal-crystalline-test-facts}(3). As
$\ell$ is a nonzerodivisor of $A^\sharp$, the map $A^\sharp\to A^\sharp[1/\ell]$ is injective, which gives
uniqueness. For orthogonal idempotents $e,e'$ of $A$ one has $[ee']=[e][e']$ and $[e+e']=[e]+[e']$ in
$W(A)$, since $W(A)$ is the product of the rings $W(eA)$, $W(e'A)$ and $W((1-e-e')A)$. So the
$e_ae_b$ for $a\neq b$ and $\sum_ae_a$ are idempotents with images $0$ and $1$, and uniqueness gives
$e_ae_b=0$ and $\sum_ae_a=1$.

(2) If $f$ and $f'$ have idempotents $e_a$ and $e'_b$, then $f+f'$ and $ff'$ have the idempotents
$\sum_{a+b=c}e_ae'_b$ and $\sum_{ab=c}e_ae'_b$. With the identities for Teichm\"uller representatives of
orthogonal idempotents, this gives $\lambda_A(f+f')=\lambda_A(f)+\lambda_A(f')$,
$\lambda_A(ff')=\lambda_A(f)\lambda_A(f')$, and $\lambda_A(1)=[1]=1$. The Frobenius of $W(A)$ sends
$[e]$ to $[e^\ell]=[e]$. The images $h(e_a)$ are orthogonal idempotents with sum $1$, and
$g(\theta([e_a]))=\theta([h(e_a)])$, so a prime $\mathfrak p$ of $B^\sharp[1/\ell]$ does not contain $\theta([h(e_a)])$ if
and only if $g^{-1}(\mathfrak p)$ does not contain $\theta([e_a])$, that is, if and only if $f(g^{-1}(\mathfrak p))=a$. So
the $h(e_a)$ are the idempotents of $f\circ\operatorname{Spec}g$, and $\lambda_B(f\circ\operatorname{Spec}g)=\sum_a\tilde
a[h(e_a)]=h(\lambda_A(f))$, as $W(h)([e_a])=[h(e_a)]$. For a map $h$ of test algebras, $W(h)$ fixes
$\xi=\ell-[\ell^\flat]$, as $h(\ell^\flat)=\ell^\flat$, and the map $g$ induced by $h$ commutes with $\theta$. The
idempotent of the
reduction of $f$ at $b\in\mathbf Z/\ell^\nu$ is the sum of the $e_a$ with $a$ reducing to $b$.
\end{proof}

\begin{definition}\label{gal-realization-eta}
Let $\ell$ and $C$ be as in Definition~\ref{gal-crystalline-test-algebras}, $\bar K$ the algebraic closure of
$\mathbf Q_\ell$ in $C$, $K$ a field with $\mathbf Q_\ell\subset K\subset\bar K$ (for example a finite extension of
$\mathbf Q_\ell$ in $C$, as in Definition~\ref{gal-realization-split-diagrams}(3); $\bar K$ is also the algebraic closure
of $K$ in $C$), $\mathcal O_K=K\cap\mathcal O_C$ its valuation ring, $X$ an $\mathcal O_K$-scheme, with $X^\sharp$ formed
for the inclusion $\mathcal O_K\subset\mathcal O_C$, $X_{\bar K}=X\times_{\mathcal O_K}\bar K$, and $\nu\geq1$. For an object
$(A,y)$ of $\mathrm{El}(X^\sharp)$ put $U_{(A,y)}=\operatorname{Spec}A^\sharp[1/\ell]$, and let
$y_{(A,y)}:U_{(A,y)}\to X_{\bar K}$ have the components $U_{(A,y)}\to\operatorname{Spec}A^\sharp\to X$,
the second map being $y$, and $U_{(A,y)}\to\operatorname{Spec}\bar K$, from $\bar K\subset
C=\mathcal O_C[1/\ell]\to A^\sharp[1/\ell]$. By Lemma~\ref{gal-crystalline-test-facts}(2) this is a split
diagram over $\mathrm{El}(X^\sharp)$ with points in $X_{\bar K}$. The \emph{comparison map}
\[
 \eta_\nu=\lambda_*\circ\varepsilon_y:H^k(X_{\bar K},\mathbf Z/\ell^\nu)\to
 H^k(X^\sharp,A_{\mathrm{cris}}/\ell^\nu)
\]
is the composite of the map $\varepsilon_y$ of Lemma~\ref{gal-realization-edge-map}(1) for $q=\ell^\nu$
with the map induced by the homomorphisms $\lambda_A:L_{\ell^\nu}(A,y)\to A_{\mathrm{cris}}(A)/\ell^\nu$ of
Lemma~\ref{gal-realization-locally-constant}, which are natural in $(A,y)$ by that lemma.
\end{definition}
\leanref{Development/WeilBound/CoreStatements/PadicHodgeComparison/EtaleCohomologyComparison/KummerCupProducts/MapUgCToUg.lean}{CrystallineSemistableComparison.E.geomPt}
\leanref{Development/WeilBound/CoreStatements/PadicHodgeComparison/EtaleCohomologyComparison/KummerCupProducts/MapUgCToUg.lean}{CrystallineSemistableComparison.E.probe}
\leanref{Development/WeilBound/CoreStatements/PadicHodgeComparison/EtaleCohomologyComparison/KummerCupProducts/MapUgCToUg.lean}{CrystallineSemistableComparison.E.pts}
\leanref{Development/WeilBound/CoreStatements/PadicHodgeComparison/EtaleCohomologyComparison/KummerCupProducts/MapUgCToUg.lean}{CrystallineSemistableComparison.E.lcMor}
\leanref{Development/WeilBound/CoreStatements/PadicHodgeComparison/EtaleCohomologyComparison/KummerCupProducts/MapUgCToUg.lean}{CrystallineSemistableComparison.E.eM}
\leanref{Development/WeilBound/CoreStatements/PadicHodgeComparison/EtaleCohomologyComparison/KummerCupProducts/EtaleComparisonMapOfFacts.lean}{CrystallineSemistableComparison.E.eFam}

\begin{lemma}\label{gal-realization-duality}
Let $\ell$ be a prime, $K$ a finite extension of $\mathbf Q_\ell$ with valuation ring $\mathcal O_K$,
algebraic closure $\bar K$ and $G_K=\operatorname{Gal}(\bar K/K)$, and $Y$ a smooth projective
$\mathcal O_K$-scheme. For $j\geq0$ write $T^j=H^j(Y_{\bar K},\mathbf Z_\ell)$ and $V^j=T^j\otimes
\mathbf Q_\ell$, with the cup product and the action of $G_K$.
\begin{enumerate}
\item Each $V^j$ is finite-dimensional.
\item Suppose that the generic fibre $Y_K$ is connected of dimension $d$. Then $V^j=0$ for $j>2d$. For
  $i+j=2d$ there are $g_1,\ldots,g_r$ spanning $V^{2d}$ such that for all $s,s'$ there are $\sigma\in G_K$
  and $z\in\mathbf Q_\ell^\times$ with $\sigma g_s=zg_{s'}$, and such that the $\mathbf Q_\ell$-linear forms
  $\mu$ on $V^i$ for which there are $s$ and $y\in T^j$ with $v\cup y=\mu(v)g_s$ in $V^{2d}$ for all
  $v\in T^i$ span the dual of $V^i$.
\item There are smooth projective $\mathcal O_K$-schemes $Y_1,\ldots,Y_r$ with connected generic fibres
  and $\mathcal O_K$-morphisms $u_s:Y_s\to Y$ such that $v\mapsto(u_s^*v)_s$ is injective on every $V^j$.
\end{enumerate}
These follow from the finiteness theorem and Poincar\'e duality for \'etale cohomology
\cite[Expos\'e~XVIII]{AGV73}, \cite[Chapter~VI, \S11]{Mil80}. In (2), $G_K$ permutes the connected
components of $Y_{\bar K}$ transitively, $V^{2d}$ is spanned by their fundamental classes $g_s$, and the
cup product into each line $\mathbf Q_\ell g_s$ is a perfect pairing on the cohomology of the
corresponding component. In (3) the $Y_s$ are the connected components of $Y$.
\end{lemma}
\leanref{Development/WeilBound/CoreStatements/PadicHodgeComparison/EtaleCohomologyComparison/BocksteinBaseChange/EtSide.lean}{CrystallineSemistableComparison.EtSide}
\leanref{Development/WeilBound/CoreStatements/PadicHodgeComparison/EtaleCohomologyComparison/BocksteinBaseChange/EtSide.lean}{CrystallineSemistableComparison.Model.Dual}
\leanref{Development/AutomorphicArithmetic/Foundations/GaloisRepresentations/KunnethLefschetz/EtSideGoodHolds.lean}{UpperTheory.Gal.KunnethLefschetz.Top.etSide_sm_holds}

\begin{definition}\label{gal-realization-period-datum}
Let $\ell$, $C$, $A_{\mathrm{cris}}$, $K$, $\mathcal O_K$, $\bar K$ and $G_K$ be as in
Definition~\ref{gal-realization-split-diagrams}(3), $t\in A_{\mathrm{cris}}$ as in
Theorem~\ref{gal-crystalline-realization}, and $B_{\mathrm{dR}}^+\supset A_{\mathrm{cris}}$ and its
fraction field $B_{\mathrm{dR}}$ Fontaine's rings \cite{Fon94a}, with their actions of $G_K$. Let $Y$ be a
smooth projective $\mathcal O_K$-scheme whose generic fibre is connected of dimension $d$. For $j\geq0$
write $T^j=H^j(Y_{\bar K},\mathbf Z_\ell)$, let $E^j=\varprojlim_\nu H^j(Y^\sharp,A_{\mathrm{cris}}/\ell^\nu)$,
with the cup product and the action of $G_K$ of Definition~\ref{gal-realization-split-diagrams}, and let
$\eta:T^j\to E^j$ send $v$ to $(\eta_\nu(v))_\nu$ for the comparison maps $\eta_\nu$ of
Definition~\ref{gal-realization-eta}. A \emph{period datum} for $Y$ consists of:
\begin{itemize}
\item $A_{\mathrm{cris}}$-modules $H^j$, $j\geq0$, with products $H^i\times H^j\to H^{i+j}$,
  $(x,y)\mapsto xy$, an element $1\in H^0$, and for each $\sigma\in G_K$ additive maps
  $\sigma:H^j\to H^j$;
\item $G_K$-equivariant $A_{\mathrm{cris}}$-linear maps $\beta:H^j\to E^j$ with $\beta(xy)=\beta(x)\cup
  \beta(y)$ and $\beta(1)=1$;
\item finite-dimensional $K$-vector spaces $D^j$ with $K$-bilinear products $D^i\times D^j\to D^{i+j}$,
  $(b,b')\mapsto b\cdot b'$;
\item additive maps $\kappa$ from $H^j$ to the set of functions $\operatorname{Hom}_K(D^j,K)\to
  B_{\mathrm{dR}}$;
\end{itemize}
such that:
\begin{enumerate}
\item for each $j$ there are $m,n\geq0$ with $t^m\ker\beta=0$ on $H^j$ and $t^nE^j\subset\beta(H^j)$;
\item $\kappa(cx)(\psi)=c\,\kappa(x)(\psi)$, $\kappa(x)(a\psi)=a\,\kappa(x)(\psi)$ and
  $\kappa(\sigma x)(\psi)=\sigma(\kappa(x)(\psi))$ for $c\in A_{\mathrm{cris}}$, $a\in K$ and $\sigma\in
  G_K$;
\item for $x\in H^i$, $x'\in H^j$, bases $(b_s)$ of $D^i$ and $(b'_{s'})$ of $D^j$ with dual bases
  $(b_s^*)$ and $(b'^*_{s'})$, and $\psi\in\operatorname{Hom}_K(D^{i+j},K)$,
  \[
   \kappa(xx')(\psi)=\sum_{s,s'}\kappa(x)(b_s^*)\,\kappa(x')(b'^*_{s'})\,\psi(b_s\cdot b'_{s'});
  \]
\item there are $x\in H^2$, $v\in T^2$ and $\psi\in\operatorname{Hom}_K(D^{2d},K)$ with $\beta(x)=t\,\eta(v)$
  and $\kappa(x^d)(\psi)$ a unit of $B_{\mathrm{dR}}^+$, where $x^0=1$ and $x^{k+1}=x^kx$.
\end{enumerate}
\end{definition}
\leanref{Development/WeilBound/CoreStatements/PadicHodgeComparison/EtaleCohomologyComparison/BocksteinBaseChange/ThetaSetupAlgebraMaps.lean}{CrystallineSemistableComparison.Per.Facts}
\leanref{Development/WeilBound/CoreStatements/PadicHodgeComparison/EtaleCohomologyComparison/BocksteinBaseChange/ThetaSetupAlgebraMaps.lean}{CrystallineSemistableComparison.Dat}
\leanref{Development/WeilBound/CoreStatements/PadicHodgeComparison/EtaleCohomologyComparison/BocksteinBaseChange/ThetaSetupAlgebraMaps.lean}{CrystallineSemistableComparison.Dat.Laws}
\leanref{Development/WeilBound/CoreStatements/PadicHodgeComparison/EtaleCohomologyComparison/BocksteinBaseChange/ThetaSetupAlgebraMaps.lean}{CrystallineSemistableComparison.EMap}
\leanref{Development/WeilBound/CoreStatements/PadicHodgeComparison/EtaleCohomologyComparison/BocksteinBaseChange/ThetaSetupAlgebraMaps.lean}{CrystallineSemistableComparison.BMap}
\leanref{Development/WeilBound/CoreStatements/PadicHodgeComparison/EtaleCohomologyComparison/BocksteinBaseChange/ThetaSetupAlgebraMaps.lean}{CrystallineSemistableComparison.PMap}
\leanref{Development/WeilBound/CoreStatements/PadicHodgeComparison/EtaleCohomologyComparison/BocksteinBaseChange/BcrisRegularPeriodRing.lean}{CrystallineSemistableComparison.per}
\leanref{Development/WeilBound/CoreStatements/PadicHodgeComparison/EtaleCohomologyComparison/BocksteinBaseChange/DatPull.lean}{CrystallineSemistableComparison.dat}

\begin{lemma}\label{gal-realization-injectivity}
Let $\ell$, $A_{\mathrm{cris}}$, $t$, $K$, $\mathcal O_K$, $G_K$ and $B_{\mathrm{dR}}$ be as in
Definition~\ref{gal-realization-period-datum}. Let $Y$ be a smooth projective $\mathcal O_K$-scheme whose
generic fibre is connected of dimension $d$, let $T^j$, $E^j$ and $\eta$ be as in
Definition~\ref{gal-realization-period-datum}, let $(H,\beta,D,\kappa)$
be a period datum for $Y$, and let $V^j=T^j\otimes\mathbf Q_\ell$. Suppose that $\eta$ is
$G_K$-equivariant, that $\eta(v\cup v')=\eta(v)\cup\eta(v')$, and that $\eta(w)=1$ for some $w\in T^0$.
Let $v_1,\ldots,v_r\in T^j$, $c_1,\ldots,c_r\in A_{\mathrm{cris}}$ and $u\geq0$ satisfy
$t^u\sum_kc_k\eta(v_k)=0$ in $E^j$. Then $\sum_k\mu(v_k)c_k=0$ in $B_{\mathrm{dR}}$ for every
$\mathbf Q_\ell$-linear form $\mu$ on $V^j$.
\end{lemma}
\leanref{Development/AutomorphicArithmetic/CommutativeAlgebra/EtaleCohomology/VInduction.lean}{UpperTheory.Ru.crystallineComparison.IsVal}
\leanref{Development/AutomorphicArithmetic/CommutativeAlgebra/EtaleCohomology/VInduction.lean}{UpperTheory.Ru.crystallineComparison.exists_val}
\leanref{Development/AutomorphicArithmetic/CommutativeAlgebra/EtaleCohomology/VInduction.lean}{UpperTheory.Ru.crystallineComparison.inj_conn}

\begin{proof}
The map $\eta$ is $\mathbf Z_\ell$-linear, since each $\eta_\nu$ is additive and its target is killed by
$\ell^\nu$. In $B_{\mathrm{dR}}$ the element $t$ is nonzero. The numbers (1) to (4) below refer to the
conditions of Definition~\ref{gal-realization-period-datum}.

\emph{Periods of classes.} Fix $j$, and let $m,n$ be as in (1). For $v\in T^j$ choose $h_v\in H^j$ with
$\beta(h_v)=t^n\eta(v)$, and put $C(v)(\psi)=t^{-n}\kappa(h_v)(\psi)$ for $\psi\in\operatorname{Hom}_K(D^j,K)$.
\begin{enumerate}
\item[(a)] If $h\in H^j$ and $c\in A_{\mathrm{cris}}$ satisfy $\beta(h)=c\,\eta(v)$, then
  $\kappa(h)(\psi)=c\,C(v)(\psi)$. Indeed $\beta(t^nh-ch_v)=0$, so $t^m(t^nh-ch_v)=0$, and (2) gives
  $t^{m+n}\kappa(h)(\psi)=t^mc\,\kappa(h_v)(\psi)$ in the field $B_{\mathrm{dR}}$.
\item[(b)] Applying (a) to $h_v+h_{v'}$ and to $zh_v$ for $z\in\mathbf Z_\ell$ shows that $C$ is
  $\mathbf Z_\ell$-linear in $v$. So $C$ extends to a $\mathbf Q_\ell$-linear map on $V^j$.
\item[(c)] If $t^u\sum_kc_k\eta(v_k)=0$, then $\beta(t^u\sum_kc_kh_{v_k})=0$, so
  $t^{m+u}\sum_kc_kh_{v_k}=0$, and applying $\kappa$ gives $\sum_kc_kC(v_k)(\psi)=0$ for every $\psi$.
\item[(d)] For $\sigma\in G_K$, $\beta(\sigma h_v)=\sigma(t^n)\eta(\sigma v)$, so by (a) and (2),
  $\sigma(t^n)\,C(\sigma v)(\psi)=\kappa(\sigma h_v)(\psi)=\sigma(t^n)\,\sigma(C(v)(\psi))$. Hence
  $C(\sigma w)(\psi)=\sigma(C(w)(\psi))$ for $w\in V^j$, as $\sigma$ fixes $\mathbf Q_\ell$.
\item[(e)] For $v\in T^i$, $v'\in T^j$, with $\beta(h_v)=t^n\eta(v)$ and $\beta(h_{v'})=t^{n'}\eta(v')$,
  one has $\beta(h_vh_{v'})=t^{n+n'}\eta(v\cup v')$. By (a) and (3), for bases as in (3),
  $C(v\cup v')(\psi)=\sum_{s,s'}C(v)(b_s^*)\,C(v')(b'^*_{s'})\,\psi(b_s\cdot b'_{s'})$.
\item[(f)] Take $x$, $v$ and $\psi$ as in (4), and put $w_0=w$ and $w_{k+1}=w_k\cup v$. Since $\beta$
  and $\eta$ are multiplicative and $\beta(1)=1=\eta(w)$, induction gives $\beta(x^k)=t^k\eta(w_k)$. By
  (a), $\kappa(x^d)(\psi)=t^dC(w_d)(\psi)$, which is a unit, so $C(w_d)(\psi)\neq0$.
\end{enumerate}

\emph{Duality.} If $j>2d$ then $V^j=0$ by Lemma~\ref{gal-realization-duality}(2), and there is nothing to
prove. Otherwise let $j'=2d-j$ and take $g_1,\ldots,g_r$ as in Lemma~\ref{gal-realization-duality}(2) for
$(j,j')$, and $\psi$ as in (f). Since the $g_s$ span $V^{2d}$ and $C(w_d)(\psi)\neq0$, some
$C(g_{s_0})(\psi)$ is nonzero. For every $s$ there are $\sigma$ and $z\neq0$ with $\sigma g_s=zg_{s_0}$,
and (d) gives $\sigma(C(g_s)(\psi))=z\,C(g_{s_0})(\psi)\neq0$; so $C(g_s)(\psi)\neq0$ for all $s$. The
conclusion is linear in $\mu$, so by Lemma~\ref{gal-realization-duality}(2) it suffices to take $\mu$
with $v\cup y=\mu(v)g_s$ for some $s$, some $y\in T^{j'}$ and all $v\in T^j$. For such $\mu$, (e) with
bases of $D^j$ and $D^{j'}$ and then (c) give
\[
 \Bigl(\sum_kc_k\mu(v_k)\Bigr)C(g_s)(\psi)=\sum_kc_kC(v_k\cup y)(\psi)=\sum_{s_1,s_2}
 \Bigl(\sum_kc_kC(v_k)(b_{s_1}^*)\Bigr)C(y)(b'^*_{s_2})\,\psi(b_{s_1}\cdot b'_{s_2})=0.
\]
As $C(g_s)(\psi)\neq0$, $\sum_k\mu(v_k)c_k=0$.
\end{proof}

\begin{lemma}\label{gal-realization-descent}
Let $\ell$ be a prime and $W$ a domain in which every nonzero element is $\ell^m$ times a unit for some
$m\geq0$. Let $A$ be a $W$-algebra, $t\in A$, and $c:A\to B$ a ring homomorphism such that $\ell$ and
$c(t)$ are units of $B$, every element of $B$ is $c(a)c(t)^{-r}$ for some $a\in A$ and $r\geq0$, and every
$a$ with $c(a)=0$ is killed by a power of $t$. Let a group $\Gamma$ act on $A$ and on $B$ by ring
homomorphisms $\sigma_A$ and $\sigma_B$ with $\sigma_Bc=c\sigma_A$, fixing the image of $W$ in $A$, with
$\sigma_A(t)=\chi(\sigma)t$ for elements $\chi(\sigma)\in A$, and such that for every $\Gamma$-invariant
$x\in B$ there are $a\geq0$ and $w\in W$ with $\ell^ax=c(w)$. Let $L$ be a $W$-module, with $\Gamma$ acting
on $A\otimes_WL$ by $\sigma_A\otimes1$. Let $z\in A\otimes_WL$ and $s\geq0$ be such that for every
$\sigma\in\Gamma$ there are $a,b\geq0$ with $\ell^at^b(\sigma_A\otimes1)(z)=\ell^at^b\chi(\sigma)^sz$.
Then there are $a,b\geq0$ and $y\in L$ with $\ell^at^bz=t^{b+s}(1\otimes y)$.
\end{lemma}
\leanref{Development/AutomorphicArithmetic/EtaleIdempotents/GaloisEquivariantTensorDescentOfPrimePowerUnits.lean}{UpperTheory.galois_equivariant_tensor_descent_of_prime_power_units}

\begin{proof}
The hypotheses on $c$ say that $B=A[1/t]$. Nonzero elements of $W$ map to units of $B$, so the fraction
field $K_0$ of $W$ maps to $B$, and $\Gamma$ fixes it. Choose a basis $(b_i)$ of the $K_0$-vector space
$K_0\otimes_WL$. Then $B\otimes_WL$ is free over $B$ with basis $(b_i)$; let $\Phi_i:A\otimes_WL\to B$ be
the $i$-th coordinate of the image, so that $\Phi_i(a\otimes y)=c(a)\cdot(\text{the }i\text{-th coordinate
of }1\otimes y)$ and $\Phi_i((\sigma_A\otimes1)z)=\sigma_B(\Phi_i(z))$. Cancelling the unit
$c(\ell^at^b)$ in the hypothesis gives $\sigma_B(\Phi_i(z))=c(\chi(\sigma))^s\Phi_i(z)$, so
$c(t)^{-s}\Phi_i(z)$ is $\Gamma$-invariant and $\ell^{a_i}c(t)^{-s}\Phi_i(z)=c(w_i)$ with $a_i\geq0$ and
$w_i\in W$. Only finitely many $\Phi_i(z)$ are nonzero; let $N$ be the largest of their $a_i$. The vector
$\sum_i\ell^{N-a_i}w_ib_i$ of $K_0\otimes_WL$ is $(1\otimes y_0)/w$ for some $y_0\in L$ and nonzero $w\in W$,
and writing $w=\ell^{m}\beta$ with $\beta$ a unit, $y=\beta^{-1}y_0$ satisfies $1\otimes y=\ell^m\sum_i
\ell^{N-a_i}w_ib_i$. Then $\ell^{m+N}z-t^s(1\otimes y)$ has all coordinates $\Phi_i$ equal to $0$, so it
maps to $0$ in $B\otimes_WL=(A\otimes_WL)[1/t]$ and is killed by some $t^b$. Hence
$\ell^{m+N}t^bz=t^{b+s}(1\otimes y)$.
\end{proof}

\begin{lemma}\label{gal-realization-invariants}
Let $\ell$, $C$, $A_{\mathrm{cris}}$, $t$, $K$ and $G_K$ be as in
Definition~\ref{gal-realization-period-datum}, let $\mathbf F_q$ be the residue field of $K$, and let
$W(\mathbf F_q)\to A_{\mathrm{cris}}$ be induced by the embedding of $\mathbf F_q$ in $\mathcal O_C^\flat$
that lifts its embedding in the residue field of $C$. Let $B_{\mathrm{cris}}=A_{\mathrm{cris}}[1/t]$. Then
$G_K$ fixes the image of $W(\mathbf F_q)$, and for every $G_K$-invariant $x\in B_{\mathrm{cris}}$ there are
$a\geq0$ and $w\in W(\mathbf F_q)$ with $\ell^ax=w$. This is Fontaine's computation
$B_{\mathrm{cris}}^{G_K}=W(\mathbf F_q)[1/\ell]$ \cite{Fon94a}.
\end{lemma}
\leanref{Development/AutomorphicArithmetic/CommutativeAlgebra/EtaleCohomology/AbsInvOf.lean}{UpperTheory.Ru.crystallineComparison.InvC}
\leanref{Development/AutomorphicArithmetic/CommutativeAlgebra/EtaleCohomology/PullbackFinRelation.lean}{UpperTheory.Ru.crystallineComparison.pbInv}

\begin{definition}\label{gal-realization-frobenius-datum}
Let $\ell$, $C$, $A_{\mathrm{cris}}$, $t$, $K$, $\mathcal O_K$, $\bar K$ and $G_K$ be as in
Definition~\ref{gal-realization-period-datum}, let $\mathbf F_q$, $q=\ell^f$, be the residue field of $K$,
$W=W(\mathbf F_q)\to A_{\mathrm{cris}}$ as in Lemma~\ref{gal-realization-invariants}, $g:\mathcal O_K\to
\mathcal O_C/\ell$ the map sending $x$ to the Teichm\"uller representative of its residue, modulo $\ell$,
and $N$ a multiple of $f$ with $(x\bmod\ell)^{\ell^N}=g(x)$ for all $x\in\mathcal O_K$. Let $Y$ be a
smooth projective $\mathcal O_K$-scheme, $F:Y_g\to Y_g$ the $q$-power Frobenius and $\rho_N:Y^\sharp\to
Y_g$ the transformation of Definition~\ref{gal-crystalline-test-algebras}(4), and $j\geq0$. A
\emph{Frobenius datum} of level $N$ for $Y$ in degree $j$ consists of:
\begin{itemize}
\item an $A_{\mathrm{cris}}$-module $H$ with a $\phi$-semilinear endomorphism $\phi_H$ and an action of
  $G_K$, semilinear over its action on $A_{\mathrm{cris}}$;
\item $A_{\mathrm{cris}}$-linear maps $b_\nu:H\to H^j(Y^\sharp,A_{\mathrm{cris}}/\ell^\nu)$, $\nu\geq1$,
  with $b_\nu\phi_H=\phi b_\nu$ and $b_\nu\sigma=\sigma b_\nu$ for $\sigma\in G_K$
  (Definition~\ref{gal-realization-split-diagrams}(3));
\item a $W$-module $L$, an integer $e\geq0$, an additive map $\alpha:H\to A_{\mathrm{cris}}\otimes_WL$ with
  $\alpha(cx)=\phi^N(c)\alpha(x)$ and $\alpha(\sigma x)=(\sigma\otimes1)\alpha(x)$ for $c\in A_{\mathrm{cris}}$
  and $\sigma\in G_K$, and an $A_{\mathrm{cris}}$-linear map $\beta:A_{\mathrm{cris}}\otimes_WL\to H$ with
  $\beta(\alpha(x))=\ell^e\phi_H^N(x)$;
\end{itemize}
such that:
\begin{enumerate}
\item there is $u\geq0$ such that for every family $(y_\nu)$ in $\varprojlim_\nu
  H^j(Y^\sharp,A_{\mathrm{cris}}/\ell^\nu)$ there is $x\in H$ with $b_\nu(x)=t^uy_\nu$ for all $\nu$, and
  $t^ux=0$ for every $x\in H$ with $b_\nu(x)=0$ for all $\nu$;
\item for $y\in L$ and $\nu\geq1$ there is $X_\nu\in H^j(Y_g,A_{\mathrm{cris}}/\ell^\nu)$ with
  $\phi^fX_\nu=F^*X_\nu$ and $b_\nu(\beta(1\otimes y))=\rho_N^*X_\nu$.
\end{enumerate}
\end{definition}
\leanref{Development/WeilBound/CoreStatements/PadicHodgeComparison/EtaleCohomologyComparison/PerfectoidPullbacks/FrobeniusWittActions/PhiTTmul.lean}{CrystallineSemistableComparison.Pb.Fac}
\leanref{Development/AutomorphicArithmetic/CommutativeAlgebra/EtaleCohomology/AbsInvOf.lean}{UpperTheory.Ru.crystallineComparison.BIn}
\leanref{Development/AutomorphicArithmetic/CommutativeAlgebra/EtaleCohomology/AbsInvOf.lean}{UpperTheory.Ru.crystallineComparison.Pin}
\leanref{Development/AutomorphicArithmetic/CommutativeAlgebra/EtaleCohomology/AbsInvOf.lean}{UpperTheory.Ru.crystallineComparison.AbsCl}
\leanref{Development/AutomorphicArithmetic/CommutativeAlgebra/EtaleCohomology/PhiAIter.lean}{UpperTheory.Ru.crystallineComparison.keyStep}

\begin{lemma}\label{gal-realization-twist}
Let $\ell$, $C$, $K$, $\mathcal O_K$, $\bar K$ and $G_K$ be as in
Definition~\ref{gal-realization-split-diagrams}(3), let $\sigma\in G_K$ and $\nu\geq1$, and let $A$ be a test
algebra, with $A^\sigma$ as there.
\begin{enumerate}
\item The identity of $W(A)$ induces a ring isomorphism $\tau_A:(A^\sigma)^\sharp\to A^\sharp$, and the
  identity of $W(A)[1/\ell]$ induces a ring homomorphism $\tau_A:A_{\mathrm{cris}}(A^\sigma)/\ell^\nu\to
  A_{\mathrm{cris}}(A)/\ell^\nu$. The first carries the map $\mathcal O_C\to(A^\sigma)^\sharp$ to $\sigma$
  followed by $\mathcal O_C\to A^\sharp$, and after inverting $\ell$ it carries $\bar K\to
  (A^\sigma)^\sharp[1/\ell]$ to $\sigma$ followed by $\bar K\to A^\sharp[1/\ell]$. The second satisfies
  $\tau_A(cx)=\sigma(c)\tau_A(x)$ for $c\in A_{\mathrm{cris}}$. A map $h:A\to B$ of test algebras is also a
  map $A^\sigma\to B^\sigma$, and both maps $\tau$ commute with the maps induced by $h$.
\item For a locally constant function $f$ on $\operatorname{Spec}(A^\sigma)^\sharp[1/\ell]$ with values in
  $\mathbf Z/\ell^\nu$, $\lambda_A(f\circ\operatorname{Spec}\tau_A)=\tau_A(\lambda_{A^\sigma}(f))$, with
  $\lambda$ as in Lemma~\ref{gal-realization-locally-constant}.
\item Let $X$ be an $\mathcal O_K$-scheme. For $y\in X^\sharp(A)$, $y^\sigma=y\circ(\operatorname{Spec}
  \tau_A)^{-1}$ lies in $X^\sharp(A^\sigma)$, and $t_\sigma(A,y)=(A^\sigma,y^\sigma)$ defines a functor
  $t_\sigma$ from $\mathrm{El}(X^\sharp)$ to itself. The action of $\sigma$ on
  $H^k(X^\sharp,A_{\mathrm{cris}}/\ell^\nu)$ in Definition~\ref{gal-realization-split-diagrams}(3) is the map
  induced by $t_\sigma$ and the homomorphisms $\tau_A$, as in
  Definition~\ref{gal-realization-split-diagrams}(1). For the points $y_{(A,y)}$ of
  Definition~\ref{gal-realization-eta},
  \[
   (1_X\times\operatorname{Spec}\sigma)\circ y_{(A,y)}=y_{(A^\sigma,y^\sigma)}\circ\operatorname{Spec}\tau_A
  \]
  as morphisms $\operatorname{Spec}A^\sharp[1/\ell]\to X_{\bar K}$, where $\operatorname{Spec}\sigma$ is the
  automorphism of $\operatorname{Spec}\bar K$ induced by $\sigma$.
\end{enumerate}
\end{lemma}
\leanref{Development/WeilBound/CoreStatements/PadicHodgeComparison/EtaleCohomologyComparison/AcrisKummerUnits/GaloisTwistIdentitySection.lean}{CrystallineSemistableComparison.PAlg.twU}
\leanref{Development/WeilBound/CoreStatements/PadicHodgeComparison/EtaleCohomologyComparison/AcrisKummerUnits/GaloisTwistIdentitySection.lean}{CrystallineSemistableComparison.PAlg.twU_bijective}
\leanref{Development/WeilBound/CoreStatements/PadicHodgeComparison/EtaleCohomologyComparison/AcrisKummerUnits/GaloisTwistIdentitySection.lean}{CrystallineSemistableComparison.PAlg.twU_oC}
\leanref{Development/WeilBound/CoreStatements/PadicHodgeComparison/EtaleCohomologyComparison/AcrisKummerUnits/GaloisTwistIdentitySection.lean}{CrystallineSemistableComparison.PAlg.twAcr_cris}
\leanref{Development/WeilBound/CoreStatements/PadicHodgeComparison/EtaleCohomologyComparison/KummerCupProducts/EtaleComparisonMapOfFacts.lean}{CrystallineSemistableComparison.E.twUg}
\leanref{Development/AutomorphicSpectrum/Proofs/BlockDecomposition/TwUgCompAlgebraMap.lean}{GlobalTheory.CompanionTheory.TwUg.twUg_prop}
\leanref{Development/AutomorphicSpectrum/Proofs/BlockDecomposition/PullPullOfPullC.lean}{GlobalTheory.CompanionTheory.FEquiProof.twAcr_lcA}
\leanref{Development/WeilBound/CoreStatements/PadicHodgeComparison/EtaleCohomologyComparison/AcrisKummerUnits/GaloisTwistIdentitySection.lean}{CrystallineSemistableComparison.pxTw}
\leanref{Development/WeilBound/CoreStatements/PadicHodgeComparison/EtaleCohomologyComparison/AcrisKummerUnits/GaloisTwistIdentitySection.lean}{CrystallineSemistableComparison.elTw}
\leanref{Development/WeilBound/CoreStatements/PadicHodgeComparison/EtaleCohomologyComparison/AcrisKummerUnits/GaloisTwistIdentitySection.lean}{CrystallineSemistableComparison.galMor}
\leanref{Development/AbelianVarietyTheory/Definitions/GaloisGeometry/IsMatrixForm.lean}{GlobalTheory.Et.galMap}

\begin{proof}
(1) Let $\xi_0=\ell-[\ell^\flat]\in W(\mathcal O_C^\flat)$, and let $\theta_0:W(\mathcal O_C^\flat)\to\mathcal O_C$
be Fontaine's map, which commutes with the action of $\sigma$. For $A$ the element $\xi\in W(A)$ is the
image of $\xi_0$, and for $A^\sigma$ it is the image $\xi_\sigma$ of $\sigma(\xi_0)$. The map
$\mathcal O_C\to A^\sharp$ is induced by $\theta_0$, so $\theta(\xi_\sigma)$ is the image of
$\theta_0(\sigma(\xi_0))=\sigma(\theta_0(\xi_0))=0$. In the same way, $\xi$ is the image of
$\sigma^{-1}(\xi_0)$ under the structure map of $A^\sigma$, so its image in $(A^\sigma)^\sharp$ is the image
of $\theta_0(\sigma^{-1}(\xi_0))=0$. Hence $\xi W(A)=\xi_\sigma W(A)$, and the identity of
$W(A)$ induces an isomorphism $\tau_A:(A^\sigma)^\sharp\to A^\sharp$. For $x=\theta_0(v)\in\mathcal O_C$,
the image of $x$ in $(A^\sigma)^\sharp$ is the class of the image of $\sigma(v)$ in $W(A)$, which $\tau_A$
sends to the image of $\theta_0(\sigma(v))=\sigma(x)$ in $A^\sharp$. Inverting $\ell$ and restricting to
$\bar K\subset C=\mathcal O_C[1/\ell]$ gives the statement for $\bar K$. The element $\sigma(\xi_0)$ lies in the
kernel of $\theta_0$, so its divided powers lie in $A_{\mathrm{cris}}$, and their images
$\xi_\sigma^i/i!$ lie in $A_{\mathrm{cris}}(A)$. So the identity of $W(A)[1/\ell]$ maps
$A_{\mathrm{cris}}(A^\sigma)$ into $A_{\mathrm{cris}}(A)$. An element $c\in A_{\mathrm{cris}}$ acts on
$A_{\mathrm{cris}}(A^\sigma)$ through $\sigma$ followed by the map to $W(A)[1/\ell]$, which gives
$\tau_A(cx)=\sigma(c)\tau_A(x)$. All these maps are induced by the identity of $W(A)$ or of $W(A)[1/\ell]$,
so they commute with the maps induced by $W(h)$.

(2) Let $e_a$ be the idempotents of $f$ given by Lemma~\ref{gal-realization-locally-constant}(1) for
$A^\sigma$; they are idempotents of the ring $A$. As $\tau_A$ is induced by the identity of $W(A)$, it
sends $\theta([e_a])$ in $(A^\sigma)^\sharp$ to $\theta([e_a])$ in $A^\sharp$. So a prime $\mathfrak p$ of
$A^\sharp[1/\ell]$ does not contain $\theta([e_a])$ if and only if $\tau_A^{-1}(\mathfrak p)$ does not
contain $\theta([e_a])$, that is, if and only if $f(\tau_A^{-1}(\mathfrak p))=a$. Hence the $e_a$ are also
the idempotents of $f\circ\operatorname{Spec}\tau_A$, and $\lambda_A(f\circ\operatorname{Spec}\tau_A)=
\sum_a\tilde a\,[e_a]=\tau_A(\lambda_{A^\sigma}(f))$.

(3) The $\mathcal O_K$-algebra structure of $(A^\sigma)^\sharp$ comes from $\mathcal O_K\subset\mathcal O_C\to
(A^\sigma)^\sharp$. By (1), $\tau_A$ carries it to $\sigma$ followed by $\mathcal O_K\to A^\sharp$, which is
the $\mathcal O_K$-algebra structure of $A^\sharp$ because $\sigma$ fixes $K$. So $y^\sigma$ is an
$\mathcal O_K$-morphism. For a map $h:(A,y)\to(B,y')$ in $\mathrm{El}(X^\sharp)$, the naturality in (1) gives
$(\operatorname{Spec}\tau_A)^{-1}\circ\operatorname{Spec}h=\operatorname{Spec}h\circ(\operatorname{Spec}
\tau_B)^{-1}$ on $\operatorname{Spec}(B^\sigma)^\sharp$, so $h$ is a map $(A^\sigma,y^\sigma)\to
(B^\sigma,y'^\sigma)$, and $t_\sigma$ is a functor. The action of Definition~\ref{gal-realization-split-diagrams}(3)
is the transport of structure along $t_\sigma$ and the $\tau_A$. Finally, both sides of the displayed
identity have the same component in $X$, namely $\operatorname{Spec}A^\sharp[1/\ell]\to\operatorname{Spec}
A^\sharp\xrightarrow{y}X$, since $1_X\times\operatorname{Spec}\sigma$ does not change it and
$y^\sigma\circ\operatorname{Spec}\tau_A=y$ on $\operatorname{Spec}A^\sharp$. Their components in
$\operatorname{Spec}\bar K$ are induced by $\bar K\xrightarrow{\sigma}\bar K\to A^\sharp[1/\ell]$ and by
$\tau_A$ composed with $\bar K\to(A^\sigma)^\sharp[1/\ell]$, which agree by (1).
\end{proof}

\begin{lemma}\label{gal-realization-cup-staircase}
Let $\mathcal C$, $(U_c)$, $(y_c)$, $X$, $q$ and $I^\bullet$ be as in Lemma~\ref{gal-realization-edge-map}.
Let $\mu$ be bilinear maps $\Gamma(U,I^s)\times\Gamma(U,I^t)\to\Gamma(U,I^{s+t})$, for schemes $U$ and
$s,t\geq0$, that commute with restriction along morphisms of schemes and satisfy
$d\mu(a,b)=\mu(da,b)+(-1)^s\mu(a,db)$. Suppose that $\mu(1,1)=1$ in $\Gamma(U_c,I^0)$ for every $c$, where
$1\in L_q(c)\subset\Gamma(U_c,I^0)$, and that for $a\in\Gamma(X,I^i)$ and $b\in\Gamma(X,I^j)$ with $da=0$
and $db=0$ the cup product of their classes in $H^*(X,\mathbf Z/q)$ is the class of $\mu(a,b)$. Then
\[
 \varepsilon_y(x\cup x')=\varepsilon_y(x)\cup\varepsilon_y(x')
\]
for $x\in H^i(X,\mathbf Z/q)$ and $x'\in H^j(X,\mathbf Z/q)$, with the cup product of
Definition~\ref{gal-realization-split-diagrams}(1) on the right.
\end{lemma}
\leanref{Development/WeilBound/CoreStatements/PadicHodgeComparison/EtaleCohomologyComparison/KummerCupProducts/EtaleChernClassEqKummerClass.lean}{CrystallineSemistableComparison.E.Main.etCup}
\leanref{Development/WeilBound/CoreStatements/PadicHodgeComparison/EtaleCohomologyComparison/KummerCupProducts/EtaleChernClassEqKummerClass.lean}{CrystallineSemistableComparison.E.Main.edgeCup}
\leanref{Development/WeilBound/CoreStatements/PadicHodgeComparison/EtaleCohomologyComparison/KummerCupProducts/EtaleChernClassEqKummerClass.lean}{CrystallineSemistableComparison.E.Main.stTop}
\leanref{Development/WeilBound/CoreStatements/PadicHodgeComparison/EtaleCohomologyComparison/KummerCupProducts/EtaleChernClassEqKummerClass.lean}{CrystallineSemistableComparison.E.Main.stCupA}
\leanref{Development/WeilBound/CoreStatements/PadicHodgeComparison/EtaleCohomologyComparison/KummerCupProducts/EtaleChernClassEqKummerClass.lean}{CrystallineSemistableComparison.E.Main.stCupB}
\leanref{Development/WeilBound/CoreStatements/PadicHodgeComparison/EtaleCohomologyComparison/KummerCupProducts/EtaleChernClassEqKummerClass.lean}{CrystallineSemistableComparison.E.Main.kapMul}
\leanref{Development/WeilBound/CoreStatements/PadicHodgeComparison/EtaleCohomologyComparison/KummerCupProducts/CupLeib.lean}{CrystallineSemistableComparison.E.cupLeib}
\leanref{Development/AbelianVarietyTheory/Definitions/Divisors/PullHlComp.lean}{GlobalTheory.EtC.Pair}

\begin{proof}
Let $T^{k,s}$, $C^k\subset T^{k,0}$ and staircases be as in the proof of
Lemma~\ref{gal-realization-edge-map}(1). For $e\in T^{k,s}$ and $g\in T^{k',t}$ let $e\cup g\in
T^{k+k',s+t}$ be the cup product formed with $\mu$,
\[
 (e\cup g)(c_0\to\cdots\to c_{k+k'})=\mu\bigl(e(c_0\to\cdots\to c_k),\,g(c_k\to\cdots\to c_{k+k'})\bigr),
\]
with the first argument restricted to $U_{c_{k+k'}}$. The Leibniz rule for $\mu$ gives
$d(e\cup g)=de\cup g+(-1)^se\cup dg$. Since $\mu$ commutes with restriction, the computation for singular
cochains \cite[Lemma~3.6]{Hat02} gives $\delta(e\cup g)=\delta e\cup g+(-1)^ke\cup\delta g$.

For $f,f'\in L_q(c)$ one has $\mu(f,f')=ff'$. Indeed, $U_c$ is the disjoint union of the open and closed
subschemes $W$ on which $f$ and $f'$ are constant, with values $\alpha_W$ and $\alpha'_W$, and a section of
$I^0$ over $U_c$ is determined by its restrictions to the $W$. On $W$, $\mu(f,f')$ restricts to
$\alpha_W\alpha'_W\mu(1,1)=\alpha_W\alpha'_W$, the restriction of $ff'$. So on cochains with values in
$L_q$ the product $\cup$ is the cup product of Definition~\ref{gal-realization-split-diagrams}(1), and
$z\cup z'\in C^{k+k'}$ for $z\in C^k$ and $z'\in C^{k'}$. Three facts are proved by induction.
\begin{enumerate}
\item[(a)] If there is a staircase from $w\in T^{k,s}$ to a cocycle $z$, then $dw=0$ and $\delta w=0$. For
  $s=0$, $w=z$ lies in $C^k$, so $dw=0$, and $\delta w=\delta z=0$. For $s\geq1$, $w=de$ for the first
  step $e$, so $dw=0$, and the induction hypothesis for $w'=(-1)^k\delta e$ gives $dw'=0$, so
  $\delta w=d\delta e=(-1)^kdw'=0$.
\item[(b)] Let $z\in C^i$ with $\delta z=0$, and suppose there is a staircase from $g\in T^{k,t}$ to
  $z'\in C^j$. Then there is a staircase from $z\cup g\in T^{i+k,t}$ to $z\cup z'$. For $t=0$, $g=z'$. For
  $t\geq1$, let $e$ be the first step, with $de=g$ and a staircase from $(-1)^k\delta e$ to $z'$. As $dz=0$,
  $d(z\cup e)=z\cup g$, and as $\delta z=0$, $(-1)^{i+k}\delta(z\cup e)=z\cup((-1)^k\delta e)$. The
  induction hypothesis gives a staircase from $z\cup((-1)^k\delta e)$ to $z\cup z'$, and $z\cup e$ is the
  first step of a staircase from $z\cup g$ to $z\cup z'$.
\item[(c)] Let $g\in T^{0,t}$ with $dg=0$ and $\delta g=0$, let $z\in C^i$, and suppose there is a
  staircase from $z\cup g\in T^{i,t}$ to some $u\in C^{i+t}$. If there is a staircase from $w\in T^{k,s}$
  to $z$, then there is one from $w\cup g\in T^{k,s+t}$ to $u$. For $s=0$, $w=z$. For $s\geq1$, let $e$ be
  the first step, with $de=w$ and a staircase from $(-1)^k\delta e$ to $z$. As $dg=0$, $d(e\cup g)=w\cup g$,
  and as $\delta g=0$, $(-1)^k\delta(e\cup g)=((-1)^k\delta e)\cup g$. The induction hypothesis gives a
  staircase from $((-1)^k\delta e)\cup g$ to $u$, and $e\cup g$ is the first step of a staircase from
  $w\cup g$ to $u$.
\end{enumerate}

Now let $a$ and $b$ represent $x$ and $x'$, and take staircases from $w=(y_c^*a)_c$ to $z$ and from
$w'=(y_c^*b)_c$ to $z'$ as in Lemma~\ref{gal-realization-edge-map}(1), so that $\varepsilon_y(x)$ and
$\varepsilon_y(x')$ are the classes of $z$ and $z'$. By (a), $dw'=0$ and $\delta w'=0$. By (b) with $g=w'$
there is a staircase from $z\cup w'$ to $z\cup z'$, and by (c) with $g=w'$ there is one from $w\cup w'$ to
$z\cup z'$. Since $\mu$ commutes with restriction, $w\cup w'=(y_c^*\mu(a,b))_c$, and $\mu(a,b)$ represents
$x\cup x'$. By Lemma~\ref{gal-realization-edge-map}(1), $\varepsilon_y(x\cup x')$ is the class of
$z\cup z'$, which is $\varepsilon_y(x)\cup\varepsilon_y(x')$.
\end{proof}

\begin{lemma}\label{gal-realization-period-prism}
Let $\sigma:R\to R'$ be a ring homomorphism, $\mathcal C$ and $\mathcal C'$ categories, $F$ a functor from
$\mathcal C$ to $R$-modules and $F'$ a functor from $\mathcal C'$ to $R'$-modules. Let $G,G':\mathcal C'\to\mathcal
C$ be functors with $\sigma$-semilinear maps $\tau_c:F(G(c))\to F'(c)$ and $\tau'_c:F(G'(c))\to F'(c)$, natural in
$c\in\mathcal C'$, and let $\eta:G\to G'$ be a natural transformation with $\tau'_c\circ F(\eta_c)=\tau_c$ for every
$c$. Then $(G,\tau)$ and $(G',\tau')$ induce the same maps $H^k(\mathcal C,F)\to H^k(\mathcal C',F')$
(Definition~\ref{gal-realization-split-diagrams}(1)).
\end{lemma}
\leanref{Development/WeilBound/CoreStatements/PadicHodgeComparison/EtaleCohomologyComparison/CechPairings/FaceMapCompSuccAboveCastSucc.lean}{CrystallineSemistableComparison.P.prism}
\leanref{Development/WeilBound/CoreStatements/PadicHodgeComparison/EtaleCohomologyComparison/CechPairings/FaceMapCompSuccAboveCastSucc.lean}{CrystallineSemistableComparison.Auxfc.prism_main}
\leanref{Development/WeilBound/CoreStatements/PadicHodgeComparison/EtaleCohomologyComparison/CechPairings/FaceMapCompSuccAboveCastSucc.lean}{CrystallineSemistableComparison.Auxfc.hh}
\leanref{Development/WeilBound/CoreStatements/PadicHodgeComparison/EtaleCohomologyComparison/CechPairings/FaceMapCompSuccAboveCastSucc.lean}{CrystallineSemistableComparison.Auxfc.htpy_succ}
\leanref{Development/WeilBound/CoreStatements/PadicHodgeComparison/EtaleCohomologyComparison/BocksteinBaseChange/CrSEqOne.lean}{CrystallineSemistableComparison.cr_pull_eq}

\begin{proof}
Write $G^*f$ and $G'^*f$ for the induced maps on cochains. For a $(k+1)$-cochain $f$ of $\mathcal C$ with values
in $F$, let $hf$ be the $k$-cochain of $\mathcal C'$ whose value on $c_0\to\cdots\to c_k$ is
\[
 \sum_{i=0}^k(-1)^i\,\tau'_{c_k}f\bigl(G(c_0)\to\cdots\to G(c_i)\xrightarrow{\eta_{c_i}}G'(c_i)\to\cdots\to
 G'(c_k)\bigr).
\]
This is the prism operator of the nerves, and a computation with the face maps gives $G'^*f-G^*f=d(hf)$ for every
cocycle $f$ of positive degree. A $0$-cocycle $f$ satisfies $f(G'(c))=F(\eta_c)f(G(c))$, so
$G'^*f(c)=\tau'_c(F(\eta_c)f(G(c)))=\tau_c(f(G(c)))=G^*f(c)$.
\end{proof}

\subsection{Thickening cohomology and Chern classes}\label{subsection-gal-thickening-cohomology}

\begin{definition}\label{gal-realization-period-functors}
Let $\mathcal O$ be a ring.
\begin{enumerate}
\item A \emph{point functor} over $\mathcal O$ consists of a functor $F$ from rings to sets and, for every ring $B$, a
  relation between elements $y\in F(B)$ and ring maps $f:\mathcal O\to B$, read ``$y$ lies over $f$'', such that each $y$
  lies over at most one $f$, and $F(\varphi)(y)$ lies over $\varphi\circ f$ whenever $y$ lies over $f$ and $\varphi:B\to
  B'$ is a ring map. A \emph{map} $u:F\to F'$ of point functors over $\mathcal O$ is a natural transformation such that
  $u(y)$ lies over $f$ whenever $y$ lies over $f$.
\item An $\mathcal O$-scheme $X$ gives the point functor $F_X$ over $\mathcal O$: $F_X(B)$ is the set of morphisms
  $y:\operatorname{Spec}B\to X$, $F_X(\varphi)(y)=y\circ\operatorname{Spec}\varphi$, and $y$ lies over $f$ if $y$
  followed by the structure morphism of $X$ is $\operatorname{Spec}f$. An $\mathcal O$-morphism $h:X\to X'$ gives the
  map $h_*:F_X\to F_{X'}$, $y\mapsto h\circ y$.
\item Let $\ell$, $C$ and the test algebras be as in Definition~\ref{gal-crystalline-test-algebras}, $F$ a point functor
  over $\mathcal O$ and $\kappa:\mathcal O\to\mathcal O_C$ a ring map. The functor $F^\sharp$ on test algebras sends $A$
  to the set of $y\in F(A^\sharp)$ that lie over $\mathcal O\xrightarrow{\kappa}\mathcal O_C\to A^\sharp$, and a map
  $h:A\to A'$ to $F(h^\sharp)$, where $h^\sharp:A^\sharp\to A'^\sharp$ is induced by $W(h)$. For an $\mathcal O$-scheme
  $X$, $(F_X)^\sharp$ is the functor $X^\sharp$ of Definition~\ref{gal-crystalline-test-algebras}(4). A map $u:F\to F'$
  restricts to a natural transformation $u:F^\sharp\to F'^\sharp$.
\item Let moreover $\bar{\mathbf Q}_\ell$ be the algebraic closure of $\mathbf Q_\ell$ in $C$, $K$ a field with
  $\mathbf Q_\ell\subset K\subset\bar{\mathbf Q}_\ell$, and $\sigma\in\operatorname{Gal}(\bar{\mathbf Q}_\ell/K)$ with
  $\sigma(\kappa(x))=\kappa(x)$ for $x\in\mathcal O$; $\sigma$ acts on $C$, $\mathcal O_C^\flat$ and $A_{\mathrm{cris}}$
  as an element of $G_{\mathbf Q_\ell}$ (Definition~\ref{gal-realization-split-diagrams}(3)). For a test algebra $A$ let
  $A^\sigma$ and the maps $\tau_A$ be as in Lemma~\ref{gal-realization-twist}(1), for $\mathbf Q_\ell$ in place of $K$.
  As $\tau_A$ carries $\mathcal O_C\to(A^\sigma)^\sharp$ to $\sigma$ followed by $\mathcal O_C\to A^\sharp$, and $\sigma$
  fixes $\kappa(\mathcal O)$, $y^\sigma=F(\tau_A^{-1})(y)$ lies in $F^\sharp(A^\sigma)$ for $y\in F^\sharp(A)$. As the
  maps $\tau$ commute with the maps of test algebras, $t_\sigma(A,y)=(A^\sigma,y^\sigma)$ is a functor from
  $\mathrm{El}(F^\sharp)$ to itself. Let $\sigma$ also denote the map of $H^k(F^\sharp,A_{\mathrm{cris}}/\ell^\nu)$
  induced by $t_\sigma$ and the maps $\tau_A:A_{\mathrm{cris}}(A^\sigma)/\ell^\nu\to A_{\mathrm{cris}}(A)/\ell^\nu$
  (Definition~\ref{gal-realization-split-diagrams}(1)); it is semilinear over the action of $\sigma$ on
  $A_{\mathrm{cris}}$. If $K$ is finite over $\mathbf Q_\ell$, $\mathcal O$ is its valuation ring $\mathcal O_K$,
  $\kappa$ is the inclusion and $F=F_X$, then $y^\sigma=y\circ(\operatorname{Spec}\tau_A)^{-1}$, and by
  Lemma~\ref{gal-realization-twist}(3) this map is the action of $\sigma$ of
  Definition~\ref{gal-realization-split-diagrams}(3).
\end{enumerate}
\end{definition}
\leanref{Development/WeilBound/CoreStatements/PadicHodgeComparison/EtaleCohomologyComparison/ZariskiEtaleLifting/SubThickRees.lean}{CrystallineSemistableComparison.B.Fn}
\leanref{Development/WeilBound/CoreStatements/PadicHodgeComparison/EtaleCohomologyComparison/ZariskiEtaleLifting/MapAcrFromBase.lean}{CrystallineSemistableComparison.B.Fn.Map}
\leanref{Development/WeilBound/CoreStatements/PadicHodgeComparison/EtaleCohomologyComparison/ZariskiEtaleLifting/SubThickRees.lean}{CrystallineSemistableComparison.B.schFn}
\leanref{Development/WeilBound/CoreStatements/PadicHodgeComparison/EtaleCohomologyComparison/ZariskiEtaleLifting/MapAcrFromBase.lean}{CrystallineSemistableComparison.B.schMap}
\leanref{Development/WeilBound/CoreStatements/PadicHodgeComparison/EtaleCohomologyComparison/ZariskiEtaleLifting/MapAcrFromBase.lean}{CrystallineSemistableComparison.B.Fn.pt}
\leanref{Development/WeilBound/CoreStatements/PadicHodgeComparison/EtaleCohomologyComparison/ZariskiEtaleLifting/MapAcrFromBase.lean}{CrystallineSemistableComparison.B.schFn_pt}
\leanref{Development/WeilBound/CoreStatements/PadicHodgeComparison/EtaleCohomologyComparison/ZariskiEtaleLifting/MapAcrFromBase.lean}{CrystallineSemistableComparison.B.Fn.Map.pt}
\leanref{Development/WeilBound/CoreStatements/PadicHodgeComparison/EtaleCohomologyComparison/ZariskiEtaleLifting/MapAcrFromBase.lean}{CrystallineSemistableComparison.B.schMap_pt}
\leanref{Development/WeilBound/CoreStatements/PadicHodgeComparison/EtaleCohomologyComparison/ZariskiEtaleLifting/SchFnPtTw.lean}{CrystallineSemistableComparison.B.ptTw}
\leanref{Development/WeilBound/CoreStatements/PadicHodgeComparison/EtaleCohomologyComparison/ZariskiEtaleLifting/SchFnPtTw.lean}{CrystallineSemistableComparison.B.schFn_ptTw}
\leanref{Development/WeilBound/CoreStatements/PadicHodgeComparison/EtaleCohomologyComparison/AcrisKummerUnits/GaloisTwistIdentitySection.lean}{CrystallineSemistableComparison.elTw}
\leanref{Development/WeilBound/CoreStatements/PadicHodgeComparison/EtaleCohomologyComparison/AcrisKummerUnits/GaloisTwistIdentitySection.lean}{CrystallineSemistableComparison.galMor}

\begin{definition}\label{gal-realization-period-thickenings}
Let $\ell$, $C$, $\mathcal O_C^\flat$, $\xi$, $\theta$ and $A_{\mathrm{cris}}$ be as in
Definition~\ref{gal-crystalline-test-algebras}, and $\nu\geq1$.
\begin{enumerate}
\item Let $A$ be a test algebra or $A=\mathcal O_C^\flat$. Let $I(A)$ be the ideal of $A_{\mathrm{cris}}(A)$
  generated by the $\xi^k/k!$, $k\geq1$, and $I'(A)=\ell A_{\mathrm{cris}}(A)+I(A)$. Their images $I_\nu(A)$ and
  $I'_\nu(A)$ in $A_{\mathrm{cris}}(A)/\ell^\nu$ carry the divided powers $\delta$ with $\delta_k(\bar x)=
  \overline{x^k/k!}$ for $x$ in $I(A)$, resp.\ $I'(A)$, where $x^k/k!$ is formed in $W(A)[1/\ell]$. The inclusion
  $W(A)\subset A_{\mathrm{cris}}(A)$ induces $u_A:A^\sharp=W(A)/\xi\to(A_{\mathrm{cris}}(A)/\ell^\nu)/I_\nu(A)$. For
  $A=\mathcal O_C^\flat$ write $A_\nu=A_{\mathrm{cris}}/\ell^\nu$, $I_\nu$, $I'_\nu$, $\gamma$ for the divided powers
  and $u:\mathcal O_C\to A_\nu/I_\nu$.
\item Let $\mathcal O$ be a ring, $F$ a point functor over $\mathcal O$
  (Definition~\ref{gal-realization-period-functors}) and $\kappa:\mathcal O\to\mathcal O_C$ a ring map, and let $I$ be
  $I_\nu$ or $I'_\nu$. A \emph{thickening} of $F$ of level $\nu$ for $I$ consists of a ring $B$, a ring map
  $i:A_\nu\to B$, an ideal $J$ of $B$ with divided powers $\delta$ such that $i(I)\subset J$ and
  $\delta_k(i(x))=i(\gamma_k(x))$ for $x\in I$, and an element $y\in F(B/J)$ that lies over $\mathcal
  O\xrightarrow{\kappa}\mathcal O_C\xrightarrow{u}A_\nu/I_\nu\to A_\nu/I\xrightarrow{i}B/J$. A map
  $(B,i,J,\delta,y)\to(B',i',J',\delta',y')$ of thickenings is a ring map $\varphi:B\to B'$ with $\varphi\circ i=i'$,
  $\varphi(J)\subset J'$, $\delta'_k\circ\varphi=\varphi\circ\delta_k$ on $J$ and $y'=F(\bar\varphi)(y)$, where
  $\bar\varphi:B/J\to B'/J'$ is induced by $\varphi$. Let $\mathrm{Th}_\nu(F)$ and $\mathrm{Th}'_\nu(F)$ be the
  categories of thickenings for $I_\nu$ and for $I'_\nu$. With the functor $\mathcal O$ that sends a thickening to the
  ring $B$, an $A_{\mathrm{cris}}$-algebra through $A_{\mathrm{cris}}\to A_\nu\xrightarrow{i}B$, put
  \[
   H^j_\nu(F)=H^j(\mathrm{Th}_\nu(F),\mathcal O),\qquad H'^j_\nu(F)=H^j(\mathrm{Th}'_\nu(F),\mathcal O),
  \]
  with the cup products of Definition~\ref{gal-realization-split-diagrams}(1) and the unit $1$, the class of the
  constant $0$-cochain $1$. For an $\mathcal O$-scheme $X$, a thickening of $F_X$ has an $\mathcal O$-morphism
  $y:\operatorname{Spec}B/J\to X$, where $B/J$ is an $\mathcal O$-algebra through the composite above, and the condition
  on a map is $y'=y\circ\operatorname{Spec}\bar\varphi$. Write $\mathrm{Th}_\nu(X)$, $\mathrm{Th}'_\nu(X)$,
  $H^j_\nu(X)$ and $H'^j_\nu(X)$ for $\mathrm{Th}_\nu(F_X)$, $\mathrm{Th}'_\nu(F_X)$, $H^j_\nu(F_X)$ and
  $H'^j_\nu(F_X)$, and call their objects thickenings of $X$.
\item Each of the following functors, with the identity maps of the rings $B$, induces maps of these groups as
  in Definition~\ref{gal-realization-split-diagrams}(1).
  \begin{enumerate}
  \item[(a)] For $n\leq m$, a thickening of level $n$ is one of level $m$ through $A_m\to A_n$. This gives the
    \emph{reductions} $H^j_m(F)\to H^j_n(F)$ and $H'^j_m(F)\to H'^j_n(F)$.
  \item[(b)] Let $K$ and $\sigma\in\operatorname{Gal}(\bar{\mathbf Q}_\ell/K)$, with $\sigma(\kappa(x))=\kappa(x)$
    for $x\in\mathcal O$, be as in Definition~\ref{gal-realization-period-functors}(4). Sending $(B,i,J,\delta,y)$
    to $(B,i\circ\sigma,J,\delta,y)$ gives $\sigma$-semilinear maps $\sigma$ of $H^j_\nu(F)$ and of $H'^j_\nu(F)$.
  \item[(c)] On $\mathrm{Th}'_\nu(F)$, sending $(B,i,J,\delta,y)$ to $(B,i\circ\phi,J,\delta,F(F_{B/J})(y))$, where
    $F_{B/J}$ is the $\ell$-th power map of $B/J$, gives a $\phi$-semilinear map $\phi$ of $H'^j_\nu(F)$; for $F=F_X$,
    $F(F_{B/J})(y)=y\circ\operatorname{Spec}F_{B/J}$.
  \item[(d)] A thickening for $I'_\nu$ is one for $I_\nu$. This gives the \emph{comparison}
    $\iota_\nu:H^j_\nu(F)\to H'^j_\nu(F)$.
  \item[(e)] A map $u:F\to F'$ of point functors over $\mathcal O$ sends $(B,i,J,\delta,y)$ to $(B,i,J,\delta,u(y))$,
    which gives $u^*:H^j_\nu(F')\to H^j_\nu(F)$ and $u^*:H'^j_\nu(F')\to H'^j_\nu(F)$. For an $\mathcal O$-morphism
    $h:X\to X'$ write $h^*=(h_*)^*$; the functor sends $(B,i,J,\delta,y)$ to $(B,i,J,\delta,h\circ y)$.
  \end{enumerate}
\item For a test algebra $A$ and $y\in F^\sharp(A)$ (Definition~\ref{gal-realization-period-functors}(3), for
  $\kappa$), let $g_\nu(A,y)$ be the thickening $(A_{\mathrm{cris}}(A)/\ell^\nu,i_A,I_\nu(A),\delta,F(u_A)(y))$ for
  $I_\nu$, where $i_A$ is induced by $\mathcal O_C^\flat\to A$, and let $g'_\nu(A,y)$ be the thickening for $I'_\nu$
  with $I'_\nu(A)$ in place of $I_\nu(A)$ and $u_A$ followed by the projection. For $F=F_X$,
  $F(u_A)(y)=y\circ\operatorname{Spec}u_A$. Maps of test algebras induce maps of these thickenings, so $g_\nu$ and
  $g'_\nu$ are functors on $\mathrm{El}(F^\sharp)$. With the identity maps of the rings $A_{\mathrm{cris}}(A)/\ell^\nu$
  they induce
  \[
   \beta_\nu:H^j_\nu(F)\to H^j(F^\sharp,A_{\mathrm{cris}}/\ell^\nu),\qquad
   \beta'_\nu:H'^j_\nu(F)\to H^j(F^\sharp,A_{\mathrm{cris}}/\ell^\nu).
  \]
  For $n\leq m$, the identity functor of $\mathrm{El}(F^\sharp)$ with the projections $A_{\mathrm{cris}}(A)/\ell^m\to
  A_{\mathrm{cris}}(A)/\ell^n$ induces the \emph{reduction} $H^j(F^\sharp,A_{\mathrm{cris}}/\ell^m)\to
  H^j(F^\sharp,A_{\mathrm{cris}}/\ell^n)$.
\item By Lemma~\ref{gal-realization-period-compatible} below, the maps of (3) and (4) commute with the
  reductions, and the reductions preserve cup products and units. Put $H^j_\infty(F)=\varprojlim_\nu H^j_\nu(F)$,
  $H'^j_\infty(F)=\varprojlim_\nu H'^j_\nu(F)$ and $E^j(F)=\varprojlim_\nu H^j(F^\sharp,A_{\mathrm{cris}}/\ell^\nu)$,
  with the cup products, units and maps $\sigma$, $\phi$ and $u^*$ formed componentwise, and let
  $\beta:H^j_\infty(F)\to E^j(F)$ and $\beta':H'^j_\infty(F)\to E^j(F)$ be the limits of the $\beta_\nu$ and
  $\beta'_\nu$. For an $\mathcal O$-scheme $X$ write $H^j_\infty(X)$, $H'^j_\infty(X)$ and $E^j(X)$ for these groups for
  $F_X$.
\item For a thickening $T=(B,i,J,\delta,y)$ and $r\geq0$, let $J_T^{[r]}$ be the ideal of $B$ generated by the products
  $\delta_{k_1}(x_1)\cdots\delta_{k_n}(x_n)$ with $x_i\in J$ and $k_1+\cdots+k_n\geq r$. A map of thickenings commutes with
  the divided powers, so it carries $J_T^{[r]}$ into $J_{T'}^{[r]}$, and $T\mapsto J_T^{[r]}$ is a subfunctor
  $\mathcal O^{[r]}$ of $\mathcal O$. The \emph{divided power filtration} $\mathrm{Fil}^rH^j_\nu(F)$ is the image of the
  map $H^j(\mathrm{Th}_\nu(F),\mathcal O^{[r]})\to H^j_\nu(F)$ induced by the inclusion, and $\mathrm{Fil}^rH'^j_\nu(F)$
  is formed in the same way on $\mathrm{Th}'_\nu(F)$.
\end{enumerate}
\end{definition}
\leanref{Development/WeilBound/CoreStatements/PadicHodgeComparison/EtaleCohomologyComparison/AcrisKummerUnits/MapUgAlgebraMap.lean}{CrystallineSemistableComparison.PAlg.Fil}
\leanref{Development/WeilBound/CoreStatements/PadicHodgeComparison/EtaleCohomologyComparison/AcrisKummerUnits/MapUgAlgebraMap.lean}{CrystallineSemistableComparison.PAlg.PFil}
\leanref{Development/WeilBound/CoreStatements/PadicHodgeComparison/EtaleCohomologyComparison/AcrisKummerUnits/FiltrationSpanDividedPowers.lean}{CrystallineSemistableComparison.RL.pd_ex}
\leanref{Development/WeilBound/CoreStatements/PadicHodgeComparison/EtaleCohomologyComparison/ZariskiEtaleLifting/MapAcrFromBase.lean}{CrystallineSemistableComparison.B.JJ}
\leanref{Development/WeilBound/CoreStatements/PadicHodgeComparison/EtaleCohomologyComparison/ZariskiEtaleLifting/MapAcrFromBase.lean}{CrystallineSemistableComparison.B.DD}
\leanref{Development/WeilBound/CoreStatements/PadicHodgeComparison/EtaleCohomologyComparison/ZariskiEtaleLifting/MapAcrFromBase.lean}{CrystallineSemistableComparison.B.uJ}
\leanref{Development/WeilBound/CoreStatements/PadicHodgeComparison/EtaleCohomologyComparison/ZariskiEtaleLifting/MapAcrFromBase.lean}{CrystallineSemistableComparison.B.bs}
\leanref{Development/WeilBound/CoreStatements/PadicHodgeComparison/EtaleCohomologyComparison/ZariskiEtaleLifting/MapAcrFromBase.lean}{CrystallineSemistableComparison.B.κb}
\leanref{Development/AbelianVarietyTheory/Definitions/GaloisGeometry/IsDPMorphismBot.lean}{GlobalTheory.LogCris.Thick}
\leanref{Development/AbelianVarietyTheory/Definitions/GaloisGeometry/IsDPMorphismBot.lean}{GlobalTheory.LogCris.Thick.Hom}
\leanref{Development/WeilBound/CoreStatements/PadicHodgeComparison/EtaleCohomologyComparison/ZariskiEtaleLifting/SubThickRees.lean}{CrystallineSemistableComparison.B.Fn.cr}
\leanref{Development/WeilBound/CoreStatements/PadicHodgeComparison/EtaleCohomologyComparison/ZariskiEtaleLifting/MapAcrFromBase.lean}{CrystallineSemistableComparison.B.Lc}
\leanref{Development/WeilBound/CoreStatements/PadicHodgeComparison/EtaleCohomologyComparison/ZariskiEtaleLifting/MapAcrFromBase.lean}{CrystallineSemistableComparison.B.OcA}
\leanref{Development/WeilBound/CoreStatements/PadicHodgeComparison/EtaleCohomologyComparison/ZariskiEtaleLifting/MapAcrFromBase.lean}{CrystallineSemistableComparison.B.Hc}
\leanref{Development/WeilBound/CoreStatements/PadicHodgeComparison/EtaleCohomologyComparison/ZariskiEtaleLifting/SchFnPtTw.lean}{CrystallineSemistableComparison.B.redHc}
\leanref{Development/WeilBound/CoreStatements/PadicHodgeComparison/EtaleCohomologyComparison/ZariskiEtaleLifting/SchFnPtTw.lean}{CrystallineSemistableComparison.B.galHc}
\leanref{Development/WeilBound/CoreStatements/PadicHodgeComparison/EtaleCohomologyComparison/ZariskiEtaleLifting/SchFnPtTw.lean}{CrystallineSemistableComparison.B.phiHc}
\leanref{Development/WeilBound/CoreStatements/PadicHodgeComparison/EtaleCohomologyComparison/ZariskiEtaleLifting/SchFnPtTw.lean}{CrystallineSemistableComparison.B.cmpMor}
\leanref{Development/WeilBound/CoreStatements/PadicHodgeComparison/EtaleCohomologyComparison/ZariskiEtaleLifting/SchFnPtTw.lean}{CrystallineSemistableComparison.B.pullHc}
\leanref{Development/WeilBound/CoreStatements/PadicHodgeComparison/EtaleCohomologyComparison/ZariskiEtaleLifting/MapAcrFromBase.lean}{CrystallineSemistableComparison.B.gObj}
\leanref{Development/WeilBound/CoreStatements/PadicHodgeComparison/EtaleCohomologyComparison/ZariskiEtaleLifting/MapAcrFromBase.lean}{CrystallineSemistableComparison.B.gE}
\leanref{Development/WeilBound/CoreStatements/PadicHodgeComparison/EtaleCohomologyComparison/ZariskiEtaleLifting/MapAcrFromBase.lean}{CrystallineSemistableComparison.B.βMor}
\leanref{Development/WeilBound/CoreStatements/PadicHodgeComparison/EtaleCohomologyComparison/ZariskiEtaleLifting/MapAcrFromBase.lean}{CrystallineSemistableComparison.B.β}
\leanref{Development/WeilBound/CoreStatements/PadicHodgeComparison/EtaleCohomologyComparison/ZariskiEtaleLifting/SchFnPtTw.lean}{CrystallineSemistableComparison.B.HcSys}
\leanref{Development/WeilBound/CoreStatements/PadicHodgeComparison/EtaleCohomologyComparison/BocksteinBaseChange/DatPull.lean}{CrystallineSemistableComparison.Model.HA}
\leanref{Development/WeilBound/CoreStatements/PadicHodgeComparison/EtaleCohomologyComparison/BocksteinBaseChange/AtlasGroupActions/AtlasOf.lean}{CrystallineSemistableComparison.GlueB'.βS}
\leanref{Development/WeilBound/CoreStatements/PadicHodgeComparison/EtaleCohomologyComparison/BocksteinBaseChange/AtlasGroupActions/AtlasOf.lean}{CrystallineSemistableComparison.GlueB'.bβ}
\leanref{Development/WeilBound/CoreStatements/PadicHodgeComparison/EtaleCohomologyComparison/BocksteinBaseChange/AtlasGroupActions/AtlasOf.lean}{CrystallineSemistableComparison.GlueB'.bβ'}
\leanref{Development/AbelianVarietyTheory/Definitions/GaloisGeometry/IsDPMorphismBot.lean}{GlobalTheory.LogCris.Thick.pdPow}
\leanref{Development/WeilBound/CoreStatements/PadicHodgeComparison/EtaleCohomologyComparison/ZariskiEtaleLifting/SubThickRees.lean}{CrystallineSemistableComparison.B.pdS}
\leanref{Development/WeilBound/CoreStatements/PadicHodgeComparison/EtaleCohomologyComparison/ZariskiEtaleLifting/MapAcrFromBase.lean}{CrystallineSemistableComparison.B.HcFil}
\leanref{Development/WeilBound/CoreStatements/PadicHodgeComparison/EtaleCohomologyComparison/AcrisKummerUnits/GaloisTwistIdentitySection.lean}{CrystallineSemistableComparison.redEl}

\begin{lemma}\label{gal-realization-period-compatible}
Let $\mathcal O$ be a ring, $F$ and $F'$ point functors over $\mathcal O$, $u:F\to F'$ a map of point functors and
$\kappa:\mathcal O\to\mathcal O_C$ a ring map (Definition~\ref{gal-realization-period-functors}). Let $K$ be a field with
$\mathbf Q_\ell\subset K\subset\bar{\mathbf Q}_\ell$ and $\sigma\in\operatorname{Gal}(\bar{\mathbf Q}_\ell/K)$ with
$\sigma(\kappa(x))=\kappa(x)$ for $x\in\mathcal O$, and let $\nu\geq1$ and $1\leq n\leq m$. Use the notation of
Definition~\ref{gal-realization-period-thickenings}.
\begin{enumerate}
\item The reductions $H^j_m(F)\to H^j_n(F)$ and $H'^j_m(F)\to H'^j_n(F)$ commute with the maps $\sigma$, $\phi$,
  $\iota$ and $u^*$, and preserve cup products and units.
\item The maps $\beta_\nu$ and $\beta'_\nu$ commute with the reductions and preserve cup products and units.
  Moreover $\beta'_\nu\circ\iota_\nu=\beta_\nu$, $\beta'_\nu\circ\phi=\phi\circ\beta'_\nu$, $\beta_\nu\circ u^*=
  u^*\circ\beta_\nu$, $\beta'_\nu\circ u^*=u^*\circ\beta'_\nu$, $\beta_\nu\circ\sigma=\sigma\circ\beta_\nu$ and
  $\beta'_\nu\circ\sigma=\sigma\circ\beta'_\nu$, with the maps $u^*$ and $\sigma$ of
  Definitions~\ref{gal-crystalline-test-algebras}(2) and~\ref{gal-realization-period-functors}(4) on the right.
\item The reductions carry $\mathrm{Fil}^rH^j_m(F)$ into $\mathrm{Fil}^rH^j_n(F)$ and $\mathrm{Fil}^rH'^j_m(F)$ into
  $\mathrm{Fil}^rH'^j_n(F)$ for $r\geq0$. If $a\in A_{\mathrm{cris}}$ has image in $I_\nu$, resp.\ $I'_\nu$, then
  $a\,\mathrm{Fil}^rH^j_\nu(F)\subset\mathrm{Fil}^{r+1}H^j_\nu(F)$, resp.\ $a\,\mathrm{Fil}^rH'^j_\nu(F)\subset
  \mathrm{Fil}^{r+1}H'^j_\nu(F)$.
\end{enumerate}
For an $\mathcal O$-scheme $X$ and an $\mathcal O$-morphism $h:X\to X'$ this applies to $F_X$ and $h_*$; if $\mathcal
O=\mathcal O_K$ for a finite extension $K$ of $\mathbf Q_\ell$ and $\kappa$ is the inclusion, the map $\sigma$ on the
right of (2) is the action of Definition~\ref{gal-realization-split-diagrams}(3)
(Definition~\ref{gal-realization-period-functors}(4)).
\end{lemma}
\leanref{Development/WeilBound/CoreStatements/PadicHodgeComparison/EtaleCohomologyComparison/ZariskiEtaleLifting/CompatH.lean}{CrystallineSemistableComparison.B.CompatH}
\leanref{Development/WeilBound/CoreStatements/PadicHodgeComparison/EtaleCohomologyComparison/ZariskiEtaleLifting/FirstChernCompatibleL.lean}{CrystallineSemistableComparison.B.compatH}
\leanref{Development/WeilBound/CoreStatements/PadicHodgeComparison/EtaleCohomologyComparison/ZariskiEtaleLifting/SchFnPtTw.lean}{CrystallineSemistableComparison.B.Compat}
\leanref{Development/WeilBound/CoreStatements/PadicHodgeComparison/EtaleCohomologyComparison/ZariskiEtaleLifting/FirstChernCompatibleL.lean}{CrystallineSemistableComparison.B.compat}
\leanref{Development/WeilBound/CoreStatements/PadicHodgeComparison/EtaleCohomologyComparison/BocksteinBaseChange/CrSEqOne.lean}{CrystallineSemistableComparison.cr_square}
\leanref{Development/WeilBound/CoreStatements/PadicHodgeComparison/EtaleCohomologyComparison/CechPairings/TotPairing.lean}{CrystallineSemistableComparison.P.cup_pull}
\leanref{Development/WeilBound/CoreStatements/PadicHodgeComparison/EtaleCohomologyComparison/CechPairings/FaceMapCompSuccAboveCastSucc.lean}{CrystallineSemistableComparison.P.pull_comp}
\leanref{Development/AutomorphicSpectrum/Proofs/BlockDecomposition/SquarePull.lean}{GlobalTheory.CompanionTheory.BetaRed.beta_red}
\leanref{Development/AutomorphicSpectrum/Proofs/BlockDecomposition/TwAcrFromBase.lean}{GlobalTheory.CompanionTheory.BetaGal.etaGal}
\leanref{Development/AutomorphicSpectrum/Proofs/BlockDecomposition/TwAcrFromBase.lean}{GlobalTheory.CompanionTheory.BetaGal.β_gal}
\leanref{Development/AutomorphicSpectrum/Proofs/BlockDecomposition/PullFrob.lean}{GlobalTheory.CompanionTheory.CrysPhi.etaPhi}
\leanref{Development/AutomorphicSpectrum/Proofs/BlockDecomposition/PullFrob.lean}{GlobalTheory.CompanionTheory.CrysPhi.phi_beta}
\leanref{Development/AutomorphicSpectrum/Proofs/BlockDecomposition/QuotientMapFactor.lean}{GlobalTheory.CompanionTheory.BetaCmp.β_cmp}
\leanref{Development/AutomorphicSpectrum/Proofs/BlockDecomposition/UnitsPullbackEmptyProductOne.lean}{GlobalTheory.CompanionTheory.BetaPull.beta_pull}
\leanref{Development/AutomorphicSpectrum/Proofs/BlockDecomposition/PullCup.lean}{GlobalTheory.CompanionTheory.CupBeta.beta_cup}
\leanref{Development/AutomorphicSpectrum/Proofs/BlockDecomposition/MapHClsEq.lean}{GlobalTheory.CompanionTheory.CrisOne.β_one}

\begin{proof}
Every map in the statement is induced, as in Definition~\ref{gal-realization-split-diagrams}(1), by a functor and
maps of coefficients, and a composite of two such maps is induced by the composite functor with the composite
maps of coefficients. So each identity compares two functors with maps of coefficients, and in each case below
they are related as in Lemma~\ref{gal-realization-period-prism}, which gives the identity.

(1) For the commutation with the reductions, the two composite functors, from $\mathrm{Th}_n$ or $\mathrm{Th}'_n$ of
$F$ or $F'$ to $\mathrm{Th}_m$ or $\mathrm{Th}'_m$ of $F$, send a thickening to thickenings with the same ring, ideal
and divided powers. Their maps from $A_m$ agree because $A_m\to A_n$ commutes with $\sigma$ and with $\phi$, and their
points agree, so the identity of the ring is a natural map between them. The cup product of
cochains is formed from the products of the rings $B$, and the maps of coefficients of the reductions are
identities, hence ring homomorphisms; so the reductions preserve cup products, and they send the class of the
constant cochain $1$ to the class of $1$.

(2) The maps of coefficients of $\beta_\nu$ and $\beta'_\nu$ are identities, so $\beta_\nu$ and $\beta'_\nu$
preserve cup products and units as in (1). For each identity, the two composite functors on
$\mathrm{El}(F^\sharp)$ are related by the following natural maps of thickenings.
\begin{itemize}
\item Reductions: the reduction $A_{\mathrm{cris}}(A)/\ell^m\to A_{\mathrm{cris}}(A)/\ell^n$ is a map from $g_m(A,y)$
  to $g_n(A,y)$, viewed at level $m$, and similarly for $g'$; its composite with the identity of
  $A_{\mathrm{cris}}(A)/\ell^n$ is the map of coefficients of the reduction of $H^j(F^\sharp,A_{\mathrm{cris}}/\ell^m)$.
\item $\iota_\nu$: the identity of $A_{\mathrm{cris}}(A)/\ell^\nu$ is a map from $g_\nu(A,y)$ to $g'_\nu(A,y)$, viewed
  as a thickening for $I_\nu$, since $I_\nu(A)\subset I'_\nu(A)$ with the same divided powers.
\item $\phi$: the Frobenius $\phi$ of $A_{\mathrm{cris}}(A)/\ell^\nu$ commutes with the maps from $A_\nu$, preserves
  $I'_\nu(A)$ and its divided powers, and induces the $\ell$-th power map modulo $I'_\nu(A)$. So it is a map from
  $g'_\nu(A,y)$ to the image of $g'_\nu(A,y)$ under the functor of
  Definition~\ref{gal-realization-period-thickenings}(3)(c).
\item $u^*$: the two composite functors agree, as $u(F(u_A)(y))=F'(u_A)(u(y))$ by the naturality of $u$.
\item $\sigma$: use the notation of Definition~\ref{gal-realization-period-functors}(4). There,
  $\sigma\circ\beta_\nu$ is induced by $(A,y)\mapsto g_\nu(A^\sigma,y^\sigma)$ with the maps $\tau_A$ on
  coefficients, and $\beta_\nu\circ\sigma$ by the image of $g_\nu(A,y)$ under the functor of
  Definition~\ref{gal-realization-period-thickenings}(3)(b), with identity maps. Since $\xi W(A)=\xi_\sigma W(A)$
  (proof of Lemma~\ref{gal-realization-twist}(1)), the identity of $W(A)[1/\ell]$ carries $I(A^\sigma)$ onto $I(A)$,
  compatibly with $x\mapsto x^k/k!$, and by Lemma~\ref{gal-realization-twist}(1) $\tau_A$ composed with the map from
  $A_\nu$ for $A^\sigma$ is $\sigma$ followed by the map from $A_\nu$ for $A$. The map $\bar\tau_A$ that $\tau_A$
  induces modulo the ideals satisfies $\bar\tau_A\circ u_{A^\sigma}=u_A\circ\tau_A$ on $(A^\sigma)^\sharp$, as both
  send the class of $w\in W(A)$ to the class of $w$. So $F(\bar\tau_A)$ carries the point
  $F(u_{A^\sigma})(y^\sigma)=F(u_{A^\sigma}\circ\tau_A^{-1})(y)$ to $F(u_A)(y)$, and $\tau_A$ is a map of thickenings
  from the first to the second, natural in $(A,y)$. The same holds for $g'$.
\end{itemize}

(3) The functor of a reduction keeps rings, ideals and divided powers, so it carries $J^{[r]}$ to $J^{[r]}$, and its
maps of coefficients are identities. So it carries a cocycle with values in $\mathcal O^{[r]}$ to a cocycle with
values in $\mathcal O^{[r]}$: with these maps of coefficients the reduction is also defined on $H^j(\cdot,\mathcal
O^{[r]})$, compatibly with the inclusions into $\mathcal O$, and it carries $\mathrm{Fil}^r$ into $\mathrm{Fil}^r$. Let
$a\in A_{\mathrm{cris}}$ have image $\bar a$ in $I_\nu$, resp.\ $I'_\nu$. For every thickening $T=(B,i,J,\delta,y)$,
$a$ acts on $B$ as $i(\bar a)\in J$, and $zx\in J^{[r+1]}$ for $z\in J$ and $x\in J^{[r]}$: for a generator
$x=\delta_{k_1}(x_1)\cdots\delta_{k_n}(x_n)$, $zx=\delta_1(z)\,\delta_{k_1}(x_1)\cdots\delta_{k_n}(x_n)$ is a generator
of $J^{[r+1]}$. So multiplication by $a$ is a natural map $\mathcal O^{[r]}\to\mathcal O^{[r+1]}$, compatible with the
multiplication by $a$ on $\mathcal O$. The map it induces carries the class of a cocycle $c$ with values in $\mathcal
O^{[r]}$ to that of $ac$, so $a$ times the image of the class of $c$ in $H^j_\nu(F)$, resp.\ $H'^j_\nu(F)$, lies in
$\mathrm{Fil}^{r+1}$.
% Lean: (3) is Compat.red_fil = CrisFil.red_fil (omMor_img, through pull_img and subMor for the subfunctors pdS) and
% Compat.smul_fil = BetaFil.smul_fil (pdS_smul, through img_smul and mul_pdPow). The archive states every item for all
% n <= m in N, including n = 0, where A_0 = 0 and all the groups vanish.
\end{proof}
\leanref{Development/AutomorphicSpectrum/Proofs/BlockDecomposition/OmMorImg.lean}{GlobalTheory.CompanionTheory.CrisFil.red_fil}
\leanref{Development/AutomorphicSpectrum/Proofs/BlockDecomposition/OmMorImg.lean}{GlobalTheory.CompanionTheory.CrisFil.omMor_img}
\leanref{Development/AutomorphicSpectrum/Proofs/BlockDecomposition/SmulFil.lean}{GlobalTheory.CompanionTheory.BetaFil.smul_fil}
\leanref{Development/AutomorphicSpectrum/Proofs/BlockDecomposition/SmulFil.lean}{GlobalTheory.CompanionTheory.BetaFil.pdS_smul}
\leanref{Development/WeilBound/CoreStatements/PadicHodgeComparison/EtaleCohomologyComparison/ZariskiEtaleLifting/SubThickRees.lean}{CrystallineSemistableComparison.B.mul_pdPow}

\begin{lemma}\label{gal-realization-period-limit}
Let $R$ be a ring, $s\in R$, $(M_\nu)_{\nu\geq1}$ and $(N_\nu)_{\nu\geq1}$ inverse systems of $R$-modules and
$f_\nu:M_\nu\to N_\nu$ $R$-linear maps compatible with the transition maps. Suppose that for every $\nu$, $s$ kills
the kernel of $f_\nu$ and $sN_\nu\subset f_\nu(M_\nu)$. Then $f=\varprojlim_\nu f_\nu$ satisfies $s^2x=0$ whenever
$f(x)=0$, and $s^2\varprojlim_\nu N_\nu\subset f(\varprojlim_\nu M_\nu)$.
\end{lemma}
\leanref{Development/WeilBound/CoreStatements/PadicHodgeComparison/EtaleCohomologyComparison/CechPairings/MapHComp.lean}{CrystallineSemistableComparison.P.lim_upTo}
\leanref{Development/WeilBound/CoreStatements/PadicHodgeComparison/EtaleCohomologyComparison/BocksteinBaseChange/AtlasGroupActions/LimSurjFil.lean}{CrystallineSemistableComparison.GlueBL.lim_inj}
\leanref{Development/WeilBound/CoreStatements/PadicHodgeComparison/EtaleCohomologyComparison/BocksteinBaseChange/AtlasGroupActions/LimSurjFil.lean}{CrystallineSemistableComparison.GlueBL.lim_surj}

\begin{proof}
If $f(x)=0$, then $f_\nu(x_\nu)=0$, so $sx_\nu=0$ for every $\nu$, and $s^2x=0$. Let $y=(y_\nu)$ be compatible, and
choose $x_\nu\in M_\nu$ with $f_\nu(x_\nu)=sy_\nu$. Let $r$ denote the transition maps. Then $f_\nu(r(x_{\nu+1})-
x_\nu)=s\,r(y_{\nu+1})-sy_\nu=0$, so $s\,r(x_{\nu+1})=sx_\nu$. Hence $(sx_\nu)_\nu$ is compatible, and
$f((sx_\nu)_\nu)=s^2y$.
\end{proof}

\begin{lemma}\label{gal-realization-period-t-mu}
Let $\ell$, $\mathcal O_C^\flat$, $\theta$ and $A_{\mathrm{cris}}$ be as in
Definition~\ref{gal-crystalline-test-algebras}, and let $\varepsilon\in\mathcal O_C^\flat$ satisfy $\theta([\varepsilon])=1$
and $\theta([\varepsilon^{1/\ell}])\neq1$. Let $\mu=[\varepsilon]-1$ and $t=\log[\varepsilon]=\sum_{n\geq1}
(-1)^{n-1}(n-1)!\,\mu^n/n!$ in $A_{\mathrm{cris}}$. Then $t=\mu a$ for a unit $a$ of $A_{\mathrm{cris}}$. This is a
computation: $\mu^{\ell-1}=\ell c$ with $c^\ell\in\ell A_{\mathrm{cris}}$, so $a=\sum_{n\geq0}(-1)^n\mu^n/(n+1)$
converges $\ell$-adically in $A_{\mathrm{cris}}$, and $a-1$ lies in the Jacobson radical of $A_{\mathrm{cris}}$.
\end{lemma}
\leanref{GroundTheory/PAdicHodge/FontaineLemmaPeriod.lean}{GroundTheory.PAdicHodge.ACris.exists_periodT_eq_mu_mul_unit}
\leanref{Development/WeilBound/CoreStatements/PadicHodgeComparison/EtaleCohomologyComparison/BocksteinBaseChange/AtlasGroupActions/ExistsEpsPeriodTUnit.lean}{CrystallineSemistableComparison.GlueBE.exists_eps_period_t_unit}
\leanref{Development/WeilBound/CoreStatements/PadicHodgeComparison/EtaleCohomologyComparison/ZariskiEtaleLifting/MapUSharp.lean}{CrystallineSemistableComparison.B.Eps}
\leanref{Development/WeilBound/CoreStatements/PadicHodgeComparison/EtaleCohomologyComparison/ZariskiEtaleLifting/MapUSharp.lean}{CrystallineSemistableComparison.B.Eps.μ}

\begin{definition}\label{gal-realization-period-de-rham}
Let $K$ be a field and $X$ a $K$-scheme. A \emph{nilpotent thickening} of $X$ consists of a $K$-algebra $B$, a
nilpotent ideal $J$ of $B$ with divided powers $\delta$, and a $K$-morphism $y:\operatorname{Spec}B/J\to X$. A map
$(B,J,\delta,y)\to(B',J',\delta',y')$ is a $K$-algebra map $\varphi:B\to B'$ with $\varphi(J)\subset J'$,
$\delta'_k\circ\varphi=\varphi\circ\delta_k$ on $J$ and $y'=y\circ\operatorname{Spec}\bar\varphi$. Let
$\mathrm{Th}_{\mathrm{dR}}(X)$ be the category of nilpotent thickenings, and let the \emph{de Rham cohomology}
$H^j_{\mathrm{dR}}(X)=H^j(\mathrm{Th}_{\mathrm{dR}}(X),\mathcal O)$, with $\mathcal O$ sending a thickening to the
$K$-algebra $B$, be the $K$-vector space with the cup products and unit of
Definition~\ref{gal-realization-split-diagrams}(1). A $K$-morphism $h:X\to X'$ sends $(B,J,\delta,y)$ to
$(B,J,\delta,h\circ y)$, which gives $h^*:H^j_{\mathrm{dR}}(X')\to H^j_{\mathrm{dR}}(X)$.
\end{definition}
\leanref{Development/WeilBound/CoreStatements/PadicHodgeComparison/EtaleCohomologyComparison/HarishChandraFormulaStatement.lean}{CrystallineSemistableComparison.infBase}
\leanref{Development/WeilBound/CoreStatements/PadicHodgeComparison/EtaleCohomologyComparison/HarishChandraFormulaStatement.lean}{CrystallineSemistableComparison.smFun}
\leanref{Development/WeilBound/CoreStatements/PadicHodgeComparison/EtaleCohomologyComparison/HarishChandraFormulaStatement.lean}{CrystallineSemistableComparison.HdR}
\leanref{Development/WeilBound/CoreStatements/PadicHodgeComparison/EtaleCohomologyComparison/HarishChandraFormulaStatement.lean}{CrystallineSemistableComparison.pullR}
\leanref{Development/WeilBound/CoreStatements/PadicHodgeComparison/EtaleCohomologyComparison/PerfectoidPullbacks/InclStd.lean}{CrystallineSemistableComparison.Pb.cupK}
\leanref{Development/WeilBound/CoreStatements/PadicHodgeComparison/EtaleCohomologyComparison/PerfectoidPullbacks/InclStd.lean}{CrystallineSemistableComparison.Pb.oneK}
\leanref{Development/WeilBound/CoreStatements/PadicHodgeComparison/EtaleCohomologyComparison/PerfectoidPullbacks/InclStd.lean}{CrystallineSemistableComparison.Pb.powK}

\begin{definition}\label{gal-realization-period-chern}
Let $X$ be a scheme.
\begin{enumerate}
\item A \emph{line datum} $\mathcal L$ on $X$ consists of a smooth surjective morphism $f:X_{\mathcal L}\to X$ and
  units $u(a,b)\in\Gamma(V,\mathcal O_V)^\times$ for all schemes $V$ and morphisms $a,b:V\to X_{\mathcal L}$ with
  $fa=fb$, such that $u(ah,bh)=h^*u(a,b)$ for $h:V'\to V$ and $u(a,b)u(b,b')=u(a,b')$. Its \emph{base change} along
  $X'\to X$ is the line datum on $X'$ given by $X_{\mathcal L}\times_XX'\to X'$ and the units $u(\mathrm{pr}_1a,
  \mathrm{pr}_1b)$.
\item Let $R$ be a ring, $X$ an $R$-scheme and $i:X\to\mathbf P^M_R$ a closed immersion over $R$, and let
  $U_k\subset\mathbf P^M_R$ be the open subset where the coordinate $x_k$ is invertible. A line datum $\mathcal L$ on
  $X$ is \emph{tautological for $i$} if $f$ is \'etale and there are $X$-morphisms $s_k:i^{-1}(U_k)\to X_{\mathcal
  L}$ with $u(s_ka,s_{k'}b)=(ia)^*(x_{k'}/x_k)$ for all morphisms $a:V\to i^{-1}(U_k)$ and $b:V\to i^{-1}(U_{k'})$
  with the same composite to $X$. It is \emph{tautological} if it is tautological for some such $i$.
\item Let $\mathcal C$ be a category and $G$ a functor from $\mathcal C$ to $\mathrm{Th}_\nu(X)$ or to
  $\mathrm{Th}'_\nu(X)$ (Definition~\ref{gal-realization-period-thickenings}), or to $\mathrm{Th}_{\mathrm{dR}}(X)$
  (Definition~\ref{gal-realization-period-de-rham}), with $G(c)=(B_c,J_c,\delta_c,y_c,\ldots)$. A \emph{lift} of
  $\mathcal L$ along $G$ consists of morphisms $p_c:\operatorname{Spec}B_c/J_c\to X_{\mathcal L}$ with $fp_c=y_c$
  and units $w_\varphi\in B_{c'}^\times$ for the maps $\varphi:c\to c'$ of $\mathcal C$ whose images in
  $B_{c'}/J_{c'}$ are $u(p_c\circ\operatorname{Spec}\overline{G(\varphi)},p_{c'})$. For a chain
  $s=(c_0\xrightarrow{\varphi_1}c_1\xrightarrow{\varphi_2}c_2)$ put $r(s)=G(\varphi_2)(w_{\varphi_1})\,w_{\varphi_2}\,
  w_{\varphi_2\varphi_1}^{-1}\in B_{c_2}^\times$. Then $r(s)-1\in J_{c_2}$. Indeed, put
  $a=p_{c_0}\circ\operatorname{Spec}\overline{G(\varphi_2\varphi_1)}$, $b=p_{c_1}\circ\operatorname{Spec}
  \overline{G(\varphi_2)}$ and $c=p_{c_2}$. As $G(\varphi_2\varphi_1)=G(\varphi_2)G(\varphi_1)$, the naturality of
  $u$ shows that the images of $G(\varphi_2)(w_{\varphi_1})$, $w_{\varphi_2}$ and $w_{\varphi_2\varphi_1}$ in
  $B_{c_2}/J_{c_2}$ are $u(a,b)$, $u(b,c)$ and $u(a,c)$, and $u(a,b)u(b,c)=u(a,c)$.
\item A class $x\in H^2_\nu(X)$, resp.\ $H'^2_\nu(X)$, is a \emph{first Chern class} of $\mathcal L$ if for every
  category $\mathcal C$, functor $G$ to $\mathrm{Th}_\nu(X)$, resp.\ $\mathrm{Th}'_\nu(X)$, and lift of $\mathcal L$
  along $G$, the $2$-cochain $s\mapsto\sum_{k=0}^{\ell\nu-1}(-1)^kk!\,\delta_{c_2,k+1}(r(s)-1)$ of $\mathcal C$, with
  values in $c\mapsto B_c$, is a cocycle whose class is $G^*x$. A family $(x_\nu)_{\nu\geq1}$ of first Chern classes
  $x_\nu$ of $\mathcal L$ is a \emph{first Chern family} of $\mathcal L$.
\item Let $X$ be an $\mathcal O$-scheme and $\kappa:\mathcal O\to\mathcal O_C$ as in
  Definition~\ref{gal-realization-period-thickenings}. A family $(y_\nu)_{\nu\geq1}$ with $y_\nu\in H^2(X^\sharp,
  A_{\mathrm{cris}}/\ell^\nu)$ is a \emph{first Chern class} of $\mathcal L$ if the following holds for every category
  $\mathcal C$ and functor $G:\mathcal C\to\mathrm{El}(X^\sharp)$, $G(c)=(A_c,y_c)$. Here a lift of $\mathcal L$ along
  $G$ consists of morphisms $p_c:\operatorname{Spec}A_c^\sharp\to X_{\mathcal L}$ with $fp_c=y_c$ and units
  $w_\varphi\in A_{c'}^\times$ with $\theta([w_\varphi])=u(p_c\circ\operatorname{Spec}G(\varphi)^\sharp,p_{c'})$ in
  $A_{c'}^\sharp$, and $r(s)\in A_{c_2}^\times$ is defined as in (3). For every lift and every $\nu$ there is a
  $2$-cocycle $z$ of $\mathcal C$ with values in $c\mapsto A_{\mathrm{cris}}(A_c)/\ell^\nu$ whose class is $G^*y_\nu$,
  such that for every chain $s$ the elements $([r(s)]-1)^k/k!$ of $W(A_{c_2})[1/\ell]$ lie in
  $A_{\mathrm{cris}}(A_{c_2})$ and $z(s)$ is the image of $\sum_{k=0}^{\ell\nu-1}(-1)^kk!\,([r(s)]-1)^{k+1}/(k+1)!$.
\item Let $X$ be a scheme over a field $K$. A class $\omega\in H^2_{\mathrm{dR}}(X)$ is a \emph{first Chern class} of
  $\mathcal L$ if the following holds for every $m\geq0$, category $\mathcal C$ and functor $G:\mathcal
  C\to\mathrm{Th}_{\mathrm{dR}}(X)$ such that for every $c$, $J_c^{[m]}=0$ and $y_c$ factors through an open subset
  of $X$ over which $f$ has a section. Here $J^{[m]}$ is the ideal generated by the products
  $\delta_{k_1}(x_1)\cdots\delta_{k_n}(x_n)$ with $x_i\in J$ and $k_1+\cdots+k_n\geq m$. For every lift of $\mathcal
  L$ along $G$, the $2$-cochain $s\mapsto\sum_{k=0}^{m-1}(-1)^kk!\,\delta_{c_2,k+1}(r(s)-1)$ is a cocycle whose class
  is $G^*\omega$.
\end{enumerate}
\end{definition}
\leanref{Development/WeilBound/CoreStatements/PadicHodgeComparison/EtaleCohomologyComparison/KummerCupProducts/IsKumCoc.lean}{CrystallineSemistableComparison.E.LineData}
\leanref{Development/WeilBound/CoreStatements/PadicHodgeComparison/EtaleCohomologyComparison/PerfectoidPullbacks/ConeSectionRestriction/IsTautologicalLineAt.lean}{CrystallineSemistableComparison.E.LineData.pb}
\leanref{Development/WeilBound/CoreStatements/PadicHodgeComparison/EtaleCohomologyComparison/PerfectoidPullbacks/ConeSectionRestriction/IsTautologicalLineAt.lean}{CrystallineSemistableComparison.Pb.IsTautologicalLineAt}
\leanref{Development/WeilBound/CoreStatements/PadicHodgeComparison/EtaleCohomologyComparison/PerfectoidPullbacks/ConeSectionRestriction/IsTautologicalLineAt.lean}{CrystallineSemistableComparison.Pb.IsTautologicalLine}
\leanref{Development/WeilBound/CoreStatements/PadicHodgeComparison/EtaleCohomologyComparison/PerfectoidPullbacks/ConeSectionRestriction/IsTautologicalLineAt.lean}{CrystallineSemistableComparison.Pb.ratio}
\leanref{Development/WeilBound/CoreStatements/PadicHodgeComparison/EtaleCohomologyComparison/ZariskiEtaleLifting/IsFirstChernClassOfLineBundle.lean}{CrystallineSemistableComparison.B.CLift}
\leanref{Development/WeilBound/CoreStatements/PadicHodgeComparison/EtaleCohomologyComparison/ZariskiEtaleLifting/IsFirstChernClassOfLineBundle.lean}{CrystallineSemistableComparison.B.CLift.r}
\leanref{Development/WeilBound/CoreStatements/PadicHodgeComparison/EtaleCohomologyComparison/ZariskiEtaleLifting/IsFirstChernClassOfLineBundle.lean}{CrystallineSemistableComparison.B.dlog}
\leanref{Development/WeilBound/CoreStatements/PadicHodgeComparison/EtaleCohomologyComparison/ZariskiEtaleLifting/IsFirstChernClassOfLineBundle.lean}{CrystallineSemistableComparison.B.IsFirstChernClassOfLineBundle}
\leanref{Development/WeilBound/CoreStatements/PadicHodgeComparison/EtaleCohomologyComparison/ZariskiEtaleLifting/IsFirstChernClassOfLineBundle.lean}{CrystallineSemistableComparison.B.IsFirstChernFamily}
\leanref{Development/WeilBound/CoreStatements/PadicHodgeComparison/EtaleCohomologyComparison/KummerCupProducts/IsFirstChernClassOfLineBundle.lean}{CrystallineSemistableComparison.E.Lift}
\leanref{Development/WeilBound/CoreStatements/PadicHodgeComparison/EtaleCohomologyComparison/KummerCupProducts/IsFirstChernClassOfLineBundle.lean}{CrystallineSemistableComparison.E.Lift.r}
\leanref{Development/WeilBound/CoreStatements/PadicHodgeComparison/EtaleCohomologyComparison/KummerCupProducts/IsFirstChernClassOfLineBundle.lean}{CrystallineSemistableComparison.E.IsLog}
\leanref{Development/WeilBound/CoreStatements/PadicHodgeComparison/EtaleCohomologyComparison/KummerCupProducts/IsFirstChernClassOfLineBundle.lean}{CrystallineSemistableComparison.E.Lift.IsCoc}
\leanref{Development/WeilBound/CoreStatements/PadicHodgeComparison/EtaleCohomologyComparison/KummerCupProducts/IsFirstChernClassOfLineBundle.lean}{CrystallineSemistableComparison.E.IsFirstChernClassOfLineBundle}
\leanref{Development/WeilBound/CoreStatements/PadicHodgeComparison/EtaleCohomologyComparison/PerfectoidPullbacks/ConeSectionRestriction/IsTautologicalLineAt.lean}{CrystallineSemistableComparison.Pb.CLiftK}
\leanref{Development/WeilBound/CoreStatements/PadicHodgeComparison/EtaleCohomologyComparison/PerfectoidPullbacks/ConeSectionRestriction/IsTautologicalLineAt.lean}{CrystallineSemistableComparison.Pb.plogK}
\leanref{Development/WeilBound/CoreStatements/PadicHodgeComparison/EtaleCohomologyComparison/PerfectoidPullbacks/ConeSectionRestriction/IsTautologicalLineAt.lean}{CrystallineSemistableComparison.Pb.IsFirstChernDeRham}
\leanref{Development/AutomorphicSpectrum/Proofs/BlockDecomposition/CrystallineCocycleClosed.lean}{GlobalTheory.CompanionTheory.CocCls.red_Rf}
\leanref{Development/AutomorphicSpectrum/Proofs/BlockDecomposition/CrystallineCocycleClosed.lean}{GlobalTheory.CompanionTheory.CocCls.sub_one_mem}
\leanref{Development/AutomorphicSpectrum/Proofs/BlockDecomposition/ConeLiftRatioMemIdeal.lean}{GlobalTheory.CompanionTheory.ConeUnit.r_sub_mem}
\leanref{Development/WeilBound/CoreStatements/PadicHodgeComparison/EtaleCohomologyComparison/PerfectoidPullbacks/ConeSectionRestriction/RSubOneMem.lean}{CrystallineSemistableComparison.Pb.r_sub_one_mem}

\begin{definition}\label{gal-realization-frobenius-special-fibre}
Let $\ell$, $C$, $A_{\mathrm{cris}}$, $K$ and $\mathcal O_K$ be as in Definition~\ref{gal-realization-period-datum},
$\mathbf F_q$, $q=\ell^f$, the residue field of $K$, $W=W(\mathbf F_q)\to A_{\mathrm{cris}}$ as in
Lemma~\ref{gal-realization-invariants}, and $\nu\geq1$. Let $W_\nu=W/\ell^\nu$, with the divided powers
$\gamma_k(\ell x)=\ell^kx^k/k!$ on $\ell W_\nu$, and let $W_\nu\to A_\nu$ be induced by $W\to A_{\mathrm{cris}}$
(Definition~\ref{gal-realization-period-thickenings}(1)). Let $Y$ be an $\mathcal O_K$-scheme.
\begin{enumerate}
\item A \emph{thickening of $Y$ over $W_\nu$} consists of a ring $B$, a ring map $i:W_\nu\to B$, an ideal $J$ of $B$
  with divided powers $\delta$ such that $i(\ell W_\nu)\subset J$ and $\delta_k(i(x))=i(\gamma_k(x))$ for
  $x\in\ell W_\nu$, and an $\mathcal O_K$-morphism $y:\operatorname{Spec}B/J\to Y$, where $B/J$ is an $\mathcal
  O_K$-algebra through $\mathcal O_K\to\mathbf F_q=W_\nu/\ell\xrightarrow{i}B/J$. Maps are as in
  Definition~\ref{gal-realization-period-thickenings}(2). Let $\mathrm{Th}_\nu(Y/W)$ be the category of these
  thickenings and $H^j_\nu(Y/W)=H^j(\mathrm{Th}_\nu(Y/W),\mathcal O)$, with $\mathcal O$ sending a thickening to the
  $W$-algebra $B$. A thickening over $W_\nu$ is one over $W_{\nu+1}$ through $W_{\nu+1}\to W_\nu$, and
  $L^j(Y)=\varprojlim_\nu H^j_\nu(Y/W)$ for the induced maps.
\item Let $N\geq0$. A \emph{Frobenius twist} of level $N$ is a functor $T\mapsto T^\dagger$ from
  $\mathrm{Th}'_\nu(Y)$ (Definition~\ref{gal-realization-period-thickenings}, for $\mathcal O=\mathcal O_K$ and its
  inclusion into $\mathcal O_C$) to $\mathrm{Th}_\nu(Y/W)$ together with ring isomorphisms $\tau_T:B_{T^\dagger}\to
  B_T$, natural in $T$, with the following properties: $\tau_T$ carries $J_{T^\dagger}$ onto $J_T$ and commutes with
  the divided powers; $\tau_T\circ i_{T^\dagger}$ is the composite of $W_\nu\to A_\nu$ with $i_T$; and
  $y_{T^\dagger}\circ\operatorname{Spec}\bar\tau_T=y_T\circ\operatorname{Spec}F_T^N$, where $F_T$ is the $\ell$-th
  power map of $B_T/J_T$. With the maps $\tau_T$ on coefficients it induces a map $m:H^j_\nu(Y/W)\to H'^j_\nu(Y)$
  (Definition~\ref{gal-realization-split-diagrams}(1)), semilinear for $W\to A_{\mathrm{cris}}$.
\end{enumerate}
\end{definition}
\leanref{Development/WeilBound/CoreStatements/PadicHodgeComparison/EtaleCohomologyComparison/BocksteinBaseChange/IsUnitWitt.lean}{CrystallineSemistableComparison.Res}
\leanref{Development/WeilBound/CoreStatements/PadicHodgeComparison/EtaleCohomologyComparison/BocksteinBaseChange/ResCoeff.lean}{CrystallineSemistableComparison.Wc.Wm}
\leanref{Development/WeilBound/CoreStatements/PadicHodgeComparison/EtaleCohomologyComparison/BocksteinBaseChange/ResCoeff.lean}{CrystallineSemistableComparison.Wc.pd}
\leanref{Development/WeilBound/CoreStatements/PadicHodgeComparison/EtaleCohomologyComparison/BocksteinBaseChange/ResCoeff.lean}{CrystallineSemistableComparison.Wc.bs}
\leanref{Development/WeilBound/CoreStatements/PadicHodgeComparison/EtaleCohomologyComparison/BocksteinBaseChange/ResCoeff.lean}{CrystallineSemistableComparison.Wc.L}
\leanref{Development/WeilBound/CoreStatements/PadicHodgeComparison/EtaleCohomologyComparison/BocksteinBaseChange/ResCoeff.lean}{CrystallineSemistableComparison.Wc.OM}
\leanref{Development/WeilBound/CoreStatements/PadicHodgeComparison/EtaleCohomologyComparison/BocksteinBaseChange/ResCoeff.lean}{CrystallineSemistableComparison.Wc.redMor}
\leanref{Development/WeilBound/CoreStatements/PadicHodgeComparison/EtaleCohomologyComparison/BocksteinBaseChange/ResCoeff.lean}{CrystallineSemistableComparison.Wc.HS}
\leanref{Development/WeilBound/CoreStatements/PadicHodgeComparison/EtaleCohomologyComparison/PerfectoidPullbacks/FrobeniusWittActions/PhiTTmul.lean}{CrystallineSemistableComparison.Pb.LW}
\leanref{Development/WeilBound/CoreStatements/PadicHodgeComparison/EtaleCohomologyComparison/PerfectoidPullbacks/FrobeniusWittActions/TeichmullerCongruence.lean}{CrystallineSemistableComparison.Pb.Mid.wn}
\leanref{Development/AutomorphicArithmetic/CommutativeAlgebra/EtaleCohomology/AbsInvOf.lean}{UpperTheory.Ru.crystallineComparison.MuMor}
\leanref{Development/AutomorphicArithmetic/CommutativeAlgebra/EtaleCohomology/AbsInvOf.lean}{UpperTheory.Ru.crystallineComparison.IsMu}

\begin{lemma}\label{gal-realization-frobenius-fixed}
Let $\ell$, $C$, $A_{\mathrm{cris}}$, $K$, $\mathcal O_K$, $\mathbf F_q$, $q=\ell^f$, and $g:\mathcal O_K\to\mathcal
O_C/\ell$ be as in Definition~\ref{gal-realization-frobenius-datum}, and $W$ as in
Definition~\ref{gal-realization-frobenius-special-fibre}. Let $Y$ be an $\mathcal O_K$-scheme, $F:Y_g\to Y_g$ the
$q$-power Frobenius, $N\geq0$ with $(x\bmod\ell)^{\ell^N}=g(x)$ for all $x\in\mathcal O_K$, and $\rho_N:Y^\sharp\to
Y_g$ as in Definition~\ref{gal-crystalline-test-algebras}(4). Let $\nu\geq1$, $j\geq0$, and let $m:H^j_\nu(Y/W)\to
H'^j_\nu(Y)$ be the map of a Frobenius twist of level $N$. Then for every $y\in H^j_\nu(Y/W)$ there is $X\in
H^j(Y_g,A_{\mathrm{cris}}/\ell^\nu)$ with $\phi^fX=F^*X$ and $\beta'_\nu(m(y))=\rho_N^*X$.
\end{lemma}
\leanref{Development/AutomorphicArithmetic/CommutativeAlgebra/EtaleCohomology/AbsInvOf.lean}{UpperTheory.Ru.crystallineComparison.AbsCl}
\leanref{Development/AutomorphicArithmetic/CommutativeAlgebra/EtaleCohomology/PhiAZero.lean}{UpperTheory.Ru.crystallineComparison.absCl}
\leanref{Development/AutomorphicArithmetic/CommutativeAlgebra/EtaleCohomology/PhiAZero.lean}{UpperTheory.Ru.crystallineComparison.qq}
\leanref{Development/AutomorphicArithmetic/CommutativeAlgebra/EtaleCohomology/PhiAZero.lean}{UpperTheory.Ru.crystallineComparison.gObj}
\leanref{Development/AutomorphicArithmetic/CommutativeAlgebra/EtaleCohomology/PhiAZero.lean}{UpperTheory.Ru.crystallineComparison.gam}
\leanref{Development/AutomorphicArithmetic/CommutativeAlgebra/EtaleCohomology/PhiAZero.lean}{UpperTheory.Ru.crystallineComparison.gK}
\leanref{Development/AutomorphicArithmetic/CommutativeAlgebra/EtaleCohomology/PhiAZero.lean}{UpperTheory.Ru.crystallineComparison.gK_iter}
\leanref{Development/AutomorphicArithmetic/CommutativeAlgebra/EtaleCohomology/PhiAZero.lean}{UpperTheory.Ru.crystallineComparison.phiA_fix}
\leanref{Development/AutomorphicArithmetic/CommutativeAlgebra/EtaleCohomology/PhiAZero.lean}{UpperTheory.Ru.crystallineComparison.absCl_fr}
\leanref{Development/AutomorphicArithmetic/CommutativeAlgebra/EtaleCohomology/PhiAZero.lean}{UpperTheory.Ru.crystallineComparison.psi}
\leanref{Development/AutomorphicArithmetic/CommutativeAlgebra/EtaleCohomology/PhiAZero.lean}{UpperTheory.Ru.crystallineComparison.qq_psi}
\leanref{Development/AutomorphicArithmetic/CommutativeAlgebra/EtaleCohomology/PhiAZero.lean}{UpperTheory.Ru.crystallineComparison.absCl_rd}

\begin{proof}
For a test algebra $A$ write $A_\nu(A)=A_{\mathrm{cris}}(A)/\ell^\nu$. The map $W(A)\to A_\nu(A)/I'_\nu(A)$ kills
$\ell$ and $\xi$, hence $[\ell^\flat]=\ell-\xi$. Since $W(A)/\ell W(A)=A$ through the first Witt component, it
induces a ring map $q_A:A/\ell^\flat\to A_\nu(A)/I'_\nu(A)$ with $q_A(x_0)=\bar x$ for $x\in W(A)$.

For $(A,z)$ in $\mathrm{El}(Y_g)$ let $\Gamma(A,z)$ be the thickening of $Y$ over $W_\nu$ given by $A_\nu(A)$, the
map $W_\nu\to A_\nu\to A_\nu(A)$, the ideal $I'_\nu(A)$ with its divided powers, and the point $z\circ
\operatorname{Spec}q_A$. This map from $W_\nu$ sends $\ell W_\nu$ into $I'_\nu(A)$ compatibly with the divided
powers, which on both sides are induced by $x\mapsto x^k/k!$. The point is an $\mathcal O_K$-morphism: for
$x\in\mathcal O_K$ with residue $\bar x$, both structure maps send $x$ to the class of $[s(\bar x)]\in W(A)$, where
$s:\mathbf F_q\to\mathcal O_C^\flat$ is the embedding of Lemma~\ref{gal-realization-invariants}. Maps of test
algebras give maps of these
thickenings, so $\Gamma$ is a functor from $\mathrm{El}(Y_g)$ to $\mathrm{Th}_\nu(Y/W)$. For $k\geq0$ let $X_k\in
H^j(Y_g,A_{\mathrm{cris}}/\ell^\nu)$ be the image of $y$ under the map induced by $\Gamma$ with the maps $\phi^k$ of
the rings $A_\nu(A)$ on coefficients, and put $X=X_0$.

\emph{Step 1.} $\phi X_k=X_{k+1}$, since $\phi$ of $H^j(Y_g,A_{\mathrm{cris}}/\ell^\nu)$ is induced by the identity
functor with the maps $\phi$ on coefficients. Hence $\phi^fX=X_f$.

\emph{Step 2.} $X_f=F^*X$. The residue field has $q=\ell^f$ elements, so $\phi^f$ is the identity on $W$, and the
map $\phi^f$ of $A_\nu(A)$ commutes with the maps from $W_\nu$. It preserves $I'_\nu(A)$ and its divided powers and
induces the $q$-power map modulo $I'_\nu(A)$, and $q_A$ commutes with $q$-th powers. As $F(A,z)$ is $z$ composed
with the $q$-power map of $A/\ell^\flat$, the map $\phi^f$ is a map of thickenings $\Gamma(A,z)\to\Gamma(F(A,z))$,
natural in $(A,z)$. Lemma~\ref{gal-realization-period-prism}, with $\Gamma$ and the maps $\phi^f$ on one side and
$\Gamma\circ F$ and identity maps on the other, gives $X_f=F^*X$.

\emph{Step 3.} $\beta'_\nu(m(y))=\rho_N^*X$. The left side is induced by $(A,y')\mapsto g'_\nu(A,y')^\dagger$
with the maps $\tau$ on coefficients, and the right side by $(A,y')\mapsto\Gamma(A,\rho_N(y'))$ with identity maps.
For $T=g'_\nu(A,y')$, the isomorphism $\tau_T$ is a map of thickenings $T^\dagger\to\Gamma(A,\rho_N(y'))$. Its
compatibility with the maps from $W_\nu$ and with the divided powers is part of the definition of a Frobenius
twist. For the points, let $u'_A:A^\sharp\to A_\nu(A)/I'_\nu(A)$ be as in
Definition~\ref{gal-realization-period-thickenings}(4) and $\psi_N:A^\sharp\to A/\ell^\flat$, $\theta(x)\mapsto
x_0^{\ell^N}$. Then $F_T^N\circ u'_A=q_A\circ\psi_N$, since both send $\theta(x)$ to $\bar x^{\ell^N}$. Hence
$y_{T^\dagger}\circ\operatorname{Spec}\bar\tau_T=y'\circ\operatorname{Spec}u'_A\circ\operatorname{Spec}F_T^N=
y'\circ\operatorname{Spec}\psi_N\circ\operatorname{Spec}q_A=\rho_N(y')\circ\operatorname{Spec}q_A$. The maps
$\tau_T$ are natural in $(A,y')$, and Lemma~\ref{gal-realization-period-prism} gives the claim.
\end{proof}

\begin{definition}\label{gal-realization-ingredients-cech-line}
Let $X$ be a scheme, $(U_k)_{k\in\Lambda}$ a family of open subschemes that cover $X$, and $g_{jk}$, for
$j,k\in\Lambda$, units of $\Gamma(U_j\cap U_k,\mathcal O_X)$ with $g_{kk}=1$ and $g_{ij}g_{jk}=g_{ik}$ on $U_i\cap
U_j\cap U_k$. Let $X_g=\coprod_kU_k$, and let $f:X_g\to X$ be the inclusion on each summand. For a scheme $V$,
morphisms $a,b:V\to X_g$ with $fa=fb$ and $j,k\in\Lambda$, let $V_{jk}\subset V$ be the open subscheme of the
points that $a$ maps into the summand $U_j$ and $b$ into the summand $U_k$. The $V_{jk}$ are disjoint and cover $V$,
and $fa$ maps $V_{jk}$ into $U_j\cap U_k$. Let $u(a,b)$ be the unit of $\Gamma(V,\mathcal O_V)$ whose restriction to
each $V_{jk}$ is $(fa)^*g_{jk}$. The \emph{line datum of $g$} consists of $f$ and the units $u(a,b)$.
\end{definition}
\leanref{ReciprocityTheory/ReductiveGroups/ArithmeticGeometry/SemiabelianSchemes/CubicalStructures/CubeSymmetricSwapFirst.lean}{ReciprocityTheory.ChartRestriction.G.Semi.LB}
\leanref{ReciprocityTheory/ReductiveGroups/ArithmeticGeometry/SemiabelianSchemes/CubicalStructures/CubeSymmetricSwapFirst.lean}{ReciprocityTheory.ChartRestriction.G.Semi.LB.IsLB}
\leanref{Development/AutomorphicSpectrum/Proofs/BlockDecomposition/ExistsLineDataTwistingSheaf.lean}{GlobalTheory.CompanionTheory.LineCech.exists_u}
\leanref{Development/AutomorphicSpectrum/Proofs/BlockDecomposition/ExistsLineDataTwistingSheaf.lean}{GlobalTheory.CompanionTheory.LineCech.uc}
\leanref{Development/AutomorphicSpectrum/Proofs/BlockDecomposition/ExistsLineDataTwistingSheaf.lean}{GlobalTheory.CompanionTheory.LineCech.toLine}

\begin{lemma}\label{gal-realization-ingredients-tautological}
Let $R$ be a ring, and use the notation of Definition~\ref{gal-realization-period-chern}(1) and~(2).
\begin{enumerate}
\item For a family $g$ as in Definition~\ref{gal-realization-ingredients-cech-line}, the line datum of $g$ is a line
  datum on $X$, and its map $f$ is \'etale.
\item Let $X$ be an $R$-scheme and $i:X\to\mathbf P^M_R$ a closed immersion over $R$. The line datum of the family
  $(i^{-1}(U_k))_{0\leq k\leq M}$ with $g_{jk}=i^*(x_k/x_j)$ is tautological for $i$.
\item Let $X$ and $i$ be as in (2), let $\mathcal L$ be a line datum on $X$ that is tautological for $i$, let
  $R\to R'$ be a ring map, $X'=X\times_{\operatorname{Spec}R}\operatorname{Spec}R'$, and $i':X'\to\mathbf P^M_{R'}$
  the base change of $i$. The base change of $\mathcal L$ along $X'\to X$ is tautological for $i'$.
\item If $\mathcal L$ is a tautological line datum on an $R$-scheme $X$, then $X$ is the union of the open
  subschemes $V$ over which $f$ has a section, that is, a morphism $\sigma:V\to X_{\mathcal L}$ with $f\sigma$ the
  inclusion of $V$.
\end{enumerate}
\end{lemma}
\leanref{Development/AutomorphicSpectrum/Proofs/BlockDecomposition/ExistsLineDataTwistingSheaf.lean}{GlobalTheory.CompanionTheory.LineCech.uc_nat}
\leanref{Development/AutomorphicSpectrum/Proofs/BlockDecomposition/ExistsLineDataTwistingSheaf.lean}{GlobalTheory.CompanionTheory.LineCech.uc_coc}
\leanref{Development/AutomorphicSpectrum/Proofs/BlockDecomposition/ExistsLineDataTwistingSheaf.lean}{GlobalTheory.CompanionTheory.LineCech.pullC}
\leanref{Development/AbelianVarietyTheory/Definitions/ThetaMetrics/TwistingSheafOneChartEqPreimageStandardChart.lean}{GlobalTheory.ProjSpace.O1}
\leanref{Development/WeilBound/CoreStatements/PadicHodgeComparison/EtaleCohomologyComparison/PerfectoidPullbacks/ConeSectionRestriction/IsFirstChernDeRhamUnique.lean}{CrystallineSemistableComparison.Pb.isO1At_exists}
\leanref{Development/AutomorphicSpectrum/Proofs/BlockDecomposition/ExistsTwistingLineDataProjective.lean}{GlobalTheory.CompanionTheory.ConeTwistingSheafMap.main}
\leanref{Development/WeilBound/CoreStatements/PadicHodgeComparison/EtaleCohomologyComparison/PerfectoidPullbacks/ConeSectionRestriction/IsFirstChernDeRhamUnique.lean}{CrystallineSemistableComparison.Pb.isTwistingSheaf_pullback}
\leanref{Development/AutomorphicSpectrum/Proofs/BlockDecomposition/TwistingSheafBaseChange.lean}{GlobalTheory.CompanionTheory.ConeTwistingSheaf.isTwistingSheaf_pullback}
\leanref{Development/WeilBound/CoreStatements/PadicHodgeComparison/EtaleCohomologyComparison/PerfectoidPullbacks/ConeSectionRestriction/Precomp.lean}{CrystallineSemistableComparison.Pb.𝒱D}
\leanref{Development/WeilBound/CoreStatements/PadicHodgeComparison/EtaleCohomologyComparison/PerfectoidPullbacks/ConeSectionRestriction/RSubOneMem.lean}{CrystallineSemistableComparison.Pb.sSup_𝒱D}
\leanref{Development/AutomorphicSpectrum/Proofs/BlockDecomposition/SupremumTwistingOpensEqTop.lean}{GlobalTheory.CompanionTheory.ConeCover.main}

\begin{proof}
(1) On each summand $f$ is the inclusion of an open subscheme, so $f$ is \'etale, hence smooth, and it is surjective
because the $U_k$ cover $X$. The units $u(a,b)$ exist and are unique because the $V_{jk}$ are open, disjoint and
cover $V$. Let $h:V'\to V$. The open subschemes of $V'$ attached to $ah$ and $bh$ are the $h^{-1}(V_{jk})$, and
$(fah)^*g_{jk}=h^*((fa)^*g_{jk})$, so $u(ah,bh)=h^*u(a,b)$. Let $a,b,b':V\to X_g$ with $fa=fb=fb'$, and let
$V_{ijk}$ be the open subscheme of the points that $a$, $b$ and $b'$ map into the summands $U_i$, $U_j$ and $U_k$.
These cover $V$, and on $V_{ijk}$ the identity $g_{ij}g_{jk}=g_{ik}$ gives
$u(a,b)u(b,b')=(fa)^*(g_{ij}g_{jk})=(fa)^*g_{ik}=u(a,b')$.

(2) The $U_k$ cover $\mathbf P^M_R$, so the $i^{-1}(U_k)$ cover $X$, and the units $i^*(x_k/x_j)$ satisfy $g_{kk}=1$
and the cocycle identity. By (1) the line datum of $g$ is a line datum whose map $f$ is \'etale. Let
$s_k:i^{-1}(U_k)\to X_g$ be the inclusion of the $k$-th summand, an $X$-morphism. Let $a:V\to i^{-1}(U_k)$ and
$b:V\to i^{-1}(U_{k'})$ have the same composite to $X$. Then $s_ka$ and $s_{k'}b$ map all of $V$ into the summands
$U_k$ and $U_{k'}$, so $u(s_ka,s_{k'}b)=(fs_ka)^*g_{kk'}=(ia)^*(x_{k'}/x_k)$.

(3) Let $\pi:\mathbf P^M_{R'}\to\mathbf P^M_R$ and $\pi_X:X'\to X$ be the projections, so that $i\pi_X=\pi i'$. The
morphism $i'$ is a closed immersion, as a base change of $i$, and the map $f':X_{\mathcal L}\times_XX'\to X'$ of the
base change of $\mathcal L$ is \'etale, as a base change of $f$. Let $U'_k\subset\mathbf P^M_{R'}$ be the open
subscheme where $x_k$ is invertible. Then $\pi$ maps $U'_k$ into $U_k$, $\pi^*(x_{k'}/x_k)=x_{k'}/x_k$, and $\pi_X$
induces $m_k:i'^{-1}(U'_k)\to i^{-1}(U_k)$. Let $s_k$ be the sections of
Definition~\ref{gal-realization-period-chern}(2) for $i$, and let $s'_k:i'^{-1}(U'_k)\to X_{\mathcal L}\times_XX'$
have the components $s_km_k$ and the inclusion. For $a:V\to i'^{-1}(U'_k)$ and $b:V\to i'^{-1}(U'_{k'})$ with the
same composite to $X'$, the units of the base change give
\[
 u'(s'_ka,s'_{k'}b)=u(s_km_ka,s_{k'}m_{k'}b)=(im_ka)^*(x_{k'}/x_k)=(i'a)^*\pi^*(x_{k'}/x_k)=(i'a)^*(x_{k'}/x_k).
\]

(4) Let $i$ and the $s_k$ be as in Definition~\ref{gal-realization-period-chern}(2). Each $s_k$ is a section of $f$
over $i^{-1}(U_k)$, and these open subschemes cover $X$.
\end{proof}

\begin{definition}\label{gal-realization-ingredients-atlas}
Let $\mathcal O$ be a ring, $F$ a point functor over $\mathcal O$ (Definition~\ref{gal-realization-period-functors}) and
$d\geq0$.
\begin{enumerate}
\item A \emph{subfunctor} $V$ of $F$ consists of subsets $V(B)\subset F(B)$ for all rings $B$ with $F(\varphi)(V(B))
  \subset V(B')$ for ring maps $\varphi:B\to B'$; with the relation of $F$ it is a point functor over $\mathcal O$. For
  subfunctors $V_\alpha$, $\alpha\in S$, $\bigcap_{\alpha\in S}V_\alpha$ is the subfunctor $B\mapsto\bigcap_{\alpha\in
  S}V_\alpha(B)$.
\item A \emph{small presentation} of dimension $d$ of $F$ consists of a ring $R$, a ring map $a:\mathcal O\to R$, units
  $T_1,\ldots,T_d$ of $R$ such that the ring map $\mathcal O[x_1,\ldots,x_d]\to R$ extending $a$ with $x_k\mapsto T_k$ is
  \'etale, and an element $y_0\in F(R)$ that lies over $a$, such that for every ring $B$ and every $y\in F(B)$ that lies
  over some ring map $\mathcal O\to B$ there is exactly one ring map $g:R\to B$ with $F(g)(y_0)=y$.
\item A \emph{Zariski atlas} of dimension $d$ of $F$ consists of a finite set $E$ and subfunctors $V_\alpha$ of $F$,
  $\alpha\in E$, with the following properties.
  \begin{enumerate}
  \item[(a)] For every ring $B$ and $y\in F(B)$ there are $s_1,\ldots,s_m\in B$ that generate the unit ideal, such that
    for every $k$ and every ring map $\varphi:B\to B'$ with $\varphi(s_k)$ a unit, $F(\varphi)(y)\in V_\alpha(B')$ for
    some $\alpha$.
  \item[(b)] Let $\alpha\in E$, $\varphi:B\to B'$ a surjective ring map whose kernel lies in the Jacobson radical of
    $B$, and $y\in F(B)$. If $F(\varphi)(y)\in V_\alpha(B')$, then $y\in V_\alpha(B)$.
  \item[(c)] For every nonempty $S\subset E$, $\bigcap_{\alpha\in S}V_\alpha$ has a small presentation of some
    dimension $d'\leq d$.
  \end{enumerate}
\end{enumerate}
\end{definition}
\leanref{Development/WeilBound/CoreStatements/PadicHodgeComparison/EtaleCohomologyComparison/ZariskiEtaleLifting/SubThickRees.lean}{CrystallineSemistableComparison.B.Fn.Sub}
\leanref{Development/WeilBound/CoreStatements/PadicHodgeComparison/EtaleCohomologyComparison/ZariskiEtaleLifting/SubThickRees.lean}{CrystallineSemistableComparison.B.Fn.Sub.fn}
\leanref{Development/WeilBound/CoreStatements/PadicHodgeComparison/EtaleCohomologyComparison/ZariskiEtaleLifting/SubThickRees.lean}{CrystallineSemistableComparison.B.Fn.Sub.inter}
\leanref{Development/WeilBound/CoreStatements/PadicHodgeComparison/EtaleCohomologyComparison/ZariskiEtaleLifting/SubThickRees.lean}{CrystallineSemistableComparison.B.Fn.Small}
\leanref{Development/WeilBound/CoreStatements/PadicHodgeComparison/EtaleCohomologyComparison/ZariskiEtaleLifting/SubThickRees.lean}{CrystallineSemistableComparison.B.Fn.Atlas}

\begin{lemma}\label{gal-realization-ingredients-atlas-scheme}
Let $\mathcal O$ be a ring and $X$ a smooth, separated and quasi-compact $\mathcal O$-scheme. Then $F_X$
(Definition~\ref{gal-realization-period-functors}(2)) has a Zariski atlas of some dimension
(Definition~\ref{gal-realization-ingredients-atlas}).
\end{lemma}
\leanref{Development/WeilBound/CoreStatements/PadicHodgeComparison/EtaleCohomologyComparison/ZariskiEtaleLifting/PropsA.lean}{CrystallineSemistableComparison.B.P.atlas}
\leanref{Development/AutomorphicSpectrum/Proofs/BlockDecomposition/PreimageOpenTopOfJacobsonSurjection.lean}{GlobalTheory.CompanionTheory.SchAtlas.exists_atlas}

\begin{proof}
\emph{Step 1: charts with units.} Let $x\in X$. As $X$ is smooth over $\mathcal O$, $x$ has an affine open
neighbourhood $V$ with $t_1,\ldots,t_n\in\Gamma(V,\mathcal O_X)$ such that the $\mathcal O$-algebra map $\mathcal
O[x_1,\ldots,x_n]\to\Gamma(V,\mathcal O_X)$, $x_k\mapsto t_k$, is \'etale \cite[Corollaire~17.11.4]{Gro67}. Let $P$ be
the prime ideal of $\Gamma(V,\mathcal O_X)$ at $x$. For each $k$ put $e_k=0$ if $t_k\notin P$ and $e_k=1$ otherwise; then
$t'_k=t_k+e_k\notin P$. The map $x_k\mapsto t'_k$ is the composite of the automorphism $x_k\mapsto x_k+e_k$ of $\mathcal
O[x_1,\ldots,x_n]$ with the \'etale map above, so it is \'etale. Let $h=\prod_kt'_k$. Then $h\notin P$, so the basic
open $U=D(h)\subset V$ is an affine open neighbourhood of $x$, the restrictions $T_k$ of the $t'_k$ to $U$ are units, and
$\mathcal O[x_1,\ldots,x_n]\to\Gamma(U,\mathcal O_X)$, $x_k\mapsto T_k$, is \'etale by
Lemma~\ref{gal-realization-thickening-cover}(1).

\emph{Step 2: the atlas.} As $X$ is quasi-compact, finitely many of the opens $U$ of Step 1, say $U_\alpha$ with units
$T_{\alpha,1},\ldots,T_{\alpha,n_\alpha}$ for $\alpha\in E$, cover $X$. Let $d=\max_\alpha n_\alpha$, and let
$V_\alpha(B)$ be the set of morphisms $\operatorname{Spec}B\to X$ with image in $U_\alpha$, a subfunctor of $F_X$.

(a) Let $y:\operatorname{Spec}B\to X$, and let $S$ be the set of $s\in B$ such that $y$ maps the basic open $D(s)$ into
some $U_\alpha$. If $S$ generated a proper ideal, it would lie in a maximal ideal $M$. Then $y(M)\in U_\alpha$ for some
$\alpha$, and as $y^{-1}(U_\alpha)$ is open there is $r\notin M$ with $D(r)\subset y^{-1}(U_\alpha)$; so $r\in S\subset
M$, a contradiction. Hence finitely many $s_1,\ldots,s_m\in S$ generate the unit ideal. If $\varphi(s_k)$ is a unit of
$B'$, every prime ideal of $B'$ pulls back to a prime ideal of $B$ that does not contain $s_k$, so $y\circ
\operatorname{Spec}\varphi$ maps into $U_\alpha$ when $y$ maps $D(s_k)$ into $U_\alpha$.

(b) Let $\varphi:B\to B'$ be surjective with kernel in the Jacobson radical, and let $y\circ\operatorname{Spec}\varphi$
have image in $U_\alpha$. Let $\mathfrak p$ be a prime ideal of $B$ and $M\supset\mathfrak p$ a maximal ideal. Then $M$
contains the kernel of $\varphi$, so $M=\varphi^{-1}(\varphi(M))$ with $\varphi(M)$ a prime ideal of $B'$, and $y(M)\in
U_\alpha$. As $\mathfrak p\subset M$, $y(\mathfrak p)$ is a generization of $y(M)$, and it lies in the open set
$U_\alpha$.

(c) Let $S\subset E$ be nonempty and $\alpha\in S$. As $X$ is separated, $W=\bigcap_{\beta\in S}U_\beta$ is affine (proof
of Lemma~\ref{gal-realization-thickening-cover}(2)), and $W\subset U_\alpha$. The restrictions of the $T_{\alpha,k}$ to
$W$ are units, and the map $\mathcal O[x_1,\ldots,x_{n_\alpha}]\to R=\Gamma(W,\mathcal O_X)$ they define is \'etale by
Lemma~\ref{gal-realization-thickening-cover}(1). The subfunctor $\bigcap_{\beta\in S}V_\beta$ sends $B$ to the set of
morphisms $\operatorname{Spec}B\to X$ with image in $W$. Let $a:\mathcal O\to R$ be the structure map and $y_0$ the
composite of the isomorphism $\operatorname{Spec}R\cong W$ with the inclusion; it lies over $a$. A morphism
$y:\operatorname{Spec}B\to X$ with image in $W$ factors through the open immersion $W\to X$ in exactly one way, so
$y=y_0\circ\operatorname{Spec}g$ for exactly one ring map $g:R\to B$. This is a small presentation of dimension
$n_\alpha\leq d$.
% Lean: SchAtlas.exists_atlas, with exists_unit_chart (Step 1: exists_etale_chart, et_shift, et_res), the cover from
% isCompact_univ.elim_nhds_subcover, subOf for V_alpha, exists_unimodular_cover_refining_opens for (a),
% preimage_open_top_of_jacobson_surjection for (b), and small_of_affine with IsAffineOpen.biInf for (c).
\end{proof}
\leanref{Development/AutomorphicSpectrum/Proofs/BlockDecomposition/PreimageOpenTopOfJacobsonSurjection.lean}{GlobalTheory.CompanionTheory.SchAtlas.exists_unit_chart}
\leanref{Development/AutomorphicSpectrum/Proofs/BlockDecomposition/PreimageOpenTopOfJacobsonSurjection.lean}{GlobalTheory.CompanionTheory.SchAtlas.exists_unimodular_cover_refining_opens}
\leanref{Development/AutomorphicSpectrum/Proofs/BlockDecomposition/PreimageOpenTopOfJacobsonSurjection.lean}{GlobalTheory.CompanionTheory.SchAtlas.preimage_open_top_of_jacobson_surjection}
\leanref{Development/AutomorphicSpectrum/Proofs/BlockDecomposition/PreimageOpenTopOfJacobsonSurjection.lean}{GlobalTheory.CompanionTheory.SchAtlas.small_of_affine}

\begin{lemma}\label{gal-realization-ingredients-bounds}
Let $\ell$, $C$ and $A_{\mathrm{cris}}$ be as in Definition~\ref{gal-crystalline-test-algebras}. For all integers
$d,j\geq0$ there is $N\geq0$ with the following property. Let $\mathcal O$ be a ring, $F$ a point functor over
$\mathcal O$ with a Zariski atlas of dimension $d$ (Definition~\ref{gal-realization-ingredients-atlas}), $\kappa:\mathcal
O\to\mathcal O_C$ a ring map, $\varepsilon$ and $\mu$ as in Lemma~\ref{gal-realization-period-t-mu}, and $\nu\geq1$, and
use the notation of Definition~\ref{gal-realization-period-thickenings}. Then $\mu^N$ kills the kernels of $\beta_\nu$ on
$H^j_\nu(F)$ and of $\beta'_\nu$ on $H'^j_\nu(F)$; $\mu^NH^j(F^\sharp,A_{\mathrm{cris}}/\ell^\nu)$ lies in the image of
$\beta'_\nu$; and for every $y\in H^j(F^\sharp,A_{\mathrm{cris}}/\ell^\nu)$ there is $x\in\mathrm{Fil}^NH^j_\nu(F)$ with
$\beta_\nu(x)=\mu^Ny$.
\end{lemma}
\leanref{Development/WeilBound/CoreStatements/PadicHodgeComparison/EtaleCohomologyComparison/ZariskiEtaleLifting/MapUSharp.lean}{CrystallineSemistableComparison.B.Main}
\leanref{Development/WeilBound/CoreStatements/PadicHodgeComparison/EtaleCohomologyComparison/ZariskiEtaleLifting/FirstChernCompatibleL.lean}{CrystallineSemistableComparison.B.main}
\leanref{Development/WeilBound/CoreStatements/PadicHodgeComparison/EtaleCohomologyComparison/ZariskiEtaleLifting/SubsingletonH.lean}{CrystallineSemistableComparison.B.main_of}
\leanref{Development/WeilBound/CoreStatements/PadicHodgeComparison/EtaleCohomologyComparison/CechPairings/InstAddCommGroup.lean}{CrystallineSemistableComparison.UpTo}

\begin{proof}
For a smooth $\mathcal O$-scheme $X$, the published comparison concerns the $\ell$-adic completion $\mathfrak X$ of
$X\otimes_{\mathcal O}\mathcal O_C$, a smooth formal scheme over $\mathcal O_C$, and its fibre $\mathfrak X_0$ over
$\mathcal O_C/\ell$. Bhatt, Morrow and Scholze define $A\Omega_{\mathfrak X}=L\eta_\mu(R\pi_*A_{\mathrm{inf},\mathfrak
X})$, where $\pi$ is the projection from the pro-\'etale site of the generic fibre of $\mathfrak X$ to the Zariski site
of $\mathfrak X$ \cite[Definition~9.1]{BMS18}, and prove that $A\Omega_{\mathfrak X}\hat\otimes_{A_{\mathrm{inf}}}
A_{\mathrm{cris}}$ computes the crystalline cohomology of $\mathfrak X_0$ over $A_{\mathrm{cris}}$
\cite[Theorem~12.1]{BMS18}. They also bound how far $L\eta_\mu$ is from the identity: for a complex $M$ in $D^{\geq0}$
with $H^0(M)$ free of $\mu$-torsion and $j\geq0$, there are natural maps from $\mu^j\tau^{\leq j}M$ to $\tau^{\leq
j}L\eta_\mu M$ and from $L\eta_\mu M$ to $M$ whose composites, in either order, are multiplication by $\mu^j$ in the
degrees up to $j$ \cite[Lemmas~6.9 and~6.10]{BMS18}. Bhatt and Scholze prove the crystalline comparison for prismatic
cohomology \cite[Theorem~1.8(1) and Theorem~5.2]{BS22}, and for an $\ell$-completely smooth $\mathcal O_C$-algebra they
identify $A\Omega$ with the Frobenius pullback of its prismatic cohomology over $A_{\mathrm{inf}}$
\cite[Theorem~17.2]{BS22}. Alpöge proves the lemma first for point functors with a small presentation of dimension at
most $d$, then glues over the atlas, and then passes to every level $\nu$.
\openstep{The proof of the lemma is not yet written out; it goes beyond the published results above, as its groups are
the cohomology of the categories of thickenings and of test algebras, its $N$ depends only on $d$ and $j$, $F$ is any
point functor with a Zariski atlas of dimension $d$, and for $\beta_\nu$ the preimages lie in $\mathrm{Fil}^N$.}
% Lean: B.main = main_of (D.out l). main_of is global_of (local_of hD) followed by level_all: local_of is the case
% F.Small d' (B.Local, for the Rees module of the filtration flt; it uses P.glue, P.elift, P.tower_ex, P.thick_act,
% P.theta_ex, P.local_ni, P.line_univ and P.desc hD), global_of glues it over the atlas with P.glob, with exponent
% (max over d' <= d of N(d', j+1)) (2j+3), and level_all passes to each level (level_zero, level_pos). Locators checked
% against arXiv:1602.03148v3 (BMS18) and arXiv:1905.08229v4 (BS22). Size (tools/lean_size.py): 17,950 lines for B.main_of
% and D.out together (--union); local_of 11,527 (with D.Out as a hypothesis), D.out 5,810, global_of 2,874 and level_all
% 628, which overlap.
\end{proof}

\begin{lemma}\label{gal-realization-ingredients-crystalline-chern}
Let $(R,I,\gamma,p)$, $\mathcal O$, $X$ and $\kappa$ be as in Definition~\ref{gal-realization-thickening-nilpotent}, with
$p\neq0$, let $N\geq0$ be an integer with $N!=0$ in $R$, and let $\mathcal L$ be a line datum on $X$
(Definition~\ref{gal-realization-period-chern}) such that $X$ is the union of the open subschemes over which $f$ has a
section. Then there is exactly one $x\in H^2(\mathrm{Th}(X/R),\mathcal O)$ with the following property: for every category
$\mathcal C$, functor $G:\mathcal C\to\mathrm{Th}(X/R)$ and lift of $\mathcal L$ along $G$, defined as in
Definition~\ref{gal-realization-period-chern}(3), the $2$-cochain $s\mapsto\sum_{k=0}^{N-1}(-1)^kk!\,\delta_{c_2,k+1}(r(s)-1)$
is a cocycle whose class is $G^*x$. For $\mathcal O$, $X$, $\kappa$ and $\nu$ as in
Definition~\ref{gal-realization-period-thickenings}, the base $(A_\nu,I_\nu,\gamma,\ell)$ and the map $u\circ\kappa$ to
$A_\nu/I_\nu$, $\mathrm{Th}(X/A_\nu)=\mathrm{Th}_\nu(X)$, and $(\ell\nu)!=0$ in $A_\nu$ as $\ell^\nu$ divides
$(\ell\nu)!$; with $N=\ell\nu$ the class $x$ is then the unique first Chern class of $\mathcal L$ in $H^2_\nu(X)$
(Definition~\ref{gal-realization-period-chern}(4)). This is the crystalline first Chern class of Berthelot and Illusie
\cite{BI70}, the image of the class of the line bundle under the logarithm
$\log(1+x)=\sum_{k\geq0}(-1)^kk!\,\delta_{k+1}(x)$ on the divided power ideal, written as a cocycle on every family of
thickenings with a lift of $\mathcal L$. The sum stops at $k=N-1$ because $k!=0$ in $R$ for $k\geq N$.
\end{lemma}
\leanref{Development/AutomorphicSpectrum/Proofs/BlockDecomposition/ExistsUniqueChernClassNilpotentBaseStatement.lean}{GlobalTheory.T.existsUnique_chern_class_nilpotent_base}
\leanref{Development/AutomorphicSpectrum/Proofs/BlockDecomposition/ExistsUniqueChernClassNilpotentBase.lean}{GlobalTheory.CompanionTheory.CrystallineChernClass.existsUnique_crystalline_chern_class}
\leanref{Development/AutomorphicSpectrum/Proofs/BlockDecomposition/ChernClassReductionCompatible.lean}{GlobalTheory.CompanionTheory.FirstChernClassFactorial.pow_dvd_factorial_mul}
\leanref{Development/AutomorphicSpectrum/Proofs/BlockDecomposition/ChernClassReductionCompatible.lean}{GlobalTheory.CompanionTheory.FirstChernClassFactorial.factorial_acr}

\begin{lemma}\label{gal-realization-ingredients-chern-reduction}
Let $\mathcal O$, $X$ and $\kappa$ be as in Definition~\ref{gal-realization-period-thickenings}, let $\mathcal L$ be
a line datum on $X$, and $\nu\geq1$. The reduction $H^2_{\nu+1}(X)\to H^2_\nu(X)$ carries every first Chern class of
$\mathcal L$ to a first Chern class of $\mathcal L$ (Definition~\ref{gal-realization-period-chern}(4)).
\end{lemma}
\leanref{Development/AutomorphicSpectrum/Proofs/BlockDecomposition/ChernClassReductionCompatible.lean}{GlobalTheory.CompanionTheory.FirstChernClassFactorial.chern_class_reduction_compatible}
\leanref{Development/AutomorphicSpectrum/Proofs/BlockDecomposition/ChernClassReductionCompatible.lean}{GlobalTheory.CompanionTheory.FirstChernClassFactorial.isFirstChern_base_change}
\leanref{Development/AutomorphicSpectrum/Proofs/BlockDecomposition/ChernClassReductionCompatible.lean}{GlobalTheory.CompanionTheory.FirstChernClassFactorial.dlog_eq_of_factorial}

\begin{proof}
Let $x$ be a first Chern class of $\mathcal L$ in $H^2_{\nu+1}(X)$, let $G:\mathcal C\to\mathrm{Th}_\nu(X)$ be a
functor from a category $\mathcal C$, and $(p_c,w_\varphi)$ a lift of $\mathcal L$ along $G$. Let
$\rho:\mathrm{Th}_\nu(X)\to\mathrm{Th}_{\nu+1}(X)$ be the functor of
Definition~\ref{gal-realization-period-thickenings}(3)(a), which keeps the ring, the ideal, the divided powers and
the point of a thickening. Then $(p_c,w_\varphi)$ is also a lift of $\mathcal L$ along $\rho G$,
with the same elements $r(s)$. As $x$ is a first Chern class, the cochain
$s\mapsto\sum_{k=0}^{\ell(\nu+1)-1}(-1)^kk!\,\delta_{c_2,k+1}(r(s)-1)$ is a cocycle whose class is $(\rho G)^*x$.
Each $B_c$ is an algebra over $A_\nu=A_{\mathrm{cris}}/\ell^\nu$, and $\ell^\nu$ divides $(\ell\nu)!$, hence $k!$ for
$k\geq\ell\nu$; so the terms with $k\geq\ell\nu$ vanish, and the cochain is the one of
Definition~\ref{gal-realization-period-chern}(4) at level $\nu$. The reduction is induced by $\rho$ with the identity
maps of the rings $B$, so $G^*$ of the reduction of $x$ is $(\rho G)^*x$. Hence the reduction of $x$ is a first
Chern class.
\end{proof}

\begin{lemma}\label{gal-realization-ingredients-sharp-units}
Let $A$ be a test algebra, with untilt $A^\sharp$ and $\theta:W(A)\to A^\sharp$
(Definition~\ref{gal-crystalline-test-algebras}). Every unit of $A^\sharp$ is $\theta([w])$ for a unit $w$ of $A$.
\end{lemma}
\leanref{Development/WeilBound/CoreStatements/PadicHodgeComparison/EtaleCohomologyComparison/AcrisKummerUnits/ValuationOfValuationModSquare.lean}{CrystallineSemistableComparison.PAlg.IsTest.sharp_unit}
\leanref{Development/WeilBound/CoreStatements/PadicHodgeComparison/EtaleCohomologyComparison/AcrisKummerUnits/ValuationOfValuationModSquare.lean}{CrystallineSemistableComparison.PAlg.IsTest.sharp_surjective}
\leanref{Development/WeilBound/CoreStatements/PadicHodgeComparison/EtaleCohomologyComparison/AcrisKummerUnits/FiltrationSpanDividedPowers.lean}{CrystallineSemistableComparison.RL.tilt_ex}
\leanref{Development/AutomorphicSpectrum/Proofs/BlockDecomposition/CompatibleRootsLiftToTilt.lean}{GlobalTheory.CompanionTheory.Untilt.exists_compatible_roots_of_sharp}
\leanref{Development/WeilBound/CoreStatements/PadicHodgeComparison/EtaleCohomologyComparison/AcrisKummerUnits/FiltrationSpanDividedPowers.lean}{CrystallineSemistableComparison.RL.unit_lift}
\leanref{Development/AutomorphicSpectrum/Proofs/BlockDecomposition/IsUnitOfIsUnitSharp.lean}{GlobalTheory.CompanionTheory.Untilt.isUnit_of_isUnit_sharp}

\begin{proof}
Let $u$ be a unit of $A^\sharp$. By Lemma~\ref{gal-passage-untilt}(3) and Lemma~\ref{gal-passage-lift-square},
$A^\sharp$ is absolutely integrally closed, so there are $u_0=u,u_1,u_2,\ldots$ in $A^\sharp$ with
$u_{i+1}^\ell=u_i$. By Lemma~\ref{gal-passage-untilt}(2) there are $b_i\in A$ with $\theta([b_i])\equiv u_i$ modulo
$\ell$, and then $b_{i+1}^\ell-b_i\in\ell^\flat A$. As $A$ is perfect of characteristic $\ell$,
$b_{i+n+1}^{\ell^{n+1}}-b_{i+n}^{\ell^n}=(b_{i+n+1}^\ell-b_{i+n})^{\ell^n}$ lies in $(\ell^\flat)^nA$, so
$a_i=\lim_nb_{i+n}^{\ell^n}$ exists in the $\ell^\flat$-adically complete ring $A$. Since $A$ is $\ell^\flat$-adically
separated, $a_{i+1}^\ell=a_i$ and $a_i\equiv b_i$ modulo $\ell^\flat$, so $\theta([a_i])\equiv u_i$ modulo $\ell$.
Hence $\theta([a_0])=\theta([a_n])^{\ell^n}$ and $u=u_n^{\ell^n}$ agree modulo $\ell^{n+1}$ for every $n$, and
$\theta([a_0])=u$ because $A^\sharp$ is $\ell$-adically separated (Lemma~\ref{gal-passage-untilt}(1)). Finally, the
image of $a_0$ in $A/\ell^\flat A=A^\sharp/\ell A^\sharp$ is a unit, so $a_0c\in1+\ell^\flat A$ for some $c\in A$.
As $A$ is $\ell^\flat$-adically complete, $\ell^\flat$ lies in its Jacobson radical, $a_0c$ is a unit, and so is
$w=a_0$.
\end{proof}

\begin{lemma}\label{gal-realization-ingredients-roots}
Let $q\geq1$, let $U$ be a scheme on which $q$ is invertible, and let $\zeta\in\Gamma(U,\mathcal O_U)$ be a root of the
cyclotomic polynomial $\Phi_q$, a \emph{primitive $q$-th root of unity} on $U$.
\begin{enumerate}
\item $\zeta^q=1$, $\zeta^a-\zeta^b$ is a unit for all integers $a\not\equiv b\pmod q$, and
  $r^q-1=\prod_{a\in\mathbf Z/q}(r-\zeta^a)$ for every $r\in\Gamma(U,\mathcal O_U)$.
\item For every $r\in\Gamma(U,\mathcal O_U)$ with $r^q=1$ there is a locally constant function $f:U\to\mathbf Z/q$ such
  that $r=\zeta^a$ on the open and closed subscheme $f^{-1}(a)$ for every $a\in\mathbf Z/q$. Write $r=\zeta^f$.
\end{enumerate}
Item (1) is a computation with the cyclotomic polynomial.
\end{lemma}
\leanref{Development/WeilBound/CoreStatements/PadicHodgeComparison/EtaleCohomologyComparison/KummerCupProducts/FirstChernCompatibleSystem.lean}{CrystallineSemistableComparison.E.PrimAlg}
\leanref{Development/AutomorphicSpectrum/Proofs/BlockDecomposition/IsUnitOneSubPow.lean}{GlobalTheory.CompanionTheory.PrimAlgProof.primAlg}
\leanref{Development/WeilBound/CoreStatements/PadicHodgeComparison/EtaleCohomologyComparison/KummerCupProducts/FirstChernCompatibleSystem.lean}{CrystallineSemistableComparison.E.RootPow}
\leanref{Development/AutomorphicSpectrum/Proofs/BlockDecomposition/MemBasicOpenOfIsUnitSub.lean}{GlobalTheory.CompanionTheory.RootPowLC.rootPow}
\leanref{Development/AutomorphicSpectrum/Proofs/BlockDecomposition/MemBasicOpenOfIsUnitSub.lean}{GlobalTheory.CompanionTheory.RootPowLC.exists_isPow}
\leanref{Development/WeilBound/CoreStatements/PadicHodgeComparison/EtaleCohomologyComparison/KummerCupProducts/IsKumCoc.lean}{CrystallineSemistableComparison.E.IsPow}
\leanref{Development/AbelianVarietyTheory/Definitions/Divisors/PullHlComp.lean}{GlobalTheory.EtC.Prim}

\begin{proof}
(2) By (1), $\prod_a(r-\zeta^a)=r^q-1=0$. Let $u\in U$. The factors $r-\zeta^a$ are not all invertible in the local
ring $\mathcal O_{U,u}$, since their product is $0$. At most one of them is not invertible, since the difference
$\zeta^b-\zeta^a$ of two of them is a unit and the non-units of a local ring form an ideal. Let $f(u)$ be the $a$ for which
$r-\zeta^a$ is not invertible at $u$. Then $f^{-1}(a)$ is the intersection of the finitely many open subsets on which
the $r-\zeta^b$, $b\neq a$, are invertible, so $f$ is locally constant. On $f^{-1}(a)$ these $r-\zeta^b$ are units, and
the product formula gives $r-\zeta^a=0$ there.
\end{proof}

\begin{definition}\label{gal-realization-ingredients-kummer-maps}
Let $q\geq1$. The sheaves below are sheaves on the big \'etale site of $\mathbf Z[1/q]$-schemes.
\begin{enumerate}
\item For $i\in\mathbf Z$ let $(\mathbf Z/q)(i)=\mu_q^{\otimes i}$, and $(\mathbf Z/q)(i)=\operatorname{Hom}(\mu_q^{\otimes
  -i},\mathbf Z/q)$ for $i<0$, so that $(\mathbf Z/q)(0)=\mathbf Z/q$ and $(\mathbf Z/q)(1)=\mu_q$. A primitive $q$-th
  root of unity $\zeta$ on $U$ (Lemma~\ref{gal-realization-ingredients-roots}) gives a basis $\zeta^{\otimes i}$ of
  $(\mathbf Z/q)(i)$ over $U$ (for $i<0$, the dual basis of $\zeta^{\otimes-i}$), with $(\zeta^a)^{\otimes
  i}=a^i\zeta^{\otimes i}$ for $a\in(\mathbf Z/q)^\times$. The tensor product gives products $(\mathbf
  Z/q)(i)\times(\mathbf Z/q)(i')\to(\mathbf Z/q)(i+i')$ with $\zeta^{\otimes i}\zeta^{\otimes i'}=\zeta^{\otimes(i+i')}$,
  and for $q\mid Q$ the reductions $(\mathbf Z/Q)(i)\to(\mathbf Z/q)(i)$ send $\zeta^{\otimes i}$ to
  $(\zeta^{Q/q})^{\otimes i}$.
\item Let $\mu_q\to I^\bullet$ be an injective resolution. A \emph{Kummer map} for $I^\bullet$ consists of maps
  $\operatorname{kum}:\Gamma(U,\mathcal O_U)^\times\to\Gamma(U,I^0)$ for all $\mathbf Z[1/q]$-schemes $U$, compatible
  with pullback along morphisms $V\to U$, such that $\operatorname{kum}(\zeta v)=\operatorname{kum}(v)+\zeta$ for every
  unit $v$ and every primitive $q$-th root of unity $\zeta$ on $U$, where $\zeta\in\mu_q(U)\subset\Gamma(U,I^0)$.
  Kummer maps exist, as $I^0$ is injective.
\end{enumerate}
\end{definition}
\leanref{Development/AbelianVarietyTheory/Definitions/Divisors/PullHlComp.lean}{GlobalTheory.EtC.tw}
\leanref{Development/AbelianVarietyTheory/Definitions/Divisors/PullHlComp.lean}{GlobalTheory.EtC.twGen}
\leanref{Development/AbelianVarietyTheory/Definitions/Divisors/PullHlComp.lean}{GlobalTheory.EtC.IsKum}
\leanref{Development/AbelianVarietyTheory/Definitions/Divisors/PullHlComp.lean}{GlobalTheory.EtC.kumT}
\leanref{Development/AbelianVarietyTheory/Proofs/DivisorProofs/IsKummerTorsorSplittingKummerTorsorSplitting.lean}{GlobalTheory.EtC.KummerSplit.exists_isKum}
\leanref{Development/WeilBound/CoreStatements/PadicHodgeComparison/EtaleCohomologyComparison/KummerCupProducts/IsKumCoc.lean}{CrystallineSemistableComparison.E.InKum}

\begin{lemma}\label{gal-realization-ingredients-twisted-edge}
Let $\mathcal C$ be a category, $(U_c)$ a split diagram over $\mathcal C$ with points $(y_c)$ in a scheme $X$
(Definition~\ref{gal-realization-split-diagrams}), and $q\geq1$ an integer invertible on $X$ and on every $U_c$. Let
$\zeta_c$ be primitive $q$-th roots of unity on the $U_c$ with $U_\varphi^*\zeta_c=\zeta_{c'}$ for $\varphi:c\to c'$, let
$i\in\mathbf Z$, and let $(\mathbf Z/q)(i)\to I^\bullet$ be an injective resolution
(Definition~\ref{gal-realization-ingredients-kummer-maps}).
\begin{enumerate}
\item The maps $L_q(c)\to\Gamma(U_c,(\mathbf Z/q)(i))$, $f\mapsto f\zeta_c^{\otimes i}$, are isomorphisms compatible with
  the $U_\varphi$. With them in place of $L_q(c)=\Gamma(U_c,\mathbf Z/q)$, Lemma~\ref{gal-realization-edge-map}(1) holds
  for $H^j(X,(\mathbf Z/q)(i))$. Write $\varepsilon^{(i)}_y:H^j(X,(\mathbf Z/q)(i))\to H^j(\mathcal C,L_q)$ for the
  homomorphism it gives. For $i=0$ it is $\varepsilon_y$.
\item In the situation of Lemma~\ref{gal-realization-edge-map}(2), with primitive $q$-th roots of unity $\zeta'_c$ on the
  $U'_c$ such that $g_c^*\zeta_{m(c)}=\zeta'_c$, one has $m^*\circ\varepsilon^{(i)}_y=\varepsilon^{(i)}_{y'}\circ u^*$.
\item $\varepsilon^{(i+i')}_y(x\cup x')=\varepsilon^{(i)}_y(x)\cup\varepsilon^{(i')}_y(x')$ for $x\in H^j(X,(\mathbf
  Z/q)(i))$ and $x'\in H^{j'}(X,(\mathbf Z/q)(i'))$, with the products of
  Definition~\ref{gal-realization-ingredients-kummer-maps}(1).
\item If $\zeta$ is a primitive $q$-th root of unity on $X$ with $y_c^*\zeta=\zeta_c$ for all $c$, then the class $w\in
  H^0(X,(\mathbf Z/q)(i))$ of $\zeta^{\otimes i}$ satisfies $\varepsilon^{(i)}_y(w)=1$.
\end{enumerate}
\end{lemma}
\leanref{Development/WeilBound/CoreStatements/PadicHodgeComparison/EtaleCohomologyComparison/KummerCupProducts/IsKumCoc.lean}{CrystallineSemistableComparison.E.kap}
\leanref{Development/WeilBound/CoreStatements/PadicHodgeComparison/EtaleCohomologyComparison/KummerCupProducts/IsKumCoc.lean}{CrystallineSemistableComparison.E.Probe}
\leanref{Development/WeilBound/CoreStatements/PadicHodgeComparison/EtaleCohomologyComparison/KummerCupProducts/IsKumCoc.lean}{CrystallineSemistableComparison.E.Probe.res}
\leanref{Development/WeilBound/CoreStatements/PadicHodgeComparison/EtaleCohomologyComparison/KummerCupProducts/FirstChernCompatibleSystem.lean}{CrystallineSemistableComparison.E.KapSpec}
\leanref{Development/WeilBound/CoreStatements/PadicHodgeComparison/EtaleCohomologyComparison/KummerCupProducts/FirstChernCompatibleSystem.lean}{CrystallineSemistableComparison.E.EtSpec}
\leanref{Development/WeilBound/CoreStatements/PadicHodgeComparison/EtaleCohomologyComparison/KummerCupProducts/FirstChernCompatibleSystem.lean}{CrystallineSemistableComparison.E.EtNat}
\leanref{Development/WeilBound/CoreStatements/PadicHodgeComparison/EtaleCohomologyComparison/KummerCupProducts/FirstChernCompatibleSystem.lean}{CrystallineSemistableComparison.E.EtCup}
\leanref{Development/WeilBound/CoreStatements/PadicHodgeComparison/EtaleCohomologyComparison/KummerCupProducts/FirstChernCompatibleSystem.lean}{CrystallineSemistableComparison.E.KapMul}
\leanref{Development/WeilBound/CoreStatements/PadicHodgeComparison/EtaleCohomologyComparison/KummerCupProducts/FirstChernCompatibleSystem.lean}{CrystallineSemistableComparison.E.EtRoot}
\leanref{Development/WeilBound/CoreStatements/PadicHodgeComparison/EtaleCohomologyComparison/KummerCupProducts/EtaleChernClassEqKummerClass.lean}{CrystallineSemistableComparison.E.Main.kapSpec}
\leanref{Development/WeilBound/CoreStatements/PadicHodgeComparison/EtaleCohomologyComparison/KummerCupProducts/EtaleChernClassEqKummerClass.lean}{CrystallineSemistableComparison.E.Main.kapMul}
\leanref{Development/WeilBound/CoreStatements/PadicHodgeComparison/EtaleCohomologyComparison/KummerCupProducts/EtaleChernClassEqKummerClass.lean}{CrystallineSemistableComparison.E.Main.etCup}
\leanref{Development/WeilBound/CoreStatements/PadicHodgeComparison/EtaleCohomologyComparison/KummerCupProducts/EtaleChernClassEqKummerClass.lean}{CrystallineSemistableComparison.E.Main.etRoot}

\begin{proof}
(1) Over $U_c$ the sheaf $(\mathbf Z/q)(i)$ is free with basis $\zeta_c^{\otimes i}$, so its sections are the
$f\zeta_c^{\otimes i}$ with $f$ locally constant, and $U_\varphi^*\zeta_c^{\otimes i}=\zeta_{c'}^{\otimes i}$. The proof
of Lemma~\ref{gal-realization-edge-map}(1) uses of $L_q(c)$ only that it is the kernel of $d:\Gamma(U_c,I^0)\to
\Gamma(U_c,I^1)$, compatibly with the $U_\varphi$, and that $\Gamma(U_c,-)$ keeps the resolution exact because every
\'etale surjection onto $U_c$ has a section. Both hold here, so that proof applies. For $i=0$, $\zeta_c^{\otimes0}=1$,
and the restriction of an injective resolution of $\mathbf Z/q$ on all schemes to $\mathbf Z[1/q]$-schemes is an
injective resolution \cite[Tag 03F3]{Sta}; by the uniqueness in Lemma~\ref{gal-realization-edge-map}(1),
$\varepsilon^{(0)}_y=\varepsilon_y$.

(2) The proof of Lemma~\ref{gal-realization-edge-map}(2) applies, since pullback along $g_c$ carries
$f\zeta_{m(c)}^{\otimes i}$ to $(g_c^*f)\zeta'^{\otimes i}_c$.

(3) The proof of Lemma~\ref{gal-realization-cup-staircase} applies to a pairing $\mu$ of the resolutions of
$(\mathbf Z/q)(i)$ and $(\mathbf Z/q)(i')$ into the resolution of $(\mathbf Z/q)(i+i')$ that satisfies the Leibniz rule,
induces the product of Definition~\ref{gal-realization-ingredients-kummer-maps}(1) in degree $0$ and computes the cup
product. The identity $\mu(f,f')=ff'$ for $f,f'\in L_q(c)$ used there becomes $\mu(f\zeta_c^{\otimes i},
f'\zeta_c^{\otimes i'})=ff'\zeta_c^{\otimes(i+i')}$.

(4) For $j=0$ the conditions of Lemma~\ref{gal-realization-edge-map}(1) say that $\varepsilon^{(i)}_y(w)$ is the class
of the $0$-cocycle $(y_c^*\zeta^{\otimes i})_c=(\zeta_c^{\otimes i})_c$, which corresponds to the constant function $1$.
\end{proof}

\begin{lemma}\label{gal-realization-ingredients-kummer-class}
Let $\ell$ be a prime, $X$ a scheme on which $\ell$ is invertible, and $\mathcal L$ a line datum on $X$, with
$f:X_{\mathcal L}\to X$ and units $u(b,c)$ (Definition~\ref{gal-realization-period-chern}(1)). For $n\geq1$ let
$\mu_{\ell^n}\to I_n^\bullet$ be an injective resolution and $\operatorname{kum}_n$ a Kummer map for it
(Definition~\ref{gal-realization-ingredients-kummer-maps}). There are classes $x_n\in H^2(X,\mu_{\ell^n})$, $n\geq1$,
such that the reduction $\mu_{\ell^{n+1}}\to\mu_{\ell^n}$, $v\mapsto v^\ell$, carries $x_{n+1}$ to $x_n$, and such that
for every $n$, with $q=\ell^n$, $I^\bullet=I_n^\bullet$ and $\operatorname{kum}=\operatorname{kum}_n$, there are
\begin{itemize}
\item $a\in\Gamma(X,I^2)$ with $da=0$ representing $x_n$,
\item $h_b\in\Gamma(V,I^1)$ for all schemes $V$ and morphisms $b:V\to X_{\mathcal L}$,
\item $k_{b,c}\in\Gamma(V,I^0)$ for all morphisms $b,c:V\to X_{\mathcal L}$ with $fb=fc$,
\end{itemize}
compatible with pullback along morphisms $V'\to V$, with the following properties, where $v$ and $w$ are units of
$\Gamma(V,\mathcal O_V)$:
\begin{enumerate}
\item $dh_b=(fb)^*a$;
\item $h_c-h_b=d(k_{b,c}-\operatorname{kum}(v))$ if $v^q=u(b,c)$;
\item $k_{c,e}-k_{b,e}+k_{b,c}=\operatorname{kum}(v)+\operatorname{kum}(w)-\operatorname{kum}(vw)$ if $fb=fc=fe$,
  $v^q=u(b,c)$ and $w^q=u(c,e)$.
\end{enumerate}
\end{lemma}
\leanref{Development/WeilBound/CoreStatements/PadicHodgeComparison/EtaleCohomologyComparison/KummerCupProducts/FirstChernCompatibleSystem.lean}{CrystallineSemistableComparison.E.FirstChernCoboundary}
\leanref{Development/AutomorphicSpectrum/Proofs/BlockDecomposition/CechTripleSubfaceZero.lean}{GlobalTheory.CompanionTheory.KummerPrimitive.main}
\leanref{Development/WeilBound/CoreStatements/PadicHodgeComparison/EtaleCohomologyComparison/KummerCupProducts/FirstChernCompatibleSystem.lean}{CrystallineSemistableComparison.E.FirstChernCompatibleSystem}
\leanref{Development/AutomorphicSpectrum/Proofs/BlockDecomposition/ChernClassCompatibleSystem.lean}{GlobalTheory.CompanionTheory.ChernSys.chern_class_compatible_system}
\leanref{Development/WeilBound/CoreStatements/PadicHodgeComparison/EtaleCohomologyComparison/KummerCupProducts/FirstChernCompatibleSystem.lean}{CrystallineSemistableComparison.E.LineMap}
\leanref{Development/AutomorphicSpectrum/Proofs/BlockDecomposition/LineMapExists.lean}{GlobalTheory.CompanionTheory.CechLine.lineMap_exists}
\leanref{Development/AbelianVarietyTheory/Definitions/Divisors/PullHlComp.lean}{GlobalTheory.EtC.c1}
\leanref{Development/AbelianVarietyTheory/Proofs/DivisorProofs/ExistsEtaleCocycleFirstChernClassTateTwist.lean}{GlobalTheory.EtC.firstChernClassTateTwist_cocycle_representative}
\leanref{Development/AbelianVarietyTheory/Definitions/Divisors/DivisorSupportCondition.lean}{GlobalTheory.EtC.A2}
\leanref{Development/AbelianVarietyTheory/Proofs/DivisorProofs/SmoothSurjectiveCechDescentCohomology.lean}{GlobalTheory.EtC.smooth_surjective_cech_descent_cohomology}
\leanref{Development/AbelianVarietyTheory/Proofs/DivisorProofs/CoefficientReductionFirstChernClassTateTwist.lean}{GlobalTheory.EtC.RedChern.coefficientReduction_firstChernClass_tateTwist}

\begin{proof}
Fix $n$ and write $q$, $I^\bullet$ and $\operatorname{kum}$. For $p\geq0$ let $X_p$ be the fibre product of $p+1$
copies of $X_{\mathcal L}$ over $X$, so that a morphism $V\to X_p$ is a tuple $(b_0,\ldots,b_p)$ of morphisms $V\to
X_{\mathcal L}$ with $fb_0=\cdots=fb_p$. Let $T^{p,s}=\Gamma(X_p,I^s)$; an element is a family of sections
$\gamma(b_0,\ldots,b_p)\in\Gamma(V,I^s)$, compatible with pullback. Let $\delta:T^{p,s}\to T^{p+1,s}$ be the alternating
sum of the face maps, and give the total complex the differential $D=\delta+(-1)^pd$ on $T^{p,s}$. Sending $a\in
\Gamma(X,I^s)$ to $((fb)^*a)_b\in T^{0,s}$ is a map of complexes $e$ from $\Gamma(X,I^\bullet)$ to the total complex,
since $(fb)^*a=(fc)^*a$ when $fb=fc$. It is a quasi-isomorphism. Indeed, the restriction of $I^s$ to the big \'etale
site of $X$ is injective \cite[Tag 03F3]{Sta}, and as $f$ is smooth and surjective, the proof of
\cite[Tag 06XJ]{Sta}, through \cite[Tags 06XE, 06XF and 06XH]{Sta}, shows that $0\to\Gamma(X,I^s)\to T^{0,s}\to
T^{1,s}\to\cdots$ is exact. The acyclic assembly lemma \cite[Lemma~2.7.3]{Wei94}, applied to the double complex with
the column $\Gamma(X,I^\bullet)$ added, then shows that the cone of $e$ is acyclic.

\emph{A cocycle.} Let $(b,c)$ be a morphism $V\to X_1$. \'Etale locally on $V$, $u(b,c)$ has a $q$-th root $v$, by the
Kummer sequence \cite[Tag 03PL]{Sta}, and there is a primitive $q$-th root of unity $\zeta$, since $\mathcal
O_V[T]/(\Phi_q(T))$ is finite \'etale and surjective over $V$. Two $q$-th roots of $u(b,c)$ differ by some $r$ with
$r^q=1$, which is locally a power $\zeta^a$ by Lemma~\ref{gal-realization-ingredients-roots}(2), and
$\operatorname{kum}(\zeta^av)=\operatorname{kum}(v)+a\zeta$. As $d$ kills $\zeta\in\mu_q(V)$, the local sections
$d\operatorname{kum}(v)$ do not depend on $v$ and glue to $\gamma^{1,1}(b,c)\in\Gamma(V,I^1)$. In the same way, for
$(b,c,e):V\to X_2$ and local $q$-th roots $v$ of $u(b,c)$ and $w$ of $u(c,e)$, the local sections
$-(\operatorname{kum}(v)+\operatorname{kum}(w)-\operatorname{kum}(vw))$ do not depend on the roots, since replacing $v$
by $\zeta^av$ adds $a\zeta$ to $\operatorname{kum}(v)$ and to $\operatorname{kum}(vw)$, and similarly for $w$; they
glue to $\gamma^{2,0}(b,c,e)\in\Gamma(V,I^0)$. With $\gamma^{0,2}=0$ this gives $\gamma=(\gamma^{0,2},\gamma^{1,1},
\gamma^{2,0})$ of total degree $2$, and $D\gamma=0$. The component in $T^{1,2}$ is
$-dd\operatorname{kum}(v)=0$. The component in $T^{2,1}$ is $\delta\gamma^{1,1}+d\gamma^{2,0}$, which vanishes at
$(b,c,e)$ because $vw$ is a $q$-th root of $u(b,e)=u(b,c)u(c,e)$. The component in $T^{3,0}$ vanishes on expanding its
four terms. Let $x_n$ be the class with $e(x_n)=[\gamma]$.

\emph{The sections.} Let $a$ represent $x_n$. Then $\gamma-e(a)=D\beta$ for some $\beta=(\beta^{0,1},\beta^{1,0})$ of
total degree $1$. Put $h_b=-\beta^{0,1}(b)$ and $k_{b,c}=-\beta^{1,0}(b,c)$. The components of $D\beta=\gamma-e(a)$ in
$T^{0,2}$, $T^{1,1}$ and $T^{2,0}$ are $d\beta^{0,1}=-e(a)$, $\delta\beta^{0,1}-d\beta^{1,0}=\gamma^{1,1}$ and
$\delta\beta^{1,0}=\gamma^{2,0}$. Evaluated at $b$, at $(b,c)$ with $v^q=u(b,c)$ (where $\gamma^{1,1}(b,c)=
d\operatorname{kum}(v)$) and at $(b,c,e)$, these are (1), (2) and (3).

\emph{Reduction.} Let $\rho:I_{n+1}^\bullet\to I_n^\bullet$ lift the reduction $\mu_{\ell^{n+1}}\to\mu_{\ell^n}$, and
let $\gamma_n$ be the cocycle above for $n$. For $(b,c):V\to X_1$ and a local $\ell^{n+1}$-th root $v$ of $u(b,c)$ put
$m(b,c)=\rho(\operatorname{kum}_{n+1}(v))-\operatorname{kum}_n(v^\ell)$. Replacing $v$ by $\zeta v$, for a primitive
$\ell^{n+1}$-th root of unity $\zeta$, adds $\rho(\zeta)=\zeta^\ell$ to both terms, as $\zeta^\ell$ is a primitive
$\ell^n$-th root of unity; so these sections glue to $m\in T^{1,0}$ for $I_n^\bullet$. As $v^\ell$ is an $\ell^n$-th
root of $u(b,c)$, the components of $\rho(\gamma_{n+1})-\gamma_n$ are $0$, $dm$ and $-\delta m$, which are those of
$D(0,-m)$. As $\rho$ commutes with $e$, the reduction of $x_{n+1}$ is $x_n$.
\end{proof}

\begin{lemma}\label{gal-realization-ingredients-kummer-cocycle}
Let $\ell$, $X$, $\mathcal L$, $f$, $u$ and $n$ be as in Lemma~\ref{gal-realization-ingredients-kummer-class}, with
$q=\ell^n$ and the class $x_n$ given there. Let $\mathcal C$, $(U_c)$, $(y_c)$ and $(\zeta_c)$ be as in
Lemma~\ref{gal-realization-ingredients-twisted-edge}, with points in $X$. Let $p_c:U_c\to X_{\mathcal L}$ be morphisms
with $fp_c=y_c$, let $r_\varphi\in\Gamma(U_{c'},\mathcal O)^\times$ for $\varphi:c\to c'$ satisfy
$r_\varphi^q=u(p_cU_\varphi,p_{c'})$, and let $z$ be a $2$-cochain of $\mathcal C$ with values in $L_q$ such that
$U_{\varphi_2}^*(r_{\varphi_1})\,r_{\varphi_2}\,r_{\varphi_2\varphi_1}^{-1}=\zeta_{c_2}^{z(s)}$
(Lemma~\ref{gal-realization-ingredients-roots}(2)) for every chain $s=(c_0\xrightarrow{\varphi_1}c_1
\xrightarrow{\varphi_2}c_2)$. Then $z$ is a cocycle, and $\varepsilon^{(1)}_y(x_n)$ is its class.
\end{lemma}
\leanref{Development/WeilBound/CoreStatements/PadicHodgeComparison/EtaleCohomologyComparison/KummerCupProducts/FirstChernCompatibleSystem.lean}{CrystallineSemistableComparison.E.EtaleFirstChernKummer}
\leanref{Development/AutomorphicSpectrum/Proofs/BlockDecomposition/EtaleChernClassEqKummerClass.lean}{GlobalTheory.CompanionTheory.EtaleChernClass.etale_chern_class_eq_kummer_class}
\leanref{Development/WeilBound/CoreStatements/PadicHodgeComparison/EtaleCohomologyComparison/KummerCupProducts/IsKumCoc.lean}{CrystallineSemistableComparison.E.IsKumCoc}
\leanref{Development/WeilBound/CoreStatements/PadicHodgeComparison/EtaleCohomologyComparison/KummerCupProducts/FirstChernCompatibleSystem.lean}{CrystallineSemistableComparison.E.KumPow}
\leanref{Development/AutomorphicSpectrum/Proofs/BlockDecomposition/IsKumPow.lean}{GlobalTheory.CompanionTheory.KumPowProof.kumPow}

\begin{proof}
Let $a$, $h$ and $k$ be as in Lemma~\ref{gal-realization-ingredients-kummer-class} for $x_n$, with the resolution
$I^\bullet$ and the Kummer map $\operatorname{kum}$, and for $g\in L_q(c)$ let $\iota(g)\in\Gamma(U_c,I^0)$ be the
image of $g\zeta_c$, as in Lemma~\ref{gal-realization-ingredients-twisted-edge}(1) for $i=1$. If $\rho=\zeta_c^g$ and
$v$ is a unit on $U_c$, then $\operatorname{kum}(\rho v)=\operatorname{kum}(v)+\iota(g)$: on $g^{-1}(a)$, with $0\leq a<q$,
this is the defining property of $\operatorname{kum}$ applied $a$ times, and these open subschemes cover $U_c$.

Let $b_0$ be the $0$-cochain $c\mapsto h_{p_c}$ and $b_1$ the $1$-cochain $(\varphi:c\to c')\mapsto
k_{p_cU_\varphi,p_{c'}}-\operatorname{kum}(r_\varphi)$. By (1), $db_0=(y_c^*a)_c$. By the compatibility of $h$ with
pullback and by (2) for $v=r_\varphi$, $\delta b_0(\varphi)=h_{p_{c'}}-h_{p_cU_\varphi}=db_1(\varphi)$. For a chain $s$
as above put $\varphi_3=\varphi_2\varphi_1$, $\alpha=p_{c_0}U_{\varphi_3}=p_{c_0}U_{\varphi_1}U_{\varphi_2}$,
$\alpha'=p_{c_1}U_{\varphi_2}$, $\alpha''=p_{c_2}$, $v=U_{\varphi_2}^*(r_{\varphi_1})$ and $w=r_{\varphi_2}$. By the
naturality of $u$, $v^q=u(\alpha,\alpha')$ and $w^q=u(\alpha',\alpha'')$, and by the compatibility of $k$ and
$\operatorname{kum}$ with pullback and (3),
\[
 \delta b_1(s)=b_1(\varphi_2)-b_1(\varphi_3)+U_{\varphi_2}^*b_1(\varphi_1)
 =(k_{\alpha',\alpha''}-k_{\alpha,\alpha''}+k_{\alpha,\alpha'})-\operatorname{kum}(w)+\operatorname{kum}(r_{\varphi_3})
 -\operatorname{kum}(v)=\operatorname{kum}(r_{\varphi_3})-\operatorname{kum}(vw).
\]
By hypothesis $vw=\zeta_{c_2}^{z(s)}r_{\varphi_3}$, so $\delta b_1(s)=-\iota(z(s))$. Hence $b_0,b_1$ and $z$ satisfy
the conditions of Lemma~\ref{gal-realization-edge-map}(1), in the form of
Lemma~\ref{gal-realization-ingredients-twisted-edge}(1), for the representative $a$ of $x_n$. As $\iota$ is injective
and commutes with $\delta$, and $\iota(\delta z)=-\delta\delta b_1=0$, $z$ is a cocycle, and $\varepsilon^{(1)}_y(x_n)$
is its class.
\end{proof}

\begin{lemma}\label{gal-realization-ingredients-logarithm}
Let $A$ be a test algebra (Definition~\ref{gal-crystalline-test-algebras}), $\nu\geq1$, and $\varepsilon$ and
$t=\log[\varepsilon]$ as in Lemma~\ref{gal-realization-period-t-mu}, and let $\varepsilon_A$ be the image of
$\varepsilon$ in $A$.
\begin{enumerate}
\item If $w\in W(A)$ and $\theta(w)=1$, the elements $(w-1)^k/k!$ of $W(A)[1/\ell]$, $k\geq1$, lie in
  $A_{\mathrm{cris}}(A)$. Let $\log_\nu(w)\in A_{\mathrm{cris}}(A)/\ell^\nu$ be the image of
  $\sum_{k=0}^{\ell\nu-1}(-1)^kk!\,(w-1)^{k+1}/(k+1)!$.
\item If $\theta(w)=\theta(w')=1$, then $\log_\nu(ww')=\log_\nu(w)+\log_\nu(w')$. In particular $\log_\nu(1)=0$ and
  $\log_\nu(w^m)=m\log_\nu(w)$ for $m\geq0$.
\item If $\theta(w)=\theta(w')=1$ and $e\in W(A)$ is an idempotent with $ew=ew'$, then $e\log_\nu(w)=e\log_\nu(w')$.
\item $\theta([\varepsilon_A])=1$, and $\log_\nu([\varepsilon_A])$ is the image of $t$.
\end{enumerate}
Item (2) is a computation: the truncations after $\ell\nu$ terms of the two sides of $\log((1+x)(1+y))=\log(1+x)+
\log(1+y)$ differ by $(\ell\nu)!$ times an element of the divided power algebra, and $\ell^\nu$ divides $(\ell\nu)!$.
Item (4) is a computation with the series defining $t$, whose terms of index at least $\ell\nu$ lie in $\ell^\nu
A_{\mathrm{cris}}$.
\end{lemma}
\leanref{Development/WeilBound/CoreStatements/PadicHodgeComparison/EtaleCohomologyComparison/KummerCupProducts/IsFirstChernClassOfLineBundle.lean}{CrystallineSemistableComparison.E.IsLog}
\leanref{Development/WeilBound/CoreStatements/PadicHodgeComparison/EtaleCohomologyComparison/KummerCupProducts/MapUgCToUg.lean}{CrystallineSemistableComparison.E.LogEx}
\leanref{Development/AutomorphicSpectrum/Proofs/BlockDecomposition/IsLogUnique.lean}{GlobalTheory.CompanionTheory.CrisLog.logEx}
\leanref{Development/WeilBound/CoreStatements/PadicHodgeComparison/EtaleCohomologyComparison/KummerCupProducts/MapUgCToUg.lean}{CrystallineSemistableComparison.E.LogMul}
\leanref{Development/AutomorphicSpectrum/Proofs/BlockDecomposition/FactorialAcr.lean}{GlobalTheory.CompanionTheory.LogMulPf.logMul}
\leanref{Development/WeilBound/CoreStatements/PadicHodgeComparison/EtaleCohomologyComparison/KummerCupProducts/EtaleComparisonMapOfFacts.lean}{CrystallineSemistableComparison.E.LogAlg}
\leanref{Development/AutomorphicSpectrum/Proofs/BlockDecomposition/NatDegreeLgLe.lean}{GlobalTheory.CompanionTheory.LogAlgPf.logAlg}
\leanref{Development/WeilBound/CoreStatements/PadicHodgeComparison/EtaleCohomologyComparison/KummerCupProducts/MapUgCToUg.lean}{CrystallineSemistableComparison.E.LogEps}
\leanref{Development/AutomorphicSpectrum/Proofs/BlockDecomposition/AcrisZeroImageDividedPower.lean}{GlobalTheory.CompanionTheory.LogEpsP.logEps}

\begin{proof}
(1) As $\theta(w-1)=0$, $w-1=v\xi$ with $v\in W(A)$, and $(w-1)^k/k!=v^k\xi^k/k!$. (2) Apply the first claim to
$1\cdot1$ and to $w^m\cdot w$. (3) Since $e$ is idempotent and $e(w-1)=e(w'-1)$, $e(w-1)^k=(e(w-1))^k=e(w'-1)^k$ for
$k\geq1$. (4) $\theta([\varepsilon_A])$ is the image of $\theta([\varepsilon])=1$ under $\mathcal O_C\to A^\sharp$.
\end{proof}

\begin{lemma}\label{gal-realization-ingredients-log-roots}
Let $A$ be a test algebra, $n\geq1$ and $q=\ell^n$, and let $\varepsilon$, $t$ and $\log_n$ be as in
Lemma~\ref{gal-realization-ingredients-logarithm}. Let $\varepsilon^{1/q}$ be the $q$-th root of $\varepsilon$ in the
perfect ring $\mathcal O_C^\flat$; then $\theta([\varepsilon^{1/q}])\in\mathcal O_C$ is a primitive $q$-th root of
unity, as its $q$-th power is $\theta([\varepsilon])=1$ and its $(q/\ell)$-th power is
$\theta([\varepsilon^{1/\ell}])\neq1$. Let $\zeta_A$ be its image in
$A^\sharp[1/\ell]$. Let $u\in A$, and let $f$ be a locally constant function on $\operatorname{Spec}A^\sharp[1/\ell]$
with values in $\mathbf Z/q$ such that the image of $\theta([u])$ in $A^\sharp[1/\ell]$ is $\zeta_A^f$
(Lemma~\ref{gal-realization-ingredients-roots}(2)). Then $\theta([u^q])=1$, and $\log_n([u^q])=t\,\lambda_A(f)$ in
$A_{\mathrm{cris}}(A)/\ell^n$, with $\lambda_A$ as in Lemma~\ref{gal-realization-locally-constant}.
\end{lemma}
\leanref{Development/WeilBound/CoreStatements/PadicHodgeComparison/EtaleCohomologyComparison/KummerCupProducts/MapUgCToUg.lean}{CrystallineSemistableComparison.E.LcLog}
\leanref{Development/AutomorphicSpectrum/Proofs/BlockDecomposition/IdemMulEqZero.lean}{GlobalTheory.CompanionTheory.LcLogProof.lcLog}
\leanref{Development/AutomorphicSpectrum/Proofs/BlockDecomposition/IdemMulEqZero.lean}{GlobalTheory.CompanionTheory.LcLogProof.log_on_part}
\leanref{Development/WeilBound/CoreStatements/PadicHodgeComparison/EtaleCohomologyComparison/KummerCupProducts/MapUgCToUg.lean}{CrystallineSemistableComparison.E.EpsRoot}
\leanref{Development/AutomorphicSpectrum/Proofs/BlockDecomposition/UntiltIterateSymmEpsilon.lean}{GlobalTheory.CompanionTheory.EpsRootProof.epsRoot}

\begin{proof}
Let $\varepsilon_q$ be the image of $\varepsilon^{1/q}$ in $A$. As $\mathcal O_C\to A^\sharp$ is induced by
$W(\mathcal O_C^\flat)\to W(A)$, $\theta([\varepsilon_q])$ maps to $\zeta_A$, and $\theta([\varepsilon_q])^q=
\theta([\varepsilon_A])=1$. Let $(e_a)_{a\in\mathbf Z/q}$ be the idempotents of $f$
(Lemma~\ref{gal-realization-locally-constant}(1)) and $E_a=[e_a]$. On the open and closed subset of
$\operatorname{Spec}A^\sharp[1/\ell]$ where $\theta(E_a)$ is invertible, $f=a$ and $\theta([u])=\theta([\varepsilon_q])^a$;
on its complement $\theta(E_a)=0$. So $\theta(E_a)(\theta([u])-\theta([\varepsilon_q])^a)$ vanishes in
$A^\sharp[1/\ell]$, hence in $A^\sharp$, where $\ell$ is a nonzerodivisor (Lemma~\ref{gal-crystalline-test-facts}(1)).
As $\sum_aE_a=1$ (proof of Lemma~\ref{gal-realization-locally-constant}), $\theta([u^q])=\sum_a\theta(E_a)\theta([u])^q=
\sum_a\theta(E_a)\theta([\varepsilon_q])^{aq}=1$.

Fix $a$, represented by $0\leq a<q$, and put $m=(q-1)a$ and $w=E_a[u][\varepsilon_q]^m+1-E_a$. Then $\theta(w)=
\theta(E_a)\theta([\varepsilon_q])^{a+m}+1-\theta(E_a)=1$, since $q$ divides $a+m$. By
Lemma~\ref{gal-realization-ingredients-logarithm}(2), $\log_n(w^q)=q\log_n(w)=0$ in $A_{\mathrm{cris}}(A)/\ell^n$. As
$E_a$ is idempotent, $E_aw^q=E_a[u]^q[\varepsilon_q]^{mq}=E_a[u^q][\varepsilon_A]^m$, so by
Lemma~\ref{gal-realization-ingredients-logarithm}(3), (2) and (4), $0=E_a\log_n(w^q)=E_a(\log_n([u^q])+mt)$. Since
$a+m=qa$ and $q=0$ in $A_{\mathrm{cris}}(A)/\ell^n$, $E_a\log_n([u^q])=aE_at$. Summing over $a$ gives
$\log_n([u^q])=\sum_aa[e_a]\,t=t\,\lambda_A(f)$.
\end{proof}

\subsection{Approximation systems and truncated periods}\label{subsection-gal-truncated-periods}

\begin{definition}\label{gal-realization-kappa-approximations}
Let $\ell$, $C$, $A_{\mathrm{cris}}$, $K$, $\mathcal O_K$, $G_K$, $t$ and $B_{\mathrm{dR}}^+\subset B_{\mathrm{dR}}$ be as in
Definition~\ref{gal-realization-period-datum}, so that $K\subset B_{\mathrm{dR}}^+$ and $A_{\mathrm{cris}}\subset
B_{\mathrm{dR}}^+$ \cite{Fon94a}, and let $\xi$ and $\theta:W(\mathcal O_C^\flat)\to\mathcal O_C$ be as in
Definition~\ref{gal-crystalline-test-algebras}. For an ideal $I$ of a ring with divided powers $\gamma$ and $r\geq0$, let
$I^{[r]}$ be the ideal generated by the products $\gamma_{k_1}(x_1)\cdots\gamma_{k_n}(x_n)$ with $x_i\in I$ and
$k_1+\cdots+k_n\geq r$. A \emph{divided power map} is a ring map that carries an ideal with divided powers into another
and commutes with the divided powers.
\begin{enumerate}
\item Let $m\geq1$. An \emph{approximation of height $m$} consists of a $K$-algebra $B$ with a ring map
  $\tau:B_{\mathrm{dR}}^+\to B$ and an ideal $J$ with divided powers $\delta$, an $\mathcal O_K$-algebra $A$ with an
  ideal $I$ with divided powers $\gamma$, ring maps $\iota:A\to B$ and $\lambda:A_{\mathrm{cris}}\to A$, and for each
  $\sigma\in G_K$ ring endomorphisms $\sigma_A$ of $A$ and $\sigma_B$ of $B$, with the following properties.
  \begin{enumerate}
  \item[(a)] $\tau$ is surjective with kernel $t^mB_{\mathrm{dR}}^+$, $\tau(a)=a$ for $a\in K$, and $J^{[r]}=\tau(t)^rB$
    for every $r\geq0$.
  \item[(b)] $\iota$ is injective and a divided power map, $I=\iota^{-1}(J)$, $J=\iota(I)B$, $\iota(a)=a$ for
    $a\in\mathcal O_K$, and $\iota(\lambda(c))=\tau(c)$ for $c\in A_{\mathrm{cris}}$.
  \item[(c)] $\sigma_B\circ\tau=\tau\circ\sigma$, $\iota\circ\sigma_A=\sigma_B\circ\iota$ and
    $\sigma_A\circ\lambda=\lambda\circ\sigma$; $\sigma_A$ fixes $\mathcal O_K$, $\sigma_B$ fixes $K$, and both are
    divided power maps.
  \item[(d)] $I^{[m]}=0$; $\ell x\in I^{[r]}$ implies $x\in I^{[r]}$; for every integer $n\geq1$ some power of $\ell$
    lies in $nA$; each $I^{[r]}$ is finitely generated; $I$ is prime; for all $x,y\in A$, $x\in yA+I$ or
    $y\in xA+I$; for $x\notin I$ some power of $\ell$ lies in $xA+I$; $A$ is $\ell$-adically complete and
    separated; and every $x\in A$ that lies in $\ell^kA+I^{[r]}$ for all $k$ lies in $I^{[r]}$.
  \item[(e)] $\lambda(\xi)\in I$, and $a-\lambda(w)\in I$ for $a\in\mathcal O_K$ and $w\in W(\mathcal O_C^\flat)$ with
    $\theta(w)=a$.
  \end{enumerate}
\item An \emph{approximation system} consists of approximations $(B_m,\tau_m,J_m,\delta_m,A_m,I_m,\gamma_m,\iota_m,
  \lambda_m,\sigma_{A_m},\sigma_{B_m})$ of height $m+1$ for $m\geq0$, and ring maps $d_A:A_{m+1}\to A_m$ and
  $d_B:B_{m+1}\to B_m$, such that $d_B\circ\tau_{m+1}=\tau_m$, $d_B\circ\iota_{m+1}=\iota_m\circ d_A$ and
  $d_A\circ\lambda_{m+1}=\lambda_m$; $d_A$ fixes $\mathcal O_K$ and $d_B$ fixes $K$; $d_A$ and $d_B$ are divided
  power maps that commute with the $\sigma_A$ and the $\sigma_B$; and for every family $(y_m)_{m\geq0}$ with
  $y_m\in B_m$ and $d_B(y_{m+1})=y_m$ there is exactly one $z\in B_{\mathrm{dR}}^+$ with $\tau_m(z)=y_m$ for all $m$.
\end{enumerate}
\end{definition}
\leanref{Development/WeilBound/CoreStatements/PadicHodgeComparison/EtaleCohomologyComparison/PerfectoidPullbacks/LevelChainValuations/ToBSurj.lean}{CrystallineSemistableComparison.Pb.Kap.Lev}
\leanref{Development/WeilBound/CoreStatements/PadicHodgeComparison/EtaleCohomologyComparison/PerfectoidPullbacks/LevelChainValuations/ToBSurj.lean}{CrystallineSemistableComparison.Pb.Kap.Lev.Good}
\leanref{Development/WeilBound/CoreStatements/PadicHodgeComparison/EtaleCohomologyComparison/PerfectoidPullbacks/LevelChainValuations/ChainOf.lean}{CrystallineSemistableComparison.Pb.Kap.Lev.CrisOK}
\leanref{Development/WeilBound/CoreStatements/PadicHodgeComparison/EtaleCohomologyComparison/PerfectoidPullbacks/LevelChainValuations/ToBSurj.lean}{CrystallineSemistableComparison.Pb.Kap.LevSys}
\leanref{Development/WeilBound/CoreStatements/PadicHodgeComparison/EtaleCohomologyComparison/PerfectoidPullbacks/LevelChainValuations/ToBSurj.lean}{CrystallineSemistableComparison.Pb.Kap.LevSys.Good}
\leanref{Development/WeilBound/CoreStatements/PadicHodgeComparison/EtaleCohomologyComparison/PerfectoidPullbacks/SplittingFieldLattices/GrSplit.lean}{CrystallineSemistableComparison.Pb.LatBase}
\leanref{Development/WeilBound/CoreStatements/PadicHodgeComparison/EtaleCohomologyComparison/PerfectoidPullbacks/InclStd.lean}{CrystallineSemistableComparison.Pb.TFBase}
\leanref{Development/WeilBound/CoreStatements/PadicHodgeComparison/EtaleCohomologyComparison/PerfectoidPullbacks/SEqOne.lean}{CrystallineSemistableComparison.Pb.pdP}

\begin{definition}\label{gal-realization-kappa-lattice}
Let $W$ be a ring, $B$ a $\mathbf Q$-algebra, $f:W\to B$ a ring map, $\xi\in W$, and $\ell\geq1$ and $M\geq0$ integers.
Put $\eta=\ell^{-M}f(\xi)$ and $e_i=\eta^i/i!$ for $i\geq0$. The \emph{lattice ring} $A(f,\xi,M)$ is the subring of $B$
generated by $f(W)$ and the $e_i$. The ideal $f(\xi)B$ has the divided powers $\delta_k(x)=x^k/k!$; put
$I(f,\xi,M)=A(f,\xi,M)\cap f(\xi)B$. For an integer $m\geq1$, the triple $(f,\xi,\ell)$ has \emph{height} $m$ if:
\begin{enumerate}
\item $f(\xi)^m=0$;
\item if $0\leq i<m$, $a\in W$, $b\in B$ and $f(a)f(\xi)^i=f(\xi)^{i+1}b$, then $\xi$ divides $a$;
\item $W/\xi W$ is a domain, and of any two elements of $W/\xi W$ one divides the other;
\item every nonzero element of $W/\xi W$ divides a power of $\ell$, the image of $\ell$ in $W/\xi W$ is not $0$, and
  $W/\xi W$ is $\ell$-adically complete and separated;
\item for every integer $n\geq1$ there are $k\geq0$ and $u\in W$ with $nu=\ell^k$.
\end{enumerate}
\end{definition}
\leanref{Development/WeilBound/CoreStatements/PadicHodgeComparison/EtaleCohomologyComparison/PerfectoidPullbacks/LevelChainValuations/LatDefs.lean}{CrystallineSemistableComparison.Pb.Lat.dpE}
\leanref{Development/WeilBound/CoreStatements/PadicHodgeComparison/EtaleCohomologyComparison/PerfectoidPullbacks/LevelChainValuations/LatDefs.lean}{CrystallineSemistableComparison.Pb.Lat.sub}
\leanref{Development/WeilBound/CoreStatements/PadicHodgeComparison/EtaleCohomologyComparison/PerfectoidPullbacks/LevelChainValuations/LatDefs.lean}{CrystallineSemistableComparison.Pb.Lat.I}
\leanref{Development/WeilBound/CoreStatements/PadicHodgeComparison/EtaleCohomologyComparison/PerfectoidPullbacks/LevelChainValuations/LatDefs.lean}{CrystallineSemistableComparison.Pb.Lat.Hyp}

\begin{lemma}\label{gal-realization-kappa-lattice-ring}
Let $W$, $B$, $f$, $\xi$, $\ell$, $M$ and the $e_i$ be as in Definition~\ref{gal-realization-kappa-lattice}, with
$A=A(f,\xi,M)$ and $I=I(f,\xi,M)$, and suppose that $(f,\xi,\ell)$ has height $m$.
\begin{enumerate}
\item $e_0=1$, $e_ie_j=\binom{i+j}{i}e_{i+j}$, $f(\xi)e_i=\ell^M(i+1)e_{i+1}$, and $e_i=0$ for $i\geq m$. The elements of
  $A$ are the sums $\sum_{i<m}f(c_i)e_i$ with $c_i\in W$, and for $r\geq0$ the elements of $A\cap f(\xi)^rB$ are the sums
  $\sum_{r\leq i<m}f(c_i)e_i$.
\item If $x\in I$ and $k\geq0$, then $x^k/k!\in A$. So $\gamma_k(x)=x^k/k!$ are divided powers on $I$, and the
  inclusion $A\to B$ is a divided power map.
\item For $r\geq0$, $I^{[r]}=A\cap f(\xi)^rB$ (with $I^{[r]}$ as in Definition~\ref{gal-realization-kappa-approximations}),
  and this ideal is generated by $e_r,\ldots,e_{m-1}$. In particular $I^{[m]}=0$.
\item If $x\in A$ and $\ell x\in I^{[r]}$, then $x\in I^{[r]}$. For every integer $n\geq1$ some power of $\ell$ lies in
  $nA$.
\item For every $x\in A$ there is $a\in W$ with $x-f(a)\in I$, and $f(a)\in I$ exactly when $\xi$ divides $a$. Hence $I$
  is prime; for all $x,y\in A$, $x\in yA+I$ or $y\in xA+I$; and for $x\notin I$ some power of $\ell$ lies in $xA+I$.
\item $A$ is $\ell$-adically complete and separated, and every $x\in A$ that lies in $\ell^kA+I^{[r]}$ for every $k$
  lies in $I^{[r]}$.
\end{enumerate}
\end{lemma}
\leanref{Development/AutomorphicSpectrum/Proofs/BlockDecomposition/MemLatticeSubIffSumStatement.lean}{GlobalTheory.T.mem_lattice_sub_iff_sum}
\leanref{Development/AutomorphicSpectrum/Proofs/BlockDecomposition/MemLatticeSubIffSum.lean}{GlobalTheory.CompanionTheory.LatSpan.dpE_mul}
\leanref{Development/AutomorphicSpectrum/Proofs/BlockDecomposition/LatticeXiMulDpowSuccStatement.lean}{GlobalTheory.T.lattice_xi_mul_dpow_succ}
\leanref{Development/AutomorphicSpectrum/Proofs/BlockDecomposition/MemLatticeIdealIffSumStatement.lean}{GlobalTheory.T.mem_lattice_ideal_iff_sum}
\leanref{Development/AutomorphicSpectrum/Proofs/BlockDecomposition/LatticeFiltrationComapSpanPow.lean}{GlobalTheory.CompanionTheory.LatPdpEq.exists_rep_high}
\leanref{Development/AutomorphicSpectrum/Proofs/BlockDecomposition/LatticeDpowClosedOfHyp.lean}{GlobalTheory.CompanionTheory.LatDP.dpClosed}
\leanref{Development/AutomorphicSpectrum/Proofs/BlockDecomposition/LatticeDpowClosedOfHyp.lean}{GlobalTheory.CompanionTheory.LatDP.inv_fact_smul_dpE_pow}
\leanref{Development/WeilBound/CoreStatements/PadicHodgeComparison/EtaleCohomologyComparison/PerfectoidPullbacks/LevelChainValuations/LatDefs.lean}{CrystallineSemistableComparison.Pb.Lat.γ}
\leanref{Development/AutomorphicSpectrum/Proofs/BlockDecomposition/DividedPowerFiltrationSpanPowStatement.lean}{GlobalTheory.T.divided_power_filtration_span_pow}
\leanref{Development/AutomorphicSpectrum/Proofs/BlockDecomposition/LatticeFiltrationComapSpanPowStatement.lean}{GlobalTheory.T.lattice_filtration_comap_span_pow}
\leanref{Development/AutomorphicSpectrum/Proofs/BlockDecomposition/LatticeFiltrationFg.lean}{GlobalTheory.CompanionTheory.LatFG.dividedPowerFiltration_fg}
\leanref{Development/AutomorphicSpectrum/Proofs/BlockDecomposition/LatticeFiltrationTopEqBot.lean}{GlobalTheory.CompanionTheory.LatNil.dividedPowerFiltration_eq_bot}
\leanref{Development/AutomorphicSpectrum/Proofs/BlockDecomposition/LatticeFiltrationLTorsionFreeStatement.lean}{GlobalTheory.T.lattice_filtration_l_torsion_free}
\leanref{Development/AutomorphicSpectrum/Proofs/BlockDecomposition/LatticeNatDividesLPow.lean}{GlobalTheory.CompanionTheory.LatticeUnit.natCast_dvd_pow}
\leanref{Development/AutomorphicSpectrum/Proofs/BlockDecomposition/LatticeIdealIsPrime.lean}{GlobalTheory.CompanionTheory.LatPrime.I_isPrime}
\leanref{Development/AutomorphicSpectrum/Proofs/BlockDecomposition/LatticeIdealQuotientChain.lean}{GlobalTheory.CompanionTheory.LatVal.val}
\leanref{Development/AutomorphicSpectrum/Proofs/BlockDecomposition/LatticeLPowMemSpanOfNotMem.lean}{GlobalTheory.CompanionTheory.Lat.pow_mem_span_sup_I}
\leanref{Development/AutomorphicSpectrum/Proofs/BlockDecomposition/LatticeIsAdicComplete.lean}{GlobalTheory.CompanionTheory.LatCpl.isAdicComplete_A}
\leanref{Development/AutomorphicSpectrum/Proofs/BlockDecomposition/LatticeFiltrationLAdicClosed.lean}{GlobalTheory.CompanionTheory.LatClosed.closed}

\begin{proof}
(1) The first identities follow from $e_i=\eta^i/i!$ and $f(\xi)=\ell^M\eta$, and $e_i=0$ for $i\geq m$ because $e_i$ is
a rational multiple of $f(\xi)^i$. By the product formula and $e_0=1$, the sums $\sum_if(c_i)e_i$ form a subring of $B$
that contains $f(W)$ and the $e_i$ and lies in $A$; so it is $A$. Let $x\in A\cap f(\xi)^rB$. By induction on $j\leq
\min(r,m)$, $x=\sum_{j\leq i<m}f(c_i)e_i$ for some $c_i$. For $j=0$ this was just shown. Suppose it holds for some
$j<\min(r,m)$. The terms with $i>j$ lie in $f(\xi)^{j+1}B$, and so does $x$; hence $f(c_j)e_j\in f(\xi)^{j+1}B$. As
$e_j$ is a nonzero rational multiple of $f(\xi)^j$, condition (2) of Definition~\ref{gal-realization-kappa-lattice}
gives $c_j=\xi d$ with $d\in W$, and then $f(c_j)e_j=f(\ell^M(j+1)d)e_{j+1}$, which can be added to the term of index
$j+1$. The converse inclusion holds as $e_i\in f(\xi)^iB$.

(2) Let $x\in I$. By (1) for $r=1$, $x=\sum_{1\leq j<m}f(c_j)e_j$. By the binomial theorem, $(z+w)^k/k!=\sum_{a+b=k}
(z^a/a!)(w^b/b!)$, so it suffices that $(f(c)e_j)^k/k!\in A$ for $c\in W$ and $j\geq1$. This element is
$f(c)^kN_{k,j}e_{kj}$, where $N_{k,j}=(kj)!/(k!\,(j!)^k)$ is an integer, the number of partitions of a set of $kj$
elements into $k$ blocks of $j$ elements. For $k\geq1$, $x^k/k!$ also lies in $f(\xi)B$. So the divided powers $\delta$ of
$f(\xi)B$ restrict to divided powers $\gamma$ on $I$.

(3) A divided power map carries $I^{[r]}$ into $(f(\xi)B)^{[r]}$, which is $f(\xi)^rB$: a product
$\delta_{k_1}(y_1)\cdots\delta_{k_s}(y_s)$ with $y_i\in f(\xi)B$ lies in $f(\xi)^{k_1+\cdots+k_s}B$, and
$f(\xi)^r=\delta_1(f(\xi))^r$. So $I^{[r]}\subset A\cap f(\xi)^rB$. Conversely, $e_1\in I$ and $\gamma_i(e_1)=e_i$, so
$e_i\in I^{[i]}\subset I^{[r]}$ for $i\geq r$, and (1) gives the reverse inclusion and the generators. As $f(\xi)^m=0$,
$I^{[m]}=0$.

(4) As $\ell$ is a unit of $B$, $\ell x\in f(\xi)^rB$ implies $x\in f(\xi)^rB$; now use (3). For $n\geq1$, condition (5)
gives $\ell^k=nu$, so $\ell^k=nf(u)\in nA$.

(5) The elements $x$ of $B$ with $x-f(a)\in f(\xi)B$ for some $a\in W$ form a subring that contains $f(W)$ and the $e_i$,
as $e_0=1$ and $e_i\in f(\xi)B$ for $i\geq1$; so it contains $A$. If $\xi$ divides $a$, then $f(a)\in A\cap f(\xi)B=I$.
If $f(a)=f(\xi)b$, condition (2) for $i=0$ gives $\xi\mid a$. So $W\to A/I$ is surjective with kernel $\xi W$, and
$A/I\cong W/\xi W$. The three claims follow from conditions (3) and (4): for instance, if $x\notin I$ and $x-f(a)\in
I$, then $a\notin\xi W$, so $au-\ell^n\in\xi W$ for some $u$ and $n$, and $\ell^n\in xA+I$.

(6) Call an abelian group $P$ \emph{complete} if $\bigcap_n\ell^nP=0$ and every sequence $(s_n)$ in $P$ with
$s_{n+1}-s_n\in\ell^nP$ for all $n$ has an $x$ with $x-s_n\in\ell^nP$ for all $n$; for a ring this means
$\ell$-adically complete and separated. If $N\subset P$ is a subgroup with $N$ and $P/N$ complete and $P/N$ without
$\ell$-torsion, then $P$ is complete. Indeed, if $x\in\ell^nP$ for all $n$, its image in $P/N$ vanishes, so $x\in N$;
writing $x=\ell^ny_n$, the image of $y_n$ is killed by $\ell^n$, so $y_n\in N$ and $x\in\bigcap_n\ell^nN=0$. Given
$(s_n)$ with $s_{n+1}-s_n=\ell^ny_n$, let $\bar p$ be a limit of the images in $P/N$, with $\bar p-\bar s_n=\ell^n\bar
z_n$, and lift $\bar p$ and $\bar z_n$ to $p$ and $w_n$. Then $t_n=s_n+\ell^nw_n-p\in N$, and $t_{n+1}-t_n=\ell^n(y_n+
\ell w_{n+1}-w_n)$ lies in $N\cap\ell^nP=\ell^nN$, as $P/N$ has no $\ell$-torsion. A limit $T$ of $(t_n)$ in $N$ gives
the limit $T+p$ of $(s_n)$.

Put $F_i=A\cap f(\xi)^iB$. Then $F_m=0$, and for $i<m$ the map $c\mapsto f(c)e_i$ induces $W/\xi W\cong F_i/F_{i+1}$:
it is onto by (1), and its kernel is $\xi W$ by condition (2) and $f(\xi d)e_i=f(\ell^M(i+1)d)e_{i+1}$. By
conditions (3) and (4), $W/\xi W$ is complete without $\ell$-torsion. Descending induction on $i$ shows that every
$F_i$ is complete, and $F_0=A$.

For the last claim, induct on $r$; for $r=0$, $I^{[0]}=A$. Let $x\in\ell^kA+I^{[r+1]}$ for every $k$. Then $x\in
I^{[r]}$ by induction. If $r\geq m$, $I^{[r]}=0$. Otherwise, by (1) and (3), $x=f(c)e_r+t$ with $t\in I^{[r+1]}$. For
each $n$ write $x=\ell^na+y$ with $y\in I^{[r+1]}$. Then $\ell^na\in I^{[r]}$, so $a\in I^{[r]}$ by (4), and
$a=f(d)e_r+t'$ with $t'\in I^{[r+1]}$. So $f(c-\ell^nd)e_r=\ell^nt'+y-t\in f(\xi)^{r+1}B$, and condition (2) gives
$c-\ell^nd\in\xi W$. The image of $c$ in $W/\xi W$ lies in every $\ell^n(W/\xi W)$, so it is $0$: $c=\xi e$, and
$f(c)e_r=f(\ell^M(r+1)e)e_{r+1}\in I^{[r+1]}$. Hence $x\in I^{[r+1]}$.
\end{proof}
\leanref{Development/AutomorphicSpectrum/Proofs/BlockDecomposition/LatticeIsAdicComplete.lean}{GlobalTheory.CompanionTheory.LatCpl.Cpl.ext}
\leanref{Development/AutomorphicSpectrum/Proofs/BlockDecomposition/LatticeIsAdicComplete.lean}{GlobalTheory.CompanionTheory.LatCpl.cpl_step}

\begin{lemma}\label{gal-realization-kappa-lattice-maps}
Let $(W,B,f,\xi)$ and $(W',B',f',\xi')$ be as in Definition~\ref{gal-realization-kappa-lattice}, for the same $\ell$ and
$M$, with lattice rings $A$ and $A'$ and ideals $I$ and $I'$. Let $g:B\to B'$ be a ring map.
\begin{enumerate}
\item If $s:W\to W'$ is a ring map with $g\circ f=f'\circ s$, and $s(\xi)=u\xi'$ for some $u\in W'$, then
  $g(A)\subset A'$.
\item Suppose that $x^k/k!\in A$ for $x\in I$ and $k\geq0$, and likewise for $A'$ and $I'$, and give $I$ and $I'$ the
  divided powers $x\mapsto x^k/k!$. If $g(A)\subset A'$ and $g(f(\xi))=f'(u\xi')$ for some $u\in W'$, then the
  restriction $A\to A'$ of $g$ is a divided power map.
\end{enumerate}
\end{lemma}
\leanref{Development/AutomorphicSpectrum/Proofs/BlockDecomposition/GaloisPreservesLatticeSub.lean}{GlobalTheory.CompanionTheory.LatGal.map_mem_sub_of_mul}
\leanref{Development/AutomorphicSpectrum/Proofs/BlockDecomposition/LatticeSubReductionMem.lean}{GlobalTheory.CompanionTheory.LatLevel.dk_mem_sub}
\leanref{Development/AutomorphicSpectrum/Proofs/BlockDecomposition/LatticeIsDPMorphismOfCompat.lean}{GlobalTheory.CompanionTheory.LatDP.isDPMorphism_map}
\leanref{Development/AutomorphicSpectrum/Proofs/BlockDecomposition/LevelSystemTransitionIsDPMorphism.lean}{GlobalTheory.CompanionTheory.LatC.map_isDPMorphism}

\begin{proof}
(1) As $g$ is $\mathbf Q$-linear, $g(\eta)=\ell^{-M}f'(u\xi')=f'(u)\eta'$, so $g(e_i)=f'(u)^ie'_i\in A'$; and $g(f(w))=
f'(s(w))\in A'$. These elements generate $A$. (2) As $g(f(\xi)B)\subset f'(\xi')B'$, $g(I)\subset I'$, and
$g(x^k/k!)=g(x)^k/k!$.
\end{proof}

\begin{lemma}\label{gal-realization-kappa-newton}
Let $R$ be a ring and $E\in R[X]$, with derivative $E'$.
\begin{enumerate}
\item Let $a,T,y,v,e,w\in R$ and $m\geq0$ with $y^m=0$, $E'(a)v=T(1+Tyw)$ and $E(a)=T^2ye$. Then $E(a+Tyz)=0$ for some
  $z\in R$.
\item Let $a\in R$ with $E'(a)$ a unit, and let $\alpha,\beta\in R$ be roots of $E$ with $\alpha-a$ and $\beta-a$
  nilpotent. Then $\alpha=\beta$.
\end{enumerate}
\end{lemma}
\leanref{Development/AutomorphicSpectrum/Proofs/BlockDecomposition/IntegerMemLatticeSubBound.lean}{GlobalTheory.CompanionTheory.LatOKElem.exists_root}
\leanref{Development/AutomorphicSpectrum/Proofs/BlockDecomposition/IntegerMemLatticeSubBound.lean}{GlobalTheory.CompanionTheory.LatOKElem.newton_step}
\leanref{Development/AutomorphicSpectrum/Proofs/BlockDecomposition/IntegerMemLatticeSubBound.lean}{GlobalTheory.CompanionTheory.LatOKElem.root_unique}

\begin{proof}
For $b,h\in R$ there are $k,k'\in R$ with $E(b+h)=E(b)+E'(b)h+kh^2$ and $E'(b+h)=E'(b)+k'h$.

(1) As $Tyw$ is nilpotent, $1+Tyw$ is a unit; replacing $v$ by $v(1+Tyw)^{-1}$, assume $E'(a)v=T$. By induction on $n$
there are $z_n,r_n$ with $E(a+Tyz_n)=T^2y^{n+1}r_n$: take $z_0=0$ and $r_0=e$, and given $z_n$ and $r_n$ put
$h=-Tvy^{n+1}r_n$ and $z_{n+1}=z_n-vy^nr_n$, so that $a+Tyz_{n+1}=a+Tyz_n+h$. With $E'(a+Tyz_n)=E'(a)+k'Tyz_n$, the terms
$T^2y^{n+1}r_n$ and $E'(a)h=-T^2y^{n+1}r_n$ cancel in $E(a+Tyz_n+h)$, and the remaining terms are
$T^2y^{n+2}(-k'z_nvr_n+kv^2y^nr_n^2)$. For $n=m$, $y^{m+1}=0$.

(2) $0=E(\alpha)-E(\beta)=(\alpha-\beta)(E'(\beta)+k(\alpha-\beta))$ and $E'(\beta)=E'(a)+k'(\beta-a)$. As $\alpha-\beta$
and $\beta-a$ are nilpotent, $E'(a)+k'(\beta-a)+k(\alpha-\beta)$ is a unit, so $\alpha=\beta$.
\end{proof}

\begin{lemma}\label{gal-realization-kappa-fontaine}
Let $\ell$, $C$, $\mathcal O_C^\flat$, $\xi$, $\theta$ and $A_{\mathrm{cris}}$ be as in
Definition~\ref{gal-crystalline-test-algebras}, $K$, $\bar K$ and $G_K$ as in
Definition~\ref{gal-realization-split-diagrams}(3), and $t$ and $B_{\mathrm{dR}}^+$ as in
Definition~\ref{gal-realization-period-datum}. Put $A_{\mathrm{inf}}=W(\mathcal O_C^\flat)$, and let
$A_{\mathrm{cris}}(\mathcal O_C^\flat)\subset A_{\mathrm{inf}}[1/\ell]$ be the $A_{\mathrm{inf}}$-subalgebra generated by the
$\xi^n/n!$, whose $\ell$-adic completion is $A_{\mathrm{cris}}$.
\begin{enumerate}
\item $\theta:A_{\mathrm{inf}}\to\mathcal O_C$ is surjective with kernel $\xi A_{\mathrm{inf}}$, and commutes with $G_K$;
  so $\sigma(\xi)\in\xi A_{\mathrm{inf}}$ for $\sigma\in G_K$.
\item $B_{\mathrm{dR}}^+$ is a complete discrete valuation ring with residue field $C$, and the map $A_{\mathrm{inf}}\to
  B_{\mathrm{dR}}^+$ commutes with $G_K$. The maximal ideal of $B_{\mathrm{dR}}^+$ is the kernel of the map
  $\theta:B_{\mathrm{dR}}^+\to C$ that extends $\theta$, and it is generated by $t$ and also by the image of $\xi$.
\item $\bar K\subset B_{\mathrm{dR}}^+$, compatibly with the actions of $G_K$, and $\theta(x)=x$ for $x\in\bar K$. On
  $\mathbf Z_\ell=W(\mathbf F_\ell)$ the inclusion agrees with the map through $A_{\mathrm{inf}}$.
\item $A_{\mathrm{cris}}\subset B_{\mathrm{dR}}^+$, compatibly with $A_{\mathrm{inf}}\to B_{\mathrm{dR}}^+$ and with the
  actions of $G_K$, with $t$ mapping to $t$. For $c\in A_{\mathrm{cris}}$ and $m\geq0$ there is $s\in
  A_{\mathrm{cris}}(\mathcal O_C^\flat)$ with $c-s\in t^mB_{\mathrm{dR}}^+$.
\end{enumerate}
These are results of Fontaine \cite{Fon94a}.
\end{lemma}
\leanref{Development/WeilBound/CoreStatements/PadicHodgeComparison/EtaleCohomologyComparison/AcrisKummerUnits/MapUgAlgebraMap.lean}{CrystallineSemistableComparison.ker_theta0}
\leanref{GroundTheory/PAdicHodge/PerfectoidField.lean}{GroundTheory.PAdicHodge.IsPerfectoidField.surjective_fontaineTheta}
\leanref{GroundTheory/PAdicHodge/BdRGalois.lean}{GroundTheory.PAdicHodge.gal_smul_mem_ker_fontaineTheta}
\leanref{Development/AutomorphicSpectrum/Proofs/BlockDecomposition/TwistFrobeniusFiltrationCompat.lean}{GlobalTheory.CompanionTheory.TwFil.exists_gal_xi}
\leanref{GroundTheory/PAdicHodge/ThetaPlus.lean}{GroundTheory.PAdicHodge.IsPerfectoidField.maximalIdeal_bDeRhamPlus}
\leanref{GroundTheory/PAdicHodge/ThetaPlus.lean}{GroundTheory.PAdicHodge.IsPerfectoidField.maximalIdeal_bDeRhamPlus_eq_span}
\leanref{GroundTheory/PAdicHodge/ThetaPlus.lean}{GroundTheory.PAdicHodge.IsPerfectoidField.irreducible_wittToBDeRhamPlus}
\leanref{GroundTheory/PAdicHodge/UntiltFrobeniusEquivSymmEpsilonNeOne.lean}{GroundTheory.PAdicHodge.maximalIdeal_eq_span_tBdR}
\leanref{GroundTheory/PAdicHodge/UntiltFrobeniusEquivSymmEpsilonNeOne.lean}{GroundTheory.PAdicHodge.IsPerfectoidField.instIsAdicCompleteMaximalIdealBdR}
\leanref{GroundTheory/PAdicHodge/BdRGalois.lean}{GroundTheory.PAdicHodge.gal_smul_wittToBDeRhamPlus}
\leanref{GroundTheory/PAdicHodge/BdRAlgCl.lean}{GroundTheory.PAdicHodge.algClToBdR}
\leanref{GroundTheory/PAdicHodge/BdRAlgCl.lean}{GroundTheory.PAdicHodge.thetaAlgHom_algClToBdR}
\leanref{GroundTheory/PAdicHodge/BdRAlgCl.lean}{GroundTheory.PAdicHodge.gal_smul_algClToBdR}
\leanref{GroundTheory/PAdicHodge/BdRAlg.lean}{GroundTheory.PAdicHodge.IsPerfectoidField.algebraMap_padic_coe}
\leanref{GroundTheory/PAdicHodge/ACrisBdR.lean}{GroundTheory.PAdicHodge.ACris.iota}
\leanref{GroundTheory/PAdicHodge/ACrisBdR.lean}{GroundTheory.PAdicHodge.ACris.embedding₂}
\leanref{GroundTheory/PAdicHodge/ACrisBdR.lean}{GroundTheory.PAdicHodge.ACris.evalₐ_iota}
\leanref{GroundTheory/PAdicHodge/ACrisBdRGalois.lean}{GroundTheory.PAdicHodge.ACris.iota_gal_smul}
\leanref{GroundTheory/PAdicHodge/ACrisBdRGalois.lean}{GroundTheory.PAdicHodge.ACris.iota_periodT}

\begin{lemma}\label{gal-realization-kappa-truncations}
Let $\ell$, $\xi$, $K$, $\mathcal O_K$, $G_K$, $t$, $B_{\mathrm{dR}}^+$ and $A_{\mathrm{inf}}$ be as in
Lemma~\ref{gal-realization-kappa-fontaine}. For $m\geq1$ let $Q_m=B_{\mathrm{dR}}^+/t^mB_{\mathrm{dR}}^+$, a
$K$-algebra through $K\subset B_{\mathrm{dR}}^+$, and $f_m:A_{\mathrm{inf}}\to Q_m$ the induced map, and for $M\geq0$ let
$\Lambda_m(M)$ be the lattice ring of $(f_m,\xi,M)$ (Definition~\ref{gal-realization-kappa-lattice}).
\begin{enumerate}
\item $(f_m,\xi,\ell)$ has height $m$.
\item The ideal $t^mB_{\mathrm{dR}}^+$ is stable under $G_K$, and for the induced action on $Q_m$,
  $\sigma(\Lambda_m(M))\subset\Lambda_m(M)$ for $\sigma\in G_K$. The projection $Q_{m+1}\to Q_m$ carries
  $\Lambda_{m+1}(M)$ into $\Lambda_m(M)$, and the image of $A_{\mathrm{cris}}$ in $Q_m$ lies in $\Lambda_m(M)$.
\item There is $M\geq0$ such that the image of $\mathcal O_K$ in $Q_m$ lies in $\Lambda_m(M)$ for every $m\geq1$.
\end{enumerate}
\end{lemma}
\leanref{Development/AutomorphicSpectrum/Proofs/BlockDecomposition/IsAdicCompleteOfEquiv.lean}{GlobalTheory.CompanionTheory.LatCHyp.hyp_k}
\leanref{Development/AutomorphicSpectrum/Proofs/BlockDecomposition/GaloisPreservesTPowerSpanStatement.lean}{GlobalTheory.T.galois_preserves_t_power_span}
\leanref{Development/AutomorphicSpectrum/Proofs/BlockDecomposition/GaloisPreservesLatticeSubStatement.lean}{GlobalTheory.T.galois_preserves_lattice_sub}
\leanref{Development/AutomorphicSpectrum/Proofs/BlockDecomposition/LatticeSubReductionMemStatement.lean}{GlobalTheory.T.lattice_sub_reduction_mem}
\leanref{Development/AutomorphicSpectrum/Proofs/BlockDecomposition/AcrisMemLatticeSubStatement.lean}{GlobalTheory.T.acris_mem_lattice_sub}
\leanref{Development/AutomorphicSpectrum/Proofs/BlockDecomposition/ExistsLatticeBoundIntegersStatement.lean}{GlobalTheory.T.exists_lattice_bound_integers}
\leanref{Development/AutomorphicSpectrum/Proofs/BlockDecomposition/IntegerMemLatticeSubBoundStatement.lean}{GlobalTheory.T.integer_mem_lattice_sub_bound}

\begin{proof}
(1) By Lemma~\ref{gal-realization-kappa-fontaine}(2), $\xi$ and $t$ generate the same ideal of $B_{\mathrm{dR}}^+$, so
$f_m(\xi)^m=0$. Let $i<m$, $a\in A_{\mathrm{inf}}$ and $b\in B_{\mathrm{dR}}^+$ with $a\xi^i-\xi^{i+1}b\in
t^mB_{\mathrm{dR}}^+$, writing $a$ also for its image in $B_{\mathrm{dR}}^+$. As $\xi^{i+1}$ divides $t^m$, it divides
$a\xi^i$. As $B_{\mathrm{dR}}^+$ is a domain and $\xi\neq0$ there, $\xi$ divides $a$ in $B_{\mathrm{dR}}^+$, so
$\theta(a)=0$ and $a\in\xi A_{\mathrm{inf}}$ by Lemma~\ref{gal-realization-kappa-fontaine}(1). By the same statement $\theta$ identifies
$A_{\mathrm{inf}}/\xi$ with $\mathcal O_C$, which is a domain, in which of two elements the one of smaller absolute value
is divisible by the other, a nonzero $x$ divides $\ell^n$ once $|\ell|^n\leq|x|$, $\ell\neq0$, and which is
$\ell$-adically complete and separated. For condition (5), $n=\ell^kv$ with $v\in\mathbf Z_\ell^\times$, so
$nv^{-1}=\ell^k$, and $\mathbf Z_\ell\subset A_{\mathrm{inf}}$.

(2) Each $\sigma$ is an automorphism of the local ring $B_{\mathrm{dR}}^+$, so it preserves the maximal ideal $(t)$ and
its powers. Lemma~\ref{gal-realization-kappa-lattice-maps}(1) applies to $\sigma$ on $Q_m$ and on $A_{\mathrm{inf}}$,
with $\sigma(\xi)=u\xi$ (Lemma~\ref{gal-realization-kappa-fontaine}(1)), and to the projection $Q_{m+1}\to Q_m$ with the
identity of $A_{\mathrm{inf}}$ and $u=1$. For $c\in A_{\mathrm{cris}}$, Lemma~\ref{gal-realization-kappa-fontaine}(4)
gives $s\in A_{\mathrm{cris}}(\mathcal O_C^\flat)$ with the same image in $Q_m$. Now $s$ is a sum of terms $w\xi^n/n!$
with $w\in A_{\mathrm{inf}}$, and $f_m(\xi)^n/n!=\ell^{Mn}e_n$ lies in $\Lambda_m(M)$.

(3) As $\mathcal O_K$ is a finitely generated $\mathbf Z_\ell$-module and $\mathbf Z_\ell$ maps into $f_m(A_{\mathrm{inf}})$
(Lemma~\ref{gal-realization-kappa-fontaine}(3)), it suffices that for each $x\in\mathcal O_K$ there is $c\geq0$ with $x\in
\Lambda_m(M)$ for all $m\geq1$ and $M\geq2c$; then take $M=2\max c$ over a set of generators. Let $E_0$ be the minimal
polynomial of $x$ over $\mathbf Q_\ell$. As $|x|\leq1$, $E_0\in\mathbf Z_\ell[X]$, and $E_0'(x)\neq0$ as $E_0$ is
separable. Let $E$ be its image in $A_{\mathrm{inf}}[X]$. Choose $a\in A_{\mathrm{inf}}$ with $\theta(a)=x$. Then
$\theta(E(a))=E_0(x)=0$, so $E(a)=\xi e$ for some $e$. Choose $c$ with $|\ell|^c<|E_0'(x)|$ and $v\in A_{\mathrm{inf}}$
with $\theta(v)=\ell^c/E_0'(x)$. Then $\theta(E'(a)v-\ell^c)=0$, so $E'(a)v=\ell^c+\xi w$ for some $w$.

Let $M\geq2c$ and $m\geq1$, write $A=\Lambda_m(M)$ and drop $f_m$ from the notation. Put $T=\ell^c$ and
$y=\ell^{-2c}\xi=\ell^{M-2c}e_1\in A$. Then $y^m=0$, $E'(a)v=T(1+Tyw)$ and $E(a)=T^2ye$, so
Lemma~\ref{gal-realization-kappa-newton}(1) gives a root $b=a+Tyz\in A$ of $E$. The image of $x$ in $Q_m$ is also a root
of $E$, and $x-a$ lies in the nilpotent ideal $\xi Q_m$, by Lemma~\ref{gal-realization-kappa-fontaine}(3) and~(2), as
does $b-a$. The element $E'(a)v=\ell^c+\xi w$ is a unit of $Q_m$, hence so is $E'(a)$. By
Lemma~\ref{gal-realization-kappa-newton}(2), $x=b\in A$.
\end{proof}
\leanref{Development/AutomorphicSpectrum/Proofs/BlockDecomposition/IsAdicCompleteOfEquiv.lean}{GlobalTheory.CompanionTheory.LatCHyp.gr_k}
\leanref{Development/AutomorphicSpectrum/Proofs/BlockDecomposition/GaloisPreservesLatticeSub.lean}{GlobalTheory.CompanionTheory.LatGal.gal_mem_sub}
\leanref{Development/AutomorphicSpectrum/Proofs/BlockDecomposition/AcrisMemLatticeSub.lean}{GlobalTheory.CompanionTheory.LatCris.cris_in}
\leanref{Development/AutomorphicSpectrum/Proofs/BlockDecomposition/ExistsLatticeBoundIntegers.lean}{GlobalTheory.CompanionTheory.LatCOK.exists_uniform_mem_sub}
\leanref{Development/AutomorphicSpectrum/Proofs/BlockDecomposition/ExistsFinsetGenerators.lean}{GlobalTheory.CompanionTheory.LocalIntegers.exists_finset_generators}
\leanref{Development/AutomorphicSpectrum/Proofs/BlockDecomposition/PadicIntMemSub.lean}{GlobalTheory.CompanionTheory.LatCZp.padicInt_mem_sub}
\leanref{Development/AutomorphicSpectrum/Proofs/BlockDecomposition/IntegerMemLatticeSubBound.lean}{GlobalTheory.CompanionTheory.LatOKElem.exists_data}
\leanref{Development/AutomorphicSpectrum/Proofs/BlockDecomposition/IntegerMemLatticeSubBound.lean}{GlobalTheory.CompanionTheory.LatOKElem.mem_sub_of_root}

\begin{lemma}\label{gal-realization-kappa-approximation-system}
Let $\ell$, $K$, $G_K$ and $B_{\mathrm{dR}}^+$ be as in Definition~\ref{gal-realization-period-datum}. There is an
approximation system (Definition~\ref{gal-realization-kappa-approximations}(2)).
\end{lemma}
\leanref{Development/WeilBound/CoreStatements/PadicHodgeComparison/EtaleCohomologyComparison/HodgeTateResidues/LatLink.lean}{CrystallineSemistableComparison.HT.levSys_ex}
\leanref{Development/WeilBound/CoreStatements/PadicHodgeComparison/EtaleCohomologyComparison/PerfectoidPullbacks/LevelChainValuations/CrystallineDeRhamFirstChernExists.lean}{CrystallineSemistableComparison.Pb.levSys_ex'}

\begin{proof}
Let $A_{\mathrm{inf}}$, the $Q_m$, $f_m$ and $\Lambda_m(M)$ be as in Lemma~\ref{gal-realization-kappa-truncations}, and
fix $M$ as in its part (3). For $m\geq0$ let the approximation of height $m+1$ consist of:
\begin{itemize}
\item $B_m=Q_{m+1}$, with the projection $\tau_m:B_{\mathrm{dR}}^+\to Q_{m+1}$, and $J_m=f_{m+1}(\xi)Q_{m+1}$ with the
  divided powers $\delta_k(x)=x^k/k!$;
\item $A_m=\Lambda_{m+1}(M)$, an $\mathcal O_K$-algebra by Lemma~\ref{gal-realization-kappa-truncations}(3), and
  $I_m=A_m\cap J_m$ with the divided powers $\gamma_k(x)=x^k/k!$ of Lemma~\ref{gal-realization-kappa-lattice-ring}(2),
  which applies by Lemma~\ref{gal-realization-kappa-truncations}(1);
\item $\iota_m:A_m\to B_m$ the inclusion, and $\lambda_m:A_{\mathrm{cris}}\to A_m$ the map to $Q_{m+1}$, which lands
  in $A_m$ by Lemma~\ref{gal-realization-kappa-truncations}(2);
\item $\sigma_{B_m}$ the map of $Q_{m+1}$ induced by $\sigma\in G_K$, and $\sigma_{A_m}$ its restriction to $A_m$
  (Lemma~\ref{gal-realization-kappa-truncations}(2)).
\end{itemize}
Let $d_B:B_{m+1}\to B_m$ be the projection $Q_{m+2}\to Q_{m+1}$ and $d_A:A_{m+1}\to A_m$ its restriction
(Lemma~\ref{gal-realization-kappa-truncations}(2)). Fix $m$ and drop it from the notation in (a) to (e).

\emph{Condition (a).} $\tau$ is onto with kernel $t^{m+1}B_{\mathrm{dR}}^+$, and $\tau(a)=a$ for $a\in K$ by the
definition of the $K$-algebra $Q_{m+1}$. By Lemma~\ref{gal-realization-kappa-fontaine}(2), $J=\tau(t)B$, and as in the
proof of Lemma~\ref{gal-realization-kappa-lattice-ring}(3), $J^{[r]}=\tau(t)^rB$.

\emph{Condition (b).} $\iota$ is injective, and it is a divided power map with $I=\iota^{-1}(J)$ by construction. As
$f_{m+1}(\xi)\in I$, $J=\iota(I)B$. By definition, $\iota(a)=\tau(a)=a$ for $a\in\mathcal O_K$ and
$\iota(\lambda(c))=\tau(c)$ for $c\in A_{\mathrm{cris}}$.

\emph{Condition (c).} The identities $\sigma_B\circ\tau=\tau\circ\sigma$ and $\iota\circ\sigma_A=\sigma_B\circ\iota$
hold by definition, and $\sigma_A\circ\lambda=\lambda\circ\sigma$ as $A_{\mathrm{cris}}\to B_{\mathrm{dR}}^+$ commutes
with $G_K$ (Lemma~\ref{gal-realization-kappa-fontaine}(4)). As $G_K$ fixes $K\subset B_{\mathrm{dR}}^+$
(Lemma~\ref{gal-realization-kappa-fontaine}(3)), $\sigma_B$ fixes $K$ and $\sigma_A$ fixes $\mathcal O_K$. By
Lemma~\ref{gal-realization-kappa-fontaine}(1) and~(2), $\sigma(f_{m+1}(\xi))=f_{m+1}(u\xi)$ for some $u\in
A_{\mathrm{inf}}$, so $\sigma_A$ is a divided power map by Lemma~\ref{gal-realization-kappa-lattice-maps}(2). Then
$\sigma_B(J)=\sigma_B(\iota(I)B)\subset\iota(\sigma_A(I))B\subset J$, and $\sigma_B(x^k/k!)=\sigma_B(x)^k/k!$, so
$\sigma_B$ is a divided power map.

\emph{Condition (d).} This is Lemma~\ref{gal-realization-kappa-lattice-ring}(3) to~(6), for the height $m+1$.

\emph{Condition (e).} $\iota(\lambda(\xi))=f_{m+1}(\xi)\in J$, so $\lambda(\xi)\in I$. Let $a\in\mathcal O_K$ and $w\in
A_{\mathrm{inf}}$ with $\theta(w)=a$. By Lemma~\ref{gal-realization-kappa-fontaine}(3), $\theta(a-w)=0$ in $C$, so $a-w$
lies in the maximal ideal $\xi B_{\mathrm{dR}}^+$, and $\iota(a-\lambda(w))=\tau(a-w)\in J$.

\emph{The system.} The identities $d_B\circ\tau_{m+1}=\tau_m$, $d_B\circ\iota_{m+1}=\iota_m\circ d_A$ and
$d_A\circ\lambda_{m+1}=\lambda_m$ hold because all these maps are induced by maps to $B_{\mathrm{dR}}^+$ followed by
projections; for the same reason $d_A$ fixes $\mathcal O_K$, $d_B$ fixes $K$, and $d_A$ and $d_B$ commute with the
$\sigma_A$ and the $\sigma_B$. As $d_B(f_{m+2}(\xi))=f_{m+1}(\xi)$, $d_A$ is a divided power map by
Lemma~\ref{gal-realization-kappa-lattice-maps}(2), and $d_B$ is one by the argument given for $\sigma_B$. Finally let
$y_m\in B_m$ with $d_B(y_{m+1})=y_m$ for all $m$, and lift each $y_m$ to $z_m\in B_{\mathrm{dR}}^+$. Then
$z_{m+1}-z_m\in t^{m+1}B_{\mathrm{dR}}^+$. As $B_{\mathrm{dR}}^+$ is $t$-adically complete and separated
(Lemma~\ref{gal-realization-kappa-fontaine}(2)), the $z_m$ converge to some $z$ with $\tau_m(z)=y_m$ for all $m$, and $z$
is unique as $\bigcap_mt^mB_{\mathrm{dR}}^+=0$.
\end{proof}
\leanref{Development/WeilBound/CoreStatements/PadicHodgeComparison/EtaleCohomologyComparison/PerfectoidPullbacks/LevelChainValuations/LatCDefs.lean}{CrystallineSemistableComparison.Pb.Lat.LevHyp}
\leanref{Development/WeilBound/CoreStatements/PadicHodgeComparison/EtaleCohomologyComparison/PerfectoidPullbacks/LevelChainValuations/LatCDefs.lean}{CrystallineSemistableComparison.Pb.Lat.lev}
\leanref{Development/WeilBound/CoreStatements/PadicHodgeComparison/EtaleCohomologyComparison/PerfectoidPullbacks/LevelChainValuations/LatCDefs.lean}{CrystallineSemistableComparison.Pb.Lat.levSys}
\leanref{Development/WeilBound/CoreStatements/PadicHodgeComparison/EtaleCohomologyComparison/HodgeTateResidues/LatLink.lean}{CrystallineSemistableComparison.HT.levHyp_ex}
\leanref{Development/WeilBound/CoreStatements/PadicHodgeComparison/EtaleCohomologyComparison/HodgeTateResidues/LatLink.lean}{CrystallineSemistableComparison.HT.lev_good}
\leanref{Development/WeilBound/CoreStatements/PadicHodgeComparison/EtaleCohomologyComparison/HodgeTateResidues/LatLink.lean}{CrystallineSemistableComparison.HT.levSys_good}
\leanref{Development/WeilBound/CoreStatements/PadicHodgeComparison/EtaleCohomologyComparison/HodgeTateResidues/LatLink.lean}{CrystallineSemistableComparison.HT.dp_of_char}
\leanref{Development/AutomorphicSpectrum/Proofs/BlockDecomposition/FontaineLatticeDpowClosedStatement.lean}{GlobalTheory.T.fontaine_lattice_dpow_closed}
\leanref{Development/AutomorphicSpectrum/Proofs/BlockDecomposition/LevelLatticeBaseStatement.lean}{GlobalTheory.T.level_lattice_base}
\leanref{Development/AutomorphicSpectrum/Proofs/BlockDecomposition/LevelFiltrationSpanTPowStatement.lean}{GlobalTheory.T.level_filtration_span_t_pow}
\leanref{Development/AutomorphicSpectrum/Proofs/BlockDecomposition/LevelGaloisIsDPMorphismStatement.lean}{GlobalTheory.T.level_galois_isDPMorphism}
\leanref{Development/AutomorphicSpectrum/Proofs/BlockDecomposition/LevelLamGaloisEquivariantStatement.lean}{GlobalTheory.T.level_lam_galois_equivariant}
\leanref{Development/AutomorphicSpectrum/Proofs/BlockDecomposition/LevelSystemTransitionIsDPMorphismStatement.lean}{GlobalTheory.T.level_system_transition_isDPMorphism}
\leanref{Development/AutomorphicSpectrum/Proofs/BlockDecomposition/ExistsUniquePeriodOfCompatibleStatement.lean}{GlobalTheory.T.existsUnique_period_of_compatible}
\leanref{Development/AutomorphicSpectrum/Proofs/BlockDecomposition/LevelXiMemIdealAndThetaCompatStatement.lean}{GlobalTheory.T.level_xi_mem_ideal_and_theta_compat}

\begin{definition}\label{gal-realization-kappa-truncated-periods}
Let $\ell$, $A_{\mathrm{cris}}$, $K$, $\mathcal O_K$ and $G_K$ be as in Definition~\ref{gal-realization-period-datum},
and let the approximations $(B_m,\tau_m,\ldots)$ of height $m+1$, with $d_A$ and $d_B$, form an approximation system
(Definition~\ref{gal-realization-kappa-approximations}). For a smooth projective $\mathcal O_K$-scheme $Y$ with generic
fibre $Y_K$, let $H^j_\infty(Y)$ be as in Definition~\ref{gal-realization-period-thickenings}(5), for $\mathcal
O=\mathcal O_K$ and its inclusion into $\mathcal O_C$, and $H^j_{\mathrm{dR}}(Y_K)$ as in
Definition~\ref{gal-realization-period-de-rham}. A \emph{system of truncated periods} for the approximation system
consists of relations $\rho_m(x,f)$, one for every smooth projective $\mathcal O_K$-scheme $Y$, tautological line datum
$\mathcal L$ on $Y$ (Definition~\ref{gal-realization-period-chern}(2)) and $j,m\geq0$, between elements $x\in
H^j_\infty(Y)$ and additive maps $f:\operatorname{Hom}_K(H^j_{\mathrm{dR}}(Y_K),K)\to B_m$, with the following properties
for every $Y$ and $\mathcal L$.
\begin{enumerate}
\item For all $j$, $m$ and $x$ there is exactly one $f$ with $\rho_m(x,f)$.
\item If $\rho_m(x,f)$, then $\rho_m(cx,\iota_m(\lambda_m(c))f)$ for $c\in A_{\mathrm{cris}}$ and $f(a\psi)=af(\psi)$
  for $a\in K$; if moreover $\rho_m(x',f')$, then $\rho_m(x+x',f+f')$.
\item If $\rho_{m+1}(x,f)$, then $\rho_m(x,d_B\circ f)$.
\item If $\rho_m(x,f)$ and $\sigma\in G_K$, then $\rho_m(\sigma x,\sigma_{B_m}\circ f)$.
\item Let $x\in H^i_\infty(Y)$ and $x'\in H^j_\infty(Y)$, with $\rho_m(x,f)$, $\rho_m(x',f')$ and $\rho_m(x\cup x',h)$.
  For bases $(b_s)$ of $H^i_{\mathrm{dR}}(Y_K)$ and $(b'_{s'})$ of $H^j_{\mathrm{dR}}(Y_K)$ with dual bases $(b_s^*)$
  and $(b'^*_{s'})$, and $\psi\in\operatorname{Hom}_K(H^{i+j}_{\mathrm{dR}}(Y_K),K)$,
  \[
   h(\psi)=\sum_{s,s'}f(b_s^*)\,f'(b'^*_{s'})\,\psi(b_s\cup b'_{s'}).
  \]
\item If $\rho_m(1,f)$ for the unit $1\in H^0_\infty(Y)$, then $f(\psi)=\psi(1)$ for every $\psi$.
\item If the components of $x\in H^2_\infty(Y)$ form a first Chern family of $\mathcal L$, $\omega\in
  H^2_{\mathrm{dR}}(Y_K)$ is a first Chern class of the base change of $\mathcal L$ to $Y_K$
  (Definition~\ref{gal-realization-period-chern}(4) and~(6)), and $\rho_m(x,f)$, then $f(\psi)=\psi(\omega)$ for every
  $\psi$.
\item Let $Y'$ be a smooth projective $\mathcal O_K$-scheme with a tautological line datum $\mathcal L'$, $h:Y\to Y'$ an $\mathcal
  O_K$-morphism with generic fibre $h_K:Y_K\to Y'_K$, and $h^*:H^j_\infty(Y')\to H^j_\infty(Y)$ and $h_K^*:H^j_{\mathrm
  {dR}}(Y'_K)\to H^j_{\mathrm{dR}}(Y_K)$ the pullbacks (Definitions~\ref{gal-realization-period-thickenings}(3)(e)
  and~\ref{gal-realization-period-de-rham}). If $x'\in H^j_\infty(Y')$, $\rho_m(x',f')$ for $(Y',\mathcal L')$ and
  $\rho_m(h^*x',f)$ for $(Y,\mathcal L)$, then $f(\psi)=f'(\psi\circ h_K^*)$ for every $\psi$.
\item Let $n\leq m$ and $r\geq0$, and for a thickening $T$ with ideal $J_T$ and divided powers $\delta_T$ let
  $J_T^{[k]}$ be as in Definition~\ref{gal-realization-kappa-approximations}. Suppose that $\rho_m(x,f)$, that for
  every $\nu$ the component $x_\nu\in H^j_\nu(Y)$ of $x$ is the class of a cocycle whose value on every chain
  $T_0\to\cdots\to T_j$ lies in $J_{T_j}^{[n]}$, and that $\psi$ vanishes on the classes in $H^j_{\mathrm{dR}}(Y_K)$ of
  the cocycles whose value on every chain $T_0\to\cdots\to T_j$ lies in $J_{T_j}^{[r+1]}$. Then $f(\psi)$ lies in the
  ideal $J_m^{[\max(n-r,0)]}$ of $B_m$, formed for $\delta_m$.
\end{enumerate}
Conditions (8) and (9) are used only for (3)(d) and (3)(e) of Lemma~\ref{gal-realization-period-ingredients}.
\end{definition}
\leanref{Development/WeilBound/CoreStatements/PadicHodgeComparison/EtaleCohomologyComparison/PerfectoidPullbacks/LevelChainValuations/FinDefs.lean}{CrystallineSemistableComparison.Pb.Kap.RhoSys}

\begin{definition}\label{gal-realization-kappa-truncated-periods-exist-diagram}
Let $K$ be a field and $\Lambda$, $A$ and $B$ rings.
\begin{enumerate}
\item A \emph{comparison diagram} $\mathcal D$ over $(\Lambda,A,B,K)$ consists of an integer $\ell$; for $n\geq1$ and
  $j\geq0$, a $\Lambda$-module $S^j_n$ and an $A$-module $Z^j_n$; for $j\geq0$, an $A$-module $Z^j$, a $B$-module $W^j$
  and a $K$-vector space $V^j$; ring maps $\lambda:\Lambda\to A$, $\iota:A\to B$ and $o:K\to B$; and, for all $n\geq1$
  and $j\geq0$, a $\Lambda$-linear map $t_S:S^j_{n+1}\to S^j_n$, an $A$-linear map $t_Z:Z^j_{n+1}\to Z^j_n$, a
  $\lambda$-semilinear map $r_n:S^j_n\to Z^j_n$, an $A$-linear map $\mathrm{red}_n:Z^j\to Z^j_n$, an $\iota$-semilinear
  map $b:Z^j\to W^j$ and an $o$-semilinear map $\mu:V^j\to W^j$. A family $x=(x_n)_{n\geq1}$ with $x_n\in S^j_n$ is
  \emph{compatible} if $t_S(x_{n+1})=x_n$ for every $n$.
\item $\mathcal D$ is \emph{effective} if the following hold for every $j$.
  \begin{enumerate}
  \item[(a)] The image of $\ell$ in $B$ is a unit, and $V^j$ is finite-dimensional.
  \item[(b)] For every basis $(v_t)$ of $V^j$ and every $w\in W^j$ there is exactly one family $(g_t)$ in $B$ with
    $w=\sum_tg_t\,\mu(v_t)$.
  \item[(c)] $t_Z\circ r_{n+1}=r_n\circ t_S$ and $t_Z\circ\mathrm{red}_{n+1}=\mathrm{red}_n$ for every $n$.
  \item[(d)] For every compatible family $x$ with $x_n\in S^j_n$ there are $e\geq0$ and $z\in Z^j$ with
    $\mathrm{red}_n(z)=\ell^er_n(x_n)$ for every $n$.
  \item[(e)] If $z\in Z^j$ and $\mathrm{red}_n(z)=0$ for every $n$, then $b(z)=0$.
  \end{enumerate}
\item Let $x=(x_n)_{n\geq1}$ be a family with $x_n\in S^j_n$. An additive map $f:\operatorname{Hom}_K(V^j,K)\to B$ is a
  \emph{period} of $x$ in $\mathcal D$ if $f(a\psi)=o(a)f(\psi)$ for $a\in K$, and there are $e\geq0$ and $z\in Z^j$ with
  $\mathrm{red}_n(z)=\ell^er_n(x_n)$ for every $n$ and
  \[
   \ell^e\sum_tf(v_t^*)\,\mu(v_t)=b(z)
  \]
  for every finite basis $(v_t)$ of $V^j$, with dual basis $(v_t^*)$.
\item Let $\mathcal D'$ be a comparison diagram over $(\Lambda',A',B',K)$, with data $\ell',S'^j_n,\ldots,\mu'$. A
  \emph{map} $h:\mathcal D\to\mathcal D'$ consists of ring maps $s_\Lambda:\Lambda\to\Lambda'$, $s_A:A\to A'$ and
  $s_B:B\to B'$, and for all $n$ and $j$ an $s_\Lambda$-semilinear map $h_S:S^j_n\to S'^j_n$, $s_A$-semilinear maps
  $h_Z:Z^j_n\to Z'^j_n$ and $h_Z:Z^j\to Z'^j$, an $s_B$-semilinear map $h_W:W^j\to W'^j$ and a $K$-linear map
  $h_V:V^j\to V'^j$, such that $\ell'=\ell$, $s_A\circ\lambda=\lambda'\circ s_\Lambda$, $s_B\circ\iota=\iota'\circ s_A$,
  $s_B\circ o=o'$, $h_S\circ t_S=t'_S\circ h_S$, $h_Z\circ r_n=r'_n\circ h_S$, $h_Z\circ\mathrm{red}_n=\mathrm{red}'_n\circ
  h_Z$, $h_W\circ b=b'\circ h_Z$ and $h_W\circ\mu=\mu'\circ h_V$. For an additive map $f:\operatorname{Hom}_K(V^j,K)\to
  B$, let $h_*f:\operatorname{Hom}_K(V'^j,K)\to B'$ send $\psi'$ to $s_B(f(\psi'\circ h_V))$.
\item A \emph{product structure} on $\mathcal D$ consists of a $\Lambda$-bilinear map $S^i_n\times S^j_n\to S^{i+j}_n$,
  $A$-bilinear maps $Z^i_n\times Z^j_n\to Z^{i+j}_n$ and $Z^i\times Z^j\to Z^{i+j}$, a $B$-bilinear map $W^i\times
  W^j\to W^{i+j}$ and a $K$-bilinear map $V^i\times V^j\to V^{i+j}$, for all $i$, $j$ and $n$, each written
  $(x,y)\mapsto xy$, and elements $1_n\in S^0_n$, $1_Z\in Z^0$ and $1_V\in V^0$, such that $t_S$, the $r_n$, the
  $\mathrm{red}_n$, $b$ and $\mu$ preserve products, $t_S(1_{n+1})=1_n$, $\mathrm{red}_n(1_Z)=r_n(1_n)$ and
  $b(1_Z)=\mu(1_V)$.
\end{enumerate}
\end{definition}
\leanref{Development/WeilBound/CoreStatements/PadicHodgeComparison/EtaleCohomologyComparison/PerfectoidPullbacks/LevelChainValuations/OneSTr.lean}{CrystallineSemistableComparison.Pb.Kap.Chain}
\leanref{Development/WeilBound/CoreStatements/PadicHodgeComparison/EtaleCohomologyComparison/PerfectoidPullbacks/LevelChainValuations/OneSTr.lean}{CrystallineSemistableComparison.Pb.Kap.Chain.Comp}
\leanref{Development/WeilBound/CoreStatements/PadicHodgeComparison/EtaleCohomologyComparison/PerfectoidPullbacks/LevelChainValuations/OneSTr.lean}{CrystallineSemistableComparison.Pb.Kap.Chain.Good}
\leanref{Development/WeilBound/CoreStatements/PadicHodgeComparison/EtaleCohomologyComparison/PerfectoidPullbacks/LevelChainValuations/OneSTr.lean}{CrystallineSemistableComparison.Pb.Kap.Chain.Rho}
\leanref{Development/WeilBound/CoreStatements/PadicHodgeComparison/EtaleCohomologyComparison/PerfectoidPullbacks/LevelChainValuations/OneSTr.lean}{CrystallineSemistableComparison.Pb.Kap.Hom}
\leanref{Development/WeilBound/CoreStatements/PadicHodgeComparison/EtaleCohomologyComparison/PerfectoidPullbacks/LevelChainValuations/OneSTr.lean}{CrystallineSemistableComparison.Pb.Kap.Hom.Good}
\leanref{Development/WeilBound/CoreStatements/PadicHodgeComparison/EtaleCohomologyComparison/PerfectoidPullbacks/LevelChainValuations/OneSTr.lean}{CrystallineSemistableComparison.Pb.Kap.Hom.tr}
\leanref{Development/WeilBound/CoreStatements/PadicHodgeComparison/EtaleCohomologyComparison/PerfectoidPullbacks/LevelChainValuations/OneSTr.lean}{CrystallineSemistableComparison.Pb.Kap.Cup}
\leanref{Development/WeilBound/CoreStatements/PadicHodgeComparison/EtaleCohomologyComparison/PerfectoidPullbacks/LevelChainValuations/OneSTr.lean}{CrystallineSemistableComparison.Pb.Kap.Cup.Good}

\begin{lemma}\label{gal-realization-kappa-truncated-periods-exist-periods}
Let $\mathcal D$ be a comparison diagram over $(\Lambda,A,B,K)$, with data as in
Definition~\ref{gal-realization-kappa-truncated-periods-exist-diagram}(1), let $j\geq0$, and let $x=(x_n)$ and
$y=(y_n)$ be families with $x_n,y_n\in S^j_n$.
\begin{enumerate}
\item If $\mathcal D$ is effective and $x$ is compatible, then $x$ has a period in $\mathcal D$.
\item If $\mathcal D$ is effective, then $x$ has at most one period in $\mathcal D$.
\item If $f$ is a period of $x$ and $g$ one of $y$, then $f+g$ is a period of $(x_n+y_n)_n$.
\item If $f$ is a period of $x$ and $c\in\Lambda$, then $\psi\mapsto\iota(\lambda(c))f(\psi)$ is a period of
  $(cx_n)_n$.
\item Let $h:\mathcal D\to\mathcal D'$ be a map of effective comparison diagrams. If $f$ is a period of $x$ in $\mathcal
  D$, then $h_*f$ is a period of $(h_S(x_n))_n$ in $\mathcal D'$.
\item Suppose that $\mathcal D$ is effective and has a product structure. Let $f$ be a period of $x$, $x'=(x'_n)$ a
  family with $x'_n\in S^{j'}_n$ with a period $f'$, and $g$ a period of $(x_nx'_n)_n$. Then
  \[
   g(\psi)=\sum_{s,s'}f(v_s^*)\,f'(v'^*_{s'})\,o\bigl(\psi(v_sv'_{s'})\bigr)
  \]
  for bases $(v_s)$ of $V^j$ and $(v'_{s'})$ of $V^{j'}$, with dual bases $(v_s^*)$ and $(v'^*_{s'})$, and
  $\psi\in\operatorname{Hom}_K(V^{j+j'},K)$.
\item Suppose that $\mathcal D$ is effective and has a product structure. If $g$ is a period of $(1_n)_n$, then
  $g(\psi)=o(\psi(1_V))$ for every $\psi\in\operatorname{Hom}_K(V^0,K)$.
\end{enumerate}
\end{lemma}
\leanref{Development/AutomorphicSpectrum/Proofs/BlockDecomposition/ExistsRealizationOfCompatibleStatement.lean}{GlobalTheory.T.exists_realization_of_compatible}
\leanref{Development/AutomorphicSpectrum/Proofs/BlockDecomposition/RealizationUniqueStatement.lean}{GlobalTheory.T.realization_unique}
\leanref{Development/AutomorphicSpectrum/Proofs/BlockDecomposition/RealizationAddStatement.lean}{GlobalTheory.T.realization_add}
\leanref{Development/AutomorphicSpectrum/Proofs/BlockDecomposition/RealizationSmulStatement.lean}{GlobalTheory.T.realization_smul}
\leanref{Development/AutomorphicSpectrum/Proofs/BlockDecomposition/RealizationFunctorialStatement.lean}{GlobalTheory.T.realization_functorial}
\leanref{Development/AutomorphicSpectrum/Proofs/BlockDecomposition/RealizationCupProductStatement.lean}{GlobalTheory.T.realization_cup_product}
\leanref{Development/AutomorphicSpectrum/Proofs/BlockDecomposition/RealizationUnitClassStatement.lean}{GlobalTheory.T.realization_unit_class}

\begin{proof}
Let $f:\operatorname{Hom}_K(V^j,K)\to B$ be additive with $f(a\psi)=o(a)f(\psi)$ for $a\in K$, and $(v_t)$ a finite basis
of $V^j$. Every $\psi$ is $\sum_t\psi(v_t)v_t^*$, and every $v\in V^j$ is $\sum_tv_t^*(v)v_t$, so
\[
 f(\psi)=\sum_tf(v_t^*)\,o(\psi(v_t)),\qquad \mu(v)=\sum_to(v_t^*(v))\,\mu(v_t).\tag{$*$}
\]
In (1), (2), (6) and (7), let $u$ be the inverse of $\ell$ in $B$, and recall that $b(az)=\iota(a)b(z)$ for $a\in A$.

(1) By (d) there are $e$ and $z$ with $\mathrm{red}_n(z)=\ell^er_n(x_n)$ for every $n$. Fix a basis $(v_t)$ of $V^j$,
which is finite by (a), and let $(g_t)$ be the family of (b) with $u^eb(z)=\sum_tg_t\,\mu(v_t)$. Put
$f(\psi)=\sum_to(\psi(v_t))g_t$; it is additive, and $f(a\psi)=o(a)f(\psi)$. For a finite basis $(v'_s)$ of $V^j$,
$f(v'^*_s)=\sum_to(v'^*_s(v_t))g_t$, so by $(*)$
\[
 \sum_sf(v'^*_s)\,\mu(v'_s)=\sum_tg_t\sum_so(v'^*_s(v_t))\,\mu(v'_s)=\sum_tg_t\,\mu(v_t)=u^eb(z),
\]
and $\ell^e$ times this is $b(z)$.

(2) Let $f$, with $e$ and $z$, and $f'$, with $e'$ and $z'$, be periods of $x$. Then $\mathrm{red}_n(\ell^ez'-\ell^{e'}z)=
\ell^{e+e'}r_n(x_n)-\ell^{e'+e}r_n(x_n)=0$ for every $n$, so $\ell^eb(z')=\ell^{e'}b(z)$ by (e). Fix a basis $(v_t)$ of
$V^j$. Then $\ell^{e+e'}\sum_tf'(v_t^*)\,\mu(v_t)=\ell^eb(z')=\ell^{e'}b(z)=\ell^{e'+e}\sum_tf(v_t^*)\,\mu(v_t)$. As $\ell$
is a unit, $\sum_t(f(v_t^*)-f'(v_t^*))\,\mu(v_t)=0=\sum_t0\cdot\mu(v_t)$, and (b) gives $f(v_t^*)=f'(v_t^*)$ for every $t$.
By $(*)$, $f=f'$.

(3) Let $(e_1,z_1)$ and $(e_2,z_2)$ be as in Definition~\ref{gal-realization-kappa-truncated-periods-exist-diagram}(3)
for $f$ and $g$. Then $e=e_1+e_2$ and $z=\ell^{e_2}z_1+\ell^{e_1}z_2$ have $\mathrm{red}_n(z)=\ell^er_n(x_n+y_n)$ and
$b(z)=\ell^e\sum_t(f+g)(v_t^*)\,\mu(v_t)$ for every finite basis $(v_t)$, and $f+g$ is semilinear as $f$ and $g$ are.

(4) Let $(e,z)$ be as there for $f$. As $r_n$ is $\lambda$-semilinear, $\mathrm{red}_n(\lambda(c)z)=\lambda(c)\ell^er_n(x_n)=
\ell^er_n(cx_n)$, and $b(\lambda(c)z)=\iota(\lambda(c))b(z)=\ell^e\sum_t\iota(\lambda(c))f(v_t^*)\,\mu(v_t)$.

(5) Let $(e,z)$ be as there for $f$, and $z'=h_Z(z)$. As $s_A(\ell)=\ell=\ell'$,
$\mathrm{red}'_n(z')=h_Z(\mathrm{red}_n(z))=\ell^eh_Z(r_n(x_n))=\ell^er'_n(h_S(x_n))$. For $a\in K$,
$h_*f(a\psi')=s_B(o(a)f(\psi'\circ h_V))=o'(a)\,h_*f(\psi')$. Fix a basis $(v_s)$ of $V^j$, and let $(v'_t)$ be a finite
basis of $V'^j$. Applying $h_W$ to $\ell^e\sum_sf(v_s^*)\,\mu(v_s)=b(z)$ gives $\ell^e\sum_ss_B(f(v_s^*))\,
\mu'(h_V(v_s))=b'(z')$. By $(*)$ for $\mathcal D'$, $\mu'(h_V(v_s))=\sum_to'(v'^*_t(h_V(v_s)))\,\mu'(v'_t)$, and by $(*)$
for $f$ and $\psi=v'^*_t\circ h_V$, with $s_B\circ o=o'$, $h_*f(v'^*_t)=\sum_ss_B(f(v_s^*))\,o'(v'^*_t(h_V(v_s)))$. So
$\sum_th_*f(v'^*_t)\,\mu'(v'_t)=\sum_ss_B(f(v_s^*))\,\mu'(h_V(v_s))$, and $\ell^e$ times it is $b'(z')$.

(6) Let $(e,z)$, $(e',z')$ and $(e'',z'')$ be as there for $f$, $f'$ and $g$. As $\mathrm{red}_n$ and $r_n$ preserve
products, $\mathrm{red}_n(zz')=\ell^{e+e'}r_n(x_nx'_n)$, so $\mathrm{red}_n(\ell^{e+e'}z''-\ell^{e''}zz')=0$ for every $n$,
and (e) gives $\ell^{e+e'}b(z'')=\ell^{e''}b(z)b(z')$. Let $(c_t)$ be a basis of $V^{j+j'}$ and
$C_t=\sum_{s,s'}f(v_s^*)f'(v'^*_{s'})\,o(c_t^*(v_sv'_{s'}))$. As $\mu$ preserves products, $(*)$ gives
$\mu(v_s)\mu(v'_{s'})=\sum_to(c_t^*(v_sv'_{s'}))\,\mu(c_t)$, so $b(z)b(z')=\ell^{e+e'}\sum_tC_t\,\mu(c_t)$. Hence
$\ell^{e+e'+e''}\sum_tg(c_t^*)\,\mu(c_t)=\ell^{e''+e+e'}\sum_tC_t\,\mu(c_t)$; as $\ell$ is a unit, (b) gives
$g(c_t^*)=C_t$. By $(*)$ for $g$, $g(\psi)=\sum_tC_t\,o(\psi(c_t))$, which is the claimed sum, as
$\sum_tc_t^*(v)\psi(c_t)=\psi(v)$ for $v\in V^{j+j'}$.

(7) Let $(e,z)$ be as there for $g$. Then $\mathrm{red}_n(z-\ell^e1_Z)=\ell^er_n(1_n)-\ell^er_n(1_n)=0$ for every $n$, so by
(e) $b(z)=\ell^eb(1_Z)=\ell^e\mu(1_V)$. For a basis $(v_t)$ of $V^0$, $\ell^e\sum_tg(v_t^*)\,\mu(v_t)=b(z)=\ell^e\mu(1_V)$,
so $\sum_tg(v_t^*)\,\mu(v_t)=\mu(1_V)=\sum_to(v_t^*(1_V))\,\mu(v_t)$ by $(*)$, and (b) gives $g(v_t^*)=o(v_t^*(1_V))$. By
$(*)$, $g(\psi)=\sum_to(v_t^*(1_V)\psi(v_t))=o(\psi(1_V))$.
% Lean: (1) is KapRho.exists_rho (with Module.finBasis for the basis); (2) is KapUniq.rho_unique through
% eq_of_unit_smul_combination; (3) is KapRho.rho_add; (4) is KapChain.rho_smul; (5) is KapRhoHom.rho_hom (with
% dual_expand for (*)); (6) is KapRho.rho_cup (with semilinear_eq_sum_coord for (*)); (7) is KapUnit.rho_one_eval.
\end{proof}
\leanref{Development/AutomorphicSpectrum/Proofs/BlockDecomposition/ExistsRealizationOfCompatible.lean}{GlobalTheory.CompanionTheory.KapRho.exists_rho}
\leanref{Development/AutomorphicSpectrum/Proofs/BlockDecomposition/RealizationUnique.lean}{GlobalTheory.CompanionTheory.KapUniq.rho_unique}
\leanref{Development/AutomorphicSpectrum/Proofs/BlockDecomposition/RealizationUnique.lean}{GlobalTheory.CompanionTheory.KapUniq.eq_of_unit_smul_combination}
\leanref{Development/AutomorphicSpectrum/Proofs/BlockDecomposition/RealizationAdd.lean}{GlobalTheory.CompanionTheory.KapRho.rho_add}
\leanref{Development/AutomorphicSpectrum/Proofs/BlockDecomposition/RealizationSmul.lean}{GlobalTheory.CompanionTheory.KapChain.rho_smul}
\leanref{Development/AutomorphicSpectrum/Proofs/BlockDecomposition/RealizationFunctorial.lean}{GlobalTheory.CompanionTheory.KapRhoHom.rho_hom}
\leanref{Development/AutomorphicSpectrum/Proofs/BlockDecomposition/RealizationCupProduct.lean}{GlobalTheory.CompanionTheory.KapRho.rho_cup}
\leanref{Development/AutomorphicSpectrum/Proofs/BlockDecomposition/RealizationUnitClass.lean}{GlobalTheory.CompanionTheory.KapUnit.rho_one_eval}

\begin{lemma}\label{gal-realization-kappa-truncated-periods-exist-reduction}
Let $\ell$, $\mathcal O_C^\flat$, $\xi$, $\theta$, $A_{\mathrm{cris}}$, $A_\nu$, $I_\nu$, $\gamma$ and $u:\mathcal
O_C\to A_\nu/I_\nu$ be as in Definition~\ref{gal-realization-period-thickenings}(1), $A_{\mathrm{inf}}=W(\mathcal
O_C^\flat)$, and $K$ and $\mathcal O_K$ as in Definition~\ref{gal-realization-period-datum}. Let $A$ be an $\mathcal
O_K$-algebra with divided powers $\gamma_A$ on an ideal $I$, $s\geq1$, and $\lambda:A_{\mathrm{cris}}\to A$ a ring map,
such that $I^{[s]}=0$ (Definition~\ref{gal-realization-kappa-approximations}), $\ell x\in I^{[r]}$ implies $x\in I^{[r]}$
for $r\geq0$ and $x\in A$, for every integer $n\geq1$ some power of $\ell$ lies in $nA$, $\lambda(\xi)\in I$, and
$a-\lambda(w)\in I$ for $a\in\mathcal O_K$ and $w\in A_{\mathrm{inf}}$ with $\theta(w)=a$. Let $\nu\geq1$, and write
$\bar y$ for the class in $A/\ell^\nu A$ of $y\in A$. Then $\lambda$ induces a ring map $\rho_\nu:A_\nu\to A/\ell^\nu A$
that sends the class of $c\in A_{\mathrm{cris}}$ to $\overline{\lambda(c)}$, and:
\begin{enumerate}
\item for $x\in I_\nu$ there is $y\in I$ with $\rho_\nu(x)=\bar y$; if $x\in I_\nu$ and $y\in I$ satisfy
  $\rho_\nu(x)=\bar y$, then $\rho_\nu(\gamma_k(x))=\overline{\gamma_{A,k}(y)}$ for every $k\geq0$;
\item if $a\in\mathcal O_K$ and $x\in A_\nu$ satisfy $u(a)=x+I_\nu$, then $\rho_\nu(x)=\overline{a+y}$ for some
  $y\in I$.
\end{enumerate}
\end{lemma}
\leanref{Development/WeilBound/CoreStatements/PadicHodgeComparison/EtaleCohomologyComparison/PerfectoidPullbacks/DividedPowerReduction/JEqPdPow.lean}{CrystallineSemistableComparison.Pb.RedIn}
\leanref{Development/WeilBound/CoreStatements/PadicHodgeComparison/EtaleCohomologyComparison/PerfectoidPullbacks/DividedPowerReduction/RedInEx.lean}{CrystallineSemistableComparison.Pb.redIn_ex}

\begin{proof}
The projection $A_{\mathrm{cris}}\to A_\nu$ is onto with kernel $\ell^\nu A_{\mathrm{cris}}$, which $\lambda$ maps into
$\ell^\nu A$; so $\rho_\nu$ exists. Write $A^0=A_{\mathrm{cris}}(\mathcal O_C^\flat)$ and $I^0=I(\mathcal O_C^\flat)$ as in
Definition~\ref{gal-realization-period-thickenings}(1), and $\lambda^0:A^0\to A$ for $\lambda$ composed with $A^0\to
A_{\mathrm{cris}}$. Then $\rho_\nu$ sends the class of $y\in A^0$ to $\overline{\lambda^0(y)}$, $I_\nu$ is the image of
$I^0$, and $\gamma_k$ sends the class of $x\in I^0$ to the class of $x^k/k!\in A^0$.

\emph{Cancelling.} (i) If $\ell^Nx\in I$, then $x\in I$, by the hypothesis for $r=1$, as $I^{[1]}=I$, applied $N$ times.
(ii) If $\ell^Nx=0$, then $x=0$, by the hypothesis for $r=s$, as $I^{[s]}=0$. (iii) For $k\geq0$ there are $N\geq0$ and
$v\in A$ with $k!\,v=\ell^N$; so $k!\,x\in I$ implies $x\in I$, and $k!\,x=k!\,x'$ implies $x=x'$, by (i) and (ii).
(iv) If $x,x'\in I$ and $x-x'\in\ell^\nu A$, then $\gamma_{A,k}(x)-\gamma_{A,k}(x')\in\ell^\nu A$ for $k\geq0$. Indeed,
$x=x'+\ell^\nu a$ with $a\in I$ by (i), and $\gamma_{A,k}(x)=\sum_{i+i'=k}\gamma_{A,i}(x')\gamma_{A,i'}(\ell^\nu a)$, where
$\gamma_{A,i'}(\ell^\nu a)=\ell^{\nu i'}\gamma_{A,i'}(a)$ lies in $\ell^\nu A$ for $i'\geq1$.

\emph{The map $\lambda^0$ on $I^0$.} For $k\geq1$, $k!\,\lambda^0(\xi^k/k!)=\lambda(\xi)^k\in I$, so
$\lambda^0(\xi^k/k!)\in I$ by (iii); hence $\lambda^0(I^0)\subset I$. For $x\in I^0$ and $k\geq0$,
$k!\,\lambda^0(x^k/k!)=\lambda^0(x)^k=k!\,\gamma_{A,k}(\lambda^0(x))$, so $\lambda^0(x^k/k!)=\gamma_{A,k}(\lambda^0(x))$ by
(iii).

(1) Let $x$ be the class of $x_0\in I^0$. Then $\rho_\nu(x)=\overline{\lambda^0(x_0)}$ with $\lambda^0(x_0)\in I$. If also
$\rho_\nu(x)=\bar y$ with $y\in I$, then $\rho_\nu(\gamma_k(x))=\overline{\lambda^0(x_0^k/k!)}=
\overline{\gamma_{A,k}(\lambda^0(x_0))}=\overline{\gamma_{A,k}(y)}$, by (iv).

(2) Choose $w\in A_{\mathrm{inf}}$ with $\theta(w)=a$ (Lemma~\ref{gal-realization-kappa-fontaine}(1)). Then $u(a)$ is the
class of $w$, by the definition of $u$, so $x$ minus the class of $w$ lies in $I_\nu$; let it be the class of $x_0\in
I^0$. Then $\rho_\nu(x)=\overline{\lambda(w)+\lambda^0(x_0)}=\overline{a+y}$ with $y=\lambda^0(x_0)-(a-\lambda(w))$, which
lies in $I$.
% Lean: RedInEx.redIn_ex, with ρ = RedInEx.rho (liftOfSurjective through RL.base_acr), rho_cris and rho_mk; A^0 is
% (base l).Acr0 = Acris0 l through e0, λ^0 is lam0, and I^0 is Fil0 1, whose image is the ideal of B.DD false (isPd_pdFil).
% (i) to (iv) are mem_of_pow_mul, eq_zero_of_pow_mul, mem_of_fact_mul with eq_of_fact_mul, and dpow_congr; the map λ^0 on
% I^0 is lam0_mem_of_mem_fil0_one and lam0_dpow (x^k/k! from RL.pdi through exists_unit_one_add_generator); (2) is the hκ
% field, with surjective_fontaineTheta, B.uJ_th and the hypothesis hO, which is Lev.CrisOK.
\end{proof}
\leanref{Development/AutomorphicSpectrum/Proofs/BlockDecomposition/DividedPowersOnFiltration.lean}{GlobalTheory.CompanionTheory.RedInEx.redIn_ex}
\leanref{Development/AutomorphicSpectrum/Proofs/BlockDecomposition/DividedPowersOnFiltration.lean}{GlobalTheory.CompanionTheory.RedInEx.rho}
\leanref{Development/AutomorphicSpectrum/Proofs/BlockDecomposition/DividedPowersOnFiltration.lean}{GlobalTheory.CompanionTheory.RedInEx.dpow_congr}
\leanref{Development/AutomorphicSpectrum/Proofs/BlockDecomposition/DividedPowersOnFiltration.lean}{GlobalTheory.CompanionTheory.RedInEx.lam0_mem_of_mem_fil0_one}
\leanref{Development/AutomorphicSpectrum/Proofs/BlockDecomposition/DividedPowersOnFiltration.lean}{GlobalTheory.CompanionTheory.RedInEx.lam0_dpow}

\begin{definition}\label{gal-realization-kappa-truncated-periods-exist-categories}
Let $\ell$, $A_{\mathrm{cris}}$, $K$, $\mathcal O_K$ and $G_K$ be as in Definition~\ref{gal-realization-period-datum}, and
let the approximations $(B_m,\tau_m,J_m,\delta_m,A_m,I_m,\gamma_m,\iota_m,\lambda_m,\sigma_{A_m},\sigma_{B_m})$ of height
$m+1$, with $d_A$ and $d_B$, form an approximation system (Definition~\ref{gal-realization-kappa-approximations}). For
$m\geq0$ and $\nu\geq1$ let $\rho_{m,\nu}:A_\nu\to A_m/\ell^\nu A_m$ be the ring map of
Lemma~\ref{gal-realization-kappa-truncated-periods-exist-reduction} for $A=A_m$, $\lambda=\lambda_m$ and $s=m+1$, whose
hypotheses are conditions (d) and (e) of Definition~\ref{gal-realization-kappa-approximations}(1). Let $Y$ be a smooth
projective $\mathcal O_K$-scheme with generic fibre $Y_K=Y\times_{\mathcal O_K}K$ and projection $\mathrm{pr}:Y_K\to Y$,
and let $m\geq0$ and $s\geq1$.
\begin{enumerate}
\item $\mathcal Q_m(Y,s)$ is the full subcategory of $\mathrm{Th}^{(s)}(Y/A_m)$, formed for $(A_m,I_m,\gamma_m,0)$,
  $\mathcal O=\mathcal O_K$ and the structure map $\mathcal O_K\to A_m/I_m$
  (Definition~\ref{gal-realization-kappa-de-rham-truncated}(1)), of the thickenings $(B,i,J,\delta,y)$ such that $\ell
  x\in J^{[r]}$ implies $x\in J^{[r]}$ for all $r\geq0$ and $x\in B$. $\mathcal Q'_m(Y,s)$ is
  $\mathrm{Th}^{(s)}(Y_K/B_m)$, formed for $(B_m,J_m,\delta_m,0)$, $\mathcal O=K$ and the structure map $K\to B_m/J_m$.
\item Let $T=(B,i,J,\delta,y)$ be in $\mathcal Q_m(Y,s)$, $\nu\geq1$, $\bar B=B/\ell^\nu B$ and $\pi:B\to\bar B$ the
  projection. If $\ell^\nu b\in J$, then $b\in J$, by the condition of (1) for $r=1$, as $J^{[1]}=J$; and
  $\delta_k(\ell^\nu b)=\ell^{\nu k}\delta_k(b)\in\ell^\nu B$ for $k\geq1$. So $\delta$ induces divided powers $\bar\delta$
  on $\bar J=\pi(J)$ for which $\pi$ is a divided power map \cite[Tag 07H2]{Sta}. Let $R_\nu(T)=(\bar
  B,\bar\imath\rho_{m,\nu},\bar J,\bar\delta,\bar y)$, where $\bar\imath:A_m/\ell^\nu A_m\to\bar B$ is induced by $i$ and
  $\bar y$ is $\operatorname{Spec}\bar B/\bar J\to\operatorname{Spec}B/J$ followed by $y$. By
  Lemma~\ref{gal-realization-kappa-truncated-periods-exist-reduction}(1), $\bar\imath\rho_{m,\nu}$ carries $I_\nu$ into
  $\bar J$ compatibly with the divided powers, and by
  Lemma~\ref{gal-realization-kappa-truncated-periods-exist-reduction}(2) the $\mathcal O_K$-algebra structures of $\bar
  B/\bar J$ through $u$ and through $A_m/I_m$ agree. So $R_\nu(T)$ is a thickening of $Y$ of level $\nu$ for $I_\nu$
  (Definition~\ref{gal-realization-period-thickenings}(2)). Maps $T\to T'$ induce maps $R_\nu(T)\to R_\nu(T')$, and
  $R_\nu$ is a functor $\mathcal Q_m(Y,s)\to\mathrm{Th}_\nu(Y)$.
\item For $T'=(B',i',J',\delta',y')$ in $\mathcal Q'_m(Y,s)$, $N(T')=(B',i'\iota_m,J',\delta',\mathrm{pr}\circ y')$ lies
  in $\mathcal Q_m(Y,s)$: $\iota_m$ is a divided power map, $\iota_m(a)=a$ for $a\in\mathcal O_K$, $I_m=\iota_m^{-1}(J_m)$,
  and $\ell$ is a unit of the $K$-algebra $B'$. This gives a functor $N:\mathcal Q'_m(Y,s)\to\mathcal Q_m(Y,s)$.
  Restriction of scalars along $K\to B_m$ (Definition~\ref{gal-realization-thickening-nilpotent}(4)) is a functor
  $E:\mathcal Q'_m(Y,s)\to\mathrm{Th}^{(s)}_{\mathrm{dR}}(Y_K)$ (Definition~\ref{gal-realization-kappa-de-rham-truncated}(2)).
\item The comparison diagram $\mathcal D_m(Y,s)$ over $(A_{\mathrm{cris}},A_m,B_m,K)$
  (Definition~\ref{gal-realization-kappa-truncated-periods-exist-diagram}) has the integer $\ell$;
  $S^j_n=H^j_n(Y)$ (Definition~\ref{gal-realization-period-thickenings}), with $t_S$ the reduction;
  $Z^j_n=H^j(\mathcal Q_m(Y,s),\mathcal O/\ell^n\mathcal O)$ and $Z^j=H^j(\mathcal Q_m(Y,s),\mathcal O)$
  (Definition~\ref{gal-realization-thickening-nilpotent}(3)), with $t_Z$ and $\mathrm{red}_n$ induced by the projections;
  $W^j=H^j(\mathcal Q'_m(Y,s),\mathcal O)$; $V^j=H^j_{\mathrm{dR}}(Y_K)$
  (Definition~\ref{gal-realization-period-de-rham}); $\lambda=\lambda_m$, $\iota=\iota_m$ and $o:K\to B_m$ the structure
  map; $r_n=R_n^*$ and $b=N^*$, with the identity maps of the rings $\bar B$ and $B'$
  (Definition~\ref{gal-realization-split-diagrams}(1)); and $\mu$ the restriction $H^j_{\mathrm{dR}}(Y_K)\to
  H^j(\mathrm{Th}^{(s)}_{\mathrm{dR}}(Y_K),\mathcal O)$ followed by $E^*$, with identity maps. Its \emph{cup data} are
  the cup products of Definitions~\ref{gal-realization-split-diagrams}(1), \ref{gal-realization-period-thickenings}(2)
  and~\ref{gal-realization-period-de-rham}, and the classes $1_n$, $1_Z$ and $1_V$ of the constant cochains $1$.
\end{enumerate}
\end{definition}
\leanref{Development/WeilBound/CoreStatements/PadicHodgeComparison/EtaleCohomologyComparison/PerfectoidPullbacks/LevelChainValuations/RhoFirstChernLaw.lean}{CrystallineSemistableComparison.Pb.Kap.dOf}
\leanref{Development/WeilBound/CoreStatements/PadicHodgeComparison/EtaleCohomologyComparison/PerfectoidPullbacks/SplittingFieldLattices/GrSplit.lean}{CrystallineSemistableComparison.Pb.TS}
\leanref{Development/WeilBound/CoreStatements/PadicHodgeComparison/EtaleCohomologyComparison/PerfectoidPullbacks/SEqOne.lean}{CrystallineSemistableComparison.Pb.TF}
\leanref{Development/WeilBound/CoreStatements/PadicHodgeComparison/EtaleCohomologyComparison/PerfectoidPullbacks/LevelChainValuations/ChainOf.lean}{CrystallineSemistableComparison.Pb.Kap.QA}
\leanref{Development/WeilBound/CoreStatements/PadicHodgeComparison/EtaleCohomologyComparison/PerfectoidPullbacks/LevelChainValuations/ChainOf.lean}{CrystallineSemistableComparison.Pb.Kap.QB}
\leanref{Development/WeilBound/CoreStatements/PadicHodgeComparison/EtaleCohomologyComparison/PerfectoidPullbacks/DividedPowerReduction/JEqPdPow.lean}{CrystallineSemistableComparison.Pb.subDP}
\leanref{Development/WeilBound/CoreStatements/PadicHodgeComparison/EtaleCohomologyComparison/PerfectoidPullbacks/DividedPowerReduction/JEqPdPow.lean}{CrystallineSemistableComparison.Pb.redObj}
\leanref{Development/WeilBound/CoreStatements/PadicHodgeComparison/EtaleCohomologyComparison/PerfectoidPullbacks/DividedPowerReduction/JEqPdPow.lean}{CrystallineSemistableComparison.Pb.redG}
\leanref{Development/WeilBound/CoreStatements/PadicHodgeComparison/EtaleCohomologyComparison/PerfectoidPullbacks/DividedPowerReduction/JEqPdPow.lean}{CrystallineSemistableComparison.Pb.redMor}
\leanref{Development/WeilBound/CoreStatements/PadicHodgeComparison/EtaleCohomologyComparison/PerfectoidPullbacks/LevelChainValuations/ChainOf.lean}{CrystallineSemistableComparison.Pb.Kap.nmB}
\leanref{Development/WeilBound/CoreStatements/PadicHodgeComparison/EtaleCohomologyComparison/PerfectoidPullbacks/LevelChainValuations/ValuationChainProofs/QbLePf.lean}{CrystallineSemistableComparison.Pb.Kap.qb_le_pf}
\leanref{Development/WeilBound/CoreStatements/PadicHodgeComparison/EtaleCohomologyComparison/PerfectoidPullbacks/LevelChainValuations/ChainOf.lean}{CrystallineSemistableComparison.Pb.Kap.bMap}
\leanref{Development/WeilBound/CoreStatements/PadicHodgeComparison/EtaleCohomologyComparison/PerfectoidPullbacks/SplittingFieldLattices/GrSplit.lean}{CrystallineSemistableComparison.Pb.muB}
\leanref{Development/WeilBound/CoreStatements/PadicHodgeComparison/EtaleCohomologyComparison/PerfectoidPullbacks/SplittingFieldLattices/GrSplit.lean}{CrystallineSemistableComparison.Pb.redM}
\leanref{Development/WeilBound/CoreStatements/PadicHodgeComparison/EtaleCohomologyComparison/PerfectoidPullbacks/SplittingFieldLattices/GrSplit.lean}{CrystallineSemistableComparison.Pb.trM}
\leanref{Development/WeilBound/CoreStatements/PadicHodgeComparison/EtaleCohomologyComparison/PerfectoidPullbacks/LevelChainValuations/RHof.lean}{CrystallineSemistableComparison.Pb.Kap.rHof}
\leanref{Development/WeilBound/CoreStatements/PadicHodgeComparison/EtaleCohomologyComparison/PerfectoidPullbacks/LevelChainValuations/ChainOf.lean}{CrystallineSemistableComparison.Pb.Kap.chainOf}
\leanref{Development/WeilBound/CoreStatements/PadicHodgeComparison/EtaleCohomologyComparison/PerfectoidPullbacks/LevelChainValuations/RhoFirstChernLaw.lean}{CrystallineSemistableComparison.Pb.Kap.chK}
\leanref{Development/WeilBound/CoreStatements/PadicHodgeComparison/EtaleCohomologyComparison/PerfectoidPullbacks/LevelChainValuations/GaloisCupProducts/CupZnApply.lean}{CrystallineSemistableComparison.Pb.Kap.cupOf}

\begin{lemma}\label{gal-realization-kappa-truncated-periods-exist-base-change-zariski}
Let $(R,I,\gamma,0)$, $\mathcal O$, $X$ and $\kappa$ be as in Definition~\ref{gal-realization-kappa-de-rham-truncated}, and
$s,\ell\geq1$ integers. Let $\mathrm{Th}^{(s)}_\ell(X/R)$ be the full subcategory of $\mathrm{Th}^{(s)}(X/R)$ of the
thickenings $(B,i,J,\delta,y)$ such that $\ell x\in J^{[r]}$ implies $x\in J^{[r]}$ for all $r\geq0$ and $x\in B$; thus
$\mathrm{Th}^{(s)}_1(X/R)=\mathrm{Th}^{(s)}(X/R)$. Then (1) and (2) of Lemma~\ref{gal-realization-thickening-zariski} hold
for $p=0$ and $\mathcal Q_0=\mathrm{Th}^{(s)}_\ell(X/R)$.
\end{lemma}
\leanref{Development/AutomorphicSpectrum/Proofs/BlockDecomposition/CechRefinementCohomologyBijective.lean}{GlobalTheory.CompanionTheory.NlZar.used_locN}
\leanref{Development/AutomorphicSpectrum/Proofs/BlockDecomposition/CechRefinementCohomologyBijective.lean}{GlobalTheory.CompanionTheory.NlZar.tf_loc}

\begin{proof}
The proof of Lemma~\ref{gal-realization-kappa-de-rham-zariski} applies, once the localization $T'_g$ of a thickening
$T'=(B,i,J,\delta,y)$ of $\mathrm{Th}^{(s)}_\ell(X/R)$ at an element $g\in B$ is known to lie in
$\mathrm{Th}^{(s)}_\ell(X/R)$. Let $\lambda:B\to B_g$ be the localization map. By
Lemma~\ref{gal-realization-thickening-zariski-localization}(2), $(JB_g)^{[r]}=\lambda(J^{[r]})B_g$ for $r\geq0$, which is
$0$ for $r=s$. Let $r\geq0$ and $x\in B_g$ with $\ell x\in\lambda(J^{[r]})B_g$, and write $x=\lambda(b)/g^n$ with $b\in B$.
Then $\lambda(\ell b)=g^n\ell x$ lies in $\lambda(J^{[r]})B_g$, so $g^k\ell b\in J^{[r]}$ for some $k\geq0$. By the condition
on $T'$, $g^kb\in J^{[r]}$, and $x=\lambda(g^kb)/g^{n+k}$ lies in $\lambda(J^{[r]})B_g$.
% Lean: NlZar.zariski_descent (GlobalTheory.T.cech_refinement_cohomology_bijective) for Used Q₀ with p = 0 and Q₀ = TF l s,
% the book's Th^(s)_l(X/R); used_locN gives the closure under localization, through pdPow_locN and tf_loc.
\end{proof}

\begin{definition}\label{gal-realization-kappa-truncated-periods-exist-base-change-splitting}
Let $\Lambda$ be a ring, $A$ a $\Lambda$-algebra with divided powers $\gamma$ on an ideal $I$, and $m\geq0$ an integer with
$I^{[m]}=0$ (Definition~\ref{gal-realization-kappa-approximations}). A \emph{graded splitting} of $(A,\gamma)$ of length $m$
over $\Lambda$ consists of flat $\Lambda$-modules $N_0,\ldots,N_{m-1}$ and $\Lambda$-linear maps $\nu_i:N_i\to A$ such that,
for every $i<m$, $\nu_i(N_i)\subset I^{[i]}$, $\nu_i(n)\in I^{[i+1]}$ only for $n=0$, and every $x\in I^{[i]}$ has an $n\in
N_i$ with $x-\nu_i(n)\in I^{[i+1]}$.
\end{definition}
\leanref{Development/WeilBound/CoreStatements/PadicHodgeComparison/EtaleCohomologyComparison/PerfectoidPullbacks/SplittingFieldLattices/GrSplit.lean}{CrystallineSemistableComparison.Pb.GrSplit}

\begin{lemma}\label{gal-realization-kappa-truncated-periods-exist-base-change-splitting-field}
Let $\Lambda$ be a ring, $A$ a $\Lambda$-algebra with divided powers $\gamma$ on an ideal $I$, and $m\geq0$ with
$I^{[m]}=0$.
\begin{enumerate}
\item If $\Lambda$ is a field, $(A,\gamma)$ has a graded splitting of length $m$ over $\Lambda$
  (Definition~\ref{gal-realization-kappa-truncated-periods-exist-base-change-splitting}).
\item Let $(N_i,\nu_i)$ be a graded splitting of length $m$ over $\Lambda$. Every $x\in A$ is $\sum_{i<m}\nu_i(n_i)$ for
  exactly one family $(n_i)$ with $n_i\in N_i$, and for $0\leq k\leq m$, $x\in I^{[k]}$ if and only if $n_i=0$ for $i<k$.
\end{enumerate}
\end{lemma}
\leanref{Development/AutomorphicSpectrum/Proofs/BlockDecomposition/GradedSplittingOverFieldStatement.lean}{GlobalTheory.T.graded_splitting_over_field}
\leanref{Development/AutomorphicSpectrum/Proofs/BlockDecomposition/BasisBOfBcOn.lean}{GlobalTheory.CompanionTheory.BasisBase.gr_decomp}

\begin{proof}
The ideals $I^{[k]}$ decrease in $k$, and $I^{[0]}=A$.

(1) For $i<m$ choose a $\Lambda$-subspace $N_i$ of $I^{[i]}$ with $I^{[i]}=N_i\oplus I^{[i+1]}$, and let $\nu_i$ be the
inclusion. Vector spaces are flat.

(2) Existence, for $x\in I^{[k]}$ with a family that has $n_i=0$ for $i<k$, by descending induction on $k$: for $k=m$,
$x\in I^{[m]}=0$ and all $n_i=0$; for $k<m$, choose $n_k$ with $x-\nu_k(n_k)\in I^{[k+1]}$ and apply the induction to
$x-\nu_k(n_k)$. For $k=0$ this covers every $x\in A$. Uniqueness: if $\sum_i\nu_i(n_i)=0$, then $n_i=0$ for $i<k$, by
induction on $k$; given this for $i<k$, $\nu_k(n_k)=-\sum_{i>k}\nu_i(n_i)\in I^{[k+1]}$, so $n_k=0$. Hence the family of
$x\in I^{[k]}$ is the one constructed, with $n_i=0$ for $i<k$; conversely such a sum lies in $I^{[k]}$.
% Lean: (1) is GrSplitField.nonempty_grSplit (exists_step chooses the complements); (2) is BasisBase.gr_decomp (key, by
% induction on m - k, for existence, and the induction on k for uniqueness). The statement on I^[k] is the case Γ = Λ of
% GrTens.main (graded_splitting_tensor_filtration).
\end{proof}
\leanref{Development/AutomorphicSpectrum/Proofs/BlockDecomposition/GradedSplittingOverField.lean}{GlobalTheory.CompanionTheory.GrSplitField.nonempty_grSplit}

\begin{lemma}\label{gal-realization-kappa-truncated-periods-exist-base-change-chart}
Let $\mathcal O$ be a ring, $R$ an $\mathcal O$-algebra with divided powers $\gamma$ on an ideal $I$, and $s,\ell\geq1$
integers such that $I^{[s]}=0$ (Definition~\ref{gal-realization-kappa-approximations}), $\ell x\in I^{[r]}$ implies $x\in
I^{[r]}$ for $r\geq0$ and $x\in R$, and for every integer $n\geq1$ some power of $\ell$ lies in $nR$. Let $X$ be an
$\mathcal O$-scheme and $\kappa:\mathcal O\to R/I$ the structure map followed by the projection.
\begin{enumerate}
\item Let $t_1,\ldots,t_d$ be coordinates on an affine open $V\subset X$
  (Definition~\ref{gal-realization-thickening-coordinates}(1)), $\Gamma=\Gamma(V,\mathcal O_X)$, $B=R\otimes_{\mathcal
  O}\Gamma$ and $i:R\to B$, $a\mapsto a\otimes1$. Then $\Gamma$ is flat over $\mathcal O$, and there is a divided power
  chart $D$ with $t'_1,\ldots,t'_d$ of dimension $d$ on $V$ over $(R,I,\gamma,0)$
  (Definition~\ref{gal-realization-thickening-coordinates}(3)) with $B_D=B$, $i_D=i$ and $t'_k=1\otimes t_k$, which lies
  in $\mathrm{Th}^{(s)}_\ell(X/R)$ (Lemma~\ref{gal-realization-kappa-truncated-periods-exist-base-change-zariski}), such
  that the ring map $\Gamma\to B/J_D$ that corresponds to $y_D$ sends $f$ to the class of $1\otimes f$.
\item For a set $\mathcal V$ of open subsets of $X$ and $\mathcal Q=\mathrm{Th}^{(s)}_\ell(X/R)_{\mathcal V}$,
  $\Phi_{\mathcal Q}$ (Definition~\ref{gal-realization-thickening-torsor}(1)) has a divided power polynomial structure
  that respects, in every degree $r\leq s$, the sequence $(J^{[r']})_{r'}$ and, for every $c\in R$, the sequence
  $(J^{[r']}+cB)_{r'}$.
\end{enumerate}
\end{lemma}
\leanref{Development/AutomorphicSpectrum/Proofs/BlockDecomposition/ExistsChartFromCoordinates.lean}{GlobalTheory.CompanionTheory.ChartAlg.chart_exists}
\leanref{Development/WeilBound/CoreStatements/PadicHodgeComparison/EtaleCohomologyComparison/PerfectoidPullbacks/ChartGrConst.lean}{CrystallineSemistableComparison.Pb.poly_tf}

\begin{proof}
(1) The ring $\Gamma$ is \'etale over $\mathcal O[x_1,\ldots,x_d]$, which is free over $\mathcal O$, so $\Gamma$ is flat over
$\mathcal O$ \cite[Tags 00U2 and 00HC]{Sta}, and $B$ is flat over $R$. Steps 2 and 3 of the proof of
Lemma~\ref{gal-realization-kappa-de-rham-chart} do not use that the integers are invertible; they apply here and give the
chart, once $J=i(I)B$ carries divided powers $\delta$ for which $i$ is a divided power map, with $J^{[s]}=0$ and $J$
nilpotent. The next three paragraphs show this, and Step 1 there refers to them.

\emph{Torsion.} Multiplication by $\ell$ is injective on $R$: if $\ell x=0$, then $\ell x\in I^{[s]}$, so $x\in
I^{[s]}=0$. As $B$ is flat over $R$, it is injective on $B$. For $n\geq1$ write $\ell^k=nu$ with $u\in R$; if $nx=0$ in
$B$, then $\ell^kx=0$, so $x=0$. So $B$ has no $\mathbf Z$-torsion.

\emph{The divided powers.} For $x\in J$ there are $y_n\in B$, $n\geq0$, with $y_0=1$, $y_n\in J$ for $n\geq1$ and
$n!\,y_n=x^n$: for $x=i(a)b$ with $a\in I$ take $y_n=i(\gamma_n(a))b^n$; if $(y_n)$ serves $x$ and $(y'_n)$ serves $x'$,
then $(\sum_{a+b=n}y_ay'_b)$ serves $x+x'$, by the binomial formula, and $(b^ny_n)$ serves $bx$. As $B$ has no
$\mathbf Z$-torsion, $y_n$ is determined by $x$; put $\delta_n(x)=y_n$. Each axiom of divided powers for $\delta$ is an
identity that holds after multiplication by a nonzero integer, for example $m!\,n!\,\delta_m(x)\delta_n(x)=x^{m+n}=
m!\,n!\binom{m+n}m\delta_{m+n}(x)$, so it holds. For $a\in I$, $n!\,\delta_n(i(a))=i(a)^n=n!\,i(\gamma_n(a))$, so $i$ is a
divided power map. By the proof of Lemma~\ref{gal-realization-thickening-chart}(2), $J^{[r]}=i(I^{[r]})B$, the image of
$I^{[r]}\otimes_{\mathcal O}\Gamma$, for $r\geq0$. So $J^{[s]}=0$, and $J^s=0$, as a product $x_1\cdots x_s=\delta_1(x_1)
\cdots\delta_1(x_s)$ of elements of $J$ lies in $J^{[s]}$.

\emph{The condition on $\ell$.} As $\Gamma$ is flat, $B/J^{[r]}=(R/I^{[r]})\otimes_{\mathcal O}\Gamma$, and multiplication
by $\ell$, which is injective on $R/I^{[r]}$ by hypothesis, stays injective after tensoring with $\Gamma$. So $\ell x\in
J^{[r]}$ implies $x\in J^{[r]}$, and $D$ lies in $\mathrm{Th}^{(s)}_\ell(X/R)$.

(2) This is Lemma~\ref{gal-realization-kappa-de-rham-poly}, which is stated for $\mathrm{Th}^{(s)}_\ell(X/R)_{\mathcal
V}$; its proof checks that the thickening $P(T)$ constructed there satisfies the condition on $\ell$.
% Lean: (1) is ChartAlg.chart_exists for TFBase γ l s, as is Lemma gal-realization-kappa-de-rham-chart. Torsion is
% hl, hlk and hn; the divided powers are exists_dp (hasDP_of_mem, fact_sum, dpExt with cancel), with
% dividedPowerFiltration_eq_map for J^[r]; flatness is flat_of_etale and the condition on l is tf_tensor. (2) is
% Pb.poly_tf = Site.poly_tf' with polyTf, whose membership in TS γ q l s 𝒱 is polyTr_mem with SiteTr.tr_tf; the condition
% on l is the second part of tr_tf, through the filtration criterion (divided_power_polynomial_filtration).
\end{proof}
\leanref{Development/AutomorphicSpectrum/Proofs/BlockDecomposition/ExistsChartFromCoordinates.lean}{GlobalTheory.CompanionTheory.ChartAlg.exists_dp}
\leanref{Development/AutomorphicSpectrum/Proofs/BlockDecomposition/ExistsChartFromCoordinates.lean}{GlobalTheory.CompanionTheory.ChartAlg.tf_tensor}
\leanref{Development/AutomorphicSpectrum/Proofs/BlockDecomposition/TruncatedEnvelopeTorsionFree.lean}{GlobalTheory.CompanionTheory.SiteTr.tr_tf}

\begin{lemma}\label{gal-realization-kappa-truncated-periods-exist-base-change-tensor}
Let $\Lambda$ be a ring, $A$ a $\Lambda$-algebra with divided powers $\gamma$ on an ideal $I$, and $s,\ell\geq1$
integers. Suppose that the following hold for $(R,I_R)=(\Lambda,0)$, with the zero divided powers, and for
$(R,I_R)=(A,I)$, with $\gamma$: $I_R^{[s]}=0$; $\ell x\in I_R^{[r]}$ implies $x\in I_R^{[r]}$ for $r\geq0$ and $x\in R$;
and for every integer $n\geq1$ some power of $\ell$ lies in $nR$. Let $(N_i,\nu_i)_{i<m}$ be a graded splitting of
$(A,\gamma)$ over $\Lambda$ (Definition~\ref{gal-realization-kappa-truncated-periods-exist-base-change-splitting}). Let
$X$ be a $\Lambda$-scheme with a coordinate cover (Definition~\ref{gal-realization-thickening-coordinates}(2)),
$\mathcal V$ a set of open subsets of $X$ with union $X$, $\mathcal Q_\Lambda=\mathrm{Th}^{(s)}_\ell(X/\Lambda)_{\mathcal
V}$, formed for $(\Lambda,0,0,0)$ and the identity of $\Lambda$, and $\mathcal Q_A=\mathrm{Th}^{(s)}_\ell(X/A)_{\mathcal V}$,
formed for $(A,I,\gamma,0)$ and $\Lambda\to A\to A/I$
(Lemma~\ref{gal-realization-kappa-truncated-periods-exist-base-change-zariski}).
\begin{enumerate}
\item The restriction of scalars $G$ along $\Lambda\to A$ maps $\mathcal Q_A$ into $\mathcal Q_\Lambda$, and is, with the
  identity maps of the rings, a morphism $\Phi_{\mathcal Q_\Lambda}\to\Phi_{\mathcal Q_A}$ over $\Lambda\to A$
  (Definition~\ref{gal-realization-thickening-torsor}(4)). For $a\in I^{[i]}$ and $r\in\mathbf Z$, $a\rho_T$ carries
  $\mathcal O^{[r-i]}(G(T))$ into $\mathcal O^{[r]}(T)$, so multiplication by $a$ maps
  $H^j(\mathcal Q_\Lambda,\mathcal O^{[r-i]}/\mathcal O^{[r'-i]})$ to $H^j(\mathcal Q_A,\mathcal O^{[r]}/\mathcal O^{[r']})$.
\item For integers $r\leq r'\leq s$ and $j\geq0$, the map
  \[
   \bigoplus_{i<m}N_i\otimes_\Lambda H^j(\mathcal Q_\Lambda,\mathcal O^{[r-i]}/\mathcal O^{[r'-i]})\to
   H^j(\mathcal Q_A,\mathcal O^{[r]}/\mathcal O^{[r']})
  \]
  that sends $n\otimes x$, for $n\in N_i$, to the image of $x$ under multiplication by $\nu_i(n)$, is bijective.
\end{enumerate}
\end{lemma}
\leanref{Development/WeilBound/CoreStatements/PadicHodgeComparison/EtaleCohomologyComparison/PerfectoidPullbacks/SplittingFieldLattices/BasisB.lean}{CrystallineSemistableComparison.Pb.BcOn}
\leanref{Development/WeilBound/CoreStatements/PadicHodgeComparison/EtaleCohomologyComparison/PerfectoidPullbacks/SplittingFieldLattices/BcOnFalse.lean}{CrystallineSemistableComparison.Pb.bcOn_false}
\leanref{Development/AutomorphicSpectrum/Proofs/BlockDecomposition/ScalarMultiplicationShiftsFiltrationStatement.lean}{GlobalTheory.T.scalar_multiplication_shifts_filtration}

\begin{proof}
(1) For $T=(B,i,J,\delta,y)$ in $\mathcal Q_A$, $G(T)=(B,i\circ(\Lambda\to A),J,\delta,y)$ is a thickening over
$(\Lambda,0,0,0)$, with the same ideals $J^{[r]}$, so it lies in $\mathcal Q_\Lambda$; the conditions of a morphism hold
as the maps of the rings are identities. For $a\in I^{[i]}$, $i(a)\in J^{[i]}$, as $i$ is a divided power map, and
$J^{[i]}J^{[r-i]}\subset J^{[r]}$, also when $r-i\leq0$, as then $i\geq r$.

(2) Write $F_i=\mathcal O^{[r-i]}/\mathcal O^{[r'-i]}$ and $F=\mathcal O^{[r]}/\mathcal O^{[r']}$, and say that \emph{base
change holds} for full subcategories $\mathcal W$ of $\mathrm{Th}^{(s)}_\ell(X/\Lambda)$ and $\mathcal W'$ of
$\mathrm{Th}^{(s)}_\ell(X/A)$ with $G(\mathcal W')\subset\mathcal W$ if the map of (2), formed for $\mathcal W$ and
$\mathcal W'$, is bijective for every $j$. For $n\in N_i$ let $\tau_i(n)_T:F_i(G(T))\to F(T)$ be $x\mapsto\nu_i(n)x$; these
maps are semilinear for $\Lambda\to A$, natural in $T$, additive in $n$, and $\tau_i(an)=a\tau_i(n)$ for $a\in\Lambda$.

\emph{Step 1: a coordinate cover.} Let $(U_\alpha)_{\alpha\in E}$ be a coordinate cover of $X$ and $\mathcal U=\{U_\alpha:
\alpha\in E\}$. Every thickening lies over $X$ (Lemma~\ref{gal-realization-thickening-over}(2)). By
Lemma~\ref{gal-realization-kappa-truncated-periods-exist-base-change-zariski} and
Lemma~\ref{gal-realization-thickening-zariski}(1), with $\mathcal W=\{X\}$ and $\mathcal V$ or $\mathcal U$ in place of
$\mathcal V$, restriction from $\mathrm{Th}^{(s)}_\ell(X/\Lambda)$ to its full subcategories over $\mathcal V$ and over
$\mathcal U$ is bijective on $H^j(\cdot,F_i)$, and likewise over $A$ on $H^j(\cdot,F)$. These restrictions commute with
multiplication by $\nu_i(n)$, as both composites are induced by the same functor with the same maps. So base change for
$(\mathcal Q_\Lambda,\mathcal Q_A)$ is equivalent to base change for the categories over $\{X\}$, which is equivalent to
base change for the categories over $\mathcal U$.

\emph{Step 2: the intersections.} By Lemma~\ref{gal-realization-thickening-over}, the thickenings over $U_\alpha$ form
cosieves of the categories over $\mathcal U$ that cover them, $G$ maps the one over $A$ into the one over $\Lambda$, and
for nonempty $S\subset E$ their intersections are the categories over $U_S=\bigcap_{\alpha\in S}U_\alpha$. By
Lemma~\ref{gal-realization-thickening-cosieves}(1), for the flat $\Lambda$-modules $N_i$ and the maps $\tau_i(n)$, it
suffices that base change holds for $\mathcal C=\mathrm{Th}^{(s)}_\ell(X/\Lambda)_{U_S}$ and $\mathcal
C'=\mathrm{Th}^{(s)}_\ell(X/A)_{U_S}$ for every nonempty $S$.

\emph{Step 3: one open with coordinates.} Fix $S$ and let $V=U_S$, an affine open with coordinates $t_1,\ldots,t_d$
restricted from some $U_\alpha$ (Lemma~\ref{gal-realization-thickening-cover}(1)), and $\Gamma=\Gamma(V,\mathcal O_X)$. By
Lemma~\ref{gal-realization-kappa-truncated-periods-exist-base-change-chart}(1), for $(\Lambda,0)$ and for $(A,I)$, there are
divided power charts $D$ in $\mathcal C$ and $D'$ in $\mathcal C'$ on $V$, with $B_D=\Lambda\otimes_\Lambda\Gamma$ and
$B_{D'}=A\otimes_\Lambda\Gamma$, coordinates the $1\otimes t_k$, and points corresponding to $f\mapsto1\otimes f$. The ring
map $e:\Lambda\otimes_\Lambda\Gamma\to A\otimes_\Lambda\Gamma$, $a\otimes f\mapsto a\otimes f$, is a map $D\to G(D')$: it
commutes with the maps from $\Lambda$, it is a divided power map as $J_D=0$, and it is compatible with the points, as
both correspond to $f\mapsto1\otimes f$. It sends $1\otimes t_k$ to $1\otimes t_k$. By
Lemma~\ref{gal-realization-thickening-chart-torsor}(1) and~(2), $D$ and $D'$ give torsor data of dimension $d$ for
$\Phi_{\mathcal C}$ and $\Phi_{\mathcal C'}$, compatible along $e$ for the morphism of (1), whose maps $\rho_T$ are
injective. By Lemma~\ref{gal-realization-kappa-truncated-periods-exist-base-change-chart}(2), $\Phi_{\mathcal C}$ and
$\Phi_{\mathcal C'}$ have divided power polynomial structures that respect $(\mathcal O^{[k]})_k$ in every degree $k\leq
s$.

Write $r=r'-c$ with $c\geq0$, and apply Lemma~\ref{gal-realization-thickening-torsor-base-change} to the modules $N_i$ and
maps $\nu_i$, the sequences $\mathfrak a_{i,t}=\mathcal O^{[t-c-i]}$ and $\mathfrak b_{i,t}=\mathcal O^{[t-i]}$ of
$\Phi_{\mathcal C}$ and $\mathfrak a'_t=\mathcal O^{[t-c]}$ and $\mathfrak b'_t=\mathcal O^{[t]}$ of $\Phi_{\mathcal C'}$,
and the integer $r'$; for $t=r'$ they give $F_i$ and $F$. The polynomial structures are ladders up to $r'$ for these
sequences, as $\mathcal O^{[t-i]}\subset\mathcal O^{[t-c-i]}$ and all the degrees involved are at most $r'\leq s$; the
conditions on multiplication hold by (1), as $\nu_i(n)\in I^{[i]}$. It remains to check the hypothesis at the charts: for
$t\leq r'$, the map $\bigoplus_iN_i\otimes_\Lambda(\mathcal O^{[t-c-i]}/\mathcal O^{[t-i]})(D)\to(\mathcal
O^{[t-c]}/\mathcal O^{[t]})(D')$, $n\otimes x\mapsto\nu_i(n)e(x)$, is bijective.

\emph{Step 4: the charts.} Let $u=\max(t-c,0)$ and $v=\max(t,0)$. As $J_D=0$, $J_D^{[k]}$ is $0$ for $k\geq1$ and $B_D$
for $k\leq0$; so $(\mathcal O^{[t-c-i]}/\mathcal O^{[t-i]})(D)$ is $B_D$, identified with $\Gamma$, if $u\leq i<v$, and $0$
otherwise. By Lemma~\ref{gal-realization-kappa-truncated-periods-exist-base-change-splitting-field}(2), $(n_i)\mapsto
\sum_i\nu_i(n_i)$ is a bijection $\bigoplus_iN_i\to A$ that maps $\bigoplus_{i\geq k}N_i$ onto $I^{[k]}$. Tensoring with
$\Gamma$ gives a bijection $g:\bigoplus_iN_i\otimes_\Lambda\Gamma\to A\otimes_\Lambda\Gamma$, $n\otimes f\mapsto\nu_i(n)
\otimes f$, which maps $\bigoplus_{i\geq k}N_i\otimes_\Lambda\Gamma$ onto the image of $I^{[k]}\otimes_\Lambda\Gamma$; and
this image is $J_{D'}^{[k]}=i_{D'}(I^{[k]})B_{D'}$ (Lemma~\ref{gal-realization-thickening-chart}(2)). So an element $g(w)$
lies in $J_{D'}^{[k]}$ if and only if $w_i=0$ for $i<k$, and $g$ induces a bijection from $\bigoplus_{u\leq
i<v}N_i\otimes_\Lambda\Gamma$ onto $J_{D'}^{[u]}/J_{D'}^{[v]}=(\mathcal O^{[t-c]}/\mathcal O^{[t]})(D')$. This is the map
of Step 3, as $\nu_i(n)e(1\otimes f)=\nu_i(n)\otimes f$.
% Lean: Pb.bcOn_false = bc_glob glob_bc_st bcOn_true. (1) is scalar_multiplication_shifts_filtration (LecShift.leC_shift)
% and GlobBc.thM/smO. Step 1 is GlobBc.bcAt_iff (cech_refinement_cohomology_bijective for TF l s over O and over A, with
% comm_top and tens_transport); Step 2 is glob_bc_st = GlobalTheory.T.stable_cover_cech_vanishing (Lemma
% gal-realization-thickening-cosieves(1)), with tM_add and tM_smul; Step 3 is bcOn_true = BcLoc.bc_single: the charts are
% exists_chart_from_coordinates, e is chart_tensor_base_compat (BcMor.chart_mor), the frames are Chart.frame with
% chart_is_good, the polynomial structures are poly_tf, and loc_bc is LocBC.local_base_change (Lemma
% gal-realization-thickening-torsor-base-change). Step 4 is BcLoc.chart_tens with graded_split_range_iso_of_interval
% (GrPiece.main), through graded_splitting_tensor_filtration (GrTens.main) for g; dividedPowerFiltration_dividedPowersBot'
% gives J_D^[k] = 0 for k >= 1 and chart_divided_powers_graded_constant gives J_{D'}^[k].
\end{proof}
\leanref{Development/AutomorphicSpectrum/Proofs/BlockDecomposition/TensTransportSymm.lean}{GlobalTheory.CompanionTheory.GlobBc.main}
\leanref{Development/AutomorphicSpectrum/Proofs/BlockDecomposition/TensTransportSymm.lean}{GlobalTheory.CompanionTheory.GlobBc.bcAt_iff}
\leanref{Development/AutomorphicSpectrum/Proofs/BlockDecomposition/PdPBotPos.lean}{GlobalTheory.CompanionTheory.BcLoc.bc_single}
\leanref{Development/AutomorphicSpectrum/Proofs/BlockDecomposition/PdPBotPos.lean}{GlobalTheory.CompanionTheory.BcLoc.chart_tens}
\leanref{Development/AutomorphicSpectrum/Proofs/BlockDecomposition/GradedSplitRangeIsoOfInterval.lean}{GlobalTheory.CompanionTheory.GrPiece.main}
\leanref{Development/AutomorphicSpectrum/Proofs/BlockDecomposition/ChartTensorBaseCompat.lean}{GlobalTheory.CompanionTheory.BcMor.chart_mor}

\begin{lemma}\label{gal-realization-kappa-truncated-periods-exist-base-change-filtration}
Let $K$ be a field of characteristic $0$, $X$ a $K$-scheme with a coordinate cover $(U_\alpha)$ of dimensions $d_\alpha$
(Definition~\ref{gal-realization-thickening-coordinates}(2)), $s\geq1$, $\mathcal V$ a set of open subsets of $X$ with
union $X$, and $\mathcal Q=\mathrm{Th}^{(s)}_{\mathrm{dR}}(X)_{\mathcal V}$
(Definition~\ref{gal-realization-kappa-de-rham-truncated}). For $0\leq u\leq s$ and $j\geq0$ let $q_u:H^j_{\mathrm{dR}}(X)
\to H^j(\mathcal Q,\mathcal O/\mathcal O^{[u]})$ be restriction to $\mathcal Q$ followed by the map induced by the
projection $\mathcal O\to\mathcal O/\mathcal O^{[u]}$.
\begin{enumerate}
\item If $u\geq1$ and $d_\alpha<u$ for every $\alpha$, then $q_u$ is bijective.
\item Suppose that $d_\alpha<s$ for every $\alpha$. A class $x\in H^j_{\mathrm{dR}}(X)$ has $q_u(x)=0$ if and only if it
  is the class of a cocycle of $\mathrm{Th}_{\mathrm{dR}}(X)$ whose value on every chain $T_0\to\cdots\to T_j$ lies in
  $J_{T_j}^{[u]}$.
\end{enumerate}
\end{lemma}
\leanref{Development/AutomorphicSpectrum/Proofs/BlockDecomposition/QuotientFiltrationBijectiveOfCoverStatement.lean}{GlobalTheory.T.quotient_filtration_bijective_of_cover}
\leanref{Development/WeilBound/CoreStatements/PadicHodgeComparison/EtaleCohomologyComparison/PerfectoidPullbacks/GlobalSerreRelations/GlobSerreSt.lean}{CrystallineSemistableComparison.Pb.KH.KIso}

\begin{proof}
Write $\mathrm{Th}=\mathrm{Th}_{\mathrm{dR}}(X)=\mathrm{Th}(X/K)$ and $F=\mathcal O/\mathcal O^{[u]}$.

(2) By Lemma~\ref{gal-realization-kappa-de-rham-restriction}, for $\ell=1$, $\mathcal Q_0=\mathrm{Th}$, $\{X\}$ in place of
$\mathcal V$, $r=0$ and $r'=u\leq s$, restriction $H^j(\mathrm{Th},F)\to H^j(\mathrm{Th}^{(s)}_{\mathrm{dR}}(X)_{\{X\}},F)$ is
bijective, and by Lemma~\ref{gal-realization-kappa-de-rham-zariski} with Lemma~\ref{gal-realization-thickening-zariski}(1),
for $\mathcal Q_0=\mathrm{Th}^{(s)}_{\mathrm{dR}}(X)$, $\mathcal V$ and $\mathcal W=\{X\}$, so is restriction to $\mathcal
Q$; every thickening lies over $X$ (Lemma~\ref{gal-realization-thickening-over}(2)). Restriction commutes with the maps
induced by $\mathcal O\to F$. So $q_u(x)=0$ if and only if the image of $x$ in $H^j(\mathrm{Th},F)$ is $0$. If $x$ is the
class of a cocycle $z$ with values in $\mathcal O^{[u]}$, the image of $z$ is $0$. Conversely, let $z$ be a cocycle with
class $x$ whose image is $0$: if $j=0$, $z$ has values in $\mathcal O^{[u]}$; if $j\geq1$, the image of $z$ is the
coboundary of a cochain $\bar y$ with values in $F$, and for a cochain $y$ lifting $\bar y$, $z-dy$ has values in $\mathcal
O^{[u]}$ and class $x$.

(1) By Lemma~\ref{gal-realization-kappa-de-rham-truncation}, with $u$ in place of $s$, restriction $H^j_{\mathrm{dR}}(X)
\to H^j(\mathrm{Th}^{(u)}_{\mathrm{dR}}(X)_{\mathcal V},\mathcal O)$ is bijective, and on
$\mathrm{Th}^{(u)}_{\mathrm{dR}}(X)_{\mathcal V}$ the projection $\mathcal O\to F$ is an isomorphism of functors, as
$\mathcal O^{[u]}=0$ there. By Lemma~\ref{gal-realization-kappa-de-rham-restriction}, with $\ell=1$, $u$ in place of $s$,
$\mathcal Q_0=\mathrm{Th}^{(s)}_{\mathrm{dR}}(X)$, $r=0$ and $r'=u$, restriction $H^j(\mathcal Q,F)\to
H^j(\mathrm{Th}^{(u)}_{\mathrm{dR}}(X)_{\mathcal V},F)$ is bijective. Its composite with $q_u$ is restriction followed by
$\mathcal O\to F$ on $\mathrm{Th}^{(u)}_{\mathrm{dR}}(X)_{\mathcal V}$, which is bijective. So $q_u$ is bijective.
% Lean: (1) is QuoK.quoK_bijective: quoK_bij_self with the first part of KIso for u (Lemma
% gal-realization-kappa-de-rham-truncation), filtered_cohomology_restriction_bijective (Lemma
% gal-realization-kappa-de-rham-restriction) and quoK_comm. (2) is the second part of KIso (KIso.k_iso), through
% chain_bij, chain_comm and chain_zero (the bijective restrictions) and fil_iff (pull_eq_zero_iff), with FilR q j u the set
% of the classes of cocycles with values in O^[u].
\end{proof}
\leanref{Development/AutomorphicSpectrum/Proofs/BlockDecomposition/QuotientFiltrationBijectiveOfCover.lean}{GlobalTheory.CompanionTheory.QuoK.quoK_bijective}
\leanref{Development/AutomorphicSpectrum/Proofs/BlockDecomposition/HodgeFiltrationViaTruncation.lean}{GlobalTheory.CompanionTheory.KIso.fil_iff}

\begin{lemma}\label{gal-realization-kappa-truncated-periods-exist-base-change}
Let $K$ be a field of characteristic $0$, $B$ a $K$-algebra with divided powers $\delta$ on an ideal $J$, and $1\leq
m\leq s$ integers with $J^{[m]}=0$ (Definition~\ref{gal-realization-kappa-approximations}). Let $X$ be a $K$-scheme with a
coordinate cover $(U_\alpha)$ of dimensions $d_\alpha$ (Definition~\ref{gal-realization-thickening-coordinates}(2)) such
that $d_\alpha+m\leq s$ for every $\alpha$, let $\mathcal V$ be a set of open subsets of $X$ with union $X$, and let
$\mathcal Q=\mathrm{Th}^{(s)}(X/B)_{\mathcal V}$, formed for $(B,J,\delta,0)$, $\mathcal O=K$ and the structure map $K\to
B/J$ (Definition~\ref{gal-realization-kappa-de-rham-truncated}). For $j\geq0$ let $\mu:H^j_{\mathrm{dR}}(X)\to
H^j(\mathcal Q,\mathcal O)$ be the restriction to $\mathrm{Th}^{(s)}_{\mathrm{dR}}(X)_{\mathcal V}$ followed by the map
induced by restriction of scalars along $K\to B$
(Definition~\ref{gal-realization-thickening-nilpotent}(4)), with identity maps, and let $(v_t)$ be a finite basis of
$H^j_{\mathrm{dR}}(X)$.
\begin{enumerate}
\item Every $w\in H^j(\mathcal Q,\mathcal O)$ is $\sum_tg_t\,\mu(v_t)$ for exactly one family $(g_t)$ in $B$.
\item Let $0\leq r\leq s$, and let $(g_t)$ be a family in $B$ such that the image of $\sum_tg_t\,\mu(v_t)$ in
  $H^j(\mathcal Q,\mathcal O/\mathcal O^{[r]})$ is $0$. Let $c\geq0$, and let $\psi\in\operatorname{Hom}_K(
  H^j_{\mathrm{dR}}(X),K)$ vanish on the classes of the cocycles of $\mathrm{Th}_{\mathrm{dR}}(X)$ whose value on every
  chain $T_0\to\cdots\to T_j$ lies in $J_{T_j}^{[c+1]}$. Then $\sum_tg_t\psi(v_t)\in J^{[r-c]}$, where $J^{[k]}=B$ for
  $k\leq0$.
\end{enumerate}
\end{lemma}
\leanref{Development/WeilBound/CoreStatements/PadicHodgeComparison/EtaleCohomologyComparison/PerfectoidPullbacks/SplittingFieldLattices/BasisB.lean}{CrystallineSemistableComparison.Pb.BasisB}
\leanref{Development/WeilBound/CoreStatements/PadicHodgeComparison/EtaleCohomologyComparison/PerfectoidPullbacks/SplittingFieldLattices/BcOnFalse.lean}{CrystallineSemistableComparison.Pb.basisB}

\begin{proof}
Write $\mathcal Q_K=\mathrm{Th}^{(s)}_{\mathrm{dR}}(X)_{\mathcal V}$, which is $\mathrm{Th}^{(s)}_1(X/K)_{\mathcal V}$ for
$(K,0,0,0)$ (Definition~\ref{gal-realization-kappa-de-rham-truncated}(2)), and note $\mathcal
Q=\mathrm{Th}^{(s)}_1(X/B)_{\mathcal V}$ (Lemma~\ref{gal-realization-kappa-truncated-periods-exist-base-change-zariski}). Let $E$ be the restriction of scalars
along $K\to B$, which maps $\mathcal Q$ into $\mathcal Q_K$, so that $\mu(x)=E^*(x|_{\mathcal Q_K})$. The hypotheses of
Lemma~\ref{gal-realization-kappa-truncated-periods-exist-base-change-tensor} hold for $\Lambda=K$, $A=B$, $\ell=1$, $s$ and
$\mathcal V$: $0^{[s]}=0$ and $J^{[s]}\subset J^{[m]}=0$, the conditions on $\ell=1$ are empty, and every integer
$n\geq1$ is invertible in $K$ and in $B$, as $K$ has characteristic $0$. By
Lemma~\ref{gal-realization-kappa-truncated-periods-exist-base-change-splitting-field}(1), $(B,\delta)$ has a graded
splitting $(N_i,\nu_i)_{i<m}$ over $K$.

For $r\leq s$, $i<m$ and $x\in H^j_{\mathrm{dR}}(X)$ let $e^r_i(x)$ be the image of $x|_{\mathcal Q_K}$ in $H^j(\mathcal
Q_K,\mathcal O/\mathcal O^{[r-i]})$; for $r-i\geq0$, $e^r_i$ is the map $q_{r-i}$ of
Lemma~\ref{gal-realization-kappa-truncated-periods-exist-base-change-filtration} for $\mathcal V$. As $\mathcal
O^{[-i]}=\mathcal O$, Lemma~\ref{gal-realization-kappa-truncated-periods-exist-base-change-tensor}(2) for $(0,r)$ gives a
bijection $\Phi_r:\bigoplus_iN_i\otimes_KH^j(\mathcal Q_K,\mathcal O/\mathcal O^{[r-i]})\to H^j(\mathcal Q,\mathcal
O/\mathcal O^{[r]})$. Let $\pi_r:H^j(\mathcal Q,\mathcal O)\to H^j(\mathcal Q,\mathcal O/\mathcal O^{[r]})$ be induced by the
projection. For $n\in N_i$,
\[
 \Phi_r(n\otimes e^r_i(x))=\pi_r(\nu_i(n)\mu(x)),\tag{$*$}
\]
as both sides are induced by the functor $E$, followed by restriction, with the maps $b\mapsto\nu_i(n)b$ followed by the
projections on coefficients.

(1) For $i<m$, $d_\alpha<s-i$, as $d_\alpha+m\leq s$; so $e^s_i$ is bijective
(Lemma~\ref{gal-realization-kappa-truncated-periods-exist-base-change-filtration}(1)), and the $e^s_i(v_t)$ form a basis of
$H^j(\mathcal Q_K,\mathcal O/\mathcal O^{[s-i]})$. As $J_T^{[s]}=0$ for $T$ in $\mathcal Q$, $\pi_s$ is bijective. Let
$(g_t)$ be a family in $B$ with $\sum_tg_t\mu(v_t)=0$, and write $g_t=\sum_i\nu_i(n_{t,i})$
(Lemma~\ref{gal-realization-kappa-truncated-periods-exist-base-change-splitting-field}(2)). By $(*)$, $\Phi_s$ maps the
family $(\sum_tn_{t,i}\otimes e^s_i(v_t))_i$ to $\pi_s(\sum_tg_t\mu(v_t))=0$, so $\sum_tn_{t,i}\otimes e^s_i(v_t)=0$ for
every $i$, and $n_{t,i}=0$ as the $e^s_i(v_t)$ form a basis. So $g_t=0$, which gives uniqueness. Let $w\in H^j(\mathcal
Q,\mathcal O)$. The components of $\Phi_s^{-1}(\pi_s(w))$ are $\sum_tn_{i,t}\otimes e^s_i(v_t)$ for some $n_{i,t}\in N_i$;
with $g_t=\sum_i\nu_i(n_{i,t})$, $(*)$ gives $\pi_s(\sum_tg_t\mu(v_t))=\pi_s(w)$, so $w=\sum_tg_t\mu(v_t)$.

(2) Let $r$, $(g_t)$, $c$ and $\psi$ be as in (2), and write $g_t=\sum_i\nu_i(n_{t,i})$. As in (1), $(*)$ and the
injectivity of $\Phi_r$ give $\sum_tn_{t,i}\otimes e^r_i(v_t)=0$ for every $i$. As $\nu_i$ is $K$-linear,
$\sum_tg_t\psi(v_t)=\sum_i\nu_i(\sum_t\psi(v_t)n_{t,i})$. For $i\geq r-c$ the $i$-th term lies in $J^{[i]}\subset
J^{[r-c]}$. Let $i<r-c$, so that $1\leq c+1\leq r-i\leq s$. If $e^r_i(x)=0$, then $x$ is the class of a cocycle with values
in $\mathcal O^{[r-i]}\subset\mathcal O^{[c+1]}$
(Lemma~\ref{gal-realization-kappa-truncated-periods-exist-base-change-filtration}(2)), so $\psi(x)=0$. Hence
$\psi=\psi'\circ e^r_i$ for a linear form $\psi'$ on $H^j(\mathcal Q_K,\mathcal O/\mathcal O^{[r-i]})$, and applying
$1\otimes\psi'$ to $\sum_tn_{t,i}\otimes e^r_i(v_t)=0$ gives $\sum_t\psi(v_t)n_{t,i}=0$. So $\sum_tg_t\psi(v_t)\in
J^{[r-c]}$.
% Lean: Pb.basisB = basis_b bcOn_false @k_iso = BasisBase.main, for BasisB stated for every set 𝒱 of opens with union X,
% with TS γ q 1 s 𝒱 = Th^(s)(X/B)_𝒱. graded_splitting_over_field gives S, tfbase_K and tfbase_B
% the hypotheses, scalar_multiplication_shifts_filtration the multiplications, and BcOn False (Lemma
% gal-realization-kappa-truncated-periods-exist-base-change-tensor) the bijections T r. (1) is core_basis, with zK
% (the maps e^r_i, through zK_transport for O^[-i] = O), quotient_filtration_bijective_of_cover (Lemma
% gal-realization-kappa-truncated-periods-exist-base-change-filtration(1)), quoM_bij for pi_s, tens_span and tens_dual;
% (2) is core_filt with the second part of KIso (Lemma ...-filtration(2)) and tens_dual. (*) is
% smooth_pull_basis_filtration (sqC_zK). The Lean also allows m = 0, where B = 0.
\end{proof}
\leanref{Development/AutomorphicSpectrum/Proofs/BlockDecomposition/BasisBOfBcOn.lean}{GlobalTheory.CompanionTheory.BasisBase.main}
\leanref{Development/AutomorphicSpectrum/Proofs/BlockDecomposition/BasisBOfBcOn.lean}{GlobalTheory.CompanionTheory.BasisBase.core_basis}
\leanref{Development/AutomorphicSpectrum/Proofs/BlockDecomposition/BasisBOfBcOn.lean}{GlobalTheory.CompanionTheory.BasisBase.core_filt}
\leanref{Development/AutomorphicSpectrum/Proofs/BlockDecomposition/SmoothPullBasisFiltrationStatement.lean}{GlobalTheory.T.smooth_pull_basis_filtration}

\begin{lemma}\label{gal-realization-kappa-truncated-periods-exist-lifting-limits}
Let $R$ be a ring, $c\in R$, and $K^\bullet$ a complex of $R$-modules on whose terms multiplication by $c$ is injective. For
$n\geq0$ let $\mathrm{red}_n:K^\bullet\to K^\bullet_n$ be surjective maps of complexes whose kernels are the $c^nK^\bullet$,
and $t_n:K^\bullet_{n+1}\to K^\bullet_n$ maps of complexes with $t_n\circ\mathrm{red}_{n+1}=\mathrm{red}_n$. Let
$j,k\geq0$, and suppose:
\begin{enumerate}
\item[(a)] every $x\in H^{j+1}(K^\bullet)$ with $c^nx=0$ for some $n$ has $c^kx=0$;
\item[(b)] every $x\in H^j(K^\bullet)$ that lies in $c^nH^j(K^\bullet)$ for every $n$ has $c^kx=0$;
\item[(c)] for every sequence $(x_n)_{n\geq0}$ in $H^j(K^\bullet)$ with $x_{n+1}-x_n\in c^nH^j(K^\bullet)$ for every $n$
  there is $z\in H^j(K^\bullet)$ with $z-c^kx_n\in c^nH^j(K^\bullet)$ for every $n$.
\end{enumerate}
Then:
\begin{enumerate}
\item for every family $(x_n)_{n\geq0}$ with $x_n\in H^j(K^\bullet_n)$ and $t_n(x_{n+1})=x_n$ for every $n$ there is $z\in
  H^j(K^\bullet)$ with $\mathrm{red}_n(z)=c^{2k}x_n$ for every $n$;
\item if $z\in H^j(K^\bullet)$ and $\mathrm{red}_n(z)=0$ for every $n$, then $c^kz=0$.
\end{enumerate}
\end{lemma}
\leanref{Development/AutomorphicSpectrum/Proofs/BlockDecomposition/InverseLimitCohomologyComparisonStatement.lean}{GlobalTheory.T.inverse_limit_cohomology_comparison}
\leanref{Development/AutomorphicSpectrum/Proofs/BlockDecomposition/InverseLimitCohomologyComparison.lean}{GlobalTheory.CompanionTheory.LimLift.lim_lift}

\begin{proof}
Write $H^j=H^j(K^\bullet)$ and $d$ for the differentials. Multiplication by $c^n$ is injective on the terms of $K^\bullet$.

\emph{Step 1: if $\mathrm{red}_n(z)=0$, then $z\in c^nH^j$.} Let $a$ be a cocycle with class $z$. If $j=0$, then
$\mathrm{red}_n(a)=0$; put $v=0$. If $j\geq1$, then $\mathrm{red}_n(a)=d\bar v$ for a cochain $\bar v$; lift it to $v$. In
both cases $\mathrm{red}_n(a-dv)=0$, so $a-dv=c^ny$. Then $c^ndy=d(a-dv)=0$, so $dy=0$, and $z$ is $c^n$ times the class
of $y$.

\emph{Step 2: $c^n$ kills $H^j(K^\bullet_n)$.} A class is the class of $\mathrm{red}_n(a)$ for a cochain $a$, and
$c^n\mathrm{red}_n(a)=\mathrm{red}_n(c^na)=0$.

\emph{Step 3: for $x\in H^j(K^\bullet_n)$ there is $y\in H^j$ with $\mathrm{red}_n(y)=c^kx$.} Let $x$ be the class of
$\mathrm{red}_n(a)$. Then $\mathrm{red}_n(da)=0$, so $da=c^nb$, and $c^ndb=dda=0$, so $db=0$. The class of $b$ in
$H^{j+1}(K^\bullet)$ is killed by $c^n$, so by (a) $c^kb=de$ for a cochain $e$. Then $c^ka-c^ne$ is a cocycle, as
$d(c^ka-c^ne)=c^{k+n}b-c^nc^kb=0$, and $\mathrm{red}_n(c^ka-c^ne)=c^k\mathrm{red}_n(a)$.

(2) By Step 1, $z\in c^nH^j$ for every $n$, so $c^kz=0$ by (b).

(1) By Step 3 choose $y_n\in H^j$ with $\mathrm{red}_n(y_n)=c^kx_n$. Then $\mathrm{red}_n(y_{n+1}-y_n)=
t_n(c^kx_{n+1})-c^kx_n=0$, so $y_{n+1}-y_n\in c^nH^j$ by Step 1, and (c) gives $z$ with $z-c^ky_n=c^nu_n$, $u_n\in H^j$.
By Step 2, $\mathrm{red}_n(z)=c^k\mathrm{red}_n(y_n)+c^n\mathrm{red}_n(u_n)=c^{2k}x_n$.
% Lean: LimLift.lim_lift; Step 1 is ker_red, Step 2 kill_H, Step 3 lift_level, and t_n(red_{n+1}(y)) = red_n(y) on
% cohomology is mapH_tr.
\end{proof}
\leanref{Development/AutomorphicSpectrum/Proofs/BlockDecomposition/InverseLimitCohomologyComparison.lean}{GlobalTheory.CompanionTheory.LimLift.ker_red}
\leanref{Development/AutomorphicSpectrum/Proofs/BlockDecomposition/InverseLimitCohomologyComparison.lean}{GlobalTheory.CompanionTheory.LimLift.lift_level}

\begin{lemma}\label{gal-realization-kappa-truncated-periods-exist-lifting-modules}
Let $A$ be a ring, $I$ a finitely generated prime ideal of $A$ with $I^N=0$ for some $N\geq0$, and $c\in A$, such that: for
all $x,y\in A$, $x\in yA+I$ or $y\in xA+I$; for $x\notin I$ some power of $c$ lies in $xA+I$; every $x\in A$ that lies in
$c^nA+I$ for every $n$ lies in $I$; and for every sequence $(a_n)_{n\geq0}$ in $A$ with $a_{n'}-a_n\in c^nA$ for $n\leq n'$
there is $a\in A$ with $a-a_n\in c^nA$ for every $n$. The condition that $x\in yA+I$ or $y\in xA+I$ is not used. Let $M$ be
an $A$-module with a finitely generated submodule $M_0$ and an integer $e\geq0$ such that $c^eM\subset M_0$.
There is $k\geq0$ with the following properties.
\begin{enumerate}
\item If $x\in M$ and $c^nx=0$ for some $n$, then $c^kx=0$.
\item If $x\in M$ lies in $c^nM$ for every $n$, then $c^kx=0$.
\item For every sequence $(x_n)_{n\geq0}$ in $M$ with $x_{n+1}-x_n\in c^nM$ for every $n$ there is $z\in M$ with
  $z-c^kx_n\in c^nM$ for every $n$.
\end{enumerate}
\end{lemma}
\leanref{Development/AutomorphicSpectrum/Proofs/BlockDecomposition/AlmostFiniteModuleBoundedTorsionStatement.lean}{GlobalTheory.T.almost_finite_module_bounded_torsion}
\leanref{Development/AutomorphicSpectrum/Proofs/BlockDecomposition/AlmostFiniteModuleBoundedTorsion.lean}{GlobalTheory.CompanionTheory.AlmostFinite.afg_separated_complete}

\begin{proof}
Call an $A$-module \emph{bounded} if some $k$ has properties (1) and (2) for it.

(a) Let $M'\subset M$ be a submodule such that $M'$ is bounded with $k'$ and $M/M'$ with $k''$. If $c^ny\in M'$ for some
$n$, then $c^{k''}y\in M'$, by (1) for $M/M'$. Let $c^nx=0$. Then $c^{k''}x\in M'$ and $c^n(c^{k''}x)=0$, so
$c^{k'+k''}x=0$. Let $x\in c^nM$ for every $n$. By (2) for $M/M'$, $c^{k''}x\in M'$. For every $n$ write $x=c^{n+k''}y$;
as $c^{k''+n+k''}y=c^{k''}x\in M'$, $c^{k''}y\in M'$, so $x\in c^nM'$. In particular $x\in M'$, and $c^{k'}x=0$ by (2)
for $M'$. So $M$ is bounded with $k'+k''$.

(b) If $M_0\subset M$ is bounded with $k$ and $c^eM\subset M_0$, then $M$ is bounded with $k+e$: apply (1) and (2) for
$M_0$ to $c^ex$, which lies in $c^nM_0$ if $x$ lies in $c^{n}M$.

(c) Let $c\notin I$, and $M=Am$ a module generated by one element and killed by $I$. If every $t$ with $tm=0$ lies in
$I$, then $k=0$ works: if $c^nam=0$, then $c^na\in I$, so $a\in I$, as $I$ is prime and $c\notin I$, and $am=0$; if
$am\in c^nM$ for every $n$, then $a-c^nb_n$ kills $m$ for some $b_n$, so $a\in c^nA+I$ for every $n$, $a\in I$ and
$am=0$. Otherwise some $t\notin I$ kills $m$, and $c^n=ut+i$ with $u\in A$ and $i\in I$ for some $n$; then $c^nM=0$, and
$k=n$ works.

(d) Let $c\notin I$. A finitely generated module killed by $I$ is bounded, by induction on the number of generators: the
submodule generated by the first one is bounded by (c), and the quotient by it is killed by $I$ and has fewer generators;
(a) applies.

(e) Let $c\notin I$. A finitely generated module $M$ killed by $I^{N'}$ is bounded, by induction on $N'$: for $N'=0$,
$M=0$; otherwise $IM$ is finitely generated, as $I$ and $M$ are, and killed by $I^{N'-1}$, and $M/IM$ is finitely
generated and killed by $I$; (a) and (d) apply.

(1) and (2). If $c\in I$, then $c^N=0$, and $M$ is bounded with $k_0=N$. Otherwise $M_0$ is bounded by (e), as $I^N=0$,
and so is $M$, by (b), with some $k_0$.

(3) Let $g_1,\ldots,g_p$ generate $M_0$, and write $x_{n+1}-x_n=c^ny_n$ and $c^ey_n=\sum_ia_{n,i}g_i$ with $a_{n,i}\in A$.
Let $S_i(n)=\sum_{n'<n}c^{n'}a_{n',i}$. Then $S_i(n')-S_i(n)\in c^nA$ for $n\leq n'$, so there is $s_i\in A$ with
$s_i-S_i(n)\in c^nA$ for every $n$. By induction on $n$, $c^ex_n=c^ex_0+\sum_iS_i(n)g_i$, so $z_0=c^ex_0+\sum_is_ig_i$ has
$z_0-c^ex_n=\sum_i(s_i-S_i(n))g_i\in c^nM$. With $k=k_0+e$, which also has (1) and (2), $z=c^{k_0}z_0$ has $z-c^kx_n\in
c^nM$ for every $n$.
% Lean: AlmostFinite.afg_separated_complete (BS is "bounded"); (a) is BS.ext, (b) BS.of_submodule, (c) BS.cyclic, (d)
% BS.of_gens, (e) BS.of_finite, the case c in I is BS.of_pow, and (3) is limit_of_cauchy. The hypothesis that x is in
% yA + I or y in xA + I for all x and y is introduced as _ and not used.
\end{proof}
\leanref{Development/AutomorphicSpectrum/Proofs/BlockDecomposition/AlmostFiniteModuleBoundedTorsion.lean}{GlobalTheory.CompanionTheory.AlmostFinite.BS.ext}
\leanref{Development/AutomorphicSpectrum/Proofs/BlockDecomposition/AlmostFiniteModuleBoundedTorsion.lean}{GlobalTheory.CompanionTheory.AlmostFinite.BS.cyclic}
\leanref{Development/AutomorphicSpectrum/Proofs/BlockDecomposition/AlmostFiniteModuleBoundedTorsion.lean}{GlobalTheory.CompanionTheory.AlmostFinite.BS.of_finite}
\leanref{Development/AutomorphicSpectrum/Proofs/BlockDecomposition/AlmostFiniteModuleBoundedTorsion.lean}{GlobalTheory.CompanionTheory.AlmostFinite.limit_of_cauchy}

\begin{lemma}\label{gal-realization-kappa-truncated-periods-exist-lifting-splitting}
Let $\Lambda$ be a discrete valuation ring and $\ell\geq1$ an integer such that every nonzero element of $\Lambda$ divides a
power of $\ell$. Let $A$ be a $\Lambda$-algebra with divided powers $\gamma$ on an ideal $I$, and $m\geq0$, such that
$I^{[m]}=0$ (Definition~\ref{gal-realization-kappa-approximations}) and: $\ell x\in I^{[r]}$ implies $x\in I^{[r]}$; for
every integer $n\geq1$ some power of $\ell$ lies in $nA$; each $I^{[r]}$ is finitely generated; $I$ is prime; for all
$x,y\in A$, $x\in yA+I$ or $y\in xA+I$; for $x\notin I$ some power of $\ell$ lies in $xA+I$; $A$ is $\ell$-adically
complete and separated; and every $x\in A$ that lies in $\ell^kA+I^{[r]}$ for every $k$ lies in $I^{[r]}$. Then $(A,\gamma)$
has a graded splitting of length $m$ over $\Lambda$
(Definition~\ref{gal-realization-kappa-truncated-periods-exist-base-change-splitting}).
\end{lemma}
\leanref{Development/AutomorphicSpectrum/Proofs/BlockDecomposition/GradedSplittingOverLatticeBaseStatement.lean}{GlobalTheory.T.graded_splitting_over_lattice_base}
\leanref{Development/AutomorphicSpectrum/Proofs/BlockDecomposition/GradedSplittingOverLatticeBase.lean}{GlobalTheory.CompanionTheory.LatSplit.nonempty_grSplit}

\begin{proof}
Let $\varpi$ be a uniformizer of $\Lambda$, $c$ the image of $\ell$ in $\Lambda$, and $L_r=I^{[r]}$, a $\Lambda$-submodule
of $A$. Then $c\neq0$: otherwise $\ell\cdot1=0\in I^{[m]}$, so $1\in I^{[m]}=0$, and $A=0$, which has no prime ideal.

(a) For $0\neq a\in\Lambda$, $x\in A$ and $r\geq0$, $ax\in L_r$ implies $x\in L_r$: $a$ divides some $\ell^k$, so $\ell^kx\in
L_r$, and $x\in L_r$ by the condition on $\ell$, applied $k$ times.

(b) If $y\in L_r$ lies in $\ell^kL_r$ for every $k$, then $y=0$, as $A$ is $\ell$-adically separated. If $(y_k)$ is a
sequence in $L_r$ with $y_{k+1}-y_k\in\ell^kL_r$ for every $k$, it has a limit $y\in A$ with $y-y_k\in\ell^kA$, as $A$ is
complete; $y$ lies in $\ell^kA+L_r$ for every $k$, so in $L_r$; and $y-y_k=\ell^kt$ with $\ell^kt\in L_r$, so $t\in L_r$
by (a), and $y-y_k\in\ell^kL_r$.

(c) \emph{A section.} Fix $r$ and let $M=A/L_r$, a torsion-free $\Lambda$-module by (a). Choose elements $e_j$ of $M$ whose
images form a basis of the $\Lambda/\varpi$-vector space $M/\varpi M$. If $\sum_jg_je_j\in\varpi M$ for a finite family
$(g_j)$ in $\Lambda$, then $\varpi$ divides every $g_j$. If $\sum_jg_je_j=0$, write $g_j=\varpi g'_j$; then
$\sum_jg'_je_j=0$, as $M$ is torsion-free, and by induction every power of $\varpi$ divides every $g_j$, so $g_j=0$. So
the $e_j$ span a free submodule $P$ with $M=P+\varpi M$, and $\varpi x\in P$ implies $x\in P$: if $\varpi
x=\sum_jf_je_j$, then $f_j=\varpi f'_j$, and $\varpi(x-\sum_jf'_je_j)=0$. As $c=u\varpi^k$ with $u$ a unit, induction on
$k$ gives $M=P+cM$, and $cx\in P$ implies $x\in P$. As $P$ is free, there is a $\Lambda$-linear $\sigma:P\to A$ whose
composite with the projection $A\to M$ is the inclusion. Let $W\subset A$ be the inverse image of $P$, which contains
$L_r$; then $A=W+cA$, and $cx\in W$ implies $x\in W$. Let $\rho_0:W\to L_r$ send $w$ to $w-\sigma(\bar w)$, where $\bar w$
is the image of $w$ in $P$.

The map $\rho_0$ extends to a $\Lambda$-linear $\rho:A\to L_r$. For $x\in A$ and $k\geq0$ there are $w\in W$ and $z\in A$
with $x=w+c^kz$, by induction on $k$. If $w,w'\in W$ and $w-w'=c^kt$, then $t\in W$ and
$\rho_0(w)-\rho_0(w')=c^k\rho_0(t)$. Choose decompositions $x=w_k+c^kz_k$ for every $k$; then
$w_{k+1}-w_k=c^k(z_k-cz_{k+1})$, so $\rho_0(w_{k+1})-\rho_0(w_k)\in c^kL_r$, and by (b) the $\rho_0(w_k)$ have a limit
$\rho(x)\in L_r$. For every decomposition $x=w+c^kz$ with $w\in W$, $\rho(x)-\rho_0(w)\in c^kL_r$; by (b) this property
determines $\rho(x)$, and it shows that $\rho$ is $\Lambda$-linear and extends $\rho_0$.

For $x\in L_r$, $\bar x=0$, so $\rho(x)=\rho_0(x)=x$. So $x\mapsto x-\rho(x)$ induces a $\Lambda$-linear map $s_r:M\to A$,
and $s_r(\bar x)-x=-\rho(x)\in L_r$ for $x\in A$.

(d) For $i<m$ let $N_i$ be the image of $L_i$ in $A/L_{i+1}$, which is torsion-free by (a), hence flat over $\Lambda$
\cite[Tag 0539]{Sta}, and $\nu_i$ the restriction of $s_{i+1}$ to $N_i$. For $x\in L_i$, $\nu_i(\bar x)-x\in L_{i+1}\subset
L_i$, so $\nu_i(N_i)\subset L_i$ and $x-\nu_i(\bar x)\in L_{i+1}$. If $\nu_i(n)\in L_{i+1}$, then $n$, the image of
$\nu_i(n)$ in $A/L_{i+1}$, is $0$.
% Lean: LatSplit.nonempty_grSplit: hc for c /= 0, ht for (a), hsec with the separation (IsHausdorff) and the limit
% (IsPrecomplete, LatBase.closed, htf) for (b), and exists_section for (c), through exists_basic (exists_basic_irreducible
% for P, and the reduction from ϖ to c) and exists_ext for ρ; (d) is the final structure, with hflat
% (flat_iff_torsion_eq_bot_of_isBezout) and hν.
\end{proof}
\leanref{Development/AutomorphicSpectrum/Proofs/BlockDecomposition/GradedSplittingOverLatticeBase.lean}{GlobalTheory.CompanionTheory.LatSplit.exists_section}
\leanref{Development/AutomorphicSpectrum/Proofs/BlockDecomposition/GradedSplittingOverLatticeBase.lean}{GlobalTheory.CompanionTheory.LatSplit.exists_ext}
\leanref{Development/AutomorphicSpectrum/Proofs/BlockDecomposition/GradedSplittingOverLatticeBase.lean}{GlobalTheory.CompanionTheory.LatSplit.exists_basic_irreducible}

\begin{lemma}\label{gal-realization-kappa-truncated-periods-exist-lifting-finite-local}
Let $\Lambda$ be a ring and $s,\ell\geq1$ integers such that multiplication by $\ell$ is injective on $\Lambda$ and for
every integer $n\geq1$ some power of $\ell$ lies in $n\Lambda$. Let $X$ be a $\Lambda$-scheme, and form
$\mathrm{Th}^{(s)}_\ell(X/\Lambda)$ for $(\Lambda,0,0,0)$ and the identity of $\Lambda$
(Lemma~\ref{gal-realization-kappa-truncated-periods-exist-base-change-zariski}), with $\mathrm{gr}_r=\mathcal
O^{[r]}/\mathcal O^{[r+1]}$. Then (1) and (2) of Lemma~\ref{gal-realization-kappa-de-rham-local} hold with $\Lambda$ in
place of $K$ and $\mathrm{Th}^{(s)}_\ell(X/\Lambda)_W$ in place of $\mathrm{Th}^{(s)}_{\mathrm{dR}}(X)_W$, and
Lemma~\ref{gal-realization-kappa-de-rham-bundle} holds with $\Lambda$ in place of $K$ and the categories
$\mathrm{Th}^{(s)}_\ell(X/\Lambda)_{U_i}$ in place of the $\mathrm{Th}^{(s)}_{\mathrm{dR}}(X)_{U_i}$.
\end{lemma}
\leanref{Development/AutomorphicSpectrum/Proofs/BlockDecomposition/GradedFormVanishOfFrame.lean}{GlobalTheory.CompanionTheory.GrForm.van}
\leanref{Development/AutomorphicSpectrum/Proofs/BlockDecomposition/GradedFormLocalOfFrame.lean}{GlobalTheory.CompanionTheory.GrForm.stLoc}
\leanref{Development/AutomorphicSpectrum/Proofs/BlockDecomposition/GradedCechCohomologyFinite.lean}{GlobalTheory.CompanionTheory.GrCech.vb_cech}

\begin{proof}
Step 1 of the proof of Lemma~\ref{gal-realization-kappa-de-rham-local} takes the chart of
Lemma~\ref{gal-realization-kappa-de-rham-chart} for $\mathcal O=R=K$ and $I=0$. Here
Lemma~\ref{gal-realization-kappa-truncated-periods-exist-base-change-chart}(1), for $\mathcal O=R=\Lambda$ and $I=0$, whose
hypotheses are those of the statement, gives a divided power chart $D$ on $W$ in $\mathrm{Th}^{(s)}_\ell(X/\Lambda)_W$ with
$B_D=\Lambda\otimes_\Lambda\Gamma(W,\mathcal O_X)$, $t'_k=1\otimes t_k$ and $J_D=0$, whose point corresponds to
$f\mapsto1\otimes f$. Step 3 there takes the divided power polynomial structure of
Lemma~\ref{gal-realization-kappa-de-rham-poly}; here
Lemma~\ref{gal-realization-kappa-truncated-periods-exist-base-change-chart}(2) gives one for
$\mathrm{Th}^{(s)}_\ell(X/\Lambda)_{\{W\}}$. The rest of that proof, through
Lemmas~\ref{gal-realization-thickening-chart-torsor}, \ref{gal-realization-thickening-chart},
\ref{gal-realization-thickening-level-one-forms} and~\ref{gal-realization-thickening-level-one-graded}, uses neither that
$K$ is a field of characteristic $0$ nor that $\ell=1$, and gives (1) and (2). The proof of
Lemma~\ref{gal-realization-kappa-de-rham-bundle} then applies, with (2) for $\Lambda$ in place of
Lemma~\ref{gal-realization-kappa-de-rham-local}(2).
% Lean: GrForm.van (StVan) and GrForm.stLoc (StLoc), through GrForm.stFrame, are stated for every ring O with TFBase
% (dividedPowersBot O) l s, and GrCech.vb_cech likewise; the book's Lemmas gal-realization-kappa-de-rham-local and
% gal-realization-kappa-de-rham-bundle state their case O = K, l = 1. FinOverO.van_coord repeats the argument of (1).
\end{proof}
\leanref{Development/AutomorphicSpectrum/Proofs/BlockDecomposition/CapOAffine.lean}{GlobalTheory.CompanionTheory.FinOverO.van_coord}

\begin{lemma}\label{gal-realization-kappa-truncated-periods-exist-lifting-finite}
Let $\Lambda$ be a Noetherian ring and $s,\ell\geq1$ integers such that multiplication by $\ell$ is injective on $\Lambda$
and for every integer $n\geq1$ some power of $\ell$ lies in $n\Lambda$. Let $X$ be a proper $\Lambda$-scheme with a
coordinate cover (Definition~\ref{gal-realization-thickening-coordinates}(2)), $\mathcal V$ a set of open subsets of $X$
with union $X$, and $\mathcal Q=\mathrm{Th}^{(s)}_\ell(X/\Lambda)_{\mathcal V}$, formed for $(\Lambda,0,0,0)$ and the
identity of $\Lambda$ (Lemma~\ref{gal-realization-kappa-truncated-periods-exist-base-change-zariski}). For integers $r\leq
r'\leq s$ and $j\geq0$, $H^j(\mathcal Q,\mathcal O^{[r]}/\mathcal O^{[r']})$ is a finitely generated $\Lambda$-module.
\end{lemma}
\leanref{Development/WeilBound/CoreStatements/PadicHodgeComparison/EtaleCohomologyComparison/PerfectoidPullbacks/SplittingFieldLattices/BasisB.lean}{CrystallineSemistableComparison.Pb.FinO}
\leanref{Development/WeilBound/CoreStatements/PadicHodgeComparison/EtaleCohomologyComparison/PerfectoidPullbacks/SplittingFieldLattices/BcOnFalse.lean}{CrystallineSemistableComparison.Pb.finO'}

\begin{proof}
Write $\mathcal Q_{\mathcal W}=\mathrm{Th}^{(s)}_\ell(X/\Lambda)_{\mathcal W}$ for a set $\mathcal W$ of open subsets of
$X$, $\mathrm{gr}_n=\mathcal O^{[n]}/\mathcal O^{[n+1]}$, and let $(U_\alpha)_{\alpha\in E}$ be a coordinate cover of $X$
and $\mathcal U=\{U_\alpha:\alpha\in E\}$. Every thickening lies over $X$
(Lemma~\ref{gal-realization-thickening-over}(2)); by
Lemma~\ref{gal-realization-kappa-truncated-periods-exist-base-change-zariski} with
Lemma~\ref{gal-realization-thickening-zariski}(1), restriction from $\mathcal Q_{\{X\}}$ to $\mathcal Q_{\mathcal U}$ and to
$\mathcal Q=\mathcal Q_{\mathcal V}$ is bijective on $H^j(\cdot,\mathcal O^{[r]}/\mathcal O^{[r']})$. So it suffices to treat
$\mathcal Q_{\mathcal U}$. Let $\mathcal S$ be the class of finitely generated $\Lambda$-modules; as $\Lambda$ is
Noetherian, it satisfies the hypothesis of Lemma~\ref{gal-realization-thickening-level-one-serre}.

(a) For $n<s$ and every $j$, $H^j(\mathcal Q_{\mathcal U},\mathrm{gr}_n)\in\mathcal S$. If $n<0$, then $\mathrm{gr}_n=0$.
Let $0\leq n<s$. By Lemma~\ref{gal-realization-thickening-over}, the thickenings over the $U_\alpha$ form cosieves of
$\mathcal Q_{\mathcal U}$ that cover it, with intersections the $\mathcal C_i=\mathcal Q_{U_i}$,
$U_i=\bigcap_kU_{i_k}$; by Lemma~\ref{gal-realization-thickening-level-one-cech} it suffices that $H^a(\check C^\bullet_b)
\in\mathcal S$ for $a+b=j$, where $\check C^a_b=\prod_{i\in E^{a+1}}H^b(\mathcal C_i,\mathrm{gr}_n)$. If $b\neq n$, then
$\check C^a_b=0$: $U_i$ is an affine open of $U_{i_0}$, which has coordinates, so $H^b(\mathcal C_i,\mathrm{gr}_n)=0$ by
Lemma~\ref{gal-realization-kappa-truncated-periods-exist-lifting-finite-local} and
Lemma~\ref{gal-realization-kappa-de-rham-local}(1). If $b=n$, then by
Lemma~\ref{gal-realization-kappa-truncated-periods-exist-lifting-finite-local} and
Lemma~\ref{gal-realization-kappa-de-rham-bundle}, $\check C^\bullet_n$ is the complex of a cosimplicial $\Lambda$-module
with data (i), (ii) and (iii) of Lemma~\ref{gal-realization-thickening-level-one-coherent} for $Z=X$, which is proper over
$\Lambda$, and the opens $U_\alpha$, which cover $X$; by that lemma $H^a(\check C^\bullet_n)$ is finitely generated.

(b) For $r\leq r'\leq s$, $H^j(\mathcal Q_{\mathcal U},\mathcal O^{[r]}/\mathcal O^{[r']})\in\mathcal S$, by induction on
$r'-r$. For $r=r'$ the functor is $0$. For $r<r'$, the long exact sequence of $0\to\mathcal O^{[r+1]}/\mathcal
O^{[r']}\to\mathcal O^{[r]}/\mathcal O^{[r']}\to\mathrm{gr}_r\to0$ gives an exact sequence $H^j(\mathcal Q_{\mathcal
U},\mathcal O^{[r+1]}/\mathcal O^{[r']})\to H^j(\mathcal Q_{\mathcal U},\mathcal O^{[r]}/\mathcal O^{[r']})\to H^j(\mathcal
Q_{\mathcal U},\mathrm{gr}_r)$, whose outer terms lie in $\mathcal S$ by induction and by (a), as $r<r'\leq s$.
% Lean: Pb.finO' = fin_o kh1 poly_tf kh2 kh3.2 = FinOverO.finO: B1 and B2 (cech_refinement_cohomology_bijective) for the
% reduction to the cover, range_fin for (b) (sq_exact and the Serre class finite_modules_serre_class), gr_range for (a):
% KH.GlobSerreSt (Lemma gal-realization-thickening-level-one-cech), van_coord for b /= n (the Λ-version of Lemma
% gal-realization-kappa-de-rham-local(1), with coordinates_restrict_affine_open and capO_affine), and KH.GrCechFin =
% gr_cech_fin for b = n (vb_cech, the Λ-version of Lemma gal-realization-kappa-de-rham-bundle, with
% VBGlue.finite_H_cx, Lemma gal-realization-thickening-level-one-coherent).
\end{proof}
\leanref{Development/AutomorphicSpectrum/Proofs/BlockDecomposition/CapOAffine.lean}{GlobalTheory.CompanionTheory.FinOverO.finO}
\leanref{Development/AutomorphicSpectrum/Proofs/BlockDecomposition/CapOAffine.lean}{GlobalTheory.CompanionTheory.FinOverO.gr_range}
\leanref{Development/AutomorphicSpectrum/Proofs/BlockDecomposition/CapOAffine.lean}{GlobalTheory.CompanionTheory.FinOverO.range_fin}
\leanref{Development/WeilBound/CoreStatements/PadicHodgeComparison/EtaleCohomologyComparison/PerfectoidPullbacks/GrCechFin.lean}{CrystallineSemistableComparison.Pb.gr_cech_fin}

\begin{lemma}\label{gal-realization-kappa-truncated-periods-exist-lifting-generation}
Let $\Lambda$ be a discrete valuation ring and $\ell\geq1$ an integer such that every nonzero element of $\Lambda$ divides a
power of $\ell$, multiplication by $\ell$ is injective on $\Lambda$, and for every integer $n\geq1$ some power of $\ell$
lies in $n\Lambda$. Let $A$, $\gamma$, $I$ and $m$ be as in
Lemma~\ref{gal-realization-kappa-truncated-periods-exist-lifting-splitting}, and $s\geq\max(m,1)$. Let $X$ be a proper
$\Lambda$-scheme with a coordinate cover (Definition~\ref{gal-realization-thickening-coordinates}(2)), $\mathcal V$ a set of
open subsets of $X$ with union $X$, and $\mathcal Q_A=\mathrm{Th}^{(s)}_\ell(X/A)_{\mathcal V}$, formed for
$(A,I,\gamma,0)$ and $\Lambda\to A\to A/I$
(Lemma~\ref{gal-realization-kappa-truncated-periods-exist-base-change-zariski}). For integers $r\leq r'\leq s$ and
$j\geq0$, $H^j(\mathcal Q_A,\mathcal O^{[r]}/\mathcal O^{[r']})$ is a finitely generated $A$-module.
\end{lemma}
\leanref{Development/WeilBound/CoreStatements/PadicHodgeComparison/EtaleCohomologyComparison/PerfectoidPullbacks/SplittingFieldLattices/BasisB.lean}{CrystallineSemistableComparison.Pb.AfgA}
\leanref{Development/WeilBound/CoreStatements/PadicHodgeComparison/EtaleCohomologyComparison/PerfectoidPullbacks/SplittingFieldLattices/BcOnFalse.lean}{CrystallineSemistableComparison.Pb.afgA}

\begin{proof}
Write $F=\mathcal O^{[r]}/\mathcal O^{[r']}$. By Lemma~\ref{gal-realization-kappa-truncated-periods-exist-lifting-splitting},
$(A,\gamma)$ has a graded splitting $(N_i,\nu_i)_{i<m}$ over $\Lambda$. The hypotheses of
Lemma~\ref{gal-realization-kappa-truncated-periods-exist-base-change-tensor} hold: for $(\Lambda,0)$, $0^{[s]}=0$ as $s\geq1$,
and the condition on $\ell$ says that multiplication by $\ell$ is injective on $\Lambda$; for $(A,I)$, $I^{[s]}\subset
I^{[m]}=0$. So, with $\mathcal Q_\Lambda=\mathrm{Th}^{(s)}_\ell(X/\Lambda)_{\mathcal V}$, the map of
Lemma~\ref{gal-realization-kappa-truncated-periods-exist-base-change-tensor}(2) from $\bigoplus_iN_i\otimes_\Lambda M_i$, where
$M_i=H^j(\mathcal Q_\Lambda,\mathcal O^{[r-i]}/\mathcal O^{[r'-i]})$, onto $H^j(\mathcal Q_A,F)$ is surjective. By
Lemma~\ref{gal-realization-kappa-truncated-periods-exist-lifting-finite}, each $M_i$ is generated by finitely many
$x_{i,1},\ldots,x_{i,q_i}$, and by hypothesis the ideal $I^{[i]}$ by finitely many $g_{i,1},\ldots,g_{i,p_i}$. For $a\in
I^{[i]}$ let $\mu_a:M_i\to H^j(\mathcal Q_A,F)$ be multiplication by $a$
(Lemma~\ref{gal-realization-kappa-truncated-periods-exist-base-change-tensor}(1)). On cochains it multiplies the values by
$a$, so $\mu_{ba+b'a'}=b\mu_a+b'\mu_{a'}$ for $b,b'\in A$, and $\mu_a(\lambda x)=\lambda\mu_a(x)$ for $\lambda\in\Lambda$.
For $n\in N_i$ and $x\in M_i$ write $\nu_i(n)=\sum_kb_kg_{i,k}$ and $x=\sum_l\lambda_lx_{i,l}$; then the image of $n\otimes x$
is $\mu_{\nu_i(n)}(x)=\sum_{k,l}b_k\lambda_l\mu_{g_{i,k}}(x_{i,l})$. So the finitely many $\mu_{g_{i,k}}(x_{i,l})$ generate
$H^j(\mathcal Q_A,F)$ over $A$.
% Lean: Pb.afgA = afg_a bcOn_false finO'.2 = LatticeFinite.afgA_of_bcOn_finO, for AfgA, whose conclusion AFG (l : A)
% (HQ (TS γ q l s 𝒱) r r' j) (an exponent n and a finitely generated submodule containing l^n times every element) is
% proved with n = 0: graded_splitting_over_lattice_base (Lemma ...-lifting-splitting), tfBase_mono, BcOn False (Lemma
% ...-base-change-tensor), FinO (Lemma ...-lifting-finite) and exists_fg_pull_sqC, with pull_sqC_comb for the formula for
% μ_{ba+b'a'}.
\end{proof}
\leanref{Development/AutomorphicSpectrum/Proofs/BlockDecomposition/ExistsFgPullSqC.lean}{GlobalTheory.CompanionTheory.LatticeFinite.afgA_of_bcOn_finO}
\leanref{Development/AutomorphicSpectrum/Proofs/BlockDecomposition/ExistsFgPullSqC.lean}{GlobalTheory.CompanionTheory.LatticeFinite.exists_fg_pull_sqC}

\begin{lemma}\label{gal-realization-kappa-truncated-periods-exist-lifting}
Let $\Lambda$ be a discrete valuation ring and $\ell$, $m$ and $s$ integers with $\ell\geq0$, $s\geq1$ and $0\leq
m\leq s$, such that every nonzero element of $\Lambda$ divides a power of $\ell$, multiplication by $\ell$ is injective
on $\Lambda$, and for every integer $n\geq1$ some power of $\ell$ lies in $n\Lambda$. Let $A$ be a $\Lambda$-algebra with
divided powers $\gamma$ on an ideal $I$ such that $I^{[m]}=0$
(Definition~\ref{gal-realization-kappa-approximations}) and: $\ell x\in I^{[r]}$ implies $x\in I^{[r]}$; for every integer
$n\geq1$ some power of $\ell$ lies in $nA$; each $I^{[r]}$ is finitely generated; $I$ is prime; for all $x,y\in A$, $x\in
yA+I$ or $y\in xA+I$; for $x\notin I$ some power of $\ell$ lies in $xA+I$; $A$ is $\ell$-adically complete and
separated; and every $x\in A$ that lies in $\ell^kA+I^{[r]}$ for every $k$ lies in $I^{[r]}$. Let $X$ be a proper
$\Lambda$-scheme with a coordinate cover (Definition~\ref{gal-realization-thickening-coordinates}(2)), $\mathcal V$ a set
of open subsets of $X$ with union $X$, $j\geq0$, and $\mathcal Q$ the full subcategory of
$\mathrm{Th}^{(s)}(X/A)_{\mathcal V}$, formed for $(A,I,\gamma,0)$, $\Lambda$ in place of $\mathcal O$ and the structure
map $\Lambda\to A/I$ (Definition~\ref{gal-realization-kappa-de-rham-truncated}(1)), of the thickenings $(B,i,J,\delta,y)$
such that $\ell
x\in J^{[r]}$ implies $x\in J^{[r]}$ for all $r\geq0$ and $x\in B$. Let
$\mathrm{red}_n:H^j(\mathcal Q,\mathcal O)\to H^j(\mathcal Q,\mathcal O/\ell^n\mathcal O)$ and $t:H^j(\mathcal Q,\mathcal
O/\ell^{n+1}\mathcal O)\to H^j(\mathcal Q,\mathcal O/\ell^n\mathcal O)$ be induced by the projections
(Definition~\ref{gal-realization-thickening-nilpotent}(3)). There is $e\geq0$ with the following properties.
\begin{enumerate}
\item For every family $(x_n)_{n\geq1}$ with $x_n\in H^j(\mathcal Q,\mathcal O/\ell^n\mathcal O)$ and $t(x_{n+1})=x_n$
  for every $n$, there is $z\in H^j(\mathcal Q,\mathcal O)$ with $\mathrm{red}_n(z)=\ell^ex_n$ for every $n$.
\item If $z\in H^j(\mathcal Q,\mathcal O)$ and $\mathrm{red}_n(z)=0$ for every $n$, then $\ell^ez=0$.
\item Let $r\leq s$, and let $z\in H^j(\mathcal Q,\mathcal O)$ have image $0$ in $H^j(\mathcal Q,\mathcal O/(\mathcal
  O^{[r]}+\ell^n\mathcal O))$ for every $n\geq1$. Then $\ell^e$ times the image of $z$ in $H^j(\mathcal Q,\mathcal
  O/\mathcal O^{[r]})$ is $0$.
\end{enumerate}
\end{lemma}
\leanref{Development/WeilBound/CoreStatements/PadicHodgeComparison/EtaleCohomologyComparison/PerfectoidPullbacks/SplittingFieldLattices/BasisB.lean}{CrystallineSemistableComparison.Pb.LiftP}
\leanref{Development/WeilBound/CoreStatements/PadicHodgeComparison/EtaleCohomologyComparison/PerfectoidPullbacks/SplittingFieldLattices/BasisB.lean}{CrystallineSemistableComparison.Pb.LiftA}
\leanref{Development/WeilBound/CoreStatements/PadicHodgeComparison/EtaleCohomologyComparison/PerfectoidPullbacks/SplittingFieldLattices/BcOnFalse.lean}{CrystallineSemistableComparison.Pb.liftA}
\leanref{Development/WeilBound/CoreStatements/PadicHodgeComparison/EtaleCohomologyComparison/PerfectoidPullbacks/LevelChainValuations/ValuationChainProofs/CgLift.lean}{CrystallineSemistableComparison.Pb.Kap.cg_lift_pf}

\begin{proof}
As $\Lambda\neq0$ and multiplication by $\ell$ is injective on $\Lambda$, $\ell\geq1$. The category $\mathcal Q$ is
$\mathrm{Th}^{(s)}_\ell(X/A)_{\mathcal V}$ (Lemma~\ref{gal-realization-kappa-truncated-periods-exist-base-change-zariski}).
For $T=(B,i,J,\delta,y)$ in $\mathcal Q$ and $r\geq0$, multiplication by $\ell$ is injective on $B/J^{[r]}$, by the
condition defining $\mathcal Q$, and on $B$, the case $r=s$, as $J^{[s]}=0$.

\emph{Step 1: bounds.} By hypothesis every nonzero element of $\Lambda$ divides a power of $\ell$, multiplication by
$\ell$ is injective on $\Lambda$, for every integer $n\geq1$ some power of $\ell$ lies in $n\Lambda$, and
$s\geq\max(m,1)$. So Lemma~\ref{gal-realization-kappa-truncated-periods-exist-lifting-generation} applies, for $\Lambda$,
$X$ and $\mathcal V$: for $0\leq t\leq s$ and every $j'$, $H^{j'}(\mathcal Q,\mathcal O/\mathcal
O^{[t]})$ is a finitely generated $A$-module. The hypotheses of
Lemma~\ref{gal-realization-kappa-truncated-periods-exist-lifting-modules} hold for $A$, $I$ and $c=\ell$: $I=I^{[1]}$ is
finitely generated and prime, $I^m\subset I^{[m]}=0$, the conditions on all $x,y$ and on $x\notin I$ are hypotheses, the
closedness is the
hypothesis for $r=1$, and limits exist as $A$ is $\ell$-adically complete. Applying it to $H^j$ and $H^{j+1}$ gives, for
each $t\leq s$, an integer $k_t$ such that (a), (b) and (c) of
Lemma~\ref{gal-realization-kappa-truncated-periods-exist-lifting-limits} hold for $c=\ell$, $k=k_t$,
$H^{j+1}(\mathcal Q,\mathcal O/\mathcal O^{[t]})$ and $H^j(\mathcal Q,\mathcal O/\mathcal O^{[t]})$.

\emph{Step 2: (1) and (2).} Let $K^\bullet$ be the complex of cochains of $\mathcal Q$ with values in $\mathcal O$, and
$K^\bullet_n$ the one with values in $\mathcal O/\ell^n\mathcal O$, with $\mathrm{red}_n$ and $t_n$ induced by the
projections; for $n=0$ the latter is $0$. Multiplication by $\ell$ is injective on the terms of $K^\bullet$, and
$\mathrm{red}_n$ is surjective with kernel $\ell^nK^\bullet$, as both hold valuewise. As $J_T^{[s]}=0$, the projection
$\mathcal O\to\mathcal O/\mathcal O^{[s]}$ is an isomorphism of functors on $\mathcal Q$, so the bounds of Step 1 for
$t=s$ hold for $H^{j+1}(K^\bullet)$ and $H^j(K^\bullet)$. Lemma~\ref{gal-realization-kappa-truncated-periods-exist-lifting-limits},
with $x_0=0$, gives (1) with $\ell^{2k_s}$ in place of $\ell^e$, and (2) with $\ell^{k_s}$.

\emph{Step 3: (3).} Let $r$ and $z$ be as in (3). If $r\leq0$, then $\mathcal O/\mathcal O^{[r]}=0$. Let $1\leq r\leq s$,
let $K^\bullet$ be the complex of cochains with values in $\mathcal O/\mathcal O^{[r]}$, and $K^\bullet_n$ the one with
values in $\mathcal O/(\mathcal O^{[r]}+\ell^n\mathcal O)$. Multiplication by $\ell$ is injective on the terms of
$K^\bullet$, and $\mathrm{red}_n$ is surjective with kernel $\ell^nK^\bullet$, as $(J^{[r]}+\ell^nB)/J^{[r]}=\ell^n(B/J^{[r]})$.
The image of $z$ in $H^j(K^\bullet)$ maps to $0$ in every $H^j(K^\bullet_n)$, so by
Lemma~\ref{gal-realization-kappa-truncated-periods-exist-lifting-limits}(2), with the bounds of Step 1 for $t=r$,
$\ell^{k_r}$ kills it.

So $e=2k_s+\max_{1\leq r\leq s}k_r$ has the three properties.
% Lean: Pb.liftA = lift_a afgA = LiftLat.liftA_of_afgA, for LiftA, whose conclusion is LiftP (TS γ q l s 𝒱) l s j e.
% The hypotheses on Λ are hdv and TFBase (dividedPowersBot Λ) l s, those on A are LatBase γ l m, and hms is m <= s;
% all are passed unchanged to AfgA. liftingProperty_of_bounds: hzero and htf are the torsion facts, hQ is AfgA (Lemma
% ...-lifting-generation, through afg_of_bij and bij_quoM for H^j(Q, O)), hb is bounds (almost_finite_module_bounded_torsion,
% Lemma ...-lifting-modules, with dividedPowerFiltration_one_eq and pow_le_pdP'); top_part is Step 2 (gen_top, gen and
% inverse_limit_cohomology_comparison, Lemma ...-lifting-limits), quo_part is Step 3 (gen_quo), and sub_HQ is the case
% r <= 0. The exponent is 2k + sup of the f(t), as here. LiftP indexes from n = 0, where the groups vanish.
\end{proof}
\leanref{Development/AutomorphicSpectrum/Proofs/BlockDecomposition/LiftPOfBounds.lean}{GlobalTheory.CompanionTheory.LiftLat.liftA_of_afgA}
\leanref{Development/AutomorphicSpectrum/Proofs/BlockDecomposition/LiftPOfBounds.lean}{GlobalTheory.CompanionTheory.LiftLat.liftingProperty_of_bounds}
\leanref{Development/AutomorphicSpectrum/Proofs/BlockDecomposition/LiftPOfBounds.lean}{GlobalTheory.CompanionTheory.LiftLat.bounds}

\begin{lemma}\label{gal-realization-kappa-truncated-periods-exist-effective}
Let $\mathcal D_m(Y,s)$ be as in Definition~\ref{gal-realization-kappa-truncated-periods-exist-categories}, for an
approximation system, a smooth projective $\mathcal O_K$-scheme $Y$, $m\geq0$ and $s\geq m+1$. Suppose that $Y_K$ has a
coordinate cover $(U_\alpha)$ of dimensions $d_\alpha$ (Definition~\ref{gal-realization-thickening-coordinates}(2)) with
$d_\alpha+m+1\leq s$ for every $\alpha$. Then $\mathcal D_m(Y,s)$ is effective
(Definition~\ref{gal-realization-kappa-truncated-periods-exist-diagram}(2)).
\end{lemma}
\leanref{Development/WeilBound/CoreStatements/PadicHodgeComparison/EtaleCohomologyComparison/PerfectoidPullbacks/LevelChainValuations/InstSt.lean}{CrystallineSemistableComparison.Pb.Kap.chain_good}
\leanref{Development/WeilBound/CoreStatements/PadicHodgeComparison/EtaleCohomologyComparison/PerfectoidPullbacks/LevelChainValuations/ValuationChainProofs/LevelSystemGood.lean}{CrystallineSemistableComparison.Pb.Kap.level_system_good}

\begin{proof}
Write $\mathcal Q=\mathcal Q_m(Y,s)$ and $\mathcal Q'=\mathcal Q'_m(Y,s)$, and fix $j$.

(a) As $\ell\neq0$ in $K$ and $B_m$ is a $K$-algebra, $\ell$ is a unit of $B_m$. As $Y_K$ is proper and smooth over $K$,
$H^j_{\mathrm{dR}}(Y_K)$ is finite-dimensional (Lemma~\ref{gal-realization-kappa-de-rham}(1)).

(b) As $J_m=\iota_m(I_m)B_m$ and $\iota_m$ is a divided power map, $J_m^{[m+1]}=\iota_m(I_m^{[m+1]})B_m$ (proof of
Lemma~\ref{gal-realization-thickening-chart}(2)), which is $0$ by (d) of
Definition~\ref{gal-realization-kappa-approximations}(1). So
Lemma~\ref{gal-realization-kappa-truncated-periods-exist-base-change}(1) applies to $B_m$, $\delta_m$, $m+1\leq s$,
$X=Y_K$ and $\mathcal V=\{Y_K\}$, with $\mathcal Q'$ and the map $\mu$ of $\mathcal D_m(Y,s)$.

(c) The maps $t_Z\circ\mathrm{red}_{n+1}$ and $\mathrm{red}_n$ are both induced by the identity functor with the
projections $B\to B/\ell^nB$. The map $t_Z\circ r_{n+1}$ is induced by $R_{n+1}$ with the projections $B/\ell^{n+1}B\to
B/\ell^nB$, and $r_n\circ t_S$ by $R_n$, viewed as a functor to $\mathrm{Th}_{n+1}(Y)$
(Definition~\ref{gal-realization-period-thickenings}(3)(a)), with identity maps. For $T$ in $\mathcal Q$ the projection
$B/\ell^{n+1}B\to B/\ell^nB$ is a map of thickenings from $R_{n+1}(T)$ to $R_n(T)$, natural in $T$: it is a divided power
map, its points agree, and it is compatible with the maps from $A_{n+1}$, as $\rho_{m,n}$ composed with $A_{n+1}\to A_n$
and $\rho_{m,n+1}$ followed by $A_m/\ell^{n+1}A_m\to A_m/\ell^nA_m$ both send the class of $c\in A_{\mathrm{cris}}$ to that
of $\lambda_m(c)$. Lemma~\ref{gal-realization-period-prism} gives $t_Z\circ r_{n+1}=r_n\circ t_S$.

(d) and (e). As $Y$ is projective over $\mathcal O_K$, it is proper \cite[Tag 01WC]{Sta}, and as it is also smooth, it
has a coordinate cover (Lemma~\ref{gal-realization-thickening-cover}(2)). The ring $\mathcal O_K$ is a discrete valuation
ring of characteristic $0$ in which $\ell$ is not a unit, so every nonzero element of $\mathcal O_K$ divides a power of
$\ell$, multiplication by $\ell$ is injective on $\mathcal O_K$, and every integer $n\geq1$, being nonzero in $\mathcal
O_K$, divides a power of $\ell$. By (d) of Definition~\ref{gal-realization-kappa-approximations}(1), and as
$I_m^{[s]}\subset I_m^{[m+1]}=0$ with $s\geq m+1\geq1$,
Lemma~\ref{gal-realization-kappa-truncated-periods-exist-lifting} applies to $\Lambda=\mathcal O_K$, $A_m$, $I_m$,
$\gamma_m$, $s$ in place of $m$, $X=Y$, $\mathcal V=\{Y\}$ and $j$, with the category $\mathcal Q$, as every thickening lies
over $Y$ (Lemma~\ref{gal-realization-thickening-over}(2)); let $e$ be as there. Let
$x$ be a compatible family with $x_n\in H^j_n(Y)$. By (c) the $r_n(x_n)$ are compatible under $t_Z$, so
Lemma~\ref{gal-realization-kappa-truncated-periods-exist-lifting}(1) gives $z$ with $\mathrm{red}_n(z)=\ell^er_n(x_n)$
for every $n$. If $\mathrm{red}_n(z)=0$ for every $n$, then $\ell^ez=0$ by
Lemma~\ref{gal-realization-kappa-truncated-periods-exist-lifting}(2), so $\ell^eb(z)=b(\ell^ez)=0$, and $b(z)=0$ as $\ell$
is a unit of $B_m$.
% Lean: chain_good = level_system_good: unit is unit_B_pf, fin is m_fin, basis is cg_basis_pf (basisB, with cg_basis_pdP
% through dividedPowerFiltration_le_map, and cg_basis_empty for a cover with no member, which the hypothesis m + 1 <= s
% makes unnecessary here), r_tr is rHof_tr_gen (Pb.red_tr = RedTrans.red_trans, with tr_comp, tr_dp and the natural
% map eta), red_tr is trM_redM, and lift and sep are cg_lift_pf (liftA with ok_dvr, sm_proper and cov_ex).
\end{proof}
\leanref{Development/WeilBound/CoreStatements/PadicHodgeComparison/EtaleCohomologyComparison/PerfectoidPullbacks/LevelChainValuations/ValuationChainProofs/unitB.lean}{CrystallineSemistableComparison.Pb.Kap.unit_B_pf}
\leanref{Development/WeilBound/CoreStatements/PadicHodgeComparison/EtaleCohomologyComparison/PerfectoidPullbacks/LevelChainValuations/ValuationChainProofs/CgBasis.lean}{CrystallineSemistableComparison.Pb.Kap.cg_basis_pf}
\leanref{Development/WeilBound/CoreStatements/PadicHodgeComparison/EtaleCohomologyComparison/PerfectoidPullbacks/LevelChainValuations/ValuationChainProofs/CgBasis.lean}{CrystallineSemistableComparison.Pb.Kap.cg_basis_pdP}
\leanref{Development/WeilBound/CoreStatements/PadicHodgeComparison/EtaleCohomologyComparison/PerfectoidPullbacks/LevelChainValuations/RhoActionCup/RHofTrGen.lean}{CrystallineSemistableComparison.Pb.Kap.rHof_tr_gen}
\leanref{Development/AutomorphicSpectrum/Proofs/BlockDecomposition/RedTrans.lean}{GlobalTheory.CompanionTheory.RedTrans.red_trans}

\begin{lemma}\label{gal-realization-kappa-truncated-periods-exist-maps}
Let the approximation system, $\mathcal Q_m(Y,s)$, $\mathcal Q'_m(Y,s)$ and $\mathcal D_m(Y,s)$ be as in
Definition~\ref{gal-realization-kappa-truncated-periods-exist-categories}, for smooth projective $\mathcal O_K$-schemes
$Y$ and $Y'$, $m\geq0$ and $s\geq1$. In (1) to (4) below, the data listed form a map of comparison diagrams
(Definition~\ref{gal-realization-kappa-truncated-periods-exist-diagram}(4)), where $h_Z$ and $h_W$ are induced by the
functors named, with the identity maps of the rings (Definition~\ref{gal-realization-split-diagrams}(1)).
\begin{enumerate}
\item For $s\leq s'$, from $\mathcal D_m(Y,s')$ to $\mathcal D_m(Y,s)$: identity ring maps, $h_S$ and $h_V$ the
  identities, and the inclusions $\mathcal Q_m(Y,s)\subset\mathcal Q_m(Y,s')$ and $\mathcal Q'_m(Y,s)\subset\mathcal
  Q'_m(Y,s')$.
\item From $\mathcal D_{m+1}(Y,s)$ to $\mathcal D_m(Y,s)$: $s_\Lambda$ the identity, $s_A=d_A$ and $s_B=d_B$, $h_S$ and
  $h_V$ the identities, and the restrictions of scalars along $d_A$ and $d_B$
  (Definition~\ref{gal-realization-thickening-nilpotent}(4)), from $\mathcal Q_m(Y,s)$ to $\mathcal Q_{m+1}(Y,s)$ and from
  $\mathcal Q'_m(Y,s)$ to $\mathcal Q'_{m+1}(Y,s)$.
\item For $\sigma\in G_K$, from $\mathcal D_m(Y,s)$ to itself: $s_\Lambda=\sigma$, $s_A=\sigma_{A_m}$ and
  $s_B=\sigma_{B_m}$, $h_S$ the map $\sigma$ of Definition~\ref{gal-realization-period-thickenings}(3)(b), $h_V$ the
  identity, and the restrictions of scalars along $\sigma_{A_m}$ and $\sigma_{B_m}$.
\item For an $\mathcal O_K$-morphism $h:Y\to Y'$ with generic fibre $h_K:Y_K\to Y'_K$, from $\mathcal D_m(Y',s)$ to
  $\mathcal D_m(Y,s)$: identity ring maps, $h_S=h^*$ (Definition~\ref{gal-realization-period-thickenings}(3)(e)),
  $h_V=h_K^*$ (Definition~\ref{gal-realization-period-de-rham}), and the functors $(B,i,J,\delta,y)\mapsto
  (B,i,J,\delta,h\circ y)$ from $\mathcal Q_m(Y,s)$ to $\mathcal Q_m(Y',s)$ and $(B,i,J,\delta,y)\mapsto
  (B,i,J,\delta,h_K\circ y)$ from $\mathcal Q'_m(Y,s)$ to $\mathcal Q'_m(Y',s)$.
\end{enumerate}
Moreover the cup data of $\mathcal D_m(Y,s)$ (Definition~\ref{gal-realization-kappa-truncated-periods-exist-categories}(4))
form a product structure on it (Definition~\ref{gal-realization-kappa-truncated-periods-exist-diagram}(5)).
\end{lemma}
\leanref{Development/WeilBound/CoreStatements/PadicHodgeComparison/EtaleCohomologyComparison/PerfectoidPullbacks/LevelChainValuations/GaloisCupProducts/HomLevGood.lean}{CrystallineSemistableComparison.Pb.Kap.homS}
\leanref{Development/WeilBound/CoreStatements/PadicHodgeComparison/EtaleCohomologyComparison/PerfectoidPullbacks/LevelChainValuations/GaloisCupProducts/HomSGood.lean}{CrystallineSemistableComparison.Pb.Kap.homS_good_of_inclHM}
\leanref{Development/WeilBound/CoreStatements/PadicHodgeComparison/EtaleCohomologyComparison/PerfectoidPullbacks/LevelChainValuations/ValuationChainProofs/RHofIncl.lean}{CrystallineSemistableComparison.Pb.Kap.rHof_incl_pf}
\leanref{Development/WeilBound/CoreStatements/PadicHodgeComparison/EtaleCohomologyComparison/PerfectoidPullbacks/LevelChainValuations/GaloisCupProducts/HomLevGood.lean}{CrystallineSemistableComparison.Pb.Kap.homLev}
\leanref{Development/WeilBound/CoreStatements/PadicHodgeComparison/EtaleCohomologyComparison/PerfectoidPullbacks/LevelChainValuations/RhoActionCup/RhoFinal.lean}{CrystallineSemistableComparison.Pb.Kap.Mo.lev_good}
\leanref{Development/WeilBound/CoreStatements/PadicHodgeComparison/EtaleCohomologyComparison/PerfectoidPullbacks/LevelChainValuations/GaloisCupProducts/HomGalGoodRHof.lean}{CrystallineSemistableComparison.Pb.Kap.homGal}
\leanref{Development/WeilBound/CoreStatements/PadicHodgeComparison/EtaleCohomologyComparison/PerfectoidPullbacks/LevelChainValuations/GaloisCupProducts/HomGalGood.lean}{CrystallineSemistableComparison.Pb.Kap.homGal_good_of_galRed}
\leanref{Development/WeilBound/CoreStatements/PadicHodgeComparison/EtaleCohomologyComparison/PerfectoidPullbacks/LevelChainValuations/GaloisCupProducts/HomGalGood.lean}{CrystallineSemistableComparison.Pb.Kap.galRed_rHof}
\leanref{Development/WeilBound/CoreStatements/PadicHodgeComparison/EtaleCohomologyComparison/PerfectoidPullbacks/LevelChainValuations/GaloisCupProducts/HomPullGoodRHof.lean}{CrystallineSemistableComparison.Pb.Kap.homPull}
\leanref{Development/WeilBound/CoreStatements/PadicHodgeComparison/EtaleCohomologyComparison/PerfectoidPullbacks/LevelChainValuations/GaloisCupProducts/HomPullRHofGood.lean}{CrystallineSemistableComparison.Pb.Kap.homPull_rHof_good}
\leanref{Development/WeilBound/CoreStatements/PadicHodgeComparison/EtaleCohomologyComparison/PerfectoidPullbacks/LevelChainValuations/GaloisCupProducts/CupGood.lean}{CrystallineSemistableComparison.Pb.Kap.cupOf_rHof_good}

\begin{proof}
\emph{The functors.} Each functor named sends a thickening to one with the same ring, ideal and divided powers, so it
keeps the conditions $J^{[s]}=0$ and $\ell x\in J^{[r]}\Rightarrow x\in J^{[r]}$; in (1), $J^{[s']}\subset J^{[s]}$. The
restrictions of scalars in (2) and (3) are defined, as $d_A$, $d_B$, $\sigma_{A_m}$ and $\sigma_{B_m}$ are divided power
maps, $d_A$ and $\sigma_{A_m}$ fix $\mathcal O_K$, and $d_B$ and $\sigma_{B_m}$ fix $K$
(Definition~\ref{gal-realization-kappa-approximations}). In (4), $h\circ y$ and $h_K\circ y$ are morphisms over $\mathcal
O_K$ and over $K$. An element $a$ of the ring of the source acts on the ring of the image of a thickening through
$s_A(a)$, resp.\ $s_B(a)$, so $h_Z$ is $s_A$-semilinear and $h_W$ is $s_B$-semilinear.

\emph{The ring maps.} The identities $\ell'=\ell$, $s_A\circ\lambda=\lambda'\circ s_\Lambda$, $s_B\circ\iota=\iota'\circ
s_A$ and $s_B\circ o=o'$ are trivial in (1) and (4). In (2) they are $d_A\circ\lambda_{m+1}=\lambda_m$, $d_B\circ
\iota_{m+1}=\iota_m\circ d_A$ and $d_B$ fixing $K$ (Definition~\ref{gal-realization-kappa-approximations}(2)); in (3) they
are $\sigma_{A_m}\circ\lambda_m=\lambda_m\circ\sigma$, $\sigma_{B_m}\circ\iota_m=\iota_m\circ\sigma_{A_m}$ and
$\sigma_{B_m}$ fixing $K$ (Definition~\ref{gal-realization-kappa-approximations}(1)(c)).

\emph{The squares.} Each remaining identity compares two maps induced by functors with maps of coefficients, and a
composite of two such maps is induced by the composite functor with the composite maps. In each case below, the two
composite functors agree, or the identity maps of their rings form a natural map between them, so
Lemma~\ref{gal-realization-period-prism} gives the identity.
\begin{itemize}
\item $t_S$: in (1) and (2) $h_S$ is the identity; in (3) and (4), $h_S$ commutes with the reductions by
  Lemma~\ref{gal-realization-period-compatible}(1).
\item $\mathrm{red}_n$: both composites are induced by the same functor, with the projections $B\to B/\ell^nB$.
\item $r_n$: in (1) and (4), $R_n$ commutes with the inclusions and with composing points with $h$. In (2), for $T$ in
  $\mathcal Q_m(Y,s)$, $R_n$ of the restriction of scalars of $T$ along $d_A$ has the map $\bar\imath\bar d_A\rho_{m+1,n}$
  from $A_n$, and $R_n(T)$ has $\bar\imath\rho_{m,n}$, where $\bar d_A$ is induced by $d_A$; these agree, as both
  $\bar d_A\rho_{m+1,n}$ and $\rho_{m,n}$ send the class of $c\in A_{\mathrm{cris}}$ to that of
  $d_A(\lambda_{m+1}(c))=\lambda_m(c)$. In (3), $R_n$ of the restriction of scalars of $T$ along $\sigma_{A_m}$ has
  $\bar\imath\bar\sigma_{A_m}\rho_{m,n}$, and the image of $R_n(T)$ under the functor of
  Definition~\ref{gal-realization-period-thickenings}(3)(b) has $\bar\imath\rho_{m,n}\sigma$; these agree, as
  $\sigma_{A_m}\circ\lambda_m=\lambda_m\circ\sigma$. The other data of the thickenings agree.
\item $b$: the two composites send $T'=(B',i',J',\delta',y')$ to thickenings with ring $B'$ and the same ideal and divided
  powers. Their maps from the ring $A$ of the source are $i'\circ\iota'\circ s_A$ and $i'\circ s_B\circ\iota$, which agree
  as $s_B\circ\iota=\iota'\circ s_A$, and their points are $\mathrm{pr}\circ y'$ in (1) to (3), and $\mathrm{pr}\circ
  h_K\circ y'=h\circ\mathrm{pr}\circ y'$ in (4).
\item $\mu$: in (1), the restrictions to $\mathrm{Th}^{(s)}_{\mathrm{dR}}(Y_K)$ and to
  $\mathrm{Th}^{(s')}_{\mathrm{dR}}(Y_K)$ are compatible with the inclusion, and $E$ commutes with it. In (2) and (3), $E$
  composed with the restriction of scalars along $d_B$ or $\sigma_{B_m}$ is $E$, as $d_B$ and $\sigma_{B_m}$ fix $K$. In
  (4), both composites $H^j_{\mathrm{dR}}(Y'_K)\to H^j(\mathcal Q'_m(Y,s),\mathcal O)$ are induced by the functor
  $T'\mapsto E(T')$ with its point composed with $h_K$, with identity maps.
\end{itemize}

\emph{Products.} The maps $t_S$, $r_n$, $\mathrm{red}_n$ and $b$, and the two maps whose composite is $\mu$, are induced
by functors with maps of coefficients that are ring homomorphisms (identities or projections). The cup product of cochains
is formed from the products of the rings, so these maps preserve cup products, and they send the class of the constant
cochain $1$ to the class of $1$. This gives $t_S(1_{n+1})=1_n$, $\mathrm{red}_n(1_Z)=r_n(1_n)$ and $b(1_Z)=\mu(1_V)$.
% Lean: (1) is homS with homS_good_of_inclHM (rHof_incl_pf through red_s_pf and P.prism, incl_redM, bMap_incl through
% pull_pull, and MuThick.muB_incl). (2) is homLev with Mo.lev_good (nmHM_nmLevA_rHof_succ through red_nm with idMor,
% nmHM_redM, nmHO_nmLevB_bMap through etaLev and P.prism, and nm_muB with pullR_id'). (3) is homGal with
% homGal_good_of_galRed (galRed_rHof through red_nm and PAlg.galBase_cris, the commutation of galH with the reductions,
% nmHO_nmGB_bMap through etaGal, nm_muB). (4) is homPull with homPull_rHof_good (rHof_pull, Model.pullH commuting with the
% reductions, nmHO_nmPullB_bMap through etaPull, nm_muB = MuNat.muB_nat). The cup data are cupOf, and their good
% conditions are cupOf_rHof_good (red_cup, redM_cup, smor_cup, muB_cup, CupAbs.om_cup and om_one). RedNat.red_nm checks
% that the two scalar actions on the groups modulo l^n agree (lam_smul).
\end{proof}
\leanref{Development/WeilBound/CoreStatements/PadicHodgeComparison/EtaleCohomologyComparison/PerfectoidPullbacks/DividedPowerReduction/CoefficientReductionFirstChernClassTateTwist.lean}{CrystallineSemistableComparison.Pb.red_nm}
\leanref{Development/AutomorphicSpectrum/Proofs/BlockDecomposition/PullRetype.lean}{GlobalTheory.CompanionTheory.RedNat.red_nm}
\leanref{Development/WeilBound/CoreStatements/PadicHodgeComparison/EtaleCohomologyComparison/PerfectoidPullbacks/DividedPowerReduction/CoefficientReductionFirstChernClassTateTwist.lean}{CrystallineSemistableComparison.Pb.nmHM_redM}
\leanref{Development/WeilBound/CoreStatements/PadicHodgeComparison/EtaleCohomologyComparison/PerfectoidPullbacks/DividedPowerReduction/ReductionChangeLevel.lean}{CrystallineSemistableComparison.Pb.nm_muB}
\leanref{Development/AutomorphicSpectrum/Proofs/BlockDecomposition/MuBNat.lean}{GlobalTheory.CompanionTheory.MuNat.muB_nat}
\leanref{Development/WeilBound/CoreStatements/PadicHodgeComparison/EtaleCohomologyComparison/PerfectoidPullbacks/DividedPowerReduction/ReductionChangeLevel.lean}{CrystallineSemistableComparison.Pb.muB_cup}
\leanref{Development/WeilBound/CoreStatements/PadicHodgeComparison/EtaleCohomologyComparison/PerfectoidPullbacks/LevelChainValuations/ValuationChainProofs/RHofLev.lean}{CrystallineSemistableComparison.Pb.Kap.nmHM_nmLevA_rHof_succ}
\leanref{Development/WeilBound/CoreStatements/PadicHodgeComparison/EtaleCohomologyComparison/PerfectoidPullbacks/LevelChainValuations/GaloisCupProducts/NmHONmPullBBMap.lean}{CrystallineSemistableComparison.Pb.Kap.nmHO_nmLevB_bMap}
\leanref{Development/WeilBound/CoreStatements/PadicHodgeComparison/EtaleCohomologyComparison/PerfectoidPullbacks/LevelChainValuations/GaloisCupProducts/HomGalGood.lean}{CrystallineSemistableComparison.Pb.Kap.nmHO_nmGB_bMap}
\leanref{Development/WeilBound/CoreStatements/PadicHodgeComparison/EtaleCohomologyComparison/PerfectoidPullbacks/LevelChainValuations/GaloisCupProducts/NmHONmPullBBMap.lean}{CrystallineSemistableComparison.Pb.Kap.nmHO_nmPullB_bMap}
\leanref{Development/WeilBound/CoreStatements/PadicHodgeComparison/EtaleCohomologyComparison/PerfectoidPullbacks/LevelChainValuations/GaloisCupProducts/BMapIncl.lean}{CrystallineSemistableComparison.Pb.Kap.bMap_incl}
\leanref{Development/WeilBound/CoreStatements/PadicHodgeComparison/EtaleCohomologyComparison/PerfectoidPullbacks/LevelChainValuations/GaloisCupProducts/HomPullRHofGood.lean}{CrystallineSemistableComparison.Pb.Kap.homPull_good_of_nmHM_pull}

\begin{lemma}\label{gal-realization-kappa-truncated-periods-exist-chern-log}
Let $\delta$ be divided powers on an ideal $J$ of a ring $B$, and $N\geq0$ an integer such that $k!\,\delta_{k+1}(x)=0$
for all $k\geq N$ and $x\in J$. Then $\log_N((1+x)(1+y)-1)=\log_N(x)+\log_N(y)$ for $x,y\in J$, with $\log_N$ as in
Lemma~\ref{gal-realization-kappa-de-rham-log}.
\end{lemma}
\leanref{Development/AutomorphicSpectrum/Proofs/BlockDecomposition/DividedPowerLogMultiplicativeStatement.lean}{GlobalTheory.T.divided_power_log_multiplicative}
\leanref{Development/AutomorphicSpectrum/Proofs/BlockDecomposition/DividedPowerLogMultiplicative.lean}{GlobalTheory.CompanionTheory.TruncLog.dlog_mul}

\begin{proof}
Let $x,y\in J$. If $N<a\leq k+1$, then $(a-1)!$ divides $k!$ and $a-1\geq N$, so $k!\,\delta_a(x)=0$. As
$x\,\delta_a(x)=(a+1)\,\delta_{a+1}(x)$, induction on $n$ gives
\[
 (1+x)\sum_{a=0}^n(-1)^aa!\,\delta_a(x)=1+(-1)^nn!\,(n+1)\,\delta_{n+1}(x)
\]
and, for $m\geq0$,
\[
 (1+x)\sum_{a=0}^n(-1)^a(a+m+1)!\,\delta_a(x)=(m+1)\sum_{a=0}^n(-1)^a(a+m)!\,\delta_a(x)+(-1)^n(n+m+1)!\,(n+1)\,
 \delta_{n+1}(x).
\]
For $n=N$ the last terms vanish, by the first remark, and induction on $m$ gives
\[
 (1+x)^{m+1}\sum_{a=0}^N(-1)^a(a+m)!\,\delta_a(x)=m!\qquad(m\geq0).\tag{$*$}
\]
Put $z=x+(1+x)y\in J$, so that $(1+x)(1+y)-1=z$. As the terms with $k\geq N$ vanish, $\log_N(z)=\sum_{k<2N}(-1)^kk!\,
\delta_{k+1}(z)$. As $\delta_b((1+x)y)=(1+x)^b\delta_b(y)$,
\[
 (-1)^kk!\,\delta_{k+1}(z)=\sum_{a+b=k+1}g(a,b),\qquad g(a,b)=(-1)^{a+b+1}(a+b-1)!\,\delta_a(x)(1+x)^b\delta_b(y).
\]
If $a>N$ or $b>N$, then $g(a,b)=0$, by the first remark for $x$ or for $y$ with $k=a+b-1$. So $\log_N(z)$ is the sum of the
$g(a,b)$ over $0\leq a,b\leq N$ with $(a,b)\neq(0,0)$. The terms with $b=0$ give $\sum_{a=1}^N(-1)^{a-1}(a-1)!\,\delta_a(x)
=\log_N(x)$. For $1\leq b\leq N$, $(*)$ with $m=b-1$ gives $\sum_{a=0}^Ng(a,b)=(-1)^{b-1}(b-1)!\,\delta_b(y)$, and these
terms add up to $\log_N(y)$.
% Lean: TruncLog.dlog_mul, with factorial_mul_dpow_eq_zero (the first remark), mul_dpow_eq_succ, one_add_mul_sum_zero and
% one_add_mul_sum_succ (the two identities), one_add_pow_mul_sum ((*)), and sum_antidiagonal_eq_sum_box for the regrouping.
\end{proof}
\leanref{Development/AutomorphicSpectrum/Proofs/BlockDecomposition/DividedPowerLogMultiplicative.lean}{GlobalTheory.CompanionTheory.TruncLog.one_add_pow_mul_sum}

\begin{lemma}\label{gal-realization-kappa-truncated-periods-exist-chern-cocycle}
Let $(R,I,\gamma,p)$, $\mathcal O$, $X$ and $\kappa$ be as in Definition~\ref{gal-realization-thickening-nilpotent},
$\mathcal L$ a line datum on $X$ with morphism $f:X_{\mathcal L}\to X$ (Definition~\ref{gal-realization-period-chern}(1)),
$\mathcal C$ a category, and $G:\mathcal C\to\mathrm{Th}(X/R)$ a functor with $G(c)=(B_c,i_c,J_c,\delta_c,y_c)$. Lifts of
$\mathcal L$ along $G$ and their elements $r(s)$ are defined as in Definition~\ref{gal-realization-period-chern}(3). For a
lift $\Lambda$ and $N\geq0$ let $z_N(\Lambda)$ be the $2$-cochain of $\mathcal C$ with values in $c\mapsto B_c$ that sends
a chain $s=(c_0\to c_1\to c_2)$ to $\log_N(r(s)-1)$, formed for $\delta_{c_2}$
(Lemma~\ref{gal-realization-kappa-de-rham-log}).
\begin{enumerate}
\item If for every $c$ there is a morphism $p_c:\operatorname{Spec}B_c/J_c\to X_{\mathcal L}$ with $fp_c=y_c$, then
  $\mathcal L$ has a lift along $G$.
\item Suppose that $k!\,\delta_{c,k+1}(x)=0$ for every $c$, every $k\geq N$ and every $x\in J_c$. Then $z_N(\Lambda)$ is a
  cocycle for every lift $\Lambda$, and $z_N(\Lambda')-z_N(\Lambda)$ is a coboundary for any two lifts $\Lambda$ and
  $\Lambda'$.
\item For a functor $h:\mathcal C'\to\mathcal C$, the morphisms $p_{h(c)}$ and units $w_{h(\varphi)}$ of a lift $\Lambda$
  form a lift $\Lambda h$ of $\mathcal L$ along $Gh$, and $z_N(\Lambda h)$ sends $c_0\to c_1\to c_2$ to
  $z_N(\Lambda)(h(c_0)\to h(c_1)\to h(c_2))$.
\end{enumerate}
\end{lemma}
\leanref{Development/AutomorphicSpectrum/Proofs/BlockDecomposition/ExistsCrystallineLiftStatement.lean}{GlobalTheory.T.exists_crystalline_lift}
\leanref{Development/AutomorphicSpectrum/Proofs/BlockDecomposition/CrystallineCocycleClosedStatement.lean}{GlobalTheory.T.crystalline_cocycle_closed}
\leanref{Development/WeilBound/CoreStatements/PadicHodgeComparison/EtaleCohomologyComparison/ZariskiEtaleLifting/IsFirstChernClassOfLineBundle.lean}{CrystallineSemistableComparison.B.CLift}
\leanref{Development/WeilBound/CoreStatements/PadicHodgeComparison/EtaleCohomologyComparison/ZariskiEtaleLifting/IsFirstChernClassOfLineBundle.lean}{CrystallineSemistableComparison.B.CLift.coc}

\begin{proof}
Each $J_c$ is nilpotent: if $p=0$ by definition, and if $p\neq0$, then $x^{pn}=(p!\,\delta_p(x))^n$ for $x\in J_c$, which is
$0$ once $p^n=0$ in $B_c$. So a unit of $B_c/J_c$ lifts to a unit of $B_c$, as in the proof of
Lemma~\ref{gal-realization-kappa-de-rham-lift}. The argument of Definition~\ref{gal-realization-period-chern}(3) shows that
$r(s)-1\in J_{c_2}$ also here.

(1) For a map $\varphi:c\to c'$, the morphisms $p_c\circ\operatorname{Spec}\overline{G(\varphi)}$ and $p_{c'}$ have the same
composite $y_{c'}$ with $f$; lift the unit $u(p_c\circ\operatorname{Spec}\overline{G(\varphi)},p_{c'})$ of $B_{c'}/J_{c'}$
to a unit $w_\varphi$ of $B_{c'}$.

(2) Under the hypothesis, Lemma~\ref{gal-realization-kappa-truncated-periods-exist-chern-log} gives
$\log_N(ab-1)=\log_N(a-1)+\log_N(b-1)$ for units $a,b$ of $B_c$ with $a-1,b-1\in J_c$. With this identity in place of
Lemma~\ref{gal-realization-kappa-de-rham-log}(3), and with the lifting of units above, the proofs of (1) and (2) of
Lemma~\ref{gal-realization-kappa-de-rham-cocycle} apply without change. They also use that the maps $G(\varphi)$ are
divided power maps, which commute with $\log_N$ (Lemma~\ref{gal-realization-kappa-de-rham-log}(1)).

(3) This holds by the definitions.
% Lean: (1) is CLiftEx.nonempty_cLift (thick_isNilpotent and CocCls.unit_lift). (2) is CocCls.coc_cls: the cocycle
% condition is closed, through rel_closed and dlog_rel (divided_power_log_multiplicative); the independence is exists_v
% (the units v_c), Ef (the inverse of the unit epsilon of the proof of Lemma gal-realization-kappa-de-rham-cocycle(2)),
% rel_ind, divided_power_log_three_term_relation and d_yc, which gives d(yc) = z_N(Λ) - z_N(Λ'). (3) is
% CrystallineFirstChernTorsionFree.precomp, which holds by rfl.
\end{proof}
\leanref{Development/AutomorphicSpectrum/Proofs/BlockDecomposition/ExistsCrystallineLift.lean}{GlobalTheory.CompanionTheory.CLiftEx.nonempty_cLift}
\leanref{Development/AutomorphicSpectrum/Proofs/BlockDecomposition/CrystallineCocycleClosed.lean}{GlobalTheory.CompanionTheory.CocCls.coc_cls}
\leanref{Development/AutomorphicSpectrum/Proofs/BlockDecomposition/CrystallineCocycleClosed.lean}{GlobalTheory.CompanionTheory.CocCls.closed}
\leanref{Development/AutomorphicSpectrum/Proofs/BlockDecomposition/CrystallineCocycleClosed.lean}{GlobalTheory.CompanionTheory.CocCls.d_yc}
\leanref{Development/AutomorphicSpectrum/Proofs/BlockDecomposition/ExistsUniqueIsFirstChernTorsionFree.lean}{GlobalTheory.CompanionTheory.CrystallineFirstChernTorsionFree.precomp}

\begin{definition}\label{gal-realization-kappa-truncated-periods-exist-chern-class}
Let $(R,I,\gamma,p)$, $\mathcal O$, $X$ and $\kappa$ be as in Definition~\ref{gal-realization-thickening-nilpotent},
$\mathcal L$ a line datum on $X$ with morphism $f:X_{\mathcal L}\to X$, $\mathcal Q$ a full subcategory of
$\mathrm{Th}(X/R)$, and $N\geq0$. A class $z\in H^2(\mathcal Q,\mathcal O)$ is a \emph{first Chern class of $\mathcal L$
of order $N$} if the following holds for every category $\mathcal C$ and functor $G:\mathcal C\to\mathcal Q$ such that
every $G(c)$ lies over an open subset of $X$ over which $f$ has a section: for every lift $\Lambda$ of $\mathcal L$ along
$G$, the cochain $z_N(\Lambda)$ of Lemma~\ref{gal-realization-kappa-truncated-periods-exist-chern-cocycle} is a cocycle whose
class is $G^*z$ (Definition~\ref{gal-realization-split-diagrams}(1), with identity maps).
\end{definition}
\leanref{Development/WeilBound/CoreStatements/PadicHodgeComparison/EtaleCohomologyComparison/PerfectoidPullbacks/DividedPowerReduction/IsFirstChernThickening.lean}{CrystallineSemistableComparison.Pb.IsFirstChernThickening}

\begin{lemma}\label{gal-realization-kappa-truncated-periods-exist-chern-unique}
Let $(R,I,\gamma,0)$, $\mathcal O$, $X$ and $\kappa$ be as in Definition~\ref{gal-realization-kappa-de-rham-truncated},
$s,\ell\geq1$ integers, and $\mathcal Q=\mathrm{Th}^{(s)}_\ell(X/R)$
(Lemma~\ref{gal-realization-kappa-truncated-periods-exist-base-change-zariski}). Let $\mathcal L$ be a line datum on $X$
such that $X$ is the union of the open subsets over which its morphism $f$ has a section. Then $\mathcal L$ has exactly one
first Chern class of order $s$ in $H^2(\mathcal Q,\mathcal O)$
(Definition~\ref{gal-realization-kappa-truncated-periods-exist-chern-class}).
\end{lemma}
\leanref{Development/AutomorphicSpectrum/Proofs/BlockDecomposition/ExistsUniqueIsFirstChernTorsionFree.lean}{GlobalTheory.CompanionTheory.CrystallineFirstChernTorsionFree.existsUnique_isFirstChernTorsionFree}

\begin{proof}
Let $\mathcal V$ be the set of open subsets of $X$ over which $f$ has a section, $\mathcal Q'=\mathcal Q_{\mathcal V}$ and
$\iota:\mathcal Q'\to\mathcal Q$ the inclusion. Every thickening lies over $X$
(Lemma~\ref{gal-realization-thickening-over}(2)), so $\mathcal Q=\mathcal Q_{\{X\}}$, and by
Lemma~\ref{gal-realization-kappa-truncated-periods-exist-base-change-zariski} with
Lemma~\ref{gal-realization-thickening-zariski}(1), for $\mathcal W=\{X\}$, $\iota^*:H^2(\mathcal Q,\mathcal O)\to
H^2(\mathcal Q',\mathcal O)$ is bijective. For $T$ in $\mathcal Q$, $k\geq s$ and $x\in J_T$, $k!\,\delta_{T,k+1}(x)$ lies in
$J_T^{[k+1]}\subset J_T^{[s]}=0$; so Lemma~\ref{gal-realization-kappa-truncated-periods-exist-chern-cocycle}(2) applies to
every functor to $\mathcal Q$, with $N=s$.

A thickening $T$ of $\mathcal Q'$ lies over some $V\in\mathcal V$, with a section $\sigma$; its point is a morphism
$y':\operatorname{Spec}B_T/J_T\to V$ followed by the inclusion, and $p_T=\sigma y'$ has $fp_T=y_T$. By
Lemma~\ref{gal-realization-kappa-truncated-periods-exist-chern-cocycle}(1) there is a lift $\Lambda_0$ of $\mathcal L$ along
$\iota$, viewed as a functor to $\mathrm{Th}(X/R)$. Then $z_s(\Lambda_0)$ is a cocycle, and there is exactly one $z\in
H^2(\mathcal Q,\mathcal O)$ such that $\iota^*z$ is its class.

\emph{Existence.} Let $G:\mathcal C\to\mathcal Q$ be a functor such that every $G(c)$ lies over a member of $\mathcal V$,
and $\Lambda$ a lift of $\mathcal L$ along $G$. Then $G=\iota G'$ for a functor $G':\mathcal C\to\mathcal Q'$, and $\Lambda_0
G'$ is a lift along $G$ with $z_s(\Lambda_0G')=G'^*z_s(\Lambda_0)$
(Lemma~\ref{gal-realization-kappa-truncated-periods-exist-chern-cocycle}(3)). By
Lemma~\ref{gal-realization-kappa-truncated-periods-exist-chern-cocycle}(2), $z_s(\Lambda)$ is a cocycle with the class of
$z_s(\Lambda_0G')$, which is $G'^*\iota^*z=G^*z$.

\emph{Uniqueness.} If $z'$ is a first Chern class of $\mathcal L$ of order $s$, then its definition, for $G=\iota$ and
$\Lambda_0$, shows that $\iota^*z'$ is the class of $z_s(\Lambda_0)$, which is $\iota^*z$. So $z'=z$.
% Lean: CrystallineFirstChernTorsionFree.existsUnique_isFirstChernTorsionFree: Q' = St (TF l s) (𝒱D D), hb from
% cech_refinement_cohomology_bijective, Λ₀ from exists_crystalline_lift with lifts_of_mem, dpow_vanish for the hypothesis
% of crystalline_cocycle_closed, and precomp for Λ₀G'.
\end{proof}
\leanref{Development/AutomorphicSpectrum/Proofs/BlockDecomposition/ExistsUniqueIsFirstChernTorsionFree.lean}{GlobalTheory.CompanionTheory.CrystallineFirstChernTorsionFree.dpow_vanish}

\begin{definition}\label{gal-realization-kappa-truncated-periods-exist-chern-reduction-functor}
Let $\mathcal O$ and $R$ be rings, $A$ an $\mathcal O$-algebra with divided powers $\gamma$ on an ideal $I$, $\delta$
divided powers on an ideal $J$ of $R$, $\kappa:\mathcal O\to R/J$ a ring map, and $\ell\geq1$ and $n\geq0$ integers. Write
$\bar y$ for the class in $A/\ell^nA$ of $y\in A$, and let $\rho:R\to A/\ell^nA$ be a ring map such that:
\begin{enumerate}
\item[(a)] for $x\in J$ there is $y\in I$ with $\rho(x)=\bar y$;
\item[(b)] if $x\in J$ and $y\in I$ satisfy $\rho(x)=\bar y$, then $\rho(\delta_k(x))=\overline{\gamma_k(y)}$ for every
  $k\geq0$;
\item[(c)] if $a\in\mathcal O$ and $r\in R$ satisfy $\kappa(a)=r+J$, then $\rho(r)=\overline{a+y}$ for some $y\in I$.
\end{enumerate}
Let $X$ be an $\mathcal O$-scheme, $s\geq1$, $\mathcal Q=\mathrm{Th}^{(s)}_\ell(X/A)$, formed for $(A,I,\gamma,0)$ and the
structure map $\mathcal O\to A/I$ (Lemma~\ref{gal-realization-kappa-truncated-periods-exist-base-change-zariski}), and
$\mathrm{Th}(X/R)$ formed for $(R,J,\delta,\ell)$ and $\kappa$ (Definition~\ref{gal-realization-thickening-nilpotent}). Let
$T=(B,i,J_T,\delta_T,y)$ be in $\mathcal Q$, $\bar B=B/\ell^nB$ and $\pi:B\to\bar B$ the projection. If $\ell^nb\in J_T$,
then $b\in J_T$, by the condition defining $\mathcal Q$ for $r=1$, applied $n$ times, as $J_T^{[1]}=J_T$; and
$\delta_{T,k}(\ell^nb)=\ell^{nk}\delta_{T,k}(b)\in\ell^nB$ for $k\geq1$. So $\delta_T$ induces divided powers $\bar\delta_T$
on $\bar J_T=\pi(J_T)$ for which $\pi$ is a divided power map \cite[Tag 07H2]{Sta}. The \emph{reduction} of $T$ along
$\rho$ is $\mathrm{Red}_\rho(T)=(\bar B,\bar\imath\rho,\bar J_T,\bar\delta_T,\bar y)$, where $\bar\imath:A/\ell^nA\to\bar B$
is induced by $i$ and $\bar y$ is $\operatorname{Spec}\bar B/\bar J_T\to\operatorname{Spec}B/J_T$ followed by $y$. By (a)
and (b), $\bar\imath\rho$ carries $J$ into $\bar J_T$ compatibly with the divided powers; by (c), the $\mathcal
O$-algebra structures of $\bar B/\bar J_T$ through $\kappa$ and through $A/I$ agree; and $\ell^n=0$ in $\bar B$. So
$\mathrm{Red}_\rho(T)$ is a thickening of $X$ over $(R,J,\delta,\ell)$. Maps $T\to T'$ induce maps
$\mathrm{Red}_\rho(T)\to\mathrm{Red}_\rho(T')$, and $\mathrm{Red}_\rho$ is a functor $\mathcal Q\to\mathrm{Th}(X/R)$.
\end{definition}
\leanref{Development/WeilBound/CoreStatements/PadicHodgeComparison/EtaleCohomologyComparison/PerfectoidPullbacks/DividedPowerReduction/JEqPdPow.lean}{CrystallineSemistableComparison.Pb.RedIn}
\leanref{Development/WeilBound/CoreStatements/PadicHodgeComparison/EtaleCohomologyComparison/PerfectoidPullbacks/DividedPowerReduction/JEqPdPow.lean}{CrystallineSemistableComparison.Pb.redObj}
\leanref{Development/WeilBound/CoreStatements/PadicHodgeComparison/EtaleCohomologyComparison/PerfectoidPullbacks/DividedPowerReduction/JEqPdPow.lean}{CrystallineSemistableComparison.Pb.redG}
\leanref{Development/WeilBound/CoreStatements/PadicHodgeComparison/EtaleCohomologyComparison/PerfectoidPullbacks/DividedPowerReduction/JEqPdPow.lean}{CrystallineSemistableComparison.Pb.red_compat}
\leanref{Development/WeilBound/CoreStatements/PadicHodgeComparison/EtaleCohomologyComparison/PerfectoidPullbacks/DividedPowerReduction/JEqPdPow.lean}{CrystallineSemistableComparison.Pb.ψBar_κ}

\begin{lemma}\label{gal-realization-kappa-truncated-periods-exist-chern-reduction}
Let $\mathcal O$, $A$, $I$, $\gamma$, $R$, $J$, $\delta$, $\kappa$, $\ell\geq1$, $n\geq0$, $\rho$, $X$, $s\geq1$, $\mathcal
Q=\mathrm{Th}^{(s)}_\ell(X/A)$, $\mathrm{Th}(X/R)$ and $\mathrm{Red}_\rho$ be as in
Definition~\ref{gal-realization-kappa-truncated-periods-exist-chern-reduction-functor}. Let $\Lambda$ be a ring and
$c:\Lambda\to R$ and $\lambda:\Lambda\to A$ ring maps with $\rho(c(a))=\overline{\lambda(a)}$ for $a\in\Lambda$. Through
$c$, $H^2(\mathrm{Th}(X/R),\mathcal O)$ is a $\Lambda$-module, and the map $\mathrm{Red}_\rho^*:H^2(\mathrm{Th}(X/R),\mathcal
O)\to H^2(\mathcal Q,\mathcal O/\ell^n\mathcal O)$ induced by $\mathrm{Red}_\rho$ with the identity maps of the rings
$B/\ell^nB$ is $\lambda$-semilinear, as $\bar\imath(\rho(c(a)))=\bar\imath(\overline{\lambda(a)})$. Let $\mathcal L$ be a line
datum on $X$ such that $X$ is the union of the open subsets over which its morphism $f$ has a section. Let $x\in
H^2(\mathrm{Th}(X/R),\mathcal O)$ be such that for every category $\mathcal C$, functor $G:\mathcal C\to\mathrm{Th}(X/R)$
and lift $\Pi$ of $\mathcal L$ along $G$, the cochain $z_{\ell n}(\Pi)$ of
Lemma~\ref{gal-realization-kappa-truncated-periods-exist-chern-cocycle} is a cocycle whose class is $G^*x$; and let $z\in
H^2(\mathcal Q,\mathcal O)$ be a first Chern class of $\mathcal L$ of order $s$
(Definition~\ref{gal-realization-kappa-truncated-periods-exist-chern-class}). Then the image of $z$ in $H^2(\mathcal
Q,\mathcal O/\ell^n\mathcal O)$ is $\mathrm{Red}_\rho^*(x)$.
\end{lemma}
\leanref{Development/AutomorphicSpectrum/Proofs/BlockDecomposition/CoefficientReductionFirstChernClassTateTwist.lean}{GlobalTheory.CompanionTheory.ReductionFirstChern.coefficientReduction_firstChernClass_tateTwist}

\begin{proof}
Let $\mathcal V$ be the set of open subsets of $X$ over which $f$ has a section, $\mathcal Q'=\mathcal Q_{\mathcal V}$ and
$\iota:\mathcal Q'\to\mathcal Q$ the inclusion. The functor $\mathcal O/\ell^n\mathcal O$ is $T\mapsto(J_T^{[0]}+\ell^nB_T)/
\ell^nB_T$ (Definition~\ref{gal-realization-thickening-nilpotent}(3)), so, as in the proof of
Lemma~\ref{gal-realization-kappa-truncated-periods-exist-chern-unique}, $\iota^*$ is injective on $H^2(\cdot,\mathcal
O/\ell^n\mathcal O)$, and there is a lift $\Pi=(p,w)$ of $\mathcal L$ along $\iota$. By the definition of $z$,
$z_s(\Pi)$ is a cocycle whose class is $\iota^*z$.

\emph{The reduced lift.} For $T=(B,i,J_T,\delta_T,y)$ in $\mathcal Q'$ let $\mathrm{Red}_\rho(T)=(\bar B,\ldots,\bar
J_T,\bar\delta_T,\bar y)$ and $\pi:B\to\bar B$ be as in
Definition~\ref{gal-realization-kappa-truncated-periods-exist-chern-reduction-functor}, and $\pi_0:\operatorname{Spec}\bar
B/\bar J_T\to\operatorname{Spec}B/J_T$ the morphism induced by $\pi$, so that $\bar y=y\pi_0$. Put $\bar p_T=p_T\pi_0$ and
$\bar w_\varphi=\pi(w_\varphi)$ for the maps $\varphi$ of $\mathcal Q'$. Then $f\bar p_T=\bar y$, and, by the naturality
of $u$, the image of $\bar w_\varphi$ in $\bar B_{T'}/\bar J_{T'}$ is $u(\bar p_T\circ\operatorname{Spec}
\overline{\mathrm{Red}_\rho(\varphi)},\bar p_{T'})$, as $\mathrm{Red}_\rho(\varphi)$ is induced by $\varphi$. So
$\bar\Pi=(\bar p,\bar w)$ is a lift of $\mathcal L$ along $\mathrm{Red}_\rho\iota:\mathcal Q'\to\mathrm{Th}(X/R)$, and its
elements are the $\bar r(\sigma)=\pi(r(\sigma))$ for the chains $\sigma$ of $\mathcal Q'$. By the hypothesis on $x$, the
cochain $z_{\ell n}(\bar\Pi)$, formed for the $\bar\delta_T$, is a cocycle whose class is
$(\mathrm{Red}_\rho\iota)^*x=\iota^*\mathrm{Red}_\rho^*x$.

\emph{The cochains.} For a chain $\sigma$ of $\mathcal Q'$ ending at $T$,
$\pi(\log_s(r(\sigma)-1))=\log_{\ell n}(\bar r(\sigma)-1)$. Let $M=\max(s,\ell n)$. As $J_T^{[s]}=0$,
$\log_s(r(\sigma)-1)=\log_M(r(\sigma)-1)$ (Lemma~\ref{gal-realization-kappa-de-rham-log}(2)); as $\pi$ is a divided power
map, $\pi(\log_M(r(\sigma)-1))=\log_M(\bar r(\sigma)-1)$ (Lemma~\ref{gal-realization-kappa-de-rham-log}(1)); and as
$\ell^n$ divides $k!$ for $k\geq\ell n$, through the factors $\ell,2\ell,\ldots,n\ell$ of $(\ell n)!$, and $\ell^n\bar B=0$,
the terms of $\log_M(\bar r(\sigma)-1)$ with $k\geq\ell n$ vanish. So the reduction of $z_s(\Pi)$ modulo $\ell^n$ is
$z_{\ell n}(\bar\Pi)$, and $\iota^*$ of the image of $z$ in $H^2(\mathcal Q,\mathcal O/\ell^n\mathcal O)$ is the class of
$z_{\ell n}(\bar\Pi)$, which is $\iota^*\mathrm{Red}_\rho^*x$. As $\iota^*$ is injective, the claim follows.
% Lean: ReductionFirstChern.coefficientReduction_firstChernClass_tateTwist, for RedIn data d (the book's rho with (a) to
% (c)), c and lam with hc, x with B.IsFirstChernClassOfLineBundle (nlBase R J δ l) D c (l * n) x and z with
% IsFirstChernThickening (TS γ q l s {⊤}) D s z; Red_rho is redG and Red_rho^* is (redMor q s d {⊤} c lam hc none).pull.
% hinj from cech_refinement_cohomology_bijective (with bij_top and stdI_top), Lf from exists_crystalline_lift, redLift
% for the reduced lift (with uq_map), coc_red with r_red and dlog_red (dlog_trunc, dlog_map and
% FirstChernClassFactorial.dlog_eq_of_factorial, with factorial_quot and pow_dvd_factorial_mul), and the comparison of
% classes through m₁, m₂ and P.mapH_cls. The Lean allows s = 0, where every thickening of Q has B = 0 and both sides
% vanish; the book's Th^(s)_l is defined for s >= 1.
\end{proof}
\leanref{Development/AutomorphicSpectrum/Proofs/BlockDecomposition/CoefficientReductionFirstChernClassTateTwist.lean}{GlobalTheory.CompanionTheory.ReductionFirstChern.redLift}
\leanref{Development/AutomorphicSpectrum/Proofs/BlockDecomposition/CoefficientReductionFirstChernClassTateTwist.lean}{GlobalTheory.CompanionTheory.ReductionFirstChern.coc_red}
\leanref{Development/AutomorphicSpectrum/Proofs/BlockDecomposition/CoefficientReductionFirstChernClassTateTwist.lean}{GlobalTheory.CompanionTheory.ReductionFirstChern.dlog_red}

\begin{lemma}\label{gal-realization-kappa-truncated-periods-exist-chern-pullback}
Let $(R,I,\gamma,p)$, $\mathcal O$, $X$ and $\kappa$, and $(R',I',\gamma',p)$, $\mathcal O'$, $X'$ and $\kappa'$, be as in
Definition~\ref{gal-realization-thickening-nilpotent}, with the same $p$. Let $\sigma:R\to R'$ be a ring map with
$\sigma(I)\subset I'$ and $\gamma'_k(\sigma(x))=\sigma(\gamma_k(x))$ for $x\in I$, $g:\mathcal O\to\mathcal O'$ a ring map
such that $\kappa'(g(a))$ is the image of $\kappa(a)$ under the map $R/I\to R'/I'$ induced by $\sigma$ for every $a\in\mathcal
O$, and $h:X'\to X$ a morphism whose composite with $X\to\operatorname{Spec}\mathcal O$ is $X'\to\operatorname{Spec}\mathcal
O'\xrightarrow{\operatorname{Spec}g}\operatorname{Spec}\mathcal O$. Sending $(B,i,J,\delta,y)$ to $(B,i\circ\sigma,J,\delta,
h\circ y)$, and a map to the same ring map, defines a functor $M:\mathrm{Th}(X'/R')\to\mathrm{Th}(X/R)$. Let $\mathcal Q$ be
a full subcategory of $\mathrm{Th}(X/R)$, and $\mathcal Q_0$ a full subcategory of $\mathrm{Th}(X'/R')$ that $M$ maps into
$\mathcal Q$, such that either $p=0$ and $\mathcal Q_0=\mathrm{Th}^{(s)}_\ell(X'/R')$ for integers $s,\ell\geq1$
(Lemma~\ref{gal-realization-kappa-truncated-periods-exist-base-change-zariski}), or $p\neq0$ and $\mathcal Q_0$ is
$\mathrm{Th}(X'/R')$ or $\mathrm{Th}^{\mathrm{fl}}(X'/R')$. Let $\mathcal L$ be a line datum on $X$ such that $X$ is the
union of the open subsets over which its morphism $f$ has a section, $\mathcal L'$ its base change along $h$
(Definition~\ref{gal-realization-period-chern}(1)), and $N\geq0$. If $z\in H^2(\mathcal Q,\mathcal O)$ is a first Chern
class of $\mathcal L$ of order $N$ and $z'\in H^2(\mathcal Q_0,\mathcal O)$ one of $\mathcal L'$ of order $N$
(Definition~\ref{gal-realization-kappa-truncated-periods-exist-chern-class}), then $M^*z=z'$, where $M^*$ is induced by $M$
with the identity maps of the rings.
\end{lemma}
\leanref{Development/AutomorphicSpectrum/Proofs/BlockDecomposition/FirstChernPullback.lean}{GlobalTheory.CompanionTheory.FirstChernPullback.firstChern_pullback}

\begin{proof}
\emph{The functor.} Let $T=(B,i,J,\delta,y)$ be in $\mathrm{Th}(X'/R')$. Then $i\sigma$ carries $I$ into $i(I')\subset J$,
and $\delta_k(i\sigma(x))=i(\gamma'_k(\sigma(x)))=i\sigma(\gamma_k(x))$ for $x\in I$. The composite of $h\circ y$ with
$X\to\operatorname{Spec}\mathcal O$ is $y$ followed by $X'\to\operatorname{Spec}\mathcal O'\to\operatorname{Spec}\mathcal
O$, which is $\operatorname{Spec}$ of $\mathcal O\xrightarrow{g}\mathcal O'\xrightarrow{\kappa'}R'/I'\to B/J$; by the
condition on $g$ this is $\mathcal O\xrightarrow{\kappa}R/I\to B/J$, through $i\sigma$, so $h\circ y$ is an $\mathcal
O$-morphism. The conditions on $p$ and $J$ concern $B$ and $J$ only. So $M(T)$ is a thickening of $X$ over $(R,I,\gamma,p)$.
A map $\varphi:T\to T'$ satisfies $\varphi\circ i\sigma=i'\sigma$ and $h\circ y'=h\circ y\circ\operatorname{Spec}
\bar\varphi$, so it is a map $M(T)\to M(T')$.

\emph{The classes.} The morphism of $\mathcal L'$ is $f':X_{\mathcal L}\times_XX'\to X'$; let $q:X_{\mathcal
L}\times_XX'\to X_{\mathcal L}$ be the projection. Let $\mathcal V$ be the set of open subsets of $X$ over which $f$ has a
section, and $\mathcal V'=\{h^{-1}(V):V\in\mathcal V\}$. A section $\tau$ of $f$ over $V$ gives the section of $f'$ over
$h^{-1}(V)$ with components $\tau\circ h$ and the inclusion, and the members of $\mathcal V'$ cover $X'$, as those of
$\mathcal V$ cover $X$. Let $\mathcal Q'_0=(\mathcal Q_0)_{\mathcal V'}$ and $\iota:\mathcal Q'_0\to\mathcal Q_0$ the
inclusion. Every thickening lies over $X'$ (Lemma~\ref{gal-realization-thickening-over}(2)), so $\mathcal Q_0=(\mathcal
Q_0)_{\{X'\}}$, and $\iota^*:H^2(\mathcal Q_0,\mathcal O)\to H^2(\mathcal Q'_0,\mathcal O)$ is bijective by
Lemma~\ref{gal-realization-thickening-zariski}(1) for $\mathcal W=\{X'\}$, which holds for $\mathcal Q_0$ if $p\neq0$ and,
by Lemma~\ref{gal-realization-kappa-truncated-periods-exist-base-change-zariski}, if $p=0$. As in the proof of
Lemma~\ref{gal-realization-kappa-truncated-periods-exist-chern-unique}, there is a lift $\Lambda'=(p',w')$ of $\mathcal L'$
along $\iota$; by the definition of $z'$, $z_N(\Lambda')$ is a cocycle whose class is $\iota^*z'$.

For $T$ in $\mathcal Q'_0$, $M(T)$ has the ring, ideal and divided powers of $T$ and the point $h\circ y_T$, and
$f\circ qp'_T=h\circ f'p'_T=h\circ y_T$. The units of $\mathcal L'$ are $u'(a,b)=u(qa,qb)$, so the $qp'_T$ and the units
$w'_\varphi$ form a lift $\Lambda$ of $\mathcal L$ along $M\iota:\mathcal Q'_0\to\mathcal Q$. Every $M(T)$ lies over a member
of $\mathcal V$: if $y_T$ factors through $h^{-1}(V)$, then $h\circ y_T$ factors through $V$. By the definition of $z$,
$z_N(\Lambda)$ is a cocycle whose class is $(M\iota)^*z=\iota^*M^*z$. The cochains $z_N(\Lambda)$ and $z_N(\Lambda')$ are
equal, as they are formed from the same units in the same rings with the same divided powers. So $\iota^*M^*z=\iota^*z'$,
and $M^*z=z'$.
% Lean: FirstChernPullback.firstChern_pullback, for an NMor m with fields hI (the condition on sigma), u (the book's h), g,
% hu and hκ; M is m.mor p (NMor.mor, whose pt field is the check of the first paragraph), M^* is
% (m.smor Q (St Q₀ {⊤}) h).o.pull, and Used Q₀ is the condition on Q₀; the category of z' is St Q₀ {⊤}, which is Q₀. 𝒱'
% is the image of 𝒱D D under preimage (sSup_preimage, mem_VD_pb), the injectivity is cech_refinement_cohomology_bijective,
% Λ' is from exists_crystalline_lift, Λ has p = Λ'.p ≫ pullback.fst and w = Λ'.w (u_congr), and hE, which holds by rfl,
% is the equality of the cochains. Used Q₀ also allows p = 0 with Q₀ = TF l s for l = 0 or s = 0; then every thickening
% of Q₀ has B = 0 and both sides vanish, and the book's Th^(s)_l is defined for s, l >= 1.
\end{proof}
\leanref{Development/AutomorphicSpectrum/Proofs/BlockDecomposition/FirstChernPullback.lean}{GlobalTheory.CompanionTheory.FirstChernPullback.mem_VD_pb}

\begin{lemma}\label{gal-realization-kappa-truncated-periods-exist-chern-de-rham}
Let $K$ be a field of characteristic $0$, $B$ a $K$-algebra with divided powers $\delta$ on an ideal $J$, $s\geq1$, $X$ a
$K$-scheme, and $\mathcal Q=\mathrm{Th}^{(s)}(X/B)$, formed for $(B,J,\delta,0)$, $\mathcal O=K$ and the structure map
$K\to B/J$ (Definition~\ref{gal-realization-kappa-de-rham-truncated}). Let $E:\mathcal
Q\to\mathrm{Th}^{(s)}_{\mathrm{dR}}(X)$ be the restriction of scalars along $K\to B$
(Definition~\ref{gal-realization-thickening-nilpotent}(4)), and $\mu:H^2_{\mathrm{dR}}(X)\to H^2(\mathcal Q,\mathcal O)$ the
restriction to $\mathrm{Th}^{(s)}_{\mathrm{dR}}(X)$ followed by $E^*$, with identity maps. Let $\mathcal L$ be a line datum
on $X$ and $\omega\in H^2_{\mathrm{dR}}(X)$ a first Chern class of $\mathcal L$
(Definition~\ref{gal-realization-period-chern}(6)). Then $\mu(\omega)$ is a first Chern class of $\mathcal L$ of order $s$
(Definition~\ref{gal-realization-kappa-truncated-periods-exist-chern-class}).
\end{lemma}
\leanref{Development/AutomorphicSpectrum/Proofs/BlockDecomposition/IsFirstChernThickeningMuBOfIsFirstChernDeRham.lean}{GlobalTheory.CompanionTheory.FirstChernRationalLift.isFirstChernThickening_muB_of_isFirstChernDeRham}

\begin{proof}
Let $G:\mathcal C\to\mathcal Q$ be a functor such that every $G(c)$ lies over an open subset over which the morphism of
$\mathcal L$ has a section, and $\Lambda$ a lift of $\mathcal L$ along $G$. The functor $E$ keeps the rings, ideals, divided
powers and points, so the $EG(c)$ have $J^{[s]}=0$ and lie over such open subsets, and $\Lambda$ is a lift of $\mathcal L$
along $EG$, viewed as a functor to $\mathrm{Th}_{\mathrm{dR}}(X)$, with the same units and elements $r$. By
Definition~\ref{gal-realization-period-chern}(6) with $m=s$, $z_s(\Lambda)$ is a cocycle whose class is $(EG)^*\omega$, and
this is $G^*\mu(\omega)$, as both are induced by the same functor with identity maps.
% Lean: FirstChernRationalLift.isFirstChernThickening_muB_of_isFirstChernDeRham: famK G is EG, liftK Λ the lift along it,
% the field of IsFirstChernDeRham at s gives the cocycle, and e2 (P.pull_comp) identifies (EG)^*ω with G^*μ(ω).
\end{proof}
\leanref{Development/AutomorphicSpectrum/Proofs/BlockDecomposition/IsFirstChernThickeningMuBOfIsFirstChernDeRham.lean}{GlobalTheory.CompanionTheory.FirstChernRationalLift.liftK}

\begin{lemma}\label{gal-realization-kappa-truncated-periods-exist-chern}
Let $\mathcal D_m(Y,s)$ be as in Definition~\ref{gal-realization-kappa-truncated-periods-exist-categories}, for an
approximation system, a smooth projective $\mathcal O_K$-scheme $Y$, $m\geq0$ and $s\geq m+1$. Suppose that $Y_K$ has a
coordinate cover $(U_\alpha)$ of dimensions $d_\alpha$ (Definition~\ref{gal-realization-thickening-coordinates}(2)) with
$d_\alpha+m+1\leq s$ for every $\alpha$. Let $\mathcal L$ be a tautological line datum on $Y$
(Definition~\ref{gal-realization-period-chern}(2)), $x\in H^2_\infty(Y)$ an element whose components form a first Chern
family of $\mathcal L$, and $\omega\in H^2_{\mathrm{dR}}(Y_K)$ a first Chern class of the base change of $\mathcal L$ to
$Y_K$ (Definition~\ref{gal-realization-period-chern}(4) and~(6)). Then $\psi\mapsto\psi(\omega)$ is a period of the
components $(x_\nu)_\nu$ of $x$ in $\mathcal D_m(Y,s)$
(Definition~\ref{gal-realization-kappa-truncated-periods-exist-diagram}(3)).
\end{lemma}
\leanref{Development/WeilBound/CoreStatements/PadicHodgeComparison/EtaleCohomologyComparison/PerfectoidPullbacks/LevelChainValuations/ValuationChainProofs/CrystallineFirstChernLevelCompat.lean}{CrystallineSemistableComparison.Pb.Kap.crystalline_firstChern_level}

\begin{proof}
The map $f:\psi\mapsto\psi(\omega)$ is additive, and $f(a\psi)=af(\psi)$ for $a\in K$. Take $e=0$ in
Definition~\ref{gal-realization-kappa-truncated-periods-exist-diagram}(3). For a finite basis $(v_t)$ of
$H^2_{\mathrm{dR}}(Y_K)$, $\omega=\sum_tv_t^*(\omega)v_t$, so $\sum_tf(v_t^*)\,\mu(v_t)=\mu(\omega)$. So it suffices to
find $z\in H^2(\mathcal Q_m(Y,s),\mathcal O)$ with $\mathrm{red}_\nu(z)=r_\nu(x_\nu)$ for every $\nu$ and $b(z)=\mu(\omega)$.

As $\mathcal L$ is tautological, $Y$ is the union of the open subschemes over which its morphism $f$ has a section
(Lemma~\ref{gal-realization-ingredients-tautological}(4)). The category $\mathcal Q_m(Y,s)$ is
$\mathrm{Th}^{(s)}_\ell(Y/A_m)$, formed for $(A_m,I_m,\gamma_m,0)$
(Lemma~\ref{gal-realization-kappa-truncated-periods-exist-base-change-zariski}), so by
Lemma~\ref{gal-realization-kappa-truncated-periods-exist-chern-unique} $\mathcal L$ has a first Chern class $z\in
H^2(\mathcal Q_m(Y,s),\mathcal O)$ of order $s$.

Let $\nu\geq1$. Lemma~\ref{gal-realization-kappa-truncated-periods-exist-chern-reduction} applies to $\mathcal O=\mathcal
O_K$, $A=A_m$ with $\gamma_m$ on $I_m$, $R=A_\nu$ with $\gamma$ on $I_\nu$, $\kappa:\mathcal O_K\to A_\nu/I_\nu$ induced by
$u$, $n=\nu$, $\rho=\rho_{m,\nu}$, $X=Y$, $\Lambda=A_{\mathrm{cris}}$, $c$ the projection $A_{\mathrm{cris}}\to A_\nu$ and
$\lambda=\lambda_m$: conditions (a) to (c) of
Definition~\ref{gal-realization-kappa-truncated-periods-exist-chern-reduction-functor} are (1) and (2) of
Lemma~\ref{gal-realization-kappa-truncated-periods-exist-reduction}, which also gives $\rho_{m,\nu}(c(a))=
\overline{\lambda_m(a)}$. For these data $\mathrm{Th}(Y/A_\nu)$, formed for $(A_\nu,I_\nu,\gamma,\ell)$, is
$\mathrm{Th}_\nu(Y)$, as $\ell^\nu=0$ in $A_\nu$, and $\mathrm{Red}_\rho$ is $R_\nu$
(Definition~\ref{gal-realization-kappa-truncated-periods-exist-categories}(2)). As $x_\nu$ is a first Chern class of
$\mathcal L$ in $H^2_\nu(Y)$ (Definition~\ref{gal-realization-period-chern}(4)), it satisfies the hypothesis of the lemma,
which gives $\mathrm{red}_\nu(z)=R_\nu^*(x_\nu)=r_\nu(x_\nu)$.

The ring $B_m$ is a $K$-algebra with the divided powers $\delta_m$ on $J_m$, and $\mathcal Q'_m(Y,s)$ and $\mu$ are as in
Lemma~\ref{gal-realization-kappa-truncated-periods-exist-chern-de-rham} for $B=B_m$ and $X=Y_K$; so $\mu(\omega)$ is a first
Chern class of order $s$ of the base change $\mathcal L_K$ of $\mathcal L$ to $Y_K$.
Lemma~\ref{gal-realization-kappa-truncated-periods-exist-chern-pullback} applies to $p=0$, $(R,I,\gamma)=(A_m,I_m,\gamma_m)$,
$\mathcal O=\mathcal O_K$, $X=Y$, $(R',I',\gamma')=(B_m,J_m,\delta_m)$, $\mathcal O'=K$, $X'=Y_K$, $\sigma=\iota_m$, $g$ the
inclusion $\mathcal O_K\to K$, $h=\mathrm{pr}$ and $N=s$: $\iota_m$ is a divided power map, and if $a\in\mathcal O_K$ and
$r\in A_m$ satisfy $a-r\in I_m$, then $a-\iota_m(r)=\iota_m(a-r)\in J_m$. The functor $M$ is then $N$, which maps $\mathcal
Q'_m(Y,s)=\mathrm{Th}^{(s)}_1(Y_K/B_m)$ into $\mathcal Q_m(Y,s)$
(Definition~\ref{gal-realization-kappa-truncated-periods-exist-categories}(3)). So $b(z)=N^*z=\mu(\omega)$.
% Lean: crystalline_firstChern_level, with z from thickening_firstChern_existsUnique L.γ M.p l s a.1 hcov (Lemma
% gal-realization-kappa-truncated-periods-exist-chern-unique), where hcov is sSup_𝒱D (the opens over which the tautological
% line datum has a section cover Y); red_ν(z) = r_ν(x_ν) is Pb.coefficientReduction_firstChernClass_tateTwist (d n) M.p s
% ((PAlg.base l).cris n) L.lam (hc n) a.1 hcov _ z (hx n) hz (Lemma ...-chern-reduction, with d = Kap.dOf and hc = Kap.hcOf
% from redIn_ex); b(z) = μ(ω) is thickening_firstChern_pullback (nmB M hL) (QA L M s) (TF 1 s) with Used from
% Or.inl ⟨rfl, 1, s, rfl⟩ and qb_le hL M s (Lemma ...-chern-pullback; nmB has hI = ι_dp, u = pullback.fst, g the
% inclusion, and hκ from ι_I and ι_O), with deRham_firstChern_to_thickening L.γB M.q s (lineK a.1) c hcK (Lemma
% ...-chern-de-rham). The hypotheses on the
% coordinate cover and on s are not used here.
\end{proof}
\leanref{Development/WeilBound/CoreStatements/PadicHodgeComparison/EtaleCohomologyComparison/PerfectoidPullbacks/DividedPowerReduction/CoefficientReductionFirstChernClassTateTwist.lean}{CrystallineSemistableComparison.Pb.thickening_firstChern_existsUnique}
\leanref{Development/WeilBound/CoreStatements/PadicHodgeComparison/EtaleCohomologyComparison/PerfectoidPullbacks/DividedPowerReduction/CoefficientReductionFirstChernClassTateTwist.lean}{CrystallineSemistableComparison.Pb.coefficientReduction_firstChernClass_tateTwist}
\leanref{Development/WeilBound/CoreStatements/PadicHodgeComparison/EtaleCohomologyComparison/PerfectoidPullbacks/DividedPowerReduction/CoefficientReductionFirstChernClassTateTwist.lean}{CrystallineSemistableComparison.Pb.thickening_firstChern_pullback}
\leanref{Development/WeilBound/CoreStatements/PadicHodgeComparison/EtaleCohomologyComparison/PerfectoidPullbacks/DividedPowerReduction/CoefficientReductionFirstChernClassTateTwist.lean}{CrystallineSemistableComparison.Pb.deRham_firstChern_to_thickening}

\begin{lemma}\label{gal-realization-kappa-truncated-periods-exist-filtration}
Let $\mathcal D_m(Y,s)$ be as in Definition~\ref{gal-realization-kappa-truncated-periods-exist-categories}, for an
approximation system, a smooth projective $\mathcal O_K$-scheme $Y$, $m\geq0$ and $s\geq m+1$. Suppose that $Y_K$ has a
coordinate cover $(U_\alpha)$ of dimensions $d_\alpha$ (Definition~\ref{gal-realization-thickening-coordinates}(2)) with
$d_\alpha+m+1\leq s$ for every $\alpha$. Let $j\geq0$, $0\leq n\leq s$ and $r\geq0$. Let $x=(x_\nu)$ be a compatible
family with $x_\nu\in H^j_\nu(Y)$, such that every $x_\nu$ is the class of a cocycle whose value on every chain
$T_0\to\cdots\to T_j$ lies in $J_{T_j}^{[n]}$; let $f$ be a period of $x$ in $\mathcal D_m(Y,s)$; and let
$\psi\in\operatorname{Hom}_K(H^j_{\mathrm{dR}}(Y_K),K)$ vanish on the classes of the cocycles of
$\mathrm{Th}_{\mathrm{dR}}(Y_K)$ whose value on every chain $T_0\to\cdots\to T_j$ lies in $J_{T_j}^{[r+1]}$. Then
$f(\psi)\in J_m^{[n-r]}$, formed for $\delta_m$, where $J_m^{[k]}=B_m$ for $k\leq0$.
\end{lemma}
\leanref{Development/WeilBound/CoreStatements/PadicHodgeComparison/EtaleCohomologyComparison/PerfectoidPullbacks/LevelChainValuations/ValuationChainProofs/fil_lev.lean}{CrystallineSemistableComparison.Pb.Kap.fil_lev_pf}

\begin{proof}
Write $\mathcal Q=\mathcal Q_m(Y,s)$ and $\mathcal Q'=\mathcal Q'_m(Y,s)$, and let $e$ and $z$ be as in
Definition~\ref{gal-realization-kappa-truncated-periods-exist-diagram}(3) for $f$. As in the proof of
Lemma~\ref{gal-realization-kappa-truncated-periods-exist-effective}, $Y$ is proper with a coordinate cover, and
Lemma~\ref{gal-realization-kappa-truncated-periods-exist-lifting} applies to $\Lambda=\mathcal O_K$, $A_m$, $I_m$,
$\gamma_m$, $s$ in place of $m$, $X=Y$, $\mathcal V=\{Y\}$ and $j$, with the category $\mathcal Q$; let $e'$ be as there.

\emph{Step 1: for every $\nu$, the image of $z$ in $H^j(\mathcal Q,\mathcal O/(\mathcal O^{[n]}+\ell^\nu\mathcal O))$ is
$0$.} As $\mathrm{red}_\nu(z)=\ell^er_\nu(x_\nu)$, it suffices that the image of $r_\nu(x_\nu)$ there is $0$. For $T$ in
$\mathcal Q$, with $R_\nu(T)=(\bar B,\ldots,\bar J,\bar\delta,\ldots)$ and $\pi:B\to\bar B$ as in
Definition~\ref{gal-realization-kappa-truncated-periods-exist-categories}(2), $\bar J^{[n]}\subset\pi(J^{[n]})$: as $\pi$ is
a surjective divided power map and $\bar J=\pi(J)$, each generator $\prod_k\bar\delta_{a_k}(\pi(b_k))$, $b_k\in J$, is
$\pi(\prod_k\delta_{a_k}(b_k))$. Let $c$ be a cocycle of $\mathrm{Th}_\nu(Y)$ with values in $\mathcal O^{[n]}$ and class
$x_\nu$. Then $r_\nu(x_\nu)$ is the class of the cocycle $T_0\to\cdots\to T_j\mapsto c(R_\nu(T_0)\to\cdots\to R_\nu(T_j))$,
whose values lie in $\pi(J_{T_j}^{[n]})$, and whose image in the cochains with values in $\mathcal O/(\mathcal
O^{[n]}+\ell^\nu\mathcal O)$ is therefore $0$.

\emph{Step 2.} By Step 1 and Lemma~\ref{gal-realization-kappa-truncated-periods-exist-lifting}(3), for $n\leq s$, $\ell^{e'}$
times the image $\bar z$ of $z$ in $H^j(\mathcal Q,\mathcal O/\mathcal O^{[n]})$ is $0$.

\emph{Step 3.} The functor $N$ keeps rings, ideals and divided powers, so with the identity maps of the rings it induces a
map $H^j(\mathcal Q,\mathcal O/\mathcal O^{[n]})\to H^j(\mathcal Q',\mathcal O/\mathcal O^{[n]})$, which carries $\bar z$
to the image of $b(z)$, as both composites are induced by $N$ with the projections. So $\ell^{e'}$ times the image of
$b(z)$ in $H^j(\mathcal Q',\mathcal O/\mathcal O^{[n]})$ is $0$. Let $(v_t)$ be a basis of $H^j_{\mathrm{dR}}(Y_K)$, which
is finite-dimensional by Lemma~\ref{gal-realization-kappa-de-rham}(1). Then $b(z)=\ell^e\sum_tf(v_t^*)\,\mu(v_t)$, and as
$\ell$ is a unit of $B_m$, the image of $\sum_tf(v_t^*)\,\mu(v_t)$ in $H^j(\mathcal Q',\mathcal O/\mathcal O^{[n]})$ is $0$.

\emph{Step 4.} By (a) of Definition~\ref{gal-realization-kappa-approximations}(1), $J_m^{[m+1]}=\tau_m(t)^{m+1}B_m$,
which is $0$ as $t^{m+1}$ lies in the kernel of $\tau_m$. So
Lemma~\ref{gal-realization-kappa-truncated-periods-exist-base-change}(2), for $B_m$, $m+1\leq s$, $X=Y_K$, $\mathcal
V=\{Y_K\}$, $r=n$,
$g_t=f(v_t^*)$ and $c=r$, gives $\sum_tf(v_t^*)\psi(v_t)\in J_m^{[n-r]}$. As $f(a\psi)=af(\psi)$ for $a\in K$ and
$\psi=\sum_t\psi(v_t)v_t^*$, $f(\psi)=\sum_tf(v_t^*)\psi(v_t)$.
% Lean: fil_lev_pf. Step 1 is h1, with toSome_redM and red_fil = RedFil.red_kills_toSome (red_kills, val_zero and
% redObj_pdPow through dividedPowerFiltration_le_map); Step 2 is the third part of cg_lift (LiftP); Step 3 is quoM_bMap,
% with unit_B; Step 4 is the second part of basisB, with pdP_B_bot (Lev.Good.pdP_B and toB_ker) for J_m^[m+1] = 0, and h8.
\end{proof}
\leanref{Development/WeilBound/CoreStatements/PadicHodgeComparison/EtaleCohomologyComparison/PerfectoidPullbacks/DividedPowerReduction/CoefficientReductionFirstChernClassTateTwist.lean}{CrystallineSemistableComparison.Pb.red_fil}
\leanref{Development/AutomorphicSpectrum/Proofs/BlockDecomposition/RedObjPdPow.lean}{GlobalTheory.CompanionTheory.RedFil.red_kills}
\leanref{Development/WeilBound/CoreStatements/PadicHodgeComparison/EtaleCohomologyComparison/PerfectoidPullbacks/LevelChainValuations/ValuationChainProofs/fil_lev.lean}{CrystallineSemistableComparison.Pb.Kap.quoM_bMap}

\begin{lemma}\label{gal-realization-kappa-truncated-periods-exist}
Let $\ell$, $K$ and $\mathcal O_K$ be as in Definition~\ref{gal-realization-period-datum}. Every approximation system
(Definition~\ref{gal-realization-kappa-approximations}(2)) has a system of truncated periods
(Definition~\ref{gal-realization-kappa-truncated-periods}).
\end{lemma}
\leanref{Development/WeilBound/CoreStatements/PadicHodgeComparison/EtaleCohomologyComparison/PerfectoidPullbacks/LevelChainValuations/RhoActionCup/RhoFinal.lean}{CrystallineSemistableComparison.Pb.Kap.Mo.rhoSys_ex}
\leanref{Development/WeilBound/CoreStatements/PadicHodgeComparison/EtaleCohomologyComparison/PerfectoidPullbacks/LevelChainValuations/CrystallineDeRhamFirstChernExists.lean}{CrystallineSemistableComparison.Pb.rhoSys_ex'}

\begin{proof}
Let the approximation system be given, and let $\mathcal D_m(Y,s)$ be as in
Definition~\ref{gal-realization-kappa-truncated-periods-exist-categories}. For a smooth projective $\mathcal O_K$-scheme
$Y$, a tautological line datum $\mathcal L$ on $Y$, $j,m\geq0$, $x\in H^j_\infty(Y)$ and an additive map
$f:\operatorname{Hom}_K(H^j_{\mathrm{dR}}(Y_K),K)\to B_m$, let $\rho_m(x,f)$ hold if there is $s_0$ such that $f$ is a
period of the components $(x_\nu)_\nu$ of $x$ in $\mathcal D_m(Y,s)$ for every $s\geq s_0$
(Definition~\ref{gal-realization-kappa-truncated-periods-exist-diagram}(3)); it does not depend on $\mathcal L$. The
components of $x$ form a compatible family, as $H^j_\infty(Y)=\varprojlim_\nu H^j_\nu(Y)$. The numbers (1) to (9) below
are those of Definition~\ref{gal-realization-kappa-truncated-periods}.

Fix $Y$. Its generic fibre $Y_K$ is projective and smooth over $K$, so it is separated, quasi-compact and smooth, and it
has a coordinate cover $(U_\alpha)$ of dimensions $d_\alpha$ (Lemma~\ref{gal-realization-thickening-cover}(2)). For
$m\geq0$ let $s_m=m+1+d$, with $d$ the largest $d_\alpha$ ($d=0$ if there is no $\alpha$). For $s\geq s_m$,
$\mathcal D_m(Y,s)$ is effective by Lemma~\ref{gal-realization-kappa-truncated-periods-exist-effective}.

\emph{Step 0: if $s_m\leq s\leq s'$ and $f$ is a period of a family $x$ in $\mathcal D_m(Y,s')$, then it is one in
$\mathcal D_m(Y,s)$.} Apply Lemma~\ref{gal-realization-kappa-truncated-periods-exist-periods}(5) to the map of
Lemma~\ref{gal-realization-kappa-truncated-periods-exist-maps}(1): its ring maps, $h_S$ and $h_V$ are identities, so
$h_*f=f$.

(1) Let $x\in H^j_\infty(Y)$. By Lemma~\ref{gal-realization-kappa-truncated-periods-exist-periods}(1), $x$ has a period
$f$ in $\mathcal D_m(Y,s_m)$. For $s\geq s_m$ it has a period $g$ in $\mathcal D_m(Y,s)$, which is a period in
$\mathcal D_m(Y,s_m)$ by Step 0, so $g=f$ by Lemma~\ref{gal-realization-kappa-truncated-periods-exist-periods}(2). So
$\rho_m(x,f)$, with $s_0=s_m$. If $\rho_m(x,f)$ and $\rho_m(x,f')$, with $s_0$ and $s'_0$, then $f$ and $f'$ are periods of
$x$ in $\mathcal D_m(Y,s)$ for $s=\max(s_0,s'_0,s_m)$, so $f=f'$ by
Lemma~\ref{gal-realization-kappa-truncated-periods-exist-periods}(2).

(2) Let $\rho_m(x,f)$ and $\rho_m(x',f')$, with $s_0$ and $s'_0$, and $s\geq\max(s_0,s'_0)$. The components of $x+x'$ and of
$cx$, for $c\in A_{\mathrm{cris}}$, are the $x_\nu+x'_\nu$ and the $cx_\nu$. By
Lemma~\ref{gal-realization-kappa-truncated-periods-exist-periods}(3) and~(4), $f+f'$ and $\psi\mapsto
\iota_m(\lambda_m(c))f(\psi)$ are periods of them in $\mathcal D_m(Y,s)$. By
Definition~\ref{gal-realization-kappa-truncated-periods-exist-diagram}(3), $f(a\psi)=af(\psi)$ for $a\in K$.

(3) Let $\rho_{m+1}(x,f)$, with $s_0$, and $s\geq\max(s_0,s_{m+1})$. Then $\mathcal D_{m+1}(Y,s)$ and $\mathcal D_m(Y,s)$ are
effective, and Lemma~\ref{gal-realization-kappa-truncated-periods-exist-periods}(5), for the map of
Lemma~\ref{gal-realization-kappa-truncated-periods-exist-maps}(2), whose $h_S$ and $h_V$ are identities, shows that
$d_B\circ f$ is a period of $(x_\nu)_\nu$ in $\mathcal D_m(Y,s)$. So $\rho_m(x,d_B\circ f)$.

(4) Let $\rho_m(x,f)$, with $s_0$, $\sigma\in G_K$ and $s\geq\max(s_0,s_m)$. By
Lemma~\ref{gal-realization-kappa-truncated-periods-exist-periods}(5) for the map of
Lemma~\ref{gal-realization-kappa-truncated-periods-exist-maps}(3), $\sigma_{B_m}\circ f$ is a period of the components
$(\sigma x_\nu)_\nu$ of $\sigma x$ in $\mathcal D_m(Y,s)$.

(5) Let $s$ be at least $s_m$ and the numbers $s_0$ of the three relations. The components of $x\cup x'$ are the
$x_\nu\cup x'_\nu$ (Definition~\ref{gal-realization-period-thickenings}(5)). By
Lemma~\ref{gal-realization-kappa-truncated-periods-exist-maps}, the cup data form a product structure on $\mathcal
D_m(Y,s)$, and Lemma~\ref{gal-realization-kappa-truncated-periods-exist-periods}(6) gives the formula.

(6) Likewise, the components of $1\in H^0_\infty(Y)$ are the units $1_\nu$ of the $H^0_\nu(Y)$, $1_V$ is the unit of
$H^0_{\mathrm{dR}}(Y_K)$, and Lemma~\ref{gal-realization-kappa-truncated-periods-exist-periods}(7) gives $f(\psi)=\psi(1)$.

(7) Let $x$ and $\omega$ be as in (7), $\rho_m(x,f)$ with $s_0$, and $s\geq\max(s_0,s_m)$. By
Lemma~\ref{gal-realization-kappa-truncated-periods-exist-chern}, $\psi\mapsto\psi(\omega)$ is a period of $(x_\nu)_\nu$ in
$\mathcal D_m(Y,s)$, so $f(\psi)=\psi(\omega)$ by Lemma~\ref{gal-realization-kappa-truncated-periods-exist-periods}(2).

(8) Let $h$, $x'$, $f'$ and $f$ be as in (8), with $s_0$ for $\rho_m(x',f')$ and $s'_0$ for $\rho_m(h^*x',f)$. Let $s$ be at
least $s_0$, $s'_0$, $s_m$ and the number $s'_m$ formed in the same way for a coordinate cover of $Y'_K$. By
Lemma~\ref{gal-realization-kappa-truncated-periods-exist-periods}(5) for the map of
Lemma~\ref{gal-realization-kappa-truncated-periods-exist-maps}(4), the map $\psi\mapsto f'(\psi\circ h_K^*)$ is a period
of the components $(h^*x'_\nu)_\nu$ of $h^*x'$ in $\mathcal D_m(Y,s)$. By
Lemma~\ref{gal-realization-kappa-truncated-periods-exist-periods}(2) it is $f$.

(9) Let $n\leq m$, $r$, $x$ and $\psi$ be as in (9), $\rho_m(x,f)$ with $s_0$, and $s\geq\max(s_0,s_m)$. Then $n\leq m<s$,
and Lemma~\ref{gal-realization-kappa-truncated-periods-exist-filtration} gives $f(\psi)\in J_m^{[n-r]}$.
% Lean: Pb.Kap.Mo.rhoSys_ex with ρ = rhoK: rhoK a j m x f is M.Sm together with an s₀ such that (chK S hS hC M m s).Rho
% holds for s ≥ s₀, where chK is the book's D_m(Y,s) (chainOf, with rHof from dOf and hcOf). s_m is the bound of chK_good
% (cov_sm_pf, chain_good, rHof_tr_gen); Step 0 is rho_mono (realization_functorial for homS). (1) is rho_ex_pf and
% rho_uniq_pf; (2) is rho_add_pf, rho_smul_pf and rho_lin'_pf; (3) is rho_lev_pf (homLev, lev_good); (4) is rho_act_of
% (homGal, homGal_good_of_galRed, galRed_rHof); (5) is rho_cup_pf (cup_genA with realization_cup_product, cupOf_rHof_good
% and cup_iff_cS_lim); (6) is rho_one_of (realization_unit_class, one_iff); (7) is rho_firstChern_formula
% (crystalline_firstChern_level, realization_unique with evB); (8) is rho_nat_of (homPull, homPull_rHof_good,
% realization_unique); (9) is rho_fil_pf (fil_lev_pf).
\end{proof}
\leanref{Development/WeilBound/CoreStatements/PadicHodgeComparison/EtaleCohomologyComparison/PerfectoidPullbacks/LevelChainValuations/RhoFirstChernLaw.lean}{CrystallineSemistableComparison.Pb.Kap.rhoK}
\leanref{Development/WeilBound/CoreStatements/PadicHodgeComparison/EtaleCohomologyComparison/PerfectoidPullbacks/LevelChainValuations/RhoActionCup/RhoFirstChernFormula.lean}{CrystallineSemistableComparison.Pb.Kap.Mo.chK_good}
\leanref{Development/WeilBound/CoreStatements/PadicHodgeComparison/EtaleCohomologyComparison/PerfectoidPullbacks/LevelChainValuations/RhoActionCup/RhoFirstChernFormula.lean}{CrystallineSemistableComparison.Pb.Kap.Mo.rho_mono}
\leanref{Development/WeilBound/CoreStatements/PadicHodgeComparison/EtaleCohomologyComparison/PerfectoidPullbacks/LevelChainValuations/RhoActionCup/RhoFirstChernFormula.lean}{CrystallineSemistableComparison.Pb.Kap.Mo.rho_ex_pf}
\leanref{Development/WeilBound/CoreStatements/PadicHodgeComparison/EtaleCohomologyComparison/PerfectoidPullbacks/LevelChainValuations/RhoActionCup/RhoFirstChernFormula.lean}{CrystallineSemistableComparison.Pb.Kap.Mo.rho_uniq_pf}
\leanref{Development/WeilBound/CoreStatements/PadicHodgeComparison/EtaleCohomologyComparison/PerfectoidPullbacks/LevelChainValuations/RhoActionCup/RhoFinal.lean}{CrystallineSemistableComparison.Pb.Kap.Mo.rho_lev_pf}
\leanref{Development/WeilBound/CoreStatements/PadicHodgeComparison/EtaleCohomologyComparison/PerfectoidPullbacks/LevelChainValuations/RhoActionCup/RhoFinal.lean}{CrystallineSemistableComparison.Pb.Kap.Mo.rho_act_of}
\leanref{Development/WeilBound/CoreStatements/PadicHodgeComparison/EtaleCohomologyComparison/PerfectoidPullbacks/LevelChainValuations/RhoActionCup/RhoFinal.lean}{CrystallineSemistableComparison.Pb.Kap.Mo.rho_cup_pf}
\leanref{Development/WeilBound/CoreStatements/PadicHodgeComparison/EtaleCohomologyComparison/PerfectoidPullbacks/LevelChainValuations/RhoActionCup/RhoFinal.lean}{CrystallineSemistableComparison.Pb.Kap.Mo.rho_one_of}
\leanref{Development/WeilBound/CoreStatements/PadicHodgeComparison/EtaleCohomologyComparison/PerfectoidPullbacks/LevelChainValuations/RhoActionCup/RhoFinal.lean}{CrystallineSemistableComparison.Pb.Kap.Mo.rho_nat_of}
\leanref{Development/WeilBound/CoreStatements/PadicHodgeComparison/EtaleCohomologyComparison/PerfectoidPullbacks/LevelChainValuations/RhoActionCup/RhoFirstChernFormula.lean}{CrystallineSemistableComparison.Pb.Kap.Mo.rho_fil_pf}
\leanref{Development/WeilBound/CoreStatements/PadicHodgeComparison/EtaleCohomologyComparison/PerfectoidPullbacks/LevelChainValuations/RhoActionCup/RhoFirstChernFormula.lean}{CrystallineSemistableComparison.Pb.Kap.Mo.rho_firstChern_formula}
\leanref{Development/WeilBound/CoreStatements/PadicHodgeComparison/EtaleCohomologyComparison/PerfectoidPullbacks/LevelChainValuations/RhoActionCup/CupGenA.lean}{CrystallineSemistableComparison.Pb.Kap.cup_genA}

\subsection{Truncated de Rham thickenings}\label{subsection-gal-de-rham-thickenings}

\begin{definition}\label{gal-realization-kappa-de-rham-truncated}
Let $(R,I,\gamma,p)$, $\mathcal O$, $X$ and $\kappa$ be as in Definition~\ref{gal-realization-thickening-nilpotent}, with
$p=0$, and let $s\geq1$.
\begin{enumerate}
\item $\mathrm{Th}^{(s)}(X/R)$ is the full subcategory of $\mathrm{Th}(X/R)$ of the thickenings $T$ with $J_T^{[s]}=0$,
  and for a set $\mathcal V$ of open subsets of $X$, $\mathrm{Th}^{(s)}(X/R)_{\mathcal V}$ is its full subcategory of the
  thickenings that lie over some $V\in\mathcal V$ (Definition~\ref{gal-realization-thickening-nilpotent}(2) and (3)).
\item Let $K$ be a field and $X$ a $K$-scheme, and take $R=\mathcal O=K$, $I=0$ with its divided powers, and $\kappa$ the
  identity. A thickening of $X$ over $(K,0,0,0)$ is a nilpotent thickening of $X$ in the sense of
  Definition~\ref{gal-realization-period-de-rham}, with $i$ its $K$-algebra structure, as the condition on $i$ is empty
  for $I=0$ and $p=0$ asks that $J$ be nilpotent; and the maps of the two definitions agree. So
  $\mathrm{Th}_{\mathrm{dR}}(X)=\mathrm{Th}(X/K)$, and $H^j_{\mathrm{dR}}(X)=H^j(\mathrm{Th}(X/K),\mathcal O)$. Write
  $\mathrm{Th}^{(s)}_{\mathrm{dR}}(X)=\mathrm{Th}^{(s)}(X/K)$ and $\mathrm{gr}_r=\mathcal O^{[r]}/\mathcal O^{[r+1]}$. In
  Definition~\ref{gal-realization-thickening-level-one-fibre}(1), $X_K=X$.
\end{enumerate}
\end{definition}
\leanref{Development/WeilBound/CoreStatements/PadicHodgeComparison/EtaleCohomologyComparison/PerfectoidPullbacks/SEqOne.lean}{CrystallineSemistableComparison.Pb.TF}
\leanref{Development/WeilBound/CoreStatements/PadicHodgeComparison/EtaleCohomologyComparison/PerfectoidPullbacks/InclStd.lean}{CrystallineSemistableComparison.Pb.KS}

\begin{lemma}\label{gal-realization-kappa-de-rham-subideal}
Let $J$ be an ideal of a ring $B$ with divided powers $\delta$, and $m\geq1$. Then $J^{[m]}\subset J$, and
$\delta_n(a)\in J^{[nm]}$ for every $a\in J^{[m]}$ and $n\geq0$ (Definition~\ref{gal-realization-kappa-approximations}).
In particular $\delta_n(a)\in J^{[m]}$ for $a\in J^{[m]}$ and $n\geq1$.
\end{lemma}
\leanref{Development/WeilBound/CoreStatements/AbelianHeightsDuality/DividedPowerAlgebras/FilIsSubDPIdeal.lean}{FubiniStudyChowHeightFormula₂.Adr.Fil.fil_isSubDPIdeal}
\leanref{Development/AbelianVarietyTheory/Proofs/DivisorProofs/DividedPowerFiltration.lean}{GlobalTheory.Cris.ExpBij.dpow_mem_dpFil_mul}
\leanref{Development/AbelianVarietyTheory/Proofs/DivisorProofs/DividedPowerFiltration.lean}{GlobalTheory.Cris.ExpBij.dpFil_le}

\begin{proof}
Each generator $\prod_j\delta_{k_j}(x_j)$ of $J^{[m]}$, with $x_j\in J$ and $\sum_jk_j\geq m\geq1$, has a factor with
$k_j\geq1$, which lies in $J$. So $J^{[m]}\subset J$.

\emph{Generators.} Let $z=\prod_{j=1}^p\delta_{k_j}(x_j)$ with $x_j\in J$, and $k=\sum_jk_j$. Then $\delta_n(z)\in
J^{[nk]}$ for every $n\geq0$, by induction on $p$. For $p=0$, $z=1$ and $J^{[0]}=B$. Let $p\geq1$ and
$z'=\prod_{j\geq2}\delta_{k_j}(x_j)$, so that $z=\delta_{k_1}(x_1)z'$ and $z'\in J^{[k-k_1]}$. If $k_1=0$, then $z=z'$, and
the induction applies. If $k_1\geq1$, then $\delta_{k_1}(x_1)\in J$, and
\[
 \delta_n(z)=z'^n\,\delta_n(\delta_{k_1}(x_1))=\frac{(nk_1)!}{n!\,(k_1!)^n}\,z'^n\,\delta_{nk_1}(x_1),
\]
where $z'^n\in J^{[n(k-k_1)]}$, as $J^{[a]}J^{[b]}\subset J^{[a+b]}$, and $\delta_{nk_1}(x_1)\in J^{[nk_1]}$. So
$\delta_n(z)\in J^{[nk]}$.

\emph{All elements.} Let $S$ be the set of $a\in J^{[m]}$ with $\delta_n(a)\in J^{[nm]}$ for every $n\geq0$. It contains
the generators $z$ of $J^{[m]}$, by the above, as $nk\geq nm$ and the ideals $J^{[r]}$ decrease in $r$. It contains $0$, as
$\delta_n(0)=0$ for $n\geq1$. For $a,a'\in S$, which lie in $J$,
$\delta_n(a+a')=\sum_{u+v=n}\delta_u(a)\delta_v(a')\in J^{[um]}J^{[vm]}\subset J^{[nm]}$, so $a+a'\in S$. For $b\in B$ and
$a\in S$, $\delta_n(ba)=b^n\delta_n(a)$, so $ba\in S$. Hence $S=J^{[m]}$. For $n\geq1$, $J^{[nm]}\subset J^{[m]}$.
% Lean: dpFil_le for the inclusion; dpow_mem_dpFil_mul, whose hgen is the induction on the generators (dpow_mul,
% dpow_comp and pow_mem_dpFil) and whose span induction is the step "All elements"; fil_isSubDPIdeal adds dpFil_anti.
\end{proof}

\begin{lemma}\label{gal-realization-kappa-de-rham-poly}
Let $(R,I,\gamma,0)$, $\mathcal O$, $X$ and $\kappa$ be as in Definition~\ref{gal-realization-kappa-de-rham-truncated},
$s,\ell\geq1$ integers, $\mathcal V$ a set of open subsets of $X$ and $\mathcal Q=\mathrm{Th}^{(s)}_\ell(X/R)_{\mathcal V}$
(Lemma~\ref{gal-realization-kappa-truncated-periods-exist-base-change-zariski}). Then $\Phi_{\mathcal Q}$
(Definition~\ref{gal-realization-thickening-torsor}(1)) has a divided power polynomial structure that respects, in every
degree $r\leq s$, the sequence $(J^{[r']})_{r'}$ and, for every $c\in R$, the sequence $(J^{[r']}+cB)_{r'}$.
\end{lemma}
\leanref{Development/WeilBound/CoreStatements/PadicHodgeComparison/EtaleCohomologyComparison/PerfectoidPullbacks/ChartGrConst.lean}{CrystallineSemistableComparison.Pb.poly_tf}
\leanref{Development/WeilBound/CoreStatements/PadicHodgeComparison/EtaleCohomologyComparison/PerfectoidPullbacks/PolynomialSiteTraces/PolyTf.lean}{CrystallineSemistableComparison.Pb.Site.poly_tf'}

\begin{proof}
Let $T=(B,i,J,\delta,y)$ be in $\mathcal Q$, and let $B\langle X\rangle$, its basis $X^{[k]}$, the ideal $J\langle X\rangle$,
the divided powers $\eta$ on it and $j:B\to B\langle X\rangle$ be as in the proof of
Lemma~\ref{gal-realization-thickening-chart-torsor}(3). Let $N=J\langle X\rangle^{[s]}$, formed for $\eta$.

\emph{The quotient.} By Lemma~\ref{gal-realization-kappa-de-rham-subideal}, for $\eta$ and $m=s$, $N\subset J\langle
X\rangle$, and $\eta_n(z)\in N$ for every $z\in N$ and $n\geq1$. So $\eta$ extends to $P=B\langle X\rangle/N$
\cite[Tag 07H2]{Sta}, uniquely \cite[Tag 07H1]{Sta}: there is exactly one divided power structure $\bar\eta$ on the image
$\bar J$ of $J\langle X\rangle$ in $P$ for which the projection $\pi$ is a divided power map. If $g:B\langle X\rangle\to C$ is a divided power map to a ring with divided
powers that kills $N$, the induced map $P\to C$ is a divided power map, as $\bar\eta_n(\pi(z))=\pi(\eta_n(z))$. For
$r\geq0$, $\bar J^{[r]}=\pi(J\langle X\rangle^{[r]})$: the right side lies in the left as $\pi$ is a divided power map,
and the generators $\prod_i\bar\eta_{k_i}(\pi(z_i))$ of the left side are the images of the $\prod_i\eta_{k_i}(z_i)$. In
particular $\bar J^{[s]}=0$, and $\bar J^s=0$, as a product $x_1\cdots x_s=\bar\eta_1(x_1)\cdots\bar\eta_1(x_s)$ of $s$
elements of $\bar J$ lies in $\bar J^{[s]}$.

\emph{The object $P(T)$.} As $N\subset J\langle X\rangle$, $P/\bar J=B\langle X\rangle/J\langle X\rangle=B/J$. Let
$P(T)=(P,\pi ji,\bar J,\bar\eta,y)$, $j_T=\pi j$ and $X_T=\pi(X)\in\bar J$. As in the proof of
Lemma~\ref{gal-realization-thickening-chart-torsor}(3), $P(T)$ is a thickening and $j_T$ a map $T\to P(T)$: the map
$\pi ji$ is a divided power map, as a composite of three, and $\bar J$ is nilpotent. It lies in $\mathcal Q$, as
$\bar J^{[s]}=0$, it lies over every open that $T$ lies over, and it satisfies the condition on $\ell$. For the last, let
$r\geq0$ and $m=\min(r,s)$. As $N\subset J\langle X\rangle^{[m]}$ and $J\langle X\rangle^{[r]}\subset N$ if $r\geq s$,
an element $\pi(z)$ lies in $\bar J^{[r]}=\pi(J\langle X\rangle^{[r]})$ if and only if $z\in J\langle X\rangle^{[m]}$.
Write $z=\sum_kc_kX^{[k]}$. If $\ell\pi(z)\in\bar J^{[r]}$, then $\ell z=\sum_k\ell c_kX^{[k]}\in J\langle X\rangle^{[m]}$, so
$\ell c_k\in J^{[m-k]}$ for every $k$, by the filtration in the proof of
Lemma~\ref{gal-realization-thickening-chart-torsor}(3); then $c_k\in J^{[m-k]}$, as $T$ is in $\mathcal Q$, so $z\in
J\langle X\rangle^{[m]}$ and $\pi(z)\in\bar J^{[r]}$.

\emph{The universal property.} Let $\varphi:T\to T'$ be a map of $\mathcal Q$ and $x\in J_{T'}$. By the proof of
Lemma~\ref{gal-realization-thickening-chart-torsor}(3) there is exactly one divided power map $F:B\langle X\rangle\to
B_{T'}$ with $Fj=\varphi$ and $F(X)=x$; it commutes with the maps from $R$ and is compatible with the points. It carries
$N$ into $J_{T'}^{[s]}=0$, so it induces a divided power map $\psi:P\to B_{T'}$, which is a map $P(T)\to T'$ of
$\mathcal Q$ with $\psi j_T=\varphi$ and $\psi(X_T)=x$. For any such $\psi$, $\psi\pi$ is a divided power map with these
two properties, so it is $F$, and $\psi$ is unique.

\emph{Spanning.} Every element of $P$ is $\pi(\sum_kc_kX^{[k]})=\sum_kj_T(c_k)\bar\eta_k(X_T)$ with $c_k\in B$.

\emph{The sequences.} Let $r\leq s$, $(c_k)_{k\geq0}$ finitely supported in $B$ and $z=\sum_kc_kX^{[k]}$. If $r\leq0$,
both sides of the two criteria hold, as $\bar J^{[r]}=P$ and $J^{[r-k]}=B$. Let $0<r\leq s$. Then $N\subset J\langle
X\rangle^{[r]}$, so $\pi(z)\in\bar J^{[r]}=\pi(J\langle X\rangle^{[r]})$ if and only if $z\in J\langle X\rangle^{[r]}$,
that is, $c_k\in J^{[r-k]}$ for every $k$, by the filtration in the proof of
Lemma~\ref{gal-realization-thickening-chart-torsor}(3). Likewise, for $c\in R$, $\bar J^{[r]}+cP=\pi(J\langle
X\rangle^{[r]}+cB\langle X\rangle)$, and $\pi(z)$ lies in it if and only if $z\in J\langle X\rangle^{[r]}+cB\langle X\rangle$,
that is, $c_k\in J^{[r-k]}+cB$ for every $k$, by the description of these ideals in that proof.
% Lean: Pb.poly_tf = Site.poly_tf' with polyTf. P(T) is polyTr T s: trB = (pq B).P / trI with trI = fil (eta δ)
% (max s 1), trD with trD_spec from Adr.Fil.exists_quot (Mathlib's DividedPowers.Quotient, the book's Stacks 07H2,
% (a) implies (b)) for the sub-pd ideal of Adr.Fil.fil_isSubDPIdeal (Lemma gal-realization-kappa-de-rham-subideal).
% Membership in Q is polyTr_mem with truncated_envelope_torsion_free (SiteTr.tr_tf:
% fil δ' s = ⊥ through fil_le_map_fil and map_fil_le, and nilpotence from J^n ⊂ J^[n]); univ is
% divided_power_envelope_universal with map_fil_le and trD_spec; span is divided_power_envelope_expansion (trSum); the
% sequences are truncated_envelope_filtration_quotient (SiteTrCrit.tr_crit, with crit_sup and PDFil.fil_pdPoly).
\end{proof}
\leanref{Development/WeilBound/CoreStatements/PadicHodgeComparison/EtaleCohomologyComparison/PerfectoidPullbacks/PolynomialSiteTraces/PolyTf.lean}{CrystallineSemistableComparison.Pb.Site.polyTf}
\leanref{Development/WeilBound/CoreStatements/PadicHodgeComparison/EtaleCohomologyComparison/PerfectoidPullbacks/PolynomialSiteTraces/AdmissibleMap.lean}{CrystallineSemistableComparison.Pb.Site.trD_spec}
\leanref{Development/WeilBound/CoreStatements/AbelianHeightsDuality/DividedPowerAlgebras/FilIsSubDPIdeal.lean}{FubiniStudyChowHeightFormula₂.Adr.Fil.exists_quot}
\leanref{Development/AutomorphicSpectrum/Proofs/BlockDecomposition/TruncatedEnvelopeTorsionFree.lean}{GlobalTheory.CompanionTheory.SiteTr.tr_tf}
\leanref{Development/AutomorphicSpectrum/Proofs/BlockDecomposition/TruncatedEnvelopeFiltrationQuotient.lean}{GlobalTheory.CompanionTheory.SiteTrCrit.tr_crit}

\begin{lemma}\label{gal-realization-kappa-de-rham-chart}
Let $\mathcal O$ be a ring, $R$ an $\mathcal O$-algebra with divided powers $\gamma$ on an ideal $I$, and $s,\ell\geq1$
integers such that $I^{[s]}=0$ (Definition~\ref{gal-realization-kappa-approximations}), $\ell x\in I^{[r]}$ implies $x\in
I^{[r]}$ for $r\geq0$ and $x\in R$, and for every integer $n\geq1$ some power of $\ell$ lies in $nR$; for $\ell=1$ these
conditions say that $I^{[s]}=0$ and that every integer $n\geq1$ is invertible in $R$. Let $X$ be an $\mathcal O$-scheme,
and $\kappa:\mathcal O\to R/I$ the structure map followed by the projection. Let $t_1,\ldots,t_d$ be coordinates on an
affine open $V\subset X$ (Definition~\ref{gal-realization-thickening-coordinates}(1)), $G=\Gamma(V,\mathcal O_X)$,
$B=R\otimes_{\mathcal O}G$ and $i:R\to B$, $a\mapsto a\otimes1$. There is a divided power chart $D$ with
$t'_1,\ldots,t'_d$ of dimension $d$ on $V$ over $(R,I,\gamma,0)$
(Definition~\ref{gal-realization-thickening-coordinates}(3)) with $B_D=B$, $i_D=i$ and $t'_k=1\otimes t_k$, which lies in
$\mathrm{Th}^{(s)}_\ell(X/R)$ (Lemma~\ref{gal-realization-kappa-truncated-periods-exist-base-change-zariski}), such that
the ring map $G\to B/J_D$ that corresponds to $y_D$ sends $f$ to the class of $1\otimes f$.
\end{lemma}
\leanref{Development/AutomorphicSpectrum/Proofs/BlockDecomposition/ExistsChartFromCoordinatesStatement.lean}{GlobalTheory.T.exists_chart_from_coordinates}
\leanref{Development/AutomorphicSpectrum/Proofs/BlockDecomposition/ExistsChartFromCoordinates.lean}{GlobalTheory.CompanionTheory.ChartAlg.chart_exists}

\begin{proof}
\emph{Step 1: the divided powers.} The ring $G$ is \'etale over $\mathcal O[x_1,\ldots,x_d]$, which is free over
$\mathcal O$, so $G$ is flat over $\mathcal O$ \cite[Tags 00U2 and 00HC]{Sta}, and $B$ is flat over $R$. With this, the
paragraphs \emph{Torsion}, \emph{The divided powers} and \emph{The condition on $\ell$} of the proof of
Lemma~\ref{gal-realization-kappa-truncated-periods-exist-base-change-chart}(1), which use only the hypotheses on $R$
and the flatness, give divided powers $\delta$ on $J=i(I)B$ for which $i$ is a divided power map, with
$J^{[r]}=i(I^{[r]})B$ for $r\geq0$, $J^{[s]}=0$ and $J^s=0$, and such that $\ell x\in J^{[r]}$ implies $x\in J^{[r]}$.
Here $\delta_n(x)$ is the element $y$ of $B$ with $n!\,y=x^n$. For $\ell=1$ every integer $n\geq1$ is invertible in $B$,
through $i$, and $\delta_n(x)=x^n/n!$; on every ideal of $B$ these maps are divided powers, as the axioms come down to
the binomial formula and the identities $\frac{x^a}{a!}\frac{x^b}{b!}=\binom{a+b}{a}\frac{x^{a+b}}{(a+b)!}$ and
$\frac1{a!}\bigl(\frac{x^b}{b!}\bigr)^a=\frac{(ab)!}{a!\,(b!)^a}\frac{x^{ab}}{(ab)!}$.

\emph{Step 2: the thickening.} Let $g:G\to B/J$ send $f$ to the class of $1\otimes f$, and let $\psi:R/I\to B/J$ be
induced by $i$. For $a\in\mathcal O$, $a\otimes1=1\otimes a$ in $B$, so $g$ and $\psi\circ\kappa$ agree on $\mathcal O$. As
$V$ is affine, $g$ corresponds to a morphism $y:\operatorname{Spec}B/J\to X$ that factors through $V$
\cite[Tag 01I1]{Sta}, and $y$ is an $\mathcal O$-morphism for the $\mathcal O$-algebra structure of $B/J$ through
$\psi\circ\kappa$. As $J$ is nilpotent, $D=(B,i,J,\delta,y)$ is a thickening of $X$ over $(R,I,\gamma,0)$
(Definition~\ref{gal-realization-thickening-nilpotent}(1)); it lies over $V$, it lies in $\mathrm{Th}^{(s)}_\ell(X/R)$ by
Step 1, and the ring map $G\to B/J$ that corresponds to $y$ is $g$.

\emph{Step 3: the chart.} The ring map $R[x_1,\ldots,x_d]\to B$ extending $i$ with $x_k\mapsto1\otimes t_k$ is the base
change of the \'etale map $\mathcal O[x_1,\ldots,x_d]\to G$, $x_k\mapsto t_k$, along $\mathcal
O[x_1,\ldots,x_d]\to R[x_1,\ldots,x_d]$, as $B=R[x_1,\ldots,x_d]\otimes_{\mathcal O[x_1,\ldots,x_d]}G$; so it is \'etale
\cite[Tag 00U2]{Sta}. And $J=i(I)B$. Let $B'$ be a ring, $s':R/I\to B'$ a ring map and $z:\operatorname{Spec}B'\to X$ an
$\mathcal O$-morphism that factors through $V$, where $B'$ is an $\mathcal O$-algebra through $s'\circ\kappa$. Then $z$
corresponds to a ring map $h:G\to B'$ \cite[Tag 01I1]{Sta}, which agrees with $s'\circ\kappa$ on $\mathcal O$, as $z$ is
an $\mathcal O$-morphism. By the universal property of the tensor product there is exactly one ring map $\Phi:B\to B'$
with $\Phi(a\otimes f)=s'(\bar a)h(f)$, where $\bar a$ is the class of $a$; it kills $J=i(I)B$, so it induces
$\varphi:B/J\to B'$ with $\varphi\circ\psi=s'$ and $\varphi\circ g=h$, that is, with $y\circ\operatorname{Spec}\varphi=z$.
If $\varphi_2:B/J\to B'$ also has $\varphi_2\circ\psi=s'$ and $y\circ\operatorname{Spec}\varphi_2=z$, then $\varphi_2\circ
g$ corresponds to $z$, so $\varphi_2\circ g=h$; as $a\otimes f=(a\otimes1)(1\otimes f)$, $\varphi_2$ and $\varphi$ agree on
the classes of the $a\otimes f$, which generate $B/J$ additively, so $\varphi_2=\varphi$. This is the defining property of
a divided power chart, for $D$ with $t'_k=1\otimes t_k$.
% Lean: ChartAlg.chart_exists, which assumes TFBase γ l s. Step 1 is flat_of_etale, the torsion facts hl, hlk and hn,
% exists_dp (dpExt: δ_n(x) is the y with n! y = x^n, through hasDP_of_mem and fact_sum),
% dividedPowerFiltration_eq_map (hP), hnil through pow_le_pdP, and tf_tensor; Step 2 is g, ψ, hcomp with fromSpec_comp,
% and mkT with
% top_le_specMap; Step 3 is etale_tensor (RingHom.Etale.isStableUnderBaseChange) and rep_tensor; the last claim of the
% lemma is evalAt_specMap. The Lean's isomorphism e is the identity of R ⊗_O Γ(X, V), and TF l s is membership in
% Th^(s)_l(X/R).
\end{proof}
\leanref{Development/AutomorphicSpectrum/Proofs/BlockDecomposition/ExistsChartFromCoordinates.lean}{GlobalTheory.CompanionTheory.ChartAlg.exists_dp}
\leanref{Development/AutomorphicSpectrum/Proofs/BlockDecomposition/ExistsChartFromCoordinates.lean}{GlobalTheory.CompanionTheory.ChartAlg.dividedPowerFiltration_eq_map}
\leanref{Development/AutomorphicSpectrum/Proofs/BlockDecomposition/ExistsChartFromCoordinates.lean}{GlobalTheory.CompanionTheory.ChartAlg.rep_tensor}
\leanref{Development/AutomorphicSpectrum/Proofs/BlockDecomposition/ExistsChartFromCoordinates.lean}{GlobalTheory.CompanionTheory.ChartAlg.evalAt_specMap}

\begin{lemma}\label{gal-realization-kappa-de-rham-local}
Let $K$ be a field of characteristic $0$, $X$ a $K$-scheme, $V\subset X$ an open with coordinates $t_1,\ldots,t_d$
(Definition~\ref{gal-realization-thickening-coordinates}(1)), $W\subset V$ an affine open, $s\geq1$, $\mathcal
Q=\mathrm{Th}^{(s)}_{\mathrm{dR}}(X)_W$ and $0\leq r<s$ (Definition~\ref{gal-realization-kappa-de-rham-truncated}).
\begin{enumerate}
\item $H^j(\mathcal Q,\mathrm{gr}_r)=0$ for $j\neq r$.
\item Let $u_k(T)$ be lifts of the $t_k$ at the thickenings $T$ over $V$
  (Definition~\ref{gal-realization-thickening-level-one-fibre}(2)). Every class of $H^r(\mathcal Q,\mathrm{gr}_r)$ is
  represented by exactly one family $(x_I)$ of elements of $\Gamma(W,\mathcal O_X)$, for the restrictions of the $t_k$ to
  $W$ and the lifts $u_k(T)$ at the thickenings $T$ of $\mathcal Q$
  (Definition~\ref{gal-realization-thickening-level-one-fibre}(4)).
\item If $d<r$, then $H^j(\mathcal Q,\mathrm{gr}_r)=0$ for every $j$.
\end{enumerate}
\end{lemma}
\leanref{Development/AutomorphicSpectrum/Proofs/BlockDecomposition/GradedFormVanishOfFrame.lean}{GlobalTheory.CompanionTheory.GrForm.van}
\leanref{Development/AutomorphicSpectrum/Proofs/BlockDecomposition/GradedFormLocalOfFrame.lean}{GlobalTheory.CompanionTheory.GrForm.stLoc}
\leanref{Development/AutomorphicSpectrum/Proofs/BlockDecomposition/HodgeFiltrationViaTruncation.lean}{GlobalTheory.CompanionTheory.KIso.loc_vanish}

\begin{proof}
The restrictions of the $t_k$ to $W$ are coordinates on $W$ (Lemma~\ref{gal-realization-thickening-cover}(1)); write
$t_k$ also for them. Definition~\ref{gal-realization-thickening-level-one-fibre} and
Lemma~\ref{gal-realization-thickening-level-one-forms} apply to the full subcategory $\mathcal Q$ of $\mathrm{Th}(X/K)$,
with $X_K=X$. In (1) and (3) choose any lifts $u_k(T)$ of the $t_k$ at the thickenings $T$ over $W$.

\emph{Step 1: a chart.} Let $G=\Gamma(W,\mathcal O_X)$. Apply Lemma~\ref{gal-realization-kappa-de-rham-chart} with $\ell=1$
to $\mathcal O=R=K$, in which every integer $n\geq1$ is invertible as $K$ has characteristic $0$, the ideal $I=0$, which has
$0^{[s]}=0$, and the coordinates $t_k$ on $W$. It gives a divided power chart $D$ with $t'_1,\ldots,t'_d$ on $W$ over
$(K,0,0,0)$, with $B_D=K\otimes_KG$, $t'_k=1\otimes t_k$ and $J_D^{[s]}=0$, such that $\mathrm{ev}_D$ sends $f$ to the
class of $1\otimes f$. It lies in $\mathcal Q$, as $J_D^{[s]}=0$ and it lies over $W$, and $J_D=i_D(0)B_D=0$. Let $\theta:G\to B_D$,
$f\mapsto1\otimes f$, a ring isomorphism with inverse $a\otimes f\mapsto af$. Then $\theta(t_k)=t'_k$, and
$\mathrm{ev}_D=\theta$, as $J_D=0$.

\emph{Step 2: the torsor datum.} By Lemma~\ref{gal-realization-thickening-chart-torsor}(1), as $D\in\mathcal Q$ and every
thickening of $\mathcal Q$ lies over $W$, $D$ with $\tau_k(\sigma,\sigma')=\sigma'(t'_k)-\sigma(t'_k)$ is a torsor datum of
dimension $d$ for $\Phi_{\mathcal Q}$. Lemma~\ref{gal-realization-thickening-chart}(1) holds here, as its proof covers
$p=0$. So, as in Step 2 of the proof of Lemma~\ref{gal-realization-thickening-level-one-local}, there are maps
$\sigma_T:D\to T$ with $\sigma_T(t'_k)=u_k(T)$, $\sigma_c(\theta(f))$ is a lift of $f$ at $c$, and the cochain
$\sum_I\omega_I(\theta(x_I))$ of Lemma~\ref{gal-realization-thickening-level-one-graded} is the sum of
Definition~\ref{gal-realization-thickening-level-one-fibre}(4) for $x$ with the lifts $b_I=\sigma_{c_r}(\theta(x_I))$.

\emph{Step 3: the hypotheses of Lemma~\ref{gal-realization-thickening-level-one-graded}.} As in Step 3 of the proof of
Lemma~\ref{gal-realization-thickening-level-one-local}, the families $\mathcal O^{[r']}$ form a filtration of
$\Phi_{\mathcal Q}$ for which $D$ is graded constant (Definition~\ref{gal-realization-thickening-level-one-graded-filtration}),
and $J_D^{[q]}=0$ for $q\geq1$ (Lemma~\ref{gal-realization-thickening-chart}(2)). By
Lemma~\ref{gal-realization-kappa-de-rham-poly}, for $R=K$, $I=0$ and $\mathcal V=\{W\}$, $\Phi_{\mathcal Q}$ has a
divided power polynomial structure that respects $(\mathcal O^{[r']})_{r'}$ in every degree $r'\leq s$, so in every
degree $r'\leq r+1$.

(1) If $j>r$, Lemma~\ref{gal-realization-thickening-level-one-graded}(1) applies. If $j<r$, let $q=r-j\geq1$. By
Lemma~\ref{gal-realization-thickening-level-one-graded}(2) every class is represented by a cocycle with values the
classes of $\sum_I\omega_I(x_I)$ with $x_I\in J_D^{[q]}=0$, so it is $0$.

(2) By Lemma~\ref{gal-realization-thickening-level-one-graded}(2) with $j=r$ and $q=0$, and by Step 2, every class is the
class of $\omega_u(x)$ for $x_I=\theta^{-1}(x'_I)$, where $x'_I\in J_D^{[0]}=B_D$ are the elements given there. If $x$
and $x'$ represent the same class, then $x-x'$ represents $0$ by Lemma~\ref{gal-realization-thickening-level-one-forms}(3),
so $\theta(x_I-x'_I)\in J_D^{[1]}=0$ by Lemma~\ref{gal-realization-thickening-level-one-graded}(2) and Step 2, and
$x=x'$.

(3) If $j>r$, Lemma~\ref{gal-realization-thickening-level-one-graded}(1) applies. If $j\leq r$, let $q=r-j$. By
Lemma~\ref{gal-realization-thickening-level-one-graded}(2) every class is $[\sum_I\omega_I(x_I)]_r$, the sum over the
strictly increasing sequences $I$ of length $j$ in $\{1,\ldots,d\}$, with $x_I\in J_D^{[q]}$. If $j>d$ there is no such
$I$, and if $j\leq d$, then $q\geq r-d\geq1$ and $x_I=0$; in both cases the class is $0$.
% Lean: GrForm.stFrame builds the torsor datum from the chart of exists_chart_from_coordinates (Lemma
% gal-realization-kappa-de-rham-chart) for R = O and I = ⊥, with θ = e ∘ (TensorProduct.lid)⁻¹ (Step 1), chart_is_good,
% chart_divided_powers_graded_constant, hσ and hval; the polynomial structure is poly_tf. (1) is GrForm.van (StVan), (2) is
% GrForm.stLoc (StLoc, with form_sub_mem and rep_congr), and (3) is KIso.loc_vanish. The Lean states (1) and (2) for any
% ring O with TFBase (dividedPowersBot O) l s, which for O = K and l = 1 holds as K has characteristic 0.
\end{proof}
\leanref{Development/AutomorphicSpectrum/Proofs/BlockDecomposition/GradedFormFrame.lean}{GlobalTheory.CompanionTheory.GrForm.stFrame}

\begin{lemma}\label{gal-realization-kappa-de-rham-bundle}
Let $K$ be a field of characteristic $0$, $X$ a $K$-scheme, $(U_\alpha)_{\alpha\in A}$ a finite family of open subsets of
$X$ with coordinates $t^\alpha_1,\ldots,t^\alpha_{d_\alpha}$ on $U_\alpha$, such that $U_i=\bigcap_kU_{i_k}$ is affine for
every $i\in A^{n+1}$, $n\geq0$, and let $0\leq r<s$. Let $\mathcal C_i=\mathrm{Th}^{(s)}_{\mathrm{dR}}(X)_{U_i}$, and $M$
the cosimplicial $K$-module with $M^n=\prod_{i\in A^{n+1}}H^r(\mathcal C_i,\mathrm{gr}_r)$ and $(\delta_bx)_i$ the
restriction of $x_{i\circ\partial_b}$ from $\mathcal C_{i\circ\partial_b}$ to $\mathcal C_i$. There are finite sets
$\kappa_\alpha$, elements $g_{\alpha\beta}(k,k')\in\Gamma(U_\alpha\cap U_\beta,\mathcal O_X)$ and bijections $e_{n,a}$
that satisfy (i), (ii) and (iii) of Lemma~\ref{gal-realization-thickening-level-one-coherent} for $\Lambda=K$, the opens
$U_\alpha$ and $M$.
\end{lemma}
\leanref{Development/AutomorphicSpectrum/Proofs/BlockDecomposition/GradedCechCohomologyFinite.lean}{GlobalTheory.CompanionTheory.GrCech.vb_cech}
\leanref{Development/AutomorphicSpectrum/Proofs/BlockDecomposition/GradedFormVectorBundle.lean}{GlobalTheory.CompanionTheory.GrVB.main}

\begin{proof}
The proof of Lemma~\ref{gal-realization-thickening-level-one-bundle} applies, with $R=K$, $X_R=X$, the categories $\mathcal
C_i$ in place of $\mathrm{Th}(X/R)_{U_i}$, and Lemma~\ref{gal-realization-kappa-de-rham-local}(2), for $V=U_\alpha$, in
place of Lemma~\ref{gal-realization-thickening-level-one-local}(2) in its Step 1. Lemma~\ref{gal-realization-thickening-level-one-forms},
used there, holds for these full subcategories of $\mathrm{Th}(X/K)$, and the proof uses only that the $U_i$ are
affine; it does not use that the $U_\alpha$ cover $X$.
% Lean: GrCech.vb_cech, for a Cov U, through exists_graded_form_on_affine_chart (GrCech.Omega.basis, the analogue of
% FibGr.basis_gen, from stLoc, free_at, exists_of_cocycle and rp_fAct) and cover_cup_action_basis (CechVB.vb_basis), with
% fHom, fAct_str, res_fAct and res_res. For a family covering an affine open the Lean uses GrVB.main (StVB, through
% exists_equiv and chg), which gives the same data by the same argument.
\end{proof}
\leanref{Development/AutomorphicSpectrum/Proofs/BlockDecomposition/ExistsGradedFormOnAffineChartStatement.lean}{GlobalTheory.T.exists_graded_form_on_affine_chart}
\leanref{Development/AutomorphicSpectrum/Proofs/BlockDecomposition/ExistsGradedFormOnAffineChart.lean}{GlobalTheory.CompanionTheory.GrCech.Omega.basis}

\begin{lemma}\label{gal-realization-kappa-de-rham-zariski}
Let $(R,I,\gamma,0)$, $\mathcal O$, $X$ and $\kappa$ be as in Definition~\ref{gal-realization-kappa-de-rham-truncated},
and $s\geq1$. Then (1) and (2) of Lemma~\ref{gal-realization-thickening-zariski} hold for $p=0$ and $\mathcal
Q_0=\mathrm{Th}^{(s)}(X/R)$.
\end{lemma}
\leanref{Development/AutomorphicSpectrum/Proofs/BlockDecomposition/CechRefinementCohomologyBijectiveStatement.lean}{GlobalTheory.T.cech_refinement_cohomology_bijective}

\begin{proof}
The proof of Lemma~\ref{gal-realization-thickening-zariski} applies. It uses $p\neq0$ and the choice of $\mathcal Q_0$
only to see that the localization $L_S=T'_{g_S}$ of a thickening $T'\in\mathcal Q_0$ lies in $\mathcal Q_0$. Here
Lemma~\ref{gal-realization-thickening-zariski-localization}(1) covers $p=0$, and $(J_{T'}B_{L_S})^{[s]}=
\lambda_S(J_{T'}^{[s]})B_{L_S}=0$ by Lemma~\ref{gal-realization-thickening-zariski-localization}(2), so $L_S\in\mathcal
Q_0$. The elements of $J_T$ are nilpotent also for $p=0$, as $J_T$ is.
% Lean: NlZar.zariski_descent, for Used Q₀ with p = 0 and Q₀ = TF l s (l = 1 here), where used_locN (through
% pdPow_locN, the book's Lemma gal-realization-thickening-zariski-localization(2)) gives the closure under localization.
\end{proof}
\leanref{Development/AutomorphicSpectrum/Proofs/BlockDecomposition/CechRefinementCohomologyBijective.lean}{GlobalTheory.CompanionTheory.NlZar.used_locN}

\begin{lemma}\label{gal-realization-kappa-de-rham-affine-basic}
Let $\mathcal O$ be a ring, $X$ an $\mathcal O$-scheme smooth over $\mathcal O$ and $V\subset X$ an affine open. For
$f\in\Gamma(V,\mathcal O_X)$ let $D(f)\subset V$ be the open where $f$ is a unit.
\begin{enumerate}
\item There are a finite set $A$, elements $f_\alpha\in\Gamma(V,\mathcal O_X)$ and integers $d_\alpha\geq0$, for
  $\alpha\in A$, such that the $D(f_\alpha)$ cover $V$ and $D(f_\alpha)$ has coordinates of dimension $d_\alpha$
  (Definition~\ref{gal-realization-thickening-coordinates}(1)). If $X$ is smooth of relative dimension $d$ over $\mathcal
  O$, they can be chosen with $d_\alpha=d$ for every $\alpha$.
\item For $f_1,\ldots,f_m\in\Gamma(V,\mathcal O_X)$, $m\geq1$, $\bigcap_kD(f_k)=D(f_1\cdots f_m)$, and it is affine.
\end{enumerate}
\end{lemma}
\leanref{Development/AutomorphicSpectrum/Proofs/BlockDecomposition/GradedFormCover.lean}{GlobalTheory.CompanionTheory.GrForm.AffCover.stCover}
\leanref{Development/AutomorphicSpectrum/Proofs/BlockDecomposition/AffineGradedQuotientVanishAboveDim.lean}{GlobalTheory.CompanionTheory.GrTop.exists_cover}
\leanref{Development/AutomorphicSpectrum/Proofs/BlockDecomposition/GradedFormCover.lean}{GlobalTheory.CompanionTheory.GrForm.AffCover.iInf_basicOpen}

\begin{proof}
(1) Let $G=\Gamma(V,\mathcal O_X)$. As $X$ is smooth over $\mathcal O$ and $V$ is affine, $G$ is a smooth $\mathcal
O$-algebra \cite[Tag 01V6]{Sta}, so it is locally standard smooth \cite[Tag 00TA]{Sta}: the $g\in G$ such that $G_g$ is a
standard smooth $\mathcal O$-algebra generate the unit ideal of $G$. So the $D(g)$ for these $g$ cover $V$.

Let $\xi\in V$, and let $f$ be one of these $g$ with $\xi\in D(f)$. Write $G_f=\mathcal
O[x_1,\ldots,x_N]/(h_1,\ldots,h_c)$ with $\det(\partial h_k/\partial x_l)_{1\leq k,l\leq c}$ invertible in $G_f$, and let
$e=N-c$. Over $P=\mathcal O[x_{c+1},\ldots,x_N]$, $G_f=P[x_1,\ldots,x_c]/(h_1,\ldots,h_c)$ is a standard smooth algebra
with as many variables as equations, so it is smooth over $P$ with $\Omega_{G_f/P}=0$ \cite[Tag 00T7]{Sta}, that is,
\'etale over $P$ \cite[Tag 00U1]{Sta}, a polynomial algebra in $e$ variables. As $V$ is affine, restriction identifies
$\Gamma(D(f),\mathcal O_X)$ with $G_f$ \cite[Tag 01HV]{Sta}. So the images of $x_{c+1},\ldots,x_N$ are coordinates of
dimension $e$ on the affine open $D(f)$. If $X$ is smooth of relative dimension $d$, then $e=d$: as $\xi\in D(f)$,
$G_f\neq0$; $\Omega_{G_f/\mathcal O}$ is free of rank $e$ \cite[Tag 00T7]{Sta}, and it is $\Gamma(D(f),\Omega_{X/\mathcal
O})$ \cite[Tag 01UT]{Sta}, where $\Omega_{X/\mathcal O}$ is locally free of rank $d$, as the fibres of $X$ have dimension
$d$ \cite[Tag 02G1]{Sta}.

As $V$ is quasi-compact, finitely many of these $D(f)$, say $D(f_\alpha)$ for $\alpha\in A$, cover $V$.

(2) A point of $V$ lies in every $D(f_k)$ if and only if every $f_k$ is a unit in its local ring, that is, if and only if
$f_1\cdots f_m$ is. So $\bigcap_kD(f_k)=D(f_1\cdots f_m)$, a basic open of the affine scheme $V$, which is affine.
% Lean: (1) is AffCover.stCover (the loc step: RingHom.smooth_iff_locally_isStandardSmooth for q.smooth_appLE,
% CovEx.exists_basicOpen, exists_etale_mvPolynomial and CovEx.coord_basicOpen, then elim_finite_subcover), and, for
% relative dimension d, GrTop.exists_cover, which gets local standard smoothness of relative dimension d directly from
% Mathlib's SmoothOfRelativeDimension (HasRingHomProperty.appLE). (2) is AffCover.iInf_basicOpen and GrTop.iInf_basicOpen
% (induction with Scheme.basicOpen_mul) and IsAffineOpen.basicOpen. About 90 lines.
\end{proof}
\leanref{Development/AutomorphicSpectrum/Proofs/BlockDecomposition/ExistsSmoothCoordinateCover.lean}{GlobalTheory.CompanionTheory.CovEx.exists_basicOpen}
\leanref{Development/AutomorphicSpectrum/Proofs/BlockDecomposition/ExistsSmoothCoordinateCover.lean}{GlobalTheory.CompanionTheory.CovEx.coord_basicOpen}

\begin{lemma}\label{gal-realization-kappa-de-rham-affine}
Let $K$ be a field of characteristic $0$, $X$ a smooth $K$-scheme, $V\subset X$ an affine open, $0\leq r<s$ and $j\neq r$.
Then $H^j(\mathrm{Th}^{(s)}_{\mathrm{dR}}(X)_V,\mathrm{gr}_r)=0$
(Definition~\ref{gal-realization-kappa-de-rham-truncated}).
\end{lemma}
\leanref{Development/WeilBound/CoreStatements/PadicHodgeComparison/EtaleCohomologyComparison/PerfectoidPullbacks/GrCechFin.lean}{CrystallineSemistableComparison.Pb.gr_aff}
\leanref{Development/AutomorphicSpectrum/Proofs/BlockDecomposition/AffineGradedCohomologyVanishingStatement.lean}{GlobalTheory.T.affine_graded_cohomology_vanishing}
\leanref{Development/AutomorphicSpectrum/Proofs/BlockDecomposition/AffineGradedCohomologyVanishing.lean}{GlobalTheory.CompanionTheory.GrForm.gr_aff}

\begin{proof}
By Lemma~\ref{gal-realization-kappa-de-rham-affine-basic}(1) there are finitely many basic opens $U_\alpha=D(f_\alpha)$,
$\alpha\in A$, of $V$ that cover $V$, each with coordinates. For $i\in A^{n+1}$, $n\geq0$, $U_i=\bigcap_kU_{i_k}$ is the
basic open $D(\prod_kf_{i_k})$ of $V$, so it is affine (Lemma~\ref{gal-realization-kappa-de-rham-affine-basic}(2)). Let
$\mathcal U=\{U_\alpha:\alpha\in A\}$,
$\mathcal Q=\mathrm{Th}^{(s)}_{\mathrm{dR}}(X)_{\mathcal U}$ and $\mathcal C_i=\mathrm{Th}^{(s)}_{\mathrm{dR}}(X)_{U_i}$.
By Lemma~\ref{gal-realization-kappa-de-rham-zariski} and Lemma~\ref{gal-realization-thickening-zariski}(1), with
$\mathcal V=\mathcal U$ and $\mathcal W=\{V\}$, restriction $H^j(\mathrm{Th}^{(s)}_{\mathrm{dR}}(X)_V,\mathrm{gr}_r)\to
H^j(\mathcal Q,\mathrm{gr}_r)$ is bijective. So it suffices to show $H^j(\mathcal Q,\mathrm{gr}_r)=0$.

By Lemma~\ref{gal-realization-thickening-over}, the categories $\mathcal C_{(\alpha)}$ are cosieves of $\mathcal Q$ that
cover it, with intersections the $\mathcal C_i$. The class of zero modules satisfies the hypothesis of
Lemma~\ref{gal-realization-thickening-level-one-serre}, so by Lemma~\ref{gal-realization-thickening-level-one-cech} it
suffices that $H^a(\check C^\bullet_b)=0$ for $a+b=j$, where $\check C^a_b=\prod_{i\in A^{a+1}}H^b(\mathcal
C_i,\mathrm{gr}_r)$.

If $b\neq r$, then $\check C^a_b=0$: $U_i$ is an affine open of $U_{i_0}$, which has coordinates, so $H^b(\mathcal
C_i,\mathrm{gr}_r)=0$ by Lemma~\ref{gal-realization-kappa-de-rham-local}(1).

If $b=r$, then $a=j-r\geq1$, as $j\neq r$. By Lemma~\ref{gal-realization-kappa-de-rham-bundle}, $\check C^\bullet_r$ is the
complex of a cosimplicial $K$-module with data (i), (ii) and (iii) of
Lemma~\ref{gal-realization-thickening-level-one-coherent} for the opens $U_\alpha$. As in the proof of that lemma, these
data glue the free modules $\mathcal O_{U_\alpha}^{\kappa_\alpha}$ along the matrices $g_{\alpha\beta}$ to a locally free
$\mathcal O_V$-module $\mathcal E$ of finite rank whose \v{C}ech complex for the cover $(U_\alpha)$ of $V$ is $\check
C^\bullet_r$. As the $U_i$ are affine, $H^a(\check C^\bullet_r)=H^a(V,\mathcal E)$ \cite[Tag 01XD]{Sta}, which vanishes
for $a\geq1$, as $V$ is affine and $\mathcal E$ is quasi-coherent \cite[Tag 01XB]{Sta}.
% Lean: GrForm.gr_aff = GrForm.main, with StCover = graded_form_cover (AffCover.stCover: V is covered by basic opens
% with coordinates, whose intersections are basic opens of V; Lemma gal-realization-kappa-de-rham-affine-basic),
% cech_refinement_cohomology_bijective (injective restriction, hz), GlobSerreSt with zeroSerre, cech_van (StVan with
% coordinates_restrict_affine_open and capT_eq) and cech_pos (StVB with affine_cover_higher_cohomology_vanishes =
% VBAff.subsingleton_H, the Lean's own proof of the vanishing, where the book cites Stacks).
\end{proof}
\leanref{Development/AutomorphicSpectrum/Proofs/BlockDecomposition/AffineGradedCohomologyVanishing.lean}{GlobalTheory.CompanionTheory.GrForm.main}
\leanref{Development/AutomorphicSpectrum/Proofs/BlockDecomposition/AffineGradedCohomologyVanishing.lean}{GlobalTheory.CompanionTheory.GrForm.cech_van}
\leanref{Development/AutomorphicSpectrum/Proofs/BlockDecomposition/AffineGradedCohomologyVanishing.lean}{GlobalTheory.CompanionTheory.GrForm.cech_pos}
\leanref{Development/AutomorphicSpectrum/Proofs/BlockDecomposition/GradedFormCover.lean}{GlobalTheory.CompanionTheory.GrForm.AffCover.stCover}
\leanref{Development/AutomorphicSpectrum/Proofs/BlockDecomposition/AffineCoverHigherCohomologyVanishesStatement.lean}{GlobalTheory.T.affine_cover_higher_cohomology_vanishes}

\begin{lemma}\label{gal-realization-kappa-de-rham-tower}
Let $\mathcal C$ be a category, $A$ a functor from $\mathcal C$ to commutative rings, and $(\mathfrak a_k)_{k\geq0}$
families of ideals $\mathfrak a_k(T)\subset A(T)$, with $A(\varphi)(\mathfrak a_k(T))\subset\mathfrak a_k(T')$ for maps
$\varphi:T\to T'$, such that $\mathfrak a_{k+1}(T)\subset\mathfrak a_k(T)$ for all $k$ and $T$, and every object $T$ has
some $k$ with $\mathfrak a_k(T)=0$. If $H^j(\mathcal C,\mathfrak a_k/\mathfrak a_{k+1})=0$ for all $k$ and $j$, then
$H^j(\mathcal C,\mathfrak a_0)=0$ for every $j$ (Definition~\ref{gal-realization-split-diagrams}(1)).
\end{lemma}
\leanref{Development/AutomorphicSpectrum/Proofs/BlockDecomposition/CohomologyVanishingAlongIdealFiltrationStatement.lean}{GlobalTheory.T.cohomology_vanishing_along_ideal_filtration}
\leanref{Development/AutomorphicSpectrum/Proofs/BlockDecomposition/CohomologyVanishingAlongIdealFiltration.lean}{GlobalTheory.CompanionTheory.TowerEnd.vanish}

\begin{proof}
For an object $T$ let $K(T)$ be some $k$ with $\mathfrak a_k(T)=0$; then $\mathfrak a_k(T)=0$ for $k\geq K(T)$. A cochain
\emph{has values in} $\mathfrak a_k$ if its value at every chain $c_0\to\cdots\to c_n$ lies in $\mathfrak a_k(c_n)$;
then so does its coboundary.

\emph{Degree $0$.} Let $z$ be a $0$-cocycle with values in $\mathfrak a_0$. If $z$ has values in $\mathfrak a_k$, the
classes of its values form a $0$-cocycle with values in $\mathfrak a_k/\mathfrak a_{k+1}$, which is $0$ as
$H^0(\mathcal C,\mathfrak a_k/\mathfrak a_{k+1})=0$ and there are no coboundaries in degree $0$; so $z$ has values in
$\mathfrak a_{k+1}$. By induction $z$ has values in every $\mathfrak a_k$, so $z(T)\in\mathfrak a_{K(T)}(T)=0$.

\emph{Degree $n+1$.} If $z'$ is an $(n+1)$-cocycle with values in $\mathfrak a_k$, the classes of its values form a
cocycle with values in $\mathfrak a_k/\mathfrak a_{k+1}$, which is the coboundary of an $n$-cochain with values in
$\mathfrak a_k/\mathfrak a_{k+1}$, as $H^{n+1}(\mathcal C,\mathfrak a_k/\mathfrak a_{k+1})=0$. Lifting its values gives an
$n$-cochain $y$ with values in $\mathfrak a_k$ such that $z'-dy$ has values in $\mathfrak a_{k+1}$. Let $z$ be an
$(n+1)$-cocycle with values in $\mathfrak a_0$. Applying this repeatedly gives $n$-cochains $y_0,y_1,\ldots$, with $y_k$
having values in $\mathfrak a_k$, such that $z-d(y_0+\cdots+y_{k-1})$ has values in $\mathfrak a_k$ for every $k$. At a
chain $s=(c_0\to\cdots\to c_n)$, $y_k(s)=0$ for $k\geq K(c_n)$, so $Y(s)=\sum_ky_k(s)$ is a finite sum, and $Y$ is an
$n$-cochain with values in $\mathfrak a_0$. Let $s=(c_0\to\cdots\to c_{n+1})$, and let $N$ be at least $K(c_{n+1})$ and
the $K$ of the last objects of the faces of $s$. The value of $z-d(y_0+\cdots+y_{N-1})$ at $s$ lies in $\mathfrak
a_N(c_{n+1})=0$, and $d(y_0+\cdots+y_{N-1})$ and $dY$ agree at $s$, as $y_0+\cdots+y_{N-1}$ and $Y$ agree at the faces of
$s$. So $z(s)=dY(s)$, and $z=dY$.
% Lean: TowerEnd.vanish, through step_zero for degree 0, and step_succ, exists_seq (the y_k) and exists_primitive (Y,
% with N the maximum over the faces) for degree n + 1.
\end{proof}
\leanref{Development/AutomorphicSpectrum/Proofs/BlockDecomposition/CohomologyVanishingAlongIdealFiltration.lean}{GlobalTheory.CompanionTheory.TowerEnd.step_zero}
\leanref{Development/AutomorphicSpectrum/Proofs/BlockDecomposition/CohomologyVanishingAlongIdealFiltration.lean}{GlobalTheory.CompanionTheory.TowerEnd.step_succ}
\leanref{Development/AutomorphicSpectrum/Proofs/BlockDecomposition/CohomologyVanishingAlongIdealFiltration.lean}{GlobalTheory.CompanionTheory.TowerEnd.exists_seq}
\leanref{Development/AutomorphicSpectrum/Proofs/BlockDecomposition/CohomologyVanishingAlongIdealFiltration.lean}{GlobalTheory.CompanionTheory.TowerEnd.exists_primitive}

\begin{lemma}\label{gal-realization-kappa-de-rham-restriction}
Let $(R,I,\gamma,0)$, $\mathcal O$, $X$ and $\kappa$ be as in Definition~\ref{gal-realization-kappa-de-rham-truncated},
$s,\ell\geq1$ integers, $\mathcal Q_0$ either $\mathrm{Th}(X/R)$, if $\ell=1$, or $\mathrm{Th}^{(s')}_\ell(X/R)$ for some
$s'\geq s$ (Lemma~\ref{gal-realization-kappa-truncated-periods-exist-base-change-zariski}), $\mathcal V$ a set of open
subsets of $X$, and $r,r'\in\mathbf Z$ with $r'\leq s$. For every $j$, restriction $H^j((\mathcal Q_0)_{\mathcal
V},\mathcal O^{[r]}/\mathcal O^{[r']})\to H^j(\mathrm{Th}^{(s)}_\ell(X/R)_{\mathcal V},\mathcal O^{[r]}/\mathcal O^{[r']})$
is bijective (Definition~\ref{gal-realization-thickening-nilpotent}(3)).
\end{lemma}
\leanref{Development/AutomorphicSpectrum/Proofs/BlockDecomposition/FilteredCohomologyRestrictionBijectiveStatement.lean}{GlobalTheory.T.filtered_cohomology_restriction_bijective}
\leanref{Development/AutomorphicSpectrum/Proofs/BlockDecomposition/FilteredCohomologyRestrictionBijective.lean}{GlobalTheory.CompanionTheory.PDTrunc.pull_incl_bijective}

\begin{proof}
Let $\mathcal Q'=(\mathcal Q_0)_{\mathcal V}$ and $\mathcal Q=\mathrm{Th}^{(s)}_\ell(X/R)_{\mathcal V}$, a full subcategory
of $\mathcal Q'$, as $J^{[s']}\subset J^{[s]}$; let $\iota$ be the inclusion and $F=\mathcal O^{[r]}/\mathcal O^{[r']}$.

\emph{The truncation.} Let $T=(B,i,J,\delta,y)$ be in $\mathcal Q'$. By
Lemma~\ref{gal-realization-kappa-de-rham-subideal}, for $\delta$ and $m=s$, $J^{[s]}\subset J$, and $\delta_n(z)\in
J^{[s]}$ for every $z\in J^{[s]}$ and $n\geq1$. So, as in the proof of Lemma~\ref{gal-realization-kappa-de-rham-poly},
with $\delta$ in place of $\eta$, there is exactly one divided power structure $\bar\delta$ on the image $\bar J$ of $J$
in $\bar B=B/J^{[s]}$ for which the projection $\lambda_T$ is a divided power map \cite[Tags 07H2 and 07H1]{Sta};
divided power maps out of $B$ that kill $J^{[s]}$ induce divided power maps out of $\bar B$; and
$\bar J^{[t]}=\lambda_T(J^{[t]})$ for $t\geq0$. So $\bar J^{[s]}=0$, $\bar J$ is nilpotent as $J$ is, and $L(T)=(\bar
B,\lambda_Ti,\bar J,\bar\delta,y)$, with $\bar B/\bar J=B/J$, is a thickening in $\mathcal Q$ that lies over the opens
$T$ lies over, and $\lambda_T$ is a map $T\to L(T)$. For $L(T)\in\mathcal Q$ it remains to check the condition on $\ell$,
which is empty if $\ell=1$; otherwise $T$ satisfies it. Let $r_0\geq0$, $m=\min(r_0,s)$ and $x\in B$ with
$\ell\lambda_T(x)\in\bar J^{[r_0]}$. Then $\lambda_T(\ell x)\in\bar J^{[m]}=\lambda_T(J^{[m]})$, and as $J^{[s]}\subset
J^{[m]}$, $\ell x\in J^{[m]}$, so $x\in J^{[m]}$. If $r_0\leq s$, this gives $\lambda_T(x)\in\bar J^{[r_0]}$; if $r_0>s$, then $x\in J^{[s]}$ and
$\lambda_T(x)=0$. A map $\varphi:T\to T'$ carries $J^{[s]}$ into $J_{T'}^{[s]}$, so it induces a map $L(\varphi):L(T)\to
L(T')$, and $L$ is a functor $\mathcal Q'\to\mathcal Q$.

\emph{The adjunction.} For $T\in\mathcal Q'$ and $T''\in\mathcal Q$, every map $\psi:T\to T''$ carries $J^{[s]}$ into
$J_{T''}^{[s]}=0$, so $\psi=\chi\lambda_T$ for exactly one ring map $\chi:\bar B\to B_{T''}$, which is a map $L(T)\to T''$:
it is a divided power map, commutes with the maps from $R$ and is compatible with the points, as $\psi$ is. This
bijection is natural in $T$ and $T''$, so $L$ is left adjoint to $\iota$, with unit $\lambda$.

\emph{The unit on $F$.} For $T\in\mathcal Q'$, $F(\lambda_T):F(T)\to F(L(T))$ is bijective. If $r'\leq0$, both sides are
$0$. Let $r'\geq1$. Then $J^{[s]}\subset J^{[r']}$, so $\bar B/\bar J^{[r']}=B/J^{[r']}$, and $F(L(T))$, the image of
$\bar J^{[r]}=\lambda_T(J^{[r]})$ in it (with $J^{[r]}=B$ if $r\leq0$), is $F(T)$.

\emph{The conclusion.} Restriction along $\iota$ followed by the map $H^j(\mathcal Q,F)\to H^j(\mathcal Q',F\circ\iota L)$
induced by $L$ is induced by $\iota L$ with the identity maps of the modules $F(L(T))$. By
Lemma~\ref{gal-realization-period-prism}, for the natural transformation $\lambda$ from the identity functor to $\iota L$,
it equals the map induced by the identity functor with the maps $F(\lambda_T)$, which is bijective, as the $F(\lambda_T)$
form an isomorphism of functors. The map induced by $L$ is bijective by
Lemma~\ref{gal-realization-thickening-torsor-adjoint}(1), as $L$ has the right adjoint $\iota$. So restriction along
$\iota$ is bijective.
% Lean: PDTrunc.pull_incl_bijective, through pull_bij_of_adj (P.adj, P.prism and coef_upTo); L is trunc with qT (the
% sub-pd ideal sub_inf from fil_isSubDPIdeal, Quotient.dividedPowers), unitH, counitH and qMap, the adjunction is
% truncAdj, F(λ_T) bijective is unit_bij (pdPow_q, mem_of_mk), and L(T) ∈ Q is q_bot and inU_q (q_tf for the Lean's
% l-torsion condition, empty for l = 1). The Lean treats s = 0 separately, where both groups vanish.
\end{proof}
\leanref{Development/AutomorphicSpectrum/Proofs/BlockDecomposition/FilteredCohomologyRestrictionBijective.lean}{GlobalTheory.CompanionTheory.PDTrunc.pull_bij_of_adj}
\leanref{Development/AutomorphicSpectrum/Proofs/BlockDecomposition/FilteredCohomologyRestrictionBijective.lean}{GlobalTheory.CompanionTheory.PDTrunc.qT}
\leanref{Development/AutomorphicSpectrum/Proofs/BlockDecomposition/FilteredCohomologyRestrictionBijective.lean}{GlobalTheory.CompanionTheory.PDTrunc.truncAdj}
\leanref{Development/AutomorphicSpectrum/Proofs/BlockDecomposition/FilteredCohomologyRestrictionBijective.lean}{GlobalTheory.CompanionTheory.PDTrunc.unit_bij}
\leanref{Development/AutomorphicSpectrum/Proofs/BlockDecomposition/FilteredCohomologyRestrictionBijective.lean}{GlobalTheory.CompanionTheory.PDTrunc.sub_inf}
\leanref{Development/AutomorphicSpectrum/Proofs/BlockDecomposition/FilteredCohomologyRestrictionBijective.lean}{GlobalTheory.CompanionTheory.PDTrunc.pdPow_q}

\begin{lemma}\label{gal-realization-kappa-de-rham-truncation}
Let $K$ be a field of characteristic $0$, $X$ a $K$-scheme with a coordinate cover $(U_\alpha)_{\alpha\in A}$ of dimensions
$d_\alpha$ (Definition~\ref{gal-realization-thickening-coordinates}(2)), $s\geq1$ with $d_\alpha<s$ for every $\alpha$,
and $\mathcal V$ a set of open subsets of $X$ with union $X$. For every $j$, restriction $H^j_{\mathrm{dR}}(X)\to
H^j(\mathrm{Th}^{(s)}_{\mathrm{dR}}(X)_{\mathcal V},\mathcal O)$ is bijective
(Definition~\ref{gal-realization-kappa-de-rham-truncated}).
\end{lemma}
\leanref{Development/WeilBound/CoreStatements/PadicHodgeComparison/EtaleCohomologyComparison/PerfectoidPullbacks/GrCechFin.lean}{CrystallineSemistableComparison.Pb.k_iso}
\leanref{Development/WeilBound/CoreStatements/PadicHodgeComparison/EtaleCohomologyComparison/PerfectoidPullbacks/GlobalSerreRelations/GlobSerreSt.lean}{CrystallineSemistableComparison.Pb.KH.KIso}
\leanref{Development/AutomorphicSpectrum/Proofs/BlockDecomposition/HodgeFiltrationViaTruncationStatement.lean}{GlobalTheory.T.hodge_filtration_via_truncation}
\leanref{Development/AutomorphicSpectrum/Proofs/BlockDecomposition/HodgeFiltrationViaTruncation.lean}{GlobalTheory.CompanionTheory.KIso.k_iso}

\begin{proof}
Write $\mathrm{Th}=\mathrm{Th}_{\mathrm{dR}}(X)=\mathrm{Th}(X/K)$ and $\mathrm{Th}^{(t)}=\mathrm{Th}^{(t)}_{\mathrm{dR}}(X)$,
and let $\mathcal U=\{U_\alpha:\alpha\in A\}$. Every thickening lies over $X$
(Lemma~\ref{gal-realization-thickening-over}(2)), so $\mathrm{Th}=\mathrm{Th}_{\{X\}}$.

\emph{Step 1: $H^j(\mathrm{Th},\mathrm{gr}_r)=0$ for every $j$ and every $r\geq1$ with $r>d_\alpha$ for all $\alpha$.} By
Lemma~\ref{gal-realization-kappa-de-rham-restriction}, for $\ell=1$, $\mathcal Q_0=\mathrm{Th}$, $s=r+1$, $\mathcal V=\{X\}$
and
$r'=r+1$, restriction $H^j(\mathrm{Th},\mathrm{gr}_r)\to H^j(\mathrm{Th}^{(r+1)}_{\{X\}},\mathrm{gr}_r)$ is bijective. By
Lemma~\ref{gal-realization-kappa-de-rham-zariski} and Lemma~\ref{gal-realization-thickening-zariski}(1), for $\mathcal
Q_0=\mathrm{Th}^{(r+1)}$, $\mathcal V=\mathcal U$ and $\mathcal W=\{X\}$, so is restriction to
$\mathrm{Th}^{(r+1)}_{\mathcal U}$. By Lemma~\ref{gal-realization-thickening-over}, the thickenings over the $U_\alpha$
form cosieves of $\mathrm{Th}^{(r+1)}_{\mathcal U}$ that cover it, with intersections $\mathrm{Th}^{(r+1)}_{U_S}$ for
nonempty $S\subset A$, where $U_S=\bigcap_{\alpha\in S}U_\alpha$. The affine open $U_S$ of $U_\alpha$, $\alpha\in S$, has
coordinates of dimension $d_\alpha<r$, restricted from $U_\alpha$, so $H^j(\mathrm{Th}^{(r+1)}_{U_S},\mathrm{gr}_r)=0$
for every $j$ by Lemma~\ref{gal-realization-kappa-de-rham-local}(3), with $s=r+1$. By
Lemma~\ref{gal-realization-thickening-cosieves}(3), $H^j(\mathrm{Th}^{(r+1)}_{\mathcal U},\mathrm{gr}_r)=0$ for every
$j$, and so $H^j(\mathrm{Th},\mathrm{gr}_r)=0$.

\emph{Step 2: $H^j(\mathrm{Th},\mathcal O^{[s]})=0$ for every $j$.} Apply
Lemma~\ref{gal-realization-kappa-de-rham-tower} to $A=\mathcal O$ and $\mathfrak a_k=\mathcal O^{[s+k]}$, which decrease
in $k$. For a thickening $T$, $J_T$ is nilpotent, say $J_T^N=0$, and $\delta_n(x)=x^n/n!\in J_T^n$ for $x\in J_T$, as
$n!$ is invertible in the $K$-algebra $B_T$; so each generator $\prod_i\delta_{k_i}(x_i)$ of $J_T^{[N]}$ lies in
$J_T^N=0$, and $\mathfrak a_k(T)=0$ for $s+k\geq N$. The quotients $\mathfrak a_k/\mathfrak a_{k+1}$ are the
$\mathrm{gr}_{s+k}$, whose cohomology vanishes by Step 1, as $s+k\geq s>d_\alpha$.

\emph{Step 3.} By Step 2 and the long exact sequence of $0\to\mathcal O^{[s]}\to\mathcal O\to\mathcal O/\mathcal
O^{[s]}\to0$, the map $H^j(\mathrm{Th},\mathcal O)\to H^j(\mathrm{Th},\mathcal O/\mathcal O^{[s]})$ is bijective. By
Lemma~\ref{gal-realization-kappa-de-rham-restriction}, for $\ell=1$, $\mathcal Q_0=\mathrm{Th}$, $\mathcal V=\{X\}$, $r=0$
and
$r'=s$, and by Lemma~\ref{gal-realization-kappa-de-rham-zariski} with
Lemma~\ref{gal-realization-thickening-zariski}(1), for $\mathcal Q_0=\mathrm{Th}^{(s)}$, $\mathcal V$ and $\mathcal
W=\{X\}$, restriction $H^j(\mathrm{Th},\mathcal O/\mathcal O^{[s]})\to H^j(\mathrm{Th}^{(s)}_{\mathcal V},\mathcal
O/\mathcal O^{[s]})$ is bijective. On $\mathrm{Th}^{(s)}_{\mathcal V}$, $\mathcal O^{[s]}=0$, so $\mathcal O\to\mathcal
O/\mathcal O^{[s]}$ is an isomorphism of functors there. Restriction commutes with the maps induced by $\mathcal
O\to\mathcal O/\mathcal O^{[s]}$, as both composites are induced by the same functor and maps. So restriction
$H^j(\mathrm{Th},\mathcal O)\to H^j(\mathrm{Th}^{(s)}_{\mathcal V},\mathcal O)$ is bijective.
% Lean: KIso.k_iso, first part. Step 1 is thick_vanish (filtered_cohomology_restriction_bijective as b1,
% cech_refinement_cohomology_bijective as b2, restrict_bij as b0, and glob_vanish through GlobVanSt and loc_vanish);
% Step 2 is tower_vanish (cohomology_vanishing_along_ideal_filtration, with pdPow_nil and
% dividedPowerFiltration_le_pow); Step 3 is toSq_bij, chain_bij, toSq_bij_bot and chain_comm. The Lean's KIso also
% characterizes the filtration FilR through quoK, which this book does not use.
\end{proof}
\leanref{Development/AutomorphicSpectrum/Proofs/BlockDecomposition/HodgeFiltrationViaTruncation.lean}{GlobalTheory.CompanionTheory.KIso.thick_vanish}
\leanref{Development/AutomorphicSpectrum/Proofs/BlockDecomposition/HodgeFiltrationViaTruncation.lean}{GlobalTheory.CompanionTheory.KIso.glob_vanish}
\leanref{Development/AutomorphicSpectrum/Proofs/BlockDecomposition/HodgeFiltrationViaTruncation.lean}{GlobalTheory.CompanionTheory.KIso.tower_vanish}
\leanref{Development/AutomorphicSpectrum/Proofs/BlockDecomposition/HodgeFiltrationViaTruncation.lean}{GlobalTheory.CompanionTheory.KIso.pdPow_nil}
\leanref{Development/AutomorphicSpectrum/Proofs/BlockDecomposition/HodgeFiltrationViaTruncation.lean}{GlobalTheory.CompanionTheory.KIso.dividedPowerFiltration_le_pow}
\leanref{Development/AutomorphicSpectrum/Proofs/BlockDecomposition/HodgeFiltrationViaTruncation.lean}{GlobalTheory.CompanionTheory.KIso.toSq_bij}
\leanref{Development/AutomorphicSpectrum/Proofs/BlockDecomposition/HodgeFiltrationViaTruncation.lean}{GlobalTheory.CompanionTheory.KIso.chain_bij}
\leanref{Development/AutomorphicSpectrum/Proofs/BlockDecomposition/HodgeFiltrationViaTruncation.lean}{GlobalTheory.CompanionTheory.KIso.toSq_bij_bot}
\leanref{Development/AutomorphicSpectrum/Proofs/BlockDecomposition/HodgeFiltrationViaTruncation.lean}{GlobalTheory.CompanionTheory.KIso.chain_comm}

\begin{lemma}\label{gal-realization-kappa-de-rham-log}
Let $\delta$ be divided powers on an ideal $J$ of a ring $B$, and $m\geq0$. For $x\in J$ let
$\log_m(x)=\sum_{k=0}^{m-1}(-1)^kk!\,\delta_{k+1}(x)$.
\begin{enumerate}
\item $\log_m(x)\in J$ for $x\in J$. If $\varphi:B\to B'$ is a divided power map to a ring $B'$ with divided powers
  $\delta'$ on an ideal $J'$, then $\varphi(\log_m(x))=\log_m(\varphi(x))$, formed for $\delta'$.
\item If $J^{[m]}=0$ (Definition~\ref{gal-realization-kappa-approximations}) and $m\leq m'$, then
  $\log_m(x)=\log_{m'}(x)$ for $x\in J$.
\item If $J^{[m]}=0$ and every integer $n\geq1$ is invertible in $B$, then $\log_m((1+x)(1+y)-1)=\log_m(x)+\log_m(y)$
  for $x,y\in J$.
\end{enumerate}
\end{lemma}
\leanref{Development/WeilBound/CoreStatements/PadicHodgeComparison/EtaleCohomologyComparison/PerfectoidPullbacks/ConeSectionRestriction/RSubOneMem.lean}{CrystallineSemistableComparison.Pb.plogK_mem}
\leanref{Development/AutomorphicSpectrum/Proofs/BlockDecomposition/ConeCocycleClosed.lean}{GlobalTheory.CompanionTheory.ConeZcoc.plogK_map}
\leanref{Development/WeilBound/CoreStatements/PadicHodgeComparison/EtaleCohomologyComparison/PerfectoidPullbacks/ConeSectionRestriction/RSubOneMem.lean}{CrystallineSemistableComparison.Pb.plogK_congr}
\leanref{Development/AutomorphicSpectrum/Proofs/BlockDecomposition/TruncatedLogMultiplicativeStatement.lean}{GlobalTheory.T.truncated_log_multiplicative}
\leanref{Development/AutomorphicSpectrum/Proofs/BlockDecomposition/TruncatedLogMultiplicative.lean}{GlobalTheory.CompanionTheory.PlogMul.plogK_mul}

\begin{proof}
(1) Each term $(-1)^kk!\,\delta_{k+1}(x)$ lies in $J$. A divided power map commutes with the $\delta_{k+1}$ and fixes the
integers.

(2) The terms of $\log_{m'}(x)$ with $k\geq m$ contain $\delta_{k+1}(x)\in J^{[k+1]}\subset J^{[m]}=0$.

(3) For $z\in J$, $(k+1)\cdot k!\,\delta_{k+1}(z)=z^{k+1}$, so $\log_m(z)=\ell(z)$ for the polynomial
$\ell(T)=\sum_{k=0}^{m-1}(-1)^kT^{k+1}/(k+1)$ of $B[T]$, whose derivative is $\ell'(T)=\sum_{k=0}^{m-1}(-T)^k$. For $a+b\geq
m$, $x^ay^b=a!\,b!\,\delta_a(x)\delta_b(y)=0$, as $\delta_a(x)\delta_b(y)\in J^{[m]}$. In $B[T]$ let $u=xT$,
$Z=(1+u)(1+y)-1$ and $F=\ell(Z)-\ell(u)-\ell(y)$. Then $u^m=0$, and $Z^m=0$, as $Z=u(1+y)+y$ and each term of the binomial
expansion of $Z^m$ contains a product $x^iy^{m-i}$. By $(1+v)\sum_{k=0}^{m-1}(-v)^k=1-(-v)^m$, $(1+u)\ell'(u)=1$ and
$(1+Z)\ell'(Z)=1$, where $1+Z=(1+u)(1+y)$. So $F'=x(1+y)\ell'(Z)-x\ell'(u)$ has $(1+u)F'=x-x=0$, and $F'=0$, as $1+u$ is
a unit. As the integers are invertible in $B$, $F$ is constant, so $F=F(0)=\ell(y)-\ell(0)-\ell(y)=0$. At $T=1$ this is
the claim, as $(1+x)(1+y)-1\in J$.
% Lean: (1) is plogK_mem and plogK_map, (2) is plogK_congr, and (3) is PlogMul.plogK_mul: hlog (log_m = logPoly), hxy,
% and eval_logPoly_mul, with derivative_logPoly, one_add_mul_geom and the constancy of G from hderiv.
\end{proof}
\leanref{Development/AutomorphicSpectrum/Proofs/BlockDecomposition/TruncatedLogMultiplicative.lean}{GlobalTheory.CompanionTheory.PlogMul.eval_logPoly_mul}
\leanref{Development/AutomorphicSpectrum/Proofs/BlockDecomposition/TruncatedLogMultiplicative.lean}{GlobalTheory.CompanionTheory.PlogMul.derivative_logPoly}

\begin{lemma}\label{gal-realization-kappa-de-rham-lift}
Let $K$ be a ring, $X$ a $K$-scheme, $\mathcal L$ a line datum on $X$ with morphism $f:X_{\mathcal L}\to X$, $\mathcal V$
the set of open subsets $V\subset X$ over which $f$ has a section (a morphism $\sigma:V\to X_{\mathcal L}$ with $f\sigma$
the inclusion of $V$), and $s\geq1$. Form $\mathrm{Th}(X/K)$ and $\mathrm{Th}^{(s)}(X/K)_{\mathcal V}$ for
$(K,0,0,0)$, with $\mathcal O=K$ and $\kappa$ the identity
(Definitions~\ref{gal-realization-thickening-nilpotent} and~\ref{gal-realization-kappa-de-rham-truncated}(1)); for a
field $K$ they are $\mathrm{Th}_{\mathrm{dR}}(X)$ and $\mathrm{Th}^{(s)}_{\mathrm{dR}}(X)_{\mathcal V}$
(Definition~\ref{gal-realization-kappa-de-rham-truncated}(2)). Then $\mathcal L$ has a lift along the inclusion of
$\mathrm{Th}^{(s)}(X/K)_{\mathcal V}$ into $\mathrm{Th}(X/K)$, defined as in
Definition~\ref{gal-realization-period-chern}(3).
\end{lemma}
\leanref{Development/AutomorphicSpectrum/Proofs/BlockDecomposition/ExistsConeLiftStatement.lean}{GlobalTheory.T.exists_cone_lift}
\leanref{Development/AutomorphicSpectrum/Proofs/BlockDecomposition/ExistsConeLift.lean}{GlobalTheory.CompanionTheory.ConeLiftEx.nonempty_cLiftK}
\leanref{Development/WeilBound/CoreStatements/PadicHodgeComparison/EtaleCohomologyComparison/PerfectoidPullbacks/ConeSectionRestriction/RSubOneMem.lean}{CrystallineSemistableComparison.Pb.lift_exists}

\begin{proof}
Let $T=(B,i,J,\delta,y)$ be in $\mathrm{Th}^{(s)}(X/K)_{\mathcal V}$. It lies over some $V\in\mathcal V$, with a
section $\sigma$; so $y$ is a morphism $y':\operatorname{Spec}B/J\to V$ followed by the inclusion, and $p_T=\sigma y'$ has
$fp_T=y$. For a map $\varphi:T\to T'$, $p_T\circ\operatorname{Spec}\bar\varphi$ and $p_{T'}$ have the same composite
$y_{T'}$ with $f$, so $u(p_T\circ\operatorname{Spec}\bar\varphi,p_{T'})$ is a unit of $B_{T'}/J_{T'}$. It lifts to a unit
$w_\varphi$ of $B_{T'}$: if $a$ and $b$ lift it and its inverse, then $ab-1\in J_{T'}$ is nilpotent, as $J_{T'}$ is, so
$ab$ is a unit, and so is $a$.
% Lean: ConeLiftEx.nonempty_cLiftK, with IsOpenImmersion.lift for y' and exists_unit_lift for w.
\end{proof}
\leanref{Development/AutomorphicSpectrum/Proofs/BlockDecomposition/ExistsConeLift.lean}{GlobalTheory.CompanionTheory.ConeLiftEx.exists_unit_lift}

\begin{lemma}\label{gal-realization-kappa-de-rham-cocycle}
Let $K$ be a field of characteristic $0$, $X$ a $K$-scheme, $\mathcal L$ a line datum on $X$, $\mathcal C$ a category,
$G:\mathcal C\to\mathrm{Th}_{\mathrm{dR}}(X)$ a functor with $G(c)=(B_c,J_c,\delta_c,y_c)$, and $m\geq0$
(Definition~\ref{gal-realization-period-de-rham}). For a lift $\Lambda=(p,w)$ of $\mathcal L$ along $G$
(Definition~\ref{gal-realization-period-chern}(3)) let $z_m(\Lambda)$ be the $2$-cochain of $\mathcal C$ with values in
$c\mapsto B_c$ that sends a chain $s=(c_0\to c_1\to c_2)$ to $\log_m(r(s)-1)$, formed for $\delta_{c_2}$
(Lemma~\ref{gal-realization-kappa-de-rham-log}).
\begin{enumerate}
\item If $J_c^{[m]}=0$ for every $c$, then $z_m(\Lambda)$ is a cocycle.
\item If $J_c^{[m]}=0$ for every $c$ and $\Lambda'$ is another lift of $\mathcal L$ along $G$, then
  $z_m(\Lambda')-z_m(\Lambda)$ is a coboundary.
\item For a functor $h:\mathcal C'\to\mathcal C$, the morphisms $p_{h(c)}$ and units $w_{h(\varphi)}$ form a lift
  $\Lambda h$ of $\mathcal L$ along $Gh$, and $z_m(\Lambda h)$ is the cochain $h^*z_m(\Lambda)$, which sends
  $c_0\to c_1\to c_2$ to $z_m(\Lambda)(h(c_0)\to h(c_1)\to h(c_2))$.
\end{enumerate}
\end{lemma}
\leanref{Development/WeilBound/CoreStatements/PadicHodgeComparison/EtaleCohomologyComparison/PerfectoidPullbacks/ConeSectionRestriction/Precomp.lean}{CrystallineSemistableComparison.Pb.zK}
\leanref{Development/AutomorphicSpectrum/Proofs/BlockDecomposition/ConeCocycleClosedStatement.lean}{GlobalTheory.T.cone_cocycle_closed}
\leanref{Development/WeilBound/CoreStatements/PadicHodgeComparison/EtaleCohomologyComparison/PerfectoidPullbacks/ConeSectionRestriction/RSubOneMem.lean}{CrystallineSemistableComparison.Pb.zK_cocycle}
\leanref{Development/AutomorphicSpectrum/Proofs/BlockDecomposition/ConeCocycleClassIndependentStatement.lean}{GlobalTheory.T.cone_cocycle_class_independent}
\leanref{Development/WeilBound/CoreStatements/PadicHodgeComparison/EtaleCohomologyComparison/PerfectoidPullbacks/ConeSectionRestriction/RSubOneMem.lean}{CrystallineSemistableComparison.Pb.zK_cls_eq}
\leanref{Development/WeilBound/CoreStatements/PadicHodgeComparison/EtaleCohomologyComparison/PerfectoidPullbacks/ConeSectionRestriction/Precomp.lean}{CrystallineSemistableComparison.Pb.CLiftK.precomp}
\leanref{Development/WeilBound/CoreStatements/PadicHodgeComparison/EtaleCohomologyComparison/PerfectoidPullbacks/ConeSectionRestriction/RSubOneMem.lean}{CrystallineSemistableComparison.Pb.zK_precomp}

\begin{proof}
Each $B_c$ is a $K$-algebra, so every integer $n\geq1$ is invertible in it. So in (1) and (2),
Lemma~\ref{gal-realization-kappa-de-rham-log}(3) applies in every $B_c$: for units $a,b$ of $B_c$ with $a-1,b-1\in J_c$,
$\log_m(ab-1)=\log_m(a-1)+\log_m(b-1)$; for $b=a^{-1}$ this gives $\log_m(a^{-1}-1)=-\log_m(a-1)$, as $\log_m(0)=0$.
Recall that $r(s)-1\in J_{c_2}$ (Definition~\ref{gal-realization-period-chern}(3)), and that the maps $G(\varphi)$ are
divided power maps, so they carry $J_c$ into $J_{c'}$ and commute with $\log_m$
(Lemma~\ref{gal-realization-kappa-de-rham-log}(1)).

(1) Let $\tau=(c_0\to c_1\to c_2\to c_3)$, with maps $\varphi_{ij}:c_i\to c_j$, and write $G_{ij}=G(\varphi_{ij})$,
$w_{ij}=w_{\varphi_{ij}}$ and $r_{ijk}=r(c_i\to c_j\to c_k)=G_{jk}(w_{ij})w_{jk}w_{ik}^{-1}$. Expanding,
\[
 G_{23}(r_{012})\,r_{023}=G_{13}(w_{01})\,G_{23}(w_{12})\,w_{23}\,w_{03}^{-1}=r_{123}\,r_{013}.
\]
The four factors minus $1$ lie in $J_{c_3}$. So
$\log_m(G_{23}(r_{012})-1)+\log_m(r_{023}-1)=\log_m(r_{123}-1)+\log_m(r_{013}-1)$, where
$\log_m(G_{23}(r_{012})-1)=G_{23}(z_m(\Lambda)(c_0\to c_1\to c_2))$. This says that the coboundary of $z_m(\Lambda)$ at
$\tau$, the alternating sum of its values at $c_1\to c_2\to c_3$, $c_0\to c_2\to c_3$, $c_0\to c_1\to c_3$ and, carried
by $G_{23}$, $c_0\to c_1\to c_2$, is $0$.

(2) Let $\Lambda'=(p',w')$. For each $c$, $u(p_c,p'_c)$ is a unit of $B_c/J_c$; as in the proof of
Lemma~\ref{gal-realization-kappa-de-rham-lift} it lifts to a unit $v_c$ of $B_c$. For $\varphi:c\to c'$ put
$\varepsilon_\varphi=w'_\varphi\,G(\varphi)(v_c)\,w_\varphi^{-1}v_{c'}^{-1}\in B_{c'}^\times$. With
$a=p_c\circ\operatorname{Spec}\overline{G(\varphi)}$ and $a'=p'_c\circ\operatorname{Spec}\overline{G(\varphi)}$, the
images in $B_{c'}/J_{c'}$ of $w_\varphi$, $w'_\varphi$, $G(\varphi)(v_c)$ and $v_{c'}$ are $u(a,p_{c'})$,
$u(a',p'_{c'})$, $u(a,a')$, by the naturality of $u$, and $u(p_{c'},p'_{c'})$. By the cocycle identity of $u$ the image
of $\varepsilon_\varphi$ is $u(a,p'_{c'})u(a,p'_{c'})^{-1}=1$, so $\varepsilon_\varphi-1\in J_{c'}$. For a chain
$s=(c_0\xrightarrow{\varphi_1}c_1\xrightarrow{\varphi_2}c_2)$, expanding the definitions gives
\[
 r'(s)=r(s)\,G(\varphi_2)(\varepsilon_{\varphi_1})\,\varepsilon_{\varphi_2}\,\varepsilon_{\varphi_2\varphi_1}^{-1},
\]
where $r'$ is formed with $\Lambda'$. Let $e$ be the $1$-cochain $(c\xrightarrow\varphi c')\mapsto
\log_m(\varepsilon_\varphi-1)$. Taking $\log_m$ of $r'(s)-1$ gives $z_m(\Lambda')(s)=z_m(\Lambda)(s)+e(\varphi_2)-
e(\varphi_2\varphi_1)+G(\varphi_2)(e(\varphi_1))$, that is, $z_m(\Lambda')=z_m(\Lambda)+de$.

(3) This holds by the definitions.
% Lean: (1) is ConeZcoc.main, through rr_coboundary, plogK_four (truncated_log_multiplicative), map_zK and
% nat_isUnit; (2) is ConeZcls.main, through exists_unit_lift (v), eps and eps_sub, rr_eq, log_eq (plog_mul, plog_inv)
% and zK_eq with the 1-cochain eC; (3) is zK_precomp, which holds by rfl.
\end{proof}
\leanref{Development/AutomorphicSpectrum/Proofs/BlockDecomposition/ConeCocycleClosed.lean}{GlobalTheory.CompanionTheory.ConeZcoc.main}
\leanref{Development/AutomorphicSpectrum/Proofs/BlockDecomposition/ConeCocycleClosed.lean}{GlobalTheory.CompanionTheory.ConeZcoc.rr_coboundary}
\leanref{Development/AutomorphicSpectrum/Proofs/BlockDecomposition/ConeCocycleClosed.lean}{GlobalTheory.CompanionTheory.ConeZcoc.plogK_four}
\leanref{Development/AutomorphicSpectrum/Proofs/BlockDecomposition/ConeCocycleClassIndependent.lean}{GlobalTheory.CompanionTheory.ConeZcls.main}
\leanref{Development/AutomorphicSpectrum/Proofs/BlockDecomposition/ConeCocycleClassIndependent.lean}{GlobalTheory.CompanionTheory.ConeZcls.eps_sub}
\leanref{Development/AutomorphicSpectrum/Proofs/BlockDecomposition/ConeCocycleClassIndependent.lean}{GlobalTheory.CompanionTheory.ConeZcls.rr_eq}
\leanref{Development/AutomorphicSpectrum/Proofs/BlockDecomposition/ConeCocycleClassIndependent.lean}{GlobalTheory.CompanionTheory.ConeZcls.zK_eq}

\subsection{Cone coverings and residues}\label{subsection-gal-cones-residues}

\begin{definition}\label{gal-realization-kappa-de-rham-cone}
Let $K$ be a field, $X$ a $K$-scheme, $i:X\to\mathbf P^M_K$ a $K$-morphism, $y_0,\ldots,y_m$ polynomials in
$K[x_0,\ldots,x_M]$, and $s\geq1$. For a polynomial $y$, $D_+(y)\subset\mathbf P^M_K$ is the open subset of the
homogeneous prime ideals that do not contain $y$, the union of the $D_+$ of the homogeneous components of $y$.
\begin{enumerate}
\item Let $U_k=i^{-1}(D_+(y_k))$, $X_j=i^{-1}(D_+(x_j))$ for $0\leq j\leq M$, and $U_S=\bigcap_{k\in S}U_k$ for a
  nonempty $S\subset\{0,\ldots,m\}$. Suppose in this item that the $y_k$ are homogeneous of one degree $e\geq1$. Through
  $i$, the section $y_k/y_{k'}$ of $\mathcal O_{\mathbf P^M_K}$ over $D_+(y_{k'})$ gives an element of $\Gamma(U_{k'},\mathcal
  O_X)$, also written $y_k/y_{k'}$. On an open contained in $U_a$, $U_b$ and $U_c$, $(y_a/y_b)(y_b/y_c)=y_a/y_c$ and
  $y_a/y_a=1$; so $y_k/y_{k'}$ is a unit on $U_k\cap U_{k'}$.
\item Let $\mathcal W_k$ be the full subcategory of $\mathrm{Th}^{(s)}_{\mathrm{dR}}(X)$
  (Definition~\ref{gal-realization-kappa-de-rham-truncated}) of the thickenings that lie over $U_k\cap X_j$ for some
  $j$. For a nonempty $S\subset\{0,\ldots,m\}$ let $\mathcal W_S$ be the full subcategory of the thickenings that lie in
  $\mathcal W_k$ for every $k\in S$, and $\mathcal W$ that of the thickenings that lie in some $\mathcal W_k$. Cohomology
  of these categories is that of the restricted functors (Definition~\ref{gal-realization-split-diagrams}(1)), and
  $\mathrm{gr}_r=\mathcal O^{[r]}/\mathcal O^{[r+1]}$.
\item Let $\mathcal Q$ be a full subcategory of $\mathrm{Th}_{\mathrm{dR}}(X)$ whose thickenings lie over an open
  $U\subset X$, and $g$ a unit of $\Gamma(U,\mathcal O_X)$. A \emph{family of lifts} of $g$ on $\mathcal Q$ is a family
  $\hat g=(\hat g_T)$ of lifts $\hat g_T\in B_T$ of $g$ at the thickenings $T$ of $\mathcal Q$
  (Definition~\ref{gal-realization-thickening-level-one-fibre}(2), for $R=K$). Each $\hat g_T$ is a unit, as $J_T$ is
  nilpotent (proof of Lemma~\ref{gal-realization-kappa-de-rham-lift}), and $\hat g_{T'}\varphi(\hat g_T)^{-1}-1\in
  J_{T'}$ for a map $\varphi:T\to T'$, as $\bar\varphi\circ\mathrm{ev}_T=\mathrm{ev}_{T'}$
  (Lemma~\ref{gal-realization-thickening-level-one-forms}(1)). Let $\mathrm{dlog}_s(\hat g)$ be the $1$-cochain of
  $\mathcal Q$ with values in $\mathcal O$ that sends $\varphi:T\to T'$ to $\log_s(\hat g_{T'}\varphi(\hat g_T)^{-1}-1)$
  (Lemma~\ref{gal-realization-kappa-de-rham-log}).
\end{enumerate}
\end{definition}
\leanref{Development/WeilBound/CoreStatements/PadicHodgeComparison/EtaleCohomologyComparison/PerfectoidPullbacks/ConeSectionRestriction/ConeSite.lean}{CrystallineSemistableComparison.Pb.uy}
\leanref{Development/WeilBound/CoreStatements/PadicHodgeComparison/EtaleCohomologyComparison/PerfectoidPullbacks/ConeSectionRestriction/ConeSite.lean}{CrystallineSemistableComparison.Pb.wch}
\leanref{Development/WeilBound/CoreStatements/PadicHodgeComparison/EtaleCohomologyComparison/PerfectoidPullbacks/ConeSectionRestriction/ConeSite.lean}{CrystallineSemistableComparison.Pb.yQ}
\leanref{Development/WeilBound/CoreStatements/PadicHodgeComparison/EtaleCohomologyComparison/PerfectoidPullbacks/ConeSectionRestriction/JacSpanLe.lean}{CrystallineSemistableComparison.Pb.yQ_mul}
\leanref{Development/WeilBound/CoreStatements/PadicHodgeComparison/EtaleCohomologyComparison/PerfectoidPullbacks/ConeSectionRestriction/JacSpanLe.lean}{CrystallineSemistableComparison.Pb.yQ_unit}
\leanref{Development/WeilBound/CoreStatements/PadicHodgeComparison/EtaleCohomologyComparison/PerfectoidPullbacks/ConeSectionRestriction/ConeSite.lean}{CrystallineSemistableComparison.Pb.wK}
\leanref{Development/WeilBound/CoreStatements/PadicHodgeComparison/EtaleCohomologyComparison/PerfectoidPullbacks/ConeSectionRestriction/ConeSite.lean}{CrystallineSemistableComparison.Pb.wU}
\leanref{Development/WeilBound/CoreStatements/PadicHodgeComparison/EtaleCohomologyComparison/PerfectoidPullbacks/CapOComp.lean}{CrystallineSemistableComparison.Pb.capS}
\leanref{Development/WeilBound/CoreStatements/PadicHodgeComparison/EtaleCohomologyComparison/PerfectoidPullbacks/CapOComp.lean}{CrystallineSemistableComparison.Pb.HU}
\leanref{Development/WeilBound/CoreStatements/PadicHodgeComparison/EtaleCohomologyComparison/PerfectoidPullbacks/ConeSectionRestriction/LocDefs.lean}{CrystallineSemistableComparison.Pb.dlogC}

\begin{definition}\label{gal-realization-kappa-de-rham-torus}
Let $K$ be a field, $R$ a $K$-algebra, $n\geq0$, $t_0,\ldots,t_n$ units of $R$ and $D_0,\ldots,D_n$ $K$-derivations of
$R$.
\begin{enumerate}
\item For $S\subset\{0,\ldots,n+1\}$, $\mathcal A(S)$ is the $K$-subalgebra of $R$ generated by the products
  $\prod_{k=0}^nt_k^{a_k}$ with $a\in\mathbf Z^{n+1}$ such that $a_k\geq0$ whenever $k+1\notin S$, and $\sum_ka_k\leq0$
  if $0\notin S$.
\item For a $K$-subalgebra $B$ of $R$, $\mathrm{Jac}(B)$ is the additive subgroup of $R$ generated by the elements
  $b\det(D_k(b_l))_{0\leq k,l\leq n}$ with $b,b_0,\ldots,b_n\in B$. For $0\leq k\leq n$, $\mathrm{Jac}_k(B)$ is the
  additive subgroup generated by the elements $b\det(D_{k'}(b_l))$, where $k'$ runs through $\{0,\ldots,n\}\smallsetminus
  \{k\}$ in increasing order and $1\leq l\leq n$, with $b,b_1,\ldots,b_n\in B$.
\end{enumerate}
\end{definition}
\leanref{Development/WeilBound/CoreStatements/PadicHodgeComparison/EtaleCohomologyComparison/PerfectoidPullbacks/ConeSectionRestriction/SeamDefs.lean}{CrystallineSemistableComparison.Pb.tpow}
\leanref{Development/WeilBound/CoreStatements/PadicHodgeComparison/EtaleCohomologyComparison/PerfectoidPullbacks/ConeSectionRestriction/SeamDefs.lean}{CrystallineSemistableComparison.Pb.expS}
\leanref{Development/WeilBound/CoreStatements/PadicHodgeComparison/EtaleCohomologyComparison/PerfectoidPullbacks/ConeSectionRestriction/SeamDefs.lean}{CrystallineSemistableComparison.Pb.torusS}
\leanref{Development/WeilBound/CoreStatements/PadicHodgeComparison/EtaleCohomologyComparison/PerfectoidPullbacks/ConeSectionRestriction/SeamDefs.lean}{CrystallineSemistableComparison.Pb.jacS}
\leanref{Development/WeilBound/CoreStatements/PadicHodgeComparison/EtaleCohomologyComparison/PerfectoidPullbacks/ConeSectionRestriction/LocDefs.lean}{CrystallineSemistableComparison.Pb.jacSpan}

\begin{lemma}\label{gal-realization-kappa-de-rham-cone-vanishing}
Let $K$ be a field.
\begin{enumerate}
\item Let $X$, $i$, $y_0,\ldots,y_m$ and $s$ be as in Definition~\ref{gal-realization-kappa-de-rham-cone}, $S$ a
  nonempty subset of $\{0,\ldots,m\}$ and $r,r'\in\mathbf Z$. Every thickening of $\mathcal W_S$ lies over $U_S$, and
  restriction $H^j(\mathrm{Th}^{(s)}_{\mathrm{dR}}(X)_{U_S},\mathcal O^{[r]}/\mathcal O^{[r']})\to H^j(\mathcal
  W_S,\mathcal O^{[r]}/\mathcal O^{[r']})$ is bijective for every $j$.
\item Let $K$ have characteristic $0$, $X$ be a $K$-scheme smooth of relative dimension $d$, $V\subset X$ an affine
  open, $s\geq1$ and $r>d$. Then $H^j(\mathrm{Th}^{(s)}_{\mathrm{dR}}(X)_V,\mathrm{gr}_r)=0$ for every $j$.
\item Let $X$ be a $K$-scheme, $\mathcal Q$ a full subcategory of $\mathrm{Th}_{\mathrm{dR}}(X)$ and $r<0$. Then
  $H^j(\mathcal Q,\mathrm{gr}_r)=0$ for every $j$.
\end{enumerate}
\end{lemma}
\leanref{Development/WeilBound/CoreStatements/PadicHodgeComparison/EtaleCohomologyComparison/PerfectoidPullbacks/ConeSectionRestriction/WchCover.lean}{CrystallineSemistableComparison.Pb.zar_capS}
\leanref{Development/AutomorphicSpectrum/Proofs/BlockDecomposition/ChartRestrictPullBijectiveStatement.lean}{GlobalTheory.T.chart_restrict_pull_bijective}
\leanref{Development/WeilBound/CoreStatements/PadicHodgeComparison/EtaleCohomologyComparison/PerfectoidPullbacks/ConeSectionRestriction/OneKNeZero.lean}{CrystallineSemistableComparison.Pb.gr_top_van}
\leanref{Development/AutomorphicSpectrum/Proofs/BlockDecomposition/AffineGradedQuotientVanishAboveDimStatement.lean}{GlobalTheory.T.affine_graded_quotient_vanish_above_dim}
\leanref{Development/WeilBound/CoreStatements/PadicHodgeComparison/EtaleCohomologyComparison/PerfectoidPullbacks/ConeSectionRestriction/GrNegVan.lean}{CrystallineSemistableComparison.Pb.gr_neg_van}
\leanref{Development/AutomorphicSpectrum/Proofs/BlockDecomposition/ThickeningGradedVanishNegativeStatement.lean}{GlobalTheory.T.thickening_graded_vanish_negative}

\begin{proof}
(1) Let $T$ be in $\mathcal W_S$. For each $k\in S$ it lies over $U_k\cap X_{c(k)}$ for some $c(k)$, so it lies over
$\bigcap_{k\in S}(U_k\cap X_{c(k)})\subset U_S$ (Lemma~\ref{gal-realization-thickening-over}(2)). Conversely, a thickening
of $\mathrm{Th}^{(s)}_{\mathrm{dR}}(X)$ that lies over such an intersection lies in $\mathcal W_S$. So $\mathcal
W_S=\mathrm{Th}^{(s)}_{\mathrm{dR}}(X)_{\mathcal V}$ for the set $\mathcal V$ of the opens $\bigcap_{k\in S}(U_k\cap
X_{c(k)})$, $c:S\to\{0,\ldots,M\}$. As the $D_+(x_j)$ cover $\mathbf P^M_K$, the $X_j$ cover $X$, and the union of
$\mathcal V$ is $U_S$. So Lemma~\ref{gal-realization-kappa-de-rham-zariski} and
Lemma~\ref{gal-realization-thickening-zariski}(1), for $\mathcal Q_0=\mathrm{Th}^{(s)}_{\mathrm{dR}}(X)$, the set $\mathcal
V$ and the set $\{U_S\}$ in place of the $\mathcal W$ there, give the bijectivity.

(2) If $s\leq r$, then $\mathrm{gr}_r=0$ on $\mathrm{Th}^{(s)}_{\mathrm{dR}}(X)$, as $J^{[r]}\subset J^{[s]}=0$. Let
$s>r$, and for $f\in\Gamma(V,\mathcal O_X)$ let $D(f)\subset V$ be the basic open where $f$ is a unit.

\emph{A cover by basic opens.} As $X$ is smooth of relative dimension $d$ over $K$, there are $f_\alpha\in\Gamma(V,\mathcal
O_X)$, for $\alpha$ in a finite set $A$, such that the $V_\alpha=D(f_\alpha)$ cover $V$ and each has coordinates of
dimension $d$ (Lemma~\ref{gal-realization-kappa-de-rham-affine-basic}(1)). Let $\mathcal U=\{V_\alpha:\alpha\in A\}$ and
$V_S=\bigcap_{\alpha\in S}V_\alpha$.

\emph{The vanishing.} For nonempty $S\subset A$, $V_S=D(\prod_{\alpha\in S}f_\alpha)$ is affine
(Lemma~\ref{gal-realization-kappa-de-rham-affine-basic}(2)), and it is an affine open of $V_\alpha$ for $\alpha\in S$. So it has coordinates of dimension $d$,
restricted from $V_\alpha$ (Lemma~\ref{gal-realization-thickening-cover}(1)), and as $d<r<s$,
$H^j(\mathrm{Th}^{(s)}_{\mathrm{dR}}(X)_{V_S},\mathrm{gr}_r)=0$ for every $j$ by
Lemma~\ref{gal-realization-kappa-de-rham-local}(3). By Lemma~\ref{gal-realization-thickening-over}, the thickenings over
the $V_\alpha$ form cosieves of $\mathrm{Th}^{(s)}_{\mathrm{dR}}(X)_{\mathcal U}$ that cover it, with intersections the
$\mathrm{Th}^{(s)}_{\mathrm{dR}}(X)_{V_S}$, so $H^j(\mathrm{Th}^{(s)}_{\mathrm{dR}}(X)_{\mathcal U},\mathrm{gr}_r)=0$ for
every $j$ (Lemma~\ref{gal-realization-thickening-cosieves}(3)). By Lemma~\ref{gal-realization-kappa-de-rham-zariski} and
Lemma~\ref{gal-realization-thickening-zariski}(1), for $\mathcal V=\mathcal U$ and $\mathcal W=\{V\}$, restriction from
$\mathrm{Th}^{(s)}_{\mathrm{dR}}(X)_V$ to $\mathrm{Th}^{(s)}_{\mathrm{dR}}(X)_{\mathcal U}$ is bijective on cohomology,
which gives the claim.

(3) As $r<r+1\leq0$, $\mathcal O^{[r]}=\mathcal O^{[r+1]}=\mathcal O$, so $\mathrm{gr}_r=0$.
% Lean: (1) is ConeZarCap.main: the category is rewritten as KS q s of the opens ⨅ k ∈ S, uy ⊓ wch (c k) (e), whose
% union contains U_S by iSup_wch, and NlZar.zariski_descent applies. (2) is GrTop.grTop, with
% twoVariableSeriesVanishesBelowDegree_one_shiftVariable for s ≤ r; the cover of V is exists_cover, by finitely many basic
% opens with coordinates (Lemma gal-realization-kappa-de-rham-affine-basic(1), through CovEx.exists_basicOpen,
% exists_etale_mvPolynomial and CovEx.coord_basicOpen); V_S = D(∏ f_α) is iInf_basicOpen with IsAffineOpen.basicOpen, its coordinates are
% coordinates_restrict_affine_open, and then loc_vanish, stable_cover_cech_vanishing and
% cech_refinement_cohomology_bijective. (3) is gr_neg_van.
\end{proof}
\leanref{Development/AutomorphicSpectrum/Proofs/BlockDecomposition/ChartRestrictPullBijective.lean}{GlobalTheory.CompanionTheory.ConeZarCap.main}
\leanref{Development/AutomorphicSpectrum/Proofs/BlockDecomposition/AffineGradedQuotientVanishAboveDim.lean}{GlobalTheory.CompanionTheory.GrTop.grTop}
\leanref{Development/AutomorphicSpectrum/Proofs/BlockDecomposition/AffineGradedQuotientVanishAboveDim.lean}{GlobalTheory.CompanionTheory.GrTop.exists_cover}
\leanref{Development/AutomorphicSpectrum/Proofs/BlockDecomposition/AffineGradedQuotientVanishAboveDim.lean}{GlobalTheory.CompanionTheory.GrTop.iInf_basicOpen}

\begin{lemma}\label{gal-realization-kappa-de-rham-setup-normalization}
Let $K$ be an infinite field, $M\geq0$, and $I\neq K[x_0,\ldots,x_M]$ an ideal that contains the homogeneous components
of its elements. There are $e\geq1$, $m\geq0$ and homogeneous polynomials $y_k$, $0\leq k<m$, of degree $e$ whose images
$\bar y_k$ in $K[x_0,\ldots,x_M]/I$ are algebraically independent over $K$, such that $K[x_0,\ldots,x_M]/I$ is a finitely
generated module over the $K$-subalgebra generated by the $\bar y_k$.
\end{lemma}
\leanref{Development/AutomorphicSpectrum/Proofs/BlockDecomposition/NoetherNormalizationHomogeneousStatement.lean}{GlobalTheory.T.noether_normalization_homogeneous}

\begin{proof}
Let $A=K[x_0,\ldots,x_M]/I$, write $\bar f$ for the image of $f$ in $A$, and call a polynomial homogeneous of degree $1$ a
\emph{linear form}; $0$ is one. For a finite family $l=(l_i)_{i\in\iota}$ of polynomials let $K[\bar l]\subset A$ be the
$K$-subalgebra generated by the $\bar l_i$. The proof gives $e=1$. As $I\neq K[x_0,\ldots,x_M]$, $A\neq0$, so $K\to A$
is injective.

\emph{Step 1: a homogeneous relation.} Let $l=(l_i)_{i\in\iota}$ be linear forms whose images are algebraically
dependent over $K$. Then there is a nonzero homogeneous $F\in K[z_i:i\in\iota]$ with $F(l)\in I$. Indeed, let $F'\neq0$
with $F'(\bar l)=0$, that is, $F'(l)\in I$, and let $F$ be a nonzero homogeneous component of $F'$, of degree $n$. As the
$l_i$ are linear forms, each homogeneous component $F'_k$ of $F'$ gives a homogeneous polynomial $F'_k(l)$ of degree
$k$; so $F(l)$ is the homogeneous component of degree $n$ of $F'(l)$, and it lies in $I$.

\emph{Step 2: integrality.} Let $R\subset S$ be $K$-algebras, $G\in K[z_i:i\in\iota]$ homogeneous of degree $N$, $a\in
K^\iota$ with $G(a)\neq0$, $v_i\in R$ and $u\in S$ with $G((v_i+a_iu)_i)=0$. Then $u$ is integral over $R$. Indeed, let
$p(T)=G((v_i+a_iT)_i)\in R[T]$. For a monomial $\prod_iz_i^{d_i}$ of $G$, $\sum_id_i=N$, and $\prod_i(v_i+a_iT)^{d_i}$ has
degree at most $N$ with coefficient $\prod_ia_i^{d_i}$ at $T^N$. So $p$ has degree at most $N$, with coefficient $G(a)$ at
$T^N$, and $G(a)^{-1}p$ is a monic polynomial over $R$ with root $u$.

\emph{Step 3: the induction.} For every finite family $l=(l_i)_{i\in\iota}$ of linear forms such that $A$ is integral over
$K[\bar l]$ there are linear forms $y_0,\ldots,y_{m-1}$ whose images are algebraically independent over $K$, such that
$A$ is integral over $K[\bar y]$. This is shown by induction on the number of elements of $\iota$. If the $\bar l_i$ are
algebraically independent, which they are if $\iota$ is empty, take the $l_i$ in some order. Otherwise $\iota$ is
nonempty; choose $j\in\iota$, and $F$ as in Step 1, homogeneous of degree $d$. Then $G=z_jF$ is nonzero, homogeneous of
degree $d+1$, and $G(l)\in I$. As $K$ is infinite and $G\neq0$, there is $a\in K^\iota$ with $G(a)\neq0$; then $a_j\neq0$.
Let $u=a_j^{-1}l_j$ and $v_i=l_i-a_iu$ for $i\in\iota$, linear forms with $v_j=0$, and let $R=K[\bar v_i:i\neq j]$, which
contains every $\bar v_i$. As $\bar l_i=\bar v_i+a_i\bar u$, $G((\bar v_i+a_i\bar u)_i)=\overline{G(l)}=0$, so $\bar u$ is
integral over $R$ by Step 2, for $S=A$. So every $\bar l_i$ is integral over $R$, hence so is $K[\bar l]$
\cite[Corollary~5.3]{AM69}, and $A$ is integral over $R$ by transitivity \cite[Corollary~5.4]{AM69}. The induction
hypothesis applies to the family $(v_i)_{i\neq j}$.

\emph{Step 4: conclusion.} The $\bar x_i$ generate $A$, so $A=K[\bar x]$ is integral over it. Step 3 for $l=(x_0,\ldots,
x_M)$ gives linear forms $y_0,\ldots,y_{m-1}$, homogeneous of degree $e=1$, with algebraically independent images, such
that $A$ is integral over $K[\bar y]$. As $A$ is generated over $K[\bar y]$ by the finitely many $\bar x_i$, which are
integral over it, $A$ is a finitely generated $K[\bar y]$-module \cite[Corollary~5.2]{AM69}.
% Lean: GlobalTheory.T.noether_normalization_homogeneous = ConeNoether.main (NoetherNormalizationHomogeneous.lean).
% Step 1 is exists_homogeneous_relation, with homogeneousComponent_aeval; Step 2 is isIntegral_of_homogeneous_relation,
% with prod_pow_lead (UpperTheory.Reformulated.B.Yra8.main) for the coefficient at T^N; Step 3 is exists_linear_of_integral
% (the case of independent forms is done, the point a from MvPolynomial.funext, u and v, hdec, hu and hA, and
% isIntegral_trans_subalgebra); Step 4 is htop and Algebra.IsIntegral.finite with Algebra.FiniteType. About 250 lines.
\end{proof}
\leanref{Development/AutomorphicSpectrum/Proofs/BlockDecomposition/NoetherNormalizationHomogeneous.lean}{GlobalTheory.CompanionTheory.ConeNoether.main}
\leanref{Development/AutomorphicSpectrum/Proofs/BlockDecomposition/NoetherNormalizationHomogeneous.lean}{GlobalTheory.CompanionTheory.ConeNoether.exists_homogeneous_relation}
\leanref{Development/AutomorphicSpectrum/Proofs/BlockDecomposition/NoetherNormalizationHomogeneous.lean}{GlobalTheory.CompanionTheory.ConeNoether.isIntegral_of_homogeneous_relation}
\leanref{Development/AutomorphicSpectrum/Proofs/BlockDecomposition/NoetherNormalizationHomogeneous.lean}{GlobalTheory.CompanionTheory.ConeNoether.exists_linear_of_integral}

\begin{lemma}\label{gal-realization-kappa-de-rham-setup-cone-relation}
Let $K$ be a field, $I$ an ideal of $K[x_0,\ldots,x_M]$ that contains the homogeneous components of its elements, $e\geq1$,
and $y_0,\ldots,y_{m-1}$ homogeneous polynomials of degree $e$ such that $K[x_0,\ldots,x_M]/I$ is a finitely generated
module over the subalgebra generated by the images of the $y_k$. Let $w\geq0$ and $F$ homogeneous of degree $ew$. There are
$s\geq1$ and $c_l\in K[z_0,\ldots,z_{m-1}]$, homogeneous of degree $(s-l)w$, for $0\leq l<s$, such that
$F^s+\sum_{l<s}c_l(y)F^l\in I$.
\end{lemma}
\leanref{Development/AutomorphicSpectrum/Proofs/BlockDecomposition/HomogeneousIntegralExpression.lean}{GlobalTheory.CompanionTheory.Geo.exists_hom_relation}

\begin{proof}
Write $\bar f$ for the image of $f$ in $A=K[x_0,\ldots,x_M]/I$, and $B$ for the subalgebra generated by the $\bar y_k$. As
$A$ is a finitely generated $B$-module, $\bar F$ is integral over $B$ \cite[Proposition~5.1]{AM69}: it is a root of a
monic polynomial over $B$, which, multiplied by $X$, has some degree $s\geq1$. Its coefficients are $C_l(\bar y)$ for
polynomials $C_l\in K[z_0,\ldots,z_{m-1}]$, so $F^s+\sum_{l<s}C_l(y)F^l\in I$. Its homogeneous component of degree $ews$
lies in $I$. It is $F^s+\sum_lc_l(y)F^l$, where $c_l$ is the homogeneous component of degree $(s-l)w$ of $C_l$: the
homogeneous component $C_{l,r}$ of degree $r$ of $C_l$ gives $C_{l,r}(y)F^l$, homogeneous of degree $er+ewl$, and as $e\geq1$
this is $ews$ only for $r=(s-l)w$.
% Lean: exists_hom_relation, through exists_monic_relation (Algebra.IsIntegral.of_finite, the factor X, and the lift of
% the coefficients through Algebra.adjoin_range_eq_range_aeval) and homogeneousComponent_aeval_mul.
\end{proof}
\leanref{Development/AutomorphicSpectrum/Proofs/BlockDecomposition/HomogeneousIntegralExpression.lean}{GlobalTheory.CompanionTheory.Geo.exists_monic_relation}
\leanref{Development/AutomorphicSpectrum/Proofs/BlockDecomposition/HomogeneousIntegralExpression.lean}{GlobalTheory.CompanionTheory.Geo.homogeneousComponent_aeval_mul}

\begin{lemma}\label{gal-realization-kappa-de-rham-setup-cone-vanish}
Let $K$ be a field, $X$ a $K$-scheme, $i:X\to\mathbf P^M_K$ a $K$-morphism, $X_j=i^{-1}(D_+(x_j))$, and $I_X$ the ideal of
$K[x_0,\ldots,x_M]$ generated by the homogeneous $F$ such that, for every $j$, the section $F/x_j^{\deg F}$ of $\mathcal
O_{\mathbf P^M_K}$ over $D_+(x_j)$, carried through $i$, vanishes on $X_j$. For $g$ homogeneous of degree $e\geq1$ and
$F$ homogeneous of degree $se$, write $F/g^s$ also for the section of $\mathcal O_X$ over $i^{-1}(D_+(g))$ carried through
$i$, and for its restrictions.
\begin{enumerate}
\item If $F\in I_X$, then $F/g^s=0$ on $i^{-1}(D_+(g))$.
\item If $X$ is integral and $F/g^s$ vanishes on a nonempty open $V\subset i^{-1}(D_+(g))$, then $F\in I_X$.
\item For $p\in X$, let $\mathfrak q_p$ be the homogeneous prime of $K[x_0,\ldots,x_M]$ of the point $i(p)$. Then
  $I_X\subset\mathfrak q_p$.
\end{enumerate}
\end{lemma}
\leanref{Development/AutomorphicSpectrum/Proofs/BlockDecomposition/EvalZeroOfMemHomogeneousIdeal.lean}{GlobalTheory.CompanionTheory.GeoVanish.homIdeal_vanish}
\leanref{Development/AutomorphicSpectrum/Proofs/BlockDecomposition/MemHomogeneousIdealOfEvalZero.lean}{GlobalTheory.CompanionTheory.GeoVan.mem_homIdeal_of_vanish}
\leanref{Development/AutomorphicSpectrum/Proofs/BlockDecomposition/BasicOpensCoverOfFinite.lean}{GlobalTheory.CompanionTheory.GeoCover.mem_of_evV_eq_zero}

\begin{proof}
On $D_+(GH)=D_+(G)\cap D_+(H)$ the sections satisfy $F/G^aH^b=(F/G^{a'}H^{b'})\cdot(G^{a'}H^{b'}/G^aH^b)$ whenever the
degrees match, as in the proof of Lemma~\ref{gal-realization-kappa-de-rham-local-exact-charts}; this is used throughout.

(1) \emph{The charts.} Fix $j$. Let $\mathcal J$ be the set of polynomials $p$ such that, for every $n$, the homogeneous
component $p_n$ of degree $n$ of $p$ has $p_n/x_j^n=0$ on $X_j$. It is closed under sums. For $a$ homogeneous of degree
$k$, $(ap)_n=ap_{n-k}$ (and $0$ for $n<k$), and $(ap_{n-k})/x_j^n=(a/x_j^k)(p_{n-k}/x_j^{n-k})$; so $\mathcal J$ is closed
under multiplication by homogeneous elements, hence by all polynomials, and it is an ideal. It contains the generators of
$I_X$, which are homogeneous, so $I_X\subset\mathcal J$. So $F/x_j^{se}=0$ on $X_j$.

\emph{The open $i^{-1}(D_+(g))$.} On $W_j=i^{-1}(D_+(g))\cap X_j=i^{-1}(D_+(gx_j))$, $F/g^s=(x_j^{se}/g^s)(F/x_j^{se})=0$. As the
$D_+(x_j)$ cover $\mathbf P^M_K$, the $W_j$ cover $i^{-1}(D_+(g))$, so $F/g^s=0$ there.

(2) Let $j$ be given. If $X_j$ is empty, $F/x_j^{se}=0$ on $X_j$. Otherwise $W=V\cap X_j$ is nonempty, as $X$ is
irreducible, and on $W$, $F/x_j^{se}=(g^s/x_j^{se})(F/g^s)=0$. As $X$ is integral, restriction from $X_j$ to its nonempty
open $W$ is injective, so $F/x_j^{se}=0$ on $X_j$. So $F$ is one of the generators of $I_X$.

(3) Let $F$ be one of the generators of $I_X$, of degree $n$, and $j$ with $i(p)\in D_+(x_j)$. If $F\notin\mathfrak q_p$,
then $i(p)\in D_+(x_jF)$, where $F/x_j^n=x_jF/x_j^{n+1}$ is a unit. So $F/x_j^n$, carried through $i$, is a unit at $p\in
X_j$, and it is not $0$ on $X_j$, which is false. So $F\in\mathfrak q_p$.
% Lean: (1) is GeoVanish.homIdeal_vanish, with chart_zero (the ideal J, by span induction and
% DirectSum.Decomposition.inductionOn) for the charts, ratio_eq, evV_awayMap and evV_res, and
% TopCat.Sheaf.eq_of_locally_eq' for the cover; (2) is GeoVan.mem_homIdeal_of_vanish, with away_factor,
% nonempty_preirreducible_inter and map_injective_of_isIntegral; (3) is GeoCover.mem_of_evV_eq_zero, through
% basicOpen_awayToSection_isLocalizationElem, and the step hle of GeoCover.iSup_preimage_eq_top.
\end{proof}
\leanref{Development/AutomorphicSpectrum/Proofs/BlockDecomposition/EvalZeroOfMemHomogeneousIdeal.lean}{GlobalTheory.CompanionTheory.GeoVanish.chart_zero}

\begin{lemma}\label{gal-realization-kappa-de-rham-setup-cone-monomial}
Let $K$ be a field, $X$ a $K$-scheme, $i:X\to\mathbf P^M_K$ a $K$-morphism, $d\geq0$, $e\geq1$, $y_0,\ldots,y_d$ homogeneous
polynomials of degree $e$, $U_k=i^{-1}(D_+(y_k))$, $\Lambda=\{0,\ldots,d\}$, $V\subset\bigcap_{k\in\Lambda}U_k$ an open,
and $t_k\in\Gamma(V,\mathcal O_X)$ the restriction of $y_{k+1}/y_0$ for $0\leq k\leq d-1$. Let $S\subset\Lambda$ be
nonempty, $y_S=\prod_{k\in S}y_k$, and $E(S)$ the set of $m\in\mathbf Z^d$ with $m_k\geq0$ whenever $k+1\notin S$ and
$\sum_km_k\leq0$ if $0\notin S$.
\begin{enumerate}
\item The $t_k$ are units. For $j\geq0$ and $\alpha\in\mathbf N^\Lambda$ with $\sum_k\alpha_k=j|S|$, the restriction of
  $\prod_ky_k^{\alpha_k}/y_S^j$ to $V$ is $\prod_kt_k^{m_k}$, where $m_k=\alpha_{k+1}-j$ if $k+1\in S$ and $m_k=\alpha_{k+1}$
  otherwise; and $m\in E(S)$.
\item Every $\prod_kt_k^{m_k}$ with $m\in E(S)$ is of this form.
\item For $C\in K[z_0,\ldots,z_d]$ homogeneous of degree $j|S|$, the restriction of $C(y)/y_S^j$ to $V$ lies in the
  $K$-subalgebra of $\Gamma(V,\mathcal O_X)$ generated by the $\prod_kt_k^{m_k}$, $m\in E(S)$.
\end{enumerate}
\end{lemma}
\leanref{Development/AutomorphicSpectrum/Proofs/BlockDecomposition/TorusMonomialEvalCompat.lean}{GlobalTheory.CompanionTheory.GeoTorus.evV_monomial}
\leanref{Development/AutomorphicSpectrum/Proofs/BlockDecomposition/TorusMonomialMemRange.lean}{GlobalTheory.CompanionTheory.GeoTpow.tpow_mem_range}
\leanref{Development/AutomorphicSpectrum/Proofs/BlockDecomposition/TorusMonomialEvalCompat.lean}{GlobalTheory.CompanionTheory.GeoTorus.aeval_mem_torusS}

\begin{proof}
(1) On $D_+(y_0y_{k+1})$, $y_{k+1}/y_0$ has the inverse $y_0/y_{k+1}$, and $V$ lies in its inverse image. Let
$\beta_k=j$ for $k\in S$ and $\beta_k=0$ otherwise, so that $y_S^j=\prod_ky_k^{\beta_k}$ and $\sum_k(\alpha_k-\beta_k)=0$.
On $V$, where all the $y_k/y_0$ are units,
\[
 \frac{\prod_ky_k^{\alpha_k}}{y_S^j}=\prod_{k\in\Lambda}\Bigl(\frac{y_k}{y_0}\Bigr)^{\alpha_k-\beta_k}=\prod_{k=0}^{d-1}t_k^{m_k},
\]
as $y_0/y_0=1$. If $k+1\notin S$, $m_k=\alpha_{k+1}\geq0$. If $0\notin S$, $\sum_km_k=-(\alpha_0-\beta_0)=-\alpha_0\leq0$.

(2) Let $m\in E(S)$, and $n\in\mathbf Z^\Lambda$ with $n_0=-\sum_km_k$ and $n_{k+1}=m_k$; so $\sum_kn_k=0$, and $n_k\geq0$
for $k\notin S$. Let $j=\sum_k|n_k|$ and $\alpha_k=n_k+\beta_k$, with $\beta$ as in (1). Then $\alpha_k\geq0$, $\sum_k\alpha_k
=j|S|$, and (1) gives $\prod_kt_k^{m_k}$.

(3) $C(y)=\sum_\alpha c_\alpha\prod_ky_k^{\alpha_k}$, the sum over $\alpha$ with $\sum_k\alpha_k=j|S|$, so by (1) the
restriction of $C(y)/y_S^j$ is $\sum_\alpha c_\alpha\prod_kt_k^{m_k(\alpha)}$ with $m(\alpha)\in E(S)$; the constants
$c_\alpha$ are carried to $\Gamma(V,\mathcal O_X)$ by the structure map, as $i$ is a $K$-morphism.
% Lean: (1) is GeoUnit.tV_isUnit (awayMap_ratio_mul) for the units and GeoTorus.evV_monomial (with mk_prod_pow,
% evV_mul_cross and bS) and mem_expS for the monomials; (2) is GeoTpow.tpow_mem_range, with exists_shift and core; (3) is
% GeoTorus.aeval_mem_torusS, with evV_const for the constants.
\end{proof}
\leanref{Development/AutomorphicSpectrum/Proofs/BlockDecomposition/TorusMonomialEvalCompat.lean}{GlobalTheory.CompanionTheory.GeoTorus.mem_expS}
\leanref{Development/AutomorphicSpectrum/Proofs/BlockDecomposition/TorusMonomialMemRange.lean}{GlobalTheory.CompanionTheory.GeoTpow.exists_shift}

\begin{lemma}\label{gal-realization-kappa-de-rham-setup-cone}
Let $K$ be a field, $X$ a $K$-scheme and $i:X\to\mathbf P^M_K$ a closed immersion over $K$, with $X_j=i^{-1}(D_+(x_j))$.
Let $I_X$ be the ideal of $K[x_0,\ldots,x_M]$ generated by the homogeneous $F$ such that, for every $j$, the section
$F/x_j^{\deg F}$ of $\mathcal O_{\mathbf P^M_K}$ over $D_+(x_j)$, carried through $i$, vanishes on $X_j$. For a homogeneous
$G$ of positive degree let $K[x_0,\ldots,x_M]_{(G)}$ be as in Lemma~\ref{gal-realization-kappa-de-rham-local-exact-charts}.
\begin{enumerate}
\item If $X$ is nonempty, then $I_X\neq K[x_0,\ldots,x_M]$, and $I_X$ contains the homogeneous components of its elements.
\end{enumerate}
Let $e\geq1$ and $y_k$, $k\in\Lambda$, a finite family of homogeneous polynomials of degree $e$, whose images in
$K[x_0,\ldots,x_M]/I_X$ are algebraically independent over $K$, such that $K[x_0,\ldots,x_M]/I_X$ is a finitely generated
module over the subalgebra they generate. Let $U_k=i^{-1}(D_+(y_k))$, and for nonempty $S\subset\Lambda$ let
$y_S=\prod_{k\in S}y_k$ and $U_S=\bigcap_{k\in S}U_k$.
\begin{enumerate}
\setcounter{enumi}{1}
\item The $U_k$ cover $X$.
\item $U_S=i^{-1}(D_+(y_S))$ is affine, and for every open $V\subset U_S$ the images in $\Gamma(V,\mathcal O_X)$ of
  $\Gamma(U_S,\mathcal O_X)$ and of $K[x_0,\ldots,x_M]_{(y_S)}$, carried through $i$, are equal.
\item If $X$ is integral and $G\notin I_X$ is homogeneous of positive degree, then $i^{-1}(D_+(G))$ is nonempty.
\end{enumerate}
Let now $\Lambda=\{0,\ldots,d\}$ with $d\geq0$, $V\subset U_\Lambda$ an open and $t_k\in\Gamma(V,\mathcal O_X)$ the
restriction of $y_{k+1}/y_0$ for $0\leq k\leq d-1$.
\begin{enumerate}
\setcounter{enumi}{4}
\item The $t_k$ are units. If $X$ is integral and $V$ is nonempty, they are algebraically independent over $K$.
\item Let $V\subset U_S$, and let $\mathcal A(S)$ be the $K$-subalgebra of $\Gamma(V,\mathcal O_X)$ generated by the
  products $\prod_kt_k^{a_k}$ with $a\in\mathbf Z^d$ such that $a_k\geq0$ whenever $k+1\notin S$, and $\sum_ka_k\leq0$ if
  $0\notin S$ (for $d=n+1$, Definition~\ref{gal-realization-kappa-de-rham-torus}(1)). Then $\mathcal A(S)$ lies in the image
  of $K[x_0,\ldots,x_M]_{(y_S)}$ in $\Gamma(V,\mathcal O_X)$, and every element of this image is integral over $\mathcal
  A(S)$.
\end{enumerate}
\end{lemma}
\leanref{Development/WeilBound/CoreStatements/PadicHodgeComparison/EtaleCohomologyComparison/PerfectoidPullbacks/ConeSectionRestriction/GeoDefs.lean}{CrystallineSemistableComparison.Pb.homIdeal}
\leanref{Development/AutomorphicSpectrum/Proofs/BlockDecomposition/HomogeneousIdealNeTop.lean}{GlobalTheory.CompanionTheory.GeoNetop.homIdeal_ne_top}
\leanref{Development/AutomorphicSpectrum/Proofs/BlockDecomposition/HomogeneousIdealComponentMem.lean}{GlobalTheory.CompanionTheory.Geo.homIdeal_homog}
\leanref{Development/AutomorphicSpectrum/Proofs/BlockDecomposition/BasicOpensCoverOfFiniteStatement.lean}{GlobalTheory.T.basic_opens_cover_of_finite}
\leanref{Development/WeilBound/CoreStatements/PadicHodgeComparison/EtaleCohomologyComparison/PerfectoidPullbacks/ConeSectionRestriction/IsAffineOpenPreimage.lean}{CrystallineSemistableComparison.Pb.range_resO_eq}
\leanref{Development/AutomorphicSpectrum/Proofs/BlockDecomposition/BasicOpenNonemptyOfNeZeroStatement.lean}{GlobalTheory.T.basic_open_nonempty_of_ne_zero}
\leanref{Development/AutomorphicSpectrum/Proofs/BlockDecomposition/TorusCoordinateIsUnitStatement.lean}{GlobalTheory.T.torus_coordinate_isUnit}
\leanref{Development/AutomorphicSpectrum/Proofs/BlockDecomposition/AlgebraicallyIndependentCoordsStatement.lean}{GlobalTheory.T.algebraically_independent_coords}
\leanref{Development/AutomorphicSpectrum/Proofs/BlockDecomposition/TorusMonomialMemRangeStatement.lean}{GlobalTheory.T.torus_monomial_mem_range}
\leanref{Development/AutomorphicSpectrum/Proofs/BlockDecomposition/NoetherNormalizationFiniteStatement.lean}{GlobalTheory.T.noether_normalization_finite}

\begin{proof}
For homogeneous $g$ of degree $e'\geq1$ and $F$ of degree $se'$ write $F/g^s$ also for the section of $\mathcal O_X$ over
$i^{-1}(D_+(g))$ carried through $i$, and for its restrictions, as in
Lemma~\ref{gal-realization-kappa-de-rham-setup-cone-vanish}, which applies to $I_X$.

(1) Let $F$ be one of the generators of $I_X$, of degree $n$. If $n\geq1$, its constant term is $0$. If $n=0$, $F=c\in K$
and $c$ vanishes on every $X_j$; as $X$ is nonempty and the $X_j$ cover it, some $\Gamma(X_j,\mathcal O_X)$ is nonzero, so
$c=0$. So $I_X$ lies in the kernel of the constant term, and $1\notin I_X$. An ideal generated by homogeneous elements
contains the homogeneous components of its elements: for $\sum_ia_iF_i$ with $F_i$ homogeneous, the component of degree
$n$ is $\sum_i(a_i)_{n-\deg F_i}F_i$.

(2) Let $p\in X$, $j$ with $i(p)\in D_+(x_j)$, and $\mathfrak q_p$ the homogeneous prime of $i(p)$, so that
$x_j\notin\mathfrak q_p$. Suppose that $p$ lies in no $U_k$, that is, $y_k\in\mathfrak q_p$ for every $k$. By
Lemma~\ref{gal-realization-kappa-de-rham-setup-cone-vanish}(3), $I_X\subset\mathfrak q_p$. Let $A=K[x_0,\ldots,x_M]/I_X$,
$B\subset A$ the subalgebra generated by the images of the $y_k$, $\varphi:A\to K[x_0,\ldots,x_M]/\mathfrak q_p$ the
projection and $Q=\ker\varphi$, a prime. As $\varphi$ kills the images of the $y_k$, it maps $B$ onto $K$, so $Q\cap B$ is
a maximal ideal of $B$. As $A$ is a finitely generated $B$-module, it is integral over $B$ \cite[Proposition~5.1]{AM69}, so
$Q$ is maximal \cite[Corollary~5.8]{AM69}. As $x_j\notin\mathfrak q_p$, the image of $x_j$ is not in $Q$, so it is
invertible modulo $Q$: there is $g$ with $gx_j-1\in\mathfrak q_p$. As $\mathfrak q_p$ is homogeneous, it contains the
homogeneous component of degree $0$ of $gx_j-1$, which is $-1$; this is impossible. So $p$ lies in some $U_k$.

(3) As $D_+(GH)=D_+(G)\cap D_+(H)$, $U_S=i^{-1}(D_+(y_S))$. As $y_S$ has positive degree, $D_+(y_S)$ is affine with ring
$K[x_0,\ldots,x_M]_{(y_S)}$ \cite[Tag 01MB]{Sta}, and as $i$ is a closed immersion, $U_S$ is affine and the ring map
$K[x_0,\ldots,x_M]_{(y_S)}\to\Gamma(U_S,\mathcal O_X)$ is surjective \cite[Tag 01IN]{Sta}. For $V\subset U_S$ the map from
$K[x_0,\ldots,x_M]_{(y_S)}$ to $\Gamma(V,\mathcal O_X)$ is this map followed by restriction, so the two images are equal.

(4) Let $G$ have degree $e'$, and suppose $i^{-1}(D_+(G))$ is empty. For every $j$, the section $G/x_j^{e'}$ over $X_j$ is
nowhere a unit, as $D_+(x_j)\cap D_+(G)=D_+(x_jG)$, where it is a unit, has empty inverse image. As $X$ is reduced, a
section that is nowhere a unit is $0$; so $G/x_j^{e'}=0$ on $X_j$ for every $j$, and $G$ is one of the generators of $I_X$,
which is false.

(5) The $t_k$ are units by Lemma~\ref{gal-realization-kappa-de-rham-setup-cone-monomial}(1). Let $P=\sum_\alpha
c_\alpha\prod_{k=1}^dz_k^{\alpha_k}\in K[z_1,\ldots,z_d]$ of total degree $D$ with $P(t_0,\ldots,t_{d-1})=0$, and
$P^h=\sum_\alpha c_\alpha z_0^{D-|\alpha|}\prod_kz_k^{\alpha_k}\in K[z_0,\ldots,z_d]$ its homogenization, where
$|\alpha|=\sum_k\alpha_k$; it is homogeneous of degree $D$, and $P^h(1,z_1,\ldots,z_d)=P$. So $F=P^h(y_0,\ldots,y_d)$ is
homogeneous of degree $De$, and on $V$
\[
 F/y_0^D=\sum_\alpha c_\alpha\prod_k(y_k/y_0)^{\alpha_k}=P(t_0,\ldots,t_{d-1})=0.
\]
As $X$ is integral and $V$ is nonempty, $F\in I_X$ by Lemma~\ref{gal-realization-kappa-de-rham-setup-cone-vanish}(2). So
$P^h$ vanishes at the images of the $y_k$, which are algebraically independent; so $P^h=0$, and $P=0$.

(6) By Lemma~\ref{gal-realization-kappa-de-rham-setup-cone-monomial}(2) every generator $\prod_kt_k^{a_k}$ of $\mathcal
A(S)$ is the restriction of some $\prod_ky_k^{\alpha_k}/y_S^j\in K[x_0,\ldots,x_M]_{(y_S)}$, so $\mathcal A(S)$ lies in
the image. An element of $K[x_0,\ldots,x_M]_{(y_S)}$ is $F/y_S^r$ with $F$ homogeneous of degree $er|S|$. By (1) and
Lemma~\ref{gal-realization-kappa-de-rham-setup-cone-relation}, for $w=r|S|$, there are $s\geq1$ and $c_l$ homogeneous of
degree $(s-l)r|S|$ with $G=F^s+\sum_{l<s}c_l(y)F^l\in I_X$. Dividing by $y_S^{rs}$, on $V$
\[
 \Bigl(\frac F{y_S^r}\Bigr)^s+\sum_{l<s}\frac{c_l(y)}{y_S^{(s-l)r}}\Bigl(\frac F{y_S^r}\Bigr)^l=\frac G{y_S^{rs}}=0,
\]
by Lemma~\ref{gal-realization-kappa-de-rham-setup-cone-vanish}(1), and the coefficients $c_l(y)/y_S^{(s-l)r}$ lie in
$\mathcal A(S)$ by Lemma~\ref{gal-realization-kappa-de-rham-setup-cone-monomial}(3). So the image of $F/y_S^r$ is integral
over $\mathcal A(S)$.
% Lean: (1) is GeoNetop.homIdeal_ne_top (constantCoeff_gen, with exists_chart and mk_C_mul_inv) and Geo.homIdeal_homog
% (Ideal.homogeneous_span); (2) is GeoCover.iSup_preimage_eq_top, with mem_of_evV_eq_zero (hle) and X_mem_of_finite (Q
% maximal through Ideal.isMaximal_of_isIntegral_of_isMaximal_comap, and the degree-0 component of g x_j - 1); (3) is
% biInf_preimage, isAffineOpen_preimage (Proj.isAffineOpen_basicOpen and IsAffineOpen.preimage for a closed immersion)
% and range_resO_eq, through evV_surjective = basic_open_eval_surjective (Proj.basicOpenIsoAway, app_surjective); (4) is
% preimage_basicOpen_nonempty_of_mk_ne_zero (basicOpen_eq_bot_iff, basicOpen_awayToSection_isLocalizationElem); (5) is
% GeoUnit.tV_isUnit and GeoAlgInd.algebraicIndependent_tV (homogP, homogenize_isHomogeneous₂, aeval_homogP, away_key,
% evV_cst and mem_homIdeal_of_vanish); (6) is GeoTpow.tpow_mem_range and GeoInteg.isIntegral_torusS (exists_hom_relation,
% rel_mem, away_relation, homIdeal_vanish and aeval_mem_torusS).
% The statement-level names are basic_opens_cover_of_finite, basic_open_nonempty_of_ne_zero, torus_coordinate_isUnit,
% algebraically_independent_coords, torus_monomial_mem_range and noether_normalization_finite. About 1,350 lines in
% BasicOpensCoverOfFinite, TorusCoordinateIsUnit, AlgebraicallyIndependentCoords, TorusMonomialMemRange,
% BasicOpenNonemptyOfNeZero, NoetherNormalizationFinite, IsAffineOpenPreimage, HomogeneousIdealNeTop,
% HomogeneousIdealComponentMem, MemHomogeneousIdealOfEvalZero, EvalZeroOfMemHomogeneousIdeal,
% HomogeneousIntegralExpression and TorusMonomialEvalCompat.
\end{proof}
\leanref{Development/AutomorphicSpectrum/Proofs/BlockDecomposition/HomogeneousIdealNeTop.lean}{GlobalTheory.CompanionTheory.GeoNetop.constantCoeff_gen}
\leanref{Development/AutomorphicSpectrum/Proofs/BlockDecomposition/BasicOpensCoverOfFinite.lean}{GlobalTheory.CompanionTheory.GeoCover.X_mem_of_finite}
\leanref{Development/AutomorphicSpectrum/Proofs/BlockDecomposition/BasicOpenEvalSurjective.lean}{GlobalTheory.CompanionTheory.Geo.evV_surjective}
\leanref{Development/AutomorphicSpectrum/Proofs/BlockDecomposition/AlgebraicallyIndependentCoords.lean}{GlobalTheory.CompanionTheory.GeoAlgInd.algebraicIndependent_tV}
\leanref{Development/AutomorphicSpectrum/Proofs/BlockDecomposition/NoetherNormalizationFinite.lean}{GlobalTheory.CompanionTheory.GeoInteg.isIntegral_torusS}

\begin{lemma}\label{gal-realization-kappa-de-rham-setup-generic}
Let $K$ be a field.
\begin{enumerate}
\item A connected $K$-scheme that is smooth of relative dimension $n$ over $K$ is integral.
\item Let $K$ have characteristic $0$, $A$ a finitely generated $K$-algebra that is a domain, and $s_1,\ldots,s_d\in A$
  algebraically independent over $K$, such that $A$ is algebraic over $K[s_1,\ldots,s_d]$. There is a nonzero $g\in A$ such
  that the $K$-algebra map $K[x_1,\ldots,x_d]\to A_g$, $x_k\mapsto s_k$, is \'etale.
\item Let $X$ be a $K$-scheme smooth of relative dimension $n$ over $K$ and $V\subset X$ a nonempty open with coordinates
  of dimension $m$ (Definition~\ref{gal-realization-thickening-coordinates}(1)). Then $m=n$.
\item If $X$ is a $K$-scheme locally of finite type and $W\subset X$ an affine open, then $\Gamma(W,\mathcal O_X)$ is a
  finitely generated $K$-algebra.
\end{enumerate}
\end{lemma}
\leanref{Development/WeilBound/CoreStatements/PadicHodgeComparison/EtaleCohomologyComparison/PerfectoidPullbacks/ConeSectionRestriction/ExistsCoordBasicOpen.lean}{CrystallineSemistableComparison.Pb.isIntegral_of_smooth}
\leanref{Development/WeilBound/CoreStatements/PadicHodgeComparison/EtaleCohomologyComparison/PerfectoidPullbacks/ConeSectionRestriction/ExistsCoordBasicOpen.lean}{CrystallineSemistableComparison.Pb.exists_etale_away}
\leanref{Development/WeilBound/CoreStatements/PadicHodgeComparison/EtaleCohomologyComparison/PerfectoidPullbacks/ConeSectionRestriction/ExistsCoordBasicOpen.lean}{CrystallineSemistableComparison.Pb.coord_card_eq}
\leanref{Development/AutomorphicSpectrum/Proofs/BlockDecomposition/AffineOpenFiniteTypeStatement.lean}{GlobalTheory.T.affine_open_finiteType}

\begin{proof}
(1) Let $X$ be such a scheme.

\emph{Domains after inverting one element.} Let $A$ be a smooth $K$-algebra and $\mathfrak p$ a prime of $A$. There is
$r\in A\smallsetminus\mathfrak p$ such that $A_r$ is a domain. Indeed, there is $g\notin\mathfrak p$ such that $A_g$ is a
standard smooth $K$-algebra \cite[Tag 00TA]{Sta}, so, as in the proof of
Lemma~\ref{gal-realization-kappa-de-rham-affine-basic}(1), there are $m\geq0$ and a $K$-algebra map
$P=K[x_1,\ldots,x_m]\to A'=A_g$ that makes $A'$ \'etale over $P$. Let $F$ be the field of fractions of $P$ and
$\iota:A'\to A'\otimes_PF$, $a\mapsto a\otimes1$. It is injective, as $A'$ is flat over $P$ \cite[Tag 00U2]{Sta} and
$P\to F$ is injective. The $F$-algebra $A'\otimes_PF$ is \'etale \cite[Tag 00U2]{Sta}, so it is isomorphic to a finite
product $\prod_{i\in E}L_i$ of finite separable field extensions of $F$ \cite[Tag 00U3]{Sta}; let $e_i$ be the idempotent
of the factor $L_i$, so that $\sum_ie_i=1$. Each $e_i$ is integral over $A'$, as $e_i^2=e_i$. As $P\to A'$ is \'etale,
hence smooth, the integral closure of $A'$ in $A'\otimes_PF$ is the image of $A'\otimes_PP'$, where $P'$ is the integral
closure of $P$ in $F$ \cite[Tag 03GG]{Sta}; and $P'=P$, as $P$ is factorial. So $e_i=\iota(a_i)$ with $a_i\in A'$. As
$\iota(\sum_ia_i)=1$ and $\iota$ is injective, $\sum_ia_i=1$, so some $a=a_{i_0}$ is not in the prime $\mathfrak pA'$ of
$A'$. Let $\phi:A'\to L_{i_0}$ be $\iota$ followed by the projection. Then $\phi(a)=1$, so $\phi$ extends to a ring map
$A'_a\to L_{i_0}$, which is injective: if $\phi(y)=0$, then $\iota(ya)=\iota(y)e_{i_0}=0$, as the component of
$\iota(y)e_{i_0}$ at $i_0$ is $\phi(y)$ and the others are $0$, so $ya=0$, and $y$ is $0$ in $A'_a$. So $A'_a$ is a
domain, as a subring of a field. Write $a=b/g^k$ with $b\in A$. Then $b\notin\mathfrak p$, and $r=gb\notin\mathfrak p$ has
$A_r=A'_a$.

\emph{Integral neighbourhoods.} Let $x\in X$, $V$ an affine open neighbourhood of $x$, $A=\Gamma(V,\mathcal O_X)$ and
$\mathfrak p$ the prime of $A$ at $x$. As $X$ is smooth over $K$, $A$ is a smooth $K$-algebra \cite[Tag 01V6]{Sta}. Let
$r\notin\mathfrak p$ be as above. Then $W=D(r)$ is an open neighbourhood of $x$, affine with ring $A_r$
\cite[Tag 01HV]{Sta}, a domain; so $W$ is an integral scheme.

\emph{Irreducibility.} As $X$ is covered by integral opens, it is reduced. It is connected, so nonempty; let $x_0\in X$
and $Z$ an irreducible component of $X$ that contains $x_0$, which is closed. It is also open: for $z\in Z$, let $W$ be an
integral, hence irreducible, open neighbourhood of $z$. As $Z$ is irreducible and meets the open $W$, $Z$ lies in the
closure of $Z\cap W$, hence in $\overline W$, which is irreducible; so $\overline W=Z$ by maximality, and $W\subset Z$. As
$X$ is connected, $Z=X$, so $X$ is irreducible, and $X$ is integral \cite[Tag 01ON]{Sta}.

(2) Let $P=K[x_1,\ldots,x_d]\to A$ be the map $x_k\mapsto s_k$, which is injective as the $s_k$ are algebraically
independent. As $A$ is finitely generated over $K$, it is finitely generated over $P$, so of finite presentation over $P$,
as $P$ is noetherian. Let $F$ and $L$ be the fields of fractions of $P$ and $A$, so that $L$ is the localization of $A$ at
its prime $0$. As $A$ is algebraic over $P$ and generated over it by finitely many elements, $L$ is a finitely generated
algebraic field extension of $F$, separable as $K$ has characteristic $0$; so $L$ is formally \'etale over $F$
\cite[Tag 090W]{Sta}. As $F$ is formally \'etale over $P$ \cite[Tag 04EG]{Sta}, and a lift along a square-zero ideal can be
made first on $F$ and then on $L$, $L$ is formally \'etale over $P$ \cite[Tag 00UQ]{Sta}. So $\Omega_{L/P}=0$
\cite[Tag 00UO]{Sta} and $H_1(\mathrm{NL}_{L/P})=0$ \cite[Tag 031J]{Sta}, where $\mathrm{NL}$ is the naive cotangent
complex. These are the localizations at $0$ of the $A$-modules $\Omega_{A/P}$ and $H_1(\mathrm{NL}_{A/P})$
\cite[Tags 00RT and 00S7]{Sta}, which are finitely generated: for a presentation $A=P[y_1,\ldots,y_N]/I$
\cite[Tag 00S1]{Sta}, $\Omega_{A/P}$ is generated by the $dy_j$, and $H_1(\mathrm{NL}_{A/P})$ is a submodule of $I/I^2$, a
finitely generated module over the noetherian ring $A$. So each generator is killed by a nonzero element of $A$, and the
product $g\neq0$ of these elements kills both modules. Then $P\to A_g$ is of finite presentation with
$\Omega_{A_g/P}=0$ and $H_1(\mathrm{NL}_{A_g/P})=0$ \cite[Tags 00RT and 00S7]{Sta}, that is, it is \'etale
\cite[Tag 00U1]{Sta}. This is the map $x_k\mapsto s_k$.

(3) Let $G=\Gamma(V,\mathcal O_X)$ and $t_1,\ldots,t_m$ the coordinates, so that $P=K[x_1,\ldots,x_m]\to G$, $x_k\mapsto
t_k$, is \'etale. As $V$ is nonempty, $G\neq0$; let $\mathfrak m$ be a maximal ideal of $G$. As $X$ is smooth of relative
dimension $n$ over $K$ and $V$ is affine, the $u\in G$ such that $G_u$ is a standard smooth $K$-algebra of relative
dimension $n$, that is, $G_u\cong K[y_1,\ldots,y_N]/(h_1,\ldots,h_c)$ with $N-c=n$ and $\det(\partial h_k/\partial
y_l)_{1\leq k,l\leq c}$ invertible, generate the unit ideal of $G$. So one of them, $u$, is not in $\mathfrak m$, and
$G_u\neq0$, as $u$ is not nilpotent. The map $P\to G_u$ is formally \'etale, as $P\to G$ is \'etale, hence formally
\'etale \cite[Tag 00UR]{Sta}, and $G\to G_u$ is formally \'etale \cite[Tag 04EG]{Sta}. So the map $G_u\otimes_P\Omega_{P/K}
\to\Omega_{G_u/K}$ is an isomorphism \cite[Tag 0H93]{Sta}, and $\Omega_{G_u/K}$ is free of rank $m$, on the $dt_k$. On the
other hand $\Omega_{G_u/K}$ is free of rank $N-c=n$ \cite[Tag 00T7]{Sta}. As $G_u\neq0$, $m=n$.

(4) As $X$ is locally of finite type over $K$ and $W$ is affine, $K\to\Gamma(W,\mathcal O_X)$ is of finite type
\cite[Tag 01T2]{Sta}.
% Lean: (1) is isIntegral_of_smooth_pf (IsIntegralOfSmoothPf.lean): exists_integral_nbhd (ExistsIntegralNbhd.lean) gives
% each point a basic open of an affine open with a domain as ring, through exists_away_isDomain_of_smooth
% (ExistsAwayIsDomainOfSmooth.lean); then IsReduced.of_openCover, irreducibleSpace_of_connected_of_locally (the
% irreducibility step) and isIntegral_of_irreducibleSpace_of_isReduced. "Domains after inverting one element" is
% exists_away_isDomain_of_smooth: A_g standard smooth (Algebra.IsSmoothAt.exists_notMem_isStandardSmooth = Stacks 00TA),
% étale over a polynomial ring (exists_etale_mvPolynomial), and exists_away_isDomain_of_etale
% (ExistsTmulOneEqOfIsIntegral.lean: hι by flatness, Algebra.FormallyEtale.iff_exists_algEquiv_prod for the product of
% fields, isIntegral_of_mul_self_eq and exists_tmul_one_eq_of_isIntegral through
% TensorProduct.toIntegralClosure_bijective_of_smooth = Stacks 03GG, then hψ for the injectivity into the factor).
% (2) is exists_etale_away_pf (ExistsEtaleAwayPf.lean): finite presentation over P, L/F algebraic and separable, hence
% formally étale (Algebra.FormallyEtale.of_isSeparable, which needs no finiteness; the book cites Stacks 090W, for the
% finitely generated L), P → F formally étale (FormallyEtale.of_isLocalization = Stacks 04EG), FormallyEtale.comp, so
% IsEtaleAt (⊥ : Ideal A), and Algebra.exists_etale_of_isEtaleAt; the book's argument for the last step is that of
% Mathlib's etaleLocus_eq_compl_support and basicOpen_subset_etaleLocus_iff. (3) is coord_card_eq =
% coord_dim_eq_relative_dim = CoordCard.main: the u are those of Mathlib's SmoothOfRelativeDimension, read on V through
% HasRingHomProperty.appLE (for the definition of Stacks 01V5 with fibres of dimension n they come from Stacks 01V7,
% 00T7 and 02G1); then card_eq_of_away (Algebra.FormallyEtale.of_isLocalization, FormallyEtale.comp) and
% card_eq_of_etale (KaehlerDifferential.tensorKaehlerEquivOfFormallyEtale = Stacks 0H93 and
% IsStandardSmoothOfRelativeDimension.rank_kaehlerDifferential). (4) is finiteType_affine = affine_open_finiteType
% (HasRingHomProperty.appLE for LocallyOfFiniteType). About 440 lines in these files.
\end{proof}
\leanref{Development/WeilBound/CoreStatements/PadicHodgeComparison/EtaleCohomologyComparison/PerfectoidPullbacks/ConeSectionRestriction/IsIntegralOfSmoothPf.lean}{CrystallineSemistableComparison.Pb.isIntegral_of_smooth_pf}
\leanref{Development/WeilBound/CoreStatements/PadicHodgeComparison/EtaleCohomologyComparison/PerfectoidPullbacks/ConeSectionRestriction/ExistsEtaleAwayPf.lean}{CrystallineSemistableComparison.Pb.exists_etale_away_pf}
\leanref{Development/AutomorphicSpectrum/Proofs/BlockDecomposition/CoordDimEqRelativeDimStatement.lean}{GlobalTheory.T.coord_dim_eq_relative_dim}
\leanref{Development/WeilBound/CoreStatements/PadicHodgeComparison/EtaleCohomologyComparison/PerfectoidPullbacks/ConeSectionRestriction/ExistsAwayIsDomainOfSmooth.lean}{CrystallineSemistableComparison.Pb.EtaleIntegralityFacts.exists_away_isDomain_of_smooth}
\leanref{Development/WeilBound/CoreStatements/PadicHodgeComparison/EtaleCohomologyComparison/PerfectoidPullbacks/ConeSectionRestriction/ExistsTmulOneEqOfIsIntegral.lean}{CrystallineSemistableComparison.Pb.EtaleIntegralityFacts.exists_away_isDomain_of_etale}
\leanref{Development/AutomorphicSpectrum/Proofs/BlockDecomposition/CoordDimEqRelativeDim.lean}{GlobalTheory.CompanionTheory.CoordCard.card_eq_of_etale}

\begin{lemma}\label{gal-realization-kappa-de-rham-setup}
Let $K$ be a field of characteristic $0$, $n\geq0$, $X$ a connected $K$-scheme smooth of relative dimension $n+1$, and
$i:X\to\mathbf P^M_K$ a closed immersion over $K$. There are $e\geq1$, homogeneous polynomials $y_0,\ldots,y_{n+1}$ of
degree $e$ in $K[x_0,\ldots,x_M]$, with $U_S$ as in Definition~\ref{gal-realization-kappa-de-rham-cone}(1), and an affine
open $V\subset U_{\{0,\ldots,n+1\}}$ with coordinates $t_0,\ldots,t_n$
(Definition~\ref{gal-realization-thickening-coordinates}(1)), with the following properties. For a nonempty
$S\subset\{0,\ldots,n+1\}$ let $B(S)$ be the image of restriction $\Gamma(U_S,\mathcal O_X)\to\Gamma(V,\mathcal O_X)$.
\begin{enumerate}
\item The $U_k$ cover $X$, and $U_S$ is affine for every nonempty $S$.
\item $\Gamma(V,\mathcal O_X)$ is a domain; $t_k$ is the restriction of $y_{k+1}/y_0$ to $V$ for $0\leq k\leq n$; the
  $t_k$ are units of $\Gamma(V,\mathcal O_X)$, algebraically independent over $K$.
\item For nonempty $S$, restriction $\Gamma(U_S,\mathcal O_X)\to\Gamma(V,\mathcal O_X)$ is injective, and $B(S)$ is a
  finitely generated $K$-algebra that contains $\mathcal A(S)$ (Definition~\ref{gal-realization-kappa-de-rham-torus}(1),
  for $R=\Gamma(V,\mathcal O_X)$ and the $t_k$) and is integral over it. $B(\{0,\ldots,n+1\}\smallsetminus\{j\})\subset
  B(\{0,\ldots,n+1\})$ for every $j$, and for every $a\in\Gamma(V,\mathcal O_X)$ there are $b,b'\in B(\{0,\ldots,n+1\})$
  with $b'\neq0$ and $ab'=b$.
\end{enumerate}
\end{lemma}
\leanref{Development/WeilBound/CoreStatements/PadicHodgeComparison/EtaleCohomologyComparison/PerfectoidPullbacks/ConeSectionRestriction/ConeHyp.lean}{CrystallineSemistableComparison.Pb.LocalCohomologySetup.GeoSetupP}
\leanref{Development/WeilBound/CoreStatements/PadicHodgeComparison/EtaleCohomologyComparison/PerfectoidPullbacks/ConeSectionRestriction/GeometricGenericComponents.lean}{CrystallineSemistableComparison.Pb.geo_setup}

\begin{proof}
Let $I_X$ be as in Lemma~\ref{gal-realization-kappa-de-rham-setup-cone}.

\emph{Step 1: the polynomials.} By Lemma~\ref{gal-realization-kappa-de-rham-setup-generic}(1), $X$ is integral, so
nonempty, and $K$ is infinite, as it has characteristic $0$. By
Lemma~\ref{gal-realization-kappa-de-rham-setup-cone}(1) and Lemma~\ref{gal-realization-kappa-de-rham-setup-normalization}
there are $e\geq1$ and homogeneous $y_k$, $0\leq k<m$, of degree $e$ whose images in $K[x_0,\ldots,x_M]/I_X$ are
algebraically independent, with $K[x_0,\ldots,x_M]/I_X$ finite over the subalgebra they generate. The $U_k$ cover $X$
(Lemma~\ref{gal-realization-kappa-de-rham-setup-cone}(2)); as $X$ is nonempty, $m\geq1$, say $m=d+1$. The $U_S$ are affine
for nonempty $S$ (Lemma~\ref{gal-realization-kappa-de-rham-setup-cone}(3)).

\emph{Step 2: the open $W$.} Let $\top=\{0,\ldots,d\}$ and $W=U_\top$, an affine open. The images of the $y_k$ in
$K[x_0,\ldots,x_M]/I_X$ are algebraically independent, so their product is not $0$, and $y_\top\notin I_X$; so $W$ is
nonempty (Lemma~\ref{gal-realization-kappa-de-rham-setup-cone}(4)). Let $t'_k\in\Gamma(W,\mathcal O_X)$, $0\leq k\leq d-1$,
be the restriction of $y_{k+1}/y_0$. The $t'_k$ are units, algebraically independent over $K$, and by
Lemma~\ref{gal-realization-kappa-de-rham-setup-cone}(3) and~(6), for $V=W$ and $S=\top$, every element of $\Gamma(W,\mathcal
O_X)$ is integral over the algebra $\mathcal A(\top)$ of the $t'_k$.

\emph{Step 3: the open $V$.} Each $t'^{-1}_k$ is a root of $t'_kx-1$, so $\mathcal A(\top)$ is algebraic over $K[t'_0,\ldots,
t'_{d-1}]$, and so is $\Gamma(W,\mathcal O_X)$. It is a domain, as $X$ is integral and $W$ is nonempty, and a finitely
generated $K$-algebra (Lemma~\ref{gal-realization-kappa-de-rham-setup-generic}(4)). By
Lemma~\ref{gal-realization-kappa-de-rham-setup-generic}(2) there is a nonzero $g\in\Gamma(W,\mathcal O_X)$ such that
$K[x_0,\ldots,x_{d-1}]\to\Gamma(W,\mathcal O_X)_g$, $x_k\mapsto t'_k$, is \'etale. Let $V=D(g)\subset W$ be the basic open
where $g$ is a unit. As $W$ is affine, $\Gamma(V,\mathcal O_X)=\Gamma(W,\mathcal O_X)_g$, so the restrictions $t_k$ of the
$t'_k$ are coordinates of dimension $d$ on the affine open $V$. As $X$ is integral and $g\neq0$, $V$ is nonempty, and
$d=n+1$ by Lemma~\ref{gal-realization-kappa-de-rham-setup-generic}(3). The $t_k$ are the restrictions of the $y_{k+1}/y_0$.

\emph{Step 4: the properties.} (1) holds by Step 1. (2): $\Gamma(V,\mathcal O_X)$ is a domain, as $X$ is integral and $V$
is nonempty, and the $t_k$ are units, algebraically independent over $K$, by
Lemma~\ref{gal-realization-kappa-de-rham-setup-cone}(5). (3): restriction from $U_S$ to the nonempty open $V$ is
injective, as $X$ is integral. By Lemma~\ref{gal-realization-kappa-de-rham-setup-cone}(3), $B(S)$ is the image of
$K[x_0,\ldots,x_M]_{(y_S)}$ in $\Gamma(V,\mathcal O_X)$, so by Lemma~\ref{gal-realization-kappa-de-rham-setup-cone}(6) it
contains $\mathcal A(S)$ and is integral over it. It is finitely generated, as an image of $\Gamma(U_S,\mathcal O_X)$, which
is finitely generated (Lemma~\ref{gal-realization-kappa-de-rham-setup-generic}(4)). As $U_\top\subset
U_{\top\smallsetminus\{j\}}$, restriction from $U_{\top\smallsetminus\{j\}}$ to $V$ factors through $U_\top$, so
$B(\top\smallsetminus\{j\})\subset B(\top)$. Finally every $a\in\Gamma(V,\mathcal O_X)=\Gamma(W,\mathcal O_X)_g$ is
$b/g^k$ with $b\in\Gamma(W,\mathcal O_X)$, so $ab'=b|_V$ with $b'=g^k|_V\neq0$, and $b|_V,b'\in B(\top)$.
% Lean: Pb.geo_setup (GeometricGenericComponents.lean). Step 1 is isIntegral_of_smooth,
% noether_normalization_homogeneous for homIdeal i (homIdeal_ne_top, homIdeal_homog), iSup_preimage_eq_top and haff
% (biInf_preimage, isAffineOpen_preimage); Step 2 is W, hWne (preimage_basicOpen_nonempty), tV_isUnit,
% algebraicIndependent_tV and hintW (range_resO_eq, isIntegral_torusS); Step 3 is isAlgebraic_of_integral_torusS
% (GeoTorus.isAlgebraic_of_integral_torusS, through isAlgebraic_unit_inv), finiteType_affine, exists_etale_away,
% exists_coord_basicOpen (GeoCoord.exists_coord_basicOpen, through IsAffineOpen.isLocalization_basicOpen) and
% coord_card_eq; Step 4 is the final refine, with map_injective_of_isIntegral, tpow_mem_range, range_resO_eq,
% Algebra.FiniteType.out.map, isIntegral_torusS and exists_frac_basicOpen. About 170 lines in
% GeometricGenericComponents.lean; the rest of the 2,800 lines is in Lemmas
% gal-realization-kappa-de-rham-setup-normalization, -setup-cone and -setup-generic.
\end{proof}
\leanref{Development/WeilBound/CoreStatements/PadicHodgeComparison/EtaleCohomologyComparison/PerfectoidPullbacks/ConeSectionRestriction/IsAlgebraicOfIntegralTorusS.lean}{CrystallineSemistableComparison.Pb.isAlgebraic_of_integral_torusS}
\leanref{Development/WeilBound/CoreStatements/PadicHodgeComparison/EtaleCohomologyComparison/PerfectoidPullbacks/ConeSectionRestriction/ExistsCoordBasicOpen.lean}{CrystallineSemistableComparison.Pb.exists_coord_basicOpen}
\leanref{Development/WeilBound/CoreStatements/PadicHodgeComparison/EtaleCohomologyComparison/PerfectoidPullbacks/ConeSectionRestriction/ExistsCoordBasicOpen.lean}{CrystallineSemistableComparison.Pb.exists_frac_basicOpen}

\begin{lemma}\label{gal-realization-kappa-de-rham-residue-torus}
Let $K$ be a field, $R$ a domain that is a $K$-algebra, $n\geq0$, and $t_0,\ldots,t_n$ units of $R$ that are
algebraically independent over $K$, with the algebras $\mathcal A(S)$ of
Definition~\ref{gal-realization-kappa-de-rham-torus}(1). Let $\top=\{0,\ldots,n+1\}$ and $S_j=\top\smallsetminus\{j\}$. For
$m\in\mathbf Z^{n+1}$ let $t^m=\prod_kt_k^{m_k}$, and let $E(S)$ be the set of the $m$ with $m_k\geq0$ whenever $k+1\notin
S$ and $\sum_km_k\leq0$ if $0\notin S$.
\begin{enumerate}
\item The $t^m$, $m\in\mathbf Z^{n+1}$, are linearly independent over $K$, and $\mathcal A(S)$ is spanned over $K$ by the
  $t^m$ with $m\in E(S)$. In particular $\mathcal A(S)\subset\mathcal A(\top)$, and $(t_0\cdots t_n)^{-1}$ and the
  $t_k^{-1}$ lie in $\mathcal A(\top)$.
\item Let $0\in S$. Then $t_k\in\mathcal A(S)$, the $K$-algebra map $K[x_0,\ldots,x_n]\to\mathcal A(S)$, $x_k\mapsto t_k$,
  is injective, and for every $a\in\mathcal A(S)$ there are polynomials $f,g$ such that $g(t)$ is a unit of $\mathcal
  A(S)$ and $a\,g(t)=f(t)$.
\item Let $u_0=t_0^{-1}$ and $u_k=t_kt_0^{-1}$ for $1\leq k\leq n$. Then the $u_k$ lie in $\mathcal A(S_0)$, and (2) holds
  for $\mathcal A(S_0)$ with the $u_k$ in place of the $t_k$.
\item For nonempty $S$ and every $k$ there is a nonzero $a'\in\mathcal A(S)$ with $t_ka'\in\mathcal A(S)$.
\item For every $j\in\top$ there is a nonzero $w\in\mathcal A(S_j)$ such that every $x\in\mathcal A(\top)$ has some
  $N\geq0$ with $w^Nx\in\mathcal A(S_j)$.
\item For $1\leq j\leq n+1$, $\mathcal A(S_j)$ lies in the span of the $t^m$ with $m_{j-1}\geq0$, and
  $t_0^{-(n+2)}\mathcal A(S_0)$ lies in the span of the $t^m$ with $\sum_km_k\leq-(n+2)$.
\end{enumerate}
\end{lemma}
\leanref{Development/WeilBound/CoreStatements/PadicHodgeComparison/EtaleCohomologyComparison/PerfectoidPullbacks/ConeSectionRestriction/SeamDefs.lean}{CrystallineSemistableComparison.Pb.tpow}
\leanref{Development/WeilBound/CoreStatements/PadicHodgeComparison/EtaleCohomologyComparison/PerfectoidPullbacks/ConeSectionRestriction/SeamDefs.lean}{CrystallineSemistableComparison.Pb.expS}
\leanref{Development/WeilBound/CoreStatements/PadicHodgeComparison/EtaleCohomologyComparison/PerfectoidPullbacks/ConeSectionRestriction/ResidResPf.lean}{CrystallineSemistableComparison.Pb.resRPf_linInd}
\leanref{Development/WeilBound/CoreStatements/PadicHodgeComparison/EtaleCohomologyComparison/PerfectoidPullbacks/ConeSectionRestriction/RaTpowSinglePow.lean}{CrystallineSemistableComparison.Pb.ra_torusS_le_span}
\leanref{Development/WeilBound/CoreStatements/PadicHodgeComparison/EtaleCohomologyComparison/PerfectoidPullbacks/ConeSectionRestriction/RaTpowSinglePow.lean}{CrystallineSemistableComparison.Pb.torusS_loc}
\leanref{Development/WeilBound/CoreStatements/PadicHodgeComparison/EtaleCohomologyComparison/PerfectoidPullbacks/ConeSectionRestriction/InfinityCoordAlgebraicIndependent.lean}{CrystallineSemistableComparison.Pb.infinityCoord_localization}
\leanref{Development/WeilBound/CoreStatements/PadicHodgeComparison/EtaleCohomologyComparison/PerfectoidPullbacks/ConeSectionRestriction/RaTpowSinglePow.lean}{CrystallineSemistableComparison.Pb.torusS_frac}
\leanref{Development/WeilBound/CoreStatements/PadicHodgeComparison/EtaleCohomologyComparison/PerfectoidPullbacks/ConeSectionRestriction/RaTpowSinglePow.lean}{CrystallineSemistableComparison.Pb.torus_scale}
\leanref{Development/WeilBound/CoreStatements/PadicHodgeComparison/EtaleCohomologyComparison/PerfectoidPullbacks/ConeSectionRestriction/RaTpowSinglePow.lean}{CrystallineSemistableComparison.Pb.torusS_succ_span}
\leanref{Development/WeilBound/CoreStatements/PadicHodgeComparison/EtaleCohomologyComparison/PerfectoidPullbacks/ConeSectionRestriction/InfinityCoordAlgebraicIndependent.lean}{CrystallineSemistableComparison.Pb.infinityCoord_span_residue}

\begin{proof}
(1) Let $\sum_mc_mt^m=0$, a finite sum, and $N\geq|m_k|$ for all $m$ with $c_m\neq0$ and all $k$. Multiplying by
$(t_0\cdots t_n)^N$ gives a polynomial relation $\sum_mc_mt^{m+(N,\ldots,N)}=0$ among the $t_k$ with distinct
exponents, so all $c_m=0$. The set $E(S)$ contains $0$ and is closed under addition, so the $t^m$, $m\in E(S)$, span a
$K$-subalgebra, which is $\mathcal A(S)$. And $E(S)\subset E(\top)=\mathbf Z^{n+1}$.

(2) For $m\in E(S)$ let $m^+$ and $m^-$ be the positive and negative parts, $m=m^+-m^-$. As $0\in S$, $-m^-\in E(S)$: its
nonzero entries are those $k$ with $m_k<0$, for which $k+1\in S$. So $g(t)=t^{m^-}$ is a unit of $\mathcal A(S)$, and
$t^mg(t)=t^{m^+}=f(t)$ for monomials $f$ and $g$. If $ag(t)=f(t)$ and $a'g'(t)=f'(t)$, then $(a+a')gg'(t)=(fg'+f'g)(t)$
and $aa'gg'(t)=ff'(t)$; by (1) this covers every $a\in\mathcal A(S)$. The map is injective as the $t_k$ are
algebraically independent, and $t_k=t^{e_k}$ with $e_k\in E(S)$.

(3) A monomial in the $u_k$ is $\prod_ku_k^{e_k}=t^{m}$ with $m_0=-\sum_ke_k$ and $m_k=e_k$ for $k\geq1$; so $u_k\in\mathcal
A(S_0)$, and $u_k^{-1}=t_0t_k^{-1}\in\mathcal A(S_0)$ for $k\geq1$. The $u_k$ are algebraically independent: if $f(u)=0$
with $f$ of total degree $N$, then $t_0^N\prod_ku_k^{\alpha_k}=t_0^{N-|\alpha|}\prod_{k\geq1}t_k^{\alpha_k}$ for
$|\alpha|=\sum_k\alpha_k\leq N$, and $\alpha\mapsto(N-|\alpha|,\alpha_1,\ldots,\alpha_n)$ is injective, so $t_0^Nf(u)=0$ is a
polynomial relation among the $t_k$ whose coefficients are those of $f$. For $m\in E(S_0)$, that is $\sum_km_k\leq0$,
\[
 t^m\prod_{k\geq1}u_k^{m_k^-}=u_0^{-\sum_km_k}\prod_{k\geq1}u_k^{m_k^+},
\]
where $\prod_{k\geq1}u_k^{m_k^-}$ is a unit of $\mathcal A(S_0)$. With (1), sums and products are treated as in (2).

(4) If $0\in S$ take $a'=1$, by (2). Otherwise $j+1\in S$ for some $j$; take $a'=t_j^{-1}$, as $-e_j$ and $e_k-e_j$ lie in
$E(S)$.

(5) By (1) and as the set of $x$ with this property is a subalgebra, it suffices to treat $x=t^m$. For $j=0$ take
$w=t_0^{-1}$ and $N\geq\sum_km_k$: $t^{m-Ne_0}\in\mathcal A(S_0)$. For $j=k+1$ take $w=t_k$ and $N\geq-m_k$:
$t^{m+Ne_k}\in\mathcal A(S_j)$, as $0\in S_j$.

(6) For $j\geq1$, every $m\in E(S_j)$ has $m_{j-1}\geq0$. For $m\in E(S_0)$, $t_0^{-(n+2)}t^m=t^{m-(n+2)e_0}$ with
exponent sum at most $-(n+2)$. Both claims follow from (1).
% Lean: (1) is resRPf_linInd and ra_torusS_le_span, with torusS_le_univ, prod_inv_mem and inv_mem_univ; (2) is
% t_mem_torusS and torusS_loc; (3) is u0 with infinityCoord_mem_torus, infinityCoord_inv_mem,
% infinityCoord_algebraicIndependent (through infinityCoord_monomial_eq) and infinityCoord_localization; (4) is
% torusS_frac; (5) is torus_scale; (6) is torusS_succ_span and infinityCoord_span_residue.
\end{proof}
\leanref{Development/WeilBound/CoreStatements/PadicHodgeComparison/EtaleCohomologyComparison/PerfectoidPullbacks/ConeSectionRestriction/InfinityCoordAlgebraicIndependent.lean}{CrystallineSemistableComparison.Pb.u0}
\leanref{Development/WeilBound/CoreStatements/PadicHodgeComparison/EtaleCohomologyComparison/PerfectoidPullbacks/ConeSectionRestriction/InfinityCoordAlgebraicIndependent.lean}{CrystallineSemistableComparison.Pb.infinityCoord_algebraicIndependent}
\leanref{Development/WeilBound/CoreStatements/PadicHodgeComparison/EtaleCohomologyComparison/PerfectoidPullbacks/ConeSectionRestriction/RaTpowSinglePow.lean}{CrystallineSemistableComparison.Pb.prod_inv_mem}

\begin{lemma}\label{gal-realization-kappa-de-rham-residue-functional}
Let $K$, $R$, $n$, the $t_k$ and $t^m$ be as in Lemma~\ref{gal-realization-kappa-de-rham-residue-torus}, and
$D_0,\ldots,D_n$ $K$-derivations of $R$ with $D_k(t_j)=1$ for $k=j$ and $0$ otherwise. There is a $K$-linear map
$\mathrm{res}:R\to K$ with $\mathrm{res}((t_0\cdots t_n)^{-1})=1$, $\mathrm{res}(t^m)=0$ for $m\neq(-1,\ldots,-1)$, and
$\mathrm{res}(D_k(x))=0$ for every $k$ and every $x\in\mathcal A(\top)$.
\end{lemma}
\leanref{Development/WeilBound/CoreStatements/PadicHodgeComparison/EtaleCohomologyComparison/PerfectoidPullbacks/ConeSectionRestriction/ResR.lean}{CrystallineSemistableComparison.Pb.res_R}

\begin{proof}
By Lemma~\ref{gal-realization-kappa-de-rham-residue-torus}(1) the $t^m$ form a basis of the subspace they span. The linear
form on it with $t^m\mapsto1$ for $m=(-1,\ldots,-1)$ and $t^m\mapsto0$ otherwise extends to a $K$-linear map
$\mathrm{res}:R\to K$, as every subspace of a $K$-vector space has a complement. As $D_k$ is a derivation,
$t^{-m}D_k(t^m)=\sum_jm_jt_j^{-1}D_k(t_j)=m_kt_k^{-1}$, so $D_k(t^m)=m_kt^{m-e_k}$. If $m-e_k=(-1,\ldots,-1)$, then
$m_k=0$; so $\mathrm{res}(D_k(t^m))=0$ for every $m$, and $\mathrm{res}\circ D_k$ vanishes on $\mathcal A(\top)$, which
is spanned by the $t^m$.
% Lean: res_R = res_R_pf, through resRPf_linInd, LinearMap.exists_extend, resRPf_logD and resRPf_Dk.
\end{proof}
\leanref{Development/WeilBound/CoreStatements/PadicHodgeComparison/EtaleCohomologyComparison/PerfectoidPullbacks/ConeSectionRestriction/ResidResPf.lean}{CrystallineSemistableComparison.Pb.res_R_pf}
\leanref{Development/WeilBound/CoreStatements/PadicHodgeComparison/EtaleCohomologyComparison/PerfectoidPullbacks/ConeSectionRestriction/ResidResPf.lean}{CrystallineSemistableComparison.Pb.resRPf_Dk}

\begin{lemma}\label{gal-realization-kappa-de-rham-residue-finite}
Let $K$, $R$, $n$ and the $t_k$ be as in Lemma~\ref{gal-realization-kappa-de-rham-residue-torus}, $L$ the field of
fractions of $R$ and $F=K(t_0,\ldots,t_n)\subset L$. Let $B$ be a finitely generated $K$-subalgebra of $R$ that contains
$\mathcal A(\top)$ and is integral over it, such that for every $a\in R$ there are $b,b'\in B$ with $b'\neq0$ and
$ab'=b$. Then $\mathcal A(\top)\subset F$, and $L$ is a finite extension of $F$.
\end{lemma}
\leanref{Development/WeilBound/CoreStatements/PadicHodgeComparison/EtaleCohomologyComparison/PerfectoidPullbacks/ConeSectionRestriction/RaTpowSinglePow.lean}{CrystallineSemistableComparison.Pb.fldF}
\leanref{Development/WeilBound/CoreStatements/PadicHodgeComparison/EtaleCohomologyComparison/PerfectoidPullbacks/ConeSectionRestriction/TorusMemFldF.lean}{CrystallineSemistableComparison.Pb.fldF_findim}

\begin{proof}
The generators $t^m$ of $\mathcal A(\top)$ lie in $F$. Let $s$ be a finite set of generators of $B$. They are integral
over $\mathcal A(\top)\subset F$, so $E=F(s)$ is a finite extension of $F$. It contains $B$, and every $a\in R$, as
$a=b/b'$ with $b,b'\in B$; so $E=L$.
% Lean: fldF_findim = fldF_findim_pf, with fldF_findim_pf_mem and IntermediateField.finiteDimensional_adjoin.
\end{proof}
\leanref{Development/WeilBound/CoreStatements/PadicHodgeComparison/EtaleCohomologyComparison/PerfectoidPullbacks/ConeSectionRestriction/ResidFinPf.lean}{CrystallineSemistableComparison.Pb.fldF_findim_pf}

\begin{lemma}\label{gal-realization-kappa-de-rham-residue-chart}
Let $K$ be a field, $R$ a domain that is a $K$-algebra, and $\mathcal A\subset\mathcal T$ and $B_0\subset B_u$
$K$-subalgebras of $R$ with $\mathcal A\subset B_0$ and $\mathcal T\subset B_u$, such that $B_0$ and $B_u$ are finitely
generated, $B_0$ is integral over $\mathcal A$, $B_u$ is integral over $\mathcal T$, and for every $a\in R$ there are
$b,b'\in B_u$ with $b'\neq0$ and $ab'=b$. Let $w\in\mathcal A$ be nonzero such that every $x\in\mathcal T$ has some
$N\geq0$ with $w^Nx\in\mathcal A$. Then there is a finitely generated $K$-subalgebra $B'\supset B_0$ of $R$, integral
over $\mathcal A$, such that for every $a\in R$ there are $b,b'\in B'$ with $b'\neq0$ and $ab'=b$.
\end{lemma}
\leanref{Development/WeilBound/CoreStatements/PadicHodgeComparison/EtaleCohomologyComparison/PerfectoidPullbacks/ConeSectionRestriction/ChartExt.lean}{CrystallineSemistableComparison.Pb.chart_ext}

\begin{proof}
\emph{Scaling.} If $g\in R$ is integral over $\mathcal T$, then $w^Ng$ is integral over $\mathcal A$ for some $N$. Indeed,
let $p=x^d+c_{d-1}x^{d-1}+\cdots+c_0$ with $c_i\in\mathcal T$ and $p(g)=0$, and $N$ with $w^Nc_i\in\mathcal A$ for all
$i$. Then $w^Ng$ is a root of $x^d+\sum_i(w^N)^{d-i}c_ix^i$, whose coefficients lie in $\mathcal A$, as $d-i\geq1$.

\emph{The algebra.} Let $t$ and $s$ be finite sets of generators of $B_0$ and $B_u$, and for $g\in s$ let $N_g$ be as in
the scaling step. Let $B'$ be the $K$-algebra generated by $t$ and the $w^{N_g}g$, $g\in s$. It contains $B_0$, hence
$w$, and it is integral over $\mathcal A$, as its generators are. Every $b\in B_u$ has some $w^Lb\in B'$, as this
holds for the generators of $B_u$ and is preserved by sums and products. For $a\in R$ with $ab'=b$, $b,b'\in B_u$, choose
$L_1,L_2$ with $w^{L_1}b,w^{L_2}b'\in B'$; then $w^{L_1+L_2}b$ and $w^{L_1+L_2}b'\neq0$ lie in $B'$, and
$a\,w^{L_1+L_2}b'=w^{L_1+L_2}b$.
% Lean: chart_ext = chart_ext_pf, with chart_ext_scale (Polynomial.scaleRoots) for the scaling step and hpow.
\end{proof}
\leanref{Development/WeilBound/CoreStatements/PadicHodgeComparison/EtaleCohomologyComparison/PerfectoidPullbacks/ConeSectionRestriction/ResidChartPf.lean}{CrystallineSemistableComparison.Pb.chart_ext_pf}
\leanref{Development/WeilBound/CoreStatements/PadicHodgeComparison/EtaleCohomologyComparison/PerfectoidPullbacks/ConeSectionRestriction/ResidChartPf.lean}{CrystallineSemistableComparison.Pb.chart_ext_scale}

\begin{lemma}\label{gal-realization-kappa-de-rham-residue-jacobian-poly}
Let $R\neq0$ be a ring in which every integer $n\geq1$ is invertible, $S$ an $R$-algebra integral over $R$,
$\mathfrak P\neq S$ an ideal of $S$ and $b\in S$. There is $G\in R[X]$ with $G(b)\in\mathfrak P$ and $G'(b)\notin\mathfrak
P$.
\end{lemma}
\leanref{Development/AutomorphicSpectrum/Proofs/BlockDecomposition/ExistsGoodPoly.lean}{GlobalTheory.CompanionTheory.TraceJac.exists_good_poly}

\begin{proof}
As $b$ is integral over $R$, some monic $G\in R[X]$ has $G(b)=0\in\mathfrak P$. Let $G$ be a monic polynomial of least
degree $n$ with $G(b)\in\mathfrak P$. Then $n\geq1$, as $1\notin\mathfrak P$. If $G'(b)\in\mathfrak P$, then $n^{-1}G'$
is monic of degree $n-1$ with value in $\mathfrak P$ at $b$, which is impossible. So $G'(b)\notin\mathfrak P$.
% Lean: exists_good_poly, with Nat.find for the least degree.
\end{proof}

\begin{lemma}\label{gal-realization-kappa-de-rham-residue-jacobian-valuation}
Let $R$ be a discrete valuation ring with uniformizer $\pi$ in which every integer $n\geq1$ is invertible, $F$ its field
of fractions, $L$ a finite separable extension of $F$ and $S$ the integral closure of $R$ in $L$. Let $\mathfrak P$ be a
nonzero prime of $S$ and $v$ the valuation of $L$ at $\mathfrak P$. Let $E$ be an additive derivation of $L$ with
$E(R)\subset R$ and $E(\pi)\in\pi S$, and $I$ an ideal of $S$ with $E(R)\subset I$. Then $v(E(x))\geq0$ for every $x\in
S$. If $I\subset\mathfrak P$, then $v(E(x))>0$ for every $x\in S$.
\end{lemma}
\leanref{Development/AutomorphicSpectrum/Proofs/BlockDecomposition/FaithfulIntegralClosure.lean}{GlobalTheory.CompanionTheory.TraceJac.valuation_bounds}
\leanref{Development/AutomorphicSpectrum/Proofs/BlockDecomposition/ExistsGoodPoly.lean}{GlobalTheory.CompanionTheory.TraceJac.core}

\begin{proof}
\emph{The ring $S$.} It is a Dedekind domain with field of fractions $L$, and a finitely generated $R$-module
\cite[Chapter~I, \S4, Propositions~8 and~9]{Ser79}. Its nonzero primes contain $\pi$: if $0\neq x\in\mathfrak P$ and $x^n+r_{n-1}x^{n-1}+\cdots
+r_0=0$ with $r_i\in R$ and $n$ least, then $r_0\neq0$, as $S$ is a domain, and $r_0\in\mathfrak P\cap R$; so $\mathfrak
P\cap R$ is a nonzero prime of $R$, which is $\pi R$. So they are the finitely many primes that divide $\pi S$, and $S$
is a principal ideal domain \cite[Chapter~I, \S3, Corollary to Proposition~7]{Ser79}. Let $\mathfrak P=\varpi S$, and normalize $v$ by $v(\varpi)=1$.
For $x\in S$, $v(x)\geq0$, and $v(x)>0$ if and only if $\varpi\mid x$. The integers $n\geq1$ are units of $R\subset S$, so
$v(n)=0$.

\emph{Step 1: a bound.} There are $\delta\neq0$ in $S$ and $m\geq0$ with $v(E(x))\geq-m$ for every $x\in S$. Indeed, let
$s_1,\ldots,s_q$ generate $S$ as an $R$-module, $E(s_i)=a_i/d_i$ with $a_i,d_i\in S$, $d_i\neq0$, and $\delta=\prod_id_i$.
For $x=\sum_ir_is_i$ with $r_i\in R$, $\delta E(x)=\sum_i(\delta E(r_i)s_i+r_i\delta E(s_i))\in S$, as $E(r_i)\in R$. Let
$m=v(\delta)$.

\emph{Step 2: the uniformizer.} Let $T\geq0$ with $v(E(x))\geq-T$ for every $x\in S$. Then $v(E(y))\geq1-T$ for every
$y\in\mathfrak P$. As $\pi\in\mathfrak P$ is nonzero, $\pi=\varpi^ew$ with $e\geq1$ and $\varpi\nmid w$, so $v(w)=0$; and
$E(\pi)=\pi s$ with $s\in S$. Then $\varpi^eE(w)+e\varpi^{e-1}wE(\varpi)=E(\pi)=\varpi^ews$, so $ewE(\varpi)=\varpi(ws-E(w))$,
and $v(E(\varpi))=1+v(ws-E(w))\geq1-T$. For $y=\varpi z$, $z\in S$, $E(y)=\varpi E(z)+zE(\varpi)$ has valuation at least
$1-T$.

\emph{Step 3: all elements.} Let $T$ be as in Step 2 and $b\in S$. Then $v(E(b))\geq\min(1-T,0)$, and $v(E(b))\geq1-T$ if
$I\subset\mathfrak P$. Indeed, by Lemma~\ref{gal-realization-kappa-de-rham-residue-jacobian-poly}, for $\mathfrak P$, there
is $G=\sum_ig_iX^i\in R[X]$ with $G(b)\in\mathfrak P$ and $v(G'(b))=0$. As $E$ is a derivation, $E(G(b))=s+G'(b)E(b)$
with $s=\sum_iE(g_i)b^i\in I$. So $v(E(b))=v(E(G(b))-s)\geq\min(v(E(G(b))),v(s))$, where $v(E(G(b)))\geq1-T$ by Step 2,
and $v(s)\geq0$, or $v(s)\geq1$ if $I\subset\mathfrak P$.

\emph{Step 4: conclusion.} By Step 1 the hypothesis of Step 2 holds for $T=m$, and if it holds for $T\geq1$, then by
Step 3 it holds for $T-1$. So it holds for $T=0$: $v(E(x))\geq0$ for $x\in S$. If $I\subset\mathfrak P$, Step 3 with
$T=0$ gives $v(E(x))\geq1$.
% Lean: valuation_bounds, with mem_prime (π ∈ 𝔓, through Ideal.comap_ne_bot_of_integral_mem), isPID_integralClosure
% (IsIntegralClosure.isDedekindDomain, finitely many primes through Ideal.finite_factors) and exists_den (Step 1), and
% core, written with a multiplicative valuation: stepA and stepB are Step 2, stepC is Step 3 (exists_good_poly and
% deriv_aeval), hbase and hdesc are Step 4.
\end{proof}
\leanref{Development/AutomorphicSpectrum/Proofs/BlockDecomposition/FaithfulIntegralClosure.lean}{GlobalTheory.CompanionTheory.TraceJac.exists_den}
\leanref{Development/AutomorphicSpectrum/Proofs/BlockDecomposition/ExistsGoodPoly.lean}{GlobalTheory.CompanionTheory.TraceJac.deriv_aeval}
\leanref{Development/AutomorphicSpectrum/Proofs/BlockDecomposition/FaithfulIntegralClosure.lean}{GlobalTheory.CompanionTheory.TraceJac.mem_prime}
\leanref{Development/AutomorphicSpectrum/Proofs/BlockDecomposition/FaithfulIntegralClosure.lean}{GlobalTheory.CompanionTheory.TraceJac.isPID_integralClosure}

\begin{lemma}\label{gal-realization-kappa-de-rham-residue-jacobian-nilpotent}
Let $R$, $\pi$, $F$, $L$ and $S$ be as in Lemma~\ref{gal-realization-kappa-de-rham-residue-jacobian-valuation}, without
the condition on the integers. If $y\in S$ and $y^N\in\pi S$ for some $N\geq1$, then $\operatorname{Tr}_{L/F}(y)\in\pi R$.
\end{lemma}
\leanref{Development/AutomorphicSpectrum/Proofs/BlockDecomposition/FaithfulIntegralClosure.lean}{GlobalTheory.CompanionTheory.TraceJac.trace_radical}

\begin{proof}
The $R$-module $S$ is finitely generated \cite[Chapter~I, \S4, Proposition~8]{Ser79} and torsion-free, so free, with a basis
$e_1,\ldots,e_r$. Every element of $L$ is $x/a$ with $x\in S$ and $0\neq a\in R$, as $L$ is algebraic over $F$; so the
$e_i$ form a basis of $L$ over $F$, and $\operatorname{Tr}_{L/F}(y)$ is the trace of the matrix $M_y$ of multiplication by
$y$ on $S$, whose entries lie in $R$. Let $y^N=\pi z$ with $z\in S$. Then $M_y^N=\pi M_z$, so the reduction of $M_y$
modulo $\pi$ is a nilpotent matrix over the field $R/\pi R$, and its trace is $0$. So $\operatorname{tr}M_y\in\pi R$.
% Lean: trace_radical, with IsIntegralClosure.finite and module_free, Algebra.trace_localization for the trace of S
% over R, and Matrix.isNilpotent_trace_of_isNilpotent.
\end{proof}

\begin{lemma}\label{gal-realization-kappa-de-rham-residue-jacobian-dvr}
Let $R$, $\pi$, $F$ and $L$ be as in Lemma~\ref{gal-realization-kappa-de-rham-residue-jacobian-valuation}, and
$E_1,\ldots,E_d$ additive derivations of $L$ with $E_k(R)\subset R$ for every $k$, such that $E_j(\pi)$ is a unit of $R$
for some $j$. Then $\operatorname{Tr}_{L/F}\bigl(c\det(E_k(b_l))_{1\leq k,l\leq d}\bigr)\in R$ for all $c,b_1,\ldots,b_d\in
L$ integral over $R$.
\end{lemma}
\leanref{Development/AutomorphicSpectrum/Proofs/BlockDecomposition/FaithfulIntegralClosure.lean}{GlobalTheory.CompanionTheory.TraceJac.trace_jac_dvr}

\begin{proof}
Let $S$ be the integral closure of $R$ in $L$, a principal ideal domain whose nonzero primes contain $\pi$ (proof of
Lemma~\ref{gal-realization-kappa-de-rham-residue-jacobian-valuation}); $c$ and the $b_l$ lie in $S$. For a nonzero prime
$\mathfrak P$ of $S$ let $v_{\mathfrak P}$ be the valuation of $L$ at $\mathfrak P$.

\emph{Step 1: the derivations.} Let $u=E_j(\pi)$ and $\alpha_k=E_k(\pi)u^{-1}\in R$. For $k\neq j$ let $E'_k=E_k-\alpha_kE_j$;
then $E'_k(\pi)=0$ and $E'_k(R)\subset R$. Let $E''=\pi E_j$; then $E''(R)\subset\pi R\subset\pi S$, and $E''(\pi)=\pi u$.

\emph{Step 2: valuations.} Let $\mathfrak P$ be a nonzero prime of $S$. By
Lemma~\ref{gal-realization-kappa-de-rham-residue-jacobian-valuation}, for $E'_k$ and $I=S$, $v_{\mathfrak P}(E'_k(x))\geq0$
for $x\in S$ and $k\neq j$. By the same lemma, for $E''$ and $I=\pi S$, which lies in $\mathfrak P$ as $\pi\in\mathfrak P$,
$v_{\mathfrak P}(\pi E_j(x))>0$ for $x\in S$.

\emph{Step 3: the matrix.} Let $M=(E_k(b_l))_{k,l}$, and let $M_2$ be the matrix whose row $j$ is $\pi$ times row $j$ of
$M$ and whose row $k\neq j$ is row $k$ of $M$ minus $\alpha_k$ times row $j$ of $M$. Then $\det M_2=\pi\det M$. The
entries of row $k\neq j$ of $M_2$ are the $E'_k(b_l)$, and those of row $j$ are the $\pi E_j(b_l)$. So for every nonzero
prime $\mathfrak P$, the entries of $M_2$ have $v_{\mathfrak P}\geq0$ and those of row $j$ have $v_{\mathfrak P}>0$, by
Step 2. Expanding along row $j$, $v_{\mathfrak P}(\det M_2)>0$, and $v_{\mathfrak P}(c\det M_2)>0$, as $c\in S$.

\emph{Step 4: a nilpotent element.} Let $w=c\det M_2=\pi c\det M$. As $S$ is a principal ideal domain, an element of $L$
with $v_{\mathfrak P}\geq0$ at every nonzero prime lies in $S$; so $w\in S$. Write $\pi S=\varpi_1^{e_1}\cdots\varpi_q^{e_q}S$
with prime elements $\varpi_i$. As $v_{\varpi_iS}(w)>0$, every $\varpi_i$ divides $w$, so $w^N\in\pi S$ for $N=\max_ie_i$.

\emph{Step 5: the trace.} By Lemma~\ref{gal-realization-kappa-de-rham-residue-jacobian-nilpotent},
$\pi\operatorname{Tr}_{L/F}(c\det M)=\operatorname{Tr}_{L/F}(w)=\pi r$ with $r\in R$. So $\operatorname{Tr}_{L/F}(c\det
M)=r\in R$.
% Lean: trace_jac_dvr. Step 1 is α, hE'π, hE'R and hE''R; Step 2 is hvb (valuation_bounds with I = ⊤ and with I =
% span {π}); Step 3 is M2, hdet (det_rowop_lemma, det_updateRow_smul), hent and hw (val_det_lt_lemma_3, through
% val_det_le_lemma_2 for the cofactors); Step 4 is mem_integers_of_valuation_le_one and hrad (Ideal.radical_eq_sInf);
% Step 5 is trace_radical.
\end{proof}

\begin{lemma}\label{gal-realization-kappa-de-rham-residue-jacobian}
Let $K$ be a field of characteristic $0$, $d\geq0$, $\mathcal A$ a domain that is a $K$-algebra, and $u_1,\ldots,u_d\in
\mathcal A$ such that the $K$-algebra map $K[x_1,\ldots,x_d]\to\mathcal A$, $x_k\mapsto u_k$, is injective and for every
$a\in\mathcal A$ there are polynomials $f,g$ with $g(u)$ a unit of $\mathcal A$ and $a\,g(u)=f(u)$. Let $B$ be a domain
that is an $\mathcal A$-algebra, finitely generated as an $\mathcal A$-module, with $\mathcal A\to B$ injective, let $F$ and
$L$ be the fields of fractions of $\mathcal A$ and $B$, and $D_1,\ldots,D_d$ $K$-derivations of $L$ with $D_k(u_j)=1$ for
$k=j$ and $0$ otherwise. Then $\operatorname{Tr}_{L/F}\bigl(c\det(D_k(b_l))_{1\leq k,l\leq d}\bigr)\in\mathcal A$ for all
$c,b_1,\ldots,b_d\in B$.
\end{lemma}
\leanref{Development/AutomorphicSpectrum/Proofs/BlockDecomposition/ExistsPrimeDvdNotUnit.lean}{GlobalTheory.CompanionTheory.TraceJac.trace_jacobian_mem}

\begin{proof}
Let $P=K[x_1,\ldots,x_d]$, identified with its image in $\mathcal A$, so that $f$ is $f(u)$; it is a factorial ring. Let
$\tau=\operatorname{Tr}_{L/F}(c\det(D_k(b_l)))$.

\emph{Step 1: derivatives of polynomials.} $D_k(f)=\partial f/\partial x_k$ for $f\in P$, as both sides are $K$-derivations
$P\to L$ that agree on the $x_j$.

\emph{Step 2: prime elements.} For every prime element $\pi$ of $P$ there is $j$ with $\pi\nmid\partial\pi/\partial x_j$.
Indeed, $\pi$ is neither $0$ nor a unit, so it is not constant, and some monomial of $\pi$ has a positive exponent at some
$x_j$. As $K$ has characteristic $0$, $\partial\pi/\partial x_j\neq0$, and its total degree is less than that of $\pi$, while
nonzero multiples of $\pi$ have total degree at least that of $\pi$.

\emph{Step 3: the fields.} By the hypothesis on $\mathcal A$, every element of $\mathcal A$ is a fraction $f/g$ of elements
of $P$, so $F$ is the field of fractions of $P$. As $B$ is a finitely generated $\mathcal A$-module, its generators are
integral over $F$, and they generate $L$ over $F$, as every element of $L$ is a fraction of elements of $B$. So $L$ is
finite over $F$, and separable, as $K$ has characteristic $0$. The elements of $B$ are integral over $\mathcal A$.

\emph{Step 4: at a prime.} Let $\pi$ be a prime element of $P$ that is not a unit of $\mathcal A$, and $R=P_{(\pi)}\subset
F$. It is a noetherian local domain with maximal ideal $\pi R\neq0$, and $\bigcap_n\pi^nR=0$ \cite[Corollary~10.20]{AM69};
so every nonzero element of $R$ is a unit times a power of $\pi$, and $R$ is a discrete valuation ring with uniformizer
$\pi$, in which every integer $n\geq1$ is invertible. It contains $\mathcal A$: for $a\in\mathcal A$ with $ag=f$, $g$ a unit
of $\mathcal A$, $\pi\nmid g$, as otherwise $\pi$ would divide a unit of $\mathcal A$; so $a=f/g\in R$. So $c$ and the $b_l$
are integral over $R$. By Step 1 and the quotient rule, $D_k(f/g)=(g\,\partial f/\partial x_k-f\,\partial g/\partial
x_k)/g^2\in R$ for $f/g\in R$, so $D_k(R)\subset R$. By Step 2, $D_j(\pi)=\partial\pi/\partial x_j$ is a unit of $R$ for some
$j$. So $\tau\in R$ by Lemma~\ref{gal-realization-kappa-de-rham-residue-jacobian-dvr}: $\tau=f/s$ with $f,s\in P$ and
$\pi\nmid s$.

\emph{Step 5: conclusion.} As $P$ is factorial, $\tau=n_0/d_0$ with coprime $n_0,d_0\in P$, $d_0\neq0$. If $d_0$ is a unit
of $\mathcal A$, then $\tau=n_0d_0^{-1}\in\mathcal A$. Otherwise, as $d_0$ is a unit of $K$ times a product of prime
elements, one of them, $\pi$, is not a unit of $\mathcal A$. By Step 4, $\tau=f/s$ with $\pi\nmid s$, so $sn_0=fd_0$. Then
$\pi$ divides $sn_0$ and not $s$, so $\pi\mid n_0$, against the coprimality of $n_0$ and $d_0$. So $d_0$ is a unit of
$\mathcal A$, and $\tau\in\mathcal A$.
% Lean: trace_jacobian_mem, through main_abs. Step 1 is deriv_aeval_mv (hE, hτ); Step 2 is exists_pderiv_not_dvd (hjac);
% Step 3 is the IsFractionRing P F instance of main_abs (from hloc') and finite_frac, with Algebra.IsSeparable from
% characteristic 0; Step 4 is key, through at_prime_abs (the localization R at span {π} is a DVR by
% IsDiscreteValuationRing.TFAE, hR by the quotient rule, u from hfj and hj) with hr and integral_of_range, and
% trace_jac_dvr (Lemma gal-realization-kappa-de-rham-residue-jacobian-dvr); Step 5 is IsFractionRing.exists_reduced_fraction
% and exists_prime_dvd_not_unit. About 1,000 lines in ExistsPrimeDvdNotUnit.lean, ExistsPderivNotDvd.lean,
% FaithfulIntegralClosure.lean and ExistsGoodPoly.lean.
\end{proof}
\leanref{Development/AutomorphicSpectrum/Proofs/BlockDecomposition/ExistsPderivNotDvd.lean}{GlobalTheory.CompanionTheory.TraceJac.at_prime_abs}
\leanref{Development/AutomorphicSpectrum/Proofs/BlockDecomposition/ExistsPderivNotDvd.lean}{GlobalTheory.CompanionTheory.TraceJac.exists_pderiv_not_dvd}
\leanref{Development/AutomorphicSpectrum/Proofs/BlockDecomposition/ExistsPrimeDvdNotUnit.lean}{GlobalTheory.CompanionTheory.TraceJac.deriv_aeval_mv}
\leanref{Development/AutomorphicSpectrum/Proofs/BlockDecomposition/ExistsPrimeDvdNotUnit.lean}{GlobalTheory.CompanionTheory.TraceJac.exists_prime_dvd_not_unit}
\leanref{Development/AutomorphicSpectrum/Proofs/BlockDecomposition/ExistsPrimeDvdNotUnit.lean}{GlobalTheory.CompanionTheory.TraceJac.main_abs}

\begin{lemma}\label{gal-realization-kappa-de-rham-residue-trace}
Let $K$ be a field of characteristic $0$, and $R$, $n$, the $t_k$, $\mathcal A(S)$, $L$ and $F$ as in
Lemmas~\ref{gal-realization-kappa-de-rham-residue-torus} and~\ref{gal-realization-kappa-de-rham-residue-finite}. Let $S$
be nonempty, and $u_0,\ldots,u_n\in\mathcal A(S)$ such that the $K$-algebra map $K[x_0,\ldots,x_n]\to\mathcal A(S)$,
$x_k\mapsto u_k$, is injective and for every $a\in\mathcal A(S)$ there are polynomials $f,g$ with $g(u)$ a unit of
$\mathcal A(S)$ and $a\,g(u)=f(u)$. Let $B'$ be a finitely generated $K$-subalgebra of $R$ that contains $\mathcal A(S)$
and is integral over it, such that for every $a\in R$ there are $b,b'\in B'$ with $b'\neq0$ and $ab'=b$, and let
$E_0,\ldots,E_n$ be $K$-derivations of $R$ with $E_k(u_j)=1$ for $k=j$ and $0$ otherwise. Let $0\leq m\leq n+1$ and
$\kappa:\{1,\ldots,m\}\to\{0,\ldots,n\}$ be injective, and let $y$ lie in the additive subgroup of $R$ generated by the
elements $b\det(E_{\kappa(k)}(b_l))_{1\leq k,l\leq m}$ with $b,b_1,\ldots,b_m\in B'$. Then $L$ is finite over $F$, and
$\operatorname{Tr}_{L/F}(y)\in\mathcal A(S)$.
\end{lemma}
\leanref{Development/WeilBound/CoreStatements/PadicHodgeComparison/EtaleCohomologyComparison/PerfectoidPullbacks/ConeSectionRestriction/TraceChart.lean}{CrystallineSemistableComparison.Pb.trace_chart}
\leanref{Development/AutomorphicSpectrum/Proofs/BlockDecomposition/FractionFieldTranscendenceBasisExtStatement.lean}{GlobalTheory.T.fraction_field_transcendence_basis_ext}

\begin{proof}
\emph{The fields.} $\mathcal A(S)\subset\mathcal A(\top)\subset F$ (Lemmas~\ref{gal-realization-kappa-de-rham-residue-torus}(1)
and~\ref{gal-realization-kappa-de-rham-residue-finite}), and $F$ is the field of fractions of $\mathcal A(S)$: by
Lemma~\ref{gal-realization-kappa-de-rham-residue-torus}(4) each $t_k$ is $(t_ka')/a'$ with $a',t_ka'\in\mathcal A(S)$,
and the fractions of elements of $\mathcal A(S)$ form a subfield. $L$ is the field of fractions of $B'$, as every element of
$R$ is a fraction of elements of $B'$. As $B'$ is finitely generated and integral over $\mathcal A(S)$, it is a finitely
generated $\mathcal A(S)$-module, so $L$ is finite over $F$.

\emph{The derivations.} Each $E_k$ extends to a $K$-derivation of $L$, by $E_k(r/s)=(sE_k(r)-rE_k(s))/s^2$.

\emph{Minors.} Let $b,b_1,\ldots,b_m\in B'$, and complete the $b_l$ to a family $(b'_j)_{0\leq j\leq n}$ with
$b'_{\kappa(l)}=b_l$ and $b'_j=u_j$ for $j$ outside the image of $\kappa$. Ordering rows and columns alike, with the
image of $\kappa$ first, the matrix $(E_k(b'_j))_{k,j}$ becomes block lower triangular, with diagonal blocks
$(E_{\kappa(k)}(b_l))_{k,l}$ and the identity, as $E_k(u_j)=0$ for $k\neq j$. So
$\det(E_k(b'_j))=\det(E_{\kappa(k)}(b_l))$, and Lemma~\ref{gal-realization-kappa-de-rham-residue-jacobian}, for $\mathcal
A=\mathcal A(S)$ with the $u_k$, $B=B'$ and $d=n+1$, gives $\operatorname{Tr}_{L/F}(b\det(E_{\kappa(k)}(b_l)))\in\mathcal
A(S)$. As the trace is additive, the claim holds for every $y$ in the subgroup.
% Lean: trace_chart, with hQ and IsFractionRing.of_field for the fields, Algebra.IsIntegral.finite, frac_deriv for the
% derivations, and fraction_field_transcendence_basis_ext = TraceJac.trace_minor_mem, which completes the minor
% (det_completed_eq_minor) and applies trace_jacobian_mem (Lemma gal-realization-kappa-de-rham-residue-jacobian).
\end{proof}
\leanref{Development/AutomorphicSpectrum/Proofs/BlockDecomposition/FractionFieldTranscendenceBasisExt.lean}{GlobalTheory.CompanionTheory.TraceJac.trace_minor_mem}
\leanref{Development/AutomorphicSpectrum/Proofs/BlockDecomposition/FractionFieldTranscendenceBasisExt.lean}{GlobalTheory.CompanionTheory.TraceJac.det_completed_eq_minor}
\leanref{Development/WeilBound/CoreStatements/PadicHodgeComparison/EtaleCohomologyComparison/PerfectoidPullbacks/ConeSectionRestriction/TorusMemFldF.lean}{CrystallineSemistableComparison.Pb.frac_deriv}

\begin{lemma}\label{gal-realization-kappa-de-rham-residue-derivation}
Let $K$ be a field of characteristic $0$, $L/F$ a finite extension of fields that contain $K$, and $D$ a $K$-derivation
of $L$ with $D(F)\subset F$. Then $\operatorname{Tr}_{L/F}(D(y))=D(\operatorname{Tr}_{L/F}(y))$ for every $y\in L$. The
hypothesis that $K$ has characteristic $0$ is not used.
\end{lemma}
\leanref{Development/AutomorphicSpectrum/Proofs/BlockDecomposition/DerivationCommutesWithTraceStatement.lean}{GlobalTheory.T.derivation_commutes_with_trace}

\begin{proof}
Let $b_1,\ldots,b_N$ be a basis of $L$ over $F$, $M$ and $M'$ the matrices of multiplication by $y$ and by $D(y)$, so
that $yb_j=\sum_iM_{ij}b_i$, and $C$ the matrix with $D(b_j)=\sum_iC_{ij}b_i$. Applying $D$ to $yb_j=\sum_iM_{ij}b_i$ and
using $D(ab)=D(a)b+aD(b)$ for $a\in F$ gives
\[
 D(y)b_j+yD(b_j)=\sum_i\bigl(D(M_{ij})b_i+M_{ij}D(b_i)\bigr),
\]
that is, $M'+MC=D(M)+CM$, where $D(M)$ is $D$ applied to the entries. Taking traces, as $\operatorname{tr}(MC)=
\operatorname{tr}(CM)$, gives $\operatorname{tr}M'=D(\operatorname{tr}M)$, which is the claim.
% Lean: GlobalTheory.T.derivation_commutes_with_trace = TraceDerivation.trace_map_derivation, with leftMulMatrix for M
% and M' (key, hmat, htr). It assumes K of characteristic 0, which the argument does not use.
\end{proof}
\leanref{Development/AutomorphicSpectrum/Proofs/BlockDecomposition/DerivationCommutesWithTrace.lean}{GlobalTheory.CompanionTheory.TraceDerivation.trace_map_derivation}

\begin{lemma}\label{gal-realization-kappa-de-rham-residue}
Let $K$ be a field of characteristic $0$, $n\geq0$, $R$ a domain that is a $K$-algebra, $t_0,\ldots,t_n$ units of $R$
that are algebraically independent over $K$, and $D_0,\ldots,D_n$ $K$-derivations of $R$ with $D_k(t_j)=1$ for $k=j$
and $D_k(t_j)=0$ for $k\neq j$. Let $\top=\{0,\ldots,n+1\}$, $S_j=\top\smallsetminus\{j\}$, and for every nonempty
$S\subset\top$ let $B(S)$ be a finitely generated $K$-subalgebra of $R$ that contains $\mathcal A(S)$
(Definition~\ref{gal-realization-kappa-de-rham-torus}) and is integral over it. Suppose that $B(S_j)\subset B(\top)$
for every $j$, and that for every $a\in R$ there are $b,b'\in B(\top)$ with $b'\neq0$ and $ab'=b$. Then
\[
 (t_0\cdots t_n)^{-1}\neq\sum_{j=0}^{n+1}x_j+\sum_{k=0}^n(-1)^kD_k(\eta_k)
\]
for all $x_j\in\mathrm{Jac}(B(S_j))$ and $\eta_k\in\mathrm{Jac}_k(B(\top))$.
\end{lemma}
\leanref{Development/WeilBound/CoreStatements/PadicHodgeComparison/EtaleCohomologyComparison/PerfectoidPullbacks/ConeSectionRestriction/ConeHyp.lean}{CrystallineSemistableComparison.Pb.LocalCohomologySetup.ResidAlgP}
\leanref{Development/WeilBound/CoreStatements/PadicHodgeComparison/EtaleCohomologyComparison/PerfectoidPullbacks/ConeSectionRestriction/ResidStmt.lean}{CrystallineSemistableComparison.Pb.resid_alg}

\begin{proof}
Suppose that $(t_0\cdots t_n)^{-1}=\sum_jx_j+\sum_k(-1)^kD_k(\eta_k)$ with $x_j\in\mathrm{Jac}(B(S_j))$ and
$\eta_k\in\mathrm{Jac}_k(B(\top))$. Write $t^m$ as in Lemma~\ref{gal-realization-kappa-de-rham-residue-torus}.

\emph{Step 1: the trace and the residue.} Let $L$ be the field of fractions of $R$, which contains $R$, and
$F=K(t_0,\ldots,t_n)\subset L$. By Lemma~\ref{gal-realization-kappa-de-rham-residue-finite}, for $B=B(\top)$, $L$ is
finite over $F$; let $\operatorname{Tr}=\operatorname{Tr}_{L/F}$. Let $\mathrm{res}:R\to K$ be as in
Lemma~\ref{gal-realization-kappa-de-rham-residue-functional}.

\emph{Step 2: the faces $S_j$, $j\geq1$.} By Lemma~\ref{gal-realization-kappa-de-rham-residue-torus}(5) and
Lemma~\ref{gal-realization-kappa-de-rham-residue-chart}, for $\mathcal A=\mathcal A(S_j)$, $\mathcal T=\mathcal A(\top)$,
$B_0=B(S_j)$ and $B_u=B(\top)$, there is a finitely generated $K$-subalgebra $B'\supset B(S_j)$ of $R$, integral over
$\mathcal A(S_j)$, whose fractions give all of $R$. As $0\in S_j$,
Lemma~\ref{gal-realization-kappa-de-rham-residue-torus}(2) and~(4) and
Lemma~\ref{gal-realization-kappa-de-rham-residue-trace}, for $S=S_j$, $u_k=t_k$, $E_k=D_k$, $m=n+1$ and $\kappa(l)=l-1$,
give $\operatorname{Tr}(x_j)=a_j\in\mathcal A(S_j)$, as $\mathrm{Jac}(B(S_j))\subset\mathrm{Jac}(B')$. By
Lemma~\ref{gal-realization-kappa-de-rham-residue-torus}(6), $a_j$ is a combination of $t^m$ with $m_{j-1}\geq0$, so
$\mathrm{res}(a_j)=0$.

\emph{Step 3: the face $S_0$.} Let $u_0=t_0^{-1}$ and $u_k=t_kt_0^{-1}$ for $k\geq1$, and let $D'_0=-t_0\sum_kt_kD_k$ and
$D'_k=t_0D_k$ for $1\leq k\leq n$. Then $D'_k(u_j)=1$ for $k=j$ and $0$ otherwise: $D_k(t_0^{-1})=-t_0^{-2}D_k(t_0)$ gives
$D'_0(u_0)=t_0t_0t_0^{-2}=1$ and, for $j\geq1$, $D'_0(u_j)=-t_0(t_jt_0^{-1}-t_0t_jt_0^{-2})=0$; for $k\geq1$,
$D'_k(u_0)=0$, and $D'_k(u_j)=t_0t_0^{-1}=1$ for $j=k$ and $0$ for $j\geq1$, $j\neq k$. For $b_0,\ldots,b_n\in R$, $(D'_k(b_l))_{k,l}=C\,(D_k(b_l))_{k,l}$, where $C$ has first
row $(-t_0t_0,\ldots,-t_0t_n)$ and $C_{kk}=t_0$, $C_{kl}=0$ otherwise, for $k\geq1$. It is upper triangular with
$\det C=-t_0^{n+2}$. So $x_0=-t_0^{-(n+2)}y'$ for some $y'$ in the additive subgroup generated by the
$b\det(D'_k(b_l))$, $b,b_l\in B(S_0)$. As in Step 2, with Lemma~\ref{gal-realization-kappa-de-rham-residue-torus}(3) for
the $u_k$ and the $E_k=D'_k$, $\operatorname{Tr}(y')=a'\in\mathcal A(S_0)$. As $c=-t_0^{-(n+2)}\in\mathcal A(\top)\subset
F$, $\operatorname{Tr}(x_0)=a_0=ca'$, and by Lemma~\ref{gal-realization-kappa-de-rham-residue-torus}(6), $a_0$ is a
combination of $t^m$ with $\sum_km_k\leq-(n+2)$, so $\mathrm{res}(a_0)=0$.

\emph{Step 4: the terms $D_k(\eta_k)$.} Lemma~\ref{gal-realization-kappa-de-rham-residue-trace}, for $S=\top$,
$B'=B(\top)$, $u_j=t_j$, $E_j=D_j$, $m=n$ and $\kappa$ the increasing enumeration of $\{0,\ldots,n\}\smallsetminus\{k\}$,
gives $\operatorname{Tr}(\eta_k)=b_k\in\mathcal A(\top)$. Extend $D_k$ to $L$ by the quotient rule. As $D_k(t_j)\in K$, the
elements of $F$ on which the extension takes values in $F$ form a subfield that contains the $t_j$, so $D_k(F)\subset F$.
By Lemma~\ref{gal-realization-kappa-de-rham-residue-derivation}, $\operatorname{Tr}(D_k(\eta_k))=D_k(b_k)$, and
$\mathrm{res}(D_k(b_k))=0$ by Lemma~\ref{gal-realization-kappa-de-rham-residue-functional}.

\emph{Step 5: the contradiction.} As $(t_0\cdots t_n)^{-1}\in\mathcal A(\top)\subset F$, its trace is
$[L:F](t_0\cdots t_n)^{-1}$. Taking traces in the supposed identity gives
$[L:F](t_0\cdots t_n)^{-1}=\sum_ja_j+\sum_k(-1)^kD_k(b_k)$, an identity in $R$, as $R\to L$ is injective. Applying
$\mathrm{res}$ gives $[L:F]=0$ in $K$, which is false, as $[L:F]\geq1$ and $K$ has characteristic $0$.
% Lean: resid_alg (ResidStmt.lean). Step 1 is fldF_findim, res_R and Φ (the trace, carried back to FractionRing R);
% Step 2 is hchart (torus_scale, chart_ext) and the succ case of hω' (torusS_loc, trace_chart, torusS_succ_span); Step 3
% is the zero case of hω' (jacobian_infinityDerivation with infinityDerivation and infinityDerivation_coord,
% infinityCoord_localization, trace_chart, infinityCoord_span_residue_ne); Step 4 is hη' (trace_chart with
% Fin.succAbove k, trace_deriv); Step 5 is hrel (Algebra.trace_algebraMap) and hres with Module.finrank_pos. About 1,500
% lines in ResidStmt.lean and the files it uses, of which about 1,000 are those of Lemma
% gal-realization-kappa-de-rham-residue-jacobian.
\end{proof}
\leanref{Development/WeilBound/CoreStatements/PadicHodgeComparison/EtaleCohomologyComparison/PerfectoidPullbacks/ConeSectionRestriction/InfinityCoordAlgebraicIndependent.lean}{CrystallineSemistableComparison.Pb.infinityDerivation}
\leanref{Development/WeilBound/CoreStatements/PadicHodgeComparison/EtaleCohomologyComparison/PerfectoidPullbacks/ConeSectionRestriction/InfinityCoordAlgebraicIndependent.lean}{CrystallineSemistableComparison.Pb.infinityDerivation_coord}
\leanref{Development/WeilBound/CoreStatements/PadicHodgeComparison/EtaleCohomologyComparison/PerfectoidPullbacks/ConeSectionRestriction/InfinityCoordAlgebraicIndependent.lean}{CrystallineSemistableComparison.Pb.jacobian_infinityDerivation}
\leanref{Development/WeilBound/CoreStatements/PadicHodgeComparison/EtaleCohomologyComparison/PerfectoidPullbacks/ConeSectionRestriction/TorusMemFldF.lean}{CrystallineSemistableComparison.Pb.trace_deriv}

\begin{lemma}\label{gal-realization-kappa-de-rham-local-exact-charts}
Let $K$ be a field, $X$ a $K$-scheme and $i:X\to\mathbf P^M_K$ a $K$-morphism. For a homogeneous $F\in
K[x_0,\ldots,x_M]$ of positive degree let $K[x_0,\ldots,x_M]_{(F)}$ be the ring of the fractions $H/F^n$ with $H$
homogeneous of degree $n\deg F$, so that $D_+(F)=\operatorname{Spec}K[x_0,\ldots,x_M]_{(F)}$ \cite[Tag 01MB]{Sta}. Let
$X_j=i^{-1}(D_+(x_j))$. For a thickening $T=(B,J,\delta,y)$ of $\mathrm{Th}_{\mathrm{dR}}(X)$ that lies over $X_j$, let
$\ell_{T,j}:K[x_0,\ldots,x_M]_{(x_j)}\to B/J$ be the ring map that corresponds to $i\circ y:\operatorname{Spec}B/J\to
D_+(x_j)$.
\begin{enumerate}
\item If $\varphi:T\to T'$ is a map of $\mathrm{Th}_{\mathrm{dR}}(X)$ and $T$ lies over $X_j$, then $T'$ lies over $X_j$
  and $\ell_{T',j}=\bar\varphi\circ\ell_{T,j}$, where $\bar\varphi:B_T/J_T\to B_{T'}/J_{T'}$ is induced by $\varphi$.
\item Let $T$ lie over $X_j\cap X_{j'}$, and let $F$ be homogeneous of degree $e\geq1$. Then
  $\ell_{T,j}(F/x_j^e)=\ell_{T,j'}(F/x_{j'}^e)\,\ell_{T,j}(x_{j'}/x_j)^e$.
\item Let $F$ and $G$ be homogeneous of degree $e\geq1$, and let $T$ lie over $X_j\cap i^{-1}(D_+(G))$. Then
  $\ell_{T,j}(G/x_j^e)$ is a unit, and $\ell_{T,j}(F/x_j^e)=\ell_{T,j}(G/x_j^e)\,\mathrm{ev}_T(F/G)$. Here
  $F/G\in\Gamma(i^{-1}(D_+(G)),\mathcal O_X)$ is the section $F/G$ of $\mathcal O_{\mathbf P^M_K}$ over $D_+(G)$ carried
  through $i$, and $\mathrm{ev}_T$ is as in Definition~\ref{gal-realization-thickening-level-one-fibre}(2), for $R=K$.
\end{enumerate}
\end{lemma}
\leanref{Development/WeilBound/CoreStatements/PadicHodgeComparison/EtaleCohomologyComparison/PerfectoidPullbacks/ConeSectionRestriction/LocExactC.lean}{CrystallineSemistableComparison.Pb.LX.lj}
\leanref{Development/WeilBound/CoreStatements/PadicHodgeComparison/EtaleCohomologyComparison/PerfectoidPullbacks/ConeSectionRestriction/LocExactC.lean}{CrystallineSemistableComparison.Pb.LX.lj_comp}
\leanref{Development/WeilBound/CoreStatements/PadicHodgeComparison/EtaleCohomologyComparison/PerfectoidPullbacks/ConeSectionRestriction/LocExactA.lean}{CrystallineSemistableComparison.Pb.LX.liftA_yx_change}
\leanref{Development/WeilBound/CoreStatements/PadicHodgeComparison/EtaleCohomologyComparison/PerfectoidPullbacks/ConeSectionRestriction/LocExactA.lean}{CrystallineSemistableComparison.Pb.LX.liftA_yx_isUnit}
\leanref{Development/WeilBound/CoreStatements/PadicHodgeComparison/EtaleCohomologyComparison/PerfectoidPullbacks/ConeSectionRestriction/LocExactA.lean}{CrystallineSemistableComparison.Pb.LX.liftA_yx_quot}

\begin{proof}
(1) As $\varphi$ is a map, $y_{T'}=y_T\circ\operatorname{Spec}\bar\varphi$, so $T'$ lies over $X_j$
(Lemma~\ref{gal-realization-thickening-over}(1)), and $i\circ y_{T'}=(i\circ y_T)\circ\operatorname{Spec}\bar\varphi$,
which corresponds to $\bar\varphi\circ\ell_{T,j}$ \cite[Tag 01I1]{Sta}.

(2) and (3). For homogeneous $F_1,F_2$ of positive degree, $D_+(F_1)\cap D_+(F_2)=D_+(F_1F_2)$, and restriction from
$D_+(F_1)$ to $D_+(F_1F_2)$ is the ring map $K[x_0,\ldots,x_M]_{(F_1)}\to K[x_0,\ldots,x_M]_{(F_1F_2)}$ with $H/F_1^n\mapsto
HF_2^n/(F_1F_2)^n$ \cite[Tag 01M4]{Sta}. As $K[x_0,\ldots,x_M]$ is a domain, these rings are subrings of its field of
fractions, and these maps are the inclusions. In (2), $i\circ y$ factors through $D_+(x_jx_{j'})$; the ring map
$\ell:K[x_0,\ldots,x_M]_{(x_jx_{j'})}\to B/J$ that corresponds to it restricts to $\ell_{T,j}$ and $\ell_{T,j'}$, and
$F/x_j^e=(F/x_{j'}^e)(x_{j'}/x_j)^e$ in the field of fractions. In (3), $i\circ y$ factors through $D_+(x_jG)$, with ring
map $\ell$ on $K[x_0,\ldots,x_M]_{(x_jG)}$. It restricts to $\ell_{T,j}$ on $K[x_0,\ldots,x_M]_{(x_j)}$, and on
$K[x_0,\ldots,x_M]_{(G)}$ to the ring map of $i\circ y$ as a morphism to $D_+(G)$, which sends $F/G$ to
$\mathrm{ev}_T(F/G)$. In $K[x_0,\ldots,x_M]_{(x_jG)}$, $G/x_j^e$ has the inverse $x_j^e/G=x_j^{e+1}/(x_jG)$, and
$F/x_j^e=(G/x_j^e)(F/G)$. Applying $\ell$ gives the claims.
% Lean: lj is liftA (pt T ≫ i) for the chart Proj.awayι of X j; (1) is inU_map_c with lj_comp (liftA_comp); (2) is
% liftA_yx_change and (3) is liftA_yx_isUnit and liftA_yx_quot (with evalAt_comp_appLE), all through liftA_awayMap
% (Proj.SpecMap_awayMap_awayι) and liftA_eval_mul, with the identities checked in Localization.
\end{proof}
\leanref{Development/WeilBound/CoreStatements/PadicHodgeComparison/EtaleCohomologyComparison/PerfectoidPullbacks/ConeSectionRestriction/LocExactA.lean}{CrystallineSemistableComparison.Pb.LX.liftA_awayMap}
\leanref{Development/WeilBound/CoreStatements/PadicHodgeComparison/EtaleCohomologyComparison/PerfectoidPullbacks/ConeSectionRestriction/LocExactA.lean}{CrystallineSemistableComparison.Pb.LX.liftA_eval_mul}
\leanref{Development/WeilBound/CoreStatements/PadicHodgeComparison/EtaleCohomologyComparison/PerfectoidPullbacks/ConeSectionRestriction/LocExactC.lean}{CrystallineSemistableComparison.Pb.LX.inU_map_c}

\begin{lemma}\label{gal-realization-kappa-de-rham-local-exact}
Let $K$ be a field of characteristic $0$, $X$ a $K$-scheme, $i:X\to\mathbf P^M_K$ a closed immersion over $K$,
$\mathcal L$ a line datum on $X$ that is tautological for $i$ (Definition~\ref{gal-realization-period-chern}(2)), and
$y_0,\ldots,y_m$, $s$ and the categories $\mathcal W_k$, $\mathcal W_S$ and $\mathcal W$ as in
Definition~\ref{gal-realization-kappa-de-rham-cone}, with $y_0,\ldots,y_m$ homogeneous of one degree $e\geq1$. Let $\Lambda$
be a lift of $\mathcal L$ along the inclusion of
$\mathcal W$ into $\mathrm{Th}_{\mathrm{dR}}(X)$, and $z_s(\Lambda)$ as in
Lemma~\ref{gal-realization-kappa-de-rham-cocycle}. There are $1$-cochains $b_k$ of $\mathcal W_k$ with values in $J$ such
that $e\,z_s(\Lambda)|_{\mathcal W_k}=db_k$ for every $k$, and such that for all $k,k'$ there is a family $\hat g$ of lifts
of $y_k/y_{k'}$ on $\mathcal W_{\{k,k'\}}$ with $b_k|_{\mathcal W_{\{k,k'\}}}-b_{k'}|_{\mathcal
W_{\{k,k'\}}}=\mathrm{dlog}_s(\hat g)$.
\end{lemma}
\leanref{Development/WeilBound/CoreStatements/PadicHodgeComparison/EtaleCohomologyComparison/PerfectoidPullbacks/ConeSectionRestriction/ConeHyp.lean}{CrystallineSemistableComparison.Pb.LocalCohomologySetup.LocExactP}
\leanref{Development/WeilBound/CoreStatements/PadicHodgeComparison/EtaleCohomologyComparison/PerfectoidPullbacks/ConeSectionRestriction/LocExact.lean}{CrystallineSemistableComparison.Pb.firstChern_local_exact}

\begin{proof}
Write $T=(B_T,J_T,\delta_T,y_T)$ for a thickening, $\bar\varphi:B_T/J_T\to B_{T'}/J_{T'}$ for the map induced by a map
$\varphi:T\to T'$, $f:X_{\mathcal L}\to X$ and $u$ for the morphism and the units of $\mathcal L$, and $\Lambda=(p,w)$. Let
$X_j$ and $\ell_{T,j}$ be as in Lemma~\ref{gal-realization-kappa-de-rham-local-exact-charts}. As $\mathcal L$ is
tautological for $i$, there are $X$-morphisms $\sigma_j:X_j\to X_{\mathcal L}$ with $u(\sigma_ja,\sigma_{j'}b)=
(ia)^*(x_{j'}/x_j)$ for morphisms $a:V\to X_j$ and $b:V\to X_{j'}$ with the same composite to $X$
(Definition~\ref{gal-realization-period-chern}(2)). Every integer $n\geq1$ is invertible in every $B_T$, and $J_T^{[s]}=0$
for $T$ in $\mathcal W$. So Lemma~\ref{gal-realization-kappa-de-rham-log}(3) applies in $B_T$: for units $a,b$ with
$a-1,b-1\in J_T$, $\log_s(ab-1)=\log_s(a-1)+\log_s(b-1)$, so also $\log_s(a^{-1}-1)=-\log_s(a-1)$ and $\log_s(a^e-1)=
e\log_s(a-1)$.

\emph{Step 1: the units $A_T$.} A thickening $T$ of $\mathcal W$ lies in some $\mathcal W_k$, so it lies over some $X_j$;
choose such a $j(T)$. Let $a_T:\operatorname{Spec}B_T/J_T\to X_{j(T)}$ be the morphism through which $y_T$ factors. Then
$f\sigma_{j(T)}a_T=y_T=fp_T$; let $A_T=u(\sigma_{j(T)}a_T,p_T)$, a unit of $B_T/J_T$.

\emph{Step 2: the units $v_k(T)$.} For $T$ in $\mathcal W_k$, $T$ lies over $U_k$
(Lemma~\ref{gal-realization-kappa-de-rham-cone-vanishing}(1)), so over $U_k\cap X_{j(T)}$
(Lemma~\ref{gal-realization-thickening-over}(2)), and
$\mu_k(T)=\ell_{T,j(T)}(y_k/x_{j(T)}^e)\,A_T^{-e}$ is a unit of $B_T/J_T$, by
Lemma~\ref{gal-realization-kappa-de-rham-local-exact-charts}(3) for $G=y_k$. As in the proof of
Lemma~\ref{gal-realization-kappa-de-rham-lift}, it lifts to a unit $v_k(T)$ of $B_T$.

\emph{Step 3: the units $\varepsilon_k(\varphi)$.} For a map $\varphi:T\to T'$ of $\mathcal W_k$ let
$\varepsilon_k(\varphi)=w_\varphi^e\,v_k(T')\,\varphi(v_k(T))^{-1}$. Then $\varepsilon_k(\varphi)-1\in J_{T'}$. Indeed,
let $j=j(T)$ and $j'=j(T')$. By Lemma~\ref{gal-realization-kappa-de-rham-local-exact-charts}(1), $T'$ lies over $X_j$, and
$a_T\circ\operatorname{Spec}\bar\varphi$ is the morphism $a'$ through which $y_{T'}$ factors into $X_j$. In
$B_{T'}/J_{T'}$ the class of $w_\varphi$ is $W=u(p_T\circ\operatorname{Spec}\bar\varphi,p_{T'})$, and by the naturality of
$u$, $\bar\varphi(A_T)=u(\sigma_ja',p_T\circ\operatorname{Spec}\bar\varphi)$. By the cocycle identity of $u$ and the
tautological property,
\[
 \bar\varphi(A_T)\,W=u(\sigma_ja',p_{T'})=u(\sigma_ja',\sigma_{j'}a_{T'})\,A_{T'}=\rho\,A_{T'},\qquad
 \rho=\ell_{T',j}(x_{j'}/x_j).
\]
By Lemma~\ref{gal-realization-kappa-de-rham-local-exact-charts}(1) and~(2),
$\bar\varphi(\ell_{T,j}(y_k/x_j^e))=\ell_{T',j}(y_k/x_j^e)=\ell_{T',j'}(y_k/x_{j'}^e)\,\rho^e$. So
\[
 \bar\varphi(\mu_k(T))=\ell_{T',j'}(y_k/x_{j'}^e)\,\rho^e\,\bar\varphi(A_T)^{-e}
 =\ell_{T',j'}(y_k/x_{j'}^e)\,W^eA_{T'}^{-e}=W^e\mu_k(T'),
\]
and the class of $\varepsilon_k(\varphi)$ is $W^e\mu_k(T')\bar\varphi(\mu_k(T))^{-1}=1$.

\emph{Step 4: the primitives.} For a map $\varphi:T\to T'$ of $\mathcal W_k$ let
$b_k(\varphi)=\log_s(\varepsilon_k(\varphi)-1)$, which lies in $J_{T'}$ by Lemma~\ref{gal-realization-kappa-de-rham-log}(1).
Let $\tau=(T_0\xrightarrow{\varphi}T_1\xrightarrow{\psi}T_2)$ be a chain of $\mathcal W_k$ and $\chi=\psi\varphi$. In the
product $\psi(\varepsilon_k(\varphi))\,\varepsilon_k(\psi)\,\varepsilon_k(\chi)^{-1}$ the factors $v_k(T_m)$ cancel, so
\[
 \psi(\varepsilon_k(\varphi))\,\varepsilon_k(\psi)\,\varepsilon_k(\chi)^{-1}=\bigl(\psi(w_\varphi)\,w_\psi\,
 w_\chi^{-1}\bigr)^e=r(\tau)^e
\]
(Definition~\ref{gal-realization-period-chern}(3)). The three factors on the left minus $1$ lie in $J_{T_2}$, as $\psi$
carries $J_{T_1}$ into $J_{T_2}$, and $r(\tau)-1\in J_{T_2}$. As $\psi$ commutes with $\log_s$
(Lemma~\ref{gal-realization-kappa-de-rham-log}(1)), taking $\log_s$ gives
\[
 b_k(\psi)-b_k(\chi)+\psi(b_k(\varphi))=e\log_s(r(\tau)-1)=e\,z_s(\Lambda)(\tau).
\]
The left side is $(db_k)(\tau)$, so $e\,z_s(\Lambda)|_{\mathcal W_k}=db_k$.

\emph{Step 5: the families of lifts.} Let $k,k'$ be given, and for $T$ in $\mathcal W_{\{k,k'\}}$ let $\hat g_T=
v_k(T)v_{k'}(T)^{-1}$. With $j=j(T)$, its class is $\mu_k(T)\mu_{k'}(T)^{-1}=\ell_{T,j}(y_k/x_j^e)\,
\ell_{T,j}(y_{k'}/x_j^e)^{-1}=\mathrm{ev}_T(y_k/y_{k'})$, by
Lemma~\ref{gal-realization-kappa-de-rham-local-exact-charts}(3) for $F=y_k$ and $G=y_{k'}$, as $T$ lies over $U_{k'}\cap
X_j$, as in Step 2. So $\hat g$ is a family of lifts of $y_k/y_{k'}$ on $\mathcal W_{\{k,k'\}}$. For a map $\varphi:T\to T'$ of
$\mathcal W_{\{k,k'\}}$ the factors $w_\varphi^e$ cancel, and $\varepsilon_k(\varphi)\varepsilon_{k'}(\varphi)^{-1}=\hat
g_{T'}\varphi(\hat g_T)^{-1}$. Taking $\log_s$ gives $b_k(\varphi)-b_{k'}(\varphi)=\log_s(\hat g_{T'}\varphi(\hat
g_T)^{-1}-1)=\mathrm{dlog}_s(\hat g)(\varphi)$.
% Lean: firstChern_local_exact (LocExact.lean). Step 1 is jc, hjc, az (with sO, the sections of IsTautologicalLineAt)
% and AU; Step 2 is μ, μ_isUnit and vv (exists_unit_lift); Step 3 is epsE and eps_mem (ub_nat, ub_mul, sO_u, eR,
% liftA_yx_change, lj_comp and the identity key); Step 4 is bE, bE_mem (mem_thFil_one), and zK_smul_eq_d through epsE_rr
% (the group identity grpE), log_rr (plog_pow, plog_mul, plog_inv, plogK_map) and the coboundary formula hd; Step 5 is the
% family vv k · (vv k')⁻¹ with vv_spec and liftA_yx_quot, and bE_sub = lE_sub through epsE_div. About 750 lines in
% LocExact.lean, LocExactA.lean, LocExactB.lean and LocExactC.lean.
\end{proof}
\leanref{Development/WeilBound/CoreStatements/PadicHodgeComparison/EtaleCohomologyComparison/PerfectoidPullbacks/ConeSectionRestriction/LocExactC.lean}{CrystallineSemistableComparison.Pb.LX.AU}
\leanref{Development/WeilBound/CoreStatements/PadicHodgeComparison/EtaleCohomologyComparison/PerfectoidPullbacks/ConeSectionRestriction/LocExactC.lean}{CrystallineSemistableComparison.Pb.LX.vv}
\leanref{Development/WeilBound/CoreStatements/PadicHodgeComparison/EtaleCohomologyComparison/PerfectoidPullbacks/ConeSectionRestriction/LocExactC.lean}{CrystallineSemistableComparison.Pb.LX.eps_mem}
\leanref{Development/WeilBound/CoreStatements/PadicHodgeComparison/EtaleCohomologyComparison/PerfectoidPullbacks/ConeSectionRestriction/LocExactB.lean}{CrystallineSemistableComparison.Pb.LX.epsE}
\leanref{Development/WeilBound/CoreStatements/PadicHodgeComparison/EtaleCohomologyComparison/PerfectoidPullbacks/ConeSectionRestriction/LocExactB.lean}{CrystallineSemistableComparison.Pb.LX.bE}
\leanref{Development/WeilBound/CoreStatements/PadicHodgeComparison/EtaleCohomologyComparison/PerfectoidPullbacks/ConeSectionRestriction/LocExactB.lean}{CrystallineSemistableComparison.Pb.LX.epsE_rr}
\leanref{Development/WeilBound/CoreStatements/PadicHodgeComparison/EtaleCohomologyComparison/PerfectoidPullbacks/ConeSectionRestriction/LocExactB.lean}{CrystallineSemistableComparison.Pb.LX.zK_smul_eq_d}
\leanref{Development/WeilBound/CoreStatements/PadicHodgeComparison/EtaleCohomologyComparison/PerfectoidPullbacks/ConeSectionRestriction/LocExactB.lean}{CrystallineSemistableComparison.Pb.LX.bE_sub}

\subsection{Transfer and the top power of a Chern class}\label{subsection-gal-top-power}

\begin{lemma}\label{gal-realization-kappa-de-rham-cup}
Let $\Phi$, given by $A$, $J$ and $\gamma$, be a divided power functor on a category $\mathcal C$ over a ring $R$ with a
filtration $(\mathfrak f_r)$ (Definition~\ref{gal-realization-thickening-level-one-graded-filtration}), $m\geq0$, and
$\ell_1,\ldots,\ell_m$ $1$-cocycles of $\mathcal C$ with values in $A$ whose values lie in $\mathfrak f_1$. Then
$\ell_1\cup\cdots\cup\ell_m$, the constant cochain $1$ for $m=0$, is a cocycle with values in $\mathfrak f_m$. So
$[\ell_1\cup\cdots\cup\ell_m]_m\in H^m(\mathcal C,\mathrm{gr}_m)$ is defined.
\end{lemma}
\leanref{Development/WeilBound/CoreStatements/PadicHodgeComparison/EtaleCohomologyComparison/PerfectoidPullbacks/ConeSectionRestriction/RepDlogLifts.lean}{CrystallineSemistableComparison.Pb.rep_cupList}
\leanref{Development/AutomorphicSpectrum/Proofs/BlockDecomposition/ExistsCupListRepresentativeStatement.lean}{GlobalTheory.T.exists_cup_list_representative}
\leanref{Development/AutomorphicSpectrum/Proofs/BlockDecomposition/ExistsCupListRepresentative.lean}{GlobalTheory.CompanionTheory.ConeRepCup.exists_rep}

\begin{proof}
By induction on $m$, with $\ell_1\cup\cdots\cup\ell_{m+1}=(\ell_1\cup\cdots\cup\ell_m)\cup\ell_{m+1}$. For $m=0$, the
coboundary of the constant cochain $1$ is $1-1=0$, and $\mathfrak f_0=A$. For the step, the Leibniz rule $d(a\cup
b)=da\cup b+(-1)^ia\cup db$ for an $i$-cochain $a$ shows that the cup product of two cocycles is a cocycle, and its
values lie in $\mathfrak f_m\mathfrak f_1\subset\mathfrak f_{m+1}$. The coboundary $0$ has values in $\mathfrak
f_{m+1}$, so $[\ell_1\cup\cdots\cup\ell_m]_m$ is defined.
% Lean: ConeRepCup.exists_rep, through cupList_cocycle (P.leibniz), cupList_mem (cup_mem) and ConeLocal.rep_exists.
\end{proof}
\leanref{Development/AutomorphicSpectrum/Proofs/BlockDecomposition/ExistsCupListRepresentative.lean}{GlobalTheory.CompanionTheory.ConeRepCup.cupList_cocycle}
\leanref{Development/AutomorphicSpectrum/Proofs/BlockDecomposition/ExistsCupListRepresentative.lean}{GlobalTheory.CompanionTheory.ConeRepCup.cupList_mem}

\begin{lemma}\label{gal-realization-kappa-de-rham-cech-stair}
Let $K$ be a double complex of abelian groups $K^{p,q}$, $p,q\geq0$, with commuting differentials $d_1:K^{p,q}\to
K^{p+1,q}$ and $d_2:K^{p,q}\to K^{p,q+1}$, and $\mathrm{Tot}(K)$ its total complex, with differential $D=d_1+(-1)^pd_2$ on
$K^{p,q}$ (Lemma~\ref{gal-realization-thickening-zariski-double}). For $x\in\mathrm{Tot}^m(K)$ write $x_{p,q}$ for its
component in $K^{p,q}$, and for $c\in K^{p,q}$ write $[c]$ for $c$ viewed in $\mathrm{Tot}^{p+q}(K)$. Let
$L_r^{p,q}\subset K^{p,q}$, $r\in\mathbf Z$, be subgroups with $L_0^{p,q}=K^{p,q}$, $L_{r'}^{p,q}\subset L_r^{p,q}$ for
$r\leq r'$, $d_1(L_r^{p,q})\subset L_r^{p+1,q}$ and $d_2(L_r^{p,q})\subset L_r^{p,q+1}$. Let $n\geq0$, and suppose that
$K^{p,q}=0$ for $p>n+1$ and that the following holds for all $r,p,q\geq0$ with $q\neq r$:
\begin{enumerate}
\item[(E)] if $x\in L_r^{p,q}$ and $d_2x\in L_{r+1}^{p,q+1}$, then $x\in L_{r+1}^{p,q}$ if $q=0$, and $x-d_2v\in
  L_{r+1}^{p,q}$ for some $v\in L_r^{p,q-1}$ if $q\geq1$.
\end{enumerate}
Let $P\in L_{n+1}^{n+1,n+1}$ and $y\in\mathrm{Tot}^{2n+1}(K)$ with $Dy=[P]$. Then there are $\omega\in L_{n+1}^{n,n+1}$
with $d_2\omega\in L_{n+2}^{n,n+2}$ and $f\in L_n^{n+1,n}$ such that $P=d_1\omega+(-1)^{n+1}d_2f$.
\end{lemma}
\leanref{Development/AutomorphicSpectrum/Proofs/BlockDecomposition/DvMemOf.lean}{GlobalTheory.CompanionTheory.ConeCech.DFil}
\leanref{Development/AutomorphicSpectrum/Proofs/BlockDecomposition/DvMemOf.lean}{GlobalTheory.CompanionTheory.ConeCech.stair}

\begin{proof}
Say that $x\in\mathrm{Tot}^m(K)$ lies in $L_r$ if all its components do.

\emph{Step 1: one column at a time.} Let $r,m,e\geq0$ and $y\in\mathrm{Tot}^{m+1}(K)$ in $L_r$ such that
$(Dy)_{p,m+2-p}\in L_{r+1}$ for $p<e$, and such that (E) holds for $r$ and $(p,m+1-p)$ for every $p<e$. Then there is
$w\in\mathrm{Tot}^m(K)$ with $w_{p,q}=0$ for $p\geq e$, such that $y-Dw$ lies in $L_r$ and $(y-Dw)_{p,q}\in L_{r+1}$ for
$p<e$. This is shown by induction on $e$; for $e=0$ take $w=0$. Let $w$ be given for $e$, and $y'=y-Dw$, so that
$Dy'=Dy$. If $e>m+1$, $y'$ has no component in column $e$, and $w$ works for $e+1$. Otherwise let $q=m+1-e$. As
$(Dy')_{e,q+1}=d_1y'_{e-1,q+1}+(-1)^ed_2y'_{e,q}$, where the first term is absent for $e=0$ and lies in $L_{r+1}$
otherwise, $d_2y'_{e,q}\in L_{r+1}$. If $q=0$, (E) gives $y'_{e,0}\in L_{r+1}$, and $w$ works for $e+1$. If $q\geq1$,
(E) gives $v\in L_r^{e,q-1}$ with $y'_{e,q}-d_2v\in L_{r+1}$. Replace $w$ by $w+(-1)^e[v]$. As
$D((-1)^e[v])=(-1)^e[d_1v]+[d_2v]$, the new $y-Dw$ has the same components as $y'$ in the columns $p<e$, the component
$y'_{e,q}-d_2v\in L_{r+1}$ in column $e$, and a component in column $e+1$ changed by an element of $L_r$.

\emph{Step 2: the filtration of $y$.} For $0\leq r\leq n$ there is $y\in\mathrm{Tot}^{2n+1}(K)$ in $L_r$ with $Dy=[P]$.
For $r=0$ take the given $y$, as $L_0=K$. For $r+1\leq n$, apply Step 1 to $y$ for $r$, with $m=2n$ and $e=n+2$:
$Dy=[P]$ lies in $L_{r+1}$, as $P\in L_{n+1}\subset L_{r+1}$, and for $p\leq n+1$, $q=2n+1-p\geq n>r$. Then $y-Dw$ lies in
$L_{r+1}$, as $K^{p,q}=0$ for $p>n+1$, and $D(y-Dw)=[P]$.

\emph{Step 3.} Let $y$ be as in Step 2 for $r=n$. Step 1 for $r=n$, $m=2n$ and $e=n+1$ applies, as $q=2n+1-p\geq n+1$
for $p\leq n$; it gives $y_1=y-Dw$ in $L_n$ with $Dy_1=[P]$, whose components in the columns $p\leq n$ lie in
$L_{n+1}$. Let $f=(y_1)_{n+1,n}\in L_n^{n+1,n}$ and $t=y_1-[f]$. As $d_1f\in K^{n+2,n}=0$, $Dt=[P-(-1)^{n+1}d_2f]$, and $t$
lies in $L_{n+1}$: its components in the columns $p\leq n$ are those of $y_1$, and it has none in the columns $p\geq n+1$.
Step 1 for $r=n+1$, $m=2n$ and $e=n$ applies, as $Dt$ vanishes in the columns $p<n$ and $q=2n+1-p\geq n+2$ for $p<n$; it
gives $w'$ with $w'_{p,q}=0$ for $p\geq n$, such that $t'=t-Dw'$ lies in $L_{n+1}$ and in $L_{n+2}$ in the columns $p<n$.
Then $t'_{n+1,n}=t_{n+1,n}=0$, as $(Dw')_{n+1,n}=d_1w'_{n,n}+(-1)^{n+1}d_2w'_{n+1,n-1}=0$. Let $\omega=t'_{n,n+1}\in
L_{n+1}$. As $Dt'=Dt$ vanishes at $(n,n+2)$, $(-1)^nd_2\omega=-d_1t'_{n-1,n+2}$ (or $0$ if $n=0$), which lies in
$L_{n+2}$. At $(n+1,n+1)$, $P-(-1)^{n+1}d_2f=(Dt')_{n+1,n+1}=d_1\omega+(-1)^{n+1}d_2t'_{n+1,n}=d_1\omega$.
% Lean: ConeCech.stair, for the filtration DFil (L, mem_h, mem_v, anti) with Ex as (E). Step 1 is kill (with dv_mem_of
% for d_2 y'_{e,q} ∈ L_{r+1} and the correction (-1)^e sgl e q v); Step 2 is stair_up; Step 3 is the rest of stair (y1,
% t = y1 - sgl (n+1) n (y1 (n+1) n), w2, ht0 and the final identity h).
\end{proof}
\leanref{Development/AutomorphicSpectrum/Proofs/BlockDecomposition/DvMemOf.lean}{GlobalTheory.CompanionTheory.ConeCech.kill}
\leanref{Development/AutomorphicSpectrum/Proofs/BlockDecomposition/DvMemOf.lean}{GlobalTheory.CompanionTheory.ConeCech.stair_up}
\leanref{Development/AutomorphicSpectrum/Proofs/BlockDecomposition/DvMemOf.lean}{GlobalTheory.CompanionTheory.ConeCech.dv_mem_of}

\begin{lemma}\label{gal-realization-kappa-de-rham-cech-zigzag}
Let $K$, $\mathrm{Tot}(K)$, $D$ and $[c]$ be as in Lemma~\ref{gal-realization-kappa-de-rham-cech-stair}, and let
$B:K^{p,q}\to K^{p+1,q+1}$ be homomorphisms with $d_1B=Bd_1$ and $d_2B=Bd_2$. Let $a_k\in K^{0,2k}$, $g_k\in K^{0,2k+1}$
and $P_k\in K^{k,k}$, $k\geq0$, with $d_2g_k=a_{k+1}$, $d_1g_k=-Ba_k$, $P_0=a_0$ and $P_{k+1}=(-1)^kBP_k$. Then for every
$k$ there is $U_k\in\mathrm{Tot}^{2k-1}(K)$, with $U_0=0$, such that $DU_k=[a_k]-[P_k]$.
\end{lemma}
\leanref{Development/AutomorphicSpectrum/Proofs/BlockDecomposition/ConeZigzagCompat.lean}{GlobalTheory.CompanionTheory.ConeZig.cone_zigzag_compat}

\begin{proof}
For $x\in\mathrm{Tot}^m(K)$ let $S(x)\in\mathrm{Tot}^{m+2}(K)$ have the components $S(x)_{p+1,q+1}=(-1)^qBx_{p,q}$, and
$0$ in row $0$ and column $0$. Then $S[c]=(-1)^q[Bc]$ for $c\in K^{p,q}$, and $DS(x)=S(Dx)$: at $(p+1,q+1)$, as $B$
commutes with $d_1$ and $d_2$,
\[
 d_1S(x)_{p,q+1}+(-1)^{p+1}d_2S(x)_{p+1,q}=(-1)^qB\bigl(d_1x_{p-1,q}+(-1)^pd_2x_{p,q-1}\bigr)=(-1)^qB(Dx)_{p,q},
\]
with the terms for $p=0$ or $q=0$ absent, and both sides vanish in row $0$ and column $0$. Let $U_0=0$, which works as
$P_0=a_0$, and $U_{k+1}=S(U_k)+[g_k]$. Then
\[
 DU_{k+1}=S([a_k]-[P_k])+[d_1g_k]+[d_2g_k]=[Ba_k]-(-1)^k[BP_k]-[Ba_k]+[a_{k+1}]=[a_{k+1}]-[P_{k+1}].
\]
% Lean: ConeZig.cone_zigzag_compat, with S = shift (shift_sgl, shift_sub and D_shift); zig_stair combines it with
% Lemma gal-realization-kappa-de-rham-cech-stair.
\end{proof}
\leanref{Development/AutomorphicSpectrum/Proofs/BlockDecomposition/ConeZigzagCompat.lean}{GlobalTheory.CompanionTheory.ConeZig.shift}
\leanref{Development/AutomorphicSpectrum/Proofs/BlockDecomposition/ConeZigzagCompat.lean}{GlobalTheory.CompanionTheory.ConeZig.D_shift}
\leanref{Development/AutomorphicSpectrum/Proofs/BlockDecomposition/ConeZigzagCompat.lean}{GlobalTheory.CompanionTheory.ConeZig.zig_stair}

\begin{lemma}\label{gal-realization-kappa-de-rham-cech}
Let $\Phi$, given by $A$, $J$ and $\gamma$, be a divided power functor on a category $\mathcal C$ over a ring $R$ with a
filtration $(\mathfrak f_r)$ (Definition~\ref{gal-realization-thickening-level-one-graded-filtration}), and let $s,n\geq0$.
Let $\mathcal W_0,\ldots,\mathcal W_{n+1}$ be full subcategories of $\mathcal C$, $\top=\{0,\ldots,n+1\}$, and for a
nonempty $S\subset\top$ let $\mathcal W_S$ be the full subcategory of the objects in $\mathcal W_\alpha$ for every
$\alpha\in S$, and $\mathcal W$ that of the objects in some $\mathcal W_\alpha$. Suppose:
\begin{enumerate}
\item[(a)] $\mathfrak f_s(T)=0$ for every $T$ in $\mathcal W$;
\item[(b)] if $T\to T'$ is a map of $\mathcal C$ with $T$ in $\mathcal W_\alpha$ and $T'$ in $\mathcal W$, then $T'$ is in
  $\mathcal W_\alpha$;
\item[(c)] for every nonempty $S\subset\top$, $r\in\mathbf Z$ and $j\geq0$ with $j\neq r$ or $r>n+1$, $H^j(\mathcal
  W_S,\mathrm{gr}_r)=0$.
\end{enumerate}
Let $z$ be a $2$-cocycle of $\mathcal W$ with values in $A$ whose values lie in $\mathfrak f_1$, such that the cup power
$z^{\cup(n+1)}$ is a coboundary, where $z^{\cup0}=1$ and $z^{\cup(k+1)}=z^{\cup k}\cup z$. Let $b_\alpha$ be $1$-cochains of
$\mathcal W_\alpha$ with values in $\mathfrak f_1$ and $db_\alpha=z|_{\mathcal W_\alpha}$, let $\ell_k=(b_k-b_{k+1})|_{\mathcal
W_\top}$ for $0\leq k\leq n$, and $c=[\ell_0\cup\cdots\cup\ell_n]_{n+1}\in H^{n+1}(\mathcal W_\top,\mathrm{gr}_{n+1})$
(Lemma~\ref{gal-realization-kappa-de-rham-cup}). Then there are classes $\omega_\alpha\in H^{n+1}(\mathcal
W_{\top\smallsetminus\{\alpha\}},\mathrm{gr}_{n+1})$, $\alpha\in\top$, and an $n$-cochain $f$ of $\mathcal W_\top$ with
values in $\mathfrak f_n$ whose coboundary has values in $\mathfrak f_{n+1}$, such that
$[df]_{n+1}=c-\sum_\alpha\omega_\alpha|_{\mathcal W_\top}$.
\end{lemma}
\leanref{Development/WeilBound/CoreStatements/PadicHodgeComparison/EtaleCohomologyComparison/PerfectoidPullbacks/ConeSectionRestriction/CechRel.lean}{CrystallineSemistableComparison.Pb.cech_top_rel}
\leanref{Development/AutomorphicSpectrum/Proofs/BlockDecomposition/GradedCoverVanishingInductionStatement.lean}{GlobalTheory.T.graded_cover_vanishing_induction}

\begin{proof}
Hypotheses (a) and (b), the case $r>n+1$ of (c) and the condition on the values of $z$ are not used. A cochain of a full
subcategory \emph{has values in} $\mathfrak f_r$ if all its values do; restriction to a smaller full subcategory is
evaluation at its chains, and commutes with coboundaries and cup products.

\emph{Step 1: the \v{C}ech double complex.} For $p\geq0$ let $I_p$ be the set of strictly increasing maps
$i:\{0,\ldots,p\}\to\top$, $\mathcal W_i=\mathcal W_{i(\{0,\ldots,p\})}$, and for $i\in I_{p+1}$ and $0\leq k\leq p+1$
let $i\circ\partial_k\in I_p$ be $i$ with $i(k)$ omitted; then $\mathcal W_i\subset\mathcal W_{i\circ\partial_k}$. Let
$K^{p,q}=\prod_{i\in I_p}C^q(\mathcal W_i,A)$, with $d_2$ the product of the coboundaries and $(d_1x)_i=\sum_{k=0}^{p+1}
(-1)^kx_{i\circ\partial_k}|_{\mathcal W_i}$. As in the proof of Lemma~\ref{gal-realization-thickening-level-one-cech},
$K$ is a double complex: $d_1d_1=0$ by the identities $\partial_k\partial_l=\partial_{l+1}\partial_k$ for $k\leq l$ of the
omissions, and $d_1$ commutes with $d_2$, as restriction commutes with coboundaries. Let $L_r^{p,q}$ be the set of the $x$ whose components have values in $\mathfrak f_r$. As the
$\mathfrak f_r$ are families of ideals of $\Phi$ that decrease in $r$, with $\mathfrak f_0=A$, the hypotheses of
Lemma~\ref{gal-realization-kappa-de-rham-cech-stair} on the $L_r$ hold. $K^{p,q}=0$ for $p>n+1$, as $I_p$ is then empty.
Condition (E) holds: let $q\neq r$, and let $x\in L_r^{p,q}$ with $d_2x\in L_{r+1}$. For $i\in I_p$ the cochain
$\bar x_i$ of the classes of the values of $x_i$ is a cocycle with values in $\mathrm{gr}_r$, whose class $[x_i]_r$
vanishes by (c), as $\mathcal W_i=\mathcal W_S$ with $S=i(\{0,\ldots,p\})$ nonempty. If $q=0$, then $\bar x_i=0$, as
there are no nonzero $0$-coboundaries, so $x_i$ has values in $\mathfrak f_{r+1}$. If $q\geq1$, $\bar x_i$ is the
coboundary of a $(q-1)$-cochain with values in $\mathrm{gr}_r$, and lifting its values to $\mathfrak f_r$ gives $v_i$
with $x_i-dv_i$ having values in $\mathfrak f_{r+1}$.

\emph{Step 2: the operator $B$.} For $i\in I_p$ and $a,c\in\{0,\ldots,p\}$ let $\ell_i(a,c)=b_{i(a)}|_{\mathcal
W_i}-b_{i(c)}|_{\mathcal W_i}$. It has values in $\mathfrak f_1$, and it is a cocycle, as $d\ell_i(a,c)=z|_{\mathcal
W_i}-z|_{\mathcal W_i}$. Let $B:K^{p,q}\to K^{p+1,q+1}$ send $x$ to $(Bx)_i=x_{i\circ\partial_{p+1}}|_{\mathcal W_i}\cup
\ell_i(p,p+1)$. Then $d_2B=Bd_2$, by the Leibniz rule $d(x\cup\ell)=dx\cup\ell+(-1)^qx\cup d\ell$. And $d_1B=Bd_1$: for
$i\in I_{p+2}$ and $k\leq p$, the term $k$ of $(d_1Bx)_i$ is $(-1)^kx_{i\circ\partial_{p+2}\circ\partial_k}|\cup
\ell_i(p+1,p+2)$, the term $k$ of $(Bd_1x)_i$; with $y=x_{i\circ\partial_{p+2}\circ\partial_{p+1}}|$, the terms $p+1$
and $p+2$ of $(d_1Bx)_i$ are $(-1)^{p+1}y\cup\ell_i(p,p+2)$ and $(-1)^py\cup\ell_i(p,p+1)$, whose sum is
$(-1)^{p+1}y\cup\ell_i(p+1,p+2)$, the term $p+1$ of $(Bd_1x)_i$, as $\ell_i(p,p+2)=\ell_i(p,p+1)+\ell_i(p+1,p+2)$.

\emph{Step 3: the zigzag.} Let $(a_k)_\alpha=z^{\cup k}|_{\mathcal W_\alpha}$ and $(g_k)_\alpha=z^{\cup
k}|_{\mathcal W_\alpha}\cup b_\alpha$ for $\alpha\in\top=I_0$, let $P_0=a_0$ and $P_{k+1}=(-1)^kBP_k$. By the Leibniz
rule $z^{\cup k}$ is a cocycle, and $d_2g_k=a_{k+1}$, as $db_\alpha=z|_{\mathcal W_\alpha}$. For $i=(\alpha_0,\alpha_1)\in
I_1$, $(d_1g_k)_i=z^{\cup k}|\cup(b_{\alpha_1}-b_{\alpha_0})|=-(Ba_k)_i$. By
Lemma~\ref{gal-realization-kappa-de-rham-cech-zigzag} there is $U\in\mathrm{Tot}^{2n+1}(K)$ with
$DU=[a_{n+1}]-[P_{n+1}]$. Let $u$ be a $(2n+1)$-cochain of $\mathcal W$ with $du=z^{\cup(n+1)}$, and $v_\alpha=u|_{\mathcal
W_\alpha}$. Then $d_2v=a_{n+1}$ and $d_1v=0$, so $y=[v]-U$ has $Dy=[P_{n+1}]$.

\emph{Step 4: the components of $P_k$.} By induction on $k$, $P_k\in L_k$, as $\mathfrak f_k\mathfrak f_1\subset\mathfrak
f_{k+1}$, and there is $e_k\in\{1,-1\}$ with $(P_k)_i=e_k\,\ell_i(0,1)\cup\ell_i(1,2)\cup\cdots\cup\ell_i(k-1,k)$ for
every $i\in I_k$: for $k=0$ this is $(P_0)_\alpha=1$ with $e_0=1$, and $(BP_k)_i=e_k(\ell_{i\circ\partial_{k+1}}(0,1)
\cup\cdots\cup\ell_{i\circ\partial_{k+1}}(k-1,k))|_{\mathcal W_i}\cup\ell_i(k,k+1)$, where
$\ell_{i\circ\partial_{k+1}}(a,a+1)|_{\mathcal W_i}=\ell_i(a,a+1)$; so $e_{k+1}=(-1)^ke_k$.

\emph{Step 5: the conclusion.} By Lemma~\ref{gal-realization-kappa-de-rham-cech-stair}, applied to $P_{n+1}\in L_{n+1}$ and
$y$, there are $\omega'\in L_{n+1}^{n,n+1}$ with $d_2\omega'\in L_{n+2}$ and $f'\in L_n^{n+1,n}$ with
$P_{n+1}=d_1\omega'+(-1)^{n+1}d_2f'$. Take the component at the identity $\iota$ of $\top$, the only element of $I_{n+1}$,
with $\mathcal W_\iota=\mathcal W_\top$; $\iota\circ\partial_\alpha$ has image $\top\smallsetminus\{\alpha\}$, and
$\ell_\iota(k,k+1)=\ell_k$. With $e=e_{n+1}$, multiplying by $e$ gives
\[
 \ell_0\cup\cdots\cup\ell_n-\sum_{\alpha\in\top}w_\alpha|_{\mathcal W_\top}=df,\qquad
 w_\alpha=e(-1)^\alpha\omega'_{\iota\circ\partial_\alpha},\quad f=e(-1)^{n+1}f'_\iota.
\]
Here $f$ has values in $\mathfrak f_n$, and each $w_\alpha$ is a cochain of $\mathcal W_{\top\smallsetminus\{\alpha\}}$
with values in $\mathfrak f_{n+1}$ and coboundary with values in $\mathfrak f_{n+2}$; let $\omega_\alpha=[w_\alpha]_{n+1}$.
The left side has values in $\mathfrak f_{n+1}$ (Lemma~\ref{gal-realization-kappa-de-rham-cup}), so $df$ does. As
$[\cdot]_{n+1}$ is additive (Lemma~\ref{gal-realization-thickening-level-one-graded-forms}(2)) and commutes with
restriction, $[df]_{n+1}=c-\sum_\alpha\omega_\alpha|_{\mathcal W_\top}$.
% Lean: GlobalTheory.T.graded_cover_vanishing_induction = ConeZig.cech_top_cup_power. Step 1 is cosK and KK (the Čech
% double complex over Idx (p + 1) (n + 2), the strictly increasing maps, with fc k i = i ∘ succAbove k), dfil, KV, and KEx
% through ex_zero and ex_succ (with capT_eq); Step 2 is LL, BB, BB_dv and BB_dh (with LL_add and neg_one_smul_cancel);
% Step 3 is aa, gg, PP, cupPow_d, gg_dv, gg_dh, and v with hv1 and hv2, fed into zig_stair; Step 4 is PP_mem and PP_eq;
% Step 5 is alg_fin and fin_step (gC, gC_d, rep_fin). About 1,000 lines in GradedCoverVanishingInduction.lean,
% ConeZigzagCompat.lean, DvMemOf.lean, NegOneSmulCancel.lean and Mk2EqZeroIffMemS.lean.
\end{proof}
\leanref{Development/AutomorphicSpectrum/Proofs/BlockDecomposition/GradedCoverVanishingInduction.lean}{GlobalTheory.CompanionTheory.ConeZig.cech_top_cup_power}
\leanref{Development/AutomorphicSpectrum/Proofs/BlockDecomposition/NegOneSmulCancel.lean}{GlobalTheory.CompanionTheory.ConeZig.cosK}
\leanref{Development/AutomorphicSpectrum/Proofs/BlockDecomposition/NegOneSmulCancel.lean}{GlobalTheory.CompanionTheory.ConeZig.dfil}
\leanref{Development/AutomorphicSpectrum/Proofs/BlockDecomposition/GradedCoverVanishingInduction.lean}{GlobalTheory.CompanionTheory.ConeZig.KEx}
\leanref{Development/AutomorphicSpectrum/Proofs/BlockDecomposition/Mk2EqZeroIffMemS.lean}{GlobalTheory.CompanionTheory.ConeZig.ex_succ}
\leanref{Development/AutomorphicSpectrum/Proofs/BlockDecomposition/NegOneSmulCancel.lean}{GlobalTheory.CompanionTheory.ConeZig.BB}
\leanref{Development/AutomorphicSpectrum/Proofs/BlockDecomposition/NegOneSmulCancel.lean}{GlobalTheory.CompanionTheory.ConeZig.BB_dh}
\leanref{Development/AutomorphicSpectrum/Proofs/BlockDecomposition/NegOneSmulCancel.lean}{GlobalTheory.CompanionTheory.ConeZig.BB_dv}
\leanref{Development/AutomorphicSpectrum/Proofs/BlockDecomposition/NegOneSmulCancel.lean}{GlobalTheory.CompanionTheory.ConeZig.gg_dv}
\leanref{Development/AutomorphicSpectrum/Proofs/BlockDecomposition/NegOneSmulCancel.lean}{GlobalTheory.CompanionTheory.ConeZig.gg_dh}
\leanref{Development/AutomorphicSpectrum/Proofs/BlockDecomposition/GradedCoverVanishingInduction.lean}{GlobalTheory.CompanionTheory.ConeZig.PP_mem}
\leanref{Development/AutomorphicSpectrum/Proofs/BlockDecomposition/GradedCoverVanishingInduction.lean}{GlobalTheory.CompanionTheory.ConeZig.PP_eq}
\leanref{Development/AutomorphicSpectrum/Proofs/BlockDecomposition/GradedCoverVanishingInduction.lean}{GlobalTheory.CompanionTheory.ConeZig.fin_step}

\begin{lemma}\label{gal-realization-kappa-de-rham-transfer-lifts}
Let $K$ be a field of characteristic $0$, $X$ a $K$-scheme, $s\geq1$ and $\mathcal Q$ a full subcategory of
$\mathrm{Th}_{\mathrm{dR}}(X)$ with $J_T^{[s]}=0$ for every $T$ in $\mathcal Q$. Call a family $\hat g=(\hat g_T)$ of units
$\hat g_T\in B_T$, for the thickenings $T$ of $\mathcal Q$, \emph{compatible} if $\hat g_{T'}-\varphi(\hat g_T)\in J_{T'}$ for
every map $\varphi:T\to T'$ of $\mathcal Q$, and let $\mathrm{dlog}_s(\hat g)$ be as in
Definition~\ref{gal-realization-kappa-de-rham-cone}(3).
\begin{enumerate}
\item For a compatible family $\hat g$, $\mathrm{dlog}_s(\hat g)$ is a $1$-cocycle of $\mathcal Q$ with values in $J$.
\item Let $m\geq0$, $\hat g_1,\ldots,\hat g_m$ compatible families, and $\hat g'_1,\ldots,\hat g'_m$ families of elements
  $\hat g'_k(T)\in B_T$ with $\hat g_k(T)-\hat g'_k(T)\in J_T$ for all $k$ and $T$; they are units, as $J_T$ is nilpotent.
  Then $z'=\mathrm{dlog}_s(\hat g'_1)\cup\cdots\cup\mathrm{dlog}_s(\hat g'_m)$ has values in $J^{[m]}$ and coboundary with
  values in $J^{[m+1]}$, and in $H^m(\mathcal Q,\mathrm{gr}_m)$ (Lemma~\ref{gal-realization-kappa-de-rham-cup})
  \[
   [\mathrm{dlog}_s(\hat g_1)\cup\cdots\cup\mathrm{dlog}_s(\hat g_m)]_m=[z']_m.
  \]
\end{enumerate}
\end{lemma}
\leanref{Development/AutomorphicSpectrum/Proofs/BlockDecomposition/UnitLogRepresentativeCompat.lean}{GlobalTheory.CompanionTheory.RepLifts.d_dlogC}
\leanref{Development/AutomorphicSpectrum/Proofs/BlockDecomposition/UnitLogRepresentativeCompat.lean}{GlobalTheory.CompanionTheory.RepLifts.main}

\begin{proof}
(1) For $\varphi:T\to T'$, $\hat g_{T'}\varphi(\hat g_T)^{-1}-1=(\hat g_{T'}-\varphi(\hat g_T))\varphi(\hat g_T)^{-1}\in
J_{T'}$, so the values of $\mathrm{dlog}_s(\hat g)$ lie in $J$ (Lemma~\ref{gal-realization-kappa-de-rham-log}(1)). Let
$T\xrightarrow{\varphi}T'\xrightarrow{\psi}T''$ be a chain of $\mathcal Q$, and $r(\varphi)=\hat
g_{T'}\varphi(\hat g_T)^{-1}$. Then $r(\psi\varphi)=\psi(r(\varphi))\,r(\psi)$. As every integer $n\geq1$ is invertible in
$B_{T''}$ and $J_{T''}^{[s]}=0$, Lemma~\ref{gal-realization-kappa-de-rham-log}(3) and~(1) give
$\log_s(r(\psi\varphi)-1)=\psi(\log_s(r(\varphi)-1))+\log_s(r(\psi)-1)$, which says that the coboundary of
$\mathrm{dlog}_s(\hat g)$ vanishes at the chain.

(2) The families $\mathcal O^{[r]}$ form a filtration of $\Phi_{\mathcal Q}$, by the definition of the ideals $J^{[r]}$
(Definition~\ref{gal-realization-thickening-nilpotent}(3)), with $\mathcal O^{[1]}=J$. As in (1), $\log_s(ab-1)=
\log_s(a-1)+\log_s(b-1)$ for units $a,b$ of a $B_T$ with $a-1,b-1\in J_T$; so also $\log_s(a^{-1}-1)=-\log_s(a-1)$, as
$a^{-1}-1=-a^{-1}(a-1)\in J_T$. Let $\ell_k=\mathrm{dlog}_s(\hat g_k)$, a $1$-cocycle with values in $J$ by (1).

\emph{Step 1: the units $e_k(T)$.} Let $x=\hat g_k(T)^{-1}(\hat g'_k(T)-\hat g_k(T))\in J_T$. As $x$ is nilpotent,
$e_k(T)=1+x$ is a unit, and $\hat g'_k(T)=\hat g_k(T)\,e_k(T)$. Let $a_k(T)=\log_s(e_k(T)-1)$, which lies in $J_T$
(Lemma~\ref{gal-realization-kappa-de-rham-log}(1)), and $\alpha_k$ the $0$-cochain with $\alpha_k(T)=a_k(T)$.

\emph{Step 2: $\mathrm{dlog}_s(\hat g'_k)=\ell_k+d\alpha_k$.} Let $\varphi:T\to T'$ be a map of $\mathcal Q$ and
$r(\varphi)=\hat g_k(T')\varphi(\hat g_k(T))^{-1}$, so that $r(\varphi)-1\in J_{T'}$, as in (1). Then
\[
 \hat g'_k(T')\,\varphi(\hat g'_k(T))^{-1}=r(\varphi)\,e_k(T')\,\varphi(e_k(T))^{-1},
\]
a product of units congruent to $1$ modulo $J_{T'}$, as $\varphi(e_k(T))-1=\varphi(e_k(T)-1)$. Taking $\log_s$, and as
$\varphi$ commutes with $\log_s$ (Lemma~\ref{gal-realization-kappa-de-rham-log}(1)), $\mathrm{dlog}_s(\hat
g'_k)(\varphi)=\ell_k(\varphi)+a_k(T')-\varphi(a_k(T))=\ell_k(\varphi)+(d\alpha_k)(\varphi)$.

\emph{Step 3: the cup products.} Let $\ell'_k=\ell_k+d\alpha_k$, a $1$-cocycle with values in $J$, and for $0\leq j\leq m$
let $L_j=\ell_1\cup\cdots\cup\ell_j$ and $L'_j=\ell'_1\cup\cdots\cup\ell'_j$, the constant cochain $1$ for $j=0$. By
Lemma~\ref{gal-realization-kappa-de-rham-cup}, $L_j$ is a cocycle with values in $J^{[j]}$. Then $L'_j=L_j$ for $j=0$, and
$L'_j-L_j=dH_j$ for a $(j-1)$-cochain $H_j$ with values in $J^{[j]}$ for $j\geq1$. Indeed, by induction on $j$,
\[
 L'_{j+1}-L_{j+1}=(L'_j-L_j)\cup\ell'_{j+1}+L_j\cup d\alpha_{j+1}.
\]
By the Leibniz rule $d(a\cup b)=da\cup b+(-1)^ia\cup db$ for an $i$-cochain $a$: the first term is $0$ for $j=0$ and
$d(H_j\cup\ell'_{j+1})$ for $j\geq1$, as $d\ell'_{j+1}=0$; and the second is $d((-1)^jL_j\cup\alpha_{j+1})$, as $dL_j=0$.
Both cochains in parentheses have values in $J^{[j]}J\subset J^{[j+1]}$.

\emph{Step 4: the classes.} By Step 2, $z'=L'_m$. By Step 3 it is $L_m$, or $L_m+dH_m$ with $H_m$ with values in
$J^{[m]}$, so its values lie in $J^{[m]}$, and $dz'=dL_m=0$. For $m=0$, $z'=L_0$. For $m\geq1$, by
Lemma~\ref{gal-realization-thickening-level-one-graded-forms}(2), $[z']_m=[L_m]_m+[dH_m]_m$, and $[dH_m]_m$ is the class
of the coboundary of the cochain of the classes of the values of $H_m$ in $\mathrm{gr}_m$, so it is $0$.
% Lean: (1) is RepLifts.d_dlogC with rU_sub_mem and lg_comp (plog_mul, plog_map_unit), and dlog_mem
% (ProjectiveChartAffineCoords.lean); (2) is RepLifts.main (UnitLogRepresentativeCompat.lean), with isFil_sFil for the
% filtration. Step 1 is he (hun, through J_nil) and α = zeroC of a = plogK (e - 1); Step 2 is key, through lg_shift
% (grp2, plog_mul, plog_inv, plog_map_unit), dlogC_eq and d_zeroC; Step 3 is cupList_diff, with bd_cup and cup_d (both
% through P.leibniz), d_cupList and cupList_mem, the hypotheses hm1, hd and hα; Step 4 is rep_of_bd (rep_add, rep_d). About
% 200 of the 330 lines of UnitLogRepresentativeCompat.lean.
\end{proof}
\leanref{Development/AutomorphicSpectrum/Proofs/BlockDecomposition/UnitLogRepresentativeCompat.lean}{GlobalTheory.CompanionTheory.RepLifts.lg_shift}
\leanref{Development/AutomorphicSpectrum/Proofs/BlockDecomposition/UnitLogRepresentativeCompat.lean}{GlobalTheory.CompanionTheory.RepLifts.cupList_diff}
\leanref{Development/AutomorphicSpectrum/Proofs/BlockDecomposition/UnitLogRepresentativeCompat.lean}{GlobalTheory.CompanionTheory.RepLifts.rep_of_bd}
\leanref{Development/AutomorphicSpectrum/Proofs/BlockDecomposition/UnitLogRepresentativeCompat.lean}{GlobalTheory.CompanionTheory.RepLifts.lg_comp}
\leanref{Development/AutomorphicSpectrum/Proofs/BlockDecomposition/UnitLogRepresentativeCompat.lean}{GlobalTheory.CompanionTheory.RepLifts.rU_sub_mem}

\begin{lemma}\label{gal-realization-kappa-de-rham-transfer-extend}
Let $\Phi$, given by $A$, $J$ and $\gamma$, be a divided power functor on a category $\mathcal C$ over a ring $R$ with a
filtration $(\mathfrak f_r)$ (Definition~\ref{gal-realization-thickening-level-one-graded-filtration}), $\mathcal
V\subset\mathcal U$ full subcategories and $n\geq0$, such that restriction $H^n(\mathcal U,\mathrm{gr}_n)\to H^n(\mathcal
V,\mathrm{gr}_n)$ is surjective. Let $f$ be an $n$-cochain of $\mathcal V$ with values in $A$ whose values lie in
$\mathfrak f_n$ and whose coboundary has values in $\mathfrak f_{n+1}$. Then there is an $n$-cochain $f'$ of $\mathcal U$
with values in $\mathfrak f_n$ and coboundary with values in $\mathfrak f_{n+1}$ such that
$[df]_{n+1}=[df']_{n+1}|_{\mathcal V}$.
\end{lemma}
\leanref{Development/AutomorphicSpectrum/Proofs/BlockDecomposition/FilteredCochainExtendOfSurjective.lean}{GlobalTheory.CompanionTheory.ConeRepExact.exact_lift}

\begin{proof}
For a cochain $z$ whose values lie in $\mathfrak f_r$, let $\bar z$ be the cochain with values in $\mathrm{gr}_r$ whose
values are the classes of those of $z$; then $\overline{dz}=d\bar z$, and restriction to $\mathcal V$ commutes with
$z\mapsto\bar z$ and with $d$. A cochain with values in $\mathrm{gr}_r$ is $\bar z$ for some $z$ with values in $\mathfrak
f_r$, by choosing an element of $\mathfrak f_r$ in each class.

\emph{Step 1: a class on $\mathcal U$.} As $df$ has values in $\mathfrak f_{n+1}$, $d\bar f=\overline{df}=0$, so $\bar f$
is a cocycle of $\mathcal V$ with values in $\mathrm{gr}_n$. By hypothesis its class is the restriction of a class of
$H^n(\mathcal U,\mathrm{gr}_n)$; let $y'$ be a cocycle that represents it.

\emph{Step 2: the cochain $f'$.} Let $f'$ be an $n$-cochain of $\mathcal U$ with values in $\mathfrak f_n$ and
$\bar f'=y'$. Then $\overline{df'}=dy'=0$, so $df'$ has values in $\mathfrak f_{n+1}$.

\emph{Step 3: a correction.} Let $z=f'|_{\mathcal V}-f$. Then $\bar z=y'|_{\mathcal V}-\bar f$, a coboundary, as both
cocycles represent the same class. There is an $n$-cochain $e$ of $\mathcal V$ with values in $\mathfrak f_{n+1}$ and
$de=dz$. If $n=0$, then $\bar z=0$, so $z$ has values in $\mathfrak f_1$; take $e=z$. If $n\geq1$, then $\bar z=d\bar g$
for an $(n-1)$-cochain $g$ with values in $\mathfrak f_n$. Then $\overline{z-dg}=0$, so $e=z-dg$ has values in
$\mathfrak f_{n+1}$, and $de=dz$, as $dd=0$.

\emph{Step 4: the classes.} As restriction commutes with $d$, $d(f'|_{\mathcal V})-df=dz=de$. By
Lemma~\ref{gal-realization-thickening-level-one-graded-forms}(2),
\[
 [df']_{n+1}|_{\mathcal V}-[df]_{n+1}=[d(f'|_{\mathcal V})-df]_{n+1}=[de]_{n+1},
\]
which is the class of $\overline{de}=d\bar e$, a coboundary, where $\bar e$ is formed in $\mathrm{gr}_{n+1}$. So it is
$0$.
% Lean: ConeRepExact.exact_lift. z ↦ z̄ is the relation Rp (Rp.d, Rp.sub, Rp.res, Rp.unique); Step 1 is exists_rp and
% hdy, and the class c' from the surjectivity hs; Step 2 is exists_lift and hf'.d; Step 3 is sub_mem_B (hB) and
% exists_of_mem_B (cases n = 0 and n + 1, with exists_lift for the primitive); Step 4 is e1 (d_pullC), hde, exists_rp for e
% and cls_eq_of_sub_mem. About 150 lines.
\end{proof}
\leanref{Development/AutomorphicSpectrum/Proofs/BlockDecomposition/FilteredCochainExtendOfSurjective.lean}{GlobalTheory.CompanionTheory.ConeRepExact.exists_lift}
\leanref{Development/AutomorphicSpectrum/Proofs/BlockDecomposition/FilteredCochainExtendOfSurjective.lean}{GlobalTheory.CompanionTheory.ConeRepExact.exists_of_mem_B}

\begin{lemma}\label{gal-realization-kappa-de-rham-transfer-local-near}
Let $\Phi$, given by $A$, $J$ and $\gamma$, be a divided power functor on a category $\mathcal C$ over a ring,
$(\mathfrak f_r)$ a filtration of $\Phi$, $(D,\tau)$ a torsor datum of dimension $d$ for $\Phi$ that is graded constant
for $(\mathfrak f_r)$, and $\sigma_T$, $a_k$, $\omega_I$ and $[z]_r$ as in
Definition~\ref{gal-realization-thickening-level-one-graded-filtration}. Cochains are cochains of $\mathcal C$ with values in
$A$, and for $w\in A(D)$ let $\tilde w$ be the $0$-cochain $T\mapsto A(\sigma_T)(w)$. Let $m\geq0$, $\ell_1,\ldots,\ell_m$
$1$-cochains, $c_{k,i}\in A(D)$ for $1\leq k\leq m$ and $1\leq i\leq d$, and $h_1,\ldots,h_m$ $0$-cochains with values in
$\mathfrak f_1$, such that $\ell_k-\sum_i\omega_{(i)}(c_{k,i})-dh_k$ has values in $\mathfrak f_2$ for every $k$. Let
$L=\ell_1\cup\cdots\cup\ell_m$, the constant cochain $1$ for $m=0$, and for $x\in A(D)$ let
$\Omega(x)=\sum_I\omega_I(x\prod_{k=1}^mc_{k,I_k})$, the sum over all sequences $I=(I_1,\ldots,I_m)$ in $\{1,\ldots,d\}$.
\begin{enumerate}
\item There is an $m$-cochain $g$, with $g=0$ if $m=0$ and $g=dH$ for an $(m-1)$-cochain $H$ with values in $\mathfrak f_m$
  if $m\geq1$, such that $L-\Omega(1)-g$ has values in $\mathfrak f_{m+1}$.
\item For $w\in A(D)$, if $[L\cup\tilde w]_m$ is defined, then $[L\cup\tilde w]_m=[\Omega(w)]_m$.
\end{enumerate}
\end{lemma}
\leanref{Development/AutomorphicSpectrum/Proofs/BlockDecomposition/CupListRepresentativeOfForms.lean}{GlobalTheory.CompanionTheory.CupRep.near}
\leanref{Development/AutomorphicSpectrum/Proofs/BlockDecomposition/CupListRepresentativeOfForms.lean}{GlobalTheory.CompanionTheory.CupRep.cup_rep}
\leanref{Development/AutomorphicSpectrum/Proofs/BlockDecomposition/CupListRepresentativeOfFormsStatement.lean}{GlobalTheory.T.cup_list_representative_of_forms}

\begin{proof}
Maps carry each $\mathfrak f_r$ into itself, so the values of a cup product of cochains with values in $\mathfrak f_r$ and
$\mathfrak f_{r'}$ lie in $\mathfrak f_{r+r'}$. For a sequence $I$ and $x,w\in A(D)$, $\omega_I(wx)=\omega_I(x)\cup\tilde
w$, as $A(\sigma_T)$ is multiplicative; and $\omega_I(x)$ has values in $\mathfrak f_j$ and coboundary with values in
$\mathfrak f_{j+1}$, where $j$ is the length of $I$ (Lemma~\ref{gal-realization-thickening-level-one-graded-forms}(3), for
$q=0$). The Leibniz rule $d(a\cup b)=da\cup b+(-1)^ia\cup db$ holds for an $i$-cochain $a$.

\emph{Step 1: two products.} (a) For a sequence $I$ of length $j$, $i\in\{1,\ldots,d\}$ and $x,y\in A(D)$,
$\omega_I(x)\cup\omega_{(i)}(y)-\omega_{(I,i)}(xy)$ has values in $\mathfrak f_{j+2}$, where $(I,i)=(I_1,\ldots,I_j,i)$.
Indeed, at a chain $c_0\to\cdots\to c_{j+1}$, with all elements carried to $A(c_{j+1})$, let $P$ be the product of the
$a_{I_m}(c_{m-1}\to c_m)$, $m\leq j$, and of $a_i(c_j\to c_{j+1})$, which lies in $\mathfrak f_{j+1}$. The two values are
$PWY$ and $PXY$, with $W=A(c_j\to c_{j+1})(A(\sigma_{c_j})(x))$, $X=A(\sigma_{c_{j+1}})(x)$ and
$Y=A(\sigma_{c_{j+1}})(y)$, and $W-X\in\mathfrak f_1$, as $(D,\tau)$ is graded constant. (b) For $w\in A(D)$ and a
$j$-cochain $H$ with values in $\mathfrak f_{j+1}$, $dH\cup\tilde w-d(H\cup\tilde w)$ has values in $\mathfrak f_{j+2}$.
Indeed, at a chain ending at $c_{j+1}$, the term of this difference for the face that omits $c_i$ is $\pm$ the value of
$H$ at that face, carried to $c_{j+1}$, times $A(\sigma_{c_{j+1}})(w)-A(\varphi\sigma_{c'})(w)$, where $c'$ is the last
object of the face and $\varphi:c'\to c_{j+1}$ the map of the chain; this difference lies in $\mathfrak f_1$, as $(D,\tau)$
is graded constant.

\emph{Step 2: proof of (1).} By induction on $m$. For $m=0$, $L=1$ and $\Omega(1)=\omega_\emptyset(1)=\tilde1=1$; take
$g=0$. Suppose that (1) holds for $\ell_1,\ldots,\ell_m$, with $L$, $B=\Omega(1)$ and $g$ formed for them, and let
$l=\ell_{m+1}$, $b=\sum_i\omega_{(i)}(c_{m+1,i})$ and $t=dh_{m+1}$. Then $b$ and $t$ have values in $\mathfrak f_1$, so
$l$ does, as $l-b-t$ has values in $\mathfrak f_2$; and $db$ has values in $\mathfrak f_2$ and $dt=0$, so $dl=d(l-b-t)+db$
has values in $\mathfrak f_2$. Similarly $B$ has values in $\mathfrak f_m$ and $dB$ in $\mathfrak f_{m+1}$. Let $B'$ be
$\Omega(1)$ formed for $\ell_1,\ldots,\ell_{m+1}$. Let $H_1=H\cup l$ if $m\geq1$ and $H_1=0$ if $m=0$, and
$H_2=(-1)^mB\cup h_{m+1}$; both have values in $\mathfrak f_{m+1}$, and
\[
 L\cup l-B'-d(H_1+H_2)=(L-B-g)\cup l+(g\cup l-dH_1)+(B\cup b-B')+(B\cup t-dH_2)+B\cup(l-b-t).
\]
Each term has values in $\mathfrak f_{m+2}$. The first lies in $\mathfrak f_{m+1}\mathfrak f_1$ and the last in
$\mathfrak f_m\mathfrak f_2$. By the Leibniz rule, $g\cup l-dH_1$ is $0$ for $m=0$ and $-(-1)^{m-1}H\cup dl$ for $m\geq1$,
and $B\cup t-dH_2=-(-1)^mdB\cup h_{m+1}$. With $x_I=\prod_{k\leq m}c_{k,I_k}$,
$B\cup b-B'=\sum_I\sum_i\bigl(\omega_I(x_I)\cup\omega_{(i)}(c_{m+1,i})-\omega_{(I,i)}(x_Ic_{m+1,i})\bigr)$, which is
Step 1(a). As $L\cup l=\ell_1\cup\cdots\cup\ell_{m+1}$, (1) holds for $m+1$, with $g=d(H_1+H_2)$.

\emph{Step 3: proof of (2).} By the remark above, $\Omega(w)=\Omega(1)\cup\tilde w$, and $[\Omega(w)]_m$ is defined. If
$m=0$, then $L\cup\tilde w-\Omega(w)=(L-\Omega(1))\cup\tilde w$ has values in $\mathfrak f_1$. If $m\geq1$, let
$H'=H\cup\tilde w$, with values in $\mathfrak f_m$; then
\[
 L\cup\tilde w-\Omega(w)-dH'=(L-\Omega(1)-dH)\cup\tilde w+\bigl(dH\cup\tilde w-d(H\cup\tilde w)\bigr)
\]
has values in $\mathfrak f_{m+1}$, by (1) and Step 1(b), and $[dH']_m$ is the class of the coboundary of the cochain of
the classes of the values of $H'$, so it is $0$. In both cases $[L\cup\tilde w]_m=[\Omega(w)]_m$, by
Lemma~\ref{gal-realization-thickening-level-one-graded-forms}(2).
% Lean: (1) is CupRep.near (CupListRepresentativeOfForms.lean), by induction on m: the identity is bil, with cup_mem,
% FC_anti, bd_cup (H_1, through P.leibniz), cup_fc (Step 1(a), through graded constancy) and b_cup_dh (H_2); (2) is
% CupRep.cup_rep, with form_mul, d_mul (Step 1(b)), rep_add and rep_d in the cases m = 0 and m = n + 1. Statement form:
% GlobalTheory.T.cup_list_representative_of_forms. About 400 lines.
\end{proof}
\leanref{Development/AutomorphicSpectrum/Proofs/BlockDecomposition/CupListRepresentativeOfForms.lean}{GlobalTheory.CompanionTheory.CupRep.cup_fc}
\leanref{Development/AutomorphicSpectrum/Proofs/BlockDecomposition/CupListRepresentativeOfForms.lean}{GlobalTheory.CompanionTheory.CupRep.d_mul}
\leanref{Development/AutomorphicSpectrum/Proofs/BlockDecomposition/CupListRepresentativeOfForms.lean}{GlobalTheory.CompanionTheory.CupRep.bd_cup}
\leanref{Development/AutomorphicSpectrum/Proofs/BlockDecomposition/CupListRepresentativeOfForms.lean}{GlobalTheory.CompanionTheory.CupRep.b_cup_dh}

\begin{lemma}\label{gal-realization-kappa-de-rham-transfer-local-alternating}
Let $\Phi$, given by $A$, $J$ and $\gamma$, be a divided power functor on a category $\mathcal C$ over a ring,
$(\mathfrak f_r)$ a filtration of $\Phi$, $(D,\tau)$ a torsor datum of dimension $d$ for $\Phi$ that is graded constant
for $(\mathfrak f_r)$, and $\sigma_T$, $\omega_I$ and $[z]_r$ as in
Definition~\ref{gal-realization-thickening-level-one-graded-filtration}.
\begin{enumerate}
\item Let $j,q\geq0$, $I$ a sequence of length $j$ in $\{1,\ldots,d\}$ and $x\in\mathfrak f_q(D)$. If $I_a=I_b$ for some
  $a\neq b$, then $[\omega_I(x)]_{j+q}=0$. For a permutation $\pi$ of $\{1,\ldots,j\}$, $[\omega_{I\circ\pi}(x)]_{j+q}=
  \operatorname{sgn}(\pi)[\omega_I(x)]_{j+q}$, where $I\circ\pi=(I_{\pi(1)},\ldots,I_{\pi(j)})$.
\item For a sequence $J$ and a strictly increasing sequence $I$ in $\{1,\ldots,d\}$, both of length $j$, let
  $\varepsilon(I,J)=\operatorname{sgn}(\pi)$ if $J=I\circ\pi$ for a permutation $\pi$, and $\varepsilon(I,J)=0$
  otherwise. Let $E$ be a finite set, and for $e\in E$ let $J_e$ be a sequence of length $j$ in $\{1,\ldots,d\}$ and
  $w_e\in A(D)$. Then
  \[
   \Bigl[\sum_{e\in E}\omega_{J_e}(w_e)\Bigr]_j=\Bigl[\sum_I\omega_I\Bigl(\sum_{e\in E}\varepsilon(I,J_e)w_e\Bigr)\Bigr]_j,
  \]
  the second sum over the strictly increasing sequences $I$ of length $j$ in $\{1,\ldots,d\}$.
\end{enumerate}
\end{lemma}
\leanref{Development/AutomorphicSpectrum/Proofs/BlockDecomposition/FormAlternating.lean}{GlobalTheory.CompanionTheory.AltForm.form_alternating}
\leanref{Development/AutomorphicSpectrum/Proofs/BlockDecomposition/FormSumRepresentativeLift.lean}{GlobalTheory.CompanionTheory.ConeLocal.rep_collapse}
\leanref{Development/AutomorphicSpectrum/Proofs/BlockDecomposition/FormSumRepresentativeLiftStatement.lean}{GlobalTheory.T.form_sum_representative_lift}

\begin{proof}
The classes are defined by Lemma~\ref{gal-realization-thickening-level-one-graded-forms}(3), and $\varepsilon(I,J)$ is well
defined, as $I$ is injective. Cochains are cochains of $\mathcal C$ with values in $A$. For $r\in\mathbf Z$, say that a
$j$-cochain $z$ is a \emph{coboundary modulo $\mathfrak f_{r+1}$} if $z$ has values in $\mathfrak f_{r+1}$, when $j=0$, or
if $z-d\beta$ has values in $\mathfrak f_{r+1}$ for a $(j-1)$-cochain $\beta$ with values in $\mathfrak f_r$, when
$j\geq1$. Sums of such cochains are such cochains.

\emph{Step 1: general facts.} (i) If $[z]_r$ is defined and $z'-z$ is a coboundary modulo $\mathfrak f_{r+1}$, then
$[z']_r$ is defined and equals $[z]_r$. Indeed, $z'=z+d\beta+e$, or $z'=z+e$, with $\beta$ as above and $e$ with values in
$\mathfrak f_{r+1}$; so $z'$ has values in $\mathfrak f_r$ and $dz'=dz+de$ in $\mathfrak f_{r+1}$, and $[d\beta]_r$ and
$[e]_r$ are $0$ (Lemma~\ref{gal-realization-thickening-level-one-graded-forms}(2)). (ii) If a $j$-cochain $z$ is a
coboundary modulo $\mathfrak f_{r+1}$ and $y$ is a cochain with values in $\mathfrak f_q$ whose coboundary has values in
$\mathfrak f_{q+1}$, then $z\cup y$ is a coboundary modulo $\mathfrak f_{r+q+1}$. For $j=0$ this is clear; for $j\geq1$,
by the Leibniz rule, $z\cup y=(z-d\beta)\cup y+d(\beta\cup y)-(-1)^{j-1}\beta\cup dy$, where $\beta\cup y$ has values in
$\mathfrak f_{r+q}$ and the other two terms in $\mathfrak f_{r+q+1}$. (iii) Write $a_i$ also for the $1$-cochain
$\varphi\mapsto a_i(\varphi)$. It is a cocycle with values in $J=\mathfrak f_1$, by
Lemma~\ref{gal-realization-thickening-level-one-graded-forms}(1). For $x\in\mathfrak f_q(D)$ let $\tilde x$ be the
$0$-cochain $T\mapsto A(\sigma_T)(x)$, with values in $\mathfrak f_q$ and coboundary with values in $\mathfrak f_{q+1}$, as
$(D,\tau)$ is graded constant. By the definitions of $\omega_I$ and of the cup product
(Definition~\ref{gal-realization-split-diagrams}(1)), $\omega_I(x)=a_{I_1}\cup\cdots\cup a_{I_j}\cup\tilde x$. (iv) For
$1$-cocycles $a,b$ with values in $J$, let $ab$ and $\gamma_2(a)$ be the $1$-cochains $\varphi\mapsto a(\varphi)b(\varphi)$
and $\varphi\mapsto\gamma_2(a(\varphi))$, with values in $\mathfrak f_2$. Then $d(ab)=-(a\cup b+b\cup a)$ and
$d(\gamma_2(a))=-a\cup a$. Indeed, at a chain $c_0\to c_1\to c_2$ let $\alpha_0$, $\alpha_1$ and $\alpha_2$ be the values of
$a$ at $c_1\to c_2$, $c_0\to c_2$ and $c_0\to c_1$, the last carried to $c_2$, and $\beta_0,\beta_1,\beta_2$ those of $b$.
Then $\alpha_1=\alpha_0+\alpha_2$ and $\beta_1=\beta_0+\beta_2$, as $a$ and $b$ are cocycles; the value of $d(ab)$ is
$\alpha_0\beta_0-\alpha_1\beta_1+\alpha_2\beta_2=-(\alpha_2\beta_0+\beta_2\alpha_0)$, and that of $a\cup b+b\cup a$ is
$\alpha_2\beta_0+\beta_2\alpha_0$. As maps commute with $\gamma_2$, the value of $d(\gamma_2(a))$ is
$\gamma_2(\alpha_0)-\gamma_2(\alpha_0+\alpha_2)+\gamma_2(\alpha_2)=-\alpha_0\alpha_2$.

\emph{Step 2: products of cocycles.} Let $j\geq1$, $c_1,\ldots,c_{j+1}$ $1$-cocycles with values in $\mathfrak f_1$,
$1\leq k\leq j$, and $c'$ the sequence $c$ with $c_k$ and $c_{k+1}$ exchanged. Then $c_1\cup\cdots\cup c_{j+1}+c'_1\cup
\cdots\cup c'_{j+1}$ is a coboundary modulo $\mathfrak f_{j+2}$, and so is $c_1\cup\cdots\cup c_{j+1}$ if $c_k=c_{k+1}$.
This is shown by induction on $j$. If $k=j$, let $u=c_1\cup\cdots\cup c_{j-1}$, the constant cochain $1$ for $j=1$, a
cocycle with values in $\mathfrak f_{j-1}$ (Lemma~\ref{gal-realization-kappa-de-rham-cup}), $a=c_j$ and $b=c_{j+1}$. By
associativity of the cup product, (iv) and the Leibniz rule, $(u\cup a)\cup b+(u\cup b)\cup a=-u\cup d(ab)=
(-1)^jd(u\cup ab)$ and $(u\cup a)\cup a=(-1)^jd(u\cup\gamma_2(a))$, where $u\cup ab$ and $u\cup\gamma_2(a)$ have values in
$\mathfrak f_{j+1}$. If $k<j$, the products are $(c_1\cup\cdots\cup c_j)\cup c_{j+1}$ and $(c'_1\cup\cdots\cup c'_j)\cup
c_{j+1}$, and the claims follow from the case $j-1$ and (ii), for $y=c_{j+1}$, whose coboundary is $0$.

\emph{Step 3: proof of (1).} Let $1\leq k<j$ and $\tau$ the transposition of $k$ and $k+1$. By Step 2 for $c_m=a_{I_m}$, (ii)
for $y=\tilde x$ and (iii), $\omega_I(x)+\omega_{I\circ\tau}(x)$ is a coboundary modulo $\mathfrak f_{j+q+1}$, and so is
$\omega_I(x)$ if $I_k=I_{k+1}$. By (i), applied to $z=0$, and
Lemma~\ref{gal-realization-thickening-level-one-graded-forms}(2), $[\omega_{I\circ\tau}(x)]_{j+q}=-[\omega_I(x)]_{j+q}$,
and $[\omega_I(x)]_{j+q}=0$ if $I_k=I_{k+1}$. As every permutation is a product of such transpositions, and $(I\circ\pi)
\circ\pi'=I\circ(\pi\pi')$, the formula for $\pi$ follows. If $I_a=I_b$ with $a<b$, let $\pi$ be the transposition of $a+1$
and $b$, the identity if $b=a+1$; then $(I\circ\pi)_a=I_a=I_b=(I\circ\pi)_{a+1}$, so $[\omega_{I\circ\pi}(x)]_{j+q}=0$, and
$[\omega_I(x)]_{j+q}=\operatorname{sgn}(\pi)[\omega_{I\circ\pi}(x)]_{j+q}=0$.

\emph{Step 4: proof of (2).} Let $e\in E$. If $J_e$ has a repeated entry, then $[\omega_{J_e}(w_e)]_j=0$ by (1), and
$\varepsilon(I,J_e)=0$ for every $I$. Otherwise there are exactly one strictly increasing $I$ and one permutation $\pi$
with $J_e=I\circ\pi$, obtained by ordering the entries of $J_e$; then $[\omega_{J_e}(w_e)]_j=\operatorname{sgn}(\pi)
[\omega_I(w_e)]_j=[\omega_I(\varepsilon(I,J_e)w_e)]_j$ by (1), and $\varepsilon(I',J_e)=0$ for $I'\neq I$. As $\omega_I$ is
additive in its argument, summing over $e$ gives (2), by
Lemma~\ref{gal-realization-thickening-level-one-graded-forms}(2).
% Lean: (1) is the first two parts of AltForm.form_alternating (FormAlternating.lean). "Coboundary modulo f_{r+1}" is Cob,
% with Cob.add; (i) is Rep.of_cob; (ii) is Cob.cp_right; (iii) is ca_coc (through a_comp), cx_mem, cx_dd and form_eq (with
% wd_apply); (iv) is dd_pw and dd_gm; Step 2 is wd_swap and wd_sq, through cob_swap_top, cob_sq_top, cob_top, wd_mem and
% CupAssocC.cup_assocC for associativity; Step 3 is rep_swap, rep_sq, rep_perm (through
% Equiv.Perm.mclosure_swap_castSucc_succ) and the end of form_alternating (with Equiv.swap k.succ b). Its third part, the
% independence of the maps σ_T (wd_change, ca_sub, cx_sub, rep_change), is not used here. (2) is ConeLocal.rep_collapse
% (FormSumRepresentativeLift.lean), through rep_collapse_one (exists_increasing with Tuple.sort, sgn_comp_self,
% sgn_comp_of_ne, sgn_eq_zero_of_not_injective) and rep_finset_sum; ε(I, J) is ConeLocal.sgn. About 740 lines in
% FormAlternating.lean and 130 in FormSumRepresentativeLift.lean.
\end{proof}
\leanref{Development/AutomorphicSpectrum/Proofs/BlockDecomposition/FormAlternating.lean}{GlobalTheory.CompanionTheory.AltForm.Rep.of_cob}
\leanref{Development/AutomorphicSpectrum/Proofs/BlockDecomposition/FormAlternating.lean}{GlobalTheory.CompanionTheory.AltForm.dd_pw}
\leanref{Development/AutomorphicSpectrum/Proofs/BlockDecomposition/FormAlternating.lean}{GlobalTheory.CompanionTheory.AltForm.dd_gm}
\leanref{Development/AutomorphicSpectrum/Proofs/BlockDecomposition/FormAlternating.lean}{GlobalTheory.CompanionTheory.AltForm.wd_swap}
\leanref{Development/AutomorphicSpectrum/Proofs/BlockDecomposition/FormAlternating.lean}{GlobalTheory.CompanionTheory.AltForm.wd_sq}
\leanref{Development/AutomorphicSpectrum/Proofs/BlockDecomposition/FormAlternating.lean}{GlobalTheory.CompanionTheory.AltForm.rep_perm}

\begin{lemma}\label{gal-realization-kappa-de-rham-transfer-local-connecting}
Let $\Phi$, given by $A$, $J$ and $\gamma$, be a divided power functor on a category $\mathcal C$ over a ring,
$(\mathfrak f_r)$ a filtration of $\Phi$, $(D,\tau)$ a torsor datum of dimension $d$ for $\Phi$ that is graded constant
for $(\mathfrak f_r)$, and $\sigma_T$, $a_k$, $\omega_I$ and $[z]_r$ as in
Definition~\ref{gal-realization-thickening-level-one-graded-filtration}. Let $p_1,\ldots,p_d$ be maps $A(D)\to A(D)$ such
that
\[
 A(u)(x)-A(v)(x)-\sum_{k=1}^d\tau_k(v,u)\,A(v)\bigl(p_k(x)\bigr)\in\mathfrak f_2(T)
\]
for every object $T$, all maps $u,v:D\to T$ and every $x\in A(D)$.
\begin{enumerate}
\item For a sequence $I$ of length $j$ in $\{1,\ldots,d\}$ and $x\in A(D)$, $d\omega_I(x)-(-1)^j\sum_k\omega_{(I,k)}(p_k(x))$
  has values in $\mathfrak f_{j+2}$, where $(I,k)=(I_1,\ldots,I_j,k)$.
\item Let $n\geq0$, $f$ an $n$-cochain of $\mathcal C$ with values in $A$ whose values lie in $\mathfrak f_n$ and whose
  coboundary has values in $\mathfrak f_{n+1}$, and $x_I\in A(D)$, for the strictly increasing sequences $I$ of length $n$
  in $\{1,\ldots,d\}$, such that $[f]_n=[\sum_I\omega_I(x_I)]_n$. Then
  \[
   [df]_{n+1}=\Bigl[\sum_I\sum_{k=1}^d\omega_{(I,k)}\bigl((-1)^np_k(x_I)\bigr)\Bigr]_{n+1}.
  \]
\end{enumerate}
\end{lemma}
\leanref{Development/AutomorphicSpectrum/Proofs/BlockDecomposition/TaylorDifferentialRepresentative.lean}{GlobalTheory.CompanionTheory.ConeLocal.d_form}
\leanref{Development/AutomorphicSpectrum/Proofs/BlockDecomposition/TaylorDifferentialRepresentative.lean}{GlobalTheory.CompanionTheory.ConeLocal.connecting_form}
\leanref{Development/AutomorphicSpectrum/Proofs/BlockDecomposition/TaylorDifferentialRepresentativeStatement.lean}{GlobalTheory.T.taylor_differential_representative}

\begin{proof}
(1) Let $s=(c_0\to\cdots\to c_{j+1})$ be a chain, $T=c_{j+1}$, and carry all elements to $A(T)$ along the chain. By the
computation in the proof of Lemma~\ref{gal-realization-thickening-level-one-graded-forms}(3), the value of $d\omega_I(x)$
at $s$ is $(-1)^jP(Y-W)$, where $P$ is the product of the images of the $a_{I_m}(c_{m-1}\to c_m)$, $m\leq j$,
$Y=A(\sigma_T)(x)$ and $W=A(v)(x)$ for $v=(c_j\to T)\sigma_{c_j}$. By the hypothesis, for $u=\sigma_T$ and $v$,
$Y-W-\sum_k\tau_k(v,\sigma_T)A(v)(p_k(x))\in\mathfrak f_2(T)$, where $\tau_k(v,\sigma_T)=a_k(c_j\to T)\in J(T)$. As
$(D,\tau)$ is graded constant, $A(v)(p_k(x))-A(\sigma_T)(p_k(x))\in\mathfrak f_1(T)$, so $Y-W-\sum_ka_k(c_j\to T)
A(\sigma_T)(p_k(x))\in\mathfrak f_2(T)$. Multiplying by $(-1)^jP\in\mathfrak f_j(T)$ gives the claim, as
$P\,a_k(c_j\to T)A(\sigma_T)(p_k(x))$ is the value of $\omega_{(I,k)}(p_k(x))$ at $s$.

(2) Let $g=\sum_I\omega_I(x_I)$. As $[f]_n=[g]_n$, the cochain of the classes in $\mathrm{gr}_n$ of the values of $f-g$ is
a coboundary. If $n=0$, it is $0$, and $e=f-g$ has values in $\mathfrak f_1$. If $n\geq1$, it is the coboundary of an
$(n-1)$-cochain with values in $\mathrm{gr}_n$; lifting its values to $\mathfrak f_n$ gives an $(n-1)$-cochain $\beta$
with values in $\mathfrak f_n$ such that $e=f-g-d\beta$ has values in $\mathfrak f_{n+1}$. In both cases $df=dg+de$. The
classes $[dg]_{n+1}$ and $[de]_{n+1}$ are defined, and $[de]_{n+1}$ is the class of the coboundary of the cochain of the
classes of the values of $e$, so it is $0$; so $[df]_{n+1}=[dg]_{n+1}$ by
Lemma~\ref{gal-realization-thickening-level-one-graded-forms}(2). By (1), as $\omega_{(I,k)}$ is additive in its argument,
$dg-\sum_I\sum_k\omega_{(I,k)}((-1)^np_k(x_I))$ has values in $\mathfrak f_{n+2}$, and the same lemma, with its part (3),
gives the claim.
% Lean: (1) is ConeLocal.d_form (TaylorDifferentialRepresentative.lean), through alt_sum_form (the value (−1)^j P (Y − W),
% with headP for P and the telescoping identity CocProd.tele), the Taylor hypothesis hT for u = σ_T and v (hin), ga_last
% for τ_k(v, σ_T) = a_k, hdiff for the graded constancy, prod_mem and form_snoc; (2) is connecting_form, through
% d_eq_of_rep (the cochain e), rep_d, rep_add and form_neg_pow. The hypothesis on the p_k is ConeLocal.Taylor
% (RepUnique.lean). Statement form: GlobalTheory.T.taylor_differential_representative. About 270 lines.
\end{proof}
\leanref{Development/AutomorphicSpectrum/Proofs/BlockDecomposition/RepUnique.lean}{GlobalTheory.CompanionTheory.ConeLocal.Taylor}

\begin{lemma}\label{gal-realization-kappa-de-rham-transfer-local-dlog}
Let $\Phi$, given by $A$, $J$ and $\gamma$, be a divided power functor on a category $\mathcal C$ over a ring,
$(\mathfrak f_r)$ a filtration of $\Phi$, $(D,\tau)$ a torsor datum of dimension $d$ for $\Phi$ that is graded constant
for $(\mathfrak f_r)$, $\sigma_T$ and $a_k$ as in Definition~\ref{gal-realization-thickening-level-one-graded-filtration},
and $p_1,\ldots,p_d$ as in Lemma~\ref{gal-realization-kappa-de-rham-transfer-local-connecting}. Suppose that $1+x$ is a unit
of $A(T)$ for every object $T$ and $x\in J(T)$. Let $G$ be a unit of $A(D)$, and $\hat g_T\in A(T)$, for the objects $T$,
with $\hat g_T-A(\sigma_T)(G)\in J(T)$; let $h(T)=\hat g_TA(\sigma_T)(G)^{-1}-1$. Then $h(T)\in J(T)$, $\hat g_T$ is a unit,
and for every map $\varphi:T\to T'$, $r(\varphi)=\hat g_{T'}A(\varphi)(\hat g_T)^{-1}$ has $r(\varphi)-1\in J(T')$ and
\[
 r(\varphi)-1-\sum_{i=1}^da_i(\varphi)\,A(\sigma_{T'})\bigl(G^{-1}p_i(G)\bigr)-\bigl(h(T')-A(\varphi)(h(T))\bigr)
 \in\mathfrak f_2(T').
\]
\end{lemma}
\leanref{Development/AutomorphicSpectrum/Proofs/BlockDecomposition/LocalFrameTaylorClassEq.lean}{GlobalTheory.CompanionTheory.ConeLocal.DlogApprox.frame_main}

\begin{proof}
As $J(T')J(T')\subset\mathfrak f_2(T')$, a sum of products each with at least two factors in $J(T')$ lies in $\mathfrak
f_2(T')$. First, $h(T)=(\hat g_T-A(\sigma_T)(G))A(\sigma_T)(G)^{-1}\in J(T)$, and $\hat g_T=(1+h(T))A(\sigma_T)(G)$ is a
unit. Let $v=\varphi\sigma_T$, $Y_0=A(v)(G)$, $I_0=A(v)(G^{-1})$, $\Delta=A(\sigma_{T'})(G)-Y_0$, which lies in $J(T')$ as
$(D,\tau)$ is graded constant, $H=h(T')$ and $H_0=A(\varphi)(h(T))$, both in $J(T')$. Then $\hat g_{T'}=(1+H)(Y_0+\Delta)$
and $A(\varphi)(\hat g_T)=(1+H_0)Y_0$, so
\[
 r(\varphi)=(1+H)(1+\Delta I_0)(1+H_0)^{-1},\qquad (1+H_0)^{-1}=1-H_0+H_0^2(1+H_0)^{-1},
\]
and $r(\varphi)-1-(\Delta I_0+H-H_0)\in\mathfrak f_2(T')$; in particular $r(\varphi)-1\in J(T')$. By the hypothesis on the
$p_i$, for $u=\sigma_{T'}$, $v$ and $x=G$, $\Delta-\sum_i\tau_i(v,\sigma_{T'})A(v)(p_i(G))\in\mathfrak f_2(T')$, where
$\tau_i(v,\sigma_{T'})=a_i(\varphi)$; multiplying by $I_0$ gives $\Delta I_0-\sum_ia_i(\varphi)A(v)(G^{-1}p_i(G))\in
\mathfrak f_2(T')$. Finally $a_i(\varphi)\bigl(A(v)(G^{-1}p_i(G))-A(\sigma_{T'})(G^{-1}p_i(G))\bigr)$ lies in $J(T')J(T')$,
as $(D,\tau)$ is graded constant. Adding up gives the claim.
% Lean: DlogApprox.frame_main (LocalFrameTaylorClassEq.lean): hh for h(T) ∈ J, e1, e2, e3, hU and hW for the units,
% and frame_core: alg_core for r(φ) − 1 − (ΔI_0 + H − H_0) (with mul2 for J J ⊂ f_2), the Taylor hypothesis ht
% multiplied by I_0 (h2), and h3 for the change from v to σ_{T'} (const1). About 150 lines.
\end{proof}
\leanref{Development/AutomorphicSpectrum/Proofs/BlockDecomposition/LocalFrameTaylorClassEq.lean}{GlobalTheory.CompanionTheory.ConeLocal.DlogApprox.frame_core}
\leanref{Development/AutomorphicSpectrum/Proofs/BlockDecomposition/LocalFrameTaylorClassEq.lean}{GlobalTheory.CompanionTheory.ConeLocal.DlogApprox.alg_core}

\begin{lemma}\label{gal-realization-kappa-de-rham-transfer-local-taylor}
Let $K$ be a ring, $G$ a $K$-algebra and $t_1,\ldots,t_d\in G$ such that $G$ is formally unramified over
$K[x_1,\ldots,x_d]$ for the $K$-algebra map $x_k\mapsto t_k$, and $D_1,\ldots,D_d$ $K$-derivations of $G$ with $D_k(t_j)=1$
for $k=j$ and $0$ otherwise. Let $B$ be a $K$-algebra, $N$ and $P$ ideals of $B$ with $N^2\subset P$, and $a,b:G\to B$
maps of $K$-algebras with $a(x)-b(x)\in N$ for every $x\in G$. Then for every $x\in G$
\[
 a(x)-b(x)-\sum_{k=1}^d\bigl(a(t_k)-b(t_k)\bigr)\,b\bigl(D_k(x)\bigr)\in P.
\]
\end{lemma}
\leanref{Development/AutomorphicSpectrum/Proofs/BlockDecomposition/TaylorOfLocalFrame.lean}{GlobalTheory.CompanionTheory.ConeLocal.taylor_alg}

\begin{proof}
Let $e_k=a(t_k)-b(t_k)\in N$ and $S(x)=\sum_ke_kb(D_k(x))\in N$. Then $S$ is additive, kills the image of $K$, has
$S(t_j)=e_j$, and $S(xy)=b(x)S(y)+b(y)S(x)$, as the $D_k$ are derivations. Let $\pi:B\to B/P$ be the projection and
$c(x)=\pi(b(x)+S(x))$. Then $c$ is additive, $c(1)=1$, and $c(xy)-c(x)c(y)=-\pi(S(x)S(y))=0$, as $S(x)S(y)\in N^2\subset
P$; so $c$ is a map of $K$-algebras. It agrees with $\pi\circ a$ on the image of $K[x_1,\ldots,x_d]$, as both are maps of
$K$-algebras and $c(t_j)=\pi(b(t_j)+e_j)=\pi(a(t_j))$. So $c$ and $\pi\circ a$ are maps of $K[x_1,\ldots,x_d]$-algebras
$G\to B/P$, and $c(x)-\pi(a(x))=\pi(b(x)-a(x)+S(x))$ lies in the ideal $\pi(N)$, whose square is $0$. As $G$ is formally
unramified over $K[x_1,\ldots,x_d]$, $c=\pi\circ a$ \cite[Tag 00UN]{Sta}. So $a(x)-b(x)-S(x)\in P$.
% Lean: ConeLocal.taylor_alg (TaylorOfLocalFrame.lean): S with hSadd, hSmul, hSone and hSzero, the ring map c, hcf on the
% polynomial algebra (MvPolynomial.ringHom_ext), and Algebra.FormallyUnramified.ext for the ideal N.map π of square 0
% (hnil). About 90 lines.
\end{proof}

\begin{lemma}\label{gal-realization-kappa-de-rham-transfer-local-chart}
Let $K$ be a field of characteristic $0$, $X$ a $K$-scheme, $W\subset X$ an affine open with coordinates
$t_1,\ldots,t_d$ (Definition~\ref{gal-realization-thickening-coordinates}(1)), $s\geq1$,
$\mathcal Q=\mathrm{Th}^{(s)}_{\mathrm{dR}}(X)_W$ (Definition~\ref{gal-realization-kappa-de-rham-truncated}), with the
divided power functor $\Phi_{\mathcal Q}$ (Definition~\ref{gal-realization-thickening-torsor}(1)), and
$G=\Gamma(W,\mathcal O_X)$. There are a thickening $D$ of $\mathcal Q$ and a ring isomorphism $\theta:G\to B_D$ with the
following properties, where $t'_k=\theta(t_k)$.
\begin{enumerate}
\item $\theta$ is a map of $K$-algebras, and $J_D=0$.
\item For every map $\sigma:D\to T$ of $\mathcal Q$ and $f\in G$, $\sigma(\theta(f))$ is a lift of $f$ at $T$
  (Definition~\ref{gal-realization-thickening-level-one-fibre}(2)).
\item The families $\mathcal O^{[r]}$ form a filtration of $\Phi_{\mathcal Q}$, and with $\tau_k(\sigma,\sigma')=
  \sigma'(t'_k)-\sigma(t'_k)$, $(D,\tau)$ is a torsor datum of dimension $d$ for $\Phi_{\mathcal Q}$ that is graded constant
  for it (Definitions~\ref{gal-realization-thickening-torsor}(2)
  and~\ref{gal-realization-thickening-level-one-graded-filtration}).
\item Choose maps $\sigma_T:D\to T$, and let $\omega_I$ and $[z]_r$ be as in
  Definition~\ref{gal-realization-thickening-level-one-graded-filtration}, for this filtration. For $0\leq m<s$, every class
  of $H^m(\mathcal Q,\mathrm{gr}_m)$ is $[\sum_I\omega_I(\theta(y_I))]_m$ for exactly one family $(y_I)$ of elements of
  $G$, indexed by the strictly increasing sequences $I$ of length $m$ in $\{1,\ldots,d\}$.
\item Let $D_1,\ldots,D_d$ be $K$-derivations of $G$ with $D_k(t_j)=1$ for $k=j$ and $0$ otherwise, and
  $p_k=\theta D_k\theta^{-1}$. Then $u(x)-v(x)-\sum_k\tau_k(v,u)\,v(p_k(x))\in J_T^{[2]}$ for every thickening $T$ of
  $\mathcal Q$, all maps $u,v:D\to T$ and every $x\in B_D$.
\end{enumerate}
\end{lemma}
\leanref{Development/AutomorphicSpectrum/Proofs/BlockDecomposition/ExistsLocalFrameOfCoordsStatement.lean}{GlobalTheory.T.exists_local_frame_of_coords}
\leanref{Development/AutomorphicSpectrum/Proofs/BlockDecomposition/RepUnique.lean}{GlobalTheory.CompanionTheory.ConeLocal.IsLoc}
\leanref{Development/AutomorphicSpectrum/Proofs/BlockDecomposition/LocalGradedClassRepresentativeStatement.lean}{GlobalTheory.T.local_graded_class_representative}
\leanref{Development/AutomorphicSpectrum/Proofs/BlockDecomposition/TaylorOfLocalFrameStatement.lean}{GlobalTheory.T.taylor_of_local_frame}

\begin{proof}
Steps 1 to 3 of the proof of Lemma~\ref{gal-realization-kappa-de-rham-local}, for $V=W$ with its coordinates, give $D$ and
$\theta$, $f\mapsto1\otimes f$, with $\theta(t_k)=t'_k$, $J_D=0$ and $\mathrm{ev}_D=\theta$; they give (3), and a divided
power polynomial structure on $\Phi_{\mathcal Q}$ that respects $(\mathcal O^{[r]})_r$ in every degree $r\leq s$.

(1) $\theta$ is $K$-linear. (2) By Lemma~\ref{gal-realization-thickening-level-one-forms}(1),
$\mathrm{ev}_T=\bar\sigma\circ\mathrm{ev}_D=\bar\sigma\circ\theta$, so $\sigma(\theta(f))$ has class $\mathrm{ev}_T(f)$.

(4) As $m+1\leq s$, Lemma~\ref{gal-realization-thickening-level-one-graded}(2), with $j=r=m$ and $q=0$, shows that every
class is $[\sum_I\omega_I(x_I)]_m$ with $x_I\in\mathcal O^{[0]}(D)=B_D=\theta(G)$. If two families $(y_I)$ and $(y'_I)$
give the same class, then $[\sum_I\omega_I(\theta(y_I-y'_I))]_m=0$, as $\omega_I$ is additive in its argument
(Lemma~\ref{gal-realization-thickening-level-one-graded-forms}(2)); so $\theta(y_I-y'_I)\in J_D^{[1]}=0$ by the same part
of Lemma~\ref{gal-realization-thickening-level-one-graded}, and $y_I=y'_I$.

(5) The maps $a=u\theta$ and $b=v\theta$ from $G$ to $B_T$ are maps of $K$-algebras by (1), and $a(f)-b(f)\in J_T$ for
$f\in G$, by (2). The ideals $N=J_T$ and $P=J_T^{[2]}$ have $N^2\subset P$, and $G$ is formally unramified over
$K[x_1,\ldots,x_d]$, as it is \'etale over it \cite[Tags 00UR and 00UQ]{Sta}. As $\tau_k(v,u)=a(t_k)-b(t_k)$,
Lemma~\ref{gal-realization-kappa-de-rham-transfer-local-taylor} gives the claim for $x=\theta(f)$, as
$v(p_k(\theta(f)))=b(D_k(f))$.
% Lean: (1) to (3) are ConeLocal.exists_loc (ExistsLocalFrameOfCoords.lean), through ChartAlg.chart_exists (Lemma
% gal-realization-kappa-de-rham-chart for O = R = K and I = 0, with tfBase_dpBot), ChartMaps.good, ChartFil.chart_grConst
% and c'.frame, with θ = e ∘ (TensorProduct.lid)⁻¹; ConeLocal.IsLoc (RepUnique.lean) records them, (2) as
% ev T g = class of u(θ g). (4) is ConeLocal.loc_coef (LocalGradedClassRepresentative.lean), through
% GlobalTheory.T.polynomial_transfer (Lemma gal-realization-kappa-de-rham-poly) and local_graded_vanishing (Lemma
% gal-realization-thickening-level-one-graded). (5) is ConeLocal.taylor_loc (TaylorOfLocalFrame.lean), with
% Algebra.FormallyUnramified from the étale coordinates c.etale. About 140 lines.
\end{proof}
\leanref{Development/AutomorphicSpectrum/Proofs/BlockDecomposition/ExistsLocalFrameOfCoords.lean}{GlobalTheory.CompanionTheory.ConeLocal.exists_loc}
\leanref{Development/AutomorphicSpectrum/Proofs/BlockDecomposition/LocalGradedClassRepresentative.lean}{GlobalTheory.CompanionTheory.ConeLocal.loc_coef}
\leanref{Development/AutomorphicSpectrum/Proofs/BlockDecomposition/TaylorOfLocalFrame.lean}{GlobalTheory.CompanionTheory.ConeLocal.taylor_loc}

\begin{lemma}\label{gal-realization-kappa-de-rham-transfer-local-compare}
Let $K$ be a field of characteristic $0$, $X$ a $K$-scheme, $s\geq1$, and $W\subset U$ affine opens of $X$ with
coordinates $t_1,\ldots,t_d$ on $W$ and $u_1,\ldots,u_e$ on $U$; write $f|_W$ for the restriction of $f\in\Gamma(U,\mathcal
O_X)$ to $W$. Let $p_1,\ldots,p_d$ be $K$-derivations of $\Gamma(W,\mathcal O_X)$ with $p_k(t_j)=1$ for $k=j$ and $0$
otherwise. Let $\mathcal Q(W)\subset\mathcal Q(U)$ and $(D_W,\theta_W)$ and $(D_U,\theta_U)$ be as in
Lemma~\ref{gal-realization-kappa-de-rham-transfer-local-chart} for $W$ and for $U$, with maps $\sigma^W_T$ and $\sigma^U_T$
and forms $\omega^W_I$ and $\omega^U_L$. Let $0\leq m<s$, $\xi\in H^m(\mathcal Q(W),\mathrm{gr}_m)$,
$y_I\in\Gamma(W,\mathcal O_X)$ for the strictly increasing sequences $I$ of length $m$ in $\{1,\ldots,d\}$ and
$z_L\in\Gamma(U,\mathcal O_X)$ for those $L$ in $\{1,\ldots,e\}$, such that $\xi=[\sum_I\omega^W_I(\theta_W(y_I))]_m$ and
$\xi$ is the restriction to $\mathcal Q(W)$ of $[\sum_L\omega^U_L(\theta_U(z_L))]_m$. Then for every $I$
\[
 y_I=\sum_Lz_L|_W\det\bigl(p_{I_k}(u_{L_l}|_W)\bigr)_{1\leq k,l\leq m}.
\]
\end{lemma}
\leanref{Development/AutomorphicSpectrum/Proofs/BlockDecomposition/BasicOpenRsLe.lean}{GlobalTheory.CompanionTheory.ConeLocal.CoefMemCmp.cmp_coef}

\begin{proof}
A lift of $f\in\Gamma(U,\mathcal O_X)$ at a thickening of $\mathcal Q(W)$ is a lift of $f|_W$
(Lemma~\ref{gal-realization-thickening-level-one-forms}(1)). Let $\tau^W$ and $\tau^U$ be the torsor data of the two charts
and $a^W_i$ and $a^U_k$ the elements of Definition~\ref{gal-realization-thickening-level-one-graded-filtration}(3) for
them. For $x\in B_{D_W}$ let $\tilde x$ be the $0$-cochain $T\mapsto\sigma^W_T(x)$ of $\mathcal Q(W)$.

\emph{Step 1: one form.} Let $L$ be strictly increasing of length $m$ in $\{1,\ldots,e\}$, $z\in\Gamma(U,\mathcal O_X)$ and
$\omega'$ the restriction of $\omega^U_L(\theta_U(z))$ to $\mathcal Q(W)$. Then
\[
 [\omega']_m=\Bigl[\sum_J\omega^W_J\Bigl(\theta_W\Bigl(z|_W\prod_{k=1}^mp_{J_k}(u_{L_k}|_W)\Bigr)\Bigr)\Bigr]_m,
\]
the sum over all sequences $J$ of length $m$ in $\{1,\ldots,d\}$. Indeed, let $\ell_k$ be the $1$-cochain of $\mathcal Q(W)$
that sends $\varphi:T\to T'$ to $a^U_{L_k}(\varphi)=\sigma^U_{T'}(\theta_U(u_{L_k}))-\varphi\sigma^U_T(\theta_U(u_{L_k}))$,
let $c_{k,i}=\theta_W(p_i(u_{L_k}|_W))$, and let $h_k$ be the $0$-cochain $T\mapsto\sigma^U_T(\theta_U(u_{L_k}))-
\sigma^W_T(\theta_W(u_{L_k}|_W))$, with values in $J$, as both terms are lifts of $u_{L_k}|_W$
(Lemma~\ref{gal-realization-kappa-de-rham-transfer-local-chart}(2)). With $x=\theta_W(u_{L_k}|_W)$,
$\ell_k(\varphi)-(dh_k)(\varphi)=\sigma^W_{T'}(x)-\varphi\sigma^W_T(x)$. By
Lemma~\ref{gal-realization-kappa-de-rham-transfer-local-chart}(5) for $W$, with $u=\sigma^W_{T'}$ and $v=\varphi\sigma^W_T$,
this differs from $\sum_ia^W_i(\varphi)\,\varphi\sigma^W_T(p_i(x))$ by an element of $J^{[2]}_{T'}$, as
$\tau^W_i(\varphi\sigma^W_T,\sigma^W_{T'})=a^W_i(\varphi)$, and $\varphi\sigma^W_T(p_i(x))-\sigma^W_{T'}(p_i(x))\in J_{T'}$,
as $(D_W,\tau^W)$ is graded constant. As $\omega^W_{(i)}(c)(\varphi)=a^W_i(\varphi)\sigma^W_{T'}(c)$, the cochain
$\ell_k-\sum_i\omega^W_{(i)}(c_{k,i})-dh_k$ has values in $J^{[2]}$. By the definitions, $\omega'$ is
$\ell_1\cup\cdots\cup\ell_m$, with values in $J^{[m]}$, times the $0$-cochain $T\mapsto\sigma^U_T(\theta_U(z))$, which
differs from $\widetilde{\theta_W(z|_W)}$ by values in $J$, as both are lifts of $z|_W$. So $\omega'-(\ell_1\cup\cdots\cup
\ell_m)\cup\widetilde{\theta_W(z|_W)}$ has values in $J^{[m+1]}$, and the two classes are defined and equal
(Lemma~\ref{gal-realization-thickening-level-one-graded-forms}(2) and (3)).
Lemma~\ref{gal-realization-kappa-de-rham-transfer-local-near}(2), for $w=\theta_W(z|_W)$, gives the claim.

\emph{Step 2: the coefficients.} By Step 1 and Lemma~\ref{gal-realization-thickening-level-one-graded-forms}(2),
$\xi=[\sum_L\sum_J\omega^W_J(\theta_W(z_L|_W\prod_kp_{J_k}(u_{L_k}|_W)))]_m$. By
Lemma~\ref{gal-realization-kappa-de-rham-transfer-local-alternating}(2), $\xi=[\sum_I\omega^W_I(\theta_W(y'_I))]_m$ with
$y'_I=\sum_Lz_L|_W\sum_J\varepsilon(I,J)\prod_kp_{J_k}(u_{L_k}|_W)$. As $J=I\circ\pi$ runs through the sequences with
$\varepsilon(I,J)\neq0$, $\sum_J\varepsilon(I,J)\prod_kp_{J_k}(u_{L_k}|_W)=\sum_\pi\operatorname{sgn}(\pi)\prod_k
p_{I_{\pi(k)}}(u_{L_k}|_W)=\det(p_{I_k}(u_{L_l}|_W))_{k,l}$. By the uniqueness in
Lemma~\ref{gal-realization-kappa-de-rham-transfer-local-chart}(4), $y_I=y'_I$.
% Lean: CoefMemCmp.cmp_coef (BasicOpenRsLe.lean). Step 1 is CoefMemCmp.rep_one, through CupRep.cup_rep with ℓ_k = lc
% (lc_val, cup_lc, form_cup), the h_k of lift_diff (through GrForm.ev_res) and taylor_form (the Taylor property with the
% graded constancy); the factor σ^U(θ_U z) is replaced by σ^W(θ_W(z|_W)) through lift_diff and CocProd.prod_mem. Step 2 is
% rep_sum, rep_collapse (h4), sum_sgn (Matrix.det_apply) and the uniqueness of loc_coef (h7). About 240 lines.
\end{proof}
\leanref{Development/AutomorphicSpectrum/Proofs/BlockDecomposition/BasicOpenRsLe.lean}{GlobalTheory.CompanionTheory.ConeLocal.CoefMemCmp.rep_one}
\leanref{Development/AutomorphicSpectrum/Proofs/BlockDecomposition/BasicOpenRsLe.lean}{GlobalTheory.CompanionTheory.ConeLocal.CoefMemCmp.taylor_form}
\leanref{Development/AutomorphicSpectrum/Proofs/BlockDecomposition/BasicOpenRsLe.lean}{GlobalTheory.CompanionTheory.ConeLocal.CoefMemCmp.lift_diff}
\leanref{Development/AutomorphicSpectrum/Proofs/BlockDecomposition/BasicOpenRsLe.lean}{GlobalTheory.CompanionTheory.ConeLocal.CoefMemCmp.sum_sgn}

\begin{lemma}\label{gal-realization-kappa-de-rham-transfer-local-clear}
Let $\rho:\mathcal A\to\mathcal B$, $\lambda:\mathcal B\to\mathcal B_1$, $\iota:\mathcal A\to\mathcal A_1$ and
$\mu:\mathcal A_1\to\mathcal B_1$ be maps of commutative rings with $\mu\iota=\lambda\rho$, and $f\in\mathcal A$, such that
for every $a\in\mathcal A_1$ there are $b\in\mathcal A$ and $N\geq0$ with $\iota(f)^Na=\iota(b)$, and for every
$x\in\mathcal B$ with $\lambda(x)=0$ there is $N\geq0$ with $\rho(f)^Nx=0$. Let $D_1,\ldots,D_d$ be additive maps
$\mathcal B\to\mathcal B$ and $p_1,\ldots,p_d$ additive maps $\mathcal B_1\to\mathcal B_1$ with $p_k(xy)=x\,p_k(y)+
y\,p_k(x)$ and $p_k\lambda=\lambda D_k$. Let $\kappa=(\kappa_1,\ldots,\kappa_m)$ be a sequence in $\{1,\ldots,d\}$, and
$M_\kappa$ the additive subgroup of $\mathcal B$ generated by the elements $\rho(b)\det(D_{\kappa_k}(\rho(b_l)))_{1\leq
k,l\leq m}$ with $b,b_1,\ldots,b_m\in\mathcal A$. Let $x\in\mathcal B$, $\Lambda$ a finite set, and $z_L$ and
$u_{L,1},\ldots,u_{L,m}$ elements of $\mathcal A_1$, for $L\in\Lambda$, with
\[
 \lambda(x)=\sum_{L\in\Lambda}\mu(z_L)\det\bigl(p_{\kappa_k}(\mu(u_{L,l}))\bigr)_{1\leq k,l\leq m}.
\]
Then $\rho(f)^Mx\in M_\kappa$ for some $M\geq0$.
\end{lemma}
\leanref{Development/AutomorphicSpectrum/Proofs/BlockDecomposition/MinorGroupMul.lean}{GlobalTheory.CompanionTheory.ConeLocal.JacAlg.clear_sum}

\begin{proof}
$M_\kappa$ is stable under multiplication by the elements of $\rho(\mathcal A)$, as $\rho(a)\rho(b)=\rho(ab)$. Let
$F=\lambda(\rho(f))$.

\emph{Step 1: a column.} Let $a\in\mathcal A_1$, $b\in\mathcal A$ and $N\geq0$ with $\iota(f)^Na=\iota(b)$. Then for every
$k$
\[
 \lambda\bigl(\rho(f)^ND_k(\rho(b))-N\rho(f)^{N-1}\rho(b)D_k(\rho(f))\bigr)=F^{2N}p_k(\mu(a)),
\]
where the second term is $0$ for $N=0$. Indeed, $F^N\mu(a)=\mu(\iota(f)^Na)=\lambda(\rho(b))$, so
$\lambda(D_k(\rho(b)))=p_k(F^N\mu(a))=F^Np_k(\mu(a))+NF^{N-1}\mu(a)\lambda(D_k(\rho(f)))$ by the Leibniz rule, as
$p_k(F)=\lambda(D_k(\rho(f)))$; multiply by $F^N$.

\emph{Step 2: a determinant.} For $u_1,\ldots,u_m\in\mathcal A_1$ there are $M\geq0$ and $n\in M_\kappa$ with
$F^M\det(p_{\kappa_k}(\mu(u_l)))_{k,l}=\lambda(n)$. Indeed, there are $N$ and $b_l\in\mathcal A$ with
$\iota(f)^Nu_l=\iota(b_l)$ for every $l$, multiplying by powers of $f$ to make the exponents equal. By Step 1,
$F^{2Nm}\det(p_{\kappa_k}(\mu(u_l)))=\lambda(n)$, where $n$ is the determinant of the matrix with entries
$\rho(f^N)D_{\kappa_k}(\rho(b_l))-\rho(Nf^{N-1}b_l)D_{\kappa_k}(\rho(f))$. As the determinant is additive in each column,
$n$ is a sum of determinants whose column $l$ is $\rho(f^N)$ times $(D_{\kappa_k}(\rho(b_l)))_k$ or $-\rho(Nf^{N-1}b_l)$
times $(D_{\kappa_k}(\rho(f)))_k$; taking the factors out of the columns, each is $\rho(c)\det(D_{\kappa_k}(\rho(b'_l)))$
with $c\in\mathcal A$ and $b'_l\in\{b_l,f\}$, an element of $M_\kappa$.

\emph{Step 3: conclusion.} For $L\in\Lambda$ let $M_L$ and $n_L$ be as in Step 2 for $(u_{L,l})_l$, and $b_L\in\mathcal A$
and $N_L$ with $\iota(f)^{N_L}z_L=\iota(b_L)$, and let $K'\geq M_L+N_L$ for every $L$. Then
\[
 F^{K'}\lambda(x)=\sum_LF^{K'-M_L-N_L}\cdot F^{N_L}\mu(z_L)\cdot F^{M_L}\det\bigl(p_{\kappa_k}(\mu(u_{L,l}))\bigr)=\lambda(y),
 \qquad y=\sum_L\rho(f^{K'-M_L-N_L}b_L)\,n_L\in M_\kappa.
\]
So $\lambda(\rho(f)^{K'}x-y)=0$, and there is $N'$ with $\rho(f)^{N'}(\rho(f)^{K'}x-y)=0$; so $\rho(f)^{N'+K'}x=
\rho(f^{N'})y\in M_\kappa$.
% Lean: JacAlg.clear_sum (MinorGroupMul.lean), with minorGroup for M_κ and minorGroup_mul; Step 1 is column_clear (with
% leibniz_pow), Step 2 is det_clear with det_mem (Matrix.detRowAlternating.map_add_univ and Matrix.det_mul_column), and
% Step 3 is the end of clear_sum (hterm, hymem, h0). About 170 lines.
\end{proof}
\leanref{Development/AutomorphicSpectrum/Proofs/BlockDecomposition/MinorGroupMul.lean}{GlobalTheory.CompanionTheory.ConeLocal.JacAlg.column_clear}
\leanref{Development/AutomorphicSpectrum/Proofs/BlockDecomposition/MinorGroupMul.lean}{GlobalTheory.CompanionTheory.ConeLocal.JacAlg.det_clear}

\begin{lemma}\label{gal-realization-kappa-de-rham-transfer-local-minors}
Let $K$ be a field of characteristic $0$, $X$ a $K$-scheme smooth over $K$, $s\geq1$, $V\subset V'$ opens of $X$ with $V'$
affine and coordinates $t_1,\ldots,t_d$ on $V$ (Definition~\ref{gal-realization-thickening-coordinates}(1)),
$D_1,\ldots,D_d$ $K$-derivations of $\Gamma(V,\mathcal O_X)$ with $D_k(t_j)=1$ for $k=j$ and $0$ otherwise, and $B'$ the
image of restriction $\Gamma(V',\mathcal O_X)\to\Gamma(V,\mathcal O_X)$. For an open $U$ let $\mathcal
Q(U)=\mathrm{Th}^{(s)}_{\mathrm{dR}}(X)_U$. Let $(D,\theta)$ be as in
Lemma~\ref{gal-realization-kappa-de-rham-transfer-local-chart} for $V$, with maps $\sigma_T$ and forms $\omega_I$, let
$0\leq m<s$ and $\xi\in H^m(\mathcal Q(V'),\mathrm{gr}_m)$, and let $y_I\in\Gamma(V,\mathcal O_X)$, for the strictly
increasing sequences $I$ of length $m$ in $\{1,\ldots,d\}$, be such that the restriction of $\xi$ to $\mathcal Q(V)$ is
$[\sum_I\omega_I(\theta(y_I))]_m$. Then for every $I$, $y_I$ lies in the additive subgroup $\mathrm{Jac}_I(B')$ of
$\Gamma(V,\mathcal O_X)$ generated by the elements $b\det(D_{I_k}(b_l))_{1\leq k,l\leq m}$ with $b,b_1,\ldots,b_m\in B'$.
\end{lemma}
\leanref{Development/AutomorphicSpectrum/Proofs/BlockDecomposition/LocalFrameDerivationCompatStatement.lean}{GlobalTheory.T.local_frame_derivation_compat}
\leanref{Development/AutomorphicSpectrum/Proofs/BlockDecomposition/LocalFrameDerivationCompat.lean}{GlobalTheory.CompanionTheory.ConeLocal.coef_mem_jacSpan}

\begin{proof}
Fix $I$. Let $\mathfrak c$ be the set of $a\in\Gamma(V',\mathcal O_X)$ with $a|_Vy_I\in\mathrm{Jac}_I(B')$. As
$\mathrm{Jac}_I(B')$ is stable under multiplication by the elements of $B'$, $\mathfrak c$ is an ideal.

\emph{Step 1: local elements of $\mathfrak c$.} For every $x\in V'$ there are $f\in\Gamma(V',\mathcal O_X)$ with $x$ in the
open $U$ of $V'$ where $f$ is a unit, and $M\geq0$ with $f^M\in\mathfrak c$. Indeed, by
Lemma~\ref{gal-realization-kappa-de-rham-affine-basic}(1), for $V'$, there is such an $f$ for which $U$ has coordinates
$u_1,\ldots,u_e$; $U$ is affine with $\Gamma(U,\mathcal O_X)=\Gamma(V',\mathcal O_X)_f$ \cite[Tag 01HV]{Sta}. Let $W\subset
V$ be the open where $f|_V$ is a unit. As $V$ is affine, $W$ is affine with $\Gamma(W,\mathcal O_X)=\Gamma(V,\mathcal
O_X)_{f|_V}$, and $W\subset U$. The restrictions of the $t_k$ are coordinates on $W$, as a composite of \'etale maps is
\'etale \cite[Tag 00U2]{Sta}. Each $D_k$ extends to a $K$-derivation $p_k$ of $\Gamma(W,\mathcal O_X)$ with
$p_k(g|_W)=D_k(g)|_W$: the map $g\mapsto g|_W+D_k(g)|_W\varepsilon$ from $\Gamma(V,\mathcal O_X)$ to $\Gamma(W,\mathcal
O_X)[\varepsilon]/(\varepsilon^2)$ is a ring map that sends $f|_V$ to a unit, so it extends to $\Gamma(W,\mathcal O_X)$,
and $p_k$ is the coefficient of $\varepsilon$ of the extension. Then $p_k(t_j|_W)=1$ for $k=j$ and $0$ otherwise. Let
$(D_W,\theta_W)$ and $(D_U,\theta_U)$ be as in Lemma~\ref{gal-realization-kappa-de-rham-transfer-local-chart} for $W$ and
$U$. By its part (4), the restrictions of $\xi$ to $\mathcal Q(U)$ and to $\mathcal Q(W)$ are
$[\sum_L\omega^U_L(\theta_U(z_L))]_m$ and $[\sum_I\omega^W_I(\theta_W(y^W_I))]_m$ for some $z_L\in\Gamma(U,\mathcal O_X)$
and $y^W_I\in\Gamma(W,\mathcal O_X)$, and restriction from $\mathcal Q(V')$ to $\mathcal Q(W)$ is restriction to $\mathcal
Q(V)$, or to $\mathcal Q(U)$, followed by restriction to $\mathcal Q(W)$. By
Lemma~\ref{gal-realization-kappa-de-rham-transfer-local-compare} for $W\subset V$, $y^W_I=\sum_Ly_L|_W\det(p_{I_k}(t_{L_l}|_W))
=y_I|_W$, as $p_{I_k}(t_{L_l}|_W)$ is $1$ if $I_k=L_l$ and $0$ otherwise, and this matrix has determinant $1$ for $L=I$
and a zero row for $L\neq I$. By the same lemma for $W\subset U$, $y^W_I=\sum_Lz_L|_W\det(p_{I_k}(u_{L_l}|_W))$. So
\[
 y_I|_W=\sum_Lz_L|_W\det\bigl(p_{I_k}(u_{L_l}|_W)\bigr)_{k,l}.
\]
Apply Lemma~\ref{gal-realization-kappa-de-rham-transfer-local-clear} to $\mathcal A=\Gamma(V',\mathcal O_X)$, $\mathcal
B=\Gamma(V,\mathcal O_X)$, $\mathcal A_1=\Gamma(U,\mathcal O_X)$, $\mathcal B_1=\Gamma(W,\mathcal O_X)$, the restriction
maps, $f$, the $D_k$, the $p_k$, $\kappa=I$ and $x=y_I$: every element of $\Gamma(U,\mathcal O_X)=\Gamma(V',\mathcal
O_X)_f$ is $b/f^N$, an element of $\Gamma(V,\mathcal O_X)$ with restriction $0$ to $W$ is killed by a power of $f|_V$, and
$M_I=\mathrm{Jac}_I(B')$. So $f|_V^My_I\in\mathrm{Jac}_I(B')$ for some $M$, that is, $f^M\in\mathfrak c$.

\emph{Step 2: conclusion.} If $\mathfrak c\neq\Gamma(V',\mathcal O_X)$, it lies in a maximal ideal $\mathfrak m$; let
$x\in V'$ be the point of $\mathfrak m$, as $V'$ is affine. With $f$ and $M$ as in Step 1, $f\notin\mathfrak m$, as $f$ is
a unit at $x$, while $f^M\in\mathfrak c\subset\mathfrak m$, which is impossible. So $1\in\mathfrak c$, that is,
$y_I\in\mathrm{Jac}_I(B')$.
% Lean: GlobalTheory.T.local_frame_derivation_compat = ConeLocal.coef_mem_jacSpan (LocalFrameDerivationCompat.lean):
% c is coefIdeal (rs_mul_mem_jacSpan), Step 2 is ideal_eq_top_of_local. Step 1 is coord_near (CoordNear.lean, as in Lemma
% gal-realization-kappa-de-rham-affine-basic(1)) and CoefMemLoc.main (BasicOpenRsLe.lean): W = X.basicOpen (rs hVV f),
% CoordRes.main (coordinates on W, through etale_rs), DerivExt.main (the p_k, through exists_derivation_ext with
% TrivSqZeroExt), exists_loc for W and U, loc_coef, cmp_coef twice (e1, e2), sum_delta and det_delta (e3),
% exists_eq_pow_mul_of_isAffineOpen, exists_pow_mul_eq_zero_of_res_basicOpen_eq_zero_of_isAffineOpen and JacAlg.clear_sum.
% About 200 lines beyond Lemmas gal-realization-kappa-de-rham-transfer-local-compare and -clear.
\end{proof}
\leanref{Development/AutomorphicSpectrum/Proofs/BlockDecomposition/BasicOpenRsLe.lean}{GlobalTheory.CompanionTheory.ConeLocal.CoefMemLoc.main}
\leanref{Development/AutomorphicSpectrum/Proofs/BlockDecomposition/LocalFrameDerivationCompat.lean}{GlobalTheory.CompanionTheory.ConeLocal.ideal_eq_top_of_local}
\leanref{Development/AutomorphicSpectrum/Proofs/BlockDecomposition/ExistsDerivationExt.lean}{GlobalTheory.CompanionTheory.ConeLocal.DerivExt.exists_derivation_ext}
\leanref{Development/AutomorphicSpectrum/Proofs/BlockDecomposition/CoordNear.lean}{GlobalTheory.CompanionTheory.ConeLocal.coord_near}

\begin{lemma}\label{gal-realization-kappa-de-rham-transfer-local}
Let $K$ be a field of characteristic $0$, $n\geq0$, $X$ a $K$-scheme smooth of relative dimension $n+1$, $s>n+1$, $V\subset
X$ an open with coordinates $t_0,\ldots,t_n$ (Definition~\ref{gal-realization-thickening-coordinates}(1)) and
$D_0,\ldots,D_n$ $K$-derivations of $\Gamma(V,\mathcal O_X)$ with $D_k(t_j)=1$ for $k=j$ and $0$ otherwise. For an open
$U$ let $\mathcal Q(U)=\mathrm{Th}^{(s)}_{\mathrm{dR}}(X)_U$ (Definition~\ref{gal-realization-kappa-de-rham-truncated}).
For every thickening $T$ of $\mathcal Q(V)$ and $a\in\Gamma(V,\mathcal O_X)$ fix a lift $\hat a(T)$ of $a$ at $T$
(Definition~\ref{gal-realization-thickening-level-one-fibre}(2)). For $x\in\Gamma(V,\mathcal O_X)$ let $\theta(x)$ be the
$(n+1)$-cochain of $\mathcal Q(V)$ that sends $c_0\to\cdots\to c_{n+1}$ to
$\hat x(c_0)\prod_{k=0}^n(\hat t_k(c_{k+1})-\hat t_k(c_k))$, with all factors carried to $B_{c_{n+1}}$; it is the cup
product of the $0$-cochain $\hat x$ and the coboundaries of the $0$-cochains $\hat t_k$. As the coboundaries of $\hat x$
and of the $\hat t_k$ have values in $J$ (Lemma~\ref{gal-realization-thickening-level-one-forms}(1)), and a product of $m$
elements of $J$ lies in $J^{[m]}$, the values of $\theta(x)$ lie in $J^{[n+1]}$ and those of its coboundary in
$J^{[n+2]}$; say that $x$ \emph{represents} $[\theta(x)]_{n+1}\in H^{n+1}(\mathcal Q(V),\mathrm{gr}_{n+1})$. Let $V\subset V'$ with $V'$ affine, and $B'$ the image of restriction
$\Gamma(V',\mathcal O_X)\to\Gamma(V,\mathcal O_X)$, with $\mathrm{Jac}$ and $\mathrm{Jac}_k$ as in
Definition~\ref{gal-realization-kappa-de-rham-torus}(2), for $R=\Gamma(V,\mathcal O_X)$.
\begin{enumerate}
\item Every class of $H^{n+1}(\mathcal Q(V),\mathrm{gr}_{n+1})$ is represented by exactly one $x\in\Gamma(V,\mathcal
  O_X)$, and if $x$ and $x'$ represent $h$ and $h'$, then $x+x'$ represents $h+h'$.
\item Let $g_0,\ldots,g_n$ be units of $\Gamma(V',\mathcal O_X)$, $\hat g_k$ a family of lifts of $g_k$ on $\mathcal
  Q(V')$ (Definition~\ref{gal-realization-kappa-de-rham-cone}(3)), and $c'=[\mathrm{dlog}_s(\hat g_0)\cup\cdots\cup
  \mathrm{dlog}_s(\hat g_n)]_{n+1}$ (Lemmas~\ref{gal-realization-kappa-de-rham-transfer-lifts}(1)
  and~\ref{gal-realization-kappa-de-rham-cup}). Then there is $\varepsilon\in\{1,-1\}$ such that the restriction of $c'$
  to $\mathcal Q(V)$ is represented by $\varepsilon(\prod_kg_k|_V)^{-1}\det(D_k(g_l|_V))_{0\leq k,l\leq n}$.
\item For every $\omega\in H^{n+1}(\mathcal Q(V'),\mathrm{gr}_{n+1})$, the restriction of $\omega$ to $\mathcal Q(V)$ is
  represented by an element of $\mathrm{Jac}(B')$.
\item Let $f$ be an $n$-cochain of $\mathcal Q(V')$ with values in $J^{[n]}$ whose coboundary has values in $J^{[n+1]}$.
  Then there are $\eta_k\in\mathrm{Jac}_k(B')$, $0\leq k\leq n$, such that the restriction of $[df]_{n+1}$ to $\mathcal
  Q(V)$ is represented by $\sum_k(-1)^kD_k(\eta_k)$.
\end{enumerate}
\end{lemma}
\leanref{Development/WeilBound/CoreStatements/PadicHodgeComparison/EtaleCohomologyComparison/PerfectoidPullbacks/ConeSectionRestriction/ZeroC.lean}{CrystallineSemistableComparison.Pb.xdtC}
\leanref{Development/AutomorphicSpectrum/Proofs/BlockDecomposition/ExistsUniqueTopGradedRep.lean}{GlobalTheory.CompanionTheory.ConeTopChart.top_chart}
\leanref{Development/AutomorphicSpectrum/Proofs/BlockDecomposition/TopGradedRepRestrict.lean}{GlobalTheory.CompanionTheory.ConeLdlog.main}
\leanref{Development/AutomorphicSpectrum/Proofs/BlockDecomposition/TopGradedRepDerivationCompat.lean}{GlobalTheory.CompanionTheory.ConeLRestr.main}
\leanref{Development/AutomorphicSpectrum/Proofs/BlockDecomposition/TopGradedRepOpenCompat.lean}{GlobalTheory.CompanionTheory.ConeLocal.LExact.main}

\begin{proof}
Let $\mathcal Q=\mathcal Q(V)$ and $G=\Gamma(V,\mathcal O_X)$. The families $\mathcal O^{[r]}$ form a filtration of
$\Phi_{\mathcal Q}$, by the definition of the ideals $J^{[r]}$ (Definition~\ref{gal-realization-thickening-nilpotent}(3)).

(1) \emph{Step 1: the chart.} Let $(D,\theta)$ be as in Lemma~\ref{gal-realization-kappa-de-rham-transfer-local-chart}
for $V$, which is affine as it has coordinates, with the coordinates indexed by $k\in\{0,\ldots,n\}$, $t'_k=\theta(t_k)$
and $\tau_k(\sigma,\sigma')=\sigma'(t'_k)-\sigma(t'_k)$. Choose maps $\sigma_T:D\to T$, and let $a_k$ and $\omega_I$ be as
in Definition~\ref{gal-realization-thickening-level-one-graded-filtration}(3), for the filtration $(\mathcal O^{[r]})$.

\emph{Step 2: the forms.} Let $I_0=(0,1,\ldots,n)$, the only strictly increasing sequence of length $n+1$ in
$\{0,\ldots,n\}$. As $n+1<s$, every class $h$ of $H^{n+1}(\mathcal Q,\mathrm{gr}_{n+1})$ is
$[\omega_{I_0}(\theta(y))]_{n+1}$ for exactly one $y\in G$
(Lemma~\ref{gal-realization-kappa-de-rham-transfer-local-chart}(4)).

\emph{Step 3: comparison of $\theta(x)$ with $\omega_{I_0}(\theta(x))$.} Let $x\in G$. There is an $n$-cochain $G_x$ with
values in $J^{[n+1]}$ such that $\theta(x)-\omega_{I_0}(\theta(x))-dG_x$ has values in $J^{[n+2]}$. Indeed, let $\ell_k$ be
the coboundary of the $0$-cochain $\hat t_k$, and $h_k(T)=\hat t_k(T)-\sigma_T(t'_k)$, which lies in $J_T$, as both terms
are lifts of $t_k$. For $\varphi:T\to T'$, $a_k(\varphi)=\sigma_{T'}(t'_k)-\varphi(\sigma_T(t'_k))$, so
$\ell_k=\omega_{(k)}(1)+dh_k$, as $\omega_{(i)}(c)(\varphi)=a_i(\varphi)\sigma_{T'}(c)$. By
Lemma~\ref{gal-realization-kappa-de-rham-transfer-local-near}(1), for $c_{k,i}=1$ if $i=k$ and $0$ otherwise, there is an
$n$-cochain $H$ with values in $J^{[n+1]}$ such that $\ell_0\cup\cdots\cup\ell_n-\omega_{I_0}(1)-dH$ has values in
$J^{[n+2]}$, as $\prod_kc_{k,I_k}$ is $1$ for $I=I_0$ and $0$ otherwise. Also $\omega_{I_0}(\theta(x))$ is
$\omega_{I_0}(1)$ multiplied, at a chain $c_0\to\cdots\to c_{n+1}$, by $Y=\sigma_{c_{n+1}}(\theta(x))$. The value of
$\theta(x)$ at the chain is $\hat x(c_0)$,
carried to $c_{n+1}$, times the value of $\ell_0\cup\cdots\cup\ell_n$, which lies in $J^{[n+1]}$; and $\hat x(c_0)$ and $Y$
are lifts of $x$ at $c_{n+1}$, so they differ by an element of $J$. So, modulo $J^{[n+2]}$, $\theta(x)$ is
$\omega_{I_0}(\theta(x))+Y\,dH$. Let $G_x$ send $c_0\to\cdots\to c_n$ to $H(c_0\to\cdots\to c_n)\,\sigma_{c_n}(\theta(x))$.
At a chain ending at $c_{n+1}$, the terms of $dG_x$ are those of $dH$ multiplied by $Y$, except the term of the face that
omits $c_{n+1}$, which is multiplied by the image of $\sigma_{c_n}(\theta(x))$ instead; these two images of $\theta(x)$
under maps $D\to c_{n+1}$ differ by an element of $J$, as $D$ is graded constant. So $Y\,dH-dG_x$ has values in
$J^{[n+1]}J\subset J^{[n+2]}$.

\emph{Step 4: conclusion.} By Step 3 and
Lemma~\ref{gal-realization-thickening-level-one-graded-forms}(2) and~(3), $[\theta(x)]_{n+1}=[\omega_{I_0}(\theta(x))]_{n+1}
+[dG_x]_{n+1}$, and $[dG_x]_{n+1}$ is the class of the coboundary of the cochain of the classes of the values of $G_x$, so
it is $0$. So $x$ represents $h$ if and only if $[\omega_{I_0}(\theta(x))]_{n+1}=h$, and by Step 2 every class is
represented by exactly one $x$. For $x,x'\in G$, $\widehat{x+x'}(c_0)-\hat x(c_0)-\hat x'(c_0)\in J$, so
$\theta(x+x')-\theta(x)-\theta(x')$ has values in $J\,J^{[n+1]}\subset J^{[n+2]}$, and $x+x'$ represents $h+h'$ by
Lemma~\ref{gal-realization-thickening-level-one-graded-forms}(2).

(2) \emph{Step 1: the derivations.} Let $p_k=\theta D_k\theta^{-1}$, a map $B_D\to B_D$. By
Lemma~\ref{gal-realization-kappa-de-rham-transfer-local-chart}(5), the $p_k$ satisfy the hypothesis of
Lemma~\ref{gal-realization-kappa-de-rham-transfer-local-connecting} for $\Phi_{\mathcal Q}$, the filtration $(\mathcal
O^{[r]})$ and $(D,\tau)$. Every $1+y$ with $y\in J_T$ is a unit, as $J_T$ is nilpotent (proof of
Lemma~\ref{gal-realization-kappa-de-rham-lift}).

\emph{Step 2: the logarithmic derivatives.} Fix $k$, let $g=g_k|_V$, a unit of $G$, and $\upsilon=\theta(g)$, a unit of
$B_D$. The thickenings $T$ of $\mathcal Q$ lie in $\mathcal Q(V')$, and $\hat g_k(T)$ and $\sigma_T(\upsilon)$ are lifts
of $g$ at $T$ (Lemmas~\ref{gal-realization-thickening-level-one-forms}(1)
and~\ref{gal-realization-kappa-de-rham-transfer-local-chart}(2)), so $\hat g_k(T)-\sigma_T(\upsilon)\in J_T$. By
Lemma~\ref{gal-realization-kappa-de-rham-transfer-local-dlog}, applied to the unit $\upsilon$ of $B_D$ and the elements
$\hat g_k(T)$, there is a $0$-cochain $h_k$ of $\mathcal Q$ with values in $J$ such that for every $\varphi:T\to T'$,
$r(\varphi)=\hat g_k(T')\varphi(\hat g_k(T))^{-1}$ has $r(\varphi)-1\in J_{T'}$ and
$r(\varphi)-1-\sum_i\omega_{(i)}(c_{k,i})(\varphi)-(dh_k)(\varphi)\in J^{[2]}_{T'}$, where
$c_{k,i}=\upsilon^{-1}p_i(\upsilon)=\theta(g^{-1}D_i(g))$, as $\omega_{(i)}(c)(\varphi)=a_i(\varphi)\sigma_{T'}(c)$. For $y\in
J_{T'}$, $\log_s(y)-y=\sum_{1\leq j<s}(-1)^jj!\,\gamma_{j+1}(y)\in J^{[2]}_{T'}$
(Lemma~\ref{gal-realization-kappa-de-rham-log}). So the restriction $\ell_k$ of $\mathrm{dlog}_s(\hat g_k)$ to $\mathcal
Q$, which sends $\varphi$ to $\log_s(r(\varphi)-1)$, has $\ell_k-\sum_i\omega_{(i)}(c_{k,i})-dh_k$ with values in
$J^{[2]}$.

\emph{Step 3: the cup product.} Restriction from $\mathcal Q(V')$ to $\mathcal Q$ commutes with cup products, so the
restriction of $c'$ is $[\ell_0\cup\cdots\cup\ell_n]_{n+1}=[(\ell_0\cup\cdots\cup\ell_n)\cup\tilde1]_{n+1}$, where
$\tilde1$ is the constant $0$-cochain $1$. By Lemma~\ref{gal-realization-kappa-de-rham-transfer-local-near}(2), for $w=1$
and the $c_{k,i}$ and $h_k$ of Step 2, it is $[\sum_J\omega_J(\prod_kc_{k,J_k})]_{n+1}$, the sum over all sequences $J$ of
length $n+1$ in $\{0,\ldots,n\}$.

\emph{Step 4: the determinant.} As $I_0$ is the only strictly increasing sequence of length $n+1$,
Lemma~\ref{gal-realization-kappa-de-rham-transfer-local-alternating}(2) shows that the restriction of $c'$ is
$[\omega_{I_0}(\sum_J\varepsilon(I_0,J)\prod_kc_{k,J_k})]_{n+1}$. Here $\varepsilon(I_0,J)=\operatorname{sgn}(\pi)$ if
$J=(\pi(0),\ldots,\pi(n))$ for a permutation $\pi$, and $0$ otherwise, so $\sum_J\varepsilon(I_0,J)\prod_kc_{k,J_k}=
\det(c_{k,i})_{k,i}=\theta\bigl(\det(g_k^{-1}D_i(g_k))_{k,i}\bigr)$, with the $g_k$ restricted to $V$. Taking the factor
$g_k^{-1}$ out of row $k$ and transposing, $\det(g_k^{-1}D_i(g_k))_{k,i}=(\prod_kg_k)^{-1}\det(D_k(g_l))_{k,l}$.

\emph{Step 5: conclusion.} By Step 4 and Step 4 of the proof of (1), $(\prod_kg_k|_V)^{-1}\det(D_k(g_l|_V))$ represents
the restriction of $c'$ to $\mathcal Q$. So (2) holds, with $\varepsilon=1$.

(3) Let $\omega\in H^{n+1}(\mathcal Q(V'),\mathrm{gr}_{n+1})$. By Step 2 of the proof of (1), its restriction to $\mathcal Q$
is $[\omega_{I_0}(\theta(y))]_{n+1}$ for some $y\in G$. By Lemma~\ref{gal-realization-kappa-de-rham-transfer-local-minors},
for $m=n+1<s$, $y\in\mathrm{Jac}_{I_0}(B')$, which is $\mathrm{Jac}(B')$
(Definition~\ref{gal-realization-kappa-de-rham-torus}(2)). By Step 4 of the proof of (1), $y$ represents the restriction
of $\omega$.

(4) The restriction $f'$ of $f$ to $\mathcal Q$ has values in $J^{[n]}$ and coboundary with values in $J^{[n+1]}$, and the
restriction of $[df]_{n+1}$ is $[df']_{n+1}$. By Lemma~\ref{gal-realization-kappa-de-rham-transfer-local-chart}(4), for
$m=n<s$, $[f']_n=[\sum_I\omega_I(\theta(y_I))]_n$ for some $y_I\in G$, indexed by the strictly increasing sequences $I$ of
length $n$ in $\{0,\ldots,n\}$; $[f']_n$ is the restriction of $[f]_n\in H^n(\mathcal Q(V'),\mathrm{gr}_n)$, so
$y_I\in\mathrm{Jac}_I(B')$ by Lemma~\ref{gal-realization-kappa-de-rham-transfer-local-minors}, for $m=n$. By
Lemma~\ref{gal-realization-kappa-de-rham-transfer-local-connecting}(2), with the $p_k$ of Step 1 of (2),
$[df']_{n+1}=[\sum_I\sum_k\omega_{(I,k)}((-1)^n\theta(D_k(y_I)))]_{n+1}$. By
Lemma~\ref{gal-realization-kappa-de-rham-transfer-local-alternating}(2), this is $[\omega_{I_0}(\theta(z))]_{n+1}$ with
$z=\sum_I\sum_k\varepsilon(I_0,(I,k))(-1)^nD_k(y_I)$. For strictly increasing $I$, $(I,k)$ is a permutation of
$(0,\ldots,n)$ only if $I$ is the sequence $I_k$ of the elements of $\{0,\ldots,n\}\smallsetminus\{k\}$ in increasing
order; let $\varepsilon_k=\varepsilon(I_0,(I_k,k))\in\{1,-1\}$. So $z=\sum_k(-1)^kD_k(\eta_k)$ with
$\eta_k=(-1)^{k+n}\varepsilon_ky_{I_k}$, which lies in $\mathrm{Jac}_{I_k}(B')=\mathrm{Jac}_k(B')$
(Definition~\ref{gal-realization-kappa-de-rham-torus}(2)). By Step 4 of the proof of (1), $z$ represents the restriction
of $[df]_{n+1}$.
% Lean for (1): ConeTopChart.top_chart. Step 1 is ConeLocal.exists_loc (Lemma
% gal-realization-kappa-de-rham-transfer-local-chart(1) to (3)); Step 2 is ConeLocal.loc_coef (part (4) of that lemma)
% with top_index_eq and uniqForm. Step 3 is key: addC_near (ℓ_k = ω_{(k)}(1) + d h_k), CupRep.near (Lemma
% gal-realization-kappa-de-rham-transfer-local-near(1), here with c the identity matrix), hsum, form_mul, lift_sub and
% CupRep.d_mul. Step 4 is toForm with rep_shift, and the additivity is ConeLocal.rep_xdtC_add (TopRepresentativeAdd.lean).
% Lean for (2): ConeLdlog.main (TopGradedRepRestrict.lean). Step 1 is ConeLocal.taylor_loc (hT) and hun; Step 2 is
% DlogApprox.main (LocalFrameTaylorClassEq.lean: hgl through GrForm.ev_res, frame_main, and plogK_sub_mem for
% log_s(y) − y); Step 3 is rep_res, cup_resQ (h2) and CupRep.cup_rep (h3); Step 4 is ConeLocal.rep_collapse (h4), rep_one
% (h5), sgn_det, inv_prod, Matrix.det_mul_row and map_det (e); Step 5 is xdt_rep, which redoes Step 3 of the proof of (1)
% through rep_back and CupRep.near. The Lean's conclusion is the disjunction of the two signs, and it proves the left one.
% Lean for (3): ConeLRestr.main (TopGradedRepDerivationCompat.lean): local_graded_class_representative = loc_coef,
% local_frame_derivation_compat (Lemma gal-realization-kappa-de-rham-transfer-local-minors), rep_top, and rep_xdt, which
% redoes Step 3 of the proof of (1) through rep_cup and rep_zero_cup.
% Lean for (4): ConeLocal.LExact.main (TopGradedRepOpenCompat.lean): rep_exists and rep_res for [f]_n, loc_coef for m = n,
% local_frame_derivation_compat, d_resQ, taylor_differential_representative (Lemma
% gal-realization-kappa-de-rham-transfer-local-connecting(2)), form_sum_representative_lift = rep_collapse, rep_one, the
% sequences Ik k = succAbove k with sgn_snoc (InvProd.lean, through SnocPerm.eq_succAbove_of_snoc_eq_perm), and rep_xdt,
% again a version of Step 3 of the proof of (1). About 800 lines in the four main files, the rest in the files of Lemmas
% gal-realization-kappa-de-rham-transfer-local-near to -minors.
% Lean overview: (1) is ConeTopChart.top_chart (ExistsUniqueTopGradedRep.lean, through ConeLocal.exists_loc and loc_coef) with
% rep_add' = ConeLocal.rep_xdtC_add (TopRepresentativeAdd.lean); (2) is ConeLdlog.main (TopGradedRepRestrict.lean); (3)
% is ConeLRestr.main (TopGradedRepDerivationCompat.lean); (4) is ConeLocal.LExact.main (TopGradedRepOpenCompat.lean). All
% take KH.LocGr, KH.PolyTf and KH.GrAff (Lemmas gal-realization-thickening-level-one-graded,
% gal-realization-kappa-de-rham-poly and gal-realization-kappa-de-rham-affine) as inputs. (1) also follows from Lemma
% gal-realization-kappa-de-rham-local(2), with W = V and r = n + 1, as θ(x) and the cochain of Definition
% gal-realization-thickening-level-one-fibre(4) for the lifts û = t̂ differ by values in J^[n+2]; the Lean does not argue
% this way.
\end{proof}
\leanref{Development/AutomorphicSpectrum/Proofs/BlockDecomposition/ExistsUniqueTopGradedRep.lean}{GlobalTheory.CompanionTheory.ConeTopChart.key}
\leanref{Development/AutomorphicSpectrum/Proofs/BlockDecomposition/ExistsLocalFrameOfCoords.lean}{GlobalTheory.CompanionTheory.ConeLocal.exists_loc}
\leanref{Development/AutomorphicSpectrum/Proofs/BlockDecomposition/LocalGradedClassRepresentative.lean}{GlobalTheory.CompanionTheory.ConeLocal.loc_coef}
\leanref{Development/AutomorphicSpectrum/Proofs/BlockDecomposition/CupListRepresentativeOfForms.lean}{GlobalTheory.CompanionTheory.CupRep.near}
\leanref{Development/AutomorphicSpectrum/Proofs/BlockDecomposition/CupListRepresentativeOfForms.lean}{GlobalTheory.CompanionTheory.CupRep.d_mul}
\leanref{Development/AutomorphicSpectrum/Proofs/BlockDecomposition/TopRepresentativeAdd.lean}{GlobalTheory.CompanionTheory.ConeLocal.rep_xdtC_add}
\leanref{Development/AutomorphicSpectrum/Proofs/BlockDecomposition/LocalFrameTaylorClassEq.lean}{GlobalTheory.CompanionTheory.ConeLocal.DlogApprox.main}
\leanref{Development/AutomorphicSpectrum/Proofs/BlockDecomposition/TopGradedRepRestrict.lean}{GlobalTheory.CompanionTheory.ConeLdlog.sgn_det}
\leanref{Development/AutomorphicSpectrum/Proofs/BlockDecomposition/TopGradedRepRestrict.lean}{GlobalTheory.CompanionTheory.ConeLdlog.xdt_rep}
\leanref{Development/AutomorphicSpectrum/Proofs/BlockDecomposition/InvProd.lean}{GlobalTheory.CompanionTheory.ConeLocal.sgn_snoc}

\begin{lemma}\label{gal-realization-kappa-de-rham-transfer}
Let $K$ be a field of characteristic $0$, $n\geq0$, $X$ a $K$-scheme smooth of relative dimension $n+1$, and $i$ and
$y_0,\ldots,y_{n+1}$ as in Definition~\ref{gal-realization-kappa-de-rham-cone}, with $s=n+2$, such that $U_S$ is affine for
every nonempty $S\subset\top=\{0,\ldots,n+1\}$; let $S_j=\top\smallsetminus\{j\}$. Let $V\subset U_\top$ be an open with
coordinates $t_0,\ldots,t_n$ (Definition~\ref{gal-realization-thickening-coordinates}(1)), $D_0,\ldots,D_n$
$K$-derivations of $\Gamma(V,\mathcal O_X)$ with $D_k(t_j)=1$ for $k=j$ and $0$ otherwise, and $B'(S)$ the image of
restriction $\Gamma(U_S,\mathcal O_X)\to\Gamma(V,\mathcal O_X)$. Let $g_0,\ldots,g_n$ be units of $\Gamma(U_\top,\mathcal
O_X)$, $\hat g_k$ a family of lifts of $g_k$ on $\mathcal W_\top$, and $c_0\in H^{n+1}(\mathcal W_\top,\mathrm{gr}_{n+1})$
such that $z_0=\mathrm{dlog}_s(\hat g_0)\cup\cdots\cup\mathrm{dlog}_s(\hat g_n)$ has values in $J^{[n+1]}$ and
coboundary with values in $J^{[n+2]}$, and $c_0=[z_0]_{n+1}$
(Definition~\ref{gal-realization-thickening-level-one-graded-filtration}(4), for the filtration
$(\mathcal O^{[r]})$). Let $\omega_j\in H^{n+1}(\mathcal W_{S_j},\mathrm{gr}_{n+1})$ for $j\in\top$, and $f$ an
$n$-cochain of $\mathcal W_\top$ with values in $J^{[n]}$ whose coboundary has values in $J^{[n+1]}$, such that
$[df]_{n+1}=c_0-\sum_j\omega_j|_{\mathcal W_\top}$. Then there are $x_j\in\mathrm{Jac}(B'(S_j))$ for $j\in\top$,
$\eta_k\in\mathrm{Jac}_k(B'(\top))$ for $0\leq k\leq n$ (Definition~\ref{gal-realization-kappa-de-rham-torus}(2), for
$R=\Gamma(V,\mathcal O_X)$) and $\varepsilon\in\{1,-1\}$ such that
\[
 \varepsilon\Bigl(\prod_{k=0}^ng_k|_V\Bigr)^{-1}\det\bigl(D_k(g_l|_V)\bigr)_{0\leq k,l\leq n}
 =\sum_{j\in\top}x_j+\sum_{k=0}^n(-1)^kD_k(\eta_k).
\]
\end{lemma}
\leanref{Development/WeilBound/CoreStatements/PadicHodgeComparison/EtaleCohomologyComparison/PerfectoidPullbacks/ConeSectionRestriction/ConeTransfer.lean}{CrystallineSemistableComparison.Pb.cone_transfer}
\leanref{Development/AutomorphicSpectrum/Proofs/BlockDecomposition/ProjectiveChartAffineCoordsStatement.lean}{GlobalTheory.T.projective_chart_affine_coords}

\begin{proof}
Let $\mathcal Q(U)=\mathrm{Th}^{(s)}_{\mathrm{dR}}(X)_U$ for an open $U$, and $\mathrm{gr}=\mathrm{gr}_{n+1}$. Every
thickening of $\mathcal W_S$ lies in $\mathcal Q(U_S)$ (Lemma~\ref{gal-realization-kappa-de-rham-cone-vanishing}(1)), and
$\mathcal Q(V)\subset\mathcal Q(U_\top)\subset\mathcal Q(U_{S_j})$.

\emph{Step 1: comparison.} By Lemma~\ref{gal-realization-kappa-de-rham-cone-vanishing}(1), restriction from $\mathcal
Q(U_S)$ to $\mathcal W_S$ is bijective on $H^{n+1}(-,\mathrm{gr})$ for $S=\top$ and $S=S_j$, and on
$H^n(-,\mathrm{gr}_n)$ for $S=\top$.

\emph{Step 2: the class $c'$.} For a thickening $T$ of $\mathcal Q(U_\top)$ choose lifts $\hat g'_k(T)$ of the $g_k$. They
are units, as $J_T$ is nilpotent, and the families $\hat g'_k$ are compatible in the sense of
Lemma~\ref{gal-realization-kappa-de-rham-transfer-lifts}, by
Lemma~\ref{gal-realization-thickening-level-one-forms}(1). By
Lemma~\ref{gal-realization-kappa-de-rham-transfer-lifts}(1) and Lemma~\ref{gal-realization-kappa-de-rham-cup}, for the
filtration $(\mathcal O^{[r]})$,
$c'=[\mathrm{dlog}_s(\hat g'_0)\cup\cdots\cup\mathrm{dlog}_s(\hat g'_n)]_{n+1}\in H^{n+1}(\mathcal Q(U_\top),\mathrm{gr})$ is
defined.

\emph{Step 3: $c'$ restricts to $c_0$.} Restriction to $\mathcal W_\top$ commutes with cup products and with
$\mathrm{dlog}_s$, so the restriction of $c'$ is the class of the cup product of the $\mathrm{dlog}_s$ of the restricted
families. On $\mathcal W_\top$ these and the $\hat g_k$ are compatible families of lifts of the same $g_k$, which differ by
elements of $J$; so by Lemma~\ref{gal-realization-kappa-de-rham-transfer-lifts}(2), $c'|_{\mathcal W_\top}=c_0$.

\emph{Step 4: lifting $\omega_j$ and $f$.} By Step 1 there are $\omega'_j\in H^{n+1}(\mathcal Q(U_{S_j}),\mathrm{gr})$ with
$\omega'_j|_{\mathcal W_{S_j}}=\omega_j$. By Step 1 and Lemma~\ref{gal-realization-kappa-de-rham-transfer-extend}, for
$\mathcal V=\mathcal W_\top\subset\mathcal U=\mathcal Q(U_\top)$, there is an $n$-cochain $f'$ of $\mathcal Q(U_\top)$ with
values in $J^{[n]}$ and coboundary with values in $J^{[n+1]}$, such that $x'=[df']_{n+1}$ has $x'|_{\mathcal
W_\top}=[df]_{n+1}$.

\emph{Step 5: the relation on $U_\top$.} As restrictions compose, $\omega'_j|_{\mathcal Q(U_\top)}$ restricts to
$\omega_j|_{\mathcal W_\top}$ on $\mathcal W_\top$. So $c'-\sum_j\omega'_j|_{\mathcal Q(U_\top)}-x'$ restricts to
$c_0-\sum_j\omega_j|_{\mathcal W_\top}-[df]_{n+1}=0$. Restriction from $\mathcal Q(U_\top)$ to $\mathcal W_\top$ is
injective (Step 1), so $c'=\sum_j\omega'_j|_{\mathcal Q(U_\top)}+x'$, and restricting to $\mathcal Q(V)$,
\[
 c'|_{\mathcal Q(V)}=\sum_j\omega'_j|_{\mathcal Q(V)}+x'|_{\mathcal Q(V)}.
\]

\emph{Step 6: the identity.} Apply Lemma~\ref{gal-realization-kappa-de-rham-transfer-local}, with $s=n+2$ and the $t_k$
and $D_k$. By its part (2), for $V'=U_\top$, $c'|_{\mathcal Q(V)}$ is represented by
$\varepsilon(\prod_kg_k|_V)^{-1}\det(D_k(g_l|_V))$ for some $\varepsilon\in\{1,-1\}$. By its part (3), for $V'=U_{S_j}$,
$\omega'_j|_{\mathcal Q(V)}$ is represented by some $x_j\in\mathrm{Jac}(B'(S_j))$. By its part (4), for $V'=U_\top$ and
$f'$, $x'|_{\mathcal Q(V)}$ is represented by $\sum_k(-1)^kD_k(\eta_k)$ for some $\eta_k\in\mathrm{Jac}_k(B'(\top))$. By its
part (1), the right side of Step 5 is represented by $\sum_jx_j+\sum_k(-1)^kD_k(\eta_k)$, and as representatives are
unique, this is $\varepsilon(\prod_kg_k|_V)^{-1}\det(D_k(g_l|_V))$.
% Lean: GlobalTheory.T.projective_chart_affine_coords = ConeTransfer.main (ProjectiveChartAffineCoords.lean). Step 1 is
% ConeZarCap.main (h1, hb1, hb0, hjl, hjb); Step 2 is liftV, hlift, hunit, hcomp (RepLifts.lift_compat) and
% ConeRepCup.exists_rep with dlog_cocycle and dlog_mem; Step 3 is ConeRep.rep_resB, ConeLocal.cup_resQ, RepLifts.main and
% rep_uniqueU (e0); Step 4 is hω' and ConeRepExact.exact_lift with ConeLocal.rep_exists (e2); Step 5 is glue (e3) and
% glue_res (e4); Step 6 is ConeLdlog.main (hd), ConeLRestr.main (hr), ConeLocal.LExact.main, rep_sum, rep_add',
% rep_zero_xdt and the uniqueness of ConeTopChart.top_chart. About 250 lines in ProjectiveChartAffineCoords.lean; the
% rest of the 5,700 lines is in Lemmas gal-realization-kappa-de-rham-transfer-lifts(2), -transfer-extend and
% -transfer-local.
\end{proof}
\leanref{Development/AutomorphicSpectrum/Proofs/BlockDecomposition/ProjectiveChartAffineCoords.lean}{GlobalTheory.CompanionTheory.ConeTransfer.main}
\leanref{Development/AutomorphicSpectrum/Proofs/BlockDecomposition/ProjectiveChartAffineCoords.lean}{GlobalTheory.CompanionTheory.ConeTransfer.glue}
\leanref{Development/AutomorphicSpectrum/Proofs/BlockDecomposition/ProjectiveChartAffineCoords.lean}{GlobalTheory.CompanionTheory.ConeTransfer.dlog_cocycle}
\leanref{Development/AutomorphicSpectrum/Proofs/BlockDecomposition/ProjectiveChartAffineCoords.lean}{GlobalTheory.CompanionTheory.ConeTransfer.rep_zero_xdt}

\begin{lemma}\label{gal-realization-kappa-de-rham}
Let $K$ be a field of characteristic $0$ and $X$ a $K$-scheme, with the de Rham cohomology of
Definition~\ref{gal-realization-period-de-rham} and the first Chern classes of
Definition~\ref{gal-realization-period-chern}(6).
\begin{enumerate}
\item If $X$ is proper and smooth over $K$, each $H^j_{\mathrm{dR}}(X)$ is finite-dimensional over $K$.
\item If $X$ is smooth of relative dimension $d$ over $K$ and $\mathcal L$ is a tautological line datum on $X$
  (Definition~\ref{gal-realization-period-chern}(2), for $R=K$), then $\mathcal L$ has a first Chern class in
  $H^2_{\mathrm{dR}}(X)$.
\item If moreover $X$ is connected, then $\omega^d\neq0$ for every first Chern class $\omega\in H^2_{\mathrm{dR}}(X)$ of
  $\mathcal L$, where $\omega^0=1$ and $\omega^{k+1}=\omega^k\cup\omega$.
\end{enumerate}
\end{lemma}
\leanref{Development/WeilBound/CoreStatements/PadicHodgeComparison/EtaleCohomologyComparison/PerfectoidPullbacks/GrCechFin.lean}{CrystallineSemistableComparison.Pb.k_fin}
\leanref{Development/WeilBound/CoreStatements/PadicHodgeComparison/EtaleCohomologyComparison/PerfectoidPullbacks/IsProperOfProj.lean}{CrystallineSemistableComparison.Pb.m_fin}
\leanref{Development/WeilBound/CoreStatements/PadicHodgeComparison/EtaleCohomologyComparison/PerfectoidPullbacks/ConeSectionRestriction/IsFirstChernDeRhamUnique.lean}{CrystallineSemistableComparison.Pb.isFirstChernDeRham_exists}
\leanref{Development/WeilBound/CoreStatements/PadicHodgeComparison/EtaleCohomologyComparison/PerfectoidPullbacks/ConeSectionRestriction/IsFirstChernDeRhamUnique.lean}{CrystallineSemistableComparison.Pb.powK_ne_zero}

\begin{proof}
(1) Let $X$ be proper and smooth over $K$, and $j\geq0$. As $X$ is separated, quasi-compact and smooth over $K$, it has a
coordinate cover $(U_\alpha)_{\alpha\in A}$ of dimensions $d_\alpha$ (Lemma~\ref{gal-realization-thickening-cover}(2)). Let
$\mathcal U=\{U_\alpha:\alpha\in A\}$, $s=\max_\alpha d_\alpha+j+1$, $\mathcal Q=\mathrm{Th}^{(s)}_{\mathrm{dR}}(X)_{\mathcal
U}$ and $\mathcal C_i=\mathrm{Th}^{(s)}_{\mathrm{dR}}(X)_{U_i}$ for $i\in A^{n+1}$, with $U_i=\bigcap_kU_{i_k}$
(Definition~\ref{gal-realization-kappa-de-rham-truncated}). Let $\mathcal S$ be the class of finite-dimensional
$K$-vector spaces; it satisfies the hypothesis of Lemma~\ref{gal-realization-thickening-level-one-serre}.

(a) Restriction $H^j_{\mathrm{dR}}(X)\to H^j(\mathcal Q,\mathcal O)$ is bijective, by
Lemma~\ref{gal-realization-kappa-de-rham-truncation} with $\mathcal V=\mathcal U$, as $s>d_\alpha$ for every $\alpha$.

(b) On $\mathcal Q$, $\mathcal O^{[s]}=0$, so $\mathcal O\to\mathcal O/\mathcal O^{[s]}$ is an isomorphism of functors and
$H^j(\mathcal Q,\mathcal O)\to H^j(\mathcal Q,\mathcal O/\mathcal O^{[s]})$ is bijective.

(c) For $0\leq r<s$, $H^j(\mathcal Q,\mathrm{gr}_r)\in\mathcal S$. By Lemma~\ref{gal-realization-thickening-over}, the
$\mathcal C_{(\alpha)}$ are cosieves of $\mathcal Q$ that cover it, with intersections the $\mathcal C_i$, so by
Lemma~\ref{gal-realization-thickening-level-one-cech} it suffices that $H^a(\check C^\bullet_b)\in\mathcal S$ for
$a+b=j$, where $\check C^a_b=\prod_{i\in A^{a+1}}H^b(\mathcal C_i,\mathrm{gr}_r)$. If $b\neq r$, then $\check C^a_b=0$, as
$H^b(\mathcal C_i,\mathrm{gr}_r)=0$ by Lemma~\ref{gal-realization-kappa-de-rham-affine} for the affine open $U_i$. If
$b=r$, Lemma~\ref{gal-realization-kappa-de-rham-bundle} shows that $\check C^\bullet_r$ is the complex of a cosimplicial
$K$-module with data as in Lemma~\ref{gal-realization-thickening-level-one-coherent} for $\Lambda=K$, $Z=X$, which is
proper over $K$, and the opens $U_\alpha$, which cover $X$; by that lemma $H^a(\check C^\bullet_r)$ is finite-dimensional.

(d) For $0\leq r\leq s$, $H^j(\mathcal Q,\mathcal O^{[r]}/\mathcal O^{[s]})\in\mathcal S$, by descending induction on $r$.
For $r=s$ the functor is $0$. For $r<s$, the long exact sequence of $0\to\mathcal O^{[r+1]}/\mathcal O^{[s]}\to\mathcal
O^{[r]}/\mathcal O^{[s]}\to\mathrm{gr}_r\to0$ gives an exact sequence $H^j(\mathcal Q,\mathcal O^{[r+1]}/\mathcal
O^{[s]})\to H^j(\mathcal Q,\mathcal O^{[r]}/\mathcal O^{[s]})\to H^j(\mathcal Q,\mathrm{gr}_r)$, whose outer terms lie in
$\mathcal S$ by induction and (c).

By (d) for $r=0$, $H^j(\mathcal Q,\mathcal O/\mathcal O^{[s]})$ is finite-dimensional, and by (a) and (b) so is
$H^j_{\mathrm{dR}}(X)$.
% Lean: Pb.k_fin = k_fin_pf = de_rham_finite_and_filtration_bounded = KFin.main, first part: the cover is
% exists_smooth_coordinate_cover, S is finite_modules_serre_class, (a) is KIso (Lemma gal-realization-kappa-de-rham-
% truncation), (b) is toSq_bijective with fil_bot, (c) is gr_fin (GlobSerreSt, which is CosieveRel.serre_union, with
% capT_eq; GrAff for b /= r; GrCechFin = GrCech.main_nl, vb_cech with affine_cover_cohomology_finite = VBGlue.finite_H_cx,
% for b = r), and (d) is dev with sq_les. The second part of k_fin, FilR q j (j + 1) = 0, is not used in this book.

(2) Let $X$ be smooth of relative dimension $d$ over $K$, and $\mathcal L$ tautological for a closed immersion
$i:X\to\mathbf P^M_K$, with morphism $f:X_{\mathcal L}\to X$. Let $\mathcal V$ be the set of open subsets of $X$ over
which $f$ has a section. For $t\geq1$ let $\mathcal Q_t=\mathrm{Th}^{(t)}_{\mathrm{dR}}(X)_{\mathcal V}$, with inclusion
$\iota_t$ into $\mathrm{Th}_{\mathrm{dR}}(X)$, and $\mathrm{res}_t=\iota_t^*:H^2_{\mathrm{dR}}(X)\to H^2(\mathcal
Q_t,\mathcal O)$ (Definition~\ref{gal-realization-split-diagrams}(1), with identity maps). Lifts along $\iota_t$ exist by
Lemma~\ref{gal-realization-kappa-de-rham-lift}. Write $z_m(\Lambda)$ as in
Lemma~\ref{gal-realization-kappa-de-rham-cocycle}.

\emph{Step 1: $\mathrm{res}_t$ is bijective for $t>d$.} As $X$ is a closed subscheme of $\mathbf P^M_K$, it is proper
over $K$ \cite[Tag 01WC]{Sta}, so separated and quasi-compact, and it has a coordinate cover with $d_\alpha=d$ for every
$\alpha$ (Lemma~\ref{gal-realization-thickening-cover}(2)). The union of the members of $\mathcal V$ is $X$
(Lemma~\ref{gal-realization-ingredients-tautological}(4)). So Lemma~\ref{gal-realization-kappa-de-rham-truncation}, for
$s=t$ and $\mathcal V$, applies.

\emph{Step 2: the classes of lifts.} For $t\geq1$ and a lift $\Lambda$ along $\iota_t$, $z_t(\Lambda)$ is a cocycle, as
$J^{[t]}=0$ on $\mathcal Q_t$ (Lemma~\ref{gal-realization-kappa-de-rham-cocycle}(1)); let $\mathrm{cl}_t(\Lambda)\in
H^2(\mathcal Q_t,\mathcal O)$ be its class. It does not depend on $\Lambda$, by
Lemma~\ref{gal-realization-kappa-de-rham-cocycle}(2). For $t\leq t'$, $\mathcal Q_t\subset\mathcal Q_{t'}$, and
restriction $H^2(\mathcal Q_{t'},\mathcal O)\to H^2(\mathcal Q_t,\mathcal O)$ sends $\mathrm{cl}_{t'}(\Lambda')$ to
$\mathrm{cl}_t(\Lambda)$: by Lemma~\ref{gal-realization-kappa-de-rham-cocycle}(3) it sends it to the class of
$z_{t'}(\Lambda'|_{\mathcal Q_t})$, which is $z_t(\Lambda'|_{\mathcal Q_t})$ by
Lemma~\ref{gal-realization-kappa-de-rham-log}(2), as $J^{[t]}=0$ on $\mathcal Q_t$ and $r(s)-1\in J$.

\emph{Step 3: the class.} Let $\Lambda_0$ be a lift along $\iota_{d+1}$, and by Step 1 let $\omega\in H^2_{\mathrm{dR}}(X)$
with $\mathrm{res}_{d+1}(\omega)=\mathrm{cl}_{d+1}(\Lambda_0)$. Then $\mathrm{res}_t(\omega)=\mathrm{cl}_t(\Lambda)$ for
every $t\geq d+1$ and every lift $\Lambda$ along $\iota_t$. Indeed, by Step 1 there is $\omega'$ with
$\mathrm{res}_t(\omega')=\mathrm{cl}_t(\Lambda)$. As $\mathrm{res}_{d+1}$ is $\mathrm{res}_t$ followed by restriction to
$\mathcal Q_{d+1}$, Step 2 gives $\mathrm{res}_{d+1}(\omega')=\mathrm{cl}_{d+1}(\Lambda_0)=\mathrm{res}_{d+1}(\omega)$, and
$\omega'=\omega$ by Step 1.

\emph{Step 4: $\omega$ is a first Chern class of $\mathcal L$.} Let $m\geq0$, $\mathcal C$ and $G:\mathcal
C\to\mathrm{Th}_{\mathrm{dR}}(X)$ be as in Definition~\ref{gal-realization-period-chern}(6), and $\Lambda$ a lift of
$\mathcal L$ along $G$. Its cochain is $z_m(\Lambda)$, a cocycle by Lemma~\ref{gal-realization-kappa-de-rham-cocycle}(1).
Let $t=\max(m,d+1)$. For every $c$, $J_c^{[t]}\subset J_c^{[m]}=0$ and $G(c)$ lies over a member of $\mathcal V$, so
$G=\iota_tG'$ for a functor $G':\mathcal C\to\mathcal Q_t$. Let $\Lambda_t$ be a lift along $\iota_t$. Then
$G^*\omega=G'^*\mathrm{res}_t(\omega)=G'^*\mathrm{cl}_t(\Lambda_t)$ by Step 3. By
Lemma~\ref{gal-realization-kappa-de-rham-cocycle}(3) this is the class of $z_t(\Lambda_tG')$, which is
$z_m(\Lambda_tG')$ by Lemma~\ref{gal-realization-kappa-de-rham-log}(2), as $J_c^{[m]}=0$. As $\Lambda_tG'$ and $\Lambda$
are lifts along $G$, its class is that of $z_m(\Lambda)$, by Lemma~\ref{gal-realization-kappa-de-rham-cocycle}(2).
% Lean: Pb.isFirstChernDeRham_exists_of. Step 1 is restriction_bijective_above_dimension (isProper_of_isO1, cov_ex,
% sSup_𝒱D and k_iso); Step 2 is clsK with clsK_res (pull_cls_zK, plogK_congr and zK_cls_eq); Step 3 is x and hA (with
% resK_res); Step 4 is zK_cocycle and test_factor (g', pull_restrict_comp, lift_exists, plogK_congr and zK_cls_eq).

(3) Let $X$, $d$, $\mathcal L$, $i$, $f$ and $\mathcal V$ be as in the proof of (2), with $X$ connected, and let $\omega$
be a first Chern class of $\mathcal L$.

\emph{The case $d=0$.} Here $\omega^0=1$. As $X$ is connected, it is nonempty; let $\xi$ be a point of $X$ with residue
field $\kappa(\xi)$. Then $(\kappa(\xi),0,0,y_\xi)$, with $y_\xi:\operatorname{Spec}\kappa(\xi)\to X$ the canonical
morphism, is a nilpotent thickening of $X$. There are no nonzero $0$-coboundaries, so if $1=0$ in
$H^0_{\mathrm{dR}}(X)$, the constant $0$-cochain $1$ is $0$, and $1=0$ in $\kappa(\xi)$, which is false.

Let now $d=n+1$ with $n\geq0$, and suppose that $\omega^{n+1}=0$. Let $\top=\{0,\ldots,n+1\}$ and
$S_j=\top\smallsetminus\{j\}$.

\emph{Step 1: the setup.} Lemma~\ref{gal-realization-kappa-de-rham-setup} gives $e$, $y_0,\ldots,y_{n+1}$, $V$, $t_0,\ldots,
t_n$ and the algebras $B(S)$. Let $U_S$ and $X_j$ be as in Definition~\ref{gal-realization-kappa-de-rham-cone}(1), $s=n+2$,
and $\mathcal W_k$, $\mathcal W_S$ and $\mathcal W$ as in Definition~\ref{gal-realization-kappa-de-rham-cone}(2). As
$K[x_0,\ldots,x_n]\to\Gamma(V,\mathcal O_X)$, $x_k\mapsto t_k$, is \'etale, $\Omega_{\Gamma(V,\mathcal O_X)/K}$ is free on
$dt_0,\ldots,dt_n$ \cite[Tags 00U1 and 06B6]{Sta}, so there are $K$-derivations $D_0,\ldots,D_n$ of $\Gamma(V,\mathcal O_X)$
with $D_k(t_j)=1$ for $k=j$ and $0$ otherwise.

\emph{Step 2: the cocycle.} A thickening $T$ of $\mathcal W$ has $J_T^{[s]}=0$ and lies over some $X_j$, over which $f$
has a section, as $\mathcal L$ is tautological for $i$ (Definition~\ref{gal-realization-period-chern}(2)). So $\mathcal W$ is a full subcategory of
$\mathrm{Th}^{(s)}_{\mathrm{dR}}(X)_{\mathcal V}$, and a lift along the inclusion of the latter
(Lemma~\ref{gal-realization-kappa-de-rham-lift}), composed with the inclusion of $\mathcal W$, is a lift $\Lambda$ of
$\mathcal L$ along the inclusion of $\mathcal W$ (Lemma~\ref{gal-realization-kappa-de-rham-cocycle}(3)). As $\omega$ is a
first Chern class (Definition~\ref{gal-realization-period-chern}(6), for $m=s$), $z=z_s(\Lambda)$ is a cocycle whose class
is the restriction $\omega|_{\mathcal W}$.

\emph{Step 3: the cup power.} Restriction to $\mathcal W$ preserves cup products and units, and the cup product of the
classes of two cocycles is the class of their cup product (Definition~\ref{gal-realization-split-diagrams}(1)). So by
induction $\omega^k|_{\mathcal W}$ is the class of $z^{\cup k}$, where $z^{\cup0}=1$ and $z^{\cup(k+1)}=z^{\cup k}\cup
z$. As $\omega^{n+1}=0$, $z^{\cup(n+1)}$ is a coboundary, and so is $(ez)^{\cup(n+1)}=e^{n+1}z^{\cup(n+1)}$.

\emph{Step 4: the local primitives.} By Lemma~\ref{gal-realization-kappa-de-rham-local-exact} there are $1$-cochains
$b_k$ of $\mathcal W_k$ with values in $J$ and $db_k=ez|_{\mathcal W_k}$, and for $0\leq k\leq n$ families $\hat g'_k$ of
lifts of $y_k/y_{k+1}$ on $\mathcal W_{\{k,k+1\}}$ with $b_k-b_{k+1}=\mathrm{dlog}_s(\hat g'_k)$ on $\mathcal
W_{\{k,k+1\}}$. Let $g_k\in\Gamma(U_\top,\mathcal O_X)$ be the restriction of $y_k/y_{k+1}$, a unit
(Definition~\ref{gal-realization-kappa-de-rham-cone}(1)), and $\hat g_k$ the restriction of $\hat g'_k$ to $\mathcal
W_\top$, a family of lifts of $g_k$ by Lemma~\ref{gal-realization-thickening-level-one-forms}(1). Let
$\ell_k=(b_k-b_{k+1})|_{\mathcal W_\top}=\mathrm{dlog}_s(\hat g_k)$. Then $d\ell_k=ez-ez=0$ on $\mathcal W_\top$, and the
values of $\ell_k$ lie in $J$. By Lemma~\ref{gal-realization-kappa-de-rham-cup}, for the filtration
$(\mathcal O^{[r]})$ of $\Phi_{\mathrm{Th}_{\mathrm{dR}}(X)}$ (Step 3 of the proof of
Lemma~\ref{gal-realization-thickening-level-one-local}), $c_0=[\ell_0\cup\cdots\cup\ell_n]_{n+1}\in H^{n+1}(\mathcal
W_\top,\mathrm{gr}_{n+1})$ is defined.

\emph{Step 5: the relation.} Lemma~\ref{gal-realization-kappa-de-rham-cech} applies to $\Phi_{\mathrm{Th}_{\mathrm{dR}}(X)}$
with the filtration $(\mathcal O^{[r]})$, the categories $\mathcal W_0,\ldots,\mathcal W_{n+1}$, $s$, the cocycle $ez$,
whose values lie in $J$ by Lemma~\ref{gal-realization-kappa-de-rham-log}(1), and the $b_k$. Its hypotheses hold: (a) as
$J^{[s]}=0$ on $\mathcal W$; (b) as for a map $T\to T'$ with $T$ over $U_k\cap X_j$, $T'$ lies over $U_k\cap X_j$
(Lemma~\ref{gal-realization-thickening-over}(1)); and (c) by
Lemma~\ref{gal-realization-kappa-de-rham-cone-vanishing}(1), with $r'=r+1$, which identifies $H^j(\mathcal
W_S,\mathrm{gr}_r)$ with $H^j(\mathrm{Th}^{(s)}_{\mathrm{dR}}(X)_{U_S},\mathrm{gr}_r)$ for the affine $U_S$. The latter
vanishes for $r<0$ by Lemma~\ref{gal-realization-kappa-de-rham-cone-vanishing}(3), for $0\leq r\leq n+1<s$ and $j\neq r$
by Lemma~\ref{gal-realization-kappa-de-rham-affine}, and for $r>n+1=d$ by
Lemma~\ref{gal-realization-kappa-de-rham-cone-vanishing}(2). So there are classes $\omega_j\in H^{n+1}(\mathcal
W_{S_j},\mathrm{gr}_{n+1})$ and an $n$-cochain $f$ of $\mathcal W_\top$ with values in $J^{[n]}$ whose coboundary has values
in $J^{[n+1]}$, such that $[df]_{n+1}=c_0-\sum_j\omega_j|_{\mathcal W_\top}$.

\emph{Step 6: the transfer.} By Lemma~\ref{gal-realization-kappa-de-rham-transfer}, applied with $V$, the $t_k$, the $D_k$,
the $g_k$ and $\hat g_k$, $c_0$, whose hypotheses hold by Step 4 as $\ell_k=\mathrm{dlog}_s(\hat g_k)$, the $\omega_j$ and
$f$, there are $x_j\in\mathrm{Jac}(B(S_j))$, $\eta_k\in
\mathrm{Jac}_k(B(\top))$ and $\varepsilon\in\{1,-1\}$ with
\[
 \varepsilon\Bigl(\prod_kg_k|_V\Bigr)^{-1}\det\bigl(D_k(g_l|_V)\bigr)=\sum_jx_j+\sum_k(-1)^kD_k(\eta_k);
\]
the algebras $B'(S)$ there are the $B(S)$ of Lemma~\ref{gal-realization-kappa-de-rham-setup}.

\emph{Step 7: the computation.} On $V$, $(y_a/y_b)(y_b/y_c)=y_a/y_c$, $y_a/y_a=1$ and $t_k=y_{k+1}/y_0$, so
$g_0|_V=t_0^{-1}$ and $g_l|_V=t_{l-1}t_l^{-1}$ for $1\leq l\leq n$. Let $u_k=t_k^{-1}$; then $D_k(u_j)=-u_j^2$ for $k=j$
and $0$ otherwise. It follows that $D_k(g_l|_V)=u_k\beta_{kl}\,g_l|_V$, where $\beta_{kk}=-1$, $\beta_{k,k+1}=1$ and
$\beta_{kl}=0$ otherwise: for $l=0$ this is $D_k(u_0)$, and for $l\geq1$, $D_k(t_{l-1}u_l)$ is $u_l=u_{l-1}t_{l-1}u_l$ for
$k=l-1$, $-t_{l-1}u_l^2$ for $k=l$, and $0$ otherwise. The matrix $\beta$ is upper triangular with determinant
$(-1)^{n+1}$, so $(\prod_kg_k|_V)^{-1}\det(D_k(g_l|_V))=(-1)^{n+1}u_0\cdots u_n=(-1)^{n+1}(t_0\cdots t_n)^{-1}$. Replacing
the $x_j$ and $\eta_k$ by their negatives if $\varepsilon(-1)^{n+1}=-1$, which keeps them in the subgroups
$\mathrm{Jac}(B(S_j))$ and $\mathrm{Jac}_k(B(\top))$, gives
\[
 (t_0\cdots t_n)^{-1}=\sum_{j\in\top}x_j+\sum_{k=0}^n(-1)^kD_k(\eta_k).
\]

\emph{Step 8.} This contradicts Lemma~\ref{gal-realization-kappa-de-rham-residue}, for $R=\Gamma(V,\mathcal O_X)$, a
domain, the units $t_k$, which are algebraically independent over $K$, the $D_k$ and the algebras $B(S)$
(Lemma~\ref{gal-realization-kappa-de-rham-setup}(2) and~(3)). So $\omega^{n+1}\neq0$.
% Lean: Pb.powK_ne_zero_main = powK_ne_zero_of geo_setup resid_alg firstChern_local_exact. The case d = 0 is
% oneK_ne_zero (ConeOne.oneK_ne_zero, at the thickening resT q x). Step 1 is hG and
% exists_dual_derivations_of_coords (ConeDerivs.exists_dual_derivations, through
% EtaleCoordJacobian.exists_basis_of_etale, where the book cites Stacks); Step 2 is hsQ, hVQ (wch_mem), Λ (lift_exists
% precomposed along wU_le) and hc; Step 3 is powK_pull with cupPowC_smul (hzero); Step 4 is hL, ℓ, hℓcoc, hℓ1, hĝ, hℓeq
% and rep_cupList (hc0'); Step 5 is cech_top_rel with hsU, wK_rel, hvan1 and hvan2 (zar_capS, gr_neg_van, gr_aff,
% gr_top_van); Step 6 is cone_transfer; Step 7 is g_rel, dlog_ratios (Cone.dlog_ratio_det), jacSpan_le and hneg; Step 8
% is hR = resid_alg. Size of (3): about 29,400 lines beyond the properties it uses, of which about 12,000 lie beyond
% the lemmas the book already states; the open steps are in Lemmas gal-realization-kappa-de-rham-setup (about 2,800),
% -residue (about 1,500), -local-exact (about 750), -cech (about 1,000) and -transfer (about 5,700).
\end{proof}
\leanref{Development/WeilBound/CoreStatements/PadicHodgeComparison/EtaleCohomologyComparison/PerfectoidPullbacks/ConeSectionRestriction/RestrictionBijectiveAboveDimension.lean}{CrystallineSemistableComparison.Pb.isFirstChernDeRham_exists_of}
\leanref{Development/WeilBound/CoreStatements/PadicHodgeComparison/EtaleCohomologyComparison/PerfectoidPullbacks/ConeSectionRestriction/RestrictionBijectiveAboveDimension.lean}{CrystallineSemistableComparison.Pb.restriction_bijective_above_dimension}
\leanref{Development/WeilBound/CoreStatements/PadicHodgeComparison/EtaleCohomologyComparison/PerfectoidPullbacks/ConeSectionRestriction/RestrictionBijectiveAboveDimension.lean}{CrystallineSemistableComparison.Pb.isProper_of_isO1}
\leanref{Development/WeilBound/CoreStatements/PadicHodgeComparison/EtaleCohomologyComparison/PerfectoidPullbacks/ConeSectionRestriction/RestrictionBijectiveAboveDimension.lean}{CrystallineSemistableComparison.Pb.clsK}
\leanref{Development/WeilBound/CoreStatements/PadicHodgeComparison/EtaleCohomologyComparison/PerfectoidPullbacks/ConeSectionRestriction/RestrictionBijectiveAboveDimension.lean}{CrystallineSemistableComparison.Pb.clsK_res}
\leanref{Development/WeilBound/CoreStatements/PadicHodgeComparison/EtaleCohomologyComparison/PerfectoidPullbacks/ConeSectionRestriction/RestrictionBijectiveAboveDimension.lean}{CrystallineSemistableComparison.Pb.test_factor}
\leanref{Development/WeilBound/CoreStatements/PadicHodgeComparison/EtaleCohomologyComparison/PerfectoidPullbacks/ConeSectionRestriction/PowTop.lean}{CrystallineSemistableComparison.Pb.powK_ne_zero_main}
\leanref{Development/WeilBound/CoreStatements/PadicHodgeComparison/EtaleCohomologyComparison/PerfectoidPullbacks/ConeSectionRestriction/PowKNeZeroOf.lean}{CrystallineSemistableComparison.Pb.powK_ne_zero_of}
\leanref{Development/WeilBound/CoreStatements/PadicHodgeComparison/EtaleCohomologyComparison/PerfectoidPullbacks/ConeSectionRestriction/OneKNeZero.lean}{CrystallineSemistableComparison.Pb.oneK_ne_zero}
\leanref{Development/AutomorphicSpectrum/Proofs/BlockDecomposition/UnitClassNeZeroOfNonempty.lean}{GlobalTheory.CompanionTheory.ConeOne.oneK_ne_zero}
\leanref{Development/AutomorphicSpectrum/Proofs/BlockDecomposition/ExistsDualDerivationsOfCoordsStatement.lean}{GlobalTheory.T.exists_dual_derivations_of_coords}
\leanref{Development/WeilBound/CoreStatements/PadicHodgeComparison/EtaleCohomologyComparison/PerfectoidPullbacks/ConeSectionRestriction/WchCover.lean}{CrystallineSemistableComparison.Pb.wch_mem}
\leanref{Development/WeilBound/CoreStatements/PadicHodgeComparison/EtaleCohomologyComparison/PerfectoidPullbacks/ConeSectionRestriction/WchCover.lean}{CrystallineSemistableComparison.Pb.wU_le}
\leanref{Development/WeilBound/CoreStatements/PadicHodgeComparison/EtaleCohomologyComparison/PerfectoidPullbacks/ConeSectionRestriction/WchCover.lean}{CrystallineSemistableComparison.Pb.wK_rel}
\leanref{Development/WeilBound/CoreStatements/PadicHodgeComparison/EtaleCohomologyComparison/PerfectoidPullbacks/ConeSectionRestriction/PowKPull.lean}{CrystallineSemistableComparison.Pb.powK_pull}
\leanref{Development/WeilBound/CoreStatements/PadicHodgeComparison/EtaleCohomologyComparison/PerfectoidPullbacks/ConeSectionRestriction/OneKNeZero.lean}{CrystallineSemistableComparison.Pb.cupPowC_smul}
\leanref{Development/WeilBound/CoreStatements/PadicHodgeComparison/EtaleCohomologyComparison/PerfectoidPullbacks/ConeSectionRestriction/PowKNeZeroOf.lean}{CrystallineSemistableComparison.Pb.g_rel}
\leanref{Development/WeilBound/CoreStatements/PadicHodgeComparison/EtaleCohomologyComparison/PerfectoidPullbacks/ConeSectionRestriction/RepDlogLifts.lean}{CrystallineSemistableComparison.Pb.dlog_ratios}
\leanref{Development/AutomorphicSpectrum/Proofs/BlockDecomposition/TorusJacobianDet.lean}{GlobalTheory.CompanionTheory.Cone.dlog_ratio_det}
\leanref{Development/WeilBound/CoreStatements/PadicHodgeComparison/EtaleCohomologyComparison/PerfectoidPullbacks/ConeSectionRestriction/JacSpanLe.lean}{CrystallineSemistableComparison.Pb.jacSpan_le}

\subsection{Thickenings, torsor data and Zariski descent}\label{subsection-gal-torsor-data}

\begin{lemma}\label{gal-realization-frobenius-congruence}
Let $\ell$, $C$, $A_{\mathrm{cris}}$, $K$, $\mathcal O_K$, $\mathbf F_q$, $q=\ell^f$, and $g$ be as in
Definition~\ref{gal-realization-frobenius-datum}, and let $W\to A_{\mathrm{cris}}$ be as in
Lemma~\ref{gal-realization-invariants}, induced by the embedding $s:\mathbf F_q\to\mathcal O_C^\flat$ given there. For
$x\in\mathcal O_K$ with residue $\bar x$ let $\tilde x=\theta([s(\bar x)])\in\mathcal O_C$ be the Teichm\"uller
representative of $\bar x$. Let $A_\nu$, $I_\nu$, $I'_\nu$ and $u:\mathcal O_C\to A_\nu/I_\nu$ be as in
Definition~\ref{gal-realization-period-thickenings}(1), and $W_\nu=W/\ell^\nu$. For every integer $f_1\geq1$ there is a
multiple $N$ of $f_1$ with the following properties.
\begin{enumerate}
\item $x^{\ell^N}-\tilde x\in\ell\mathcal O_C$ for every $x\in\mathcal O_K$; that is, $(x\bmod\ell)^{\ell^N}=g(x)$.
\item For every $\nu\geq1$, the map $\mathcal O_K\to A_\nu/I'_\nu$ that sends $x$ to the $\ell^N$-th power of the image of
  $u(x)$ is the composite of the residue map $\mathcal O_K\to\mathbf F_q=W_\nu/\ell$ with the map $W_\nu/\ell\to
  A_\nu/I'_\nu$ induced by $W\to A_{\mathrm{cris}}$.
\end{enumerate}
\end{lemma}
\leanref{Development/WeilBound/CoreStatements/PadicHodgeComparison/EtaleCohomologyComparison/TeichmullerRootsUnity/RootsUnityPullback/TeichDvd.lean}{CrystallineSemistableComparison.crystallineComparison.PinP.teich_dvd}
\leanref{Development/WeilBound/CoreStatements/PadicHodgeComparison/EtaleCohomologyComparison/TeichmullerRootsUnity/RootsUnityPullback/TeichDvd.lean}{CrystallineSemistableComparison.crystallineComparison.PinP.close_dvd}
\leanref{Development/WeilBound/CoreStatements/PadicHodgeComparison/EtaleCohomologyComparison/PerfectoidPullbacks/FrobeniusWittActions/TeichmullerCongruence.lean}{CrystallineSemistableComparison.Pb.Mid.TeichmullerCongruence}
\leanref{Development/AutomorphicArithmetic/EtaleIdempotents/IntegersResidueEqThetaCoeffPowOfTeichmuller.lean}{UpperTheory.integersResidue_eq_theta_coeff_pow_of_teichmuller}

\begin{proof}
Absolute values are those of $C$. The residue field has $\ell^f$ elements, so $z^{\ell^f}=z$ for $z\in\mathbf F_q$. As
$s$ and $y\mapsto\theta([y])$ are multiplicative, $\tilde x^{\ell^{fk}}=\tilde x$ for $x\in\mathcal O_K$ and $k\geq0$.

\emph{Step 1.} If $a,b\in\mathcal O_C$, $|a-b|\leq c$ and $k\geq0$, then $|a^{\ell^k}-b^{\ell^k}|\leq\max(|\ell|,c^{\ell^k})$.
By the binomial theorem, $a^\ell-b^\ell=(a-b)^\ell+\ell b(a-b)r$ with $r\in\mathcal O_C$, so $|a^\ell-b^\ell|\leq
\max(|a-b|^\ell,|\ell|)$. Induction on $k$ gives the claim, since $|\ell|^\ell\leq|\ell|$.

\emph{Step 2.} There is $c<1$ with $|x-\tilde x|\leq c$ for every $x\in\mathcal O_K$. Let $c_0<1$ bound the absolute
values of the elements of the maximal ideal of $\mathcal O_K$, and put $c=\max(c_0,|\ell|)$. Suppose that $v=x-\tilde x$
has $|v|>c$. As $x$ and $\tilde x$ have the same residue, $|v|<1$. Step 1 for $c=|v|$ and $k=f$ gives
$|x^{\ell^f}-\tilde x|\leq\max(|\ell|,|v|^{\ell^f})<|v|$. As $x^{\ell^f}-x$ lies in the maximal ideal of $\mathcal O_K$,
$|x^{\ell^f}-x|\leq c_0<|v|$. Since $v=(x^{\ell^f}-\tilde x)-(x^{\ell^f}-x)$, this gives $|v|<|v|$, which is absurd.

\emph{Step 3.} Choose $j\geq1$ with $c^j<|\ell|$, and put $N=ff_1j$, a multiple of $f$ and of $f_1$. Then $\ell^N\geq
j$, so $c^{\ell^N}\leq c^j<|\ell|$, and Step 1 with $\tilde x^{\ell^N}=\tilde x$ gives $|x^{\ell^N}-\tilde x|\leq|\ell|$
for every $x\in\mathcal O_K$. This
is (1), as $g(x)$ is the class of $\tilde x$.

\emph{Step 4.} Let $\nu\geq1$. As $\ell\in I'_\nu$, elements of $\mathcal O_C$ that agree modulo $\ell$ have the same image
in $A_\nu/I'_\nu$. As $u(\theta(z))$ is the class of $z$ for $z\in W(\mathcal O_C^\flat)$, the image of $\tilde x$ is the
image of $[s(\bar x)]$, which is the image of the Teichm\"uller representative of $\bar x$ under $W\to A_{\mathrm{cris}}$.
By (1), the $\ell^N$-th power of the image of $u(x)$, which is the image of $x^{\ell^N}$, equals the image of $\tilde x$.
This is (2).
\end{proof}
\leanref{Development/WeilBound/CoreStatements/PadicHodgeComparison/EtaleCohomologyComparison/PerfectoidPullbacks/FrobeniusWittActions/FrobeniusNormBounds/ResiduePadicIntegersToIntegers.lean}{CrystallineSemistableComparison.Pb.G3.res_frob}
\leanref{Development/WeilBound/CoreStatements/PadicHodgeComparison/EtaleCohomologyComparison/PerfectoidPullbacks/FrobeniusWittActions/FrobeniusNormBounds/ResiduePadicIntegersToIntegers.lean}{CrystallineSemistableComparison.Pb.G3.norm_pow_sub_pow_le_bis}
\leanref{Development/WeilBound/CoreStatements/PadicHodgeComparison/EtaleCohomologyComparison/PerfectoidPullbacks/FrobeniusWittActions/FrobeniusNormBounds/ResiduePadicIntegersToIntegers.lean}{CrystallineSemistableComparison.Pb.G3.pow_close}
\leanref{Development/WeilBound/CoreStatements/PadicHodgeComparison/EtaleCohomologyComparison/PerfectoidPullbacks/FrobeniusWittActions/FrobeniusNormBounds/ResiduePadicIntegersToIntegers.lean}{CrystallineSemistableComparison.Pb.G3.disc}
\leanref{Development/WeilBound/CoreStatements/PadicHodgeComparison/EtaleCohomologyComparison/PerfectoidPullbacks/FrobeniusWittActions/FrobeniusNormBounds/ResiduePadicIntegersToIntegers.lean}{CrystallineSemistableComparison.Pb.G3.exists_teichmuller_close}
\leanref{Development/WeilBound/CoreStatements/PadicHodgeComparison/EtaleCohomologyComparison/PerfectoidPullbacks/FrobeniusWittActions/FrobeniusNormBounds/ResiduePadicIntegersToIntegers.lean}{CrystallineSemistableComparison.Pb.G3.cong}
\leanref{Development/WeilBound/CoreStatements/PadicHodgeComparison/EtaleCohomologyComparison/PerfectoidPullbacks/FrobeniusWittActions/FrobeniusNormBounds/ResiduePadicIntegersToIntegers.lean}{CrystallineSemistableComparison.Pb.G3.struct}

\begin{definition}\label{gal-realization-frobenius-teichmuller}
Let $\ell$, $C$, $A_{\mathrm{cris}}$, $K$, $\mathcal O_K$, $G_K$, $\mathbf F_q$, $W$, $W_\nu$, $W_\nu\to A_\nu$ and
$\mathrm{Th}_\nu(Y/W)$ be as in Definition~\ref{gal-realization-frobenius-special-fibre}, and $I'_\nu$,
$\mathrm{Th}'_\nu(Y)$ and $H'^j_\nu(Y)$ as in Definition~\ref{gal-realization-period-thickenings}, for $\mathcal
O=\mathcal O_K$ and its inclusion into $\mathcal O_C$. Let $Y$ be an $\mathcal O_K$-scheme, $\nu\geq1$ and $j\geq0$.
\begin{enumerate}
\item A \emph{Teichm\"uller thickening} of $Y$ of level $\nu$ is defined as a thickening for $I'_\nu$ in
  Definition~\ref{gal-realization-period-thickenings}(2), except that $B/J$ is an $\mathcal O_K$-algebra through
  $\mathcal O_K\to\mathbf F_q=W_\nu/\ell\to A_\nu/I'_\nu\xrightarrow{i}B/J$, the second map being induced by
  $W_\nu\to A_\nu$. With the maps of that definition they form a category $\mathrm{Th}^{\mathrm T}_\nu(Y)$. Put
  $H^{\mathrm T,j}_\nu(Y)=H^j(\mathrm{Th}^{\mathrm T}_\nu(Y),\mathcal O)$, an $A_{\mathrm{cris}}$-module, where
  $\mathcal O$ sends a Teichm\"uller thickening to its ring $B$.
\item The following functors, with the identity maps of the rings $B$, induce maps as in
  Definition~\ref{gal-realization-split-diagrams}(1). In (b) and (c), $N\geq0$ is such that the conclusion of
  Lemma~\ref{gal-realization-frobenius-congruence}(2) holds for $\nu$, and $F_T$ is the $\ell$-th power map of $B_T/J_T$
  for $T=(B_T,i_T,J_T,\delta_T,y_T)$.
  \begin{enumerate}
  \item[(a)] $\Gamma_W:\mathrm{Th}^{\mathrm T}_\nu(Y)\to\mathrm{Th}_\nu(Y/W)$, $(B,i,J,\delta,y)\mapsto(B,i\circ
    (W_\nu\to A_\nu),J,\delta,y)$, gives a map $b_\nu:H^j_\nu(Y/W)\to H^{\mathrm T,j}_\nu(Y)$, semilinear for
    $W\to A_{\mathrm{cris}}$.
  \item[(b)] $M_N:\mathrm{Th}'_\nu(Y)\to\mathrm{Th}^{\mathrm T}_\nu(Y)$, $T\mapsto(B_T,i_T,J_T,\delta_T,y_T\circ
    \operatorname{Spec}F_T^N)$, gives an $A_{\mathrm{cris}}$-linear map $\mu_\nu:H^{\mathrm T,j}_\nu(Y)\to H'^j_\nu(Y)$.
  \item[(c)] $V_N:\mathrm{Th}^{\mathrm T}_\nu(Y)\to\mathrm{Th}'_\nu(Y)$, $(B,i,J,\delta,y)\mapsto(B,i\circ\phi^N,J,\delta,
    y)$, gives a $\phi^N$-semilinear map $\nu_\nu:H'^j_\nu(Y)\to H^{\mathrm T,j}_\nu(Y)$.
  \item[(d)] For $\sigma\in G_K$, $(B,i,J,\delta,y)\mapsto(B,i\circ\sigma,J,\delta,y)$ on $\mathrm{Th}^{\mathrm
    T}_\nu(Y)$ gives a $\sigma$-semilinear map $\sigma$ of $H^{\mathrm T,j}_\nu(Y)$.
  \item[(e)] A Teichm\"uller thickening of level $\nu$ is one of level $\nu+1$ through $A_{\nu+1}\to A_\nu$. This gives
    the reductions $H^{\mathrm T,j}_{\nu+1}(Y)\to H^{\mathrm T,j}_\nu(Y)$.
  \end{enumerate}
\end{enumerate}
These functors take values in the stated categories. The element $\ell$ lies in $I'_\nu$, and the divided powers of
$\ell W_\nu$ and of $I'_\nu$ are both induced by $x\mapsto x^k/k!$. The maps $\phi$ and $\sigma$ preserve $I'_\nu$ and its
divided powers, $\phi$ induces the $\ell$-th power map of $A_\nu/I'_\nu$, and $\sigma$ fixes the image of $W$
(Lemma~\ref{gal-realization-invariants}). In (b) and (c) the new points are $\mathcal O_K$-morphisms by the choice of
$N$.
\end{definition}
\leanref{Development/WeilBound/CoreStatements/PadicHodgeComparison/EtaleCohomologyComparison/PerfectoidPullbacks/FrobeniusWittActions/TeichmullerCongruence.lean}{CrystallineSemistableComparison.Pb.Mid.κm}
\leanref{Development/WeilBound/CoreStatements/PadicHodgeComparison/EtaleCohomologyComparison/PerfectoidPullbacks/FrobeniusWittActions/TeichmullerCongruence.lean}{CrystallineSemistableComparison.Pb.Mid.Lm}
\leanref{Development/WeilBound/CoreStatements/PadicHodgeComparison/EtaleCohomologyComparison/PerfectoidPullbacks/FrobeniusWittActions/MuMor.lean}{CrystallineSemistableComparison.Pb.Mid.HM}
\leanref{Development/WeilBound/CoreStatements/PadicHodgeComparison/EtaleCohomologyComparison/PerfectoidPullbacks/FrobeniusFactorisation/PullT.lean}{CrystallineSemistableComparison.Pb.bWMor}
\leanref{Development/WeilBound/CoreStatements/PadicHodgeComparison/EtaleCohomologyComparison/PerfectoidPullbacks/FrobeniusFactorisation/PullT.lean}{CrystallineSemistableComparison.Pb.bW}
\leanref{Development/WeilBound/CoreStatements/PadicHodgeComparison/EtaleCohomologyComparison/PerfectoidPullbacks/FrobeniusWittActions/TeichmullerCongruence.lean}{CrystallineSemistableComparison.Pb.Mid.mIn}
\leanref{Development/WeilBound/CoreStatements/PadicHodgeComparison/EtaleCohomologyComparison/PerfectoidPullbacks/FrobeniusWittActions/MuMor.lean}{CrystallineSemistableComparison.Pb.Mid.muMor}
\leanref{Development/WeilBound/CoreStatements/PadicHodgeComparison/EtaleCohomologyComparison/PerfectoidPullbacks/FrobeniusWittActions/MuMor.lean}{CrystallineSemistableComparison.Pb.Mid.mu}
\leanref{Development/WeilBound/CoreStatements/PadicHodgeComparison/EtaleCohomologyComparison/PerfectoidPullbacks/FrobeniusWittActions/TeichmullerCongruence.lean}{CrystallineSemistableComparison.Pb.Mid.mOut}
\leanref{Development/WeilBound/CoreStatements/PadicHodgeComparison/EtaleCohomologyComparison/PerfectoidPullbacks/FrobeniusWittActions/MuMor.lean}{CrystallineSemistableComparison.Pb.Mid.nuMor}
\leanref{Development/WeilBound/CoreStatements/PadicHodgeComparison/EtaleCohomologyComparison/PerfectoidPullbacks/FrobeniusWittActions/MuMor.lean}{CrystallineSemistableComparison.Pb.Mid.nu}
\leanref{Development/WeilBound/CoreStatements/PadicHodgeComparison/EtaleCohomologyComparison/PerfectoidPullbacks/FrobeniusFactorisation/PullT.lean}{CrystallineSemistableComparison.Pb.galT'}
\leanref{Development/WeilBound/CoreStatements/PadicHodgeComparison/EtaleCohomologyComparison/PerfectoidPullbacks/FrobeniusFactorisation/PullT.lean}{CrystallineSemistableComparison.Pb.galM}

\begin{lemma}\label{gal-realization-frobenius-teichmuller-maps}
Let $\ell$, $K$, $\mathcal O_K$ and $G_K$ be as in Definition~\ref{gal-realization-frobenius-teichmuller}, $Y$ an
$\mathcal O_K$-scheme, $\nu\geq1$, $j\geq0$, $N$ as in Lemma~\ref{gal-realization-frobenius-congruence}, and $b_\nu$,
$\mu_\nu$, $\nu_\nu$ and the maps $\sigma$ of $H^{\mathrm T,j}_\nu(Y)$ as in
Definition~\ref{gal-realization-frobenius-teichmuller}. Let $\phi$ and $\sigma$ be the maps of $H'^j_\nu(Y)$ of
Definition~\ref{gal-realization-period-thickenings}(3).
\begin{enumerate}
\item The maps $b_\nu$, $\mu_\nu$ and $\nu_\nu$ commute with the reductions.
\item $\mu_\nu\circ\nu_\nu=\phi^N$ on $H'^j_\nu(Y)$.
\item $\sigma\circ b_\nu=b_\nu$ and $\nu_\nu\circ\sigma=\sigma\circ\nu_\nu$ for $\sigma\in G_K$.
\item The functor $\Gamma_W\circ M_N:\mathrm{Th}'_\nu(Y)\to\mathrm{Th}_\nu(Y/W)$, with the identity maps $\tau_T$ of the
  rings $B_T$, is a Frobenius twist of level $N$ (Definition~\ref{gal-realization-frobenius-special-fibre}(2)), and its
  map $m:H^j_\nu(Y/W)\to H'^j_\nu(Y)$ is $\mu_\nu\circ b_\nu$.
\end{enumerate}
\end{lemma}
\leanref{Development/WeilBound/CoreStatements/PadicHodgeComparison/EtaleCohomologyComparison/PerfectoidPullbacks/FrobeniusFactorisation/NuPull.lean}{CrystallineSemistableComparison.Pb.bW_red}
\leanref{Development/WeilBound/CoreStatements/PadicHodgeComparison/EtaleCohomologyComparison/PerfectoidPullbacks/FrobeniusFactorisation/NuPull.lean}{CrystallineSemistableComparison.Pb.mu_red}
\leanref{Development/WeilBound/CoreStatements/PadicHodgeComparison/EtaleCohomologyComparison/PerfectoidPullbacks/FrobeniusFactorisation/NuPull.lean}{CrystallineSemistableComparison.Pb.nu_red}
\leanref{Development/WeilBound/CoreStatements/PadicHodgeComparison/EtaleCohomologyComparison/PerfectoidPullbacks/FrobeniusWittActions/MuNu.lean}{CrystallineSemistableComparison.Pb.Mid.mu_nu}
\leanref{Development/WeilBound/CoreStatements/PadicHodgeComparison/EtaleCohomologyComparison/PerfectoidPullbacks/FrobeniusFactorisation/NuPull.lean}{CrystallineSemistableComparison.Pb.bW_gal}
\leanref{Development/WeilBound/CoreStatements/PadicHodgeComparison/EtaleCohomologyComparison/PerfectoidPullbacks/FrobeniusFactorisation/NuPull.lean}{CrystallineSemistableComparison.Pb.nu_gal}
\leanref{Development/WeilBound/CoreStatements/PadicHodgeComparison/EtaleCohomologyComparison/TeichmullerRootsUnity/RootsUnityPullback/MmPull.lean}{CrystallineSemistableComparison.crystallineComparison.PinP.mm}
\leanref{Development/WeilBound/CoreStatements/PadicHodgeComparison/EtaleCohomologyComparison/TeichmullerRootsUnity/RootsUnityPullback/IsMuMmG.lean}{CrystallineSemistableComparison.crystallineComparison.PinP.isMu_mmG}
\leanref{Development/WeilBound/CoreStatements/PadicHodgeComparison/EtaleCohomologyComparison/TeichmullerRootsUnity/RootsUnityPullback/MmPull.lean}{CrystallineSemistableComparison.crystallineComparison.PinP.mm_pull}

\begin{proof}
Each map in the statement is induced, as in Definition~\ref{gal-realization-split-diagrams}(1), by a functor with the
identity maps of the rings, and a composite of two such maps is induced by the composite functor in the opposite order.
So each identity follows once the two composite functors agree.

(1) For $b_\nu$, both composites send a Teichm\"uller thickening of level $\nu$ to the same ring, ideal, divided powers
and point, with the map $W_{\nu+1}\to W_\nu\to A_\nu\xrightarrow{i}B$, which is also $W_{\nu+1}\to A_{\nu+1}\to
A_\nu\xrightarrow{i}B$. For $\mu_\nu$ and $\nu_\nu$ the composites agree because $A_{\nu+1}\to A_\nu$ commutes with
$\phi^N$, and the $\ell$-th power maps of the rings $B/J$ do not depend on the level.

(2) The map $\nu_\nu$ is induced by $V_N$ and $\mu_\nu$ by $M_N$, so $\mu_\nu\circ\nu_\nu$ is induced by $V_N\circ M_N$,
which sends $T$ to $(B_T,i_T\circ\phi^N,J_T,\delta_T,y_T\circ\operatorname{Spec}F_T^N)$. This is the $N$-th power of the
functor of Definition~\ref{gal-realization-period-thickenings}(3)(c) that induces $\phi$, so $\mu_\nu\circ\nu_\nu=
\phi^N$.

(3) The map $\sigma\circ b_\nu$ is induced by $\Gamma_W$ after the functor of
Definition~\ref{gal-realization-frobenius-teichmuller}(2)(d). This composite sends $(B,i,J,\delta,y)$ to $(B,i\circ\sigma
\circ(W_\nu\to A_\nu),J,\delta,y)$, which is $\Gamma_W(B,i,J,\delta,y)$ since $\sigma$ fixes the image of $W$
(Lemma~\ref{gal-realization-invariants}). Both $\nu_\nu\circ\sigma$ and $\sigma\circ\nu_\nu$ are induced by functors
that send $(B,i,J,\delta,y)$ to $(B,i\circ\phi^N\circ\sigma,J,\delta,y)$ and to $(B,i\circ\sigma\circ\phi^N,J,\delta,y)$,
which agree since $\phi$ and $\sigma$ commute on $A_{\mathrm{cris}}$.

(4) The functor $\Gamma_W\circ M_N$ sends $T$ to $T^\dagger=(B_T,i_T\circ(W_\nu\to A_\nu),J_T,\delta_T,y_T\circ
\operatorname{Spec}F_T^N)$. The identity of $B_T$ carries $J_{T^\dagger}$ onto $J_T$ with the same divided powers, its
composite with $i_{T^\dagger}$ is $W_\nu\to A_\nu$ followed by $i_T$, and $y_{T^\dagger}=y_T\circ\operatorname{Spec}
F_T^N$. These are the conditions of a Frobenius twist of level $N$. Its map is induced by $\Gamma_W\circ M_N$, so it is
$\mu_\nu\circ b_\nu$.
\end{proof}

\begin{definition}\label{gal-realization-thickening-nilpotent}
Let $R$ be a ring, $\gamma$ divided powers on an ideal $I$ of $R$, $p\geq0$ an integer, $\mathcal O$ a ring, $X$ an
$\mathcal O$-scheme and $\kappa:\mathcal O\to R/I$ a ring map.
\begin{enumerate}
\item A \emph{thickening} of $X$ over $(R,I,\gamma,p)$ consists of a ring $B$, a ring map $i:R\to B$, an ideal $J$ of $B$
  with divided powers $\delta$ such that $i(I)\subset J$ and $\delta_k(i(x))=i(\gamma_k(x))$ for $x\in I$, and an
  $\mathcal O$-morphism $y:\operatorname{Spec}B/J\to X$, where $B/J$ is an $\mathcal O$-algebra through $\mathcal
  O\xrightarrow{\kappa}R/I\xrightarrow{i}B/J$; moreover $p$ is nilpotent in $B$ if $p\neq0$, and $J$ is nilpotent if
  $p=0$. Maps are as in Definition~\ref{gal-realization-period-thickenings}(2). Let $\mathrm{Th}(X/R)$ be the category of
  these thickenings, $\mathcal O$ the functor sending a thickening to the $R$-module $B$, and
  $\mathrm{Th}^{\mathrm{fl}}(X/R)$ the full subcategory of the thickenings with $B$ flat over $R$ through $i$.
\item A thickening \emph{lies over} an open subset $V$ of $X$ if $y$ factors through $V$. For a full subcategory
  $\mathcal Q$ of $\mathrm{Th}(X/R)$ and a set $\mathcal V$ of open subsets of $X$, $\mathcal Q_{\mathcal V}$ is the full
  subcategory of the thickenings in $\mathcal Q$ that lie over some $V\in\mathcal V$, and $\mathcal Q_V=\mathcal
  Q_{\{V\}}$.
\item For a thickening $T=(B,i,J,\delta,y)$ and $r\in\mathbf Z$, $J^{[r]}$ is the ideal of
  Definition~\ref{gal-realization-kappa-approximations} for $\delta$ if $r\geq0$, and $J^{[r]}=B$ if $r\leq0$; for $c\in
  R$, $cB$ denotes $i(c)B$. For two of the families of ideals $T\mapsto J_T^{[r]}+cB_T$ and $T\mapsto cB_T$ ($r\in\mathbf
  Z$, $c\in R$), say $\mathfrak a$ and $\mathfrak b$, $\mathfrak a/\mathfrak b$ is the functor sending $T$ to the image of
  $\mathfrak a_T$ in $B_T/\mathfrak b_T$, an $R$-module. Thus $\mathcal O=J^{[0]}/0B$, and $\mathcal O^{[r]}$ denotes the
  functor $T\mapsto J_T^{[r]}$.
\item Let $\gamma'$ be divided powers on an ideal $I'$ of a ring $R'$, and $\sigma:R\to R'$ a ring map with
  $\sigma(I)\subset I'$ and $\gamma'_k(\sigma(x))=\sigma(\gamma_k(x))$ for $x\in I$. Let $\kappa'$ be $\kappa$ followed by
  the map $R/I\to R'/I'$ induced by $\sigma$, and form $\mathrm{Th}(X/R')$ for $(R',I',\gamma',p)$ and $\kappa'$. The
  \emph{restriction of scalars} along $\sigma$ is the functor $\mathrm{Th}(X/R')\to\mathrm{Th}(X/R)$,
  $(B,i,J,\delta,y)\mapsto(B,i\circ\sigma,J,\delta,y)$. With the identity maps of the rings it induces maps
  $H^j(\mathcal Q,\mathcal O)\to H^j(\mathcal Q',\mathcal O)$ for full subcategories $\mathcal Q'$ of
  $\mathrm{Th}(X/R')$ and $\mathcal Q$ of $\mathrm{Th}(X/R)$ with $\mathcal Q'$ mapped into $\mathcal Q$
  (Definition~\ref{gal-realization-split-diagrams}(1)).
\end{enumerate}
For $\nu\geq1$, the category $\mathrm{Th}_\nu(Y/W)$ of Definition~\ref{gal-realization-frobenius-special-fibre} is
$\mathrm{Th}(Y/W_\nu)$ for $(W_\nu,\ell W_\nu,\gamma,\ell)$ and $\mathcal O_K\to\mathbf F_q=W_\nu/\ell$, as $\ell^\nu=0$ in
$W_\nu$; write $\mathrm{Th}^{\mathrm{fl}}_\nu(Y/W)$ for $\mathrm{Th}^{\mathrm{fl}}(Y/W_\nu)$. The category
$\mathrm{Th}^{\mathrm T}_\nu(Y)$ of Definition~\ref{gal-realization-frobenius-teichmuller} is $\mathrm{Th}(Y/A_\nu)$ for
$(A_\nu,I'_\nu,\gamma,\ell)$ and $\mathcal O_K\to\mathbf F_q\to A_\nu/I'_\nu$, and its functor $\Gamma_W$ is the
restriction of scalars along $W_\nu\to A_\nu$.
\end{definition}
\leanref{Development/WeilBound/CoreStatements/PadicHodgeComparison/EtaleCohomologyComparison/PerfectoidPullbacks/SEqOne.lean}{CrystallineSemistableComparison.Pb.nlFun}
\leanref{Development/AbelianVarietyTheory/Definitions/GaloisGeometry/IsDPMorphismBot.lean}{GlobalTheory.LogCris.Thick}
\leanref{Development/AbelianVarietyTheory/Definitions/GaloisGeometry/IsDPMorphismBot.lean}{GlobalTheory.LogCris.Thick.Hom}
\leanref{Development/WeilBound/CoreStatements/PadicHodgeComparison/EtaleCohomologyComparison/PerfectoidPullbacks/SEqOne.lean}{CrystallineSemistableComparison.Pb.InU}
\leanref{Development/WeilBound/CoreStatements/PadicHodgeComparison/EtaleCohomologyComparison/PerfectoidPullbacks/SEqOne.lean}{CrystallineSemistableComparison.Pb.St}
\leanref{Development/WeilBound/CoreStatements/PadicHodgeComparison/EtaleCohomologyComparison/PerfectoidPullbacks/InclStd.lean}{CrystallineSemistableComparison.Pb.FlatT}
\leanref{Development/WeilBound/CoreStatements/PadicHodgeComparison/EtaleCohomologyComparison/PerfectoidPullbacks/SEqOne.lean}{CrystallineSemistableComparison.Pb.thFil}
\leanref{Development/WeilBound/CoreStatements/PadicHodgeComparison/EtaleCohomologyComparison/PerfectoidPullbacks/InclStd.lean}{CrystallineSemistableComparison.Pb.stdI}
\leanref{Development/WeilBound/CoreStatements/PadicHodgeComparison/EtaleCohomologyComparison/PerfectoidPullbacks/DpZeroRight.lean}{CrystallineSemistableComparison.Pb.sq}
\leanref{Development/WeilBound/CoreStatements/PadicHodgeComparison/EtaleCohomologyComparison/PerfectoidPullbacks/SEqOne.lean}{CrystallineSemistableComparison.Pb.NMor}
\leanref{Development/WeilBound/CoreStatements/PadicHodgeComparison/EtaleCohomologyComparison/PerfectoidPullbacks/SEqOne.lean}{CrystallineSemistableComparison.Pb.NMor.mor}

\begin{lemma}\label{gal-realization-thickening-over}
Let $(R,I,\gamma,p)$, $\mathcal O$, $X$ and $\kappa$ be as in Definition~\ref{gal-realization-thickening-nilpotent}.
\begin{enumerate}
\item If $T\to T'$ is a map in $\mathrm{Th}(X/R)$ and $T$ lies over an open $V\subset X$, then $T'$ lies over $V$.
\item A thickening lies over $V_1\cap\cdots\cap V_n$, for open subsets $V_1,\ldots,V_n$ of $X$ with $n\geq1$, if and only
  if it lies over every $V_k$. Every thickening lies over $X$.
\end{enumerate}
\end{lemma}
\leanref{Development/AutomorphicSpectrum/Proofs/BlockDecomposition/InOpenOfThickeningMorphismStatement.lean}{GlobalTheory.T.inOpen_of_thickening_morphism}
\leanref{Development/WeilBound/CoreStatements/PadicHodgeComparison/EtaleCohomologyComparison/PerfectoidPullbacks/PolynomialSiteTraces/AdmissibleMap.lean}{CrystallineSemistableComparison.Pb.Site.inU_hom}
\leanref{Development/AutomorphicSpectrum/Proofs/BlockDecomposition/InOpenFiniteInfIffStatement.lean}{GlobalTheory.T.inOpen_finite_inf_iff}

\begin{proof}
(1) For a map $\varphi:T\to T'$, the point of $T'$ is $\operatorname{Spec}\bar\varphi:\operatorname{Spec}B_{T'}/J_{T'}\to
\operatorname{Spec}B_T/J_T$ followed by the point of $T$, so it factors through $V$ if the point of $T$ does.

(2) A morphism to $X$ factors through an intersection of open subschemes if and only if its image lies in each of them.
\end{proof}
\leanref{Development/AutomorphicSpectrum/Proofs/BlockDecomposition/InOpenFiniteInfIff.lean}{GlobalTheory.CompanionTheory.Thick.inU_iInf_finset}

\begin{definition}\label{gal-realization-thickening-coordinates}
Let $\mathcal O$ be a ring, $X$ an $\mathcal O$-scheme and $d\geq0$.
\begin{enumerate}
\item \emph{Coordinates} of dimension $d$ on an affine open subset $V$ of $X$ are elements $t_1,\ldots,t_d$ of
  $\Gamma(V,\mathcal O_X)$ such that the $\mathcal O$-algebra map $\mathcal O[x_1,\ldots,x_d]\to\Gamma(V,\mathcal O_X)$,
  $x_k\mapsto t_k$, is \'etale.
\item A \emph{coordinate cover} of $X$ is a finite family $(U_\alpha)_{\alpha\in A}$ of open subsets with coordinates of
  dimension $d_\alpha$ on $U_\alpha$, such that $X=\bigcup_\alpha U_\alpha$ and $U_S=\bigcap_{\alpha\in S}U_\alpha$ is
  affine for every nonempty $S\subset A$.
\item Let $(R,I,\gamma,p)$ and $\kappa$ be as in Definition~\ref{gal-realization-thickening-nilpotent}. A \emph{divided
  power chart} of dimension $d$ on an open $V\subset X$ over $(R,I,\gamma,p)$ consists of a thickening
  $D=(B_D,i_D,J_D,\delta_D,y_D)$ of $X$ over $(R,I,\gamma,p)$ that lies over $V$ and elements $t_1,\ldots,t_d\in B_D$,
  such that the ring map $R[x_1,\ldots,x_d]\to B_D$ extending $i_D$ with $x_k\mapsto t_k$ is \'etale, $J_D=i_D(I)B_D$,
  and the following holds. For every ring $B$, ring map $s:R/I\to B$ and $\mathcal O$-morphism
  $z:\operatorname{Spec}B\to X$ that factors through $V$, where $B$ is an $\mathcal O$-algebra through $s\circ\kappa$,
  there is exactly one ring map $\varphi:B_D/J_D\to B$ with $\varphi\circ\bar\imath_D=s$ and
  $y_D\circ\operatorname{Spec}\varphi=z$, where $\bar\imath_D:R/I\to B_D/J_D$ is induced by $i_D$.
\end{enumerate}
\end{definition}
\leanref{Development/WeilBound/CoreStatements/PadicHodgeComparison/EtaleCohomologyComparison/PerfectoidPullbacks/SEqOne.lean}{CrystallineSemistableComparison.Pb.Coord}
\leanref{Development/WeilBound/CoreStatements/PadicHodgeComparison/EtaleCohomologyComparison/PerfectoidPullbacks/SEqOne.lean}{CrystallineSemistableComparison.Pb.Cov}
\leanref{Development/WeilBound/CoreStatements/PadicHodgeComparison/EtaleCohomologyComparison/PerfectoidPullbacks/SEqOne.lean}{CrystallineSemistableComparison.Pb.Chart}

\begin{lemma}\label{gal-realization-thickening-cover}
Let $\mathcal O$ be a ring and $X$ an $\mathcal O$-scheme.
\begin{enumerate}
\item If $t_1,\ldots,t_d$ are coordinates on an open $V\subset X$ and $W\subset V$ is an affine open, the restrictions of
  the $t_k$ to $W$ are coordinates on $W$.
\item If $X$ is separated, quasi-compact and smooth over $\mathcal O$, it has a coordinate cover, and if $X$ is smooth of
  relative dimension $d$ over $\mathcal O$, it has one with $d_\alpha=d$ for every $\alpha$.
\end{enumerate}
\end{lemma}
\leanref{Development/AutomorphicSpectrum/Proofs/BlockDecomposition/CoordinatesRestrictAffineOpenStatement.lean}{GlobalTheory.T.coordinates_restrict_affine_open}
\leanref{Development/AutomorphicSpectrum/Proofs/BlockDecomposition/ExistsSmoothCoordinateCoverStatement.lean}{GlobalTheory.T.exists_smooth_coordinate_cover}

\begin{proof}
(1) The open immersion $W\to V$ of affine schemes is \'etale, so the restriction map $\Gamma(V,\mathcal
O_X)\to\Gamma(W,\mathcal O_X)$ is \'etale. Its composite with the \'etale map $\mathcal O[x_1,\ldots,x_d]\to\Gamma(V,
\mathcal O_X)$ is the map for the restricted $t_k$, and composites of \'etale maps are \'etale.

(2) Let $x\in X$. As $X$ is smooth over $\mathcal O$, $x$ has an affine open neighbourhood $W$ with an \'etale $\mathcal
O$-algebra map $\mathcal O[x_1,\ldots,x_n]\to\Gamma(W,\mathcal O_X)$, where $n=d$ if $X$ is smooth of relative dimension
$d$ \cite[Corollaire~17.11.4]{Gro67}; the images of the $x_k$ are coordinates on $W$. As $X$ is quasi-compact, finitely
many such opens cover $X$. As $X$ is separated over $\mathcal O$, the intersection of two affine opens $U$ and $U'$ is
the inverse image of $U\times_{\mathcal O}U'$ under the diagonal, which is a closed immersion, so it is affine. By
induction every nonempty finite intersection of the chosen opens is affine.
\end{proof}
\leanref{Development/AutomorphicSpectrum/Proofs/BlockDecomposition/CoordinatesRestrictAffineOpen.lean}{GlobalTheory.CompanionTheory.CoordRes.main}
\leanref{Development/AutomorphicSpectrum/Proofs/BlockDecomposition/CoordinatesRestrictAffineOpen.lean}{GlobalTheory.CompanionTheory.CoordRes.etale_rs}
\leanref{Development/AutomorphicSpectrum/Proofs/BlockDecomposition/ExistsSmoothCoordinateCover.lean}{GlobalTheory.CompanionTheory.CovEx.exists_cov}
\leanref{Development/AutomorphicSpectrum/Proofs/BlockDecomposition/ExistsSmoothCoordinateCover.lean}{GlobalTheory.CompanionTheory.CovEx.cov_of_local}
\leanref{Development/AutomorphicSpectrum/Proofs/BlockDecomposition/ExistsSmoothCoordinateCover.lean}{GlobalTheory.CompanionTheory.CovEx.inter_affine}
\leanref{Development/AutomorphicSpectrum/Proofs/BlockDecomposition/ExistsSmoothCoordinateCover.lean}{GlobalTheory.CompanionTheory.CovEx.coord_basicOpen}

\begin{lemma}\label{gal-realization-thickening-chart}
Let $(R,I,\gamma,p)$, $\mathcal O$, $X$ and $\kappa$ be as in Definition~\ref{gal-realization-thickening-nilpotent}, $V$
an open subset of $X$ and $d\geq0$.
\begin{enumerate}
\item Let $D$ with $t_1,\ldots,t_d$ be a divided power chart on $V$
  (Definition~\ref{gal-realization-thickening-coordinates}(3)). A thickening that receives a map from $D$ lies over $V$,
  and every thickening $T$ that lies over $V$ receives one. For maps $\sigma,\sigma':D\to T$ and $b\in B_D$,
  $\sigma'(b)-\sigma(b)\in J_T$. For $\sigma:D\to T$ and $y_1,\ldots,y_d\in J_T$ there is exactly one map $\sigma':D\to T$
  with $\sigma'(t_k)-\sigma(t_k)=y_k$ for every $k$.
\item For a divided power chart $D$ and $r\geq0$, $J_D^{[r]}=i_D(I^{[r]})B_D$, where $I^{[r]}$ is formed for $\gamma$ as
  in Definition~\ref{gal-realization-kappa-approximations}.
\item Let $p\neq0$ be nilpotent in $R$, and $t_1,\ldots,t_d$ coordinates on $V$. There is a divided power chart $D$ with
  $t'_1,\ldots,t'_d$ on $V$ such that $B_D$ is flat over $R$ and the ring map $\Gamma(V,\mathcal O_X)\to B_D/J_D$ induced
  by $y_D$ sends each $t_k$ to the class of $t'_k$.
\item Let $\sigma:R\to R'$, $(R',I',\gamma')$ and $\kappa'$ be as in
  Definition~\ref{gal-realization-thickening-nilpotent}(4), let $p\neq0$, and let $D$ with $t_1,\ldots,t_d$ be a divided
  power chart on $V$ over $(R,I,\gamma,p)$. There are a divided power chart $D'$ with $t'_1,\ldots,t'_d$ on $V$ over
  $(R',I',\gamma',p)$ and a map $e$ from $D$ to the restriction of scalars of $D'$ along $\sigma$ with $e(t_k)=t'_k$ for
  every $k$, such that $B_{D'}$ is flat over $R'$ if $B_D$ is flat over $R$, and the ring map $R'\otimes_RB_D\to B_{D'}$,
  $a\otimes b\mapsto i_{D'}(a)e(b)$, is bijective.
\end{enumerate}
\end{lemma}
\leanref{Development/AutomorphicSpectrum/Proofs/BlockDecomposition/ChartIsGood.lean}{GlobalTheory.CompanionTheory.ChartMaps.good}
\leanref{Development/AutomorphicSpectrum/Proofs/BlockDecomposition/ChartIsGood.lean}{GlobalTheory.CompanionTheory.ChartMaps.inU_of_hom}
\leanref{Development/AutomorphicSpectrum/Proofs/BlockDecomposition/ChartDividedPowersGradedConstant.lean}{GlobalTheory.CompanionTheory.ChartFil.sub_mem_J}
\leanref{Development/AutomorphicSpectrum/Proofs/BlockDecomposition/ChartDividedPowersGradedConstant.lean}{GlobalTheory.CompanionTheory.ChartFil.chart_grConst}
\leanref{Development/AutomorphicSpectrum/Proofs/BlockDecomposition/FlatOfEtalePoly.lean}{GlobalTheory.CompanionTheory.FlatChartP.exists_flat_chart}
\leanref{Development/AutomorphicSpectrum/Proofs/BlockDecomposition/ChartBcOfNe.lean}{GlobalTheory.CompanionTheory.ChartBc.chart_bc_of_ne}

\begin{proof}
(1) Elements of $J_T$ are nilpotent: if $p=0$, $J_T$ is nilpotent, and if $p\neq0$, $x^p=p!\,\delta_p(x)$ for $x\in J_T$,
and $p$ is nilpotent in $B_T$. So for an \'etale ring map $P\to E$ and ring maps $f:P\to B_T$ and $g:E\to B_T/J_T$ that
agree on $P$, there is exactly one ring map $h:E\to B_T$ that extends $f$ and lifts $g$, by the equivalence between
\'etale schemes over $\operatorname{Spec}B_T$ and over $\operatorname{Spec}B_T/J_T$ \cite[Th\'eor\`eme~18.1.2]{Gro67},
applied to sections of $\operatorname{Spec}(E\otimes_PB_T)$. Applied to $R[x_1,\ldots,x_d]\to B_D$: for a ring map
$g:B_D/J_D\to B_T/J_T$ with $g\circ\bar\imath_D=\bar\imath_T$ and lifts $y'_k\in B_T$ of the $g(t_k)$, there is exactly
one ring map $f:B_D\to B_T$ with $f\circ i_D=i_T$ and $f(t_k)=y'_k$ that lifts $g$. Such an $f$ is a divided power map:
$f(J_D)\subset J_T$, and the set of $x\in J_D$ with $f(\delta_{D,n}(x))=\delta_{T,n}(f(x))$ for all $n$ contains
$i_D(I)$, as $i_D$ and $i_T$ are divided power maps, and is closed under sums and under multiplication by $B_D$, by
$\delta_n(x+x')=\sum_{a+b=n}\delta_a(x)\delta_b(x')$ and $\delta_n(cx)=c^n\delta_n(x)$. If moreover
$y_D\circ\operatorname{Spec}g=y_T$, then $f$ is a map $D\to T$.

A map $\sigma:D\to T$ gives $y_T=y_D\circ\operatorname{Spec}\bar\sigma$, which factors through $V$. Let $T$ lie over
$V$. The defining property of $D$, for $B=B_T/J_T$, $s=\bar\imath_T$ and $z=y_T$, gives exactly one ring map
$g:B_D/J_D\to B_T/J_T$ with $g\circ\bar\imath_D=\bar\imath_T$ and $y_D\circ\operatorname{Spec}g=y_T$. Lifting the
$g(t_k)$ arbitrarily gives a map $D\to T$. Every map $\sigma:D\to T$ has $\bar\sigma=g$, by this uniqueness, so
$\sigma'(b)-\sigma(b)\in J_T$ for two maps $\sigma,\sigma'$. Given $\sigma$ and the $y_k$, the lift $f$ of $g$ with
$f(t_k)=\sigma(t_k)+y_k$ is a map $\sigma'$ as required. Any other such $\sigma'$ also lifts $g$ and has the same values
on the $t_k$, so it is $f$.

(2) As $i_D$ is a divided power map, $i_D(I^{[r]})\subset J_D^{[r]}$. Conversely, the set of $x\in J_D=i_D(I)B_D$ with
$\delta_{D,n}(x)\in i_D(I^{[n]})B_D$ for all $n$ contains the $i_D(z)$, $z\in I$, as $\delta_{D,n}(i_D(z))=
i_D(\gamma_n(z))$, and is closed under sums and under multiplication by $B_D$, by the two formulas in (1) and
$I^{[a]}I^{[b]}\subset I^{[a+b]}$; so it is $J_D$. A generator $\prod_j\delta_{D,k_j}(x_j)$ of $J_D^{[r]}$, with
$\sum_jk_j\geq r$, therefore lies in $i_D(I^{[\sum_jk_j]})B_D\subset i_D(I^{[r]})B_D$.

(3) Let $G=\Gamma(V,\mathcal O_X)$, an $\mathcal O$-algebra, and $\rho:\mathcal O[x_1,\ldots,x_d]\to G$ the \'etale map
with $x_k\mapsto t_k$. Let $\pi:R[x_1,\ldots,x_d]\to(R/I)[x_1,\ldots,x_d]$ be the reduction, a surjection with kernel
$IR[x_1,\ldots,x_d]$, and $\mu:\mathcal O[x_1,\ldots,x_d]\to(R/I)[x_1,\ldots,x_d]$ the map induced by $\kappa$.

\emph{Step 1: an \'etale lift.} The base change $\bar E=(R/I)[x_1,\ldots,x_d]\otimes_{\mathcal O[x_1,\ldots,x_d]}G$ of
$\rho$ along $\mu$ is \'etale over $(R/I)[x_1,\ldots,x_d]$ \cite[Tag 00U2]{Sta}. \'Etale algebras lift along the
surjection $\pi$ \cite[Tag 04D1]{Sta}: there is an \'etale $R[x_1,\ldots,x_d]$-algebra $E$ with $E/IE\cong\bar E$ as
$(R/I)[x_1,\ldots,x_d]$-algebras. Let $i:R\to E$ be the structure map, $t'_k\in E$ the image of $x_k$, and $j:G\to
E/IE$ the map to the second factor of $\bar E$. Then $R[x_1,\ldots,x_d]\to E$, $x_k\mapsto t'_k$, is \'etale,
$j(t_k)$ is the class of $t'_k$, and $j$ and the map $R/I\to E/IE$ induced by $i$ agree on $\mathcal O$, through
$\mathcal O\to G$ and $\kappa$. By the universal property of the tensor product, for every ring $B$, ring maps
$E/IE\to B$ correspond to pairs of ring maps $s:R/I\to B$ and $g:G\to B$ that agree on $\mathcal O$: the map
$(R/I)[x_1,\ldots,x_d]\to B$ is then determined by $s$ and $x_k\mapsto g(t_k)$.

\emph{Step 2: the chart.} The ring $E$ is flat over $R$, as it is flat over $R[x_1,\ldots,x_d]$ \cite[Tag 00U2]{Sta},
which is free over $R$ \cite[Tag 00HC]{Sta}. So $\gamma$ extends to $E$ \cite[Tag 07H1]{Sta}: there are divided powers
$\delta$ on $IE$ for which $i$ is a divided power map. As $V$ is affine, ring maps $G\to C$ correspond to morphisms
$\operatorname{Spec}C\to X$ that factor through $V$ \cite[Tag 01I1]{Sta}; let $y:\operatorname{Spec}E/IE\to X$
correspond to $j$. It is an $\mathcal O$-morphism for the structure of $E/IE$ through $\kappa$ and $i$, as $j$ and
$i$ agree on $\mathcal O$. As $p$ is nilpotent in $R$, it is nilpotent in $E$. So $D=(E,i,IE,\delta,y)$ is a thickening
of $X$ over $(R,I,\gamma,p)$ that lies over $V$, with $J_D=i(I)E$, $R[x_1,\ldots,x_d]\to E$ \'etale and $E$ flat over
$R$. The ring map $G\to E/IE$ induced by $y$ is $j$, which sends $t_k$ to the class of $t'_k$.

It remains to check the property of a divided power chart. Let $B$ be a ring, $s:R/I\to B$ a ring map and
$z:\operatorname{Spec}B\to X$ an $\mathcal O$-morphism that factors through $V$, with $B$ an $\mathcal O$-algebra through
$s\circ\kappa$. Then $z$ corresponds to a ring map $g:G\to B$, which agrees with $s$ on $\mathcal O$, as $z$ is an
$\mathcal O$-morphism. By Step 1 there is exactly one $\varphi:E/IE\to B$ that restricts to $s$ on $R/I$ and to $g$ on
$G$; the second condition says $y\circ\operatorname{Spec}\varphi=z$. This is the defining property of the chart.
% Lean: GlobalTheory.CompanionTheory.FlatChartP.exists_flat_chart. Step 1 is EtaleLift.nonempty_lift, the structure
% B.Lift, through exists_lift_hom, whose lift of the étale algebra is exists_etale_lift_of_surjective, proved through a
% standard smooth presentation and the localization at its Jacobian (tensorAwayEquiv), and tensor_existsUnique_lift. Step
% 2 is flat_of_etale_poly, JLiftDP.dpFlat (the Lean's proof of Stacks 07H1(3)), mkT, top_le_specMap and rep_lift, with
% evalAt_specMap, evalAt_base and eq_specMap for the correspondence between ring maps out of Γ(X, V) and morphisms
% into V.

(4) Let $B'=R'\otimes_RB_D$, with $i':R'\to B'$, $a\mapsto a\otimes1$, $e:B_D\to B'$, $b\mapsto1\otimes b$, and
$t'_k=e(t_k)$. Then $e\circ i_D=i'\circ\sigma$.

The map $R'[x_1,\ldots,x_d]\to B'$ extending $i'$ with $x_k\mapsto t'_k$ is the base change of the \'etale map
$R[x_1,\ldots,x_d]\to B_D$ along $R[x_1,\ldots,x_d]\to R'[x_1,\ldots,x_d]$, as $B'=R'[x_1,\ldots,x_d]
\otimes_{R[x_1,\ldots,x_d]}B_D$; so it is \'etale, and $B'$ is flat over $R'[x_1,\ldots,x_d]$ and therefore over $R'$
\cite[Tags 00U2 and 00HC]{Sta}. So $\gamma'$ extends to $B'$ \cite[Tag 07H1]{Sta}: there are divided powers $\delta'$ on
$J'=i'(I')B'$ for which $i'$ is a divided power map. As $p\neq0$ is nilpotent in $B_D$, it is nilpotent in $B'$.

The map $e$ carries $J_D=i_D(I)B_D$ into $J'$, as $e(i_D(x))=i'(\sigma(x))$ and $\sigma(I)\subset I'$. It is a divided
power map: the set of $x\in J_D$ with $\delta'_n(e(x))=e(\delta_{D,n}(x))$ for all $n$ contains the $i_D(x)$, $x\in I$,
as $\delta'_n(i'(\sigma(x)))=i'(\gamma'_n(\sigma(x)))=i'(\sigma(\gamma_n(x)))=e(i_D(\gamma_n(x)))$, and it is closed
under sums and under multiplication by $B_D$, by the two formulas in (1). Let $h:B_D/J_D\to B'/J'$ and $\psi:R'/I'\to
B'/J'$ be induced by $e$ and $i'$, and $\bar\sigma:R/I\to R'/I'$ by $\sigma$; then $h\circ\bar\imath_D=\psi\circ
\bar\sigma$. Let $y'=y_D\circ\operatorname{Spec}h$. It is an $\mathcal O$-morphism for the structure of $B'/J'$ through
$\psi\circ\kappa'=\psi\circ\bar\sigma\circ\kappa=h\circ\bar\imath_D\circ\kappa$. So $D'=(B',i',J',\delta',y')$ is a
thickening of $X$ over $(R',I',\gamma',p)$. It lies over $V$, as $y_D$ factors through $V$, and $e$ is a map from $D$ to
the restriction of scalars $(B',i'\circ\sigma,J',\delta',y')$ of $D'$ along $\sigma$, with $e(t_k)=t'_k$. The map
$R'\otimes_RB_D\to B'$, $a\otimes b\mapsto i'(a)e(b)$, is the identity, and $B'$ is flat over $R'$ in every case.

It remains to check the property of a divided power chart for $D'$. By the universal property of the tensor product of
$R$-algebras, ring maps $B'\to B$ correspond to pairs of ring maps $R'\to B$ and $B_D\to B$ that agree on $R$. Such a map
kills $J'$ if and only if its first component kills $I'$, and then the second kills $J_D=i_D(I)B_D$, as $\sigma(I)\subset
I'$. So ring maps $B'/J'\to B$ correspond to pairs of ring maps $a:R'/I'\to B$ and $b:B_D/J_D\to B$ with $a\circ
\bar\sigma=b\circ\bar\imath_D$, through composition with $\psi$ and $h$. Let $B$, $s:R'/I'\to B$ and $z:\operatorname{Spec}
B\to X$ be as in Definition~\ref{gal-realization-thickening-coordinates}(3) for $D'$, with $z$ an $\mathcal O$-morphism for
$s\circ\kappa'=s\circ\bar\sigma\circ\kappa$ that factors through $V$. The chart property of $D$, for $s\circ\bar\sigma$ and
$z$, gives exactly one $b:B_D/J_D\to B$ with $b\circ\bar\imath_D=s\circ\bar\sigma$ and $y_D\circ\operatorname{Spec}b=z$.
The pair $(s,b)$ gives $\varphi:B'/J'\to B$ with $\varphi\circ\psi=s$ and $\varphi\circ h=b$, so $y'\circ
\operatorname{Spec}\varphi=y_D\circ\operatorname{Spec}b=z$. If $\varphi_2$ also has $\varphi_2\circ\psi=s$ and $y'\circ
\operatorname{Spec}\varphi_2=z$, then $\varphi_2\circ h$ has the two properties of $b$, as $\varphi_2\circ h\circ
\bar\imath_D=\varphi_2\circ\psi\circ\bar\sigma=s\circ\bar\sigma$, so $\varphi_2\circ h=b$ and $\varphi_2=\varphi$.
% Lean: GlobalTheory.CompanionTheory.ChartBc.chart_bc (LiftUnique.lean) for an NMor m with u and g identities: the
% étale map is ChartAlg.etale_tensor (RingHom.Etale.isStableUnderBaseChange), the flatness ChartAlg.flat_of_etale, the
% divided powers JLiftDP.dpFlat, the divided power map e is dp_right (IsDPMorphism.on_span), the correspondence of ring
% maps is lift_unique, and the chart property is c'.rep from c.rep. The hypothesis p ≠ 0 enters through hnil.
\end{proof}
\leanref{Development/AutomorphicSpectrum/Proofs/BlockDecomposition/ChartIsGood.lean}{GlobalTheory.CompanionTheory.ChartMaps.nil_J}
\leanref{Development/AutomorphicSpectrum/Proofs/BlockDecomposition/EtaleLiftingOfNilIdealStatement.lean}{GlobalTheory.T.etale_lifting_of_nil_ideal}
\leanref{Development/AutomorphicSpectrum/Proofs/BlockDecomposition/ChartIsGood.lean}{GlobalTheory.CompanionTheory.ChartMaps.exists_unique_lift}
\leanref{Development/AutomorphicSpectrum/Proofs/BlockDecomposition/ChartIsGood.lean}{GlobalTheory.CompanionTheory.ChartMaps.isDP}
\leanref{Development/AutomorphicSpectrum/Proofs/BlockDecomposition/ChartDividedPowersGradedConstant.lean}{GlobalTheory.CompanionTheory.ChartFil.dividedPowerFiltration_eq_map}
\leanref{Development/AutomorphicSpectrum/Proofs/BlockDecomposition/ExistsEtaleLiftOfSurjective.lean}{GlobalTheory.CompanionTheory.EtaleLift.nonempty_lift}
\leanref{Development/AutomorphicSpectrum/Proofs/BlockDecomposition/ExistsEtaleLiftOfSurjective.lean}{GlobalTheory.CompanionTheory.EtaleLift.exists_lift_hom}
\leanref{Development/AutomorphicSpectrum/Proofs/BlockDecomposition/ExistsEtaleLiftOfSurjective.lean}{GlobalTheory.CompanionTheory.EtaleLift.exists_etale_lift_of_surjective}
\leanref{Development/AutomorphicSpectrum/Proofs/BlockDecomposition/ExistsDpOfCompOnGens.lean}{GlobalTheory.CompanionTheory.JLiftDP.dpFlat}
\leanref{Development/AutomorphicSpectrum/Proofs/BlockDecomposition/FlatOfEtalePoly.lean}{GlobalTheory.CompanionTheory.FlatChartP.flat_of_etale_poly}
\leanref{Development/AutomorphicSpectrum/Proofs/BlockDecomposition/FlatOfEtalePoly.lean}{GlobalTheory.CompanionTheory.FlatChartP.rep_lift}
\leanref{Development/AutomorphicSpectrum/Proofs/BlockDecomposition/ExistsChartFromCoordinates.lean}{GlobalTheory.CompanionTheory.ChartAlg.etale_tensor}
\leanref{Development/AutomorphicSpectrum/Proofs/BlockDecomposition/ExistsChartFromCoordinates.lean}{GlobalTheory.CompanionTheory.ChartAlg.flat_of_etale}
\leanref{Development/AutomorphicSpectrum/Proofs/BlockDecomposition/LiftUnique.lean}{GlobalTheory.CompanionTheory.ChartBc.dp_right}
\leanref{Development/AutomorphicSpectrum/Proofs/BlockDecomposition/LiftUnique.lean}{GlobalTheory.CompanionTheory.ChartBc.lift_unique}
\leanref{Development/AutomorphicSpectrum/Proofs/BlockDecomposition/LiftUnique.lean}{GlobalTheory.CompanionTheory.ChartBc.chart_bc}

\begin{definition}\label{gal-realization-thickening-torsor}
Let $R$ be a ring and $\mathcal C$ a category.
\begin{enumerate}
\item A \emph{divided power functor} $\Phi$ on $\mathcal C$ over $R$ consists of a functor $A$ from $\mathcal C$ to
  commutative $R$-algebras, ideals $J(T)\subset A(T)$ with $A(\varphi)(J(T))\subset J(T')$ for $\varphi:T\to T'$, and maps
  $\gamma_k:A(T)\to A(T)$, $k\geq0$, such that $\gamma_k(0)=0$ for $k\geq1$ and, for $x,x'\in J(T)$, $\gamma_0(x)=1$,
  $\gamma_1(x)=x$, $\gamma_k(x+x')=\sum_{a+b=k}\gamma_a(x)\gamma_b(x')$ and $A(\varphi)(\gamma_k(x))=
  \gamma_k(A(\varphi)(x))$. A \emph{family of ideals} $\mathfrak a$ of $\Phi$ consists of ideals $\mathfrak a(T)\subset
  A(T)$ with $A(\varphi)(\mathfrak a(T))\subset\mathfrak a(T')$; for two of them, $\mathfrak a/\mathfrak b$ is the functor
  sending $T$ to the image of $\mathfrak a(T)$ in $A(T)/\mathfrak b(T)$. For a full subcategory $\mathcal Q$ of a category
  $\mathrm{Th}(X/R)$ of Definition~\ref{gal-realization-thickening-nilpotent}, $\Phi_{\mathcal Q}$ is the divided power
  functor on $\mathcal Q$ given by the $R$-algebras $B_T$, the ideals $J_T$ and their divided powers.
\item A \emph{torsor datum} of dimension $d$ for $\Phi$ consists of an object $D$ that has a map to every object, and
  elements $\tau_k(\sigma,\sigma')\in J(T)$, for $1\leq k\leq d$ and maps $\sigma,\sigma':D\to T$, such that
  $\tau_k(\sigma,\sigma'')=\tau_k(\sigma,\sigma')+\tau_k(\sigma',\sigma'')$, $\tau_k(\varphi\sigma,\varphi\sigma')=
  A(\varphi)(\tau_k(\sigma,\sigma'))$ for $\varphi:T\to T'$, and for every $\sigma:D\to T$ and $y_1,\ldots,y_d\in J(T)$
  there is exactly one $\sigma':D\to T$ with $\tau_k(\sigma,\sigma')=y_k$ for every $k$.
\item A \emph{divided power polynomial structure} on $\Phi$ assigns to each object $T$ an object $P(T)$, a map
  $j_T:T\to P(T)$ and an element $X_T\in J(P(T))$, such that for $\varphi:T\to T'$ and $x\in J(T')$ there is exactly one
  $\psi:P(T)\to T'$ with $\psi j_T=\varphi$ and $A(\psi)(X_T)=x$, and every element of $A(P(T))$ is a finite sum
  $\sum_kA(j_T)(c_k)\gamma_k(X_T)$ with $c_k\in A(T)$. It \emph{respects} a sequence $(\mathfrak a_r)_{r\in\mathbf Z}$ of
  families of ideals in degree $r$ if for every $T$ and every finitely supported sequence $(c_k)_{k\geq0}$ in $A(T)$, the
  sum $\sum_kA(j_T)(c_k)\gamma_k(X_T)$ lies in $\mathfrak a_r(P(T))$ if and only if $c_k\in\mathfrak a_{r-k}(T)$ for
  every $k$. It is a \emph{ladder up to} $r$ for sequences $(\mathfrak a_r)$ and $(\mathfrak b_r)$ if $\mathfrak
  b_{r'}(T)\subset\mathfrak a_{r'}(T)$ for all $r'$ and $T$ and it respects both sequences in every degree $r'\leq r$.
\item Let $\sigma:R\to R'$ be a ring map and $\Phi'$, given by $A'$, $J'$ and $\gamma'$, a divided power functor on a
  category $\mathcal C'$ over $R'$. A \emph{morphism} $\Phi\to\Phi'$ over $\sigma$ consists of a functor
  $G:\mathcal C'\to\mathcal C$ and ring maps $\rho_T:A(G(T))\to A'(T)$, natural in $T$, with $\rho_T(a\cdot1)=
  \sigma(a)\cdot1$ for $a\in R$, $\rho_T(J(G(T)))\subset J'(T)$ and $\rho_T\gamma_k=\gamma'_k\rho_T$ on $J(G(T))$. Torsor
  data $(D,\tau)$ for $\Phi$ and $(D',\tau')$ for $\Phi'$ of the same dimension are \emph{compatible} along a map
  $e:D\to G(D')$ if $\rho_T(\tau_k(G(s)e,G(s')e))=\tau'_k(s,s')$ for all $s,s':D'\to T$. For $c\in R'$ and families
  $\mathfrak a,\mathfrak b$ of $\Phi$ and $\mathfrak a',\mathfrak b'$ of $\Phi'$ with $c\rho_T(\mathfrak a(G(T)))\subset
  \mathfrak a'(T)$ and $c\rho_T(\mathfrak b(G(T)))\subset\mathfrak b'(T)$ for all $T$, the maps $x\mapsto c\rho_T(x)$ and
  $G$ induce $\sigma$-semilinear maps $H^j(\mathcal C,\mathfrak a/\mathfrak b)\to H^j(\mathcal C',\mathfrak a'/\mathfrak
  b')$ (Definition~\ref{gal-realization-split-diagrams}(1)), the maps of \emph{multiplication by} $c$.
\end{enumerate}
\end{definition}
\leanref{Development/WeilBound/CoreStatements/PadicHodgeComparison/EtaleCohomologyComparison/PerfectoidPullbacks/DpZeroRight.lean}{CrystallineSemistableComparison.Pb.PDF}
\leanref{Development/WeilBound/CoreStatements/PadicHodgeComparison/EtaleCohomologyComparison/PerfectoidPullbacks/DpZeroRight.lean}{CrystallineSemistableComparison.Pb.IdF}
\leanref{Development/WeilBound/CoreStatements/PadicHodgeComparison/EtaleCohomologyComparison/PerfectoidPullbacks/SEqOne.lean}{CrystallineSemistableComparison.Pb.thP}
\leanref{Development/WeilBound/CoreStatements/PadicHodgeComparison/EtaleCohomologyComparison/PerfectoidPullbacks/DpZeroRight.lean}{CrystallineSemistableComparison.Pb.Frame}
\leanref{Development/WeilBound/CoreStatements/PadicHodgeComparison/EtaleCohomologyComparison/PerfectoidPullbacks/DpZeroRight.lean}{CrystallineSemistableComparison.Pb.Poly}
\leanref{Development/WeilBound/CoreStatements/PadicHodgeComparison/EtaleCohomologyComparison/PerfectoidPullbacks/DpZeroRight.lean}{CrystallineSemistableComparison.Pb.Poly.Crit}
\leanref{Development/WeilBound/CoreStatements/PadicHodgeComparison/EtaleCohomologyComparison/PerfectoidPullbacks/DpZeroRight.lean}{CrystallineSemistableComparison.Pb.Poly.Ladder}
\leanref{Development/WeilBound/CoreStatements/PadicHodgeComparison/EtaleCohomologyComparison/PerfectoidPullbacks/DpZeroRight.lean}{CrystallineSemistableComparison.Pb.SMor}
\leanref{Development/WeilBound/CoreStatements/PadicHodgeComparison/EtaleCohomologyComparison/PerfectoidPullbacks/DpZeroRight.lean}{CrystallineSemistableComparison.Pb.SMor.Fr}
\leanref{Development/WeilBound/CoreStatements/PadicHodgeComparison/EtaleCohomologyComparison/PerfectoidPullbacks/CapOComp.lean}{CrystallineSemistableComparison.Pb.SMor.LeC}
\leanref{Development/WeilBound/CoreStatements/PadicHodgeComparison/EtaleCohomologyComparison/PerfectoidPullbacks/CapOComp.lean}{CrystallineSemistableComparison.Pb.SMor.sqC}

\begin{lemma}\label{gal-realization-thickening-chart-torsor}
Let $(R,I,\gamma,p)$, $\mathcal O$, $X$ and $\kappa$ be as in Definition~\ref{gal-realization-thickening-nilpotent}, and
$\mathcal Q$ a full subcategory of $\mathrm{Th}(X/R)$.
\begin{enumerate}
\item Let $D$ with $t_1,\ldots,t_d$ be a divided power chart on an open $V$
  (Definition~\ref{gal-realization-thickening-coordinates}(3)) such that $D$ lies in $\mathcal Q$ and every thickening in
  $\mathcal Q$ lies over $V$. Then $D$ with $\tau_k(\sigma,\sigma')=\sigma'(t_k)-\sigma(t_k)$ is a torsor datum of
  dimension $d$ for $\Phi_{\mathcal Q}$ (Definition~\ref{gal-realization-thickening-torsor}).
\item Let $\sigma:R\to R'$, $(R',I',\gamma')$ and $\kappa'$ be as in
  Definition~\ref{gal-realization-thickening-nilpotent}(4), and $\mathcal Q'$ a full subcategory of $\mathrm{Th}(X/R')$
  that the restriction of scalars $G$ along $\sigma$ maps into $\mathcal Q$. Then $G$ with the identity maps of the rings
  is a morphism $\Phi_{\mathcal Q}\to\Phi_{\mathcal Q'}$ over $\sigma$. If $D$ is as in (1) for $\mathcal Q$, $D'$ with
  $t'_1,\ldots,t'_d$ is as in (1) for $\mathcal Q'$, and $e:D\to G(D')$ is a map with $e(t_k)=t'_k$ for every $k$, the
  torsor data of (1) are compatible along $e$.
\item Let $p\neq0$, and let $\mathcal Q$ be $\mathrm{Th}(X/R)_{\mathcal V}$ or $\mathrm{Th}^{\mathrm{fl}}(X/R)_{\mathcal
  V}$ for a set $\mathcal V$ of open subsets of $X$. The following is a divided power polynomial structure on
  $\Phi_{\mathcal Q}$: $P(T)$ is the thickening whose ring is the divided power polynomial algebra $B_T\langle X\rangle$,
  with the ideal generated by $J_T$ and the $X^{[k]}$, $k\geq1$, the divided powers on it that extend $\delta_T$ and send
  $X$ to the $X^{[k]}$, and the point of $T$ through the map from $B_T/J_T$ to the quotient by this ideal; $j_T$ is the
  inclusion of $B_T$; and $X_T=X$. It respects, in every degree,
  the sequences $(J^{[r]})_r$ and $(0)_r$, and for every $c\in R$ the sequences $(J^{[r]}+cB)_r$ and $(cB)_r$.
\end{enumerate}
\end{lemma}
\leanref{Development/WeilBound/CoreStatements/PadicHodgeComparison/EtaleCohomologyComparison/PerfectoidPullbacks/SEqOne.lean}{CrystallineSemistableComparison.Pb.Chart.frame}
\leanref{Development/WeilBound/CoreStatements/PadicHodgeComparison/EtaleCohomologyComparison/PerfectoidPullbacks/SEqOne.lean}{CrystallineSemistableComparison.Pb.NMor.smor}
\leanref{Development/AutomorphicSpectrum/Proofs/BlockDecomposition/SmorFrFrame.lean}{GlobalTheory.CompanionTheory.ChartFrames.smor_Fr_frame}
\leanref{Development/WeilBound/CoreStatements/PadicHodgeComparison/EtaleCohomologyComparison/PerfectoidPullbacks/ChartGrConst.lean}{CrystallineSemistableComparison.Pb.poly_all}
\leanref{Development/WeilBound/CoreStatements/PadicHodgeComparison/EtaleCohomologyComparison/PerfectoidPullbacks/PolynomialSiteTraces/PolyAll.lean}{CrystallineSemistableComparison.Pb.Site.poly_all'}
\leanref{Development/WeilBound/CoreStatements/PadicHodgeComparison/EtaleCohomologyComparison/PerfectoidPullbacks/PolynomialSiteTraces/PolyAll.lean}{CrystallineSemistableComparison.Pb.Site.polyAll}
\leanref{Development/AutomorphicSpectrum/Proofs/BlockDecomposition/StdSomeZero.lean}{GlobalTheory.CompanionTheory.ChartLocBc.polyS}

\begin{proof}
(1) By Lemma~\ref{gal-realization-thickening-chart}(1), as every $T\in\mathcal Q$ lies over $V$, $D$ has a map to $T$,
and $\tau_k(\sigma,\sigma')\in J_T$; maps of $\mathcal Q$ are maps of thickenings. The identity for
$\tau_k(\sigma,\sigma'')$ holds in $B_T$, and $\varphi(\sigma'(t_k))-\varphi(\sigma(t_k))=\varphi(\tau_k(\sigma,
\sigma'))$. The last condition is the last claim of Lemma~\ref{gal-realization-thickening-chart}(1).

(2) The ring of $G(T)$ is $B_T$ with $i_T\circ\sigma$, so its identity map sends $a\cdot1$ to $\sigma(a)\cdot1$, and it
preserves the ideals and their divided powers. For $s,s':D'\to T$, $\tau_k(G(s)e,G(s')e)=s'(e(t_k))-s(e(t_k))=
s'(t'_k)-s(t'_k)=\tau'_k(s,s')$.

(3) Let $T=(B,i,J,\delta,y)$ be in $\mathcal Q$. The divided power polynomial algebra $B\langle X\rangle$ is the free
$B$-module with basis $X^{[k]}$, $k\geq0$, with $X^{[0]}=1$ and $X^{[a]}X^{[b]}=\binom{a+b}{a}X^{[a+b]}$, and $X=X^{[1]}$.
Let $J\langle X\rangle=JB\langle X\rangle+B\langle X\rangle_+$, where $B\langle X\rangle_+$ is the ideal generated by
the $X^{[k]}$, $k\geq1$. There is exactly one divided power structure $\eta$ on $J\langle X\rangle$ such that
$\eta_k(X)=X^{[k]}$ for every $k$ and $j:B\to B\langle X\rangle$ is a divided power map; and for every ring $C$ with
divided powers $\varepsilon$ on an ideal $K$, every divided power map $f:B\to C$ and every $x\in K$, there is exactly one
divided power map $F:B\langle X\rangle\to C$ with $F|_B=f$ and $F(X)=x$ \cite[Tag 07H5]{Sta}.

\emph{The object $P(T)$.} Let $P(T)=(B\langle X\rangle,i_P,J\langle X\rangle,\eta,y_P)$, where $i_P$ is $i$ followed by
$j:B\to B\langle X\rangle$, and $y_P$ is the map $\operatorname{Spec}B\langle X\rangle/J\langle X\rangle\to
\operatorname{Spec}B/J$ induced by $j$ followed by $y$. As $p\neq0$ is nilpotent in $B$, it is nilpotent in $B\langle
X\rangle$, and $i_P$ is a divided power map as a composite of two; so $P(T)$ is a thickening, and $j$ is a map $j_T:T\to
P(T)$. It lies in $\mathcal Q$: it lies over every open that $T$ lies over, by
Lemma~\ref{gal-realization-thickening-over}(1), and if $B$ is flat over $R$, so is $B\langle X\rangle$, which is free over
$B$ \cite[Tag 00HC]{Sta}. And $X_T=X\in J\langle X\rangle$.

\emph{The universal property.} Let $\varphi:T\to T'$ be a map of $\mathcal Q$ and $x\in J_{T'}$. As $\varphi$ is a divided
power map, there is exactly one divided power map $F:B\langle X\rangle\to B_{T'}$ with $F\circ j=\varphi$ and $F(X)=x$.
It commutes with the maps from $R$, as $\varphi$ does, and carries $J\langle X\rangle$ into $J_{T'}$. The point of $T'$ is
the point of $T$ composed with the map induced by $\varphi=F\circ j$, so it is the point of $P(T)$ composed with the map
induced by $F$. So $F$ is a map $\psi:P(T)\to T'$ of $\mathcal Q$, which is full, with $\psi j_T=\varphi$ and
$\psi(X_T)=x$. Every such $\psi$ is a divided power map with these two properties, so it is $F$.

\emph{Spanning.} Every element of $B\langle X\rangle$ is a finite sum $\sum_kc_kX^{[k]}=\sum_kj(c_k)\eta_k(X)$ with
$c_k\in B$.

\emph{The filtration.} Write $J^{[s]}=B$ for $s\leq0$. For $r\geq0$ let $M_r$ be the set of the $\sum_kc_kX^{[k]}$ with
$c_k\in J^{[r-k]}$ for every $k$. We show $J\langle X\rangle^{[r]}=M_r$, the left side formed for $\eta$. The $M_r$
decrease in $r$, $M_0=B\langle X\rangle$, and $M_aM_b\subset M_{a+b}$, as $cX^{[k]}\cdot c'X^{[l]}=\binom{k+l}{k}cc'
X^{[k+l]}$ and $J^{[a-k]}J^{[b-l]}\subset J^{[a+b-k-l]}$; so each $M_r$ is an ideal. The set of $z\in J\langle X\rangle$
with $\eta_n(z)\in M_n$ for every $n$ contains $J$, as $\eta_n(z)=\delta_n(z)\in J^{[n]}$ for $z\in J$, and contains every
$X^{[k+1]}$, as
\[
 \eta_n(X^{[k+1]})=\eta_n(\eta_{k+1}(X))=\frac{(n(k+1))!}{n!\,((k+1)!)^n}\,X^{[n(k+1)]}\in M_{n(k+1)}\subset M_n.
\]
It is closed under sums and under multiplication by $B\langle X\rangle$, by $\eta_n(z+z')=\sum_{a+b=n}\eta_a(z)\eta_b(z')$,
$\eta_n(cz)=c^n\eta_n(z)$ and the properties of the $M_r$. So it is $J\langle X\rangle$, and each generator
$\prod_j\eta_{k_j}(z_j)$ of $J\langle X\rangle^{[r]}$, with $\sum_jk_j\geq r$, lies in $M_{\sum_jk_j}\subset M_r$.
Conversely, for $c\in J^{[r-k]}$, $j(c)\in J\langle X\rangle^{[r-k]}$ as $j$ is a divided power map (or as $r-k\leq0$),
and $X^{[k]}=\eta_k(X)\in J\langle X\rangle^{[k]}$, so $cX^{[k]}\in J\langle X\rangle^{[r]}$.

\emph{The sequences.} Let $T\in\mathcal Q$ and $(c_k)_{k\geq0}$ finitely supported in $B$, and $z=\sum_kc_kX^{[k]}$. As
the $X^{[k]}$ form a basis, $z=0$ if and only if every $c_k=0$, which is the claim for $(0)_r$. For $r\leq0$ both
$J\langle X\rangle^{[r]}$ and every $J^{[r-k]}$ are the whole ring; for $r\geq0$, $z\in J\langle X\rangle^{[r]}$ if and
only if $c_k\in J^{[r-k]}$ for every $k$, by the filtration. For $c\in R$, the ideal $cB\langle X\rangle$, with $cB$
meaning $i(c)B$, consists of the $\sum_kd_kX^{[k]}$ with every $d_k\in cB$; adding it to $J\langle X\rangle^{[r]}$, also
described coefficientwise, gives the $\sum_kd_kX^{[k]}$ with every $d_k\in J^{[r-k]}+cB$. These are the claims for
$(cB)_r$ and $(J^{[r]}+cB)_r$.
% Lean: Pb.Site.poly_all' with polyAll. P(T) is polyT (ofMap T (algebraMap B (pq B).P) ((pq B).ext J) (eta δ)) for the
% model pq B = PDA.Model.pdPoly B, a PDPoly (basis e with e 0 = 1 and the binomial product); eta is chosen from
% FubiniStudyChowHeightFormula₂.Adr.exists_std, the existence part of Stacks 07H5. Membership in Q is polyT_mem with
% flat_comp_envelope_algebraMap (SitePD.pdPoly_flat_comp). The universal property is divided_power_envelope_universal
% (PDUniv.pdPoly_universal, the Lean's own proof: F is the B-linear map e n ↦ ε_n(x), multiplicative by mul_dpow, a
% divided power map by IsDPMorphism.on_span) with polyDesc = descOf and hom_ext; spanning is
% divided_power_envelope_expansion (PDSpanning.exists_sum_dpow); the filtration is
% divided_power_polynomial_filtration (PDFil.fil_pdPoly: coefFil is M_r, dpow_mem_coefFil and coefFil_le_fil the two
% inclusions); the sequences are divided_power_envelope_filtration (PDCrit.coeff_criteria, with repr_sum,
% mem_span_iff and mem_fil_sup_iff).
\end{proof}
\leanref{Development/AutomorphicSpectrum/Proofs/BlockDecomposition/StdSomeZero.lean}{GlobalTheory.CompanionTheory.ChartLocBc.ladder}
\leanref{Development/WeilBound/CoreStatements/PadicHodgeComparison/EtaleCohomologyComparison/PerfectoidPullbacks/PolynomialSiteTraces/AdmissibleMap.lean}{CrystallineSemistableComparison.Pb.Site.polyT}
\leanref{Development/WeilBound/CoreStatements/PadicHodgeComparison/EtaleCohomologyComparison/PerfectoidPullbacks/PolynomialSiteTraces/PolyAll.lean}{CrystallineSemistableComparison.Pb.Site.polyT_mem}
\leanref{Development/WeilBound/CoreStatements/PadicHodgeComparison/EtaleCohomologyComparison/PerfectoidPullbacks/PolynomialSiteTraces/PolyAll.lean}{CrystallineSemistableComparison.Pb.Site.polyDesc}
\leanref{Development/WeilBound/CoreStatements/AbelianHeightsDuality/DividedPowerAlgebras/ExistsDividedPowersAug.lean}{FubiniStudyChowHeightFormula₂.Adr.exists_std}
\leanref{Development/AutomorphicSpectrum/Proofs/BlockDecomposition/DividedPowerEnvelopeUniversalStatement.lean}{GlobalTheory.T.divided_power_envelope_universal}
\leanref{Development/AutomorphicSpectrum/Proofs/BlockDecomposition/DividedPowerEnvelopeExpansionStatement.lean}{GlobalTheory.T.divided_power_envelope_expansion}
\leanref{Development/AutomorphicSpectrum/Proofs/BlockDecomposition/DividedPowerEnvelopeFiltrationStatement.lean}{GlobalTheory.T.divided_power_envelope_filtration}
\leanref{Development/AutomorphicSpectrum/Proofs/BlockDecomposition/DividedPowerPolynomialFiltrationStatement.lean}{GlobalTheory.T.divided_power_polynomial_filtration}
\leanref{Development/AutomorphicSpectrum/Proofs/BlockDecomposition/FlatCompEnvelopeAlgebraMapStatement.lean}{GlobalTheory.T.flat_comp_envelope_algebraMap}
\leanref{Development/AutomorphicSpectrum/Proofs/BlockDecomposition/DividedPowerPolynomialFiltration.lean}{GlobalTheory.CompanionTheory.PDFil.fil_pdPoly}
\leanref{Development/AutomorphicSpectrum/Proofs/BlockDecomposition/DividedPowerEnvelopeFiltration.lean}{GlobalTheory.CompanionTheory.PDCrit.coeff_criteria}

\begin{definition}\label{gal-realization-thickening-torsor-cut}
Let $\Phi$, given by $A$, $J$ and $\gamma$, be a divided power functor on a category $\mathcal C$ over a ring $R$, $d\geq0$,
and $(D,\tau)$ a torsor datum of dimension $d+1$ for $\Phi$ (Definition~\ref{gal-realization-thickening-torsor}). Maps
are composed as functions.
\begin{enumerate}
\item The \emph{pointed category} $\mathcal C_D$ has as objects the pairs $(T,s)$ of an object $T$ of $\mathcal C$ and a
  map $s:D\to T$, and as maps $(T,s)\to(T',s')$ the maps $T\to T'$ of $\mathcal C$. Let $U:\mathcal C_D\to\mathcal C$ be
  the functor $(T,s)\mapsto T$. The \emph{cut category} $\mathcal C^0_D$ has the same objects, and its maps
  $(T,s)\to(T',s')$ are the maps $\varphi:T\to T'$ with $\tau_1(\varphi s,s')=0$. Let $\iota:\mathcal C^0_D\to\mathcal
  C_D$ be the inclusion and $U^0=U\iota$.
\item Let $\Phi^0$ be the divided power functor on $\mathcal C^0_D$ given by $A\circ U^0$, $J\circ U^0$ and $\gamma$; a
  family of ideals $\mathfrak a$ of $\Phi$ gives the family $\mathfrak a^0=\mathfrak a\circ U^0$ of $\Phi^0$. For
  $1\leq k\leq d$ and maps $\sigma,\sigma':(D,\mathrm{id}_D)\to(T,s)$ of $\mathcal C^0_D$, let
  $\tau^0_k(\sigma,\sigma')=\tau_{k+1}(\sigma,\sigma')$.
\item Let $P$ be a divided power polynomial structure on $\Phi$, with maps $j_T:T\to P(T)$ and elements $X_T$. Let $P^0$
  assign to $(T,s)$ the object $(P(T),j_Ts)$, the map $j_T$ and the element $X_T$. For $(T,s)$ in $\mathcal C_D$, let
  $s^+:D\to P(T)$ be the map with $\tau_1(j_Ts,s^+)=X_T$ and $\tau_k(j_Ts,s^+)=0$ for $2\leq k\leq d+1$, and
  $P_D(T,s)=(P(T),s^+)$. For a map $\varphi:(T,s)\to(T',s')$ of $\mathcal C_D$, let $P_D(\varphi):P(T)\to P(T')$ be the
  map $\psi$ with $\psi j_T=j_{T'}\varphi$ and $A(\psi)(X_T)=X_{T'}+A(j_{T'})(\tau_1(\varphi s,s'))$.
\end{enumerate}
The maps $s^+$ and $P_D(\varphi)$ exist and are unique by Definition~\ref{gal-realization-thickening-torsor}(2) and (3),
as $X_T$ lies in $J(P(T))$ and $X_{T'}+A(j_{T'})(\tau_1(\varphi s,s'))$ in $J(P(T'))$.
\end{definition}
\leanref{Development/WeilBound/CoreStatements/PadicHodgeComparison/EtaleCohomologyComparison/PerfectoidPullbacks/LocalisationLadders/TruncationHomExt.lean}{CrystallineSemistableComparison.Pb.Pt}
\leanref{Development/WeilBound/CoreStatements/PadicHodgeComparison/EtaleCohomologyComparison/PerfectoidPullbacks/LocalisationLadders/TruncationHomExt.lean}{CrystallineSemistableComparison.Pb.Pt.U}
\leanref{Development/WeilBound/CoreStatements/PadicHodgeComparison/EtaleCohomologyComparison/PerfectoidPullbacks/LocalisationLadders/TruncationHomExt.lean}{CrystallineSemistableComparison.Pb.Cut}
\leanref{Development/WeilBound/CoreStatements/PadicHodgeComparison/EtaleCohomologyComparison/PerfectoidPullbacks/LocalisationLadders/TruncationHomExt.lean}{CrystallineSemistableComparison.Pb.Cut.U}
\leanref{Development/WeilBound/CoreStatements/PadicHodgeComparison/EtaleCohomologyComparison/PerfectoidPullbacks/LocalisationLadders/TruncationHomExt.lean}{CrystallineSemistableComparison.Pb.Cut.inc}
\leanref{Development/WeilBound/CoreStatements/PadicHodgeComparison/EtaleCohomologyComparison/PerfectoidPullbacks/LocalisationLadders/CriterionTruncation.lean}{CrystallineSemistableComparison.Pb.Cut.Φ1}
\leanref{Development/WeilBound/CoreStatements/PadicHodgeComparison/EtaleCohomologyComparison/PerfectoidPullbacks/LocalisationLadders/CriterionTruncation.lean}{CrystallineSemistableComparison.Pb.Cut.E1}
\leanref{Development/WeilBound/CoreStatements/PadicHodgeComparison/EtaleCohomologyComparison/PerfectoidPullbacks/LocalisationLadders/CriterionTruncation.lean}{CrystallineSemistableComparison.Pb.Cut.Q1}
\leanref{Development/WeilBound/CoreStatements/PadicHodgeComparison/EtaleCohomologyComparison/PerfectoidPullbacks/LocalisationLadders/TruncationHomExt.lean}{CrystallineSemistableComparison.Pb.Cut.sec}
\leanref{Development/WeilBound/CoreStatements/PadicHodgeComparison/EtaleCohomologyComparison/PerfectoidPullbacks/LocalisationLadders/TruncationHomExt.lean}{CrystallineSemistableComparison.Pb.Cut.Pmap}
\leanref{Development/WeilBound/CoreStatements/PadicHodgeComparison/EtaleCohomologyComparison/PerfectoidPullbacks/LocalisationLadders/TruncationHomExt.lean}{CrystallineSemistableComparison.Pb.Cut.P0}

\begin{lemma}\label{gal-realization-thickening-torsor-cut-adjoint}
Let $\Phi$, $(D,\tau)$ and $P$ be as in Definition~\ref{gal-realization-thickening-torsor-cut}, with $\mathcal C_D$,
$\mathcal C^0_D$, $U$, $\iota$, $U^0$, $\Phi^0$, $\tau^0$, $P^0$ and $P_D$ as there.
\begin{enumerate}
\item $\mathcal C^0_D$ is a category, $(D,\mathrm{id}_D)$ with $\tau^0$ is a torsor datum of dimension $d$ for $\Phi^0$,
  and $P^0$ is a divided power polynomial structure on $\Phi^0$. If $P$ respects a sequence $(\mathfrak a_r)_r$ of
  families of ideals of $\Phi$ in degree $r$, then $P^0$ respects $(\mathfrak a^0_r)_r$ in degree $r$; if $P$ is a ladder
  up to $r$ for $(\mathfrak a_r)_r$ and $(\mathfrak b_r)_r$, then $P^0$ is one for $(\mathfrak a^0_r)_r$ and
  $(\mathfrak b^0_r)_r$.
\item For every family of maps $\sigma_T:D\to T$, $T\in\mathcal C$, the functor $S_\sigma:\mathcal C\to\mathcal C_D$,
  $T\mapsto(T,\sigma_T)$, $\varphi\mapsto\varphi$, is left adjoint and right adjoint to $U$.
\item $P_D$ is a functor $\mathcal C_D\to\mathcal C^0_D$, and it is left adjoint to $\iota$: for $(T,s)$ in $\mathcal C_D$
  and $(T'',s'')$ in $\mathcal C^0_D$, $\psi\mapsto\psi j_T$ is a bijection from the maps $P_D(T,s)\to(T'',s'')$ of
  $\mathcal C^0_D$ to the maps $T\to T''$ of $\mathcal C$. A map $\psi:P(T)\to T''$ of $\mathcal C$ is a map
  $P_D(T,s)\to(T'',s'')$ of $\mathcal C^0_D$ if and only if $A(\psi)(X_T)=\tau_1(\psi j_Ts,s'')$. The maps $j_T$ form a
  natural transformation $j:U\to U^0P_D$.
\end{enumerate}
\end{lemma}
\leanref{Development/WeilBound/CoreStatements/PadicHodgeComparison/EtaleCohomologyComparison/PerfectoidPullbacks/LocalisationLadders/CriterionTruncation.lean}{CrystallineSemistableComparison.Pb.Cut.criterion_truncation}
\leanref{Development/WeilBound/CoreStatements/PadicHodgeComparison/EtaleCohomologyComparison/PerfectoidPullbacks/LocalisationLadders/CriterionTruncation.lean}{CrystallineSemistableComparison.Pb.Cut.ladder_truncation}
\leanref{Development/WeilBound/CoreStatements/PadicHodgeComparison/EtaleCohomologyComparison/PerfectoidPullbacks/LocalisationLadders/TruncationHomExt.lean}{CrystallineSemistableComparison.Pb.Pt.adj}
\leanref{Development/WeilBound/CoreStatements/PadicHodgeComparison/EtaleCohomologyComparison/PerfectoidPullbacks/LocalisationLadders/CriterionTruncation.lean}{CrystallineSemistableComparison.Pb.Pt.adjσ}
\leanref{Development/WeilBound/CoreStatements/PadicHodgeComparison/EtaleCohomologyComparison/PerfectoidPullbacks/LocalisationLadders/TruncationHomExt.lean}{CrystallineSemistableComparison.Pb.Cut.adj}
\leanref{Development/WeilBound/CoreStatements/PadicHodgeComparison/EtaleCohomologyComparison/PerfectoidPullbacks/LocalisationLadders/TruncationHomExt.lean}{CrystallineSemistableComparison.Pb.Cut.jNat}

\begin{proof}
Taking $\sigma=\sigma'=\sigma''$, and then $\sigma''=\sigma$, in the identity
$\tau_k(\sigma,\sigma'')=\tau_k(\sigma,\sigma')+\tau_k(\sigma',\sigma'')$ gives $\tau_k(\sigma,\sigma)=0$ and
$\tau_k(\sigma',\sigma)=-\tau_k(\sigma,\sigma')$. For maps $s:D\to T$, $\varphi:T\to T'$, $s':D\to T'$, $\psi:T'\to T''$
and $s'':D\to T''$,
\[
 \tau_1(\psi\varphi s,s'')=\tau_1(\psi\varphi s,\psi s')+\tau_1(\psi s',s'')=A(\psi)(\tau_1(\varphi
 s,s'))+\tau_1(\psi s',s'').
\]

(1) The identity of $(T,s)$ is a map of $\mathcal C^0_D$, as $\tau_1(s,s)=0$, and by the displayed identity a composite
of maps of $\mathcal C^0_D$ is one. The conditions on $\tau^0$ other than the last are those on
$\tau_2,\ldots,\tau_{d+1}$, and $(D,\mathrm{id}_D)$ has the map $s$ to every $(T,s)$. Let $\sigma:(D,\mathrm{id}_D)\to(T,s)$
be a map of $\mathcal C^0_D$, so that $\tau_1(\sigma,s)=0$, and $y_1,\ldots,y_d\in J(T)$. There is exactly one
$\sigma':D\to T$ with $\tau_1(\sigma,\sigma')=0$ and $\tau_{k+1}(\sigma,\sigma')=y_k$ for $1\leq k\leq d$. It is a map of
$\mathcal C^0_D$, as $\tau_1(\sigma',s)=\tau_1(\sigma',\sigma)+\tau_1(\sigma,s)=0$. A map $\sigma'$ of $\mathcal C^0_D$
with $\tau^0_k(\sigma,\sigma')=y_k$ for every $k$ has $\tau_1(\sigma,\sigma')=\tau_1(\sigma,s)-\tau_1(\sigma',s)=0$, so
it is this one.

The map $j_T$ is a map $(T,s)\to(P(T),j_Ts)$ of $\mathcal C^0_D$, as $\tau_1(j_Ts,j_Ts)=0$. Let $\varphi:(T,s)\to(T',s')$
be a map of $\mathcal C^0_D$ and $x\in J(T')$. The map $\psi:P(T)\to T'$ with $\psi j_T=\varphi$ and $A(\psi)(X_T)=x$ is a
map $(P(T),j_Ts)\to(T',s')$ of $\mathcal C^0_D$, as $\tau_1(\psi j_Ts,s')=\tau_1(\varphi s,s')=0$, and a map of $\mathcal
C^0_D$ with these two properties is a map of $\mathcal C$ with them, so it is $\psi$. The spanning property, the condition
for respecting a sequence and those of a ladder concern $A(P(T))$ and $A(T)$ for one object $T$ at a time, and they are
the same for $P^0$ at $(T,s)$.

(2) A map $S_\sigma(T)\to(T',s')$ of $\mathcal C_D$ is a map $T\to U(T',s')$ of $\mathcal C$, and a map $(T,s)\to
S_\sigma(T')$ is a map $U(T,s)\to T'$. These identifications are natural, as maps are composed in $\mathcal C$ on both
sides.

(3) Let $\psi:P(T)\to T''$ be a map and $s'':D\to T''$. As $\tau_1(\psi s^+,\psi j_Ts)=A(\psi)(\tau_1(s^+,j_Ts))=
-A(\psi)(X_T)$,
\[
 \tau_1(\psi s^+,s'')=\tau_1(\psi s^+,\psi j_Ts)+\tau_1(\psi j_Ts,s'')=-A(\psi)(X_T)+\tau_1(\psi j_Ts,s'').
\]
So $\psi$ is a map $P_D(T,s)\to(T'',s'')$ of $\mathcal C^0_D$ if and only if $A(\psi)(X_T)=\tau_1(\psi j_Ts,s'')$. For a
map $\varphi:T\to T''$ there is exactly one $\psi:P(T)\to T''$ with $\psi j_T=\varphi$ and $A(\psi)(X_T)=\tau_1(\varphi
s,s'')$, an element of $J(T'')$. So $\psi\mapsto\psi j_T$ is bijective.

Let $\varphi:(T,s)\to(T',s')$ be a map of $\mathcal C_D$ and $\psi=P_D(\varphi)$. Then $P_D(\varphi)$ is a map of
$\mathcal C^0_D$, as
\[
 \tau_1(\psi j_Ts,s'^+)=\tau_1(j_{T'}\varphi s,j_{T'}s')+\tau_1(j_{T'}s',s'^+)=A(j_{T'})(\tau_1(\varphi
 s,s'))+X_{T'}=A(\psi)(X_T).
\]
The identity of $P(T)$ has the two properties that define $P_D(\mathrm{id})$, as $\tau_1(s,s)=0$. For a second map
$\varphi':(T',s')\to(T'',s'')$ of $\mathcal C_D$, $P_D(\varphi')P_D(\varphi)j_T=j_{T''}\varphi'\varphi$, and by the
displayed identity at the start of the proof
\[
 A(P_D(\varphi'))\bigl(X_{T'}+A(j_{T'})(\tau_1(\varphi s,s'))\bigr)=X_{T''}+A(j_{T''})\bigl(\tau_1(\varphi's',s'')+
 A(\varphi')(\tau_1(\varphi s,s'))\bigr)=X_{T''}+A(j_{T''})(\tau_1(\varphi'\varphi s,s'')).
\]
So $P_D(\varphi'\varphi)=P_D(\varphi')P_D(\varphi)$, $P_D$ is a functor, and $P_D(\varphi)j_T=j_{T'}\varphi$ says that
$j$ is natural. The bijection $\psi\mapsto\psi j_T$ is natural in $(T,s)$, as $P_D(\varphi)j_T=j_{T'}\varphi$, and in
$(T'',s'')$, as maps are composed in $\mathcal C$ on both sides; so $P_D$ is left adjoint to $\iota$.
\end{proof}
\leanref{Development/WeilBound/CoreStatements/PadicHodgeComparison/EtaleCohomologyComparison/PerfectoidPullbacks/LocalisationLadders/TruncationHomExt.lean}{CrystallineSemistableComparison.Pb.Frame.τ_self}
\leanref{Development/WeilBound/CoreStatements/PadicHodgeComparison/EtaleCohomologyComparison/PerfectoidPullbacks/LocalisationLadders/TruncationHomExt.lean}{CrystallineSemistableComparison.Pb.Frame.τ_swap}
\leanref{Development/WeilBound/CoreStatements/PadicHodgeComparison/EtaleCohomologyComparison/PerfectoidPullbacks/LocalisationLadders/TruncationHomExt.lean}{CrystallineSemistableComparison.Pb.Cut.τ_sec}
\leanref{Development/WeilBound/CoreStatements/PadicHodgeComparison/EtaleCohomologyComparison/PerfectoidPullbacks/LocalisationLadders/TruncationHomExt.lean}{CrystallineSemistableComparison.Pb.Cut.Pmap_cut}
\leanref{Development/WeilBound/CoreStatements/PadicHodgeComparison/EtaleCohomologyComparison/PerfectoidPullbacks/LocalisationLadders/TruncationHomExt.lean}{CrystallineSemistableComparison.Pb.Cut.val_comp}
\leanref{Development/WeilBound/CoreStatements/PadicHodgeComparison/EtaleCohomologyComparison/PerfectoidPullbacks/LocalisationLadders/TruncationHomExt.lean}{CrystallineSemistableComparison.Pb.Cut.adjInv}
\leanref{Development/WeilBound/CoreStatements/PadicHodgeComparison/EtaleCohomologyComparison/PerfectoidPullbacks/LocalisationLadders/TruncationHomExt.lean}{CrystallineSemistableComparison.Pb.Cut.truncation_hom_ext}

\begin{lemma}\label{gal-realization-thickening-torsor-adjoint}
Let $\Lambda$ be a ring, $\mathcal C$ a category and $F$ a functor from $\mathcal C$ to $\Lambda$-modules; cohomology is
that of Definition~\ref{gal-realization-split-diagrams}(1).
\begin{enumerate}
\item Let $L:\mathcal C'\to\mathcal C$ be a functor that has a right adjoint. For every $j$, the map $H^j(\mathcal
  C,F)\to H^j(\mathcal C',F\circ L)$ induced by $L$ and the identity maps of the modules $F(L(c))$ is bijective.
\item Let $D$ be an initial object of $\mathcal C$, and for $x\in F(D)$ let $\varepsilon(x)$ be the $0$-cochain $c\mapsto
  F(D\to c)(x)$. Then $x\mapsto\varepsilon(x)$ is a bijection from $F(D)$ to $H^0(\mathcal C,F)$, and $H^j(\mathcal
  C,F)=0$ for $j\geq1$.
\end{enumerate}
\end{lemma}
\leanref{Development/WeilBound/CoreStatements/PadicHodgeComparison/EtaleCohomologyComparison/CechPairings/FaceMapCompSuccAboveCastSucc.lean}{CrystallineSemistableComparison.P.adj}
\leanref{Development/WeilBound/CoreStatements/PadicHodgeComparison/EtaleCohomologyComparison/CechPairings/FaceMapCompSuccAboveCastSucc.lean}{CrystallineSemistableComparison.P.initial}
\leanref{Development/AutomorphicSpectrum/Proofs/BlockDecomposition/InitialObjectAugmentedResolutionStatement.lean}{GlobalTheory.T.initial_object_augmented_resolution}
\leanref{Development/AutomorphicSpectrum/Proofs/BlockDecomposition/ToH0BijectiveOfInitial.lean}{GlobalTheory.CompanionTheory.LocBC.toH0_bijective_of_initial}
\leanref{Development/AutomorphicSpectrum/Proofs/BlockDecomposition/ToH0BijectiveOfInitial.lean}{GlobalTheory.CompanionTheory.LocBC.H_succ_sub}

\begin{proof}
(1) Let $U:\mathcal C\to\mathcal C'$ be a right adjoint of $L$, with unit $\eta:\mathrm{id}\to UL$ and counit
$\epsilon:LU\to\mathrm{id}$. The functor $U$ with the maps $F(\epsilon_c):F(L(U(c)))\to F(c)$ induces a map
$m:H^j(\mathcal C',F\circ L)\to H^j(\mathcal C,F)$. A composite of two induced maps is induced by the composite functor
with the composite maps. So $m$ after the map of the statement is induced by $LU$ with the maps $F(\epsilon_c)$, and
Lemma~\ref{gal-realization-period-prism}, for the natural transformation $\epsilon$ to the identity functor with the
identity maps, shows that it is the identity. The map of the statement after $m$ is induced by $UL$ with the maps
$F(\epsilon_{L(c')}):F(LUL(c'))\to F(L(c'))$; as $F(\epsilon_{L(c')})F(L(\eta_{c'}))$ is the identity,
Lemma~\ref{gal-realization-period-prism}, for $\eta$ from the identity functor with the identity maps, shows that it is
the identity too.

(2) For a cochain $y$ of degree $j+1$, let $hy$ be the cochain of degree $j$ with $hy(c_0\to\cdots\to c_j)=y(D\to
c_0\to\cdots\to c_j)$. The face $0$ of $D\to c_0\to\cdots\to c_{j+1}$ is $c_0\to\cdots\to c_{j+1}$, and its face $i+1$ is
$D$ followed by the face $i$ of $c_0\to\cdots\to c_{j+1}$. So the value of $dy$ on $D\to c_0\to\cdots\to c_{j+1}$ is
$y(c_0\to\cdots\to c_{j+1})-d(hy)(c_0\to\cdots\to c_{j+1})$, and $y=d(hy)$ if $dy=0$. Hence $H^j(\mathcal C,F)=0$ for
$j\geq1$. For $x\in F(D)$, $\varepsilon(x)$ is a cocycle, as $F(c\to c')F(D\to c)=F(D\to c')$, and
$\varepsilon(x)(D)=x$, as the only map $D\to D$ is the identity. A $0$-cocycle $f$ has $f(c)=F(D\to c)(f(D))$, as $df$
vanishes on $D\to c$; so $f=\varepsilon(f(D))$. As $H^0(\mathcal C,F)$ is the module of $0$-cocycles, $\varepsilon$ is
bijective.
\end{proof}
\leanref{Development/AutomorphicSpectrum/Proofs/BlockDecomposition/InitialObjectAugmentedResolution.lean}{GlobalTheory.CompanionTheory.InitialCoh.con}
\leanref{Development/AutomorphicSpectrum/Proofs/BlockDecomposition/InitialObjectAugmentedResolution.lean}{GlobalTheory.CompanionTheory.InitialCoh.d_con}

\begin{lemma}\label{gal-realization-thickening-torsor-exact}
Let $\Phi$, given by $A$, $J$ and $\gamma$, $(D,\tau)$ and $P$, given by $j_T$ and $X_T$, be as in
Definition~\ref{gal-realization-thickening-torsor-cut}, with $\mathcal C_D$, $U$, $U^0$ and $P_D$ as there. Let
$(\mathfrak a_{r'})_{r'\in\mathbf Z}$ and $(\mathfrak b_{r'})_{r'\in\mathbf Z}$ be sequences of families of ideals of
$\Phi$ with $\mathfrak b_{r'}(T)\subset\mathfrak a_{r'}(T)$ for all $r'$ and $T$, and $r\in\mathbf Z$, such that $P$
respects both sequences in degrees $r$ and $r-1$. For $r'\in\{r,r-1\}$ let $F_{r'}=(\mathfrak a_{r'}/\mathfrak
b_{r'})\circ U^0P_D$, the functor on $\mathcal C_D$ with $F_{r'}(T,s)=(\mathfrak a_{r'}/\mathfrak b_{r'})(P(T))$. For a
finitely supported sequence $c=(c_k)_{k\geq0}$ in $A(T)$ write $\langle c\rangle=\sum_kA(j_T)(c_k)\gamma_k(X_T)$ and
$c^+=(c_{k+1})_{k\geq0}$.
\begin{enumerate}
\item For $(T,s)$ in $\mathcal C_D$ there is exactly one map $\beta:F_r(T,s)\to F_{r-1}(T,s)$ that sends the class of
  $\langle c\rangle$ to the class of $\langle c^+\rangle$ for every $c$ with $\langle c\rangle\in\mathfrak a_r(P(T))$. It
  is $R$-linear and natural on $\mathcal C_D$.
\item With the maps $(\mathfrak a_r/\mathfrak b_r)(j_T)$, the sequences
  \[
   0\to(\mathfrak a_r/\mathfrak b_r)(T)\to F_r(T,s)\xrightarrow{\beta}F_{r-1}(T,s)\to0
  \]
  are exact and natural on $\mathcal C_D$.
\end{enumerate}
\end{lemma}
\leanref{Development/AutomorphicSpectrum/Proofs/BlockDecomposition/FiltrationCriticalComparisonMapStatement.lean}{GlobalTheory.T.filtration_critical_comparison_map}
\leanref{Development/AutomorphicSpectrum/Proofs/BlockDecomposition/FiltrationCriticalComparisonMap.lean}{GlobalTheory.CompanionTheory.PolySes.ses}
\leanref{Development/AutomorphicSpectrum/Proofs/BlockDecomposition/FiltrationCriticalComparisonMap.lean}{GlobalTheory.CompanionTheory.PolySes.dX}

\begin{proof}
By the spanning property every element of $A(P(T))$ is some $\langle c\rangle$, and $c\mapsto\langle c\rangle$ is
additive and $R$-linear.

(1) If $\langle c\rangle\in\mathfrak a_r(P(T))$, then $c_k\in\mathfrak a_{r-k}(T)$ for every $k$, as $P$ respects
$(\mathfrak a_{r'})$ in degree $r$; so $c_{k+1}\in\mathfrak a_{r-1-k}(T)$, and $\langle c^+\rangle\in\mathfrak
a_{r-1}(P(T))$, as $P$ respects $(\mathfrak a_{r'})$ in degree $r-1$. If $\langle c\rangle$ and $\langle c'\rangle$ have
the same class modulo $\mathfrak b_r(P(T))$, the same argument for $(\mathfrak b_{r'})$ and $c-c'$ shows that
$\langle c^+\rangle-\langle c'^+\rangle\in\mathfrak b_{r-1}(P(T))$. So $\beta$ is well defined, and it is $R$-linear.

Let $\varphi:(T,s)\to(T',s')$ be a map of $\mathcal C_D$, $\psi=P_D(\varphi)$ and $t=\tau_1(\varphi s,s')\in J(T')$.
Then $\psi j_T=j_{T'}\varphi$, and $A(\psi)(X_T)=X_{T'}+A(j_{T'})(t)$ is a sum of two elements of $J(P(T'))$. By the
axioms of Definition~\ref{gal-realization-thickening-torsor}(1),
\[
 A(\psi)(\langle c\rangle)=\sum_kA(j_{T'})(A(\varphi)(c_k))\sum_{a+b=k}\gamma_a(X_{T'})A(j_{T'})(\gamma_b(t))=\langle
 c^\varphi\rangle,\qquad c^\varphi_a=\sum_{b\geq0}A(\varphi)(c_{a+b})\gamma_b(t).
\]
As $(c^+)^\varphi_a=\sum_bA(\varphi)(c_{a+1+b})\gamma_b(t)=c^\varphi_{a+1}$, $(c^\varphi)^+=(c^+)^\varphi$, so $\beta$
commutes with $F_r(\varphi)$ and $F_{r-1}(\varphi)$.

(2) The maps $(\mathfrak a_r/\mathfrak b_r)(j_T)$ are natural on $\mathcal C_D$, as $j$ is
(Lemma~\ref{gal-realization-thickening-torsor-cut-adjoint}(3)). For $w\in A(T)$, $A(j_T)(w)=\langle(w,0,0,\ldots)\rangle$,
as $\gamma_0(X_T)=1$. If $w\in\mathfrak a_r(T)$ and $A(j_T)(w)\in\mathfrak b_r(P(T))$, then $w\in\mathfrak b_r(T)$, as $P$
respects $(\mathfrak b_{r'})$ in degree $r$; so the first map is injective. The composite is $0$, as
$(w,0,0,\ldots)^+=0$. Let $\langle c\rangle\in\mathfrak a_r(P(T))$ with $\langle c^+\rangle\in\mathfrak b_{r-1}(P(T))$.
Then $c_{k+1}\in\mathfrak b_{r-1-k}(T)$ for every $k$, so $\langle c-(c_0,0,0,\ldots)\rangle\in\mathfrak b_r(P(T))$, and
$c_0\in\mathfrak a_r(T)$; so the class of $\langle c\rangle$ is the image of the class of $c_0$. Finally, for $\langle
c'\rangle\in\mathfrak a_{r-1}(P(T))$, let $c_0=0$ and $c_{k+1}=c'_k$. Then $c_{k+1}\in\mathfrak a_{r-1-k}(T)$ for every
$k$, so $\langle c\rangle\in\mathfrak a_r(P(T))$, and $c^+=c'$; so $\beta$ is surjective.
\end{proof}
\leanref{Development/AutomorphicSpectrum/Proofs/BlockDecomposition/FiltrationCriticalComparisonMap.lean}{GlobalTheory.CompanionTheory.PolySes.wd}
\leanref{Development/AutomorphicSpectrum/Proofs/BlockDecomposition/FiltrationCriticalComparisonMap.lean}{GlobalTheory.CompanionTheory.PolySes.ev_sh_mem}
\leanref{Development/AutomorphicSpectrum/Proofs/BlockDecomposition/FiltrationCriticalComparisonMap.lean}{GlobalTheory.CompanionTheory.PolySes.ev_tr}
\leanref{Development/AutomorphicSpectrum/Proofs/BlockDecomposition/FiltrationCriticalComparisonMap.lean}{GlobalTheory.CompanionTheory.PolySes.sh_tr}
\leanref{Development/AutomorphicSpectrum/Proofs/BlockDecomposition/FiltrationCriticalComparisonMap.lean}{GlobalTheory.CompanionTheory.PolySes.dX_nat}
\leanref{Development/AutomorphicSpectrum/Proofs/BlockDecomposition/FiltrationCriticalComparisonMap.lean}{GlobalTheory.CompanionTheory.PolySes.sqMap_j_inj}
\leanref{Development/AutomorphicSpectrum/Proofs/BlockDecomposition/FiltrationCriticalComparisonMap.lean}{GlobalTheory.CompanionTheory.PolySes.dX_exact}
\leanref{Development/AutomorphicSpectrum/Proofs/BlockDecomposition/FiltrationCriticalComparisonMap.lean}{GlobalTheory.CompanionTheory.PolySes.dX_surj}

\begin{lemma}\label{gal-realization-thickening-torsor-cut-morphism}
Let $\sigma:R\to R'$ be a ring map, $\Phi$ and $\Phi'$, given by $A,J,\gamma$ and $A',J',\gamma'$, divided power
functors on categories $\mathcal C$ over $R$ and $\mathcal C'$ over $R'$, and $(G,\rho)$ a morphism $\Phi\to\Phi'$ over
$\sigma$ with every $\rho_T$ injective. Let $P$ and $P'$ be divided power polynomial structures on $\Phi$ and $\Phi'$,
given by $j_T,X_T$ and $j'_T,X'_T$, and $(D,\tau)$ and $(D',\tau')$ torsor data of dimension $d+1$ for $\Phi$ and
$\Phi'$, compatible along a map $e:D\to G(D')$ (Definition~\ref{gal-realization-thickening-torsor}). Form
$\mathcal C_D$, $\mathcal C^0_D$ and $P_D$, and $\mathcal C'_{D'}$, $\mathcal C'^0_{D'}$ and $P'_{D'}$
(Definition~\ref{gal-realization-thickening-torsor-cut}), and write $P'_{D'}(T,s)=(P'(T),s^+)$.
\begin{enumerate}
\item The functor $G_D:\mathcal C'_{D'}\to\mathcal C_D$, $(T,s)\mapsto(G(T),G(s)e)$, $\varphi\mapsto G(\varphi)$, maps
  $\mathcal C'^0_{D'}$ into $\mathcal C^0_D$; let $G^0:\mathcal C'^0_{D'}\to\mathcal C^0_D$ be the restricted functor.
\item For $(T,s)$ in $\mathcal C'_{D'}$ let $\eta_{(T,s)}:P_DG_D(T,s)\to G^0P'_{D'}(T,s)$ be the map of $\mathcal C^0_D$
  with $\eta_{(T,s)}j_{G(T)}=G(j'_T)$ (Lemma~\ref{gal-realization-thickening-torsor-cut-adjoint}(3)). Then $\eta$ is a
  natural transformation $P_DG_D\to G^0P'_{D'}$, $A(\eta_{(T,s)})(X_{G(T)})=\tau_1(G(j'_Ts)e,G(s^+)e)$, and
  $\rho_{P'(T)}(A(\eta_{(T,s)})(X_{G(T)}))=X'_T$.
\item Let $a\in R'$, let $j:T\to T''$ be a map of $\mathcal C'$, $\eta':P(G(T))\to G(T'')$ a map of $\mathcal C$ with
  $\eta'j_{G(T)}=G(j)$, and $X''\in A'(T'')$ with $\rho_{T''}(A(\eta')(X_{G(T)}))=X''$. Then for every finitely
  supported sequence $(c_k)_{k\geq0}$ in $A(G(T))$,
  \[
   a\,\rho_{T''}\Bigl(A(\eta')\Bigl(\sum_kA(j_{G(T)})(c_k)\gamma_k(X_{G(T)})\Bigr)\Bigr)
   =\sum_kA'(j)\bigl(a\rho_T(c_k)\bigr)\gamma'_k(X'').
  \]
  By (2), this applies to $j=j'_T$, $\eta'=\eta_{(T,s)}$ and $X''=X'_T$.
\end{enumerate}
\end{lemma}
\leanref{Development/AutomorphicSpectrum/Proofs/BlockDecomposition/EtaNat.lean}{GlobalTheory.CompanionTheory.LocBC.Gp}
\leanref{Development/AutomorphicSpectrum/Proofs/BlockDecomposition/EtaNat.lean}{GlobalTheory.CompanionTheory.LocBC.G1}
\leanref{Development/AutomorphicSpectrum/Proofs/BlockDecomposition/EtaNat.lean}{GlobalTheory.CompanionTheory.LocBC.etaC}
\leanref{Development/AutomorphicSpectrum/Proofs/BlockDecomposition/EtaNat.lean}{GlobalTheory.CompanionTheory.LocBC.eta}
\leanref{Development/AutomorphicSpectrum/Proofs/BlockDecomposition/EtaNat.lean}{GlobalTheory.CompanionTheory.LocBC.j_eta}
\leanref{Development/AutomorphicSpectrum/Proofs/BlockDecomposition/EtaNat.lean}{GlobalTheory.CompanionTheory.LocBC.rho_eta}
\leanref{Development/AutomorphicSpectrum/Proofs/BlockDecomposition/PullRestrictFirstPieceCompat.lean}{GlobalTheory.CompanionTheory.LocBC.rho_finsupp_sum_dpow}

\begin{proof}
(1) $G_D$ is a functor, as $G$ is. If $\varphi:(T,s)\to(T',s')$ is a map of $\mathcal C'^0_{D'}$, then
$\tau'_1(\varphi s,s')=0$. By compatibility, $\rho_{T'}(\tau_1(G(\varphi s)e,G(s')e))=\tau'_1(\varphi s,s')=0$, and as
$\rho_{T'}$ is injective, $\tau_1(G(\varphi)G(s)e,G(s')e)=\tau_1(G(\varphi s)e,G(s')e)=0$. So $G(\varphi)$ is a map
$G_D(T,s)\to G_D(T',s')$ of $\mathcal C^0_D$.

(2) By Lemma~\ref{gal-realization-thickening-torsor-cut-adjoint}(3), $A(\eta_{(T,s)})(X_{G(T)})=\tau_1(\eta_{(T,s)}
j_{G(T)}G(s)e,G(s^+)e)=\tau_1(G(j'_Ts)e,G(s^+)e)$, and by compatibility its image under $\rho_{P'(T)}$ is
$\tau'_1(j'_Ts,s^+)=X'_T$. Let $\varphi:(T,s)\to(T',s')$ be a map of $\mathcal C'_{D'}$. Both
$\eta_{(T',s')}P_D(G(\varphi))$ and $G(P'_{D'}(\varphi))\eta_{(T,s)}$ are maps $P_DG_D(T,s)\to G^0P'_{D'}(T',s')$ of
$\mathcal C^0_D$, by (1) for the second, and composed with $j_{G(T)}$ both give $G(j'_{T'}\varphi)$, as
$P_D(G(\varphi))j_{G(T)}=j_{G(T')}G(\varphi)$ and $P'_{D'}(\varphi)j'_T=j'_{T'}\varphi$. By the bijection of
Lemma~\ref{gal-realization-thickening-torsor-cut-adjoint}(3) they are equal.

(3) As $\eta'j_{G(T)}=G(j)$ and $\rho$ is natural, $\rho_{T''}(A(\eta')(A(j_{G(T)})(c_k)))=\rho_{T''}(A(G(j))(c_k))=
A'(j)(\rho_T(c_k))$. As $X_{G(T)}\in J(P(G(T)))$, $A(\eta')(X_{G(T)})\in J(G(T''))$, and $\rho_{T''}$ and $A(\eta')$
commute with the divided powers on these ideals, so $\rho_{T''}(A(\eta')(\gamma_k(X_{G(T)})))=\gamma'_k(X'')$. As
$\rho_{T''}$ and $A(\eta')$ are ring maps and $A'(j)$ is $R'$-linear, the claim follows.
\end{proof}
\leanref{Development/AutomorphicSpectrum/Proofs/BlockDecomposition/EtaNat.lean}{GlobalTheory.CompanionTheory.LocBC.cut_mem}
\leanref{Development/AutomorphicSpectrum/Proofs/BlockDecomposition/EtaNat.lean}{GlobalTheory.CompanionTheory.LocBC.eta_nat}

\begin{lemma}\label{gal-realization-thickening-torsor-base-change}
Let $\sigma:R\to R'$ be a ring map, $(G,\rho)$ a morphism over $\sigma$ from a divided power functor $\Phi$ on a category
$\mathcal C$ to a divided power functor $\Phi'$ on $\mathcal C'$ with every $\rho_T$ injective, $P$ and $P'$ divided power
polynomial structures on $\Phi$ and $\Phi'$, and $(D,\tau)$ and $(D',\tau')$ torsor data of dimension $d$ for $\Phi$ and
$\Phi'$, compatible along $e:D\to G(D')$ (Definition~\ref{gal-realization-thickening-torsor}). Let $(N_i)_{i\in\iota}$ be
a finite family of flat $R$-modules with $R$-linear maps $\nu_i:N_i\to R'$, $(\mathfrak a_{i,r})_r$ and $(\mathfrak
b_{i,r})_r$ sequences of families of ideals of $\Phi$ for each $i$, $(\mathfrak a'_r)_r$ and $(\mathfrak b'_r)_r$
sequences of families of ideals of $\Phi'$, and $r\in\mathbf Z$. Suppose that $P$ is a ladder up to $r$ for $(\mathfrak
a_{i,r'})_{r'}$ and $(\mathfrak b_{i,r'})_{r'}$ for every $i$, that $P'$ is one for $(\mathfrak a'_{r'})_{r'}$ and
$(\mathfrak b'_{r'})_{r'}$, and that $c\rho_T(\mathfrak a_{i,r'}(G(T)))\subset\mathfrak a'_{r'}(T)$ and
$c\rho_T(\mathfrak b_{i,r'}(G(T)))\subset\mathfrak b'_{r'}(T)$ for all $i$, $r'$, $T$ and $c\in\nu_i(N_i)$. If for every
$r'$ with $r-d\leq r'\leq r$ the map
\[
 \bigoplus_iN_i\otimes_R(\mathfrak a_{i,r'}/\mathfrak b_{i,r'})(D)\to(\mathfrak a'_{r'}/\mathfrak b'_{r'})(D'),\qquad
 n\otimes x\mapsto\nu_i(n)\rho_{D'}(A(e)(x)),
\]
is bijective, then for every $j$ the map $\bigoplus_iN_i\otimes_RH^j(\mathcal C,\mathfrak a_{i,r}/\mathfrak b_{i,r})\to
H^j(\mathcal C',\mathfrak a'_r/\mathfrak b'_r)$ sending $n\otimes x$ to the image of $x$ under multiplication by
$\nu_i(n)$ is bijective.
\end{lemma}
\leanref{Development/WeilBound/CoreStatements/PadicHodgeComparison/EtaleCohomologyComparison/PerfectoidPullbacks/LocVan.lean}{CrystallineSemistableComparison.Pb.loc_bc}
\leanref{Development/AutomorphicSpectrum/Proofs/BlockDecomposition/ToH0BijectiveOfInitial.lean}{GlobalTheory.CompanionTheory.LocBC.local_base_change}

\begin{proof}
The proof is by induction on $d$, for all data as in the statement. For $n\in N_i$ and $x\in A(G(T))$ write
$\mu_{i,n}(x)=\nu_i(n)\rho_T(x)$.

\emph{The case $d=0$.} By the last condition of Definition~\ref{gal-realization-thickening-torsor}(2), with no elements
$y_k$, $D$ has exactly one map to every object, so it is an initial object of $\mathcal C$; likewise $D'$ is an initial
object of $\mathcal C'$. By Lemma~\ref{gal-realization-thickening-torsor-adjoint}(2), the groups $H^j$ of the statement
with $j\geq1$ vanish, and $\varepsilon$ identifies $F_i(D)$ with $H^0(\mathcal C,F_i)$ and $F'(D')$ with $H^0(\mathcal
C',F')$, where $F_i=\mathfrak a_{i,r}/\mathfrak b_{i,r}$ and $F'=\mathfrak a'_r/\mathfrak b'_r$. For $T\in\mathcal C'$ the
map $D\to G(T)$ is $G(D'\to T)e$, as $D$ is initial. So multiplication by $\nu_i(n)$ sends $\varepsilon(x)$ to the cochain
$T\mapsto\mu_{i,n}(F_i(G(D'\to T))(F_i(e)(x)))=F'(D'\to T)(\mu_{i,n}(F_i(e)(x)))$, by the naturality of $\rho$, which is
$\varepsilon(\nu_i(n)\rho_{D'}(A(e)(x)))$. Under these identifications the map in degree $0$ is the map of the hypothesis
for $r'=r$, which is bijective.

\emph{The step.} Let the data have dimension $d+1$, and suppose the claim holds for all data of dimension $d$. Form
$\mathcal C_D$, $\mathcal C^0_D$, $U$, $U^0$, $P^0$ and $P_D$ for $(D,\tau)$ and $P$, and the same for $(D',\tau')$ and
$P'$ (Definition~\ref{gal-realization-thickening-torsor-cut}). As every $\rho_T$ is injective and the torsor data are
compatible, Lemma~\ref{gal-realization-thickening-torsor-cut-morphism} gives the functors $G_D$ and $G^0$ and the natural
transformation $\eta$.

Step 1: the cut categories. The functor $G^0$ with the maps $\rho_T$ at $(T,s)$ is a morphism $\Phi^0\to\Phi'^0$ over
$\sigma$ with injective maps, as its conditions at $(T,s)$ are those of $(G,\rho)$ at $T$. The torsor data
$(D,\mathrm{id}_D)$ and $(D',\mathrm{id}_{D'})$ of dimension $d$ of
Lemma~\ref{gal-realization-thickening-torsor-cut-adjoint}(1) are compatible along $e$, which is a map
$(D,\mathrm{id}_D)\to G^0(D',\mathrm{id}_{D'})=(G(D'),e)$ of $\mathcal C^0_D$ as $\tau_1(e,e)=0$: compatibility for
$\tau^0_k$ is compatibility for $\tau_{k+1}$. Let $r_0\in\{r,r-1\}$. By
Lemma~\ref{gal-realization-thickening-torsor-cut-adjoint}(1), $P^0$ and $P'^0$ are ladders up to $r_0$ for the
sequences $(\mathfrak a^0_{i,r'})_{r'}$, $(\mathfrak b^0_{i,r'})_{r'}$ and $(\mathfrak a'^0_{r'})_{r'}$, $(\mathfrak
b'^0_{r'})_{r'}$, and these satisfy the multiplication conditions, which are those at $T$. For $r_0-d\leq r'\leq r_0$ we
have $r-(d+1)\leq r'\leq r$, and the map of the hypothesis at $(D,\mathrm{id}_D)$ and $(D',\mathrm{id}_{D'})$ is the one
at $D$ and $D'$, so it is bijective. By the claim in dimension $d$, the maps
\[
 \bigoplus_iN_i\otimes_RH^j(\mathcal C^0_D,\mathfrak a^0_{i,r_0}/\mathfrak b^0_{i,r_0})\to H^j(\mathcal
 C'^0_{D'},\mathfrak a'^0_{r_0}/\mathfrak b'^0_{r_0})
\]
given by multiplication by the $\nu_i(n)$ are bijective for $r_0\in\{r,r-1\}$ and every $j$.

Step 2: transfer along $P_D$. Let $F_{i,r_0}=(\mathfrak a_{i,r_0}/\mathfrak b_{i,r_0})\circ U^0P_D$ on $\mathcal C_D$ and
$F'_{r_0}=(\mathfrak a'_{r_0}/\mathfrak b'_{r_0})\circ U'^0P'_{D'}$ on $\mathcal C'_{D'}$. For $n\in N_i$, the functor
$G_D$ with the maps $x\mapsto\mu_{i,n}(A(\eta_{(T,s)})(x))$ from $F_{i,r_0}(G_D(T,s))$ to $F'_{r_0}(T,s)$ induces maps
$H^j(\mathcal C_D,F_{i,r_0})\to H^j(\mathcal C'_{D'},F'_{r_0})$: the maps $A(\eta_{(T,s)})$ preserve the families of
ideals, and they are natural as $\eta$ is. Restriction along $P_D$, from $H^j(\mathcal C^0_D,\mathfrak
a^0_{i,r_0}/\mathfrak b^0_{i,r_0})$ to $H^j(\mathcal C_D,F_{i,r_0})$, is bijective by
Lemma~\ref{gal-realization-thickening-torsor-adjoint}(1), as $P_D$ is left adjoint to $\iota$
(Lemma~\ref{gal-realization-thickening-torsor-cut-adjoint}(3)); likewise for restriction along $P'_{D'}$. Restriction
along $P_D$ followed by the map just defined, and multiplication by $\nu_i(n)$ on the cut categories followed by
restriction along $P'_{D'}$, are induced by the functors $P_DG_D$ and $G^0P'_{D'}$ with the maps
$\mu_{i,n}A(\eta_{(T,s)})$ and $\mu_{i,n}$, which $\eta$ relates as in Lemma~\ref{gal-realization-period-prism}; so they
are equal. By Step 1, the maps $\bigoplus_iN_i\otimes_RH^j(\mathcal C_D,F_{i,r_0})\to H^j(\mathcal C'_{D'},F'_{r_0})$ are
bijective for $r_0\in\{r,r-1\}$ and every $j$.

Step 3: the exact sequences. As $P$ is a ladder up to $r$, Lemma~\ref{gal-realization-thickening-torsor-exact} gives for
each $i$ an exact sequence $0\to(\mathfrak a_{i,r}/\mathfrak b_{i,r})\circ U\to F_{i,r}\xrightarrow{\beta}F_{i,r-1}\to0$
of functors on $\mathcal C_D$, and likewise on $\mathcal C'_{D'}$. The maps given by $G_D$, with $\mu_{i,n}$ on the first
terms and $\mu_{i,n}A(\eta_{(T,s)})$ on the other two, commute with the maps of these sequences. For the first map this
holds as $\eta_{(T,s)}j_{G(T)}=G(j'_T)$ and $\rho$ is natural. For $\beta$, let $x$ be the class of
$\sum_kA(j_{G(T)})(c_k)\gamma_k(X_{G(T)})\in\mathfrak a_{i,r}(P(G(T)))$. By
Lemma~\ref{gal-realization-thickening-torsor-cut-morphism}(3), with $a=\nu_i(n)$, $j=j'_T$, $\eta'=\eta_{(T,s)}$ and
$X''=X'_T$, $\mu_{i,n}A(\eta_{(T,s)})(x)$ is the class of $\sum_kA'(j'_T)(\mu_{i,n}(c_k))\gamma'_k(X'_T)$, an element of
$\mathfrak a'_r(P'(T))$ by the multiplication condition; so both composites send $x$ to the class of
$\sum_kA'(j'_T)(\mu_{i,n}(c_{k+1}))\gamma'_k(X'_T)$. On cochains these sequences are exact in every degree, as they are
exact at every object, so they give long exact sequences of cohomology on $\mathcal C_D$ and on $\mathcal C'_{D'}$, and
the induced maps commute with all their maps, the connecting maps included. Tensor the sequence for each $i$ with the
flat $R$-module $N_i$ and take the direct sum over $i$. By Step 2 and the five lemma, as in Step 2 of the proof of
Lemma~\ref{gal-realization-thickening-cosieves}(1), the maps $\bigoplus_iN_i\otimes_RH^j(\mathcal C_D,(\mathfrak
a_{i,r}/\mathfrak b_{i,r})\circ U)\to H^j(\mathcal C'_{D'},(\mathfrak a'_r/\mathfrak b'_r)\circ U')$ are bijective for
every $j$.

Step 4: restriction along $U$. As $D$ has a map to every object, there is a family of maps $\sigma_T:D\to T$, and $U$ is
left adjoint to $S_\sigma$ (Lemma~\ref{gal-realization-thickening-torsor-cut-adjoint}(2)). So restriction along $U$, from
$H^j(\mathcal C,\mathfrak a_{i,r}/\mathfrak b_{i,r})$ to $H^j(\mathcal C_D,(\mathfrak a_{i,r}/\mathfrak b_{i,r})\circ U)$,
is bijective by Lemma~\ref{gal-realization-thickening-torsor-adjoint}(1), and likewise along $U'$. Restriction along $U$
followed by the map of Step 3, and multiplication by $\nu_i(n)$ followed by restriction along $U'$, are induced by the
same functor $UG_D=GU'$ with the same maps. So Step 3 gives the claim.
% Lean: LocBC.main, by induction on d. The case d = 0 is base_case (Frame.initial, toH0 and toH0_bijective_of_initial,
% H_succ_sub, tens_transfer, tens_zero). The step is LocBC.step: Step 1 is the induction hypothesis for m1, Cut.Q1, Cut.E1,
% e1 and fr1, with ladder_mono, Cut.ladder_truncation and le_cover_restrict_compat; Step 2 is hP, through P.adj for
% Cut.adj and pull_restrict_first_piece_compat (P.prism for eta); Step 3 is filtration_critical_comparison_map,
% P.ses_cochains and tens_les, with alpha_hom, beta_τ (rho_sum) and hadd_Pt; Step 4 is P.adj for Pt.adj and pull_U.
\end{proof}
\leanref{Development/AutomorphicSpectrum/Proofs/BlockDecomposition/ToH0BijectiveOfInitial.lean}{GlobalTheory.CompanionTheory.LocBC.main}
\leanref{Development/AutomorphicSpectrum/Proofs/BlockDecomposition/ToH0BijectiveOfInitial.lean}{GlobalTheory.CompanionTheory.LocBC.base_case}
\leanref{Development/AutomorphicSpectrum/Proofs/BlockDecomposition/ToH0BijectiveOfInitial.lean}{GlobalTheory.CompanionTheory.LocBC.step}
\leanref{Development/WeilBound/CoreStatements/PadicHodgeComparison/EtaleCohomologyComparison/PerfectoidPullbacks/LocalisationLadders/TruncationHomExt.lean}{CrystallineSemistableComparison.Pb.Frame.initial}
\leanref{Development/AutomorphicSpectrum/Proofs/BlockDecomposition/EtaNat.lean}{GlobalTheory.CompanionTheory.LocBC.m1}
\leanref{Development/AutomorphicSpectrum/Proofs/BlockDecomposition/EtaNat.lean}{GlobalTheory.CompanionTheory.LocBC.mP}
\leanref{Development/AutomorphicSpectrum/Proofs/BlockDecomposition/EtaNat.lean}{GlobalTheory.CompanionTheory.LocBC.mPt}
\leanref{Development/AutomorphicSpectrum/Proofs/BlockDecomposition/PullRestrictFirstPieceCompat.lean}{GlobalTheory.CompanionTheory.LocBC.pull_restrict_first_piece_compat}
\leanref{Development/AutomorphicSpectrum/Proofs/BlockDecomposition/PullRestrictFirstPieceCompat.lean}{GlobalTheory.CompanionTheory.LocBC.alpha_hom}
\leanref{Development/AutomorphicSpectrum/Proofs/BlockDecomposition/PullRestrictFirstPieceCompat.lean}{GlobalTheory.CompanionTheory.LocBC.beta_τ}
\leanref{Development/AutomorphicSpectrum/Proofs/BlockDecomposition/ClsCongr.lean}{GlobalTheory.CompanionTheory.LocBC.tens_les}
\leanref{Development/AutomorphicSpectrum/Proofs/BlockDecomposition/SqCSmul.lean}{GlobalTheory.CompanionTheory.LocBC.pull_U}
\leanref{Development/WeilBound/CoreStatements/PadicHodgeComparison/EtaleCohomologyComparison/CechPairings/MemBoundariesZeroIff.lean}{CrystallineSemistableComparison.P.ses_cochains}

\begin{lemma}\label{gal-realization-thickening-local-base-change}
Let $\sigma:R\to R'$, $(R,I,\gamma)$, $(R',I',\gamma')$, $\mathcal O$, $X$, $\kappa$ and $\kappa'$ be as in
Definition~\ref{gal-realization-thickening-nilpotent}(4), with $R'$ flat over $R$, let $p\neq0$ be an integer nilpotent in
$R$, let $V\subset X$ be an open with coordinates, and $\Gamma$ the restriction of scalars along $\sigma$. For $j\geq0$,
the map $R'\otimes_RH^j(\mathrm{Th}(X/R)_V,\mathcal O)\to H^j(\mathrm{Th}(X/R')_V,\mathcal O)$, $a\otimes x\mapsto
a\,\Gamma^*(x)$, is bijective.
\end{lemma}
\leanref{Development/AutomorphicSpectrum/Proofs/BlockDecomposition/StdSomeZero.lean}{GlobalTheory.CompanionTheory.ChartLocBc.local_bc}

\begin{proof}
Let $\mathcal Q=\mathrm{Th}(X/R)_V$ and $\mathcal Q'=\mathrm{Th}(X/R')_V$; $\Gamma$ maps $\mathcal Q'$ into $\mathcal Q$,
as it does not change the point. By Lemma~\ref{gal-realization-thickening-chart}(3) there is a divided power chart $D$ on
$V$ over $(R,I,\gamma,p)$, and by Lemma~\ref{gal-realization-thickening-chart}(4) a divided power chart $D'$ on $V$ over
$(R',I',\gamma',p)$ with a map $e:D\to\Gamma(D')$ such that $e(t_k)=t'_k$ and $R'\otimes_RB_D\to B_{D'}$, $a\otimes
b\mapsto i_{D'}(a)e(b)$, is bijective. By Lemma~\ref{gal-realization-thickening-chart-torsor}(1) and (2), $D$ and $D'$
give torsor data for $\Phi_{\mathcal Q}$ and $\Phi_{\mathcal Q'}$, compatible along $e$ for the morphism given by
$\Gamma$. By Lemma~\ref{gal-realization-thickening-chart-torsor}(3) there are divided power polynomial structures on
$\Phi_{\mathcal Q}$ and $\Phi_{\mathcal Q'}$ that respect $(J^{[r]})_r$ and $(0)_r$ in every degree; as $0\subset
J^{[r]}$, they are ladders up to $0$ for these sequences. Multiplication by $a\in R'$ carries $J^{[r]}$ into $J^{[r]}$
and $0$ into $0$.

Apply Lemma~\ref{gal-realization-thickening-torsor-base-change} with one module $N=R'$, flat over $R$, $\nu$ its identity,
$\mathfrak a_r=\mathfrak a'_r=J^{[r]}$, $\mathfrak b_r=\mathfrak b'_r=0$ and $r=0$. For $r'\leq0$, $(J^{[r']}/0)(D)=B_D$
and $(J^{[r']}/0)(D')=B_{D'}$, and the map of its hypothesis is $a\otimes b\mapsto i_{D'}(a)e(b)$, which is bijective.
Its conclusion, for $J^{[0]}/0=\mathcal O$, is the claim.
\end{proof}
\leanref{Development/AutomorphicSpectrum/Proofs/BlockDecomposition/StdSomeZero.lean}{GlobalTheory.CompanionTheory.ChartLocBc.tens_of_ring}

\begin{lemma}\label{gal-realization-thickening-cosieves}
Let $\mathcal C$ be a category. A \emph{cosieve} of $\mathcal C$ is a full subcategory $\mathcal U$ such that $T'\in
\mathcal U$ whenever $T\in\mathcal U$ and there is a map $T\to T'$. For a family $(\mathcal U_\alpha)_{\alpha\in A}$ of
cosieves and a nonempty $S\subset A$, let $\mathcal U_S=\bigcap_{\alpha\in S}\mathcal U_\alpha$; cohomology is that of
Definition~\ref{gal-realization-split-diagrams}(1).
\begin{enumerate}
\item Let $\sigma:\Lambda\to\Lambda'$ be a ring map, $G:\mathcal C'\to\mathcal C$ a functor, $(N_i)_{i\in\iota}$ a finite
  family of flat $\Lambda$-modules, $F_i$ functors from $\mathcal C$ to $\Lambda$-modules and $F'$ a functor from
  $\mathcal C'$ to $\Lambda'$-modules. For $n\in N_i$ let $\tau_i(n)_T:F_i(G(T))\to F'(T)$ be $\sigma$-semilinear maps,
  natural in $T$, with $\tau_i(n+n')=\tau_i(n)+\tau_i(n')$ and $\tau_i(an)=\sigma(a)\tau_i(n)$ for $a\in\Lambda$. Let $A$
  be finite, and $(\mathcal U_\alpha)_{\alpha\in A}$ and $(\mathcal U'_\alpha)_{\alpha\in A}$ families of cosieves of
  $\mathcal C$ and of $\mathcal C'$ such that every object of $\mathcal C$ lies in some $\mathcal U_\alpha$, every object
  of $\mathcal C'$ lies in some $\mathcal U'_\alpha$, and $G$ maps $\mathcal U'_\alpha$ into $\mathcal U_\alpha$. If for
  every nonempty $S\subset A$ and every $j$ the map $\bigoplus_iN_i\otimes_\Lambda H^j(\mathcal U_S,F_i)\to
  H^j(\mathcal U'_S,F')$, $n\otimes x\mapsto\tau_i(n)^*(x)$, is bijective, then so is
  $\bigoplus_iN_i\otimes_\Lambda H^j(\mathcal C,F_i)\to H^j(\mathcal C',F')$ for every $j$.
\item Let $G:\mathcal C'\to\mathcal C$ be a functor, $F$ and $F'$ functors from $\mathcal C$ and $\mathcal C'$ to
  $\Lambda$-modules, $F(G(T))\to F'(T)$ $\Lambda$-linear maps natural in $T$, and $(\mathcal U_\alpha)$ and $(\mathcal
  U'_\alpha)$ as in (1). If the induced maps $H^j(\mathcal U_S,F)\to H^j(\mathcal U'_S,F')$ are bijective for every
  nonempty $S$ and every $j$, so are the maps $H^j(\mathcal C,F)\to H^j(\mathcal C',F')$.
\item Let $F$ be a functor from $\mathcal C$ to $\Lambda$-modules, $(\mathcal U_\alpha)_{\alpha\in A}$ cosieves with every
  object in some $\mathcal U_\alpha$, $A$ finite with $N$ elements, and $s\geq0$. If $H^j(\mathcal U_S,F)=0$ for every
  nonempty $S$ and every $j<s$, then $H^j(\mathcal C,F)=0$ for $j<s$. If $H^j(\mathcal U_S,F)=0$ for every nonempty $S$
  and every $j$, then $H^j(\mathcal C,F)=0$ for every $j$. If $H^j(\mathcal U_S,F)=0$ for every nonempty $S$ and every
  $j>s$, then $H^j(\mathcal C,F)=0$ for $j\geq s+N$.
\end{enumerate}
\end{lemma}
\leanref{Development/WeilBound/CoreStatements/PadicHodgeComparison/EtaleCohomologyComparison/PerfectoidPullbacks/CapOComp.lean}{CrystallineSemistableComparison.Pb.Cosieve}
\leanref{Development/AutomorphicSpectrum/Proofs/BlockDecomposition/FlatBaseChangeCohomology.lean}{GlobalTheory.CompanionTheory.CosieveBC.base_change_of_cover}
\leanref{Development/AutomorphicSpectrum/Proofs/BlockDecomposition/SmulMPull.lean}{GlobalTheory.CompanionTheory.DisoHm.cover_bij}
\leanref{Development/AutomorphicSpectrum/Proofs/BlockDecomposition/CechVanishingFromCoverStatement.lean}{GlobalTheory.T.cech_vanishing_from_cover}

\begin{proof}
A chain $c_0\to\cdots\to c_k$ of $\mathcal C$ \emph{starts in} a full subcategory $\mathcal W$ if $c_0\in\mathcal W$. If
$\mathcal W$ is a cosieve, such a chain lies in $\mathcal W$, so the chains of $\mathcal W$ are the chains of $\mathcal C$
that start in $\mathcal W$; and every face of such a chain starts in $\mathcal W$, as it starts at $c_0$ or at $c_1$.
Unions and intersections of cosieves are cosieves. For a family $(\mathcal U_\alpha)_{\alpha\in A}$ of cosieves and
finite subsets $A',S\subset A$, let $\mathcal W_{A',S}$ be the full subcategory of the objects that lie in some $\mathcal
U_\alpha$ with $\alpha\in A'$ and in every $\mathcal U_\alpha$ with $\alpha\in S$, a cosieve. Then $\mathcal
W_{\emptyset,S}$ is empty, $\mathcal W_{A,\emptyset}=\mathcal C$ if $A$ is finite and every object lies in some
$\mathcal U_\alpha$, and for $a\in A\smallsetminus A'$
\[
 \mathcal W_{A'\cup\{a\},S}=\mathcal W_{A',S}\cup\mathcal U_{S\cup\{a\}},\qquad \mathcal W_{A',S}\cap\mathcal
 U_{S\cup\{a\}}=\mathcal W_{A',S\cup\{a\}}.
\]

(1) Form $\mathcal W'_{A',S}$ from the $\mathcal U'_\alpha$ in the same way; $G$ maps $\mathcal W'_{A',S}$ into
$\mathcal W_{A',S}$, as it maps each $\mathcal U'_\alpha$ into $\mathcal U_\alpha$. For full subcategories $\mathcal
W\subset\mathcal C$ and $\mathcal W'\subset\mathcal C'$ with $G(\mathcal W')\subset\mathcal W$, say that \emph{base
change holds} for $(\mathcal W,\mathcal W')$ if for every $j$ the map $\bigoplus_iN_i\otimes_\Lambda H^j(\mathcal
W,F_i)\to H^j(\mathcal W',F')$, $n\otimes x\mapsto\tau_i(n)^*(x)$, is bijective. This map is well defined: $\tau_i(n)^*$
is additive in $n$ and in $x$, and $\tau_i(an)^*(x)=\sigma(a)\tau_i(n)^*(x)=\tau_i(n)^*(ax)$, as $\tau_i(n)$ is
$\sigma$-semilinear. The claim is that base change holds for $(\mathcal W_{A,\emptyset},\mathcal W'_{A,\emptyset})=
(\mathcal C,\mathcal C')$.

\emph{Step 1.} Let $\mathcal U$ and $\mathcal V$ be cosieves of a category $\mathcal D$ and $E$ a functor from
$\mathcal D$ to $\Lambda$-modules, and write $C^k(\mathcal W,E)$ for the $k$-cochains of a full subcategory $\mathcal
W$. For every $k$ the sequence
\[
 0\to C^k(\mathcal U\cup\mathcal V,E)\to C^k(\mathcal U,E)\oplus C^k(\mathcal V,E)\to C^k(\mathcal U\cap\mathcal V,E)\to0,
\]
with the maps $x\mapsto(x|_{\mathcal U},x|_{\mathcal V})$ and $(a,b)\mapsto a|_{\mathcal U\cap\mathcal
V}-b|_{\mathcal U\cap\mathcal V}$ given by restriction of cochains, is exact, and its maps commute with the coboundaries.
Restriction commutes with the coboundaries, as the faces of a chain of a full subcategory are chains of it. The first
map is injective: a chain of $\mathcal U\cup\mathcal V$ starts in $\mathcal U$ or in $\mathcal V$, and is then a chain of
$\mathcal U$ or of $\mathcal V$. The composite is $0$. If $a$ and $b$ agree on the chains of $\mathcal U\cap\mathcal V$,
the cochain of $\mathcal U\cup\mathcal V$ equal to $a$ on the chains that start in $\mathcal U$ and to $b$ on the others
restricts to $a$ and to $b$, as a chain of $\mathcal V$ that starts in $\mathcal U$ is a chain of $\mathcal
U\cap\mathcal V$. The last map is surjective: a cochain $y$ of $\mathcal U\cap\mathcal V$ is the image of $(y',0)$,
where $y'$ is $y$ on the chains of $\mathcal U$ that start in $\mathcal V$ and $0$ on the others. So there is a long
exact sequence
\[
 0\to H^0(\mathcal U\cup\mathcal V,E)\to\cdots\to H^{j-1}(\mathcal U\cap\mathcal V,E)\xrightarrow{\partial}
 H^j(\mathcal U\cup\mathcal V,E)\to H^j(\mathcal U,E)\oplus H^j(\mathcal V,E)\to H^j(\mathcal U\cap\mathcal V,E)
 \xrightarrow{\partial}\cdots.
\]

\emph{Step 2.} Let $\mathcal U,\mathcal V$ be cosieves of $\mathcal C$ and $\mathcal U',\mathcal V'$ cosieves of
$\mathcal C'$ with $G(\mathcal U')\subset\mathcal U$ and $G(\mathcal V')\subset\mathcal V$. If base change holds for
$(\mathcal U,\mathcal U')$, $(\mathcal V,\mathcal V')$ and $(\mathcal U\cap\mathcal V,\mathcal U'\cap\mathcal V')$, it
holds for $(\mathcal U\cup\mathcal V,\mathcal U'\cup\mathcal V')$. For each $i$ and $n\in N_i$, the maps on cochains
$x\mapsto\bigl(c'_0\to\cdots\to c'_k\mapsto\tau_i(n)(x(G(c'_0)\to\cdots\to G(c'_k)))\bigr)$ form a map from the
sequence of Step 1 for $F_i$ to the one for $F'$: they commute with restriction and with the coboundaries. The maps they
induce on cohomology therefore commute with the maps of the two long exact sequences, the connecting maps included.
Tensor the long exact sequence for $F_i$ with the flat module $N_i$ and take the direct sum over $i$. The result is
exact, and $n\otimes x\mapsto\tau_i(n)^*(x)$ maps it to the long exact sequence for $F'$, with commuting squares. On
$H^j(\mathcal U)\oplus H^j(\mathcal V)$ this map is the direct sum of the maps for $(\mathcal U,\mathcal U')$ and
$(\mathcal V,\mathcal V')$, so it is bijective for every $j$, and so is the map on $H^j(\mathcal U\cap\mathcal V)$. The
five lemma, applied to the terms from $H^{j-1}(\mathcal U)\oplus H^{j-1}(\mathcal V)$ to $H^j(\mathcal U\cap\mathcal
V)$, with zeros in place of the first two when $j=0$, shows that the map on $H^j(\mathcal U\cup\mathcal V)$ is bijective.

\emph{Step 3.} By induction on $A'$, base change holds for $(\mathcal W_{A',S},\mathcal W'_{A',S})$ for all finite
$A',S\subset A$. For $A'=\emptyset$ both categories are empty, so all their cohomology is $0$. For $a\notin A'$, apply
Step 2 with $\mathcal U=\mathcal W_{A',S}$, $\mathcal V=\mathcal U_{S\cup\{a\}}$, $\mathcal U'=\mathcal W'_{A',S}$ and
$\mathcal V'=\mathcal U'_{S\cup\{a\}}$: base change holds for the first pair by induction, for the second by hypothesis,
as $S\cup\{a\}$ is nonempty, and for the intersections $\mathcal W_{A',S\cup\{a\}}$ and $\mathcal W'_{A',S\cup\{a\}}$ by
induction. The unions are $\mathcal W_{A'\cup\{a\},S}$ and $\mathcal W'_{A'\cup\{a\},S}$. The case $A'=A$,
$S=\emptyset$ is the claim.

(2) Apply (1) with $\Lambda'=\Lambda$, $\sigma$ the identity, one module $N=\Lambda$, $F_1=F$ and $\tau(a)$ equal to $a$
times the given maps. The map $\Lambda\otimes_\Lambda M\to M'$, $a\otimes x\mapsto af(x)$, is bijective if and only if
the linear map $f:M\to M'$ is, as it is $f$ composed with $\Lambda\otimes_\Lambda M\cong M$. So the hypotheses of (1)
hold, and its conclusion is the claim.

(3) For a cosieve $\mathcal W$ and $j\geq0$, say that $\mathcal W$ is \emph{$j$-acyclic} if every $j$-cochain $x$ of
$\mathcal C$ whose coboundary $dx$ vanishes on the $(j+1)$-chains that start in $\mathcal W$ satisfies: $x$ vanishes on
the objects of $\mathcal W$ if $j=0$, and there is a $(j-1)$-cochain $y$ of $\mathcal C$ with $x(s)=dy(s)$ for every
$j$-chain $s$ that starts in $\mathcal W$ if $j\geq1$. An empty cosieve is $j$-acyclic for every $j$.

\emph{Step 1.} If $H^j(\mathcal W,F)=0$, then $\mathcal W$ is $j$-acyclic. Let $x$ be as above. As the chains of
$\mathcal W$ are the chains of $\mathcal C$ that start in $\mathcal W$, $x$ restricts to a cocycle $x_{\mathcal W}$ of
$\mathcal W$. If $j=0$, $x_{\mathcal W}=0$. If $j\geq1$, $x_{\mathcal W}=dy'$ for a $(j-1)$-cochain $y'$ of $\mathcal
W$. Let $y$ be the cochain of $\mathcal C$ equal to $y'$ on the chains that start in $\mathcal W$ and to $0$ on the
others. Its restriction to $\mathcal W$ is $y'$, and restriction commutes with $d$, so $dy$ and $x$ agree on the chains
that start in $\mathcal W$.

\emph{Step 2.} Let $\mathcal U$ and $\mathcal V$ be $j$-acyclic cosieves, and if $j\geq1$ let $\mathcal U\cap\mathcal V$
be $(j-1)$-acyclic. Then $\mathcal U\cup\mathcal V$ is $j$-acyclic. Let $x$ be a $j$-cochain with $dx=0$ on the chains
that start in $\mathcal U\cup\mathcal V$. If $j=0$, $x$ vanishes on the objects of $\mathcal U$ and of $\mathcal V$. Let
$j\geq1$. There are $(j-1)$-cochains $a$ and $b$ with $x=da$ on the $j$-chains that start in $\mathcal U$ and $x=db$ on
those that start in $\mathcal V$. Then $d(a-b)$ vanishes on the $j$-chains that start in $\mathcal U\cap\mathcal V$.
There is a $(j-1)$-cochain $a'$ with $x=da'$ on the $j$-chains that start in $\mathcal U$ and $a'=b$ on the
$(j-1)$-chains that start in $\mathcal U\cap\mathcal V$: if $j=1$, $a-b$ vanishes on the objects of $\mathcal
U\cap\mathcal V$, and $a'=a$; if $j\geq2$, $a-b=dc$ on the $(j-1)$-chains that start in $\mathcal U\cap\mathcal V$, and
$a'=a-dc$, as $ddc=0$. Let $z$ be the $(j-1)$-cochain equal to $a'$ on the chains that start in $\mathcal U$ and to $b$ on
the others. For a $j$-chain $s$ that starts in $\mathcal U$, every face of $s$ starts in $\mathcal U$, so
$dz(s)=da'(s)=x(s)$. For a $j$-chain $s$ that starts in $\mathcal V$ but not in $\mathcal U$, every face of $s$ starts
in $\mathcal V$, and on a face that starts in $\mathcal U$ also $a'=b$; so $dz(s)=db(s)=x(s)$.

\emph{Step 3.} If the class of all objects is $j$-acyclic, then $H^j(\mathcal C,F)=0$: every class is the class of a
cocycle $x$, and $x=0$ if $j=0$, and $x=dy$ if $j\geq1$.

\emph{Step 4.} Let $A$ have $N$ elements, and consider the three hypotheses in turn. By induction on $A'$, every
$\mathcal W_{A',S}$ is $j$-acyclic, for all $j<s$, for all $j$, and for all $j\geq s+|A'|$ respectively. For
$A'=\emptyset$, $\mathcal W_{\emptyset,S}$ is empty. For $a\notin A'$ and $j$ in the given range for $A'\cup\{a\}$,
$\mathcal W_{A',S}$ is $j$-acyclic by induction, as $j$ is in the range for $A'$; $\mathcal U_{S\cup\{a\}}$ is $j$-acyclic
by Step 1, as $j<s$, any $j$, or $j\geq s+|A'|+1>s$; and if $j\geq1$, $\mathcal W_{A',S\cup\{a\}}$ is $(j-1)$-acyclic by
induction, as $j-1$ is in the range for $A'$. By Step 2 and the identities above, $\mathcal W_{A'\cup\{a\},S}$ is
$j$-acyclic. For $A'=A$ and $S=\emptyset$, $\mathcal W_{A,\emptyset}$ is the class of all objects, and Step 3 gives
$H^j(\mathcal C,F)=0$ for $j<s$, for every $j$, and for $j\geq s+N$ respectively.
\end{proof}
\leanref{Development/AutomorphicSpectrum/Proofs/BlockDecomposition/SmulMPull.lean}{GlobalTheory.CompanionTheory.DisoHm.tens_iff}
\leanref{Development/AutomorphicSpectrum/Proofs/BlockDecomposition/MvFGExact.lean}{GlobalTheory.CompanionTheory.CosieveBC.mvFG_exact}
\leanref{Development/AutomorphicSpectrum/Proofs/BlockDecomposition/MvFGExact.lean}{GlobalTheory.CompanionTheory.CosieveBC.mvF_inj}
\leanref{Development/AutomorphicSpectrum/Proofs/BlockDecomposition/MvFGExact.lean}{GlobalTheory.CompanionTheory.CosieveBC.mvG_surj}
\leanref{Development/AutomorphicSpectrum/Proofs/BlockDecomposition/MvFGExact.lean}{GlobalTheory.CompanionTheory.CosieveBC.bc_sup}
\leanref{Development/AutomorphicSpectrum/Proofs/BlockDecomposition/SubsingletonTensor.lean}{GlobalTheory.CompanionTheory.CosieveBC.les_pair}
\leanref{Development/AutomorphicSpectrum/Proofs/BlockDecomposition/SubsingletonTensor.lean}{GlobalTheory.CompanionTheory.CosieveBC.bc_ses}
\leanref{Development/AutomorphicSpectrum/Proofs/BlockDecomposition/SubsingletonTensor.lean}{GlobalTheory.CompanionTheory.CosieveBC.bc_prod}
\leanref{Development/AutomorphicSpectrum/Proofs/BlockDecomposition/BalancedLiftTmul.lean}{GlobalTheory.CompanionTheory.CosieveBC.tens_five}
\leanref{Development/AutomorphicSpectrum/Proofs/BlockDecomposition/FlatBaseChangeCohomology.lean}{GlobalTheory.CompanionTheory.CosieveBC.bc_Wc}
\leanref{Development/AutomorphicSpectrum/Proofs/BlockDecomposition/FlatBaseChangeCohomology.lean}{GlobalTheory.CompanionTheory.CosieveBC.tens_top}
\leanref{Development/AutomorphicSpectrum/Proofs/BlockDecomposition/CechVanishingFromCover.lean}{GlobalTheory.CompanionTheory.CosieveMV.van_of_subsingleton}
\leanref{Development/AutomorphicSpectrum/Proofs/BlockDecomposition/CechVanishingFromCover.lean}{GlobalTheory.CompanionTheory.CosieveMV.van_sup}
\leanref{Development/AutomorphicSpectrum/Proofs/BlockDecomposition/CechVanishingFromCover.lean}{GlobalTheory.CompanionTheory.CosieveMV.van_Wc}
\leanref{Development/AutomorphicSpectrum/Proofs/BlockDecomposition/CechVanishingFromCover.lean}{GlobalTheory.CompanionTheory.CosieveMV.subsingleton_of_van}

\begin{definition}\label{gal-realization-thickening-zariski-class}
Let $\mathcal C$ be a category, $R$ a ring, $F$ a functor from $\mathcal C$ to $R$-modules and $\mathcal P$ a class of
objects of $\mathcal C$. Maps are composed as functions: $\chi\psi$ is $\psi$ followed by $\chi$. Then $\mathcal P$ is a
\emph{Zariski class} for $F$ if the following hold.
\begin{enumerate}
\item If $T\in\mathcal P$ and there is a map $T\to T'$, then $T'\in\mathcal P$.
\item For every object $T$ there are an integer $m\geq0$ and classes $M_1,\ldots,M_m$ of maps with source $T$ such that
  $\chi\psi\in M_k$ whenever $\psi\in M_k$ and $\chi$ is a map out of the target of $\psi$, the target of every map in
  $M_k$ lies in $\mathcal P$, and the following hold for every map $\varphi:T\to T'$.
  \begin{enumerate}
  \item For every $S\subset\{1,\ldots,m\}$ there is a map $\lambda_S:T'\to L_S$ such that $\lambda_S\varphi\in M_k$ for
    every $k\in S$, and every map $\psi:T'\to T''$ with $\psi\varphi\in M_k$ for every $k\in S$ is $\chi\lambda_S$ for
    exactly one map $\chi:L_S\to T''$.
  \item There are a commutative ring $\Lambda$, elements $f_1,\ldots,f_m$ that generate the unit ideal of $\Lambda$, and
    ring maps from $\Lambda$ to the rings of $R$-linear endomorphisms of $F(T')$ and of every $F(L_S)$ with the following
    properties. The maps $F(\lambda_S)$, and $F(t)$ for every map $t:L_S\to L_{S'}$ with $t\lambda_S=\lambda_{S'}$, are
    $\Lambda$-linear. Every $f_k$ with $k\in S$ acts bijectively on $F(L_S)$. With $f_S=\prod_{k\in S}f_k$, for every
    $z\in F(L_S)$ some $f_S^nz$ lies in the image of $F(\lambda_S)$, and every $x\in F(T')$ with $F(\lambda_S)(x)=0$ has
    $f_S^nx=0$ for some $n\geq0$.
  \end{enumerate}
\end{enumerate}
\end{definition}
\leanref{Development/WeilBound/CoreStatements/PadicHodgeComparison/EtaleCohomologyComparison/ZariskiEtaleLifting/XSmulSucc.lean}{CrystallineSemistableComparison.B.ZarCov}
\leanref{Development/WeilBound/CoreStatements/PadicHodgeComparison/EtaleCohomologyComparison/ZariskiEtaleLifting/XSmulSucc.lean}{CrystallineSemistableComparison.B.LocMod}

\begin{definition}\label{gal-realization-thickening-zariski-double-up-to}
Let $\Lambda$ be a commutative ring and $s\in\Lambda$. A $\Lambda$-linear map $f:M\to N$ is \emph{injective up to
$s$} if $sx=0$ for every $x\in M$ with $f(x)=0$, \emph{surjective up to $s$} if $sy\in f(M)$ for every $y\in N$, and an
\emph{isomorphism up to $s$} if it is both. Linear maps $f:M\to N$ and $g:N\to P$ are \emph{exact up to $s$} if
$sy\in f(M)$ for every $y\in N$ with $g(y)=0$. Let $C^0\to C^1\to\cdots$ be a complex of $\Lambda$-modules with
differential $d$, and $\epsilon:M\to C^0$ a linear map with $d\epsilon=0$. Then $C$ is a \emph{resolution up to $s$} of
$M$ by $\epsilon$ if $\epsilon$ is injective up to $s$, $\epsilon$ and $d:C^0\to C^1$ are exact up to $s$, and
$d:C^j\to C^{j+1}$ and $d:C^{j+1}\to C^{j+2}$ are exact up to $s$ for every $j\geq0$. For $s=1$ these are injectivity,
surjectivity, isomorphism, exactness, and exactness of $0\to M\to C^0\to C^1\to\cdots$.
\end{definition}
\leanref{Development/WeilBound/CoreStatements/PadicHodgeComparison/EtaleCohomologyComparison/CechPairings/InstAddCommGroup.lean}{CrystallineSemistableComparison.InjUpTo}
\leanref{Development/WeilBound/CoreStatements/PadicHodgeComparison/EtaleCohomologyComparison/CechPairings/InstAddCommGroup.lean}{CrystallineSemistableComparison.SurjUpTo}
\leanref{Development/WeilBound/CoreStatements/PadicHodgeComparison/EtaleCohomologyComparison/CechPairings/InstAddCommGroup.lean}{CrystallineSemistableComparison.UpTo}
\leanref{Development/WeilBound/CoreStatements/PadicHodgeComparison/EtaleCohomologyComparison/CechPairings/InstAddCommGroup.lean}{CrystallineSemistableComparison.ExactUpTo}
\leanref{Development/WeilBound/CoreStatements/PadicHodgeComparison/EtaleCohomologyComparison/CechPairings/InstAddCommGroup.lean}{CrystallineSemistableComparison.Cx.ResUpTo}

\begin{lemma}\label{gal-realization-thickening-zariski-double}
Let $\Lambda$ be a commutative ring, $s\in\Lambda$, and $K$ a double complex of $\Lambda$-modules $K^{p,q}$, $p,q\geq0$,
with commuting differentials $d_1:K^{p,q}\to K^{p+1,q}$ and $d_2:K^{p,q}\to K^{p,q+1}$, and let $\mathrm{Tot}(K)$ be its
total complex, with $\mathrm{Tot}^n(K)=\bigoplus_{p+q=n}K^{p,q}$ and differential $d_1+(-1)^pd_2$ on $K^{p,q}$. Use the
terms of Definition~\ref{gal-realization-thickening-zariski-double-up-to}.
\begin{enumerate}
\item Let $A$ be a complex of $\Lambda$-modules and $\epsilon^q:A^q\to K^{0,q}$ maps with $d_1\epsilon^q=0$ and
  $d_2\epsilon^q=\epsilon^{q+1}d_A$, such that for every $q$ the row $K^{0,q}\to K^{1,q}\to\cdots$ is a resolution up
  to $s$ of $A^q$ by $\epsilon^q$. Then for every $n$ the map $H^n(A)\to H^n(\mathrm{Tot}(K))$ induced by the maps
  $\epsilon^q$ is an isomorphism up to $s^{n+2}$.
\item Let $n\geq0$ and let $\varphi:K\to K'$ be a map of such double complexes that induces isomorphisms up to $s$
  $H^q(K^{p,\bullet})\to H^q(K'^{p,\bullet})$ on the cohomology of the columns for all $p,q$ with $p+q\leq n+1$. Then
  the map $H^n(\mathrm{Tot}(K))\to H^n(\mathrm{Tot}(K'))$ induced by $\varphi$ is an isomorphism up to $s^{2n+3}$.
\end{enumerate}
For $s=1$, (1) says that the $\epsilon^q$ induce isomorphisms $H^n(A)\to H^n(\mathrm{Tot}(K))$ when every sequence
$0\to A^q\to K^{0,q}\to K^{1,q}\to\cdots$ is exact, and (2) that $\varphi$ induces an isomorphism on $H^n$ of the total
complexes when it induces isomorphisms on the cohomology of the columns for $p+q\leq n+1$.
\end{lemma}
\leanref{Development/WeilBound/CoreStatements/PadicHodgeComparison/EtaleCohomologyComparison/CechPairings/DcColSurj.lean}{CrystallineSemistableComparison.P.dc_row}
\leanref{Development/WeilBound/CoreStatements/PadicHodgeComparison/EtaleCohomologyComparison/CechPairings/DoubleComplexMapMemTotal.lean}{CrystallineSemistableComparison.P.dc_map_col}

\begin{proof}
Write $D$ for the differentials of the total complexes. For $x\in\mathrm{Tot}^m(K)$ write $x_{p,q}$ for its component in
$K^{p,q}$, and for $c\in K^{p,q}$ write $[c]$ for $c$ viewed in $\mathrm{Tot}^{p+q}(K)$, so that
$D[c]=[d_1c]+(-1)^p[d_2c]$. The same notation is used for $K'$.

(1) Let $K^t$ be the double complex with $(K^t)^{p,q}=K^{q,p}$, $d_1^t=d_2$ and $d_2^t=d_1$. The map
$\mathrm{Tot}(K)\to\mathrm{Tot}(K^t)$ that multiplies the component in $K^{p,q}$ by $(-1)^{pq}$ is an isomorphism of
complexes: for $c\in K^{p,q}$ both ways around the square send $[c]$ to $(-1)^{pq}[d_2c]+(-1)^{pq+q}[d_1c]$, with
$[d_2c]$ in $(K^t)^{q+1,p}$ and $[d_1c]$ in $(K^t)^{q,p+1}$. It is the identity on $K^{0,n}=(K^t)^{n,0}$, so it carries
the map induced by the $\epsilon^q$ to the map induced by $\eta^p=\epsilon^p:A^p\to(K^t)^{p,0}$. Since composing with an
isomorphism does not change being an isomorphism up to $s^{n+2}$, it suffices to prove the following, for $K^t$ in
place of $K$. Let $\eta^p:A^p\to K^{p,0}$ be maps with $d_2\eta^p=0$ and $d_1\eta^p=\eta^{p+1}d_A$, such that every
column $K^{p,\bullet}$ is a resolution up to $s$ of $A^p$ by $\eta^p$, and let $\eta:A\to\mathrm{Tot}(K)$ be the map
$a\mapsto[\eta^p(a)]$ on $A^p$; it commutes with the differentials. Then $\eta$ induces isomorphisms up to $s^{n+2}$ on
$H^n$.

\emph{A staircase.} Let $N\geq0$ and $x\in\mathrm{Tot}^N(K)$ with $(Dx)_{p,q}=0$ for all $p$ and all $q\geq1$. We show
that there is $y\in\mathrm{Tot}^{N-1}(K)$, with $y=0$ if $N=0$, such that $s^Nx-Dy=[z]$ for some $z\in K^{N,0}$. By
induction on $k\leq N$ we find $y_k$ with $(s^kx-Dy_k)_{p,q}=0$ for $q>N-k$; take $y_0=0$. Given $y_k$ with $k<N$, put
$x'=s^kx-Dy_k$ and $r=N-k-1$. As $x'_{k-1,r+2}=0$, the component of $Dx'=s^kDx$ in $K^{k,r+2}$ is
$(-1)^kd_2x'_{k,r+1}$, which vanishes; so $d_2x'_{k,r+1}=0$, and as the column $K^{k,\bullet}$ is a resolution up to
$s$, $sx'_{k,r+1}=d_2w$ for some $w\in K^{k,r}$. Put $y_{k+1}=sy_k+(-1)^k[w]$. Then
$s^{k+1}x-Dy_{k+1}=sx'-(-1)^k[d_1w]-[sx'_{k,r+1}]$. The only component of $x'$ in the rows $q\geq r+1$ is
$x'_{k,r+1}$, since $x'$ has degree $N$, and $[d_1w]$ lies in row $r$; so $y_{k+1}$ does for $k+1$. For $k=N$ the only
component of $s^Nx-Dy_N$ is in $K^{N,0}$.

\emph{Surjectivity.} Let $x\in\mathrm{Tot}^n(K)$ with $Dx=0$, and let $y$ and $z$ be given by the staircase. Then
$D[z]=0$, so $d_1z=0$ and $d_2z=0$. As $\eta^n$ and $d_2$ are exact up to $s$, $\eta^n(a)=sz$ for some $a\in A^n$. Then
$\eta^{n+1}(d_Aa)=d_1\eta^n(a)=sd_1z=0$, so $sd_Aa=0$ because $\eta^{n+1}$ is injective up to $s$. So $sa$ is a cocycle,
and $\eta(sa)=s^2[z]=s^{n+2}x-D(s^2y)$: the class of $sa$ maps to $s^{n+2}$ times the class of $x$.

\emph{Injectivity.} Let $a\in A^n$ with $d_Aa=0$ and $\eta(a)=DY$ for some $Y\in\mathrm{Tot}^{n-1}(K)$, with $Y=0$ if
$n=0$. If $n=0$, then $\eta^0(a)=0$, so $sa=0$ and $s^2a=0$. Let $n=m+1$. Then $DY=[\eta^n(a)]$ has no components in
the rows $q\geq1$, and the staircase gives $y\in\mathrm{Tot}^{m-1}(K)$ and $w\in K^{m,0}$ with $s^mY-Dy=[w]$. So
$D[w]=s^m[\eta^n(a)]$, that is, $d_2w=0$ and $d_1w=s^m\eta^n(a)$. As $\eta^m$ and $d_2$ are exact up to $s$,
$\eta^m(b)=sw$ for some $b\in A^m$. Then $\eta^n(d_Ab-s^{m+1}a)=d_1\eta^m(b)-s^{m+1}\eta^n(a)=0$, so
$sd_Ab=s^{m+2}a$ because $\eta^n$ is injective up to $s$, and $s^{n+2}a=d_A(s^2b)$ is a coboundary.
% Lean: P.dc_row is P.dc_col for the flip (AugRow -> AugCol of K.flip) composed with the isomorphism P.dc_flip
% (GlobalTheory.T.total_complex_flip_isomorphism, the sign (-1)^(pq)) through upTo_alg; P.dc_col is dc_col_surj and
% dc_col_inj, with the staircase Auxfb.stair and stair_top.

(2) The proof is a diagram chase in the total complexes. Write $\varphi$ also for the map
$\mathrm{Tot}(K)\to\mathrm{Tot}(K')$, which commutes with $D$. Say that $x\in\mathrm{Tot}^m(K)$ \emph{starts in column}
$k$ if $x_{p,q}=0$ for $p<k$; then $x_{k,m-k}$ is its only component in column $k$. Two facts are used throughout.
\begin{enumerate}
\item[(a)] If $x\in\mathrm{Tot}^m(K)$ starts in column $k$, then $(Dx)_{k,m-k+1}=(-1)^kd_2x_{k,m-k}$, as
  $x_{k-1,m-k+1}=0$. In particular $d_2x_{k,m-k}=0$ if $Dx$ starts in column $k+1$.
\item[(b)] If $c\in K'^{k,q}$ is $d_2v$ for some $v\in K'^{k,q-1}$, the coboundary $w=D((-1)^k[v])=(-1)^k[d_1v]+[c]$
  has the same components as $[c]$ in the columns $p\leq k$. If $q=0$ and $c=0$, the same holds for $w=0$.
\end{enumerate}

\emph{Surjectivity.} Let $y\in\mathrm{Tot}^n(K')$ with $Dy=0$. We show by induction on $k$, for $0\leq k\leq n+1$, that
there are $a\in\mathrm{Tot}^n(K)$ and a coboundary $w\in\mathrm{Tot}^n(K')$ such that $s^{2k}y-\varphi(a)-w$ and $Da$
start in column $k$. For $k=0$ take $a=0$ and $w=0$. Let $k\leq n$, let $a$ and $w$ be given for $k$, put $q=n-k$ and
$r=s^{2k}y-\varphi(a)-w$. Then $Dr=-\varphi(Da)$, as $Dy=0=Dw$.

First, let $b=(Da)_{k,q+1}$. By (a) for $Da$, which has $D(Da)=0$, $d_2b=0$. By (a) for $r$, $\varphi(b)=-(Dr)_{k,q+1}=
-(-1)^kd_2r_{k,q}$. As $H^{q+1}(K^{k,\bullet})\to H^{q+1}(K'^{k,\bullet})$ is injective up to $s$, $sb=d_2c$ for some
$c\in K^{k,q}$. Put $a_1=sa-(-1)^k[c]$. Then $Da_1=sDa-(-1)^k[d_1c]-[sb]$ starts in column $k+1$, as $b$ is the only
component of $Da$ in column $k$.

Second, $r_1=s^{2k+1}y-\varphi(a_1)-sw=sr+(-1)^k[\varphi(c)]$ starts in column $k$, and $Dr_1=-\varphi(Da_1)$ starts in
column $k+1$, so $d_2(r_1)_{k,q}=0$ by (a). As $H^q(K^{k,\bullet})\to H^q(K'^{k,\bullet})$ is surjective up to $s$,
there is $c'\in K^{k,q}$ with $d_2c'=0$ such that $\varphi(c')-s(r_1)_{k,q}$ is $d_2v$ for some $v\in K'^{k,q-1}$, or
is $0$ if $q=0$. Let $w_0$ be the coboundary of (b) for this element, and put $a_2=sa_1+[c']$ and $w_2=s^2w-w_0$. Then
$s^{2k+2}y-\varphi(a_2)-w_2=sr_1-[\varphi(c')]+w_0$ starts in column $k$, and its component in column $k$ is
$s(r_1)_{k,q}-\varphi(c')+(\varphi(c')-s(r_1)_{k,q})=0$. And $Da_2=sDa_1+[d_1c']$ starts in column $k+1$. So $a_2$ and
$w_2$ do for $k+1$.

For $k=n+1$, $s^{2n+2}y-\varphi(a)-w$ lies in degree $n$ and starts in column $n+1$, so it is $0$. And $Da$ lies in
degree $n+1$ and starts in column $n+1$, so $Da=[b]$ with $b\in K^{n+1,0}$. By (a), $d_2b=0$, and $\varphi(b)$ is the
only component of $\varphi(Da)=D(s^{2n+2}y-w)=0$. As $H^0(K^{n+1,\bullet})$, the kernel of $d_2$ on $K^{n+1,0}$, maps
injectively up to $s$ to $H^0(K'^{n+1,\bullet})$, $sb=0$. So $sa$ is a cocycle, and $\varphi(sa)=s^{2n+3}y-sw$: the class
of $sa$ maps to $s^{2n+3}$ times the class of $y$.

\emph{Injectivity.} Let $x\in\mathrm{Tot}^n(K)$ with $Dx=0$ and $\varphi(x)=Dz$, where $z\in\mathrm{Tot}^{n-1}(K')$, and
$z=0$ if $n=0$. We show by induction on $k$, for $0\leq k\leq n$, that there are a coboundary $g\in\mathrm{Tot}^n(K)$ and
$t\in\mathrm{Tot}^{n-1}(K')$, with $t=0$ if $n=0$, such that $s^{2k}x-g$ starts in column $k$, $t$ starts in column
$k$, and $\varphi(s^{2k}x-g)=Dt$. For $k=0$ take $g=0$ and $t=z$. Let $k<n$, let $g$ and $t$ be given for $k$, put
$j=n-1-k$ and $e=s^{2k}x-g$, a cocycle.

First, by (a) $d_2e_{k,j+1}=0$, and $\varphi(e_{k,j+1})=(Dt)_{k,j+1}=(-1)^kd_2t_{k,j}$. As $H^{j+1}(K^{k,\bullet})\to
H^{j+1}(K'^{k,\bullet})$ is injective up to $s$, $se_{k,j+1}=d_2c$ for some $c\in K^{k,j}$. Put
$g_1=sg+D((-1)^k[c])$ and $t_1=st-(-1)^k[\varphi(c)]$. Then $s^{2k+1}x-g_1=se-(-1)^k[d_1c]-[se_{k,j+1}]$ starts in
column $k+1$, $t_1$ starts in column $k$, and $\varphi(s^{2k+1}x-g_1)=Dt_1$.

Second, $Dt_1=\varphi(s^{2k+1}x-g_1)$ starts in column $k+1$, so $d_2(t_1)_{k,j}=0$ by (a). As
$H^j(K^{k,\bullet})\to H^j(K'^{k,\bullet})$ is surjective up to $s$, there is $c'\in K^{k,j}$ with $d_2c'=0$ such that
$\varphi(c')-s(t_1)_{k,j}$ is $d_2v$ for some $v\in K'^{k,j-1}$, or is $0$ if $j=0$. Let $w_0$ be the coboundary of (b)
for this element, and put $g_2=sg_1+D[c']$ and $t_2=st_1-[\varphi(c')]+w_0$. Then
$s^{2k+2}x-g_2=s(s^{2k+1}x-g_1)-[d_1c']$ starts in column $k+1$, and
$\varphi(s^{2k+2}x-g_2)=sDt_1-D[\varphi(c')]=Dt_2$, as $Dw_0=0$. And $t_2$ starts in column $k$, with component
$s(t_1)_{k,j}-\varphi(c')+(\varphi(c')-s(t_1)_{k,j})=0$ in column $k$. So $g_2$ and $t_2$ do for $k+1$.

For $k=n$, $t$ lies in degree $n-1$ and starts in column $n$, so $t=0$ and $\varphi(s^{2n}x-g)=0$. And $s^{2n}x-g$ lies
in degree $n$ and starts in column $n$, so $s^{2n}x-g=[e]$ with $e\in K^{n,0}$, $d_2e=0$ by (a), and $\varphi(e)=0$. As
$H^0(K^{n,\bullet})\to H^0(K'^{n,\bullet})$ is injective up to $s$, $se=0$. So $s^{2n+1}x=sg$, and
$s^{2n+3}x=s^2(sg)$ is a coboundary.

The surjectivity used the hypothesis for $p+q=n$ and $p+q=n+1$, and the injectivity for $p+q=n$ and $p+q=n-1$.
% Lean: dc_map_col_surj (surj_main, surj_stage, surj_step) and dc_map_col_inj (inj_main, inj_stage, inj_step,
% inj_last); (a) is D_at with sign_eq_zero, (b) is colB, starting in column k is Sup, and [c] is sgl. Each step
% multiplies by s twice, once for each of the two column hypotheses, which gives s^(2n+2) after n+1 steps and
% s^(2n+3) with the last step.
\end{proof}

\begin{remark}\label{gal-realization-thickening-zariski-double-filtration}
If the hypothesis of Lemma~\ref{gal-realization-thickening-zariski-double}(2) holds with $s=1$ for all $p$ and $q$, its
conclusion for all $n$ also follows from the filtration by columns. For $p\geq0$ let $F^p\subset\mathrm{Tot}(K)$ be the subcomplex
$\bigoplus_{p'\geq p}K^{p',\bullet}$, and likewise $F'^p\subset\mathrm{Tot}(K')$. Then $F^p/F^{p+1}$ is the column
$K^{p,\bullet}$ shifted by $p$. By induction on $p$, the map $\mathrm{Tot}(K)/F^p\to\mathrm{Tot}(K')/F'^p$ induces
isomorphisms on cohomology: for $p=0$ both complexes are $0$, and the step is the five lemma for the long exact
sequences of $0\to F^p/F^{p+1}\to\mathrm{Tot}(K)/F^{p+1}\to\mathrm{Tot}(K)/F^p\to0$ and of its analogue for $K'$. As
$F^{n+2}$ and $F'^{n+2}$ vanish in degrees at most $n+1$, $H^n(\mathrm{Tot}(K))=H^n(\mathrm{Tot}(K)/F^{n+2})$, and
likewise for $K'$.
\end{remark}

\begin{lemma}\label{gal-realization-thickening-zariski-fibre}
Let $\Lambda$ be a ring, $G:\mathcal C'\to\mathcal C$ a functor, $F$ and $F'$ functors from $\mathcal C$ and $\mathcal C'$
to $\Lambda$-modules, and $\tau_c:F(G(c))\to F'(c)$ maps natural in $c\in\mathcal C'$. For an object $T$ of $\mathcal C$
let $T/G$ be the category of the pairs $(c,u)$ of an object $c$ of $\mathcal C'$ and a map $u:T\to G(c)$, whose maps
$(c,u)\to(c',u')$ are the maps $v:c\to c'$ with $G(v)u=u'$; write $F'$ also for the functor $(c,u)\mapsto F'(c)$ on it.
Let $\varepsilon_T:F(T)\to C^0(T/G,F')$ send $x$ to the cochain $(c,u)\mapsto\tau_c(F(u)(x))$. Suppose that for every
object $T$ of $\mathcal C$ the sequence
\[
 0\to F(T)\xrightarrow{\varepsilon_T}C^0(T/G,F')\xrightarrow{d}C^1(T/G,F')\xrightarrow{d}\cdots
\]
of cochains (Definition~\ref{gal-realization-split-diagrams}(1)) is exact. Then the maps $H^j(\mathcal C,F)\to
H^j(\mathcal C',F')$ induced by $G$ and $\tau$ are bijective for every $j$.
\end{lemma}
\leanref{Development/WeilBound/CoreStatements/PadicHodgeComparison/EtaleCohomologyComparison/ZariskiEtaleLifting/FibInit.lean}{CrystallineSemistableComparison.B.P.fib}
\leanref{Development/WeilBound/CoreStatements/PadicHodgeComparison/EtaleCohomologyComparison/CechPairings/TerminalRes.lean}{CrystallineSemistableComparison.P.rel}

\begin{proof}
For an object $T$ of $\mathcal C$ write $T/\mathcal C$ for the category $T/G$ formed for the identity functor of
$\mathcal C$, with the functor $F$. A map $a:T\to T'$ gives the functor $T'/G\to T/G$, $(c,u)\mapsto(c,ua)$, and so maps
$C^q(T/G,F')\to C^q(T'/G,F')$ that commute with $d$. This makes $K'^q:T\mapsto C^q(T/G,F')$ a functor on $\mathcal C$,
and likewise $K^q:T\mapsto C^q(T/\mathcal C,F)$. Let $K^{p,q}=C^p(\mathcal C,K^q)$ and $K'^{p,q}=C^p(\mathcal C,K'^q)$,
double complexes with the differentials of the cochains of $\mathcal C$ and of the categories $T/\mathcal C$ and $T/G$.
The functors $T/G\to T/\mathcal C$, $(c,u)\mapsto(G(c),u)$, with the maps $\tau_c$, give maps $K^q(T)\to K'^q(T)$,
natural in $T$ and compatible with $d$, and so a map of double complexes $K\to K'$.

\emph{The rows.} Pulling back along $(c,u)\mapsto c$ gives maps $C^q(\mathcal C',F')\to C^0(\mathcal C,K'^q)$, which
commute with the differentials. An element of $K'^{p,q}$ assigns to a chain $T_0\to\cdots\to T_p$ of $\mathcal C$ and a
chain $(c_0,u_0)\to\cdots\to(c_q,u_q)$ of $T_p/G$ an element of $F'(c_q)$. Such a chain of $T_p/G$ is a chain
$c_0\to\cdots\to c_q$ of $\mathcal C'$ with a map $u_0:T_p\to G(c_0)$, which determines the other $u_i$. So $K'^{p,q}$ is
the product, over the chains $c_0\to\cdots\to c_q$ of $\mathcal C'$, of the $p$-cochains of the category $\mathcal
C/G(c_0)$ of objects over $G(c_0)$ with values in the constant functor $F'(c_q)$, and $d_1$ is their differential. The
category $\mathcal C/G(c_0)$ has the terminal object $(G(c_0),\mathrm{id})$. For a category $\mathcal D$ with a terminal
object $Z$ and a module $N$, the functor from $\mathcal D$ to the category with one object and one map, and the functor
back to $Z$, induce inverse maps between the cohomology of the constant functors $N$, by
Lemma~\ref{gal-realization-period-prism} for the maps $c\to Z$; and by
Lemma~\ref{gal-realization-thickening-torsor-adjoint}(2) the cohomology on the category with one object is $N$ in degree
$0$, given by the constant cochains, and $0$ in the other degrees. So every row $0\to C^q(\mathcal C',F')\to
K'^{0,q}\to K'^{1,q}\to\cdots$ is exact, and likewise for $K$ and $C^q(\mathcal C,F)$. By
Lemma~\ref{gal-realization-thickening-zariski-double}(1) with $s=1$, these maps induce isomorphisms $H^n(\mathcal C,F)\to
H^n(\mathrm{Tot}(K))$ and $H^n(\mathcal C',F')\to H^n(\mathrm{Tot}(K'))$. They carry the map induced by $G$ and $\tau$ to
the map induced by $K\to K'$: both send an $n$-cochain $f$ of $\mathcal C$ to the element of $K'^{0,n}$ whose value at
$T$ and $(c_0,u_0)\to\cdots\to(c_n,u_n)$ is $\tau_{c_n}(f(G(c_0)\to\cdots\to G(c_n)))$.

\emph{The columns.} The column $K^{p,\bullet}$ is the product, over the chains $T_0\to\cdots\to T_p$ of $\mathcal C$, of
the complexes $C^\bullet(T_p/\mathcal C,F)$, and likewise $K'^{p,\bullet}$ with $C^\bullet(T_p/G,F')$, and $K\to K'$ is
the product of the maps of complexes defined above. The category $T_p/\mathcal C$ has the initial object
$(T_p,\mathrm{id})$, so by Lemma~\ref{gal-realization-thickening-torsor-adjoint}(2) the cohomology of
$C^\bullet(T_p/\mathcal C,F)$ is $F(T_p)$ in degree $0$, embedded by $x\mapsto((c,u)\mapsto F(u)(x))$, and $0$ in the
other degrees. The composite of this embedding with the map to $C^0(T_p/G,F')$ is $\varepsilon_{T_p}$, so by the
hypothesis the maps $H^q(T_p/\mathcal C,F)\to H^q(T_p/G,F')$ are bijective for every $q$. So $K\to K'$ induces
isomorphisms on the cohomology of the columns, and by Lemma~\ref{gal-realization-thickening-zariski-double}(2) with $s=1$ on the
cohomology of the total complexes. With the rows, this gives the claim.
\end{proof}
\leanref{Development/WeilBound/CoreStatements/PadicHodgeComparison/EtaleCohomologyComparison/CechPairings/TerminalRes.lean}{CrystallineSemistableComparison.RelativeAugmentedComplex.relCx}
\leanref{Development/WeilBound/CoreStatements/PadicHodgeComparison/EtaleCohomologyComparison/CechPairings/TerminalRes.lean}{CrystallineSemistableComparison.RelativeAugmentedComplex.relAugRow}
\leanref{Development/WeilBound/CoreStatements/PadicHodgeComparison/EtaleCohomologyComparison/CechPairings/TerminalRes.lean}{CrystallineSemistableComparison.RelativeAugmentedComplex.relHom}
\leanref{Development/WeilBound/CoreStatements/PadicHodgeComparison/EtaleCohomologyComparison/CechPairings/TerminalRes.lean}{CrystallineSemistableComparison.RelativeAugmentedComplex.cmp}
\leanref{Development/WeilBound/CoreStatements/PadicHodgeComparison/EtaleCohomologyComparison/CechPairings/TerminalRes.lean}{CrystallineSemistableComparison.RelativeAugmentedComplex.terminal_res}
\leanref{Development/WeilBound/CoreStatements/PadicHodgeComparison/EtaleCohomologyComparison/CechPairings/TerminalRes.lean}{CrystallineSemistableComparison.RelativeAugmentedComplex.row_res}
\leanref{Development/WeilBound/CoreStatements/PadicHodgeComparison/EtaleCohomologyComparison/CechPairings/TerminalRes.lean}{CrystallineSemistableComparison.RelativeAugmentedComplex.col_upTo}
\leanref{Development/WeilBound/CoreStatements/PadicHodgeComparison/EtaleCohomologyComparison/ZariskiEtaleLifting/FibInit.lean}{CrystallineSemistableComparison.B.fib_init}
\leanref{Development/WeilBound/CoreStatements/PadicHodgeComparison/EtaleCohomologyComparison/ZariskiEtaleLifting/FibInit.lean}{CrystallineSemistableComparison.B.fib_aug}
\leanref{Development/WeilBound/CoreStatements/PadicHodgeComparison/EtaleCohomologyComparison/ZariskiEtaleLifting/FResLe.lean}{CrystallineSemistableComparison.B.P.res_le}

\begin{lemma}\label{gal-realization-thickening-zariski-cech}
Let $m\geq0$, and let $\mathcal N_m$ be the set of nonempty subsets of $\{1,\ldots,m\}$, ordered by inclusion and viewed
as a category. Let $R$ be a ring, $G$ a functor from $\mathcal N_m$ to $R$-modules, $M$ an $R$-module and
$\epsilon_S:M\to G(S)$, $S\in\mathcal N_m$, $R$-linear maps with $G(S\subset S')\epsilon_S=\epsilon_{S'}$. Let $\Lambda$
be a commutative ring with elements $f_1,\ldots,f_m$ that generate the unit ideal, and let ring maps from $\Lambda$ to the
rings of $R$-linear endomorphisms of $M$ and of every $G(S)$ be given, for which the maps $\epsilon_S$ and $G(S\subset
S')$ are $\Lambda$-linear. Write $f_S=\prod_{k\in S}f_k$. Suppose that every $f_k$ with $k\in S$ acts bijectively on
$G(S)$, that for every $y\in G(S)$ some $f_S^ny$ lies in $\epsilon_S(M)$, and that every $x\in M$ with $\epsilon_S(x)=0$
has $f_S^nx=0$ for some $n\geq0$. Then the sequence
\[
 0\to M\xrightarrow{\epsilon}C^0(\mathcal N_m,G)\xrightarrow{d}C^1(\mathcal N_m,G)\xrightarrow{d}\cdots,
\]
with $\epsilon(x)=(S\mapsto\epsilon_S(x))$, is exact.
\end{lemma}
\leanref{Development/WeilBound/CoreStatements/PadicHodgeComparison/EtaleCohomologyComparison/ZariskiEtaleLifting/FCechLoc.lean}{CrystallineSemistableComparison.B.P.cech_loc}
\leanref{Development/WeilBound/CoreStatements/PadicHodgeComparison/EtaleCohomologyComparison/ZariskiEtaleLifting/FCechLoc.lean}{CrystallineSemistableComparison.B.ClD}
\leanref{Development/WeilBound/CoreStatements/PadicHodgeComparison/EtaleCohomologyComparison/ZariskiEtaleLifting/FCechLoc.lean}{CrystallineSemistableComparison.B.ClD.res}

\begin{proof}
The ring $\Lambda$ acts on cochains value by value, and this action commutes with $d$ and with $\epsilon$, as the maps
$G(S\subset S')$ and $\epsilon_S$ are $\Lambda$-linear; and $d\epsilon=0$, as $G(S\subset S')\epsilon_S=\epsilon_{S'}$.
Let $V$ be a $\Lambda$-stable submodule of $M$ or of a module of cochains, and $z$ an element such that for every $k$
some $f_k^nz$ lies in $V$. Then $z\in V$: the $a\in\Lambda$ with $az\in V$ form an ideal that contains a power of each
$f_k$, and these powers generate the unit ideal, as the $f_k$ do.

Fix $k$, and for $S\in\mathcal N_m$ let $\alpha_S=G(S\subset S\cup\{k\})$. Write $f_kf_S=cf_{S\cup\{k\}}$, with $c=f_k$
if $k\in S$ and $c=1$ otherwise. (a) For $y\in G(S\cup\{k\})$ there are $n\geq0$ and $x\in G(S)$ with
$f_k^ny=\alpha_S(x)$. Indeed, choose $n$ and $x_0\in M$ with $f_{S\cup\{k\}}^ny=\epsilon_{S\cup\{k\}}(x_0)$, and, as
$f_S^n$ acts bijectively on $G(S)$, $x\in G(S)$ with $f_S^nx=\epsilon_S(c^nx_0)$. Then
$f_S^n\alpha_S(x)=\epsilon_{S\cup\{k\}}(c^nx_0)=c^nf_{S\cup\{k\}}^ny=f_S^nf_k^ny$, and $f_S^n$ acts bijectively on
$G(S\cup\{k\})$, so $\alpha_S(x)=f_k^ny$. (b) If $x\in G(S)$ has $\alpha_S(x)=0$, then $f_k^Nx=0$ for some $N$. Indeed,
choose $n$ and $x_0\in M$ with $f_S^nx=\epsilon_S(x_0)$. Then $\epsilon_{S\cup\{k\}}(x_0)=\alpha_S(f_S^nx)=0$, so
$f_{S\cup\{k\}}^Nx_0=0$ for some $N$, and $f_S^{N+n}f_k^Nx=c^Nf_{S\cup\{k\}}^N\epsilon_S(x_0)=0$; as $f_S$ acts bijectively
on $G(S)$, $f_k^Nx=0$. If (a) holds for $n$, it holds for every larger exponent, and likewise for (b). As there are
finitely many chains of each length in $\mathcal N_m$, (a) and (b) hold with one exponent for all the values of a
cochain.

Let $\mathcal K_k\subset\mathcal N_m$ be the full subcategory of the sets that contain $k$, with initial object $\{k\}$,
and $G_k$ the functor $S\mapsto G(S\cup\{k\})$ on $\mathcal N_m$. The functor $S\mapsto S\cup\{k\}$ from $\mathcal N_m$ to
$\mathcal K_k$ is left adjoint to the inclusion, as $S\cup\{k\}\subset S'$ if and only if $S\subset S'$, for $S'\ni k$. By
Lemma~\ref{gal-realization-thickening-torsor-adjoint}(1) and (2), the complex $C^\bullet(\mathcal N_m,G_k)$ has no
cohomology in degrees at least $1$, and its $0$-cocycles are the cochains $S\mapsto G(\{k\}\subset S\cup\{k\})(y)$ with
$y\in G(\{k\})$. The maps $\alpha_S$ form a natural transformation $G\to G_k$, which gives a map of complexes $\alpha$
from $C^\bullet(\mathcal N_m,G)$ to $C^\bullet(\mathcal N_m,G_k)$ that commutes with $\Lambda$.

Let $z$ be a cocycle of degree $j+1\geq1$. Then $\alpha(z)$ is a cocycle, so $\alpha(z)=dy$ for a cochain $y$ of
$C^\bullet(\mathcal N_m,G_k)$ of degree $j$. By (a), $f_k^{N_1}y=\alpha(y'')$ for some $N_1$ and some cochain $y''$ of
$G$. Then $\alpha(f_k^{N_1}z-dy'')=f_k^{N_1}dy-d\alpha(y'')=0$, so by (b) $f_k^{N_2}(f_k^{N_1}z-dy'')=0$ for some $N_2$,
and $f_k^{N_1+N_2}z=d(f_k^{N_2}y'')$. As this holds for every $k$, $z$ is a coboundary.

Let $z$ be a $0$-cocycle. Then $\alpha(z)$ is a $0$-cocycle of $G_k$, so $\alpha(z)(S)=G(\{k\}\subset S\cup\{k\})(y)$
for some $y\in G(\{k\})$, and $f_k^ny=\epsilon_{\{k\}}(x)$ for some $n$ and $x\in M$. Then $\alpha(f_k^nz-\epsilon(x))=0$,
so by (b) $f_k^N(f_k^nz-\epsilon(x))=0$ for some $N$, and $f_k^{N+n}z=\epsilon(f_k^Nx)$. As this holds for every $k$,
$z\in\epsilon(M)$. Finally, if $\epsilon(x)=0$, then $\epsilon_{\{k\}}(x)=0$, so $f_k^nx=0$ for some $n$, for every $k$,
and $x=0$.
\end{proof}
\leanref{Development/WeilBound/CoreStatements/PadicHodgeComparison/EtaleCohomologyComparison/ZariskiEtaleLifting/FCechLoc.lean}{CrystallineSemistableComparison.B.cl_mem_of_pow}
\leanref{Development/WeilBound/CoreStatements/PadicHodgeComparison/EtaleCohomologyComparison/ZariskiEtaleLifting/FCechLoc.lean}{CrystallineSemistableComparison.B.ClD.e1}
\leanref{Development/WeilBound/CoreStatements/PadicHodgeComparison/EtaleCohomologyComparison/ZariskiEtaleLifting/FCechLoc.lean}{CrystallineSemistableComparison.B.ClD.e2}
\leanref{Development/WeilBound/CoreStatements/PadicHodgeComparison/EtaleCohomologyComparison/ZariskiEtaleLifting/FCechLoc.lean}{CrystallineSemistableComparison.B.cl_gc}
\leanref{Development/WeilBound/CoreStatements/PadicHodgeComparison/EtaleCohomologyComparison/ZariskiEtaleLifting/FCechLoc.lean}{CrystallineSemistableComparison.B.cl_res}
\leanref{Development/WeilBound/CoreStatements/PadicHodgeComparison/EtaleCohomologyComparison/ZariskiEtaleLifting/FCechLoc.lean}{CrystallineSemistableComparison.B.ClD.hi}
\leanref{Development/WeilBound/CoreStatements/PadicHodgeComparison/EtaleCohomologyComparison/ZariskiEtaleLifting/FCechLoc.lean}{CrystallineSemistableComparison.B.ClD.h0}

\begin{lemma}\label{gal-realization-thickening-zariski-refinement}
Let $\mathcal C$ be a category, $R$ a ring and $F$ a functor from $\mathcal C$ to $R$-modules; maps are composed as
functions. Let $\varphi:T\to T'$ be a map of $\mathcal C$, $m\geq0$, and $M_1,\ldots,M_m$ classes of maps with source $T$
such that $\chi\psi\in M_k$ whenever $\psi\in M_k$ and $\chi$ is a map out of the target of $\psi$. Let maps
$\lambda_S:T'\to L_S$, for $S\subset\{1,\ldots,m\}$, and a ring $\Lambda$ with elements $f_1,\ldots,f_m$ and actions on
$F(T')$ and the $F(L_S)$ satisfy (2)(a) and (2)(b) of Definition~\ref{gal-realization-thickening-zariski-class} for
$\varphi$. Let $\Theta:\mathcal E\to\mathcal C$ be a faithful functor and $\psi_e:T'\to\Theta(e)$, $e\in\mathcal E$, maps
with $\Theta(v)\psi_e=\psi_{e'}$ for every map $v:e\to e'$. Suppose that for every $e$, $\psi_e\varphi\in M_k$ for some
$k$; that every map $\chi:\Theta(e)\to\Theta(e')$ with $\chi\psi_e=\psi_{e'}$ is $\Theta(v)$ for a map $v:e\to e'$; and
that for every nonempty $S$ there is an object $e_S$ with $\Theta(e_S)=L_S$ and $\psi_{e_S}=\lambda_S$. Then the sequence
$0\to F(T')\to C^0(\mathcal E,F\circ\Theta)\xrightarrow{d}C^1(\mathcal E,F\circ\Theta)\xrightarrow{d}\cdots$, whose first
map sends $x$ to the cochain $e\mapsto F(\psi_e)(x)$, is exact.
\end{lemma}
\leanref{Development/WeilBound/CoreStatements/PadicHodgeComparison/EtaleCohomologyComparison/ZariskiEtaleLifting/ZariskiCocycleRefinement.lean}{CrystallineSemistableComparison.B.zariski_cocycle_refinement}

\begin{proof}
For $e\in\mathcal E$ let $V(e)$ be the set of the $k$ with $\psi_e\varphi\in M_k$, which is nonempty. For a map $v:e\to
e'$, $\psi_{e'}\varphi=\Theta(v)\psi_e\varphi$, so $V(e)\subset V(e')$, and $V$ is a functor $\mathcal E\to\mathcal N_m$
(Lemma~\ref{gal-realization-thickening-zariski-cech}). Let $S$ be nonempty and $e\in\mathcal E$. There is at most one
map $e_S\to e$: for two maps $v$ and $v'$, $\Theta(v)\lambda_S=\psi_e=\Theta(v')\lambda_S$ and $\psi_e\varphi=
\Theta(v)\lambda_S\varphi\in M_k$ for $k\in S$, so $\Theta(v)=\Theta(v')$ by the uniqueness in (2)(a), and $v=v'$ as
$\Theta$ is faithful. There is one if and only if $S\subset V(e)$: a map $v$ gives $\psi_e\varphi=\Theta(v)\lambda_S
\varphi\in M_k$ for $k\in S$; conversely, if $S\subset V(e)$, (2)(a) gives $\chi:L_S\to\Theta(e)$ with
$\chi\lambda_S=\psi_e$, and $\chi=\Theta(v)$ for a map $v:e_S\to e$ by hypothesis. As $S'\subset V(e_{S'})$ for
nonempty $S'$, there is exactly one map $e_S\to e_{S'}$ when $S\subset S'$. So $\Xi:S\mapsto e_S$ is a functor $\mathcal
N_m\to\mathcal E$, left adjoint to $V$, and by Lemma~\ref{gal-realization-thickening-torsor-adjoint}(1) restriction
along $\Xi$, from $H^j(\mathcal E,F\circ\Theta)$ to $H^j(\mathcal N_m,F\circ\Theta\circ\Xi)$, is bijective for every $j$.

The functor $F\circ\Theta\circ\Xi$ sends $S$ to $F(L_S)$, and $S\subset S'$ to $F(t)$ for the map $t=\Theta(e_S\to
e_{S'})$, which has $t\lambda_S=\lambda_{S'}$. With $M=F(T')$, $\epsilon_S=F(\lambda_S)$ and the actions of (2)(b), the
hypotheses of Lemma~\ref{gal-realization-thickening-zariski-cech} hold, so the sequence $0\to F(T')\to C^0(\mathcal
N_m,F\circ\Theta\circ\Xi)\to\cdots$ is exact. Its first map is the first map of the statement followed by restriction
along $\Xi$, as $\psi_{e_S}=\lambda_S$. The first map of the statement followed by $d$ is $0$, as
$F(\Theta(v))F(\psi_e)=F(\psi_{e'})$. As restriction along $\Xi$ is bijective on $H^0$, the module of $0$-cocycles, and
on every $H^j$ with $j\geq1$, the first map of the statement is injective with image the $0$-cocycles, and the
cohomology of $C^\bullet(\mathcal E,F\circ\Theta)$ in degrees at least $1$ vanishes.
\end{proof}
\leanref{Development/WeilBound/CoreStatements/PadicHodgeComparison/EtaleCohomologyComparison/ZariskiEtaleLifting/ZariskiCocycleRefinement.lean}{CrystallineSemistableComparison.B.zar_d_aug}
\leanref{Development/WeilBound/CoreStatements/PadicHodgeComparison/EtaleCohomologyComparison/CechPairings/MemBoundariesZeroIff.lean}{CrystallineSemistableComparison.P.res_transfer}

\begin{lemma}\label{gal-realization-thickening-zariski-descent}
Let $\mathcal C$ be a category, $R$ a ring, $F$ a functor from $\mathcal C$ to $R$-modules, $\mathcal P$ a Zariski class
for $F$ (Definition~\ref{gal-realization-thickening-zariski-class}), and $h:\mathcal C'\to\mathcal C$ a fully faithful
functor such that $h(T')\in\mathcal P$ for every object $T'$ of $\mathcal C'$ and every object of $\mathcal P$ is isomorphic
to $h(T')$ for some $T'$. For every $j\geq0$, restriction along $h$, $H^j(\mathcal C,F)\to H^j(\mathcal C',F\circ h)$, is
bijective (Definition~\ref{gal-realization-split-diagrams}(1)). In particular this holds for the inclusion of the full
subcategory $\mathcal C_{\mathcal P}$ of $\mathcal C$ on $\mathcal P$.
\end{lemma}
\leanref{Development/WeilBound/CoreStatements/PadicHodgeComparison/EtaleCohomologyComparison/ZariskiEtaleLifting/ZariskiCocycleRefinement.lean}{CrystallineSemistableComparison.B.P.zar}
\leanref{Development/WeilBound/CoreStatements/PadicHodgeComparison/EtaleCohomologyComparison/ZariskiEtaleLifting/XSmulSucc.lean}{CrystallineSemistableComparison.B.Onto}

\begin{proof}
Maps are composed as functions. Let $\iota:\mathcal C_{\mathcal P}\to\mathcal C$ be the inclusion. For an object $T$ of
$\mathcal C$, the category $T/\iota$ of Lemma~\ref{gal-realization-thickening-zariski-fibre} has as objects the pairs
$(T'',u)$ of $T''\in\mathcal P$ and a map $u:T\to T''$; write $F$ also for the functor $(T'',u)\mapsto F(T'')$ on it. By
Lemma~\ref{gal-realization-thickening-zariski-fibre}, for $G=\iota$, $F'=F\circ\iota$ and the identity maps, it suffices
to show that for every $T$ the sequence $0\to F(T)\to C^0(T/\iota,F)\to C^1(T/\iota,F)\to\cdots$, whose first map sends
$x$ to $(T'',u)\mapsto F(u)(x)$, is exact.

Fix $T$, and let $m$ and $M_1,\ldots,M_m$ be as in (2) of Definition~\ref{gal-realization-thickening-zariski-class}. Let
$\mathcal Q_T$ be the full subcategory of $T/\iota$ of the pairs $(T'',u)$ with $u\in M_k$ for some $k$, and
$\kappa:\mathcal Q_T\to T/\iota$ the inclusion.

\emph{Step 1: restriction along $\kappa$ is bijective on $H^j(-,F)$ for every $j$.} By
Lemma~\ref{gal-realization-thickening-zariski-fibre} for $\kappa$, it suffices that for every object $x=(T',\varphi)$ of
$T/\iota$ the sequence $0\to F(T')\to C^0(x/\kappa,F)\to\cdots$ of that lemma is exact. The objects of $x/\kappa$ are the
pairs $(y,v)$ of an object $y=(T'',u)$ of $\mathcal Q_T$ and a map $v:T'\to T''$ with $v\varphi=u$, a map $(y,v)\to
(y',v')$ is a map $w:T''\to T'''$ with $wv=v'$, and the first map of the sequence sends $z$ to $(y,v)\mapsto F(v)(z)$.
Apply Lemma~\ref{gal-realization-thickening-zariski-refinement} to $\varphi:T\to T'$ with the data of (2)(a) and (2)(b)
of Definition~\ref{gal-realization-thickening-zariski-class} for $\varphi$, $\mathcal E=x/\kappa$, $\Theta(y,v)=T''$ and
$\psi_{(y,v)}=v$. The functor $\Theta$ is faithful, and $\psi_{(y,v)}\varphi=u\in M_k$ for some $k$. A map $\chi:T''\to
T'''$ with $\chi v=v'$ has $\chi u=\chi v\varphi=v'\varphi=u'$, so it is a map $y\to y'$ of $\mathcal Q_T$ and a map
$(y,v)\to(y',v')$. For nonempty $S$, $\lambda_S\varphi\in M_k$ for $k\in S$, so $L_S\in\mathcal P$, and
$((L_S,\lambda_S\varphi),\lambda_S)$ is an object $e_S$ of $x/\kappa$ with $\Theta(e_S)=L_S$ and $\psi_{e_S}=\lambda_S$.
So the sequence is exact.

\emph{Step 2.} The sequence $0\to F(T)\to C^0(\mathcal Q_T,F)\to C^1(\mathcal Q_T,F)\to\cdots$, whose first map sends $x$
to $(T'',u)\mapsto F(u)(x)$, is exact, by Lemma~\ref{gal-realization-thickening-zariski-refinement} for
$\mathrm{id}_T$, with the data of (2)(a) and (2)(b) of Definition~\ref{gal-realization-thickening-zariski-class} for
$\mathrm{id}_T$, $\mathcal E=\mathcal Q_T$, $\Theta(T'',u)=T''$ and $\psi_{(T'',u)}=u$. Its hypotheses hold: $u\in M_k$
for some $k$ by the definition of $\mathcal Q_T$; a map $\chi:T''\to T'''$ with $\chi u=u'$ is a map
$(T'',u)\to(T''',u')$; and for nonempty $S$, $\lambda_S\in M_k$ for $k\in S$, so $L_S\in\mathcal P$ and $(L_S,\lambda_S)$
is an object of $\mathcal Q_T$.

\emph{Step 3.} The first map of the sequence for $T/\iota$, followed by restriction along $\kappa$, is the first map of
the sequence of Step 2, and the first map for $T/\iota$ followed by $d$ is $0$. By Step 1 restriction along $\kappa$ is
bijective on $H^0$, the module of $0$-cocycles, and on every $H^j$ with $j\geq1$. So, as at the end of the proof of
Lemma~\ref{gal-realization-thickening-zariski-refinement}, the sequence for $T/\iota$ is exact. This proves the lemma for
$\iota$.

\emph{Step 4: the functor $h$.} Let $L:\mathcal C'\to\mathcal C_{\mathcal P}$ send $T'$ to $h(T')$ and a map to its image
under $h$, so that $h=\iota L$. The functor $L$ is fully faithful, as $h$ is and $\mathcal C_{\mathcal P}$ is a full
subcategory, and every object of $\mathcal C_{\mathcal P}$ is isomorphic to some $L(T')$; so $L$ is an equivalence and has a
right adjoint, its quasi-inverse. Restriction along $h$ is restriction along $\iota$ followed by restriction along $L$, as
a composite of two induced maps is induced by the composite functor. The first is bijective by Steps 1 to 3, and the
second by Lemma~\ref{gal-realization-thickening-torsor-adjoint}(1).
% Lean: B.zar_main. key T is Step 3, through P.res_transfer: Step 1 is hup, B.P.fib for Mor.restrict FK incl, with z2
% (zariski_cocycle_refinement for StructuredArrow x incl), and Step 2 is hK' (zariski_cocycle_refinement for P'.FullSubcategory);
% the conclusion is B.P.fib for Mor.restrict F ιQ. Step 4 is B.P.zar: L is ObjectProperty.lift with IsEquivalence from
% essential surjectivity, ri_calc (the equivalence case, through P.adj for G.asEquivalence.toAdjunction) for L, its
% composition rule for ιQ and L, and its isomorphism rule for liftCompιIso. The condition (1) of a Zariski class (hQm) is
% not used.
\end{proof}
\leanref{Development/WeilBound/CoreStatements/PadicHodgeComparison/EtaleCohomologyComparison/ZariskiEtaleLifting/ZariskiCocycleRefinement.lean}{CrystallineSemistableComparison.B.zar_main}
\leanref{Development/WeilBound/CoreStatements/PadicHodgeComparison/EtaleCohomologyComparison/ZariskiEtaleLifting/FRiCalc.lean}{CrystallineSemistableComparison.B.P.ri_calc}

\begin{lemma}\label{gal-realization-thickening-zariski-localization}
Let $(R,I,\gamma,p)$, $\mathcal O$, $X$ and $\kappa$ be as in Definition~\ref{gal-realization-thickening-nilpotent}, let
$T=(B,i,J,\delta,y)$ be a thickening of $X$ over $(R,I,\gamma,p)$, $g\in B$, and $\lambda:B\to B_g=B[1/g]$ the
localization map.
\begin{enumerate}
\item There are divided powers $\delta_g$ on $JB_g$, and only one such, for which $\lambda$ is a divided power map. With
  them, $T_g=(B_g,\lambda\circ i,JB_g,\delta_g,y_g)$ is a thickening, where $y_g$ is the map $\operatorname{Spec}
  B_g/JB_g\to\operatorname{Spec}B/J$ induced by $\lambda$ followed by $y$, and $\lambda$ is a map $T\to T_g$. If $B$ is
  flat over $R$, so is $B_g$.
\item For $r\geq0$, the ideal $(JB_g)^{[r]}$ formed for $\delta_g$ is $\lambda(J^{[r]})B_g$.
\item For every map $\psi:T\to T''$ of thickenings such that $\psi(g)$ is a unit of $B_{T''}$, there is exactly one map
  $\chi:T_g\to T''$ with $\chi\circ\lambda=\psi$.
\end{enumerate}
\end{lemma}
\leanref{Development/WeilBound/CoreStatements/PadicHodgeComparison/EtaleCohomologyComparison/PerfectoidPullbacks/PolynomialSiteTraces/AdmissibleMap.lean}{CrystallineSemistableComparison.Pb.Site.locT}
\leanref{Development/WeilBound/CoreStatements/PadicHodgeComparison/EtaleCohomologyComparison/PerfectoidPullbacks/PolynomialSiteTraces/AdmissibleMap.lean}{CrystallineSemistableComparison.Pb.Site.locH}
\leanref{Development/WeilBound/CoreStatements/PadicHodgeComparison/EtaleCohomologyComparison/PerfectoidPullbacks/PolynomialSiteTraces/AdmissibleMap.lean}{CrystallineSemistableComparison.Pb.Site.locD_dp}
\leanref{Development/AutomorphicSpectrum/Proofs/BlockDecomposition/CohomologyRestrictionBetweenStableFamilies.lean}{GlobalTheory.CompanionTheory.ZarCosieve.pdPow_locT}
\leanref{Development/AutomorphicSpectrum/Proofs/BlockDecomposition/CohomologyRestrictionBetweenStableFamilies.lean}{GlobalTheory.CompanionTheory.ZarCosieve.locT_univ}
\leanref{Development/AutomorphicSpectrum/Proofs/BlockDecomposition/CohomologyRestrictionBetweenStableFamilies.lean}{GlobalTheory.CompanionTheory.ZarCosieve.Q_locT}

\begin{proof}
(1) The ring $B_g$ is flat over $B$ \cite[Tag 00HT]{Sta}, so $\delta$ extends to $B_g$, and uniquely \cite[Tag
07H1]{Sta}: there is exactly one divided power structure $\delta_g$ on $JB_g$ for which $\lambda$ is a divided power
map. If
$p\neq0$, $p$ is nilpotent in $B$ and so in $B_g$; if $p=0$, $J$ is nilpotent and so is $JB_g$. The map $\lambda\circ i$
is a divided power map from $(R,I,\gamma)$, as a composite of two. The ring map $B/J\to B_g/JB_g$ induced by $\lambda$
commutes with the maps from $R/I$, so $y_g$ is an $\mathcal O$-morphism and $T_g$ is a thickening; $\lambda$ commutes
with the maps from $R$, carries $J$ into $JB_g$ with its divided powers, and induces the point of $T_g$ from that of $T$.
If $B$ is flat over $R$, then $B_g$, flat over $B$, is flat over $R$ \cite[Tag 00HC]{Sta}.

(2) This is the argument of Lemma~\ref{gal-realization-thickening-chart}(2), with $\lambda$ in place of $i_D$. As $\lambda$
is a divided power map, $\lambda(J^{[r]})\subset(JB_g)^{[r]}$. The set of $x\in JB_g$ with $\delta_{g,n}(x)\in
\lambda(J^{[n]})B_g$ for all $n$ contains $\lambda(J)$ and is closed under sums and under multiplication by $B_g$, by the
formulas $\delta_n(x+x')=\sum_{a+b=n}\delta_a(x)\delta_b(x')$ and $\delta_n(cx)=c^n\delta_n(x)$ and $J^{[a]}J^{[b]}\subset
J^{[a+b]}$; so it is $JB_g$. A generator $\prod_j\delta_{g,k_j}(x_j)$ of $(JB_g)^{[r]}$, with $\sum_jk_j\geq r$, therefore
lies in $\lambda(J^{[\sum_jk_j]})B_g\subset\lambda(J^{[r]})B_g$.

(3) By the universal property of localization there is exactly one ring map $\chi:B_g\to B_{T''}$ with
$\chi\circ\lambda=\psi$. It commutes with the maps from $R$ and carries $JB_g$ into $J_{T''}$, as $\psi(J)\subset
J_{T''}$. It is a divided power map: the set of $x\in JB_g$ with $\chi(\delta_{g,n}(x))=\delta_{T'',n}(\chi(x))$ for
all $n$ contains $\lambda(J)$, as $\lambda$ and $\psi$ are divided power maps, and is closed under sums and under
multiplication by $B_g$ by the same two formulas. The point of $T''$ is the point of $T$ composed with the map
$\operatorname{Spec}B_{T''}/J_{T''}\to\operatorname{Spec}B/J$ induced by $\psi$, which factors through
$\operatorname{Spec}B_g/JB_g$; so $\chi$ is a map $T_g\to T''$. A map of thickenings is determined by its ring map, so
$\chi$ is the only one.
\end{proof}
\leanref{Development/WeilBound/CoreStatements/AbelianHeightsDuality/DividedPowerAlgebras/ExistsDividedPowersAug.lean}{FubiniStudyChowHeightFormula₂.Adr.flatExt}
\leanref{Development/AutomorphicSpectrum/Proofs/BlockDecomposition/ExtendedDividedPowers.lean}{GlobalTheory.CompanionTheory.PDMap.extended_dividedPowers}
\leanref{Development/WeilBound/CoreStatements/PadicHodgeComparison/EtaleCohomologyComparison/PerfectoidPullbacks/PolynomialSiteTraces/AdmissibleMap.lean}{CrystallineSemistableComparison.Pb.Site.admissible_map}

\begin{lemma}\label{gal-realization-thickening-zariski}
Let $(R,I,\gamma,p)$, $\mathcal O$, $X$ and $\kappa$ be as in Definition~\ref{gal-realization-thickening-nilpotent}. Let
$\mathcal Q_0$ be $\mathrm{Th}(X/R)$ or $\mathrm{Th}^{\mathrm{fl}}(X/R)$ if $p\neq0$, and $\mathrm{Th}^{(s)}_\ell(X/R)$ for
integers $s,\ell\geq1$ (Lemma~\ref{gal-realization-kappa-truncated-periods-exist-base-change-zariski}) if $p=0$. Let
$\mathcal V$ be a set of
open subsets of $X$ with union $V_0$, and let $F$ be $\mathcal O$ or $\mathfrak a/\mathfrak b$ for two of the families
$J^{[r]}+cB$ and $cB$ of Definition~\ref{gal-realization-thickening-nilpotent}(3).
\begin{enumerate}
\item If $\mathcal W$ is a set of open subsets of $V_0$ with $(\mathcal Q_0)_{\mathcal V}\subset(\mathcal Q_0)_{\mathcal
  W}$, then restriction $H^j((\mathcal Q_0)_{\mathcal W},F)\to H^j((\mathcal Q_0)_{\mathcal V},F)$ is bijective for every
  $j$.
\item Let $\mathcal Q'$ be a full subcategory with $(\mathcal Q_0)_{\mathcal V}\subset\mathcal Q'\subset(\mathcal
  Q_0)_{V_0}$ such that $T'\in\mathcal Q'$ whenever $T\in\mathcal Q'$, $T'\in\mathcal Q_0$ and there is a map $T\to T'$.
  Then restriction $H^j(\mathcal Q',F)\to H^j((\mathcal Q_0)_{\mathcal V},F)$ is bijective for every $j$.
\end{enumerate}
\end{lemma}
\leanref{Development/WeilBound/CoreStatements/PadicHodgeComparison/EtaleCohomologyComparison/PerfectoidPullbacks/ChartGrConst.lean}{CrystallineSemistableComparison.Pb.zar}
\leanref{Development/AutomorphicSpectrum/Proofs/BlockDecomposition/CechRefinementCohomologyBijective.lean}{GlobalTheory.CompanionTheory.NlZar.zariski_descent}
\leanref{Development/AutomorphicSpectrum/Proofs/BlockDecomposition/CohomologyRestrictionBetweenStableFamilies.lean}{GlobalTheory.CompanionTheory.ZarCosieve.main}

\begin{proof}
(1) The category $\mathcal Q'=(\mathcal Q_0)_{\mathcal W}$ satisfies the hypotheses of (2). It contains $(\mathcal
Q_0)_{\mathcal V}$ by assumption. A thickening that lies over some $W\in\mathcal W$ lies over $V_0\supset W$, so
$\mathcal Q'\subset(\mathcal Q_0)_{V_0}$. If $T\in\mathcal Q'$ lies over $W\in\mathcal W$, $T'\in\mathcal Q_0$ and there is
a map $T\to T'$, then $T'$ lies over $W$ by Lemma~\ref{gal-realization-thickening-over}(1), so $T'\in\mathcal Q'$. Thus
(1) is (2) for this $\mathcal Q'$.
% Lean: NlZar.zariski_descent runs the argument of (2) below for Q' = St Q₀ 𝒲: zarCov_of is zarCov, with the closure of
% St Q₀ 𝒲 under localization (used_locN, ThickPt.inU_map) in place of hc, and locMod_sq, locMod_o and locI_std are
% locSq, locO and locI_std of ZarCosieve; both conclude with B.P.zar.

(2) Let $\mathcal Q_1=(\mathcal Q_0)_{\mathcal V}$, a class of objects of $\mathcal Q'$. Its full subcategory is
$(\mathcal Q_0)_{\mathcal V}$, so by Lemma~\ref{gal-realization-thickening-zariski-descent}, for the category $\mathcal
Q'$ and the functor $F$, it suffices to show that $\mathcal Q_1$ is a Zariski class for $F$
(Definition~\ref{gal-realization-thickening-zariski-class}). Maps are composed as functions.

Condition (1). Let $T\to T'$ be a map of $\mathcal Q'$ with $T\in\mathcal Q_1$. Then $T'\in\mathcal Q'\subset\mathcal
Q_0$, and $T'$ lies over the open $V\in\mathcal V$ that $T$ lies over, by Lemma~\ref{gal-realization-thickening-over}(1).
So $T'\in\mathcal Q_1$.

Condition (2). Let $T=(B,i,J,\delta,y)$ be in $\mathcal Q'$. As $T$ lies over $V_0=\bigcup\mathcal V$, the inverse images
$y^{-1}(V)$, $V\in\mathcal V$, cover $\operatorname{Spec}B/J$. Basic opens form a basis of its topology, so every point
lies in some $D(s)\subset y^{-1}(V)$ with $s\in B/J$ and $V\in\mathcal V$. These $s$ lie in no prime ideal, so they
generate the unit ideal, and finitely many of them do: there are $s_1,\ldots,s_m\in B/J$ generating the unit ideal and
$V_1,\ldots,V_m\in\mathcal V$ with $D(s_k)\subset y^{-1}(V_k)$. Choose $f_k\in B$ lifting $s_k$. The $f_k$ generate the
unit ideal of $B$: some $\sum_ka_kf_k$ lies in $1+J$, and it is a unit, as the elements of $J$ are nilpotent (proof of
Lemma~\ref{gal-realization-thickening-chart}(1) if $p\neq0$; if $p=0$, $J$ is nilpotent). Let $M_k$ be the class of maps
$\psi:T\to T_1$ of $\mathcal Q'$ such
that $\psi(f_k)$ is a unit of $B_{T_1}$. A ring map carries units to units, so $\chi\psi\in M_k$ for $\psi\in M_k$. For
$\psi\in M_k$, the ring map $\bar\psi:B/J\to B_{T_1}/J_{T_1}$ sends $s_k$ to a unit, so $\operatorname{Spec}\bar\psi$
maps every prime ideal into $D(s_k)$, and the point $y\circ\operatorname{Spec}\bar\psi$ of $T_1$ factors through $V_k$.
As $T_1\in\mathcal Q'\subset\mathcal Q_0$, $T_1\in\mathcal Q_1$.

Let $\varphi:T\to T'$ be a map of $\mathcal Q'$. For $S\subset\{1,\ldots,m\}$ let $g_S=\prod_{k\in S}\varphi(f_k)\in
B_{T'}$, $L_S=T'_{g_S}$ and $\lambda_S:T'\to L_S$ the map of
Lemma~\ref{gal-realization-thickening-zariski-localization}(1). Then $L_S\in\mathcal Q_0$, as $T'\in\mathcal Q_0$: if
$\mathcal Q_0$ is $\mathrm{Th}(X/R)$ or $\mathrm{Th}^{\mathrm{fl}}(X/R)$, by that lemma; if $\mathcal
Q_0=\mathrm{Th}^{(s)}_\ell(X/R)$, by the proof of
Lemma~\ref{gal-realization-kappa-truncated-periods-exist-base-change-zariski}, which uses
Lemma~\ref{gal-realization-thickening-zariski-localization}(2) for $J^{[s]}$ and checks the condition on $\ell$. By the
closure condition on $\mathcal Q'$, applied to $\lambda_S$, $L_S\in\mathcal Q'$. For $k\in S$, $\lambda_S(\varphi(f_k))$
divides the unit
$\lambda_S(g_S)$, so $\lambda_S\varphi\in M_k$. If $\psi:T'\to T''$ is a map of $\mathcal Q'$ with $\psi\varphi\in M_k$
for every $k\in S$, then $\psi(g_S)$ is a unit, and by Lemma~\ref{gal-realization-thickening-zariski-localization}(3)
there is exactly one map $\chi:L_S\to T''$ with $\chi\lambda_S=\psi$; it is a map of $\mathcal Q'$, which is full. This
is (2)(a) of Definition~\ref{gal-realization-thickening-zariski-class}.

For (2)(b) take $\Lambda=B_{T'}$ and the elements $\varphi(f_1),\ldots,\varphi(f_m)$, which generate the unit ideal as the
$f_k$ do. For every thickening $T_2$, $F(T_2)$ is a $B_{T_2}$-module: for $F=\mathcal O$ it is $B_{T_2}$, and for
$F=\mathfrak a/\mathfrak b$ it is the image of $\mathfrak a_{T_2}$ in $B_{T_2}/\mathfrak b_{T_2}$, a submodule. For a map
$t$, $F(t)(cx)=t(c)F(t)(x)$. Let $B_{T'}$ act on $F(T')$ by this structure and on $F(L_S)$ through $\lambda_S$. Then
$F(\lambda_S)$ is $B_{T'}$-linear, and so is $F(t)$ for every map $t:L_S\to L_{S'}$ with $t\lambda_S=\lambda_{S'}$, as
$t(\lambda_S(c))=\lambda_{S'}(c)$. For $k\in S$, $\varphi(f_k)$ acts on $F(L_S)$ as multiplication by the unit
$\lambda_S(\varphi(f_k))$, so bijectively. For $F=\mathcal O$, every element of $B_{L_S}=(B_{T'})_{g_S}$ is
$\lambda_S(x)/g_S^n$ for some $x\in B_{T'}$ and $n$, so $g_S^n$ times it is $\lambda_S(x)$; and $\lambda_S(x)=0$ implies
$g_S^nx=0$ for some $n$. For $F=\mathfrak a/\mathfrak b$, $\mathfrak a_{L_S}=\lambda_S(\mathfrak a_{T'})B_{L_S}$ and
$\mathfrak b_{L_S}=\lambda_S(\mathfrak b_{T'})B_{L_S}$: for $J^{[r]}$ with $r\geq0$ by
Lemma~\ref{gal-realization-thickening-zariski-localization}(2), for $J^{[r]}=B$ with $r\leq0$ and for $cB$ directly, and
the extension of a sum of ideals is the sum of their extensions. So an element of $F(L_S)$ is the class of
$\lambda_S(a)/g_S^n$ for some $a\in\mathfrak a_{T'}$ and $n$, and $g_S^n$ times it is the image of the class of $a$
under $F(\lambda_S)$. If the class of $a\in\mathfrak a_{T'}$ maps to $0$, then $\lambda_S(a)\in\lambda_S(\mathfrak
b_{T'})B_{L_S}$, so $g_S^na\in\mathfrak b_{T'}$ for some $n$, and $g_S^n$ kills the class of $a$. So $\mathcal Q_1$ is a
Zariski class for $F$.
\end{proof}
\leanref{Development/AutomorphicSpectrum/Proofs/BlockDecomposition/CohomologyRestrictionBetweenStableFamilies.lean}{GlobalTheory.CompanionTheory.ZarCosieve.zarCov}
\leanref{Development/AutomorphicSpectrum/Proofs/BlockDecomposition/CohomologyRestrictionBetweenStableFamilies.lean}{GlobalTheory.CompanionTheory.ZarCosieve.exists_cover}
\leanref{Development/AutomorphicSpectrum/Proofs/BlockDecomposition/CohomologyRestrictionBetweenStableFamilies.lean}{GlobalTheory.CompanionTheory.ZarCosieve.span_lift}
\leanref{Development/AutomorphicSpectrum/Proofs/BlockDecomposition/CohomologyRestrictionBetweenStableFamilies.lean}{GlobalTheory.CompanionTheory.ZarCosieve.mkLoc}
\leanref{Development/AutomorphicSpectrum/Proofs/BlockDecomposition/CohomologyRestrictionBetweenStableFamilies.lean}{GlobalTheory.CompanionTheory.ZarCosieve.locSq}
\leanref{Development/AutomorphicSpectrum/Proofs/BlockDecomposition/CohomologyRestrictionBetweenStableFamilies.lean}{GlobalTheory.CompanionTheory.ZarCosieve.locO}
\leanref{Development/AutomorphicSpectrum/Proofs/BlockDecomposition/CohomologyRestrictionBetweenStableFamilies.lean}{GlobalTheory.CompanionTheory.ZarCosieve.locI_std}
\leanref{Development/AutomorphicSpectrum/Proofs/BlockDecomposition/CechRefinementCohomologyBijective.lean}{GlobalTheory.CompanionTheory.NlZar.zarCov_of}

\begin{lemma}\label{gal-realization-thickening-base-change}
Let $\sigma:R\to R'$, $(R,I,\gamma)$, $(R',I',\gamma')$, $\mathcal O$, $X$, $\kappa$ and $\kappa'$ be as in
Definition~\ref{gal-realization-thickening-nilpotent}(4), with $R'$ flat over $R$, let $p\neq0$ be an integer nilpotent in
$R$, and suppose that $X$ has a coordinate cover. Let $\Gamma$ be the restriction of scalars along $\sigma$. For $j\geq0$,
the map $R'\otimes_RH^j(\mathrm{Th}(X/R),\mathcal O)\to H^j(\mathrm{Th}(X/R'),\mathcal O)$, $a\otimes x\mapsto
a\,\Gamma^*(x)$, is bijective.
\end{lemma}
\leanref{Development/AutomorphicSpectrum/Proofs/BlockDecomposition/NilpotentIdentityMorphismCompatStatement.lean}{GlobalTheory.T.nilpotent_identity_morphism_compat}

\begin{proof}
Let $(U_\alpha)_{\alpha\in A}$ be a coordinate cover and $\mathcal U=\{U_\alpha:\alpha\in A\}$. Every thickening lies over
$X=\bigcup_\alpha U_\alpha$ (Lemma~\ref{gal-realization-thickening-over}(2)), so $\mathrm{Th}(X/R)=\mathrm{Th}(X/R)_X$,
and by Lemma~\ref{gal-realization-thickening-zariski}(1), with $\mathcal V=\mathcal U$ and $\mathcal W=\{X\}$,
restriction $H^j(\mathrm{Th}(X/R),\mathcal O)\to H^j(\mathrm{Th}(X/R)_{\mathcal U},\mathcal O)$ is bijective; likewise
over $R'$. These restrictions commute with $\Gamma^*$, as both composites are induced by the same functor with the
identity maps of the rings, and with multiplication by elements of $R'$. So it suffices to prove the claim for
$\mathrm{Th}(X/R)_{\mathcal U}$ and $\mathrm{Th}(X/R')_{\mathcal U}$.

By Lemma~\ref{gal-realization-thickening-over}, the thickenings lying over $U_\alpha$ form cosieves of these categories
that cover them, $\Gamma$ maps the one over $R'$ into the one over $R$, and for nonempty $S\subset A$ their intersections
are $\mathrm{Th}(X/R)_{U_S}$ and $\mathrm{Th}(X/R')_{U_S}$, with $U_S=\bigcap_{\alpha\in S}U_\alpha$. The affine open
$U_S$ has coordinates, restricted from $U_\alpha$ for any $\alpha\in S$
(Lemma~\ref{gal-realization-thickening-cover}(1)), so by Lemma~\ref{gal-realization-thickening-local-base-change} the
claim holds for $\mathrm{Th}(X/R)_{U_S}$ and $\mathrm{Th}(X/R')_{U_S}$. Lemma~\ref{gal-realization-thickening-cosieves}(1),
with $\Lambda=R$, $\Lambda'=R'$, one module $N=R'$, $F=F'=\mathcal O$ and $\tau(a)$ the multiplication by $a$, gives the
claim for $\mathrm{Th}(X/R)_{\mathcal U}$ and $\mathrm{Th}(X/R')_{\mathcal U}$.
\end{proof}
\leanref{Development/AutomorphicSpectrum/Proofs/BlockDecomposition/NilpotentIdentityMorphismCompat.lean}{GlobalTheory.CompanionTheory.BcCov.bc_cover}
\leanref{Development/AutomorphicSpectrum/Proofs/BlockDecomposition/NilpotentIdentityMorphismCompat.lean}{GlobalTheory.CompanionTheory.BcCov.st_obj}
\leanref{Development/AutomorphicSpectrum/Proofs/BlockDecomposition/NilpotentIdentityMorphismCompat.lean}{GlobalTheory.CompanionTheory.BcCov.pull_square}
\leanref{Development/AutomorphicSpectrum/Proofs/BlockDecomposition/NilpotentIdentityMorphismCompat.lean}{GlobalTheory.CompanionTheory.BcCov.tens_transport}
\leanref{Development/AutomorphicSpectrum/Proofs/BlockDecomposition/StableCoverCechVanishingStatement.lean}{GlobalTheory.T.stable_cover_cech_vanishing}

\subsection{Bockstein sequences and flat thickenings}\label{subsection-gal-flat-thickenings}

\begin{lemma}\label{gal-realization-frobenius-base-change}
Let $K$, $\mathcal O_K$ and $W$ be as in Definition~\ref{gal-realization-frobenius-teichmuller}, $Y$ a smooth projective
$\mathcal O_K$-scheme, $j\geq0$ and $\nu\geq1$, and $b_\nu$ as there. The $A_{\mathrm{cris}}$-linear map
$A_{\mathrm{cris}}\otimes_WH^j_\nu(Y/W)\to H^{\mathrm T,j}_\nu(Y)$, $a\otimes y\mapsto a\,b_\nu(y)$, is bijective.
\end{lemma}
\leanref{Development/WeilBound/CoreStatements/PadicHodgeComparison/EtaleCohomologyComparison/PerfectoidPullbacks/PeriodBaseChange/BcTens.lean}{CrystallineSemistableComparison.Pb.bc_tens}

\begin{proof}
As $H^j_\nu(Y/W)$ is a $W_\nu$-module and $A_{\mathrm{cris}}/\ell^\nu=A_\nu$, the map in question is the map
$A_\nu\otimes_{W_\nu}H^j_\nu(Y/W)\to H^{\mathrm T,j}_\nu(Y)$, $a\otimes y\mapsto a\,b_\nu(y)$. The ring $A_\nu$ is flat
over $W_\nu$: $A_{\mathrm{cris}}(\mathcal O_C^\flat)\subset W(\mathcal O_C^\flat)[1/\ell]$
(Definition~\ref{gal-crystalline-test-algebras}) has no $\ell$-torsion, so it is flat over the discrete valuation ring
$W$, and $A_\nu$ is its reduction modulo $\ell^\nu$.

By Definition~\ref{gal-realization-thickening-nilpotent}, $\mathrm{Th}_\nu(Y/W)=\mathrm{Th}(Y/W_\nu)$ and
$\mathrm{Th}^{\mathrm T}_\nu(Y)=\mathrm{Th}(Y/A_\nu)$, and $b_\nu$ is induced by $\Gamma_W$, the restriction of scalars
along $W_\nu\to A_\nu$. This map carries $\ell W_\nu$ into $I'_\nu$ and commutes with the divided powers
(Definition~\ref{gal-realization-frobenius-teichmuller}), and $\mathcal O_K\to A_\nu/I'_\nu$ is $\mathcal O_K\to
W_\nu/\ell$ followed by the map it induces. As $Y$ is smooth and projective over $\mathcal O_K$, it is separated,
quasi-compact and smooth, so it has a coordinate cover by Lemma~\ref{gal-realization-thickening-cover}(2). As $\ell\neq0$
is nilpotent in $W_\nu$, Lemma~\ref{gal-realization-thickening-base-change}, for $W_\nu\to A_\nu$ and $p=\ell$, shows that
$a\otimes y\mapsto a\,b_\nu(y)$ is bijective.
% Lean: the first paragraph is tens_change with cris_surjective and cris_ker, and the flatness is
% GlobalTheory.T.flat_acris_truncated_witt (CrisFlat.flat_acr, from flat_acris_zero). The second is
% nilpotent_identity_morphism_compat for BcAcris.nm r M n (u and g identities, hI = witt_to_acris_isDPMorphism) and the
% cover cov_of_sm; toQ_pull_bijective and sqC_toQ identify the cohomology of sq (oT Q) (oB Q) with that of O.
\end{proof}
\leanref{Development/AutomorphicSpectrum/Proofs/BlockDecomposition/SmoothModelTensorWittComparison.lean}{GlobalTheory.CompanionTheory.BcAcris.bc_tens}
\leanref{Development/AutomorphicSpectrum/Proofs/BlockDecomposition/CrisZeroSurjective.lean}{GlobalTheory.CompanionTheory.BcAcris.tens_change}
\leanref{Development/AutomorphicSpectrum/Proofs/BlockDecomposition/FlatAcrisTruncatedWittStatement.lean}{GlobalTheory.T.flat_acris_truncated_witt}
\leanref{Development/AutomorphicSpectrum/Proofs/BlockDecomposition/FlatAcrisTruncatedWitt.lean}{GlobalTheory.CompanionTheory.CrisFlat.flat_acris_zero}
\leanref{Development/AutomorphicSpectrum/Proofs/BlockDecomposition/NilpotentIdentityMorphismCompatStatement.lean}{GlobalTheory.T.nilpotent_identity_morphism_compat}
\leanref{Development/AutomorphicSpectrum/Proofs/BlockDecomposition/CrisZeroSurjective.lean}{GlobalTheory.CompanionTheory.BcAcris.nm}
\leanref{Development/AutomorphicSpectrum/Proofs/BlockDecomposition/WittToAcrisIsDPMorphismStatement.lean}{GlobalTheory.T.witt_to_acris_isDPMorphism}
\leanref{Development/AutomorphicSpectrum/Proofs/BlockDecomposition/CrisZeroSurjective.lean}{GlobalTheory.CompanionTheory.BcAcris.cov_of_sm}
\leanref{Development/AutomorphicSpectrum/Proofs/BlockDecomposition/CrisZeroSurjective.lean}{GlobalTheory.CompanionTheory.BcAcris.toQ_pull_bijective}
\leanref{Development/AutomorphicSpectrum/Proofs/BlockDecomposition/CrisZeroSurjective.lean}{GlobalTheory.CompanionTheory.BcAcris.sqC_toQ}

\begin{definition}\label{gal-realization-frobenius-bockstein}
Let $W$ be a ring, $\ell\geq0$ an integer, and $(S_n)_{n\geq0}$ and $(S'_n)_{n\geq0}$ inverse systems of $W$-modules,
with transition maps $r_n:S_{n+1}\to S_n$ and $r'_n:S'_{n+1}\to S'_n$ and limits $\varprojlim_nS_n$ and
$\varprojlim_nS'_n$, the modules of families $(x_n)$ with $r_n(x_{n+1})=x_n$, resp.\ $r'_n(x_{n+1})=x_n$.
\emph{Bockstein sequences} for $(S,S')$ consist of $W$-linear maps $\rho_{c,b}:S_c\to S_b$, $\nu_{a,c}:S_a\to S_c$ and
$\delta_{b,a}:S_b\to S'_a$ for all $a,b,c\geq0$, such that:
\begin{enumerate}
\item $\ell^nS_n=0$ for every $n$;
\item $\rho_{c,c}$ is the identity, and $\rho_{c+1,b}=\rho_{c,b}\circ r_c$ for $b\leq c$;
\item $\nu_{a,c}\circ\rho_{c,a}$ is multiplication by $\ell^{c-a}$, and $r_c\circ\nu_{a+1,c+1}=\nu_{a,c}\circ r_a$, for
  $a\leq c$;
\item $r'_a\circ\delta_{b,a+1}=\delta_{b,a}$;
\item if $c=a+b$, the kernel of $\rho_{c,b}$ lies in the image of $\nu_{a,c}$, and the kernel of $\delta_{b,a}$ lies in
  the image of $\rho_{c,b}$.
\end{enumerate}
\end{definition}
\leanref{Development/WeilBound/CoreStatements/PadicHodgeComparison/EtaleCohomologyComparison/PerfectoidPullbacks/TwistedRepresentations/Tens.lean}{CrystallineSemistableComparison.Tw.LSeq}

\begin{lemma}\label{gal-realization-frobenius-artinian-lift}
Let $W$ be a ring and $(S_n)_{n\geq0}$ an inverse system of Artinian $W$-modules with transition maps $r_n$. Let
$N_n\subset S_n$ be submodules with $r_n(N_{n+1})\subset N_n$, and $y_n\in S_n$ elements with $r_n(y_{n+1})-y_n\in N_n$,
for every $n$. Then there is $z\in\varprojlim_nS_n$ with $z_n-y_n\in N_n$ for every $n$.
\end{lemma}
\leanref{Development/AutomorphicSpectrum/Proofs/BlockDecomposition/ArtinianLimitLiftOfCompatibleStatement.lean}{GlobalTheory.T.artinian_limit_lift_of_compatible}
\leanref{Development/AutomorphicSpectrum/Proofs/BlockDecomposition/ArtinianLimitLiftOfCompatible.lean}{GlobalTheory.CompanionTheory.ArtinTower.exists_compatible}

\begin{proof}
\emph{Step 1.} Let $M$ be an Artinian $W$-module, $K_0\supset K_1\supset\cdots$ submodules and $c_k\in M$ with
$c_{k+1}-c_k\in K_k$. Then there is $x\in M$ with $x-c_k\in K_k$ for every $k$. Indeed, $c_{k+j}-c_k\in K_k$ for all
$j$, as the $K_i$ decrease, and the chain stabilizes: $K_k=K_m$ for $k\geq m$. Put $x=c_m$. For $k\leq m$,
$x-c_k\in K_k$, and for $k\geq m$, $x-c_k=-(c_k-c_m)\in K_m=K_k$.

\emph{Step 2.} For $k,n\geq0$ let $N^{(k)}_n$ and $y^{(k)}_n$ be the images of $N_{n+k}$ and of $y_{n+k}$ in $S_n$ under
the composite of the transition maps. Then $N^{(k+1)}_n\subset N^{(k)}_n$ and $y^{(k+1)}_n-y^{(k)}_n\in N^{(k)}_n$, the
images of $r_{n+k}(N_{n+k+1})\subset N_{n+k}$ and of $r_{n+k}(y_{n+k+1})-y_{n+k}\in N_{n+k}$. Let $\Sigma_n$ be the set
of $x\in S_n$ with $x-y^{(k)}_n\in N^{(k)}_n$ for every $k$. By Step 1, $\Sigma_0$ is not empty.

\emph{Step 3.} Every $x\in\Sigma_n$ is $r_n(x')$ for some $x'\in\Sigma_{n+1}$. For every $k$,
$x-y^{(k+1)}_n\in N^{(k+1)}_n=r_n(N^{(k)}_{n+1})$ and $y^{(k+1)}_n=r_n(y^{(k)}_{n+1})$, so there is $u_k\in S_{n+1}$ with
$r_n(u_k)=x$ and $u_k-y^{(k)}_{n+1}\in N^{(k)}_{n+1}$. Then $u_{k+1}-u_k$ lies in $K_k=\ker r_n\cap N^{(k)}_{n+1}$, by
Step 2, and the $K_k$ decrease. Step 1 gives $x'$ with $x'-u_k\in K_k$ for every $k$. Then $r_n(x')=r_n(u_0)=x$, and
$x'-y^{(k)}_{n+1}=(x'-u_k)+(u_k-y^{(k)}_{n+1})\in N^{(k)}_{n+1}$.

\emph{Step 4.} Choose $z_0\in\Sigma_0$ and, by Step 3, $z_{n+1}\in\Sigma_{n+1}$ with $r_n(z_{n+1})=z_n$. Then
$z=(z_n)$ lies in the limit, and $z_n-y_n\in N^{(0)}_n=N_n$.
\end{proof}
\leanref{Development/AutomorphicSpectrum/Proofs/BlockDecomposition/ArtinianLimitLiftOfCompatible.lean}{GlobalTheory.CompanionTheory.ArtinTower.coset_chain}
\leanref{Development/AutomorphicSpectrum/Proofs/BlockDecomposition/ArtinianLimitLiftOfCompatible.lean}{GlobalTheory.CompanionTheory.ArtinTower.stable_nonempty}
\leanref{Development/AutomorphicSpectrum/Proofs/BlockDecomposition/ArtinianLimitLiftOfCompatible.lean}{GlobalTheory.CompanionTheory.ArtinTower.stable_surj}

\begin{lemma}\label{gal-realization-frobenius-bockstein-limits}
Let $W$ be a ring, $\ell\geq0$ an integer, and $(S_n)$ and $(S'_n)$ inverse systems of $W$-modules with Bockstein
sequences (Definition~\ref{gal-realization-frobenius-bockstein}). Put $L=\varprojlim_nS_n$ and $L'=\varprojlim_nS'_n$.
\begin{enumerate}
\item If $S_1$ has finite length, then every $S_n$ has finite length.
\item If $W$ is $\ell$-adically complete and separated and $S_1$ has finite length, then $L$ is a finitely generated
  $W$-module.
\item If $W$ is Noetherian, every $S_n$ is Artinian and $L'$ is finitely generated, then there is $c\geq0$ such that
  $\ell^cS_n$ lies in the image of $L\to S_n$ for every $n$.
\item If every $S_n$ is Artinian, $n\geq0$ and $x\in L$ has $x_n=0$, then $x\in\ell^nL$.
\end{enumerate}
\end{lemma}
\leanref{Development/AutomorphicSpectrum/Proofs/BlockDecomposition/FiniteLengthOfLseqFirstStatement.lean}{GlobalTheory.T.finite_length_of_lseq_first}
\leanref{Development/AutomorphicSpectrum/Proofs/BlockDecomposition/ModuleFiniteLimitOfLseqStatement.lean}{GlobalTheory.T.module_finite_limit_of_lseq}
\leanref{Development/AutomorphicSpectrum/Proofs/BlockDecomposition/BoundedCokernelLimitToLevelStatement.lean}{GlobalTheory.T.bounded_cokernel_limit_to_level}
\leanref{Development/AutomorphicSpectrum/Proofs/BlockDecomposition/LimitDivisibleOfVanishingLevelStatement.lean}{GlobalTheory.T.limit_divisible_of_vanishing_level}

\begin{proof}
(1) Induction on $n$. As $\ell^0=1$ kills $S_0$, $S_0=0$. By Definition~\ref{gal-realization-frobenius-bockstein}(5)
with $a=1$, $b=n$ and $c=n+1$, the kernel of $\rho_{n+1,n}:S_{n+1}\to S_n$ lies in the image of $\nu_{1,n+1}:S_1\to
S_{n+1}$, so it is a submodule of a quotient of $S_1$. The image of $\rho_{n+1,n}$ is a submodule of $S_n$. So $S_{n+1}$
has finite length if $S_1$ and $S_n$ do.

(4) First, $\rho_{c,b}(x_c)=x_b$ for $x\in L$ and $b\leq c$, by
Definition~\ref{gal-realization-frobenius-bockstein}(2) and induction on $c-b$. Let $x\in L$ with $x_n=0$, and $a\geq0$.
As $\rho_{n+a,n}(x_{n+a})=x_n=0$, Definition~\ref{gal-realization-frobenius-bockstein}(5) with $b=n$ and $c=n+a$ gives
$y_a\in S_a$ with $\nu_{a,n+a}(y_a)=x_{n+a}$. Let $N_a$ be the kernel of $\nu_{a,n+a}$. By
Definition~\ref{gal-realization-frobenius-bockstein}(3), $r_{n+a}\circ\nu_{a+1,n+a+1}=\nu_{a,n+a}\circ r_a$, so
$r_a(N_{a+1})\subset N_a$ and $\nu_{a,n+a}(r_a(y_{a+1})-y_a)=r_{n+a}(x_{n+a+1})-x_{n+a}=0$. Lemma~\ref{gal-realization-frobenius-artinian-lift}
gives $z\in L$ with $\nu_{a,n+a}(z_a)=\nu_{a,n+a}(y_a)=x_{n+a}$ for every $a$. By
Definition~\ref{gal-realization-frobenius-bockstein}(3), $\nu_{a,n+a}(z_a)=\nu_{a,n+a}(\rho_{n+a,a}(z_{n+a}))=
\ell^nz_{n+a}$. So $x_{n+a}=\ell^nz_{n+a}$ for every $a$, and for every $c$,
$x_c=\rho_{n+c,c}(x_{n+c})=\ell^n\rho_{n+c,c}(z_{n+c})=\ell^nz_c$. Hence $x=\ell^nz$.

(2) By (1) every $S_n$ has finite length, hence is Artinian, so (4) applies. The projection $\pi_n:L\to S_n$ kills
$\ell^nL$ by Definition~\ref{gal-realization-frobenius-bockstein}(1), so an element of every $\ell^nL$ lies in the
kernel of every $\pi_n$ and is $0$: $L$ is $\ell$-adically separated. By (4) for $n=1$, the kernel of $\pi_1$ is $\ell L$,
so $L/\ell L$ embeds into $S_1$ and has finite length. Choose $x_1,\ldots,x_r\in L$ whose images generate $L/\ell L$, and
let $f:W^r\to L$ send the standard basis to the $x_i$. For $x\in L$ there are $a_0,a_1,\ldots\in W^r$ with
$x-f(\sum_{i<k}\ell^ia_i)\in\ell^kL$ for every $k$, chosen one at a time. As $W^r$ is $\ell$-adically complete,
$a=\sum_i\ell^ia_i$ exists, and $x-f(a)\in\ell^kL$ for every $k$, so $x=f(a)$. Hence $L$ is finitely generated.

(3) As $W$ is Noetherian and $L'$ is finitely generated, the kernels of multiplication by $\ell^k$ on $L'$ form an
increasing chain that stabilizes, so there is $c$ such that $\ell^kt=0$ implies $\ell^ct=0$ for $t\in L'$. Let $x\in
S_n$. The family $\delta(x)=(\delta_{n,a}(x))_a$ lies in $L'$ by Definition~\ref{gal-realization-frobenius-bockstein}(4),
and $\ell^n\delta(x)=\delta(\ell^nx)=0$ by Definition~\ref{gal-realization-frobenius-bockstein}(1). So
$\delta_{n,a}(\ell^cx)=0$ for every $a$. It remains to show that every $x'\in S_n$ with $\delta_{n,a}(x')=0$ for all $a$
is $\pi_n(y)$ for some $y\in L$. For $m\geq n$, Definition~\ref{gal-realization-frobenius-bockstein}(5) with $b=n$,
$a=m-n$ and $c=m$ gives $u_m\in S_m$ with $\rho_{m,n}(u_m)=x'$. Put $y_m=u_m$ and $N_m=\ker\rho_{m,n}$ for $m\geq n$,
and $y_m=0$ and $N_m=S_m$ for $m<n$. As $\rho_{m+1,n}=\rho_{m,n}\circ r_m$ for $m\geq n$, $r_m(N_{m+1})\subset N_m$ and
$r_m(y_{m+1})-y_m\in N_m$ for every $m$. Lemma~\ref{gal-realization-frobenius-artinian-lift} gives $y\in L$ with
$y_m-u_m\in N_m$ for $m\geq n$. For $m=n$, $\rho_{n,n}$ is the identity, so $y_n=u_n=x'$.
\end{proof}
\leanref{Development/AutomorphicSpectrum/Proofs/BlockDecomposition/FiniteLengthOfLseqFirst.lean}{GlobalTheory.CompanionTheory.TowerLen.isFiniteLength_level}
\leanref{Development/AutomorphicSpectrum/Proofs/BlockDecomposition/LimitDivisibleOfVanishingLevel.lean}{GlobalTheory.CompanionTheory.Tower.rho_lim}
\leanref{Development/AutomorphicSpectrum/Proofs/BlockDecomposition/LimitDivisibleOfVanishingLevel.lean}{GlobalTheory.CompanionTheory.Tower.lim_eq_pow_smul}
\leanref{Development/AutomorphicSpectrum/Proofs/BlockDecomposition/ModuleFiniteLimitOfLseq.lean}{GlobalTheory.CompanionTheory.TowerLim.module_finite_lim}
\leanref{Development/AutomorphicSpectrum/Proofs/BlockDecomposition/BoundedCokernelLimitToLevel.lean}{GlobalTheory.CompanionTheory.TwCok.pow_torsion_bounded}
\leanref{Development/AutomorphicSpectrum/Proofs/BlockDecomposition/BoundedCokernelLimitToLevel.lean}{GlobalTheory.CompanionTheory.TwCok.lift_of_conn_zero}
\leanref{Development/AutomorphicSpectrum/Proofs/BlockDecomposition/BoundedCokernelLimitToLevel.lean}{GlobalTheory.CompanionTheory.TwCok.lim_cokernel_bounded}

\begin{definition}\label{gal-realization-frobenius-flat-tower}
Let $W$ be a ring and $\ell,c\geq0$ integers. A $W$-module $P$ is \emph{$\ell^c$-flat} if $\ell^cP=0$ and, whenever
$a+b=c$ and $x\in P$ with $\ell^bx=0$, $x\in\ell^aP$. A complex $K=(K^0\to K^1\to\cdots)$ of $W$-modules is
$\ell^c$-flat if every $K^i$ is; $K/\ell^b$ is the complex of the $K^i/\ell^bK^i$. A \emph{flat tower} consists of
complexes $K_n$, $n\geq0$, with $K_n$ $\ell^n$-flat, and maps of complexes $\psi_n:K_n\to K_{n+1}/\ell^n$ that induce
isomorphisms on cohomology. For $j\geq0$, its \emph{cohomology system} in degree $j$ is the inverse system
$(H^j(K_n))_{n\geq0}$ with the transition maps $r_n=\psi_n^{-1}\circ\pi_n$, where $\pi_n:H^j(K_{n+1})\to
H^j(K_{n+1}/\ell^n)$ is induced by the projection.
\end{definition}
\leanref{Development/WeilBound/CoreStatements/PadicHodgeComparison/EtaleCohomologyComparison/PerfectoidPullbacks/LocalSystemQuotients/FlatL.lean}{CrystallineSemistableComparison.Ls.FlatL}
\leanref{Development/WeilBound/CoreStatements/PadicHodgeComparison/EtaleCohomologyComparison/PerfectoidPullbacks/LocalSystemQuotients/QuoDMk.lean}{CrystallineSemistableComparison.Ls.quo}
\leanref{Development/WeilBound/CoreStatements/PadicHodgeComparison/EtaleCohomologyComparison/PerfectoidPullbacks/LocalSystemQuotients/FlatL.lean}{CrystallineSemistableComparison.Ls.tsys}

\begin{lemma}\label{gal-realization-frobenius-flat-tower-bockstein}
Let $W$ be a ring and $\ell\geq0$ an integer.
\begin{enumerate}
\item Let $K$ be an $\ell^m$-flat complex, $L$ an $\ell^{m+1}$-flat complex and $f:K\to L/\ell^m$ a map of complexes
  inducing isomorphisms on cohomology. For $b\leq m$, the induced map $K/\ell^b\to L/\ell^b$ induces isomorphisms on
  cohomology.
\item For a flat tower (Definition~\ref{gal-realization-frobenius-flat-tower}) and $j\geq0$, the cohomology systems in
  degrees $j$ and $j+1$ have Bockstein sequences for $\ell$ (Definition~\ref{gal-realization-frobenius-bockstein}).
\item If $(S_n)$ and $(S'_n)$ have Bockstein sequences and $(T_n)$ and $(T'_n)$ are inverse systems of $W$-modules
  isomorphic to them, then $(T_n)$ and $(T'_n)$ have Bockstein sequences.
\end{enumerate}
\end{lemma}
\leanref{Development/AutomorphicSpectrum/Proofs/BlockDecomposition/QuotientReductionBijectiveOfFlatStatement.lean}{GlobalTheory.T.quotient_reduction_bijective_of_flat}
\leanref{Development/AutomorphicSpectrum/Proofs/BlockDecomposition/QuotientReductionBijectiveOfFlat.lean}{GlobalTheory.CompanionTheory.QredFlat.qred_bij}
\leanref{Development/AutomorphicSpectrum/Proofs/BlockDecomposition/QuotientReductionBijectiveOfFlat.lean}{GlobalTheory.CompanionTheory.QredFlat.exact_quo}
\leanref{Development/AutomorphicSpectrum/Proofs/BlockDecomposition/LseqOfFlatTowerStatement.lean}{GlobalTheory.T.lseq_of_flat_tower}
\leanref{Development/AutomorphicSpectrum/Proofs/BlockDecomposition/LseqOfFlatTower.lean}{GlobalTheory.CompanionTheory.TowFlat.nonempty_lSeq_tsys}
\leanref{Development/AutomorphicSpectrum/Proofs/BlockDecomposition/LseqOfLevelIsosStatement.lean}{GlobalTheory.T.lseq_of_level_isos}
\leanref{Development/AutomorphicSpectrum/Proofs/BlockDecomposition/LseqOfLevelIsos.lean}{GlobalTheory.CompanionTheory.TowerSeq.lseq_transport}

\begin{proof}
(1) First, let $D$ be an $\ell^m$-flat complex with no cohomology, $c\leq m$ and $x\in D^n$ with $dx\in\ell^cD^{n+1}$;
then $x\in\ell^cD^n+dD^{n-1}$ (with $dD^{-1}=0$). Indeed, $dx=\ell^cy$ gives $d(\ell^{m-c}x)=\ell^my=0$, so
$\ell^{m-c}x=dw$ for some $w$, and $w=0$ if $n=0$. For $n=0$, $\ell^{m-c}x=0$ and flatness gives $x\in\ell^cD^0$. For
$n>0$, $dw\in\ell^{m-c}D^n$, so by induction $w=\ell^{m-c}u+v$ with $v$ a boundary; then $\ell^{m-c}(x-du)=0$, and
flatness gives $x-du\in\ell^cD^n$.

Apply this to the cone $D$ of $f$, with $D^n=K^n\oplus(L/\ell^m)^{n-1}$ and $d(x,t)=(dx,f(x)-dt)$. It has no cohomology
because $f$ is a quasi-isomorphism, and it is $\ell^m$-flat because $K$ is and $L/\ell^m$ is ($L$ being
$\ell^{m+1}$-flat). Injectivity: let $z$ be a cocycle of $K/\ell^b$ in degree $n$, lifted to $x\in K^n$ with
$dx\in\ell^bK^{n+1}$, whose image is a boundary, so that $f(x)-dt\in\ell^b(L/\ell^m)^n$ for some $t$. Then
$d(x,t)\in\ell^bD^{n+1}$, so $(x,t)\in\ell^bD^n+dD^{n-1}$, and $x\in\ell^bK^n+dK^{n-1}$: $z$ is a boundary.
Surjectivity: lift a cocycle of $L/\ell^b$ in degree $n$ to $w\in(L/\ell^m)^n$ with $dw=\ell^bp$. Then
$d(0,w)=(0,-\ell^bp)\in\ell^bD^{n+2}$, so $(0,w)=\ell^b(u_1,u_2)+d(g_1,g_2)$. The first components give $dg_1=-\ell^bu_1$,
so $g_1$ is a cocycle modulo $\ell^b$, and the second give $w\equiv f(g_1)-dg_2$ modulo $\ell^b$.

(2) For $b\leq c$ let $\theta_{b,c}:H^j(K_b)\to H^j(K_c/\ell^b)$ be the composite of the isomorphism $H^j(K_b)\cong
H^j(K_b/\ell^b)$ (as $\ell^bK_b=0$) with the maps induced by $\psi_b,\ldots,\psi_{c-1}$ reduced modulo $\ell^b$. By (1)
these are isomorphisms, so $\theta_{b,c}$ is bijective. For $c=a+b$ the sequence $0\to K_c/\ell^a\xrightarrow{\ell^b}
K_c\to K_c/\ell^b\to0$ is exact, as $K_c$ is $\ell^c$-flat; let $\partial$ be its connecting map. Put
$\rho_{c,b}=\theta_{b,c}^{-1}\circ\pi$ for $b\leq c$ and $\rho_{c,b}=0$ otherwise, with $\pi:H^j(K_c)\to
H^j(K_c/\ell^b)$ the projection; $\nu_{a,c}=\ell^{c-a}_*\circ\theta_{a,c}$ for $a\leq c$ and $0$ otherwise, where
$\ell^{c-a}:K_c/\ell^a\to K_c$ is multiplication; and $\delta_{b,a}=\theta_{a,a+b}^{-1}\circ\partial\circ\theta_{b,a+b}$.
Then $\ell^n$ kills $H^j(K_n)$, $\rho_{c,c}$ is the identity, and $\nu_{a,c}(\rho_{c,a}(x))=\ell^{c-a}_*(\pi x)=
\ell^{c-a}x$. The identities $\rho_{c+1,b}=\rho_{c,b}\circ r_c$, $r_c\circ\nu_{a+1,c+1}=\nu_{a,c}\circ r_a$ and
$r'_a\circ\delta_{b,a+1}=\delta_{b,a}$ follow from $\theta_{b,c+1}=(\psi_c\bmod\ell^b)_*\circ\theta_{b,c}$ and from the
naturality of the projections, of the multiplications by powers of $\ell$ and of $\partial$ with respect to $\psi_c$
and to the reductions $K_c/\ell^{a+1}\to K_c/\ell^a$. Condition (5) of
Definition~\ref{gal-realization-frobenius-bockstein} is the exactness of the long exact sequence of the short exact
sequence above at $H^j(K_c)$ and at $H^j(K_c/\ell^b)$, transported by the $\theta$.

(3) Compose the maps $\rho$, $\nu$ and $\delta$ with the isomorphisms.
\end{proof}

\begin{lemma}\label{gal-realization-thickening-local-flat}
Let $(R,I,\gamma)$, $\mathcal O$, $X$ and $\kappa$ be as in Definition~\ref{gal-realization-thickening-nilpotent}, let
$p\neq0$ be an integer nilpotent in $R$, and let $V\subset X$ be an open with coordinates. For $j\geq0$, restriction
$H^j(\mathrm{Th}(X/R)_V,\mathcal O)\to H^j(\mathrm{Th}^{\mathrm{fl}}(X/R)_V,\mathcal O)$ is bijective.
\end{lemma}
\leanref{Development/AutomorphicSpectrum/Proofs/BlockDecomposition/SmoothModelRestrictionBijective.lean}{GlobalTheory.CompanionTheory.FlatRes.local_bij}

\begin{proof}
Let $\mathcal Q=\mathrm{Th}(X/R)_V$ and $\mathcal Q'=\mathrm{Th}^{\mathrm{fl}}(X/R)_V$. By
Lemma~\ref{gal-realization-thickening-chart}(3) there is a divided power chart $D$ on $V$ with $B_D$ flat over $R$, so
$D$ lies in $\mathcal Q$ and in $\mathcal Q'$, and by Lemma~\ref{gal-realization-thickening-chart-torsor}(1) it gives
torsor data for $\Phi_{\mathcal Q}$ and $\Phi_{\mathcal Q'}$. The inclusion $\mathcal Q'\to\mathcal Q$ with the identity
maps of the rings is a morphism $\Phi_{\mathcal Q}\to\Phi_{\mathcal Q'}$ over the identity of $R$, and the torsor data are
compatible along the identity of $D$. By Lemma~\ref{gal-realization-thickening-chart-torsor}(3) there are divided power
polynomial structures on $\Phi_{\mathcal Q}$ and $\Phi_{\mathcal Q'}$ that respect $(0)_r$ in every degree; they respect
the constant sequence $(B)_r$, as both sides of the condition always hold, so they are ladders up to $0$ for $(B)_r$ and
$(0)_r$. Multiplication by $a\in R$ carries $B$ into $B$ and $0$ into $0$.

Apply Lemma~\ref{gal-realization-thickening-torsor-base-change} with $R'=R$, one module $N=R$, $\nu$ its identity,
$\mathfrak a_r=\mathfrak a'_r=B$ and $\mathfrak b_r=\mathfrak b'_r=0$ for every $r$, and $r=0$. The map of its hypothesis
is the bijection $R\otimes_RB_D\to B_D$. Its conclusion is that $R\otimes_RH^j(\mathcal Q,\mathcal O)\to H^j(\mathcal
Q',\mathcal O)$, $a\otimes x\mapsto a\,x|_{\mathcal Q'}$, is bijective, so restriction is bijective.
\end{proof}

\begin{lemma}\label{gal-realization-thickening-flat-restriction}
Let $(R,I,\gamma)$, $\mathcal O$, $X$ and $\kappa$ be as in Definition~\ref{gal-realization-thickening-nilpotent}, let
$p\neq0$ be an integer nilpotent in $R$, and suppose that $X$ has a coordinate cover. For $j\geq0$, restriction
$H^j(\mathrm{Th}(X/R),\mathcal O)\to H^j(\mathrm{Th}^{\mathrm{fl}}(X/R),\mathcal O)$ is bijective.
\end{lemma}
\leanref{Development/AutomorphicSpectrum/Proofs/BlockDecomposition/SmoothModelRestrictionBijective.lean}{GlobalTheory.CompanionTheory.FlatRes.all_flat}
\leanref{Development/AutomorphicSpectrum/Proofs/BlockDecomposition/SmoothModelRestrictionBijective.lean}{GlobalTheory.CompanionTheory.FlatRes.st_top_eq}

\begin{proof}
Let $(U_\alpha)_{\alpha\in A}$ be a coordinate cover and $\mathcal U=\{U_\alpha:\alpha\in A\}$. Every thickening lies over
$X=\bigcup_\alpha U_\alpha$ (Lemma~\ref{gal-realization-thickening-over}(2)), so $\mathrm{Th}(X/R)=\mathrm{Th}(X/R)_X$
and $\mathrm{Th}^{\mathrm{fl}}(X/R)=\mathrm{Th}^{\mathrm{fl}}(X/R)_X$. By
Lemma~\ref{gal-realization-thickening-zariski}(1), with $\mathcal V=\mathcal U$ and $\mathcal W=\{X\}$, restriction from
$\mathrm{Th}(X/R)$ to $\mathrm{Th}(X/R)_{\mathcal U}$ and from $\mathrm{Th}^{\mathrm{fl}}(X/R)$ to
$\mathrm{Th}^{\mathrm{fl}}(X/R)_{\mathcal U}$ are bijective on $H^j(-,\mathcal O)$. The two composites of restrictions from
$\mathrm{Th}(X/R)$ to $\mathrm{Th}^{\mathrm{fl}}(X/R)_{\mathcal U}$ agree, so it suffices that restriction from
$\mathrm{Th}(X/R)_{\mathcal U}$ to $\mathrm{Th}^{\mathrm{fl}}(X/R)_{\mathcal U}$ is bijective.

By Lemma~\ref{gal-realization-thickening-over}, the thickenings lying over $U_\alpha$ form cosieves of these two categories
that cover them, the inclusion maps one into the other, and for nonempty $S\subset A$ their intersections are
$\mathrm{Th}(X/R)_{U_S}$ and $\mathrm{Th}^{\mathrm{fl}}(X/R)_{U_S}$, with $U_S=\bigcap_{\alpha\in S}U_\alpha$. The affine
open $U_S$ has coordinates, restricted from $U_\alpha$ for any $\alpha\in S$
(Lemma~\ref{gal-realization-thickening-cover}(1)), so restriction from $\mathrm{Th}(X/R)_{U_S}$ to
$\mathrm{Th}^{\mathrm{fl}}(X/R)_{U_S}$ is bijective by Lemma~\ref{gal-realization-thickening-local-flat}.
Lemma~\ref{gal-realization-thickening-cosieves}(2) concludes.
\end{proof}
\leanref{Development/AutomorphicSpectrum/Proofs/BlockDecomposition/SmoothModelRestrictionBijective.lean}{GlobalTheory.CompanionTheory.FlatRes.glue}

\begin{definition}\label{gal-realization-thickening-torsor-quotient}
Let $\Phi$ and $\Phi'$ be divided power functors over a ring $R$ on categories $\mathcal C$ and $\mathcal C'$, given by
$A,J,\gamma$ and $A',J',\gamma'$ (Definition~\ref{gal-realization-thickening-torsor}), and $\mathfrak c$ a family of ideals
of $\Phi'$. A \emph{morphism modulo} $\mathfrak c$ from $\Phi$ to $\Phi'$ consists of a functor $G:\mathcal C'\to\mathcal
C$ and injective $R$-algebra maps $\rho_T:A(G(T))\to A'(T)/\mathfrak c(T)$ such that, whenever $x\in A(G(T))$ and $y\in
A'(T)$ have $\rho_T(x)=\bar y$, $\rho_{T'}(A(G(\varphi))(x))$ is the class of $A'(\varphi)(y)$ for every $\varphi:T\to
T'$, and if moreover $x\in J(G(T))$ and $y\in J'(T)$, then $\rho_T(\gamma_k(x))$ is the class of $\gamma'_k(y)$ for every
$k$. For a family $\mathfrak b\supset\mathfrak c$, the maps $\rho_T$ followed by $A'(T)/\mathfrak c(T)\to A'(T)/\mathfrak
b(T)$ induce, with $G$, maps $H^j(\mathcal C,A)\to H^j(\mathcal C',A'/\mathfrak b)$. Torsor data $(D,\tau)$ for $\Phi$
and $(D',\tau')$ for $\Phi'$ of the same dimension are \emph{compatible} along $e:D\to G(D')$ if
$\rho_T(\tau_k(G(s)e,G(s')e))$ is the class of $\tau'_k(s,s')$ for all $s,s':D'\to T$.
\end{definition}
\leanref{Development/WeilBound/CoreStatements/PadicHodgeComparison/EtaleCohomologyComparison/PerfectoidPullbacks/Coords.lean}{CrystallineSemistableComparison.Pb.QMor}
\leanref{Development/WeilBound/CoreStatements/PadicHodgeComparison/EtaleCohomologyComparison/PerfectoidPullbacks/Coords.lean}{CrystallineSemistableComparison.Pb.QMor.o}
\leanref{Development/WeilBound/CoreStatements/PadicHodgeComparison/EtaleCohomologyComparison/PerfectoidPullbacks/Coords.lean}{CrystallineSemistableComparison.Pb.QMor.Fr}

\begin{lemma}\label{gal-realization-thickening-torsor-cut-morphism-quotient}
Let $(G,\rho)$ be a morphism modulo $\mathfrak c$ from a divided power functor $\Phi$ on $\mathcal C$ to a divided power
functor $\Phi'$ on $\mathcal C'$, over a ring $R$ and given by $A,J,\gamma$ and $A',J',\gamma'$
(Definition~\ref{gal-realization-thickening-torsor-quotient}). Let $P$ and $P'$ be divided power polynomial structures
on $\Phi$ and $\Phi'$, given by $j_T,X_T$ and $j'_T,X'_T$, and $(D,\tau)$ and $(D',\tau')$ torsor data of dimension
$d+1$ for $\Phi$ and $\Phi'$, compatible along a map $e:D\to G(D')$. Form $\mathcal C_D$, $\mathcal C^0_D$ and $P_D$, and
$\mathcal C'_{D'}$, $\mathcal C'^0_{D'}$ and $P'_{D'}$ (Definition~\ref{gal-realization-thickening-torsor-cut}), and write
$P'_{D'}(T,s)=(P'(T),s^+)$.
\begin{enumerate}
\item The functor $G_D:\mathcal C'_{D'}\to\mathcal C_D$, $(T,s)\mapsto(G(T),G(s)e)$, $\varphi\mapsto G(\varphi)$, maps
  $\mathcal C'^0_{D'}$ into $\mathcal C^0_D$; let $G^0:\mathcal C'^0_{D'}\to\mathcal C^0_D$ be the restricted functor.
\item For $(T,s)$ in $\mathcal C'_{D'}$ let $\eta_{(T,s)}:P_DG_D(T,s)\to G^0P'_{D'}(T,s)$ be the map of $\mathcal C^0_D$
  with $\eta_{(T,s)}j_{G(T)}=G(j'_T)$ (Lemma~\ref{gal-realization-thickening-torsor-cut-adjoint}(3)). Then $\eta$ is a
  natural transformation $P_DG_D\to G^0P'_{D'}$, and $\rho_{P'(T)}(A(\eta_{(T,s)})(X_{G(T)}))$ is the class of $X'_T$.
\item If $(c_k)_{k\geq0}$ and $(y_k)_{k\geq0}$ are finitely supported sequences in $A(G(T))$ and $A'(T)$ with
  $\rho_T(c_k)$ the class of $y_k$ for every $k$, then $\rho_{P'(T)}\bigl(A(\eta_{(T,s)})(\sum_k
  A(j_{G(T)})(c_k)\gamma_k(X_{G(T)}))\bigr)$ is the class of $\sum_kA'(j'_T)(y_k)\gamma'_k(X'_T)$.
\end{enumerate}
\end{lemma}
\leanref{Development/AutomorphicSpectrum/Proofs/BlockDecomposition/MapHBijOfRes.lean}{GlobalTheory.CompanionTheory.LocCmp.Gpt}
\leanref{Development/AutomorphicSpectrum/Proofs/BlockDecomposition/MapHBijOfRes.lean}{GlobalTheory.CompanionTheory.LocCmp.Gcut}
\leanref{Development/AutomorphicSpectrum/Proofs/BlockDecomposition/MapHBijOfRes.lean}{GlobalTheory.CompanionTheory.LocCmp.eta}
\leanref{Development/AutomorphicSpectrum/Proofs/BlockDecomposition/MapHBijOfRes.lean}{GlobalTheory.CompanionTheory.LocCmp.eta_X}
\leanref{Development/AutomorphicSpectrum/Proofs/BlockDecomposition/MapHBijOfRes.lean}{GlobalTheory.CompanionTheory.LocCmp.rho_ev}

\begin{proof}
(1) $G_D$ is a functor, as $G$ is. If $\varphi:(T,s)\to(T',s')$ is a map of $\mathcal C'^0_{D'}$, then
$\tau'_1(\varphi s,s')=0$. By compatibility, $\rho_{T'}(\tau_1(G(\varphi s)e,G(s')e))$ is the class of
$\tau'_1(\varphi s,s')=0$, and as $\rho_{T'}$ is injective, $\tau_1(G(\varphi)G(s)e,G(s')e)=0$. So $G(\varphi)$ is a map
$G_D(T,s)\to G_D(T',s')$ of $\mathcal C^0_D$.

(2) The naturality of $\eta$ is proved as in Lemma~\ref{gal-realization-thickening-torsor-cut-morphism}(2), from
$\eta_{(T,s)}j_{G(T)}=G(j'_T)$ and the bijection of Lemma~\ref{gal-realization-thickening-torsor-cut-adjoint}(3), which
also gives $A(\eta_{(T,s)})(X_{G(T)})=\tau_1(G(j'_Ts)e,G(s^+)e)$. By compatibility, its image under $\rho_{P'(T)}$ is the
class of $\tau'_1(j'_Ts,s^+)=X'_T$.

(3) As $\eta_{(T,s)}j_{G(T)}=G(j'_T)$ and $\rho_T(c_k)$ is the class of $y_k$,
$\rho_{P'(T)}(A(\eta_{(T,s)})(A(j_{G(T)})(c_k)))=\rho_{P'(T)}(A(G(j'_T))(c_k))$ is the class of $A'(j'_T)(y_k)$. As
$A(\eta_{(T,s)})(X_{G(T)})$ lies in $J(G(P'(T)))$, $X'_T$ in $J'(P'(T))$, and the first maps to the class of the second by
(2), $\rho_{P'(T)}(A(\eta_{(T,s)})(\gamma_k(X_{G(T)})))=\rho_{P'(T)}(\gamma_k(A(\eta_{(T,s)})(X_{G(T)})))$ is the class
of $\gamma'_k(X'_T)$. As $\rho_{P'(T)}$ and $A(\eta_{(T,s)})$ are ring maps, the claim follows.
\end{proof}
\leanref{Development/AutomorphicSpectrum/Proofs/BlockDecomposition/MapHBijOfRes.lean}{GlobalTheory.CompanionTheory.LocCmp.τ_zero}
\leanref{Development/AutomorphicSpectrum/Proofs/BlockDecomposition/MapHBijOfRes.lean}{GlobalTheory.CompanionTheory.LocCmp.etaAt_nat}

\begin{lemma}\label{gal-realization-thickening-torsor-quotient-comparison}
Let $(G,\rho)$ be a morphism modulo $\mathfrak c$ from a divided power functor $\Phi$ on $\mathcal C$ to $\Phi'$ on
$\mathcal C'$ (Definition~\ref{gal-realization-thickening-torsor-quotient}), $P$ and $P'$ divided power polynomial
structures on $\Phi$ and $\Phi'$, and $(D,\tau)$ and $(D',\tau')$ torsor data of dimension $d$ compatible along $e:D\to
G(D')$. Let $(\mathfrak b_r)_{r\in\mathbf Z}$ be families of ideals of $\Phi'$ containing $\mathfrak c$, and $r\in\mathbf
Z$. Suppose that $P$ respects $(0)_r$ in degree $0$ and that $P'$ respects $(\mathfrak b_r)_r$ in every degree $r'\leq r$
(Definition~\ref{gal-realization-thickening-torsor}(3)). If for every $r'$ with $r-d\leq r'\leq r$ the map $A(D)\to
A'(D')/\mathfrak b_{r'}(D')$, sending $x$ to the image of $\rho_{D'}(A(e)(x))$, is bijective, then $H^j(\mathcal C,A)\to
H^j(\mathcal C',A'/\mathfrak b_r)$ is bijective for every $j$.
\end{lemma}
\leanref{Development/AutomorphicSpectrum/Proofs/BlockDecomposition/MapHBijOfRes.lean}{GlobalTheory.CompanionTheory.LocCmp.localComparison}

\begin{proof}
The proof is by induction on $d$, for all data as in the statement, and has the shape of the proof of
Lemma~\ref{gal-realization-thickening-torsor-base-change}, with the maps $\bar\rho_T:A(G(T))\to A'(T)/\mathfrak
b_{r'}(T)$, $\rho_T$ followed by the projection, in place of the multiplications.

\emph{The case $d=0$.} As there, $D$ and $D'$ are initial objects, the groups $H^j$ with $j\geq1$ vanish, and by
Lemma~\ref{gal-realization-thickening-torsor-adjoint}(2) the map in degree $0$ is, under $\varepsilon$, the map
$x\mapsto\bar\rho_{D'}(A(e)(x))$ of the hypothesis for $r'=r$: for $T\in\mathcal C'$ the map $D\to G(T)$ is $G(D'\to
T)e$, and $\bar\rho$ is natural.

\emph{The step.} Let the data have dimension $d+1$, and suppose the claim holds for all data of dimension $d$.
Lemma~\ref{gal-realization-thickening-torsor-cut-morphism-quotient} gives $G_D$, $G^0$ and $\eta$.

Step 1. The functor $G^0$ with the maps $\rho_T$ at $(T,s)$ is a morphism modulo $\mathfrak c^0$ from $\Phi^0$ to
$\Phi'^0$, and the torsor data $(D,\mathrm{id}_D)$ and $(D',\mathrm{id}_{D'})$ of dimension $d$ are compatible along
$e$, as in Step 1 of the proof of Lemma~\ref{gal-realization-thickening-torsor-base-change}. The families $\mathfrak
b^0_{r'}$ contain $\mathfrak c^0$, $P^0$ respects $(0)_{r'}$ in degree $0$, and $P'^0$ respects $(\mathfrak
b^0_{r'})_{r'}$ in every degree at most $r$ (Lemma~\ref{gal-realization-thickening-torsor-cut-adjoint}(1)). For
$r_0\in\{r,r-1\}$ and $r_0-d\leq r'\leq r_0$, the map of the hypothesis for $r'$ is bijective, as $r-(d+1)\leq r'\leq r$.
By the claim in dimension $d$, the maps $H^j(\mathcal C^0_D,A\circ U^0)\to H^j(\mathcal C'^0_{D'},(A'/\mathfrak
b_{r_0})\circ U'^0)$ are bijective for $r_0\in\{r,r-1\}$ and every $j$.

Step 2. The functor $G_D$ with the maps $x\mapsto\bar\rho_{P'(T)}(A(\eta_{(T,s)})(x))$ induces maps $H^j(\mathcal
C_D,A\circ U^0P_D)\to H^j(\mathcal C'_{D'},(A'/\mathfrak b_{r_0})\circ U'^0P'_{D'})$. As in Step 2 of the proof of
Lemma~\ref{gal-realization-thickening-torsor-base-change}, restriction along $P_D$ and along $P'_{D'}$ is bijective and
carries the maps of Step 1 to these, so they are bijective for $r_0\in\{r,r-1\}$.

Step 3. Lemma~\ref{gal-realization-thickening-torsor-exact} for the constant sequences $A$ and $0$, which $P$ respects in
degrees $0$ and $-1$ (for $A$ both sides of the condition always hold, and for $0$ the condition does not depend on the
degree), gives an exact sequence $0\to A\circ U\to A\circ U^0P_D\xrightarrow{\beta}A\circ U^0P_D\to0$ of functors on
$\mathcal C_D$. The same lemma for the constant sequence $A'$ and $(\mathfrak b_{r'})_{r'}$ gives an exact sequence
$0\to(A'/\mathfrak b_r)\circ U'\to(A'/\mathfrak b_r)\circ U'^0P'_{D'}\xrightarrow{\beta'}(A'/\mathfrak b_{r-1})\circ
U'^0P'_{D'}\to0$ on $\mathcal C'_{D'}$. The maps given by $G_D$, with $\bar\rho_T$ on the first terms and
$\bar\rho_{P'(T)}A(\eta_{(T,s)})$ on the others, commute with the maps of these sequences, as in Step 3 there, by
Lemma~\ref{gal-realization-thickening-torsor-cut-morphism-quotient}(3). The long exact sequences of
cohomology and the five lemma, with Step 2, show that the maps $H^j(\mathcal C_D,A\circ U)\to H^j(\mathcal
C'_{D'},(A'/\mathfrak b_r)\circ U')$ are bijective for every $j$.

Step 4. Restriction along $U$ and along $U'$ is bijective and carries the map of the statement to the map of Step 3, as
in Step 4 there.
% Lean: LocCmp.main, by induction on d. The case d = 0 is LocCmp.base (mapH_bij_of_res with P.initial). The step is
% LocCmp.step with hcut: Step 1 is the induction hypothesis for m1, Cut.Q1, Cut.E1, e1 and Fr1 with
% Cut.criterion_truncation; Step 2 is mid_bij (tr_bij for Cut.adj, through P.adj and P.prism for eta); Step 3 is ses0
% (filtration_critical_comparison_map for IdF.top and IdF.bot, transported by qE), filtration_critical_comparison_map for
% the target, P.ses_cochains and P.five, with comm_f and comm_g (rho_ev, eta_X); Step 4 is low_bij (tr_bij for Pt.adj).
\end{proof}
\leanref{Development/AutomorphicSpectrum/Proofs/BlockDecomposition/MapHBijOfRes.lean}{GlobalTheory.CompanionTheory.LocCmp.main}
\leanref{Development/AutomorphicSpectrum/Proofs/BlockDecomposition/MapHBijOfRes.lean}{GlobalTheory.CompanionTheory.LocCmp.base}
\leanref{Development/AutomorphicSpectrum/Proofs/BlockDecomposition/MapHBijOfRes.lean}{GlobalTheory.CompanionTheory.LocCmp.mapH_bij_of_res}
\leanref{Development/AutomorphicSpectrum/Proofs/BlockDecomposition/MapHBijOfRes.lean}{GlobalTheory.CompanionTheory.LocCmp.step}
\leanref{Development/AutomorphicSpectrum/Proofs/BlockDecomposition/MapHBijOfRes.lean}{GlobalTheory.CompanionTheory.LocCmp.m1}
\leanref{Development/AutomorphicSpectrum/Proofs/BlockDecomposition/MapHBijOfRes.lean}{GlobalTheory.CompanionTheory.LocCmp.tr_bij}
\leanref{Development/AutomorphicSpectrum/Proofs/BlockDecomposition/MapHBijOfRes.lean}{GlobalTheory.CompanionTheory.LocCmp.mid_bij}
\leanref{Development/AutomorphicSpectrum/Proofs/BlockDecomposition/MapHBijOfRes.lean}{GlobalTheory.CompanionTheory.LocCmp.low_bij}
\leanref{Development/AutomorphicSpectrum/Proofs/BlockDecomposition/MapHBijOfRes.lean}{GlobalTheory.CompanionTheory.LocCmp.ses0}
\leanref{Development/AutomorphicSpectrum/Proofs/BlockDecomposition/MapHBijOfRes.lean}{GlobalTheory.CompanionTheory.LocCmp.comm_f}
\leanref{Development/AutomorphicSpectrum/Proofs/BlockDecomposition/MapHBijOfRes.lean}{GlobalTheory.CompanionTheory.LocCmp.comm_g}
\leanref{Development/WeilBound/CoreStatements/PadicHodgeComparison/EtaleCohomologyComparison/CechPairings/MemBoundariesZeroIff.lean}{CrystallineSemistableComparison.P.five}

\begin{lemma}\label{gal-realization-thickening-chart-reduction}
Let $\ell$, $k$, $W$, $W_m$ and $\gamma$ be as in Definition~\ref{gal-realization-frobenius-flat-witt}, $\mathcal O$ a
ring, $\rho:\mathcal O\to k$ a ring map and $Y$ an $\mathcal O$-scheme, with the thickenings of $Y$ over $(W_m,\ell
W_m,\gamma,\ell)$ formed for $\mathcal O\xrightarrow{\rho}k=W_m/\ell W_m$
(Definition~\ref{gal-realization-thickening-nilpotent}). Let $n\geq1$, $V\subset Y$ an open, and $D=(B,i,J,\delta,y)$
with $t_1,\ldots,t_d$ a divided power chart on $V$ over $(W_{n+1},\ell W_{n+1},\gamma,\ell)$
(Definition~\ref{gal-realization-thickening-coordinates}(3)). Then $\ell^nB\subset J$, the divided powers $\delta$ pass
to $\bar J=J/\ell^nB$, and $\bar D=(B/\ell^nB,\bar\imath,\bar J,\bar\delta,\bar y)$, where $\bar\imath$ is induced by $i$
and $\bar y$ is the image of $y$ under the isomorphism $B/J\to(B/\ell^nB)/\bar J$, is a thickening of $Y$ over
$(W_n,\ell W_n,\gamma,\ell)$. There are a divided power chart $D'$ with $t'_1,\ldots,t'_d$ on $V$ over $(W_n,\ell
W_n,\gamma,\ell)$ and an isomorphism of thickenings $e:D'\to\bar D$ with $e(t'_i)$ the class of $t_i$ for every $i$, such
that $B_{D'}$ is flat over $W_n$ if $B/\ell^nB$ is.
\end{lemma}
\leanref{Development/AutomorphicSpectrum/Proofs/BlockDecomposition/WittChartDescendLevelStatement.lean}{GlobalTheory.T.witt_chart_descend_level}

\begin{proof}
Write $\bar B=B/\ell^nB$. As $\ell^n\in\ell W_{n+1}$, $i(\ell^n)\in J$, and $\ell^nB$ is generated by $i(\ell^n)$. For
$j\geq1$, $\delta_j(i(\ell^n))=i(\gamma_j(\ell^n))$, and $\gamma_j(\ell^n)=\ell^{nj}/j!$ lies in $\ell^nW_{n+1}$: the
$\ell$-adic valuation of $j!$ is less than $j$, so at most $n(j-1)$. Hence $\ell^nB$ is a sub-divided-power ideal of
$J$, and $\delta$ passes to $\bar J$. As $\ell^nB\subset J$, the map $B/J\to\bar B/\bar J$ is an isomorphism, and it
carries the map from $W_{n+1}/\ell W_{n+1}$ to the one from $W_n/\ell W_n$; so $\bar y$ lies over the composite that
Definition~\ref{gal-realization-thickening-nilpotent} requires, and $\bar D$ is a thickening over $(W_n,\ell
W_n,\gamma,\ell)$. Its divided powers are compatible with $\gamma$ on $\ell W_n$ because $\delta$ is on $\ell W_{n+1}$.

It remains to show that $D'=\bar D$, with $t'_i$ the class of $t_i$ and $e$ the identity, is a divided power chart;
then $B_{D'}=\bar B$, so $B_{D'}$ is flat over $W_n$ if $\bar B$ is. The kernel of $B\to\bar B$ is generated by the image of the kernel
$\ell^nW_{n+1}[x_1,\ldots,x_d]$ of $W_{n+1}[x_1,\ldots,x_d]\to W_n[x_1,\ldots,x_d]$, so $\bar B=W_n[x_1,\ldots,x_d]
\otimes_{W_{n+1}[x_1,\ldots,x_d]}B$ is \'etale over $W_n[x_1,\ldots,x_d]$, with $x_i$ going to the class of $t_i$. The
ideal $\bar J$ is the image of $J=\ell W_{n+1}B$, which is $\ell W_n\bar B$. The point of $\bar D$ is that of $D$
through the bijection $B/J\to\bar B/\bar J$; so $\bar D$ lies over $V$. As $W_{n+1}/\ell=W_n/\ell=k$, the defining
property of a divided power chart for $\bar D$ is then the one for $D$.
% Lean: chart_q = (cq q rho hn c', Iso.refl, fun i => rfl, id); the thickening is LsQ.qT, with lI_le (l^n B in J)
% and subPd (the sub-divided-power ideal). No flatness of the chart is assumed.
\end{proof}
\leanref{Development/AutomorphicSpectrum/Proofs/BlockDecomposition/WittChartDescendLevel.lean}{GlobalTheory.CompanionTheory.ChartRed.chart_q}
\leanref{Development/AutomorphicSpectrum/Proofs/BlockDecomposition/WittChartDescendLevel.lean}{GlobalTheory.CompanionTheory.ChartRed.cq}
\leanref{Development/WeilBound/CoreStatements/PadicHodgeComparison/EtaleCohomologyComparison/PerfectoidPullbacks/LocalSystemQuotients/RedMk.lean}{CrystallineSemistableComparison.Pb.LsQ.lI_le}
\leanref{Development/WeilBound/CoreStatements/PadicHodgeComparison/EtaleCohomologyComparison/PerfectoidPullbacks/LocalSystemQuotients/RedMk.lean}{CrystallineSemistableComparison.Pb.LsQ.subPd}
\leanref{Development/AutomorphicSpectrum/Proofs/BlockDecomposition/WittChartDescendLevel.lean}{GlobalTheory.CompanionTheory.ChartRed.etale_q}
\leanref{Development/AutomorphicSpectrum/Proofs/BlockDecomposition/WittChartDescendLevel.lean}{GlobalTheory.CompanionTheory.ChartRed.J_q}
\leanref{Development/AutomorphicSpectrum/Proofs/BlockDecomposition/WittChartDescendLevel.lean}{GlobalTheory.CompanionTheory.ChartRed.qe_bij}
\leanref{Development/AutomorphicSpectrum/Proofs/BlockDecomposition/WittChartDescendLevel.lean}{GlobalTheory.CompanionTheory.ChartRed.rep_abs}
\leanref{Development/WeilBound/CoreStatements/PadicHodgeComparison/EtaleCohomologyComparison/PerfectoidPullbacks/LocalSystemQuotients/InUQT.lean}{CrystallineSemistableComparison.Pb.SubquotientResolution.inU_qT}

\begin{lemma}\label{gal-realization-thickening-local-quotient}
Let $\ell$, $k$, $W$, $W_m$ and $\gamma$ be as in Definition~\ref{gal-realization-frobenius-flat-witt}, $\mathcal O$ a
ring, $\rho:\mathcal O\to k$ a ring map and $Y$ an $\mathcal O$-scheme, and for $m\geq1$ let $\mathrm{Th}^{\mathrm{fl}}_m(Y/W)
=\mathrm{Th}^{\mathrm{fl}}(Y/W_m)$, formed for $(W_m,\ell W_m,\gamma,\ell)$ and $\mathcal O\xrightarrow{\rho}k=W_m/\ell W_m$
(Definition~\ref{gal-realization-thickening-nilpotent}). Let $n\geq1$, and let $V\subset Y$ be an open with coordinates.
For $j\geq0$, the map
\[
 H^j(\mathrm{Th}^{\mathrm{fl}}_n(Y/W)_V,\mathcal O)\to H^j(\mathrm{Th}^{\mathrm{fl}}_{n+1}(Y/W)_V,\mathcal O/\ell^n
 \mathcal O)
\]
induced by $T\mapsto\bar T$ (Lemma~\ref{gal-realization-frobenius-flat-thickenings}(2) below, for $F=F_Y$) and the
identity maps of the rings $B_T/\ell^nB_T$ is bijective (Definition~\ref{gal-realization-thickening-nilpotent}).
\end{lemma}
\leanref{Development/AutomorphicSpectrum/Proofs/BlockDecomposition/WittCoordQuotientPullBijectiveStatement.lean}{GlobalTheory.T.witt_coord_quotient_pull_bijective}

\begin{proof}
Let $\mathcal Q=\mathrm{Th}^{\mathrm{fl}}_n(Y/W)_V$ and $\mathcal Q'=\mathrm{Th}^{\mathrm{fl}}_{n+1}(Y/W)_V$. The functor
$T\mapsto\bar T$ maps $\mathcal Q'$ into $\mathcal Q$, as the point of $\bar T$ is that of $T$ through the bijection
$B_T/J_T\to B_{\bar T}/J_{\bar T}$. As $\ell$ is nilpotent in $W_{n+1}$, Lemma~\ref{gal-realization-thickening-chart}(3)
gives a divided power chart $D$ with $t_1,\ldots,t_d$ on $V$ over $(W_{n+1},\ell W_{n+1},\gamma,\ell)$ with $B_D$ flat
over $W_{n+1}$. Lemma~\ref{gal-realization-thickening-chart-reduction} gives a divided power chart $D'$ with
$t'_1,\ldots,t'_d$ on $V$ over $(W_n,\ell W_n,\gamma,\ell)$ and an isomorphism $e:D'\to\bar D$ with $e(t'_i)=\bar t_i$,
the class of $t_i$. As $\bar D$ is flat by Lemma~\ref{gal-realization-frobenius-flat-thickenings}(2), so is $D'$.
By Lemma~\ref{gal-realization-thickening-chart-torsor}(1), $D'$ and $D$ give torsor data of dimension $d$ for
$\Phi_{\mathcal Q}$ and $\Phi_{\mathcal Q'}$, and by Lemma~\ref{gal-realization-thickening-chart-torsor}(3) both have
divided power polynomial structures that respect $(0)_r$ and $(\ell^nB)_r$ in every degree.

The functor $T\mapsto\bar T$ with the identity maps of the rings $B_T/\ell^nB_T$ is a morphism modulo $\ell^nB$
from $\Phi_{\mathcal Q}$ to $\Phi_{\mathcal Q'}$ (Definition~\ref{gal-realization-thickening-torsor-quotient}), as the
maps and the divided powers of the $\bar T$ are induced by those of the $T$. The torsor data are compatible along
$e:D'\to\bar D$: for $s,s':D\to T$, $\bar s'(e(t'_i))-\bar s(e(t'_i))=\bar s'(\bar t_i)-\bar s(\bar t_i)$ is the class
of $s'(t_i)-s(t_i)$. Apply Lemma~\ref{gal-realization-thickening-torsor-quotient-comparison} with $\mathfrak
b_r=\ell^nB$ for every $r$ and $r=0$: the map $B_{D'}\to B_D/\ell^nB_D$ of its hypothesis is $e$, which is bijective.
Its conclusion is the claim.
\end{proof}
\leanref{Development/AutomorphicSpectrum/Proofs/BlockDecomposition/WittCoordQuotientPullBijective.lean}{GlobalTheory.CompanionTheory.GradedDegreeZeroIsLocal.local_bij}
\leanref{Development/WeilBound/CoreStatements/PadicHodgeComparison/EtaleCohomologyComparison/PerfectoidPullbacks/LocalSystemQuotients/Functor.lean}{CrystallineSemistableComparison.Pb.LsQ.mq}
\leanref{Development/WeilBound/CoreStatements/PadicHodgeComparison/EtaleCohomologyComparison/PerfectoidPullbacks/LocalSystemQuotients/InUQT.lean}{CrystallineSemistableComparison.Pb.SubquotientResolution.mqS}

\begin{lemma}\label{gal-realization-thickening-flat-quotient}
Let $\ell$, $k$, $W$, $W_m$ and $\gamma$ be as in Definition~\ref{gal-realization-frobenius-flat-witt}, $\mathcal O$ a
ring, $\rho:\mathcal O\to k$ a ring map and $Y$ an $\mathcal O$-scheme with a coordinate cover, and for $m\geq1$ let
$\mathrm{Th}^{\mathrm{fl}}_m(Y/W)=\mathrm{Th}^{\mathrm{fl}}(Y/W_m)$, formed for $(W_m,\ell W_m,\gamma,\ell)$ and $\mathcal
O\xrightarrow{\rho}k=W_m/\ell W_m$ (Definition~\ref{gal-realization-thickening-nilpotent}). Let $n\geq1$. For $j\geq0$,
the map
$H^j(\mathrm{Th}^{\mathrm{fl}}_n(Y/W),\mathcal O)\to H^j(\mathrm{Th}^{\mathrm{fl}}_{n+1}(Y/W),\mathcal O/\ell^n\mathcal O)$
induced by $T\mapsto\bar T$ (Lemma~\ref{gal-realization-frobenius-flat-thickenings}(2) below, for $F=F_Y$) and the
identity maps of the rings $B_T/\ell^nB_T$ is bijective.
\end{lemma}
\leanref{Development/WeilBound/CoreStatements/PadicHodgeComparison/EtaleCohomologyComparison/PerfectoidPullbacks/LocalSystemQuotients/BijCov.lean}{CrystallineSemistableComparison.Pb.SubquotientResolution.bij_cov}

\begin{proof}
Let $(U_\alpha)_{\alpha\in A}$ be a coordinate cover, $\mathcal U=\{U_\alpha:\alpha\in A\}$, and $\mathcal
Q_m=\mathrm{Th}^{\mathrm{fl}}_m(Y/W)$. As in the proof of Lemma~\ref{gal-realization-thickening-local-quotient}, $T\mapsto
\bar T$ maps the thickenings in $\mathcal Q_{n+1}$ lying over $U_\alpha$ to thickenings lying over $U_\alpha$, so it maps
$(\mathcal Q_{n+1})_{\mathcal U}$ into $(\mathcal Q_n)_{\mathcal U}$ and gives the same kind of map for these categories,
compatible with the restrictions from $\mathcal Q_n$ and $\mathcal Q_{n+1}$. Every thickening lies over
$Y=\bigcup_\alpha U_\alpha$, so these restrictions are bijective by Lemma~\ref{gal-realization-thickening-zariski}: by (1),
with $\mathcal V=\mathcal U$ and $\mathcal W=\{Y\}$, for $\mathcal O$ on $\mathcal Q_n$, and by (2), with $\mathcal
Q'=\mathcal Q_{n+1}$, for $\mathcal O/\ell^n\mathcal O=J^{[0]}/\ell^nB$ on $\mathcal Q_{n+1}$. It remains to show that the
map for $(\mathcal Q_n)_{\mathcal U}$ and $(\mathcal Q_{n+1})_{\mathcal U}$ is bijective.

By Lemma~\ref{gal-realization-thickening-over}, the thickenings lying over $U_\alpha$ form cosieves of these two categories
that cover them, $T\mapsto\bar T$ maps one into the other, and for nonempty $S\subset A$ their intersections are
$(\mathcal Q_n)_{U_S}$ and $(\mathcal Q_{n+1})_{U_S}$. As $U_S$ has coordinates
(Lemma~\ref{gal-realization-thickening-cover}(1)), the map for these is bijective by
Lemma~\ref{gal-realization-thickening-local-quotient}. Lemma~\ref{gal-realization-thickening-cosieves}(2) concludes.
\end{proof}
\leanref{Development/AutomorphicSpectrum/Proofs/BlockDecomposition/WittLevelMorphismCoverCompatStatement.lean}{GlobalTheory.T.witt_level_morphism_cover_compat}
\leanref{Development/AutomorphicSpectrum/Proofs/BlockDecomposition/WittLevelMorphismCoverCompat.lean}{GlobalTheory.CompanionTheory.FlatCoverGlue.glue}
\leanref{Development/AutomorphicSpectrum/Proofs/BlockDecomposition/WittSiteRestrictBijectiveStatement.lean}{GlobalTheory.T.witt_site_restrict_bijective}
\leanref{Development/AutomorphicSpectrum/Proofs/BlockDecomposition/WittSiteQuotientRestrictBijectiveStatement.lean}{GlobalTheory.T.witt_site_quotient_restrict_bijective}
\leanref{Development/WeilBound/CoreStatements/PadicHodgeComparison/EtaleCohomologyComparison/PerfectoidPullbacks/LocalSystemQuotients/BijOfSub.lean}{CrystallineSemistableComparison.Pb.SubquotientResolution.bij_of_sub}

\begin{definition}\label{gal-realization-frobenius-flat-witt}
Let $\ell$ be a prime, $k$ a perfect field of characteristic $\ell$, $W=W(k)$, and for $m\geq0$ let $W_m=W/\ell^m$, with
$W_0=0$ and the divided powers $\gamma_k(\ell x)=\ell^kx^k/k!$ on $\ell W_m$. Let $\mathcal O$ be a ring, $F$ a point
functor over $\mathcal O$ (Definition~\ref{gal-realization-period-functors}) and $\rho:\mathcal O\to k$ a ring map. A
\emph{thickening} of $F$ over $W_m$ for $\rho$ consists of a ring $B$, a ring map $i:W_m\to B$, an ideal $J$ of $B$ with
divided powers $\delta$ such that $i(\ell W_m)\subset J$ and $\delta_k(i(x))=i(\gamma_k(x))$ for $x\in\ell W_m$, and an
element $y\in F(B/J)$ that lies over $\mathcal O\xrightarrow{\rho}k=W/\ell\to W_m/\ell W_m\xrightarrow{i}B/J$. Maps are
as in Definition~\ref{gal-realization-period-thickenings}(2). Let $\mathrm{Th}_m(F/W)$ be the category of these
thickenings and $H^j_m(F/W)=H^j(\mathrm{Th}_m(F/W),\mathcal O)$, with $\mathcal O$ sending a thickening to the $W$-module
$B$. For $n\leq m$ a thickening over $W_n$ is one over $W_m$ through $W_m\to W_n$, which gives the reduction
$H^j_m(F/W)\to H^j_n(F/W)$. For $K$, $\mathcal O_K$ and $\mathbf F_q$ as in
Definition~\ref{gal-realization-frobenius-special-fibre}, $k=\mathbf F_q$, $\rho$ the residue map and an $\mathcal
O_K$-scheme $Y$, the thickenings of $F_Y$ over $W_m$ are the thickenings of $Y$ over $W_m$ of
Definition~\ref{gal-realization-frobenius-special-fibre}(1), so $\mathrm{Th}_m(F_Y/W)=\mathrm{Th}_m(Y/W)$ and
$H^j_m(F_Y/W)=H^j_m(Y/W)$.
\end{definition}
\leanref{Development/WeilBound/CoreStatements/PadicHodgeComparison/EtaleCohomologyComparison/BocksteinBaseChange/ResCoeff.lean}{CrystallineSemistableComparison.Wc.Wm}
\leanref{Development/WeilBound/CoreStatements/PadicHodgeComparison/EtaleCohomologyComparison/BocksteinBaseChange/ResCoeff.lean}{CrystallineSemistableComparison.Wc.pd}
\leanref{Development/WeilBound/CoreStatements/PadicHodgeComparison/EtaleCohomologyComparison/BocksteinBaseChange/ResCoeff.lean}{CrystallineSemistableComparison.Wc.res}
\leanref{Development/WeilBound/CoreStatements/PadicHodgeComparison/EtaleCohomologyComparison/BocksteinBaseChange/ResCoeff.lean}{CrystallineSemistableComparison.Wc.L}
\leanref{Development/WeilBound/CoreStatements/PadicHodgeComparison/EtaleCohomologyComparison/BocksteinBaseChange/ResCoeff.lean}{CrystallineSemistableComparison.Wc.redMor}
\leanref{Development/WeilBound/CoreStatements/PadicHodgeComparison/EtaleCohomologyComparison/BocksteinBaseChange/ResCoeff.lean}{CrystallineSemistableComparison.Wc.HS}

\begin{lemma}\label{gal-realization-frobenius-flat-thickenings}
Let $\ell$, $k$, $W$, $W_m$, $\mathcal O$, $F$, $\rho$, $\mathrm{Th}_m(F/W)$ and $H^j_m(F/W)$ be as in
Definition~\ref{gal-realization-frobenius-flat-witt}, for $m\geq0$. Let $\mathrm{Th}^{\mathrm{fl}}_m(F/W)$ be the full
subcategory of the thickenings $(B,i,J,\delta,y)$ with $B$ flat over $W_m$ through $i$, let $K_m$ be the complex of
cochains of $\mathrm{Th}^{\mathrm{fl}}_m(F/W)$ with values in $\mathcal O$
(Definition~\ref{gal-realization-split-diagrams}(1)), and let $\mathrm{res}_m:H^j_m(F/W)\to H^j(K_m)$ be the
restriction.
\begin{enumerate}
\item $K_m$ is $\ell^m$-flat over $W$ (Definition~\ref{gal-realization-frobenius-flat-tower}).
\item Let $n\geq1$ and $T=(B,i,J,\delta,y)$ in $\mathrm{Th}^{\mathrm{fl}}_{n+1}(F/W)$. Then $\ell^nB\subset J$, the
  divided powers $\delta$ pass to $\bar J=J/\ell^nB$, and $\bar T=(B/\ell^nB,\bar i,\bar J,\bar\delta,\bar y)$ lies in
  $\mathrm{Th}^{\mathrm{fl}}_n(F/W)$, where $\bar y$ is the image of $y$ under the isomorphism $B/J\to(B/\ell^nB)/\bar J$;
  the projections $B\to B/\ell^nB$ are maps $T\to\bar T$ of $\mathrm{Th}_{n+1}(F/W)$, natural in $T$.
\item For $n\geq1$, the map from $K_{n+1}/\ell^n$ to the complex of cochains of $\mathrm{Th}^{\mathrm{fl}}_{n+1}(F/W)$
  with values in $T\mapsto B_T/\ell^nB_T$ is bijective. Let $\psi_n:K_n\to K_{n+1}/\ell^n$ be the pullback along
  $T\mapsto\bar T$ followed by the inverse of this bijection, and $\psi_0=0$.
\item On $H^j_{n+1}(F/W)$, $\psi_n\circ\mathrm{res}_n\circ\mathrm{red}=\pi\circ\mathrm{res}_{n+1}$, where $\mathrm{red}$
  is the reduction to $H^j_n(F/W)$ and $\pi$ is induced by the projection $K_{n+1}\to K_{n+1}/\ell^n$.
\item Let $C$, $\mathcal O_C^\flat$ and $\theta$ be as in Definition~\ref{gal-crystalline-test-algebras},
  $\bar{\mathbf Q}_\ell$ the algebraic closure of $\mathbf Q_\ell$ in $C$, $K$ a field with $\mathbf Q_\ell\subset
  K\subset\bar{\mathbf Q}_\ell$ whose valuation ring $\mathcal O_K=K\cap\mathcal O_C$ has an irreducible element, and $Y$
  a smooth projective $\mathcal O_K$-scheme. Let $\mathcal O=\mathcal O_K$ and $F=F_Y$, let $\rho$ be surjective, and let
  $s:k\to\mathcal O_C^\flat$ be a ring map such that $\theta([s(\rho(x))])$ and $x$ have the same image in the residue
  field of $\mathcal O_C$ for every $x\in\mathcal O_K$; the irreducible element, the surjectivity of $\rho$ and $s$ are
  not used. Then every $\mathrm{res}_m$ is bijective, and every $\psi_n$ induces isomorphisms on cohomology.
\end{enumerate}
% Lean form: every item is in the Lean's form; (5) is seqflat0_pf and flat_restriction, for every intermediate field K
% of PadicAlgCl l over Q_l, every residue datum r : Res l K (k, rho, s) and every model with M.Sm.
\end{lemma}
\leanref{Development/WeilBound/CoreStatements/PadicHodgeComparison/EtaleCohomologyComparison/PerfectoidPullbacks/LocalSystemQuotients/ModFSW.lean}{CrystallineSemistableComparison.Pb.LsG.fl}
\leanref{Development/WeilBound/CoreStatements/PadicHodgeComparison/EtaleCohomologyComparison/PerfectoidPullbacks/LocalSystemQuotients/ModFSW.lean}{CrystallineSemistableComparison.Pb.LsG.Kt}
\leanref{Development/WeilBound/CoreStatements/PadicHodgeComparison/EtaleCohomologyComparison/PerfectoidPullbacks/LocalSystemQuotients/ModFSW.lean}{CrystallineSemistableComparison.Pb.LsG.resF}
\leanref{Development/AutomorphicSpectrum/Proofs/BlockDecomposition/FlatWittTowerLevelStatement.lean}{GlobalTheory.T.flat_witt_tower_level}
\leanref{Development/WeilBound/CoreStatements/PadicHodgeComparison/EtaleCohomologyComparison/PerfectoidPullbacks/LocalSystemQuotients/RedMk.lean}{CrystallineSemistableComparison.Pb.LsQ.qT}
\leanref{Development/WeilBound/CoreStatements/PadicHodgeComparison/EtaleCohomologyComparison/PerfectoidPullbacks/LocalSystemQuotients/RedMk.lean}{CrystallineSemistableComparison.Pb.LsQ.subPd}
\leanref{Development/AutomorphicSpectrum/Proofs/BlockDecomposition/FlatThickeningQuotientStatement.lean}{GlobalTheory.T.flat_thickening_quotient}
\leanref{Development/AutomorphicSpectrum/Proofs/BlockDecomposition/CompareQuotientBijectiveStatement.lean}{GlobalTheory.T.compare_quotient_bijective}
\leanref{Development/WeilBound/CoreStatements/PadicHodgeComparison/EtaleCohomologyComparison/PerfectoidPullbacks/LocalSystemQuotients/QuoZeroSub.lean}{CrystallineSemistableComparison.Pb.LsQ.psiP}
\leanref{Development/WeilBound/CoreStatements/PadicHodgeComparison/EtaleCohomologyComparison/PerfectoidPullbacks/LocalSystemQuotients/QuoZeroSub.lean}{CrystallineSemistableComparison.Pb.LsQ.psiP_compat}
\leanref{Development/WeilBound/CoreStatements/PadicHodgeComparison/EtaleCohomologyComparison/PerfectoidPullbacks/LocalSystemQuotients/QuoZeroSub.lean}{CrystallineSemistableComparison.Pb.LsQ.seqflat0_pf}
\leanref{Development/AutomorphicSpectrum/Proofs/BlockDecomposition/SmoothModelRestrictionBijectiveStatement.lean}{GlobalTheory.T.smooth_model_restriction_bijective}
\leanref{Development/WeilBound/CoreStatements/PadicHodgeComparison/EtaleCohomologyComparison/PerfectoidPullbacks/LocalSystemQuotients/BijCov.lean}{CrystallineSemistableComparison.Pb.SubquotientResolution.bij}

\begin{proof}
(1) A cochain of degree $i$ assigns to each chain $s$ an element of the ring $B$ of its last thickening, which is a
$W_m$-module, so $\ell^m$ kills $K^i_m$. Let $a+b=m$ and $x\in K^i_m$ with $\ell^bx=0$. For each chain $s$, as $B$ is
flat over $W_m$ and $\ell^b\cdot\ell^a=0$ in $W_m$, the relation $\ell^bx(s)=0$ is trivial: $x(s)=\sum_ju_jv_j$ with
$u_j\in W_m$, $\ell^bu_j=0$ and $v_j\in B$. As $\ell$ is a nonzerodivisor of $W$, $\ell^bu_j=0$ in $W_m=W/\ell^m$
gives $u_j\in\ell^aW_m$, so $x(s)=\ell^az(s)$ with $z(s)\in B$, and $x=\ell^az$.

(2) $\ell^n\in\ell W_{n+1}$, so $i(\ell^n)\in J$. The ideal $\ell^nB$ is generated by $i(\ell^n)$, and for $k\geq1$,
$\delta_k(i(\ell^n))=i(\gamma_k(\ell^n))$, where $\gamma_k(\ell^n)=\ell^{nk}/k!$ lies in $\ell^nW_{n+1}$: the
$\ell$-adic valuation of $k!$ is less than $k$, so at most $n(k-1)$. Hence $\ell^nB$ is a sub-divided-power ideal of
$J$, and $\delta$ passes to the quotient. $B/\ell^nB=W_n\otimes_{W_{n+1}}B$ is flat over $W_n$. As $\ell^nB\subset J$,
$B/J\to(B/\ell^nB)/\bar J$ is an isomorphism, and it carries the map from $W_{n+1}/\ell W_{n+1}$ to the one from
$W_n/\ell W_n$; so $\bar y$ lies over the composite required in
Definition~\ref{gal-realization-frobenius-flat-witt}, and the projection carries $y$ to $\bar y$. The remaining conditions
hold by construction.

(3) Both complexes have in degree $i$ the products over chains of the rings $B$ and $B/\ell^nB$, and $\ell^n$ times a
product is the product of $\ell^n$ times the factors. Then $\psi_n$ is a map of complexes.

(4) Both sides are induced by functors from $\mathrm{Th}^{\mathrm{fl}}_{n+1}(F/W)$ to $\mathrm{Th}_{n+1}(F/W)$, the
inclusion and $T\mapsto\bar T$, with coefficient maps $B\to B/\ell^nB$ and the identity of $B/\ell^nB$. The projections
$T\to\bar T$ of (2) form a natural transformation compatible with these maps, so
Lemma~\ref{gal-realization-period-prism} gives the identity.

(5) For $m=0$ both sides vanish, as $W_0=0$. Let $m\geq1$. As $\ell^m=0$ in $W_m$, the thickenings of $F_Y$ over $W_m$
for $\rho$ (Definition~\ref{gal-realization-frobenius-flat-witt}) are the thickenings of $Y$ over $(W_m,\ell
W_m,\gamma,\ell)$ for $\mathcal O_K\xrightarrow{\rho}k=W_m/\ell W_m$, with the same maps
(Definition~\ref{gal-realization-thickening-nilpotent}). So $\mathrm{Th}_m(F_Y/W)=\mathrm{Th}(Y/W_m)$, with $\ell\neq0$
nilpotent in $W_m$, and its flat thickenings form $\mathrm{Th}^{\mathrm{fl}}(Y/W_m)$. As $Y$ is smooth and projective
over $\mathcal O_K$, it is separated, quasi-compact and smooth, so it has a coordinate cover
(Lemma~\ref{gal-realization-thickening-cover}(2)). A cochain with values in
$\mathcal O$ is the same whether the rings $B$ are viewed as $W$-modules or as $W_m$-modules. So $\mathrm{res}_m$ is
bijective by Lemma~\ref{gal-realization-thickening-flat-restriction}. For $n\geq1$, $\psi_n$ is the cochain map
inducing the map of Lemma~\ref{gal-realization-thickening-flat-quotient}, followed by the inverse of the bijection of
(3), so it induces isomorphisms on cohomology by that lemma. Finally $K_0=0$, as $W_0=0$, and $K_1/\ell^0=0$, so $\psi_0$
induces isomorphisms between zero groups.
% Lean: (1) is flat_witt_tower_level (FlatLevel.flatL_Kt), (2) is LsQ.qT with flat_thickening_quotient, (3) is
% compare_quotient_bijective and LsQ.psiP, (4) is psiP_compat; all four are stated for every ring O, F : B.Fn O and
% rho : O -> k, over every perfect field k of characteristic l. In (5),
% res_m is FlatRes.resF_bij (restrict_bij to St top {top}, scM_bij for the change of scalars from W to W_m,
% all_flat and st_top_eq), with the cover from cov_ex in flat_restriction; psi_n is LsQ.psiP, and its bijectivity is
% SubquotientResolution.bij composed with hinv of cmpq, in seqflat0_pf; psi_0 is zh, with Kt_zero_sub and quo_zero_sub.
% Both use only r.rho and, from M.Sm, projectivity (hM.1.1, for properness) and smoothness (hM.2, for cov_ex); the
% semistability clause of M.Good, r.surj and r.s are not used.
\end{proof}
\leanref{Development/AutomorphicSpectrum/Proofs/BlockDecomposition/SmoothModelRestrictionBijective.lean}{GlobalTheory.CompanionTheory.FlatRes.flat_restriction}
\leanref{Development/AutomorphicSpectrum/Proofs/BlockDecomposition/SmoothModelRestrictionBijective.lean}{GlobalTheory.CompanionTheory.FlatRes.resF_bij}
\leanref{Development/WeilBound/CoreStatements/PadicHodgeComparison/EtaleCohomologyComparison/PerfectoidPullbacks/LocalSystemQuotients/QuoZeroSub.lean}{CrystallineSemistableComparison.Pb.LsQ.Kt_zero_sub}

\subsection{Level-one finiteness and Bockstein limits}\label{subsection-gal-level-one-finiteness}

\begin{lemma}\label{gal-realization-thickening-torsor-vanishing}
Let $\Phi$ be a divided power functor on a category $\mathcal C$ over a ring, $P$ a divided power polynomial structure on
$\Phi$, $(D,\tau)$ a torsor datum of dimension $d$ for $\Phi$, $(\mathfrak a_r)_{r\in\mathbf Z}$ and $(\mathfrak
b_r)_{r\in\mathbf Z}$ sequences of families of ideals of $\Phi$, and $r\in\mathbf Z$, such that $P$ is a ladder up to $r$
for $(\mathfrak a_r)$ and $(\mathfrak b_r)$ (Definition~\ref{gal-realization-thickening-torsor}). If $\mathfrak
a_{r'}(D)\subset\mathfrak b_{r'}(D)$ for every $r'$ with $r-d\leq r'\leq r$, then $H^j(\mathcal C,\mathfrak a_r/\mathfrak
b_r)=0$ for every $j$.
\end{lemma}
\leanref{Development/WeilBound/CoreStatements/PadicHodgeComparison/EtaleCohomologyComparison/PerfectoidPullbacks/LocVan.lean}{CrystallineSemistableComparison.Pb.loc_van}
\leanref{Development/WeilBound/CoreStatements/PadicHodgeComparison/EtaleCohomologyComparison/PerfectoidPullbacks/LocalisationLadders/VanishingOfLadder.lean}{CrystallineSemistableComparison.Pb.vanishing_of_ladder}

\begin{proof}
The proof is by induction on $d$, for all data as in the statement.

\emph{The case $d=0$.} As in the proof of Lemma~\ref{gal-realization-thickening-torsor-base-change}, $D$ is an initial
object. By Lemma~\ref{gal-realization-thickening-torsor-adjoint}(2), $H^j(\mathcal C,\mathfrak a_r/\mathfrak b_r)=0$ for
$j\geq1$, and $H^0(\mathcal C,\mathfrak a_r/\mathfrak b_r)$ is isomorphic to $(\mathfrak a_r/\mathfrak b_r)(D)$, which is
$0$ as $\mathfrak a_r(D)\subset\mathfrak b_r(D)$.

\emph{The step.} Let the torsor datum have dimension $d+1$, and suppose the claim holds for all data of dimension $d$.
Choose maps $\sigma_T:D\to T$ for $T\in\mathcal C$, and form $\mathcal C_D$, $\mathcal C^0_D$, $U$, $\iota$, $U^0$,
$\Phi^0$, $P^0$ and $P_D$ (Definition~\ref{gal-realization-thickening-torsor-cut}). By
Lemma~\ref{gal-realization-thickening-torsor-cut-adjoint}(2) and (3), $S_\sigma$ is left adjoint to $U$ and $P_D$ to
$\iota$, so $P_DS_\sigma:\mathcal C\to\mathcal C^0_D$ is left adjoint to $U\iota=U^0$.

Let $r''\in\{r,r-1\}$. By Lemma~\ref{gal-realization-thickening-torsor-cut-adjoint}(1), $(D,\mathrm{id}_D)$ is a torsor
datum of dimension $d$ for $\Phi^0$, and $P^0$ is a ladder up to $r''$ for $(\mathfrak a^0_{r'})_{r'}$ and $(\mathfrak
b^0_{r'})_{r'}$, as $r''\leq r$. For $r''-d\leq r'\leq r''$ we have $r-(d+1)\leq r'\leq r$, so $\mathfrak
a^0_{r'}(D,\mathrm{id}_D)=\mathfrak a_{r'}(D)\subset\mathfrak b_{r'}(D)=\mathfrak b^0_{r'}(D,\mathrm{id}_D)$. By the claim
in dimension $d$, $H^j(\mathcal C^0_D,\mathfrak a^0_{r''}/\mathfrak b^0_{r''})=0$ for every $j$, and by
Lemma~\ref{gal-realization-thickening-torsor-adjoint}(1) for the left adjoint $P_DS_\sigma$, $H^j(\mathcal C,(\mathfrak
a_{r''}/\mathfrak b_{r''})\circ U^0P_DS_\sigma)=0$ for every $j$.

As $P$ is a ladder up to $r$, Lemma~\ref{gal-realization-thickening-torsor-exact}, composed with $S_\sigma$, gives an
exact sequence of functors on $\mathcal C$
\[
 0\to\mathfrak a_r/\mathfrak b_r\to(\mathfrak a_r/\mathfrak b_r)\circ U^0P_DS_\sigma\xrightarrow{\beta}(\mathfrak
 a_{r-1}/\mathfrak b_{r-1})\circ U^0P_DS_\sigma\to0,
\]
as $US_\sigma$ is the identity. It is exact on cochains in every degree, as it is exact at every object, and in its long
exact sequence of cohomology the groups of the second and third functors vanish. So $H^0(\mathcal C,\mathfrak
a_r/\mathfrak b_r)$, which injects into $H^0$ of the second functor, vanishes, and $H^{j+1}(\mathcal C,\mathfrak
a_r/\mathfrak b_r)$, the image of the connecting map from $H^j$ of the third functor, vanishes for $j\geq0$.
% Lean: vanishing_of_ladder, by induction on d. The case d = 0 is cohomology_vanishing_of_trivial_frame
% (FramedSite.cohomology_subsingleton_of_frame_zero, through P.initial and Frame.initial). The step: hcut is the induction
% hypothesis for Cut.Q1 and Cut.E1 with Cut.ladder_truncation, transported by P.adj for Cut.adjPσ (the composite of
% Pt.adjσ and Cut.adj); the exact sequence is filtration_critical_comparison_map whiskered by Pt.Sσ; and the conclusion
% is cohomology_vanishing_short_exact (Cohomology.subsingleton_H_of_ses, through P.ses_cochains and P.les).
\end{proof}
\leanref{Development/AutomorphicSpectrum/Proofs/BlockDecomposition/CohomologyVanishingOfTrivialFrameStatement.lean}{GlobalTheory.T.cohomology_vanishing_of_trivial_frame}
\leanref{Development/WeilBound/CoreStatements/PadicHodgeComparison/EtaleCohomologyComparison/PerfectoidPullbacks/LocalisationLadders/CriterionTruncation.lean}{CrystallineSemistableComparison.Pb.Cut.adjPσ}
\leanref{Development/AutomorphicSpectrum/Proofs/BlockDecomposition/CohomologyVanishingShortExactStatement.lean}{GlobalTheory.T.cohomology_vanishing_short_exact}
\leanref{Development/WeilBound/CoreStatements/PadicHodgeComparison/EtaleCohomologyComparison/CechPairings/MemBoundariesZeroIff.lean}{CrystallineSemistableComparison.P.les}

\begin{lemma}\label{gal-realization-thickening-level-one-serre}
Let $\Lambda$ be a ring and $\mathcal S$ a class of $\Lambda$-modules that contains the zero modules and contains $M_2$
whenever $M_1\to M_2\to M_3$ is an exact sequence of $\Lambda$-modules with $M_1,M_3\in\mathcal S$. Let $K$ be a double
complex of $\Lambda$-modules with total complex $\mathrm{Tot}(K)$, as in
Lemma~\ref{gal-realization-thickening-zariski-double}, and for $q\geq0$ let $E^{\bullet,q}$ be the complex $p\mapsto
H^q(K^{p,\bullet})$, with the maps induced by $d_1$. Let $n\geq0$. If $H^p(E^{\bullet,q})\in\mathcal S$ for all
$p,q\geq0$ with $p+q=n$, then $H^n(\mathrm{Tot}(K))\in\mathcal S$.
\end{lemma}
\leanref{Development/WeilBound/CoreStatements/PadicHodgeComparison/EtaleCohomologyComparison/PerfectoidPullbacks/CapOComp.lean}{CrystallineSemistableComparison.Pb.Serre}
\leanref{Development/AutomorphicSpectrum/Proofs/BlockDecomposition/SerreClassOfCechSpectralSequence.lean}{GlobalTheory.CompanionTheory.CechSerre.serre_tot}

\begin{proof}
Applied to $0\to M\to N$ and to $M\to N\to0$, the condition shows that $\mathcal S$ contains every submodule and every
quotient of a module in $\mathcal S$. We use the notation of the proof of
Lemma~\ref{gal-realization-thickening-zariski-double}(2): the components $x_{p,q}$, elements that start in a column, and
$[c]$. For $0\leq a\leq n+1$ let $Z_a$ be the module of the cocycles of $\mathrm{Tot}^n(K)$ that start in column $a$, and
$F_a$ its image in $H^n(\mathrm{Tot}(K))$. Then $F_0=H^n(\mathrm{Tot}(K))$, $F_{a+1}\subset F_a$, and $F_{n+1}=0$, as
$Z_{n+1}=0$. We show $F_a\in\mathcal S$ by descending induction on $a$; the case $a=0$ is the claim.

Let $a\leq n$, $q=n-a$, and suppose $F_{a+1}\in\mathcal S$. Let $z\in Z_a$. By fact (a) of that proof, $d_2z_{a,q}=0$. The
component of $Dz=0$ in $K^{a+1,q}$ is $d_1z_{a,q}+(-1)^{a+1}d_2z_{a+1,q-1}$, or $d_1z_{a,q}$ if $q=0$, so $d_1z_{a,q}$ is
a coboundary of the column $K^{a+1,\bullet}$, and the class of $z_{a,q}$ in $H^q(K^{a,\bullet})$ is a cocycle of
$E^{\bullet,q}$. Let $\psi(z)\in H^a(E^{\bullet,q})$ be its class; $\psi$ is linear.

If $\psi(z)=0$, the class of $z$ lies in $F_{a+1}$. Indeed, the class of $z_{a,q}$ is then the image under $d_1$ of a
class in $H^q(K^{a-1,\bullet})$, or $0$ if $a=0$: there are $y\in K^{a-1,q}$ with $d_2y=0$, $y=0$ if $a=0$, and $u\in
K^{a,q-1}$, $u=0$ if $q=0$, with $z_{a,q}=d_1y+d_2u$. Let $w=[y]+(-1)^a[u]$, in degree $n-1$. Then
$Dw=[d_1y]+[d_2u]+(-1)^a[d_1u]$, whose components in the columns $p\leq a$ are those of $[z_{a,q}]$. So $z-Dw$ starts in
column $a+1$, and it has the class of $z$.

The map $Z_a\to F_a/F_{a+1}$ is surjective and vanishes on the kernel of $\psi$, so $F_a/F_{a+1}$ is a quotient of
$Z_a/\ker\psi$, which $\psi$ embeds in $H^a(E^{\bullet,q})\in\mathcal S$. So $F_a/F_{a+1}\in\mathcal S$, and the exact
sequence $F_{a+1}\to F_a\to F_a/F_{a+1}$ gives $F_a\in\mathcal S$.
% Lean: serre_tot, through serre_Ff (the descending induction, with Ff_top_subsingleton for F_{n+1} = 0; Zf and Ff are
% Z_a and F_a), psi_exists (psi is leadE, through dv_lead, dh_lead and leadC_mem; the kernel statement is lower_vert
% and lower) and serre_step; serre_of_injective and serre_of_surjective are the closure under submodules and quotients.
\end{proof}
\leanref{Development/AutomorphicSpectrum/Proofs/BlockDecomposition/SerreClassOfCechSpectralSequence.lean}{GlobalTheory.CompanionTheory.CechSerre.psi_exists}
\leanref{Development/AutomorphicSpectrum/Proofs/BlockDecomposition/SerreClassOfCechSpectralSequence.lean}{GlobalTheory.CompanionTheory.CechSerre.serre_step}

\begin{lemma}\label{gal-realization-thickening-level-one-cech}
Let $\mathcal C$ be a category, $\Lambda$ a ring, $F$ a functor from $\mathcal C$ to $\Lambda$-modules, and
$(\mathcal U_\alpha)_{\alpha\in A}$ a family of cosieves of $\mathcal C$ (Lemma~\ref{gal-realization-thickening-cosieves})
such that every object of $\mathcal C$ lies in some $\mathcal U_\alpha$. For $i=(i_0,\ldots,i_p)\in A^{p+1}$ let
$\mathcal U_i=\bigcap_k\mathcal U_{i_k}$, and for $0\leq k\leq p$ let $i\circ\partial_k\in A^p$ be $i$ with $i_k$ omitted.
For $q\geq0$ let $\check C^\bullet_q$ be the complex with $\check C^p_q=\prod_{i\in A^{p+1}}H^q(\mathcal U_i,F)$ and
differential $(dx)_i=\sum_{k=0}^{p+1}(-1)^k\operatorname{res}(x_{i\circ\partial_k})$ for $i\in A^{p+2}$, where
$\operatorname{res}$ is the restriction from $\mathcal U_{i\circ\partial_k}$ to $\mathcal U_i$. Let $\mathcal S$ be a class
of $\Lambda$-modules as in Lemma~\ref{gal-realization-thickening-level-one-serre} and $n\geq0$. If
$H^p(\check C^\bullet_q)\in\mathcal S$ for all $p,q\geq0$ with $p+q=n$, then $H^n(\mathcal C,F)\in\mathcal S$.
\end{lemma}
\leanref{Development/AutomorphicSpectrum/Proofs/BlockDecomposition/SerreClassOfCechSpectralSequenceStatement.lean}{GlobalTheory.T.serre_class_of_cech_spectral_sequence}
\leanref{Development/WeilBound/CoreStatements/PadicHodgeComparison/EtaleCohomologyComparison/PerfectoidPullbacks/CapOComp.lean}{CrystallineSemistableComparison.Pb.cechH}
\leanref{Development/AutomorphicSpectrum/Proofs/BlockDecomposition/SerreClassOfUpwardClosedCover.lean}{GlobalTheory.CompanionTheory.CosieveRel.serre_union}

\begin{proof}
\emph{The double complex.} As $\mathcal U_i$ is a cosieve, its chains are the chains of $\mathcal C$ that start in it. Let
$K^{p,q}=\prod_{i\in A^{p+1}}C^q(\mathcal U_i,F)$: an element $x$ assigns to every $i\in A^{p+1}$ and every chain
$s=(c_0\to\cdots\to c_q)$ of $\mathcal U_i$ an element $x(i,s)\in F(c_q)$. Let $d_2$ be the product of the coboundaries
of the $C^\bullet(\mathcal U_i,F)$, and $(d_1x)(i,s)=\sum_{k=0}^{p+1}(-1)^kx(i\circ\partial_k,s)$ for $i\in A^{p+2}$; a
chain of $\mathcal U_i$ is a chain of $\mathcal U_{i\circ\partial_k}$. The maps $d_1$ and $d_2$ commute, so $K$ is a
double complex.

\emph{The rows.} Let $\epsilon^q:C^q(\mathcal C,F)\to K^{0,q}$ send $f$ to $(\alpha,s)\mapsto f(s)$. Then $d_1\epsilon^q=0$
and $d_2\epsilon^q=\epsilon^{q+1}d$. For a chain $s$ starting at $c_0$, let $A_s$ be the set of the $\alpha$ with
$c_0\in\mathcal U_\alpha$, which is nonempty; $s$ is a chain of $\mathcal U_i$ if and only if $i\in A_s^{p+1}$. So the
row $0\to C^q(\mathcal C,F)\to K^{0,q}\to K^{1,q}\to\cdots$ is the product over the chains $s$ of the complexes
$0\to F(c_q)\to\prod_{A_s}F(c_q)\to\prod_{A_s^2}F(c_q)\to\cdots$. Choose $\beta\in A_s$. The maps $h$ that send $x$ to
$(i_0,\ldots,i_{p-1})\mapsto x(\beta,i_0,\ldots,i_{p-1})$, and to $x(\beta)$ in degree $0$, satisfy $h\epsilon=1$ and
$hd_1+d_1h=1$ in degrees $p\geq0$, where $d_1h$ means $\epsilon h$ in degree $0$: the terms with $\beta$ in front cancel
in pairs. So these complexes and the rows are exact, and by Lemma~\ref{gal-realization-thickening-zariski-double}(1) with $s=1$ the
$\epsilon^q$ induce isomorphisms $H^n(\mathcal C,F)\to H^n(\mathrm{Tot}(K))$.

\emph{The columns.} The column $K^{p,\bullet}$ is the product of the complexes $C^\bullet(\mathcal U_i,F)$, so
$H^q(K^{p,\bullet})=\prod_iH^q(\mathcal U_i,F)=\check C^p_q$, and $d_1$, the alternating sum of restrictions of cochains,
induces the differential of $\check C^\bullet_q$. So $E^{\bullet,q}=\check C^\bullet_q$ in
Lemma~\ref{gal-realization-thickening-level-one-serre}, which gives $H^n(\mathrm{Tot}(K))\in\mathcal S$. As $\mathcal S$ is
closed under isomorphisms, $H^n(\mathcal C,F)\in\mathcal S$.
% Lean: serre_cech. The double complex is dbl U hU F, the Cech double complex (CH.cechCos) of idxπ U hU, from the
% functor T ↦ {α : U α T} to the point; tup_bij identifies it with the products over i. The rows are P.cech_desc,
% through P.cech_row (the homotopy prepends a chosen index, ext p hp) and P.dc_row; const_H identifies the cohomology
% of the elements of the point with H(C, F). The columns are colCmp_bij (pi_H_bij, tup_bij) with col_comm, and
% serre_tot concludes. The archive applies it through CosieveRel.serre_union, to the cosieves of a full subcategory,
% with cech_cohomology_full_subcategory_iso identifying the Cech complexes.
\end{proof}
\leanref{Development/AutomorphicSpectrum/Proofs/BlockDecomposition/SerreClassOfCechSpectralSequence.lean}{GlobalTheory.CompanionTheory.CechSerre.serre_cech}
\leanref{Development/WeilBound/CoreStatements/PadicHodgeComparison/EtaleCohomologyComparison/CechPairings/CechDesc.lean}{CrystallineSemistableComparison.P.cech_desc}
\leanref{Development/WeilBound/CoreStatements/PadicHodgeComparison/EtaleCohomologyComparison/CechPairings/CechDesc.lean}{CrystallineSemistableComparison.P.cech_row}
\leanref{Development/AutomorphicSpectrum/Proofs/BlockDecomposition/CechCohomologyFullSubcategoryIsoStatement.lean}{GlobalTheory.T.cech_cohomology_full_subcategory_iso}

\begin{definition}\label{gal-realization-thickening-level-one-graded-filtration}
Let $\Phi$, given by $A$, $J$ and $\gamma$, be a divided power functor on a category $\mathcal C$ over a ring $R$
(Definition~\ref{gal-realization-thickening-torsor}). Maps are composed as functions.
\begin{enumerate}
\item A \emph{filtration} of $\Phi$ is a sequence $(\mathfrak f_r)_{r\in\mathbf Z}$ of families of ideals of $\Phi$ such
  that $\mathfrak f_r(T)=A(T)$ for $r\leq0$, $\mathfrak f_1=J$, $\mathfrak f_{r'}(T)\subset\mathfrak f_r(T)$ for $r\leq
  r'$, $\mathfrak f_r(T)\mathfrak f_{r'}(T)\subset\mathfrak f_{r+r'}(T)$ for $r,r'\geq0$, and $\gamma_k(J(T))\subset
  \mathfrak f_k(T)$ for $k\geq0$. Its \emph{graded pieces} are the functors $\mathrm{gr}_r=\mathfrak f_r/\mathfrak
  f_{r+1}$, $r\in\mathbf Z$.
\item A torsor datum $(D,\tau)$ for $\Phi$ is \emph{graded constant} for a filtration $(\mathfrak f_r)$ if
  $A(\sigma)(x)-A(\sigma')(x)\in\mathfrak f_{q+1}(T)$ for all $q\geq0$, $x\in\mathfrak f_q(D)$ and maps
  $\sigma,\sigma':D\to T$.
\item Let $(D,\tau)$ be a torsor datum of dimension $d$ for $\Phi$, with maps $\sigma_T:D\to T$ chosen for the objects
  $T$. For a map $\varphi:T\to T'$ and $1\leq k\leq d$ let $a_k(\varphi)=\tau_k(\varphi\sigma_T,\sigma_{T'})\in J(T')$.
  For $j\geq0$, a sequence $I=(I_1,\ldots,I_j)$ in $\{1,\ldots,d\}$ and $x\in A(D)$, let $\omega_I(x)$ be the
  $j$-cochain of $\mathcal C$ with values in $A$ that sends a chain $c_0\to\cdots\to c_j$ to
  \[
   \prod_{m=1}^jA(c_m\to c_j)\bigl(a_{I_m}(c_{m-1}\to c_m)\bigr)\cdot A(\sigma_{c_j})(x).
  \]
\item Let $(\mathfrak f_r)$ be a filtration of $\Phi$ and $r\in\mathbf Z$. For a $j$-cochain $z$ of $\mathcal C$ with
  values in $A$ whose values lie in $\mathfrak f_r$ and whose coboundary has values in $\mathfrak f_{r+1}$, let
  $[z]_r\in H^j(\mathcal C,\mathrm{gr}_r)$ be the class of the cochain whose value at every chain is the class of the
  value of $z$. This cochain is a cocycle, as the values of its coboundary are the classes of those of the coboundary of
  $z$.
\end{enumerate}
\end{definition}
\leanref{Development/WeilBound/CoreStatements/PadicHodgeComparison/EtaleCohomologyComparison/PerfectoidPullbacks/DpZeroRight.lean}{CrystallineSemistableComparison.Pb.IsFil}
\leanref{Development/WeilBound/CoreStatements/PadicHodgeComparison/EtaleCohomologyComparison/PerfectoidPullbacks/DpZeroRight.lean}{CrystallineSemistableComparison.Pb.gr}
\leanref{Development/WeilBound/CoreStatements/PadicHodgeComparison/EtaleCohomologyComparison/PerfectoidPullbacks/DpZeroRight.lean}{CrystallineSemistableComparison.Pb.Frame.GrConst}
\leanref{Development/WeilBound/CoreStatements/PadicHodgeComparison/EtaleCohomologyComparison/PerfectoidPullbacks/DpZeroRight.lean}{CrystallineSemistableComparison.Pb.Frame.a}
\leanref{Development/WeilBound/CoreStatements/PadicHodgeComparison/EtaleCohomologyComparison/PerfectoidPullbacks/DpZeroRight.lean}{CrystallineSemistableComparison.Pb.Frame.form}
\leanref{Development/WeilBound/CoreStatements/PadicHodgeComparison/EtaleCohomologyComparison/PerfectoidPullbacks/DpZeroRight.lean}{CrystallineSemistableComparison.Pb.Rep}

\begin{lemma}\label{gal-realization-thickening-level-one-graded-forms}
Let $\Phi$, given by $A$, $J$ and $\gamma$, be a divided power functor on a category $\mathcal C$ over a ring, $(\mathfrak
f_r)$ a filtration of $\Phi$, $(D,\tau)$ a torsor datum of dimension $d$ for $\Phi$ that is graded constant for
$(\mathfrak f_r)$, and $\sigma_T$, $a_k$, $\omega_I$ and $[z]_r$ as in
Definition~\ref{gal-realization-thickening-level-one-graded-filtration}.
\begin{enumerate}
\item $a_k(\psi\varphi)=A(\psi)(a_k(\varphi))+a_k(\psi)$ for maps $\varphi:T\to T'$ and $\psi:T'\to T''$ and $1\leq k\leq
  d$.
\item $[z+z']_r=[z]_r+[z']_r$, and $[z]_r=0$ if the values of $z$ lie in $\mathfrak f_{r+1}$.
\item Let $j,q\geq0$, $I$ a sequence of length $j$ in $\{1,\ldots,d\}$ and $x\in\mathfrak f_q(D)$. The values of
  $\omega_I(x)$ lie in $\mathfrak f_{j+q}$, and its coboundary has values in $\mathfrak f_{j+q+1}$; so
  $[\omega_I(x)]_{j+q}$ is defined. If $x\in\mathfrak f_{q+1}(D)$, it is $0$.
\end{enumerate}
\end{lemma}
\leanref{Development/AutomorphicSpectrum/Proofs/BlockDecomposition/DFormMem.lean}{GlobalTheory.CompanionTheory.CocProd.a_comp}
\leanref{Development/AutomorphicSpectrum/Proofs/BlockDecomposition/RepUnique.lean}{GlobalTheory.CompanionTheory.ConeLocal.rep_add}
\leanref{Development/AutomorphicSpectrum/Proofs/BlockDecomposition/RepUnique.lean}{GlobalTheory.CompanionTheory.ConeLocal.rep_zero}
\leanref{Development/AutomorphicSpectrum/Proofs/BlockDecomposition/FrameFormRespectsFiltration.lean}{GlobalTheory.CompanionTheory.FormMem.form_mem}
\leanref{Development/AutomorphicSpectrum/Proofs/BlockDecomposition/DFormMem.lean}{GlobalTheory.CompanionTheory.CocProd.d_form_mem}
\leanref{Development/AutomorphicSpectrum/Proofs/BlockDecomposition/DFormMem.lean}{GlobalTheory.CompanionTheory.CocProd.rep_form}

\begin{proof}
A product of $n$ elements of $J(T)=\mathfrak f_1(T)$ lies in $\mathfrak f_n(T)$, by induction on $n$, as $\mathfrak
f_0(T)=A(T)$ and $\mathfrak f_1(T)\mathfrak f_n(T)\subset\mathfrak f_{n+1}(T)$.

(1) By the identities of Definition~\ref{gal-realization-thickening-torsor}(2),
\[
 a_k(\psi\varphi)=\tau_k(\psi\varphi\sigma_T,\psi\sigma_{T'})+\tau_k(\psi\sigma_{T'},\sigma_{T''})=
 A(\psi)\bigl(\tau_k(\varphi\sigma_T,\sigma_{T'})\bigr)+a_k(\psi).
\]

(2) The cochains of the classes of the values of $z$ and $z'$ add, and the one for $z$ vanishes if the values of $z$ lie
in $\mathfrak f_{r+1}$.

(3) Let $c_0\to\cdots\to c_j$ be a chain. Each $a_{I_m}(c_{m-1}\to c_m)$ lies in $J(c_m)$, and maps carry $J$ into $J$, so
the product in $\omega_I(x)(c_0\to\cdots\to c_j)$ lies in $\mathfrak f_j(c_j)$; and $A(\sigma_{c_j})(x)\in\mathfrak
f_q(c_j)$. So the value lies in $\mathfrak f_{j+q}(c_j)$. If $x\in\mathfrak f_{q+1}(D)$, the values lie in $\mathfrak
f_{j+q+1}$ for the same reason, and $[\omega_I(x)]_{j+q}=0$ by (2).

For the coboundary let $s=(c_0\to\cdots\to c_{j+1})$, and carry all elements to $A(c_{j+1})$ along the chain. For $0\leq
p\leq p'\leq j+1$ let $g_k(p,p')$ be the image of $a_k(c_p\to c_{p'})$; by (1), $g_k(p,p'')=g_k(p,p')+g_k(p',p'')$ for
$p\leq p'\leq p''$. For $1\leq m\leq j$ put $u_m=g_{I_m}(m-1,m)$ and $v_m=g_{I_m}(m,m+1)$, and for $0\leq n\leq j$ let
$P_n=\prod_{m\leq n}u_m\prod_{m>n}v_m$. Let $Y=A(\sigma_{c_{j+1}})(x)$ and $W=A(c_j\to c_{j+1})(A(\sigma_{c_j})(x))$. The
face of $s$ that omits $c_i$ is the chain through the $c_{e_i(m)}$, $0\leq m\leq j$, where $e_i$ is the increasing map
from $\{0,\ldots,j\}$ to $\{0,\ldots,j+1\}$ that omits $i$. The image of the value of $\omega_I(x)$ at it is $T_iY$ if
$i\leq j$, as this face ends at $c_{j+1}$, and $T_{j+1}W$ if $i=j+1$, where $T_i=\prod_mg_{I_m}(e_i(m-1),e_i(m))$. So
$T_0=\prod_mv_m=P_0$ and $T_{j+1}=\prod_mu_m=P_j$. For $1\leq i\leq j$ the factors of $T_i$ are $u_m$ for $m<i$, $v_m$ for
$m>i$, and $g_{I_i}(i-1,i+1)=u_i+v_i$ for $m=i$, so $T_i=P_i+P_{i-1}$. The sum telescopes, and the value of the
coboundary of $\omega_I(x)$ at $s$ is
\[
 \Bigl(P_0+\sum_{i=1}^j(-1)^i(P_i+P_{i-1})\Bigr)Y+(-1)^{j+1}P_jW=(-1)^jP_j(Y-W).
\]
Here $P_j\in\mathfrak f_j(c_{j+1})$, and $Y-W=A(\sigma_{c_{j+1}})(x)-A((c_j\to c_{j+1})\sigma_{c_j})(x)$ lies in
$\mathfrak f_{q+1}(c_{j+1})$, as $(D,\tau)$ is graded constant. So the value lies in $\mathfrak f_{j+q+1}(c_{j+1})$.
% Lean: (1) is CocProd.a_comp. (2) is ConeLocal.rep_add and rep_zero, with rep_unique for the uniqueness of the class.
% (3): the values are FormMem.form_mem (prod_fil_mem); the coboundary is CocProd.d_form_mem, with ga for g_k(p,p'),
% ga_add, u, v, P, the face values T and W through face_form, hT0, hTl and hTi for T_0, T_{j+1} and T_i, the
% telescoping identity tele, and GrConst for Y - W; rep_form forms the class.
\end{proof}
\leanref{Development/AutomorphicSpectrum/Proofs/BlockDecomposition/DFormMem.lean}{GlobalTheory.CompanionTheory.CocProd.ga}
\leanref{Development/AutomorphicSpectrum/Proofs/BlockDecomposition/DFormMem.lean}{GlobalTheory.CompanionTheory.CocProd.ga_add}
\leanref{Development/AutomorphicSpectrum/Proofs/BlockDecomposition/DFormMem.lean}{GlobalTheory.CompanionTheory.CocProd.face_form}
\leanref{Development/AutomorphicSpectrum/Proofs/BlockDecomposition/DFormMem.lean}{GlobalTheory.CompanionTheory.CocProd.tele}
\leanref{Development/AutomorphicSpectrum/Proofs/BlockDecomposition/DFormMem.lean}{GlobalTheory.CompanionTheory.CocProd.prod_mem}
\leanref{Development/AutomorphicSpectrum/Proofs/BlockDecomposition/FrameFormRespectsFiltration.lean}{GlobalTheory.CompanionTheory.FormMem.prod_fil_mem}
\leanref{Development/AutomorphicSpectrum/Proofs/BlockDecomposition/RepUnique.lean}{GlobalTheory.CompanionTheory.ConeLocal.rep_unique}

\begin{lemma}\label{gal-realization-thickening-level-one-graded-cut}
Let $\Phi$, given by $A$, $J$ and $\gamma$, be a divided power functor on a category $\mathcal C$ over a ring,
$(\mathfrak f_r)$ a filtration of $\Phi$, $(D,\tau)$ a torsor datum of dimension $d+1$ for $\Phi$ that is graded constant
for $(\mathfrak f_r)$, with maps $\sigma_T$ and the $a_k$ and $\omega_I$ of
Definition~\ref{gal-realization-thickening-level-one-graded-filtration}, and $P$ a divided power polynomial structure on
$\Phi$, given by $j_T$ and $X_T$. Form $\mathcal C_D$, $\mathcal C^0_D$, $U^0$, $\Phi^0$, $\tau^0$, $P^0$ and $P_D$
(Definition~\ref{gal-realization-thickening-torsor-cut}) and $S_\sigma$
(Lemma~\ref{gal-realization-thickening-torsor-cut-adjoint}(2)), so that $P_DS_\sigma(T)=(P(T),\sigma_T^+)$; for a map
$\varphi:T\to T'$ write $P_D(\varphi):P(T)\to P(T')$ for the map $P_DS_\sigma(\varphi)$. For a sequence
$I=(I_1,\ldots,I_j)$ in $\{1,\ldots,d\}$ let $I+1=(I_1+1,\ldots,I_j+1)$ and $(1,I+1)=(1,I_1+1,\ldots,I_j+1)$.
\begin{enumerate}
\item The families $\mathfrak f^0_r=\mathfrak f_r\circ U^0$ form a filtration of $\Phi^0$, and the torsor datum
  $(D,\mathrm{id}_D)$ of dimension $d$ with $\tau^0$ is graded constant for it. For every object $(T,s)$ of $\mathcal
  C^0_D$, $s$ is a map $(D,\mathrm{id}_D)\to(T,s)$ of $\mathcal C^0_D$; let $a^0_k$ and $\omega^0_I$ be formed with these
  maps. If $P$ respects $(\mathfrak f_r)$ in degree $r$, then $P^0$ respects $(\mathfrak f^0_r)$ in degree $r$.
\item Let $j,q\geq0$, $I$ a sequence of length $j$ in $\{1,\ldots,d\}$, $x\in\mathfrak f_q(D)$ and $s=(c_0\to\cdots\to
  c_j)$ a chain of $\mathcal C$. Then $\omega^0_I(x)(P_DS_\sigma(s))-A(j_{c_j})(\omega_{I+1}(x)(s))\in\mathfrak
  f_{j+q+1}(P(c_j))$.
\item Let $k,q\geq0$, $r=k+q+1$, $I$ a sequence of length $k$ in $\{1,\ldots,d\}$ and $x\in\mathfrak f_q(D)$, and write
  $\mathrm{gr}_r=\mathfrak f_r/\mathfrak f_{r+1}$. For a chain $s=(c_0\to\cdots\to c_k)$ of $\mathcal C$, the element
  \[
   y(s)=\bigl(X_{c_k}+A(j_{c_k})(a_1(c_0\to c_k))\bigr)\cdot A(j_{c_k})\bigl(\omega_{I+1}(x)(s)\bigr)
  \]
  lies in $\mathfrak f_r(P(c_k))$. Let $\bar y$ be the $k$-cochain of $\mathcal C$ with values in $\mathrm{gr}_r\circ
  U^0P_DS_\sigma$ whose value at $s$ is the class of $y(s)$, and $\bar w$ the $(k+1)$-cochain of $\mathcal C$ with values
  in $\mathrm{gr}_r$ whose value at a chain $s=(c_0\to\cdots\to c_{k+1})$ is the class of $\omega_{(1,I+1)}(x)(s)$. Then
  $(d\bar y)(s)$ is the image of $-\bar w(s)$ under the map $\mathrm{gr}_r(c_{k+1})\to\mathrm{gr}_r(P(c_{k+1}))$ induced
  by $A(j_{c_{k+1}})$, for every chain $s=(c_0\to\cdots\to c_{k+1})$.
\end{enumerate}
\end{lemma}
\leanref{Development/WeilBound/CoreStatements/PadicHodgeComparison/EtaleCohomologyComparison/PerfectoidPullbacks/LocalisationLadders/CriterionTruncation.lean}{CrystallineSemistableComparison.Pb.Cut.isFil_1}
\leanref{Development/WeilBound/CoreStatements/PadicHodgeComparison/EtaleCohomologyComparison/PerfectoidPullbacks/LocalisationLadders/CriterionTruncation.lean}{CrystallineSemistableComparison.Pb.Cut.grConst_1}
\leanref{Development/WeilBound/CoreStatements/PadicHodgeComparison/EtaleCohomologyComparison/PerfectoidPullbacks/LocalisationLadders/CriterionTruncation.lean}{CrystallineSemistableComparison.Pb.Cut.can}
\leanref{Development/AutomorphicSpectrum/Proofs/BlockDecomposition/FrameFormRestrictionCocycle.lean}{GlobalTheory.CompanionTheory.GrPull.form_sub_mem}
\leanref{Development/AutomorphicSpectrum/Proofs/BlockDecomposition/FrameFormCoboundaryRelation.lean}{GlobalTheory.CompanionTheory.GrDelta.gr_delta}

\begin{proof}
(1) The conditions on a filtration and on a graded constant torsor datum concern one object, or maps of $\mathcal C$ out
of $D$, at a time, and the objects and maps of $\mathcal C^0_D$ are objects and maps of $\mathcal C$ through $U^0$; a
map $(D,\mathrm{id}_D)\to(T,s)$ of $\mathcal C^0_D$ is a map $D\to T$ of $\mathcal C$, and $\mathfrak
f^0_q(D,\mathrm{id}_D)=\mathfrak f_q(D)$. By Lemma~\ref{gal-realization-thickening-torsor-cut-adjoint}(1),
$(D,\mathrm{id}_D)$ with $\tau^0$ is a torsor datum of dimension $d$ for $\Phi^0$, and $P^0$ respects $(\mathfrak f^0_r)$
in degree $r$ if $P$ respects $(\mathfrak f_r)$ in degree $r$. The map $s:D\to T$ is a map $(D,\mathrm{id}_D)\to(T,s)$
of $\mathcal C^0_D$, as $\tau_1(s,s)=0$ (proof of Lemma~\ref{gal-realization-thickening-torsor-cut-adjoint}).

(2) Let $\varphi:T\to T'$ be a map of $\mathcal C$ and $\psi=P_D(\varphi)$, so that $\psi j_T=j_{T'}\varphi$. For $1\leq
k\leq d$, $\tau_{k+1}(j_T\sigma_T,\sigma_T^+)=0$ and $\tau_{k+1}(j_{T'}\sigma_{T'},\sigma_{T'}^+)=0$
(Definition~\ref{gal-realization-thickening-torsor-cut}(3)). So, by the identities of a torsor datum, applied to the maps
$\psi\sigma_T^+$, $\psi j_T\sigma_T=j_{T'}\varphi\sigma_T$, $j_{T'}\sigma_{T'}$ and $\sigma_{T'}^+$,
\[
 \tau_{k+1}(\psi\sigma_T^+,\sigma_{T'}^+)=-A(\psi)\bigl(\tau_{k+1}(j_T\sigma_T,\sigma_T^+)\bigr)+\tau_{k+1}(j_{T'}
 \varphi\sigma_T,j_{T'}\sigma_{T'})+\tau_{k+1}(j_{T'}\sigma_{T'},\sigma_{T'}^+)=A(j_{T'})(a_{k+1}(\varphi)).
\]
As the map chosen at $P_DS_\sigma(T)$ is $\sigma_T^+$, this says $a^0_k(P_D(\varphi))=A(j_{T'})(a_{k+1}(\varphi))$. For the
chain $s$, the $m$-th factor of the product in $\omega^0_I(x)(P_DS_\sigma(s))$ is therefore
$A(P_D(c_m\to c_j))(A(j_{c_m})(a_{I_m+1}(c_{m-1}\to c_m)))=A(j_{c_j})(A(c_m\to c_j)(a_{I_m+1}(c_{m-1}\to c_m)))$. So the
product is $A(j_{c_j})(\Pi)$, where $\Pi$ is the product in $\omega_{I+1}(x)(s)$, and
\[
 \omega^0_I(x)(P_DS_\sigma(s))-A(j_{c_j})(\omega_{I+1}(x)(s))=A(j_{c_j})(\Pi)\cdot\bigl(A(\sigma_{c_j}^+)(x)-A(j_{c_j}
 \sigma_{c_j})(x)\bigr).
\]
Here $A(j_{c_j})(\Pi)\in\mathfrak f_j(P(c_j))$, a product of $j$ elements of $J$, and the second factor lies in $\mathfrak
f_{q+1}(P(c_j))$, as $(D,\tau)$ is graded constant.

(3) The first factor of $y(s)$ lies in $J(P(c_k))=\mathfrak f_1(P(c_k))$, as $X_{c_k}\in J(P(c_k))$, and the second in
$\mathfrak f_{k+q}(P(c_k))$ by Lemma~\ref{gal-realization-thickening-level-one-graded-forms}(3); so $y(s)\in\mathfrak
f_{k+q+1}(P(c_k))$.

Let $\varphi:T\to T'$ be a map of $\mathcal C$ and $\psi=P_D(\varphi)$. Then $A(\psi)(A(j_T)(u))=A(j_{T'})(A(\varphi)(u))$
for $u\in A(T)$, and for a map $\chi:T_0\to T$, by Definition~\ref{gal-realization-thickening-torsor-cut}(3) and
Lemma~\ref{gal-realization-thickening-level-one-graded-forms}(1),
\[
 A(\psi)\bigl(X_T+A(j_T)(a_1(\chi))\bigr)=X_{T'}+A(j_{T'})(a_1(\varphi))+A(j_{T'})(A(\varphi)(a_1(\chi)))=
 X_{T'}+A(j_{T'})(a_1(\varphi\chi)).
\]
Let $s=(c_0\to\cdots\to c_{k+1})$, write $j$ for $j_{c_{k+1}}$, and carry elements to $A(P(c_{k+1}))$ by the maps
$A(P_D(\ldots))$ and to $A(c_{k+1})$ by the maps of the chain. For $0\leq p\leq k+1$ let $B_p=X_{c_{k+1}}+A(j)(a_1(c_p\to
c_{k+1}))$, and for $0\leq i\leq k+1$ let $t_i$ be the image of the value of $\omega_{I+1}(x)$ at the face $s_i$ of $s$
that omits $c_i$. The face $s_i$ starts at $c_1$ if $i=0$ and at $c_0$ otherwise, so by the two identities the image of
$y(s_i)$ is $B_1A(j)(t_0)$ if $i=0$ and $B_0A(j)(t_i)$ if $i\geq1$. By
Lemma~\ref{gal-realization-thickening-level-one-graded-forms}(1), $B_0=B_1+A(j)(g)$, where $g$ is the image of
$a_1(c_0\to c_1)$. So the value of the coboundary of $y$ at $s$ is
\[
 B_1A(j)(t_0)+\sum_{i=1}^{k+1}(-1)^iB_0A(j)(t_i)=B_0\,A(j)\Bigl(\sum_{i=0}^{k+1}(-1)^it_i\Bigr)-A(j)(g\,t_0).
\]
The sum $\sum_i(-1)^it_i$ is the value of the coboundary of $\omega_{I+1}(x)$ at $s$, which lies in $\mathfrak
f_{k+q+1}(c_{k+1})$ by Lemma~\ref{gal-realization-thickening-level-one-graded-forms}(3), and $B_0\in\mathfrak
f_1(P(c_{k+1}))$; so the first term lies in $\mathfrak f_{k+q+2}(P(c_{k+1}))$. And $g\,t_0=\omega_{(1,I+1)}(x)(s)$: the
first factor of the product in $\omega_{(1,I+1)}(x)(s)$ is $g$, and its other factors and $A(\sigma_{c_{k+1}})(x)$ are
those of the value of $\omega_{I+1}(x)$ at $s_0=(c_1\to\cdots\to c_{k+1})$, which ends at $c_{k+1}$. So the value of the
coboundary of $y$ at $s$ differs from $-A(j)(\omega_{(1,I+1)}(x)(s))$ by an element of $\mathfrak f_{r+1}(P(c_{k+1}))$;
as the coboundary of $\bar y$ is formed from the values of $y$ by the same maps, passing to classes in
$\mathrm{gr}_r(P(c_{k+1}))$ gives the claim.
% Lean: (1) is Cut.isFil_1, Cut.grConst_1, Cut.E1 and Cut.can, with Cut.criterion_truncation for P^0 = Cut.Q1.
% (2) is GrPull.form_sub_mem: the identity for tau_{k+1} is tau_sec_Pmap (through Cut.τ_sec and Cut.sec_spec), the
% factors are factor, and P_0 and the last factor are f1 and f2, with GrConst for u = sec and v = σ ≫ j. (3) is
% GrDelta.gr_delta: the membership is Bv_mem with FormMem.form_mem; the two identities are Pmap_j_apply and Pmap_shift;
% B_p is Bv; the images of y(s_i) are term_val; B_0 = B_1 + j(g) is Bv_zero; the rearrangement is alt_sum; the first
% term is d_form_mem; and g t_0 = omega_{(1,I+1)}(x)(s) is form_cons.
\end{proof}
\leanref{Development/AutomorphicSpectrum/Proofs/BlockDecomposition/FrameFormRestrictionCocycle.lean}{GlobalTheory.CompanionTheory.GrPull.tau_sec_Pmap}
\leanref{Development/AutomorphicSpectrum/Proofs/BlockDecomposition/FrameFormRestrictionCocycle.lean}{GlobalTheory.CompanionTheory.GrPull.factor}
\leanref{Development/AutomorphicSpectrum/Proofs/BlockDecomposition/FrameFormCoboundaryRelation.lean}{GlobalTheory.CompanionTheory.GrDelta.Pmap_j_apply}
\leanref{Development/AutomorphicSpectrum/Proofs/BlockDecomposition/FrameFormCoboundaryRelation.lean}{GlobalTheory.CompanionTheory.GrDelta.Pmap_shift}
\leanref{Development/AutomorphicSpectrum/Proofs/BlockDecomposition/FrameFormCoboundaryRelation.lean}{GlobalTheory.CompanionTheory.GrDelta.Bv}
\leanref{Development/AutomorphicSpectrum/Proofs/BlockDecomposition/FrameFormCoboundaryRelation.lean}{GlobalTheory.CompanionTheory.GrDelta.Bv_mem}
\leanref{Development/AutomorphicSpectrum/Proofs/BlockDecomposition/FrameFormCoboundaryRelation.lean}{GlobalTheory.CompanionTheory.GrDelta.term_val}
\leanref{Development/AutomorphicSpectrum/Proofs/BlockDecomposition/FrameFormCoboundaryRelation.lean}{GlobalTheory.CompanionTheory.GrDelta.Bv_zero}
\leanref{Development/AutomorphicSpectrum/Proofs/BlockDecomposition/FrameFormCoboundaryRelation.lean}{GlobalTheory.CompanionTheory.GrDelta.alt_sum}
\leanref{Development/AutomorphicSpectrum/Proofs/BlockDecomposition/FrameFormCoboundaryRelation.lean}{GlobalTheory.CompanionTheory.GrDelta.form_cons}

\begin{lemma}\label{gal-realization-thickening-level-one-graded-sequence}
Let $\Phi$, $R$, $(\mathfrak f_r)$, $(D,\tau)$ of dimension $d+1$, $\sigma_T$, $P$, $\mathcal C^0_D$, $U^0$, $P_D$,
$S_\sigma$, $(\mathfrak f^0_r)$, $\omega_I$ and $\omega^0_I$ be as in
Lemma~\ref{gal-realization-thickening-level-one-graded-cut}, and let $\mathrm{gr}_r$ and $\mathrm{gr}^0_r=\mathrm{gr}_r
\circ U^0$ be the graded pieces of $(\mathfrak f_r)$ and $(\mathfrak f^0_r)$. Let $r\in\mathbf Z$, and suppose that $P$
respects $(\mathfrak f_{r'})_{r'}$ in every degree $r'\leq r+1$. There are $R$-linear maps
\[
 \pi_n:H^n(\mathcal C,\mathrm{gr}_r)\to H^n(\mathcal C^0_D,\mathrm{gr}^0_r),\quad\varepsilon_n:H^n(\mathcal
 C^0_D,\mathrm{gr}^0_{r-1})\to H^{n+1}(\mathcal C,\mathrm{gr}_r),\quad\theta_n:H^n(\mathcal C^0_D,\mathrm{gr}^0_r)\to
 H^n(\mathcal C^0_D,\mathrm{gr}^0_{r-1})
\]
for $n\geq0$, with the following properties.
\begin{enumerate}
\item $\pi_0$ is injective, and for every $n$ the kernel of $\pi_{n+1}$ is the image of $\varepsilon_n$, the kernel of
  $\varepsilon_n$ is the image of $\theta_n$, and $\theta_n\pi_n=0$.
\item Let $j,q\geq0$ with $r=j+q$, $I$ a sequence of length $j$ in $\{1,\ldots,d\}$ and $x\in\mathfrak f_q(D)$. Then
  $\pi_j([\omega_{I+1}(x)]_r)=[\omega^0_I(x)]_r$.
\item Let $k,q\geq0$ with $r=k+1+q$, $I$ a sequence of length $k$ in $\{1,\ldots,d\}$ and $x\in\mathfrak f_q(D)$. Then
  $\varepsilon_k([\omega^0_I(x)]_{r-1})=-[\omega_{(1,I+1)}(x)]_r$.
\end{enumerate}
The classes in (2) and (3) are defined by Lemma~\ref{gal-realization-thickening-level-one-graded-forms}(3), for the cut
data through Lemma~\ref{gal-realization-thickening-level-one-graded-cut}(1).
\end{lemma}
\leanref{Development/AutomorphicSpectrum/Proofs/BlockDecomposition/LocalGradedVanishing.lean}{GlobalTheory.CompanionTheory.LocGrStep.Pack}
\leanref{Development/AutomorphicSpectrum/Proofs/BlockDecomposition/LocalGradedVanishing.lean}{GlobalTheory.CompanionTheory.LocGrStep.pack}

\begin{proof}
By Lemma~\ref{gal-realization-thickening-torsor-cut-adjoint}(2) and (3), $P_DS_\sigma:\mathcal C\to\mathcal C^0_D$ is
left adjoint to $U^0$, as in the proof of Lemma~\ref{gal-realization-thickening-torsor-vanishing}. So by
Lemma~\ref{gal-realization-thickening-torsor-adjoint}(1), for every $r'$ and $n$ the map $\rho_{r'}:H^n(\mathcal
C^0_D,\mathrm{gr}^0_{r'})\to H^n(\mathcal C,\mathrm{gr}_{r'}\circ U^0P_DS_\sigma)$ induced by $P_DS_\sigma$ and the
identity maps is bijective.

Apply Lemma~\ref{gal-realization-thickening-torsor-exact} to $\mathfrak a_{r'}=\mathfrak f_{r'}$ and $\mathfrak
b_{r'}=\mathfrak f_{r'+1}$. Here $\mathfrak b_{r'}\subset\mathfrak a_{r'}$, $P$ respects $(\mathfrak a_{r'})$ in degrees
$r$ and $r-1$, and it respects $(\mathfrak b_{r'})$ in degrees $r$ and $r-1$, as this is to respect $(\mathfrak f_{r'})$
in degrees $r+1$ and $r$. Composed with $S_\sigma$, that lemma gives an exact sequence of functors on $\mathcal C$
\[
 0\to\mathrm{gr}_r\xrightarrow{\alpha}\mathrm{gr}_r\circ U^0P_DS_\sigma\xrightarrow{\beta}\mathrm{gr}_{r-1}\circ
 U^0P_DS_\sigma\to0,
\]
where $\alpha$ is induced by the maps $j_T$, and $\beta$ sends the class of $\sum_kA(j_T)(c_k)\gamma_k(X_T)$ to the class
of $\sum_kA(j_T)(c_{k+1})\gamma_k(X_T)$. It is exact on cochains in every degree, as it is exact at every object, so it
has a long exact sequence of cohomology, in which $\alpha_*$ is injective on $H^0$. Its connecting map
$\delta_n:H^n(\mathcal C,\mathrm{gr}_{r-1}\circ U^0P_DS_\sigma)\to H^{n+1}(\mathcal C,\mathrm{gr}_r)$ sends the class of a
cocycle $\beta(y)$, for an $n$-cochain $y$, to the class of the cocycle $z$ with $\alpha(z)=dy$. Let
\[
 \pi_n=\rho_r^{-1}\alpha_*,\qquad\varepsilon_n=\delta_n\rho_{r-1},\qquad\theta_n=\rho_{r-1}^{-1}\beta_*\rho_r.
\]

(1) These are the exactness of the long exact sequence at $H^0(\mathcal C,\mathrm{gr}_r)$, at $H^{n+1}(\mathcal
C,\mathrm{gr}_r)$ and at $H^n(\mathcal C,\mathrm{gr}_{r-1}\circ U^0P_DS_\sigma)$, and $\beta_*\alpha_*=0$, transported by
the bijections $\rho_r$ and $\rho_{r-1}$.

(2) The class $\rho_r([\omega^0_I(x)]_r)$ is the class of the cocycle whose value at a chain $s=(c_0\to\cdots\to c_j)$ of
$\mathcal C$ is the class of $\omega^0_I(x)(P_DS_\sigma(s))$, and $\alpha_*([\omega_{I+1}(x)]_r)$ is the class of the
cocycle whose value at $s$ is the class of $A(j_{c_j})(\omega_{I+1}(x)(s))$, both in $\mathrm{gr}_r(P(c_j))$. By
Lemma~\ref{gal-realization-thickening-level-one-graded-cut}(2) these cocycles are equal. So
$\alpha_*([\omega_{I+1}(x)]_r)=\rho_r([\omega^0_I(x)]_r)$, which is the claim.

(3) Here $r-1=k+q$. As in (2), $\rho_{r-1}([\omega^0_I(x)]_{r-1})$ is the class of the cocycle $g$ whose value at a chain
$s=(c_0\to\cdots\to c_k)$ is the class of $A(j_{c_k})(\omega_{I+1}(x)(s))$ in $\mathrm{gr}_{r-1}(P(c_k))$. Let $y$ be the
cochain of Lemma~\ref{gal-realization-thickening-level-one-graded-cut}(3), and $\bar y$ the $k$-cochain with values in
$\mathrm{gr}_r\circ U^0P_DS_\sigma$ whose value at $s$ is the class of $y(s)$. As $\gamma_0(X_T)=1$ and
$\gamma_1(X_T)=X_T$, $y(s)=A(j_{c_k})(ab)\gamma_0(X_{c_k})+A(j_{c_k})(b)\gamma_1(X_{c_k})$ with $a=a_1(c_0\to c_k)$ and
$b=\omega_{I+1}(x)(s)$; so the value of $\beta(\bar y)$ at $s$ is the class of $A(j_{c_k})(b)$, and $\beta(\bar y)=g$. By
Lemma~\ref{gal-realization-thickening-level-one-graded-cut}(3), $d\bar y=-\alpha(w)$, where $w$ is the cocycle whose
value at a chain $s$ is the class of $\omega_{(1,I+1)}(x)(s)$, of class $[\omega_{(1,I+1)}(x)]_r$. So
$\varepsilon_k([\omega^0_I(x)]_{r-1})=\delta_k([\beta(\bar y)])=[-w]=-[\omega_{(1,I+1)}(x)]_r$.
% Lean: LocGrStep.pack, which builds Pack. rho_r is rho with rho_bij (P.adj for Cut.adjPσ); the exact sequence is
% PolySes.ses for a = Fil and b = Fil (· + 1) (crit_shift), whiskered by Pt.Sσ, with alpha the whiskering of Cut.jσ;
% P.ses_cochains and P.les give δ and the exactness. (2) is pull_form, through GrPull.gr_pull. (3) is delta_form:
% rep_form gives g, gr_pull identifies it with rho of the cut cocycle, lift_form builds the cochain y (val_ev for the
% expansion and the property of β from PolySes.ses), and GrDelta.gr_delta gives d y = -alpha(w).
\end{proof}
\leanref{Development/AutomorphicSpectrum/Proofs/BlockDecomposition/LocalGradedVanishing.lean}{GlobalTheory.CompanionTheory.LocGrStep.rho}
\leanref{Development/AutomorphicSpectrum/Proofs/BlockDecomposition/LocalGradedVanishing.lean}{GlobalTheory.CompanionTheory.LocGrStep.rho_bij}
\leanref{Development/AutomorphicSpectrum/Proofs/BlockDecomposition/LocalGradedVanishing.lean}{GlobalTheory.CompanionTheory.LocGrStep.alpha}
\leanref{Development/AutomorphicSpectrum/Proofs/BlockDecomposition/LocalGradedVanishing.lean}{GlobalTheory.CompanionTheory.LocGrStep.pull_form}
\leanref{Development/AutomorphicSpectrum/Proofs/BlockDecomposition/LocalGradedVanishing.lean}{GlobalTheory.CompanionTheory.LocGrStep.val_ev}
\leanref{Development/AutomorphicSpectrum/Proofs/BlockDecomposition/LocalGradedVanishing.lean}{GlobalTheory.CompanionTheory.LocGrStep.lift_form}
\leanref{Development/AutomorphicSpectrum/Proofs/BlockDecomposition/LocalGradedVanishing.lean}{GlobalTheory.CompanionTheory.LocGrStep.delta_form}
\leanref{Development/AutomorphicSpectrum/Proofs/BlockDecomposition/FrameFormRestrictionCocycle.lean}{GlobalTheory.CompanionTheory.GrPull.gr_pull}
\leanref{Development/WeilBound/CoreStatements/PadicHodgeComparison/EtaleCohomologyComparison/PerfectoidPullbacks/LocalisationLadders/CriterionTruncation.lean}{CrystallineSemistableComparison.Pb.Cut.Pσ}
\leanref{Development/WeilBound/CoreStatements/PadicHodgeComparison/EtaleCohomologyComparison/PerfectoidPullbacks/LocalisationLadders/CriterionTruncation.lean}{CrystallineSemistableComparison.Pb.Cut.jσ}

\begin{lemma}\label{gal-realization-thickening-level-one-graded}
Let $\Phi$, given by $A$, $J$ and $\gamma$, be a divided power functor on a category $\mathcal C$ over a ring $R$
(Definition~\ref{gal-realization-thickening-torsor}). Let $(\mathfrak f_r)_{r\in\mathbf Z}$ be families of ideals of
$\Phi$ such that $\mathfrak f_r(T)=A(T)$ for $r\leq0$, $\mathfrak f_1=J$, $\mathfrak f_{r'}(T)\subset\mathfrak f_r(T)$ for
$r\leq r'$, $\mathfrak f_r(T)\mathfrak f_{r'}(T)\subset\mathfrak f_{r+r'}(T)$ for $r,r'\geq0$, and
$\gamma_k(J(T))\subset\mathfrak f_k(T)$ for $k\geq0$. Let $(D,\tau)$ be a torsor datum of dimension $d$ for $\Phi$, and
choose maps $\sigma_T:D\to T$ for the objects $T$. Suppose that $A(\sigma)(x)-A(\sigma')(x)\in\mathfrak f_{q+1}(T)$ for
all $q\geq0$, $x\in\mathfrak f_q(D)$ and maps $\sigma,\sigma':D\to T$. For a map $\varphi:T\to T'$ and $1\leq k\leq d$ let
$a_k(\varphi)=\tau_k(\varphi\sigma_T,\sigma_{T'})\in J(T')$. For $j\geq0$, a strictly increasing sequence
$I=(I_1,\ldots,I_j)$ in $\{1,\ldots,d\}$ and $x\in A(D)$, let $\omega_I(x)$ be the $j$-cochain of $\mathcal C$ with values
in $A$ that sends a chain $c_0\to\cdots\to c_j$ to
\[
 \prod_{m=1}^jA(c_m\to c_j)\bigl(a_{I_m}(c_{m-1}\to c_m)\bigr)\cdot A(\sigma_{c_j})(x).
\]
Let $r\in\mathbf Z$, let $P$ be a divided power polynomial structure on $\Phi$ that respects $(\mathfrak f_{r'})_{r'}$ in
every degree $r'\leq r+1$, and let $j\geq0$.
\begin{enumerate}
\item If $j>r$, then $H^j(\mathcal C,\mathfrak f_r/\mathfrak f_{r+1})=0$.
\item Let $r=j+q$ with $q\geq0$. Every class of $H^j(\mathcal C,\mathfrak f_r/\mathfrak f_{r+1})$ is the class of a cocycle
  whose value at every chain is the class of the value of $\sum_I\omega_I(x_I)$, the sum over the strictly increasing
  $I$ of length $j$, for some elements $x_I\in\mathfrak f_q(D)$. If elements $x_I\in\mathfrak f_q(D)$ give in this way a
  cocycle whose class is $0$, then $x_I\in\mathfrak f_{q+1}(D)$ for every $I$.
\end{enumerate}
\end{lemma}
\leanref{Development/WeilBound/CoreStatements/PadicHodgeComparison/EtaleCohomologyComparison/PerfectoidPullbacks/GlobalSerreRelations/GlobSerreSt.lean}{CrystallineSemistableComparison.Pb.KH.LocGr}
\leanref{Development/WeilBound/CoreStatements/PadicHodgeComparison/EtaleCohomologyComparison/PerfectoidPullbacks/DpZeroRight.lean}{CrystallineSemistableComparison.Pb.IsFil}
\leanref{Development/WeilBound/CoreStatements/PadicHodgeComparison/EtaleCohomologyComparison/PerfectoidPullbacks/DpZeroRight.lean}{CrystallineSemistableComparison.Pb.Frame.GrConst}
\leanref{Development/WeilBound/CoreStatements/PadicHodgeComparison/EtaleCohomologyComparison/PerfectoidPullbacks/DpZeroRight.lean}{CrystallineSemistableComparison.Pb.Frame.form}
\leanref{Development/AutomorphicSpectrum/Proofs/BlockDecomposition/LocalGradedVanishing.lean}{GlobalTheory.CompanionTheory.LocGrStep.loc_gr}

\begin{proof}
In the terms of Definition~\ref{gal-realization-thickening-level-one-graded-filtration}, $(\mathfrak f_r)$ is a
filtration of $\Phi$, $(D,\tau)$ is graded constant for it, and $a_k$ and $\omega_I$ are those of the definition. With
$\mathrm{gr}_r=\mathfrak f_r/\mathfrak f_{r+1}$, (2) says that every class of $H^j(\mathcal C,\mathrm{gr}_r)$ is
$[\sum_I\omega_I(x_I)]_r$ for some $x_I\in\mathfrak f_q(D)$, and that $[\sum_I\omega_I(x_I)]_r=0$ with $x_I\in\mathfrak
f_q(D)$ implies $x_I\in\mathfrak f_{q+1}(D)$ for every $I$. If $r<0$, then $\mathfrak f_{r+1}=A$, so $\mathrm{gr}_r=0$ and
$H^0(\mathcal C,\mathrm{gr}_r)=0$. The proof is by induction on $d$, for all data as in the statement.

\emph{The case $d=0$.} As in the proof of Lemma~\ref{gal-realization-thickening-torsor-base-change}, $D$ is an initial
object, and $\sigma_T$ is its map to $T$. By Lemma~\ref{gal-realization-thickening-torsor-adjoint}(2),
$H^j(\mathcal C,\mathrm{gr}_r)=0$ for $j\geq1$, and every class of $H^0(\mathcal C,\mathrm{gr}_r)$ is $\varepsilon(z)$
for some $z\in\mathrm{gr}_r(D)$. With the remark above for $j=0$ and $r<0$, this gives (1). In (2) with $j\geq1$ there is
no strictly increasing sequence of length $j$ in the empty set, so the sums are empty, and both claims hold as
$H^j(\mathcal C,\mathrm{gr}_r)=0$. Let $j=0$, so that $r=q$; the only sequence is the empty one, and $\omega_\emptyset(x)$
sends $c$ to $A(\sigma_c)(x)$. If $z$ is the class of $y\in\mathfrak f_q(D)$, the value of $\varepsilon(z)$ at $c$ is the
class of $A(\sigma_c)(y)$, so $\varepsilon(z)=[\omega_\emptyset(y)]_q$. If $[\omega_\emptyset(x)]_q=0$, the cocycle with
the classes of the $A(\sigma_c)(x)$ as values is $0$, as there are no coboundaries in degree $0$; its value at $D$, where
$\sigma_D=\mathrm{id}_D$, is the class of $x$, so $x\in\mathfrak f_{q+1}(D)$.

\emph{The step.} Let the torsor datum have dimension $d+1$, and suppose the claims hold for all data of dimension $d$.
Form the data of Lemma~\ref{gal-realization-thickening-level-one-graded-cut}. By its part (1), $\Phi^0$ with $(\mathfrak
f^0_{r'})$, the torsor datum $(D,\mathrm{id}_D)$ with $\tau^0$ of dimension $d$, the maps $s:(D,\mathrm{id}_D)\to(T,s)$
and $P^0$ are data as in the statement, on $\mathcal C^0_D$, with $P^0$ respecting $(\mathfrak f^0_{r'})$ in every degree
$r'\leq r+1$. So (1) and (2) hold for them in every degree $j$, for $\mathrm{gr}^0_r$ and for $\mathrm{gr}^0_{r-1}$; call
these (1$^0$) and (2$^0$). Let $\pi_n$, $\varepsilon_n$ and $\theta_n$ be the maps of
Lemma~\ref{gal-realization-thickening-level-one-graded-sequence}. By
Lemma~\ref{gal-realization-thickening-level-one-graded-forms}(2) and the additivity of $\pi_j$ and $\varepsilon_k$, its
parts (2) and (3) also hold for sums $\sum_I$ over the strictly increasing sequences $I$ of a given length, with
coefficients $x_I\in\mathfrak f_q(D)$. A strictly increasing sequence of length $j$ in $\{1,\ldots,d+1\}$ that does not
start with $1$ is $I+1$ for exactly one strictly increasing sequence $I$ of length $j$ in $\{1,\ldots,d\}$, and one that
starts with $1$ is $(1,I'+1)$ for exactly one strictly increasing sequence $I'$ of length $j-1$ in $\{1,\ldots,d\}$; the
empty sequence is $\emptyset+1$.

(1) Let $j>r$. If $j=0$, then $r<0$ and $H^0(\mathcal C,\mathrm{gr}_r)=0$. Let $j=k+1$ and $c\in H^{k+1}(\mathcal
C,\mathrm{gr}_r)$. By (1$^0$), $H^{k+1}(\mathcal C^0_D,\mathrm{gr}^0_r)=0$, so $\pi_{k+1}(c)=0$, and $c=\varepsilon_k(m)$
for some $m\in H^k(\mathcal C^0_D,\mathrm{gr}^0_{r-1})$ by Lemma~\ref{gal-realization-thickening-level-one-graded-sequence}(1).
By (1$^0$), as $k>r-1$, $m=0$; so $c=0$.

(2) Let $r=j+q$. Two facts hold for families $(x_I)$ indexed by the strictly increasing sequences $I$ of length $j$ in
$\{1,\ldots,d\}$.

(a) For every $c\in H^j(\mathcal C,\mathrm{gr}_r)$ there are $x_I\in\mathfrak f_q(D)$ with
$\pi_j(c-[\sum_I\omega_{I+1}(x_I)]_r)=0$. Indeed, by (2$^0$), $\pi_j(c)=[\sum_I\omega^0_I(x_I)]_r$ for some
$x_I\in\mathfrak f_q(D)$, and this is $\pi_j([\sum_I\omega_{I+1}(x_I)]_r)$ by
Lemma~\ref{gal-realization-thickening-level-one-graded-sequence}(2).

(b) If $x_I\in\mathfrak f_q(D)$ and $\pi_j([\sum_I\omega_{I+1}(x_I)]_r)=0$, then $x_I\in\mathfrak f_{q+1}(D)$ for every
$I$, and $[\sum_I\omega_{I+1}(x_I)]_r=0$. Indeed, $[\sum_I\omega^0_I(x_I)]_r=0$ by
Lemma~\ref{gal-realization-thickening-level-one-graded-sequence}(2), so $x_I\in\mathfrak f_{q+1}(D)$ by (2$^0$), and then
$[\omega_{I+1}(x_I)]_r=0$ by Lemma~\ref{gal-realization-thickening-level-one-graded-forms}(3).

Let $j=0$, so that $r=q$. For $c\in H^0(\mathcal C,\mathrm{gr}_r)$, (a) and the injectivity of $\pi_0$ give
$c=[\omega_\emptyset(x_\emptyset)]_r$ with $x_\emptyset\in\mathfrak f_q(D)$; and if $[\omega_\emptyset(x)]_r=0$ with
$x\in\mathfrak f_q(D)$, then $x\in\mathfrak f_{q+1}(D)$ by (b).

Let $j=k+1$, so that $r=k+1+q$ and $r-1=k+q$. By the remark on sequences, a family indexed by the strictly increasing
sequences of length $k+1$ in $\{1,\ldots,d+1\}$ is given by a family $(x_I)$ at the sequences $I+1$, with $I$ of length
$k+1$, and a family $(x'_{I'})$ at the sequences $(1,I'+1)$, with $I'$ of length $k$, both in $\{1,\ldots,d\}$.

Existence. Let $c\in H^{k+1}(\mathcal C,\mathrm{gr}_r)$. By (a) there are $x_I\in\mathfrak f_q(D)$ with $\pi_{k+1}(c-b)=0$
for $b=[\sum_I\omega_{I+1}(x_I)]_r$. By Lemma~\ref{gal-realization-thickening-level-one-graded-sequence}(1),
$c-b=\varepsilon_k(m)$ for some $m\in H^k(\mathcal C^0_D,\mathrm{gr}^0_{r-1})$. By (2$^0$),
$-m=[\sum_{I'}\omega^0_{I'}(x'_{I'})]_{r-1}$ for some $x'_{I'}\in\mathfrak f_q(D)$, and by
Lemma~\ref{gal-realization-thickening-level-one-graded-sequence}(3), $c-b=-\varepsilon_k(-m)=[\sum_{I'}
\omega_{(1,I'+1)}(x'_{I'})]_r$. So $c=[\sum_I\omega_{I+1}(x_I)+\sum_{I'}\omega_{(1,I'+1)}(x'_{I'})]_r$, a sum over all
strictly increasing sequences of length $k+1$ in $\{1,\ldots,d+1\}$ with coefficients in $\mathfrak f_q(D)$.

Uniqueness. Let $x_I,x'_{I'}\in\mathfrak f_q(D)$ with $b+w=0$, where $b=[\sum_I\omega_{I+1}(x_I)]_r$ and
$w=[\sum_{I'}\omega_{(1,I'+1)}(x'_{I'})]_r$, and let $a=[\sum_{I'}\omega^0_{I'}(x'_{I'})]_{r-1}$. By
Lemma~\ref{gal-realization-thickening-level-one-graded-sequence}(3), $\varepsilon_k(a)=-w=b$, so $\pi_{k+1}(b)=0$ by part
(1) of that lemma. By (b), $x_I\in\mathfrak f_{q+1}(D)$ for every $I$, and $b=0$. So $\varepsilon_k(a)=0$, and
$a=\theta_k(a_2)$ for some $a_2\in H^k(\mathcal C^0_D,\mathrm{gr}^0_r)$. As $r=k+(q+1)$, (2$^0$) gives
$a_2=[\sum_{I'}\omega^0_{I'}(x''_{I'})]_r$ for some $x''_{I'}\in\mathfrak f_{q+1}(D)$, which is
$\pi_k([\sum_{I'}\omega_{I'+1}(x''_{I'})]_r)$ by Lemma~\ref{gal-realization-thickening-level-one-graded-sequence}(2). So
$a=\theta_k\pi_k([\sum_{I'}\omega_{I'+1}(x''_{I'})]_r)=0$, and (2$^0$) for $r-1=k+q$ gives $x'_{I'}\in\mathfrak
f_{q+1}(D)$ for every $I'$.
% Lean: KH.LocGr, proved by LocGrStep.loc_gr = LocGrStep.main, by induction on d. The case d = 0 is LocGr.gr_zero
% (Frame.initial and P.initial; H_succ_subsingleton for j >= 1 and H_zero_subsingleton for r < 0; Idx (n+1) 0 is
% empty). The step is LocGrStep.step: the induction hypothesis hh for Cut.Q1, Cut.isFil_1, Cut.E1, Cut.can,
% Cut.grConst_1 and Cut.criterion_truncation, at r and r - 1 (h0, h1); the maps are pack; the sums are sum_pull and
% sum_delta (rep_sum, rep_unique, CocProd.rep_form); the sequences are ia, ib, split_bij, sum_split,
% index_zero_extend_bijective and sum_zero. (1) is van_step, (a) is exA, (b) is unA, existence is ex_step and uniqueness is
% un_step.
\end{proof}
\leanref{Development/AutomorphicSpectrum/Proofs/BlockDecomposition/GradedCohomologyOfFrame.lean}{GlobalTheory.CompanionTheory.LocGr.gr_zero}
\leanref{Development/AutomorphicSpectrum/Proofs/BlockDecomposition/GradedCohomologyOfFrame.lean}{GlobalTheory.CompanionTheory.LocGr.H_succ_subsingleton}
\leanref{Development/AutomorphicSpectrum/Proofs/BlockDecomposition/GradedCohomologyOfFrame.lean}{GlobalTheory.CompanionTheory.LocGr.H_zero_subsingleton}
\leanref{Development/AutomorphicSpectrum/Proofs/BlockDecomposition/LocalGradedVanishing.lean}{GlobalTheory.CompanionTheory.LocGrStep.main}
\leanref{Development/AutomorphicSpectrum/Proofs/BlockDecomposition/LocalGradedVanishing.lean}{GlobalTheory.CompanionTheory.LocGrStep.step}
\leanref{Development/AutomorphicSpectrum/Proofs/BlockDecomposition/LocalGradedVanishing.lean}{GlobalTheory.CompanionTheory.LocGrStep.van_step}
\leanref{Development/AutomorphicSpectrum/Proofs/BlockDecomposition/LocalGradedVanishing.lean}{GlobalTheory.CompanionTheory.LocGrStep.exA}
\leanref{Development/AutomorphicSpectrum/Proofs/BlockDecomposition/LocalGradedVanishing.lean}{GlobalTheory.CompanionTheory.LocGrStep.unA}
\leanref{Development/AutomorphicSpectrum/Proofs/BlockDecomposition/LocalGradedVanishing.lean}{GlobalTheory.CompanionTheory.LocGrStep.ex_step}
\leanref{Development/AutomorphicSpectrum/Proofs/BlockDecomposition/LocalGradedVanishing.lean}{GlobalTheory.CompanionTheory.LocGrStep.un_step}
\leanref{Development/AutomorphicSpectrum/Proofs/BlockDecomposition/LocalGradedVanishing.lean}{GlobalTheory.CompanionTheory.LocGrStep.sum_pull}
\leanref{Development/AutomorphicSpectrum/Proofs/BlockDecomposition/LocalGradedVanishing.lean}{GlobalTheory.CompanionTheory.LocGrStep.sum_delta}
\leanref{Development/AutomorphicSpectrum/Proofs/BlockDecomposition/LocalGradedVanishing.lean}{GlobalTheory.CompanionTheory.LocGrStep.split_bij}
\leanref{Development/AutomorphicSpectrum/Proofs/BlockDecomposition/LocalGradedVanishing.lean}{GlobalTheory.CompanionTheory.LocGrStep.index_zero_extend_bijective}

\begin{definition}\label{gal-realization-thickening-level-one-fibre}
Let $(R,I,\gamma,p)$, $\mathcal O$, $X$ and $\kappa$ be as in Definition~\ref{gal-realization-thickening-nilpotent}, with
$I=0$, so that $\kappa$ is a ring map $\mathcal O\to R$.
\begin{enumerate}
\item Let $X_R=X\times_{\operatorname{Spec}\mathcal O}\operatorname{Spec}R$, formed with $\kappa$, with the projection
  $\pi:X_R\to X$ and the structure map $X_R\to\operatorname{Spec}R$. For an open $V\subset X$ let $V_R=\pi^{-1}(V)$, and for
  $f\in\Gamma(V,\mathcal O_X)$ write $f$ also for its image in $\Gamma(V_R,\mathcal O_{X_R})$.
\item Let $T=(B,i,J,\delta,y)$ be a thickening that lies over an open $V\subset X$. The morphism $y$ and the morphism
  $\operatorname{Spec}B/J\to\operatorname{Spec}R$ induced by $i$ define a morphism $\operatorname{Spec}B/J\to X_R$, which
  factors through $V_R$; let $\mathrm{ev}_T:\Gamma(V_R,\mathcal O_{X_R})\to B/J$ be the ring map it induces. A \emph{lift}
  of $f\in\Gamma(V_R,\mathcal O_{X_R})$ at $T$ is an element of $B$ with class $\mathrm{ev}_T(f)$.
\item Let $\mathcal Q$ be a full subcategory of $\mathrm{Th}(X/R)$ whose thickenings lie over $V$, $r\geq0$ and $f\in
  \Gamma(V_R,\mathcal O_{X_R})$. \emph{Multiplication by} $f$ on $H^j(\mathcal Q,\mathcal O^{[r]}/\mathcal O^{[r+1]})$ is the
  map induced by the maps $J_T^{[r]}/J_T^{[r+1]}\to J_T^{[r]}/J_T^{[r+1]}$, $z\mapsto bz$, where $b$ is a lift of $f$ at
  $T$ (Lemma~\ref{gal-realization-thickening-level-one-forms}(2)).
\item Let moreover $t_1,\ldots,t_d\in\Gamma(V_R,\mathcal O_{X_R})$, and let $u_k(T)$ be a lift of $t_k$ at $T$ for every
  thickening $T$ of $\mathcal Q$. For a family $x=(x_I)$ of elements of $\Gamma(V_R,\mathcal O_{X_R})$, indexed by the
  strictly increasing sequences $I=(I_1,\ldots,I_r)$ in $\{1,\ldots,d\}$, let $\omega_u(x)$ be the $r$-cochain of $\mathcal
  Q$ with values in $\mathcal O^{[r]}/\mathcal O^{[r+1]}$ that sends a chain $c_0\to\cdots\to c_r$ to the class of
  \[
   \sum_I\prod_{m=1}^r\bigl(u_{I_m}(c_m)-u_{I_m}(c_{m-1})\bigr)\,b_I,
  \]
  where $b_I$ is a lift of $x_I$ at $c_r$ and the elements of the rings of the $c_m$ are carried to that of $c_r$ along
  the chain. A class of $H^r(\mathcal Q,\mathcal O^{[r]}/\mathcal O^{[r+1]})$ is \emph{represented by} $x$, for the $t_k$
  and the $u_k$, if it is the class of $\omega_u(x)$.
\end{enumerate}
\end{definition}
\leanref{Development/AutomorphicSpectrum/Proofs/BlockDecomposition/LiftCocycleMem.lean}{GlobalTheory.CompanionTheory.FibGr.XF}
\leanref{Development/AutomorphicSpectrum/Proofs/BlockDecomposition/LiftCocycleMem.lean}{GlobalTheory.CompanionTheory.FibGr.fib}
\leanref{Development/AutomorphicSpectrum/Proofs/BlockDecomposition/LiftCocycleMem.lean}{GlobalTheory.CompanionTheory.FibGr.ev}
\leanref{Development/AutomorphicSpectrum/Proofs/BlockDecomposition/LiftCocycleMem.lean}{GlobalTheory.CompanionTheory.FibGr.lft}
\leanref{Development/AutomorphicSpectrum/Proofs/BlockDecomposition/SubMemOfLift.lean}{GlobalTheory.CompanionTheory.FibGr.fAct}
\leanref{Development/AutomorphicSpectrum/Proofs/BlockDecomposition/LiftCocycleMem.lean}{GlobalTheory.CompanionTheory.FibGr.Rp}

\begin{lemma}\label{gal-realization-thickening-level-one-forms}
Let $(R,I,\gamma,p)$, $\mathcal O$, $X$ and $\kappa$ be as in Definition~\ref{gal-realization-thickening-level-one-fibre},
with $I=0$, $V\subset X$ an open, $\mathcal Q$ a full subcategory of $\mathrm{Th}(X/R)$ whose thickenings lie over $V$, and
$r\geq0$.
\begin{enumerate}
\item For a map $\varphi:T\to T'$ of thickenings over $V$, $\mathrm{ev}_{T'}=\bar\varphi\circ\mathrm{ev}_T$, where
  $\bar\varphi:B_T/J_T\to B_{T'}/J_{T'}$ is induced by $\varphi$. For an open $V'\subset V$ and $T$ over $V'$, $\mathrm{ev}_T$
  on $\Gamma(V_R,\mathcal O_{X_R})$ is the restriction to $V'_R$ followed by $\mathrm{ev}_T$ on $\Gamma(V'_R,\mathcal
  O_{X_R})$. The image of $c\in R$ in $\Gamma(V_R,\mathcal O_{X_R})$ has $i_T(c)$ as a lift.
\item Multiplication by $f\in\Gamma(V_R,\mathcal O_{X_R})$ on $H^j(\mathcal Q,\mathcal O^{[r]}/\mathcal O^{[r+1]})$ is well
  defined, and it makes this group a $\Gamma(V_R,\mathcal O_{X_R})$-module on which the image of $c\in R$ acts as $c$. For
  a full subcategory $\mathcal Q'\subset\mathcal Q$ whose thickenings lie over an open $V'\subset V$, restriction from
  $\mathcal Q$ to $\mathcal Q'$ carries multiplication by $f$ to multiplication by the restriction of $f$ to $V'_R$.
\item In Definition~\ref{gal-realization-thickening-level-one-fibre}(4), the value of $\omega_u(x)$ does not depend on the
  lifts $b_I$, and $\omega_u(x)$ is a cocycle. Let $x$ represent $h$ and $x'$ represent $h'$. Then $x+x'$ represents
  $h+h'$, $(fx_I)_I$ represents $f$ times $h$ for $f\in\Gamma(V_R,\mathcal O_{X_R})$, and for $\mathcal Q'$ and $V'$ as in
  (2), the restrictions of the $x_I$ to $V'_R$ represent the restriction of $h$ to $\mathcal Q'$, for the restrictions of
  the $t_k$ and the same $u_k$.
\end{enumerate}
\end{lemma}
\leanref{Development/AutomorphicSpectrum/Proofs/BlockDecomposition/LiftCocycleMem.lean}{GlobalTheory.CompanionTheory.FibGr.ev_nat}
\leanref{Development/AutomorphicSpectrum/Proofs/BlockDecomposition/LiftCocycleMem.lean}{GlobalTheory.CompanionTheory.FibGr.ev_res}
\leanref{Development/AutomorphicSpectrum/Proofs/BlockDecomposition/LiftCocycleMem.lean}{GlobalTheory.CompanionTheory.FibGr.ev_str}
\leanref{Development/AutomorphicSpectrum/Proofs/BlockDecomposition/SubMemOfLift.lean}{GlobalTheory.CompanionTheory.FibGr.fHom}
\leanref{Development/AutomorphicSpectrum/Proofs/BlockDecomposition/LiftCocycleMem.lean}{GlobalTheory.CompanionTheory.FibGr.exists_rp}
\leanref{Development/AutomorphicSpectrum/Proofs/BlockDecomposition/LiftCocycleMem.lean}{GlobalTheory.CompanionTheory.FibGr.rp_add}
\leanref{Development/AutomorphicSpectrum/Proofs/BlockDecomposition/SubMemOfLift.lean}{GlobalTheory.CompanionTheory.FibGr.rp_fAct}
\leanref{Development/AutomorphicSpectrum/Proofs/BlockDecomposition/LiftCocycleMem.lean}{GlobalTheory.CompanionTheory.FibGr.rp_res}

\begin{proof}
(1) The point of $T'$ is the point of $T$ composed with $\operatorname{Spec}\bar\varphi$, and $\varphi\circ i_T=i_{T'}$,
so the morphism $\operatorname{Spec}B_{T'}/J_{T'}\to X_R$ of $T'$ is that of $T$ composed with
$\operatorname{Spec}\bar\varphi$. The ring maps induced by morphisms into $V_R$ and into $V'_R$ are compatible with
restriction. The morphism of $T$ followed by the structure map of $X_R$ is $\operatorname{Spec}$ of the map $R\to B/J$
induced by $i_T$, so $\mathrm{ev}_T$ sends the image of $c$ to the class of $i_T(c)$.

(2) Two lifts of $f$ at $T$ differ by an element of $J_T$, and $J_TJ_T^{[r]}\subset J_T^{[r+1]}$, so the map $z\mapsto bz$
does not depend on the lift $b$. For $\varphi:T\to T'$ and a lift $b$ at $T$, $\varphi(b)$ is a lift at $T'$ by (1), so
these maps commute with $\varphi$; they form an endomorphism of the functor $\mathcal O^{[r]}/\mathcal O^{[r+1]}$, which
induces the endomorphisms of cohomology. Sums, products and $1$ lift to sums, products and $1$, and $i_T(c)$ lifts the
image of $c$ by (1), where multiplication by $i_T(c)$ is the $R$-module structure. A lift of $f$ at a thickening over
$V'$ is a lift of its restriction to $V'_R$, by (1), and restriction of cochains to $\mathcal Q'$ commutes with
multiplication by lifts.

(3) For a map $c_{m-1}\to c_m$ of $\mathcal Q$, $u_k(c_m)$ and the image of $u_k(c_{m-1})$ are both lifts of $t_k$ at
$c_m$, by (1), so their difference lies in $J_{c_m}$. A product of $r$ elements of $J$ lies in $J^{[r]}$, as
$z_1\cdots z_r=\delta_1(z_1)\cdots\delta_1(z_r)$. So each product in $\omega_u(x)$ lies in $J_{c_r}^{[r]}$, and changing
$b_I$ by an element of $J_{c_r}$ changes the value by an element of $J^{[r]}J\subset J^{[r+1]}$.

For the cocycle condition, fix $I$, let $s=(c_0\to\cdots\to c_{r+1})$, carry everything to the ring of $c_{r+1}$, and
write $v_k(m)$ for the image of $u_k(c_m)$. For $0\leq l\leq r+1$ let $\Pi_l$ be the product
$\prod_{m=1}^r(v_{I_m}(e_l(m))-v_{I_m}(e_l(m-1)))$, where $e_l$ is the increasing map from $\{0,\ldots,r\}$ to
$\{0,\ldots,r+1\}$ that omits $l$; this is the image of the product of the face of $s$ that omits $c_l$. Then
$\sum_l(-1)^l\Pi_l=0$: on the ordered set $\{0,\ldots,r+1\}$, viewed as a category, $\Pi_l$ is the value at the face
that omits $l$ of the cup product of the coboundaries of the $0$-cochains $w_m(l')=v_{I_m}(l')$ with values in this
ring, and a cup product of cocycles is a cocycle. Let $b$ be a lift of $x_I$ at $c_{r+1}$ and $b'$ one at $c_r$. The
faces with $l\leq r$ end at $c_{r+1}$, and the term of the coboundary of $\omega_u(x)$ for $I$ is
\[
 \sum_{l\leq r}(-1)^l\Pi_lb+(-1)^{r+1}\Pi_{r+1}b''=(-1)^r\Pi_{r+1}(b-b''),
\]
where $b''$ is the image of $b'$, another lift of $x_I$ at $c_{r+1}$. As $\Pi_{r+1}\in J^{[r]}$ and $b-b''\in J$, this
lies in $J^{[r+1]}$, so the coboundary of $\omega_u(x)$ vanishes.

Sums of lifts are lifts of sums, so $\omega_u(x+x')=\omega_u(x)+\omega_u(x')$. If $e$ is a lift of $f$ at $c_r$, $eb_I$ is
a lift of $fx_I$, so $\omega_u((fx_I)_I)$ is $\omega_u(x)$ multiplied by lifts of $f$, which represents $f$ times $h$ by
(2). The restriction of $\omega_u(x)$ to the chains of $\mathcal Q'$ is $\omega_u$ of the restricted family, as a lift
of $x_I$ at a thickening over $V'$ is a lift of its restriction, by (1).
% Lean: (1) is ev_nat, ev_res and ev_str. (2) is fAct (mulMor and mulT, with mulT_congr for the independence of the
% lift and evLift_nat for the naturality), fAct_add, fAct_mul, fAct_one, fHom, fAct_str and res_fAct. (3): the
% cochain is Rp, through wz (the product) and lft (a chosen lift, lft_spec), with wz_mem for J^[r]; the cocycle
% condition is GrForm.stCocycle through Coc.d_mem, whose identity alt_sum is the vanishing of the alternating sum of the
% Pi_l, and exists_rp; additivity is rp_add, multiplication is rp_fAct (with evLift_mul_sub), restriction is rp_res.
\end{proof}
\leanref{Development/AutomorphicSpectrum/Proofs/BlockDecomposition/SubMemOfMk.lean}{GlobalTheory.CompanionTheory.GrCech.mulMor}
\leanref{Development/AutomorphicSpectrum/Proofs/BlockDecomposition/GradedFormCocycle.lean}{GlobalTheory.CompanionTheory.GrForm.stCocycle}
\leanref{Development/AutomorphicSpectrum/Proofs/BlockDecomposition/GradedFormCocycle.lean}{GlobalTheory.CompanionTheory.GrForm.Coc.d_mem}

\begin{lemma}\label{gal-realization-thickening-level-one-local}
Let $(R,I,\gamma,p)$, $\mathcal O$, $X$ and $\kappa$ be as in Definition~\ref{gal-realization-thickening-level-one-fibre},
with $I=0$ and $p\neq0$ nilpotent in $R$, and let $r\geq0$.
\begin{enumerate}
\item Let $W\subset X$ be an affine open with coordinates (Definition~\ref{gal-realization-thickening-coordinates}(1))
  and $\mathcal Q=\mathrm{Th}(X/R)_W$. Then $H^j(\mathcal Q,\mathcal O^{[r]}/\mathcal O^{[r+1]})=0$ for $j\neq r$.
\item Let $V\subset X$ be an open with coordinates $t_1,\ldots,t_d$, $W\subset V$ an affine open, $\mathcal
  Q=\mathrm{Th}(X/R)_W$, and $u_k(T)$ lifts of the $t_k$ at the thickenings $T$ over $V$. Every class of $H^r(\mathcal
  Q,\mathcal O^{[r]}/\mathcal O^{[r+1]})$ is represented by exactly one family $(x_I)$ of elements of
  $\Gamma(W_R,\mathcal O_{X_R})$, for the restrictions of the $t_k$ to $W_R$ and the lifts $u_k(T)$ at the thickenings
  $T$ of $\mathcal Q$.
\end{enumerate}
\end{lemma}
\leanref{Development/AutomorphicSpectrum/Proofs/BlockDecomposition/PreimageMono.lean}{GlobalTheory.CompanionTheory.FibGr.off_at}
\leanref{Development/AutomorphicSpectrum/Proofs/BlockDecomposition/PreimageMono.lean}{GlobalTheory.CompanionTheory.FibGr.free_at}
\leanref{Development/AutomorphicSpectrum/Proofs/BlockDecomposition/SubMemOfLift.lean}{GlobalTheory.CompanionTheory.FibGr.Free}

\begin{proof}
In (1) let $t_1,\ldots,t_d$ be the coordinates on $W$, and choose any lifts $u_k(T)$ of the $t_k$ at the thickenings $T$
over $W$. In (2) the restrictions of the $t_k$ to $W$ are coordinates on $W$
(Lemma~\ref{gal-realization-thickening-cover}(1)); write $t_k$ also for them, and use the given lifts at the thickenings
over $W$, which lie over $V$.

\emph{Step 1: a chart.} By Lemma~\ref{gal-realization-thickening-chart}(3) there is a divided power chart $D$ with
$t'_1,\ldots,t'_d$ on $W$ such that the map $\Gamma(W,\mathcal O_X)\to B_D/J_D$ induced by $y_D$ sends each $t_k$ to the
class of $t'_k$. As $I=0$, $J_D=i_D(I)B_D=0$. The ring map $\mathrm{ev}_D:\Gamma(W_R,\mathcal O_{X_R})\to B_D$ is
bijective. Indeed $W_R$ is affine. Let $\theta':B_D\to\Gamma(W_R,\mathcal O_{X_R})$ be the map that the chart property
gives for $B=\Gamma(W_R,\mathcal O_{X_R})$, $s$ the structure map and $z$ the morphism $W_R\to X$. The composite
$\mathrm{ev}_D\circ\theta'$ has the two properties that characterize the identity of $B_D$ in the chart property for
$s=i_D$ and $z=y_D$, so it is the identity. And $\operatorname{Spec}(\theta'\circ\mathrm{ev}_D)$ is the morphism
$W_R\to\operatorname{Spec}B_D\to W_R$ whose components to $X$ and to $\operatorname{Spec}R$ are those of the identity, so
$\theta'\circ\mathrm{ev}_D$ is the identity. Let $\theta=\mathrm{ev}_D^{-1}$; then $\theta(t_k)=t'_k$.

\emph{Step 2: the torsor datum.} By Lemma~\ref{gal-realization-thickening-chart-torsor}(1), $D$ with
$\tau_k(\sigma,\sigma')=\sigma'(t'_k)-\sigma(t'_k)$ is a torsor datum of dimension $d$ for $\Phi_{\mathcal Q}$. For $T$ in
$\mathcal Q$ there is a map $\sigma_T:D\to T$ with $\sigma_T(t'_k)=u_k(T)$ for every $k$: for any map $\sigma_0:D\to T$
(Lemma~\ref{gal-realization-thickening-chart}(1)), $\sigma_0(t'_k)$ is a lift of $t_k$ at $T$ by
Lemma~\ref{gal-realization-thickening-level-one-forms}(1), so $u_k(T)-\sigma_0(t'_k)\in J_T$, and the last claim of
Lemma~\ref{gal-realization-thickening-chart}(1) gives $\sigma_T$. Then $a_k(\varphi)=\sigma_{c'}(t'_k)-\varphi(\sigma_c(t'_k))
=u_k(c')-\varphi(u_k(c))$ for a map $\varphi:c\to c'$, and $\sigma_c(\theta(f))$ is a lift of $f$ at $c$, by
Lemma~\ref{gal-realization-thickening-level-one-forms}(1). So the cochain $\sum_I\omega_I(\theta(x_I))$ of
Lemma~\ref{gal-realization-thickening-level-one-graded} is the sum of
Definition~\ref{gal-realization-thickening-level-one-fibre}(4) for $x$ with the lifts $b_I=\sigma_{c_r}(\theta(x_I))$.

\emph{Step 3: the hypotheses of Lemma~\ref{gal-realization-thickening-level-one-graded}.} The families $\mathfrak
f_r=\mathcal O^{[r]}$ have the properties required there, by the definition of the ideals $J^{[r]}$
(Definition~\ref{gal-realization-thickening-nilpotent}(3)). For $q\geq0$, $x\in J_D^{[q]}$ and $\sigma,\sigma':D\to T$,
$\sigma(x)-\sigma'(x)\in J_T^{[q+1]}$: by Lemma~\ref{gal-realization-thickening-chart}(2), $x$ is a sum of elements
$i_D(a)b$ with $a\in I^{[q]}$, and $\sigma(i_D(a)b)-\sigma'(i_D(a)b)=i_T(a)(\sigma(b)-\sigma'(b))\in J_T^{[q]}J_T$, by
Lemma~\ref{gal-realization-thickening-chart}(1). By Lemma~\ref{gal-realization-thickening-chart-torsor}(3), for
$\mathcal Q=\mathrm{Th}(X/R)_{\{W\}}$, $\Phi_{\mathcal Q}$ has a divided power polynomial structure that respects
$(\mathcal O^{[r']})_{r'}$ in every degree. And $J_D^{[q]}\subset J_D=0$ for $q\geq1$.

(1) If $j>r$, Lemma~\ref{gal-realization-thickening-level-one-graded}(1) applies. If $j<r$, let $q=r-j\geq1$. By
Lemma~\ref{gal-realization-thickening-level-one-graded}(2) every class is represented by a cocycle with values the
classes of $\sum_I\omega_I(x_I)$ with $x_I\in J_D^{[q]}=0$, so it is $0$.

(2) By Lemma~\ref{gal-realization-thickening-level-one-graded}(2) with $j=r$ and $q=0$, and by Step 2, every class is the
class of $\omega_u(x)$ for $x_I=\theta^{-1}(x'_I)$, where $x'_I\in J_D^{[0]}=B_D$ are the elements given there. If $x$ and
$x'$ represent the same class, then $x-x'$ represents $0$ by
Lemma~\ref{gal-realization-thickening-level-one-forms}(3), so $\theta(x_I-x'_I)\in J_D^{[1]}=0$ by
Lemma~\ref{gal-realization-thickening-level-one-graded}(2) and Step 2, and $x=x'$.
% Lean: (1) is off_at, through coordinates_restrict_affine_open (in fin_gen), frame_at with the lifts lft, poly_all'
% and off_zero. (2) is free_at: coordinates_restrict_affine_open, isLift_res, frame_at, poly_all' and free_of_frame,
% whose conclusion is Free (existence and rp_coef for uniqueness). frame_at: Step 1 is exists_chart, ChartMaps.good and
% ev_bij (with c.rep and evalAt_specMap), the vanishing J_D = 0 is c'.J_eq; Step 2 is Chart.frame, het, hσ (through
% good.1 and good.2) and hval; Step 3 is chart_grConst (GrConst, through dividedPowerFiltration_eq_map), isFil_sFil
% and poly_all'. The comparison of the two cochains is form_sub_mem with key in free_of_frame.
\end{proof}
\leanref{Development/AutomorphicSpectrum/Proofs/BlockDecomposition/PreimageMono.lean}{GlobalTheory.CompanionTheory.FibGr.frame_at}
\leanref{Development/AutomorphicSpectrum/Proofs/BlockDecomposition/PreimageMono.lean}{GlobalTheory.CompanionTheory.FibGr.ev_bij}
\leanref{Development/AutomorphicSpectrum/Proofs/BlockDecomposition/SubMemOfLift.lean}{GlobalTheory.CompanionTheory.FibGr.free_of_frame}
\leanref{Development/AutomorphicSpectrum/Proofs/BlockDecomposition/SubMemOfLift.lean}{GlobalTheory.CompanionTheory.FibGr.off_zero}
\leanref{Development/AutomorphicSpectrum/Proofs/BlockDecomposition/SheafFiltrationProperty.lean}{GlobalTheory.CompanionTheory.DPFil.isFil_sFil}

\begin{lemma}\label{gal-realization-thickening-level-one-coherent}
Let $\Lambda$ be a Noetherian ring, $Z$ a scheme proper over $\operatorname{Spec}\Lambda$, and $(V_\alpha)_{\alpha\in A}$
a finite family of opens that cover $Z$, such that $V_i=\bigcap_kV_{i_k}$ is affine for every $i=(i_0,\ldots,i_n)\in
A^{n+1}$, $n\geq0$. Write $\Gamma(V)$ for $\Gamma(V,\mathcal O_Z)$, and for $0\leq b\leq n+1$ let $\partial_b:\{0,\ldots,n\}\to
\{0,\ldots,n+1\}$ be the increasing map that omits $b$. Let $K$ be a cosimplicial $\Lambda$-module, with coface maps
$\delta_b:K^n\to K^{n+1}$, and $H^n(K)$ the cohomology of the complex with differential $\sum_b(-1)^b\delta_b$. Let
$\kappa_\alpha$ be finite sets, $g_{\alpha\beta}(k,k')\in\Gamma(V_\alpha\cap V_\beta)$ for $k\in\kappa_\alpha$ and
$k'\in\kappa_\beta$, and for $0\leq a\leq n$ let $e_{n,a}$ be additive bijections from $K^n$ to the families
$(e_{i,k})$ with $e_{i,k}\in\Gamma(V_i)$ for $i\in A^{n+1}$ and $k\in\kappa_{i_a}$, such that:
\begin{enumerate}
\item[(i)] $e_{n,a}(cx)_{i,k}=c\,e_{n,a}(x)_{i,k}$ for $c\in\Lambda$;
\item[(ii)] $e_{n+1,\partial_b(a)}(\delta_bx)_{i,k}$ is the restriction to $V_i$ of $e_{n,a}(x)_{i\circ\partial_b,k}$;
\item[(iii)] $e_{n,a}(x)_{i,k}=\sum_{k'\in\kappa_{i_{a'}}}g_{i_ai_{a'}}(k,k')\,e_{n,a'}(x)_{i,k'}$ on $V_i$, for $0\leq
  a,a'\leq n$.
\end{enumerate}
Then $H^n(K)$ is a finitely generated $\Lambda$-module for every $n$.
\end{lemma}
\leanref{Development/WeilBound/CoreStatements/PadicHodgeComparison/EtaleCohomologyComparison/PerfectoidPullbacks/CapOComp.lean}{CrystallineSemistableComparison.Pb.VB}
\leanref{Development/WeilBound/CoreStatements/PadicHodgeComparison/EtaleCohomologyComparison/PerfectoidPullbacks/GlobSerre.lean}{CrystallineSemistableComparison.Pb.vb_fin}

\begin{proof}
By (iii) for $n=0$ and the bijectivity of $e_{0,0}$, $g_{\alpha\alpha}$ is the identity matrix on $V_\alpha$, and by (iii)
for $i=(\alpha,\beta,\gamma)$ and the bijectivity of $e_{2,2}$, $g_{\alpha\beta}g_{\beta\gamma}=g_{\alpha\gamma}$ on
$V_{\alpha\beta\gamma}$. So the free modules $\mathcal O_{V_\alpha}^{\kappa_\alpha}$, glued along the matrices
$g_{\alpha\beta}$, form a locally free $\mathcal O_Z$-module $\mathcal E$ of finite rank \cite[Tag 00AK]{Sta}, with
$\mathcal E(V_i)\cong\Gamma(V_i)^{\kappa_{i_a}}$ for every $a$, these identifications differing by the $g$. Through
$e_{n,0}$, $K^n$ is the module $\prod_{i\in A^{n+1}}\mathcal E(V_i)$ of \v{C}ech cochains of $\mathcal E$, $\Lambda$-linearly
by (i), and by (ii) the coface maps are the restrictions, so $H^n(K)$ is the \v{C}ech cohomology of $\mathcal E$ for the
cover $(V_\alpha)$. As the $V_i$ are affine, this is $H^n(Z,\mathcal E)$ \cite[Tag 01XD]{Sta}, which is
$\Gamma(\operatorname{Spec}\Lambda,R^nf_*\mathcal E)$ for $f:Z\to\operatorname{Spec}\Lambda$ \cite[Tag 01XK]{Sta}. And
$R^nf_*\mathcal E$ is coherent, as $f$ is proper \cite[Th\'eor\`eme~3.2.1]{Gro61}, \cite[Tag 02O5]{Sta}.
% Lean: vb_fin = VBGlue.finite_H_cx, the Lean's own proof, through about 76,000 lines of general scheme theory in the
% AbelianVarietyTheory, GroundTheory and ReciprocityTheory libraries (AffineCoverCohomologyFinite.lean and what it uses).
\end{proof}
\leanref{Development/AutomorphicSpectrum/Proofs/BlockDecomposition/AffineCoverCohomologyFinite.lean}{GlobalTheory.CompanionTheory.VBGlue.finite_H_cx}

\begin{lemma}\label{gal-realization-thickening-level-one-bundle}
Let $(R,I,\gamma,p)$, $\mathcal O$, $X$ and $\kappa$ be as in Definition~\ref{gal-realization-thickening-level-one-fibre},
with $I=0$ and $p\neq0$ nilpotent in $R$. Let $(U_\alpha)_{\alpha\in A}$ be a coordinate cover of $X$
(Definition~\ref{gal-realization-thickening-coordinates}(2)), with coordinates $t^\alpha_1,\ldots,t^\alpha_{d_\alpha}$ on
$U_\alpha$, and $r\geq0$. For $i\in A^{n+1}$ let $U_i=\bigcap_kU_{i_k}$ and $\mathcal C_i=\mathrm{Th}(X/R)_{U_i}$. Let
$K$ be the cosimplicial $R$-module with $K^n=\prod_{i\in A^{n+1}}H^r(\mathcal C_i,\mathcal O^{[r]}/\mathcal O^{[r+1]})$
and $(\delta_bx)_i$ the restriction of $x_{i\circ\partial_b}$ from $\mathcal C_{i\circ\partial_b}$ to $\mathcal C_i$. There
are finite sets $\kappa_\alpha$, elements $g_{\alpha\beta}(k,k')$ and bijections $e_{n,a}$ that satisfy (i), (ii) and (iii)
of Lemma~\ref{gal-realization-thickening-level-one-coherent} for $\Lambda=R$, $Z=X_R$, the opens $V_\alpha=U_{\alpha,R}$
and $K$.
\end{lemma}
\leanref{Development/AutomorphicSpectrum/Proofs/BlockDecomposition/PreimageMono.lean}{GlobalTheory.CompanionTheory.FibGr.vb_cech}

\begin{proof}
For each $\alpha$ let $\kappa_\alpha$ be the set of strictly increasing sequences $I=(I_1,\ldots,I_r)$ in
$\{1,\ldots,d_\alpha\}$, choose lifts $u^\alpha_k(T)$ of the $t^\alpha_k$ at the thickenings $T$ over $U_\alpha$, and let
$\omega^\alpha_I\in H^r(\mathcal C_{(\alpha)},\mathcal O^{[r]}/\mathcal O^{[r+1]})$ be the class represented by the family
with entry $1$ at $I$ and $0$ elsewhere (Lemma~\ref{gal-realization-thickening-level-one-forms}(3)).

\emph{Step 1: bases.} Let $W\subset U_\alpha$ be an affine open. Every $h\in H^r(\mathrm{Th}(X/R)_W,\mathcal
O^{[r]}/\mathcal O^{[r+1]})$ is $\sum_If_I\omega^\alpha_I|_W$ for exactly one family $(f_I)$ in $\Gamma(W_R,\mathcal
O_{X_R})$, where $\omega^\alpha_I|_W$ is the restriction and $f_I\omega$ is multiplication by $f_I$. Indeed, by
Lemma~\ref{gal-realization-thickening-level-one-forms}(3) the sum $\sum_If_I\omega^\alpha_I|_W$ is represented by $(f_I)$,
for the restrictions of the $t^\alpha_k$ and the lifts $u^\alpha_k$; and by
Lemma~\ref{gal-realization-thickening-level-one-local}(2), for $V=U_\alpha$, $h$ is represented by exactly one family.

\emph{Step 2: the data.} The thickenings of $\mathcal C_i$ are those over $U_i$
(Lemma~\ref{gal-realization-thickening-over}(2)), and $U_i$ is an affine open of every $U_{i_a}$. Let $e_{n,a}$ send
$x\in K^n$ to the family whose entries at $i$ are the coefficients of $x_i$ in Step 1, for $\alpha=i_a$ and $W=U_i$; it is
bijective by Step 1. For $\alpha,\beta\in A$ let $g_{\alpha\beta}(k,k')$, $k\in\kappa_\alpha$, $k'\in\kappa_\beta$, be the
coefficients of $\omega^\beta_{k'}|_{U_\alpha\cap U_\beta}$ in Step 1 for $\alpha$ and the affine open $W=U_\alpha\cap
U_\beta$.

(i) The image of $c\in R$ acts as $c$ (Lemma~\ref{gal-realization-thickening-level-one-forms}(2)), so the coefficients of
$cx_i$ are $c$ times those of $x_i$. (ii) Restriction from $\mathcal C_{i\circ\partial_b}$ to $\mathcal C_i$ carries
multiplication by $f$ to multiplication by its restriction
(Lemma~\ref{gal-realization-thickening-level-one-forms}(2)) and $\omega^\alpha_k|_{U_{i\circ\partial_b}}$ to
$\omega^\alpha_k|_{U_i}$, and $(i\circ\partial_b)_a=i_{\partial_b(a)}$; so the coefficients of $(\delta_bx)_i$ for
$i_{\partial_b(a)}$ are the restrictions of those of $x_{i\circ\partial_b}$ for $(i\circ\partial_b)_a$. (iii) Restricting
the identities that define $g_{i_ai_{a'}}$ to $U_i$ expresses each $\omega^{i_{a'}}_{k'}|_{U_i}$ through the
$\omega^{i_a}_k|_{U_i}$. Substituting this into $x_i=\sum_{k'}e_{n,a'}(x)_{i,k'}\,\omega^{i_{a'}}_{k'}|_{U_i}$ writes $x_i$
with the coefficients $\sum_{k'}g_{i_ai_{a'}}(k,k')e_{n,a'}(x)_{i,k'}$, which are $e_{n,a}(x)_{i,k}$ by the uniqueness in
Step 1.
% Lean: vb_cech. Step 1 is basis_gen (omega from exists_rp with Pi.single k 1, then rp_sum, rp_fAct, rp_res, rp_class
% and rp_coef) with free_at (hfreeT, hfreeP); the restriction to C_i is through capT_eq and capO_pre. Step 2 and (i) to
% (iii) are CechVB.vb_basis, with act the action fHom, b the restricted omega, hstr = fAct_str, hres = res_fAct,
% hb = res_res and hg from g' (hg'); the affine intersections are U.inter. Size of the part from Definition
% gal-realization-thickening-level-one-fibre to here: about 2,100 lines.
\end{proof}
\leanref{Development/AutomorphicSpectrum/Proofs/BlockDecomposition/SubMemOfLift.lean}{GlobalTheory.CompanionTheory.FibGr.basis_gen}
\leanref{Development/AutomorphicSpectrum/Proofs/BlockDecomposition/CoverCupActionBasis.lean}{GlobalTheory.CompanionTheory.CechVB.vb_basis}

\begin{lemma}\label{gal-realization-thickening-level-one}
\begin{enumerate}
\item Let $(R,I,\gamma)$, $\mathcal O$, $X$ and $\kappa$ be as in Definition~\ref{gal-realization-thickening-nilpotent},
  with $I=0$, let $p\neq0$ be an integer nilpotent in $R$, and let $V\subset X$ be an open with coordinates of dimension
  $e\leq d$. Then $H^j(\mathrm{Th}(X/R)_V,\mathcal O^{[d+1]})=0$ for every $j$.
\item Let $\ell$ be a prime, $k$ a perfect field of characteristic $\ell$ and $W_1(k)=W(k)/\ell W(k)$, which is $k$. Let
  $\mathcal O$ be a ring, $X$ an $\mathcal O$-scheme with a coordinate cover $(U_\alpha)_{\alpha\in A}$
  (Definition~\ref{gal-realization-thickening-coordinates}(2)), $\rho:\mathcal O\to k$ a ring map, $\mathcal
  U=\{U_\alpha:\alpha\in A\}$ and $\mathcal C=\mathrm{Th}(X/W_1(k))_{\mathcal U}$, formed for $(W_1(k),0,\gamma,\ell)$,
  with $\gamma$ the divided powers on the zero ideal, and $\kappa=\rho$. Let $d\geq d_\alpha$ for every $\alpha$. Then
  $H^j(\mathcal C,\mathcal O^{[d+1]})=0$ for every $j$.
\item With $X$ and $\mathcal C$ as in (2), suppose that $X$ is proper over $\mathcal O$. Then $H^j(\mathcal C,\mathcal
  O^{[r]}/\mathcal O^{[r+1]})$ is a finitely generated $W_1(k)$-module for all $r,j\geq0$.
\end{enumerate}
For $k=\mathbf F_q$, $\mathcal O=\mathcal O_K$, $X=Y$ and $\rho$ the reduction map, $\mathrm{Th}(Y/W_1(\mathbf F_q))$ is
$\mathrm{Th}_1(Y/W)$ (Definition~\ref{gal-realization-thickening-nilpotent}).
\end{lemma}
\leanref{Development/AutomorphicSpectrum/Proofs/BlockDecomposition/WittLevelOneVanishAboveDim.lean}{GlobalTheory.CompanionTheory.WittLevel.Van.van_on}
\leanref{Development/AutomorphicSpectrum/Proofs/BlockDecomposition/WittLevelOneVanishAboveDimStatement.lean}{GlobalTheory.T.witt_level_one_vanish_above_dim}
\leanref{Development/AutomorphicSpectrum/Proofs/BlockDecomposition/WittLevelOneGradedFiniteStatement.lean}{GlobalTheory.T.witt_level_one_graded_finite}

\begin{proof}
(1) Let $\mathcal Q=\mathrm{Th}(X/R)_V$. By Lemma~\ref{gal-realization-thickening-chart}(3) there is a divided power
chart $D$ of dimension $e$ on $V$, which gives a torsor datum of dimension $e$ for $\Phi_{\mathcal Q}$ by
Lemma~\ref{gal-realization-thickening-chart-torsor}(1). By Lemma~\ref{gal-realization-thickening-chart-torsor}(3),
$\Phi_{\mathcal Q}$ has a divided power polynomial structure that respects $(J^{[r]})_r$ and $(0)_r$ in every degree; as
$0\subset J^{[r]}$, it is a ladder up to $d+1$ for them. As $I=0$, $I^{[r']}=0$ for $r'\geq1$:
each generator is a product containing some $\gamma_k(0)=0$ with $k\geq1$. So $J_D^{[r']}=i_D(I^{[r']})B_D=0$ for
$r'\geq1$ by Lemma~\ref{gal-realization-thickening-chart}(2), and in particular for $d+1-e\leq r'\leq d+1$. By
Lemma~\ref{gal-realization-thickening-torsor-vanishing}, $H^j(\mathcal Q,J^{[d+1]}/0)=0$ for every $j$.

(2) Write $W_1=W_1(k)=k$; so $\ell W_1=0$ and $\ell=0$ in $W_1$. By Lemma~\ref{gal-realization-thickening-over}, the
thickenings lying over $U_\alpha$ form cosieves of $\mathcal C$ that cover it, and for nonempty $S\subset A$ their
intersection is $\mathrm{Th}(X/W_1)_{U_S}$, with $U_S=\bigcap_{\alpha\in S}U_\alpha$. The affine open $U_S$ has
coordinates of dimension $d_\alpha\leq d$, restricted from $U_\alpha$ for any $\alpha\in S$
(Lemma~\ref{gal-realization-thickening-cover}(1)). So by (1), for $(W_1,0,\gamma,\ell)$, $H^j(\mathrm{Th}(X/W_1)_{U_S},
\mathcal O^{[d+1]})=0$ for every $j$, and Lemma~\ref{gal-realization-thickening-cosieves}(3) gives the claim.

(3) As in (2), write $W_1=W_1(k)=k$; $\mathcal C$ is formed for $(W_1,0,\gamma,\ell)$, with $\ell\neq0$ nilpotent (indeed
$0$) in $W_1$, and $\kappa=\rho:\mathcal O\to W_1$. Let $\mathcal S$ be the class of finitely generated
$W_1$-modules. It satisfies the hypothesis of Lemma~\ref{gal-realization-thickening-level-one-serre}, as $W_1$ is a field,
hence Noetherian: if $M_1\to M_2\to M_3$ is exact with $M_1$ and $M_3$ finitely generated, the image in $M_3$ and the image
of $M_1$ are finitely generated, and so is $M_2$. As in (2), the cosieves $\mathcal C_{(\alpha)}=\mathrm{Th}(X/W_1)_{U_\alpha}$
cover $\mathcal C$, and for $i\in A^{n+1}$ their intersection is $\mathcal C_i=\mathrm{Th}(X/W_1)_{U_i}$, with
$U_i=\bigcap_kU_{i_k}$. By Lemma~\ref{gal-realization-thickening-level-one-cech}, for the functor $\mathcal
O^{[r]}/\mathcal O^{[r+1]}$, it suffices to show that $H^a(\check C^\bullet_b)$ is finitely generated for $a+b=j$, where
$\check C^a_b=\prod_{i\in A^{a+1}}H^b(\mathcal C_i,\mathcal O^{[r]}/\mathcal O^{[r+1]})$.

If $b\neq r$, then $\check C^a_b=0$: $U_i$ is an affine open of $U_{i_0}$, which has coordinates, so $U_i$ has
coordinates (Lemma~\ref{gal-realization-thickening-cover}(1)) and $H^b(\mathcal C_i,\mathcal O^{[r]}/\mathcal
O^{[r+1]})=0$ by Lemma~\ref{gal-realization-thickening-level-one-local}(1).

If $b=r$, $\check C^\bullet_r$ is the complex of the cosimplicial module $K$ of
Lemma~\ref{gal-realization-thickening-level-one-bundle}, which gives data as in
Lemma~\ref{gal-realization-thickening-level-one-coherent} for $\Lambda=W_1$, $Z=X_{W_1}=X\times_{\operatorname{Spec}\mathcal
O}\operatorname{Spec}W_1$ and the opens $U_{\alpha,W_1}$. The scheme $Z$ is proper over $W_1$, as $X$ is proper over
$\mathcal O$ and properness is stable under base change. The $U_{\alpha,W_1}$ cover $Z$, and $\bigcap_kU_{i_k,W_1}$ is the
inverse image of the affine $U_i$ under the projection $Z\to X$, which is affine as a base change of
$\operatorname{Spec}W_1\to\operatorname{Spec}\mathcal O$; so it is affine. By
Lemma~\ref{gal-realization-thickening-level-one-coherent}, $H^a(\check C^\bullet_r)$ is finitely generated over $W_1$.
% Lean: GlobalTheory.T.witt_level_one_graded_finite = FibGr.level_one_gr_finite, for the base Wc.pd l k 1 with
% Im l k 1 = ⊥ (hI) and l = 0 in W_1 (hl, hn), through fin_gen. The class S is serre_fg (finite_modules_serre_class) with
% WittLevel.isNoetherianRing_Wm_one; the reduction to the Cech terms is serre_red (glob_st, through
% CosieveRel.serre_union, Lemma gal-realization-thickening-level-one-cech); b /= r is fin_zero with
% coordinates_restrict_affine_open, capT_eq and off_at; b = r is fin_vb (vb_fin) with vb_cech, aff_pre and capO_pre
% (affine inverse images), cover_pre and IsProper (qF hI q κ) by MorphismProperty.pullback_snd. The remaining open step
% is Lemma gal-realization-thickening-level-one-graded; Lemma gal-realization-thickening-level-one-coherent is cited.
\end{proof}
\leanref{Development/AutomorphicSpectrum/Proofs/BlockDecomposition/WittLevelOneVanishAboveDim.lean}{GlobalTheory.CompanionTheory.WittLevel.Van.witt_level_one_filtration_vanish}
\leanref{Development/AutomorphicSpectrum/Proofs/BlockDecomposition/WittLevelOneGradedFinite.lean}{GlobalTheory.CompanionTheory.FibGr.level_one_gr_finite}
\leanref{Development/AutomorphicSpectrum/Proofs/BlockDecomposition/WittLevelOneGradedFinite.lean}{GlobalTheory.CompanionTheory.FibGr.fin_gen}
\leanref{Development/AutomorphicSpectrum/Proofs/BlockDecomposition/WittLevelOneGradedFinite.lean}{GlobalTheory.CompanionTheory.FibGr.serre_red}
\leanref{Development/AutomorphicSpectrum/Proofs/BlockDecomposition/WittLevelOneGradedFinite.lean}{GlobalTheory.CompanionTheory.FibGr.fin_vb}
\leanref{Development/AutomorphicSpectrum/Proofs/BlockDecomposition/WittLevelOneGradedFinite.lean}{GlobalTheory.CompanionTheory.FibGr.fin_zero}
\leanref{Development/AutomorphicSpectrum/Proofs/BlockDecomposition/FiniteModulesSerreClassStatement.lean}{GlobalTheory.T.finite_modules_serre_class}
\leanref{Development/AutomorphicSpectrum/Proofs/BlockDecomposition/StableCoverCechVanishing.lean}{GlobalTheory.CompanionTheory.GlobSt.glob_st}
\leanref{Development/AutomorphicSpectrum/Proofs/BlockDecomposition/TruncatedWittIsNoetherian.lean}{GlobalTheory.CompanionTheory.WittLevel.isNoetherianRing_Wm_one}

\begin{lemma}\label{gal-realization-frobenius-witt-bockstein}
Let $K$, $\mathcal O_K$, $W$ and $H^j_\nu(Y/W)$ be as in Definition~\ref{gal-realization-frobenius-special-fibre}, also
for $\nu=0$, where $W_0=0$ and $H^j_0(Y/W)=0$. Let $Y$ be a smooth projective $\mathcal O_K$-scheme and $j\geq0$.
\begin{enumerate}
\item The inverse systems $(H^j_\nu(Y/W))_{\nu\geq0}$ and $(H^{j+1}_\nu(Y/W))_{\nu\geq0}$, with the maps induced by
  $W_{\nu+1}\to W_\nu$, have Bockstein sequences for $\ell$ (Definition~\ref{gal-realization-frobenius-bockstein}).
\item $H^j_1(Y/W)$ is a $W$-module of finite length.
\end{enumerate}
\end{lemma}
\leanref{Development/WeilBound/CoreStatements/PadicHodgeComparison/EtaleCohomologyComparison/PerfectoidPullbacks/LocalSystemQuotients/LsSeqEq.lean}{CrystallineSemistableComparison.Pb.LsG.ls_seq}
\leanref{Development/WeilBound/CoreStatements/PadicHodgeComparison/EtaleCohomologyComparison/PerfectoidPullbacks/LocalSystemQuotients/SmoothModelFiniteLengthLink.lean}{CrystallineSemistableComparison.Pb.smooth_model_finite_length_link}

\begin{proof}
(1) Let $K_n$, $\psi_n$ and $\mathrm{res}_n$ be as in Lemma~\ref{gal-realization-frobenius-flat-thickenings}, for
$k=\mathbf F_q$, $\rho$ the residue map and $F=F_Y$, so that $H^j_n(F_Y/W)=H^j_n(Y/W)$
(Definition~\ref{gal-realization-frobenius-flat-witt}). Part (5) applies, with a uniformizer of $\mathcal O_K$ and the
embedding $s$ of Lemma~\ref{gal-realization-invariants}. By its parts (1) and (5), $(K_n,\psi_n)$ is a flat tower (Definition~\ref{gal-realization-frobenius-flat-tower}); for $n=0$,
$K_0=0$ as $W_0=0$, and $\psi_0$ is a map between complexes without cohomology. By
Lemma~\ref{gal-realization-frobenius-flat-tower-bockstein}(2), its cohomology systems in degrees $j$ and $j+1$ have
Bockstein sequences. By Lemma~\ref{gal-realization-frobenius-flat-thickenings}(4),
$r_n\circ\mathrm{res}_{n+1}=\psi_n^{-1}\circ\pi\circ\mathrm{res}_{n+1}=\mathrm{res}_n\circ\mathrm{red}$, so by part (5)
the $\mathrm{res}_n$ are isomorphisms of inverse systems from $(H^j_n(Y/W))$ to the cohomology system, and likewise in
degree $j+1$. Lemma~\ref{gal-realization-frobenius-flat-tower-bockstein}(3) concludes.

(2) The scheme $Y$ is proper and smooth over $\mathcal O_K$. Let $(U_\alpha)$ be a finite cover of $Y$ by affine opens
with \'etale maps $U_\alpha\to\mathbf A^{d_\alpha}_{\mathcal O_K}$, whose finite intersections are affine, and put
$d=\max_\alpha d_\alpha$. Let $\mathcal C\subset\mathrm{Th}_1(Y/W)$ be the full subcategory of the thickenings whose
point factors through some $U_\alpha$, and for $r\geq0$ let $\mathcal O^{[r]}$ be the subfunctor $T\mapsto
J_T^{[r]}$ of $\mathcal O$ (with $J^{[r]}$ as in Definition~\ref{gal-realization-kappa-approximations}). The proof
rests on the following:
\begin{enumerate}
\item[(a)] restriction $H^j_1(Y/W)\to H^j(\mathcal C,\mathcal O)$ is bijective;
\item[(b)] $H^j(\mathcal C,\mathcal O^{[d+1]})=0$ for all $j$;
\item[(c)] each $H^j(\mathcal C,\mathcal O^{[r]}/\mathcal O^{[r+1]})$ is a finitely generated $W_1$-module.
\end{enumerate}
Such a cover is a coordinate cover (Definition~\ref{gal-realization-thickening-coordinates}(2)), and it exists by
Lemma~\ref{gal-realization-thickening-cover}(2), as $Y$ is separated, quasi-compact and smooth. With $\mathcal
U=\{U_\alpha\}$, $\mathcal C=\mathrm{Th}_1(Y/W)_{\mathcal U}$ (Definition~\ref{gal-realization-thickening-nilpotent}). As
every thickening lies over $Y=\bigcup_\alpha U_\alpha$, (a) is Lemma~\ref{gal-realization-thickening-zariski}(1) for
$\mathrm{Th}(Y/W_1)=\mathrm{Th}_1(Y/W)$, $\mathcal V=\mathcal U$ and $\mathcal W=\{Y\}$; (b) is
Lemma~\ref{gal-realization-thickening-level-one}(2), for $k=\mathbf F_q$, $\mathcal O=\mathcal O_K$, $X=Y$ and $\rho$ the
reduction map; and (c) is Lemma~\ref{gal-realization-thickening-level-one}(3) for the same data, as $Y$ is proper.
% Lean: the cover is Pb.cov_ex (GlobalTheory.T.exists_smooth_coordinate_cover); (a) is Pb.zar for St top (range U) <=
% St top {top}, with LevelOne.toAll (the identity of the category Wc.L l F rho 1 = nlFun (Wc.pd l k 1) l X q kappa,
% viewed as its full subcategory QAll, with the change of scalars from W to W_1); (b) is
% GlobalTheory.T.witt_level_one_vanish_above_dim; (c) is GlobalTheory.T.witt_level_one_graded_finite.
Given these, $W_1=\mathbf F_q$ is Noetherian. The long exact sequences of $0\to\mathcal O^{[r+1]}/\mathcal O^{[s]}\to
\mathcal O^{[r]}/\mathcal O^{[s]}\to\mathcal O^{[r]}/\mathcal O^{[r+1]}\to0$ and descending induction on $r$ show by (c)
that every $H^j(\mathcal C,\mathcal O/\mathcal O^{[s]})$ is finitely generated. The long exact sequence of
$0\to\mathcal O^{[d+1]}\to\mathcal O\to\mathcal O/\mathcal O^{[d+1]}\to0$ and (b) show that $H^j(\mathcal C,\mathcal
O)\to H^j(\mathcal C,\mathcal O/\mathcal O^{[d+1]})$ is bijective. So $H^j(\mathcal C,\mathcal O)$, and by (a)
$H^j_1(Y/W)$, is a finitely generated module over the field $W_1$, a quotient of $W$. Hence it has finite length over
$W$.
\end{proof}
\leanref{Development/WeilBound/CoreStatements/PadicHodgeComparison/EtaleCohomologyComparison/PerfectoidPullbacks/LocalSystemQuotients/LsSeqEq.lean}{CrystallineSemistableComparison.Pb.LsG.seq_flat}
\leanref{Development/WeilBound/CoreStatements/PadicHodgeComparison/EtaleCohomologyComparison/PerfectoidPullbacks/LocalSystemQuotients/SmoothModelFiniteLengthLink.lean}{CrystallineSemistableComparison.Pb.wittCohomology_finite_length}
\leanref{Development/AutomorphicSpectrum/Proofs/BlockDecomposition/FilteredModuleFiniteOfGradedStatement.lean}{GlobalTheory.T.filtered_module_finite_of_graded}
\leanref{Development/AutomorphicSpectrum/Proofs/BlockDecomposition/SquarePullBijectiveOfAcyclicStatement.lean}{GlobalTheory.T.square_pull_bijective_of_acyclic}
\leanref{Development/AutomorphicSpectrum/Proofs/BlockDecomposition/FiniteLengthWittCohomologyStatement.lean}{GlobalTheory.T.finite_length_witt_cohomology}
\leanref{Development/WeilBound/CoreStatements/PadicHodgeComparison/EtaleCohomologyComparison/PerfectoidPullbacks/ChartGrConst.lean}{CrystallineSemistableComparison.Pb.zar}
\leanref{Development/WeilBound/CoreStatements/PadicHodgeComparison/EtaleCohomologyComparison/PerfectoidPullbacks/ChartGrConst.lean}{CrystallineSemistableComparison.Pb.cov_ex}
\leanref{Development/AutomorphicSpectrum/Proofs/BlockDecomposition/WittLevelOneVanishAboveDimStatement.lean}{GlobalTheory.T.witt_level_one_vanish_above_dim}
\leanref{Development/AutomorphicSpectrum/Proofs/BlockDecomposition/WittLevelOneGradedFiniteStatement.lean}{GlobalTheory.T.witt_level_one_graded_finite}
\leanref{Development/AutomorphicSpectrum/Proofs/BlockDecomposition/FiniteLengthWittCohomology.lean}{GlobalTheory.CompanionTheory.LevelOne.toAll}

\begin{lemma}\label{gal-realization-frobenius-flat-limit}
Let $W$ be a discrete valuation ring in which the integer $\ell$ is irreducible, $c\geq0$, and $(S_n)_{n\geq0}$ an
inverse system of $W$-modules with transition maps $r_n$, limit $L$ and projections $\pi_n:L\to S_n$. Suppose that $L$
is finitely generated and that for every $n$: $\ell^nS_n=0$, every $x\in L$ with $\pi_n(x)=0$ lies in $\ell^nL$, and
$\ell^cS_n\subset\pi_n(L)$. Let $A$ be a flat $W$-algebra that is $\ell$-adically complete and separated. Then:
\begin{enumerate}
\item if $w\in A\otimes_WL$ and $(1\otimes\pi_n)(w)=0$ for every $n$, then $w=0$;
\item there is an $A$-linear map $\alpha$ from the limit of the $A\otimes_WS_n$, for the maps $1\otimes r_n$, to
  $A\otimes_WL$ with $(1\otimes\pi_n)(\alpha(z))=\ell^cz_n$ for every $z$ and $n$.
\end{enumerate}
\end{lemma}
\leanref{Development/AutomorphicSpectrum/Proofs/BlockDecomposition/FlatBaseChangeLimitSystemStatement.lean}{GlobalTheory.T.flat_base_change_limit_system}

\begin{proof}
Write $\pi_n$ also for $1\otimes\pi_n:A\otimes_WL\to A\otimes_WS_n$; then $\pi_n=(1\otimes r_n)\circ\pi_{n+1}$.

\emph{Step 1.} $A\otimes_WL$ is $\ell$-adically complete and separated. As $\ell$ is irreducible, it generates the maximal
ideal of $W$. By the structure theorem for finitely generated modules over a discrete valuation ring,
$L\cong W^r\oplus\bigoplus_iW/\ell^{e_i}W$, so $A\otimes_WL\cong A^r\oplus\bigoplus_iA/\ell^{e_i}A$. The first summand is
$\ell$-adically complete and separated because $A$ is, and the second because it is killed by a power of $\ell$.

\emph{Step 2.} If $w\in A\otimes_WL$ and $\pi_n(w)=0$, then $w\in\ell^n(A\otimes_WL)$. Indeed, the sequence
$L\xrightarrow{\ell^n}L\xrightarrow{\pi_n}S_n$ is exact, as $\ell^nS_n=0$ and the kernel of $\pi_n$ is $\ell^nL$. As $A$
is flat over $W$, it stays exact after applying $A\otimes_W-$.

\emph{Step 3.} If $\pi_n(w)=0$ for every $n$, then $w$ lies in every $\ell^n(A\otimes_WL)$ by Step 2, so $w=0$ by Step 1.
This is (1).

\emph{Step 4.} Every $t\in A\otimes_WS_n$ has $\ell^ct=\pi_n(w)$ for some $w\in A\otimes_WL$: for $t=a\otimes s$ choose
$y\in L$ with $\pi_n(y)=\ell^cs$ and take $w=a\otimes y$, and sums follow. Let $z=(z_n)$ be compatible, with $z_n\in
A\otimes_WS_n$, and choose $f_n$ with $\pi_n(f_n)=\ell^cz_n$. Then $\pi_n(f_{n+1})=(1\otimes r_n)(\ell^cz_{n+1})=
\ell^cz_n=\pi_n(f_n)$, so $f_{n+1}-f_n\in\ell^n(A\otimes_WL)$ by Step 2. By Step 1 the $f_n$ converge to some $w$ with
$w-f_n\in\ell^n(A\otimes_WL)$ for every $n$. As $\ell^n$ kills $A\otimes_WS_n$, $\pi_n(w)=\pi_n(f_n)=\ell^cz_n$. Put
$\alpha(z)=w$, which is unique by (1). By (1) again, $\alpha$ is $A$-linear, since $\alpha(z+z')-\alpha(z)-\alpha(z')$
and $\alpha(az)-a\alpha(z)$ have image $0$ under every $\pi_n$. This is (2).
\end{proof}
\leanref{Development/AutomorphicSpectrum/Proofs/BlockDecomposition/FlatBaseChangeLimitSystem.lean}{GlobalTheory.CompanionTheory.TowerExt.limit_extension_upTo}
\leanref{Development/AutomorphicSpectrum/Proofs/BlockDecomposition/AdicCompleteTensorFiniteStatement.lean}{GlobalTheory.T.adic_complete_tensor_finite}
\leanref{Development/AutomorphicSpectrum/Proofs/BlockDecomposition/IsAdicCompleteOfPowSmulEqZero.lean}{GlobalTheory.CompanionTheory.TensorComplete.isAdicComplete_tensor}
\leanref{Development/AutomorphicSpectrum/Proofs/BlockDecomposition/FlatBaseChangeLimitSystem.lean}{GlobalTheory.CompanionTheory.TowerExt.exists_eq_pow_smul}
\leanref{Development/AutomorphicSpectrum/Proofs/BlockDecomposition/FlatBaseChangeLimitSystem.lean}{GlobalTheory.CompanionTheory.TowerExt.eq_zero_of_baseChange_pr}
\leanref{Development/AutomorphicSpectrum/Proofs/BlockDecomposition/FlatBaseChangeLimitSystem.lean}{GlobalTheory.CompanionTheory.TowerExt.exists_baseChange_pr_eq}
\leanref{Development/AutomorphicSpectrum/Proofs/BlockDecomposition/FlatBaseChangeLimitSystem.lean}{GlobalTheory.CompanionTheory.TowerExt.exists_lift}

\begin{lemma}\label{gal-realization-frobenius-alpha}
Let $W$ be a discrete valuation ring in which the integer $\ell$ is irreducible, $c\geq0$, $A$ a flat $W$-algebra that is
$\ell$-adically complete and separated, and $\psi$ a ring endomorphism of $A$. Let $(S_n)_{n\geq0}$ be an inverse system
of $W$-modules, with limit $L$, as in Lemma~\ref{gal-realization-frobenius-flat-limit} for $c$, and let $(T_n)_{n\geq0}$
and $(H_n)_{n\geq0}$ be inverse systems of $A$-modules. Let $g_n:S_n\to T_n$ be additive maps with $g_n(wy)=wg_n(y)$ for
$w\in W$, compatible with the transition maps, such that $A\otimes_WS_n\to T_n$, $a\otimes y\mapsto ag_n(y)$, is bijective
for every $n$, and let $\nu_n:H_n\to T_n$ be $\psi$-semilinear maps compatible with the transition maps. Let
$E_n:A\otimes_WL\to T_n$ be the $A$-linear map with $E_n(a\otimes y)=ag_n(y_n)$. Then:
\begin{enumerate}
\item if $w\in A\otimes_WL$ and $E_n(w)=0$ for every $n$, then $w=0$;
\item there is an additive map $\alpha:\varprojlim_nH_n\to A\otimes_WL$ with $E_n(\alpha(x))=\ell^c\nu_n(x_n)$ for every
  $x$ and $n$.
\end{enumerate}
\end{lemma}
\leanref{Development/WeilBound/CoreStatements/PadicHodgeComparison/EtaleCohomologyComparison/PerfectoidPullbacks/FrobeniusFactorisation/ExistsAlpha.lean}{CrystallineSemistableComparison.Pb.Fa.pin_ex}
\leanref{Development/WeilBound/CoreStatements/PadicHodgeComparison/EtaleCohomologyComparison/PerfectoidPullbacks/FrobeniusFactorisation/ExistsAlpha.lean}{CrystallineSemistableComparison.Pb.Fa.Pin}
\leanref{Development/WeilBound/CoreStatements/PadicHodgeComparison/EtaleCohomologyComparison/PerfectoidPullbacks/FrobeniusFactorisation/ExistsAlpha.lean}{CrystallineSemistableComparison.Pb.Fa.E}

\begin{proof}
Let $\pi_n:L\to S_n$ be the projections, and $e_n:A\otimes_WS_n\to T_n$ the bijections $a\otimes y\mapsto ag_n(y)$, so
that $E_n=e_n\circ(1\otimes\pi_n)$. As the $g_n$ are compatible with the transition maps, so are the $e_n$.

(1) If $E_n(w)=0$ for every $n$, then $(1\otimes\pi_n)(w)=0$ for every $n$, as $e_n$ is injective, and $w=0$ by
Lemma~\ref{gal-realization-frobenius-flat-limit}(1).

(2) First let $(z_n)$ be a compatible family with $z_n\in T_n$, and put $u_n=e_n^{-1}(z_n)$. Then $e_n((1\otimes
r_n)(u_{n+1}))$ is the image of $z_{n+1}$ in $T_n$, which is $z_n$, so $(1\otimes r_n)(u_{n+1})=u_n$. With $\alpha$ as in
Lemma~\ref{gal-realization-frobenius-flat-limit}(2), $w=\alpha((u_n)_n)$ satisfies $E_n(w)=e_n(\ell^cu_n)=\ell^cz_n$ for
every $n$.

Now let $x\in\varprojlim_nH_n$. The family $(\nu_n(x_n))_n$ is compatible, so there is $w_x\in A\otimes_WL$ with
$E_n(w_x)=\ell^c\nu_n(x_n)$ for every $n$, and $w_x$ is unique by (1). For $x$ and $x'$, $E_n(w_{x+x'}-w_x-w_{x'})=0$ for
every $n$, as the $\nu_n$ are additive, so $w_{x+x'}=w_x+w_{x'}$ by (1). Put $\alpha(x)=w_x$.
\end{proof}
\leanref{Development/WeilBound/CoreStatements/PadicHodgeComparison/EtaleCohomologyComparison/PerfectoidPullbacks/FrobeniusFactorisation/ExistsAlpha.lean}{CrystallineSemistableComparison.Pb.Fa.core}
\leanref{Development/WeilBound/CoreStatements/PadicHodgeComparison/EtaleCohomologyComparison/PerfectoidPullbacks/FrobeniusFactorisation/ExistsAlpha.lean}{CrystallineSemistableComparison.Pb.Fa.exists_alpha}

\subsection{The crystalline realization and the Eichler--Shimura relation}\label{subsection-gal-crystalline-realization}

\begin{lemma}\label{gal-realization-period-ingredients}
Let $\ell$, $C$, $A_{\mathrm{cris}}$, $K$, $\mathcal O_K$, $\bar K$, $G_K$, $t$ and $B_{\mathrm{dR}}$ be as in
Definition~\ref{gal-realization-period-datum}, $\mathbf F_q$, $q=\ell^f$, the residue field of $K$, $g$ as in
Definition~\ref{gal-realization-frobenius-datum}, and $W$ and $L^j$ as in
Definition~\ref{gal-realization-frobenius-special-fibre}. Where no other ring is named, the thickenings below are those
of Definition~\ref{gal-realization-period-thickenings} for $\mathcal O=\mathcal O_K$ and its inclusion into $\mathcal
O_C$, and $\eta_\nu$ are the comparison maps of Definition~\ref{gal-realization-eta}.
\begin{enumerate}
\item (Bounds.) For all integers $d,j\geq0$ there is $N\geq0$ with the following property. Let $\mathcal O$ be a ring,
  $F$ a point functor over $\mathcal O$ with a Zariski atlas of dimension $d$
  (Definition~\ref{gal-realization-ingredients-atlas}), $\kappa:\mathcal O\to\mathcal O_C$ a ring map, $\varepsilon$
  and $\mu$ as in Lemma~\ref{gal-realization-period-t-mu}, and $\nu\geq1$. Then $\mu^N$ kills the kernels of
  $\beta_\nu$ on $H^j_\nu(F)$ and of $\beta'_\nu$ on $H'^j_\nu(F)$, $\mu^NH^j(F^\sharp,A_{\mathrm{cris}}/\ell^\nu)$ lies
  in the image of $\beta'_\nu$, and for every $y\in H^j(F^\sharp,A_{\mathrm{cris}}/\ell^\nu)$ there is
  $x\in\mathrm{Fil}^NH^j_\nu(F)$ with $\beta_\nu(x)=\mu^Ny$.
\item (Chern classes; Definition~\ref{gal-realization-period-chern}.)
  \begin{enumerate}
  \item[(a)] Let $R$ be a ring and $X$ an $R$-scheme with a closed immersion into some $\mathbf P^M_R$ over $R$. There
    is a tautological line datum on $X$, and for every ring map $R\to R'$ the base change to $X\times_{\operatorname{Spec}
    R}\operatorname{Spec}R'$ of a tautological line datum on $X$ is tautological.
  \item[(b)] Let $K'$ be a field with $\mathbf Q_\ell\subset K'\subset\bar K$ and $\mathcal O_{K'}=K'\cap\mathcal O_C$
    its valuation ring, with the thickenings for $\mathcal O_{K'}$ and its inclusion into $\mathcal O_C$. Let $Y$ be an
    $\mathcal O_{K'}$-scheme and $\mathcal L$ a line datum on $Y$ such that $Y$ is the union of the open subschemes over
    which $f$ has a section. There is $x\in H^2_\infty(Y)$ whose components form a first Chern family of $\mathcal L$.
  \item[(c)] Let $\mathcal O$ be a ring, $X$ an $\mathcal O$-scheme, $\kappa:\mathcal O\to\mathcal O_C$ a ring map
    and $\mathcal L$ a line datum on $X$, with the thickenings for $\mathcal O$ and $\kappa$. If $(x_\nu)_{\nu\geq1}$,
    with $x_\nu\in H^2_\nu(X)$, resp.\ $x_\nu\in H'^2_\nu(X)$, is a first Chern family of $\mathcal L$, then
    $(\beta_\nu(x_\nu))_\nu$, resp.\ $(\beta'_\nu(x_\nu))_\nu$, is a first Chern class of $\mathcal L$.
  \item[(d)] Let $\mathcal O$, $X$, $\kappa$ and $\mathcal L$ be as in (c), with $f$ \'etale. Then $\mathcal L$ has at
    most one first Chern class $(y_\nu)_{\nu\geq1}$ with $y_\nu\in H^2(X^\sharp,A_{\mathrm{cris}}/\ell^\nu)$.
  \item[(e)] Let $K'$ and $\mathcal O_{K'}$ be as in (b), with $Y^\sharp$ formed for the inclusion $\mathcal
    O_{K'}\subset\mathcal O_C$, and let $\eta_\nu$ be the comparison maps of Definition~\ref{gal-realization-eta} for
    $K'$. For every $\mathcal O_{K'}$-scheme $Y$ and every line datum $\mathcal L$ on $Y$ there is $v\in H^2(Y_{\bar
    K},\mathbf Z_\ell)$, where $Y_{\bar K}=Y\times_{\mathcal O_{K'}}\bar K$, such that $(t\,\eta_\nu(v))_\nu$ is a first
    Chern class of $\mathcal L$.
  \end{enumerate}
\item (Periods.) For every smooth projective $\mathcal O_K$-scheme $Y$, with generic fibre $Y_K$, each
  $H^j_{\mathrm{dR}}(Y_K)$ (Definition~\ref{gal-realization-period-de-rham}) is finite-dimensional over $K$. There are
  additive maps $\kappa_{Y,\mathcal L}$ from $H^j_\infty(Y)$ to the group of functions
  $\operatorname{Hom}_K(H^j_{\mathrm{dR}}(Y_K),K)\to B_{\mathrm{dR}}$, one for every smooth projective $\mathcal
  O_K$-scheme $Y$, tautological line datum $\mathcal L$ on $Y$ and $j\geq0$, with the following properties, where
  $\kappa=\kappa_{Y,\mathcal L}$.
  \begin{enumerate}
  \item[(a)] $\kappa(cx)(\psi)=c\,\kappa(x)(\psi)$, $\kappa(x)(a\psi)=a\,\kappa(x)(\psi)$ and $\kappa(\sigma
    x)(\psi)=\sigma(\kappa(x)(\psi))$ for $c\in A_{\mathrm{cris}}$, $a\in K$ and $\sigma\in G_K$.
  \item[(b)] For $x\in H^i_\infty(Y)$, $x'\in H^j_\infty(Y)$, bases $(b_s)$ of $H^i_{\mathrm{dR}}(Y_K)$ and $(b'_{s'})$
    of $H^j_{\mathrm{dR}}(Y_K)$ with dual bases $(b_s^*)$ and $(b'^*_{s'})$, and $\psi\in\operatorname{Hom}_K(
    H^{i+j}_{\mathrm{dR}}(Y_K),K)$,
    \[
     \kappa(x\cup x')(\psi)=\sum_{s,s'}\kappa(x)(b_s^*)\,\kappa(x')(b'^*_{s'})\,\psi(b_s\cup b'_{s'}),
    \]
    and $\kappa(1)(\psi)=\psi(1)$.
  \item[(c)] If the components of $x\in H^2_\infty(Y)$ form a first Chern family of $\mathcal L$ and $\omega\in
    H^2_{\mathrm{dR}}(Y_K)$ is a first Chern class of the base change of $\mathcal L$ to $Y_K$, then
    $\kappa(x)(\psi)=\psi(\omega)$ for every $\psi$.
  \item[(d)] Let $Y'$ be a smooth projective $\mathcal O_K$-scheme with a tautological line datum $\mathcal L'$, and
    $h:Y\to Y'$ an $\mathcal O_K$-morphism with generic fibre $h_K:Y_K\to Y'_K$, and let $h^*:H^j_\infty(Y')\to
    H^j_\infty(Y)$ and $h_K^*:H^j_{\mathrm{dR}}(Y'_K)\to H^j_{\mathrm{dR}}(Y_K)$ be the pullbacks
    (Definitions~\ref{gal-realization-period-thickenings}(3)(e) and~\ref{gal-realization-period-de-rham}). Then
    $\kappa_{Y,\mathcal L}(h^*x')(\psi)=\kappa_{Y',\mathcal L'}(x')(\psi\circ h_K^*)$ for $x'\in H^j_\infty(Y')$ and
    $\psi\in\operatorname{Hom}_K(H^j_{\mathrm{dR}}(Y_K),K)$.
  \item[(e)] Let $n,r\geq0$, let $x\in H^j_\infty(Y)$ have every component $x_\nu$ in $\mathrm{Fil}^nH^j_\nu(Y)$
    (Definition~\ref{gal-realization-period-thickenings}(6)), and let $\psi$ vanish on the classes in
    $H^j_{\mathrm{dR}}(Y_K)$ of the cocycles whose value on every chain $T_0\to\cdots\to T_j$ lies in
    $J_{T_j}^{[r+1]}$. Then $t^r\kappa(x)(\psi)=t^ny$ for some $y\in B_{\mathrm{dR}}^+$.
  \end{enumerate}
  If $Y$ is a smooth projective $\mathcal O_K$-scheme with $Y_K$ connected of dimension $d$, and $\mathcal L'$ is a
  tautological line datum on $Y_K$, then $\mathcal L'$
  has a first Chern class $\omega\in H^2_{\mathrm{dR}}(Y_K)$, and $\omega^d\neq0$ for each such $\omega$, where
  $\omega^0=1$ and $\omega^{k+1}=\omega^k\cup\omega$.
\item (Frobenius.) For every integer $f_1\geq1$ there is a multiple $N$ of $f_1$ with $(x\bmod\ell)^{\ell^N}=g(x)$ for all
  $x\in\mathcal O_K$ such that for every smooth projective $\mathcal O_K$-scheme $Y$ and $j\geq0$ there are $e\geq0$, an
  additive map
  $\alpha:H'^j_\infty(Y)\to A_{\mathrm{cris}}\otimes_WL^j(Y)$ and an $A_{\mathrm{cris}}$-linear map
  $\beta^\circ:A_{\mathrm{cris}}\otimes_WL^j(Y)\to H'^j_\infty(Y)$ with the following properties.
  \begin{enumerate}
  \item[(a)] $\alpha(cx)=\phi^N(c)\alpha(x)$, $\alpha(\sigma x)=(\sigma\otimes1)(\alpha(x))$ and
    $\beta^\circ(\alpha(x))=\ell^e\phi^N(x)$ for $c\in A_{\mathrm{cris}}$, $\sigma\in G_K$ and $x\in H'^j_\infty(Y)$.
  \item[(b)] For every $\nu\geq1$ there is a Frobenius twist of level $N$ whose map $m$ satisfies
    $\beta^\circ(1\otimes y)_\nu=m(y_\nu)$ for every $y=(y_\nu)_\nu\in L^j(Y)$.
  \end{enumerate}
\end{enumerate}
% Lean form: every item is in the Lean's form. (2)(b) (Pb.firstChern_acris_family_exists) and (2)(e)
% (E.FirstChernClassFamily) are stated for every intermediate field K' of PadicAlgCl l over Q_l and every Model l K',
% and (4) (PinEx) for every f_1 > 0 (the archive's f), with K finite as PinEx requires.
\end{lemma}
\leanref{Development/WeilBound/CoreStatements/PadicHodgeComparison/EtaleCohomologyComparison/ZariskiEtaleLifting/MapUSharp.lean}{CrystallineSemistableComparison.B.Main}
\leanref{Development/WeilBound/CoreStatements/PadicHodgeComparison/EtaleCohomologyComparison/ZariskiEtaleLifting/FirstChernCompatibleL.lean}{CrystallineSemistableComparison.B.main}
\leanref{Development/WeilBound/CoreStatements/PadicHodgeComparison/EtaleCohomologyComparison/BocksteinBaseChange/AtlasGroupActions/AtlasOf.lean}{CrystallineSemistableComparison.GlueB'.atlas_of}
\leanref{Development/WeilBound/CoreStatements/PadicHodgeComparison/EtaleCohomologyComparison/PerfectoidPullbacks/ConeSectionRestriction/IsFirstChernDeRhamUnique.lean}{CrystallineSemistableComparison.Pb.isO1_exists}
\leanref{Development/WeilBound/CoreStatements/PadicHodgeComparison/EtaleCohomologyComparison/PerfectoidPullbacks/ConeSectionRestriction/IsFirstChernDeRhamUnique.lean}{CrystallineSemistableComparison.Pb.isTwistingSheaf_pullback}
\leanref{Development/WeilBound/CoreStatements/PadicHodgeComparison/EtaleCohomologyComparison/PerfectoidPullbacks/LogarithmicDerivatives/CrystallineFirstChernExistsUnique.lean}{CrystallineSemistableComparison.Pb.firstChern_acris_family_exists}
\leanref{Development/WeilBound/CoreStatements/PadicHodgeComparison/EtaleCohomologyComparison/PerfectoidPullbacks/ConeSectionRestriction/RSubOneMem.lean}{CrystallineSemistableComparison.Pb.sSup_𝒱D}
\leanref{Development/WeilBound/CoreStatements/PadicHodgeComparison/EtaleCohomologyComparison/ZariskiEtaleLifting/IsFirstChernClassOfLineBundle.lean}{CrystallineSemistableComparison.B.FirstChernCompatible}
\leanref{Development/WeilBound/CoreStatements/PadicHodgeComparison/EtaleCohomologyComparison/ZariskiEtaleLifting/FirstChernCompatibleL.lean}{CrystallineSemistableComparison.B.firstChern_compatible_l}
\leanref{Development/WeilBound/CoreStatements/PadicHodgeComparison/EtaleCohomologyComparison/KummerCupProducts/IsFirstChernClassOfLineBundle.lean}{CrystallineSemistableComparison.E.FirstChernUniqueEtale}
\leanref{Development/WeilBound/CoreStatements/PadicHodgeComparison/EtaleCohomologyComparison/KummerCupProducts/EtaleChernClassEqKummerClass.lean}{CrystallineSemistableComparison.E.Main.firstChern_unique}
\leanref{Development/WeilBound/CoreStatements/PadicHodgeComparison/EtaleCohomologyComparison/KummerCupProducts/EtaleComparisonMapOfFacts.lean}{CrystallineSemistableComparison.E.FirstChernClassFamily}
\leanref{Development/WeilBound/CoreStatements/PadicHodgeComparison/EtaleCohomologyComparison/KummerCupProducts/EtaleChernClassEqKummerClass.lean}{CrystallineSemistableComparison.E.Main.firstChern_family}
\leanref{Development/WeilBound/CoreStatements/PadicHodgeComparison/EtaleCohomologyComparison/PerfectoidPullbacks/Lin.lean}{CrystallineSemistableComparison.Pb.KapL}
\leanref{Development/WeilBound/CoreStatements/PadicHodgeComparison/EtaleCohomologyComparison/PerfectoidPullbacks/LieComplementChartContDiffAt.lean}{CrystallineSemistableComparison.Pb.LieComplementChartContDiffAt}
\leanref{Development/WeilBound/CoreStatements/PadicHodgeComparison/EtaleCohomologyComparison/PerfectoidPullbacks/LevelChainValuations/CrystallineDeRhamFirstChernExists.lean}{CrystallineSemistableComparison.Pb.crystalline_deRham_firstChern_exists}
\leanref{Development/WeilBound/CoreStatements/PadicHodgeComparison/EtaleCohomologyComparison/PerfectoidPullbacks/IsProperOfProj.lean}{CrystallineSemistableComparison.Pb.m_fin}
\leanref{Development/WeilBound/CoreStatements/PadicHodgeComparison/EtaleCohomologyComparison/PerfectoidPullbacks/IsProperOfProj.lean}{CrystallineSemistableComparison.Pb.cupR_eq}
\leanref{Development/WeilBound/CoreStatements/PadicHodgeComparison/EtaleCohomologyComparison/PerfectoidPullbacks/ConeSectionRestriction/IsFirstChernDeRhamUnique.lean}{CrystallineSemistableComparison.Pb.isFirstChernDeRham_exists}
\leanref{Development/WeilBound/CoreStatements/PadicHodgeComparison/EtaleCohomologyComparison/PerfectoidPullbacks/ConeSectionRestriction/IsFirstChernDeRhamUnique.lean}{CrystallineSemistableComparison.Pb.powK_ne_zero}
\leanref{Development/WeilBound/CoreStatements/PadicHodgeComparison/EtaleCohomologyComparison/TeichmullerRootsUnity/PinFixed.lean}{CrystallineSemistableComparison.crystallineComparison.pinEx}
\leanref{Development/WeilBound/CoreStatements/PadicHodgeComparison/EtaleCohomologyComparison/PerfectoidPullbacks/FrobeniusWittActions/PhiTTmul.lean}{CrystallineSemistableComparison.Pb.Fac}
\leanref{Development/WeilBound/CoreStatements/PadicHodgeComparison/EtaleCohomologyComparison/PerfectoidPullbacks/FrobeniusWittActions/TeichmullerCongruence.lean}{CrystallineSemistableComparison.Pb.Mid.TeichmullerCongruence}
\leanref{Development/AutomorphicArithmetic/CommutativeAlgebra/EtaleCohomology/AbsInvOf.lean}{UpperTheory.Ru.crystallineComparison.Pin}
\leanref{Development/WeilBound/CoreStatements/PadicHodgeComparison/EtaleCohomologyComparison/KummerCupProducts/FirstChernClassFamily.lean}{CrystallineSemistableComparison.E.firstChern_family_of_lifts}
\leanref{Development/AutomorphicSpectrum/Proofs/BlockDecomposition/LineBC.lean}{GlobalTheory.CompanionTheory.LineBase.lineBC}
\leanref{Development/AutomorphicSpectrum/Proofs/BlockDecomposition/ExistsRoot.lean}{GlobalTheory.CompanionTheory.FKum.fKum}
\leanref{Development/AutomorphicSpectrum/Proofs/BlockDecomposition/VnRed.lean}{GlobalTheory.CompanionTheory.FUntwPf.main}
\leanref{Development/AutomorphicSpectrum/Proofs/BlockDecomposition/PullSquare.lean}{GlobalTheory.CompanionTheory.FResProof.fRes}

\begin{proof}
(1) This is Lemma~\ref{gal-realization-ingredients-bounds}.

(2)(a) Let $i:X\to\mathbf P^M_R$ be a closed immersion over $R$. Lemma~\ref{gal-realization-ingredients-tautological}(2)
gives a line datum on $X$ that is tautological for $i$. If a line datum on $X$ is tautological for some such $i$, its
base change to $X\times_{\operatorname{Spec}R}\operatorname{Spec}R'$ is tautological for the base change of $i$, by
Lemma~\ref{gal-realization-ingredients-tautological}(3) for $R\to R'$.

(2)(b) By hypothesis $Y$ is the union of the open subschemes over which $f$ has a section, so by
Lemma~\ref{gal-realization-ingredients-crystalline-chern}, for the base $(A_\nu,I_\nu,\gamma,\ell)$ and
$N=\ell\nu$, $\mathcal L$ has exactly one first Chern class $x_\nu\in H^2_\nu(Y)$ for every $\nu$. By
Lemma~\ref{gal-realization-ingredients-chern-reduction}, the
reduction of $x_{\nu+1}$ is a first Chern class in $H^2_\nu(Y)$, so it is $x_\nu$. Hence $x=(x_\nu)_\nu$ lies in
$H^2_\infty(Y)$, and its components form a first Chern family of $\mathcal L$.

(2)(c) Consider first $x_\nu\in H^2_\nu(X)$, and let the $x_\nu$ form a first Chern family of $\mathcal L$. Fix
$\nu$, a functor $G:\mathcal C\to\mathrm{El}(X^\sharp)$ from a category $\mathcal C$, $G(c)=(A_c,y_c)$, and a lift
$(p_c,w_\varphi)$ of $\mathcal L$ along $G$ as in Definition~\ref{gal-realization-period-chern}(5). For
$\varphi:c\to c'$ write $h_\varphi=G(\varphi):A_c\to A_{c'}$.

\emph{Step 1.} For a chain $s=(c_0\xrightarrow{\varphi_1}c_1\xrightarrow{\varphi_2}c_2)$, $\theta([r(s)])=1$. Put
$a=p_{c_0}\circ\operatorname{Spec}h_{\varphi_2\varphi_1}^\sharp$, $b=p_{c_1}\circ\operatorname{Spec}
h_{\varphi_2}^\sharp$ and $e=p_{c_2}$. As $h_{\varphi_2\varphi_1}=h_{\varphi_2}h_{\varphi_1}$ and $\theta$ and the
Teichm\"uller map commute with the maps $W(h)$, the naturality of $u$ gives
$\theta([h_{\varphi_2}(w_{\varphi_1})])=u(a,b)$, and $\theta([w_{\varphi_2}])=u(b,e)$,
$\theta([w_{\varphi_2\varphi_1}])=u(a,e)$. By the cocycle identity, $\theta([r(s)])=u(a,b)u(b,e)u(a,e)^{-1}=1$. So
$[r(s)]-1=v_s\xi$ for some $v_s\in W(A_{c_2})$, and $([r(s)]-1)^k/k!=v_s^k\xi^k/k!$ lies in
$A_{\mathrm{cris}}(A_{c_2})$ for every $k$.

\emph{Step 2.} A lift along $G_\nu=g_\nu\circ G$, with $g_\nu$ as in
Definition~\ref{gal-realization-period-thickenings}(4). The thickening $G_\nu(c)$ has the ring
$A_{\mathrm{cris}}(A_c)/\ell^\nu$, the ideal $I_\nu(A_c)$ and the point $y_c\circ\operatorname{Spec}u_{A_c}$. Put
$p'_c=p_c\circ\operatorname{Spec}u_{A_c}$, so that $fp'_c=y_c\circ\operatorname{Spec}u_{A_c}$, and let $w'_\varphi$
be the image of $[w_\varphi]$ in $A_{\mathrm{cris}}(A_{c'})/\ell^\nu$, a unit. The map $\bar h_\varphi$ induced by
$G_\nu(\varphi)$ modulo the ideals satisfies $\bar h_\varphi\circ u_{A_c}=u_{A_{c'}}\circ h_\varphi^\sharp$, and
$u_A(\theta(z))$ is the class of $z$ for $z\in W(A)$. So, by the naturality of $u$, the class of $w'_\varphi$ modulo
$I_\nu(A_{c'})$ is the pullback along $\operatorname{Spec}u_{A_{c'}}$ of $\theta([w_\varphi])=u(p_c\circ
\operatorname{Spec}h_\varphi^\sharp,p_{c'})$, which is $u(p'_c\circ\operatorname{Spec}\bar h_\varphi,p'_{c'})$. Hence
$(p'_c,w'_\varphi)$ is a lift of $\mathcal L$ along $G_\nu$. As the Teichm\"uller map is multiplicative and commutes
with the maps $W(h_\varphi)$, its elements $r'(s)$ are the images of the $[r(s)]$.

\emph{Step 3.} As $x_\nu$ is a first Chern class, $z(s)=\sum_{k=0}^{\ell\nu-1}(-1)^kk!\,\delta_{k+1}(r'(s)-1)$ is a
$2$-cocycle of $\mathcal C$ whose class is $G_\nu^*x_\nu$. The map $\beta_\nu$ is induced by $g_\nu$ with the
identity maps of the rings $A_{\mathrm{cris}}(A)/\ell^\nu$, so $G^*\beta_\nu(x_\nu)=G_\nu^*x_\nu$ is the class of $z$.
By Step~1, $r'(s)-1$ is the image of $v_s\xi\in I(A_{c_2})$, so by
Definition~\ref{gal-realization-period-thickenings}(1) $\delta_{k+1}(r'(s)-1)$ is the image of
$(v_s\xi)^{k+1}/(k+1)!=([r(s)]-1)^{k+1}/(k+1)!$. So $z(s)$ is the image of
$\sum_{k=0}^{\ell\nu-1}(-1)^kk!\,([r(s)]-1)^{k+1}/(k+1)!$, as Definition~\ref{gal-realization-period-chern}(5)
requires. This holds for every $\nu$, $G$ and lift, so $(\beta_\nu(x_\nu))_\nu$ is a first Chern class of
$\mathcal L$.

For $x_\nu\in H'^2_\nu(X)$ the same argument applies with $g'_\nu$, $I'_\nu(A_c)$ and $u_{A_c}$ followed by the
projection in place of $g_\nu$, $I_\nu(A_c)$ and $u_{A_c}$: the class of $w'_\varphi$ modulo $I'_\nu(A_{c'})$ is the image
of its class modulo $I_\nu(A_{c'})$, and by Definition~\ref{gal-realization-period-thickenings}(1) the divided powers
of $I'_\nu(A_{c_2})$ on the image of $v_s\xi\in I(A_{c_2})$ are again the images of the $(v_s\xi)^{k+1}/(k+1)!$.

(2)(d) Let $y=(y_\nu)_\nu$ and $y'=(y'_\nu)_\nu$, with $y_\nu,y'_\nu\in H^2(X^\sharp,A_{\mathrm{cris}}/\ell^\nu)$, be
first Chern classes of $\mathcal L$. Take $\mathcal C=\mathrm{El}(X^\sharp)$ and $G$ the identity functor. For
$(A,y_A)$ in $\mathcal C$ the base change of $f$ along $y_A:\operatorname{Spec}A^\sharp\to X$ is \'etale and surjective
onto $\operatorname{Spec}A^\sharp$, so it has a section by Lemma~\ref{gal-crystalline-test-facts}(2), which gives
$p_{(A,y_A)}:\operatorname{Spec}A^\sharp\to X_{\mathcal L}$ with $fp_{(A,y_A)}=y_A$. For a map $\varphi:(A,y_A)\to
(A',y_{A'})$, the unit $u(p_{(A,y_A)}\circ\operatorname{Spec}\varphi^\sharp,p_{(A',y_{A'})})$ of $A'^\sharp$ is
$\theta([w_\varphi])$ for a unit $w_\varphi$ of $A'$ by Lemma~\ref{gal-realization-ingredients-sharp-units}. So
$\mathcal L$ has a lift along $G$. For every $\nu$, Definition~\ref{gal-realization-period-chern}(5) gives
$2$-cocycles $z$ and $z'$ whose classes are $G^*y_\nu=y_\nu$ and $G^*y'_\nu=y'_\nu$, since the identity functor with
identity maps induces the identity. For every chain $s$, both $z(s)$ and $z'(s)$ are the image of the same element
$\sum_{k=0}^{\ell\nu-1}(-1)^kk!\,([r(s)]-1)^{k+1}/(k+1)!$. Hence $z=z'$ and $y_\nu=y'_\nu$.

(2)(e) Let $\mathcal L$ be a line datum on $Y$, put $X=Y_{\bar K}$, and let $\mathcal L'$ be the base change of
$\mathcal L$ along $X\to Y$ (Definition~\ref{gal-realization-period-chern}(1)), with $f':X_{\mathcal L'}=X_{\mathcal
L}\times_YX\to X$, the projection $\pi:X_{\mathcal L'}\to X_{\mathcal L}$ and the units $u'(b,c)=u(\pi b,\pi c)$. Let
$\varepsilon$ be as in Lemma~\ref{gal-realization-period-t-mu}, with $t=\log[\varepsilon]$. For $n\geq1$ let
$\zeta_n=\theta([\varepsilon^{1/\ell^n}])$. By Lemma~\ref{gal-realization-ingredients-log-roots} it is a primitive
$\ell^n$-th root of unity, so it lies in $\bar K$ and is a function on $X$, and $\zeta_{n+1}^\ell=\zeta_n$. For the
split diagram $(U_{(A,y)})$ over $\mathrm{El}(Y^\sharp)$ with points $y_{(A,y)}$ in $X$ of
Definition~\ref{gal-realization-eta}, the images $\zeta_{n,(A,y)}$ of $\zeta_n$ in the rings $A^\sharp[1/\ell]$ are
primitive $\ell^n$-th roots of unity, compatible with the maps of the diagram, and $y_{(A,y)}^*\zeta_n=\zeta_{n,(A,y)}$.
Fix injective resolutions and Kummer maps as in Lemma~\ref{gal-realization-ingredients-kummer-class}.

\emph{Step 1: the class $v$.} Let $x_n\in H^2(X,\mu_{\ell^n})$ be the classes of
Lemma~\ref{gal-realization-ingredients-kummer-class} for $X$ and $\mathcal L'$, and let $w_n\in H^0(X,(\mathbf
Z/\ell^n)(-1))$ be the class of $\zeta_n^{\otimes-1}$. The reduction carries $w_{n+1}$ to $w_n$, as
$\zeta_{n+1}^\ell=\zeta_n$. Put $v_n=x_n\cup w_n\in H^2(X,\mathbf Z/\ell^n)$. The reductions are compatible with cup
products, so $v=(v_n)_n$ lies in $H^2(X,\mathbf Z_\ell)$. By Lemma~\ref{gal-realization-ingredients-twisted-edge}(1),
(3) and~(4), $\varepsilon_y(v_n)=\varepsilon^{(1)}_y(x_n)\cup\varepsilon^{(-1)}_y(w_n)=\varepsilon^{(1)}_y(x_n)$, so
$\eta_n(v_n)=\lambda_*\varepsilon^{(1)}_y(x_n)$.

\emph{Step 2: a Kummer cocycle.} Let $n\geq1$ and $q=\ell^n$, let $G:\mathcal C\to\mathrm{El}(Y^\sharp)$ be a functor
from a category $\mathcal C$, $G(c)=(A_c,y_c)$, and let $(p_c,w_\varphi)$ be a lift of $\mathcal L$ along $G$ as in
Definition~\ref{gal-realization-period-chern}(5); write $h_\varphi=G(\varphi)$. Let $U_c=\operatorname{Spec}
A_c^\sharp[1/\ell]$, $j_c:U_c\to\operatorname{Spec}A_c^\sharp$, $y'_c=y_{G(c)}$ and $\zeta_c=\zeta_{n,G(c)}$. With the
maps $U_\varphi$ induced by the $h_\varphi$, this is a split diagram over $\mathcal C$ with points $y'_c$ in $X$. As
$fp_cj_c=y_cj_c$ is the component of $y'_c$ in $Y$, there are $p'_c:U_c\to X_{\mathcal L'}$ with $\pi p'_c=p_cj_c$ and
$f'p'_c=y'_c$. As $A_{c'}$ is perfect, $w_\varphi$ has a $q$-th root $\omega_\varphi\in A_{c'}$, which is a unit. Let
$r_\varphi$ be the image of $\theta([\omega_\varphi])$ on $U_{c'}$. By the naturality of $u$, $r_\varphi^q$, the image
of $\theta([w_\varphi])=u(p_c\circ\operatorname{Spec}h_\varphi^\sharp,p_{c'})$, is
$u(p_cj_cU_\varphi,p_{c'}j_{c'})=u'(p'_cU_\varphi,p'_{c'})$. For a chain $s=(c_0\xrightarrow{\varphi_1}c_1
\xrightarrow{\varphi_2}c_2)$ let $\rho(s)=h_{\varphi_2}(\omega_{\varphi_1})\,\omega_{\varphi_2}\,
\omega_{\varphi_2\varphi_1}^{-1}\in A_{c_2}^\times$. The image of $\theta([\rho(s)])$ on $U_{c_2}$ is
$U_{\varphi_2}^*(r_{\varphi_1})\,r_{\varphi_2}\,r_{\varphi_2\varphi_1}^{-1}$, whose $q$-th power is $1$ by the
naturality and the cocycle identity of $u'$. By Lemma~\ref{gal-realization-ingredients-roots}(2) it is
$\zeta_{c_2}^{z(s)}$ for some $z(s)\in L_q(c_2)$. By Lemma~\ref{gal-realization-ingredients-kummer-cocycle}, $z$ is a
cocycle and $\varepsilon^{(1)}_{y'}(x_n)$ is its class.

\emph{Step 3: the first Chern class.} By Lemma~\ref{gal-realization-ingredients-twisted-edge}(2) for the functor $G$,
the identity maps of the $U_c$ and the identity of $X$, $G^*\varepsilon^{(1)}_y(x_n)=\varepsilon^{(1)}_{y'}(x_n)$. The
maps $G^*$ and $\lambda_*$ commute, as both composites are induced by $G$ with the maps $\lambda_{A_c}$ on
coefficients. By Step~1, $G^*(t\,\eta_n(v_n))$ is therefore the class of the cocycle $Z$ with
$Z(s)=t\,\lambda_{A_{c_2}}(z(s))$. The elements of the lift are $r(s)=h_{\varphi_2}(w_{\varphi_1})\,w_{\varphi_2}\,
w_{\varphi_2\varphi_1}^{-1}=\rho(s)^q$. By Lemma~\ref{gal-realization-ingredients-log-roots} for $A_{c_2}$, $u=\rho(s)$
and $f=z(s)$, $\theta([r(s)])=1$ and $\log_n([r(s)])=Z(s)$. By
Lemma~\ref{gal-realization-ingredients-logarithm}(1), the elements $([r(s)]-1)^k/k!$ lie in
$A_{\mathrm{cris}}(A_{c_2})$, and $Z(s)$ is the image of $\sum_{k=0}^{\ell n-1}(-1)^kk!\,([r(s)]-1)^{k+1}/(k+1)!$. As
$n$, $G$ and the lift were arbitrary, $(t\,\eta_\nu(v))_\nu$ is a first Chern class of $\mathcal L$.

(3) Let $Y$ be a smooth projective $\mathcal O_K$-scheme. Its generic fibre $Y_K$ is proper and smooth over $K$, so each
$H^j_{\mathrm{dR}}(Y_K)$ is finite-dimensional by Lemma~\ref{gal-realization-kappa-de-rham}(1). Let the approximations
$(B_m,\tau_m,\ldots)$, with $d_A$ and $d_B$, form an approximation system
(Lemma~\ref{gal-realization-kappa-approximation-system}), and let $\rho$ be a system of truncated periods for it
(Lemma~\ref{gal-realization-kappa-truncated-periods-exist}); both are fixed for all $Y$. Let $\mathcal L$ be a
tautological line datum on $Y$. The construction below of $\kappa=\kappa_{Y,\mathcal L}$ uses only $\rho$ for $Y$ and
$\mathcal L$. The numbers (P1) to (P9) refer to the items of
Definition~\ref{gal-realization-kappa-truncated-periods}. If $z,z'\in B_{\mathrm{dR}}^+$ satisfy
$\tau_m(z)=\tau_m(z')$ for every $m$, then $z=z'$, by the uniqueness in
Definition~\ref{gal-realization-kappa-approximations}(2) for the family $(\tau_m(z))_m$, which is compatible since
$d_B\circ\tau_{m+1}=\tau_m$. This is used in each step below.

\emph{Construction.} Let $j\geq0$ and $x\in H^j_\infty(Y)$. By (P1), for every $m$ there is exactly one $f^x_m$ with
$\rho_m(x,f^x_m)$, and by (P3) and (P1), $d_B\circ f^x_{m+1}=f^x_m$. By
Definition~\ref{gal-realization-kappa-approximations}(2), for every $\psi\in\operatorname{Hom}_K(H^j_{\mathrm{dR}}(Y_K),K)$
there is exactly one $\kappa(x)(\psi)\in B_{\mathrm{dR}}^+\subset B_{\mathrm{dR}}$ with $\tau_m(\kappa(x)(\psi))=
f^x_m(\psi)$ for every $m$. Each $f^x_m$ is additive, and $f^{x+x'}_m=f^x_m+f^{x'}_m$ by (P2) and (P1). So
$\kappa(x)$ and $\kappa$ are additive.

\emph{Item (a).} Let $c\in A_{\mathrm{cris}}$. By (P2), (P1) and
Definition~\ref{gal-realization-kappa-approximations}(1)(b), $f^{cx}_m=\iota_m(\lambda_m(c))f^x_m=\tau_m(c)f^x_m$, so
$\tau_m(\kappa(cx)(\psi))=\tau_m(c\,\kappa(x)(\psi))$ for every $m$, and $\kappa(cx)(\psi)=c\,\kappa(x)(\psi)$. Let
$a\in K$. By (P2) and Definition~\ref{gal-realization-kappa-approximations}(1)(a), $\tau_m(\kappa(x)(a\psi))=
af^x_m(\psi)=\tau_m(a\,\kappa(x)(\psi))$. Let $\sigma\in G_K$. By (P4) and (P1), $f^{\sigma x}_m=\sigma_{B_m}\circ
f^x_m$. As $\sigma$ preserves $B_{\mathrm{dR}}^+$, Definition~\ref{gal-realization-kappa-approximations}(1)(c) gives
$\tau_m(\kappa(\sigma x)(\psi))=\sigma_{B_m}(\tau_m(\kappa(x)(\psi)))=\tau_m(\sigma(\kappa(x)(\psi)))$.

\emph{Item (b).} For $x$, $x'$, bases and $\psi$ as in (b), (P5) applied to $f^x_m$, $f^{x'}_m$ and $f^{x\cup x'}_m$ gives
$\tau_m(\kappa(x\cup x')(\psi))=\sum_{s,s'}f^x_m(b_s^*)\,f^{x'}_m(b'^*_{s'})\,\psi(b_s\cup b'_{s'})$. As $\tau_m$ is a
ring map with $\tau_m(a)=a$ for $a\in K$, this is $\tau_m$ of the right side of (b). By (P6),
$\tau_m(\kappa(1)(\psi))=\psi(1)=\tau_m(\psi(1))$.

\emph{Item (c).} By (P7), $\tau_m(\kappa(x)(\psi))=\psi(\omega)=\tau_m(\psi(\omega))$.

\emph{Item (d).} Let $f'_m$ be the map with $\rho_m(x',f'_m)$ for $(Y',\mathcal L')$ and $f_m$ the one with
$\rho_m(h^*x',f_m)$ for $(Y,\mathcal L)$. By (P8), $f_m(\psi)=f'_m(\psi\circ h_K^*)$, so
$\tau_m(\kappa_{Y,\mathcal L}(h^*x')(\psi))=\tau_m(\kappa_{Y',\mathcal L'}(x')(\psi\circ h_K^*))$ for every $m$.

\emph{Item (e).} If $n<r$, then $y=t^{r-n}\kappa(x)(\psi)$ lies in $B_{\mathrm{dR}}^+$ and $t^r\kappa(x)(\psi)=t^ny$. Let
$n\geq r$. Every $x_\nu\in\mathrm{Fil}^nH^j_\nu(Y)$ is the class of a cocycle with values in $\mathcal O^{[n]}$, so
(P9) for $m=n$ gives $f^x_n(\psi)\in J_n^{[n-r]}$, which is $\tau_n(t)^{n-r}B_n$ by
Definition~\ref{gal-realization-kappa-approximations}(1)(a). As $\tau_n$ is surjective, $f^x_n(\psi)=\tau_n(y_0t^{n-r})$
for some $y_0\in B_{\mathrm{dR}}^+$, so $\kappa(x)(\psi)-y_0t^{n-r}$ lies in the kernel $t^{n+1}B_{\mathrm{dR}}^+$ of
$\tau_n$; write it as $t^{n+1}z$. Then $t^r\kappa(x)(\psi)=t^n(y_0+t^{r+1}z)$.

If $Y_K$ is connected of dimension $d$, it is smooth of relative dimension $d$ over $K$, and
Lemma~\ref{gal-realization-kappa-de-rham}(2) and~(3) give the last sentence of (3).

(4) Let $f_1\geq1$, and let $N$ be as in Lemma~\ref{gal-realization-frobenius-congruence} for $f_1$. It is a multiple of
$f_1$ with $(x\bmod\ell)^{\ell^N}=g(x)$ for $x\in\mathcal O_K$, and the maps $\mu_\nu$ and $\nu_\nu$ of
Definition~\ref{gal-realization-frobenius-teichmuller} are defined for every $\nu$. Let $Y$ be a smooth projective
$\mathcal O_K$-scheme and $j\geq0$. Index by $\nu\geq0$, where the groups for $\nu=0$, defined over $W_0=0$ and
$A_{\mathrm{cris}}/\ell^0=0$, are $0$. Write $S_\nu=H^j_\nu(Y/W)$, $L=L^j(Y)=\varprojlim_\nu S_\nu$,
$T_\nu=H^{\mathrm T,j}_\nu(Y)$ and $H=H'^j_\infty(Y)=\varprojlim_\nu H'^j_\nu(Y)$. By
Lemma~\ref{gal-realization-frobenius-teichmuller-maps}(1), the $b_\nu$, $\mu_\nu$ and $\nu_\nu$ are compatible with the
reductions.

\emph{Step 1: the tower.} By Lemma~\ref{gal-realization-frobenius-witt-bockstein}, applied to $j$ and to $j+1$, the
pairs $(S,H^{j+1}_\bullet(Y/W))$ and $(H^{j+1}_\bullet(Y/W),H^{j+2}_\bullet(Y/W))$ have Bockstein sequences, and
$S_1$ and $H^{j+1}_1(Y/W)$ have finite length. In particular $\ell^\nu S_\nu=0$. By
Lemma~\ref{gal-realization-frobenius-bockstein-limits}(1) every $S_\nu$ has finite length, hence is Artinian. As $W$ is
$\ell$-adically complete, Lemma~\ref{gal-realization-frobenius-bockstein-limits}(2) shows that $L$ and $L^{j+1}(Y)$ are
finitely generated. As $W$ is Noetherian, Lemma~\ref{gal-realization-frobenius-bockstein-limits}(3) gives $c\geq0$ with
$\ell^cS_\nu$ in the image of $L\to S_\nu$ for every $\nu$, and by
Lemma~\ref{gal-realization-frobenius-bockstein-limits}(4) every $x\in L$ with $x_\nu=0$ lies in $\ell^\nu L$.

\emph{Step 2: the map $\alpha$.} The ring $W$ is a discrete valuation ring in which $\ell$ is irreducible, and
$A_{\mathrm{cris}}$ is $\ell$-torsion-free, hence flat over $W$, and $\ell$-adically complete and separated
\cite{Fon94}. By Lemma~\ref{gal-realization-frobenius-base-change}, $A_{\mathrm{cris}}\otimes_WS_\nu\to T_\nu$,
$a\otimes y\mapsto a\,b_\nu(y)$, is bijective. Lemma~\ref{gal-realization-frobenius-alpha}, for $A=A_{\mathrm{cris}}$,
$\psi=\phi^N$, $g_\nu=b_\nu$, $H_\nu=H'^j_\nu(Y)$ and the $\nu_\nu$, gives an additive map $\alpha:H\to
A_{\mathrm{cris}}\otimes_WL$ with $E_\nu(\alpha(x))=\ell^c\nu_\nu(x_\nu)$, where $E_\nu(a\otimes y)=a\,b_\nu(y_\nu)$,
and shows that $w\in A_{\mathrm{cris}}\otimes_WL$ is $0$ if $E_\nu(w)=0$ for every $\nu$. Put $e=c$.

\emph{Step 3: the map $\beta^\circ$.} Put $\beta^\circ(w)=(\mu_\nu(E_\nu(w)))_\nu$. As $E_\nu$ and $\mu_\nu$ are
$A_{\mathrm{cris}}$-linear and compatible with the reductions, this is an $A_{\mathrm{cris}}$-linear map
$A_{\mathrm{cris}}\otimes_WL\to H$.

\emph{Step 4: item (a).} Let $c'\in A_{\mathrm{cris}}$ and $x\in H$. As $\nu_\nu$ is $\phi^N$-semilinear,
$E_\nu(\alpha(c'x))=\ell^e\phi^N(c')\nu_\nu(x_\nu)=E_\nu(\phi^N(c')\alpha(x))$ for every $\nu$, so
$\alpha(c'x)=\phi^N(c')\alpha(x)$. Let $\sigma\in G_K$. As $\sigma$ fixes the image of $W$
(Lemma~\ref{gal-realization-invariants}), $\sigma\otimes1$ is defined on $A_{\mathrm{cris}}\otimes_WL$, and by
Lemma~\ref{gal-realization-frobenius-teichmuller-maps}(3), $E_\nu((\sigma\otimes1)(a\otimes y))=\sigma(a)\,b_\nu(y_\nu)=
\sigma(a\,b_\nu(y_\nu))$; so $E_\nu\circ(\sigma\otimes1)=\sigma\circ E_\nu$. By the same lemma,
$E_\nu(\alpha(\sigma x))=\ell^e\nu_\nu(\sigma x_\nu)=\sigma(\ell^e\nu_\nu(x_\nu))=E_\nu((\sigma\otimes1)(\alpha(x)))$, so
$\alpha(\sigma x)=(\sigma\otimes1)(\alpha(x))$. Finally, by Lemma~\ref{gal-realization-frobenius-teichmuller-maps}(2),
$\beta^\circ(\alpha(x))_\nu=\mu_\nu(\ell^e\nu_\nu(x_\nu))=\ell^e\phi^N(x_\nu)$ for every $\nu$, so
$\beta^\circ(\alpha(x))=\ell^e\phi^N(x)$.

\emph{Step 5: item (b).} Let $\nu\geq1$ and $y\in L$. By Lemma~\ref{gal-realization-frobenius-teichmuller-maps}(4),
$\Gamma_W\circ M_N$ is a Frobenius twist of level $N$ whose map $m$ is $\mu_\nu\circ b_\nu$, and
$\beta^\circ(1\otimes y)_\nu=\mu_\nu(b_\nu(y_\nu))=m(y_\nu)$.
% The open parts of (3) and (4) are the open steps of Lemmas gal-realization-kappa-approximation-system,
% gal-realization-kappa-truncated-periods-exist, gal-realization-kappa-de-rham, gal-realization-frobenius-base-change
% and gal-realization-frobenius-witt-bockstein. Neither (3) nor (4) is a published theorem: their statements concern H_infty,
% H'_infty, HdR and L^j, which the archive does not identify with crystalline or de Rham cohomology.
% Lean: (1) to (2)(e) are written out above. (1) is B.Main l, proved as B.main = main_of (D.out l) (global_of,
% local_of, level_all); it is Lemma gal-realization-ingredients-bounds, whose proof is an open step there. The main
% theorem applies it to a
% model through GlueB'.atlas_of (B.P.atlas, Lemma gal-realization-ingredients-atlas-scheme).
% (2)(a) is Pb.isO1_exists and Pb.isTwistingSheaf_pullback; (2)(b) is Pb.firstChern_acris_family_exists
% (FirstChernClassFactorial.exists_C1H) with Pb.sSup_𝒱D; (2)(c) is B.firstChern_compatible_l =
% FirstChernComparison.main (st1, st2, st3); (2)(d) is E.Main.firstChern_unique
% (EtChern.chern_class_unique_of_lift_and_log with LineLift.liftEx and CrisLog.logEx).
% (2)(e) is E.Main.firstChern_family = firstChern_family_of_lifts lineBC lineMap chern_class_compatible_system
% etale_chern_class_eq_kummer_class fKum fRes fUntw, the field C1 of eFam; fKum = FKum.fKum rootPow lcLog and
% etale_chern_class_eq_kummer_class (etSpec, resGood, stCol, firstChern_primitive_kummer, kumPow) are the core.
% (3) is Pb.crystalline_deRham_firstChern_exists = Kap.kap_abs (facts l K) hS P, for S from levSys_ex' and P from
% rhoSys_ex', giving KapL L with LieComplementChartContDiffAt L. The proof above follows kap_abs: exists_lim (the
% limit, through hS.cpl), ext_P (uniqueness), and the fields lin (P.smul, iota_lam), lin' (P.lin', toB_K), act (P.act,
% sigmaB_toB, Per.Facts.sigma_mem), nat (P.nat), fil (P.fil at level n, Lev.Good.pdP_B for J_n^[n-r], toB_surj and
% toB_ker, with the case n < r separately), cup (P.cup, toB_K), one (P.one) and c1 (P.c1). The archive's kappa is one
% family over all models and tautological line data (dat Pb.Amp), with Mo u hu a b = PUnit, so (d) asks nothing of h
% and the line data. Finiteness is Pb.m_fin (k_fin);
% cupR = cupK is Pb.cupR_eq; the last sentence is Pb.isFirstChernDeRham_exists and Pb.powK_ne_zero.
% (4) is crystallineComparison.pinEx = PinP.pin_fixed: N from PinP.teich_dvd; F from Pb.fac_ex, whose pin_all gives
% e and alpha (Step 1: LsG.ls_seq, smooth_model_finite_length_link, finite_length_of_lseq_first, lw_fin,
% bounded_cokernel_limit_to_level, limit_divisible_of_vanishing_level; Step 2: Fa.pin_ex with bc_tens, flatW and
% adicA); beta = betaP = Fa.beta (Step 3); lin = Fa.Pin.lin, fac = Fa.Pin.fac with Mid.mu_nu, gal = Fa.Pin.op with
% bW_gal and nu_gal, and galT_tmul (Step 4); (b) is PinP.pin_of with PinP.pin_lev, the twist being PinP.mm (Step 5).
% The congruence of (4) is Lemma gal-realization-frobenius-congruence(1).
\end{proof}
\leanref{Development/WeilBound/CoreStatements/PadicHodgeComparison/EtaleCohomologyComparison/PerfectoidPullbacks/LevelChainValuations/ExistsLim.lean}{CrystallineSemistableComparison.Pb.Kap.kap_abs}
\leanref{Development/WeilBound/CoreStatements/PadicHodgeComparison/EtaleCohomologyComparison/PerfectoidPullbacks/LevelChainValuations/ExistsLim.lean}{CrystallineSemistableComparison.Pb.Kap.exists_lim}
\leanref{Development/WeilBound/CoreStatements/PadicHodgeComparison/EtaleCohomologyComparison/PerfectoidPullbacks/LevelChainValuations/ExistsLim.lean}{CrystallineSemistableComparison.Pb.Kap.ext_P}
\leanref{Development/WeilBound/CoreStatements/PadicHodgeComparison/EtaleCohomologyComparison/TeichmullerRootsUnity/PinFixed.lean}{CrystallineSemistableComparison.crystallineComparison.PinP.pin_fixed}
\leanref{Development/WeilBound/CoreStatements/PadicHodgeComparison/EtaleCohomologyComparison/PerfectoidPullbacks/FrobeniusFactorisation/PhiAIter.lean}{CrystallineSemistableComparison.Pb.fac_ex}
\leanref{Development/WeilBound/CoreStatements/PadicHodgeComparison/EtaleCohomologyComparison/PerfectoidPullbacks/FrobeniusFactorisation/PhiAIter.lean}{CrystallineSemistableComparison.Pb.pin_all}
\leanref{Development/WeilBound/CoreStatements/PadicHodgeComparison/EtaleCohomologyComparison/PerfectoidPullbacks/FrobeniusFactorisation/ExistsAlpha.lean}{CrystallineSemistableComparison.Pb.Fa.beta}
\leanref{Development/WeilBound/CoreStatements/PadicHodgeComparison/EtaleCohomologyComparison/PerfectoidPullbacks/FrobeniusFactorisation/ExistsAlpha.lean}{CrystallineSemistableComparison.Pb.Fa.Pin.lin}
\leanref{Development/WeilBound/CoreStatements/PadicHodgeComparison/EtaleCohomologyComparison/PerfectoidPullbacks/FrobeniusFactorisation/ExistsAlpha.lean}{CrystallineSemistableComparison.Pb.Fa.Pin.fac}
\leanref{Development/WeilBound/CoreStatements/PadicHodgeComparison/EtaleCohomologyComparison/PerfectoidPullbacks/FrobeniusFactorisation/ExistsAlpha.lean}{CrystallineSemistableComparison.Pb.Fa.Pin.op}
\leanref{Development/WeilBound/CoreStatements/PadicHodgeComparison/EtaleCohomologyComparison/TeichmullerRootsUnity/PinFixed.lean}{CrystallineSemistableComparison.crystallineComparison.PinP.pin_of}
\leanref{Development/WeilBound/CoreStatements/PadicHodgeComparison/EtaleCohomologyComparison/TeichmullerRootsUnity/PinFixed.lean}{CrystallineSemistableComparison.crystallineComparison.PinP.pin_lev}

\begin{lemma}\label{gal-realization-inputs}
Let $\ell$, $C$, $A_{\mathrm{cris}}$, $K$, $\mathcal O_K$, $\bar K$ and $G_K$ be as in
Definition~\ref{gal-realization-period-datum}, let $\mathbf F_q$, $q=\ell^f$, be the residue field of $K$,
and $g$ as in Definition~\ref{gal-realization-frobenius-datum}.
\begin{enumerate}
\item For every $\mathcal O_K$-scheme $X$, the comparison maps $\eta_\nu$ of
  Definition~\ref{gal-realization-eta} commute with the actions of $G_K$, satisfy
  $\eta_\nu(x\cup x')=\eta_\nu(x)\cup\eta_\nu(x')$, and there is $w\in H^0(X_{\bar K},\mathbf Z_\ell)$ with
  $\eta_\nu(w)=1$ for all $\nu$.
\item Every smooth projective $\mathcal O_K$-scheme with connected generic fibre has a period datum
  (Definition~\ref{gal-realization-period-datum}).
\item There is a multiple $N$ of $f$ with $(x\bmod\ell)^{\ell^N}=g(x)$ for all $x\in\mathcal O_K$ such
  that every smooth projective $\mathcal O_K$-scheme has a Frobenius datum of level $N$ in every degree
  (Definition~\ref{gal-realization-frobenius-datum}).
\end{enumerate}
\end{lemma}
\leanref{Development/WeilBound/CoreStatements/PadicHodgeComparison/EtaleCohomologyComparison/KummerCupProducts/EtaleChernClassEqKummerClass.lean}{CrystallineSemistableComparison.E.Main.fEqui}
\leanref{Development/WeilBound/CoreStatements/PadicHodgeComparison/EtaleCohomologyComparison/KummerCupProducts/EtaleChernClassEqKummerClass.lean}{CrystallineSemistableComparison.E.Main.fCup}
\leanref{Development/WeilBound/CoreStatements/PadicHodgeComparison/EtaleCohomologyComparison/KummerCupProducts/EtaleChernClassEqKummerClass.lean}{CrystallineSemistableComparison.E.Main.fRoot}
\leanref{Development/WeilBound/CoreStatements/PadicHodgeComparison/EtaleCohomologyComparison/KummerCupProducts/EtaleChernClassEqKummerClass.lean}{CrystallineSemistableComparison.E.Main.firstChern_family}
\leanref{Development/AutomorphicArithmetic/CommutativeAlgebra/EtaleCohomology/EPin.lean}{UpperTheory.Ru.crystallineComparison.ePin}
\leanref{Development/WeilBound/CoreStatements/PadicHodgeComparison/EtaleCohomologyComparison/TeichmullerRootsUnity/EndEx.lean}{CrystallineSemistableComparison.crystallineComparison.endEx}
\leanref{Development/WeilBound/CoreStatements/PadicHodgeComparison/EtaleCohomologyComparison/TeichmullerRootsUnity/PinFixed.lean}{CrystallineSemistableComparison.crystallineComparison.pinEx}
\leanref{Development/WeilBound/CoreStatements/PadicHodgeComparison/EtaleCohomologyComparison/TeichmullerRootsUnity/BInEx.lean}{CrystallineSemistableComparison.crystallineComparison.bInEx}
\leanref{Development/AutomorphicArithmetic/CommutativeAlgebra/EtaleCohomology/CIn.lean}{UpperTheory.Ru.crystallineComparison.cIn}
\leanref{Development/AutomorphicSpectrum/Proofs/BlockDecomposition/PullPullOfPullC.lean}{GlobalTheory.CompanionTheory.FEquiProof.equi_at}
\leanref{Development/AutomorphicSpectrum/Proofs/BlockDecomposition/LcMorCup.lean}{GlobalTheory.CompanionTheory.EtCupF.eG_cup}
\leanref{Development/WeilBound/CoreStatements/PadicHodgeComparison/EtaleCohomologyComparison/KummerCupProducts/CupLeib.lean}{CrystallineSemistableComparison.E.morCup}
\leanref{Development/AbelianVarietyTheory/Proofs/DivisorProofs/TateTwistResolutionPairingSpecAndCupUnique.lean}{GlobalTheory.EtC.TwPairUniq.main}
\leanref{Development/AutomorphicSpectrum/Proofs/BlockDecomposition/PrimResZetaX.lean}{GlobalTheory.CompanionTheory.FRootCls.rootSys}
\leanref{Development/AutomorphicSpectrum/Proofs/BlockDecomposition/PrimResZetaX.lean}{GlobalTheory.CompanionTheory.FRootCls.eM_rootSys}
\leanref{Development/WeilBound/CoreStatements/PadicHodgeComparison/EtaleCohomologyComparison/KummerCupProducts/EtaleChernClassEqKummerClass.lean}{CrystallineSemistableComparison.E.Main.etRoot}
\leanref{Development/WeilBound/CoreStatements/PadicHodgeComparison/EtaleCohomologyComparison/TeichmullerRootsUnity/EndExOf.lean}{CrystallineSemistableComparison.crystallineComparison.endEx_of}
\leanref{Development/WeilBound/CoreStatements/PadicHodgeComparison/EtaleCohomologyComparison/PerfectoidPullbacks/SumCoord.lean}{CrystallineSemistableComparison.Pb.pSeam}
\leanref{Development/WeilBound/CoreStatements/PadicHodgeComparison/EtaleCohomologyComparison/PerfectoidPullbacks/SumCoord.lean}{CrystallineSemistableComparison.Pb.pmap}
\leanref{Development/WeilBound/CoreStatements/PadicHodgeComparison/EtaleCohomologyComparison/PerfectoidPullbacks/SumCoord.lean}{CrystallineSemistableComparison.Pb.kap_pow}
\leanref{Development/WeilBound/CoreStatements/PadicHodgeComparison/EtaleCohomologyComparison/PerfectoidPullbacks/SumCoord.lean}{CrystallineSemistableComparison.Pb.pow_ex}
\leanref{Development/WeilBound/CoreStatements/PadicHodgeComparison/EtaleCohomologyComparison/BocksteinBaseChange/AtlasGroupActions/FirstChernCompatGlue.lean}{CrystallineSemistableComparison.glueB}
\leanref{Development/WeilBound/CoreStatements/PadicHodgeComparison/EtaleCohomologyComparison/BocksteinBaseChange/AtlasGroupActions/FirstChernCompatGlue.lean}{CrystallineSemistableComparison.GlueB'.bmap}
\leanref{Development/WeilBound/CoreStatements/PadicHodgeComparison/EtaleCohomologyComparison/BocksteinBaseChange/AtlasGroupActions/AtlasOf.lean}{CrystallineSemistableComparison.GlueB'.f_inj}
\leanref{Development/WeilBound/CoreStatements/PadicHodgeComparison/EtaleCohomologyComparison/BocksteinBaseChange/AtlasGroupActions/AtlasOf.lean}{CrystallineSemistableComparison.GlueB'.f_surj}
\leanref{Development/AutomorphicArithmetic/CommutativeAlgebra/EtaleCohomology/PhiAIter.lean}{UpperTheory.Ru.crystallineComparison.keyStep}

\begin{proof}
(1) Fix $X$ and $\nu$, let $\mathbf Z/\ell^\nu\to I^\bullet$ and $\varepsilon_y$ be as in
Lemma~\ref{gal-realization-edge-map} for the split diagram $(U_{(A,y)})$ of
Definition~\ref{gal-realization-eta}, so that $\eta_\nu=\lambda_*\circ\varepsilon_y$.

\emph{Galois actions.} Let $\sigma\in G_K$; it acts on $H^k(X_{\bar K},\mathbf Z/\ell^\nu)$ by pullback along
$u=1_X\times\operatorname{Spec}\sigma$. Apply Lemma~\ref{gal-realization-edge-map}(2) with this split
diagram and its points on both sides, $m=t_\sigma$, $g_{(A,y)}=\operatorname{Spec}\tau_A:U_{(A,y)}\to
U_{(A^\sigma,y^\sigma)}$ and $u$, in the notation of Lemma~\ref{gal-realization-twist}. The condition
$U_{t_\sigma(\varphi)}g_{c'}=g_cU_\varphi$ for maps $\varphi:c\to c'$ is the naturality in
Lemma~\ref{gal-realization-twist}(1), and $uy_c=y_{t_\sigma(c)}g_c$ is the identity of
Lemma~\ref{gal-realization-twist}(3). So $t_\sigma^*\circ\varepsilon_y=\varepsilon_y\circ u^*$, where
$t_\sigma^*$ is induced by $t_\sigma$ and the pullbacks of locally constant functions along the
$\operatorname{Spec}\tau_A$. By Lemma~\ref{gal-realization-twist}(2), applying the $\lambda_A$ to the values
of a cochain turns $t_\sigma^*$ into the map induced by $t_\sigma$ and the $\tau_A$, which is the action of
$\sigma$ by Lemma~\ref{gal-realization-twist}(3). Hence $\eta_\nu(u^*x)=\sigma(\eta_\nu(x))$.

\emph{Cup products.} The cup product of $H^*(X_{\bar K},\mathbf Z/\ell^\nu)$ is induced by bilinear maps
$\mu$ as in Lemma~\ref{gal-realization-cup-staircase}, for which $\mu(1,1)=1$ over every scheme with a
primitive $\ell^\nu$-th root of unity; such $\mu$ exist, and all of them induce the same product on
cohomology. Each $U_{(A,y)}$ receives $\bar K$, so it has such a root, and
Lemma~\ref{gal-realization-cup-staircase} gives $\varepsilon_y(x\cup x')=\varepsilon_y(x)\cup
\varepsilon_y(x')$. The $\lambda_A$ are ring homomorphisms and commute with the maps of test algebras
(Lemma~\ref{gal-realization-locally-constant}(2)), so applying them to the values of cochains commutes with
the cup product of Definition~\ref{gal-realization-split-diagrams}(1). Hence
$\eta_\nu(x\cup x')=\eta_\nu(x)\cup\eta_\nu(x')$.

\emph{Units.} Let $1_\nu\in H^0(X_{\bar K},\mathbf Z/\ell^\nu)$ be the class of the constant function $1$.
These classes are compatible under reduction, so they form an element $w\in H^0(X_{\bar K},\mathbf Z_\ell)$.
The class $1_\nu$ is represented by $1\in\Gamma(X_{\bar K},\mathbf Z/\ell^\nu)\subset\Gamma(X_{\bar K},I^0)$,
so for $j=0$ Lemma~\ref{gal-realization-edge-map}(1) gives $\varepsilon_y(1_\nu)$ as the class of
$(y_c^*1)_c=(1)_c$, the unit of $H^0(\mathrm{El}(X^\sharp),L_{\ell^\nu})$. Since $\lambda_A(1)=1$,
$\eta_\nu(w)=1$.

For (2) and (3), use the notation of Definition~\ref{gal-realization-period-thickenings} for $\mathcal
O=\mathcal O_K$ and its inclusion into $\mathcal O_C$. The compatible system $\varepsilon$ of $\ell$-power roots of
unity with $t=\log[\varepsilon]$ satisfies $\theta([\varepsilon])=1$ and $\theta([\varepsilon^{1/\ell}])\neq1$; let
$\mu=[\varepsilon]-1$, and let $a$ be a unit of $A_{\mathrm{cris}}$ with $t=\mu a$
(Lemma~\ref{gal-realization-period-t-mu}). Lemma~\ref{gal-realization-period-ingredients} is referred to as (I).

(2) Let $Y$ be a smooth projective $\mathcal O_K$-scheme whose generic fibre $Y_K$ is connected of dimension $d$.
As $Y$ is projective over $\mathcal O_K$, (I)(2)(a) for $R=\mathcal O_K$ gives a tautological line datum $\mathcal L$ on
$Y$; its map $f$ is \'etale. Put $H^j=H^j_\infty(Y)$, $H=\bigoplus_jH^j$, $\beta=\varprojlim_\nu\beta_\nu$ and
$D^j=H^j_{\mathrm{dR}}(Y_K)$ with the cup products, and let $\kappa=\kappa_{Y,\mathcal L}$ be the maps of (I)(3). By
Lemma~\ref{gal-realization-period-compatible}(1), the
reductions commute with the maps $\sigma$ and preserve cup products and units, so the $H^j$ are
$A_{\mathrm{cris}}$-modules with the cup products, the unit $1$ and additive maps $\sigma$ formed componentwise
(Definition~\ref{gal-realization-period-thickenings}(5)). By Lemma~\ref{gal-realization-period-compatible}(2),
$\beta$ is $A_{\mathrm{cris}}$-linear, $G_K$-equivariant and multiplicative, and $\beta(1)=1$; here
$E^j=E^j(Y)$. The $D^j$ are finite-dimensional by (I)(3), and the cup products are $K$-bilinear.

Condition (1). Fix $j$. As $Y$ is smooth, separated and quasi-compact over $\mathcal O_K$, $F_Y$ has a Zariski atlas of
some dimension $d'$ (Lemma~\ref{gal-realization-ingredients-atlas-scheme}); let $N$ be as in (I)(1) for $d'$ and $j$,
applied to $F_Y$, whose groups are $H^j_\nu(Y)$ and $H^j(Y^\sharp,A_{\mathrm{cris}}/\ell^\nu)$. Then $\mu^N$ kills the
kernel of $\beta_\nu$ and $\mu^NH^j(Y^\sharp,A_{\mathrm{cris}}/\ell^\nu)\subset\beta_\nu(H^j_\nu(Y))$ for every $\nu$.
Lemma~\ref{gal-realization-period-limit} for the $\beta_\nu$ and
$s=\mu^N$ shows that $\mu^{2N}$ kills the kernel of $\beta$ on $H^j$ and that $\mu^{2N}E^j\subset\beta(H^j)$. Since
$t^{2N}=a^{2N}\mu^{2N}$, the kernel is killed by $t^{2N}$, and if $\beta(x)=\mu^{2N}y$ then $\beta(a^{2N}x)=t^{2N}y$.
So (1) holds with $m=n=2N$.

Conditions (2) and (3) are (I)(3)(a) and (I)(3)(b).

Condition (4). By Lemma~\ref{gal-realization-ingredients-tautological}(4), $Y$ is the union of the open subschemes over
which $f$ has a section, so by (I)(2)(b) there is $x\in H^2$ whose components form a first Chern family of $\mathcal L$,
and by (I)(2)(c) $\beta(x)$ is a first Chern class of $\mathcal L$. By (I)(2)(e) there is $v\in T^2$ such that
$t\,\eta(v)=(t\,\eta_\nu(v))_\nu$ is a first Chern class of $\mathcal L$. As $f$ is \'etale, (I)(2)(d) gives
$\beta(x)=t\,\eta(v)$. The base change $\mathcal L_K$ of $\mathcal L$ to $Y_K$ is tautological by (I)(2)(a) for
$\mathcal O_K\to K$, so by
(I)(3) it has a first Chern class $\omega\in D^2$ with $\omega^d\neq0$, and $\kappa(x)(\psi)=\psi(\omega)$ for all
$\psi$ by (I)(3)(c). Let $x^0=1$ and $x^{k+1}=x^k\cup x$. Induction on $k$ gives $\kappa(x^k)(\psi)=\psi(\omega^k)$
for all $\psi$: for $k=0$ this is the last part of (I)(3)(b), and for the step, (I)(3)(b) with bases $(b_s)$ of
$D^{2k}$ and $(b'_{s'})$ of $D^2$ gives
\[
 \kappa(x^{k+1})(\psi)=\sum_{s,s'}b_s^*(\omega^k)\,b'^*_{s'}(\omega)\,\psi(b_s\cup b'_{s'})=\psi(\omega^k\cup\omega),
\]
by bilinearity. Choose $\psi$ with $\psi(\omega^d)=1$. Then $\kappa(x^d)(\psi)=1$, a unit of $B_{\mathrm{dR}}^+$.

(3) Let $N$ be as in (I)(4) for $f_1=f$. Then $N$ is a multiple of $f$ and $(x\bmod\ell)^{\ell^N}=g(x)$ for
$x\in\mathcal O_K$, so
$\rho_N$ is defined (Definition~\ref{gal-crystalline-test-algebras}(4)). Let $Y$ be a smooth projective $\mathcal
O_K$-scheme and $j\geq0$. Put $H=H'^j_\infty(Y)$ with $\phi_H=\varprojlim_\nu\phi$ and the action of $G_K$, let
$b_\nu$ be the projection to $H'^j_\nu(Y)$ followed by $\beta'_\nu$, let $L=L^j(Y)$, and let $e$, $\alpha$ and
$\beta=\beta^\circ$ be as in (I)(4) for $Y$ and $j$. Since $G_K$ fixes the image of $W$
(Lemma~\ref{gal-realization-invariants}), $\sigma\otimes1$ is defined on $A_{\mathrm{cris}}\otimes_WL$. The maps
$b_\nu$ are $A_{\mathrm{cris}}$-linear, and $b_\nu\phi_H=\phi b_\nu$ and $b_\nu\sigma=\sigma b_\nu$ by
Lemma~\ref{gal-realization-period-compatible}(2). The conditions on $\alpha$ and $\beta$ are (I)(4)(a).

Condition (1). Let $N_1$ be as in (I)(1) for $j$ and the dimension of a Zariski atlas of $F_Y$
(Lemma~\ref{gal-realization-ingredients-atlas-scheme}), applied to $F_Y$. Lemma~\ref{gal-realization-period-limit} for the
$\beta'_\nu$ and $s=\mu^{N_1}$ gives, for every compatible family $(y_\nu)$, some $x\in H$ with $b_\nu(x)=
\mu^{2N_1}y_\nu$ for all $\nu$, so that $b_\nu(a^{2N_1}x)=t^{2N_1}y_\nu$; and if $b_\nu(x)=0$ for all $\nu$, then
$\mu^{2N_1}x=0$, so $t^{2N_1}x=0$. So (1) holds with $u=2N_1$.

Condition (2). Let $y\in L$ and $\nu\geq1$. By (I)(4)(b) there is a Frobenius twist of level $N$ whose map $m$
satisfies $b_\nu(\beta(1\otimes y))=\beta'_\nu(m(y_\nu))$. By Lemma~\ref{gal-realization-frobenius-fixed} there is
$X_\nu\in H^j(Y_g,A_{\mathrm{cris}}/\ell^\nu)$ with $\phi^fX_\nu=F^*X_\nu$ and $\beta'_\nu(m(y_\nu))=\rho_N^*X_\nu$.
\end{proof}

\begin{definition}\label{gal-crystalline-relation-package}
Let $\ell$ be a prime, $C$ and $A_{\mathrm{cris}}$ as in Definition~\ref{gal-crystalline-test-algebras}, with its
Frobenius $\phi$ and the map $\mathbf Z_\ell=W(\mathbf F_\ell)\to A_{\mathrm{cris}}$, and
$t=\log[\varepsilon]\in A_{\mathrm{cris}}$ for a compatible system $\varepsilon$ of $\ell$-power roots of unity.
Let $K$ be a finite extension of $\mathbf Q_\ell$ in $C$, $\bar K$ the algebraic closure of $K$ in $C$, and
$G_K=\operatorname{Gal}(\bar K/K)$, acting on $A_{\mathrm{cris}}$. A \emph{crystalline period package} for $K$ is a
commutative $\mathbf Q_\ell$-algebra $B$ with a ring map $c:A_{\mathrm{cris}}\to B$, a $\mathbf Q_\ell$-algebra
endomorphism $\phi_B$ of $B$, and $\mathbf Q_\ell$-algebra endomorphisms $\sigma_B(g)$ of $B$ for $g\in G_K$, such
that:
\begin{enumerate}
\item $c(t)$ is a unit of $B$;
\item for every $x\in B$ there are $r\geq0$ and $a\in A_{\mathrm{cris}}$ with $c(t)^rx=c(a)$;
\item every $a\in A_{\mathrm{cris}}$ with $c(a)=0$ satisfies $t^ra=0$ for some $r\geq0$;
\item $\phi_B(c(a))=c(\phi(a))$ and $\sigma_B(g)(c(a))=c(g(a))$ for $a\in A_{\mathrm{cris}}$ and $g\in G_K$;
\item the structure map $\mathbf Q_\ell\to B$ sends every $z\in\mathbf Z_\ell$ to $c(z)$.
\end{enumerate}
By (1) and (2), two ring maps from $B$ that agree on $c(A_{\mathrm{cris}})$ are equal; so by (4) $\sigma_B$ is an
action of $G_K$ on $B$. Fontaine's ring $B_{\mathrm{cris}}=A_{\mathrm{cris}}[1/t]$, with the localization map, its
Frobenius and its action of $G_K$, is a crystalline period package for $K$.
\end{definition}
\leanref{Development/AutomorphicArithmetic/CommutativeAlgebra/EtaleCohomology/VecZero.lean}{UpperTheory.Ru.crystallineComparison.PBg}
\leanref{Development/AutomorphicArithmetic/CommutativeAlgebra/EtaleCohomology/VecZero.lean}{UpperTheory.Ru.crystallineComparison.PB}
\leanref{Development/AutomorphicArithmetic/CommutativeAlgebra/EtaleCohomology/VecZero.lean}{UpperTheory.Ru.crystallineComparison.IsInv}

\begin{lemma}\label{gal-crystalline-relation-package-embedding}
Let $B$ be a crystalline period package for $K$ (Definition~\ref{gal-crystalline-relation-package}), and
$B_{\mathrm{dR}}\supset A_{\mathrm{cris}}$ Fontaine's field of de Rham periods, with its action of $G_K$. Then
$c(a)=c(a')$ implies $a=a'$, and there is an injective ring map $\iota:B\to B_{\mathrm{dR}}$ with $\iota(c(a))=a$
for $a\in A_{\mathrm{cris}}$, $\iota(z)=z$ for $z\in\mathbf Q_\ell$, and $\iota(\sigma_B(g)x)=g(\iota(x))$ for
$g\in G_K$ and $x\in B$.
\end{lemma}
\leanref{Development/AutomorphicArithmetic/CommutativeAlgebra/EtaleCohomology/GalAwayCongr.lean}{UpperTheory.Ru.crystallineComparison.pbIota}

\begin{proof}
The inclusion $A_{\mathrm{cris}}\subset B_{\mathrm{dR}}$ is injective and $t\neq0$ in the field $B_{\mathrm{dR}}$
\cite{Fon94a}. By (1), (2) and (3) of Definition~\ref{gal-crystalline-relation-package}, $c$ makes $B$ the
localization of $A_{\mathrm{cris}}$ at the powers of $t$: $c(t)$ is a unit, every element is $c(a)c(t)^{-r}$, and
$c(a)=c(a')$ if and only if $t^r(a-a')=0$ for some $r$. As $t$ is a unit of $B_{\mathrm{dR}}$, the inclusion extends
to a ring map $\iota:B\to B_{\mathrm{dR}}$ with $\iota(c(a))=a$; it sends $x$ with $c(t)^rx=c(a)$ to $t^{-r}a$. If
$c(a)=c(a')$, then $a=\iota(c(a))=\iota(c(a'))=a'$. If $\iota(x)=0$, write $c(t)^rx=c(a)$; then $a=t^r\iota(x)=0$,
so $c(t)^rx=0$ and $x=0$. Let $g\in G_K$ and $c(t)^rx=c(a)$. Applying $\iota\circ\sigma_B(g)$ gives
$g(t)^r\iota(\sigma_B(g)x)=g(a)$ by (4), and applying $g$ to $t^r\iota(x)=a$ gives $g(t)^rg(\iota(x))=g(a)$; as
$g(t)\neq0$, $\iota(\sigma_B(g)x)=g(\iota(x))$. Finally, the ring maps $z\mapsto\iota(z)$ and $z\mapsto z$ from
$\mathbf Q_\ell$ to $B_{\mathrm{dR}}$ agree on $\mathbf Z_\ell$ by (5), hence on its fraction field $\mathbf Q_\ell$.
\end{proof}
% Lean: pbIota, with the private exists_lift (IsLocalization.lift for the powers of t, injective as above), lift_comm
% for the Galois actions and q_ext for Q_l; the map A_cris -> B_dR is DRPst.iC, injective with iC t a unit.

\begin{lemma}\label{gal-crystalline-relation-package-invariants}
Let $B$ be a crystalline period package for $K$ (Definition~\ref{gal-crystalline-relation-package}), and let
$\mathbf F_q$ be the residue field of $K$ and $W(\mathbf F_q)\to A_{\mathrm{cris}}$ as in
Lemma~\ref{gal-realization-invariants}. For every $x\in B$ with $\sigma_B(g)x=x$ for all $g\in G_K$ there are $a\geq0$
and $w\in W(\mathbf F_q)$ with $\ell^ax=c(w)$.
\end{lemma}
\leanref{Development/AutomorphicArithmetic/CommutativeAlgebra/EtaleCohomology/AbsInvOf.lean}{UpperTheory.Ru.crystallineComparison.InvC}
\leanref{Development/AutomorphicArithmetic/CommutativeAlgebra/EtaleCohomology/PullbackFinRelation.lean}{UpperTheory.Ru.crystallineComparison.pbInv}
\leanref{Development/AutomorphicArithmetic/CommutativeAlgebra/EtaleCohomology/PullbackFinRelation.lean}{UpperTheory.Ru.crystallineComparison.PBI.mainC}

\begin{proof}
For $g\in G_K$ let $\chi(g)\in\mathbf Z_\ell^\times$ be the cyclotomic character, so that $g(t)=\chi(g)t$. Write
$c(t)^nx=c(a)$ with $a\in A_{\mathrm{cris}}$. Applying $\sigma_B(g)$ and using $\sigma_B(g)x=x$ and (4) and (5) of
Definition~\ref{gal-crystalline-relation-package} gives $c(g(a))=\chi(g)^nc(t)^nx=c(\chi(g)^na)$, so
$g(a)=\chi(g)^na$ by Lemma~\ref{gal-crystalline-relation-package-embedding}. Hence $y=t^{-n}a\in B_{\mathrm{cris}}$ is
fixed by $G_K$, and by Lemma~\ref{gal-realization-invariants} there are $e\geq0$ and $w\in W(\mathbf F_q)$ with
$\ell^ey=w$, that is, $\ell^ea=t^nw$ in $B_{\mathrm{cris}}$, and so in $A_{\mathrm{cris}}$, as $A_{\mathrm{cris}}\to
B_{\mathrm{cris}}$ is injective by Lemma~\ref{gal-crystalline-relation-package-embedding} for $B_{\mathrm{cris}}$. Then $c(t)^n\ell^ex=c(\ell^ea)=c(t)^nc(w)$, and as $c(t)$ is a unit,
$\ell^ex=c(w)$.
\end{proof}
% Lean: PBI.mainC: hcinj, fix_aux for the invariance of fB(a) T^n in the Lean's B_cris, bCris_fixed_mem_maxUnramified
% (Lemma gal-realization-invariants, Fontaine's computation, proved in the Lean) with the choice of e and hK for w, and
% pullback_fin_relation with iC_inj to return to A_cris.

\begin{theorem}\label{gal-crystalline-realization}
Let $\ell$ be a prime, $C$ and $A_{\mathrm{cris}}$ as in Definition~\ref{gal-crystalline-test-algebras},
$t=\log[\varepsilon]\in A_{\mathrm{cris}}$ for a compatible system $\varepsilon$ of $\ell$-power roots of
unity, and $B_{\mathrm{cris}}=A_{\mathrm{cris}}[1/t]$, with Frobenius $\phi$ and $\phi(t)=\ell t$. Let $K$
be a finite extension of $\mathbf Q_\ell$ in $C$, with valuation ring $\mathcal O_K$, uniformizer
$\varpi$ and residue field $\mathbf F_q$, $q=\ell^f$, let $\kappa:\mathcal O_K\to\mathcal O_C$ be the
inclusion and $\bar K$ the algebraic closure of $K$ in $C$, and let $k\geq0$. Let $Y$ be a smooth
projective $\mathcal O_K$-scheme and $V=H^k(Y_{\bar K},\mathbf Q_\ell)$.
\begin{enumerate}
\item For every $\mathcal O_K$-scheme $X$ and $\nu\geq1$ there are additive maps
  $\eta_\nu:H^k(X_{\bar K},\mathbf Z/\ell^\nu)\to H^k(X^\sharp,A_{\mathrm{cris}}/\ell^\nu)$, compatible with
  reduction from $\ell^{\nu+1}$ to $\ell^\nu$, with $\phi\circ\eta_\nu=\eta_\nu$, and natural in $X$:
  $\eta_\nu\circ u^*=(u_*)^*\circ\eta_\nu$ for $\mathcal O_K$-morphisms $u$. For
  $v\in H^k(X_{\bar K},\mathbf Z_\ell)$ write $\eta_\nu(v)$ for $\eta_\nu$ of its reduction.
\item Let $B$ be a crystalline period package for $K$, with $c:A_{\mathrm{cris}}\to B$
  (Definition~\ref{gal-crystalline-relation-package}). Let $v_\alpha\in H^k(Y_{\bar K},\mathbf Z_\ell)$ and
  $c_\alpha\in A_{\mathrm{cris}}$ be finitely many classes and elements, and $u\geq0$. If
  $t^u\sum_\alpha c_\alpha\eta_\nu(v_\alpha)=0$ for every $\nu$, then $\sum_\alpha v_\alpha\otimes c(c_\alpha)=0$
  in $V\otimes_{\mathbf Q_\ell}B$.
\item Let $g:\mathcal O_K\to\mathcal O_C/\ell$ send $x$ to the Teichm\"uller representative of its
  residue, modulo $\ell$, and let $F:Y_g\to Y_g$ be the $q$-power Frobenius. For a multiple $m$ of $f$
  with $\varpi^{\ell^m}\in\ell\mathcal O_K$ one has $(x\bmod\ell)^{\ell^m}=g(x)$, so
  $\rho_m:Y^\sharp\to Y_g$ is defined. Let $B$ be a crystalline period package for $K$, with $c$,
  $\phi_B$ and $\sigma_B$. There are such an $m$ and $N_1\geq0$ with the following
  property. Let $v_\alpha\in H^k(Y_{\bar K},\mathbf Z_\ell)$, $c_\alpha\in A_{\mathrm{cris}}$ and $r\geq0$
  be such that $x=\sum_\alpha v_\alpha\otimes c(c_\alpha)c(t)^{-r}\in V\otimes_{\mathbf Q_\ell}B$ satisfies
  $(\sigma\otimes\sigma_B(\sigma))(x)=x$ for every $\sigma\in\operatorname{Gal}(\bar K/K)$. Then there are
  $a,u\geq0$ and classes $X_\nu\in H^k(Y_g,A_{\mathrm{cris}}/\ell^\nu)$ with $\phi^fX_\nu=F^*X_\nu$ and
  \[
   \ell^at^u\sum_\alpha\phi^{fN_1}(c_\alpha)\,\eta_\nu(v_\alpha)=t^{u+r}\rho_m^*X_\nu
  \]
  for every $\nu$.
\end{enumerate}
\end{theorem}
\leanref{Development/WeilBound/CoreStatements/PadicHodgeComparison/EtaleCohomologyComparison/KummerCupProducts/IsFirstChernClassOfLineBundle.lean}{CrystallineSemistableComparison.E.EtFam}
\leanref{Development/WeilBound/CoreStatements/PadicHodgeComparison/EtaleCohomologyComparison/KummerCupProducts/IsFirstChernClassOfLineBundle.lean}{CrystallineSemistableComparison.E.EtFam.Good}
\leanref{Development/WeilBound/CoreStatements/PadicHodgeComparison/EtaleCohomologyComparison/KummerCupProducts/FirstChernUnique.lean}{CrystallineSemistableComparison.E.etMap}
\leanref{Development/AutomorphicArithmetic/CommutativeAlgebra/EtaleCohomology/VecZero.lean}{UpperTheory.Ru.crystallineComparison.InjE}
\leanref{Development/AutomorphicArithmetic/CommutativeAlgebra/EtaleCohomology/VecZero.lean}{UpperTheory.Ru.crystallineComparison.AbsInv}
\leanref{Development/AutomorphicArithmetic/CommutativeAlgebra/EtaleCohomology/EqZeroOfCoord.lean}{UpperTheory.Ru.crystallineComparison.injOf}
\leanref{Development/AutomorphicArithmetic/CommutativeAlgebra/EtaleCohomology/PhiAIter.lean}{UpperTheory.Ru.crystallineComparison.absInvOf}
\leanref{Development/AutomorphicArithmetic/CommutativeAlgebra/EtaleCohomology/CIn.lean}{UpperTheory.Ru.crystallineComparison.cIn}

\begin{proof}
Let $\eta_\nu$ be the comparison maps of Definition~\ref{gal-realization-eta}.

(1) Each $\eta_\nu$ is additive, as a composite of homomorphisms. It is compatible with reduction by
Lemma~\ref{gal-realization-edge-map}(3) for $\ell^\nu\mid\ell^{\nu+1}$ and the compatibility of the
$\lambda_A$ with reduction (Lemma~\ref{gal-realization-locally-constant}(2)). The Frobenius $\phi$ of
$H^k(X^\sharp,A_{\mathrm{cris}}/\ell^\nu)$ is induced by $\phi$ on the coefficients, and
$\phi\circ\lambda_A=\lambda_A$, so $\phi\circ\eta_\nu=\eta_\nu$. Let $u:X\to X'$ be an
$\mathcal O_K$-morphism. The functor $m:\mathrm{El}(X^\sharp)\to\mathrm{El}(X'^\sharp)$, $(A,y)\mapsto(A,uy)$,
the identity maps $g_{(A,y)}$ of $\operatorname{Spec}A^\sharp[1/\ell]$ and $u_{\bar K}:X_{\bar K}\to
X'_{\bar K}$ satisfy the hypotheses of Lemma~\ref{gal-realization-edge-map}(2), which gives
$m^*\circ\varepsilon_{y'}=\varepsilon_y\circ u^*$ for the points $y'$ of $X'_{\bar K}$. The maps
$\lambda_A$ depend only on $A$, so $m^*$ commutes with $\lambda_*$, and on
$H^k(X'^\sharp,A_{\mathrm{cris}}/\ell^\nu)$ it is $(u_*)^*$. Hence $\eta_\nu\circ u^*=(u_*)^*\circ\eta_\nu$.

(2) Write $\eta(v)=(\eta_\nu(v))_\nu$ in $\varprojlim_\nu H^k(Y^\sharp,A_{\mathrm{cris}}/\ell^\nu)$, so that
$t^u\sum_\alpha c_\alpha\eta(v_\alpha)=0$. Let $\mu$ be a $\mathbf Q_\ell$-linear form on $V$. Take $Y_s$ and
$u_s:Y_s\to Y$, $1\le s\le r$, as in Lemma~\ref{gal-realization-duality}(3). A left inverse of the
injective map $v\mapsto(u_s^*v)_s$ writes $\mu=\sum_s\mu_s\circ u_s^*$ with linear forms $\mu_s$ on
$H^k(Y_{s,\bar K},\mathbf Q_\ell)$. By (1), $t^u\sum_\alpha c_\alpha\eta(u_s^*v_\alpha)=((u_s)_*)^*\bigl(t^u
\sum_\alpha c_\alpha\eta(v_\alpha)\bigr)=0$. By Lemma~\ref{gal-realization-inputs}, $Y_s$ has a period
datum and $\eta$ is $G_K$-equivariant, multiplicative and takes the value $1$, so
Lemma~\ref{gal-realization-injectivity} gives $\sum_\alpha\mu_s(u_s^*v_\alpha)c_\alpha=0$ in
$B_{\mathrm{dR}}$. Summing over $s$, $\sum_\alpha\mu(v_\alpha)c_\alpha=0$ in $B_{\mathrm{dR}}$. The injective map
$\iota:B\to B_{\mathrm{dR}}$ of Lemma~\ref{gal-crystalline-relation-package-embedding} fixes $\mathbf Q_\ell$ and sends
$c(c_\alpha)$ to $c_\alpha$, so $\sum_\alpha\mu(v_\alpha)c(c_\alpha)=0$ in $B$. Taking for $\mu$ the coordinate forms
of a basis of $V$ gives $\sum_\alpha v_\alpha\otimes c(c_\alpha)=0$ in $V\otimes_{\mathbf Q_\ell}B$.

(3) Take $N$ as in Lemma~\ref{gal-realization-inputs}(3), and let $(H,\phi_H,b_\nu,L,e,\alpha,\beta)$ be a
Frobenius datum of level $N$ for $Y$ in degree $k$, with $u$ as in
Definition~\ref{gal-realization-frobenius-datum}(1). Put $m=N$ and $N_1=N/f$. Then
$(x\bmod\ell)^{\ell^m}=g(x)$ for all $x\in\mathcal O_K$, so $\rho_m$ is defined, and
$\varpi^{\ell^m}\in\ell\mathcal O_C\cap\mathcal O_K=\ell\mathcal O_K$ because $g(\varpi)=0$. Let
$v_\alpha$, $c_\alpha$ and $r$ be such that $x=\sum_\alpha v_\alpha\otimes c(c_\alpha)c(t)^{-r}$ is fixed by every
$\sigma\otimes\sigma_B(\sigma)$, and put $E_\nu=\sum_\alpha c_\alpha\eta_\nu(v_\alpha)$. For $\sigma\in G_K$ let
$\chi(\sigma)\in\mathbf Z_\ell^\times$ be the cyclotomic character, so that $\sigma(t)=\chi(\sigma)t$; recall
$\phi(t)=\ell t$ and that $\phi$ fixes $\mathbf Z_\ell$.

\emph{Step 1.} By (1) the $E_\nu$ are compatible under reduction, so by
Definition~\ref{gal-realization-frobenius-datum}(1) there is $h\in H$ with $b_\nu(h)=t^uE_\nu$ for all
$\nu$.

\emph{Step 2.} Let $\sigma\in G_K$. There is $b\geq0$ with $t^b\sigma(E_\nu)=t^b\chi(\sigma)^rE_\nu$ for
all $\nu$. Indeed, as $\sigma_B(\sigma)c(t)=\chi(\sigma)c(t)$, the identity $(\sigma\otimes\sigma_B(\sigma))x=x$ says that
$\sum_\alpha\sigma v_\alpha\otimes c(\sigma(c_\alpha))c(t)^{-r}-\sum_\alpha v_\alpha\otimes\chi(\sigma)^rc(c_\alpha)c(t)^{-r}$
vanishes in $V\otimes_{\mathbf Q_\ell}B=T\otimes_{\mathbf Z_\ell}B$, where $T=H^k(Y_{\bar K},\mathbf Z_\ell)$. As $c$
makes $B$ the localization of $A_{\mathrm{cris}}$ at the powers of $t$ (proof of
Lemma~\ref{gal-crystalline-relation-package-embedding}), $T\otimes_{\mathbf Z_\ell}B=(T\otimes_{\mathbf
Z_\ell}A_{\mathrm{cris}})[1/t]$. So
$t^b\bigl(\sum_\alpha\sigma v_\alpha\otimes\sigma(c_\alpha)-\sum_\alpha v_\alpha\otimes\chi(\sigma)^rc_\alpha
\bigr)=0$ in $T\otimes_{\mathbf Z_\ell}A_{\mathrm{cris}}$ for some $b$. Apply the map $v\otimes c\mapsto
c\,\eta_\nu(v)$, which is well defined because $\eta_\nu(zv)=z\eta_\nu(v)$ for $z\in\mathbf Z_\ell$. Since
$\eta_\nu(\sigma v)=\sigma(\eta_\nu(v))$ by Lemma~\ref{gal-realization-inputs}(1) and $\sigma$ acts
semilinearly, the first sum goes to $\sigma(E_\nu)$, which gives the claim.

\emph{Step 3.} Let $\sigma$ and $b$ be as in Step~2, and put $\xi=t^b\sigma(h)-t^b\chi(\sigma)^{u+r}h$. Then
$b_\nu(\sigma h)=\sigma(t^uE_\nu)=\chi(\sigma)^ut^u\sigma(E_\nu)$, so
$b_\nu(\xi)=\chi(\sigma)^ut^u\bigl(t^b\sigma(E_\nu)-t^b\chi(\sigma)^rE_\nu\bigr)=0$ for all $\nu$, and
$t^u\xi=0$ by Definition~\ref{gal-realization-frobenius-datum}(1). Applying $\alpha$, which is
$\phi^N$-semilinear and carries $\sigma$ to $\sigma\otimes1$, and using $\phi^N(t)=\ell^Nt$ and
$\phi^N(\chi(\sigma))=\chi(\sigma)$,
\[
 \ell^{N(u+b)}t^{u+b}(\sigma\otimes1)(\alpha(h))=\ell^{N(u+b)}t^{u+b}\chi(\sigma)^{u+r}\alpha(h).
\]

\emph{Step 4.} Apply Lemma~\ref{gal-realization-descent} to $W=W(\mathbf F_q)$, in which every nonzero
element is a power of $\ell$ times a unit, $A=A_{\mathrm{cris}}$, the package $B$ with $c$ and the action
$\sigma_B$ (Definition~\ref{gal-crystalline-relation-package}), $\Gamma=G_K$, $z=\alpha(h)$ and $s=u+r$; $G_K$ fixes
the image of $W$ by Lemma~\ref{gal-realization-invariants}, the hypothesis on invariants is
Lemma~\ref{gal-crystalline-relation-package-invariants}, and the hypothesis on $z$ is Step~3. It gives
$a',b'\geq0$ and $y\in L$ with
$\ell^{a'}t^{b'}\alpha(h)=t^{b'+u+r}(1\otimes y)$.

\emph{Step 5.} Apply $b_\nu\circ\beta$ to this identity. On the left, $\beta(\alpha(h))=\ell^e\phi_H^N(h)$
and $b_\nu(\phi_H^Nh)=\phi^N(t^uE_\nu)=\ell^{Nu}t^u\phi^N(E_\nu)$. As $\phi$ is semilinear and
$\phi\circ\eta_\nu=\eta_\nu$ by (1), $\phi^N(E_\nu)=\sum_\alpha\phi^N(c_\alpha)\eta_\nu(v_\alpha)$. Hence
\[
 \ell^{a'+e+Nu}t^{b'+u}\sum_\alpha\phi^{fN_1}(c_\alpha)\,\eta_\nu(v_\alpha)=t^{b'+u+r}\,b_\nu(\beta(1\otimes y)).
\]
By Definition~\ref{gal-realization-frobenius-datum}(2) there are $X_\nu\in H^k(Y_g,A_{\mathrm{cris}}/\ell^\nu)$
with $\phi^fX_\nu=F^*X_\nu$ and $b_\nu(\beta(1\otimes y))=\rho_m^*X_\nu$. This is the claim with
$a=a'+e+Nu$ and $b'+u$ in place of $u$.
\end{proof}
\leanref{Development/AutomorphicSpectrum/Proofs/BlockDecomposition/PowDivPowSucc.lean}{GlobalTheory.CompanionTheory.EtReduction.fred_of_etRed}
\leanref{Development/AutomorphicSpectrum/Proofs/BlockDecomposition/FrobIdem.lean}{GlobalTheory.CompanionTheory.FrobLc.fPhi}
\leanref{Development/AutomorphicSpectrum/Proofs/BlockDecomposition/GeomPtGeomMap.lean}{GlobalTheory.CompanionTheory.FNatProof.nat_step}
\leanref{Development/AutomorphicArithmetic/CommutativeAlgebra/EtaleCohomology/VInduction.lean}{UpperTheory.Ru.crystallineComparison.injAbs}
\leanref{Development/AutomorphicArithmetic/CommutativeAlgebra/EtaleCohomology/EqZeroOfCoord.lean}{UpperTheory.Ru.crystallineComparison.eq_zero_of_coord}
\leanref{Development/AutomorphicArithmetic/CommutativeAlgebra/EtaleCohomology/GalAwayCongr.lean}{UpperTheory.Ru.crystallineComparison.pbIota}
\leanref{Development/AutomorphicArithmetic/EtaleIdempotents/EvalEtaleFamilyMemLimitOfReduced.lean}{UpperTheory.evalEtaleFamily_mem_limit_of_reduced}
\leanref{Development/AutomorphicArithmetic/CommutativeAlgebra/EtaleCohomology/TorsionOfVec.lean}{UpperTheory.Ru.crystallineComparison.vecZero}
\leanref{Development/AutomorphicArithmetic/EtaleIdempotents/GaloisActionEvalEtaleFamilyExistsDenominator.lean}{UpperTheory.galoisAction_evalEtaleFamily_exists_denominator}
\leanref{Development/AutomorphicArithmetic/CommutativeAlgebra/EtaleCohomology/PhiAIter.lean}{UpperTheory.Ru.crystallineComparison.chain1}
\leanref{Development/AutomorphicArithmetic/CommutativeAlgebra/EtaleCohomology/PhiAIter.lean}{UpperTheory.Ru.crystallineComparison.chain2}
\leanref{Development/AutomorphicArithmetic/EtaleIdempotents/CrisFrobeniusIterateEvalEtaleFamily.lean}{UpperTheory.crisFrobenius_iterate_evalEtaleFamily}

\begin{lemma}\label{gal-crystalline-splitting-cover}
Let $\ell$ be a prime, $Y$ a scheme and $T_j=(W_j,p_j,r_j)$, $j\in J$, finitely many pairs of morphisms
$p_j,r_j:W_j\to Y$ with $p_j$ finite and \'etale over $Y[1/\ell]$ and with $d_j$ geometric points over
every geometric point of $Y$ at which $\ell$ is invertible. There are a scheme $Y'$, a morphism
$\pi:Y'\to Y$ and morphisms $s_{j,i}:Y'\to W_j$, $1\leq i\leq d_j$, with $p_js_{j,i}=\pi$, such that:
\begin{enumerate}
\item for every $j$ the $s_{j,i}$ induce an isomorphism $\coprod_iY'[1/\ell]\to Y'[1/\ell]\times_YW_j$;
\item for every ring $U$ in which $\ell$ is a nonzerodivisor, every $y\in Y(U)$ and all $z_{j,i}\in
  W_j(U)$ with $p_jz_{j,i}=y$ that induce isomorphisms
  $\coprod_i\operatorname{Spec}U[1/\ell]\to\operatorname{Spec}U[1/\ell]\times_YW_j$, there is a unique
  $y'\in Y'(U)$ with $\pi y'=y$ and $s_{j,i}y'=z_{j,i}$;
\item let $Y$ be an $\mathcal O$-scheme for a ring $\mathcal O$, let $\kappa:\mathcal O\to\mathcal O_C$ be
  a ring map, give $Y'$ the structure map through $\pi$, and let $G=\prod_jS_{d_j}$ act on $Y'^\sharp$
  (Definition~\ref{gal-crystalline-test-algebras}) by permuting the $s_{j,i}$ through (2). Then
  $\pi_*:Y'^\sharp\to Y^\sharp$ and the action of $G$ satisfy the hypotheses of
  Lemma~\ref{gal-crystalline-test-cohomology}(3).
\end{enumerate}
\end{lemma}
\leanref{Development/AutomorphicArithmetic/CommutativeAlgebra/EtaleCohomology/SplitRules.lean}{UpperTheory.Ru.crystallineComparison.WDat}
\leanref{Development/AutomorphicArithmetic/CommutativeAlgebra/EtaleCohomology/SplitRules.lean}{UpperTheory.Ru.crystallineComparison.WEx}
\leanref{Development/AutomorphicArithmetic/CommutativeAlgebra/EtaleCohomology/SplitOnOfDisj.lean}{UpperTheory.Ru.crystallineComparison.wEx}
\leanref{Development/AutomorphicArithmetic/CommutativeAlgebra/EtaleCohomology/TestFacts.lean}{UpperTheory.Ru.crystallineComparison.WTrans}
\leanref{Development/AutomorphicArithmetic/CommutativeAlgebra/EtaleCohomology/SplitPointGeneric.lean}{UpperTheory.Ru.crystallineComparison.wTrans}

\begin{proof}
Let $P$ be the fibre product over $Y$ of $d_j$ copies of $W_j$ for every $j$, with projections $c_{j,i}$.
For $i\neq i'$, the locus of $P[1/\ell]$ where $c_{j,i}=c_{j,i'}$ is open and closed, since $p_j$ is
\'etale and separated over $Y[1/\ell]$. Let $D\subset P[1/\ell]$ be the open and closed complement of
the union of these loci, let $Y'$ be the closure of $D$ in $P$, and let $\pi$ and the $s_{j,i}$ be the
restrictions of the structure map and of the $c_{j,i}$. The points of $Y'$ at which $\ell$ is
invertible lie in $D$, where the $s_{j,i}$ are $d_j$ pairwise disjoint sections of the finite \'etale map
$p_j$ of degree $d_j$; this gives (1). In (2), the point $(y,(z_{j,i}))$ of $P$ maps
$\operatorname{Spec}U[1/\ell]$ into $D$, and $\operatorname{Spec}U[1/\ell]$ is schematically dense in
$\operatorname{Spec}U$, so the point factors uniquely through the closure $Y'$.

For (3), let $A$ be a test algebra and $y\in Y^\sharp(A)$. By Lemma~\ref{gal-crystalline-test-facts}(2)
the finite \'etale covers $W_j\times_Y\operatorname{Spec}A^\sharp[1/\ell]$ split, and their sections
extend to $A^\sharp$-points of $W_j$, since $p_j$ is finite and $A^\sharp$ is integrally closed in
$A^\sharp[1/\ell]$ (Lemma~\ref{gal-crystalline-test-facts}(1)); by (2) they come from a point of
$Y'^\sharp(A)$ over $y$. Two points of $Y'^\sharp(A)$ over $y$ give two splittings of the same covers,
which agree up to a permutation on the pieces of a decomposition of $A^\sharp[1/\ell]$ by idempotents.
These idempotents lie in $A^\sharp$, which is integrally closed in $A^\sharp[1/\ell]$, and come from a
decomposition of $A$ by Lemma~\ref{gal-crystalline-test-facts}(3) and~(4); on each piece the two points
differ by an element of $G$, by the uniqueness in (2), so their $G$-orbits agree as multisets.
\end{proof}

\begin{lemma}\label{gal-crystalline-norm-forms}
For a ring $B$ and $c=(c_1,\ldots,c_m)\in B^m$ write $\Phi_c=X-\sum_kc_kT_k\in B[T_1,\ldots,T_m][X]$.
Let $\ell$ be a prime and $R$ a ring in which $\ell$ is a nonzerodivisor and which is integrally closed
in $R[1/\ell]$. Let $S$ be a finite $R$-algebra such that $S[1/\ell]$ is \'etale over $R[1/\ell]$ and
has $d$ homomorphisms of $R$-algebras to every algebraically closed field in which $\ell\neq0$, and let
$s_1,\ldots,s_m\in S$. There is a polynomial $N\in R[T_1,\ldots,T_m][X]$ such that:
\begin{enumerate}
\item for every $R$-algebra $B$ in which $\ell$ is a nonzerodivisor and all $R$-algebra maps
  $z_1,\ldots,z_d:S\to B$ for which $S\otimes_RB[1/\ell]\to B[1/\ell]^d$, $s\otimes b\mapsto(z_i(s)b)_i$,
  is bijective, the image of $N$ in $B[T][X]$ is $\prod_i\Phi_{(z_i(s_k))_k}$;
\item for every valuation ring $V$ with algebraically closed fraction field that is an $R$-algebra
  with $\ell\neq0$ in $V$, there are finitely many $R$-algebra maps $S\to V$, and the image of $N$ in
  $V[T][X]$ is the product of $\Phi_{(z(s_k))_k}$ over all of them.
\end{enumerate}
\end{lemma}
\leanref{Development/AutomorphicArithmetic/Foundations/GaloisRepresentations/AnnihilatorModels/GenericMulOfProd.lean}{UpperTheory.Gal.crystallineComparison.form}
\leanref{Development/AutomorphicArithmetic/Foundations/GaloisRepresentations/AnnihilatorModels/GenericMulOfProd.lean}{UpperTheory.Gal.crystallineComparison.SameNorm}
\leanref{Development/AutomorphicArithmetic/Foundations/GaloisRepresentations/AnnihilatorModels/SameNormBasic.lean}{UpperTheory.Gal.crystallineComparison.NormRead}
\leanref{Development/AutomorphicArithmetic/Foundations/GaloisRepresentations/AnnihilatorModels/PNormRead.lean}{UpperTheory.Gal.crystallineComparison.normRead_of}

\begin{proof}
Since $S[1/\ell]$ is finite \'etale over $R[1/\ell]$, it is finite projective; let $N$ be the
characteristic polynomial of multiplication by $\sum_ks_kT_k$ on $S[1/\ell][T]$ over $R[1/\ell][T]$.
After a base change over which $S$ splits, multiplication by $\sum_ks_kT_k$ is diagonal with entries
$\sum_kz_i(s_k)T_k$, so $N$ maps to $\prod_i\Phi_{(z_i(s_k))_k}$; this holds in $B[1/\ell]$, and $B$
embeds in $B[1/\ell]$. The coefficients of $N$ are integral over $R$, as one sees after a faithfully
flat base change that splits $S[1/\ell]$, where they are polynomials in the images of the $s_k$, which
are integral over $R$. They lie in $R[1/\ell]$, hence in $R$. This gives (1). In (2), the $R$-algebra
maps $S\to V$ are the maps from $S$ to the fraction field of $V$, since $V$ is integrally closed, and
there are $d$ of them; they split $S$ over $V[1/\ell]$, and (2) follows as (1).
\end{proof}

\begin{lemma}\label{gal-crystalline-norm-identity}
Let $\mathcal O$ be a discrete valuation ring with uniformizer $\varpi$, $\ell$ a prime that is neither
zero nor a unit in $\mathcal O$, $q=\ell^f$ with $f\geq1$, and $Y$ a smooth $\mathcal O$-scheme. Let
$T_j=(W_j,p_j,r_j)$, $0\leq j\leq n$, be finite correspondences of $Y$
(Definition~\ref{gal-tower-reduction}(2), with $\ell$ for $p$), with $d_j$ geometric points in the
fibres of $p_j$ over $Y[1/\ell]$, that satisfy the Eichler--Shimura relation at the generic points of
$Y\otimes\mathcal O/\varpi$ of Definition~\ref{gal-tower-reduction}(3) with this $q$. Let $U$ and $U'$
be affine opens of $Y$ with $p_j^{-1}(U)\subset r_j^{-1}(U')$ for every $j$, and
$x_1,\ldots,x_m\in\Gamma(U',\mathcal O_Y)$. Let $B$ be a $\Gamma(U,\mathcal O_Y)$-algebra in which
$\ell$ is a nonzerodivisor, let $z_{j,i}:\Gamma(p_j^{-1}U,\mathcal O_{W_j})\to B$, $1\leq i\leq d_j$, be
maps of $\Gamma(U,\mathcal O_Y)$-algebras that split $\Gamma(p_j^{-1}U,\mathcal O_{W_j})$ over
$B[1/\ell]$ as in Lemma~\ref{gal-crystalline-norm-forms}(1), and let $\beta:B\to B'$ be a ring map with
$\beta(\varpi)=0$. Put $c_{j,i}=\bigl(\beta(z_{j,i}(r_j^*x_k))^{q^{n-j}}\bigr)_k\in B'^m$. Then, with the
notation of Lemma~\ref{gal-crystalline-norm-forms},
\[
 \prod_{j\text{ even}}\prod_i\Phi_{c_{j,i}}^{\,q^{j(j-1)/2}}=\prod_{j\text{ odd}}\prod_i\Phi_{c_{j,i}}^{\,q^{j(j-1)/2}}
 \quad\text{in }B'[T_1,\ldots,T_m][X].
\]
\end{lemma}
\leanref{Development/AutomorphicArithmetic/Foundations/GaloisRepresentations/AnnihilatorModels/GenericMulOfProd.lean}{UpperTheory.Gal.crystallineComparison.NormOn}
\leanref{Development/AutomorphicArithmetic/Foundations/GaloisRepresentations/AnnihilatorModels/GenericMulOfProd.lean}{UpperTheory.Gal.crystallineComparison.NormOfPt}
\leanref{Development/AutomorphicArithmetic/Foundations/GaloisRepresentations/AnnihilatorModels/PNormPt.lean}{UpperTheory.Gal.crystallineComparison.normPt_of}

\begin{proof}
Let $R=\Gamma(U,\mathcal O_Y)$, $S_j=\Gamma(p_j^{-1}U,\mathcal O_{W_j})$ and
$s_{j,k}=(r_j^*x_k)^{q^{n-j}}\in S_j$. Since $Y$ is smooth over $\mathcal O$, $R$ is flat over
$\mathcal O$ and normal, so $\ell$ is a nonzerodivisor in $R$ and $R$ is integrally closed in
$R[1/\ell]$; the special fibre is smooth over $\mathcal O/\varpi$, so the ideal $\varpi R$ is radical.
Each $S_j$ satisfies the hypotheses of Lemma~\ref{gal-crystalline-norm-forms} with $d=d_j$; let
$N_j$ be its polynomial for $s_{j,1},\ldots,s_{j,m}$, and put
$P_{\mathrm{ev}}=\prod_{j\text{ even}}N_j^{q^{j(j-1)/2}}$ and
$P_{\mathrm{od}}=\prod_{j\text{ odd}}N_j^{q^{j(j-1)/2}}$. By Lemma~\ref{gal-crystalline-norm-forms}(1),
applied to $B$ and followed by $\beta$, the two sides of the claim are the images of $P_{\mathrm{ev}}$
and $P_{\mathrm{od}}$. Since $\varpi R$ is radical, it suffices to show that $P_{\mathrm{ev}}$ and
$P_{\mathrm{od}}$ agree modulo every minimal prime $\mathfrak p$ over $\varpi R$.

Let $\mathfrak p$ be such a prime. There is a valuation ring $V$ with algebraically closed fraction
field and a map $R\to V$ with $\ell\neq0$ in $V$ whose maximal ideal pulls back to $\mathfrak p$; then
$R/\mathfrak p$ embeds in the residue field $\kappa_V$. The corresponding point
$y:\operatorname{Spec}V\to U\subset Y$ maps the closed point to the generic point $\mathfrak p$ of
the special fibre. The $V$-points $z$ of $W_j$ over $y$ are the $R$-algebra maps $S_j\to V$, and $r_jz$
lies in $U'$. The $\kappa_V$-point $\phi^{n-j}(\bar r_j(z))$ of Definition~\ref{gal-tower-reduction}(3)
lies in $U'$, and its coordinates $x_k$ are the residues of the $z(s_{j,k})$. The Eichler--Shimura
relation at $y$ says that the multisets of these points, with the weights $q^{j(j-1)/2}$ for even and
for odd $j$, agree, so the multisets of their coordinate vectors agree. By
Lemma~\ref{gal-crystalline-norm-forms}(2), $P_{\mathrm{ev}}$ and $P_{\mathrm{od}}$ then have the same
image in $\kappa_V[T][X]$, hence modulo $\mathfrak p$.
\end{proof}
\leanref{Development/AutomorphicArithmetic/Foundations/GaloisRepresentations/AnnihilatorModels/GenericMulOfProd.lean}{UpperTheory.Gal.crystallineComparison.SameNormRing}
\leanref{Development/AutomorphicArithmetic/Foundations/GaloisRepresentations/AnnihilatorModels/PSameNormRing.lean}{UpperTheory.Gal.crystallineComparison.sameNormRing_of}
\leanref{Development/AutomorphicArithmetic/Foundations/GaloisRepresentations/AnnihilatorModels/SameNormBasic.lean}{UpperTheory.Gal.crystallineComparison.CongrRad}
\leanref{Development/AutomorphicArithmetic/Foundations/GaloisRepresentations/AnnihilatorModels/SameNormBasic.lean}{UpperTheory.Gal.crystallineComparison.ValLift}
\leanref{Development/AutomorphicArithmetic/Foundations/GaloisRepresentations/AnnihilatorModels/SameNormBasic.lean}{UpperTheory.Gal.crystallineComparison.BaseOfChart}

\begin{definition}\label{gal-crystalline-matching-products}
Let $B$ be a ring, $m\geq0$, and $\Phi_c=X-\sum_kc_kT_k\in B[T_1,\ldots,T_m][X]$ for $c\in B^m$, as in
Lemma~\ref{gal-crystalline-norm-forms}. Let $u$ and $u'$ be finite multisets of elements of $B^m$ and
$a\in B$.
\begin{enumerate}
\item $\Phi_u=\prod_{c\in u}\Phi_c$. For a ring map $\rho:B\to B'$, $\rho(u)$ is the multiset of the
  vectors $(\rho(c_k))_k$ for the members $c$ of $u$.
\item $\Phi_u\equiv\Phi_{u'}\bmod a$ means that the coefficients of $\Phi_u-\Phi_{u'}$ lie in $aB$, that
  is, $\Phi_{\rho(u)}=\Phi_{\rho(u')}$ for the quotient map $\rho:B\to B/aB$.
\item $u\equiv u'\bmod a$ means that $\rho(u)=\rho(u')$ for the quotient map $\rho:B\to B/aB$.
\item For $x,y\in B^m$, $a\mid x-y$ means that $a$ divides $x_k-y_k$ for every $k$.
\end{enumerate}
\end{definition}
\leanref{Development/AutomorphicArithmetic/Foundations/GaloisRepresentations/AnnihilatorModels/GenericMulOfProd.lean}{UpperTheory.Gal.crystallineComparison.redMod}

\begin{lemma}\label{gal-crystalline-matching-forms}
Let $B$ be a ring, $m\geq0$, and $u,u'$ finite multisets of elements of $B^m$, with the notation of
Definition~\ref{gal-crystalline-matching-products}.
\begin{enumerate}
\item If $\Phi_u=\Phi_{u'}$, then $\Phi_{\rho(u)}=\Phi_{\rho(u')}$ for every ring map $\rho:B\to B'$.
\item If $x\in B^m$ and $\Phi_x\Phi_u=\Phi_x\Phi_{u'}$, then $\Phi_u=\Phi_{u'}$.
\item If $B\neq0$ and $\Phi_u=\Phi_{u'}$, then $u$ and $u'$ have the same number of members.
\item If $\Phi_u=\Phi_{u'}$ and $x\in u$, then $\prod_{y\in u'}\sum_k(x_k-y_k)T_k=0$ in
  $B[T_1,\ldots,T_m]$.
\end{enumerate}
\end{lemma}
\leanref{Development/AutomorphicArithmetic/Foundations/GaloisRepresentations/AnnihilatorModels/SameNormBasic.lean}{UpperTheory.Gal.crystallineComparison.SameNormBasic}
\leanref{Development/AutomorphicArithmetic/Foundations/GaloisRepresentations/AnnihilatorModels/PSnBasic.lean}{UpperTheory.Gal.crystallineComparison.snBasic_of}

\begin{proof}
Applying $\rho$ to the coefficients maps $\Phi_c$ to $\Phi_{(\rho(c_k))_k}$, which gives (1). The
polynomial $\Phi_x$ is monic in $X$, hence a nonzerodivisor, which gives (2). If $B\neq0$, a product of
$n$ polynomials that are monic of degree one in $X$ is monic of degree $n$, which gives (3). For (4),
substitute $X=\sum_kx_kT_k$. Then $\Phi_y$ becomes $\sum_k(x_k-y_k)T_k$, and $\Phi_u$ becomes $0$, since
$\Phi_x$ is one of its factors.
\end{proof}

\begin{lemma}\label{gal-crystalline-matching-gauss}
Let $A$ be a ring, $a\in A$, and $p_1,\ldots,p_n\in A[T_1,\ldots,T_m]$ polynomials each of which has a
coefficient equal to $1$. If $a\,p_1\cdots p_n=0$, then $a=0$.
\end{lemma}
\leanref{Development/AutomorphicArithmetic/Foundations/GaloisRepresentations/AnnihilatorModels/SameNormBasic.lean}{UpperTheory.Gal.crystallineComparison.GaussLin}
\leanref{Development/AutomorphicArithmetic/Foundations/GaloisRepresentations/AnnihilatorModels/PGaussLin.lean}{UpperTheory.Gal.crystallineComparison.gaussLin_of}

\begin{proof}
Suppose $a\neq0$. Then the annihilator of $a$ is a proper ideal, so it lies in a maximal ideal
$\mathfrak m$. Every coefficient of $p_1\cdots p_n$ is annihilated by $a$, so it lies in $\mathfrak m$,
and the image of $p_1\cdots p_n$ in $(A/\mathfrak m)[T_1,\ldots,T_m]$ is zero. This ring is a domain, so
the image of some $p_i$ is zero. That is impossible, since $p_i$ has a coefficient equal to $1$.
\end{proof}

\begin{lemma}\label{gal-crystalline-matching-factor}
Let $B$ be a ring, $a\in B$ and $d\in B^m$. Suppose that some coordinate $d_k$ divides $a$, and that for
all $k,k'$ with $d_k\mid a$ one of $d_k$ and $d_{k'}$ divides the other. Then there are $b\in B$ with
$b\mid a$ and $v\in B^m$ with $v_{k_0}=1$ for some $k_0$, such that $d=bv$.
\end{lemma}
\leanref{Development/AutomorphicArithmetic/Foundations/GaloisRepresentations/AnnihilatorModels/SameNormBasic.lean}{UpperTheory.Gal.crystallineComparison.LinFactor}
\leanref{Development/AutomorphicArithmetic/Foundations/GaloisRepresentations/AnnihilatorModels/PLinFactor.lean}{UpperTheory.Gal.crystallineComparison.linFactor_of}

\begin{proof}
The coordinates $d_k$ that divide $a$ are pairwise comparable for divisibility. By induction on their
number, one of them, $d_{k_0}$, divides all the others. Then $d_{k_0}$ divides every $d_k$: either
$d_{k_0}\mid d_k$, or $d_k\mid d_{k_0}$, and in the second case $d_k$ divides $a$, so $d_{k_0}\mid d_k$
by the choice of $k_0$. Put $b=d_{k_0}$ and $v_{k_0}=1$, and for $k\neq k_0$ choose $v_k$ with
$d_k=bv_k$.
\end{proof}

\begin{lemma}\label{gal-crystalline-matching-step}
Let $B$ be a ring, $w\in B$ a nonzerodivisor that is not a unit, $N\geq0$ and $r\geq1$. Let $u$ and $u'$
be finite multisets of elements of $B^m$, $u'$ with at most $N$ members, such that
$\Phi_u\equiv\Phi_{u'}\bmod w^{r(N+1)}$ (Definition~\ref{gal-crystalline-matching-products}). Let $x\in
u$ be such that for every $y\in u'$:
\begin{enumerate}
\item[(a)] for every $k$, $x_k-y_k$ divides $w^r$ or $w^r$ divides $x_k-y_k$;
\item[(b)] for all $k,k'$ with $x_k-y_k\mid w^r$, one of $x_k-y_k$ and $x_{k'}-y_{k'}$ divides the
  other.
\end{enumerate}
Then there is $y\in u'$ with $w^r\mid x-y$.
\end{lemma}
\leanref{Development/AutomorphicArithmetic/Foundations/GaloisRepresentations/AnnihilatorModels/SameNormBasic.lean}{UpperTheory.Gal.crystallineComparison.MatchStep}
\leanref{Development/AutomorphicArithmetic/Foundations/GaloisRepresentations/AnnihilatorModels/PMatchStep.lean}{UpperTheory.Gal.crystallineComparison.matchStep_of}

\begin{proof}
Suppose that there is no such $y$, and let $y\in u'$. Some coordinate $x_k-y_k$ is not divisible by
$w^r$, so it divides $w^r$ by (a). By (b) and Lemma~\ref{gal-crystalline-matching-factor} with $a=w^r$,
$x-y=b_yv_y$ with $b_y\mid w^r$ and $v_y\in B^m$ a vector with a coordinate equal to $1$. Let
$\bar B=B/w^{r(N+1)}B$, let bars denote images in $\bar B$, and let $b=\prod_{y\in u'}b_y$. By
Lemma~\ref{gal-crystalline-matching-forms}(4) over $\bar B$,
\[
 0=\prod_{y\in u'}\sum_k(\bar x_k-\bar y_k)T_k=\bar b\prod_{y\in u'}\sum_k\bar v_{y,k}T_k
 \quad\text{in }\bar B[T_1,\ldots,T_m].
\]
Each factor $\sum_k\bar v_{y,k}T_k$ has a coefficient equal to $1$, so $\bar b=0$ by
Lemma~\ref{gal-crystalline-matching-gauss}; that is, $w^{r(N+1)}$ divides $b$. On the other hand $b$
divides $w^{rs}$, where $s\leq N$ is the number of members of $u'$. Hence
$w^{rs}=w^{r(N+1)}t=w^{rs}\cdot w^{r(N+1-s)}t$ for some $t\in B$. Since $w^{rs}$ is a nonzerodivisor,
$w^{r(N+1-s)}t=1$. As $r(N+1-s)\geq1$, $w$ is a unit, a contradiction.
\end{proof}

\begin{lemma}\label{gal-crystalline-matching-induction}
Let $B$ be a ring, $w\in B$ a nonzerodivisor that is not a unit, and $N\geq0$. Let $u$ and $u'$ be
finite multisets of elements of $B^m$ with at most $N$ members each, such that
$\Phi_u\equiv\Phi_{u'}\bmod w^{(N+1)^N}$ (Definition~\ref{gal-crystalline-matching-products}). Suppose
that for all $x\in u$ and $y\in u'$:
\begin{enumerate}
\item[(a)] for every $k$ and every $0\leq j\leq N$, $x_k-y_k$ divides $w^{(N+1)^j}$ or $w^{(N+1)^j}$
  divides $x_k-y_k$;
\item[(b)] for all $k,k'$ with $x_k-y_k\mid w^{(N+1)^N}$, one of $x_k-y_k$ and $x_{k'}-y_{k'}$ divides
  the other.
\end{enumerate}
Then $u\equiv u'\bmod w$.
\end{lemma}
\leanref{Development/AutomorphicArithmetic/Foundations/GaloisRepresentations/AnnihilatorModels/GenericMulOfProd.lean}{UpperTheory.Gal.crystallineComparison.Matching}
\leanref{Development/AutomorphicArithmetic/Foundations/GaloisRepresentations/AnnihilatorModels/PMatching.lean}{UpperTheory.Gal.crystallineComparison.matching_of}

\begin{proof}
A congruence $\Phi_u\equiv\Phi_{u'}\bmod a$ implies the same congruence modulo every divisor of $a$, by
Lemma~\ref{gal-crystalline-matching-forms}(1). We show by induction on $c\leq N$: if $u$ has $c$
members, $u'$ has at most $N$ members, (a) and (b) hold, and $\Phi_u\equiv\Phi_{u'}\bmod w^{(N+1)^c}$,
then $u\equiv u'\bmod w$. Since $w^{(N+1)^c}$ divides $w^{(N+1)^N}$, this gives the lemma.

If $c=0$, then $u$ is empty and $\Phi_{u'}\equiv1\bmod w$. Since $w$ is not a unit, $B/wB\neq0$, so
$u'$ is empty by Lemma~\ref{gal-crystalline-matching-forms}(3) over $B/wB$.

Let $u$ have $c+1\leq N$ members. Choose $x\in u$ and write $u=\{x\}+u_0$. Put $r=(N+1)^c$, so that
$r(N+1)=(N+1)^{c+1}$. Hypothesis (a) of Lemma~\ref{gal-crystalline-matching-step} for $x$ is (a) with
$j=c$. Its hypothesis (b) follows from (b), since $c<N$ and so $w^r$ divides $w^{(N+1)^N}$. That lemma
gives $y\in u'$ with $w^r\mid x-y$; write $u'=\{y\}+u'_0$. Modulo $w^r$ we have $\Phi_x\equiv\Phi_y$ and
$\Phi_u\equiv\Phi_{u'}$, so $\Phi_{u_0}\equiv\Phi_{u'_0}\bmod w^r$ by
Lemma~\ref{gal-crystalline-matching-forms}(2) over $B/w^rB$. The multisets $u_0$, with $c$ members, and
$u'_0$ satisfy (a) and (b), so $u_0\equiv u'_0\bmod w$ by induction. Since $w$ divides $w^r$, $x\equiv
y\bmod w$, and so $u\equiv u'\bmod w$.
\end{proof}

\begin{lemma}\label{gal-crystalline-matching-piece}
Let $B$ be a ring, $w\in B$ a nonzerodivisor, $N\geq0$, and $c_1,\ldots,c_N,c'_1,\ldots,c'_N\in B^m$
with $\prod_i\Phi_{c_i}\equiv\prod_i\Phi_{c'_i}\bmod w^{(N+1)^N}$
(Definition~\ref{gal-crystalline-matching-products}). Let $e\in B$ be an idempotent such that for all
$i,i'$ and $k$, with $d=c_{i,k}-c'_{i',k}$:
\begin{enumerate}
\item[(a)] for every $0\leq n\leq(N+1)^N$, $d\mid ew^n$ or $w^n\mid ed$;
\item[(b)] if $d\mid ew^{(N+1)^N}$, then for every $k'$, with $d'=c_{i,k'}-c'_{i',k'}$, $d\mid ed'$ or
  $d'\mid ed$.
\end{enumerate}
Then there is a permutation $\sigma\in S_N$ with $w\mid e(c_i-c'_{\sigma(i)})$ for every $i$.
\end{lemma}
\leanref{Development/AutomorphicArithmetic/Foundations/GaloisRepresentations/AnnihilatorModels/SameNormBasic.lean}{UpperTheory.Gal.crystallineComparison.MatchPart}
\leanref{Development/AutomorphicArithmetic/Foundations/GaloisRepresentations/AnnihilatorModels/PMatchPart.lean}{UpperTheory.Gal.crystallineComparison.matchPart_of}
\leanref{Development/AutomorphicArithmetic/Foundations/GaloisRepresentations/AnnihilatorModels/SameNormBasic.lean}{UpperTheory.Gal.crystallineComparison.PermOfMap}
\leanref{Development/AutomorphicArithmetic/Foundations/GaloisRepresentations/AnnihilatorModels/PPermMap.lean}{UpperTheory.Gal.crystallineComparison.permMap_of}

\begin{proof}
Let $\rho:B\to\bar B=B/(1-e)B$ be the quotient map, so that $\rho(e)=1$. If $\rho(a)=\rho(b)$, then
$a-b\in(1-e)B$, so $ea=eb$. Hence $\rho(a)\mid\rho(b)$ implies $a\mid eb$: if $\rho(b)=\rho(at)$, then
$eb=eat$. Conversely, $a\mid eb$ implies $\rho(a)\mid\rho(b)$, since $\rho(eb)=\rho(b)$.

If $\rho(w)$ is a unit of $\bar B$, it divides every element, and every $\sigma$ satisfies the claim.
Otherwise $\rho(w)$ is a nonzerodivisor: if $\rho(z)\rho(w)=0$, then $ezw=0$, so $ez=0$ and
$\rho(z)=\rho(ez)=0$. Let $u$ and $u'$ be the multisets of the $\rho(c_i)$ and of the $\rho(c'_i)$. Then
$\Phi_u\equiv\Phi_{u'}\bmod\rho(w)^{(N+1)^N}$ by Lemma~\ref{gal-crystalline-matching-forms}(1). By the
second implication, (a) with $n=(N+1)^j$ and (b) give hypotheses (a) and (b) of
Lemma~\ref{gal-crystalline-matching-induction} for $u$, $u'$ and $\rho(w)$. So $u\equiv u'\bmod\rho(w)$,
and there is a permutation $\sigma$ with $\rho(w)\mid\rho(c_i)-\rho(c'_{\sigma(i)})$ for every $i$. By
the first implication, $w\mid e(c_i-c'_{\sigma(i)})$.
\end{proof}

\begin{lemma}\label{gal-crystalline-matching-local}
Let $B$ be a ring and $w\in B$ such that, whenever an ideal $(b_1,b_2)$ of $B$ contains a power of $w$,
there is an idempotent $e$ with $eb_1\in b_2B$ and $(1-e)b_2\in b_1B$. Let $\mathfrak p$ be a prime
ideal of $B$, $D\subset B$ a finite set and $M\geq0$. Then there is an idempotent $e\notin\mathfrak p$
such that:
\begin{enumerate}
\item for all $d\in D$ and $0\leq n\leq M$, $d\mid ew^n$ or $w^n\mid ed$;
\item for all $d,d'\in D$ with $d\mid ew^M$, $d\mid ed'$ or $d'\mid ed$.
\end{enumerate}
\end{lemma}
\leanref{Development/AutomorphicArithmetic/Foundations/GaloisRepresentations/AnnihilatorModels/GenericMulOfProd.lean}{UpperTheory.Gal.crystallineComparison.ValLike}
\leanref{Development/AutomorphicArithmetic/Foundations/GaloisRepresentations/AnnihilatorModels/SameNormBasic.lean}{UpperTheory.Gal.crystallineComparison.ValLocal}
\leanref{Development/AutomorphicArithmetic/Foundations/GaloisRepresentations/AnnihilatorModels/PValLocal.lean}{UpperTheory.Gal.crystallineComparison.valLocal_of}

\begin{proof}
For an idempotent $e$, one of $e$ and $1-e$ is not in $\mathfrak p$, since their sum is $1$. A product
of idempotents not in $\mathfrak p$ is an idempotent not in $\mathfrak p$, and a divisibility $a\mid eb$
persists when $e$ is replaced by $e'e$. So it suffices to find idempotents not in $\mathfrak p$ of two
kinds: for each $d\in D$ and $n\leq M$, an $e_{d,n}$ with $d\mid e_{d,n}w^n$ or $w^n\mid e_{d,n}d$; and
for all $d,d'\in D$, an $e_{d,d'}$ with $d\mid e_{d,d'}d'$ or $d'\mid e_{d,d'}d$ whenever $d\mid e_0w^M$
for some idempotent $e_0\notin\mathfrak p$. The product $e$ of all of them satisfies (1), and it
satisfies (2) by taking $e_0=e$.

Let $d\in D$ and $n\leq M$. The ideal $(d,w^n)$ contains $w^n$, so there is an idempotent $e_1$ with
$e_1d\in w^nB$ and $(1-e_1)w^n\in dB$. Take $e_{d,n}=e_1$ if $e_1\notin\mathfrak p$, and
$e_{d,n}=1-e_1$ otherwise.

Let $d,d'\in D$. If there is no idempotent $e_0\notin\mathfrak p$ with $d\mid e_0w^M$, take
$e_{d,d'}=1$. Otherwise write $e_0w^M=db$ for such an $e_0$, and put $f=e_0d+1-e_0$. Then $e_0f=e_0d$,
and
\[
 w^M=\bigl(e_0b+(1-e_0)w^M\bigr)f,
\]
since $e_0bd=e_0w^M$ and $e_0(1-e_0)=0$. So the ideal $(f,d')$ contains $w^M$, and there is an
idempotent $e_1$ with $d'\mid e_1f$ and $f\mid(1-e_1)d'$. Multiplying by $e_0$ gives
$d'\mid e_0e_1f=e_0e_1d$, and $d\mid e_0(1-e_1)d'$, because $e_0(1-e_1)d'$ is a multiple of
$e_0f=e_0d$. Take $e_{d,d'}=e_0e_1$ if $e_1\notin\mathfrak p$, and $e_{d,d'}=e_0(1-e_1)$ otherwise.
\end{proof}

\begin{lemma}\label{gal-crystalline-matching-cover}
Let $B$ be a ring and $(J_\lambda)_{\lambda\in\Lambda}$ a finite family of ideals of $B$ such that for
every maximal ideal $\mathfrak m$ of $B$ there are $\lambda\in\Lambda$ and an idempotent $e\in J_\lambda$
with $e\notin\mathfrak m$. Then there are orthogonal idempotents $e_\lambda\in J_\lambda$,
$\lambda\in\Lambda$, with $\sum_\lambda e_\lambda=1$.
\end{lemma}
\leanref{Development/AutomorphicArithmetic/Foundations/GaloisRepresentations/AnnihilatorModels/SameNormBasic.lean}{UpperTheory.Gal.crystallineComparison.IdemCover}
\leanref{Development/AutomorphicArithmetic/Foundations/GaloisRepresentations/AnnihilatorModels/PIdemCover.lean}{UpperTheory.Gal.crystallineComparison.idemCover_of}

\begin{proof}
The idempotents that lie in some $J_\lambda$ lie in no common maximal ideal, so they generate the unit
ideal, and a finite set $E$ of them generates it. For $\lambda\in\Lambda$ let
$f_\lambda=1-\prod_{e\in E\cap J_\lambda}(1-e)$. It is an idempotent, and it lies in $J_\lambda$, since
$1-\prod_e(1-e)$ is a combination of the $e$. The product $g=\prod_\lambda(1-f_\lambda)$ is divisible by
$1-e$ for every $e\in E$, and $e(1-e)=0$, so $eg=0$ for every $e\in E$. As $E$ generates the unit ideal,
$g=0$. Order $\Lambda$ and put $e_\lambda=f_\lambda\prod_{\mu<\lambda}(1-f_\mu)$. These are idempotents
in $J_\lambda$. For $\mu<\lambda$, $e_\lambda e_\mu=0$, since $e_\lambda$ is a multiple of $1-f_\mu$ and
$e_\mu$ of $f_\mu$. Finally $\sum_\lambda e_\lambda=1-\prod_\lambda(1-f_\lambda)=1-g=1$.
\end{proof}

\begin{lemma}\label{gal-crystalline-matching}
Let $\Phi_c$ be as in Lemma~\ref{gal-crystalline-norm-forms} and $N\geq0$.
\begin{enumerate}
\item Let $B$ be a ring and $w\in B$ a nonzerodivisor such that for $b_1,b_2\in B$ with
  $w^M\in(b_1,b_2)$ for some $M$ there is an idempotent $e$ with $eb_1\in b_2B$ and
  $(1-e)b_2\in b_1B$. Let $c_1,\ldots,c_N,c'_1,\ldots,c'_N\in B^m$ with
  $\prod_i\Phi_{c_i}\equiv\prod_i\Phi_{c'_i}\bmod w^{(N+1)^N}$. Then there are orthogonal idempotents
  $e_\sigma$, $\sigma\in S_N$, with sum $1$, such that $e_\sigma(c_i-c'_{\sigma(i)})\in wB^m$ for all
  $\sigma$ and $i$.
\item Let $A$ be a test algebra (Definition~\ref{gal-crystalline-test-algebras}), $a\geq0$ with
  $(N+1)^N\leq\ell^a$, and $c_1,\ldots,c_N,c'_1,\ldots,c'_N\in(A/\ell^\flat)^m$ with
  $\prod_i\Phi_{c_i}=\prod_i\Phi_{c'_i}$. Then there are orthogonal idempotents $e_\sigma$ of $A$,
  $\sigma\in S_N$, with sum $1$, such that $e_\sigma c_i^{\ell^a}=e_\sigma c'^{\,\ell^a}_{\sigma(i)}$ in
  $(A/\ell^\flat)^m$ for all $\sigma$ and $i$, with powers taken coordinatewise.
\end{enumerate}
\end{lemma}
\leanref{Development/AutomorphicArithmetic/Foundations/GaloisRepresentations/AnnihilatorModels/GenericMulOfProd.lean}{UpperTheory.Gal.crystallineComparison.MatchingVal}
\leanref{Development/AutomorphicArithmetic/Foundations/GaloisRepresentations/AnnihilatorModels/PMatchVal.lean}{UpperTheory.Gal.crystallineComparison.matchVal_of}
\leanref{Development/AutomorphicArithmetic/EtaleIdempotents/ExistsOrthogonalIdempotentsMatchingOfSameNorm.lean}{UpperTheory.exists_orthogonalIdempotents_matching_of_sameNorm}

\begin{proof}
(1) For $\sigma\in S_N$ let $J_\sigma$ be the ideal of the $b\in B$ with $w\mid b(c_i-c'_{\sigma(i)})$
for every $i$. Let $\mathfrak m$ be a maximal ideal of $B$, and $D$ the finite set of the differences
$c_{i,k}-c'_{i',k}$. Lemma~\ref{gal-crystalline-matching-local}, applied to $\mathfrak m$, $D$ and
$M=(N+1)^N$, gives an idempotent $e\notin\mathfrak m$ that satisfies hypotheses (a) and (b) of
Lemma~\ref{gal-crystalline-matching-piece}. By that lemma there is $\sigma\in S_N$ with $e\in J_\sigma$.
Lemma~\ref{gal-crystalline-matching-cover} then gives orthogonal idempotents $e_\sigma\in J_\sigma$ with
sum $1$, which is the claim.

(2) Since $A$ is perfect, there is $w\in A$ with $w^{\ell^a}=\ell^\flat$. It is a nonzerodivisor, and
$A$ has the property of (1) for $w$, since $w^M\in(b_1,b_2)$ implies $(\ell^\flat)^M\in(b_1,b_2)$. As
$w^{(N+1)^N}$ divides $\ell^\flat$, lifts $\tilde c_i,\tilde c'_i\in A^m$ satisfy the hypothesis of (1),
and there are orthogonal idempotents $e_\sigma$ with sum $1$ and
$e_\sigma(\tilde c_i-\tilde c'_{\sigma(i)})\in wA^m$. The $\ell^a$-power map is additive on $A$ and fixes
idempotents, so $e_\sigma(\tilde c_i^{\ell^a}-\tilde c'^{\,\ell^a}_{\sigma(i)})\in w^{\ell^a}A^m=\ell^\flat A^m$.
\end{proof}

\begin{lemma}\label{gal-crystalline-points-relation}
Let $\mathcal O$, $\varpi$, $\ell$, $q=\ell^f$ and the $T_j$, $0\leq j\leq n$, be as in
Lemma~\ref{gal-crystalline-norm-identity}, with $Y$ smooth and quasi-compact over $\mathcal O$ and every
finite set of points of $Y$ contained in an affine open, and with $p_j$ and $r_j$ over the same point
of $\operatorname{Spec}\mathcal O$. Let $\kappa:\mathcal O\to\mathcal O_C$ and
$g:\mathcal O\to\mathcal O_C/\ell$ be ring maps with $g(\varpi)=0$ and $g(x)^q=g(x)$, and $m\geq0$
with $(\kappa(x)\bmod\ell)^{\ell^m}=g(x)$ for all $x$, so that $F$ and $\rho_m$ of
Definition~\ref{gal-crystalline-test-algebras}(4) are defined. Let $\pi:Y'\to Y$ and $s_{j,i}$ be as in
Lemma~\ref{gal-crystalline-splitting-cover}, let $\tau_{j,i}=(r_js_{j,i})_*:Y'^\sharp\to Y^\sharp$, let
$D=\sum_{j\text{ even}}q^{j(j-1)/2}d_j$ and let $N_0\geq0$ with $(D+1)^D\leq q^{N_0}$. Then the sum
\[
 \sum_{j=0}^n\sum_{i=1}^{d_j}(-1)^jq^{j(j-1)/2}\,F^{n-j+N_0}\circ\rho_m\circ\tau_{j,i}
\]
of natural transformations $Y'^\sharp\to Y_g$ vanishes locally.
\end{lemma}
\leanref{Development/AutomorphicArithmetic/CommutativeAlgebra/EtaleCohomology/TestFacts.lean}{UpperTheory.Ru.crystallineComparison.WZero}
\leanref{Development/AutomorphicArithmetic/CommutativeAlgebra/EtaleCohomology/MapUSharp.lean}{UpperTheory.Ru.crystallineComparison.wZero}

\begin{proof}
Let $A$ be a test algebra and $y\in Y'^\sharp(A)$. Every point $y_0$ of $Y$ has an affine
neighbourhood $U$ with an affine open $U'$ such that $p_j^{-1}(U)\subset r_j^{-1}(U')$ for every $j$:
take $U'$ affine containing the finite set $\{y_0\}\cup\bigcup_jr_j(p_j^{-1}(y_0))$, and $U$ an affine
neighbourhood of $y_0$ outside the closed set $\bigcup_jp_j(W_j\setminus r_j^{-1}(U'))$. By
Lemma~\ref{gal-crystalline-test-facts}(4) we may replace $A$ by the pieces of a decomposition and assume
that $\pi y$ factors through such a $U$. Let $x_1,\ldots,x_m$ generate $\Gamma(U',\mathcal O_Y)$ as an
$\mathcal O$-algebra; two $\mathcal O$-morphisms from an affine scheme to $U'$ that agree on the $x_k$
are equal.

Let $\alpha_{j,i}:\Gamma(U',\mathcal O_Y)\to A^\sharp$ correspond to $r_js_{j,i}y$. The maps
$\Gamma(p_j^{-1}U,\mathcal O_{W_j})\to A^\sharp$ given by the points $s_{j,i}y$ split over
$A^\sharp[1/\ell]$ by Lemma~\ref{gal-crystalline-splitting-cover}(1), $\ell$ is a nonzerodivisor in
$A^\sharp$ (Lemma~\ref{gal-crystalline-test-facts}(1)), and $\psi_m$ kills $\varpi$ because
$\psi_m\circ\kappa=g$. Lemma~\ref{gal-crystalline-norm-identity} with $B=A^\sharp$ and $\beta=\psi_m$
gives equal products of the forms of the vectors
$c_{j,i}=\bigl(\psi_m(\alpha_{j,i}(x_k))^{q^{n-j}}\bigr)_k\in(A/\ell^\flat)^m$, with multiplicities
$q^{j(j-1)/2}$, over even and over odd $j$. If $A\neq0$, comparing degrees in $X$ shows that both
families have $D$ members. Lemma~\ref{gal-crystalline-matching}(2), with $\ell^a=q^{N_0}$, and
Lemma~\ref{gal-crystalline-test-facts}(4) give a decomposition of $A$ on each piece of which the
$q^{N_0}$-th powers of the even family agree with those of the odd family after a bijection. Now
$F^{n-j+N_0}\rho_m\tau_{j,i}(y)$ is the point of $U'$ with coordinates
$\psi_m(\alpha_{j,i}(x_k))^{q^{n-j+N_0}}$, the entries of $c_{j,i}^{q^{N_0}}$. So on each piece the
weighted multisets of these points over even and over odd $j$ agree, which is the claim.
\end{proof}
\leanref{Development/AutomorphicArithmetic/EtaleIdempotents/CorrespondenceNormBoundOfAffineOpens.lean}{UpperTheory.correspondence_norm_bound_of_affine_opens}
\leanref{Development/AutomorphicArithmetic/Foundations/GaloisRepresentations/AnnihilatorModels/PCharts.lean}{UpperTheory.Gal.crystallineComparison.charts_of}

\begin{lemma}\label{gal-crystalline-relation}
Let $\ell$ be a prime, $K$ a finite extension of $\mathbf Q_\ell$ inside an algebraic closure $\bar K$,
with valuation ring $\mathcal O_K$, a uniformizer $\varpi$ and residue field with $q=\ell^f$ elements,
and let $Y$ be a smooth projective $\mathcal O_K$-scheme. Let $T_j=(W_j,p_j,r_j)$, $0\leq j\leq n$, be
finite correspondences of $Y$ (Definition~\ref{gal-tower-reduction}(2)), with $p_j$ and $r_j$ over the
same point of $\operatorname{Spec}\mathcal O_K$ and with $d_j$ geometric points in the fibres of $p_j$
over $Y[1/\ell]$, that satisfy the Eichler--Shimura relation at the generic points of the special
fibre of Definition~\ref{gal-tower-reduction}(3) for $\varpi$ and $q$. Let $k\geq0$,
$V=H^k(Y_{\bar K},\mathbf Q_\ell)$, and let $t_0,\ldots,t_n$ be endomorphisms of $V$ with the following
property: for every scheme $X$ with a morphism $w:X\to Y$, an $\mathcal O_K$-scheme through $w$, and
all morphisms $z_1,\ldots,z_{d_j}:X\to W_j$ with $p_jz_i=w$ that induce an isomorphism
$\coprod_iX[1/\ell]\to X[1/\ell]\times_YW_j$, one has $w^*\circ t_j=\sum_i(r_jz_i)^*$ as maps
$V\to H^k(X_{\bar K},\mathbf Q_\ell)$. Let $C$ be the completion of $\bar K$, and $B$ a crystalline period package
for $K$, with $c$, $\phi_B$ and $\sigma_B$ (Definition~\ref{gal-crystalline-relation-package}). Then there is $N\geq0$
such that
\[
 \sum_{j=0}^n(-1)^jq^{j(j-1)/2}\,\bigl(t_j\otimes\phi_B^{f(n-j+N)}\bigr)(x)=0
\]
for every $x\in V\otimes_{\mathbf Q_\ell}B$ with $(\sigma\otimes\sigma_B(\sigma))(x)=x$ for all
$\sigma\in\operatorname{Gal}(\bar K/K)$. For $B=B_{\mathrm{cris}}$ these $x$ form
$D_{\mathrm{cris}}(V)=(V\otimes_{\mathbf Q_\ell}B_{\mathrm{cris}})^{\operatorname{Gal}(\bar K/K)}$, and $\phi_B$ is
the Frobenius $\phi$ of $B_{\mathrm{cris}}$.
\end{lemma}
\leanref{Development/AutomorphicArithmetic/Foundations/GaloisRepresentations/AnnihilatorModels/LocFacts.lean}{UpperTheory.Gal.crystallineComparison.Br.FinAll}
\leanref{Development/AutomorphicArithmetic/CommutativeAlgebra/EtaleCohomology/VecZero.lean}{UpperTheory.Ru.crystallineComparison.IsHk}
\leanref{Development/AutomorphicArithmetic/CommutativeAlgebra/EtaleCohomology/SmoothModelCrystalline.lean}{UpperTheory.Ru.crystallineComparison.finAll}

\begin{proof}
Let $A_{\mathrm{cris}}$, $t$, $\kappa$, $\eta_\nu$, $g$ and $F$ be as in
Theorem~\ref{gal-crystalline-realization}, and $m$, $N_1$ and $\rho=\rho_m$ as in
Theorem~\ref{gal-crystalline-realization}(3) for $B$. Write $\sigma_j=(-1)^jq^{j(j-1)/2}$. Let $\pi:Y'\to Y$ and
$s_{j,i}$ be the splitting cover of the $T_j$ (Lemma~\ref{gal-crystalline-splitting-cover}), with the
group $G=\prod_jS_{d_j}$, and put $\tau_{j,i}=(r_js_{j,i})_*:Y'^\sharp\to Y^\sharp$. By
Lemma~\ref{gal-crystalline-points-relation} there is $N_0$ such that
$\sum_{j,i}\sigma_jF^{n-j+N_0}\rho\tau_{j,i}$ vanishes locally. Put $N=N_0+N_1$ and $E=f(n+N_0)$.

Let $x\in V\otimes_{\mathbf Q_\ell}B$ be fixed by every $\sigma\otimes\sigma_B(\sigma)$. Every element of $B$ is
$c(a)c(t)^{-r}$ with $a\in A_{\mathrm{cris}}$ (Definition~\ref{gal-crystalline-relation-package}(2)), and $\ell$ is
invertible in $B$, so $x=\sum_\alpha v_\alpha\otimes c(c_\alpha)c(t)^{-r}$ with
$v_\alpha\in H^k(Y_{\bar K},\mathbf Z_\ell)$, $c_\alpha\in A_{\mathrm{cris}}$ and $r\geq0$. By
Theorem~\ref{gal-crystalline-realization}(3) there are
$a,u\geq0$ and classes $X_\nu$ with $\phi^fX_\nu=F^*X_\nu$ and
\[
 \ell^at^u\sum_\alpha\phi^{fN_1}(c_\alpha)\eta_\nu(v_\alpha)=t^{u+r}Z_\nu,\qquad Z_\nu=\rho^*X_\nu,
 \tag{a}
\]
for every $\nu$. Fix $\nu$.

Pullbacks commute with $\phi$ (Lemma~\ref{gal-crystalline-test-cohomology}(1)), so
$\phi^{fe}X_\nu=(F^e)^*X_\nu$ for $e\geq0$, and
$\tau_{j,i}^*\phi^{f(n-j+N_0)}Z_\nu=(F^{n-j+N_0}\rho\tau_{j,i})^*X_\nu$. By
Lemma~\ref{gal-crystalline-test-cohomology}(2),
\[
 \sum_{j,i}\sigma_j\tau_{j,i}^*\phi^{f(n-j+N_0)}Z_\nu=0\quad\text{in }H^k(Y'^\sharp,A_{\mathrm{cris}}/\ell^\nu).
 \tag{b}
\]
Apply $\phi^{f(n-j+N_0)}$ to (a): $\phi$ is semilinear, fixes the $\eta_\nu(v_\alpha)$ and satisfies
$\phi(t)=\ell t$. Multiplying by $\ell^{fj(u+r)}$ gives
\[
 \ell^{a+Eu}t^u\,\ell^{fjr}\sum_\alpha\phi^{f(n-j+N)}(c_\alpha)\eta_\nu(v_\alpha)
 =\ell^{E(u+r)}t^{u+r}\phi^{f(n-j+N_0)}Z_\nu.
 \tag{c}
\]

Apply the hypothesis on $t_j$ to $w=\pi$ and $z_i=s_{j,i}$, which is allowed by
Lemma~\ref{gal-crystalline-splitting-cover}(1): $\pi^*t_j(v)=\sum_i(r_js_{j,i})^*v$ for $v\in V$. Every
element of $V$ has a multiple by a power of $\ell$ in the image of $H^k(Y_{\bar K},\mathbf Z_\ell)$, and
a class of $H^k(Y'_{\bar K},\mathbf Z_\ell)$ with image $0$ in $H^k(Y'_{\bar K},\mathbf Q_\ell)$ is
killed by a power of $\ell$. So there are $b_1,b_2\geq0$ and classes
$y_{j,\alpha}\in H^k(Y_{\bar K},\mathbf Z_\ell)$ with image $\ell^{b_1}t_j(v_\alpha)$ in $V$ and
$\ell^{b_2}\pi^*y_{j,\alpha}=\ell^{b_1+b_2}\sum_i(r_js_{j,i})^*v_\alpha$. By the naturality of
$\eta_\nu$,
\[
 \ell^{b_2}\pi_*^*\eta_\nu(y_{j,\alpha})=\ell^{b_1+b_2}\sum_i\tau_{j,i}^*\eta_\nu(v_\alpha).
 \tag{d}
\]
Let $L=\ell^{a+Eu}t^u$ and $Q_\nu=\sum_{j,\alpha}\sigma_j\ell^{fjr}\phi^{f(n-j+N)}(c_\alpha)
\eta_\nu(y_{j,\alpha})$. By (d), (c) and (b),
\[
 \ell^{b_2}L\,\pi_*^*Q_\nu=\ell^{b_1+b_2}\sum_j\sigma_j\sum_i\tau_{j,i}^*\Bigl(L\ell^{fjr}\sum_\alpha
 \phi^{f(n-j+N)}(c_\alpha)\eta_\nu(v_\alpha)\Bigr)
 =\ell^{b_1+b_2+E(u+r)}t^{u+r}\sum_{j,i}\sigma_j\tau_{j,i}^*\phi^{f(n-j+N_0)}Z_\nu=0.
\]
By Lemma~\ref{gal-crystalline-splitting-cover}(3) and Lemma~\ref{gal-crystalline-test-cohomology}(3),
$|G|\ell^{b_2}LQ_\nu=0$. With the integer $\lambda=|G|\ell^{a+Eu+b_2}$, this says
$t^u\sum_{j,\alpha}\lambda\sigma_j\ell^{fjr}\phi^{f(n-j+N)}(c_\alpha)\eta_\nu(y_{j,\alpha})=0$, for
every $\nu$.

By Theorem~\ref{gal-crystalline-realization}(2),
$\sum_{j,\alpha}y_{j,\alpha}\otimes c(\lambda\sigma_j\ell^{fjr}\phi^{f(n-j+N)}(c_\alpha))=0$ in
$V\otimes_{\mathbf Q_\ell}B$, that is,
\[
 \lambda\ell^{b_1}\sum_j\sigma_j\ell^{fjr}\sum_\alpha t_j(v_\alpha)\otimes c(\phi^{f(n-j+N)}(c_\alpha))=0.
\]
Since $\phi_B(c(t))=c(\phi(t))=\ell c(t)$ and $\phi_B\circ c=c\circ\phi$, $(t_j\otimes\phi_B^{f(n-j+N)})(x)=
\ell^{-f(n-j+N)r}c(t)^{-r}\sum_\alpha t_j(v_\alpha)\otimes c(\phi^{f(n-j+N)}(c_\alpha))$. So
$\sum_j\sigma_j(t_j\otimes\phi_B^{f(n-j+N)})(x)$ is $(\lambda\ell^{b_1+f(n+N)r}c(t)^r)^{-1}$ times the left side of
the last display, and it vanishes.
\end{proof}
\leanref{Development/AutomorphicArithmetic/CommutativeAlgebra/EtaleCohomology/QuasiCompactOfProj.lean}{UpperTheory.Ru.crystallineComparison.final}
\leanref{Development/AutomorphicArithmetic/CommutativeAlgebra/EtaleCohomology/PhiHIterSmul.lean}{UpperTheory.Ru.crystallineComparison.top}
\leanref{Development/AutomorphicArithmetic/CommutativeAlgebra/EtaleCohomology/VecZero.lean}{UpperTheory.Ru.crystallineComparison.SiteAnn}
\leanref{Development/AutomorphicArithmetic/CommutativeAlgebra/EtaleCohomology/PhiIterPull.lean}{UpperTheory.Ru.crystallineComparison.siteAnn}
\leanref{Development/AutomorphicArithmetic/EtaleIdempotents/ExistsCrystallineLatticeExpansion.lean}{UpperTheory.exists_crystalline_lattice_expansion}
\leanref{Development/AutomorphicArithmetic/EtaleIdempotents/HeckeCorrespondenceProduct.lean}{UpperTheory.hecke_correspondence_product}
\leanref{Development/AutomorphicArithmetic/EtaleIdempotents/CrystallineExpansionRelationOfScaledEndomorphism.lean}{UpperTheory.crystalline_expansion_relation_of_scaled_endomorphisms}

\subsection{Packs, plans and the congruence relation}\label{subsection-gal-packs-plans}

\begin{proposition}\label{gal-tower-cohomology}
Let $E$ be a number field, $\Omega$ an algebraic closure, $\ell$ a prime,
$\iota:\bar{\mathbf Q}_\ell\to\mathbf C$ a field isomorphism, $H$ a smooth projective Hecke tower over
$E$ and $\theta:\Gamma\to\bar{\mathbf Q}_\ell$.
\begin{enumerate}
\item $H$ has a piece $\rho$ at $\theta$ that is continuous and semisimple.
\item If $H$ has a good model at $w$, it has a good surjective model at $w$; if it has a model at $w$
  satisfying a congruence relation, it has a surjective one satisfying the same relation.
\item Let $\rho$ be a continuous piece at $\theta$ and $w$ a place not over $\ell$ at which $H$ has a
  surjective model satisfying a congruence relation of degree $n$ whose Hecke polynomial is
  $\vartheta_0\prod_{z\in R}(X-\iota^{-1}(z))$ with $\vartheta_0\neq0$. Then $\rho$ is unramified at
  $w$ and $\prod_{z\in R}(\rho(\varphi)-\iota^{-1}(z))=0$ for every geometric Frobenius element
  $\varphi$ at $w$.
\item Let $\rho$ be a continuous piece at $\theta$ and $w$ a place over $\ell$ at which $H$ has a
  surjective model. Then $\rho|_{D_A}$ is crystalline for every valuation ring $A$ over $w$. If the
  model is also good and $\iota\circ\theta$ occurs in Hodge degree $j$ along an embedding $\tau$
  inducing $w$, then $j$ is a Hodge--Tate weight of $\rho|_{D_A}$ at $\tau$.
\item In the setting of (4), if the model satisfies a congruence relation with Hecke polynomial
  $\vartheta_0\prod_{z\in R}(X-\iota^{-1}(z))$, then for some geometric Frobenius element $\varphi\in W_A$
  the characteristic polynomial of $\mathrm{WD}_\iota(\rho|_{D_A})(\varphi)$ divides a power of
  $\prod_{z\in R}(X-z)$.
\end{enumerate}
\end{proposition}
\leanref{Development/AutomorphicArithmetic/Foundations/GaloisRepresentations/IntegralStructures/GeometricFlagBases/HodgeTrAt.lean}{UpperTheory.Gal.Int.Geo.RepEx}
\leanref{Development/AutomorphicArithmetic/Foundations/GaloisRepresentations/IntegralStructures/GeometricFlagBases/HodgeTrAt.lean}{UpperTheory.Gal.Int.Geo.ESStep}
\leanref{Development/AutomorphicArithmetic/Foundations/GaloisRepresentations/IntegralStructures/GeometricFlagBases/HodgeTrAt.lean}{UpperTheory.Gal.Int.Geo.CrysStep}
\leanref{Development/AutomorphicArithmetic/Foundations/GaloisRepresentations/IntegralStructures/GeometricFlagBases/HodgeTrAt.lean}{UpperTheory.Gal.Int.Geo.JumpStep}
\leanref{Development/AutomorphicArithmetic/Foundations/GaloisRepresentations/IntegralStructures/GeometricFlagBases/HodgeTrAt.lean}{UpperTheory.Gal.Int.Geo.PhiStep}
\leanref{Development/AutomorphicArithmetic/Foundations/GaloisRepresentations/IntegralStructures/GeometricFlagBases/HodgeTrAt.lean}{UpperTheory.Gal.Int.Geo.CutStep}
\leanref{Development/AutomorphicArithmetic/Foundations/GaloisRepresentations/IntegralStructures/GeometricFlagBases/MapInvMul.lean}{UpperTheory.Gal.Int.Geo.Steps}
\leanref{Development/AutomorphicArithmetic/Foundations/GaloisRepresentations/IntegralStructures/GeometricFlagBases/GeometricRealizationOfWeightOne.lean}{UpperTheory.Gal.Int.Geo.steps_holds}
\leanref{Development/AutomorphicArithmetic/Foundations/GaloisRepresentations/IntegralStructures/GeometricFlagBases/RepMain.lean}{UpperTheory.Gal.Int.Geo.repEx_holds}
\leanref{Development/AutomorphicArithmetic/Foundations/GaloisRepresentations/BasicOpenSections/PEsStep.lean}{UpperTheory.Gal.Mod.esStep}
\leanref{Development/AutomorphicArithmetic/Foundations/GaloisRepresentations/CrystallineTowers/CrysStep.lean}{UpperTheory.Gal.Seam.crysStep}
\leanref{Development/AutomorphicArithmetic/Foundations/GaloisRepresentations/PAdicWittLattices/TraceSplitHolds.lean}{UpperTheory.Gal.G3.jumpStep_of}
\leanref{Development/AutomorphicArithmetic/CommutativeAlgebra/EtaleCohomology/PhiStepOf.lean}{UpperTheory.Ru.crystallineComparison.phiStep_of}
\leanref{Development/AutomorphicArithmetic/Foundations/GaloisRepresentations/IntegralStructures/GeometricFlagBases/PGeoSide.lean}{UpperTheory.Gal.Int.Geo.geoSide_of}

\begin{proof}
Write $Z$ for the base of $H$ and $H^*$ for the total cohomology
$\bigoplus_nH^n(Z_\Omega,\mathbf Q_\ell)$, and $M\subset\bar{\mathbf Q}_\ell\otimes H^*$ for the
$\theta$-part.

(1) By Theorem~\ref{gal-etale-foundations}(1), $H^*$ is finite-dimensional, and
$\operatorname{Gal}(\Omega/E)$ acts on it continuously. Each Hecke operator $T_\gamma$ is defined by
morphisms over $E$, so it commutes with the Galois action, and $M$ is Galois-stable. Choose a
composition series of $M$ as a Galois representation and a basis adapted to it: the matrices of the
Galois action are block upper triangular, and their block diagonal part $\rho$ is continuous and
semisimple. This is a piece of $H$ at $\theta$.

(2) Given a model $Z\to Y$, replace $Y$ by the scheme-theoretic image of $Z$, which is open and closed
in $Y$ because $Z\to Y_E$ is finite \'etale and $Y$ is smooth over $\mathcal O_w$. The correspondences
witnessing goodness or the congruence relation restrict to it.

(3) Let $Z\to Y$ be the surjective model, with group $\Delta$. Since $Z\to Y_E$ is a finite \'etale
Galois covering with group $\Delta$ and $\theta=1$ on the Hecke correspondences given by $\Delta$, the
$\theta$-part lies in the $\Delta$-invariants, which pullback identifies with
$\bar{\mathbf Q}_\ell\otimes H^*(Y_\Omega,\mathbf Q_\ell)$. Since $T_j$ lies over $\gamma_j$ with
multiplicity $m_j$, the Hecke operator $T_{\gamma_j}$ acts on these invariants as $m_j$ times the
operator $t_j$ of the correspondence $T_j$ of $Y$, so $t_j$ acts on $M$ by the scalar
$\theta(\gamma_j)/m_j=\vartheta_j$. Let $A$ be a valuation ring over $w$ and $\sigma\in D_A$. If
$\sigma\in I_A$, Lemma~\ref{gal-eichler-shimura-cohomology}(1) shows that $\sigma$ acts trivially on
$M$. If $\sigma$ is an arithmetic Frobenius element, Lemma~\ref{gal-eichler-shimura-cohomology}(2),
applied to the model over the valuation ring of $E_w$, gives
$\sum_j(-1)^jq^{j(j-1)/2}\vartheta_j\sigma^{-(n-j)}=0$ on $M$, that is, $P_\vartheta(\sigma^{-1})=0$
on $M$. Writing $P_\vartheta=\vartheta_0\prod_{z\in
R}(X-\iota^{-1}(z))$ with $\vartheta_0\neq0$ gives the claim for the matrix of the Galois action on $M$,
and a polynomial that kills a block upper triangular matrix kills its block diagonal part, so it holds
for $\rho$.

(4) Let $A$ be a valuation ring over $w$, $K$ the closure of $E$ in the completion of $\Omega$ at $A$,
and $\mathcal O_K$ its valuation ring. Base change along $E\to K$ turns $H$ into a smooth projective
Hecke tower over $K$ whose $\theta$-part carries the restriction to $D_A$ of the Galois action on $M$,
and turns the surjective model into a surjective model $Z_K\to Y_{\mathcal O_K}$ with
$Y_{\mathcal O_K}$ smooth and projective over $\mathcal O_K$. As in (3), pullback embeds $M$
equivariantly into $\bar{\mathbf Q}_\ell\otimes\bigoplus_{n<N}H^n(Y_{\bar K},\mathbf Q_\ell)$, which
is crystalline by Theorem~\ref{gal-padic-comparison}(1). A subrepresentation of a crystalline
representation is crystalline, and so is a subquotient, hence so is the block diagonal part
$\rho|_{D_A}$. For the Hodge--Tate weights, suppose the model is good and $\iota\circ\theta$ occurs in
Hodge degree $j$ along $\tau$: there is a class in $H^b(Z,\Omega^j)$ on which every $s_{\gamma*}
t_\gamma^*$ acts by $\iota\theta(\gamma)$. It is $\Delta$-invariant, so it comes from
$H^b(Y_K,\Omega^j)$, and since the model is good the Hecke correspondences come from finite \'etale
correspondences of $Y$. The decomposition of Theorem~\ref{gal-padic-comparison}(2) for
$H^{b+j}(Y_{\bar K})$ is natural for these correspondences, so the $\theta$-part of
$\mathbf C_K\otimes H^{b+j}(Y_{\bar K})$ contains a nonzero vector of $\mathbf C_K(-j)\otimes
H^b(Y_K,\Omega^j)$ at $\tau$. Hence $j$ is a Hodge--Tate weight of $M$ at $\tau$, and the multiset of
Hodge--Tate weights of $M$ is that of its block diagonal part $\rho|_{D_A}$.

(5) Let $K$ be as in (4); it is the completion of $E$ at $w$, whose residue field has $q=\ell^f$
elements. Base change along $\mathcal O_w\to\mathcal O_K$ carries the congruence relation of the model,
with its correspondences $T_j$, to the model over $\mathcal O_K$. As in (3) and (4), pullback embeds
$M$ equivariantly into $\bar{\mathbf Q}_\ell\otimes\bigoplus_{k<N}H^k(Y_{\bar K},\mathbf Q_\ell)$, and
the operator $t_j$ of $T_j$ on $H^k(Y_{\bar K},\mathbf Q_\ell)$ acts on $M$ by $\vartheta_j$. Each
$H^k(Y_{\bar K},\mathbf Q_\ell)$ is crystalline by Theorem~\ref{gal-padic-comparison}(1). The trace
along $p_j$ is the sum over the points of the fibres of $p_j$ (Lemma~\ref{gal-weighted-trace}(1)), so
pulling $t_j$ back along a morphism $w$ over which $p_j$ splits into sections $z_i$ gives
$\sum_i(r_jz_i)^*$; this is the hypothesis on $t_j$ in Lemma~\ref{gal-crystalline-relation}. By that
lemma, applied to the model over $\mathcal O_K$, the $T_j$ and the crystalline period package $B_{\mathrm{cris}}$
(Definition~\ref{gal-crystalline-relation-package}), there is $N'\geq0$ such that on
$D_{\mathrm{cris}}(H^k(Y_{\bar K},\mathbf Q_\ell))$, with its crystalline Frobenius $\phi$,
\[
 \sum_{j=0}^n(-1)^jq^{j(j-1)/2}\,t_j\,\phi^{f(n-j+N')}=0.
\]
Since $\phi$ is bijective on $D_{\mathrm{cris}}$, the identity gives $P_\vartheta(\phi^f)=0$ on
$D_{\mathrm{cris}}(M)$, and hence on the $D_{\mathrm{cris}}$ of every subquotient of $M$, in particular
of the block diagonal part $\rho|_{D_A}$, which is crystalline by (4). The Weil--Deligne representation
of a crystalline representation is unramified and monodromy-free, and a geometric Frobenius element
acts on it by $\phi^f$ \cite{Fon94}. Hence the characteristic polynomial of
$\mathrm{WD}_\iota(\rho|_{D_A})(\varphi)$, for $\varphi$ a geometric Frobenius element, divides a
power of $\iota(P_\vartheta)=\iota(\vartheta_0)\prod_{z\in R}(X-z)$.
\end{proof}

\begin{lemma}\label{gal-pack-seed}
In Situation~\ref{gal-seed-situation}, let $P$ be a pack for $U$ (Definition~\ref{gal-pack}), and let
$v$ be a place over $\ell$ at which $P$ has Hodge data and at which some weight of $U$ at an embedding
$\tau$ with $\iota^{-1}\circ\tau$ inducing $v$ is lacunary in the sense of hypothesis~(5) of
Theorem~\ref{gal-seed-representations}. Then there is a seed representation $\rho$ for $U$
(Definition~\ref{gal-seed}) such that:
\begin{enumerate}
\item at every place $w$ not over $\ell$ at which $P$ is congruent, $\rho$ satisfies the condition of
  item~(3)(a) of Definition~\ref{gal-seed} at $w$;
\item at every place $w$ over $\ell$ at which $P$ has Hodge data, $\rho|_{D_A}$ is crystalline for every
  $A$ over $w$, and $\mathrm{HT}_\tau(\rho|_{D_A})$ contains $\{k-\lambda_k:0\leq k<a\}$ for every
  $\tau:E\to\bar{\mathbf Q}_\ell$ extending continuously to $E_w$ and every weight $\lambda$ of $U$ at
  $\iota\circ\tau$;
\item at every place $w$ over $\ell$ at which $P$ is congruent, for every $B$ such that $U$ has Satake
  parameters $B$ at $w$ and every $A$ over $w$, $\rho|_{D_A}$ is crystalline and for some geometric
  Frobenius element $\varphi\in W_A$ the characteristic polynomial of
  $\mathrm{WD}_\iota(\rho|_{D_A})(\varphi)$ divides a power of $\prod_{b\in B}(X-q_w^{(a-1)/2}b)$.
\end{enumerate}
\end{lemma}
\leanref{Development/AutomorphicArithmetic/Foundations/GaloisRepresentations/IntegralStructures/GeometricFlagBases/MapInvMul.lean}{UpperTheory.Gal.Int.Geo.Pack.rep}

\begin{proof}
Let $\rho_0$ be a continuous semisimple piece of the tower at $\theta$
(Proposition~\ref{gal-tower-cohomology}(1)), and let $\chi_\kappa$ be the $\ell$-adic character
attached to the algebraic Hecke character $\kappa$. Put $\rho=\rho_0\otimes\chi_\kappa^{-1}$, which is
continuous and semisimple. At a place $w\nmid\ell$ where $P$ is congruent, $\kappa$ is unramified and
Proposition~\ref{gal-tower-cohomology}(3) applies with $R=\{q_w^{(a-1)/2}b\,\kappa(\varpi_w)^{-1}\}$;
twisting by $\chi_\kappa^{-1}$, whose value at a geometric Frobenius is $\iota^{-1}(\kappa(\varpi_w))$,
multiplies the roots by $\kappa(\varpi_w)$; with Proposition~\ref{gal-tower-cohomology}(2) to make the
model surjective, this gives (1). Since the pack is congruent at all but finitely
many $d$-split places, this gives item~(3)(a) of Definition~\ref{gal-seed}. At a place $w\mid\ell$ with
Hodge data, Proposition~\ref{gal-tower-cohomology}(2) and~(4) give crystallinity and the Hodge--Tate weights
$k-\lambda_k+m_\tau$ of $\rho_0$, and the twist by $\chi_\kappa^{-1}$, crystalline at $w$ because
$\kappa$ is unramified there, shifts them by $-m_\tau$; this is (2). At $v$, (2) gives a lacunary set of
$a$ Hodge--Tate weights, which is item~(3)(b) of Definition~\ref{gal-seed}. Item (3) follows in the same
way from Proposition~\ref{gal-tower-cohomology}(2) and~(5).
\end{proof}

\begin{definition}\label{gal-pel-datum}
Let $p$ be a prime and $E$ a CM field with complex conjugation $c$. A \emph{unitary PEL datum of
simple type} for $E$ at $p$ consists of:
\begin{enumerate}
\item an element $\omega\in\mathcal O_E$ generating an imaginary quadratic subfield
  $K_q=\mathbf Q(\omega)\subset E$, with $\omega+c(\omega)=t$, $\omega c(\omega)=m$, $t^2-4m<0$,
  $p\nmid t^2-4m$ and $p$ split in $K_q$;
\item an odd integer $n\geq1$ and an $\mathcal O_E$-module $\Lambda$, free of rank $n[E:\mathbf Q]$ over
  $\mathbf Z$, with a
  nondegenerate alternating form $\langle\,,\rangle:\Lambda\times\Lambda\to\mathbf Z$ such that
  $\langle bx,y\rangle=\langle x,c(b)y\rangle$ for $b\in\mathcal O_E$, perfect up to an integer prime
  to $p$, and positive for some complex structure on $\Lambda\otimes\mathbf R$;
\item one embedding $\tau_1:E\to\mathbf C$, with $\tau_1(\omega)$ a fixed root $\rho_c$ of
  $X^2-tX+m$.
\end{enumerate}
The associated group $G$ is the group of $E$-linear similitudes of $(\Lambda\otimes\mathbf
Q,\langle\,,\rangle)$; at the infinite place of $E$ induced by $\tau_1$ its unitary part has signature
$(1,n-1)$, and at the other infinite places signature $(0,n)$. The associated \emph{Shimura variety} of
level $N$, for $N\geq1$, is the moduli space, over $\mathcal O_E$-algebras in which the integers prime
to $p$ are invertible, of abelian schemes $A$ with an action of $\mathcal O_E$, a polarization whose
Rosati involution induces $c$, the Kottwitz condition that $\mathcal O_E$ acts on the part of
$\operatorname{Lie}A$ where $\omega$ acts by $\rho_c$ through $\tau_1$ with multiplicity one, and a
level-$N$ structure. Since $[E:\mathbf Q]>2$, the unitary group is definite at some infinite place,
hence anisotropic, and this Shimura variety is proper \cite{Kot92}; it is smooth of dimension $n-1$
over $E$.
\end{definition}
\leanref{Development/AbelianVarietyTheory/Definitions/Divisors/NotCompactOfMEqOne.lean}{GlobalTheory.UPel.HT.Datum}
\leanref{Development/AbelianVarietyTheory/Definitions/Divisors/IsCyclotomicExtensionQI.lean}{GlobalTheory.LKTrip.HTGroup}
\leanref{Development/AbelianVarietyTheory/Definitions/Divisors/HasDivPlace.lean}{GlobalTheory.LKGeo.Link}
\leanref{Development/AutomorphicArithmetic/Foundations/GaloisRepresentations/IntegralStructures/GeometricFlagBases/IntegOfMem.lean}{UpperTheory.Gal.Int.Geo.Sc}
\leanref{Development/AbelianVarietyTheory/Definitions/Divisors/CompactBIffCompact.lean}{GlobalTheory.UPel.HT.Datum.CompactB}

\begin{theorem}\label{gal-holomorphic-lefschetz}
Let $Z$ be a smooth projective complex variety of dimension $d$, and let $c_1,c_2:Y\to Z$ be finite
\'etale morphisms. Suppose that the set $F$ of points $y\in Y$ with $c_1(y)=c_2(y)$ is finite and that
at each $y\in F$ the endomorphism $A_y$ of the tangent space $T_{c_1(y)}Z$ induced by
$c_2\circ c_1^{-1}$ near $y$ satisfies $\det(1-A_y)\neq0$. Then for every $a\geq0$ the spaces
$H^b(Z,\Omega^a)$ are finite-dimensional and
\[
 \sum_{b=0}^d(-1)^b\operatorname{tr}\bigl(c_{1*}c_2^*\mid H^b(Z,\Omega^a)\bigr)
 =\sum_{y\in F}\frac{\operatorname{tr}\Lambda^aA_y}{\det(1-A_y)}.
\]
This is the holomorphic Lefschetz formula of Atiyah and Bott \cite{AB68} for the Dolbeault complexes of
$\Omega^a$, applied to a correspondence.
\end{theorem}
\leanref{Development/AutomorphicArithmetic/Foundations/GaloisRepresentations/LefschetzSigmaSums/HodgeTraceFromLefschetzAndGroupSide.lean}{UpperTheory.Gal.LocalSmoothExtension.HolomorphicLefschetzFormula}
\leanref{Development/AutomorphicArithmetic/Foundations/GaloisRepresentations/IntegralStructures/GeometricFlagBases/WeightOneGeometricRealization.lean}{UpperTheory.Gal.Int.Geo.weight_one_geometric_realization}

\begin{definition}\label{gal-plan}
In Situation~\ref{gal-seed-situation}, a \emph{plan} for $U$ consists of a prime $p$, a unitary PEL
datum of simple type for $E$ at $p$ (Definition~\ref{gal-pel-datum}) of rank $n=a$ with group $G$, a
level $N\geq3$ prime to $p$, a family $(g_j)_{j\in J}$ of Hecke operators of $G$ of level $N$ with values
$\theta:J\to\bar{\mathbf Q}_\ell$, an algebraic Hecke character $\kappa$ of $E$ of infinity type $m$, a
dominant weight $\lambda$ of $G$, and a test function $f$ in the Hecke algebra of level $N$ supported
on integral elements, such that:
\begin{enumerate}
\item \emph{(spectral condition)} the operators $g_j$ commute, and there is an irreducible algebraic
  representation $\xi$ of $G$ of weight $\lambda$ such that for every polynomial $s$ in the $g_j$ with
  $s(\iota\theta)\neq0$ there is a polynomial $t$ in the $g_j$ for which $(ts)(g)\ast f$ is elliptic,
  as a test function for the subgroup $K_N$ of level $N$, with $G$ presented as in
  Situation~\ref{gal-count-situation} with $F=E$ (Definition~\ref{gal-count-pairs}(6): it vanishes
  outside the regular elements, and its orbital sums at the pairs with a marking and reducible minimal
  polynomial and its orbital integrals at the elements of groups of definite forms whose minimal
  polynomial is not irreducible, separable and of degree $n$ vanish), and the trace of $(ts)(g)\ast f$
  on automorphic forms of weight $\xi$ on the definite unitary group with the same finite adelic
  points as $G$ is nonzero;
\item \emph{(congruences)} for all but finitely many $d$-split places $w\nmid\ell$ and every $B$ such
  that $U$ has Satake parameters $B$ at $w$, $\kappa$ is unramified at $w$, and the Hecke operators of
  $G$ at the place of $E$ under $w$ take values $\vartheta$ under $\theta$ whose Hecke polynomial of
  degree $a$ is $\vartheta_0\prod_{b\in B}(X-\iota^{-1}(q_w^{(a-1)/2}b\,\kappa(\varpi_w)^{-1}))$ with
  $\vartheta_0\neq0$.
\end{enumerate}
The plan has \emph{Hodge data at} a place $v$ over $\ell$ if $\kappa$ is unramified at $v$, the Hecke
operators at $v$ are good, and for every $\tau$ with $\iota^{-1}\circ\tau$ inducing $v$ there is an
automorphism $\sigma$ of $\mathbf C$ such that each $k-\lambda'_k+m_\tau$, for $\lambda'$ a weight of $U$
at $\tau$ and $k<a$, is a Hodge exponent of $\sigma\circ\lambda$; it is \emph{congruent at} $w$ if
(2) holds at $w$.
\end{definition}
\leanref{Development/AutomorphicArithmetic/Foundations/GaloisRepresentations/IntegralStructures/GeometricFlagBases/PackExOf.lean}{UpperTheory.Gal.Int.Geo.Plan}
\leanref{Development/AutomorphicArithmetic/Foundations/GaloisRepresentations/IntegralStructures/GeometricFlagBases/PackExOf.lean}{UpperTheory.Gal.Int.Geo.Plan.HTFor}
\leanref{Development/AutomorphicArithmetic/Foundations/GaloisRepresentations/IntegralStructures/GeometricFlagBases/PackExOf.lean}{UpperTheory.Gal.Int.Geo.Plan.CongAll}
\leanref{Development/AutomorphicArithmetic/Foundations/GaloisRepresentations/LefschetzSigmaSums/HeckeTraceSemi.lean}{UpperTheory.Gal.LocalSmoothExtension.Spec}

\begin{situation}\label{gal-powers-situation}
Let $p$ be a prime, let $E$ be a CM field, and take a unitary PEL datum of simple type for $E$ at $p$
(Definition~\ref{gal-pel-datum}), with $\omega$, $n$ and $\Lambda$ as there. Let $\mathcal O$ be a
discrete valuation ring of characteristic zero with maximal ideal $\mathfrak m$, and let
$\iota:\mathcal O_E\to\mathcal O$ be a ring homomorphism. Suppose that $p\in\mathfrak m$, that the
residue field of $\mathcal O$ has $q=p^f$ elements, and that the prime $w=\iota^{-1}(\mathfrak m)$ of
$E$ has norm $q$; write $k_w=\mathcal O_E/w$. Let $r\in\mathbf Z$ with $\omega-r\in w$, which exists
because $p$ splits in $\mathbf Q(\omega)$, and let
$\mathfrak u\subset w$ be the ideal of $\mathcal O_E$ generated by $p$ and $\omega-r$. Let $N\geq3$ be
an integer prime to $p$, and let $M$ be a smooth quasi-projective $\mathcal O$-scheme of relative
dimension $n-1$ that represents, on $\mathcal O$-schemes, the moduli problem of
Definition~\ref{gal-pel-datum} at level $N$, with the Kottwitz condition read along $\iota$:
$\mathcal O_E$ acts on the part of $\operatorname{Lie}A$ where $\omega$ acts by $\iota(\omega)$
through $\iota$, with multiplicity one. Let $\mathcal A\to M$ be the universal abelian scheme. Let
$\mathfrak a$ be an ideal of $\mathcal O_E$ with $\mathfrak a+w=\mathcal O_E$ that contains a power of
$\mathfrak u$. A \emph{level-$\mathfrak a$ structure} on an object $A$ of the moduli problem over an
$\mathcal O$-scheme $S$ is an $\mathcal O_E$-linear map $\eta:\Lambda\to A(S)$ that vanishes on
$\mathfrak a\Lambda$ and induces, at every geometric point of $S$, a bijection from
$\Lambda/\mathfrak a\Lambda$ onto the points killed by $\mathfrak a$. Let $M'$ be an $\mathcal
O$-scheme that represents the objects with a level-$\mathfrak a$ structure, with universal structure
$\eta'$, such that the map $\pi:M'\to M$ forgetting the structure is flat, and let
$\mathcal A'=\pi^*\mathcal A$.
\end{situation}
\leanref{Development/AutomorphicArithmetic/Foundations/GaloisRepresentations/AnnihilatorModels/FrobKill.lean}{UpperTheory.Gal.crystallineComparison.Mod.Std}
\leanref{Development/AbelianVarietyTheory/Definitions/Divisors/NotCompactOfMEqOne.lean}{GlobalTheory.UPel.HT.Base}
\leanref{Development/AbelianVarietyTheory/Definitions/Divisors/NotCompactOfMEqOne.lean}{GlobalTheory.UPel.HT.Model}
\leanref{Development/AbelianVarietyTheory/Definitions/Divisors/ExistsTjCorrDiagModel.lean}{GlobalTheory.UPel.HT.ES.uIdeal}
\leanref{Development/AbelianVarietyTheory/Definitions/Divisors/IsEtPlace.lean}{GlobalTheory.UPel.HT.EtLev.IsEtLev}
\leanref{Development/AbelianVarietyTheory/Definitions/Divisors/IsEtPlace.lean}{GlobalTheory.UPel.HT.EtLev.EtModel}

\begin{definition}\label{gal-powers-hecke}
In Situation~\ref{gal-powers-situation}, let $0\leq j\leq n$.
\begin{enumerate}
\item Let $S$ be an $\mathcal O$-scheme and let $A$, $A'$ be objects of the moduli problem over $S$,
  with polarizations $\lambda$, $\lambda'$. An \emph{isogeny of type $j$ at $w$} from $A$ to $A'$ is a
  homomorphism $\psi:A\to A'$ that commutes with $\mathcal O_E$, carries the level-$N$ structure of $A$
  to that of $A'$ and satisfies $\psi^\vee\lambda'\psi=q^j\lambda$, for which there is
  $\psi':A'\to A$ with $\psi'\psi=q^j$ and $\psi\psi'=q^j$, and such that at every geometric point of
  $S$ of characteristic zero the points of $\ker\psi$ killed by a power of $\mathfrak u$ are killed by
  $w$ and are $q^j$ in number. The \emph{$\mathfrak u$-kernel} of $\psi$ is the subgroup scheme of
  points of $\ker\psi$ killed by a power of $\mathfrak u$.
\item A \emph{Hecke correspondence of type $j$ at $w$} on $M$ consists of a scheme $Y_j$, two
  $\mathcal O$-morphisms $s_j,t_j:Y_j\to M$ with $s_j$ finite and \'etale over $M[1/p]$, and an isogeny
  $\psi_j$ of type $j$ at $w$ from $s_j^*\mathcal A$ to $t_j^*\mathcal A$, such that for every
  $\mathcal O$-scheme $S$ on which $p$ is invertible, every $x\in M(S)$ and every isogeny $\psi$ of type
  $j$ at $w$ from $x^*\mathcal A$, there is a unique $y\in Y_j(S)$ with $s_jy=x$ such that $\psi$ is a
  pullback of $\psi_j$ along $y$. For $y\in Y_j(S)$ and $a\in\mathcal A(S)$ over $s_jy$, write
  $\psi_y(a)\in\mathcal A(S)$ for the image of $a$ under $y^*\psi_j$, a point over $t_jy$.
\end{enumerate}
\end{definition}
\leanref{Development/AbelianVarietyTheory/Definitions/Divisors/ExistsTjCorrDiagModel.lean}{GlobalTheory.UPel.HT.ES.IsPIsog}
\leanref{Development/AbelianVarietyTheory/Definitions/Divisors/ExistsTjCorrDiagModel.lean}{GlobalTheory.UPel.HT.ES.IsTj}
\leanref{Development/AutomorphicArithmetic/Foundations/GaloisRepresentations/AnnihilatorModels/FrobKill.lean}{UpperTheory.Gal.crystallineComparison.Mod.UKer}
\leanref{Development/AutomorphicArithmetic/Foundations/GaloisRepresentations/AnnihilatorModels/FrobKill.lean}{UpperTheory.Gal.crystallineComparison.Mod.Univ}
\leanref{Development/AutomorphicArithmetic/Foundations/GaloisRepresentations/AnnihilatorModels/FrobKill.lean}{UpperTheory.Gal.crystallineComparison.Mod.TjD}
\leanref{Development/AutomorphicArithmetic/Foundations/GaloisRepresentations/AnnihilatorModels/FrobKill.lean}{UpperTheory.Gal.crystallineComparison.Mod.TjD.img}

\begin{lemma}\label{gal-powers-hecke-exist}
In Situation~\ref{gal-powers-situation}, let $0\leq j\leq n$.
\begin{enumerate}
\item There is a Hecke correspondence $(Y_j,s_j,t_j,\psi_j)$ of type $j$ at $w$ on $M$
  (Definition~\ref{gal-powers-hecke}) with $Y_j$ flat over $\mathcal O$.
\item Let $(Y_j,s_j,t_j,\psi_j)$ be such a correspondence, and let $Z_j=M'\times_MY_j$ with
  projections $s'_j:Z_j\to M'$ and $g_j:Z_j\to Y_j$. There is $t'_j:Z_j\to M'$ with
  $\pi t'_j=t_jg_j$ such that $\psi_j$ carries the pullback of $\eta'$ along $s'_j$ to the pullback of
  $\eta'$ along $t'_j$.
\item Let $(Y_j,s_j,t_j,\psi_j)$ be such a correspondence and $h:Y\to M$ a morphism. Then the
  projection $Y\times_MY_j\to Y$ is finite, \'etale over $Y[1/p]$, and has
  $|\operatorname{Gr}_j(k_w^n)|$ geometric points over every geometric point of $Y$ at which $p$ is
  invertible, where $\operatorname{Gr}_j(k_w^n)$ is the set of $j$-dimensional subspaces of $k_w^n$.
\end{enumerate}
\end{lemma}
\leanref{Development/AutomorphicArithmetic/Foundations/GaloisRepresentations/AnnihilatorModels/IdealSplit.lean}{UpperTheory.Gal.crystallineComparison.Mod.TjExF}
\leanref{Development/AutomorphicArithmetic/Foundations/GaloisRepresentations/AnnihilatorModels/IdealSplitHolds.lean}{UpperTheory.Gal.crystallineComparison.Mod.tjExF_holds}
\leanref{Development/AutomorphicArithmetic/Foundations/GaloisRepresentations/AnnihilatorModels/IdealSplit.lean}{UpperTheory.Gal.crystallineComparison.Mod.TjE}
\leanref{Development/AutomorphicArithmetic/Foundations/GaloisRepresentations/AnnihilatorModels/IdealSplit.lean}{UpperTheory.Gal.crystallineComparison.Mod.TjEEx}
\leanref{Development/AutomorphicArithmetic/Foundations/GaloisRepresentations/AnnihilatorModels/IdealSplitHolds.lean}{UpperTheory.Gal.crystallineComparison.Mod.tjEEx_holds}
\leanref{Development/AutomorphicArithmetic/Foundations/GaloisRepresentations/AnnihilatorModels/IdealSplit.lean}{UpperTheory.Gal.crystallineComparison.Mod.CorrBC}
\leanref{Development/AutomorphicArithmetic/Foundations/GaloisRepresentations/AnnihilatorModels/IdealSplitHolds.lean}{UpperTheory.Gal.crystallineComparison.Mod.corrBC_holds}

\begin{proof}
(1) Over $M[1/p]$, an isogeny of type $j$ at $w$ from $\mathcal A$ is determined up to isomorphism by
its $\mathfrak u$-kernel, an $\mathcal O_E$-stable finite \'etale subgroup scheme of $\mathcal A[w]$ of
order $q^j$, and every such subgroup scheme is the $\mathfrak u$-kernel of one. So these isogenies are
classified by a scheme $Y^0_j$ finite and \'etale over $M[1/p]$. Let $Y_j$ be the normalization of $M$
in $Y^0_j$. It is finite over $M$ and normal, it is flat over $\mathcal O$ because $Y^0_j$ is dense in
it, and $Y_j[1/p]=Y^0_j$. The target of the universal isogeny over $Y^0_j$ is isogenous to the source,
so it has good reduction along every discrete valuation ring of $Y_j$ \cite{ST68}; since $Y_j$ is
normal and $M$ is quasi-projective, its classifying map extends to a finite morphism $t_j:Y_j\to M$.
The isogeny extends over the normal scheme $Y_j$ \cite[Ch.~I, Prop.~2.7]{FC90}, and the extension is
of type $j$ at $w$, since the identities between homomorphisms extend from the dense open $Y^0_j$ and
the condition at points of characteristic zero only concerns $Y^0_j$. The universal property holds
because every point of $Y_j$ with values in a scheme on which $p$ is invertible factors through
$Y^0_j$.

(2) Points killed by $\mathfrak a$ are killed by a power of $\mathfrak u$, and the $\mathfrak u$-kernel
of $\psi_j$ is killed by $w$, as it is at the dense set of points of characteristic zero. Since
$\mathfrak a+w=\mathcal O_E$, $\psi_j$ is injective on the points killed by $\mathfrak a$, and it
carries a level-$\mathfrak a$ structure to a level-$\mathfrak a$ structure. As $M'$ represents objects
with such a structure, the target of $\psi_j$ over $Z_j$ with the transported structure defines $t'_j$.

(3) The projection is the base change of $s_j$, which is finite and \'etale over $M[1/p]$. A geometric
point of $Y$ at which $p$ is invertible has characteristic zero, and by the universal property the
geometric points over it correspond to the isogenies of type $j$ at $w$ from the fibre $A$ of
$\mathcal A$ there, hence to their $\mathfrak u$-kernels. These are the $j$-dimensional
$k_w$-subspaces of the points of $A$ killed by $w$, which form a $k_w$-vector space of dimension $n$
because the Tate module of $A$ is isomorphic to $\Lambda\otimes\hat{\mathbf Z}$.
\end{proof}
\leanref{Development/AutomorphicArithmetic/Foundations/GaloisRepresentations/AnnihilatorModels/FrobKill.lean}{UpperTheory.Gal.crystallineComparison.Mod.EsGen}
\leanref{Development/AutomorphicArithmetic/Foundations/GaloisRepresentations/AnnihilatorModels/FrobKill.lean}{UpperTheory.Gal.crystallineComparison.Mod.EsNorm}
\leanref{Development/AutomorphicArithmetic/Foundations/GaloisRepresentations/AnnihilatorModels/MEsNorm.lean}{UpperTheory.Gal.crystallineComparison.Mod.esNorm_of}
\leanref{Development/AutomorphicArithmetic/Foundations/EtaleClopenDescent/ExistsHeckeCorrespondenceOnIntegralShimuraModel.lean}{UpperTheory.Cb.exists_heckeCorrespondence_on_integral_shimuraModel}
\leanref{Development/AbelianVarietyTheory/Proofs/DivisorProofs/FlagModelValueIndependent.lean}{GlobalTheory.ExtT.exists_finite_extension}
\leanref{Development/AutomorphicArithmetic/Foundations/GaloisRepresentations/AnnihilatorModels/FrobKill.lean}{UpperTheory.Gal.crystallineComparison.Mod.QuotObj}
\leanref{Development/AutomorphicArithmetic/Foundations/GaloisRepresentations/AnnihilatorModels/IdealSplit.lean}{UpperTheory.Gal.crystallineComparison.Mod.UKerW}
\leanref{Development/AutomorphicArithmetic/Foundations/GaloisRepresentations/AnnihilatorModels/IdealSplit.lean}{UpperTheory.Gal.crystallineComparison.Mod.LevPush}

\begin{lemma}\label{gal-powers-subspace-identity}
Let $k$ be a finite field with $q$ elements, $n\geq1$, $H$ a $k$-vector space of dimension $n$ and
$C\subset H$ a line, and write $\bar D$ for the image in $H/C$ of a subspace $D\subset H$. For
$0\leq j\leq n$ let $\operatorname{Gr}_j(H)$ be the set of $j$-dimensional subspaces of $H$. Let $P$
be a set, $F:P\to P$ a map, and $v_j:\operatorname{Gr}_j(H)\to P$, $0\leq j\leq n$, maps such that:
\begin{enumerate}
\item if $D,D'\in\operatorname{Gr}_j(H)$ do not contain $C$ and $\bar D=\bar D'$, then
  $v_j(D)=v_j(D')$;
\item if $D\in\operatorname{Gr}_{j+1}(H)$ contains $C$, $D'\in\operatorname{Gr}_j(H)$ does not contain
  $C$, and $\bar D=\bar D'$, then $v_{j+1}(D)=F(v_j(D'))$.
\end{enumerate}
Then the two multisets of elements of $P$
\[
 \sum_{j\text{ even}}q^{j(j-1)/2}\bigl\{F^{n-j}(v_j(D))\bigr\}_{D\in\operatorname{Gr}_j(H)},\qquad
 \sum_{j\text{ odd}}q^{j(j-1)/2}\bigl\{F^{n-j}(v_j(D))\bigr\}_{D\in\operatorname{Gr}_j(H)}
\]
are equal.
\end{lemma}
\leanref{Development/AutomorphicArithmetic/Foundations/GaloisRepresentations/AnnihilatorModels/GenericMulOfProd.lean}{UpperTheory.Gal.crystallineComparison.OrdRead}
\leanref{Development/AutomorphicArithmetic/Foundations/GaloisRepresentations/AnnihilatorModels/GenericMulOfProd.lean}{UpperTheory.Gal.crystallineComparison.AltId}
\leanref{Development/AutomorphicArithmetic/Foundations/GaloisRepresentations/AnnihilatorModels/GenericMulOfProd.lean}{UpperTheory.Gal.crystallineComparison.OrdCount}
\leanref{Development/AutomorphicArithmetic/Foundations/GaloisRepresentations/AnnihilatorModels/GenericMulOfProdHolds.lean}{UpperTheory.Gal.crystallineComparison.ordCount_holds}
\leanref{Development/AutomorphicArithmetic/Foundations/GaloisRepresentations/AnnihilatorModels/OrdReadOf.lean}{UpperTheory.Gal.crystallineComparison.Mod.OrdReadOf}
\leanref{Development/AutomorphicArithmetic/Foundations/GaloisRepresentations/AnnihilatorModels/IdealSplitHolds.lean}{UpperTheory.Gal.crystallineComparison.Mod.ordReadOf_holds}

\begin{proof}
For a subspace $B\subset H/C$ of dimension $m\leq n-1$, its inverse image in $H$ has dimension $m+1$
and contains $C$; choose a complement $D_B$ of $C$ in it, so that $D_B\in\operatorname{Gr}_m(H)$ does
not contain $C$ and $\bar D_B=B$, and put $A(B)=v_m(D_B)$. A subspace $D$ not containing $C$ meets
$C$ in $0$, so $\dim\bar D=\dim D$, and (1) gives $v_j(D)=A(\bar D)$ for every
$D\in\operatorname{Gr}_j(H)$ not containing $C$. A subspace $D\in\operatorname{Gr}_{j+1}(H)$
containing $C$ has $\dim\bar D=j$ and $D=D_{\bar D}+C$, so (2) gives $v_{j+1}(D)=F(A(\bar D))$.

The map $D\mapsto\bar D$ is a bijection from the subspaces of $H$ containing $C$ onto the subspaces of
$H/C$, lowering the dimension by one. For $B\subset H/C$ of dimension $m$, the subspaces $D$ not
containing $C$ with $\bar D=B$ are the complements of $C$ in the inverse image of $B$, and there are
$q^m$ of them. Let $A_m$ be the multiset $\{A(B)\}$ over the $m$-dimensional subspaces $B$ of $H/C$,
with $A_{-1}=A_n=0$. Then the multiset $\{v_j(D)\}_{D\in\operatorname{Gr}_j(H)}$ is
$F(A_{j-1})+q^jA_j$, the image of $A_{j-1}$ under $F$ plus $q^j$ copies of $A_j$.

Put $S_m=F^{n-m}(A_m)$. The multiset $\{F^{n-j}(v_j(D))\}_D$ is $S_{j-1}+q^jS_j$. So the first
multiset of the statement is $\sum_{j\text{ even}}q^{j(j-1)/2}S_{j-1}+\sum_{j\text{
even}}q^{j(j+1)/2}S_j$ and the second is $\sum_{j\text{ odd}}q^{j(j-1)/2}S_{j-1}+\sum_{j\text{
odd}}q^{j(j+1)/2}S_j$. Substituting $m=j-1$, the first sum on each side is
$\sum_mq^{m(m+1)/2}S_m$ over the $m$ of the other parity, which is the second sum on the other side.
\end{proof}
\leanref{Development/AutomorphicArithmetic/Foundations/GaloisRepresentations/AnnihilatorModels/SameNormBasic.lean}{UpperTheory.Gal.crystallineComparison.CountIn}
\leanref{Development/AutomorphicArithmetic/Foundations/GaloisRepresentations/AnnihilatorModels/SameNormBasic.lean}{UpperTheory.Gal.crystallineComparison.CountOut}
\leanref{Development/AutomorphicArithmetic/Foundations/GaloisRepresentations/AnnihilatorModels/SameNormBasic.lean}{UpperTheory.Gal.crystallineComparison.OrdDecomp}
\leanref{Development/AutomorphicArithmetic/Foundations/GaloisRepresentations/AnnihilatorModels/SameNormBasic.lean}{UpperTheory.Gal.crystallineComparison.Telescope}

\begin{theorem}\label{gal-powers-generic-ordinary}
In Situation~\ref{gal-powers-situation}, let $\kappa$ be an algebraically closed field and
$x:\operatorname{Spec}\kappa\to M$ a morphism whose image is a generic point of the special fibre of
$M$ over $\mathcal O$. Then $q^{n-1}$ points of $\mathcal A_x(\kappa)$ are killed by $w$; such a
geometric point $x$ is called \emph{$w$-ordinary}. This follows from the stratification of the special
fibre by the $\mathcal O_{E_w}$-height $h$ of the \'etale quotient of the $p$-divisible group
$\mathcal A[w^\infty]$ at each point, whose stratum of height $h$ is smooth of pure dimension $h$
(Harris and Taylor
\cite[Cor.~III.4.4]{HT01}): the special fibre has pure dimension $n-1$, so its generic points lie in the
stratum of height $n-1$, where the points killed by $w$ are those of the \'etale quotient. Wedhorn
proves the density of the $\mu$-ordinary locus for Shimura varieties of PEL type in general
\cite{Wed99}.
\end{theorem}
\leanref{Development/AbelianVarietyTheory/Definitions/Divisors/ExistsTjCorrDiagModel.lean}{GlobalTheory.UPel.HT.ES.IsWOrd}
\leanref{Development/AutomorphicArithmetic/Foundations/GaloisRepresentations/AnnihilatorModels/FrobKill.lean}{UpperTheory.Gal.crystallineComparison.Mod.OrdGen}
\leanref{Development/AutomorphicArithmetic/Foundations/GaloisRepresentations/RadicalStalks/OrdGen.lean}{UpperTheory.Gal.Ord.ordGen}

\begin{lemma}\label{gal-powers-special-isogenies}
In Situation~\ref{gal-powers-situation}, let $(Y_j,s_j,t_j,\psi_j)$, $0\leq j\leq n$, be Hecke
correspondences of type $j$ at $w$ on $M$ (Definition~\ref{gal-powers-hecke}). Let $\kappa$ be an
algebraically closed field of characteristic $p$ with $q$-power Frobenius $\phi$, let
$x:\operatorname{Spec}\kappa\to M$ lie over the closed point of $\operatorname{Spec}\mathcal O$, and
let $A=\mathcal A_x$. For a $\kappa$-point $y$ of $Y_j$ with $s_jy=x$, let $K_y\subset A$ be the
$\mathfrak u$-kernel of $y^*\psi_j$.
\begin{enumerate}
\item Let $y,y'$ be $\kappa$-points of $Y_j$ over $x$ such that $K_y$ and $K_{y'}$ have the same order
  $d$ and the same $\kappa$-points, and $K_y$ has $d$ $\kappa$-points. Then $t_jy=t_jy'$, and
  $\psi_y(a)=\psi_{y'}(a)$ for every $a\in A(\kappa)$.
\item Let $y$ be a $\kappa$-point of $Y_{j+1}$ and $y'$ a $\kappa$-point of $Y_j$, both over $x$, such
  that $K_y$ has order $q^{j+1}$, $K_{y'}$ has order $q^j$ and $q^j$ $\kappa$-points, and $K_y$ and
  $K_{y'}$ have the same $\kappa$-points. Then $t_{j+1}y=(t_jy')\circ\operatorname{Spec}\phi$, and
  $\psi_y(a)=\psi_{y'}(a)\circ\operatorname{Spec}\phi$ for every $a\in A(\kappa)$.
\end{enumerate}
\end{lemma}
\leanref{Development/AutomorphicArithmetic/Foundations/GaloisRepresentations/AnnihilatorModels/PulledRead.lean}{UpperTheory.Gal.crystallineComparison.Mod.CaseE}
\leanref{Development/AutomorphicArithmetic/Foundations/GaloisRepresentations/AnnihilatorModels/IdealSplitHolds.lean}{UpperTheory.Gal.crystallineComparison.Mod.caseE_holds}
\leanref{Development/AutomorphicArithmetic/Foundations/GaloisRepresentations/AnnihilatorModels/PulledRead.lean}{UpperTheory.Gal.crystallineComparison.Mod.CaseW}
\leanref{Development/AutomorphicArithmetic/Foundations/GaloisRepresentations/AnnihilatorModels/IdealSplitHolds.lean}{UpperTheory.Gal.crystallineComparison.Mod.caseW_holds}

\begin{proof}
The $\mathfrak u$-kernel $K_y$ is a finite subgroup scheme of $A$. The argument uses four facts about
homomorphisms $\psi$ from $A$ that satisfy the conditions of Definition~\ref{gal-powers-hecke}(1) with
a power $c$ of $p$ in place of $q^j$; over $\kappa$ the condition at points of characteristic zero is
empty.
\begin{enumerate}
\item[(a)] A homomorphism kills a finite subgroup scheme $K$ of $A$ as soon as it kills the
  $\kappa$-points of $K$ and its \emph{infinitesimal points}, the points with values in local
  $\kappa$-algebras that reduce to the origin.
\item[(b)] If $K$ has as many $\kappa$-points as its order, its only infinitesimal point is the origin.
\item[(c)] Two such homomorphisms with the same $c$ and the same $\mathfrak u$-kernel have the same
  kernel: a point of the kernel is the sum of a point killed by a power of $\mathfrak u$ and a point
  killed by a power of $\bar{\mathfrak u}$, and the polarizations, whose Rosati involution induces
  complex conjugation on $\mathcal O_E$, show that the second part is killed by both.
\item[(d)] Two such homomorphisms with the same $c$ and the same kernel have targets that are
  isomorphic compatibly with the homomorphisms, the $\mathcal O_E$-actions, the polarizations and the
  level structures. Since $M$ represents the moduli problem and $N\geq3$, isomorphic objects define the
  same point of $M$, and the isomorphism is the identity of the fibre of $\mathcal A$ there.
\end{enumerate}

(1) By (b), the only infinitesimal point of $K_y$ is the origin, and the $\kappa$-points of $K_y$ are
those of $K_{y'}$, which $y'^*\psi_j$ kills; by (a), $K_y\subset K_{y'}$. Since $K_{y'}$ also has $d$
$\kappa$-points, the same argument gives $K_{y'}\subset K_y$. By (c) and (d), $t_jy=t_jy'$ and
$\psi_y(a)=\psi_{y'}(a)$.

(2) Let $A'$ be the target of $y'^*\psi_j$ and $x'=t_jy'$. Since $\phi$ fixes the residue field of
$\mathcal O$, the point $x'\circ\operatorname{Spec}\phi$ of $M$ also lies over the closed point, and
the relative $q$-power Frobenius $F$ from $A'$ to the fibre of $\mathcal A$ at
$x'\circ\operatorname{Spec}\phi$ satisfies the conditions of (c) and (d) with $c=q$. Its kernel
consists of the points killed by the $q$-power Frobenius, and $F(b)=b\circ\operatorname{Spec}\phi$ for
$b\in A'(\kappa)$. Then $F\circ y'^*\psi_j$ satisfies these conditions with $c=q^{j+1}$; let $K''$ be
its $\mathfrak u$-kernel. No nonzero $\kappa$-point is killed by the Frobenius, so the $\kappa$-points of
$K''$ are those of $K_{y'}$, which are those of $K_y$. Both $K_y$ and $K''$ lie in the
$\mathfrak u$-part of $A[p^{f(j+1)}]$, whose tangent space at the origin is the part of
$\operatorname{Lie}A$ on which $\omega$ acts by $\iota(\omega)$; by the Kottwitz condition it has
dimension one. The connected component of $K_y$ has order $q$, so it is the subscheme of the
connected component of this $\mathfrak u$-part cut out by the $q$-power Frobenius, and the
infinitesimal points of $K_y$ are the infinitesimal points of the $\mathfrak u$-part killed by the
$q$-power Frobenius. Since $K_{y'}$ is \'etale, an infinitesimal point $g$ of the $\mathfrak u$-part
of $A$ is killed by the $q$-power Frobenius if and only if $\psi_{y'}(g)$ is, that is, if and only if
$g\in K''$. So $K_y$ and $K''$ have the same $\kappa$-points and the same infinitesimal points, and (a)
gives $K_y=K''$. By (c) and (d) applied to $y^*\psi_j$ and $F\circ y'^*\psi_j$, we get
$t_{j+1}y=x'\circ\operatorname{Spec}\phi$ and $\psi_y(a)=F(\psi_{y'}(a))=\psi_{y'}(a)\circ
\operatorname{Spec}\phi$.
\end{proof}
\leanref{Development/AutomorphicArithmetic/Foundations/GaloisRepresentations/AnnihilatorModels/PairCalc.lean}{UpperTheory.Gal.crystallineComparison.Mod.KerU}
\leanref{Development/AutomorphicArithmetic/Foundations/GaloisRepresentations/AnnihilatorModels/OrdReadOf.lean}{UpperTheory.Gal.crystallineComparison.Mod.KillLocal}
\leanref{Development/AutomorphicArithmetic/Foundations/GaloisRepresentations/AnnihilatorModels/OrdReadOf.lean}{UpperTheory.Gal.crystallineComparison.Mod.InfE}
\leanref{Development/AutomorphicArithmetic/Foundations/GaloisRepresentations/AnnihilatorModels/PairCalc.lean}{UpperTheory.Gal.crystallineComparison.Mod.KerFromU}
\leanref{Development/AutomorphicArithmetic/Foundations/GaloisRepresentations/AnnihilatorModels/PairCalc.lean}{UpperTheory.Gal.crystallineComparison.Mod.IsoOfKer}
\leanref{Development/AutomorphicArithmetic/Foundations/GaloisRepresentations/AnnihilatorModels/PairCalc.lean}{UpperTheory.Gal.crystallineComparison.Mod.ImgEq}
\leanref{Development/AutomorphicArithmetic/Foundations/GaloisRepresentations/AnnihilatorModels/PairCalc.lean}{UpperTheory.Gal.crystallineComparison.Mod.SameE}
\leanref{Development/AutomorphicArithmetic/Foundations/GaloisRepresentations/AnnihilatorModels/PairCalc.lean}{UpperTheory.Gal.crystallineComparison.Mod.FrobIsog}
\leanref{Development/AutomorphicArithmetic/Foundations/GaloisRepresentations/AnnihilatorModels/PairCalc.lean}{UpperTheory.Gal.crystallineComparison.Mod.IsogComp}
\leanref{Development/AutomorphicArithmetic/Foundations/GaloisRepresentations/AnnihilatorModels/PairCalc.lean}{UpperTheory.Gal.crystallineComparison.Mod.TanOne}
\leanref{Development/AutomorphicArithmetic/Foundations/GaloisRepresentations/AnnihilatorModels/OrdReadOf.lean}{UpperTheory.Gal.crystallineComparison.Mod.InfW}
\leanref{Development/AutomorphicArithmetic/Foundations/GaloisRepresentations/AnnihilatorModels/PairCalc.lean}{UpperTheory.Gal.crystallineComparison.Mod.FrobPass}
\leanref{Development/AutomorphicArithmetic/Foundations/GaloisRepresentations/AnnihilatorModels/PairCalc.lean}{UpperTheory.Gal.crystallineComparison.Mod.SameW}

\begin{proposition}\label{gal-powers-reading}
In Situation~\ref{gal-powers-situation}, let $(Y_j,s_j,t_j,\psi_j)$, $0\leq j\leq n$, be Hecke
correspondences of type $j$ at $w$ on $M$ with $Y_j$ flat over $\mathcal O$
(Definition~\ref{gal-powers-hecke}). Let $V$ be a valuation ring with algebraically closed fraction
field in which $p\neq0$; its residue field $\kappa_V$ is algebraically closed, and let $\phi$ be the
$q$-power Frobenius of $\kappa_V$. Let $x:\operatorname{Spec}V\to M$ map the closed point over the
closed point of $\operatorname{Spec}\mathcal O$, with $w$-ordinary reduction $\bar x$
(Theorem~\ref{gal-powers-generic-ordinary}), and write $\bar u$ for the reduction of a $V$-point $u$.
Let $H\subset\mathcal A_x(V)$ be the subgroup of points killed by $w$.
\begin{enumerate}
\item $H$ is a $k_w$-vector space of dimension $n$, and the $h\in H$ with $\bar h=0$ form a line $C$.
\item For each $j$, $z\mapsto e_j(z)=H\cap\ker(z^*\psi_j)$ is a bijection from the $V$-points $z$ of
  $Y_j$ with $s_jz=x$ onto $\operatorname{Gr}_j(H)$.
\item Let $z,z'$ be $V$-points of $Y_j$ over $x$ such that $e_j(z)$ and $e_j(z')$ do not contain $C$
  and $e_j(z)+C=e_j(z')+C$. Then $t_j\bar z=t_j\bar z'$, and $\psi_{\bar z}(a)=\psi_{\bar z'}(a)$ for
  every $a\in\mathcal A_{\bar x}(\kappa_V)$.
\item Let $z$ be a $V$-point of $Y_{j+1}$ and $z'$ a $V$-point of $Y_j$, both over $x$, such that
  $e_{j+1}(z)$ contains $C$, $e_j(z')$ does not, and $e_{j+1}(z)=e_j(z')+C$. Then
  $t_{j+1}\bar z=(t_j\bar z')\circ\operatorname{Spec}\phi$, and $\psi_{\bar z}(a)=\psi_{\bar
  z'}(a)\circ\operatorname{Spec}\phi$ for every $a\in\mathcal A_{\bar x}(\kappa_V)$.
\end{enumerate}
\end{proposition}
\leanref{Development/AutomorphicArithmetic/Foundations/GaloisRepresentations/AnnihilatorModels/FrobKill.lean}{UpperTheory.Gal.crystallineComparison.Mod.EW}
\leanref{Development/AutomorphicArithmetic/Foundations/GaloisRepresentations/AnnihilatorModels/IdealSplitHolds.lean}{UpperTheory.Gal.crystallineComparison.Mod.eW_holds}
\leanref{Development/AutomorphicArithmetic/Foundations/GaloisRepresentations/AnnihilatorModels/FrobKill.lean}{UpperTheory.Gal.crystallineComparison.Mod.WMod}
\leanref{Development/AutomorphicArithmetic/Foundations/GaloisRepresentations/AnnihilatorModels/PulledRead.lean}{UpperTheory.Gal.crystallineComparison.Mod.KerAt}
\leanref{Development/AutomorphicArithmetic/Foundations/GaloisRepresentations/AnnihilatorModels/OrdReadOf.lean}{UpperTheory.Gal.crystallineComparison.Mod.ResClosed}

\begin{proof}
(1) $\mathcal A_x$ is an abelian scheme over $V$, so its $V$-points are the points of its generic
fibre, over an algebraically closed field of characteristic zero; there the points killed by $w$ form a
$k_w$-vector space of dimension $n$, as in the proof of Lemma~\ref{gal-powers-hecke-exist}(3). The
subgroup scheme $\mathcal A[w]$ is finite and flat over $M$, so every $\kappa_V$-point of
$\mathcal A_{\bar x}$ killed by $w$ is the reduction of a $V$-point. The reduction map
$H\to\mathcal A_{\bar x}(\kappa_V)$ is $k_w$-linear with image of $q^{n-1}$ elements, since $\bar x$
is $w$-ordinary, so its kernel $C$ is a line.

(2) Since $Y_j$ is finite over $M$, the $V$-points of $Y_j$ over $x$ are its points over the generic
point of $\operatorname{Spec}V$, where $p$ is invertible. As in the proof of
Lemma~\ref{gal-powers-hecke-exist}(3), they correspond to the $\mathfrak u$-kernels of the isogenies of
type $j$ at $w$ from the generic fibre of $\mathcal A_x$, which are the $j$-dimensional subspaces of
$H$.

(3) and (4). Let $z$ be a $V$-point of $Y_j$ over $x$ and $D=e_j(z)$. The reduction of $z^*\psi_j$ is
$\bar z^*\psi_j$; let $K_{\bar z}$ be its $\mathfrak u$-kernel. The $\mathfrak u$-kernel of $z^*\psi_j$
is finite and flat over $V$ of order $q^j$, with generic fibre $D$; so $K_{\bar z}$ has order $q^j$,
and its $\kappa_V$-points are the reductions of the points of $D$. Since the reduction map has kernel
$C$, these reductions form a set of $q^{\dim((D+C)/C)}$ points that depends only on $D+C$. In (3),
$K_{\bar z}$ and $K_{\bar z'}$ therefore have order $q^j$ and the same $\kappa_V$-points, $q^j$ of them
because $e_j(z)$ meets $C$ in $0$; Lemma~\ref{gal-powers-special-isogenies}(1) gives the claim. In
(4), $K_{\bar z}$ has order $q^{j+1}$, $K_{\bar z'}$ has order $q^j$ and $q^j$ points, and they have the
same points; Lemma~\ref{gal-powers-special-isogenies}(2) gives the claim.
\end{proof}
\leanref{Development/AutomorphicArithmetic/Foundations/GaloisRepresentations/AnnihilatorModels/PulledRead.lean}{UpperTheory.Gal.crystallineComparison.Mod.RankN}
\leanref{Development/AutomorphicArithmetic/Foundations/GaloisRepresentations/AnnihilatorModels/PulledRead.lean}{UpperTheory.Gal.crystallineComparison.Mod.LineC}
\leanref{Development/AutomorphicArithmetic/Foundations/GaloisRepresentations/AnnihilatorModels/PulledRead.lean}{UpperTheory.Gal.crystallineComparison.Mod.SecEquiv}
\leanref{Development/AutomorphicArithmetic/Foundations/GaloisRepresentations/AnnihilatorModels/FrobKill.lean}{UpperTheory.Gal.crystallineComparison.Mod.Pulled}
\leanref{Development/AutomorphicArithmetic/Foundations/GaloisRepresentations/AnnihilatorModels/PulledRead.lean}{UpperTheory.Gal.crystallineComparison.Mod.PulledRead}
\leanref{Development/AutomorphicArithmetic/Foundations/GaloisRepresentations/AnnihilatorModels/PulledRead.lean}{UpperTheory.Gal.crystallineComparison.Mod.KerRead}
\leanref{Development/AutomorphicArithmetic/Foundations/GaloisRepresentations/AnnihilatorModels/OrdReadOf.lean}{UpperTheory.Gal.crystallineComparison.Mod.RedPts}

\begin{proposition}\label{gal-powers-eichler-shimura}
In Situation~\ref{gal-powers-situation}, let $k\geq0$, let $\mathcal Y$ be the $k$-fold fibre power of
$\mathcal A'$ over $M'$, and let $\varpi$ be a uniformizer of $\mathcal O$. For $0\leq j\leq n$ let
$(Y_j,s_j,t_j,\psi_j)$ be a Hecke correspondence of type $j$ at $w$ on $M$ with $Y_j$ flat over
$\mathcal O$ (Lemma~\ref{gal-powers-hecke-exist}(1)), and let $Z_j$, $s'_j$, $g_j$ and $t'_j$ be as in
Lemma~\ref{gal-powers-hecke-exist}(2). Let $W_j$ be the $k$-fold fibre power over $Z_j$ of the pullback
of $\mathcal A'$ along $s'_j$, let $p_j:W_j\to\mathcal Y$ be the projection over $s'_j$, and let
$r_j:W_j\to\mathcal Y$ be the map $(a_1,\ldots,a_k)\mapsto(\psi_y(a_1),\ldots,\psi_y(a_k))$ over
$t'_j$, where $y$ is the image in $Y_j$. Then:
\begin{enumerate}
\item each $(W_j,p_j,r_j)$ is a finite correspondence of $\mathcal Y$
  (Definition~\ref{gal-tower-reduction}(2)) with $p_j$ and $r_j$ over the same point of
  $\operatorname{Spec}\mathcal O$, and $p_j$ has $|\operatorname{Gr}_j(k_w^n)|$ geometric points over
  each geometric point of $\mathcal Y$ at which $p$ is invertible;
\item $(W_j,p_j,r_j)_{0\leq j\leq n}$ satisfy the Eichler--Shimura relation of degree $n$ for $q$ at
  the generic points of the special fibre $\varpi=0$ of $\mathcal Y$
  (Definition~\ref{gal-tower-reduction}(3)).
\end{enumerate}
\end{proposition}
\leanref{Development/AutomorphicArithmetic/Foundations/GaloisRepresentations/AnnihilatorModels/CongKSOf.lean}{UpperTheory.Gal.crystallineComparison.Mod.CongKS}
\leanref{Development/AutomorphicArithmetic/Foundations/GaloisRepresentations/AnnihilatorModels/IdealSplitHolds.lean}{UpperTheory.Gal.crystallineComparison.Mod.congKS_holds}
\leanref{Development/AutomorphicArithmetic/Foundations/GaloisRepresentations/AnnihilatorModels/QuotObjLink.lean}{UpperTheory.Gal.crystallineComparison.Mod.congKS_link}
\leanref{Development/AutomorphicArithmetic/Foundations/GaloisRepresentations/AnnihilatorModels/IdealSplit.lean}{UpperTheory.Gal.crystallineComparison.Mod.OrdKSE}
\leanref{Development/AutomorphicArithmetic/Foundations/GaloisRepresentations/AnnihilatorModels/IdealSplitHolds.lean}{UpperTheory.Gal.crystallineComparison.Mod.ordKSE_holds}
\leanref{Development/AutomorphicArithmetic/Foundations/GaloisRepresentations/AnnihilatorModels/IdealSplit.lean}{UpperTheory.Gal.crystallineComparison.Mod.KSOKE}
\leanref{Development/AutomorphicArithmetic/Foundations/GaloisRepresentations/AnnihilatorModels/IdealSplit.lean}{UpperTheory.Gal.crystallineComparison.Mod.KSOKE.corr}
\leanref{Development/AutomorphicArithmetic/Foundations/GaloisRepresentations/AnnihilatorModels/GenericMulOfProd.lean}{UpperTheory.Gal.crystallineComparison.PtId}

\begin{proof}
Let $h:\mathcal Y\to M$ be the projection to $M'$ followed by $\pi$. It is flat, since $\mathcal Y$ is
smooth over $M'$ and $\pi$ is flat.

(1) The square formed by $p_j$, the map $W_j\to Y_j$, $h$ and $s_j$ is cartesian, so
Lemma~\ref{gal-powers-hecke-exist}(3) applies to $p_j$. The maps $p_j$ and $r_j$ lie over the same
point of $\operatorname{Spec}\mathcal O$ because $s_j$ and $t_j$ do.

(2) A point $u$ of $\mathcal Y$ with values in a scheme $X$ is determined by $h(u)\in M(X)$, by its
coordinates $a_1(u),\ldots,a_k(u)\in\mathcal A(X)$ and by the points $\eta'(\lambda)(u)\in\mathcal A(X)$
for $\lambda\in\Lambda$. Indeed, the level-$\mathfrak a$ structure determines the point of $M'$ over
$h(u)$, because $M'$ represents objects with such a structure, and the coordinates determine the point
of the fibre power over it. For a point $z$ of $W_j$ with image $y$ in $Y_j$, these data for $r_j(z)$
are $h(r_jz)=t_jy$, $a_i(r_jz)=\psi_y(a_i(p_jz))$ and $\eta'(\lambda)(r_jz)=\psi_y(\eta'(\lambda)(p_jz))$,
the last by Lemma~\ref{gal-powers-hecke-exist}(2).

Let $V$ be a valuation ring with algebraically closed fraction field in which $p\neq0$, let
$u:\operatorname{Spec}V\to\mathcal Y$ map the closed point to a generic point of the special fibre, and
let $\phi$ be the $q$-power Frobenius of the residue field $\kappa_V$. Since $h$ is flat, generizations
lift along $h$, so $x=h\circ u$ maps the closed point to a generic point of the special fibre of $M$,
and $\bar x$ is $w$-ordinary by Theorem~\ref{gal-powers-generic-ordinary}. The square above is
cartesian, so the $V$-points of $W_j$ over $u$ correspond to the $V$-points of $Y_j$ over $x$. Let $H$,
$C$ and $e_j$ be as in Proposition~\ref{gal-powers-reading}, and for $D\in\operatorname{Gr}_j(H)$ let
$v_j(D)\in\mathcal Y(\kappa_V)$ be the reduction of $r_j(z)$, where $z$ is the $V$-point of $W_j$ over
$u$ whose image $y$ in $Y_j$ has $e_j(y)=D$.

By the description of points above, $v_j(D)$ is determined by $t_j\bar y$ and by the points
$\psi_{\bar y}(b)$ for $b$ among the reductions of the $a_i(u)$ and the $\eta'(\lambda)(u)$. If
$D,D'\in\operatorname{Gr}_j(H)$ do not contain $C$ and $D+C=D'+C$, then
Proposition~\ref{gal-powers-reading}(3) shows that these data agree for $D$ and $D'$, so
$v_j(D)=v_j(D')$. If $D\in\operatorname{Gr}_{j+1}(H)$ contains $C$, $D'\in\operatorname{Gr}_j(H)$ does
not and $D=D'+C$, then Proposition~\ref{gal-powers-reading}(4) shows in the same way that
$v_{j+1}(D)=v_j(D')\circ\operatorname{Spec}\phi$. Lemma~\ref{gal-powers-subspace-identity}, with
$P=\mathcal Y(\kappa_V)$ and $F$ the composition with $\operatorname{Spec}\phi$, gives the equality of
multisets that is the Eichler--Shimura relation at $u$.
\end{proof}
\leanref{Development/AutomorphicArithmetic/Foundations/GaloisRepresentations/AnnihilatorModels/IdealSplit.lean}{UpperTheory.Gal.crystallineComparison.Mod.OrdEmb}
\leanref{Development/AutomorphicArithmetic/Foundations/GaloisRepresentations/AnnihilatorModels/IdealSplitHolds.lean}{UpperTheory.Gal.crystallineComparison.Mod.ordEmb_holds}
\leanref{Development/AutomorphicArithmetic/Foundations/GaloisRepresentations/AnnihilatorModels/IdealSplit.lean}{UpperTheory.Gal.crystallineComparison.Mod.EtInj}
\leanref{Development/AutomorphicArithmetic/Foundations/GaloisRepresentations/AnnihilatorModels/PulledRead.lean}{UpperTheory.Gal.crystallineComparison.Mod.PowExt}
\leanref{Development/AutomorphicArithmetic/Foundations/GaloisRepresentations/AnnihilatorModels/OrdReadOf.lean}{UpperTheory.Gal.crystallineComparison.Mod.GenDown}
\leanref{Development/AutomorphicArithmetic/Foundations/GaloisRepresentations/AnnihilatorModels/GenericMulOfProd.lean}{UpperTheory.Gal.crystallineComparison.OrdPt}
\leanref{Development/AutomorphicArithmetic/Foundations/GaloisRepresentations/AnnihilatorModels/GenericMulOfProd.lean}{UpperTheory.Gal.crystallineComparison.PtIdOfOrd}
\leanref{Development/AutomorphicArithmetic/Foundations/GaloisRepresentations/AnnihilatorModels/GenericMulOfProdHolds.lean}{UpperTheory.Gal.crystallineComparison.ptIdOfOrd_holds}

\subsection{Forms, pairs and orbital sums}\label{subsection-gal-forms-pairs}

\begin{situation}\label{gal-count-situation}
Let $p$ be a prime and $F$ a CM field with complex conjugation $c$ that contains an imaginary
quadratic field $K_0$ in which $p$ splits, and fix a place $u$ of $K_0$ over $p$. Let $n\geq1$ and
let $C$ be a central simple $F$-algebra of dimension $n^2$ with an involution $x\mapsto x^*$ that
restricts to $c$ on $F$. For each infinite place $w$ of $F$, with complex embedding also written
$w$, let $\psi_w:C\to M_n(\mathbf C)$ be a ring homomorphism extending $w$, and $J_w$ an invertible
Hermitian matrix with $\psi_w(x^*)=J_w^{-1}\psi_w(x)^\dagger J_w$, where $\dagger$ is the conjugate
transpose. Suppose that $J_v$ has signature $(1,n-1)$ for exactly one infinite place $v$, and $J_w$
has signature $(0,n)$ for the others. Let $w_1,\ldots,w_m$ be the places of $F$ over $u$, and
$\rho_i:C\to M_n(F_{w_i})$ ring homomorphisms extending $F\to F_{w_i}$. Fix ring homomorphisms
$j_p:\bar{\mathbf Q}\to\bar{\mathbf Q}_p$ and $j_{\mathbf C}:\bar{\mathbf Q}\to\mathbf C$; they are not used in this
section. For a commutative
$\mathbf Q$-algebra $R$, let $G(R)$ be the group of pairs $(g,s)\in(R\otimes_{\mathbf Q}C)^\times\times
R^\times$ with $gg^*=s$. Let $\mathbf A_f^p$ be the ring of finite adeles of $\mathbf Q$ away from
$p$, and put $G_p=\mathbf Q_p^\times\times\prod_i\operatorname{GL}_n(F_{w_i})$ and
$K_p=\mathbf Z_p^\times\times\prod_i\operatorname{GL}_n(\mathcal O_{w_i})$. An element of
$G(\mathbf A_f^p)$ is \emph{regular} if, for some prime $\ell\neq p$, the elements of
$\mathbf Q_\ell\otimes C$ that commute with its component at $\ell$ commute with each other. For an
open compact subgroup $K$ of $G(\mathbf A_f^p)$, a \emph{test function} for $K$ is a function
$f:G(\mathbf A_f^p)\to\mathbf C$ with $f(k^{-1}gk)=f(g)$ for $k\in K$ that vanishes outside finitely
many double cosets $KsK$. For a subgroup $K'$ of $G_p$ commensurable with $K_p$, a test function for
$K'$ is a function on $G_p$ invariant under conjugation by $K'$ that vanishes outside finitely many
double cosets $K_psK_p$.
\end{situation}
\leanref{Development/AbelianVarietyTheory/Definitions/Divisors/IsCyclotomicExtensionQI.lean}{GlobalTheory.LKTrip.HTGroup}
\leanref{Development/AbelianVarietyTheory/Definitions/Divisors/IsCyclotomicExtensionQI.lean}{GlobalTheory.LKTrip.HTGroup.a}
\leanref{Development/AbelianVarietyTheory/Definitions/Divisors/IsCyclotomicExtensionQI.lean}{GlobalTheory.LKTrip.UPts}
\leanref{Development/AbelianVarietyTheory/Definitions/Divisors/ExistsRootSelectionUnitaryRankOne.lean}{GlobalTheory.LKMid.Gr}
\leanref{Development/AbelianVarietyTheory/Definitions/Divisors/TwOrbMEqTwOrbUOfNormal.lean}{GlobalTheory.LKRes.KHT}
\leanref{Development/AbelianVarietyTheory/Definitions/Divisors/PointsNonemptyUnitaryRankOne.lean}{GlobalTheory.LKCx.RegAf}
\leanref{Development/AbelianVarietyTheory/Definitions/Divisors/ExistsRootSelectionUnitaryRankOne.lean}{GlobalTheory.LKMid.IsTest}

\begin{definition}\label{gal-count-forms}
In Situation~\ref{gal-count-situation}:
\begin{enumerate}
\item A \emph{form} is an element $\beta\in C^\times$ with $\beta^*=\beta$ for which there are
  $x\in(\mathbf A_f^p\otimes C)^\times$, $t\in(\mathbf A_f^p)^\times$ and
  $z\in(\mathbf R\otimes C)^\times$ with $x\beta x^*=t$ and $z\beta z^*=1$. Forms $\beta$ and
  $\beta'$ are \emph{equivalent} if $\beta'=sy\beta y^*$ for some $y\in C^\times$ and some rational
  $s>0$.
\item The \emph{group} of a form $\beta$ is $G_\beta=\{(y,s)\in C^\times\times\mathbf Q^\times:
  y\beta y^*=s\beta,\ s>0\}$. With $x$ and $t$ as in (1) fixed, it maps to $G(\mathbf A_f^p)$ by
  $\iota(y,s)=(xyx^{-1},s)$ and to $G_p$ by $\iota_p(y,s)=(s,(\rho_i(y))_i)$. The \emph{minimal
  polynomial} $P_y$ of $(y,s)$ is the minimal polynomial of $y$ over $F$.
\item Put $H_\beta=J_v\psi_v(\beta^{-1})$. A vector $z\in\mathbf C^n$ is \emph{positive} for $\beta$
  if $\operatorname{Re}(z^\dagger H_\beta z)>0$. An element $(y,s)$ of $G_\beta$ \emph{fixes} a
  nonzero $z\in\mathbf C^n$ if $\psi_v(y)z=\mu z$ for some $\mu\in\mathbf C$, its
  \emph{eigenvalue} $\mu_y(z)$.
\item Let $K\subset G(\mathbf A_f^p)$ be an open compact subgroup and $K'\subset G_p$ a subgroup. A
  \emph{point} of $\beta$ of level $(K,K')$ is a class of triples $(g,g_p,z)$, with
  $g\in G(\mathbf A_f^p)$, $g_p\in G_p$ and $z$ positive for $\beta$, under the relation
  $(g,g_p,z)\sim(\iota(y)gk,\iota_p(y)g_pk',\lambda\psi_v(y)z)$ for $(y,s)\in G_\beta$, $k\in K$,
  $k'\in K'$ and $\lambda\in\mathbf C^\times$. The \emph{stabilizer} of $(g,g_p,z)$ is the set of
  $(y,s)\in G_\beta$ that fix $z$ and satisfy $g^{-1}\iota(y)g\in K$ and
  $g_p^{-1}\iota_p(y)g_p\in K'$.
\item Let $f$ and $f_p$ be test functions for $K$ and $K'$, and $\eta:F[X]\times\mathbf C\to\mathbf C$
  a function. The \emph{complex fixed-point sum} is
  \[
   T_v(\eta;f,f_p)=\sum_{[\beta]}\ \sum_{[g,g_p,z]}\frac1{|\operatorname{Stab}(g,g_p,z)|}
   \sum_{(y,s)}f(g^{-1}\iota(y)g)\,f_p(g_p^{-1}\iota_p(y)g_p)\,\eta(P_y,\mu_y(z)),
  \]
  where $[\beta]$ runs over the equivalence classes of forms, with a representative $\beta$ (and
  $x$, $t$) chosen in each, $[g,g_p,z]$ runs over the points of $\beta$ of level $(K,K')$, with a
  representative chosen in each, and $(y,s)$ runs over the elements of $G_\beta$ that fix $z$.
\end{enumerate}
In this definition and the next, a sum over an infinite set is the sum of its nonzero terms if there
are finitely many and $0$ otherwise, and $|S|^{-1}$ is $0$ for an infinite set $S$.
\end{definition}
\leanref{Development/AbelianVarietyTheory/Definitions/Divisors/PointsNonemptyUnitaryRankOne.lean}{GlobalTheory.LKCx.Form}
\leanref{Development/AbelianVarietyTheory/Definitions/Divisors/PointsNonemptyUnitaryRankOne.lean}{GlobalTheory.LKCx.grp}
\leanref{Development/AbelianVarietyTheory/Definitions/Divisors/PointsNonemptyUnitaryRankOne.lean}{GlobalTheory.LKCx.toAfF}
\leanref{Development/AbelianVarietyTheory/Definitions/Divisors/PointsNonemptyUnitaryRankOne.lean}{GlobalTheory.LKCx.toP1}
\leanref{Development/AbelianVarietyTheory/Definitions/Divisors/PointsNonemptyUnitaryRankOne.lean}{GlobalTheory.LKCx.herm}
\leanref{Development/AbelianVarietyTheory/Definitions/Divisors/PointsNonemptyUnitaryRankOne.lean}{GlobalTheory.LKCx.FixL}
\leanref{Development/AbelianVarietyTheory/Definitions/Divisors/PointsNonemptyUnitaryRankOne.lean}{GlobalTheory.LKCx.Pts}
\leanref{Development/AbelianVarietyTheory/Definitions/Divisors/PointsNonemptyUnitaryRankOne.lean}{GlobalTheory.LKCx.Stab}
\leanref{Development/AbelianVarietyTheory/Statements/PlaceShiftStatements/MarkWeightedFixedPointEqDefinitePairOdd.lean}{GlobalTheory.RootSumWeights.cxTotalM}

\begin{definition}\label{gal-count-pairs}
In Situation~\ref{gal-count-situation}, let $K\subset G(\mathbf A_f^p)$ be an open compact subgroup,
$K'\subset G_p$ a subgroup, and $f$ and $f_p$ test functions for $K$ and $K'$.
\begin{enumerate}
\item A \emph{pair} $P=(\beta,\gamma,c)$ consists of a form $\beta$, an element $\gamma\in C^\times$
  and $c\in\mathbf Q^\times$ with $\gamma\beta\gamma^*=c\beta$, such that the minimal polynomial
  $P_\gamma$ of $\gamma$ over $F$ is separable of degree $n$. Pairs $(\beta,\gamma,c)$ and
  $(\beta',\gamma',c')$ are \emph{equivalent} if $\beta'=sy\beta y^*$ and $\gamma'=y\gamma y^{-1}$
  for some $y\in C^\times$ and some rational $s>0$.
\item A \emph{marking} of a pair $P$ is a number $\mu\in\mathbf C$ such that $\psi_v(\gamma)$ has an
  eigenvector with eigenvalue $\mu$ that is positive for $\beta$.
\item Let $Z_P$ be the group of $(y,s)\in C^\times\times\mathbf Q^\times$ with $y\gamma=\gamma y$ and
  $y\beta y^*=s\beta$, mapped to $G(\mathbf A_f^p)$ and $G_p$ by $\iota$ and $\iota_p$ as in
  Definition~\ref{gal-count-forms}(2). The \emph{orbital sum} of $P$ is
  \[
   O_P(f,f_p)=\sum_{(g,g_p)}\frac{f(g^{-1}\iota(\gamma,c)g)\,f_p(g_p^{-1}\iota_p(\gamma,c)g_p)}
   {|Z_P\cap(g,g_p)(K\times K')(g,g_p)^{-1}|},
  \]
  where $(g,g_p)$ runs over representatives of $Z_P\backslash(G(\mathbf A_f^p)\times G_p)/(K\times
  K')$, and $Z_P$ acts through $(\iota,\iota_p)$.
\item A \emph{definite form} is an element $\beta\in C^\times$ with $\beta^*=\beta$ and $x\beta x^*=t$
  for some $x\in(\mathbf A_f^p\otimes C)^\times$ and $t\in(\mathbf A_f^p)^\times$, such that for every
  infinite place $w$ the Hermitian matrix $J_w\psi_w(\beta^{-1})$ is positive definite if $w=v$ and
  negative definite otherwise. A \emph{definite pair} is a triple $(\beta,\gamma,c)$ as in (1) with
  $\beta$ a definite form. Equivalence, the group $G_\beta$, the maps $\iota$, $\iota_p$ and the
  orbital sums are defined for them as for forms and pairs.
\item Suppose that $f$ is invariant under $K$ on both sides and $f_p$ under $K'$ on both sides. Let
  $\xi$ be a representation of $C^\times\times\mathbf Q^\times$ on a finite-dimensional complex vector
  space $V$, and $\beta$ a definite form. Let $\mathcal A_\beta(\xi)$ be the space of functions
  $\varphi:G(\mathbf A_f^p)\times G_p\to V$ with $\varphi((\iota,\iota_p)(y)g)=\xi(y)\varphi(g)$ for
  $y\in G_\beta$ and $\varphi(gk)=\varphi(g)$ for $k\in K\times K'$, with the Hecke operator
  $\varphi\mapsto\bigl(g\mapsto\sum_{(h,h_p)}f(h)f_p(h_p)\varphi(g(h,h_p))\bigr)$, the sum over
  $(G(\mathbf A_f^p)\times G_p)/(K\times K')$. The \emph{definite trace}
  $\operatorname{tr}_{\mathrm{def}}(\xi;f,f_p)$ is the sum, over the equivalence classes of definite
  forms, of the traces of these Hecke operators.
\item The test function $f$ is \emph{elliptic} if it vanishes outside the regular elements,
  $O_P(f,\mathbf 1_{K_p})=0$ for every pair $P$ that has a marking and has reducible $P_\gamma$, and,
  for every definite form $\beta$ and every $y\in G_\beta$ whose minimal polynomial is not
  irreducible, separable and of degree $n$, the orbital integral
  $\sum_{g}\operatorname{vol}(Z\cap g(K\times K_p)g^{-1})^{-1}(f\otimes\mathbf 1_{K_p})(g^{-1}\delta g)$
  vanishes, where $\delta=(\iota,\iota_p)(y)$, $Z$ is the centralizer of $\delta$ in
  $G(\mathbf A_f^p)\times G_p$, $g$ runs over $Z\backslash(G(\mathbf A_f^p)\times G_p)/(K\times K_p)$,
  and $\operatorname{vol}(Z\cap(K\times K_p))=1$.
\end{enumerate}
\end{definition}
\leanref{Development/AbelianVarietyTheory/Definitions/Divisors/ExistsRootSelectionUnitaryRankOne.lean}{GlobalTheory.LKMid.Pair}
\leanref{Development/AbelianVarietyTheory/Definitions/Divisors/ExistsRootSelectionUnitaryRankOne.lean}{GlobalTheory.LKMid.Pair.Mark}
\leanref{Development/AbelianVarietyTheory/Definitions/Divisors/ExistsRootSelectionUnitaryRankOne.lean}{GlobalTheory.LKMid.Pair.term}
\leanref{Development/AbelianVarietyTheory/Definitions/Divisors/PointSide.lean}{GlobalTheory.LKCount.pointSum}
\leanref{Development/AbelianVarietyTheory/Definitions/Divisors/DiagHermitianTrivialOfSumFourSquares.lean}{GlobalTheory.LKDef.DPair}
\leanref{Development/AbelianVarietyTheory/Statements/PlaceShiftStatements/TraceDefinitePairsEqAutomorphicTrace.lean}{GlobalTheory.LKDef.DForm}
\leanref{Development/AbelianVarietyTheory/Statements/PlaceShiftStatements/TraceDefinitePairsEqAutomorphicTrace.lean}{GlobalTheory.LKDef.DForm.trace}
\leanref{Development/AbelianVarietyTheory/Statements/PlaceShiftStatements/TraceDefinitePairsEqAutomorphicTrace.lean}{GlobalTheory.LKDef.dTrace}
\leanref{Development/AbelianVarietyTheory/Definitions/Divisors/PointSide.lean}{GlobalTheory.LKCount.orbK}
\leanref{Development/AutomorphicArithmetic/Foundations/GaloisRepresentations/LefschetzSigmaSums/InvPol.lean}{UpperTheory.Gal.LocalSmoothExtension.Good}

\begin{lemma}\label{gal-count-positive-line}
In Situation~\ref{gal-count-situation}, let $\beta$ be a form and $(y,s)\in G_\beta$ with $P_y$
separable of degree $n$. If $z$ and $z'$ are positive for $\beta$ and both are fixed by $(y,s)$, then
$z'\in\mathbf Cz$. In particular, a pair has at most one marking.
\end{lemma}
\leanref{Development/AbelianVarietyTheory/Proofs/DivisorProofs/SeparableFormGroupElementFixesUniquePositiveLineStatement.lean}{GlobalTheory.T.separable_formGroup_element_fixes_unique_positive_line}
\leanref{Development/AbelianVarietyTheory/Proofs/DivisorProofs/DotConjTransposeMulMul.lean}{GlobalTheory.LKMid.Pair.mark_unique}

\begin{proof}
Write $Y=\psi_v(y)$ and $H=H_\beta$. Since $\beta^*=\beta$, $H$ is Hermitian, and applying $\psi_v$ to
$\beta^{-1}=s^{-1}y^*\beta^{-1}y$ gives $Y^\dagger HY=sH$. If $z\beta z^*=1$ with
$z\in(\mathbf R\otimes C)^\times$, then $H=\psi_v(z)^\dagger J_v\psi_v(z)$, so $H$ has signature
$(1,n-1)$. Since $P_y$ is separable of degree $n$, $Y$ has $n$ distinct eigenvalues and the vectors
it fixes are its eigenvectors. Let $e,e'$ be eigenvectors with eigenvalues $\mu\neq\mu'$, both
positive. From $Y^\dagger HY=sH$, $|\mu|^2e^\dagger He=se^\dagger He$, so $|\mu|^2=s=|\mu'|^2$, and
$\bar\mu\mu'e^\dagger He'=se^\dagger He'$. If $e^\dagger He'\neq0$ then $\bar\mu\mu'=s=\bar\mu\mu$, so
$\mu'=\mu$; hence $e^\dagger He'=0$ and $H$ is positive definite on the span of $e$ and $e'$, which
contradicts the signature. For a pair $(\beta,\gamma,c)$ with a marking $\mu$ the same computation
gives $c=|\mu|^2>0$, so $(\gamma,c)\in G_\beta$, and two markings have proportional eigenvectors,
hence are equal.
\end{proof}

\begin{lemma}\label{gal-complex-majorant}
Let $n\geq1$, let $H$ be an invertible Hermitian $n\times n$ matrix of signature $(1,n-1)$, and let
$z\in\mathbf C^n$ with $a=z^\dagger Hz>0$. Put $P=2a^{-1}(Hz)(Hz)^\dagger-H$. Then $P$ is positive
definite, and $Y^\dagger PY=sP$ for every $Y\in M_n(\mathbf C)$, real $s>0$ and $\mu\in\mathbf C$ with
$Y^\dagger HY=sH$ and $Yz=\mu z$.
\end{lemma}
\leanref{Development/AbelianVarietyTheory/Proofs/DivisorProofs/ExistsPosDefMajorantOfPositiveIndexOneStatement.lean}{GlobalTheory.T.exists_posDef_majorant_of_positiveIndex_one}

\begin{proof}
Put $u=Hz$. Every vector of $\mathbf C^n$ is $\lambda z+w$ with $\lambda\in\mathbf C$ and $u^\dagger w=0$.
The form $H$ is negative definite on the hyperplane $u^\perp$. Indeed, $z^\dagger Hw=u^\dagger w=0$ for
$w\in u^\perp$, so a nonzero $w\in u^\perp$ with $w^\dagger Hw\geq0$ would span with $z$ a plane on which
$H$ is positive semidefinite, and this plane would meet the $(n-1)$-dimensional subspace on which $H$ is
negative definite. Since $Pz=u$ and $Pw=-Hw$ for $w\in u^\perp$,
\[
 (\lambda z+w)^\dagger P(\lambda z+w)=|\lambda|^2a-w^\dagger Hw,
\]
which is positive unless $\lambda z+w=0$. Now let $Y^\dagger HY=sH$ and $Yz=\mu z$. Then
$|\mu|^2a=z^\dagger Y^\dagger HYz=sa$, so $|\mu|^2=s$ and $\mu\neq0$. Also $\mu Y^\dagger u=Y^\dagger HYz=su$,
so $Y^\dagger u=(s/\mu)u$. Hence $Y^\dagger PY=2a^{-1}(s^2/|\mu|^2)uu^\dagger-sH=sP$.
\end{proof}

\begin{lemma}\label{gal-complex-archimedean}
In Situation~\ref{gal-count-situation}, let $\beta$ be a form and $z\in\mathbf C^n$ positive for $\beta$.
Then there is a real number $B$ such that $|\psi_w(y)_{ij}|^2\leq Bs$ for every $(y,s)\in G_\beta$ that
fixes $z$, every infinite place $w$ of $F$ and all $i,j$.
\end{lemma}
\leanref{Development/AbelianVarietyTheory/Proofs/DivisorProofs/FixedVectorSquaredBound.lean}{GlobalTheory.T.fixed_vector_squared_bound}
\leanref{Development/AbelianVarietyTheory/Proofs/DivisorProofs/LogResolutionFubiniStudyPositive.lean}{GlobalTheory.LKCx.entries_sq_le_of_fixL}
\leanref{Development/AbelianVarietyTheory/Proofs/DivisorProofs/HermitianSimilitudeEntrySqBoundStatement.lean}{GlobalTheory.T.hermitian_similitude_entry_sq_bound}

\begin{proof}
For an infinite place $w$ put $H_w=J_w\psi_w(\beta^{-1})$ and $Y_w=\psi_w(y)$. Applying $\psi_w$ to
$\beta^{-1}=s^{-1}y^*\beta^{-1}y$ gives $Y_w^\dagger H_wY_w=sH_w$. If $z'\beta z'^*=1$ with
$z'\in(\mathbf R\otimes C)^\times$, then $H_w=\psi_w(z')^\dagger J_w\psi_w(z')$, so $H_w$ is Hermitian with
the signature of $J_w$, as in the proof of Lemma~\ref{gal-count-positive-line}. For $w\neq v$ the matrix
$P_w=-H_w$ is therefore positive definite, and $Y_w^\dagger P_wY_w=sP_w$. For $w=v$, the matrix
$H_v=H_\beta$ has signature $(1,n-1)$, $z^\dagger H_vz>0$ and $Y_vz=\mu_y(z)z$, so
Lemma~\ref{gal-complex-majorant} gives a positive definite $P_v$, depending only on $\beta$ and $z$, with
$Y_v^\dagger P_vY_v=sP_v$. Choose $0<m_w\leq M_w$ with $m_w|e|^2\leq e^\dagger P_we\leq M_w|e|^2$ for
$e\in\mathbf C^n$. Then
\[
 |Y_we|^2\leq m_w^{-1}(Y_we)^\dagger P_w(Y_we)=m_w^{-1}s\,e^\dagger P_we\leq(M_w/m_w)\,s\,|e|^2,
\]
so every entry of $Y_w$ has square at most $(M_w/m_w)s$. Take $B=\max_wM_w/m_w$.
\end{proof}

\begin{lemma}\label{gal-complex-lattice}
\begin{enumerate}
\item Let $V$ be a $\mathbf Q$-vector space, $E$ a finite-dimensional normed real vector space and
  $\Phi:V\to E$ an additive map that sends every finite family in $V$ that is linearly independent over
  $\mathbf Q$ to a family linearly independent over $\mathbf R$. Let $S\subset V$ be a finite set, $Y$ a subset
  of the subgroup generated by $S$, and $R$ a real number with $\|\Phi(y)\|\leq R$ for all $y\in Y$. Then $Y$
  is finite.
\item In Situation~\ref{gal-count-situation}, let $\Phi:C\to\prod_wM_n(\mathbf C)$ be
  $c\mapsto(\psi_w(c))_w$, with $w$ running over the infinite places of $F$. Then $\Phi$ sends every finite
  family in $C$ that is linearly independent over $\mathbf Q$ to a family linearly independent over
  $\mathbf R$.
\end{enumerate}
\end{lemma}
\leanref{Development/AbelianVarietyTheory/Proofs/DivisorProofs/FiniteOfBoundedImageInFinitelyGeneratedSubgroupStatement.lean}{GlobalTheory.T.finite_of_bounded_image_in_finitely_generated_subgroup}
\leanref{Development/AbelianVarietyTheory/Proofs/DivisorProofs/EmbeddingsRealIndependent.lean}{GlobalTheory.T.embeddings_real_independent}
\leanref{Development/AbelianVarietyTheory/Proofs/DivisorProofs/GenericIntersectionEqualsPointsIntegral.lean}{GlobalTheory.LKCx.psi_linearIndependent_real}

\begin{proof}
(1) Let $M\subset V$ be the subgroup generated by $S$. It is finitely generated and torsion free, since $V$ is
a $\mathbf Q$-vector space, so it is free, with a basis $e_1,\ldots,e_k$. This basis is linearly independent
over $\mathbf Z$, hence over $\mathbf Q$, so the $\Phi(e_j)$ are linearly independent over $\mathbf R$. So the
map $T:\mathbf R^k\to E$, $T(a)=\sum_ja_j\Phi(e_j)$, is injective, and $\|a\|\leq\kappa\|T(a)\|$ for some
$\kappa$, with the maximum norm on $\mathbf R^k$. For $y\in Y$ with coordinates $c\in\mathbf Z^k$ in the basis
$(e_j)$, $\Phi(y)=T(c)$, so $|c_j|\leq\kappa R$ for every $j$. So $y\mapsto c$ maps $Y$ injectively into a
finite set, and $Y$ is finite.

(2) The $\mathbf R$-linear extension of $\Phi$ to $\mathbf R\otimes_{\mathbf Q}C$ is injective. Indeed, let
$\omega_1,\ldots,\omega_r$ be a $\mathbf Z$-basis of
$\mathcal O_F$ and $c_1,\ldots,c_{n^2}$ a basis of $C$ over $F$, so that the $\omega_ic_j$ form a basis of
$C$ over $\mathbf Q$. If $\sum_{i,j}r_{ij}\Phi(\omega_ic_j)=0$ with $r_{ij}\in\mathbf R$, then
$\sum_j\bigl(\sum_ir_{ij}w(\omega_i)\bigr)\psi_w(c_j)=0$ for each $w$. The matrices $\psi_w(c_j)$ are
linearly independent over $\mathbf C$, since $\psi_w$ induces a homomorphism
$\mathbf C\otimes_{F,w}C\to M_n(\mathbf C)$ whose source is simple. So $\sum_ir_{ij}w(\omega_i)=0$ for all
$w$ and $j$. Since $F$ is totally complex, the vectors $(w(\omega_i))_w$ form a basis of $\prod_w\mathbf C$
over $\mathbf R$ \cite[Chapter~I, \S5]{Neu99}, so all $r_{ij}$ vanish.

If $e_1,\ldots,e_k\in C$ are linearly independent over $\mathbf Q$, the elements $1\otimes e_j$ of
$\mathbf R\otimes_{\mathbf Q}C$ are linearly independent over $\mathbf R$, and so are their images
$\Phi(e_j)$ under the injective extension.
\end{proof}

\begin{lemma}\label{gal-complex-fixed-finite}
In Situation~\ref{gal-count-situation}, let $K\subset G(\mathbf A_f^p)$ be an open compact subgroup, $f$ a
test function for $K$, $K'\subset G_p$ a subgroup commensurable with $K_p$, and $f_p$ a test function for
$K'$. Let $\beta$ be a form, with $x$ and $t$ as in Definition~\ref{gal-count-forms}(1), let
$(g,g_p)\in G(\mathbf A_f^p)\times G_p$, and let $z$ be positive for $\beta$. Then only finitely many
$(y,s)\in G_\beta$ fix $z$ and satisfy
\[
 f(g^{-1}\iota(y,s)g)\,f_p(g_p^{-1}\iota_p(y,s)g_p)\neq0.
\]
\end{lemma}
\leanref{Development/AbelianVarietyTheory/Proofs/DivisorProofs/ExistsIntOfDenLe.lean}{GlobalTheory.FinFix.finite_fix_core}
\leanref{Development/AbelianVarietyTheory/Proofs/DivisorProofs/ExistsIntOfDenLe.lean}{GlobalTheory.FinFix.grp_val_injective}

\begin{proof}
Fix a basis $e_1,\ldots,e_N$ of $C$ over $\mathbf Q$, let $D$ and $k$ be given by
Lemma~\ref{gal-definite-form-denominators} for $h=g$ and $h_p=g_p$, and put $M=Dp^k$. Let $(y,s)$ be as
in the statement. The coordinates of $y$ in the basis $(e_j)$ lie in $M^{-1}\mathbf Z$, so $y$ lies in the
subgroup $\Lambda$ generated by the $M^{-1}e_j$. The number $Ms^{-1}$ is an integer, positive since $s>0$,
so $s\leq M$. With $B$ as in Lemma~\ref{gal-complex-archimedean}, $|\psi_w(y)_{ij}|^2\leq BM$ for all
$w$, $i$ and $j$. So only finitely many $y$ occur, by Lemma~\ref{gal-complex-lattice}(1) for the map $\Phi$
of Lemma~\ref{gal-complex-lattice}(2), with the maximum of the absolute values of the entries as the norm,
$S$ the set of the $M^{-1}e_j$ and $R=\sqrt{BM}$. Finally, $y$
determines $s$, since $y\beta y^*=s\beta$ and $\beta$ is invertible.
\end{proof}

\begin{lemma}\label{gal-complex-separable}
In Situation~\ref{gal-count-situation}, let $\beta$ be a form, with $x$ and $t$, let $(y,s)\in G_\beta$ fix
a vector positive for $\beta$, and let $g\in G(\mathbf A_f^p)$ be such that $g^{-1}\iota(y,s)g$ is regular.
Then any two elements of $C$ that commute with $y$ commute with each other, and $P_y$ is separable of
degree $n$.
\end{lemma}
\leanref{Development/AbelianVarietyTheory/Proofs/DivisorProofs/MinpolyOfRegAfOfFixL.lean}{GlobalTheory.LKCx.minpoly_of_regAf_of_fixL}
\leanref{Development/AbelianVarietyTheory/Proofs/DivisorProofs/ClassSummandEqTermOfAOne.lean}{GlobalTheory.LKCx.CxReg.sep_of_regAf}

\begin{proof}
Write $g=(g_1,\nu)$, and let $\ell\neq p$ be a prime as in the definition of regularity. The component at
$\ell$ of the first coordinate of $g^{-1}\iota(y,s)g$ is $u(1\otimes y)u^{-1}$, where $u$ is the component
at $\ell$ of $g_1^{-1}x$. Conjugating by $u$, any two elements of $\mathbf Q_\ell\otimes C$ that commute
with $1\otimes y$ commute with each other. As $C\to\mathbf Q_\ell\otimes C$ is injective, the same holds in
$C$ for $y$.

Let $Y=\psi_v(y)$. As in the proof of Lemma~\ref{gal-complex-archimedean}, $Y^\dagger PY=sP$ for a positive
definite $P$, so $s^{-1/2}Y$ is unitary for the inner product $(e,e')\mapsto e^\dagger Pe'$. For $q\in F[X]$
let $q^v$ be its image under $v$. Then $\psi_v(q(y))=q^v(Y)$ is normal for this inner product, so
$q^v(Y)^2=0$ implies $q^v(Y)=0$; as $\psi_v$ is injective, $C$ being simple, $q(y)^2=0$ implies $q(y)=0$.
If $P_y$ had a repeated irreducible factor $\pi$, then $q=P_y/\pi$ would satisfy $q(y)^2=0$ and
$q(y)\neq0$. So $P_y$ is separable.

Let $\bar F$ be an algebraic closure of $F$; then $\bar F\otimes_FC\cong M_n(\bar F)$, and let $N$ be the
image of $1\otimes y$. Its minimal polynomial $Q$ over $\bar F$ divides $P_y$, so it is separable, and
$\deg Q\geq\deg P_y$, since $1,y,\ldots,y^{\deg P_y-1}$ stay linearly independent over $\bar F$ in
$\bar F\otimes_FC$. So $\deg Q=\deg P_y$. Tensoring with $\bar F$ the exact sequence $0\to Z\to C\to C$,
where $Z$ is the centralizer of $y$ and the last map is $e\mapsto ey-ye$, shows that the centralizer of
$1\otimes y$ in $\bar F\otimes_FC$ is $\bar F\otimes_FZ$, which is commutative by the first paragraph. Now
$\deg Q\leq n$, as $Q$ divides the characteristic polynomial of $N$. Conversely, since $Q$ is separable,
$\bar F^n$ is the direct sum of the eigenspaces $V_\mu$ of $N$ for the $\deg Q$ roots $\mu$ of $Q$. Suppose
that some $V_\mu$ has dimension at least $2$, and put $q=Q/(X-\mu)$, so that $q(\mu)\neq0$, $q(N)$ maps
$\bar F^n$ into $V_\mu$, and $q(N)$ is $q(\mu)$ on $V_\mu$. For an endomorphism $A$ of $V_\mu$, the
endomorphism $A\,q(N)$ of $\bar F^n$ commutes with $N$, as both $A\,q(N)N$ and $NA\,q(N)$ are $\mu A\,q(N)$.
Choosing endomorphisms $A_1,A_2$ of $V_\mu$ that do not commute, $A_1q(N)$ and $A_2q(N)$ restrict on $V_\mu$
to $q(\mu)^2A_1A_2$ and $q(\mu)^2A_2A_1$, so they do not commute, which contradicts the commutativity of the
centralizer of $N$. So every $V_\mu$ is a line, $n\leq\deg Q$, and $\deg P_y=n$.
\end{proof}

\begin{remark}\label{gal-complex-separable-degree}
The equality $\deg P_y=n$ in Lemma~\ref{gal-complex-separable} also follows from the structure of the
centralizer of a diagonalizable matrix. Over $\bar F$ the eigenvalues of $1\otimes y$ have multiplicities
$m_1,\ldots,m_k$ adding up to $n$, its centralizer is isomorphic to $\prod_iM_{m_i}(\bar F)$, and this
algebra is commutative only if all $m_i$ are $1$.
\end{remark}

\begin{lemma}\label{gal-complex-polynomials}
In Situation~\ref{gal-count-situation}, let $K$, $f$, $K'$ and $f_p$ be as in
Lemma~\ref{gal-complex-fixed-finite}, and let $\beta$ be a form, with $x$ and $t$. Consider the
$(y,s)\in G_\beta$ for which there are $(g,g_p)\in G(\mathbf A_f^p)\times G_p$ and a vector $z$ positive for
$\beta$ and fixed by $(y,s)$ with $f(g^{-1}\iota(y,s)g)\,f_p(g_p^{-1}\iota_p(y,s)g_p)\neq0$. Their
minimal polynomials $P_y$ form a finite set.
\end{lemma}
\leanref{Development/AbelianVarietyTheory/Proofs/DivisorProofs/OrbitalTermRealTrivialisation.lean}{GlobalTheory.T.orbital_term_real_trivialisation}
\leanref{Development/AbelianVarietyTheory/Proofs/DivisorProofs/HenselianQuasiFiniteFinitePart.lean}{GlobalTheory.LKCx.finite_minpoly_fix_a_one}
\leanref{Development/AbelianVarietyTheory/Proofs/DivisorProofs/MinpolyBoundedDenominators.lean}{GlobalTheory.T.minpoly_bounded_denominators}
\leanref{Development/AbelianVarietyTheory/Proofs/DivisorProofs/MinpolyPadicBounded.lean}{GlobalTheory.T.minpoly_padic_bounded}
\leanref{Development/AbelianVarietyTheory/Proofs/DivisorProofs/SimilitudeFactorFormGroupBoundedOfTestSupportStatement.lean}{GlobalTheory.T.similitudeFactor_formGroup_bounded_of_test_support}
\leanref{Development/AbelianVarietyTheory/Proofs/DivisorProofs/FiniteMonicBoundedPolynomials.lean}{GlobalTheory.T.finite_monic_bounded_polynomials}

\begin{proof}
Let $(y,s)$, $(g,g_p)$ and $z$ be as in the statement; the bounds below do not depend on them. Since $f$
and $f_p$ vanish outside finitely many double cosets $K\sigma K$ and $K_p\sigma'K_p$, the elements
$g^{-1}\iota(y,s)g$ and $g_p^{-1}\iota_p(y,s)g_p$ lie in fixed compact sets. Let $d=\deg P_y\leq n^2$.

(a) Denominators away from $p$. For $h=(h_1,\nu)\in G(\mathbf A_f^p)$, the coefficients of the
characteristic polynomial of left multiplication by $h_1$ on $\mathbf A_f^p\otimes C$ are continuous in
$h$. At $h=g^{-1}\iota(y,s)g$, whose first coordinate is conjugate to $1\otimes y$, they are the
coefficients of the characteristic polynomial $\chi_y\in\mathbf Q[X]$ of left multiplication by $y$ on
$C$. So these rational numbers lie in a compact subset of $\mathbf A_f^p$, and as in the proof of
Lemma~\ref{gal-definite-form-denominators}(1) there is an integer $D'\geq1$ such that $D'$ times each of
them is integral at every prime $\ell\neq p$. Then for some $m$ the element $p^mD'y$ has a characteristic
polynomial with integer coefficients, so it is integral over $\mathbf Z$, and the coefficients of its
minimal polynomial over $F$ lie in $\mathcal O_F$. The coefficient of $X^j$ in that polynomial is
$(p^mD')^{d-j}$ times that of $P_y$, so every coefficient $\alpha$ of $P_y$ satisfies
$p^{m'}D\alpha\in\mathcal O_F$ for $D=D'^{n^2}$ and some $m'$.

(b) Bounds at $p$. The roots of $P_y$ under an embedding of $F$ into an algebraic closure of
$\mathbf Q_p$ are eigenvalues of $\rho_i(y)$, for some $i$, if the embedding lies over $u$, and of
$\rho_i(y^*)=s\rho_i(\beta)^{-1}\rho_i(y)^{-1}\rho_i(\beta)$ if it lies over $\bar u$. Since
$g_p^{-1}\iota_p(y,s)g_p$ and its inverse lie in compact sets, these eigenvalues, and so the coefficients
of $P_y$, are bounded at every such embedding, independently of $(y,s)$. With (a), the valuations of
$D\alpha$ at the places of $F$ over $p$ are bounded below, so there is $m_0$ with
$p^{m_0}D\alpha\in\mathcal O_F$ for every coefficient $\alpha$ of $P_y$.

(c) The similitude factor. The number $s$ is rational and positive. For $\ell\neq p$, its $\ell$-adic
valuation is that of the second coordinate of $g^{-1}\iota(y,s)g$, which is constant on each double coset
$K\sigma K$, since the second coordinates of the elements of the compact group $K$ are units. In the same
way $|s|_p$ is determined by the double coset $K_p\sigma'K_p$ containing $g_p^{-1}\iota_p(y,s)g_p$. So $s$
takes finitely many values, and $s\leq S$ for some $S$.

(d) Bounds at the infinite places. Let $w$ be an infinite place and $\mu$ an eigenvalue of $\psi_w(y)$,
with eigenvector $e$. Let $P_w$ be as in the proof of Lemma~\ref{gal-complex-archimedean}, which applies
since $(y,s)$ fixes $z$. Then $|\mu|^2e^\dagger P_we=e^\dagger\psi_w(y)^\dagger P_w\psi_w(y)e=se^\dagger
P_we$, so $|\mu|^2=s$. The roots of $P_y$ under the two complex embeddings at $w$ are the eigenvalues of
$\psi_w(y)$ and their complex conjugates, so they have absolute value $s^{1/2}\leq S^{1/2}$, and the
coefficients of $P_y$ have absolute value at most $2^{n^2}\max(1,S)^{n^2/2}$ under every complex embedding
of $F$.

(e) Since $\mathcal O_F$ is a lattice in $\mathbf R\otimes F$, only finitely many monic polynomials over $F$
of degree at most $n^2$ have their coefficients in $(p^{m_0}D)^{-1}\mathcal O_F$ and bounded as in (d). So
$P_y$ lies in a finite set.
\end{proof}

\begin{lemma}\label{gal-complex-cosets}
In Situation~\ref{gal-count-situation}, let $K$, $f$, $K'$ and $f_p$ be as in
Lemma~\ref{gal-complex-fixed-finite}, let $\beta$ be a form, with $x$ and $t$, and let
$(y_0,s_0)\in G_\beta$ with $P_{y_0}$ separable of degree $n$. Put $\delta=\iota(y_0,s_0)$ and
$\delta_p=\iota_p(y_0,s_0)$. Then there are finite sets $R\subset G(\mathbf A_f^p)$ and $R_p\subset G_p$
such that:
\begin{enumerate}
\item every $h\in G(\mathbf A_f^p)$ with $f(h^{-1}\delta h)\neq0$ is $crk$ with $c\delta=\delta c$, $r\in R$
  and $k\in K$;
\item every $h\in G_p$ with $f_p(h^{-1}\delta_ph)\neq0$ is $crk$ with $c\delta_p=\delta_pc$, $r\in R_p$ and
  $k\in K'$.
\end{enumerate}
\end{lemma}
\leanref{Development/AbelianVarietyTheory/Proofs/DivisorProofs/FiniteCosetsAndFixedPoints.lean}{GlobalTheory.FinPtsSep.fin_cosets_away_sep}
\leanref{Development/AbelianVarietyTheory/Proofs/DivisorProofs/FiniteCosetsAndFixedPoints.lean}{GlobalTheory.FinPtsSep.exists_finset_doubleCoset_cover_of_separable}

\begin{proof}
(1) For a prime $\ell\neq p$, $G(\mathbf Q_\ell)$ is the group of $(h,\nu)\in(\mathbf Q_\ell\otimes
C)^\times\times\mathbf Q_\ell^\times$ with $hh^*=\nu$, and $\delta_\ell$ denotes the component of $\delta$ at
$\ell$. Fix a basis of $C$ over $\mathbf Q$. For almost all $\ell$ its $\mathbf Z_\ell$-span $\mathcal O_\ell$
in $\mathbf Q_\ell\otimes C$ is a ring stable under the involution; let $G(\mathbf Z_\ell)$ be the group of
$(h,\nu)\in G(\mathbf Q_\ell)$ with $h,h^{-1}\in\mathcal O_\ell$ and $\nu\in\mathbf Z_\ell^\times$. The
elements of $G(\mathbf A_f^p)$ have their components in $G(\mathbf Z_\ell)$ for almost all $\ell$, and so do
the elements of a compact subset uniformly. The support of $f$ lies in a compact set $\Omega$, whose
components lie in $G(\mathbf Z_\ell)$ for $\ell$ outside a finite set $T$.

Conjugation by the component $x_\ell$ of $x$ identifies $G(\mathbf Q_\ell)$ with
$\mathrm{GU}_\beta(\mathbf Q_\ell)$ (Definition~\ref{gal-definite-form-similitude}), carries
$\delta_\ell$ to the rational element $(y_0,s_0)$, which is semisimple since $P_{y_0}$ is separable, and
for almost all $\ell$ matches $G(\mathbf Z_\ell)$ with the corresponding integral points. So by Kottwitz's
lemma \cite[Proposition~7.1]{Kot86b}, after enlarging $T$, every $h_\ell\in G(\mathbf Q_\ell)$ with
$h_\ell^{-1}\delta_\ell h_\ell\in G(\mathbf Z_\ell)$ lies in $Z(\delta_\ell)G(\mathbf Z_\ell)$, where
$Z(\delta_\ell)$ is the centralizer of $\delta_\ell$; put $D_\ell=G(\mathbf Z_\ell)$ for $\ell\notin T$. For
$\ell\in T$, the conjugacy class of the semisimple element $\delta_\ell$ is closed, and $h\mapsto
h^{-1}\delta_\ell h$ induces a homeomorphism of $Z(\delta_\ell)\backslash G(\mathbf Q_\ell)$ onto it
\cite{HC70}; so there is a compact $D_\ell\subset G(\mathbf Q_\ell)$ with $h_\ell\in Z(\delta_\ell)D_\ell$
whenever $h_\ell^{-1}\delta_\ell h_\ell$ lies in the component of $\Omega$ at $\ell$.

Now let $f(h^{-1}\delta h)\neq0$, and write $h_\ell=c_\ell d_\ell$ with $c_\ell\in Z(\delta_\ell)$ and
$d_\ell\in D_\ell$. For almost all $\ell$ both $h_\ell$ and $d_\ell$ lie in $G(\mathbf Z_\ell)$, hence so
does $c_\ell$. So $c=(c_\ell)_\ell$ is an element of $G(\mathbf A_f^p)$ commuting with $\delta$, and $h=cd$
with $d$ in the compact set $\prod_\ell D_\ell$. Finitely many cosets $rK$ cover this set, and their
representatives $r$ form $R$.

(2) The group $G_p$ is $\mathbf Q_p^\times\times\prod_i\operatorname{GL}_n(F_{w_i})$, with
$K_p=\mathbf Z_p^\times\times\prod_i\operatorname{GL}_n(\mathcal O_{w_i})$. The component $\rho_i(y_0)$ of
$\delta_p$ has minimal polynomial the image of $P_{y_0}$, of degree $n$, so its characteristic polynomial
is separable. If $f_p(h^{-1}\delta_ph)\neq0$, then $h^{-1}\delta_ph$ lies in one of finitely many double
cosets $K_p\sigma K_p$, so the component $h_i$ of $h$ satisfies
$h_i^{-1}\rho_i(y_0)h_i\in\operatorname{GL}_n(\mathcal O_{w_i})\sigma_i\operatorname{GL}_n(\mathcal O_{w_i})$
for one of finitely many $\sigma_i$. By the same properness in $\operatorname{GL}_n(F_{w_i})$
\cite{HC70}, $h_i=a_ir_ik_i$ with $a_i$ commuting with $\rho_i(y_0)$, $r_i$ in a finite set and
$k_i\in\operatorname{GL}_n(\mathcal O_{w_i})$. The component of $h$ in $\mathbf Q_p^\times$ commutes with
$\delta_p$. So $h=crk$ with $c$ commuting with $\delta_p$, $r$ in a finite set and $k\in K_p$. Since $K'$
is commensurable with $K_p$, $K_p$ is the union of finitely many cosets $t(K'\cap K_p)$, and $R_p$ consists
of the products $rt$.
\end{proof}

\begin{lemma}\label{gal-complex-torus}
In Situation~\ref{gal-count-situation}, let $K\subset G(\mathbf A_f^p)$ be an open compact subgroup and
$K'\subset G_p$ a subgroup commensurable with $K_p$. Let $\beta$ be a form, with $x$ and $t$, and let
$(y_0,s_0)\in G_\beta$ with $P_{y_0}$ separable of degree $n$, with centralizer $Z$ in $G_\beta$. Put
$\delta=\iota(y_0,s_0)$ and $\delta_p=\iota_p(y_0,s_0)$, and let $r\in G(\mathbf A_f^p)\times G_p$. Then
there is a finite set $R_r\subset G(\mathbf A_f^p)\times G_p$ such that every
$c\in G(\mathbf A_f^p)\times G_p$ commuting with $(\delta,\delta_p)$ satisfies
$cr\in(\iota,\iota_p)(Z)\,R_r\,(K\times K')$.
\end{lemma}
\leanref{Development/AbelianVarietyTheory/Proofs/DivisorProofs/FiniteClassNumberCentraliserTorusSeparableStatement.lean}{GlobalTheory.T.finite_classNumber_centraliserTorus_separable}
\leanref{Development/AbelianVarietyTheory/Proofs/DivisorProofs/FiniteClassNumberCentraliserTorusSeparable.lean}{GlobalTheory.CentCls.main}

\begin{proof}
Let $T$ be the linear algebraic group over $\mathbf Q$ whose points in a commutative $\mathbf Q$-algebra $A$
are the $(e,\sigma)\in\mathrm{GU}_\beta(A)$ (Definition~\ref{gal-definite-form-similitude}) with
$e(1\otimes y_0)=(1\otimes y_0)e$. Then $Z$ is the set of elements of $T(\mathbf Q)$ with $\sigma>0$. Let
$a:\mathrm{GU}_\beta(\mathbf A_f^p)\to G(\mathbf A_f^p)$ be $(e,\sigma)\mapsto(xex^{-1},\sigma)$ and
$b:\mathrm{GU}_\beta(\mathbf Q_p)\to G_p$ be $(e,\sigma)\mapsto(\sigma,(\rho_i(e))_i)$. Both are bijective
homomorphisms: for $a$ this is shown in the proof of Lemma~\ref{gal-definite-form-classes}; $b$ is onto by
Lemma~\ref{gal-definite-form-padic}(3), and one to one by Lemma~\ref{gal-definite-form-padic}(1), since
$\rho_i(e^*)=\sigma\rho_i(\beta)^{-1}\rho_i(e)^{-1}\rho_i(\beta)$. On $G_\beta$ they restrict to $\iota$ and
$\iota_p$, so $a\times b$ carries $T(\mathbf A_f^p)\times T(\mathbf Q_p)$ onto the centralizer of
$(\delta,\delta_p)$.

The subgroup $U$ of $T(\mathbf A_f^p)\times T(\mathbf Q_p)$ of the elements that $a\times b$ maps into
$r(K\times K_p)r^{-1}$ is open, so it is the intersection of $T(\mathbf A_f^p)\times T(\mathbf Q_p)$ with an
open neighbourhood of $1$ in the group $H_C$ of Lemma~\ref{gal-definite-form-class-number} for $D=C$. That
lemma, Borel's finiteness of class numbers, applied to $C$ with its involution, $\beta$ and $y=y_0$, whose
group is $T$, gives a finite set $S\subset H_C$ such that every element of $T(\mathbf A_f^p)\times
T(\mathbf Q_p)$ is $\gamma su$ with $\gamma\in T(\mathbf Q)$, $s\in S$ and $u\in U$; since then
$s=\gamma^{-1}(\gamma su)u^{-1}$, we may take $S\subset T(\mathbf A_f^p)\times T(\mathbf Q_p)$. If some
$\gamma_-\in T(\mathbf Q)$ has $\sigma<0$, replace $S$ by $S\cup\gamma_-S$; then every element is $\gamma su$ with
$\gamma\in Z$, $s\in S$ and $u\in U$. Let $c=(a\times b)(\gamma su)$ commute with $(\delta,\delta_p)$. Then
\[
 cr=(\iota,\iota_p)(\gamma)\,(a\times b)(s)\,r\,\bigl(r^{-1}(a\times b)(u)r\bigr),
\]
and the last factor lies in $K\times K_p$. Finally $K_p$ is the union of finitely many cosets
$t(K'\cap K_p)$, so $R_r$ can be taken to be the set of the elements $(a\times b)(s)\,r\,(1,t)$.
\end{proof}

\begin{lemma}\label{gal-complex-points-conjugate}
In Situation~\ref{gal-count-situation}, let $K$, $f$, $K'$ and $f_p$ be as in
Lemma~\ref{gal-complex-fixed-finite}, and let $\beta$ be a form, with $x$ and $t$. Suppose that $P_y$ is
separable of degree $n$ whenever $(y,s)\in G_\beta$ fixes a vector positive for $\beta$ and
$f(h^{-1}\iota(y,s)h)\neq0$ for some $h\in G(\mathbf A_f^p)$. Let $(y_0,s_0)\in G_\beta$. Then only
finitely many points of $\beta$ of level $(K,K')$ have a representative $(g,g_p,z)$ for which some
conjugate $(y,s)$ of $(y_0,s_0)$ in $G_\beta$ fixes $z$ and satisfies
$f(g^{-1}\iota(y,s)g)\,f_p(g_p^{-1}\iota_p(y,s)g_p)\neq0$.
\end{lemma}
\leanref{Development/AbelianVarietyTheory/Proofs/DivisorProofs/OrbitalIntegralTestFunctionSecond.lean}{GlobalTheory.T.orbital_integral_test_function_second}
\leanref{Development/AbelianVarietyTheory/Proofs/DivisorProofs/FiniteCosetsAndFixedPoints.lean}{GlobalTheory.FinPtsSep.finite_pts_conj_fixL_sep}

\begin{proof}
For $\gamma=(e,\sigma)\in G_\beta$ write $\psi_v(\gamma)=\psi_v(e)$. As in the proof of
Lemma~\ref{gal-count-positive-line}, $\psi_v(\gamma)^\dagger H_\beta\psi_v(\gamma)=\sigma H_\beta$ with
$\sigma>0$, so $\psi_v(\gamma)$ maps vectors positive for $\beta$ to vectors positive for $\beta$; and if
$\gamma(y_0,s_0)\gamma^{-1}$ fixes $z$, then $(y_0,s_0)$ fixes $\psi_v(\gamma)^{-1}z$.

Suppose that some point as in the statement exists, with a representative $(g_1,g_{p,1},z_1)$ and
$\gamma_1\in G_\beta$ such that $\gamma_1(y_0,s_0)\gamma_1^{-1}$ fixes $z_1$. By hypothesis the minimal
polynomial of this conjugate, which is $P_{y_0}$, is separable of degree $n$, and
$z_0=\psi_v(\gamma_1)^{-1}z_1$ is positive and fixed by $(y_0,s_0)$. Let $R$ and $R_p$ be given by
Lemma~\ref{gal-complex-cosets}, for each $r\in R\times R_p$ let $R_r$ be given by
Lemma~\ref{gal-complex-torus}, and let $\Sigma$ be the union of the $R_r$, a finite set.

Let $(g,g_p,z)$ represent a point as in the statement, with $\gamma\in G_\beta$ such that
$\gamma(y_0,s_0)\gamma^{-1}$ fixes $z$ and the product of the values of $f$ and $f_p$ is nonzero. Since
$f\bigl((\iota(\gamma)^{-1}g)^{-1}\delta\,\iota(\gamma)^{-1}g\bigr)=f(g^{-1}\iota(\gamma(y_0,s_0)\gamma^{-1})g)$,
and in the same way at $p$, Lemma~\ref{gal-complex-cosets} gives
$(\iota(\gamma)^{-1}g,\iota_p(\gamma)^{-1}g_p)=crk$ with $c$ commuting with $(\delta,\delta_p)$,
$r\in R\times R_p$ and $k\in K\times K'$. By Lemma~\ref{gal-complex-torus}, $cr=(\iota,\iota_p)(\eta)\,q\,k'$
with $\eta\in Z$, $q\in\Sigma$ and $k'\in K\times K'$. Hence $(g,g_p)=(\iota,\iota_p)(\gamma\eta)\,q\,k'k$.
The vector $\psi_v(\gamma\eta)^{-1}z$ is positive and fixed by $(y_0,s_0)$, since $\eta$ commutes with
$(y_0,s_0)$, so by Lemma~\ref{gal-count-positive-line} it is $\lambda z_0$ for some $\lambda\in\mathbf
C^\times$. Thus $z=\lambda\psi_v(\gamma\eta)z_0$, and $(g,g_p,z)$ is equivalent to $(q,z_0)$. So the point
is one of the finitely many classes of the triples $(q,z_0)$ with $q\in\Sigma$.
\end{proof}

\begin{lemma}\label{gal-complex-classes-local}
In Situation~\ref{gal-count-situation}, let $K\subset G(\mathbf A_f^p)$ be an open compact subgroup and $f$
a test function for $K$, let $\beta$ be a form, with $x$ and $t$, and let $(y_0,s_0)\in G_\beta$ with
$P_{y_0}$ separable of degree $n$. Then there is a finite set $S$ of primes, containing $p$, with the
following property. Let $(y,s)\in G_\beta$ and $g\in G(\mathbf A_f^p)$ with $f(g^{-1}\iota(y,s)g)\neq0$, and
let $c\in C^\times$ with $y=cy_0c^{-1}$. Then for every prime $\ell\notin S$ there are
$u\in\mathbf Q_\ell^\times$ and $z\in(\mathbf Q_\ell\otimes C)^\times$ commuting with $1\otimes y_0$ such
that, in $\mathbf Q_\ell\otimes C$,
\[
 c^{-1}\beta(c^{-1})^*\beta^{-1}=u\,z\beta z^*\beta^{-1}.
\]
\end{lemma}
\leanref{Development/AbelianVarietyTheory/Proofs/DivisorProofs/HarrisTaylorRegularConjugateLocalScalarNormStatement.lean}{GlobalTheory.T.harrisTaylor_regular_conjugate_local_scalar_norm}
\leanref{Development/AbelianVarietyTheory/Proofs/DivisorProofs/HarrisTaylorRegularConjugateLocalScalarNorm.lean}{GlobalTheory.CxRegLoc.local_class_scalar_norm_sep}
\leanref{Development/AbelianVarietyTheory/Proofs/DivisorProofs/EventuallyKottwitzLatticeSeparableMinpolyStatement.lean}{GlobalTheory.T.eventually_kottwitz_lattice_separable_minpoly}
\leanref{Development/AbelianVarietyTheory/Proofs/DivisorProofs/EventuallyTauFixedUnitCentralizerIsNormStatement.lean}{GlobalTheory.T.eventually_tauFixed_unit_centralizer_is_norm}
\leanref{Development/AbelianVarietyTheory/Proofs/DivisorProofs/TwistedNormUnitCommutation.lean}{GlobalTheory.twisted_norm_local_step}

\begin{proof}
Write $\tau(e)=\beta e^*\beta^{-1}$, an anti-automorphism of $C$ and of each $\mathbf Q_\ell\otimes C$ with
$\tau^2=1$, as $\beta^*=\beta$. Let $E$ be the centralizer of $y_0$ in $C$. Over an algebraic closure of
$F$, $y_0$ is diagonalizable with $n$ distinct eigenvalues, so $E=F[y_0]$ is commutative and
$\mathbf Q_\ell\otimes E$ is the centralizer of $1\otimes y_0$. Since $y_0\beta y_0^*=s_0\beta$,
$\tau(y_0)=s_0y_0^{-1}$, so $\tau$ preserves these centralizers. Fix a basis of $C$ over $\mathbf Q$ and
let $\mathcal O_\ell$ be its $\mathbf Z_\ell$-span in $\mathbf Q_\ell\otimes C$. Choose $S$ such that for
$\ell\notin S$:
\begin{enumerate}
\item[(i)] $\mathcal O_\ell$ is a ring stable under the involution and contains $\beta$ and $\beta^{-1}$,
  and the components at $\ell$ of the elements of the support of $f$ lie in $G(\mathbf Z_\ell)$, the group
  of $(h,\nu)\in G(\mathbf Q_\ell)$ with $h,h^{-1}\in\mathcal O_\ell$ and $\nu\in\mathbf Z_\ell^\times$;
\item[(ii)] if $k\in(\mathbf Q_\ell\otimes C)^\times$ and $k(1\otimes y_0)k^{-1}$ and
  $k(1\otimes y_0)^{-1}k^{-1}$ lie in $\mathcal O_\ell$, then $k^{-1}=zm$ with $z$ commuting with
  $1\otimes y_0$ and $m,m^{-1}\in\mathcal O_\ell$;
\item[(iii)] if $m'\in(\mathbf Q_\ell\otimes C)^\times$ commutes with $1\otimes y_0$, $\tau(m')=m'$ and
  $m',m'^{-1}\in\mathcal O_\ell$, then $m'=v\tau(v)$ for some $v$ commuting with $1\otimes y_0$.
\end{enumerate}
Condition (i) holds for almost all $\ell$, since the support of $f$ is compact. Condition (ii) is
Kottwitz's lemma \cite[Proposition~7.1]{Kot86b} for the semisimple element $y_0$ of $C^\times$. For (iii): $E$ is a
product of number fields permuted by $\tau$, and for almost all $\ell$, $\mathcal O_\ell\cap(\mathbf
Q_\ell\otimes E)$ is the maximal order of $\mathbf Q_\ell\otimes E$, which is unramified over the
subalgebra fixed by $\tau$. On a pair of factors exchanged by $\tau$, every fixed element $(\alpha,\tau(\alpha))$
is $v\tau(v)$ with $v=(\alpha,1)$. On the other factors, a unit of the fixed subalgebra is the norm of a
unit, since the norm maps the units of an unramified extension of local fields onto the units
\cite[Chapter~V, \S2]{Ser79}.

Now let $(y,s)$, $g$ and $c$ be as in the statement and $\ell\notin S$. Write $g=(g_1,\nu)$ and let $k$ be
the component at $\ell$ of $g_1^{-1}xc$. Then $k(1\otimes y_0)k^{-1}$ is the component at $\ell$ of the
first coordinate of $g^{-1}\iota(y,s)g$, which lies in $G(\mathbf Z_\ell)$ by (i), so (ii) gives
$k^{-1}=zm$. Put $a=c^{-1}\tau(c^{-1})=c^{-1}\beta(c^{-1})^*\beta^{-1}$. Since $x\beta x^*=t$ and
$g_1g_1^*=\nu$, the element $ka\beta k^*$ is the component at $\ell$ of
$g_1^{-1}x\beta x^*(g_1^*)^{-1}=t\nu^{-1}$, a scalar $\lambda\in\mathbf Q_\ell^\times$. Hence
$a\beta=\lambda zmm^*z^*$, that is, $a=\lambda zm'\tau(z)$ with $m'=mm^*\beta^{-1}$.

The element $a$ commutes with $y_0$. Indeed, comparing determinants of left multiplication on $C$ in
$y\beta y^*=s\beta$ and $y_0\beta y_0^*=s_0\beta$, with $y=cy_0c^{-1}$ and $y^*=(c^*)^{-1}y_0^*c^*$, gives
$s^N=s_0^N$ for $N=\dim_{\mathbf Q}C$, so $s=s_0$. Then $\tau(y)=sy^{-1}$ and $\tau(y_0)=s_0y_0^{-1}$
give $\tau(c^{-1})y_0=s_0\tau(y_0^{-1}c^{-1})=s_0\tau(c^{-1}y^{-1})=y\,\tau(c^{-1})$, so
$ay_0=c^{-1}y\tau(c^{-1})=y_0a$. So $m'=\lambda^{-1}z^{-1}a\tau(z)^{-1}$ commutes with $1\otimes y_0$. Also
$\tau(m')=\beta\beta^{-1}mm^*\beta^{-1}=m'$, and $m'$ and $m'^{-1}=\beta(m^*)^{-1}m^{-1}$ lie in $\mathcal
O_\ell$ by (i). By (iii), $m'=v\tau(v)$, and since the centralizer of $1\otimes y_0$ is commutative and
$\tau$ reverses products, $a=\lambda(zv)\tau(zv)$. This is the claim, with $u=\lambda$ and $zv$ in place of
$z$.
\end{proof}

\begin{lemma}\label{gal-complex-classes-units}
In Situation~\ref{gal-count-situation}, let $\beta$ be a form and $(y_0,s_0)\in G_\beta$ with $P_{y_0}$
separable of degree $n$. Let $E$ be the centralizer of $y_0$ in $C$, write $\tau(e)=\beta e^*\beta^{-1}$,
and let $S$ be a finite set of primes. Then there is a finite set $A\subset E$ with the following property.
Let $a\in E^\times$ with $\tau(a)=a$ be such that for every prime $\ell\notin S$ there are
$u\in\mathbf Q_\ell^\times$ and $z\in(\mathbf Q_\ell\otimes C)^\times$ commuting with $1\otimes y_0$ with
$a=u\,z\tau(z)$. Then $a=r\,ea_0\tau(e)$ for some $r\in\mathbf Q^\times$, $e\in E^\times$ and $a_0\in A$.
\end{lemma}
\leanref{Development/AbelianVarietyTheory/Proofs/DivisorProofs/FiniteFixedUnitClassesCentralizerRegularStatement.lean}{GlobalTheory.T.finite_fixed_unit_classes_centralizer_regular}
\leanref{Development/AbelianVarietyTheory/Proofs/DivisorProofs/FiniteFixedUnitClassesCentralizerRegular.lean}{GlobalTheory.LKCx.ClsReg.finite_fixed_classes_regular}
\leanref{Development/AbelianVarietyTheory/Proofs/DivisorProofs/FiniteClassesFixedUnitsModRatNormsStatement.lean}{GlobalTheory.T.finite_classes_fixedUnits_mod_rat_norms}
\leanref{Development/AbelianVarietyTheory/Proofs/DivisorProofs/ExistsGlobalNormApproxOfLocalPadicNormStatement.lean}{GlobalTheory.T.exists_global_norm_approx_of_local_padic_norm}
\leanref{Development/AbelianVarietyTheory/Proofs/DivisorProofs/InvolutionNormClassRemoveRationalFactorStatement.lean}{GlobalTheory.T.involution_norm_class_remove_rational_factor}
\leanref{Development/AbelianVarietyTheory/Proofs/DivisorProofs/NormClassFiniteOfReducedRatAlgebraStatement.lean}{GlobalTheory.T.normClassFinite_of_reduced_rat_algebra}
\leanref{Development/AbelianVarietyTheory/Proofs/DivisorProofs/ExistsFiniteClassesModuloInvolutionNormsStatement.lean}{GlobalTheory.T.exists_finite_classes_modulo_involutionNorms}

\begin{proof}
As in the proof of Lemma~\ref{gal-complex-classes-local}, $E$ is commutative, $\tau$ restricts to an
automorphism of $E$ with $\tau^2=1$, and $\mathbf Q_\ell\otimes E$ is the centralizer of $1\otimes y_0$. Over
an algebraic closure of $F$, $E$ becomes the algebra of diagonal matrices, so $E$ is reduced, a finite
product of number fields. Say that an element $b$ of $E^\times$ is an $\ell$-unit if $b$ and $b^{-1}$ are
integral at every prime above $\ell$. Let $a$ be as in the statement.

Step 1. Let $\ell\notin S$, with $a=u\,z\tau(z)$. Since $\mathbf Q$ is dense in $\mathbf Q_\ell$ and $E$ in
$\mathbf Q_\ell\otimes E$, there are $c_\ell\in\mathbf Q^\times$ and $b_\ell\in E^\times$ so close to $u$ and
$z$ that $u/c_\ell$ and $z/b_\ell$ are units of the maximal order of $\mathbf Q_\ell\otimes E$. Then
$a/(c_\ell b_\ell\tau(b_\ell))$ is an $\ell$-unit.

Step 2. Let $T$ be the finite set of primes $\ell\notin S$ at which $a$ is not an $\ell$-unit, put
$q=\prod_{\ell\in T}\ell^{v_\ell(c_\ell)}$, and $a'=q^{-1}a$. For every $\ell\notin S$ there is
$b\in E^\times$ with $a'/(b\tau(b))$ an $\ell$-unit: for $\ell\in T$ take $b=b_\ell$, as $c_\ell/q$ is an
$\ell$-adic unit; for $\ell\notin S\cup T$ take $b=1$, as $a$ and $q$ are $\ell$-units. It suffices to find a
finite set $A$ such that every such $a'$ is $ea_0\tau(e)$ with $e\in E^\times$ and $a_0\in A$.

Step 3. Write $E$ as a product of number fields, which $\tau$ permutes. On a pair of factors exchanged by
$\tau$, every element fixed by $\tau$ is $(\alpha,\tau(\alpha))=v\tau(v)$ with $v=(\alpha,1)$. The
conditions of Step 2 hold factor by factor. So it suffices to prove the following for a number field $L$
with an automorphism $\tau$ such that $\tau^2=1$: the $a'\in L^\times$ with $\tau(a')=a'$ such that for
every $\ell\notin S$ there is $b_\ell\in L^\times$ with $a'/(b_\ell\tau(b_\ell))$ an $\ell$-unit lie in
finitely many classes for the relation $a'\sim ea'\tau(e)$, $e\in L^\times$.

Step 4. Let $a'$ be such an element. For the almost all $\ell$ at which $a'$ is an $\ell$-unit we may take
$b_\ell=1$. For a prime $\mathfrak P$ of $L$ over $\ell\notin S$,
$v_{\mathfrak P}(a')=v_{\mathfrak P}(b_\ell)+v_{\tau^{-1}\mathfrak P}(b_\ell)$, so the part of the divisor of
$a'$ outside $S$ is $\mathfrak D+\tau(\mathfrak D)$, where $\mathfrak D$ is the sum, over the finitely
many $\ell\notin S$ with $b_\ell\neq1$, of the part of the divisor of $b_\ell$ above $\ell$. The part of the
divisor of $a'$ above $S$ is fixed by $\tau$. The divisors above $S$ fixed by $\tau$, modulo those of the
form $\mathfrak E+\tau(\mathfrak E)$ with $\mathfrak E$ above $S$, form a finite group: a prime fixed by
$\tau$ contributes $\mathbf Z/2$, and a pair of primes exchanged by $\tau$ contributes nothing. Choose
representatives $\mathfrak F_1,\ldots,\mathfrak F_k$ of this group and divisors $\mathfrak D_1,\ldots,
\mathfrak D_h$ representing the ideal classes of $L$, finite in number \cite[Chapter~I, (6.3)]{Neu99}. Then
the divisor of $a'$ is $(\mathfrak D+\mathfrak E)+\tau(\mathfrak D+\mathfrak E)+\mathfrak F_j$ for some
$\mathfrak E$ above $S$ and some $j$, and $\mathfrak D+\mathfrak E=\mathfrak D_i+\operatorname{div}(e)$ for
some $i$ and $e\in L^\times$. So $a'/(e\tau(e))$, which is fixed by $\tau$, has the divisor
$\mathfrak D_i+\tau(\mathfrak D_i)+\mathfrak F_j$, one of finitely many. Two elements fixed by $\tau$ with
the same divisor differ by a unit $\varepsilon$ of $\mathcal O_L$ fixed by $\tau$. The unit group
$\mathcal O_L^\times$ is finitely generated \cite[Chapter~I, (7.4)]{Neu99}, so the units fixed by $\tau$
modulo the squares of such units form a finite group, and $\eta^2=\eta\tau(\eta)$ for $\eta$ fixed by
$\tau$. Hence $a'$ lies in one of finitely many classes.
\end{proof}

\begin{lemma}\label{gal-complex-classes}
In Situation~\ref{gal-count-situation}, let $K$, $f$, $K'$ and $f_p$ be as in
Lemma~\ref{gal-complex-fixed-finite}, and let $\beta$ be a form, with $x$ and $t$. Suppose that $P_y$ is
separable of degree $n$ whenever $(y,s)\in G_\beta$ fixes a vector positive for $\beta$ and
$f(h^{-1}\iota(y,s)h)\neq0$ for some $h\in G(\mathbf A_f^p)$. Let $P\in F[X]$. Then there are finitely
many elements of $G_\beta$ such that every $(y,s)\in G_\beta$ with $P_y=P$, for which there are
$(g,g_p)\in G(\mathbf A_f^p)\times G_p$ and a vector $z$ positive for $\beta$ and fixed by $(y,s)$ with
$f(g^{-1}\iota(y,s)g)\,f_p(g_p^{-1}\iota_p(y,s)g_p)\neq0$, is conjugate in $G_\beta$ to one of them.
\end{lemma}
\leanref{Development/AbelianVarietyTheory/Proofs/DivisorProofs/OrbitalIntegralTestFunctionFirst.lean}{GlobalTheory.T.orbital_integral_test_function_first}
\leanref{Development/AbelianVarietyTheory/Proofs/DivisorProofs/PadicEmbeddingValuationSubring.lean}{GlobalTheory.LKCx.ClsReg.classes_finite}
\leanref{Development/AbelianVarietyTheory/Proofs/DivisorProofs/ConjugatePairsSameSimilitude.lean}{GlobalTheory.grp_factor_eq_of_conj}
\leanref{Development/AbelianVarietyTheory/Proofs/DivisorProofs/TwistedConjugatesAlgebraRelation.lean}{GlobalTheory.Div.conj_of_twist}

\begin{proof}
We may assume that some element $(y_0,s_0)$ as in the statement exists. By hypothesis $P=P_{y_0}$ is
separable of degree $n$. Let $S$ be given by Lemma~\ref{gal-complex-classes-local} and then $A$ by
Lemma~\ref{gal-complex-classes-units}, for $(y_0,s_0)$; let $E$ be the centralizer of $y_0$ in $C$ and
$\tau(e)=\beta e^*\beta^{-1}$.

Let $(y,s)$ be as in the statement. The elements $y$ and $y_0$ are conjugate in $C^\times$: right
multiplication by $y$, or by $y_0$, makes $C$ a module over $C\otimes_FF[X]/(P)$, a semisimple algebra
whose simple factors each have one simple module, and the two modules are isomorphic, since their
dimensions on each factor agree after extending scalars to an algebraic closure of $F$; an isomorphism is
right multiplication by some $c$ with $yc=cy_0$. So $y=cy_0c^{-1}$ with $c\in C^\times$. As in the proof of
Lemma~\ref{gal-complex-classes-local}, $s=s_0$, and $a(c)=c^{-1}\tau(c^{-1})$ is an element of $E^\times$
with $\tau(a(c))=a(c)$. By Lemma~\ref{gal-complex-classes-local} it satisfies the hypothesis of
Lemma~\ref{gal-complex-classes-units}, so $a(c)=r\,ea_0\tau(e)$ with $r\in\mathbf Q^\times$,
$e\in E^\times$ and $a_0\in A$.

Let $(y',s')$ be a second such element, with $y'=c'y_0c'^{-1}$ and $a(c')=r'e'a_0\tau(e')$ for the same
$a_0$ and with $r'/r>0$. Put $w=e'e^{-1}\in E^\times$ and $\rho=r'/r$, so that $a(c')=\rho\,wa(c)\tau(w)$, and
put $\gamma=c'wc^{-1}$. Then $\gamma y\gamma^{-1}=c'wy_0w^{-1}c'^{-1}=y'$. Since $a(c)\beta=c^{-1}\beta(c^{-1})^*$
and $\beta w^*=\tau(w)\beta$,
\[
 \gamma\beta\gamma^*=c'wa(c)\tau(w)\beta c'^*=\rho^{-1}c'a(c')\beta c'^*=\rho^{-1}\beta.
\]
So $(\gamma,\rho^{-1})\in G_\beta$, and it conjugates $(y,s)$ to $(y',s')$, as $s=s'=s_0$. Hence the
conjugacy class of $(y,s)$ is determined by $a_0\in A$ and the sign of $r$, and there are at most $2|A|$
classes.
\end{proof}

\begin{lemma}\label{gal-complex-points-finite}
In Situation~\ref{gal-count-situation}, let $K\subset G(\mathbf A_f^p)$ be an open compact subgroup, $f$ a
test function for $K$ that vanishes outside the regular elements, $K'\subset G_p$ a subgroup commensurable
with $K_p$, and $f_p$ a test function for $K'$. Let $\beta$ be a form, with $x$ and $t$. Then only finitely
many points of $\beta$ of level $(K,K')$ have a representative $(g,g_p,z)$ for which some
$(y,s)\in G_\beta$ fixing $z$ satisfies $f(g^{-1}\iota(y,s)g)\,f_p(g_p^{-1}\iota_p(y,s)g_p)\neq0$.
\end{lemma}
\leanref{Development/AbelianVarietyTheory/Proofs/DivisorProofs/ClassSummandEqTermOfAOne.lean}{GlobalTheory.LKCx.CxReg.finite_pts_reg}

\begin{proof}
If $(y,s)\in G_\beta$ fixes a vector positive for $\beta$ and $f(h^{-1}\iota(y,s)h)\neq0$, then
$h^{-1}\iota(y,s)h$ is regular, so $P_y$ is separable of degree $n$ by
Lemma~\ref{gal-complex-separable}. So Lemmas~\ref{gal-complex-classes}
and~\ref{gal-complex-points-conjugate} apply. Let $(g,g_p,z)$ and $(y,s)$ be as in the statement. By
Lemma~\ref{gal-complex-polynomials}, $P_y$ lies in a finite set of polynomials. For each of them, by
Lemma~\ref{gal-complex-classes}, $(y,s)$ is conjugate in $G_\beta$ to one of finitely many elements
$(y_0,s_0)$. For each $(y_0,s_0)$, only finitely many points occur, by
Lemma~\ref{gal-complex-points-conjugate}.
\end{proof}

\begin{lemma}\label{gal-count-complex}
In Situation~\ref{gal-count-situation}, let $K\subset G(\mathbf A_f^p)$ be an open compact subgroup,
$f$ a test function for $K$ that vanishes outside the regular elements, $K'\subset G_p$ a subgroup
commensurable with $K_p$, and $f_p$ a test function for $K'$. Then for every function
$\eta:F[X]\times\mathbf C\to\mathbf C$,
\[
 T_v(\eta;f,f_p)=\sum_{[P]}\eta(P_\gamma,\mu_P)\,O_P(f,f_p),
\]
where $[P]$ runs over the equivalence classes of pairs that have a marking, with a representative $P$
chosen in each, and $\mu_P$ is the marking of $P$ (Lemma~\ref{gal-count-positive-line}).
\end{lemma}
\leanref{Development/AbelianVarietyTheory/Statements/PlaceShiftStatements/MarkWeightedFixedPointEqDefinitePairOdd.lean}{GlobalTheory.RootSumWeights.CxCountRegMark}
\leanref{Development/AutomorphicArithmetic/Foundations/EtaleClopenDescent/ComplexCountRegularMarkedPoints.lean}{UpperTheory.Cb.complex_count_regular_marked_points}

\begin{proof}
Fix a form $\beta$, with $x$ and $t$. The group $G_\beta$ acts on the triples $(g,g_p,z)$ by
$\delta\cdot(g,g_p,z)=(\iota(\delta)g,\iota_p(\delta)g_p,\psi_v(\delta)z)$, and the summand for
$(\delta y\delta^{-1},\delta\cdot(g,g_p,z))$ equals the summand for $(y,(g,g_p,z))$, because $f$ and
$f_p$ are invariant under conjugation by $K$ and $K'$ and the eigenvalue does not change. Grouping the
pairs $(y,(g,g_p,z))$ by the conjugacy class of $y$ in $G_\beta$, the contribution of $\beta$ to
$T_v(\eta;f,f_p)$ becomes
\[
 \sum_{[y]}\ \sum_{(g,g_p,z)}\frac{f(g^{-1}\iota(y)g)\,f_p(g_p^{-1}\iota_p(y)g_p)\,\eta(P_y,\mu_y(z))}
 {|\operatorname{Stab}(g,g_p,z)\cap Z_y|},
\]
where $Z_y$ is the centralizer of $y$ in $G_\beta$ and $(g,g_p,z)$ runs over the triples with $z$
fixed by $y$, modulo $Z_y$, $K\times K'$ and scalars. Both arrangements are sums with finitely many
nonzero terms: for a positive $z$ and $(g,g_p)$, only finitely many $y$ fixing $z$ have a nonzero
summand (Lemma~\ref{gal-complex-fixed-finite}), and only finitely many points carry such a $y$
(Lemma~\ref{gal-complex-points-finite}).
If $f(g^{-1}\iota(y)g)\neq0$, then
$g^{-1}\iota(y)g$ is regular, so $P_y$ is separable of degree $n$
(Lemma~\ref{gal-complex-separable}), and by
Lemma~\ref{gal-count-positive-line} the positive vectors fixed by $y$ form one line, on which $y$ has
one eigenvalue $\mu_y$. The inner sum is then $\eta(P_y,\mu_y)$ times the orbital sum of the pair
$(\beta,y,s)$ (Definition~\ref{gal-count-pairs}(3)). Finally, a pair $(\beta',\gamma,c)$ with a
marking has $c>0$ by Lemma~\ref{gal-count-positive-line}, and sending its class to the class of
$\beta'$ together with the conjugacy class in $G_\beta$ of $(\gamma,c)$, transported to the chosen
representative $\beta$ of that class, is a bijection from the classes of pairs with a marking onto
the pairs consisting of a class of forms and a conjugacy class of its group that fixes a positive
vector. Under it the marking of the pair is $\mu_y$, which gives the formula.
\end{proof}

\subsection{Ratios of forms and the weighted count}\label{subsection-gal-weighted-count}

\begin{definition}\label{gal-ratio-integral}
In Situation~\ref{gal-count-situation}, fix a basis $b_1,\ldots,b_N$ of $C$ over $\mathbf Q$. For a prime
$\ell$, let $C_\ell=\mathbf Q_\ell\otimes_{\mathbf Q}C$, with the involution $x\mapsto x^*$ extended
$\mathbf Q_\ell$-linearly, and regard $C$ as a subring of $C_\ell$. An element of $C_\ell$ is
\emph{$\ell$-integral} if its coordinates in the basis $b_1,\ldots,b_N$ lie in $\mathbf Z_\ell$, and
$\mathcal O_\ell\subset C_\ell$ is the set of $\ell$-integral elements. For $\beta\in C^\times$ with
$\beta^*=\beta$, put $\tau_\beta(x)=\beta x^*\beta^{-1}$ for $x$ in $C$ or in $C_\ell$.
\end{definition}
\leanref{Development/AbelianVarietyTheory/Proofs/DivisorProofs/IntCZero.lean}{GlobalTheory.FinClsLoc.IntC}
\leanref{Development/AbelianVarietyTheory/Proofs/DivisorProofs/CoordL.lean}{GlobalTheory.AwayOrb.coordL}
\leanref{Development/AbelianVarietyTheory/Proofs/DivisorProofs/CommutingUnitsIntegralCofinite.lean}{GlobalTheory.LocNormTau.tauR}

\begin{lemma}\label{gal-ratio-centralizer}
In Situation~\ref{gal-count-situation}, let $\gamma\in C$ have separable minimal polynomial $P_\gamma$ of
degree $n$ over $F$, and let $E\subset C$ be the centralizer of $\gamma$.
\begin{enumerate}
\item $E=F[\gamma]$. In particular $E$ is commutative and has no nonzero nilpotent elements, and
  $E\cong F[X]/(P_\gamma)$ is a finite product of number fields.
\item For every field $k$ of characteristic $0$, the map $k\otimes_{\mathbf Q}E\to k\otimes_{\mathbf Q}C$
  is injective, and its image is the centralizer of $1\otimes\gamma$.
\item Let $\beta\in C^\times$ with $\beta^*=\beta$ and $\gamma\beta\gamma^*=s\beta$ for some
  $s\in\mathbf Q$, and put $\tau_\beta(x)=\beta x^*\beta^{-1}$. Then $\tau_\beta$ maps $E$ into $E$, and
  its restriction is an automorphism of the $\mathbf Q$-algebra $E$ with $\tau_\beta\circ\tau_\beta=1$.
  Let $L\subset E$ be the subalgebra fixed by $\tau_\beta$. For every field $k$ of characteristic $0$,
  the elements of $k\otimes_{\mathbf Q}C$ that commute with $1\otimes\gamma$ and are fixed by the
  $k$-linear extension of $\tau_\beta$ are the elements of $k\otimes_{\mathbf Q}L$.
\end{enumerate}
\end{lemma}
\leanref{Development/AbelianVarietyTheory/Proofs/DivisorProofs/NatDegreeMinpolyLeBaseChange.lean}{GlobalTheory.CentSep.centralizer_eq_adjoin}
\leanref{Development/AbelianVarietyTheory/Proofs/DivisorProofs/UnitsConjOfMinpolyEqOfSeparable.lean}{GlobalTheory.LKCx.ClsReg.commute_of_commute_of_separable}
\leanref{Development/AbelianVarietyTheory/Proofs/DivisorProofs/EqZeroOfNilpotentCommuteSeparableMinpolyStatement.lean}{GlobalTheory.T.eq_zero_of_nilpotent_commute_separable_minpoly}
\leanref{Development/AbelianVarietyTheory/Proofs/DivisorProofs/TensorCentralizerInjectiveCommutantEqStatement.lean}{GlobalTheory.T.tensor_centralizer_injective_commutant_eq}
\leanref{Development/AbelianVarietyTheory/Proofs/DivisorProofs/FixedAlgebraExists.lean}{GlobalTheory.FixedCentraliserReg.fixed_algebra_exists}
\leanref{Development/AbelianVarietyTheory/Proofs/DivisorProofs/BaseChangeFixedCentraliserAlg.lean}{GlobalTheory.FinClsLoc.baseChange_fixedCentraliser_alg}

\begin{proof}
(1) Let $\bar F$ be an algebraic closure of $F$. Then $\bar F\otimes_FC\cong M_n(\bar F)$, and the powers
$1,\gamma,\ldots,\gamma^{n-1}$ stay linearly independent over $\bar F$, so the image $\bar\gamma$ of
$\gamma$ is a matrix whose minimal polynomial has degree $n$ and divides $P_\gamma$. Hence the
characteristic polynomial of $\bar\gamma$ is $P_\gamma$, which has $n$ distinct roots. In a basis of
eigenvectors of $\bar\gamma$, the matrices commuting with $\bar\gamma$ are the diagonal ones, and these
are the polynomials in $\bar\gamma$. So the centralizer of $\bar\gamma$ is
$\bar F[\bar\gamma]=\bar F\otimes_FF[\gamma]$. The centralizer of $\gamma$ is the kernel of the $F$-linear
map $y\mapsto y\gamma-\gamma y$, and kernels are compatible with extension of scalars, so
$\bar F\otimes_FE=\bar F\otimes_FF[\gamma]$, and $E=F[\gamma]$ because $F[\gamma]\subset E$. Since
$P_\gamma$ is the minimal polynomial of $\gamma$, $F[\gamma]\cong F[X]/(P_\gamma)$, which is a product of
fields because $P_\gamma$ is separable.

(2) The centralizer of $1\otimes\gamma$ is the kernel of the $k$-linear extension of the map
$y\mapsto y\gamma-\gamma y$ on $C$, whose kernel is $E$, and $k$ is flat over $\mathbf Q$.

(3) Since $\gamma$ and $\beta$ are units, so is $s\beta$, and $s\neq0$. From $\gamma\beta\gamma^*=s\beta$,
$\tau_\beta(\gamma)=\beta\gamma^*\beta^{-1}=s\gamma^{-1}$. The map $\tau_\beta$ is $\mathbf Q$-linear,
fixes $1$ and reverses products, and $\tau_\beta(\tau_\beta(y))=y$ because $\beta^*=\beta$. If $y$
commutes with $\gamma$, then $\tau_\beta(y)$ commutes with $\tau_\beta(\gamma)=s\gamma^{-1}$, hence with
$\gamma$. So $\tau_\beta$ maps $E$ into $E$, and on the commutative ring $E$ it preserves products. The
elements of $k\otimes_{\mathbf Q}C$ in question form the kernel of the $k$-linear extension of
$y\mapsto(y\gamma-\gamma y,\tau_\beta(y)-y)$, whose kernel on $C$ is $L$, so they form
$k\otimes_{\mathbf Q}L$ as in (2).
\end{proof}

\begin{lemma}\label{gal-ratio-order}
In Situation~\ref{gal-count-situation}, with a basis of $C$ fixed as in
Definition~\ref{gal-ratio-integral}, let $\beta\in C^\times$. For all but finitely many primes $\ell$,
products of $\ell$-integral elements of $C_\ell$ are $\ell$-integral, $x^*$ is $\ell$-integral whenever
$x$ is, $\beta$ and $\beta^{-1}$ are $\ell$-integral, and every $\ell$-integral element of $C_\ell$ is a
root of a monic polynomial with coefficients in $\mathbf Z_\ell$.
\end{lemma}
\leanref{Development/AbelianVarietyTheory/Proofs/DivisorProofs/IntegralConjugationStableStructure.lean}{GlobalTheory.FinClsLoc.intC_stab}
\leanref{Development/AbelianVarietyTheory/Proofs/DivisorProofs/EventuallyIsIntegralOfIntegralCoordsLocalTensorStatement.lean}{GlobalTheory.T.eventually_isIntegral_of_integralCoordsLocalTensor}

\begin{proof}
The coordinates of the products $b_ib_j$, of the elements $b_i^*$, and of $1$, $\beta$ and $\beta^{-1}$
in the basis are finitely many rational numbers, and they lie in $\mathbf Z_\ell$ for all but finitely
many $\ell$. For such $\ell$ the first three statements follow by expanding in the basis. Moreover
$\mathcal O_\ell$ is then a subring of $C_\ell$ that is finitely generated as a $\mathbf Z_\ell$-module,
so each of its elements is integral over $\mathbf Z_\ell$, by the Cayley--Hamilton theorem for
multiplication by the element on $\mathcal O_\ell$.
\end{proof}

\begin{lemma}\label{gal-ratio-split}
In Situation~\ref{gal-count-situation}, with a basis of $C$ fixed as in
Definition~\ref{gal-ratio-integral}, let $\gamma\in C$ have separable minimal polynomial of degree $n$
over $F$. For all but finitely many primes $\ell$ there are a commutative ring $A$, a subring $R$ of $A$
with finitely many maximal ideals, and a ring isomorphism $\Phi:M_n(A)\to C_\ell$, such that an element
$y$ of $C_\ell$ is $\ell$-integral if and only if the entries of $\Phi^{-1}(y)$ lie in $R$, and such that
the characteristic polynomial of $\Phi^{-1}(\gamma)$ is a polynomial $Q$ with coefficients in $R$ that
is separable over $R$, that is, $aQ+bQ'=1$ for some polynomials $a$ and $b$ over $R$.
\end{lemma}
\leanref{Development/AbelianVarietyTheory/Proofs/DivisorProofs/EventuallyCompletedTensorMatrixSplitSemilocalStatement.lean}{GlobalTheory.T.eventually_completedTensor_matrix_split_semilocal}
\leanref{Development/AbelianVarietyTheory/Proofs/DivisorProofs/EventuallyCompletedTensorMatrixSplitSemilocal.lean}{GlobalTheory.OrbFinAf.split_eventually}
\leanref{Development/AbelianVarietyTheory/Proofs/DivisorProofs/EventuallyIntegralMatrixUnitsAtPlacesStatement.lean}{GlobalTheory.T.eventually_integral_matrix_units_at_places}
\leanref{Development/AbelianVarietyTheory/Proofs/DivisorProofs/AdjoinLocalOrderTensorPadicMatrixAlmostAllStatement.lean}{GlobalTheory.T.adjoinLocalOrderTensorPadic_matrix_almost_all}

\begin{proof}
Let $F_\ell=\mathbf Q_\ell\otimes_{\mathbf Q}F$, the centre of $C_\ell$, and let $\Lambda$ be an order
of $C$, that is, a subring that is a finitely generated $\mathbf Z$-module spanning $C$ over $\mathbf Q$;
then $\mathcal O_\ell=\mathbf Z_\ell\Lambda$ for all but finitely many $\ell$. The order $\Lambda$ is
maximal at all but finitely many primes, $C$ is split at all but finitely many places of $F$, and a
maximal order in a matrix algebra over a local field is conjugate to the matrices over its valuation ring
\cite[Theorem~8.7, Corollary~27.6]{Rei75}. Hence for all but finitely many $\ell$, $\mathbf Z_\ell\Lambda$ is isomorphic to the
matrices over $\mathbf Z_\ell\otimes_{\mathbf Z}\mathcal O_F$, where $\mathcal O_F$ is the ring of
integers of $F$, and the images of the elementary matrices are $\ell$-integral matrix units
$e_{ij}\in C_\ell$, for $1\leq i,j\leq n$, with $e_{ij}e_{kl}=\delta_{jk}e_{il}$ and
$\sum_ie_{ii}=1$. They give a ring isomorphism $\Phi:M_n(F_\ell)\to C_\ell$,
$(y_{ij})\mapsto\sum_{i,j}y_{ij}e_{ij}$, that is the identity on $F_\ell$. Exclude also the
primes at which Lemma~\ref{gal-ratio-order} fails for $\beta=1$, take $A=F_\ell$, and let
$R=F_\ell\cap\mathcal O_\ell$. Then $R$ is a subring of $F_\ell$ that is finitely generated as a
$\mathbf Z_\ell$-module, so it has finitely many maximal ideals. The entries of $\Phi^{-1}(y)$ are
$y_{ij}=\sum_ke_{ki}\,y\,e_{jk}$, which lie in $R$ when $y$ is $\ell$-integral; conversely
$y=\sum y_{ij}e_{ij}$ is $\ell$-integral when the $y_{ij}$ are. Let $P_\gamma$ be the minimal
polynomial of $\gamma$, and $a$, $b$ polynomials over $F$ with $aP_\gamma+bP_\gamma'=1$. Excluding
finitely many more primes, the coefficients of $P_\gamma$, $a$ and $b$ are $\ell$-integral, so they lie
in $R$. Finally $P_\gamma$ is the characteristic polynomial $\chi$ of $\Phi^{-1}(\gamma)$: both are monic
of degree $n$ and vanish at $\gamma$, and $1,\gamma,\ldots,\gamma^{n-1}$ are linearly independent over
$F_\ell$, so $\chi-P_\gamma=0$. Take $Q=P_\gamma$.
\end{proof}

\begin{lemma}\label{gal-ratio-semilocal}
Let $R$ be a commutative ring with finitely many maximal ideals, and let $\Gamma,\Delta\in M_n(R)$ have
the same characteristic polynomial $Q$, with $aQ+bQ'=1$ for some polynomials $a$ and $b$ over $R$. Then
$U\Delta=\Gamma U$ for some $U\in\operatorname{GL}_n(R)$.
\end{lemma}
\leanref{Development/AbelianVarietyTheory/Proofs/DivisorProofs/SemilocalConjugateSameCharpoly.lean}{GlobalTheory.T.semilocal_conjugate_same_charpoly}
\leanref{Development/AbelianVarietyTheory/Proofs/DivisorProofs/FiniteOpenHopfOrder.lean}{GlobalTheory.SemilocalConj.conj_of_charpoly_eq_of_separable}
\leanref{Development/AbelianVarietyTheory/Proofs/DivisorProofs/FiniteOpenHopfOrder.lean}{GlobalTheory.SemilocalConj.exists_cyclic_of_finite_maximal}

\begin{proof}
For $v\in R^n$ let $K_\Gamma(v)$ be the matrix with columns $v,\Gamma v,\ldots,\Gamma^{n-1}v$. By the
Cayley--Hamilton theorem, $\Gamma K_\Gamma(v)=K_\Gamma(v)C_Q$, where $C_Q$ is the companion matrix of $Q$.
For each maximal ideal $\mathfrak m$ of $R$, the reduction of $\Gamma$ modulo $\mathfrak m$ has the
separable characteristic polynomial $Q\bmod\mathfrak m$, so it has $n$ distinct eigenvalues over an
algebraic closure of $R/\mathfrak m$, and it has a cyclic vector: some $v_{\mathfrak m}\in(R/\mathfrak
m)^n$ with $\det K_{\Gamma}(v_{\mathfrak m})\neq0$ modulo $\mathfrak m$. By the Chinese remainder theorem
there is $v\in R^n$ that reduces to $v_{\mathfrak m}$ for every $\mathfrak m$. Then $\det K_\Gamma(v)$
lies in no maximal ideal, so $K_\Gamma(v)$ is invertible. Choose $w$ in the same way for $\Delta$, and put
$U=K_\Gamma(v)K_\Delta(w)^{-1}$. Then $U\Delta=K_\Gamma(v)C_QK_\Delta(w)^{-1}=\Gamma U$.
\end{proof}

\begin{lemma}\label{gal-ratio-lattice}
In Situation~\ref{gal-count-situation}, with a basis of $C$ fixed as in
Definition~\ref{gal-ratio-integral}, let $\gamma\in C^\times$ have separable minimal polynomial of degree
$n$ over $F$. For all but finitely many primes $\ell$ the following holds. If $k\in C_\ell^\times$ and
$k\gamma k^{-1}$ and $k\gamma^{-1}k^{-1}$ are $\ell$-integral, then $k^{-1}=zm$ with $z\in C_\ell^\times$
commuting with $\gamma$ and $m\in C_\ell^\times$ such that $m$ and $m^{-1}$ are $\ell$-integral. The hypothesis that
$k\gamma^{-1}k^{-1}$ is $\ell$-integral is not used.
\end{lemma}
\leanref{Development/AbelianVarietyTheory/Proofs/DivisorProofs/EventuallyKottwitzLatticeSeparableMinpolyStatement.lean}{GlobalTheory.T.eventually_kottwitz_lattice_separable_minpoly}
\leanref{Development/AbelianVarietyTheory/Proofs/DivisorProofs/EventuallyKottwitzLatticeSeparableMinpoly.lean}{GlobalTheory.FinClsLoc.SepLat.kottwitz_lattice_sep}
\leanref{Development/AbelianVarietyTheory/Proofs/DivisorProofs/SiegelMeanValueNonnegative.lean}{GlobalTheory.OrbFinAf.KottLat.conj_integral}

\begin{proof}
Exclude the finitely many primes at which Lemma~\ref{gal-ratio-split} fails for $\gamma$ or at which
$\gamma$ is not $\ell$-integral, and let $A$, $R$, $\Phi$ and $Q$ be as in
Lemma~\ref{gal-ratio-split}. Let $k$ be as in the statement. The matrices $\Gamma=\Phi^{-1}(\gamma)$ and
$\Delta=\Phi^{-1}(k\gamma k^{-1})$ have entries in $R$, and both have characteristic polynomial $Q$,
since $\Delta$ is conjugate to $\Gamma$ in $M_n(A)$. By Lemma~\ref{gal-ratio-semilocal} there is
$U\in\operatorname{GL}_n(R)$ with $U\Delta=\Gamma U$. Put $m=\Phi(U)$ and $z=k^{-1}m^{-1}$. Then $m$ and
$m^{-1}$ are $\ell$-integral and $k^{-1}=zm$, and $m\,k\gamma k^{-1}=\gamma m$ gives
$z\gamma z^{-1}=k^{-1}m^{-1}\gamma mk=\gamma$.
\end{proof}

\begin{lemma}\label{gal-ratio-unit-norm-finite}
Let $S$ be a finite commutative ring and $c$ a ring endomorphism of $S$ with $c(c(x))=x$ for all $x$, and
suppose that $\omega-c(\omega)$ is a unit for some $\omega\in S$.
\begin{enumerate}
\item If $y\in S$ satisfies $y\,c(y)=1$, then $x=y\,c(x)$ for some unit $x$ of $S$.
\item Every unit $u$ of $S$ with $c(u)=u$ is $x\,c(x)$ for some $x\in S$.
\end{enumerate}
\end{lemma}
\leanref{Development/AbelianVarietyTheory/Proofs/SatakeParameters/FiniteInvolutionFixedUnitEqMulConj.lean}{GlobalTheory.finite_involution_hilbert90}
\leanref{Development/AbelianVarietyTheory/Proofs/SatakeParameters/FiniteInvolutionFixedUnitEqMulConj.lean}{GlobalTheory.finite_involution_fixed_unit_isNorm}
\leanref{Development/AbelianVarietyTheory/Proofs/SatakeParameters/FiniteInvolutionFixedUnitEqMulConj.lean}{GlobalTheory.exists_pos_pow_mul_pow_eq_pow_of_finite}

\begin{proof}
(1) The element $y$ is a unit with inverse $c(y)$. Put $z=(1+y)\,c(1+y)$. Since $S$ is finite, $e=z^M$ is
idempotent for some $M\geq1$, and $c(e)=e$. Put $x=e(1+y)+(1-e)(\omega+y\,c(\omega))$. Using $y\,c(y)=1$,
\[
 y\,c(x)=e(y+1)+(1-e)(y\,c(\omega)+\omega)=x.
\]
Suppose that $x$ lies in a maximal ideal $\mathfrak m$. Since $y\,c(1+y)=1+y$, $z\in\mathfrak m$ implies
$1+y\in\mathfrak m$. If $e\in\mathfrak m$, then $z\in\mathfrak m$, so $1+y\in\mathfrak m$ and $1-e\notin\mathfrak
m$; then $\omega+y\,c(\omega)\in\mathfrak m$, and $y(\omega-c(\omega))=\omega(1+y)-(\omega+y\,c(\omega))$
lies in $\mathfrak m$, which is impossible as it is a unit. If $e\notin\mathfrak m$, then $e(1-e)=0$ gives
$1-e\in\mathfrak m$, so $e(1+y)\in\mathfrak m$ and $1+y\in\mathfrak m$; then $z\in\mathfrak m$ and
$e\in\mathfrak m$, a contradiction. So $x$ is a unit.

(2) Let $N$ and $\psi$ be the endomorphisms $x\mapsto x\,c(x)$ and $x\mapsto x\,c(x)^{-1}$ of the finite
abelian group $S^\times$. The kernel of $\psi$ is the group of units fixed by $c$, and it contains the
image of $N$, since $c(x\,c(x))=x\,c(x)$. By (1), the kernel of $N$ is contained in the image of $\psi$.
Hence $|\ker\psi|=|S^\times|/|\operatorname{im}\psi|\leq|S^\times|/|\ker N|=|\operatorname{im}N|$, so the
image of $N$ is the kernel of $\psi$.
\end{proof}

\begin{lemma}\label{gal-ratio-unit-norm-hensel}
Let $R$ be a commutative ring and $I$ an ideal of $R$ such that $(R,I)$ is a \emph{Henselian pair}: $I$
is contained in every maximal ideal of $R$, and for every monic $f\in R[X]$ and $a_0\in R$ with $f(a_0)\in
I$ and $f'(a_0)$ a unit modulo $I$ there is $a\in R$ with $f(a)=0$ and $a-a_0\in I$. Suppose that $R/I$ is
finite, and let $c$ be a ring endomorphism of $R$ with $c(c(x))=x$ for all $x$ and $c(I)\subset I$, such
that $\omega-c(\omega)$ is a unit for some $\omega\in R$. Then every unit $u$ of $R$ with $c(u)=u$ is
$x\,c(x)$ for some $x\in R$.
\end{lemma}
\leanref{Development/AbelianVarietyTheory/Proofs/SatakeParameters/ExistsNormEqFixedUnitHenselian.lean}{GlobalTheory.HenselianNorm.exists_norm_eq}
\leanref{Development/AbelianVarietyTheory/Proofs/SatakeParameters/ExistsNormEqFixedUnitHenselian.lean}{GlobalTheory.HenselianNorm.exists_norm_eq_one_add}
\leanref{Development/AbelianVarietyTheory/Proofs/SatakeParameters/ExistsNormEqFixedUnitHenselian.lean}{GlobalTheory.HenselianNorm.exists_mem_add_mul_sq_eq}

\begin{proof}
Put $D=\omega-c(\omega)$ and $w=\omega D^{-1}$. Since $c(D)=-D$, $w+c(w)=(\omega-c(\omega))D^{-1}=1$; put
$n=w\,c(w)$, so that $c(n)=n$.

\emph{Elements $1+e$.} Let $e\in I$ with $c(e)=e$. The polynomial $X^2+X-ne$ at $a_0=0$ has value in $I$
and derivative $1$, so it has a root $s\in I$. Then $1+s$ is a unit, as $s$ lies in every maximal ideal,
and $t=e(1+s)^{-1}\in I$ satisfies $t+nt^2=e$: as $ne=s(1+s)$,
$t+nt^2=(e(1+s)+e\cdot ne)(1+s)^{-2}=e(1+s)^2(1+s)^{-2}=e$. Applying $c$,
$c(t)$ satisfies the same equation, so $(t-c(t))(1+n(t+c(t)))=0$; the second factor is a unit, since
$n(t+c(t))\in I$, so $c(t)=t$. Then $x=1+wt$ satisfies
$x\,c(x)=1+(w+c(w))t+nt^2=1+t+nt^2=1+e$.

\emph{The general case.} The endomorphism $c$ induces an endomorphism $\bar c$ of the finite ring $R/I$
with $\bar c\circ\bar c=1$, for which the image of $\omega-c(\omega)$ is a unit. By
Lemma~\ref{gal-ratio-unit-norm-finite}(2) there is $x_0\in R$ with $x_0c(x_0)-u\in I$. Put
$a=(x_0c(x_0)-u)u^{-1}$, which lies in $I$ and is fixed by $c$, as $c(u^{-1})=u^{-1}$. Then $V=1+a$ is a
unit fixed by $c$, with $x_0c(x_0)=uV$, and $V^{-1}=1+e$ with $e=-V^{-1}a\in I$ fixed by $c$. By the first
step $V^{-1}=y\,c(y)$ for some $y\in R$, and $x=x_0y$ satisfies $x\,c(x)=uVV^{-1}=u$.
\end{proof}

\begin{lemma}\label{gal-ratio-unit-norm}
Let $O$ be a Henselian local ring with finite residue field in which $2$ is a unit, $B$ a commutative
$O$-algebra, and $m$ and $d$ units of $B$ such that $m$, $m^{-1}$, $d$ and $d^{-1}$ are integral over
$O$. Then $m=a^2-db^2$ for some $a,b\in B$.
\end{lemma}
\leanref{Development/AbelianVarietyTheory/Proofs/DivisorProofs/IntegralMonicPolynomialOfUnitRelation.lean}{GlobalTheory.OddLocalNorm.exists_eq_of_isIntegral}
\leanref{Development/AbelianVarietyTheory/Proofs/DivisorProofs/IntegralMonicPolynomialOfUnitRelation.lean}{GlobalTheory.OddLocalNorm.exists_eq_of_finite}
\leanref{Development/AbelianVarietyTheory/Proofs/DivisorProofs/ExistsEqSqSubMulSq.lean}{GlobalTheory.OddLocalNormEtale.exists_eq_sq_sub_mul_sq}
\leanref{Development/AbelianVarietyTheory/Proofs/FoundationProofs/FiniteMaximalSpectrumOfFiniteOfIsLocalRing.lean}{GlobalTheory.henselianRing_of_finite}

\begin{proof}
Replacing $B$ by its $O$-subalgebra generated by $m$, $m^{-1}$, $d$ and $d^{-1}$, which is finite over
$O$ since these elements are integral, we may assume that $B$ is finite over $O$. Let $T=B[t]/(t^2-d)$,
with the involution $c(a+bt)=a-bt$ for $a,b\in B$, and let $I=\mathfrak m_OT$, where $\mathfrak m_O$ is the
maximal ideal of $O$. Then $T$ is finite over $O$, so $(T,I)$ is a Henselian pair in the sense of
Lemma~\ref{gal-ratio-unit-norm-hensel} \cite[Tags 09XI and 09XK]{Sta}; in particular every maximal ideal
of $T$ contains $I$. The ring $T/I$ is a finite-dimensional algebra over the finite residue field of $O$,
so it is finite, and $c(I)\subset I$ as $c$ is $O$-linear. For $\omega=t$, $\omega-c(\omega)=2t$, whose
square $4d$ is a unit, so $2t$ is a unit. By Lemma~\ref{gal-ratio-unit-norm-hensel}, applied to the unit
$m$ of $T$, which is fixed by $c$, there is $x=a+bt$ with $m=x\,c(x)=a^2-db^2$.
\end{proof}

\begin{remark}\label{gal-ratio-unit-norm-count}
Lemma~\ref{gal-ratio-unit-norm} also has a direct proof. With $B$ finite over $O$, $B$ is a finite product
of Henselian local rings, and it suffices to treat one factor $B_1$, with maximal ideal $\mathfrak n$ and
residue field $k_1$, a finite field of odd characteristic. The sets $\{x^2:x\in k_1\}$ and
$\{\bar m+\bar dy^2:y\in k_1\}$ each have $(|k_1|+1)/2$ elements, so they meet, and
$\bar m=\bar a_0^2-\bar d\bar b_0^2$ for some $\bar a_0,\bar b_0\in k_1$. Lift these to $a_0,b_0\in B_1$.
Then $w=a_0^2-db_0^2$ is a unit and $mw^{-1}\in1+\mathfrak n$. Since $2$ is a unit, Hensel's lemma for
$X^2-mw^{-1}$ gives $x\in B_1$ with $x^2=mw^{-1}$, and $m=(xa_0)^2-d(xb_0)^2$.
\end{remark}

\begin{lemma}\label{gal-ratio-local-norm}
In Situation~\ref{gal-count-situation}, with a basis of $C$ fixed as in
Definition~\ref{gal-ratio-integral}, let $\beta\in C^\times$ with $\beta^*=\beta$, and let
$\gamma\in C^\times$ have separable minimal polynomial of degree $n$ over $F$ and satisfy
$\gamma\beta\gamma^*=s\beta$ for some $s\in\mathbf Q$. For all but finitely many primes $\ell$ the
following holds. If $m\in C_\ell^\times$ commutes with $\gamma$, $m$ and $m^{-1}$ are $\ell$-integral,
and $\tau_\beta(m)=m$, then $m=v\,\tau_\beta(v)$ for some $v\in C_\ell^\times$ commuting with $\gamma$.
\end{lemma}
\leanref{Development/AbelianVarietyTheory/Proofs/DivisorProofs/EventuallyTauFixedUnitCentralizerIsNormStatement.lean}{GlobalTheory.T.eventually_tauFixed_unit_centralizer_is_norm}
\leanref{Development/AbelianVarietyTheory/Proofs/DivisorProofs/EventuallyTauFixedUnitCentralizerIsNorm.lean}{GlobalTheory.LocNormTau.exists_norm_of_tau_fixed_sep}

\begin{proof}
Since $F$ is a CM field, there is $j\in F^\times$ with $j^*=-j$. It is central in $C$, and
$\tau_\beta(j)=\beta j^*\beta^{-1}=-j$. Let $E$ be the centralizer of $\gamma$ and $L\subset E$ the
subalgebra fixed by $\tau_\beta$, as in Lemma~\ref{gal-ratio-centralizer}; then $j^2$ and $j^{-2}$ lie
in $L$. Exclude $\ell=2$, the primes at which Lemma~\ref{gal-ratio-order} fails for $\beta$, and the
primes at which $j^2$ or $j^{-2}$ is not $\ell$-integral. Let $m$ be as in the statement. Then $m^{-1}$
also commutes with $\gamma$ and is fixed by $\tau_\beta$, so $m$ and $m^{-1}$ lie in
$\mathbf Q_\ell\otimes_{\mathbf Q}L$ by Lemma~\ref{gal-ratio-centralizer}(3). By
Lemma~\ref{gal-ratio-order}, $m$, $m^{-1}$, $j^2$ and $j^{-2}$ are integral over $\mathbf Z_\ell$.
Lemma~\ref{gal-ratio-unit-norm}, with $O=\mathbf Z_\ell$, $B=\mathbf Q_\ell\otimes_{\mathbf Q}L$ and
$d=j^2$, gives $a,b\in\mathbf Q_\ell\otimes_{\mathbf Q}L$ with $m=a^2-j^2b^2$. Put $v=a+bj$. Since $a$
and $b$ commute with each other and with $\gamma$ and are fixed by $\tau_\beta$, and $j$ is central with
$\tau_\beta(j)=-j$, the element $v$ commutes with $\gamma$, $\tau_\beta(v)=a-bj$, and
$v\,\tau_\beta(v)=a^2-j^2b^2=m$. In particular $v$ is a unit, with inverse $\tau_\beta(v)m^{-1}$.
\end{proof}

\begin{lemma}\label{gal-ratio-local}
In Situation~\ref{gal-count-situation}, let $K\subset G(\mathbf A_f^p)$ be an open compact subgroup, $f$
a test function for $K$ that vanishes outside the regular elements, and $(\beta_1,\gamma,c)$ a pair
(Definition~\ref{gal-count-pairs}(1)). There is a finite set $S$ of primes, containing $p$, with the
following property. Let $(\beta,\gamma,c)$ be a pair with the same $\gamma$ and $c$, let
$x\in(\mathbf A_f^p\otimes C)^\times$ and $t\in(\mathbf A_f^p)^\times$ satisfy $x\beta x^*=t$, and let
$g\in G(\mathbf A_f^p)$ satisfy $f(g^{-1}(x\gamma x^{-1},c)g)\neq0$. Then for every prime $\ell\notin S$
there are $u\in\mathbf Q_\ell^\times$ and $z\in(\mathbf Q_\ell\otimes_{\mathbf Q}C)^\times$ commuting
with $\gamma$ such that $\beta\beta_1^{-1}=u\,z\beta_1z^*\beta_1^{-1}$ in
$\mathbf Q_\ell\otimes_{\mathbf Q}C$. The hypothesis that $f$ vanishes outside the regular elements is not used.
\end{lemma}
\leanref{Development/AbelianVarietyTheory/Proofs/DivisorProofs/KottwitzPairFormRatioLocalScalarNormStatement.lean}{GlobalTheory.T.kottwitzPair_form_ratio_local_scalar_norm}
\leanref{Development/AbelianVarietyTheory/Proofs/DivisorProofs/KottwitzPairFormRatioLocalScalarNorm.lean}{GlobalTheory.CxRegPair.local_pair_scalar_norm}
\leanref{Development/AbelianVarietyTheory/Proofs/DivisorProofs/KottwitzPairFormRatioLocalScalarNorm.lean}{GlobalTheory.CxRegPair.glue_pair}
\leanref{Development/AbelianVarietyTheory/Proofs/DivisorProofs/KottwitzPairFormRatioLocalScalarNorm.lean}{GlobalTheory.CxRegPair.commute_mul_inv_of_sim}
\leanref{Development/AbelianVarietyTheory/Proofs/DivisorProofs/TwistedNormUnitCommutation.lean}{GlobalTheory.twisted_norm_local_step}
\leanref{Development/AbelianVarietyTheory/Proofs/DivisorProofs/ExistsCompactSupportTwistedHeckeTestFunctionStatement.lean}{GlobalTheory.T.exists_compact_support_twistedHeckeTestFunction}
\leanref{Development/AbelianVarietyTheory/Proofs/DivisorProofs/ExistsHodgeLatticeGermSystem.lean}{GlobalTheory.AwayOrb.eventually_compL_mem_intPts}

\begin{proof}
Fix a basis of $C$ as in Definition~\ref{gal-ratio-integral}. The function $f$ vanishes outside finitely
many double cosets $KsK$, each compact, so it vanishes outside a compact set
$\Omega\subset G(\mathbf A_f^p)$. The elements of $G(\mathbf A_f^p)$ are pairs $(h,\nu)$ with
$h\in(\mathbf A_f^p\otimes C)^\times$, $\nu\in(\mathbf A_f^p)^\times$ and $hh^*=\nu$. Since $\Omega$ is
compact, for all but finitely many primes $\ell\neq p$ and every $(h,\nu)\in\Omega$, the
$\ell$-components $h_\ell$ and $h_\ell^{-1}$ are $\ell$-integral and $\nu_\ell\in\mathbf Z_\ell^\times$.
Let $S$ consist of $p$, of these finitely many exceptions, and of the primes at which
Lemma~\ref{gal-ratio-order} fails for $\beta_1$, Lemma~\ref{gal-ratio-lattice} fails for $\gamma$, or
Lemma~\ref{gal-ratio-local-norm} fails for $\beta_1$, $\gamma$ and $s=c$; these lemmas apply, since
$\beta_1^*=\beta_1$, $\gamma\beta_1\gamma^*=c\beta_1$ and $P_\gamma$ is separable of degree $n$.

Let $(\beta,\gamma,c)$, $x$, $t$ and $g$ be as in the statement, and let $\ell\notin S$. Write
$(h,\nu)$ for the $\ell$-component of $g$, and $x_\ell$ and $t_\ell$ for those of $x$ and $t$, and put
$k=h^{-1}x_\ell\in C_\ell^\times$ and $\lambda=t_\ell\nu^{-1}$. Then $k\gamma k^{-1}$ and
$k\gamma^{-1}k^{-1}$ are the first components of the $\ell$-components of $g^{-1}(x\gamma x^{-1},c)g$,
which lies in $\Omega$, and of its inverse, so they are $\ell$-integral. Since $hh^*=\nu$ is central,
$h^*h=\nu$, and
\[
 k\beta k^*=h^{-1}(x_\ell\beta x_\ell^*)(h^*)^{-1}=t_\ell(h^*h)^{-1}=\lambda.
\]
By Lemma~\ref{gal-ratio-lattice}, $k^{-1}=zm$ with $z$ commuting with $\gamma$ and $m$, $m^{-1}$
$\ell$-integral.

Put $a=\beta\beta_1^{-1}$ and $\tau=\tau_{\beta_1}$. The ratio $a$ commutes with $\gamma$, since
$\gamma a(\beta_1\gamma^*)=\gamma\beta\gamma^*=c\beta=a\,c\beta_1=a\gamma(\beta_1\gamma^*)$ and
$\beta_1\gamma^*$ is a unit. From $k\beta k^*=\lambda$,
\[
 a\beta_1=\beta=\lambda\,k^{-1}(k^*)^{-1}=\lambda\,z\,mm^*\,z^*,
\]
so $a=\lambda\,z\,m'\,\tau(z)$ with $m'=mm^*\beta_1^{-1}$, because $z^*\beta_1^{-1}=\beta_1^{-1}\tau(z)$.
Then $m'=\lambda^{-1}z^{-1}a\,\tau(z)^{-1}$ commutes with $\gamma$, because $\tau$ preserves the
centralizer of $\gamma$ in $C_\ell$ (Lemma~\ref{gal-ratio-centralizer}(2) and (3)). Also
$\tau(m')=\beta_1(\beta_1^{-1}mm^*)\beta_1^{-1}=m'$ because $\beta_1^*=\beta_1$, and $m'$ and
$m'^{-1}=\beta_1(m^*)^{-1}m^{-1}$ are $\ell$-integral by Lemma~\ref{gal-ratio-order}. By
Lemma~\ref{gal-ratio-local-norm}, $m'=v\,\tau(v)$ with $v\in C_\ell^\times$ commuting with $\gamma$.
Since $\tau$ reverses products, $z'=zv$ commutes with $\gamma$ and
\[
 a=\lambda\,zv\,\tau(v)\tau(z)=\lambda\,z'\tau(z')=\lambda\,z'\beta_1z'^*\beta_1^{-1}.
\]
Take $u=\lambda$.
\end{proof}

\begin{definition}\label{gal-count-finite-integral-degree}
In Situation~\ref{gal-count-situation}, regard the fields $F_{w_i}$ as subfields of an algebraic closure
$\bar{\mathbf Q}_p$ of $\mathbf Q_p$, and let $\sigma$ be an automorphism of $\bar{\mathbf Q}_p$ over
$\mathbf Q_p$ with $|\sigma(x)-x^p|<1$ for every $x$ with $|x|\leq1$. Let $r\geq1$.
\begin{enumerate}
\item Let $\mathbf Q_{p^r}\subset\bar{\mathbf Q}_p$ be the field generated over $\mathbf Q_p$ by the roots of
  unity of order dividing $p^r-1$, and $\mathbf Z_{p^r}$ the ring of its elements of absolute value at most
  $1$; $\sigma$ restricts to an automorphism of $\mathbf Q_{p^r}$. For each $i$ let
  $F_{w_i,r}=F_{w_i}\otimes_{\mathbf Q_p}\mathbf Q_{p^r}$, with the automorphism $1\otimes\sigma$, and let
  $\mathcal O_{w_i,r}$ be its subring generated by $\mathcal O_{w_i}\otimes1$ and $1\otimes\mathbf Z_{p^r}$. Put
  \[
   G_{p,r}=\mathbf Q_{p^r}^\times\times\prod_i\operatorname{GL}_n(F_{w_i,r}),\qquad
   K_{p,r}=\mathbf Z_{p^r}^\times\times\prod_i\operatorname{GL}_n(\mathcal O_{w_i,r}),
  \]
  let $\sigma$ act on $G_{p,r}$ through its action on the entries, and put
  $N_r(\delta)=\delta\,\sigma(\delta)\cdots\sigma^{r-1}(\delta)$ for $\delta\in G_{p,r}$. Let $G_p\to G_{p,r}$,
  written $h\mapsto h$, be the inclusion given by $\mathbf Q_p\subset\mathbf Q_{p^r}$ and
  $F_{w_i}\to F_{w_i,r}$, $a\mapsto a\otimes1$.
\item Let $K'\subset G_{p,r}$ be a subgroup. A \emph{test function of degree $r$} for $K'$ is a function
  $\varphi:G_{p,r}\to\mathbf C$ such that $K'$ is commensurable with $K_{p,r}$, $\sigma(K')\subset K'$,
  $\varphi(k^{-1}g\sigma(k))=\varphi(g)$ for $k\in K'$ and $g\in G_{p,r}$, and $\varphi$ vanishes outside
  finitely many double cosets $K_{p,r}sK_{p,r}$.
\item Let $K\subset G(\mathbf A_f^p)$ be an open compact subgroup, $f$ a test function for $K$,
  $K'\subset G_{p,r}$ a subgroup and $\varphi$ a test function of degree $r$ for $K'$. Let
  $P=(\beta,\gamma,c)$ be a pair (Definition~\ref{gal-count-pairs}(1)), with $x$ and $t$ as in
  Definition~\ref{gal-count-forms}(1) fixed, and $Z_P$, $\iota$ and $\iota_p$ as in
  Definition~\ref{gal-count-pairs}(3). If there are $\delta,z\in G_{p,r}$ with
  $N_r(\delta)=z^{-1}\iota_p(\gamma,c)z$, fix such $\delta$ and $z$, let $Z_P$ act on $G_{p,r}$ through
  $y\mapsto z^{-1}\iota_p(y)z$, and let the \emph{orbital sum of degree $r$} of $P$ be
  \[
   O^{(r)}_P(f,\varphi)=\sum_{(g,g')}\frac{f(g^{-1}\iota(\gamma,c)g)\,\varphi(g'^{-1}\delta\,\sigma(g'))}
   {|Z_P\cap(g,g')(K\times K')(g,g')^{-1}|},
  \]
  where $(g,g')$ runs over representatives of $Z_P\backslash(G(\mathbf A_f^p)\times G_{p,r})/(K\times K')$,
  with the conventions of Definition~\ref{gal-count-forms} for infinite sums and infinite groups. If there
  are no such $\delta$ and $z$, put $O^{(r)}_P(f,\varphi)=0$.
\end{enumerate}
For $r=1$ we have $\mathbf Q_{p^1}=\mathbf Q_p$, $G_{p,1}=G_p$ and $K_{p,1}=K_p$, and $\sigma$ acts trivially.
So a test function of degree $1$ for $K'$ is a test function for $K'$ in the sense of
Situation~\ref{gal-count-situation}, and $O^{(1)}_P(f,f_p)=O_P(f,f_p)$: here
$\delta=z^{-1}\iota_p(\gamma,c)z$, and $g'\mapsto zg'$ matches the terms of the two sums.
\end{definition}
\leanref{ReciprocityTheory/ClassFieldTheory/LocalClassField/UnrDefs.lean}{ReciprocityTheory.DiscreteGroups.LocalClassField.Kn}
\leanref{ReciprocityTheory/ClassFieldTheory/LocalClassField/ExistsIntEKMulVbar.lean}{ReciprocityTheory.DiscreteGroups.LocalClassField.frobK}
\leanref{Development/AbelianVarietyTheory/Definitions/Divisors/TwOrbMEqTwOrbUOfNormal.lean}{GlobalTheory.LKRes.GHT}
\leanref{Development/AbelianVarietyTheory/Definitions/Divisors/TwOrbMEqTwOrbUOfNormal.lean}{GlobalTheory.LKRes.sigHT}
\leanref{Development/AbelianVarietyTheory/Definitions/Divisors/TwOrbMEqTwOrbUOfNormal.lean}{GlobalTheory.LKRes.KHT}
\leanref{Development/AbelianVarietyTheory/Definitions/Divisors/ExistsRootSelectionUnitaryRankOne.lean}{GlobalTheory.LKMid.nrmG}
\leanref{Development/AbelianVarietyTheory/Definitions/Divisors/ExistsRootSelectionUnitaryRankOne.lean}{GlobalTheory.LKMid.inclG}
\leanref{Development/AbelianVarietyTheory/Definitions/Divisors/ExistsRootSelectionUnitaryRankOne.lean}{GlobalTheory.LKMid.IsNrm}
\leanref{Development/AbelianVarietyTheory/Definitions/Divisors/ExistsRootSelectionUnitaryRankOne.lean}{GlobalTheory.LKMid.IsTest}
\leanref{Development/AbelianVarietyTheory/Definitions/Divisors/ExistsRootSelectionUnitaryRankOne.lean}{GlobalTheory.LKMid.Pair.toQr}
\leanref{Development/AbelianVarietyTheory/Definitions/Divisors/ExistsRootSelectionUnitaryRankOne.lean}{GlobalTheory.LKMid.Pair.termAt}
\leanref{Development/AbelianVarietyTheory/Definitions/Divisors/ExistsRootSelectionUnitaryRankOne.lean}{GlobalTheory.LKMid.Pair.term}

\begin{lemma}\label{gal-count-finite-integral}
In Situation~\ref{gal-count-situation}, let $K\subset G(\mathbf A_f^p)$ be an open compact subgroup,
$f$ a test function for $K$, $r\geq1$, $K'\subset G_{p,r}$ a subgroup, and $\varphi$ a test function of
degree $r$ for $K'$ (Definition~\ref{gal-count-finite-integral-degree}). Then there is an integer
$N\geq1$ such that $N\gamma$ and $N\gamma^{-1}$ are integral over $\mathbf Z$ for every pair
$P=(\beta,\gamma,c)$ that has a marking and satisfies $O^{(r)}_P(f,\varphi)\neq0$. The hypothesis that $P$
has a marking is not used.
\end{lemma}
\leanref{Development/AbelianVarietyTheory/Proofs/DivisorProofs/FiniteMarkedPairClassesDegreeOneTermNeZero.lean}{GlobalTheory.LKCx.FiniteClassIntegralRoots.integral_of_term_ne_zero'}
\leanref{Development/AbelianVarietyTheory/Proofs/DivisorProofs/FiniteMarkedPairClassesDegreeOneTermNeZero.lean}{GlobalTheory.LKCx.FiniteClassIntegralRoots.root_norm_le'}
\leanref{Development/AbelianVarietyTheory/Proofs/DivisorProofs/CompactConjCharpolyBoundedDenominatorsAwayPStatement.lean}{GlobalTheory.T.compact_conj_charpoly_bounded_denominators_awayP}
\leanref{Development/AbelianVarietyTheory/Proofs/DivisorProofs/TwistedOrbitNormBoundsStatement.lean}{GlobalTheory.T.twisted_orbit_norm_bounds}
\leanref{Development/AbelianVarietyTheory/Proofs/DivisorProofs/PadicCoeffDenominatorBoundOfBoundedRootsStatement.lean}{GlobalTheory.T.padic_coeff_denominator_bound_of_bounded_roots}
\leanref{Development/AbelianVarietyTheory/Proofs/DivisorProofs/IsIntegralSmulOfCharpolyCoeffMulIntStatement.lean}{GlobalTheory.T.isIntegral_natCast_smul_of_mul_charpoly_coeff_int}

\begin{proof}
For $y\in C$ let $\chi_y\in\mathbf Q[X]$ be the characteristic polynomial of left multiplication by $y$
on $C$, monic of degree $M=\dim_{\mathbf Q}C$. The function $f$ vanishes outside finitely many double
cosets $KsK$, so outside a compact set $\Omega$. As in the proof of
Lemma~\ref{gal-complex-polynomials}(a), the coefficients of the characteristic polynomial of left
multiplication on $\mathbf A_f^p\otimes C$ are continuous on $\Omega\cup\Omega^{-1}$, so there is an integer
$D\geq1$ with the following property: if $h\in\Omega\cup\Omega^{-1}$ has first coordinate
$x(1\otimes y)x^{-1}$ with $y\in C$, then every coefficient $\alpha$ of $\chi_y$ satisfies
$Dp^j\alpha\in\mathbf Z$ for some $j\geq0$.

At $p$, map $F_{w_i,r}$ to $\bar{\mathbf Q}_p$ by $a\otimes b\mapsto ab$, and $G_{p,r}$ to matrices over
$\bar{\mathbf Q}_p$ by applying these maps to the entries. Every element of $K_{p,r}$ has entries of absolute
value at most $1$, and $\sigma$ preserves $K_{p,r}$. The function $\varphi$ vanishes outside $K_{p,r}sK_{p,r}$
for $s$ in a finite set $S$. Choose $B\geq1$ bounding the absolute values of the entries of $\sigma^i(s)$ and
$\sigma^i(s^{-1})$ for $s\in S$ and $0\leq i<r$, and an integer $k\geq0$ with $B^r\leq p^k$. We claim: if
$h=(s_0,(g_i)_i)\in G_p$ and $\delta,z,y\in G_{p,r}$ satisfy $N_r(\delta)=z^{-1}hz$ and
$\varphi(y^{-1}\delta\sigma(y))\neq0$, then $|s_0|_p\leq p^k$, $|s_0^{-1}|_p\leq p^k$, and every eigenvalue of
$g_i$ or $g_i^{-1}$ in $\bar{\mathbf Q}_p$ has absolute value at most $p^k$. Indeed, write
$y^{-1}\delta\sigma(y)=k_1sk_2$ with $k_1,k_2\in K_{p,r}$ and $s\in S$. As $\sigma$ restricts to a generator of
the Galois group of $\mathbf Q_{p^r}$ over $\mathbf Q_p$, $\sigma^r$ is the identity on $G_{p,r}$, so
$N_r(y^{-1}\delta\sigma(y))=y^{-1}N_r(\delta)\sigma^r(y)=(zy)^{-1}h(zy)$. This element is the product of the $r$
elements $\sigma^i(k_1)\sigma^i(s)\sigma^i(k_2)$, whose entries have absolute value at most $B$ by the
ultrametric inequality, so its entries have absolute value at most $B^r$. Its inverse is the product of the
$\sigma^i(k_2^{-1})\sigma^i(s^{-1})\sigma^i(k_1^{-1})$ in the reverse order, with the same bound. The first
components commute, so the first component of $N_r(y^{-1}\delta\sigma(y))$ is $s_0$, and that of its inverse
is $s_0^{-1}$. For each $i$, the image of the component $g_i\otimes1$ of $h$ is $g_i$, conjugate over
$\bar{\mathbf Q}_p$ to the image of the $i$-th component of $N_r(y^{-1}\delta\sigma(y))$. An eigenvalue $\mu$ of
a matrix over $\bar{\mathbf Q}_p$ with entries of absolute value at most $b$ satisfies $|\mu|\leq b$: take an
eigenvector whose largest coordinate has absolute value $1$. This proves the claim. Put $N=Dp^{2kM}$.
% Lean: the claim is T.twisted_orbit_norm_bounds = AtPBd.main (GB.exists for B, nrmG_twConj and sig_iter_self for
% the norm of the twisted conjugate, GB.list_prod for the product, bound_of_conj with norm_root_le for the
% eigenvalues).

Let $P=(\beta,\gamma,c)$ be a pair with $O^{(r)}_P(f,\varphi)\neq0$. Then the elements $\delta$ and $z$ of
Definition~\ref{gal-count-finite-integral-degree}(3) exist, and some term of the orbital sum is nonzero, so
there are $x$ and $t$ as in Definition~\ref{gal-count-forms}(1) and $(g,g')$ with
$f(g^{-1}\iota(\gamma,c)g)\neq0$ and $\varphi(g'^{-1}\delta\sigma(g'))\neq0$. The first coordinate of
$g^{-1}\iota(\gamma,c)g\in\Omega$ is conjugate to $1\otimes\gamma$, and that of its inverse to
$1\otimes\gamma^{-1}$; so the coefficients of $\chi_\gamma$ and of $\chi_{\gamma^{-1}}$ have denominators
dividing $D$ times powers of $p$. At $p$, the claim applies to $h=\iota_p(\gamma,c)=(c,(\rho_i(\gamma))_i)$,
$\delta$, $z$ and $y=g'$.
Let $\mu\in\bar{\mathbf Q}_p$ be a root of $\chi_\gamma$. It is a root of the minimal polynomial of $\gamma$
over $\mathbf Q$, hence of the image of $P_\gamma$ under some embedding of $F$ into $\bar{\mathbf Q}_p$. If that
embedding lies over $u$, $\mu$ is an eigenvalue of some $\rho_i(\gamma)$, and $|\mu|\leq p^k$. Otherwise
its composite with $c$ lies over $u$, and $\mu$ is a root of the image of the minimal polynomial of
$\gamma^*$ under it, hence an eigenvalue of some
$\rho_i(\gamma^*)=c\rho_i(\beta)^{-1}\rho_i(\gamma)^{-1}\rho_i(\beta)$, as
$\gamma^*=c\beta^{-1}\gamma^{-1}\beta$. So $\mu=c\nu$ with $\nu$ an eigenvalue of $\rho_i(\gamma)^{-1}$, and
$|\mu|\leq p^k|c|_p\leq p^{2k}$. Since all roots of the monic polynomial $\chi_\gamma$ of degree $M$ have
absolute value at most $p^{2k}$, its coefficients $\alpha$ satisfy $|p^{2kM}\alpha|_p\leq1$. Together with
the bound away from $p$, $N\alpha\in\mathbf Z$ for every coefficient $\alpha$ of $\chi_\gamma$. Then
$N\gamma$ is a root of the monic polynomial $N^M\chi_\gamma(X/N)$, whose coefficient of $X^j$ is
$N^{M-j}$ times that of $\chi_\gamma$, an integer for $j<M$; so $N\gamma$ is integral over $\mathbf Z$.
The same argument for $(\gamma^{-1},c^{-1})$, which
satisfies $\gamma^{-1}\beta(\gamma^{-1})^*=c^{-1}\beta$, shows that $N\gamma^{-1}$ is integral.
\end{proof}

\begin{lemma}\label{gal-count-finite}
In Situation~\ref{gal-count-situation}, let $K\subset G(\mathbf A_f^p)$ be an open compact subgroup,
$f$ a test function for $K$, $K'\subset G_p$ a subgroup commensurable with $K_p$, and $f_p$ a test
function for $K'$. Then only finitely many equivalence classes of pairs that have a marking satisfy
$O_P(f,f_p)\neq0$.
\end{lemma}
\leanref{Development/AbelianVarietyTheory/Statements/PlaceShiftStatements/FinitelyManyMarkedClassesTermDegreeOne.lean}{GlobalTheory.RootSumWeights.FinitelyManyMarkedClassesTermDegreeOne}
\leanref{Development/AbelianVarietyTheory/Proofs/DivisorProofs/FinitelyManyMarkedClassesTermDegreeOne.lean}{GlobalTheory.RootSumWeights.finitely_many_marked_classes_term_degree_one}
\leanref{Development/AbelianVarietyTheory/Proofs/DivisorProofs/FiniteMarkedPairClassesDegreeOneTermNeZero.lean}{GlobalTheory.LKCx.FiniteClassIntegralRoots.finite_cls_term_one_a_one}
\leanref{Development/AbelianVarietyTheory/Proofs/DivisorProofs/FiniteKottwitzClassesPerMinpolyAEqOne.lean}{GlobalTheory.LKCx.PeriodicClassFixedPoints.finite_fixed}
\leanref{Development/AbelianVarietyTheory/Proofs/DivisorProofs/FiniteMarkedPairClassesDegreeOneTermNeZero.lean}{GlobalTheory.LKCx.FiniteClassIntegralRoots.finite_pol_marked'}
\leanref{Development/AbelianVarietyTheory/Proofs/DivisorProofs/FinitePolynomialsBoundedDenom.lean}{GlobalTheory.LKMid.Northcott.norm_le_of_isIntegral_inv}
\leanref{Development/AbelianVarietyTheory/Proofs/DivisorProofs/FiniteMonicRootsBddMulIsIntegralStatement.lean}{GlobalTheory.T.finite_monic_of_roots_bdd_of_mul_isIntegral}

\begin{proof}
For a pair $P$, every conjugate of $\iota(\gamma,c)$ is regular, since $P_\gamma$ is separable of
degree $n$; so replacing $f$ by its product with the indicator of the regular elements, again a test
function for $K$, changes no orbital sum, and we may assume that $f$ vanishes outside the regular
elements. Let $N$ be given by Lemma~\ref{gal-count-finite-integral} for $r=1$ and $\varphi=f_p$, a test
function of degree $1$ for $K'$ with $O^{(1)}_P(f,f_p)=O_P(f,f_p)$
(Definition~\ref{gal-count-finite-integral-degree}), and let $P=(\beta,\gamma,c)$ be a
pair with a marking $\mu$ and $O_P(f,f_p)\neq0$, so that $N\gamma$ and $N\gamma^{-1}$ are integral over
$\mathbf Z$. Let $\lambda$ be a root of $P_\gamma$ under a complex embedding of $F$. Replacing $\lambda$
by its complex conjugate if necessary, the embedding is the one written $w$ of an infinite place $w$, and
$\lambda$ is an eigenvalue of $\psi_w(\gamma)$. By Lemma~\ref{gal-count-positive-line}, $c>0$, so
$(\gamma,c)\in G_\beta$, and it fixes the positive eigenvector of the marking; as in step (d) of the proof
of Lemma~\ref{gal-complex-polynomials}, $|\lambda|^2=c$. The numbers $N\lambda$ and $N\lambda^{-1}$ are
eigenvalues of $\psi_w(N\gamma)$ and $\psi_w(N\gamma^{-1})$, hence algebraic integers. So $N\lambda^{-1}$
times its complex conjugate, which is $N^2/c$, is a positive rational algebraic integer, hence at least
$1$, and $|\lambda|\leq N$. Thus $P_\gamma$ is a monic polynomial of degree $n$ over $F$ whose roots under every complex
embedding have absolute value at most $N$ and become algebraic integers when multiplied by $N$. The
coefficients of $N^nP_\gamma$ are then algebraic integers of $F$ bounded under every complex embedding, so
$P_\gamma$ lies in a finite set of polynomials. Fix such a polynomial $P_0$. Two elements of $C$ with the same separable minimal polynomial of degree
$n$ are conjugate in $C^\times$, so up to equivalence $\gamma$ is a fixed $\gamma_0$ and $c=\pm c_0$.
Fix one such pair $(\beta_0,\gamma_0,c)$ with a marking. Then $c>0$ by
Lemma~\ref{gal-count-positive-line}, so $(\gamma_0,c)$ lies in the group $G_{\beta_0}$ of the form
$\beta_0$. Let $E$ be the centralizer of $\gamma_0$ in $C$ and $\tau(e)=\beta_0e^*\beta_0^{-1}$. Let $S$
be the finite set of primes of Lemma~\ref{gal-ratio-local} for $f$ and $(\beta_0,\gamma_0,c)$, and $A$
the finite set of Lemma~\ref{gal-complex-classes-units} for $\beta_0$, $(\gamma_0,c)$ and $S$. Let
$P=(\beta,\gamma_0,c)$ be another pair with $O_P(f,f_p)\neq0$, so that $f(g^{-1}(x\gamma_0x^{-1},c)g)\neq0$ for some $g$ and some $x$
and $t$ as in Definition~\ref{gal-count-forms}(1). The ratio $a=\beta\beta_0^{-1}$ is a unit commuting
with $\gamma_0$, as in the proof of Lemma~\ref{gal-ratio-local}, so $a\in E^\times$; and $\tau(a)=a$
because $\beta^*=\beta$ and $\beta_0^*=\beta_0$. By Lemma~\ref{gal-ratio-local}, for every prime
$\ell\notin S$, $a=u\,z\tau(z)$ in $\mathbf Q_\ell\otimes_{\mathbf Q}C$ for some
$u\in\mathbf Q_\ell^\times$ and some unit $z$ commuting with $\gamma_0$. By
Lemma~\ref{gal-complex-classes-units}, $a=r\,e\,a_0\,\tau(e)$ with $a_0\in A$, $r\in\mathbf Q^\times$ and
$e\in E^\times$, so $\beta=a\beta_0=r\,e\,(a_0\beta_0)\,e^*$.
% Lean: T.kottwitzPair_form_ratio_local_scalar_norm, T.finite_fixed_unit_classes_centralizer_regular,
% commute_ratio, tau_ratio and rel_of_same in LKCx.PeriodicClassFixedPoints.finite_fixed. The integer
% N is integral_of_term_ne_zero' and the finiteness of the P_gamma is finite_pol_marked'.
If a second such pair $(\beta',\gamma_0,c)$ gives the same $a_0$, with $r'$ and $e'$, and $r'/r>0$, then
$\beta'=(r'/r)\,y\beta y^*$ with $y=e'e^{-1}$, which commutes with $\gamma_0$, so the two pairs are
equivalent. So only finitely many classes remain.
\end{proof}

\begin{lemma}\label{gal-count-definite-signs}
Let $M$ be a number field with an automorphism $\sigma\neq1$ such that $\tau\circ\sigma=\bar\tau$ for
every embedding $\tau:M\to\mathbf C$, let $p$ be a prime, and let $S$ be a set of embeddings
$M\to\mathbf C$, stable under complex conjugation, whose number of elements is divisible by $4$. Then
there are $a\in M^\times$ with $\sigma(a)=a$, $\tau(a)<0$ for $\tau\in S$ and $\tau(a)>0$ for
$\tau\notin S$, and $u\in(\mathbf A_f^p\otimes M)^\times$ with $u\,\sigma(u)=a$.
\end{lemma}
\leanref{Development/AbelianVarietyTheory/Proofs/DivisorProofs/ExistsTotallyRealEvenNegPlacesLocalNormStatement.lean}{GlobalTheory.T.exists_totallyReal_even_neg_places_local_norm}
\leanref{Development/AbelianVarietyTheory/Proofs/DivisorProofs/ExistsRealSubfieldNegativeTwoPlacesNormEverywherStatement.lean}{GlobalTheory.T.exists_realSubfield_negative_twoPlaces_normEverywhere}

\begin{proof}
For disjoint sets, the products of solutions are solutions for the union, so by induction on the
number of elements it suffices to treat $S=\{\varepsilon,\bar\varepsilon,\varepsilon',\bar\varepsilon'\}$
with $\varepsilon'\neq\varepsilon,\bar\varepsilon$. Let $M^+$ be the fixed field of $\sigma$, so that
$M=M^+(\sqrt\delta)$ with $\delta\in M^+$ negative at every real place. The two real places of $M^+$
under $\varepsilon$ and $\varepsilon'$ form an even set of places at which $\delta$ is not a square,
so by the existence theorem for Hilbert symbols \cite[71:19a]{OMe63} there is $a\in M^+$ with
$(a,\delta)_w=-1$ at these two places and $(a,\delta)_w=1$ at every other place $w$ of $M^+$. At a
real place, $(a,\delta)_w=-1$ if and only if $a<0$ there, which gives the signs. At a finite place,
$(a,\delta)_w=1$ means that $a$ is a norm from $M\otimes_{M^+}M^+_w$; at the almost all places where
$a$ is a unit and $M/M^+$ is unramified, $a$ is the norm of a unit, and these local elements form $u$.
\end{proof}

\begin{lemma}\label{gal-count-definite-move}
In Situation~\ref{gal-count-situation}, suppose that $n$ is odd, and let $r$ assign to every monic
polynomial $P\in F[X]$ of degree $n$ a root $r(P)$ of $v(P)$. Then there is a bijection $\Phi$ from the
set of equivalence classes of pairs $P=(\beta,\gamma,c)$ with $P_\gamma$ irreducible and $r(P_\gamma)$ a
marking of $P$ onto the set of equivalence classes of definite pairs with irreducible $P_\gamma$, such
that $P$ and $\Phi(P)$ have the same minimal polynomial $P_\gamma$ and
$O_{\Phi(P)}(f,f_p)=O_P(f,f_p)$ for every open compact subgroup $K\subset G(\mathbf A_f^p)$, subgroup
$K'\subset G_p$, and functions $f$ and $f_p$ invariant under conjugation by $K$ and $K'$.
\end{lemma}
\leanref{Development/AbelianVarietyTheory/Proofs/DivisorProofs/MarkedFieldPairsEquivDefiniteFieldPairsOddStatement.lean}{GlobalTheory.T.markedFieldPairs_equiv_definiteFieldPairs_odd}
\leanref{Development/AbelianVarietyTheory/Proofs/DivisorProofs/ExistsPolynomialMoveDefiniteSimilitudePairMarked.lean}{GlobalTheory.LKDef.move_definite}
\leanref{Development/AbelianVarietyTheory/Proofs/DivisorProofs/KottwitzTermEqDefinitePairOfLocalIsometryStatement.lean}{GlobalTheory.T.kottwitzTerm_eq_definitePair_of_local_isometry}

\begin{proof}
Fix an irreducible monic $P_0$ of degree $n$, put $\mu=r(P_0)$, and suppose that some pair with
marking $\mu$ or some definite pair has minimal polynomial $P_0$, with $\gamma=\gamma_0$ and
$c=c_0$. Every root of every conjugate of $P_0$ has absolute value $c_0^{1/2}$. Indeed
$\psi_w(\gamma)^\dagger H_w\psi_w(\gamma)=c_0H_w$ with $H_w=J_w\psi_w(\beta^{-1})$, so eigenvectors with
eigenvalues $r,r'$ are orthogonal for $H_w$ unless $\bar rr'=c_0$. Where $H_w$ is definite this gives
$|r|^2=c_0$ for every eigenvalue. At $v$, for a pair with marking $\mu$, $|\mu|^2=c_0$, and an
eigenvalue $r$ with $|r|^2\neq c_0$ would give, with the eigenvalue $c_0/\bar r$, a plane of signature
$(1,1)$ orthogonal to the positive eigenvector of $\mu$, which contradicts the signature $(1,n-1)$.
Hence $M=F[X]/(P_0)$, a field of degree $n$ over $F$, has the automorphism $\sigma$ extending $c$ with
$\sigma(X)=c_0/X$, and $\tau\circ\sigma=\bar\tau$ for every embedding $\tau$ of $M$. Let $\varepsilon$
be the embedding of $M$ extending $v$ with $\varepsilon(X)=\mu$, and $S$ the set of embeddings
extending $v$ or $\bar v$ other than $\varepsilon$ and $\bar\varepsilon$. It has $2(n-1)$ elements, a
multiple of $4$ since $n$ is odd. Let $a$ and $u$ be as in Lemma~\ref{gal-count-definite-signs}, and
write $a=A(X)$ with $A\in F[X]$.

Let $(\beta,\gamma,c_0)$ be a definite pair with $P_\gamma=P_0$, and let $j:M\to C$ send $X$ to
$\gamma$. Put $\beta'=A(\gamma)\beta$. Since $A(\gamma)$ commutes with $\gamma$ and
$\beta j(m)^*=j(\sigma m)\beta$ for $m\in M$, $\beta'^*=\beta'$ and
$\gamma\beta'\gamma^*=c_0\beta'$; and $j(u)\beta j(u)^*=\beta'$ at the finite places away from $p$, so
$\beta'$ is trivial there. The eigenvectors $e_\tau$ of $\psi_w(\gamma)$, one for each embedding
$\tau$ of $M$ extending $w$, are orthogonal for $H_w=J_w\psi_w(\beta^{-1})$, since the eigenvalues are
distinct of absolute value $c_0^{1/2}$, and
$e_\tau^\dagger J_w\psi_w(\beta'^{-1})e_\tau=\tau(a)^{-1}e_\tau^\dagger H_we_\tau$. Hence
$J_w\psi_w(\beta'^{-1})$ is negative definite for $w\neq v$, and at $v$ it is positive on
$e_\varepsilon$, whose eigenvalue is $\mu$, and negative on the other $n-1$ eigenvectors. So it is
congruent to $J_w$ at every infinite place, $\beta'$ is trivial over $\mathbf R$, and
$(\beta',\gamma,c_0)$ is a pair with marking $\mu$. Conversely, for a pair $(\beta',\gamma,c_0)$ with
marking $\mu$ and $P_\gamma=P_0$, the same computation shows that $(A(\gamma)^{-1}\beta',\gamma,c_0)$
is a definite pair. Equivalent triples go to equivalent triples, since replacing $\gamma$ by
$y\gamma y^{-1}$ replaces $A(\gamma)$ by $yA(\gamma)y^{-1}$, and the two maps are inverse to each
other. This gives $\Phi$, which keeps $\gamma$ and $c$. The two triples have the same centralizer,
and $j(u)$, which commutes with $\gamma$, identifies their maps $\iota$ to $G(\mathbf A_f^p)$; their
maps $\iota_p$ are equal. So their orbital sums are equal.
\end{proof}

\begin{proposition}\label{gal-count-weighted}
In Situation~\ref{gal-count-situation}, suppose that $n$ is odd. Let $K\subset G(\mathbf A_f^p)$ be an
open compact subgroup, $f$ a test function for $K$ that vanishes outside the regular elements,
$K'\subset G_p$ a subgroup commensurable with $K_p$, and $f_p$ a test function for $K'$, such that
$O_P(f,f_p)=0$ for every pair $P$ that has a marking and has reducible $P_\gamma$. Then for every
function $\eta:F[X]\times\mathbf C\to\mathbf C$,
\[
 T_v(\eta;f,f_p)=\sum_{[Q]}\Bigl(\sum_{\mu}\eta(P_\gamma,\mu)\Bigr)O_Q(f,f_p),
\]
where $[Q]$ runs over the equivalence classes of definite pairs $Q=(\beta,\gamma,c)$ with $P_\gamma$
irreducible, with a representative chosen in each, and $\mu$ runs over the roots of $v(P_\gamma)$.
\end{proposition}
\leanref{Development/AbelianVarietyTheory/Statements/PlaceShiftStatements/MarkWeightedFixedPointEqDefinitePairOdd.lean}{GlobalTheory.RootSumWeights.MarkWeightedFixedPointEqDefinitePairOdd}
\leanref{Development/AutomorphicArithmetic/Foundations/GaloisRepresentations/DefiniteTraceForms/WeightedCountDefiniteUnitary.lean}{UpperTheory.Gal.Cnt.weighted_count_definite_unitary}
\leanref{Development/AbelianVarietyTheory/Proofs/DivisorProofs/MarkWeightedFixedPointDefinitePairOddOfCounts.lean}{GlobalTheory.RootSumWeights.markWeightedFixedPoint_definitePair_odd_of_counts}

\begin{proof}
By Lemma~\ref{gal-count-complex}, $T_v(\eta;f,f_p)$ is the sum of $\eta(P_\gamma,\mu_P)O_P(f,f_p)$
over the classes of pairs with a marking. Choose, for every monic $P$ of degree $n$, an enumeration
$r_1(P),\ldots,r_n(P)$ of the roots of $v(P)$ with multiplicity. Take a class $[P]$ with a marking.
If $P_\gamma$ is reducible its summand vanishes by hypothesis. If $P_\gamma$ is irreducible, its
marking is a root of $v(P_\gamma)$, and since $P_\gamma$ is separable it equals $r_i(P_\gamma)$ for
exactly one $i$. Hence
\[
 T_v(\eta;f,f_p)=\sum_{i=1}^n\ \sum_{[P]}\eta(P_\gamma,r_i(P_\gamma))\,O_P(f,f_p),
\]
the inner sum over the classes of pairs with irreducible $P_\gamma$ and marking $r_i(P_\gamma)$; the
exchange of the sums is allowed because only finitely many classes contribute
(Lemma~\ref{gal-count-finite}). Since $f$ and $f_p$ are invariant under conjugation by $K$ and $K'$,
Lemma~\ref{gal-count-definite-move} for $r=r_i$ turns the inner sum into the sum of
$\eta(P_\gamma,r_i(P_\gamma))O_Q(f,f_p)$ over the classes of definite pairs $Q$ with irreducible
$P_\gamma$, which again has finitely many nonzero terms. Summing over $i$ gives the sum over the roots
$\mu$ of $v(P_\gamma)$.
\end{proof}

\subsection{Definite forms and the group side}\label{subsection-gal-group-side}

\begin{definition}\label{gal-definite-form-similitude}
In Situation~\ref{gal-count-situation}, let $\beta\in C^\times$ with $\beta^*=\beta$. For a commutative
$\mathbf Q$-algebra $R$, the \emph{similitude group} $\mathrm{GU}_\beta(R)$ of $\beta$ is the group of
$(g,s)\in(R\otimes_{\mathbf Q}C)^\times\times R^\times$ with $g\beta g^*=s\beta$, where the involution is
extended $R$-linearly, and $s$ is the \emph{similitude factor} of $(g,s)$. For $R=\mathbf A_f^p$ and
$R=\mathbf Q_p$, the module $R\otimes_{\mathbf Q}C\cong R^{\dim_{\mathbf Q}C}$ carries the product
topology, a group of units $A^\times$ carries the topology for which $u\mapsto(u,u^{-1})$ is a
homeomorphism onto its image in $A\times A$, and $\mathrm{GU}_\beta(R)$ carries the induced topology. The
group $\mathrm{GU}_\beta(\mathbf Q)$ is embedded diagonally in $\mathrm{GU}_\beta(\mathbf A_f^p)\times
\mathrm{GU}_\beta(\mathbf Q_p)$. Thus $G=\mathrm{GU}_1$, and for a form $\beta$ the group $G_\beta$
(Definition~\ref{gal-count-forms}(2)) is the subgroup of $\mathrm{GU}_\beta(\mathbf Q)$ of the elements
with positive similitude factor.
\end{definition}
\leanref{Development/AbelianVarietyTheory/Proofs/DivisorProofs/TC.lean}{GlobalTheory.TorCls.TC}
\leanref{Development/AbelianVarietyTheory/Proofs/DivisorProofs/FormGroupToPadicPointsDegreeOneEqUnitaryPoints.lean}{GlobalTheory.TorCls.InT}
\leanref{Development/AbelianVarietyTheory/Proofs/DivisorProofs/FormGroupToPadicPointsDegreeOneEqUnitaryPoints.lean}{GlobalTheory.TorCls.dg}

\begin{lemma}\label{gal-definite-form-class-number}
Let $D$ be a finite-dimensional $\mathbf Q$-algebra, $\sigma:D\to D$ a ring antihomomorphism, $\beta,y\in D$,
$p$ a prime and $\mathbf A_f^p$ the ring of finite adeles of $\mathbf Q$ away from $p$. For a commutative
$\mathbf Q$-algebra $R$, extend $\sigma$ $R$-linearly to $R\otimes_{\mathbf Q}D$, and let $T(R)$ be the group of
$(g,s)\in(R\otimes_{\mathbf Q}D)^\times\times R^\times$ with $g(1\otimes\beta)\sigma(g)=s(1\otimes\beta)$ and
$g(1\otimes y)=(1\otimes y)g$. Let
\[
 H_D=\bigl((\mathbf A_f^p\otimes_{\mathbf Q}D)^\times\times(\mathbf A_f^p)^\times\bigr)\times
 \bigl((\mathbf Q_p\otimes_{\mathbf Q}D)^\times\times\mathbf Q_p^\times\bigr),
\]
where $R\otimes_{\mathbf Q}D\cong R^{\dim_{\mathbf Q}D}$, for $R=\mathbf A_f^p$ and $R=\mathbf Q_p$, carries the
product topology, the groups of units carry the topology of Definition~\ref{gal-definite-form-similitude}, and
$H_D$ the product topology. Embed $T(\mathbf Q)$ in $H_D$ diagonally. Let $U\subset H_D$ be a subgroup for
which there is an open neighbourhood $O$ of $1$ in $H_D$ with $O\cap\bigl(T(\mathbf A_f^p)\times
T(\mathbf Q_p)\bigr)\subset U$. Then there is a finite set $S\subset H_D$ such that every element of
$T(\mathbf A_f^p)\times T(\mathbf Q_p)$ is $\gamma su$ with $\gamma\in T(\mathbf Q)$, $s\in S$ and $u\in U$.
\end{lemma}
\leanref{Development/AbelianVarietyTheory/Proofs/DivisorProofs/FiniteDoubleCosetDefiniteHermitianFormLevel.lean}{GlobalTheory.DefClassFin.ccf_holds}
\leanref{Development/AbelianVarietyTheory/Proofs/DivisorProofs/SimilitudeCentraliserClassNumberFiniteBorelStatement.lean}{GlobalTheory.T.similitudeCentraliser_classNumber_finite_borel}
\leanref{Development/AbelianVarietyTheory/Proofs/DivisorProofs/BorelFiniteClassNumberCentraliserTorusStatement.lean}{GlobalTheory.T.borel_finite_classNumber_centraliserTorus}
\leanref{Development/AbelianVarietyTheory/Proofs/DivisorProofs/PolynomialZeroLocusCoordinatesStatement.lean}{GlobalTheory.T.exists_polynomialLocus_similitudeCentralizerTorus}
\leanref{Development/AbelianVarietyTheory/Proofs/AdelicLatticeProofs/ClassNumberFiniteHolds.lean}{GlobalTheory.classNumberFinite_holds}

\begin{proof}
In a basis of $D$ over $\mathbf Q$, $T$ is the subgroup of the units of $D\times\mathbf Q$ cut out by the
polynomial equations $g\beta\sigma(g)=s\beta$ and $gy=yg$, so it is a linear algebraic group over
$\mathbf Q$, compatibly with base change and with the topologies. Its group of points over the finite adeles
$\mathbf A_f=\mathbf A_f^p\times\mathbf Q_p$ is $T(\mathbf A_f^p)\times T(\mathbf Q_p)$, in which
$O\cap\bigl(T(\mathbf A_f^p)\times T(\mathbf Q_p)\bigr)$ is an open neighbourhood of $1$, so it contains an open
compact subgroup $U'$, and $U'\subset U$. By Borel's finiteness of class numbers \cite[Theorem~5.1]{PR94},
$T(\mathbf A_f)$ is the union of finitely many double cosets $T(\mathbf Q)\,s\,U'$, and their representatives
$s$ form $S$.
\end{proof}

\begin{lemma}\label{gal-definite-form-padic}
In Situation~\ref{gal-count-situation}, extend each $\rho_i$ to a $\mathbf Q_p$-algebra homomorphism
$\rho_i:\mathbf Q_p\otimes_{\mathbf Q}C\to M_n(F_{w_i})$.
\begin{enumerate}
\item The map $\mathbf Q_p\otimes_{\mathbf Q}C\to\prod_{i=1}^mM_n(F_{w_i})^2$,
  $a\mapsto(\rho_i(a),\rho_i(a^*))_i$, is a $\mathbf Q_p$-linear bijection.
\item The map $(g,s)\mapsto(s,(\rho_i(g))_i)$ is a bijective homomorphism $G(\mathbf Q_p)\to G_p$.
\item Let $\beta\in C^\times$ with $\beta^*=\beta$. Then every element of $G_p$ is $(s,(\rho_i(g))_i)$ for
  some $(g,s)\in\mathrm{GU}_\beta(\mathbf Q_p)$.
\end{enumerate}
\end{lemma}
\leanref{Development/AbelianVarietyTheory/Proofs/DivisorProofs/PadicPointsSimilitudeGLProductBijective.lean}{GlobalTheory.LKTrip.HTGroup.LocPBij.psiL_bijective}
\leanref{Development/AbelianVarietyTheory/Proofs/DivisorProofs/LocalMapInjectivity.lean}{GlobalTheory.HTGroupLocInj.locW_inj}
\leanref{Development/AbelianVarietyTheory/Proofs/DivisorProofs/PadicPointsSimilitudeGLProductBijectiveStatement.lean}{GlobalTheory.T.padicPoints_similitude_GL_product_bijective}
\leanref{Development/AbelianVarietyTheory/Proofs/DivisorProofs/PadicPointsSimilitudeGLProductBijective.lean}{GlobalTheory.LKTrip.HTGroup.bijective_locP}
\leanref{Development/AbelianVarietyTheory/Proofs/DivisorProofs/ExistsUnitaryCentralizerOfCommutePadicPointsSpliStatement.lean}{GlobalTheory.T.exists_unitaryCentralizer_of_commute_padicPoints_split}
\leanref{Development/AbelianVarietyTheory/Proofs/DivisorProofs/AlgebraMapQrOneBijective.lean}{GlobalTheory.TorCls.AtP.atp_of_bijective}
\leanref{Development/AbelianVarietyTheory/Proofs/DivisorProofs/AlgebraMapQrOneBijective.lean}{GlobalTheory.TorCls.AtP.exists_triv}

\begin{proof}
Since $p$ splits in $K_0$, the places of $F$ over $p$ are the $w_i$, over $u$, and their complex
conjugates $\bar w_i$, over $\bar u$, so $\mathbf Q_p\otimes F=\prod_i(F_{w_i}\times F_{\bar w_i})$.

(1) The map is a ring homomorphism to $\prod_iM_n(F_{w_i})\times\prod_iM_n(F_{w_i})^{\mathrm{op}}$, so its
kernel is a two-sided ideal of $\mathbf Q_p\otimes C=C\otimes_F(\mathbf Q_p\otimes F)$. Since $C$ is central
simple over $F$, a nonzero two-sided ideal contains a nonzero element $z$ of $\mathbf Q_p\otimes F$. For such
$z$, $\rho_i(z)$ is the scalar matrix of the component of $z$ at $w_i$, and $\rho_i(z^*)$ that of the image
under $c$ of its component at $\bar w_i$; one of them is nonzero for some $i$, so $z$ is not in the kernel.
Hence the map is injective. The local degrees $[F_{w_i}:\mathbf Q_p]$ add up to
$[F:K_0]=[F:\mathbf Q]/2$, so both sides have dimension $n^2[F:\mathbf Q]$ over $\mathbf Q_p$, and the map
is bijective.

(2) For $(g,s)\in G(\mathbf Q_p)$, $g^*=sg^{-1}$, so $\rho_i(g^*)=s\rho_i(g)^{-1}$ and the map is injective
by (1). Given $(s,(g_i)_i)\in G_p$, (1) gives $a$ with $\rho_i(a)=g_i$ and $\rho_i(a^*)=sg_i^{-1}$. Then
$aa^*$ and $s$ have the same image under the map of (1), so $aa^*=s$. So $a$ has the right inverse
$s^{-1}a^*$, hence is invertible, and $(a,s)\in G(\mathbf Q_p)$ maps to $(s,(g_i)_i)$.

(3) Let $e\in\mathbf Q_p\otimes F$ have components $1$ at the $w_i$ and $0$ at the $\bar w_i$, so that $e$
is a central idempotent with $e^*=1-e$. Put $Y=e+(1-e)\beta$. Then $Y$ is invertible with inverse
$e+(1-e)\beta^{-1}$, and $YY^*=(e+(1-e)\beta)((1-e)+\beta e)=\beta$. So $X=Y^{-1}$ satisfies $X\beta X^*=1$,
and $(g,s)\mapsto(XgX^{-1},s)$ maps $\mathrm{GU}_\beta(\mathbf Q_p)$ onto $G(\mathbf Q_p)$, with inverse
$(h,s)\mapsto(X^{-1}hX,s)$. Given $(s,(g_i)_i)\in G_p$, (2) gives $(h,s)\in G(\mathbf Q_p)$ with
$\rho_i(h)=\rho_i(X)g_i\rho_i(X)^{-1}$, and $g=X^{-1}hX$ has $\rho_i(g)=g_i$.
\end{proof}

\begin{lemma}\label{gal-definite-form-classes}
In Situation~\ref{gal-count-situation}, let $\beta$ be a definite form, with $x\in(\mathbf A_f^p\otimes
C)^\times$ and $t\in(\mathbf A_f^p)^\times$ such that $x\beta x^*=t$, which define $\iota$. Let
$K\subset G(\mathbf A_f^p)$ be an open compact subgroup and $K'\subset G_p$ a subgroup commensurable with
$K_p$, and put $H=G(\mathbf A_f^p)\times G_p$ and $U=K\times K'$. Then there is a finite set $S\subset H$
with $H=(\iota,\iota_p)(G_\beta)\,S\,U$. In particular, $G_\beta\backslash H/U$ is finite.
\end{lemma}
\leanref{Development/AbelianVarietyTheory/Proofs/DivisorProofs/FiniteDoubleCosetDefiniteHermitianFormLevelStatement.lean}{GlobalTheory.T.finite_doubleCoset_definiteHermitianForm_level}
\leanref{Development/AbelianVarietyTheory/Proofs/DivisorProofs/FiniteDoubleCosetDefiniteHermitianFormLevel.lean}{GlobalTheory.DefClassFin.finite_doubleCoset}
\leanref{Development/AbelianVarietyTheory/Proofs/DivisorProofs/FiniteDoubleCosetDefiniteHermitianFormLevel.lean}{GlobalTheory.DefClassFin.cover_KHT}
\leanref{Development/AbelianVarietyTheory/Proofs/DivisorProofs/FiniteDoubleCosetDefiniteHermitianFormLevel.lean}{GlobalTheory.DefClassFin.cover_of_relIndex}
\leanref{Development/AbelianVarietyTheory/Proofs/DivisorProofs/FiniteDoubleCosetDefiniteHermitianFormLevel.lean}{GlobalTheory.DefClassFin.finite_of_cover}

\begin{proof}
Let $a:\mathrm{GU}_\beta(\mathbf A_f^p)\to G(\mathbf A_f^p)$ be $(g,s)\mapsto(xgx^{-1},s)$, and
$b:\mathrm{GU}_\beta(\mathbf Q_p)\to G_p$ be $(g,s)\mapsto(s,(\rho_i(g))_i)$. Since $t$ is central,
$x^{-1}x^{-*}=t^{-1}\beta$, so for $(g,s)\in\mathrm{GU}_\beta(\mathbf A_f^p)$
\[
 (xgx^{-1})(xgx^{-1})^*=t^{-1}x\,g\beta g^*\,x^*=t^{-1}s\,x\beta x^*=s,
\]
and in the same way $(x^{-1}hx,s)\in\mathrm{GU}_\beta(\mathbf A_f^p)$ for $(h,s)\in G(\mathbf A_f^p)$. So $a$
is a bijective homomorphism. Both $a$ and $b$ are continuous, and on $G_\beta$ they restrict to $\iota$ and
$\iota_p$.

Apply Lemma~\ref{gal-definite-form-class-number} with $D=C$, $\sigma$ the involution and $y=1$, so that
$T=\mathrm{GU}_\beta$. The group $U_0=a^{-1}(K)\times b^{-1}(K_p)$ is open in
$\mathrm{GU}_\beta(\mathbf A_f^p)\times\mathrm{GU}_\beta(\mathbf Q_p)$, so it is the intersection of this group
with an open neighbourhood of $1$ in $H_C$. The lemma gives a finite set $S_0\subset H_C$ such that every
element of $\mathrm{GU}_\beta(\mathbf A_f^p)\times\mathrm{GU}_\beta(\mathbf Q_p)$ is $\gamma s u$ with
$\gamma\in\mathrm{GU}_\beta(\mathbf Q)$, $s\in S_0$ and $u\in U_0$; such an $s=\gamma^{-1}(\gamma su)u^{-1}$ lies in
$\mathrm{GU}_\beta(\mathbf A_f^p)\times\mathrm{GU}_\beta(\mathbf Q_p)$, so we may take $S_0$ inside this group,
and then $\mathrm{GU}_\beta(\mathbf A_f^p)\times\mathrm{GU}_\beta(\mathbf Q_p)=
\mathrm{GU}_\beta(\mathbf Q)\,S_0\,U_0$. If some $\gamma_-\in\mathrm{GU}_\beta(\mathbf Q)$ has negative
similitude factor, fix one and put $S_1=S_0\cup\gamma_-S_0$; otherwise put $S_1=S_0$. Every element of
$\mathrm{GU}_\beta(\mathbf Q)$ lies in $G_\beta$ or in $G_\beta\gamma_-$, so
$\mathrm{GU}_\beta(\mathbf A_f^p)\times\mathrm{GU}_\beta(\mathbf Q_p)=G_\beta\,S_1\,U_0$.

Let $(g,g_p)\in H$. By Lemma~\ref{gal-definite-form-padic}(3) there is $z_p\in\mathrm{GU}_\beta(\mathbf Q_p)$
with $b(z_p)=g_p$. Write $(a^{-1}(g),z_p)=\delta s_1k_0$ with $\delta\in G_\beta$, $s_1\in S_1$ and
$k_0\in U_0$. Applying $a\times b$ gives $(g,g_p)=(\iota,\iota_p)(\delta)\,(a\times b)(s_1)\,(a\times
b)(k_0)$ with $(a\times b)(k_0)\in K\times K_p$. Hence $H=(\iota,\iota_p)(G_\beta)\,S_2\,(K\times K_p)$ with
$S_2=(a\times b)(S_1)$ finite.

Since $K'$ is commensurable with $K_p$, $K'\cap K_p$ has finite index in $K_p$. Let $r_1,\ldots,r_l$
represent the cosets $k(K'\cap K_p)$ in $K_p$. Every element of $K\times K_p$ is $(1,r_j)(k,k')$ with $k\in
K$ and $k'\in K'\cap K_p$, so $H=(\iota,\iota_p)(G_\beta)\,S\,U$ for the finite set $S$ of the products
$s_2(1,r_j)$ with $s_2\in S_2$. Every double coset in $G_\beta\backslash H/U$ then contains an element of
$S$, so there are finitely many.
\end{proof}

\begin{lemma}\label{gal-definite-form-denominators}
In Situation~\ref{gal-count-situation}, fix a basis $e_1,\ldots,e_N$ of $C$ over $\mathbf Q$, and write
$\kappa_j(y)\in\mathbf Q$ for the coordinates of $y\in C$, so that $y=\sum_j\kappa_j(y)e_j$. Let $\beta$ be a
form, with $x\in(\mathbf A_f^p\otimes C)^\times$ and $t\in(\mathbf A_f^p)^\times$ such that $x\beta x^*=t$,
which define $\iota$. Let $K\subset G(\mathbf A_f^p)$ be an open compact subgroup, $f$ a test function for
$K$, $K'\subset G_p$ a subgroup commensurable with $K_p$, $f_p$ a test function for $K'$, $h\in
G(\mathbf A_f^p)$ and $h_p\in G_p$. Then there are integers $D\geq1$ and $k\geq0$ such that for every
$(y,s)\in G_\beta$:
\begin{enumerate}
\item if $f(h^{-1}\iota(y,s)h)\neq0$, then the numbers $D\kappa_j(y)$ and $Ds^{-1}$ are integral at every
  prime $\ell\neq p$;
\item if $f_p(h_p^{-1}\iota_p(y,s)h_p)\neq0$, then the numbers $p^k\kappa_j(y)$ and $p^ks^{-1}$ are
  integral at $p$.
\end{enumerate}
In particular, if both hold, then $Dp^k\kappa_j(y)$ and $Dp^ks^{-1}$ are integers.
\end{lemma}
\leanref{Development/AbelianVarietyTheory/Proofs/DivisorProofs/FormGroupTestFunctionSupportBoundedDenominatorsStatement.lean}{GlobalTheory.T.formGroup_testFunction_support_bounded_denominators}
\leanref{Development/AbelianVarietyTheory/Proofs/DivisorProofs/FormGroupTestFunctionSupportBoundedDenominators.lean}{GlobalTheory.LKCx.BoundedDen.main}
\leanref{Development/AbelianVarietyTheory/Proofs/DivisorProofs/HopfOrderPullbackAlongHom.lean}{GlobalTheory.DenAway.den_bounded_of_isCompact}
\leanref{Development/AbelianVarietyTheory/Proofs/DivisorProofs/PAdicDenominatorBound.lean}{GlobalTheory.T.p_adic_denominator_bound}
\leanref{Development/AbelianVarietyTheory/Proofs/DivisorProofs/RootFlatPushforwardGeoFun.lean}{GlobalTheory.AtP.main}
\leanref{Development/AbelianVarietyTheory/Proofs/DivisorProofs/RegularElementsBoundedPadic.lean}{GlobalTheory.T.regular_elements_bounded_padic}
\leanref{Development/AbelianVarietyTheory/Proofs/DivisorProofs/BoundedCoordinatesPPowerDenominator.lean}{GlobalTheory.T.bounded_coordinates_p_power_denominator}
\leanref{Development/AbelianVarietyTheory/Proofs/DivisorProofs/ExistsIntOfDenLe.lean}{GlobalTheory.FinFix.exists_int_of_den_le}

\begin{proof}
(1) The function $f$ vanishes outside finitely many double cosets $K\sigma K$, which are compact, so the
elements $g$ with $f(h^{-1}gh)\neq0$ lie in a compact subset of $G(\mathbf A_f^p)$. The map
$(g,s)\mapsto(x^{-1}gx,s^{-1})$ from $G(\mathbf A_f^p)$ to $(\mathbf A_f^p\otimes C)\times\mathbf A_f^p$ is
continuous and sends $\iota(y,s)=(xyx^{-1},s)$ to $(y,s^{-1})$, and the coordinates in the basis
$1\otimes e_j$ are continuous. So the numbers $\kappa_j(y)$ and $s^{-1}$, for the $(y,s)$ with
$f(h^{-1}\iota(y,s)h)\neq0$, lie in a compact subset $\Omega$ of $\mathbf A_f^p$. With $\hat{\mathbf
Z}^p=\prod_{\ell\neq p}\mathbf Z_\ell$, the open sets $m^{-1}\hat{\mathbf Z}^p$, $m\geq1$, cover
$\mathbf A_f^p$, so finitely many of them cover $\Omega$, and $\Omega\subset D^{-1}\hat{\mathbf Z}^p$ for the
product $D$ of their $m$. A rational number $q\in D^{-1}\hat{\mathbf Z}^p$ has $Dq$ integral at every
$\ell\neq p$.

(2) The function $f_p$ vanishes outside finitely many double cosets $K_p\sigma K_p$, so the elements
$\iota_p(y,s)=(s,(\rho_i(y))_i)$ with $f_p(h_p^{-1}\iota_p(y,s)h_p)\neq0$, and their inverses
$\iota_p(y^{-1},s^{-1})$, lie in compact subsets of $G_p$. Hence $s$ and $s^{-1}$ lie in a compact subset of
$\mathbf Q_p^\times$, so $p^{k_1}s^{-1}\in\mathbf Z_p$ for some $k_1$, and the entries of $\rho_i(y)$ and
$\rho_i(y^{-1})$ are bounded independently of $(y,s)$. From $y\beta y^*=s\beta$,
$y^*=s\beta^{-1}y^{-1}\beta$, so the entries of $\rho_i(y^*)=s\rho_i(\beta)^{-1}\rho_i(y^{-1})\rho_i(\beta)$
are bounded too. By Lemma~\ref{gal-definite-form-padic}(1), the coordinates of $1\otimes y$ in the basis
$1\otimes e_j$ of $\mathbf Q_p\otimes C$ are $\mathbf Q_p$-linear functions of $(\rho_i(y),\rho_i(y^*))_i$,
so they are bounded, and $p^{k_2}\kappa_j(y)\in\mathbf Z_p$ for some $k_2$. Take $k=\max(k_1,k_2)$.

A rational number that is integral at every prime is an integer, which gives the last assertion.
\end{proof}

\begin{lemma}\label{gal-definite-form-infinite}
In Situation~\ref{gal-count-situation}, fix a basis $e_1,\ldots,e_N$ of $C$ over $\mathbf Q$, write
$\kappa_j(y)\in\mathbf Q$ for the coordinates of $y\in C$, so that $y=\sum_j\kappa_j(y)e_j$, and let $\beta$
be a definite form. Then there is a real number $B_\beta$ with $|\kappa_j(y)|\leq B_\beta s^{1/2}$ for every
$(y,s)\in G_\beta$ and every $j$.
\end{lemma}
\leanref{Development/AbelianVarietyTheory/Proofs/DivisorProofs/DefiniteFormStabilizerFiniteAwayFromP.lean}{GlobalTheory.DefStab.psi_bound}
\leanref{Development/AbelianVarietyTheory/Proofs/DivisorProofs/LocalPairValueLinear.lean}{GlobalTheory.LKCx.PsiBound.bound_of_form}
\leanref{Development/AbelianVarietyTheory/Proofs/DivisorProofs/StandardAdelicCharacter.lean}{GlobalTheory.LKCx.grp_similitude_herm}
\leanref{Development/AbelianVarietyTheory/Proofs/DivisorProofs/PowerSumsVanishOfTrivialFlags.lean}{GlobalTheory.CoordArch.coord_bound}

\begin{proof}
Let $w$ be an infinite place, $H_w=J_w\psi_w(\beta^{-1})$, and $\epsilon_w=1$ if $w=v$ and $\epsilon_w=-1$
otherwise, so that $\epsilon_wH_w$ is positive definite. There are $0<m_w\leq M_w$ with
$m_w|z|^2\leq\epsilon_wz^\dagger H_wz\leq M_w|z|^2$ for $z\in\mathbf C^n$. Let $(y,s)\in G_\beta$ and
$Y=\psi_w(y)$. Applying $\psi_w$ to $\beta^{-1}=s^{-1}y^*\beta^{-1}y$ and using
$\psi_w(y^*)=J_w^{-1}Y^\dagger J_w$ gives $Y^\dagger H_wY=sH_w$. Hence
\[
 |Yz|^2\leq m_w^{-1}\epsilon_w(Yz)^\dagger H_w(Yz)=m_w^{-1}s\,\epsilon_wz^\dagger H_wz\leq(M_w/m_w)\,s\,|z|^2,
\]
so every entry of $\psi_w(y)$ has absolute value at most $(M_w/m_w)^{1/2}s^{1/2}$.

The $\mathbf R$-linear map from $\mathbf R\otimes_{\mathbf Q}C$ to $\prod_wM_n(\mathbf C)$ extending
$y\mapsto(\psi_w(y))_w$ is injective. Indeed $\mathbf R\otimes_{\mathbf Q}F$ is the product of the fields
$\mathbf C$ over the infinite places $w$, so $\mathbf R\otimes_{\mathbf Q}C$ is the product of the algebras
$C\otimes_{F,w}\mathbf C$. The extension of $\psi_w$ vanishes on the factors at the other places, and it is
injective on $C\otimes_{F,w}\mathbf C$ because this algebra is simple. So there is $B_0$ with
$|\kappa_j(y)|\leq B_0\max_{w,a,b}|\psi_w(y)_{ab}|$ for $y\in C$, and $B_\beta=B_0\max_w(M_w/m_w)^{1/2}$
will do.
\end{proof}

\begin{lemma}\label{gal-definite-form-stabilizers}
In Situation~\ref{gal-count-situation}, let $\beta$ be a definite form, with $x\in(\mathbf A_f^p\otimes
C)^\times$ and $t\in(\mathbf A_f^p)^\times$ such that $x\beta x^*=t$, which define $\iota$. Let
$K\subset G(\mathbf A_f^p)$ be an open compact subgroup, $K'\subset G_p$ a subgroup commensurable with $K_p$,
$g\in G(\mathbf A_f^p)$ and $g_p\in G_p$. Then the group of $(y,s)\in G_\beta$ with
$g^{-1}\iota(y,s)g\in K$ and $g_p^{-1}\iota_p(y,s)g_p\in K'$ is finite.
\end{lemma}
\leanref{Development/AbelianVarietyTheory/Proofs/DivisorProofs/DefiniteFormStabilizerFiniteAwayFromPStatement.lean}{GlobalTheory.T.definiteForm_stabilizer_finite_away_from_p}
\leanref{Development/AbelianVarietyTheory/Proofs/DivisorProofs/DefiniteFormStabilizerFiniteAwayFromP.lean}{GlobalTheory.DefStab.finite_locPart}
\leanref{Development/AbelianVarietyTheory/Proofs/DivisorProofs/DefiniteFormStabilizerFiniteAwayFromP.lean}{GlobalTheory.DefStab.isTest_indicator_comm}
\leanref{Development/AbelianVarietyTheory/Proofs/DivisorProofs/FiniteDenBdd.lean}{GlobalTheory.FinFix.finite_den_bdd}
\leanref{Development/AbelianVarietyTheory/Proofs/DivisorProofs/ExistsIntOfDenLe.lean}{GlobalTheory.FinFix.grp_val_injective}

\begin{proof}
The indicator function $\mathbf 1_K$ is a test function for $K$. The indicator function $\mathbf 1_{K'}$ is
invariant under conjugation by $K'$; since $K_p\cap K'$ has finite index in $K'$, the group $K'$ is the union
of finitely many cosets $r(K_p\cap K')$, each contained in the double coset $K_prK_p$, so $\mathbf 1_{K'}$ is a
test function for $K'$. Fix a basis of $C$ with coordinates $\kappa_j$, let $D$ and $k$ be given by
Lemma~\ref{gal-definite-form-denominators} for $\mathbf 1_K$, $\mathbf 1_{K'}$, $h=g$ and $h_p=g_p$, and put
$N=Dp^k$. For $(y,s)$ in the group, both functions are nonzero at the conjugates of $\iota(y,s)$ and
$\iota_p(y,s)$, so $N\kappa_j(y)$ and $Ns^{-1}$ are integers. Since $s>0$, $Ns^{-1}\geq1$, so $s\leq N$, and
Lemma~\ref{gal-definite-form-infinite} gives $|\kappa_j(y)|\leq B_\beta N^{1/2}$. So each $\kappa_j(y)$ lies
in the finite set of numbers in $N^{-1}\mathbf Z$ of absolute value at most $B_\beta N^{1/2}$, and $y$ takes
finitely many values. Since $s\beta=y\beta y^*$, $y$ determines $s$. Hence the group is finite.
\end{proof}

\begin{lemma}\label{gal-count-definite-trace}
In Situation~\ref{gal-count-situation}, let $\xi$ be a representation of $C^\times\times\mathbf Q^\times$
on a finite-dimensional complex vector space $V$ and $\eta:F[X]\to\mathbf C$ a function such that
$\operatorname{tr}\xi(y)=\eta(P_y)$ for every definite form $\beta$ and every $y\in G_\beta$ with $P_y$
separable of degree $n$. Let $K\subset G(\mathbf A_f^p)$ be an open compact subgroup, $f$ a test
function for $K$ invariant under $K$ on both sides, $K'\subset G_p$ a subgroup commensurable with $K_p$,
and $f_p$ a test function for $K'$ invariant under $K'$ on both sides. Suppose that for every definite
form $\beta$, every choice of $x$ and $t$, and every $y\in G_\beta$ whose minimal polynomial is not
irreducible, separable and of degree $n$, the orbital integral of $f\otimes f_p$ at
$(\iota,\iota_p)(y)$ (Definition~\ref{gal-count-pairs}(6), with $K'$ in place of $K_p$) vanishes.
Then
\[
 \sum_{[Q]}\eta(P_\gamma)\,O_Q(f,f_p)=\operatorname{tr}_{\mathrm{def}}(\xi;f,f_p),
\]
where $[Q]$ runs over the equivalence classes of definite pairs $Q=(\beta,\gamma,c)$ with $P_\gamma$
irreducible.
\end{lemma}
\leanref{Development/AbelianVarietyTheory/Statements/PlaceShiftStatements/TraceDefinitePairsEqAutomorphicTrace.lean}{GlobalTheory.LKDef.TraceDefinitePairsEqAutomorphicTrace}
\leanref{Development/AutomorphicArithmetic/Foundations/GaloisRepresentations/DefiniteTraceForms/WeightedCountDefiniteUnitary.lean}{UpperTheory.Gal.Cnt.trace_definite_forms_identity}
\leanref{Development/AbelianVarietyTheory/Proofs/DivisorProofs/LieModelBaseChangeMap.lean}{GlobalTheory.LKDef.trace_definite_forms_identity}
\leanref{Development/AbelianVarietyTheory/Proofs/DivisorProofs/BiInvariantTestFunctionSum.lean}{GlobalTheory.T.bi_invariant_test_function_sum}
\leanref{Development/AbelianVarietyTheory/Proofs/DivisorProofs/HeckeTraceDefiniteFormEqSumConjClassesStatement.lean}{GlobalTheory.T.heckeTrace_definiteForm_eq_sum_conjClasses}
\leanref{Development/AbelianVarietyTheory/Proofs/DivisorProofs/CentSumEqZeroOfOrbitalIntegralEqZeroStatement.lean}{GlobalTheory.T.centSum_eq_zero_of_orbitalIntegral_eq_zero}

\begin{proof}
Write $H=G(\mathbf A_f^p)\times G_p$, $U=K\times K'$ and $F_0=f\otimes f_p$. For a definite form
$\beta$ and $y\in G_\beta$, let
\[
 C_\beta(y)=\sum_{g}\frac{F_0(g^{-1}(\iota,\iota_p)(y)g)}{|Z_y\cap gUg^{-1}|},
\]
the sum over $g\in Z_y\backslash H/U$, where $Z_y$ is the centralizer of $y$ in $G_\beta$. A definite
pair $(\beta',\gamma,c)$ has $c>0$, since $J_v\psi_v(\beta'^{-1})$ is definite, so $(\gamma,c)\in
G_{\beta'}$. Sending its class to the class of $\beta'$ together with the conjugacy class of
$(\gamma,c)$, transported to the chosen representative $\beta$, is a bijection onto the pairs
consisting of a class of definite forms and a conjugacy class of its group with irreducible separable
minimal polynomial of degree $n$, and $O_Q(f,f_p)$ becomes $C_\beta(y)$. So the left side is
$\sum_{[\beta]}\sum_{[y]}\eta(P_y)C_\beta(y)$, over these conjugacy classes.

Fix one definite form $\beta$, write $\delta$ also for $(\iota,\iota_p)(\delta)$ when $\delta\in G_\beta$,
and let $g$ run over representatives of $G_\beta\backslash H/U$. This set is finite by
Lemma~\ref{gal-definite-form-classes}, and each group $\Gamma_g=G_\beta\cap gUg^{-1}$ is finite by
Lemma~\ref{gal-definite-form-stabilizers}.
A function $\varphi\in\mathcal A_\beta(\xi)$ is determined by the vectors $\varphi(g)$, and these can be
any vectors with $\varphi(g)\in V^{\Gamma_g}$, the vectors fixed by $\xi(\Gamma_g)$; let
$E_g=|\Gamma_g|^{-1}\sum_{\gamma\in\Gamma_g}\xi(\gamma)$ be the projection onto $V^{\Gamma_g}$. The cosets
$hU$ in $G_\beta g'U$ correspond to the cosets $\delta\Gamma_{g'}$ in $G_\beta$ by $h=\delta g'$, so the
Hecke operator $T$ is
\[
 (T\varphi)(g)=\sum_{h\in H/U}F_0(g^{-1}h)\varphi(h)
 =\sum_{g'}\ \sum_{\delta\in G_\beta/\Gamma_{g'}}F_0(g^{-1}\delta g')\,\xi(\delta)\varphi(g').
\]
Its trace is the sum over $g$ of the traces of the blocks $g'=g$ on $V^{\Gamma_g}$, that is, of
$\sum_{\delta\in G_\beta/\Gamma_g}F_0(g^{-1}\delta g)\operatorname{tr}(\xi(\delta)E_g)$. Since
$F_0(g^{-1}\delta\gamma g)=F_0(g^{-1}\delta g)$ for $\gamma\in\Gamma_g$, this is
$|\Gamma_g|^{-1}\sum_{y\in G_\beta}F_0(g^{-1}yg)\operatorname{tr}\xi(y)$. Group the $y$ by conjugacy
classes, $y=\delta y_0\delta^{-1}$ with $\delta\in G_\beta/Z_{y_0}$, and put $h=\delta^{-1}g$. As
$(g,\delta)$ runs over the representatives and $G_\beta/Z_{y_0}$, the double coset $Z_{y_0}hU$ runs over
$Z_{y_0}\backslash H/U$, each one $|\Gamma_g|/|Z_{y_0}\cap hUh^{-1}|$ times, so the trace of $T$ is
$\sum_{[y]}\operatorname{tr}\xi(y)C_\beta(y)$, over all conjugacy classes of $G_\beta$. For a class
whose minimal polynomial is not irreducible, separable and of degree $n$, $C_\beta(y)$ is a multiple of the vanishing orbital integral, so it vanishes. For the
other classes $\operatorname{tr}\xi(y)=\eta(P_y)$. Summing over the classes of definite forms gives the
claim.
\end{proof}

\begin{definition}\label{gal-count-weights}
Let $E$ be a CM field, $\tau_1:E\to\mathbf C$ an embedding, $\omega\in E$ with
$\tau_1(\omega)\notin\mathbf R$, and $n\geq1$. Write $\bar\tau$ for the complex conjugate of an
embedding $\tau$, put $\epsilon(\tau)=1$ if $\tau(\omega)=\overline{\tau_1(\omega)}$ and
$\epsilon(\tau)=0$ otherwise, and for $\chi\in E[X]$ let $R_\tau(\chi)$ be the multiset of roots of
$\tau(\chi)$.
\begin{enumerate}
\item For $x\in E$, $y,\mu\in\mathbf C$ and $\chi\in E[X]$, let $R'$ be $R_{\tau_1}(\chi)$ with one
  copy of $\mu$ removed, and put
  \[
   b(y;\chi,\mu)=\prod_{r\in R'}\frac{r-y\mu}{r-\mu},
  \]
  \[
   \phi_x(y;\chi,\mu)=(1-y\tau_1(x)\mu)\prod_{r\in R'}(1-\tau_1(x)r)\cdot
   (1-\overline{\tau_1(x)\mu})\prod_{r\in R'}(1-y\,\overline{\tau_1(x)r})\cdot
   \prod_{\tau\neq\tau_1,\bar\tau_1}\ \prod_{r\in R_\tau(\chi)}(1-y^{\epsilon(\tau)}\tau(x)r).
  \]
  For $k\geq0$ and $\alpha\in\mathcal O_E^k$ put
  $L_{\alpha,y}(\chi,\mu)=b(y;\chi,\mu)\prod_{s=1}^k\phi_{\alpha_s}(y;\chi,\mu)$.
\item For $P\in E[X]$ let $P^\vee=P(0)^{-1}X^{\deg P}P(1/X)$, and $P^\vee=0$ if $P(0)=0$. For an
  infinite place $v$ of $E$, the \emph{Lefschetz weight} is the function
  $w_{\alpha,y}(P,\mu)=L_{\alpha,y}(P^\vee,\mu^\vee)$ on $E[X]\times\mathbf C$, where
  $\mu^\vee=\mu^{-1}$ if $v(\omega)=\tau_1(\omega)$ and $\mu^\vee=\bar\mu^{-1}$ otherwise.
\item A \emph{weight} $\lambda$ is a family of integers $\lambda_{\tau,0}\geq\cdots\geq
  \lambda_{\tau,n-1}\geq0$, one family for each embedding $\tau$; put
  $|\lambda_\tau|=\sum_j\lambda_{\tau,j}$ and $e=\sum_{\tau\neq\tau_1,\bar\tau_1}\epsilon(\tau)
  |\lambda_\tau|$. For $\chi\in E[X]$ of degree $n$, let
  $s_\lambda(\chi)=\prod_\tau s_{\lambda_\tau}(R_\tau(\chi))$, with $s_{\lambda_\tau}$ the Schur
  polynomial, and $S_\lambda(P)=s_\lambda(P^\vee)$. A tuple $(h_0,\ldots,h_{n-1})$ is a \emph{tuple of
  Hodge exponents} of $\lambda$ if either $\lambda_{\bar\tau_1,j}=m$ for all $j$ and
  $h_i=(n-1)m+e+\lambda_{\tau_1,i}+n-1-i$ for all $i$, or $\lambda_{\tau_1,j}=m$ for all $j$ and
  $h_i=m+e+\sum_{j\neq i}\lambda_{\bar\tau_1,j}+i$ for all $i$.
\end{enumerate}
\end{definition}
\leanref{Development/AutomorphicArithmetic/Foundations/GaloisRepresentations/SphericalCoordinates/DiagAct.lean}{UpperTheory.Gal.W2.locW}
\leanref{Development/AutomorphicArithmetic/Foundations/GaloisRepresentations/LefschetzSigmaSums/InvPol.lean}{UpperTheory.Gal.LocalSmoothExtension.invPol}
\leanref{Development/AutomorphicArithmetic/Foundations/GaloisRepresentations/LefschetzSigmaSums/InvPol.lean}{UpperTheory.Gal.LocalSmoothExtension.cW}
\leanref{Development/AutomorphicArithmetic/Foundations/GaloisRepresentations/LefschetzSigmaSums/InvPol.lean}{UpperTheory.Gal.LocalSmoothExtension.schW}
\leanref{Development/AutomorphicArithmetic/Foundations/GaloisRepresentations/LefschetzSigmaSums/InvPol.lean}{UpperTheory.Gal.LocalSmoothExtension.IsHExp}

\begin{lemma}\label{gal-count-root-weights}
Let $E$, $\tau_1$, $\omega$ and $n$ be as in Definition~\ref{gal-count-weights}, $\lambda$ a weight and
$h$ a tuple of Hodge exponents of $\lambda$. Then there is $k_0$ such that for every $k\geq k_0$ there
are finitely many $\alpha\in\mathcal O_E^k$, each with all $\alpha_s$ nonzero and divisible by $2$,
and numbers $c_\alpha\in\mathbf C$, such that for every separable $\chi\in E[X]$ of degree $n$ with
$\chi(0)\neq0$ whose roots under $\tau_1$ all have the same absolute value, and every
$y\in\mathbf C$,
\[
 \sum_\alpha c_\alpha\sum_{\mu\in R_{\tau_1}(\chi)}L_{\alpha,y}(\chi,\mu)=s_\lambda(\chi)\sum_{i=0}^{n-1}y^{h_i}.
\]
\end{lemma}
\leanref{Development/AutomorphicArithmetic/Foundations/GaloisRepresentations/LefschetzSigmaSums/HeckeTraceSemi.lean}{UpperTheory.Gal.LocalSmoothExtension.RootW}
\leanref{Development/AutomorphicArithmetic/Foundations/GaloisRepresentations/LefschetzSigmaSums/HeckeTraceSemiHolds.lean}{UpperTheory.Gal.LocalSmoothExtension.RootW_holds}

\begin{proposition}\label{gal-count-group-side}
Let $p$ be a prime and take a unitary PEL datum of simple type for a CM field $E$ at $p$
(Definition~\ref{gal-pel-datum}) of rank $n$, with $\omega$, $\tau_1$ and group $G$, presented as in
Situation~\ref{gal-count-situation} with $F=E$ and $v$ the infinite place of signature $(1,n-1)$. Let
$\lambda$ be a weight, $h$ a tuple of Hodge exponents of $\lambda$, and $k\geq k_0$ and
$(c_\alpha)$ as in Lemma~\ref{gal-count-root-weights}. Let $\xi$ be a representation of
$C^\times\times\mathbf Q^\times$ on a finite-dimensional complex vector space with
$\operatorname{tr}\xi(y)=S_\lambda(P_y)$ for every definite form $\beta$ and every $y\in G_\beta$ with
$P_y$ separable of degree $n$. Let $K\subset G(\mathbf A_f^p)$ be an open compact subgroup and $f$ an
elliptic test function for $K$, invariant under $K$ on both sides. Then for every $y\in\mathbf C$,
\[
 \sum_\alpha c_\alpha\,T_v(w_{\alpha,y};f,\mathbf 1_{K_p})
 =\operatorname{tr}_{\mathrm{def}}(\xi;f,\mathbf 1_{K_p})\sum_{i=0}^{n-1}y^{h_i}.
\]
\end{proposition}
\leanref{Development/AutomorphicArithmetic/Foundations/GaloisRepresentations/LefschetzSigmaSums/HeckeTraceSemi.lean}{UpperTheory.Gal.LocalSmoothExtension.GroupTop}
\leanref{Development/AutomorphicArithmetic/Foundations/GaloisRepresentations/LefschetzSigmaSums/HeckeTraceSemiHolds.lean}{UpperTheory.Gal.LocalSmoothExtension.GroupTop_of}
\leanref{Development/AutomorphicArithmetic/Foundations/GaloisRepresentations/LefschetzSigmaSums/PGroupTop.lean}{UpperTheory.Gal.LocalSmoothExtension.groupTop_of}
\leanref{Development/AutomorphicArithmetic/Foundations/GaloisRepresentations/LefschetzSigmaSums/InvPol.lean}{UpperTheory.Gal.LocalSmoothExtension.IsWt}

\begin{proof}
The rank $n$ is odd and $v$ is the only infinite place of signature $(1,n-1)$. Since $f$ is
elliptic, Proposition~\ref{gal-count-weighted} applies to $f$ and $f_p=\mathbf 1_{K_p}$ and gives
\[
 \sum_\alpha c_\alpha T_v(w_{\alpha,y};f,\mathbf 1_{K_p})
 =\sum_{[Q]}\Bigl(\sum_\alpha c_\alpha\sum_{\mu}w_{\alpha,y}(P_\gamma,\mu)\Bigr)O_Q(f,\mathbf 1_{K_p}),
\]
over the classes of definite pairs with irreducible $P_\gamma=P$ and the roots $\mu$ of $v(P)$. For
such a pair, $P(0)\neq0$, every root $\mu$ of $v(P)$ has $|\mu|^2=c$, and $\mu\mapsto\mu^\vee$ is a
bijection from the roots of $v(P)$ onto the roots of $\tau_1(P^\vee)$, which all have absolute value
$c^{-1/2}$; and $P^\vee$ is separable of degree $n$. So Lemma~\ref{gal-count-root-weights} for
$\chi=P^\vee$ identifies the bracket with $S_\lambda(P)\sum_iy^{h_i}$. The remaining sum
$\sum_{[Q]}S_\lambda(P_\gamma)O_Q(f,\mathbf 1_{K_p})$ is
$\operatorname{tr}_{\mathrm{def}}(\xi;f,\mathbf 1_{K_p})$ by Lemma~\ref{gal-count-definite-trace} with
$\eta=S_\lambda$, whose vanishing hypothesis is part of the ellipticity of $f$.
\end{proof}

\subsection{Fixed points and packs from plans}\label{subsection-gal-fixed-points}

\begin{situation}\label{gal-fixed-situation}
Let $p$ be a prime and take a unitary PEL datum of simple type for a CM field $E$ at $p$
(Definition~\ref{gal-pel-datum}) of rank $n$, with lattice $\Lambda$, form $\langle\,,\rangle$,
element $\omega$, embedding $\tau_1$ and group $G$. Put $V=\Lambda\otimes\mathbf Q$ and
$q=n[E:\mathbf Q]/2$. Present $G$ as in Situation~\ref{gal-count-situation} with $F=E$ and $v$ the
infinite place of signature $(1,n-1)$, through an isomorphism of $C$ onto the algebra of $E$-linear
endomorphisms of $V$ under which $x^*$ is the adjoint of $x$ for $\langle\,,\rangle$. For an embedding
$\tau$ of $E$ put $\epsilon(\tau)=1$ if $\tau(\omega)=\overline{\tau_1(\omega)}$ and $\epsilon(\tau)=0$
otherwise, and let $\varrho:\mathcal O_E\to M_q(\mathbf C)$ be a ring homomorphism with
\[
 \operatorname{tr}\varrho(b)=\tau_1(b)+(n-1)\bar\tau_1(b)+n\sum_{\tau\neq\tau_1,\bar\tau_1}\epsilon(\tau)\tau(b)
\]
for $b\in\mathcal O_E$. The \emph{$\tau$-part} of $\mathbf C^q$ is the subspace on which $\varrho(b)$
acts by $\tau(b)$ for every $b$. For a form $\beta$ (Definition~\ref{gal-count-forms}), a \emph{period}
for $\beta$ is an additive map $Q:V\to\mathbf C^q$ with $Q(bw)=\varrho(b)Q(w)$ for $b\in\mathcal O_E$,
whose image spans $\mathbf C^q$ over $\mathbf R$, such that $\langle\beta^{-1}u,w\rangle$ is the
imaginary part of $Q(u)^{\mathrm T}H\,\overline{Q(w)}$ for a positive definite Hermitian matrix $H$.
For $\mu\in\mathbf C^\times$ put $\mu^\vee=\mu^{-1}$ if $v(\omega)=\tau_1(\omega)$ and
$\mu^\vee=\bar\mu^{-1}$ otherwise, as in Definition~\ref{gal-count-weights}(2). For an integer $N\geq1$
prime to $p$, $K_N\subset G(\mathbf A_f^p)$ is the open compact subgroup of level $N$ and
$\mathrm{Sh}_N$ is the Shimura variety of level $N$ of Definition~\ref{gal-pel-datum}.
\end{situation}
\leanref{Development/AutomorphicArithmetic/Foundations/GaloisRepresentations/SphericalCoordinates/IsCoordA.lean}{UpperTheory.Gal.W2.IsSigD}
\leanref{Development/AutomorphicArithmetic/Foundations/GaloisRepresentations/SphericalCoordinates/IsCoordA.lean}{UpperTheory.Gal.W2.epsD}
\leanref{Development/AutomorphicArithmetic/Foundations/GaloisRepresentations/SphericalCoordinates/IsCoordA.lean}{UpperTheory.Gal.W2.formB}
\leanref{Development/AutomorphicArithmetic/Foundations/GaloisRepresentations/SphericalCoordinates/IsCoordA.lean}{UpperTheory.Gal.W2.IsPerQ}
\leanref{Development/AutomorphicArithmetic/Foundations/GaloisRepresentations/SphericalCoordinates/LevelChange.lean}{UpperTheory.Gal.W2.RhoSig}

\begin{lemma}\label{gal-fixed-uniformization}
In Situation~\ref{gal-fixed-situation}, let $N\geq3$ be an integer prime to $p$.
\begin{enumerate}
\item For every form $\beta$ and every vector $z$ positive for $\beta$ there is a period $Q_z$ for
  $\beta$, and every period for $\beta$ is of this form. Periods $Q_z$ and $Q_{z'}$ differ by a
  $\mathbf C$-linear automorphism of $\mathbf C^q$ if and only if $z'\in\mathbf C^\times z$, and for
  $(y,s)\in G_\beta$ the map $Q_z\circ y^{-1}$ is a period $Q_{\psi_v(y)z}$. An element $(y,s)\in
  G_\beta$ fixes $z$ if and only if $Q_z\circ y=A\circ Q_z$ for a $\mathbf C$-linear map $A$ of
  $\mathbf C^q$; the $\tau_1$-part of $\mathbf C^q$ is a line, and such an $A$ acts on it by
  $1/\mu_y(z)^\vee$.
\item To a form $\beta$, with $x$ and $t$ as in Definition~\ref{gal-count-forms}(1), and to
  $(g,g_p)\in G(\mathbf A_f^p)\times G_p$ are attached an $\mathcal O_E$-stable lattice
  $M_{g,g_p}\subset V$, whose closure in $\mathbf A_f^p\otimes V$ is
  $x^{-1}g(\hat{\mathbf Z}^p\otimes\Lambda)$, a positive rational multiple $\langle\,,\rangle_{g,g_p}$
  of $\langle\beta^{-1}\cdot,\cdot\rangle$ that is integral on it, and the level-$N$ structure
  $\Lambda/N\Lambda\to M_{g,g_p}/NM_{g,g_p}$ induced by $x^{-1}g$. For $(y,s)\in G_\beta$, $y$ maps
  $M_{g',g'_p}$ onto $M_{g,g_p}$ compatibly with these forms and level structures if and only if
  $g^{-1}\iota(y)g'\in K_N$ and $g_p^{-1}\iota_p(y)g'_p\in K_p$.
\item There is a bijection $\Upsilon_N$, the complex uniformization of Shimura \cite{Shi63} and Kottwitz
  \cite[\S8]{Kot92b}, from $\mathrm{Sh}_N(\mathbf C)$ onto the disjoint union, over the equivalence
  classes of forms $\beta$, of the sets of points of $\beta$ of level $(K_N,K_p)$
  (Definition~\ref{gal-count-forms}(4)), such that the abelian variety at $\Upsilon_N^{-1}[g,g_p,z]$ is
  $\mathbf C^q/Q_z(M_{g,g_p})$, with $\mathcal O_E$ acting through $\varrho$, Riemann form
  $\langle\,,\rangle_{g,g_p}$ and the level-$N$ structure of (2).
\item Let $N'\geq3$ be a multiple of $N$ prime to $p$, $h\in G(\mathbf A_f^p)$ with
  $hK_{N'}h^{-1}\subset K_N$, and $c_1,c_2:\mathrm{Sh}_{N'}\to\mathrm{Sh}_N$ the change of level and the
  Hecke correspondence of $h$. If $\Upsilon_{N'}(x)=[g,g_p,z]$ for a form $\beta$, then
  $\Upsilon_N(c_1x)=[g,g_p,z]$, with the same abelian variety $\mathbf C^q/Q_z(M_{g,g_p})$, and
  $\Upsilon_N(c_2x)=[gh^{-1},g_p,z]$; moreover $M_{g,g_p}\subset M_{gh^{-1},g_p}$, with index prime to $p$,
  and the isogeny from the abelian variety at $x$ to the one at $c_2x$ that defines the Hecke
  correspondence is induced by the identity of $\mathbf C^q$.
\item Let $x=\Upsilon_N^{-1}[g,g_p,z]$, $M=M_{g,g_p}$, $P$ the restriction of $Q_z$ to $M$, and $H_0$ the
  positive definite Hermitian matrix with $\langle u,w\rangle_{g,g_p}=\operatorname{Im}(P(u)^{\mathrm
  T}H_0\overline{P(w)})$. Let $\mathcal C$ be the space of $C\in M_q(\mathbf C)$ with
  $\varrho(b)C=C\,\overline{\varrho(b)}$ for $b\in\mathcal O_E$ and $C^{\mathrm T}H_0$ symmetric, and
  for $C\in\mathcal C$ put $P_C(u)=P(u)+C\,\overline{P(u)}$. Then $\dim\mathcal C=n-1$. For $C$ in a
  neighbourhood of $0$ in $\mathcal C$, $P_C$ is the restriction to $M$ of a period for $\beta$; let
  $x_C$ be the point whose abelian variety is $\mathbf C^q/P_C(M)$, with the action, Riemann form and
  level structure of $M$. Then $C\mapsto x_C$ is a holomorphic chart of $\mathrm{Sh}_N(\mathbf C)$ at
  $x=x_0$, and for $k\geq0$, $(C,z_1,\ldots,z_k)\mapsto(z_s\bmod P_C(M))_{s\leq k}$, a point over
  $x_C$, gives holomorphic coordinates on the $k$-fold fibre power of the universal abelian variety
  near the fibre over $x$.
\end{enumerate}
\end{lemma}
\leanref{Development/AutomorphicArithmetic/Foundations/GaloisRepresentations/SphericalCoordinates/LevelChange.lean}{UpperTheory.Gal.W2.UnifG}
\leanref{Development/AutomorphicArithmetic/Foundations/GaloisRepresentations/SphericalCoordinates/FixedPointsIntBaseHolds.lean}{UpperTheory.Gal.W2.UnifG_holds}
\leanref{Development/AutomorphicArithmetic/Foundations/GaloisRepresentations/SphericalCoordinates/IsCoordA.lean}{UpperTheory.Gal.W2.IsCoord}
\leanref{Development/AutomorphicArithmetic/Foundations/GaloisRepresentations/SphericalCoordinates/IsCoordA.lean}{UpperTheory.Gal.W2.RDic}
\leanref{Development/AutomorphicArithmetic/Foundations/GaloisRepresentations/SphericalCoordinates/IsCoordA.lean}{UpperTheory.Gal.W2.PDic}
\leanref{Development/AutomorphicArithmetic/Foundations/GaloisRepresentations/SphericalCoordinates/LevelChange.lean}{UpperTheory.Gal.W2.FixLine}
\leanref{Development/AutomorphicArithmetic/Foundations/GaloisRepresentations/SphericalCoordinates/LevelChange.lean}{UpperTheory.Gal.W2.LatIso}
\leanref{Development/AutomorphicArithmetic/Foundations/GaloisRepresentations/SphericalCoordinates/HeckeFixedPointCMCycleClass.lean}{UpperTheory.Gal.W2.OneAt}
\leanref{Development/AutomorphicArithmetic/Foundations/GaloisRepresentations/SphericalCoordinates/HeckeFixedPointCMCycleClass.lean}{UpperTheory.Gal.W2.HeckeCoord}
\leanref{Development/AutomorphicArithmetic/Foundations/GaloisRepresentations/SphericalCoordinates/LevelChange.lean}{UpperTheory.Gal.W2.SheetOf}
\leanref{Development/AutomorphicArithmetic/Foundations/GaloisRepresentations/SphericalCoordinates/RootSumWOff.lean}{UpperTheory.Gal.W2.CspDim}
\leanref{Development/AutomorphicArithmetic/Foundations/GaloisRepresentations/SphericalCoordinates/PerCApply.lean}{UpperTheory.Gal.W2.perC}

\begin{lemma}\label{gal-fixed-stabilizer}
In Situation~\ref{gal-fixed-situation}, let $N\geq3$ be an integer prime to $p$. Then for every form
$\beta$, every point of $\beta$ of level $(K_N,K_p)$ has trivial stabilizer
(Definition~\ref{gal-count-forms}(4)).
\end{lemma}
\leanref{Development/AutomorphicArithmetic/Foundations/GaloisRepresentations/SphericalCoordinates/LevelChange.lean}{UpperTheory.Gal.W2.StabTriv}
\leanref{Development/AutomorphicArithmetic/Foundations/GaloisRepresentations/SphericalCoordinates/PStabTriv.lean}{UpperTheory.Gal.W2.stabTriv_of}
\leanref{Development/AutomorphicArithmetic/Foundations/GaloisRepresentations/SphericalCoordinates/LevelChange.lean}{UpperTheory.Gal.W2.Rigid}
\leanref{Development/AutomorphicArithmetic/Foundations/GaloisRepresentations/SphericalCoordinates/PRigid.lean}{UpperTheory.Gal.W2.rigid_of}

\begin{proof}
Let $(y,s)$ be in the stabilizer of $(g,g_p,z)$, and put $M=M_{g,g_p}$. By
Lemma~\ref{gal-fixed-uniformization}(2), with $(g',g'_p)=(g,g_p)$, $y$ restricts to an automorphism of
$M$ that preserves $\langle\,,\rangle_{g,g_p}$ and is the identity modulo $N$. By
Lemma~\ref{gal-fixed-uniformization}(1), $Q_z\circ y=A\circ Q_z$ with $A$ $\mathbf C$-linear, so $y$
commutes with the complex structure $J$ that $Q_z$ puts on $V\otimes\mathbf R$. Hence $y$ preserves the
quadratic form $u\mapsto\langle Ju,u\rangle_{g,g_p}$ on $M$, which is positive definite. An automorphism
of a lattice that preserves a positive definite quadratic form and is the identity modulo an integer
$N\geq3$ is the identity, by Serre's lemma \cite{Ser61}. So $y$ is the identity on $M$, which spans $V$;
hence $y=1$ in $C$, and $\beta=y\beta y^*=s\beta$ gives $s=1$.
\end{proof}

\begin{lemma}\label{gal-fixed-base}
In Situation~\ref{gal-fixed-situation}, let $N\mid N'$ be integers at least $3$ and prime to $p$,
$h\in G(\mathbf A_f^p)$ with $hK_{N'}h^{-1}\subset K_N$, and $c_1,c_2:\mathrm{Sh}_{N'}\to\mathrm{Sh}_N$
the change of level and the Hecke correspondence of $h$. Let $x\in\mathrm{Sh}_{N'}(\mathbf C)$ with
$\Upsilon_{N'}(x)=[g,g_p,z]$ for a form $\beta$ (Lemma~\ref{gal-fixed-uniformization}(3)). If
$c_1(x)=c_2(x)$, then some $(y,s)\in G_\beta$ fixes $z$ and satisfies $hg^{-1}\iota(y)g\in K_N$ and
$g_p^{-1}\iota_p(y)g_p\in K_p$. At most one $(y,s)\in G_\beta$ has these three properties.
\end{lemma}
\leanref{Development/AutomorphicArithmetic/Foundations/GaloisRepresentations/SphericalCoordinates/HeckeFixedPointCMCycleClass.lean}{UpperTheory.Gal.W2.FixFwd}
\leanref{Development/AutomorphicArithmetic/Foundations/GaloisRepresentations/SphericalCoordinates/PFixFwd.lean}{UpperTheory.Gal.W2.fixFwd_of}
\leanref{Development/AutomorphicArithmetic/Foundations/GaloisRepresentations/SphericalCoordinates/LevelChange.lean}{UpperTheory.Gal.W2.LatForm}
\leanref{Development/AbelianVarietyTheory/Proofs/DivisorProofs/ComplexUniformisationsSamePointDifferByIsometryStatement.lean}{GlobalTheory.T.complexUniformisations_same_point_differ_by_isometry}

\begin{proof}
Put $M=M_{g,g_p}$ and $M''=M_{gh^{-1},g_p}$. By Lemma~\ref{gal-fixed-uniformization}(4), the abelian
varieties at $c_1(x)$ and $c_2(x)$ are $\mathbf C^q/Q_z(M)$ and $\mathbf C^q/Q_z(M'')$, with the forms
and level-$N$ structures of Lemma~\ref{gal-fixed-uniformization}(2). Suppose $c_1(x)=c_2(x)$. Then the
two are isomorphic with their actions of $\mathcal O_E$, Riemann forms and level structures. An
isomorphism of complex tori $\mathbf C^q/L\to\mathbf C^q/L'$ is induced by a $\mathbf C$-linear map
carrying $L$ onto $L'$. So there are an $\mathcal O_E$-linear isomorphism $w:M\to M''$, compatible with
the Riemann forms and the level structures, and a $\mathbf C$-linear $A$ with $Q_z\circ w=A\circ Q_z$ on
$M$. Extended to $V$, $w$ is $E$-linear, so $w=y$ for some $y\in C^\times$. Since $w$ carries
$\langle\,,\rangle_{g,g_p}$ to $\langle\,,\rangle_{gh^{-1},g_p}$, and both are positive rational
multiples of $\langle\beta^{-1}\cdot,\cdot\rangle$, $y^*\beta^{-1}y=s^{-1}\beta^{-1}$ for a rational
$s>0$; that is, $y\beta y^*=s\beta$ and $(y,s)\in G_\beta$. By
Lemma~\ref{gal-fixed-uniformization}(1), $(y,s)$ fixes $z$. By Lemma~\ref{gal-fixed-uniformization}(2),
since $y$ maps $M_{g,g_p}$ onto $M_{gh^{-1},g_p}$ compatibly with the forms and level structures,
$(gh^{-1})^{-1}\iota(y)g=hg^{-1}\iota(y)g$ lies in $K_N$ and $g_p^{-1}\iota_p(y)g_p$ in $K_p$.

If $(y,s)$ and $(y',s')$ both have the three properties, then $y^{-1}y'$ fixes $z$,
$g^{-1}\iota(y^{-1}y')g=(hg^{-1}\iota(y)g)^{-1}(hg^{-1}\iota(y')g)$ lies in $K_N$, and
$g_p^{-1}\iota_p(y^{-1}y')g_p$ lies in $K_p$. So $(y,s)^{-1}(y',s')$ is in the stabilizer of
$(g,g_p,z)$ at level $(K_N,K_p)$, which is trivial by Lemma~\ref{gal-fixed-stabilizer}.
\end{proof}

\begin{lemma}\label{gal-fixed-torus}
Let $W$ be a finite-dimensional complex vector space, $\Gamma\subset W$ a discrete subgroup that spans
$W$ over $\mathbf R$, and $F:W\to W$ a $\mathbf C$-linear map with $F(\Gamma)\subset\Gamma$ and
$\det(1-F)\neq0$. Then the set of $z\in W/\Gamma$ with $F(z)=z$ is finite, with
$\det(1-F)\,\overline{\det(1-F)}$ elements.
\end{lemma}
\leanref{Development/AutomorphicArithmetic/Foundations/GaloisRepresentations/SphericalCoordinates/DetOfCharpoly.lean}{UpperTheory.Gal.W2.TorusCount}
\leanref{Development/AutomorphicArithmetic/Foundations/GaloisRepresentations/SphericalCoordinates/PTorusCount.lean}{UpperTheory.Gal.W2.torusCount_of}
\leanref{Development/AutomorphicArithmetic/Foundations/GaloisRepresentations/SphericalCoordinates/DetOfCharpoly.lean}{UpperTheory.Gal.W2.DetReal}

\begin{proof}
The fixed points are the classes of the elements of $(1-F)^{-1}\Gamma$, so they form the group
$(1-F)^{-1}\Gamma/\Gamma$, which $1-F$ maps isomorphically onto $\Gamma/(1-F)\Gamma$. The order of
this group is the absolute value of the determinant of $1-F$ as an $\mathbf R$-linear map of $W$, and
the determinant of a $\mathbf C$-linear map over $\mathbf R$ is the square of the absolute value of its
determinant over $\mathbf C$.
\end{proof}

\begin{lemma}\label{gal-fixed-eigenvalues}
In Situation~\ref{gal-fixed-situation}, let $M\subset V$ be an $\mathcal O_E$-stable lattice,
$\langle\,,\rangle_M$ a nondegenerate alternating form on $M$ with values in $\mathbf Z$, and
$P:M\to\mathbf C^q$ an injective map with $P(bu)=\varrho(b)P(u)$ for $b\in\mathcal O_E$, whose image
is discrete and spans $\mathbf C^q$ over $\mathbf R$, such that $\langle u,w\rangle_M=\operatorname{Im}
(P(u)^{\mathrm T}H_0\overline{P(w)})$ for a positive definite Hermitian matrix $H_0$. Let $\mathcal C$
be the space of $C\in M_q(\mathbf C)$ with $\varrho(b)C=C\,\overline{\varrho(b)}$ for
$b\in\mathcal O_E$ and $C^{\mathrm T}H_0$ symmetric. Let $\gamma$ be an $\mathcal O_E$-linear
endomorphism of $M$ and $A\in M_q(\mathbf C)$ with $P\circ\gamma=A\circ P$ and
$\langle\gamma u,\gamma w\rangle_M=m\langle u,w\rangle_M$ for a positive integer $m$, and suppose that
the minimal polynomial $\chi\in E[X]$ of $\gamma$ on $M\otimes\mathbf Q$ is separable of degree $n$.
Then:
\begin{enumerate}
\item the $\tau_1$-part of $\mathbf C^q$ is a line, on which $A$ acts by a root $\mu$ of $\tau_1(\chi)$;
\item $A$ is invertible, $T(C)=AC\bar A^{-1}$ defines an endomorphism $T$ of $\mathcal C$,
  $\dim\mathcal C=n-1$, and $\det(1-tT)=\prod_r(1-t\mu/r)$ for $t\in\mathbf C$, the product over the
  roots $r\neq\mu$ of $\tau_1(\chi)$; in particular $\det(1-T)\neq0$;
\item for $k\geq0$ and $\alpha\in\mathcal O_E^k$ with every $\alpha_s$ divisible by $2$, the
  endomorphism $F_\alpha(z_1,\ldots,z_k)=(A\varrho(\alpha_1)z_1,\ldots,A\varrho(\alpha_k)z_k)$ of
  $(\mathbf C^q)^k$ has $\det(1-F_\alpha)\neq0$, and for every $t\in\mathbf C$
  \[
   \frac{\det(1-tT)}{\det(1-T)}\,\det(1-tF_\alpha)\,\overline{\det(1-F_\alpha)}=L_{\alpha,t}(\chi,\mu),
  \]
  the function of Definition~\ref{gal-count-weights}(1) for $E$, $\tau_1$, $\omega$ and $n$.
\end{enumerate}
\end{lemma}
\leanref{Development/AutomorphicArithmetic/Foundations/GaloisRepresentations/SphericalCoordinates/RootSumWOff.lean}{UpperTheory.Gal.W2.EigTop}
\leanref{Development/AutomorphicArithmetic/Foundations/GaloisRepresentations/SphericalCoordinates/PEigTop.lean}{UpperTheory.Gal.W2.eigTop_of}
\leanref{Development/AutomorphicArithmetic/Foundations/GaloisRepresentations/SphericalCoordinates/DetOfCharpoly.lean}{UpperTheory.Gal.W2.AdExists}
\leanref{Development/AutomorphicArithmetic/Foundations/GaloisRepresentations/SphericalCoordinates/RootSumWOff.lean}{UpperTheory.Gal.W2.RootsFull}
\leanref{Development/AutomorphicArithmetic/Foundations/GaloisRepresentations/SphericalCoordinates/RootSumWOff.lean}{UpperTheory.Gal.W2.BallChar}
\leanref{Development/AutomorphicArithmetic/Foundations/GaloisRepresentations/SphericalCoordinates/RootSumWOff.lean}{UpperTheory.Gal.W2.FibW}
\leanref{Development/AutomorphicArithmetic/Foundations/GaloisRepresentations/SphericalCoordinates/DiagAct.lean}{UpperTheory.Gal.W2.locW}

\begin{proof}
Since $P(M)$ is discrete and spans $\mathbf C^q$ over $\mathbf R$, $P$ extends to an $\mathbf R$-linear
isomorphism $M\otimes\mathbf R\to\mathbf C^q$, and $M\otimes\mathbf C$ is identified with
$\mathbf C^q\oplus\overline{\mathbf C^q}$, where $b\in\mathcal O_E$ acts by $\varrho(b)\oplus
\overline{\varrho(b)}$ and $\gamma$ by $A\oplus\bar A$.

The identity $\langle\gamma u,\gamma w\rangle_M=m\langle u,w\rangle_M$ and $P\circ\gamma=A\circ P$ give
$A^{\mathrm T}H_0\bar A=mH_0$, so $A$ is invertible, and $A$ commutes with $\varrho(\mathcal O_E)$
because $\gamma$ is $\mathcal O_E$-linear. Then $\varrho(b)T(C)=T(C)\overline{\varrho(b)}$ for
$C\in\mathcal C$, and $T(C)^{\mathrm T}H_0=m\bar A^{-\mathrm T}(C^{\mathrm T}H_0)\bar A^{-1}$ is
symmetric, so $T$ preserves $\mathcal C$.

Since $\chi(\gamma)=0$, on the $\tau$-part the matrix $A$ is killed by the separable polynomial
$\tau(\chi)$. So $A$ and $\varrho(\mathcal O_E)$ have a common basis of eigenvectors $e_1,\ldots,e_q$:
$\varrho(b)e_i=\tau_i(b)e_i$ and $Ae_i=a_ie_i$. By the trace of $\varrho$, the number of $i$ with
$\tau_i=\tau$ is $1$ for $\tau_1$, $n-1$ for $\bar\tau_1$, and $n\epsilon(\tau)$ for the other $\tau$.
For each $\tau$, the $a_i$ with $\tau_i=\tau$ together with the conjugates of the $a_i$ with
$\tau_i=\bar\tau$ are the roots of $\tau(\chi)$, with multiplicity. Indeed, for $b\in\mathcal O_E$ the
trace of $b\gamma^j$ on $M$ is an integer and equals $\sum_\tau\tau(b)p_j(\tau)$, where $p_j(\tau)$ is
the sum of the $j$-th powers of these $n$ numbers. As this holds for every $b$, the trace form of $E$
gives $p_j(\tau)=\tau(c_j)$ for some $c_j\in E$, and Newton's identities give a monic
$\chi'\in E[X]$ of degree $n$ such that, for each $\tau$, the roots of $\tau(\chi')$ are these $n$
numbers. Then $\chi'(A)$ vanishes on every $\tau$-part, so $\chi'(\gamma)=0$, $\chi$ divides $\chi'$,
and $\chi'=\chi$.

(1) The $\tau_1$-part is the line spanned by the unique $e_{i_0}$ with $\tau_{i_0}=\tau_1$, and $A$ acts
on it by $\mu=a_{i_0}$. By the above, the roots of $\tau_1(\chi)$ are $\mu$ and the numbers $\bar a_j$
for the $n-1$ indices $j$ with $\tau_j=\bar\tau_1$; as $\chi$ is separable, these $\bar a_j$ are the
roots $r\neq\mu$.

(2) Write $C\in\mathcal C$ in the basis $(e_i)$ as $C'=S^{-1}C\bar S$, where $S$ has columns $e_i$. The
condition $\varrho(b)C=C\,\overline{\varrho(b)}$ allows a nonzero entry $C'_{ij}$ only when
$\tau_i=\bar\tau_j$. For $\tau\neq\tau_1,\bar\tau_1$ at most one of $\tau$ and $\bar\tau$ occurs among
the $\tau_i$, since $\epsilon(\tau)+\epsilon(\bar\tau)=1$, and $\tau_1$ occurs once; so the allowed
entries link $i_0$ with the $n-1$ indices $j$ of type $\bar\tau_1$, and the symmetry of
$C^{\mathrm T}H_0$ makes the entries $C'_{i_0j}$ coordinates on $\mathcal C$. Hence
$\dim\mathcal C=n-1$. In these coordinates $T$ multiplies $C'_{i_0j}$ by $a_{i_0}/\bar a_j=\mu/r$
with $r=\bar a_j$, so $\det(1-tT)=\prod_r(1-t\mu/r)$, which is nonzero at $t=1$ because $r\neq\mu$.

(3) In the basis $(e_i)$ on each factor, $\det(1-tF_\alpha)=\prod_s\prod_i(1-t\,\tau_i(\alpha_s)a_i)$.
Each $a_i$ is an eigenvalue of $\gamma$ on $M\otimes\mathbf C$, hence an algebraic integer; so
$\tau_i(\alpha_s)a_i\neq1$, since otherwise $1/2=\tau_i(\alpha_s/2)a_i$ would be an algebraic integer.
So $\det(1-F_\alpha)\neq0$. Fix $s$ and put $x=\alpha_s$. In
$\prod_i(1-t\tau_i(x)a_i)\,\overline{\prod_i(1-\tau_i(x)a_i)}$, the index $i_0$ gives
$(1-t\tau_1(x)\mu)(1-\overline{\tau_1(x)\mu})$; the indices $j$ of type $\bar\tau_1$, with $r=\bar a_j$,
give $(1-t\,\overline{\tau_1(x)r})(1-\tau_1(x)r)$; and for $\tau\neq\tau_1,\bar\tau_1$ with
$\epsilon(\tau)=1$, the $a_i$ of type $\tau$ are the $n$ roots $r$ of $\tau(\chi)$, which give
$(1-t\tau(x)r)(1-\bar\tau(x)\bar r)$, where the $\bar r$ are the roots of $\bar\tau(\chi)$ and
$\epsilon(\bar\tau)=0$. The product is $\phi_x(t;\chi,\mu)$. By (2),
$\det(1-tT)/\det(1-T)=\prod_r(r-t\mu)/(r-\mu)=b(t;\chi,\mu)$, and multiplying over $s$ gives
$L_{\alpha,t}(\chi,\mu)$.
\end{proof}

\begin{lemma}\label{gal-fixed-fibre}
In Situation~\ref{gal-fixed-situation}, let $N\mid N'$ be integers at least $3$ and prime to $p$, and
$h\in G(\mathbf A_f^p)$ with $hK_{N'}h^{-1}\subset K_N$ such that every element of $h^{-1}K_N$ is
regular. Let $k\geq0$ and $\alpha\in\mathcal O_E^k$ with every $\alpha_s$ divisible by $2$. Let $Z_N$
and $Z_{N'}$ be the $k$-fold fibre powers over $\mathbf C$ of the universal abelian schemes over
$\mathrm{Sh}_N$ and $\mathrm{Sh}_{N'}$, $c_1:Z_{N'}\to Z_N$ the change of level, and $c_2:Z_{N'}\to Z_N$
the Hecke correspondence of $h$ followed by multiplication by $\alpha_s$ on the $s$-th factor. Let
$x'\in\mathrm{Sh}_{N'}(\mathbf C)$ with $\Upsilon_{N'}(x')=[g,g_p,z]$ for a form $\beta$
(Lemma~\ref{gal-fixed-uniformization}(3)), and let $(y,s)\in G_\beta$ fix $z$ with
$hg^{-1}\iota(y)g\in K_N$ and $g_p^{-1}\iota_p(y)g_p\in K_p$. Then the set $\Phi_{x'}$ of
$x\in Z_{N'}(\mathbf C)$ over $x'$ with $c_1(x)=c_2(x)$ is finite and not empty, and at each $x\in
\Phi_{x'}$ the endomorphism $A_x$ of $T_{c_1(x)}Z_N$ induced by $c_2\circ c_1^{-1}$ near $x$ satisfies
$\det(1-A_x)\neq0$. For every $t\in\mathbf C$,
\[
 \sum_{x\in \Phi_{x'}}\frac{\det(1-tA_x)}{\det(1-A_x)}=w_{\alpha,t}(P_y,\mu_y(z)),
\]
the Lefschetz weight of Definition~\ref{gal-count-weights}(2) for $E$, $\tau_1$, $\omega$, $n$ and $v$.
\end{lemma}
\leanref{Development/AutomorphicArithmetic/Foundations/GaloisRepresentations/SphericalCoordinates/HeckeFixedPointCMCycleClass.lean}{UpperTheory.Gal.W2.FixLoc}
\leanref{Development/AutomorphicArithmetic/Foundations/GaloisRepresentations/SphericalCoordinates/PFixLoc.lean}{UpperTheory.Gal.W2.fixLoc_of}
\leanref{Development/AutomorphicArithmetic/Foundations/GaloisRepresentations/SphericalCoordinates/HeckeFixedPointCMCycleClass.lean}{UpperTheory.Gal.W2.FixGa}
\leanref{Development/AutomorphicArithmetic/Foundations/GaloisRepresentations/SphericalCoordinates/HeckeFixedPointCMCycleClass.lean}{UpperTheory.Gal.W2.FixEq}
\leanref{Development/AutomorphicArithmetic/Foundations/GaloisRepresentations/SphericalCoordinates/DetOfCharpoly.lean}{UpperTheory.Gal.W2.CorrChart}
\leanref{Development/AutomorphicArithmetic/Foundations/GaloisRepresentations/SphericalCoordinates/DetOfCharpoly.lean}{UpperTheory.Gal.W2.LocSum}
\leanref{Development/AutomorphicArithmetic/Foundations/GaloisRepresentations/SphericalCoordinates/PLocSum.lean}{UpperTheory.Gal.W2.locSum_of}
\leanref{Development/AutomorphicArithmetic/Foundations/GaloisRepresentations/SphericalCoordinates/LevelChange.lean}{UpperTheory.Gal.W2.RegSep}
\leanref{Development/AutomorphicArithmetic/Foundations/GaloisRepresentations/SphericalCoordinates/LevelChange.lean}{UpperTheory.Gal.W2.MinBridge}
\leanref{Development/AutomorphicArithmetic/Foundations/GaloisRepresentations/SphericalCoordinates/LevelChange.lean}{UpperTheory.Gal.W2.InvMin}
\leanref{Development/AutomorphicArithmetic/Foundations/GaloisRepresentations/SphericalCoordinates/LevelChange.lean}{UpperTheory.Gal.W2.IntMult}
\leanref{Development/AutomorphicArithmetic/Foundations/GaloisRepresentations/SphericalCoordinates/PIntMult.lean}{UpperTheory.Gal.W2.intMult_of}

\begin{proof}
Put $M=M_{g,g_p}$, $M''=M_{gh^{-1},g_p}$ and $P=Q_z|_M$, and let $H_0$, $\mathcal C$ and $P_C$ be as in
Lemma~\ref{gal-fixed-uniformization}(5) for the point $x'$. Write $\langle\,,\rangle_M$ for
$\langle\,,\rangle_{g,g_p}$. The fibre of $Z_{N'}$ over $x'$ is $(\mathbf C^q/P(M))^k$.

First, the element $(y,s)$ gives an endomorphism of $M$. Since $g^{-1}\iota(y^{-1})(gh^{-1})=
(hg^{-1}\iota(y)g)^{-1}$ lies in $K_N$ and $g_p^{-1}\iota_p(y^{-1})g_p$ in $K_p$,
Lemma~\ref{gal-fixed-uniformization}(2) shows that $y^{-1}$ maps $M''$ onto $M$ compatibly with the forms
and level structures. As $M\subset M''$ by Lemma~\ref{gal-fixed-uniformization}(4), $\gamma=y^{-1}|_M$
is an $\mathcal O_E$-linear endomorphism of $M$, with $\langle\gamma u,\gamma w\rangle_M=
\langle u,w\rangle_{M''}=m\langle u,w\rangle_M$ for a positive rational $m$. Here $m$ is an integer.
Indeed, the values of $\langle\,,\rangle_M$ generate an ideal $c\mathbf Z$ of $\mathbf Z$, with $c\neq0$ as
the form is nondegenerate. Writing $c=\sum_ia_i\langle u_i,w_i\rangle_M$ with $a_i\in\mathbf Z$ gives
$mc=\sum_ia_i\langle\gamma u_i,\gamma w_i\rangle_M\in c\mathbf Z$, so $m\in\mathbf Z$. The element $y^{-1}$ fixes $z$ with eigenvalue $\mu_y(z)^{-1}$, so by
Lemma~\ref{gal-fixed-uniformization}(1), $P\circ\gamma=A\circ P$ for a $\mathbf C$-linear $A$, which acts
on the $\tau_1$-part of $\mathbf C^q$ by $1/(\mu_y(z)^{-1})^\vee=\mu_y(z)^\vee$.

Second, the minimal polynomial of $\gamma$. Since $y$ fixes the positive vector $z$, $\psi_v(y)$
preserves the line of $z$ and its orthogonal complement for $H_\beta$, on which $H_\beta$ is negative
definite, so $\psi_v(y)$ is a similitude of a positive definite Hermitian form and is semisimple; and
$g^{-1}\iota(y)g$ lies in $h^{-1}K_N$, so it is regular and the centralizer of $y$ is commutative. Hence
$P_y$ is separable of degree $n$. The minimal polynomial of $\gamma$ over $E$ is that of $y^{-1}$ in $C$,
which is $P_y^\vee$, again separable of degree $n$. So Lemma~\ref{gal-fixed-eigenvalues} applies to $M$,
$\langle\,,\rangle_M$, $P$, $H_0$, $\gamma$ and $A$, with $\chi=P_y^\vee$ and, by the first step,
$\mu=\mu_y(z)^\vee$. Let $T(C)=AC\bar A^{-1}$ on $\mathcal C$ and
$F_\alpha(z_1,\ldots,z_k)=(A\varrho(\alpha_s)z_s)_s$ be as there; then $\det(1-T)\neq0$ and
$\det(1-F_\alpha)\neq0$.

Third, the correspondence in the coordinates of Lemma~\ref{gal-fixed-uniformization}(5) at $x'$, and at
$c_1(x')$, which has the same $M$, $P$ and $H_0$. For $C\in\mathcal C$ near $0$ let $x'_C$ be the point
of $\mathrm{Sh}_{N'}$ with coordinate $C$. Then $P_{T(C)}\circ\gamma=A\circ P_C$ on $M$, hence
$P_{T(C)}(y^{-1}w)=AP_C(w)$ for $w\in V$, and $AP_C(M'')=P_{T(C)}(M)$. By
Lemma~\ref{gal-fixed-uniformization}(4), applied at $x'_C$, $c_2$ sends the point $(z_s\bmod
P_C(M))_s$ over $x'_C$ to the point $(\varrho(\alpha_s)z_s\bmod P_C(M''))_s$ of the $k$-th power of
$\mathbf C^q/P_C(M'')$, and $A$ identifies this abelian variety, with its structures, with
$\mathbf C^q/P_{T(C)}(M)$, the abelian variety at $c_1(x'_{T(C)})$. So, in these coordinates,
\[
 c_2(C,z)=c_1\bigl(T(C),F_\alpha(z)\bigr).
\]

Fourth, the count. By Lemma~\ref{gal-fixed-uniformization}(4), $c_1$ is the identity in these
coordinates, so it is injective on the fibre over $x'$, and a point $(0,z)$ over $x'$ is fixed if and
only if $F_\alpha(z)\equiv z$ modulo $P(M)^k$. Since $F_\alpha$ maps $P(M)^k$ into itself,
Lemma~\ref{gal-fixed-torus} shows that $\Phi_{x'}$ is finite with $|\det(1-F_\alpha)|^2$ elements; it
contains $(0,0)$. Near each $x\in \Phi_{x'}$, in these coordinates, $c_2\circ c_1^{-1}$ is
$(C,z)\mapsto(T(C),F_\alpha(z))$ followed by a translation of $z$ by an element of $P_{T(C)}(M)^k$ that
depends holomorphically on $C$. So its differential is block triangular with diagonal blocks $T$ and
$F_\alpha$, and the characteristic polynomial of $A_x$ is the product of those of $T$ and $F_\alpha$.
Hence $\det(1-tA_x)=\det(1-tT)\det(1-tF_\alpha)$, $\det(1-A_x)\neq0$, and
\[
 \sum_{x\in \Phi_{x'}}\frac{\det(1-tA_x)}{\det(1-A_x)}
 =\frac{\det(1-tT)}{\det(1-T)}\,\det(1-tF_\alpha)\,\overline{\det(1-F_\alpha)}.
\]
By Lemma~\ref{gal-fixed-eigenvalues}(3) this is $L_{\alpha,t}(P_y^\vee,\mu_y(z)^\vee)$, which is
$w_{\alpha,t}(P_y,\mu_y(z))$.
\end{proof}

\begin{lemma}\label{gal-fixed-level}
In Situation~\ref{gal-count-situation}, let $K\subset G(\mathbf A_f^p)$ be an open compact subgroup,
$h\in G(\mathbf A_f^p)$, and $K'\subset K\cap h^{-1}Kh$ an open subgroup. Suppose that every point of
every form of level $(K,K_p)$ has trivial stabilizer, and that only finitely many points $[g,g_p,z]$
of level $(K',K_p)$ of the forms $\beta$ admit an element $(y,s)\in G_\beta$ that fixes $z$ with
$hg^{-1}\iota(y)g\in K$ and $g_p^{-1}\iota_p(y)g_p\in K_p$. Then:
\begin{enumerate}
\item only finitely many triples consist of a class of forms, a point of its chosen representative
  $\beta$ of level $(K,K_p)$, with chosen representative $(g,g_p,z)$, and an element $(y,s)\in G_\beta$
  that fixes $z$ with $g^{-1}\iota(y)g\in Kh^{-1}K$ and $g_p^{-1}\iota_p(y)g_p\in K_p$;
\item for every function $\eta:F[X]\times\mathbf C\to\mathbf C$, the sum
  $T_v(\eta;\mathbf 1_{h^{-1}K},\mathbf 1_{K_p})$ of Definition~\ref{gal-count-forms}(5) at level
  $(K',K_p)$ is $[K\cap h^{-1}Kh:K']$ times the sum $T_v(\eta;\mathbf 1_{Kh^{-1}K},\mathbf 1_{K_p})$ at
  level $(K,K_p)$.
\end{enumerate}
\end{lemma}
\leanref{Development/AutomorphicArithmetic/Foundations/GaloisRepresentations/SphericalCoordinates/LevelChange.lean}{UpperTheory.Gal.W2.LevelChange}
\leanref{Development/AutomorphicArithmetic/Foundations/GaloisRepresentations/SphericalCoordinates/PLevelChange.lean}{UpperTheory.Gal.W2.levelChange_of}
\leanref{Development/AutomorphicArithmetic/Foundations/GaloisRepresentations/SphericalCoordinates/LevelChange.lean}{UpperTheory.Gal.W2.EigSame}
\leanref{Development/AutomorphicArithmetic/Foundations/GaloisRepresentations/IsogenyEtaleTranslation/EtCancel.lean}{UpperTheory.Gal.FixZ.CxFinOf}
\leanref{Development/AutomorphicArithmetic/Foundations/GaloisRepresentations/IsogenyEtaleTranslation/PCxFinOf.lean}{UpperTheory.Gal.FixZ.cxFinOf_of}

\begin{proof}
Since $K'\subset K$, the stabilizers at level $(K',K_p)$ are trivial as well. For $u\in K'$,
$huh^{-1}$ and $u$ lie in $K$, so $\mathbf 1_{h^{-1}K}$ is invariant under conjugation by $K'$. Fix a
form $\beta$, with its $x$ and $t$. For a triple $a=(g,g_p,z)$ and a subset $S\subset G(\mathbf A_f^p)$,
let $W_S(a)$ be the sum of $\eta(P_y,\mu_y(z))$ over the $(y,s)\in G_\beta$ that fix $z$ with
$g^{-1}\iota(y)g\in S$ and $g_p^{-1}\iota_p(y)g_p\in K_p$. As the stabilizers are trivial, the term of
a point in $T_v(\eta;\mathbf 1_S,\mathbf 1_{K_p})$ is $W_S$ at its chosen representative.

(i) The value $W_{h^{-1}K}(a)$ is the same at all representatives $a$ of a point of level $(K',K_p)$.
Replacing $(g,g_p,z)$ by $(\iota(\delta)gk,\iota_p(\delta)g_pk_p,\lambda\psi_v(\delta)z)$, with $\delta\in
G_\beta$, $k\in K'$, $k_p\in K_p$ and $\lambda\in\mathbf C^\times$, replaces each $(y,s)$ by its
conjugate by $\delta$, which has the same minimal polynomial and the same eigenvalue on
$\psi_v(\delta)z$, and the conditions are kept because $\mathbf 1_{h^{-1}K}$ and $\mathbf 1_{K_p}$ are
invariant under conjugation by $K'$ and $K_p$. In the same way $W_{Kh^{-1}K}$ is constant on the points
of level $(K,K_p)$.

(ii) The sum $W_{h^{-1}K}(a)$ has at most one term: for two of its elements $(y,s)$ and $(y',s')$,
$g^{-1}\iota(y^{-1}y')g=(hg^{-1}\iota(y)g)^{-1}(hg^{-1}\iota(y')g)$ lies in $K$, so
$(y,s)^{-1}(y',s')$ is in the stabilizer of $a$ at level $(K,K_p)$, which is trivial.

(iii) Choose representatives of $K/K'$. Sending a point of level $(K,K_p)$, with chosen representative
$(g,g_p,z)$, and a coset $uK'$ to the point $[gu,g_p,z]$ of level $(K',K_p)$ is a bijection. It is onto:
a point $[g',g'_p,z']$ of level $(K',K_p)$ lies over a point of level $(K,K_p)$ whose chosen
representative $(g,g_p,z)$ is $(\iota(\delta)g'k,\iota_p(\delta)g'_pk_p,\lambda\psi_v(\delta)z')$ with
$k\in K$, and if $uK'=k^{-1}K'$ then $[g',g'_p,z']=[gu,g_p,z]$ is the image of that point and $uK'$.
It is one to one: if $[gu,g_p,z]=[gu',g_p,z]$ at level $(K',K_p)$, say $gu'=\iota(\delta)guk'$,
$g_p=\iota_p(\delta)g_pk_p$ and $z=\lambda\psi_v(\delta)z$, then $\delta$ is in the stabilizer of
$(g,g_p,z)$ at level $(K,K_p)$, so $\delta=1$ and $u'K'=uK'$.

(iv) Let $a=(g,g_p,z)$ and let $(y,s)\in G_\beta$ fix $z$, with $g_p^{-1}\iota_p(y)g_p\in K_p$; put
$\delta=g^{-1}\iota(y)g$. The number of cosets $uK'\subset K$ with $h(gu)^{-1}\iota(y)(gu)=hu^{-1}\delta
u\in K$ is $[K\cap h^{-1}Kh:K']$ if $\delta\in Kh^{-1}K$, and $0$ otherwise. Indeed, if
$hu^{-1}\delta u\in K$ then $\delta=uh^{-1}(hu^{-1}\delta u)u^{-1}\in Kh^{-1}K$. If $\delta=ah^{-1}b$
with $a,b\in K$, then $hu^{-1}\delta u=(hu^{-1}ah^{-1})(bu)$ lies in $K$ if and only if
$h(a^{-1}u)h^{-1}\in K$, that is, $a^{-1}u\in K\cap h^{-1}Kh$; so the cosets are those contained in
$a(K\cap h^{-1}Kh)$, and there are $[K\cap h^{-1}Kh:K']$ of them.

For (1), send a triple with point $a=(g,g_p,z)$ and element $(y,s)$, where $g^{-1}\iota(y)g=k_1h^{-1}k_2$
with $k_1,k_2\in K$, to the point $[gk_1,g_p,z]$ of level $(K',K_p)$. Since
$h(gk_1)^{-1}\iota(y)(gk_1)=k_2k_1\in K$, this point is one of the finitely many of the hypothesis, by
(i). The map is one to one: by (iii) the point $[gk_1,g_p,z]$ determines $a$ and the coset $k_1K'$, and
then $(y,s)$ is counted in $W_{h^{-1}K}(gk_1,g_p,z)$, which depends only on the coset $k_1K'$ because
$\mathbf 1_{h^{-1}K}$ is invariant under conjugation by $K'$, and which has at most one term by (ii).

For (2), by (iii) and (i) the left side is the sum, over the points $a=(g,g_p,z)$ of level $(K,K_p)$ and
the cosets $uK'$, of $W_{h^{-1}K}(gu,g_p,z)$. By (ii) and (iv), for fixed $a$ the sum over $u$ is
$\sum_{(y,s)}[K\cap h^{-1}Kh:K']\,\eta(P_y,\mu_y(z))$ over the $(y,s)$ fixing $z$ with
$g^{-1}\iota(y)g\in Kh^{-1}K$ and $g_p^{-1}\iota_p(y)g_p\in K_p$, that is
$[K\cap h^{-1}Kh:K']\,W_{Kh^{-1}K}(a)$. Both sides have finitely many nonzero terms, by the hypothesis
and by (1), and summing over $a$ and over the classes of forms gives (2).
\end{proof}

\begin{lemma}\label{gal-fixed-points-isolated}
Let $Y$ and $Z$ be schemes over $\mathbf C$, with $Y$ quasi-compact and locally of finite type over
$\mathbf C$ and $Z$ separated over $\mathbf C$, and let $c_1,c_2:Y\to Z$ be morphisms over $\mathbf C$.
For a $\mathbf C$-point $y$ of $Y$ with $c_1(y)=c_2(y)$ and $d\geq0$, say that a matrix $A\in M_d(\mathbf
C)$ is a \emph{tangent matrix} at $y$ if there are affine opens $U\subset Z$ and $V\subset Y$ with $y\in V$
and $c_1(V),c_2(V)\subset U$, and functions $t_1,\ldots,t_d$ on $U$, such that the differentials at $y$ of
the functions $c_1^*t_l$ on $V$ form a basis of the cotangent space $\mathfrak m_y/\mathfrak m_y^2$ of $V$
at $y$, and the differential at $y$ of $c_2^*t_j$ is $\sum_lA_{jl}\,d(c_1^*t_l)$ for every $j$. Suppose
that every $\mathbf C$-point $y$ with $c_1(y)=c_2(y)$ has a tangent matrix $A$ with $\det(1-A)\neq0$. Then
only finitely many $\mathbf C$-points $y$ of $Y$ satisfy $c_1(y)=c_2(y)$. The hypothesis that $Z$ is
separated is not used.
\end{lemma}
\leanref{Development/AutomorphicArithmetic/Foundations/GaloisRepresentations/SphericalCoordinates/LevelChange.lean}{UpperTheory.Gal.W2.FixFinite}
\leanref{Development/AutomorphicArithmetic/Foundations/GaloisRepresentations/SphericalCoordinates/PFixFinite.lean}{UpperTheory.Gal.W2.fixFinite_of}
\leanref{Development/AbelianVarietyTheory/Definitions/Divisors/IsGaloisCover.lean}{GlobalTheory.W1.IsTan}

\begin{proof}
First, each such point $y$ has an open neighbourhood that contains no other such point. Let $A$, $U$, $V$
and $t_1,\ldots,t_d$ be as in the statement, let $R$ be the ring of functions on $V$, noetherian since $V$
is of finite type over $\mathbf C$, and let $\mathfrak m\subset R$ be the maximal ideal of $y$. Put
$g_j=c_1^*t_j-c_2^*t_j\in R$. At a $\mathbf C$-point $y'$ of $V$ with $c_1(y')=c_2(y')$, both terms take
the value $t_j(c_1(y'))$, so $g_j(y')=0$; in particular $g_j\in\mathfrak m$. The differential of $g_j$ at
$y$ is $\sum_l(1-A)_{jl}\,d(c_1^*t_l)$, and $1-A$ is invertible, so the classes of the $g_j$ span
$\mathfrak m/\mathfrak m^2$. So $\mathfrak m=J+\mathfrak m^2$, where $J$ is the ideal generated by the
$g_j$. The $R$-module $\mathfrak m/J$ is finitely generated and equals $\mathfrak m\cdot(\mathfrak m/J)$,
so by Nakayama's lemma there is $r\in R$ with $r-1\in\mathfrak m$ and $r\mathfrak m\subset J$. Let $y'$ be a
$\mathbf C$-point of $V$ with $c_1(y')=c_2(y')$ and $r(y')\neq0$. Every element of $J$ vanishes at $y'$, so
$r(y')x(y')=0$, hence $x(y')=0$, for every $x\in\mathfrak m$. So $\mathfrak m$ is contained in the maximal
ideal of $y'$, and $y'=y$. The open subset of $V$ where $r$ does not vanish contains $y$, as $r(y)=1$, and
no other such point.

Second, $Y$ is quasi-compact and locally noetherian, so its underlying space is noetherian, and every
subset of it is quasi-compact. The open sets of the first step, one for each point $y$ with
$c_1(y)=c_2(y)$, cover the set of these points. Finitely many of them cover it, and each contains only one
such point.
\end{proof}

\begin{lemma}\label{gal-fixed-points}
Let $p$ be a prime and take a unitary PEL datum of simple type for a CM field $E$ at $p$
(Definition~\ref{gal-pel-datum}) of rank $n$, with $\omega$, $\tau_1$ and group $G$, presented as in
Situation~\ref{gal-count-situation} with $F=E$ and $v$ the infinite place of signature $(1,n-1)$. Let
$k\geq0$, let $N\mid N'$ be integers at least $3$ and prime to $p$, let $K_N\supset K_{N'}$ be the
open compact subgroups of $G(\mathbf A_f^p)$ of levels $N$ and $N'$, and let $Z_N$ and $Z_{N'}$ be the
$k$-fold fibre powers over $\mathbf C$ of the universal abelian schemes over the Shimura varieties of
levels $N$ and $N'$. Let $h\in G(\mathbf A_f^p)$ with $hK_{N'}h^{-1}\subset K_N$ such that every element
of $K_Nh^{-1}K_N$ is regular, and let $\alpha\in\mathcal O_E^k$ with every $\alpha_s$ nonzero and
divisible by $2$. Let $c_1:Z_{N'}\to Z_N$ be the change of level, and $c_2:Z_{N'}\to Z_N$ the Hecke
correspondence of $h$ followed by multiplication by $\alpha_s$ on the $s$-th factor. Then:
\begin{enumerate}
\item $c_2$ is finite \'etale, and the sum $T_v(w_{\alpha,y};\mathbf 1_{K_Nh^{-1}K_N},\mathbf 1_{K_p})$
  of Definition~\ref{gal-count-forms}, with $K=K_N$ and $K'=K_p$, has only finitely many nonzero
  terms;
\item the set $\Phi$ of $x\in Z_{N'}(\mathbf C)$ with $c_1(x)=c_2(x)$ is finite, and at each $x\in\Phi$
  the endomorphism $A_x$ of $T_{c_1(x)}Z_N$ induced by $c_2\circ c_1^{-1}$ near $x$ satisfies
  $\det(1-A_x)\neq0$;
\item for every $y\in\mathbf C$,
  \[
   \sum_{x\in\Phi}\frac{\det(1-yA_x)}{\det(1-A_x)}
   =[K_N\cap h^{-1}K_Nh:K_{N'}]\;T_v(w_{\alpha,y};\mathbf 1_{K_Nh^{-1}K_N},\mathbf 1_{K_p}).
  \]
\end{enumerate}
\end{lemma}
\leanref{Development/AutomorphicArithmetic/Foundations/GaloisRepresentations/LefschetzSigmaSums/HodgeTraceFromLefschetzAndGroupSide.lean}{UpperTheory.Gal.LocalSmoothExtension.FixK}
\leanref{Development/AutomorphicArithmetic/Foundations/GaloisRepresentations/IsogenyEtaleTranslation/TrFEtHolds.lean}{UpperTheory.Gal.FixZ.FixK_holds}
\leanref{Development/AutomorphicArithmetic/Foundations/GaloisRepresentations/IsogenyEtaleTranslation/PBridge.lean}{UpperTheory.Gal.FixZ.bridge_of}
\leanref{Development/AutomorphicArithmetic/Foundations/GaloisRepresentations/SphericalCoordinates/HeckeFixedPointCMCycleClass.lean}{UpperTheory.Gal.W2.FixZ}
\leanref{Development/AutomorphicArithmetic/Foundations/GaloisRepresentations/SphericalCoordinates/FixZOf.lean}{UpperTheory.Gal.W2.fixZ_of}
\leanref{Development/AutomorphicArithmetic/Foundations/GaloisRepresentations/SphericalCoordinates/HeckeFixedPointCMCycleClass.lean}{UpperTheory.Gal.W2.HeckeFixedPointCMCycleClass}
\leanref{Development/AutomorphicArithmetic/Foundations/GaloisRepresentations/SphericalCoordinates/FixedPointsIntegralBaseOfUniformisation.lean}{UpperTheory.Gal.W2.fixedPoints_integral_base_of_uniformisation}

\begin{proof}
The Hecke correspondence of $h$ is finite \'etale, and multiplication by a nonzero $\alpha_s$ is an
isogeny of abelian schemes in characteristic $0$, hence finite \'etale; so $c_2$ is finite \'etale.
The change of level $c_1$ is finite \'etale as well, so the endomorphisms $A_x$ are defined. The
presentation identifies $C$ with the algebra
of $E$-linear endomorphisms of $V=\Lambda\otimes\mathbf Q$, where $\Lambda$ is the lattice of the datum,
with $x^*$ the adjoint for its form; so Situation~\ref{gal-fixed-situation} applies, with $\varrho$ the
direct sum of the embeddings of $E$ with the multiplicities given there. Fix $t\in\mathbf C$.

The maps $c_1$ and $c_2$ lie over the change of level and the Hecke correspondence of $h$ on the Shimura
varieties, also written $c_1,c_2:\mathrm{Sh}_{N'}\to\mathrm{Sh}_N$, so every $x\in\Phi$ lies over a base
fixed point, a point $x'\in\mathrm{Sh}_{N'}(\mathbf C)$ with $c_1(x')=c_2(x')$. Write
$\Upsilon_{N'}(x')=[g,g_p,z]$ for a form $\beta$ (Lemma~\ref{gal-fixed-uniformization}(3)). By
Lemma~\ref{gal-fixed-base}, exactly one $(y,s)\in G_\beta$ fixes $z$ with $hg^{-1}\iota(y)g\in K_N$ and
$g_p^{-1}\iota_p(y)g_p\in K_p$. Conversely, if a point $[g,g_p,z]$ of level $(K_{N'},K_p)$ carries such
a $(y,s)$, then some point of $\Phi$ lies over $\Upsilon_{N'}^{-1}[g,g_p,z]$, which is therefore a base
fixed point, by Lemma~\ref{gal-fixed-fibre}; it applies because the elements of $h^{-1}K_N\subset
K_Nh^{-1}K_N$ are regular. By Lemma~\ref{gal-fixed-fibre}, every $x\in\Phi$ has $\det(1-A_x)\neq0$, and
$A_x$, written in coordinates on $Z_N$ near $c_1(x)$ pulled back by $c_1$, is a tangent matrix at $x$ in the
sense of Lemma~\ref{gal-fixed-points-isolated}. The scheme $Z_{N'}$ is quasi-compact and locally of finite
type over $\mathbf C$, and $Z_N$ is separated over $\mathbf C$. So $\Phi$ is finite by
Lemma~\ref{gal-fixed-points-isolated}. This proves (2), and the base fixed points, the images of the
points of $\Phi$, are finitely many.

By Lemma~\ref{gal-fixed-fibre},
\[
 \sum_{x\in\Phi}\frac{\det(1-tA_x)}{\det(1-A_x)}=\sum_{x'}w_{\alpha,t}(P_y,\mu_y(z)),
\]
the sum over the base fixed points $x'$, with $\Upsilon_{N'}(x')=[g,g_p,z]$ and $(y,s)$ as above. Through
$\Upsilon_{N'}$ the right side is the sum $T_v(w_{\alpha,t};\mathbf 1_{h^{-1}K_N},\mathbf 1_{K_p})$ of
Definition~\ref{gal-count-forms}(5) at level $(K_{N'},K_p)$. Indeed, the stabilizers at this level are
trivial by Lemma~\ref{gal-fixed-stabilizer}, so the term of a point is the sum of
$w_{\alpha,t}(P_y,\mu_y(z))$ over the $(y,s)$ fixing $z$ with $hg^{-1}\iota(y)g\in K_N$ and
$g_p^{-1}\iota_p(y)g_p\in K_p$. This term does not depend on the representative of the point: changing
it conjugates the elements $(y,s)$, which keeps $P_y$, the eigenvalue and the two conditions, since
$\mathbf 1_{h^{-1}K_N}$ is invariant under conjugation by $K_{N'}$. It has one summand at the base fixed
points and none elsewhere.

Finally, Lemma~\ref{gal-fixed-level} applies with $K=K_N$ and $K'=K_{N'}$, which lies in $K_N\cap
h^{-1}K_Nh$ since $N\mid N'$ and $hK_{N'}h^{-1}\subset K_N$: the stabilizers at level $(K_N,K_p)$ are
trivial by Lemma~\ref{gal-fixed-stabilizer}, and the points of level $(K_{N'},K_p)$ that carry an
element as in its hypothesis are the finitely many base fixed points. Its part (1) gives the finiteness
in (1). Its part (2) turns the sum above into $[K_N\cap h^{-1}K_Nh:K_{N'}]\,
T_v(w_{\alpha,t};\mathbf 1_{K_Nh^{-1}K_N},\mathbf 1_{K_p})$ at level $(K_N,K_p)$, which is (3) for $y=t$.
% Lean: W2.fixedPoints_integral_base_of_uniformisation assembles UnifG, StabTriv, FixFwd (with the
% uniqueness huniq), FixLoc, TanExists, FixFinite (finiteness of Phi from the isolated fixed points),
% PowSm, PwOver (hbase), SumPts and EigSame; W2.fixZ_of adds LevelChange; FixZ.bridge_of adds CxFinOf
% for the finiteness in (1) and WtEq (wtFix = cW), and moves the statement from the fibre powers to the
% base T.X.Z of the tower through the isomorphisms T.e.
\end{proof}

\begin{proposition}\label{gal-ks-towers}
Let $p$ be a prime and take a unitary PEL datum of simple type for a CM field $E$ at $p$, with group
$G$ and a dominant weight $\lambda$. Then there are $k\geq0$ and, for every level $N$ and family $g$
of Hecke operators of level $N$, a Hecke tower $H_{N,g}$ over $E$ whose base is the $k$-fold fibre
power $\mathcal A_N^k$ of the universal abelian scheme over the Shimura variety of level $N$, and whose
Hecke correspondences are those of $g$, such that:
\begin{enumerate}
\item $H_{N,g}$ is smooth projective (Definition~\ref{gal-hecke-tower});
\item at a place $w$ of $E$ over a place where the Hecke operators of $g$ are good, $H_{N,g}$ has a
  good model (Definition~\ref{gal-tower-reduction});
\item at a place $w$ where the Hecke operators of $g$ satisfy a congruence relation with coefficients
  $\vartheta$, $H_{N,g}$ has a model satisfying the congruence relation of degree $n$ with coefficients
  $\vartheta$;
\item if $\theta$ and $f$ satisfy the spectral condition of Definition~\ref{gal-plan}(1) with weight
  $\lambda$, then for every automorphism $\sigma$ of $\mathbf C$ and every Hodge exponent $h_i$ of
  $\sigma\circ\lambda$, the system $\sigma\circ\theta$ occurs in Hodge degree $h_i$ of $H_{N,g}$ along
  the embedding of $E$ given by $\sigma\circ\tau_1$.
\end{enumerate}
\end{proposition}
\leanref{Development/AutomorphicArithmetic/Foundations/GaloisRepresentations/IntegralStructures/GeometricFlagBases/PackExOf.lean}{UpperTheory.Gal.Int.Geo.KSAll}
\leanref{Development/AutomorphicArithmetic/Foundations/GaloisRepresentations/IntegralStructures/GeometricFlagBases/IntegOfMem.lean}{UpperTheory.Gal.Int.Geo.StdEx}
\leanref{Development/AutomorphicArithmetic/Foundations/GaloisRepresentations/IntegralStructures/GeometricFlagBases/IntegOfMem.lean}{UpperTheory.Gal.Int.Geo.KSEx}
\leanref{Development/AutomorphicArithmetic/Foundations/GaloisRepresentations/IntegralStructures/GeometricFlagBases/IntegOfMem.lean}{UpperTheory.Gal.Int.Geo.KSNice}
\leanref{Development/AutomorphicArithmetic/Foundations/GaloisRepresentations/IntegralStructures/GeometricFlagBases/IntegOfMem.lean}{UpperTheory.Gal.Int.Geo.KSGood}
\leanref{Development/AutomorphicArithmetic/Foundations/GaloisRepresentations/IntegralStructures/GeometricFlagBases/IntegOfMem.lean}{UpperTheory.Gal.Int.Geo.KSCong}
\leanref{Development/AutomorphicArithmetic/Foundations/GaloisRepresentations/IntegralStructures/GeometricFlagBases/PackExOf.lean}{UpperTheory.Gal.Int.Geo.KSHodge}
\leanref{Development/AutomorphicArithmetic/Foundations/GaloisRepresentations/IntegralStructures/GeometricFlagBases/GeometricRealizationOfWeightOne.lean}{UpperTheory.Gal.Int.Geo.kodaira_spencer_of_weight_one}
\leanref{Development/AutomorphicArithmetic/Foundations/GaloisRepresentations/LefschetzSigmaSums/KsHodgeOfGeo.lean}{UpperTheory.Gal.LocalSmoothExtension.ksHodge_of_geo}
\leanref{Development/AutomorphicArithmetic/Foundations/GaloisRepresentations/LefschetzSigmaSums/HodgeTraceOfLefschetzAndGroupSide.lean}{UpperTheory.Gal.LocalSmoothExtension.ksHodge_of_inputs}
\leanref{Development/AutomorphicArithmetic/Foundations/GaloisRepresentations/LefschetzSigmaSums/HodgeTraceFromLefschetzAndGroupSide.lean}{UpperTheory.Gal.LocalSmoothExtension.FixK}

\begin{proof}
Items (1) and (2) concern the moduli problem of Definition~\ref{gal-pel-datum} and its universal
abelian scheme $\mathcal A_N$. The Shimura variety of level $N$ is smooth and projective over $E$, so
$\mathcal A_N^k$ is as well, and the Hecke correspondences of level $N$, given by prime-to-$p$
isogenies compatible with the level structures, are finite \'etale on the Shimura variety and extend
to $\mathcal A_N^k$ through the isogenies. At a place $w$ over a place where the Hecke operators of $g$
are good, the moduli problem has a smooth projective integral model over which the universal abelian
scheme and these isogenies extend, the latter as finite \'etale correspondences (Kottwitz
\cite[\S5]{Kot92}); the $k$-fold fibre power of the extended abelian scheme is the good model. For
(3), let $w$ be a place at which the Hecke operators of $g$ satisfy a congruence relation with
coefficients $\vartheta$, and let $p_w$ be the prime under $w$. This relation includes that the PEL
datum is also of simple type at $p_w$; that a level $N_0\mathfrak a$ containing the level $N$ acts on
$H_{N,g}$ by deck transformations, with $N_0$ prime to $p_w$, $\mathfrak a+w=\mathcal O_E$ and
$\mathfrak a$ containing a power of the ideal $\mathfrak u$ generated by $p_w$ and $\omega-r$, where
$r\in\mathbf Z$ and $\omega-r\in w$; and that for each $0\leq i\leq n$ one of the operators,
$g_{j_i}$, is a Hecke operator of type $i$ at $w$, whose isogenies are of type $i$ at $w$ in the sense
of Definition~\ref{gal-powers-hecke}, with $\theta(g_{j_i})=m_i\vartheta_i$ for a positive integer
$m_i$, an index of level subgroups. Take Situation~\ref{gal-powers-situation} for the prime $p_w$, the
valuation ring $\mathcal O$ of $E$ at $w$, the level $N_0$ and the ideal $\mathfrak a$, with $M$ and
$M'$ the integral models of Kottwitz and of Harris and Taylor, for which $\pi$ is finite \'etale
\cite[\S5]{Kot92}, \cite[\S III.4]{HT01}. The base of $H_{N,g}$ maps to the $k$-fold fibre
power $\mathcal Y$ of $\mathcal A'$ over $M'$ by a finite \'etale covering split by the deck
transformations, and this is a model of $H_{N,g}$ at $w$. By Lemma~\ref{gal-powers-hecke-exist} and
Proposition~\ref{gal-powers-eichler-shimura}, there are finite correspondences $(W_i,p_i,r_i)$ of
$\mathcal Y$, $0\leq i\leq n$, that satisfy the Eichler--Shimura relation of degree $n$ at the generic
points of the special fibre, and the Hecke correspondence of $g_{j_i}$ lies over $(W_i,p_i,r_i)$ with
multiplicity $m_i$. So this model satisfies the congruence relation of degree $n$ with coefficients
$\vartheta$.
% Lean: items (1) and (2) are StdEx (CovMain.stdEx_holds), KSEx (G3.kSEx_holds), KSNice
% (G3.kSNice_holds) and KSGood (FiniteEtaleModels.ksGood_holds); the tower is T.tower N g for
% T : KS L gb k, with base (pow (M N).univ.A k), and GoodGrp, CongGrp say when the operators are good
% or satisfy the congruence relation at a place (CongGrp includes the operators T_i of IsTjElt with
% index the number of i-dimensional subspaces of the residue space). Item (3) is KSCong
% (CovMain.ksCong_holds = tower_covering_congruence_of_hecke_lies): SiteEx gives the data of the
% Situation (Site), KSCover the covering, MIn.cong = Mod.congKS_link the correspondences with the
% Eichler-Shimura relation, and TjLies the lying over with multiplicity.

For (4), present $G$ as in Situation~\ref{gal-count-situation} with $F=E$, and fix $\sigma$ and a
tuple $h=(h_i)$ of Hodge exponents of the weight $\sigma\circ\lambda$ (Definition~\ref{gal-count-weights}).
The spectral condition, conjugated by $\sigma$, gives a representation $\xi$ of
$C^\times\times\mathbf Q^\times$ with $\operatorname{tr}\xi(y)=S_{\sigma\circ\lambda}(P_y)$ for the
elements $y$ of the groups of definite forms with $P_y$ separable of degree $n$, and the functions
$f'=(ts)(g)\ast f$ there are elliptic test functions for $K_N$ (Definition~\ref{gal-count-pairs}(6)),
invariant under $K_N$ on both sides. Take $k$ at least the integers $k_0$ of
Lemma~\ref{gal-count-root-weights} for the finitely many weights $\sigma\circ\lambda$ and their tuples
of Hodge exponents, let $(c_\alpha)$ be as there, and let $D=\sum_\alpha c_\alpha[\alpha]^*$ on
Hodge cohomology, where $[\alpha]$ is multiplication by $\alpha_s$ on the $s$-th factor of
$Z=\mathcal A_N^k$. Write $f'=\sum_ic_i\mathbf 1_{K_Ng_i^{-1}K_N}$ with finitely many $g_i$; every
element of $K_Ng_i^{-1}K_N$ is regular. For each $i$ choose $N'$ with $N\mid N'$ and
$g_iK_{N'}g_i^{-1}\subset K_N$, and for each $\alpha$ with $c_\alpha\neq0$ let $(c_1,c_2)$ be the
correspondence of Lemma~\ref{gal-fixed-points} with $g_i$ in place of its element $h$. On
$H^b(Z,\Omega^a)$, $c_{1*}c_2^*$ is
$[K_N\cap g_i^{-1}K_Ng_i:K_{N'}]$ times the Hecke operator $T_{g_i}$ of $\mathbf 1_{K_Ng_i^{-1}K_N}$
composed with $[\alpha]^*$. By Lemma~\ref{gal-fixed-points} its fixed points are finite and
nondegenerate, and Theorem~\ref{gal-holomorphic-lefschetz}, multiplied by $(-y)^a$ and summed over $a$,
gives
\[
 \sum_a(-y)^a\sum_b(-1)^b\operatorname{tr}\bigl(T_{g_i}[\alpha]^*\mid H^b(Z,\Omega^a)\bigr)
 =T_v(w_{\alpha,y};\mathbf 1_{K_Ng_i^{-1}K_N},\mathbf 1_{K_p}),
\]
since $\sum_a(-y)^a\operatorname{tr}\Lambda^aA_x=\det(1-yA_x)$ and the index cancels. Both sides are
linear in the function, so multiplying by $c_ic_\alpha$, summing, and applying
Proposition~\ref{gal-count-group-side} gives
\[
 \sum_a(-y)^a\sum_b(-1)^b\operatorname{tr}\bigl(T_{f'}D\mid H^b(Z,\Omega^a)\bigr)
 =\operatorname{tr}_{\mathrm{def}}(\xi;f',\mathbf 1_{K_p})\sum_iy^{h_i}.
\]
This is the trace on automorphic forms of weight $\xi$ on the definite groups referred to in the
spectral condition.
% Lean: HodgeTr (hodgeTr_of: FnDecomp, HopWord, AltLin, CxLin and GeoStep), with GeoStep proved by
% geoStep_of from the Lefschetz formula, KSProj, FixK and LefSum, and GroupTop by GroupTop_of from
% weighted_count_definite_unitary and trace_definite_forms_identity; assembled in
% hodgeTrace_of_lefschetz_and_groupSide = ksHodge_of_inputs. The conjugation by sigma is conjWt and
% SpecConj in kSHodge_of.
Suppose that $\sigma\circ\theta$ did not occur in $\bigoplus_bH^b(Z,\Omega^{h_i})$. This space is
finite-dimensional and the operators $g_j$ commute on it, so there is a polynomial $s$ in the $g_j$ with
$s(\iota\theta)\neq0$ that vanishes on it. Then the alternating traces of all the operators
$(ts)(g)\ast f$ on Hodge cohomology vanish, while by the spectral condition one of the corresponding
traces on the definite group does not, a contradiction.
\end{proof}

\begin{lemma}\label{gal-plan-pack}
In Situation~\ref{gal-seed-situation}, let $P$ be a plan for $U$ (Definition~\ref{gal-plan}). Then
there is a pack $Q$ for $U$ (Definition~\ref{gal-pack}) that has Hodge data at every place where $P$
has Hodge data and is congruent at every place where $P$ is congruent.
\end{lemma}
\leanref{Development/AutomorphicArithmetic/Foundations/GaloisRepresentations/IntegralStructures/GeometricFlagBases/PackExOf.lean}{UpperTheory.Gal.Int.Geo.Plan.pack}

\begin{proof}
Let $H=H_{N,g}$ be the tower of Proposition~\ref{gal-ks-towers} for the PEL datum, level and Hecke
operators of $P$, and take $Q=(H,\theta,\kappa,X)$ with $X$ the finite set of places excluded in
Definition~\ref{gal-plan}(2). By Proposition~\ref{gal-ks-towers}(1), $H$ is smooth projective. At a
place $w$ where $P$ is congruent, $\kappa$ is unramified and the Hecke operators satisfy a congruence
relation with coefficients $\vartheta$ whose Hecke polynomial has the roots
$\iota^{-1}(q_w^{(a-1)/2}b\,\kappa(\varpi_w)^{-1})$; Proposition~\ref{gal-ks-towers}(3) gives a model
of $H$ at $w$ satisfying the congruence relation of degree $n=a$ with these coefficients, so $Q$ is
congruent at $w$, and in particular $Q$ is a pack. At a place $v$ where $P$ has Hodge data, $\kappa$
is unramified at $v$ and Proposition~\ref{gal-ks-towers}(2) gives a good model. For $\tau$ inducing $v$,
take $\sigma$ as in Definition~\ref{gal-plan}; each $k-\lambda'_k+m_\tau$ is a Hodge exponent of
$\sigma\circ\lambda$, and Proposition~\ref{gal-ks-towers}(4), applied to the spectral condition of $P$,
shows that $\iota\circ\theta$ composed with $\sigma$ occurs in that Hodge degree along $\sigma\circ\tau$.
So $Q$ has Hodge data at $v$.
\end{proof}

\subsection{Test functions and Iwahori operators}\label{subsection-gal-test-functions}

\begin{lemma}\label{gal-identification-isometry}
Let $E$ be a CM field with maximal totally real subfield $F$ and complex conjugation $c$, let
$t,m\in\mathbf Z$ with $t^2-4m\neq0$, let $\omega\in\mathcal O_E$ with $\omega+c(\omega)=t$ and
$\omega c(\omega)=m$, let $\gamma\in\mathcal O_F$ be nonzero, and let $n\geq1$ be odd. Call a prime $q$
\emph{split for $\omega$} if $q\nmid t^2-4m$ and $X^2-tX+m$ has a root modulo $q$, and let $T$ be the set
of places of $E$ over primes split for $\omega$ at which $\gamma$ is a unit. Then:
\begin{enumerate}
\item there are $x,y\in\mathbf A_{E,f}$ with $\gamma\bigl(x\,c(x)+y\,c(y)\bigr)=1$ whose components at
  the places of $T$ are integral;
\item if $R$ is a commutative ring with an involution $c$, and $x,y\in R$ and $\gamma\in R$ with
  $c(\gamma)=\gamma$ satisfy $\gamma\bigl(x\,c(x)+y\,c(y)\bigr)=1$, then the block diagonal matrix
  \[
   a=\operatorname{diag}(1,M,\ldots,M)\in\operatorname{GL}_n(R),\qquad
   M=\begin{pmatrix}x&-c(y)\\y&c(x)\end{pmatrix},
  \]
  satisfies ${}^tc(a)\,D\,a=1$ for $D=\operatorname{diag}(1,\gamma,\ldots,\gamma)$, and the entries of
  $a$ and of $a^{-1}$ lie in every subring of $R$ that contains $x$, $y$, $c(x)$, $c(y)$ and $\gamma$.
\end{enumerate}
In particular there is $a\in\operatorname{GL}_n(\mathbf A_{E,f})$ with ${}^tc(a)\,D\,a=1$ and
$a_w\in\operatorname{GL}_n(\mathcal O_w)$ for every $w\in T$.
\end{lemma}
\leanref{Development/AutomorphicArithmetic/Foundations/GaloisRepresentations/PAdicHensel/PadicUnitNormHolds.lean}{UpperTheory.Gal.G1.IsoEx_holds}
\leanref{Development/AutomorphicArithmetic/Foundations/GaloisRepresentations/PAdicHensel/PadicUnitNormHolds.lean}{UpperTheory.Gal.G1.NormTwo_holds}
\leanref{Development/AutomorphicArithmetic/Foundations/GaloisRepresentations/PAdicHensel/PadicUnitNormHolds.lean}{UpperTheory.Gal.G1.AdFour_holds}
\leanref{Development/AutomorphicArithmetic/Foundations/GaloisRepresentations/PAdicHensel/PadicUnitNormHolds.lean}{UpperTheory.Gal.G1.BlockMat_holds}

\begin{proof}
For $s,s'$ in a commutative $\mathbf Q$-algebra put $N(s,s')=s^2+tss'+ms'^2$. If $c$ fixes $s$ and $s'$,
then $(s+s'\omega)\,c(s+s'\omega)=N(s,s')$, because $\omega+c(\omega)=t$ and $\omega c(\omega)=m$. So for
(1) it suffices to find $s_0,\ldots,s_3\in\mathbf A_{F,f}$ with
$\gamma\bigl(N(s_0,s_1)+N(s_2,s_3)\bigr)=1$ that are integral at the places of $F$ under $T$, and to put
$x=s_0+s_1\omega$ and $y=s_2+s_3\omega$. Choose the components $s_{i,v}$ place by place.
\begin{itemize}
\item Let $v$ lie under a place of $T$, over the prime $q$. The square of the difference of the two roots
  of $X^2-tX+m$ is $t^2-4m$, a unit at $v$, so by Hensel's lemma there are roots
  $\rho,\rho'\in\mathcal O_v$ with $\rho-\rho'$ a unit. Then $N(s,s')=(s+\rho s')(s+\rho's')$, and
  $s'=(\gamma^{-1}-1)/(\rho-\rho')$, $s=1-\rho's'$ are integral with $N(s,s')=\gamma^{-1}$. Take
  $(s_{0,v},s_{1,v})=(s,s')$ and $s_{2,v}=s_{3,v}=0$.
\item Let $v$ be another place at which $2(t^2-4m)$ and $\gamma$ are units. Then $N$ is a binary
  quadratic form over $\mathcal O_v$ with unit discriminant and odd residue characteristic, so it represents
  the unit $\gamma^{-1}$ with $s,s'\in\mathcal O_v$ \cite[\S92]{OMe63}. Take $s_{2,v}=s_{3,v}=0$.
\item At each of the finitely many remaining places, $N\perp N$ is a nondegenerate quadratic form in four
  variables over $F_v$, so $N\perp N\perp\langle-\gamma^{-1}\rangle$, of dimension $5$, is isotropic
  \cite[63:19]{OMe63}. An
  isotropic vector with nonzero last coordinate gives a representation of $\gamma^{-1}$ by $N\perp N$; if
  the last coordinate is zero, $N\perp N$ is isotropic, so it contains a hyperbolic plane and represents
  every element of $F_v$.
\end{itemize}
Almost all components are integral, so the $s_i$ are finite adeles.

For (2), multiplying out gives ${}^tc(M)\,\gamma M=\gamma\bigl(x\,c(x)+y\,c(y)\bigr)\cdot1=1$. Since $n$
is odd, $D$ is block diagonal with blocks $1,\gamma\cdot1,\ldots,\gamma\cdot1$ matching those of $a$, so
${}^tc(a)Da=1$. Hence $a^{-1}={}^tc(a)D$, and the entries of $a$ and $a^{-1}$ are $0$, $1$ and the
elements $\pm x,\pm y,\pm c(x),\pm c(y)$, possibly multiplied by $\gamma$. The last statement follows by
applying (2) to the $x,y$ of (1) in $R=\mathbf A_{E,f}$: the entries of $a$ and $a^{-1}$ are integral
at the places of $T$ and at almost all places.
\end{proof}
% Lean: pIsoEx from NormTwo (pNormTwo AdFour_holds, place by place through SplitRep with LocRoot,
% LocUnit and LocFour) and BlockMat (pBlockMat). The Lean's integrality is stated for places over
% primes q with A.Good q, which implies q split for omega and gamma a unit at q.

\begin{lemma}\label{gal-identification}
Let $E$, $F$, $c$, $t$, $m$, $\omega$ and $\gamma$ satisfy the hypotheses of
Lemma~\ref{gal-identification-isometry}, with $t^2-4m<0$, let $\delta\in E$ be nonzero with $c(\delta)=-\delta$ and
$\operatorname{Tr}_{E/\mathbf Q}(\delta\mathcal O_E)\subset\mathbf Z$, let $n\geq1$ be odd, and let $p$ be
a prime. Call a prime $q$ \emph{good} if it is split for $\omega$ and there is an integer $e$ prime to $q$,
divisible by $\gamma$ in $\mathcal O_E$, such that $e$ times every $\mathbf Z$-linear form on $\mathcal
O_E$ is $y\mapsto\operatorname{Tr}_{E/\mathbf Q}(\delta zy)$ for some $z\in\mathcal O_E$. Suppose that $p$
is good and that $\Lambda=\mathcal O_E^n$ with the form
\[
 \langle x,y\rangle=\operatorname{Tr}_{E/\mathbf Q}\Bigl(\delta\sum_{i=1}^nh_ix_i\,c(y_i)\Bigr),\qquad
 (h_1,h_2,\ldots,h_n)=(-\gamma,-\gamma^2,\ldots,-\gamma^2),
\]
underlies a unitary PEL datum of simple type for $E$ at $p$ (Definition~\ref{gal-pel-datum}) with group
$G$. Let $\mathrm{GU}_n(\mathbf A_f)$ be the group of $(g,\nu)\in\operatorname{GL}_n(\mathbf
A_{E,f})\times\mathbf A_f^\times$ with ${}^tc(g)\,g=\nu$, let $g_w$ be the component of $g$ at a place $w$
of $E$, and let $G(\mathbf A_f^p)$ be the group of $E\otimes\mathbf A_f^p$-linear similitudes of
$\Lambda\otimes\mathbf A_f^p$, whose component at a place $w\nmid p$ acts on $E_w^n$. Then there are a
homomorphism $j:\mathrm{GU}_n(\mathbf A_f)\to G(\mathbf A_f^p)$ and matrices
$a_w\in\operatorname{GL}_n(E_w)$, one for each finite place $w$, such that:
\begin{enumerate}
\item for $w\nmid p$, $j(g,\nu)$ acts on $E_w^n$ by $a_wg_wa_w^{-1}$, and its similitude factor is the
  component of $\nu$ away from $p$;
\item $j(g,\nu)=j(g',\nu')$ if and only if $g_w=g'_w$ for every $w\nmid p$, and $j$ has a section $s$
  that is a homomorphism whose values have component $1$ at every place over $p$; in particular $j$ is
  surjective;
\item $a_w\in\operatorname{GL}_n(\mathcal O_w)$ for almost all $w$ and for every $w$ over a good prime;
\item $\theta(g,\nu)=({}^tg^{-1},\nu^{-1})$ is an automorphism of $\mathrm{GU}_n(\mathbf A_f)$;
\item the group $\Gamma$ of $(g,\nu)\in\operatorname{GL}_n(E)\times\mathbf Q^\times$ with $\nu>0$ and
  ${}^tc(g)\,g=\nu$ is the group of rational points, with positive similitude factor, of the definite
  unitary similitude group with the same finite adelic points as $G$; the resulting embedding of $\Gamma$
  into $G(\mathbf A_f^p)\times G(\mathbf Q_p)$ has first component $j$, and its second component at
  $(g,\nu)$ lies in the stabilizer $K_p$ of $\Lambda\otimes\mathbf Z_p$ if and only if
  $g\in\operatorname{GL}_n(\mathcal O_w)$ for every $w\mid p$.
\end{enumerate}
\end{lemma}
\leanref{Development/AutomorphicArithmetic/Foundations/GaloisRepresentations/IntegralStructures/GeometricFlagBases/PrimeLevelLocalization/PFrameEx.lean}{UpperTheory.Gal.Int.Geo.Pl.frameEx_holds}
\leanref{Development/AutomorphicArithmetic/Foundations/GaloisRepresentations/IntegralStructures/GeometricFlagBases/PrimeLevelLocalization/PFrameEx.lean}{UpperTheory.Gal.W2.Seed.trEx_holds}
\leanref{Development/AutomorphicArithmetic/Foundations/GaloisRepresentations/SphericalCoordinates/AlgebraEquivSetup/UptsExt.lean}{UpperTheory.Gal.W2.Seed.pTrEx}
\leanref{Development/AutomorphicArithmetic/Foundations/GaloisRepresentations/SphericalCoordinates/AlgebraEquivSetup/ApartOfNotMem.lean}{UpperTheory.Gal.W2.Seed.Tr}
\leanref{Development/AutomorphicArithmetic/Foundations/GaloisRepresentations/IntegralStructures/GeometricFlagBases/PrimeLevelLocalization/PFrameEx.lean}{UpperTheory.Gal.W2.Seed.vecCoord_holds}

\begin{proof}
Let $a\in\operatorname{GL}_n(\mathbf A_{E,f})$ be given by Lemma~\ref{gal-identification-isometry}, with
${}^tc(a)Da=1$ for $D=\operatorname{diag}(1,\gamma,\ldots,\gamma)$, and let $a_w$ be its components. A
good prime $q$ is split for $\omega$, and $\gamma$ is a unit at the places over $q$ because it divides an
integer prime to $q$. So Lemma~\ref{gal-identification-isometry} gives (3).

The Hermitian form of the datum is $h(x,y)=\sum_ih_ix_ic(y_i)$, with matrix $-\gamma D$, and
$\langle x,y\rangle=\operatorname{Tr}_{E/\mathbf Q}(\delta h(x,y))$. Since $\delta\neq0$ and the trace
form of $E$ is nondegenerate, an $E$-linear map $b$ is a similitude of $\langle\,,\rangle$ with factor
$\nu$ if and only if it is one of $h$, that is, ${}^tc(b)Db=\nu D$. For $(g,\nu)\in\mathrm{GU}_n(\mathbf
A_f)$ let $b=aga^{-1}$. Since $D={}^tc(a)^{-1}a^{-1}$,
\[
 {}^tc(b)\,D\,b={}^tc(a)^{-1}\,{}^tc(g)\,g\,a^{-1}=\nu\,{}^tc(a)^{-1}a^{-1}=\nu D.
\]
So the components of $b$ at the places prime to $p$ form an element $j(g,\nu)$ of $G(\mathbf A_f^p)$ with
factor $\nu$ away from $p$, and $j$ is a homomorphism satisfying (1).

As the $a_w$ are invertible, $j(g,\nu)=j(g',\nu')$ if and only if $g_w=g'_w$ for all $w\nmid p$; the
factor $\nu$ away from $p$ is determined by $g$. Given $h\in G(\mathbf A_f^p)$ with factor $\mu$, let
$s(h)=(g,\nu)$ with $g_w=a_w^{-1}h_wa_w$ for $w\nmid p$, $g_w=1$ for $w\mid p$, and $\nu$ equal to $\mu$
away from $p$ and to $1$ at $p$. By (3), $g_w\in\operatorname{GL}_n(\mathcal O_w)$ for almost all $w$, so
$s(h)\in\mathrm{GU}_n(\mathbf A_f)$, and $j(s(h))=h$. Two elements with component $1$ at every place over
$p$ and the same image under $j$ are equal, so $s$ is a homomorphism. This is (2).

For (4), ${}^tc(g)g=\nu$ gives $g\,{}^tc(g)=\nu$, so applying $c$ gives $c(g)\,{}^tg=\nu$, hence
${}^tg\,c(g)=\nu$, and ${}^tc({}^tg^{-1})\,{}^tg^{-1}=({}^tg\,c(g))^{-1}=\nu^{-1}$.

For (5), the standard Hermitian form is positive definite at every infinite place, and $a$ is an isometry
from it to the form with matrix $D$, which is $-\gamma^{-1}$ times that of the datum, at every finite
place. Its rational similitudes with positive factor form $\Gamma$, and the embedding defined by $a$ is
$(g,\nu)\mapsto\bigl(j(g,\nu),(a_wga_w^{-1})_{w\mid p}\bigr)$, with the factor $\nu$ at $p$. Since $p$ is
good, $a_w\in\operatorname{GL}_n(\mathcal O_w)$ for $w\mid p$ by (3), and the component at $p$ preserves
$\Lambda\otimes\mathbf Z_p=\prod_{w\mid p}\mathcal O_w^n$ if and only if $g\in\operatorname{GL}_n(\mathcal
O_w)$ for every $w\mid p$; the factor $\nu$ is then a unit at $p$.
\end{proof}
% Lean: frameEx_holds builds Pl.Frame from trEx_holds (pTrEx: AdTEx for the projection Phi away from p,
% MuEx, isoEx'_holds = IsoEx_holds, pJEx, pJEq, pJSurj), vecCoord_holds (the lattice coordinates vec,
% vec_hat, vec_act) and pStdForm2 (the DForm B with grp G B.beta = Gam0 E n, emb1 and emb2). th is
% CertificateRingHomMap.th.

\begin{definition}\label{gal-trace-vector}
In the setting of Lemma~\ref{gal-identification}, let $x\in E$, let $U$ be an automorphic representation
of $U_n$ (Definition~\ref{gal-definite-unitary}) and let $w_0$ be a place of $E$ split over $F$. Write
$u\sim w$ if the places $u$ and $w$ of $E$ lie over the same place of $\mathbf Q(x)$, and fix for each
infinite place $v$ of $E$ an embedding $\tau_v$ inducing it. A \emph{vector} for $U$ at $w_0$ consists of
the following data.
\begin{enumerate}
\item For each infinite place $v$, integers $W_{v,0}\geq\cdots\geq W_{v,n-1}$ and $m_v\geq0$ with
  $W_{v,i}+m_v\geq0$, such that $U$ has a nonzero vector of weight $(W_v)_v$ at infinity. Put
  $j_z=\sum_v\bigl(\sum_iW_{v,i}+nm_v\bigr)$, $\mathcal W=\bigotimes_v\nabla(W_v)$, and let $\xi_W$ be the
  representation of $\operatorname{GL}_n(E)\times\mathbf Q^\times$ on the dual $\mathcal W^\vee$ given by
  $(g,\nu)\mapsto\nu^{j_z}\bigotimes_v\nabla(W_v)(\tau_v(g))^\vee$.
\item A compact open subgroup $J\subset\operatorname{GL}_n(E_{w_0})$ without torsion and a character $\chi$
  of $J$ with open kernel, such that every irreducible smooth representation of
  $\operatorname{GL}_n(E_{w_0})$ with a nonzero vector on which $J$ acts by $\chi$ is cuspidal.
\item A finite set $S$ of places of $E$ containing $w_0$, and compact open subgroups
  $K_w\subset\operatorname{GL}_n(E_w)$ for $w\in S$, with $K_{w_0}=J$ and $K_w\subset\operatorname{GL}_n
  (\mathcal O_w)$ for $w\neq w_0$, such that $c(w)\notin S$ for every $w\in S$ split over $F$, and
  $w,c(w)\notin S$ for every place $w$ split over $F$, over a place of $F$ other than the one under $w_0$,
  at which $U$ has Satake parameters. The \emph{level} $K_S$ is the group of $(g,\nu)\in\mathrm{GU}_n
  (\mathbf A_f)$ with $\nu\in\hat{\mathbf Z}^\times$, $g_w\in K_w$ for $w\in S$, and
  $g_w\in\operatorname{GL}_n(\mathcal O_w)$ for $w,c(w)\notin S$.
\item A nonzero function $\Phi:\mathrm{GU}_n(\mathbf A_f)\to\mathcal W^\vee$ with
  $\Phi(\gamma g)=\xi_W(\gamma)\Phi(g)$ for $\gamma\in\Gamma$, and $\Phi(gk)=\chi(k_{w_0})\Phi(g)$ for
  $k\in K_S$ with $k_{w_0}\in J$, such that at every place $w$ split over $F$ with $w,c(w)\notin S$ at which
  $U$ has Satake parameters $B$, $\Phi$ is a spherical vector with Satake parameters $B$ for right
  translation by $\operatorname{GL}_n(E_w)$, embedded in $\mathrm{GU}_n(\mathbf A_f)$ with similitude
  factor $1$ through the components at $w$ and $c(w)$.
\item An algebraic Hecke character $\kappa$ of $E$, with infinity type determined by $x$, $W$ and $j_z$,
  that is unramified at every place $w$ with $c(u)\not\sim w$ for all $u\in S$, and such that for every such
  $w$ with moreover $c(w)\not\sim w$, $\Phi(gz)=\kappa(\varpi_w)^{-1}\Phi(g)$ for all $g$, whenever
  $z\in\mathrm{GU}_n(\mathbf A_f)$ has component the scalar $N(w)$ at the places $u$ with $c(u)\sim w$
  and $1$ at the other places.
\end{enumerate}
The \emph{weight} of the vector is $\lambda_V$, with $\lambda_V(\tau_v)=(m_v,\ldots,m_v)$ and
$\lambda_V(c\circ\tau_v)=(W_{v,i}+m_v)_{0\leq i<n}$. An element $k$ of $G(\mathbf A_f^p)$
\emph{stabilizes} the vector if $k=j(k')$ with $\theta(k')\in K_S$, $\theta(k')_{w_0}\in J$ and
$\chi(\theta(k')_{w_0})=1$. A place $w$ split over $F$ is \emph{free} for the vector if $w,c(w)\notin S$,
$c(w)\not\sim w$, and $c(u)\not\sim w$ for every $u\in S$.
\end{definition}
\leanref{Development/AutomorphicArithmetic/Foundations/GaloisRepresentations/IntegralStructures/GeometricFlagBases/PrimeLevelLocalization/PrimeProj.lean}{UpperTheory.Gal.Int.Geo.Pl.Vec}
\leanref{Development/AutomorphicArithmetic/Foundations/GaloisRepresentations/ProjectiveAutomorphicHecke/CentralEig2.lean}{UpperTheory.Gal.CertificateRingHomMap.IsForm}
\leanref{Development/AutomorphicArithmetic/Foundations/GaloisRepresentations/ProjectiveAutomorphicHecke/CentralEig2.lean}{UpperTheory.Gal.CertificateRingHomMap.IsKap}
\leanref{Development/AutomorphicArithmetic/Foundations/GaloisRepresentations/IntegralStructures/GeometricFlagBases/PrimeLevelLocalization/PrimeProj.lean}{UpperTheory.Gal.Int.Geo.Pl.Stab}
\leanref{Development/AutomorphicArithmetic/Foundations/GaloisRepresentations/IntegralStructures/GeometricFlagBases/PrimeLevelLocalization/PrimeProj.lean}{UpperTheory.Gal.Int.Geo.Pl.Free}

\begin{definition}\label{gal-trace-seen-operators}
In the setting of Lemma~\ref{gal-identification}, let $\hat\Lambda^p=\prod_{w\nmid p}\mathcal O_w^n\subset
\Lambda\otimes\mathbf A_f^p$. For $N\geq1$ prime to $p$, let $K_N\subset G(\mathbf A_f^p)$ be the group of
elements that preserve $\hat\Lambda^p$ and act trivially on $\hat\Lambda^p/N\hat\Lambda^p$. An element $h$
of $G(\mathbf A_f^p)$ is \emph{integral} if $h\hat\Lambda^p\subset\hat\Lambda^p$. For a place $w$ of $E$
split over $F$ and prime to $p$, let $\operatorname{pr}_w:G(\mathbf A_f^p)\to\operatorname{GL}_n(E_w)$ be the
homomorphism with $\operatorname{pr}_w(j(g,\nu))=g_w$, and let $\operatorname{loc}_w(y)=j(\iota_w(y))$ for
$y\in\operatorname{GL}_n(E_w)$, where $\iota_w(y)\in\mathrm{GU}_n(\mathbf A_f)$ has similitude factor $1$,
component $y$ at $w$, and component $1$ away from $w$ and $c(w)$.
\begin{enumerate}
\item A \emph{Hecke operator of level $N$} is a triple $g=(K,h,M)$ of a group $K$ with
  $K_N\subset K\subset K_1$, an integral element $h$, and a multiple $M\geq3$ of $N$ prime to $p$, such that
  $hkK_M(hk)^{-1}\subset K_N$ for every $k\in K$. Its \emph{function} is
  $f_g(y)=\#\{k\in K/K_M:hky\in K_N\}$, and its \emph{index} is
  $\operatorname{idx}(g)=[K\cap h^{-1}Kh:K_M\cap K\cap h^{-1}Kh]$.
\item Let $V$ be a vector for an automorphic representation $U$ of $U_n$ at $w_0$
  (Definition~\ref{gal-trace-vector}), and let $r_0$ be a prime. A Hecke operator $g=(K,h,M)$ of level $N$
  is \emph{seen on $U$ with value $\vartheta\in\mathbf C$ and auxiliary prime $r_0$} if either
  \begin{enumerate}
  \item $K=K_N$, $M=N$, $h\in K_1\cap K_{r_0}$, $h$ stabilizes $V$, and $\vartheta=1$; or
  \item there are a place $w$ free for $V$ over a good prime $q\neq r_0$, an integer $0\leq i\leq n$ and
    a multiset $B$ such that: $K$ is a congruence subgroup of level $N$ at the primes other than $q$ and of a
    level at $q$ adapted to $w$, and all of its elements stabilize $V$; $h=\operatorname{loc}_w(t_i)\,j(z)^i$,
    where $t_i=\operatorname{diag}(\varpi_w,\ldots,\varpi_w,1,\ldots,1)$ has $i$ entries equal to a
    uniformizer $\varpi_w$, and $z\in\mathrm{GU}_n(\mathbf A_f)$ has component the scalar $N(w)$ at the
    places $u$ with $c(u)\sim w$ and $1$ at the other places; $M=N\cdot N(w)^i$; $U$ has Satake
    parameters $B$ at $w$; and $\vartheta=\operatorname{idx}(g)\,q_w^{i(n-i)/2}e_i(B)\,\kappa(\varpi_w)^{-i}$.
  \end{enumerate}
\end{enumerate}
\end{definition}
\leanref{Development/AutomorphicArithmetic/Foundations/GaloisRepresentations/IntegralStructures/GeometricFlagBases/IntegOfMem.lean}{UpperTheory.Gal.Int.Geo.Mem}
\leanref{Development/AutomorphicArithmetic/Foundations/GaloisRepresentations/IntegralStructures/GeometricFlagBases/IntegOfMem.lean}{UpperTheory.Gal.Int.Geo.Mem.fn}
\leanref{Development/AutomorphicArithmetic/Foundations/GaloisRepresentations/IntegralStructures/GeometricFlagBases/PrimeLevelLocalization/PrimeProj.lean}{UpperTheory.Gal.Int.Geo.Pl.idx}
\leanref{Development/AutomorphicArithmetic/Foundations/GaloisRepresentations/IntegralStructures/GeometricFlagBases/PrimeLevelLocalization/PrimeProj.lean}{UpperTheory.Gal.Int.Geo.Pl.Seen}

\begin{lemma}\label{gal-trace-auxiliary-prime}
In the setting of Lemma~\ref{gal-identification}, let $x\in E$ with $\mathbf Q(x)=\mathbf Q(\omega)$, and
let $S$ be a finite set of places of $E$. For every finite set $X$ of primes there are a prime $r\notin X$
with $r\neq p$ and a nonempty finite set $W$ of places of $E$ such that:
\begin{enumerate}
\item $r$ is good and lies under no place of $S$;
\item every place in $W$ lies over $r$ and is split over $F$;
\item every place of $E$ over $r$ is $w$ or $c(w)$ for some $w\in W$;
\item distinct places in $W$ lie over distinct places of $F$.
\end{enumerate}
\end{lemma}
\leanref{Development/AutomorphicArithmetic/Foundations/GaloisRepresentations/ScalarCommutation/AuxiliaryInertOddPrime.lean}{UpperTheory.Gal.Spc.auxiliary_prime_holds}
\leanref{Development/AutomorphicArithmetic/Foundations/GaloisRepresentations/ScalarCommutation/ExistsIdempotentHeckeOperatorWithRegularSupportAlt.lean}{UpperTheory.Gal.Spc.SlopeTestRemainingChoices}
\leanref{Development/AutomorphicArithmetic/Foundations/GaloisRepresentations/IntegralStructures/GeometricFlagBases/PrimeLevelLocalization/PDatEx.lean}{UpperTheory.Gal.Int.Geo.Pl.sides_holds}

\begin{proof}
The polynomial $P(X)=X^2-tX+m$ has a root modulo infinitely many primes. Indeed $m>0$ because
$t^2-4m<0$, and if $q_1,\ldots,q_k$ were all such primes, then for $Q=q_1\cdots q_k$ and a large integer
$y$, $P(mQy)=m(mQ^2y^2-tQy+1)$, where the second factor is larger than $1$ and congruent to $1$ modulo
$Q$, so it has a prime factor outside the list. Since $\gamma$ and $\delta$ are nonzero and the trace form
of $E$ is nondegenerate, there is an integer $e\geq1$ divisible by $\gamma$ in $\mathcal O_E$ such that $e$
times every $\mathbf Z$-linear form on $\mathcal O_E$ is $y\mapsto\operatorname{Tr}_{E/\mathbf Q}(\delta
zy)$ for some $z$. Choose $r$ with a root $\rho$ of $X^2-tX+m$ modulo $r$, larger than every element of
$X$, than $p$, $e$ and $|t^2-4m|$, and than the norms of the places of $S$. Then $r$ is good, and it lies
under no place of $S$.

Since $r\nmid t^2-4m$ and $X^2-tX+m$ has the root $\rho$ modulo $r$, the prime $r$ splits in
$\mathbf Q(\omega)=\mathbf Q(x)$, into $\mathfrak r=(r,\omega-\rho)$ and $c(\mathfrak r)$. As $E=F(\omega)$
and $F\cap\mathbf Q(\omega)=\mathbf Q$, every place of $F$ over $r$ splits in $E$, into one place over
$\mathfrak r$ and its conjugate, which lies over $c(\mathfrak r)$. Let $W$ be the set of places of $E$ over
$\mathfrak r$, that is, those $w\mid r$ with $\omega-\rho\in w$. Then (2) and (3) hold, $W$ is nonempty,
and two places of $W$ over the same place of $F$ would be conjugate, which is impossible since only one of
them lies over $\mathfrak r$.
\end{proof}
% Lean: exists_root (with b exceeding X, p, e, |t^2-4m| and the norms of S), pPerfCof for e, and
% Pl.sides_holds for the description of the places over r; W is hT.toFinset and Half r W holds.

\begin{definition}\label{gal-test-function}
In the setting of Definition~\ref{gal-trace-seen-operators}, let $x\in E$ with
$\mathbf Q(x)=\mathbf Q(\omega)$, let $U$ be an automorphic representation of $U_n$, and let $V$ be a vector
for $U$ at a place $w_0$ prime to $p$ (Definition~\ref{gal-trace-vector}). Let $r$ and $W$ be given by
Lemma~\ref{gal-trace-auxiliary-prime} for the set $S$ of $V$ and some finite set of primes, and let $N\geq3$
be prime to $p$ and divisible by $r$, such that every element of $K_N$ stabilizes $V$. For $w\in W$ let
$I_w\subset\operatorname{GL}_n(\mathcal O_w)$ be the Iwahori subgroup. The \emph{type} of these data
consists of the group
\[
 K=\bigl\{k\in G(\mathbf A_f^p):k=j(k')\text{ with }\theta(k')\in K_S,\text{ and }
 \operatorname{pr}_w(k)\in I_w\text{ for }w\in W\bigr\},
\]
the group $J'=\{y\in\operatorname{GL}_n(E_{w_0}):{}^ty^{-1}\in J\}$, and its character
$\psi(y)=\chi({}^ty^{-1})$. For $h\in G(\mathbf A_f^p)$ with $\operatorname{pr}_{w_0}(h)=1$, the
\emph{test function} $e_h$ on $G(\mathbf A_f^p)$ is
\[
 e_h(g)=[K:K_N]^{-1}\,\psi(\operatorname{pr}_{w_0}(g))^{-1}\ \text{ if }g\in KhK\text{ and }
 \operatorname{pr}_{w_0}(g)\in J',\qquad e_h(g)=0\ \text{ otherwise},
\]
and $e=e_1$.
\end{definition}
\leanref{Development/AutomorphicArithmetic/Foundations/GaloisRepresentations/SphericalCoordinates/AlgebraEquivSetup/HeckeTraceEq.lean}{UpperTheory.Gal.W2.Seed.TyLev}
\leanref{Development/AutomorphicArithmetic/Foundations/GaloisRepresentations/SphericalCoordinates/AlgebraEquivSetup/HeckeTraceEq.lean}{UpperTheory.Gal.W2.Seed.TyLev.tf}
\leanref{Development/AutomorphicArithmetic/Foundations/GaloisRepresentations/ScalarCommutation/RelIndexNeZeroOfOpen.lean}{UpperTheory.Gal.Spc.tyEx}
\leanref{Development/AutomorphicArithmetic/Foundations/GaloisRepresentations/ScalarCommutation/ExistsIdempotentHeckeOperatorWithRegularSupportAlt.lean}{UpperTheory.Gal.Spc.TyOK}

\begin{lemma}\label{gal-test-function-type}
In the setting of Definition~\ref{gal-test-function}:
\begin{enumerate}
\item $K_N\subset K$ with finite index, $\operatorname{pr}_{w_0}(K)\subset J'$,
  $\operatorname{loc}_{w_0}(J')\subset K$, and $\psi(\operatorname{pr}_{w_0}(k))=1$ for $k\in K_N$;
\item for $w\in W$, $\operatorname{pr}_w(K)\subset I_w$ and $\operatorname{loc}_w(I_w)\subset K$; and for
  every place $w$ split over $F$ and prime to $p$, with $w,c(w)\notin S$ and over a place of $F$ different
  from those under $W$, $\operatorname{pr}_w(K)\subset\operatorname{GL}_n(\mathcal O_w)$ and
  $\operatorname{loc}_w(\operatorname{GL}_n(\mathcal O_w))\subset K$;
\item there is an integer $c_0\geq1$ prime to $p$ with $c_0k\hat\Lambda^p\subset\hat\Lambda^p$ for every
  $k\in K$;
\item $J'$ is compact open, $\ker\psi$ is open, and every irreducible smooth representation of
  $\operatorname{GL}_n(E_{w_0})$ with a nonzero vector on which $J'$ acts by $\psi$ is cuspidal;
\item $G(\mathbf A_f^p)$ is the product of $\operatorname{loc}_{w_0}(\operatorname{GL}_n(E_{w_0}))$ and the
  kernel of $\operatorname{pr}_{w_0}$, and $K_N$ is the product of its intersections with these two groups.
\end{enumerate}
\end{lemma}
\leanref{Development/AutomorphicArithmetic/Foundations/GaloisRepresentations/ScalarCommutation/TyFact.lean}{UpperTheory.Gal.Spc.tyFact}
\leanref{Development/AutomorphicArithmetic/Foundations/GaloisRepresentations/ScalarCommutation/ContinuousInvT.lean}{UpperTheory.Gal.Spc.tyTop}
\leanref{Development/AutomorphicArithmetic/Foundations/GaloisRepresentations/ScalarCommutation/SplitC.lean}{UpperTheory.Gal.Spc.splitC}

\begin{proof}
For $k=j(k')$, $\operatorname{pr}_{w_0}(k)=k'_{w_0}$ and $\theta(k')_{w_0}={}^tk'^{-1}_{w_0}$, so
$\operatorname{pr}_{w_0}(k)\in J'$ if and only if $\theta(k')_{w_0}\in J$, and then
$\psi(\operatorname{pr}_{w_0}(k))=\chi(\theta(k')_{w_0})$.

(1) Every $k\in K_N$ stabilizes $V$, so $k=j(k')$ with $\theta(k')\in K_S$ and
$\chi(\theta(k')_{w_0})=1$. For $w\in W$, $w$ lies over $r\mid N$ and $r$ is good, so
$a_w\in\operatorname{GL}_n(\mathcal O_w)$ by Lemma~\ref{gal-identification}(3), and
$\operatorname{pr}_w(k)=a_w^{-1}k_wa_w$ is congruent to $1$ modulo $\varpi_w$, hence lies in $I_w$. So
$K_N\subset K$, and $\psi(\operatorname{pr}_{w_0}(k))=1$. The index is finite because $K$ is compact and
$K_N$ is open. For $k\in K$, $\theta(k')\in K_S$ gives $\theta(k')_{w_0}\in K_{w_0}=J$, so
$\operatorname{pr}_{w_0}(k)\in J'$. For $y\in J'$, $\theta(\iota_{w_0}(y))$ has component ${}^ty^{-1}\in J$
at $w_0$, the corresponding component at $c(w_0)\notin S$, and $1$ elsewhere, so it lies in $K_S$; and
$\operatorname{pr}_w(\operatorname{loc}_{w_0}(y))=1$ for $w\in W$, since no place of $W$ lies over the
place of $F$ under $w_0\in S$. So $\operatorname{loc}_{w_0}(y)\in K$.

(2) The first part holds by the definition of $K$ and by the argument of (1), using that distinct places of
$W$ lie over distinct places of $F$ and that no place of $S$ lies over $r$. The second holds because $K_S$
has component $\operatorname{GL}_n(\mathcal O_w)$ at such $w$, and $\theta$ preserves
$\operatorname{GL}_n(\mathcal O_w)$.

(3) The group $K$ is compact, and its elements act at almost all $w\nmid p$ through
$a_w\operatorname{GL}_n(\mathcal O_w)a_w^{-1}=\operatorname{GL}_n(\mathcal O_w)$; so a common denominator
$c_0$ exists, and its $p$-part, a unit at the places prime to $p$, can be removed.

(4) The map $y\mapsto{}^ty^{-1}$ is a homeomorphism carrying $J'$ to $J$ and $\psi$ to $\chi$. If a
representation $\sigma$ has a nonzero vector on which $J'$ acts by $\psi$, then $\sigma$ composed with this
automorphism has one on which $J$ acts by $\chi$, so it is cuspidal; as the automorphism carries unipotent
radicals of parabolic subgroups to unipotent radicals of parabolic subgroups, $\sigma$ is cuspidal.

(5) Since $w_0$ is split over $F$, $\operatorname{loc}_{w_0}$ is a section of $\operatorname{pr}_{w_0}$
whose image commutes with the kernel of $\operatorname{pr}_{w_0}$, which gives the product decomposition;
$K_N$ is defined by conditions at each place, so it is the product of its components.
\end{proof}
% Lean: tyEx builds Ty with K = ((lev V.S V.Kw).comap F.th).map F.j meet the Iwahori conditions at W,
% K0 = S.L.KN N, J = V.J.comap invT, psi = V.chi o invT; idx by relIdx, cms by pKNIdx. tyFact gives
% (2) and (3) (the bound c through nsmul_integ_iff); tyTop gives (4) and (5) through pCuspTh and the
% factorization facEx; splitC turns (5) into the splitting Split used by specOfF.

\begin{definition}\label{gal-iwahori-operator}
Let $G$ be a group and $I\subset G$ a subgroup such that $IgI/I$ is finite for every $g\in G$. Let $A$ be a
commutative ring and $V$ an $A$-module with an action of $G$ by $A$-linear maps, and let $V^I$ be the
submodule of vectors fixed by $I$. For $g\in G$, the \emph{operator of $IgI$} is the $A$-linear map
\[
 [IgI]:V^I\to V^I,\qquad [IgI]v=\sum_{x\in IgI/I}xv,
\]
where $x$ runs over a set of representatives of the cosets in $IgI$. The sum does not depend on the
representatives because $v$ is fixed by $I$, and it is fixed by $I$ because $I$ permutes $IgI/I$. This
applies to $G=\operatorname{GL}_n(k)$, for a non-archimedean local field $k$ with ring of integers
$\mathcal O$, uniformizer $\varpi$ and residue field $\kappa$ and for $n\geq1$, and to the following
subgroups. The \emph{Iwahori subgroup} is the group of elements of $\operatorname{GL}_n(\mathcal O)$ whose
entries below the diagonal lie in $\varpi\mathcal O$. For a subgroup $C$ of $\kappa^\times$, let $I_C$ be the
group of elements of the Iwahori subgroup whose diagonal entries have residues in $C$. Thus $I_C$ is the
Iwahori subgroup when $C=\kappa^\times$, and the pro-$p$ Iwahori subgroup when $C=1$. Each $I_C$ is compact
and open, as it contains the kernel of reduction modulo $\varpi$ on $\operatorname{GL}_n(\mathcal O)$. A
diagonal matrix $t=\operatorname{diag}(t_1,\ldots,t_n)$ in
$\operatorname{GL}_n(k)$ is \emph{dominant} if $|t_1|\leq|t_2|\leq\cdots\leq|t_n|$ for the absolute value
of $k$.
\end{definition}
\leanref{Development/AutomorphicArithmetic/AutomorphicVarieties/AlgebraTools/HeckePairTop.lean}{UpperTheory.V.AlgebraTools.hk}
\leanref{Development/AutomorphicArithmetic/AutomorphicVarieties/AlgebraTools/IwHkApply.lean}{UpperTheory.V.AlgebraTools.hkApply}
\leanref{Development/AutomorphicArithmetic/AutomorphicVarieties/AlgebraTools/IwHeckePairIwC.lean}{UpperTheory.V.AlgebraTools.heckePairIwC}
\leanref{Development/AutomorphicArithmetic/AutomorphicVarieties/AlgebraTools/HeckePairTop.lean}{UpperTheory.V.AlgebraTools.iwC}
\leanref{Development/AutomorphicArithmetic/AutomorphicVarieties/AlgebraTools/IwMemIwC.lean}{UpperTheory.V.AlgebraTools.memIwC}
\leanref{Development/AutomorphicArithmetic/AutomorphicVarieties/AlgebraTools/HeckePairTop.lean}{UpperTheory.V.AlgebraTools.Dom}

\begin{lemma}\label{gal-iwahori-operator-product}
Let $G$, $I$, $A$ and $V$ be as in Definition~\ref{gal-iwahori-operator}, and let $x,y\in G$ with
$xIy\subset IxyI$ and $|IxyI/I|=|IxI/I|\cdot|IyI/I|$. Then $[IxI][IyI]=[IxyI]$ on $V^I$.
\end{lemma}
\leanref{Development/AutomorphicArithmetic/AutomorphicVarieties/AlgebraTools/IwHkMul.lean}{UpperTheory.V.AlgebraTools.hkMul}

\begin{proof}
Let $S$ and $S'$ be sets of representatives of $IxI/I$ and $IyI/I$. For $a\in S$ and $b\in S'$,
$ab\in IxI\cdot IyI=I(xIy)I\subset IxyI$. Every element of $IxyI/I$ is $cxyI$ with $c\in I$. Write $cx=ai$
with $a\in S$ and $i\in I$, and $iyI=bI$ with $b\in S'$; then $cxyI=abI$. So $(a,b)\mapsto abI$ maps
$S\times S'$ onto $IxyI/I$, and by the hypothesis on the numbers of cosets it is a bijection. Hence the
products $ab$ form a set of representatives of $IxyI/I$, and for $v\in V^I$,
$[IxI][IyI]v=\sum_{a\in S}\sum_{b\in S'}abv=[IxyI]v$.
\end{proof}

\begin{lemma}\label{gal-iwahori-operator-count}
Let $G$ be a group and $I\subset G$ a subgroup such that $IgI/I$ is finite for every $g\in G$. Let
$x,y\in G$ satisfy:
\begin{enumerate}
\item every $g\in I$ with $(xy)^{-1}gxy\in I$ satisfies $x^{-1}gx\in I$;
\item every $b\in I$ is a product $b=b'b''$ of elements $b',b''\in I$ with $xb'x^{-1}\in I$ and
  $y^{-1}b''y\in I$.
\end{enumerate}
Then $xIy\subset IxyI$ and $|IxyI/I|=|IxI/I|\cdot|IyI/I|$.
\end{lemma}
\leanref{Development/AutomorphicArithmetic/AutomorphicVarieties/AlgebraTools/IwLeviMul.lean}{UpperTheory.V.AlgebraTools.leviMul}

\begin{proof}
For $b=b'b''$ as in (2), $xby=(xb'x^{-1})\,xy\,(y^{-1}b''y)\in IxyI$. For $z\in G$, the stabilizer of
$zI$ in $I$ is $I_z=I\cap zIz^{-1}$, so $|IzI/I|=[I:I_z]$. By (1), $I_{xy}\subset I_x$, so
$|IxyI/I|=[I:I_x]\,[I_x:I_{xy}]$. As $I_{xy}$ is the stabilizer of $xyI$ in $I_x$, the index
$[I_x:I_{xy}]$ is the number of elements of the orbit of $xyI$ under $I_x$. The elements of $I_x$ are the
$xcx^{-1}$ with $c\in I$ and $xcx^{-1}\in I$, and $xcx^{-1}\cdot xyI=xcyI$. By (2), every coset $byI$ with
$b\in I$ equals $b'yI$ with $b'\in I$ and $xb'x^{-1}\in I$. So the orbit is $x\cdot(IyI/I)$, which has
$|IyI/I|$ elements.
\end{proof}

\begin{lemma}\label{gal-iwahori-operator-normalizer}
Let $G$, $I$, $A$ and $V$ be as in Definition~\ref{gal-iwahori-operator}, and let $u\in G$ normalize $I$.
Then $[IuI]v=uv$ for $v\in V^I$, $[I1I]$ is the identity, $[IuI]$ is invertible, and $[IuxI]=[IuI][IxI]$
and $[IxuI]=[IxI][IuI]$ for every $x\in G$.
\end{lemma}
\leanref{Development/AutomorphicArithmetic/AutomorphicVarieties/AlgebraTools/IwHkNorm.lean}{UpperTheory.V.AlgebraTools.hkNorm}

\begin{proof}
Let $S$ be a set of representatives of $IxI/I$. Since $uI=Iu$, $IuxI=uIxI$ and $IxuI=IxIu$, so $uS$ is a
set of representatives of $IuxI/I$, and $Su$ is one of $IxuI/I$, as $yuI=yIu$ for $y\in S$. For $x=1$ and
$S=\{1\}$ this gives $[IuI]v=uv$, and $[I1I]v=v$. For $v\in V^I$, the vector $uv$ is fixed by $I=uIu^{-1}$,
and $[IuxI]v=\sum_{y\in S}uyv=[IuI][IxI]v$ and $[IxuI]v=\sum_{y\in S}yuv=[IxI][IuI]v$. As $u^{-1}$ also
normalizes $I$, taking $x=u^{-1}$ gives $[IuI][Iu^{-1}I]=[I1I]=[Iu^{-1}I][IuI]$, so $[IuI]$ is invertible.
\end{proof}

\begin{lemma}\label{gal-iwahori-operator-factorization}
Let $k$ be a non-archimedean local field with ring of integers $\mathcal O$, uniformizer $\varpi$ and
residue field $\kappa$, let $n\geq1$, let $C$ be a subgroup of $\kappa^\times$ and let $I=I_C$
(Definition~\ref{gal-iwahori-operator}). Let $U$ and $L$ be the groups of upper and of lower unitriangular
matrices in $\operatorname{GL}_n(k)$, and $T$ the group of diagonal matrices. Then:
\begin{enumerate}
\item every $g\in I$ is a product $g=udl$ with $u\in U\cap\operatorname{GL}_n(\mathcal O)$,
  $d\in T\cap I$ and $l\in L\cap I$, and these three groups lie in $I$;
\item if $u\in U$, $d\in T$ and $l\in L$ satisfy $udl=1$, then $u=d=l=1$;
\item if $u\in U$ and $d\in T$ satisfy $ud\in I$, then $u\in I$ and $d\in I$;
\item if $t\in T$ is dominant, then $tut^{-1}\in U\cap\operatorname{GL}_n(\mathcal O)$ for
  $u\in U\cap\operatorname{GL}_n(\mathcal O)$, and $t^{-1}lt\in L\cap I$ for $l\in L\cap I$.
\end{enumerate}
\end{lemma}
\leanref{Development/AbelianVarietyTheory/Proofs/TwistedOperators/IwahoriTraceOfBorelCasselman.lean}{GlobalTheory.IwahoriTraceProof.iwahori_factorisation}
\leanref{Development/AutomorphicArithmetic/AutomorphicVarieties/AlgebraTools/IwFactTorus.lean}{UpperTheory.V.AlgebraTools.factTorus}

\begin{proof}
(1) The reduction of $g$ modulo $\varpi$ is upper triangular and invertible. For each $j$ let $A_j$ be the
matrix whose rows below the $j$-th are those of $g$ and whose other rows are those of the identity. Its
reduction is upper triangular with nonzero diagonal entries, so $A_j\in\operatorname{GL}_n(\mathcal O)$.
Let $l'$ be the matrix whose $j$-th column is the $j$-th column of $A_j^{-1}$, so that $A_jl'$ has $j$-th
column $e_j$. The rows $i\leq j$ of this identity show that $l'$ is lower unitriangular, so
$l'\in L\cap\operatorname{GL}_n(\mathcal O)$, and the rows $i>j$ show that $gl'$ is upper triangular. As
$gl'\in\operatorname{GL}_n(\mathcal O)$, its diagonal entries are units, and $gl'=du''$ with
$d\in T\cap\operatorname{GL}_n(\mathcal O)$ and $u''\in U\cap\operatorname{GL}_n(\mathcal O)$. Then
$g=udl$ with $u=du''d^{-1}\in U\cap\operatorname{GL}_n(\mathcal O)$ and $l=l'^{-1}$. The matrices $u$ and
$d$ have no nonzero entries below the diagonal, so they lie in the Iwahori subgroup, and so does
$l=(ud)^{-1}g$. The elements of $U\cap\operatorname{GL}_n(\mathcal O)$ and of $L$ have diagonal entries $1$,
so $U\cap\operatorname{GL}_n(\mathcal O)$ lies in $I$, $l\in I$, and then $d=u^{-1}gl^{-1}\in I$.

(2) The matrix $u=l^{-1}d^{-1}$ is upper unitriangular and lower triangular, so $u=1$. Then $dl=1$, and the
diagonal entries give $d=1$ and then $l=1$.

(3) Write $ud=u'd'l'$ by (1). Then $(u'^{-1}u)(dd'^{-1})(d'l'^{-1}d'^{-1})=1$ with
$d'l'^{-1}d'^{-1}\in L$, so $u=u'$ and $d=d'$ by (2), and both lie in $I$.

(4) The entry $(a,b)$ of $tut^{-1}$ is $t_at_b^{-1}u_{ab}$, and $|t_a/t_b|\leq1$ for $a\leq b$. The entry
$(a,b)$ of $t^{-1}lt$ is $t_a^{-1}t_bl_{ab}$, and $|t_b/t_a|\leq1$ for $b\leq a$.
\end{proof}

\begin{lemma}\label{gal-iwahori-operator-dominant}
Let $k$ be a non-archimedean local field with ring of integers $\mathcal O$ and residue field $\kappa$, let
$n\geq1$, let $C$ be a subgroup of $\kappa^\times$ and $I=I_C$, let $A$ be a commutative ring and $V$ an
$A$-module with an $A$-linear action of $\operatorname{GL}_n(k)$, and let $t,t'\in\operatorname{GL}_n(k)$ be
dominant diagonal matrices (Definition~\ref{gal-iwahori-operator}). Then $tt'$ is dominant and
$[ItI][It'I]=[Itt'I]$ on $V^I$.
\end{lemma}
\leanref{Development/AutomorphicArithmetic/AutomorphicVarieties/AlgebraTools/IwDomMul.lean}{UpperTheory.V.AlgebraTools.domMul}
\leanref{Development/AutomorphicArithmetic/AutomorphicVarieties/AlgebraTools/IwLeviMul.lean}{UpperTheory.V.AlgebraTools.leviMul}

\begin{proof}
The product $tt'$ is dominant since $|t_at'_a|\leq|t_bt'_b|$ for $a\leq b$. By
Lemmas~\ref{gal-iwahori-operator-product} and~\ref{gal-iwahori-operator-count}, it suffices to check
conditions (1) and (2) of Lemma~\ref{gal-iwahori-operator-count} for $x=t$ and $y=t'$. For (2), write
$b\in I$ as $udl$ by Lemma~\ref{gal-iwahori-operator-factorization}(1), and let $b'=ud$ and $b''=l$. Then
$tb't^{-1}=(tut^{-1})d$ and $t'^{-1}b''t'$ lie in $I$ by Lemma~\ref{gal-iwahori-operator-factorization}(4).
For (1), let $s=tt'$, and let $g=udl\in I$ with $s^{-1}gs\in I$. Then $s^{-1}gs=(s^{-1}us)\,d\,(s^{-1}ls)$,
and $s^{-1}ls\in I$ by Lemma~\ref{gal-iwahori-operator-factorization}(4), so $(s^{-1}us)d\in I$, and
$s^{-1}us\in I$ by Lemma~\ref{gal-iwahori-operator-factorization}(3). Hence $t^{-1}ut=t'(s^{-1}us)t'^{-1}$
lies in $I$ by Lemma~\ref{gal-iwahori-operator-factorization}(4) for $t'$, and
$t^{-1}gt=(t^{-1}ut)\,d\,(t^{-1}lt)$ lies in $I$ by Lemma~\ref{gal-iwahori-operator-factorization}(4) for
$t$.
\end{proof}

\begin{lemma}\label{gal-iwahori-operator-simple}
Let $k$ be a non-archimedean local field with ring of integers $\mathcal O$, uniformizer $\varpi$ and
residue field $\kappa$ of $q$ elements, let $n\geq2$, let $C$ be a subgroup of $\kappa^\times$ and $I=I_C$
(Definition~\ref{gal-iwahori-operator}), let $1\leq i<n$, let $s$ be the permutation matrix of the
transposition of $i$ and $i+1$, and for
$c\in\mathcal O$ let $x_c=1+cE_{i,i+1}$, where $E_{i,i+1}$ is the matrix whose only nonzero entry is $1$ at
$(i,i+1)$. Let $R\subset\mathcal O$ be a set of representatives of the residue field with $0\in R$. Then:
\begin{enumerate}
\item an element $h\in I$ satisfies $shs\in I$ if and only if $h_{i,i+1}\in\varpi\mathcal O$;
\item the cosets $x_csI$ for $c\in R$ are distinct, and they are the elements of $IsI/I$;
\item for every $g\in I$ there are $c\in R$ and $b\in I$ with $sg=bsx_c$;
\item for every commutative ring $A$ and $A$-module $V$ with an $A$-linear action of
  $\operatorname{GL}_n(k)$, the operator $[IsI]$ of $V^I$ (Definition~\ref{gal-iwahori-operator}) is
  invertible if $q$ is invertible in $A$, and $[IsI]^2=(q-1)[IsI]+q$ if $C=\kappa^\times$.
\end{enumerate}
\end{lemma}
\leanref{Development/AutomorphicArithmetic/AutomorphicVarieties/AlgebraTools/IwSimpleOrbit.lean}{UpperTheory.V.AlgebraTools.simpleOrbit}
\leanref{Development/AutomorphicArithmetic/AutomorphicVarieties/AlgebraTools/IwSimpleQuad.lean}{UpperTheory.V.AlgebraTools.simpleQuad}

\begin{proof}
(1) The entry $(a,b)$ of $shs$ is $h_{\sigma(a),\sigma(b)}$, where $\sigma$ is the transposition of $i$ and
$i+1$. For $a>b$, the position $(\sigma(a),\sigma(b))$ is again below the diagonal unless $(a,b)=(i+1,i)$,
where it is $(i,i+1)$. The diagonal entries of $shs$ are those of $h$, permuted, so their residues lie in
$C$. As $shs\in\operatorname{GL}_n(\mathcal O)$, it lies in $I$ if and only if its entries below the
diagonal lie in $\varpi\mathcal O$, which for $h\in I$ means $h_{i,i+1}\in\varpi\mathcal O$.

(2) Let $g\in I$. Since $g_{i+1,i+1}$ is a unit, there is $c\in R$ with
$g_{i,i+1}-cg_{i+1,i+1}\in\varpi\mathcal O$, and this is the entry $(i,i+1)$ of $x_c^{-1}g=x_{-c}g\in I$. By
(1), $sx_c^{-1}gs\in I$, that is, $gsI=x_csI$. As $x_c\in I$, the cosets in $IsI$ are the $x_csI$ with
$c\in R$. If $x_csI=x_{c'}sI$, then $sx_{c'-c}s=sx_c^{-1}x_{c'}s\in I$, so $c'-c\in\varpi\mathcal O$ by (1),
and $c=c'$.

(3) Let $g\in I$. There is $c\in R$ with $g_{i,i+1}-cg_{ii}\in\varpi\mathcal O$, and this is the entry
$(i,i+1)$ of $gx_c^{-1}\in I$. By (1), $b=sgx_c^{-1}s$ lies in $I$, and $sg=bsx_c$ since $s^2=1$.

(4) By (2), $[IsI]v=\sum_{c\in R}x_csv$ for $v\in V^I$, so
\[
 [IsI]^2v=\sum_{c\in R}\sum_{c'\in R}x_{c'}\,sx_cs\,v.
\]
The terms with $c=0$ give $\sum_{c'}x_{c'}v=qv$, since $s^2=1$ and $x_{c'}\in I$. Let $c\in R$ be nonzero,
so that $c$ is a unit. A computation in the rows and columns $i$ and $i+1$ gives
$sx_cs=x_{c^{-1}}\,s\,d_c\,x_{c^{-1}}$, where $d_c$ is the diagonal matrix with entries $c$ at $i$,
$-c^{-1}$ at $i+1$ and $1$ elsewhere. The matrix $d_c$ normalizes $I$, since conjugating by a diagonal
matrix with unit entries multiplies the entry $(a,b)$ by a unit and fixes the diagonal, so $[Id_cI]v=d_cv$
by Lemma~\ref{gal-iwahori-operator-normalizer}. As $x_{c^{-1}}\in I$ commutes with every $x_{c'}$ and fixes
$[IsI][Id_cI]v$, the terms with this $c$ give
$x_{c^{-1}}\sum_{c'}x_{c'}s\,d_cv=x_{c^{-1}}[IsI][Id_cI]v=[IsI][Id_cI]v$. Hence
\[
 [IsI]^2=q+\sum_{c\in R,\ c\neq0}[IsI][Id_cI].
\]
If $q$ is invertible in $A$, this shows that
$[IsI]\cdot q^{-1}\bigl([IsI]-\sum_{c\neq0}[Id_cI]\bigr)=1$. Also $sd_cs$ is diagonal with unit entries, so
$[IsI][Id_cI]=[Isd_cI]=[I(sd_cs)sI]=[Isd_csI][IsI]$ by Lemma~\ref{gal-iwahori-operator-normalizer}, and
$q^{-1}\bigl([IsI]-\sum_{c\neq0}[Isd_csI]\bigr)[IsI]=1$. So $[IsI]$ has a left and a right inverse, and is
invertible. If $C=\kappa^\times$, then $d_c\in I$, so $[Id_cI]$ is the identity, and the sum has $q-1$ terms
$[IsI]$, which gives $[IsI]^2=(q-1)[IsI]+q$.
\end{proof}

\begin{lemma}\label{gal-iwahori-operator-permutation}
Let $k$ be a non-archimedean local field with ring of integers $\mathcal O$, uniformizer $\varpi$ and
residue field $\kappa$ of $q$ elements, let $n\geq1$, let $C$ be a subgroup of $\kappa^\times$ and $I=I_C$
(Definition~\ref{gal-iwahori-operator}). For a permutation $w$ of $\{1,\ldots,n\}$ let $P_w$ be the permutation matrix with
$P_we_b=e_{w(b)}$ for the standard basis $(e_b)$, and let $\ell(w)$ be the number of pairs $a<b$ with
$w(a)>w(b)$. For $1\leq i<n$ let $s_i$ be the transposition of $i$ and $i+1$. Then:
\begin{enumerate}
\item if $w^{-1}(i)<w^{-1}(i+1)$, then $\ell(s_iw)=\ell(w)+1$; and if $w\neq1$, there is $i$ with
  $w^{-1}(i+1)<w^{-1}(i)$;
\item $|IP_wI/I|=q^{\ell(w)}$;
\item for every commutative ring $A$ and $A$-module $V$ with an $A$-linear action of
  $\operatorname{GL}_n(k)$, $[IP_{s_i}I][IP_wI]=[IP_{s_iw}I]$ on $V^I$ whenever $w^{-1}(i)<w^{-1}(i+1)$
  (Definition~\ref{gal-iwahori-operator}), and $[IP_wI]$ is invertible if $q$ is invertible in $A$.
\end{enumerate}
\end{lemma}
\leanref{Development/AutomorphicArithmetic/AutomorphicVarieties/AlgebraTools/IwLenStep.lean}{UpperTheory.V.AlgebraTools.lenStep}
\leanref{Development/AutomorphicArithmetic/AutomorphicVarieties/AlgebraTools/IwPermCard.lean}{UpperTheory.V.AlgebraTools.permCard}
\leanref{Development/AutomorphicArithmetic/AutomorphicVarieties/AlgebraTools/IwPermMul.lean}{UpperTheory.V.AlgebraTools.permMul}
\leanref{Development/AutomorphicArithmetic/AutomorphicVarieties/AlgebraTools/IwPermUnit.lean}{UpperTheory.V.AlgebraTools.permUnit}

\begin{proof}
(1) Composing with $s_i$ changes the relative order of $w(a)$ and $w(b)$ only when
$\{w(a),w(b)\}=\{i,i+1\}$. So if $w^{-1}(i)<w^{-1}(i+1)$, the pairs inverted by $s_iw$ are those inverted
by $w$ and the pair $(w^{-1}(i),w^{-1}(i+1))$. If $w^{-1}(i)<w^{-1}(i+1)$ for every $i$, then $w^{-1}$ is
increasing, and $w=1$.

(2) Let $w$ and $i$ satisfy $w^{-1}(i)<w^{-1}(i+1)$, and let $s=P_{s_i}$, $x_c$ and $R$ be as in
Lemma~\ref{gal-iwahori-operator-simple}. For $c\in\mathcal O$, $P_w^{-1}x_cP_w=1+cE_{w^{-1}(i),w^{-1}(i+1)}$
is upper unitriangular with entries in $\mathcal O$, so it lies in $I$. Consider the map
$(c,yI)\mapsto x_csyI$ on $R\times IP_wI/I$. For $y=gP_w$ with $g\in I$, write $sg=bsx_a$ with $b\in I$ by
Lemma~\ref{gal-iwahori-operator-simple}(3); then $x_csgP_wI=x_cbsP_w(P_w^{-1}x_aP_w)I=x_cbsP_wI$, so the
map takes values in $IsP_wI/I$. It is onto: for $g\in I$, Lemma~\ref{gal-iwahori-operator-simple}(2) gives
$gs=x_csb$ with $c\in R$ and $b\in I$, and $gsP_wI=x_cs(bP_w)I$. It is injective: if
$x_csgP_wI=x_{c'}sg'P_wI$ with $g,g'\in I$, then $h=g'^{-1}sx_{c'}^{-1}x_csg$ satisfies $P_w^{-1}hP_w\in I$.
The entry of $P_w^{-1}hP_w$ at $(w^{-1}(i+1),w^{-1}(i))$, which is below the diagonal, is $h_{i+1,i}$, so
$h_{i+1,i}\in\varpi\mathcal O$. Since $sx_{c'}^{-1}x_cs=1+(c-c')E_{i+1,i}$,
\[
 h_{i+1,i}=(g'^{-1}g)_{i+1,i}+(c-c')\,(g'^{-1})_{i+1,i+1}\,g_{ii}.
\]
The first term lies in $\varpi\mathcal O$ as $g'^{-1}g\in I$, and the diagonal entries of elements of $I$
are units, so $c-c'\in\varpi\mathcal O$ and $c=c'$; then $gP_wI=g'P_wI$. Hence
$|IP_{s_iw}I/I|=|IsP_wI/I|=q\,|IP_wI/I|$. Now (2) follows by induction on $\ell(w)$. For $w=1$, $I/I$ has
one element. For $w\neq1$, choose $i$ by (1) with $w^{-1}(i+1)<w^{-1}(i)$; then $w'=s_iw$ satisfies
$w'^{-1}(i)<w'^{-1}(i+1)$ and $s_iw'=w$, so $\ell(w)=\ell(w')+1$ by (1), and $|IP_wI/I|=q\,|IP_{w'}I/I|$.

(3) Let $w^{-1}(i)<w^{-1}(i+1)$ and $s=P_{s_i}$. For $g\in I$, Lemma~\ref{gal-iwahori-operator-simple}(3)
gives $sgP_w=bsP_w(P_w^{-1}x_cP_w)\in IsP_wI$, as in the proof of (2). By (1), (2) and
Lemma~\ref{gal-iwahori-operator-simple}(2), $|IsP_wI/I|=q^{\ell(w)+1}=|IsI/I|\cdot|IP_wI/I|$, so
$[IsI][IP_wI]=[IsP_wI]=[IP_{s_iw}I]$ by Lemma~\ref{gal-iwahori-operator-product}. Invertibility follows by
induction on $\ell(w)$: $[I1I]$ is the identity by Lemma~\ref{gal-iwahori-operator-normalizer}, and for
$w\neq1$, with $i$ and $w'=s_iw$ as in the proof of (2), $[IP_wI]=[IsI][IP_{w'}I]$ is a product of
invertible operators by Lemma~\ref{gal-iwahori-operator-simple}(4) and induction.
\end{proof}

\begin{lemma}\label{gal-iwahori-operator-fundamental}
Let $k$ be a non-archimedean local field with ring of integers $\mathcal O$, uniformizer $\varpi$ and
residue field $\kappa$, let $n\geq1$, let $C$ be a subgroup of $\kappa^\times$ and $I=I_C$
(Definition~\ref{gal-iwahori-operator}), and let $0\leq m\leq n$. Let $\varpi_m=\operatorname{diag}(\varpi,\ldots,\varpi,1,\ldots,1)$, with $m$ entries $\varpi$, let $w_m$ be
the permutation of $\{1,\ldots,n\}$ with $w_m(b)=b+n-m$ for $b\leq m$ and $w_m(b)=b-m$ for $b>m$, and let
$P_{w_m}$ be the permutation matrix with $P_{w_m}e_b=e_{w_m(b)}$. Then $z=P_{w_m}\varpi_m$ normalizes $I$.
\end{lemma}
\leanref{Development/AutomorphicArithmetic/AutomorphicVarieties/AlgebraTools/IwFundNorm.lean}{UpperTheory.V.AlgebraTools.fundNorm}

\begin{proof}
For indices $a,b$ let $\delta(a,b)=1$ if $b<a$ and $\delta(a,b)=0$ otherwise, and let $\epsilon_a=1$ if
$a\leq m$ and $\epsilon_a=0$ otherwise. An element $g\in\operatorname{GL}_n(k)$ lies in the Iwahori
subgroup if and only if $|g_{ab}|\leq|\varpi|^{\delta(a,b)}$ for all $a,b$ and the same holds for $g^{-1}$.
The entry of $zgz^{-1}$
at $(w_m(a),w_m(b))$ is $\varpi^{\epsilon_a-\epsilon_b}g_{ab}$. Checking the cases $a,b\leq m$, $a,b>m$,
$a\leq m<b$ and $b\leq m<a$ gives $\delta(w_m(a),w_m(b))=\epsilon_a-\epsilon_b+\delta(a,b)$. So $zgz^{-1}$
satisfies these bounds if and only if $g$ does. Applied to $g$ and to $g^{-1}$, whose conjugate is
$(zgz^{-1})^{-1}$, this shows that $zgz^{-1}$ lies in the Iwahori subgroup if and only if $g$ does. The
diagonal entry of $zgz^{-1}$ at $(w_m(a),w_m(a))$ is $g_{aa}$, so the diagonal entries of $zgz^{-1}$ are
those of $g$, permuted, and $zgz^{-1}\in I$ if and only if $g\in I$.
\end{proof}

\begin{proposition}\label{gal-iwahori-operator-invertible}
Let $k$ be a non-archimedean local field with ring of integers $\mathcal O$ and residue field $\kappa$ of
$q$ elements, let $n\geq1$, let $C$ be a subgroup of $\kappa^\times$ and $I=I_C$, let $A$ be a commutative
ring, let $V$ be an $A$-module with an $A$-linear action of $\operatorname{GL}_n(k)$, and let
$t\in\operatorname{GL}_n(k)$ be a dominant diagonal matrix (Definition~\ref{gal-iwahori-operator}).
\begin{enumerate}
\item If $q$ is invertible in $A$, the operator $[ItI]$ of $V^I$ is invertible.
\item Suppose that $C=\kappa^\times$, so that $I$ is the Iwahori subgroup, that $A$ is a field of characteristic $0$, and
  that $t=\operatorname{diag}(\varpi^{\nu_1},\ldots,\varpi^{\nu_n})$ for a uniformizer $\varpi$ and integers
  $\nu_1\geq\cdots\geq\nu_n$. Then $It^mI/I$ is finite for every $m\geq0$, and $[ItI]$ is an $A$-linear automorphism of
  $V^I$ with $[ItI]^m=[It^mI]$ for every $m\geq0$.
\item Let $C'$ be a subgroup of $\kappa^\times$ with $I_C\subset I_{C'}$. Let $\mathrm{inc}:V^{I_{C'}}\to V^{I_C}$ be the
  inclusion and $\mathrm{tr}:V^{I_C}\to V^{I_{C'}}$ the map $v\mapsto\sum_{x\in I_{C'}/I_C}xv$, over representatives of
  the cosets. Then $[I_CtI_C]\circ\mathrm{inc}=\mathrm{inc}\circ[I_{C'}tI_{C'}]$ and
  $\mathrm{tr}\circ[I_CtI_C]=[I_{C'}tI_{C'}]\circ\mathrm{tr}$.
\end{enumerate}
\end{proposition}
\leanref{Development/AutomorphicArithmetic/AutomorphicVarieties/AlgebraTools/IwDomMul.lean}{UpperTheory.V.AlgebraTools.domMul}
\leanref{Development/AbelianVarietyTheory/Proofs/DivisorProofs/ExistsTracePowRestrictNeZero.lean}{GlobalTheory.IwahoriUnit.exists_equiv_doubleCoset}

\begin{proof}
(1) Let $\varpi$ be a uniformizer, and for $a\in\mathbf Z^n$ let
$\varpi^a=\operatorname{diag}(\varpi^{a_1},\ldots,\varpi^{a_n})$, which is dominant if and only if
$a_1\geq\cdots\geq a_n$. Let $\varpi_m$, $w_m$ and $z=P_{w_m}\varpi_m$ be as in
Lemma~\ref{gal-iwahori-operator-fundamental}.

First, $[I\varpi_mI]$ is invertible for $0\leq m\leq n$. Indeed $\varpi_m=P_{w_m^{-1}}z$, and $z$
normalizes $I$, so $[I\varpi_mI]=[IP_{w_m^{-1}}I][IzI]$ by Lemma~\ref{gal-iwahori-operator-normalizer}.
Both factors are invertible, by Lemmas~\ref{gal-iwahori-operator-permutation}(3)
and~\ref{gal-iwahori-operator-normalizer}.

Next, let $c$ and $N\geq0$ be integers; by induction on $N$, $[I\varpi^aI]$ is invertible whenever
$a_1\geq\cdots\geq a_n$ lie in $[c,c+N]$. If $N=0$, then $\varpi^a$ is the scalar $\varpi^c$, which
normalizes $I$, and Lemma~\ref{gal-iwahori-operator-normalizer} applies. If $N\geq1$, let $m$ be the number
of indices $j$ with $a_j>c$; as $a$ is decreasing, these are $1,\ldots,m$. Let $a'$ be $a$ with $1$
subtracted from its first $m$ entries. Then $a'$ is decreasing with entries in $[c,c+N-1]$, and
$\varpi^a=\varpi^{a'}\varpi_m$ with both factors dominant. So $[I\varpi^aI]=[I\varpi^{a'}I][I\varpi_mI]$ by
Lemma~\ref{gal-iwahori-operator-dominant}, and both factors are invertible, by induction and by the first
step.

Finally, write $t_j=u_j\varpi^{a_j}$ with $u_j\in\mathcal O^\times$ and $a_j\in\mathbf Z$. Since $t$ is
dominant, $a_1\geq\cdots\geq a_n$. The matrix $u=\operatorname{diag}(u_1,\ldots,u_n)$ normalizes $I$,
since conjugation by it multiplies the entry $(a,b)$ by the unit $u_au_b^{-1}$ and fixes the diagonal, so
$[ItI]=[IuI][I\varpi^aI]$ by Lemma~\ref{gal-iwahori-operator-normalizer}, and both factors are
invertible.

(2) Each $It^mI/I$ is finite, as $It^mI$ is compact and $I$ is open (Definition~\ref{gal-iwahori-operator}). As $A$ has
characteristic $0$, $q$ is invertible in $A$, and $t$ is dominant as $\nu$ is decreasing; so $[ItI]$ is invertible by (1).
The identity matrix $1$ is dominant, so $[I1I][I1I]=[I1I]$ by Lemma~\ref{gal-iwahori-operator-dominant}, and $[I1I]$ is
invertible by (1); hence $[I1I]$ is the identity, which is the case $m=0$. The powers $t^m$ are dominant, so
$[ItI]^{m+1}=[ItI]^m[ItI]=[It^mI][ItI]=[It^{m+1}I]$ by induction and Lemma~\ref{gal-iwahori-operator-dominant}.

(3) Write $I=I_{C'}$, $I'=I_C$, and let $U$, $T$ and $L$ be the groups of upper unitriangular, diagonal and lower
unitriangular matrices. By Lemma~\ref{gal-iwahori-operator-factorization}(1), for $I$ and for $I'$, every $g\in I$ is
$udl$ with $u\in U\cap I$, $d\in T\cap I$ and $l\in L\cap I$, and $U\cap I=U\cap\operatorname{GL}_n(\mathcal O)$. As the
elements of $U$ and $L$ have diagonal entries $1$, $U\cap I\subset I'$ and $L\cap I\subset I'$. The residues of the
diagonal entries give a homomorphism from the Iwahori subgroup to $(\kappa^\times)^n$, as the reductions of its elements are
upper triangular; $I'$ and $I$ are the inverse images of $C^n$ and $C'^n$, so $I'$ is normal in $I$ and $I/I'$ is
finite. By Lemma~\ref{gal-iwahori-operator-factorization}(4), $t^{-1}(L\cap I)t\subset L\cap I$ and
$t^{-1}(L\cap I')t\subset L\cap I'$, and $t$ commutes with $T$, which normalizes $U$.

\emph{Representatives.} As $ItI/I$ is finite, there is a finite set $r\subset U\cap I$ such that for every $x\in U\cap I$
exactly one $u\in r$ has $(ut)^{-1}xt\in I$. For $u,x\in U\cap I$, $(ut)^{-1}xt=t^{-1}u^{-1}xt$ lies in $U$, so it lies
in $I$ if and only if it lies in $U\cap I\subset I'$: the same $u$ is the only one with $(ut)^{-1}xt\in I'$. As $I/I'$ is
finite, there is a finite set $r'\subset T\cap I$ such that for every $x\in T\cap I$ exactly one $j\in r'$ has $j^{-1}x\in
I'$.

\emph{The cosets.} Let $g=udl\in I$. Then $(ut)^{-1}gt=d\,(t^{-1}lt)\in I$, as $t$ commutes with $d$. With the $u'\in r$
that has $(u't)^{-1}ut\in I$, $gtI=u'tI$; and $u'$ is unique, by the choice of $r$. So the $ut$, $u\in r$, represent
$ItI/I$ once each. In the same way, using the factorization in $I'$ and the second property of $r$, they represent
$I'tI'/I'$ once each. Let $j\in T\cap I$. Conjugation by $j$ preserves $I'$, which is normal in $I$, and fixes $t$, so it
permutes $I'tI'/I'$, and the $j^{-1}ujt=j^{-1}(ut)j$, $u\in r$, also represent $I'tI'/I'$ once each. Finally, the $j\in
r'$ represent $I/I'$ once each: for $g=udl\in I$ and $j\in r'$ with $j^{-1}d\in I'$,
$j^{-1}g=(j^{-1}uj)(j^{-1}d)l\in I'$, as $j^{-1}uj\in U\cap I$; and $jI'=j'I'$ with $j,j'\in r'$ forces $j=j'$.

\emph{The operators.} For $v\in V^I$, $[I'tI']v=\sum_{u\in r}utv=[ItI]v$. For $v\in V^{I'}$, $\mathrm{tr}(v)=\sum_{j\in
r'}jv$, and, using for each $j$ the representatives $j^{-1}ujt$ of $I'tI'/I'$ and $jt=tj$,
\[
 \mathrm{tr}([I'tI']v)=\sum_{j\in r'}\sum_{u\in r}j\,j^{-1}ujt\,v=\sum_{u\in r}ut\sum_{j\in r'}jv=[ItI]\,\mathrm{tr}(v).
\]
% Lean: (1) is the second part of domMul (hfund, the inner induction key with exists_cut, and unit_norm). (2) is
% exists_equiv_doubleCoset: finite_cosets, hone from the first two parts of domMul for C = top, and hpow. (3) is the third
% part of domMul, through leviLevel for I = iwC C', I' = iwC C and the Levi factorizations of factTorus (the same U, M =
% the diagonal matrices, and L); leviLevel_reps gives r and r', leviLevel_cover and its variant for I' the cosets, and
% leviLevel_sum the sums sI, sI' (for j in M ⊓ I) and sJ.
\end{proof}

\begin{lemma}\label{gal-test-function-iwahori}
Let $k$ be a non-archimedean local field with ring of integers $\mathcal O$ and uniformizer $\varpi$, let
$n\geq1$, let $I\subset\operatorname{GL}_n(\mathcal O)$ be the Iwahori subgroup and let
$t=\operatorname{diag}(1,\varpi^{-1},\ldots,\varpi^{-(n-1)})$. Then:
\begin{enumerate}
\item there is a bi-$I$-invariant compactly supported function $\phi$ on $\operatorname{GL}_n(k)$ with
  $\phi\ast1_{ItI}=1_I$ for the convolution relative to $I$;
\item $ItI\cdot It^mI\subset It^{m+1}I$ for every $m\geq0$;
\item for $m\geq1$, every element of $It^mI$ has separable characteristic polynomial.
\end{enumerate}
\end{lemma}
\leanref{Development/AutomorphicArithmetic/Foundations/GaloisRepresentations/ProjectiveAutomorphicHecke/IwInvT.lean}{UpperTheory.Gal.CertificateRingHomMap.iwInvT}
\leanref{Development/AutomorphicArithmetic/IwahoriCommensurability/IwahoriDoubleCosetTorusMulOfStrictAnti.lean}{UpperTheory.iwahori_doubleCoset_torus_mul_of_strictAnti}
\leanref{Development/AbelianVarietyTheory/Proofs/DivisorProofs/IwahoriRegularSplit.lean}{GlobalTheory.IwahoriRegular.regular_split}

\begin{proof}
(1) Let $H$ be the space of bi-$I$-invariant compactly supported functions
$\operatorname{GL}_n(k)\to\mathbf C$. Let $V$ be the space of functions $f:\operatorname{GL}_n(k)\to\mathbf C$
with $f(gb)=f(g)$ for $b\in I$ that vanish outside finitely many cosets $gI$, with $\operatorname{GL}_n(k)$
acting by $(hf)(g)=f(h^{-1}g)$. Every double coset $IgI$ is a finite union of cosets $xI$, so the vectors
of $V$ fixed by $I$ are the elements of $H$. For $f\in H$, the operator of $ItI$
(Definition~\ref{gal-iwahori-operator}) is
\[
 ([ItI]f)(g)=\sum_{x\in ItI/I}f(x^{-1}g)=(1_{ItI}\ast f)(g).
\]
The exponents $0,-1,\ldots,-(n-1)$ of $t$ decrease, so $t$ is dominant, and the number of elements of the
residue field of $k$ is invertible in $\mathbf C$. By Proposition~\ref{gal-iwahori-operator-invertible}(2), for $C$
the whole multiplicative group of the residue field and $A=\mathbf C$, $f\mapsto1_{ItI}\ast f$ is a bijection of $H$. Let $\phi\in H$ with $1_{ItI}\ast\phi=1_I$. Convolution on
$H$ is associative with unit $1_I$, so $1_{ItI}\ast(\phi\ast1_{ItI})=1_I\ast1_{ItI}=1_{ItI}=1_{ItI}\ast1_I$,
and injectivity gives $\phi\ast1_{ItI}=1_I$. For a Chevalley group the corresponding invertibility of the
characteristic functions of Iwahori double cosets is part of the structure of the Iwahori--Hecke algebra
\cite{IM65}.
% Lean: iwInvT applies exists_equiv_doubleCoset (Proposition gal-iwahori-operator-invertible) to the private
% representation lrep on cfun, and the private exists_conv_inv moves the inverse to the other side
% (Final.HeckeAssoc, Final.HeckeUnit; finiteness of cosets by finite_cosets and comm_of_fin).

(2) Every element of $I$ is a product $cdl$ of an upper unitriangular $c\in I$, a diagonal $d\in I$ and a
lower unitriangular $l\in I$, whose entries below the diagonal lie in $\varpi\mathcal O$. Conjugation by
$t$ multiplies the entry $(i,j)$ by $\varpi^{j-i}$, so $tct^{-1}\in I$; and $t^{-m}lt^m\in I$, since its
entry $(i,j)$ is $\varpi^{m(i-j)}l_{ij}$. For $i_1ti_2\in ItI$ and $i_3t^mi_4\in It^mI$, write
$i_2i_3=cdl$; as $d$ commutes with $t$,
$i_1ti_2i_3t^mi_4=i_1(tct^{-1})d\cdot t^{m+1}\cdot(t^{-m}lt^m)i_4\in It^{m+1}I$.

For (3), an element of $It^mI$ is conjugate under $I$ to $xt^m$ with $x\in I$. For a set $A$ of $b$ indices,
the principal minor of $xt^m$ on $A$ is $\det(x_{A,A})\,\varpi^{-m\sum_{i\in A}(i-1)}$, and
$\det(x_{A,A})$ is a unit because the reduction of $x$ is triangular with invertible diagonal. So the
coefficient of $X^{n-b}$ in the characteristic polynomial, up to sign the sum of these minors, has valuation
$-m\sum_{i=n-b+1}^{n}(i-1)$, attained only by the last $b$ indices. The Newton polygon then has the $n$
distinct slopes $-m(n-1),\ldots,-m,0$, so the $n$ roots have distinct valuations and are distinct.
\end{proof}

\begin{lemma}\label{gal-test-function-powers}
In the setting of Definition~\ref{gal-test-function}, let $\ast$ be the convolution of functions on
$G(\mathbf A_f^p)$ relative to $K_N$, $(f\ast f')(x)=\sum_{y\in G(\mathbf A_f^p)/K_N}f(y)f'(y^{-1}x)$. There
are bi-$K_N$-invariant compactly supported functions $u$ and $u'$ such that:
\begin{enumerate}
\item $e\ast u=u=u\ast e$ and $u'\ast u=e$;
\item $u\ast e_h=e_h\ast u$ for every $h$ with $\operatorname{pr}_{w_0}(h)=1$ and $\operatorname{pr}_w(h)=1$
  for $w\in W$;
\item for every $m\geq1$ there are subsets $X_w\subset\operatorname{GL}_n(E_w)$, $w\in W$, with
  $I_wX_w\subset X_w$, whose elements have separable characteristic polynomial, such that whenever
  $(u^{\ast m}\ast e)(g)\neq0$, $g^{-1}$ is integral and $\operatorname{pr}_w(g)\in X_w$ for $w\in W$.
\end{enumerate}
\end{lemma}
\leanref{Development/AutomorphicArithmetic/Foundations/GaloisRepresentations/ScalarCommutation/ExistsIdempotentHeckeOperatorWithRegularSupport.lean}{UpperTheory.Gal.Spc.auxFn}
\leanref{Development/AutomorphicArithmetic/Foundations/GaloisRepresentations/ScalarCommutation/ExistsInvertibleHeckeFunctionAtFinitePlaces.lean}{UpperTheory.Gal.Spc.exists_invertible_hecke_function_at_finite_places}

\begin{proof}
For $w\in W$ let $t_w\in\operatorname{GL}_n(E_w)$ be the element $t$ of
Lemma~\ref{gal-test-function-iwahori}, and let $\tau=\prod_{w\in W}\operatorname{loc}_w(t_w)$. By
Lemma~\ref{gal-trace-auxiliary-prime}(4) and since no place of $W$ lies over the place of $F$ under $w_0$,
$\operatorname{pr}_{w_0}(\tau)=1$ and $\operatorname{pr}_w(\tau)=t_w$ for $w\in W$.

First, $e_\tau$ commutes with the functions $e_h$ of (2), and there is $u''$ with $u''\ast e_\tau=e$. This
follows by induction on $W$. By Lemma~\ref{gal-test-function-type}(1)--(2), at each $w\in W$ the group $K$
is the product of $\operatorname{loc}_w(I_w)$ and its intersection with the kernel of
$\operatorname{pr}_w$, and $\operatorname{loc}_w(I_w)$ commutes with $\operatorname{loc}_{w_0}(J')$. So for
$h$ with $\operatorname{pr}_{w_0}(h)=1$ and $\operatorname{pr}_w(h)=1$,
$e_{\operatorname{loc}_w(s)}\ast e_h=e_{\operatorname{loc}_w(s)h}=e_h\ast e_{\operatorname{loc}_w(s)}$, and
\[
 \bigl(e_{\operatorname{loc}_w(s)}\ast e_{\operatorname{loc}_w(s')}\bigr)(g)=
 e\bigl(g\operatorname{loc}_w(\operatorname{pr}_w(g))^{-1}\bigr)\,
 \bigl(1_{I_wsI_w}\ast1_{I_ws'I_w}\bigr)(\operatorname{pr}_w(g))
\]
with the convolution relative to $I_w$ on the right. With the function $\phi$ of
Lemma~\ref{gal-test-function-iwahori}(1) at $w$, this gives a function $v_w$ with
$v_w\ast e_{\operatorname{loc}_w(t_w)}=e$, and the claims follow.

Next, $r$ is good, so $a_w\in\operatorname{GL}_n(\mathcal O_w)$ for $w\mid r$, and there is $M\geq0$ with
$r^M\tau^{-1}$ integral. Let $c_0$ be given by Lemma~\ref{gal-test-function-type}(3) and let $D$ be the
central element of $G(\mathbf A_f^p)$ acting on $\Lambda\otimes\mathbf A_f^p$ by the scalar $c_0^2r^M$. Put
$u(g)=e_\tau(Dg)$ and $u'(g)=u''(D^{-1}g)$. Translation by the central element $D$ commutes with
convolution, so (1) follows from $e\ast e_\tau=e_\tau=e_\tau\ast e$ and $u''\ast e_\tau=e$, and (2) from
the commutation of $e_\tau$ with the $e_h$.

For (3), let $u(g)\neq0$. Then $Dg=a\tau b$ with $a,b\in K$, so
$g^{-1}=(c_0b^{-1})(r^M\tau^{-1})(c_0a^{-1})$ is integral by the choice of $c_0$ and $M$, and
$\operatorname{pr}_w(D)\operatorname{pr}_w(g)\in I_wt_wI_w$ for $w\in W$. Since $u\ast e=u$, the support of
$u^{\ast m}\ast e$ consists of products of $m$ elements of the support of $u$, so for such $g$ the inverse
$g^{-1}$ is integral and $\operatorname{pr}_w(D)^m\operatorname{pr}_w(g)\in I_wt_w^mI_w$ by
Lemma~\ref{gal-test-function-iwahori}(2). Let $X_w$ be the set of $y$ with
$\operatorname{pr}_w(D)^my\in I_wt_w^mI_w$. It is stable under left multiplication by $I_w$ because
$\operatorname{pr}_w(D)$ is a scalar, and its elements have separable characteristic polynomial by
Lemma~\ref{gal-test-function-iwahori}(3), since multiplying by a nonzero scalar preserves separability.
\end{proof}
% Lean: auxFn, with locT for Lemma gal-test-function-iwahori, exists_invertible_hecke_function_at_finite_places
% (tfMulL_holds, inv_ex) for tau and u'', exists_primePow_scalar_integral_of_good_prime for M, cen_ex
% (exists_levelSubgroup_central_scalar_of_pos) for D = j(dc) j(dr) j(dc), and pConvSupp for supports.

\subsection{The spectral condition and the existence of packs}\label{subsection-gal-spectral-condition}

\begin{definition}\label{gal-hecke-type-functions}
In the setting of Definition~\ref{gal-test-function}, let $\mathcal M$ be a complex vector space. A function
$\Xi:G(\mathbf A_f^p)\to\mathcal M$ is \emph{of type} if $\Xi(gk)=\psi(\operatorname{pr}_{w_0}(k))\,\Xi(g)$ for
all $g\in G(\mathbf A_f^p)$ and $k\in K$. For a function $\Xi$ of type and a bi-$K_N$-invariant compactly
supported function $f$ on $G(\mathbf A_f^p)$, put
\[
 (f\cdot\Xi)(g)=\sum_{y\in G(\mathbf A_f^p)/K_N}f(y)\,\Xi(gy).
\]
By Lemma~\ref{gal-test-function-type}(1), $\psi\circ\operatorname{pr}_{w_0}$ is defined on $K$ and trivial on
$K_N$, so a function of type is right invariant under $K_N$ and the sum has finitely many nonzero terms.
\end{definition}
\leanref{Development/AutomorphicArithmetic/Foundations/GaloisRepresentations/SphericalCoordinates/AlgebraEquivSetup/HeckeTraceEq.lean}{UpperTheory.Gal.W2.Seed.TyLev.IsTy}
\leanref{Formalism/SatakeDualWeyl/Statement.lean}{Formalism.UG.hAct}

\begin{lemma}\label{gal-hecke-double-cosets}
Let $X$ and $M$ be groups, and $\operatorname{pr}_w:X\to M$ and $\operatorname{loc}_w:M\to X$ homomorphisms
with $\operatorname{pr}_w\circ\operatorname{loc}_w=1$, such that every $g\in X$ with
$\operatorname{pr}_w(g)=1$ commutes with every element of $\operatorname{loc}_w(M)$. Let $H\subset X$ and
$H_w\subset M$ be subgroups with $\operatorname{loc}_w(H_w)\subset H$ and
$\operatorname{pr}_w(H)\subset H_w$, let $t\in M$, and let $\Xi$ be a function on $X$ with values in a
complex vector space, right invariant under $H$. Then
\[
 \sum_{y\in X/H}1_{H\operatorname{loc}_w(t)H}(y)\,\Xi(y)
 =\sum_{x\in M/H_w}1_{H_wtH_w}(x)\,\Xi(\operatorname{loc}_w(x)),
\]
where each side is the sum of its nonzero terms if there are finitely many, and $0$ otherwise. In the
setting of Definition~\ref{gal-trace-seen-operators}, for a place $w$ of $E$ split over $F$ and prime to
$p$, the maps $\operatorname{pr}_w$ and $\operatorname{loc}_w$ satisfy these hypotheses for
$X=G(\mathbf A_f^p)$ and $M=\operatorname{GL}_n(E_w)$.
\end{lemma}
\leanref{Development/AutomorphicArithmetic/Foundations/GaloisRepresentations/SphericalCoordinates/AlgebraEquivSetup/PDosetCosets.lean}{UpperTheory.Gal.W2.Seed.pDosetCosets}
\leanref{Development/AutomorphicArithmetic/Foundations/GaloisRepresentations/SphericalCoordinates/AlgebraEquivSetup/HeckeTraceEq.lean}{UpperTheory.Gal.W2.Seed.Loc}
\leanref{Development/AutomorphicArithmetic/Foundations/GaloisRepresentations/SphericalCoordinates/AlgebraEquivSetup/PLocComm.lean}{UpperTheory.Gal.W2.Seed.pLocComm}

\begin{proof}
For the last sentence, $\operatorname{pr}_w(\operatorname{loc}_w(y))=y$ by definition. Let $g$ satisfy
$\operatorname{pr}_w(g)=1$, and write $g=j(g',\nu)$. The relation ${}^tc(g')g'=\nu$ and $g'_w=1$ make the
component of $g'$ at $c(w)$ a scalar, and $\operatorname{loc}_w(x)$ differs from $1$ only at $w$ and
$c(w)$, so $g$ commutes with $\operatorname{loc}_w(x)$.

The map $x\mapsto\operatorname{loc}_w(x)H$ from $M/H_w$ to $X/H$ is well
defined because $\operatorname{loc}_w(H_w)\subset H$. It is injective: if
$\operatorname{loc}_w(x^{-1}x')\in H$, then $x^{-1}x'=\operatorname{pr}_w(\operatorname{loc}_w(x^{-1}x'))\in H_w$.
Applying $\operatorname{pr}_w$ and $\operatorname{loc}_w$ shows that
$\operatorname{loc}_w(x)\in H\operatorname{loc}_w(t)H$ if and only if $x\in H_wtH_w$, so the term of $x$ on
the right is the term of $\operatorname{loc}_w(x)H$ on the left. It remains to show that every coset contained
in $H\operatorname{loc}_w(t)H$ is of the form $\operatorname{loc}_w(x)H$. Let
$y=a\operatorname{loc}_w(t)b$ with $a,b\in H$, and put $a'=a\operatorname{loc}_w(\operatorname{pr}_w(a))^{-1}$.
Then $a'\in H$ and $\operatorname{pr}_w(a')=1$, so $a'$ commutes with
$\operatorname{loc}_w(\operatorname{pr}_w(a)t)$, and
\[
 y=a'\operatorname{loc}_w(\operatorname{pr}_w(a))\operatorname{loc}_w(t)\,b
 =\operatorname{loc}_w(\operatorname{pr}_w(a)t)\,a'b\in\operatorname{loc}_w(\operatorname{pr}_w(a)t)H.
\]
So the map is a bijection from the cosets with a nonzero term on the right onto those with a nonzero term
on the left, matching the terms. Hence both sides have finitely many nonzero terms or neither does, and
the two sides are equal.
\end{proof}

\begin{definition}\label{gal-hecke-test-localization}
Let $X$ and $M$ be groups. A \emph{localization} of $X$ at $M$ is a pair of homomorphisms
$\operatorname{pr}:X\to M$ and $\operatorname{loc}:M\to X$ with $\operatorname{pr}\circ\operatorname{loc}=1$, such
that every $g\in X$ with $\operatorname{pr}(g)=1$ commutes with every element of $\operatorname{loc}(M)$. A
localization $(\operatorname{pr},\operatorname{loc})$ of $X$ at $M$ and a localization
$(\operatorname{pr}',\operatorname{loc}')$ of $X$ at a group $M'$ are \emph{apart} if
$\operatorname{pr}\circ\operatorname{loc}'=1$ and $\operatorname{pr}'\circ\operatorname{loc}=1$.
\end{definition}
\leanref{Development/AutomorphicArithmetic/Foundations/GaloisRepresentations/SphericalCoordinates/AlgebraEquivSetup/HeckeTraceEq.lean}{UpperTheory.Gal.W2.Seed.Loc}
\leanref{Development/AutomorphicArithmetic/Foundations/GaloisRepresentations/SphericalCoordinates/AlgebraEquivSetup/HeckeTraceEq.lean}{UpperTheory.Gal.W2.Seed.Loc.Apart}

\begin{lemma}\label{gal-hecke-test-action}
In the setting of Definition~\ref{gal-test-function}, let $h\in G(\mathbf A_f^p)$ with
$\operatorname{pr}_{w_0}(h)=1$, and let $\Xi$ be a function of type
(Definition~\ref{gal-hecke-type-functions}). Then:
\begin{enumerate}
\item $e_h$ is bi-$K_N$-invariant with compact support, and
  $e_h(kg)=e_h(gk)=\psi(\operatorname{pr}_{w_0}(k))^{-1}e_h(g)$ for $k\in K$ and $g\in G(\mathbf A_f^p)$;
\item $e_h\ast e=e_h$, for the convolution $\ast$ of Lemma~\ref{gal-test-function-powers};
\item if $d\in G(\mathbf A_f^p)$ is central with $\operatorname{pr}_{w_0}(d)=1$, then
  $(e_d\cdot\Xi)(g)=\Xi(gd)$ for every $g$;
\item if $d\in G(\mathbf A_f^p)$ is central with $\operatorname{pr}_{w_0}(d)=1$,
  $(\operatorname{pr}',\operatorname{loc}')$ is a localization of $G(\mathbf A_f^p)$ at a group $M'$ that is
  apart from $(\operatorname{pr}_{w_0},\operatorname{loc}_{w_0})$
  (Definition~\ref{gal-hecke-test-localization}), $K_w\subset M'$ is a subgroup with
  $\operatorname{loc}'(K_w)\subset K$ and $\operatorname{pr}'(K)\subset K_w$, and $t\in M'$, then for every $g$
  \[
   (e_{\operatorname{loc}'(t)d}\cdot\Xi)(g)=\sum_{x\in M'/K_w}1_{K_wtK_w}(x)\,
   \Xi(gd\operatorname{loc}'(x)).
  \]
  The hypothesis $\operatorname{pr}'\circ\operatorname{loc}_{w_0}=1$ in the apartness is not used.
\end{enumerate}
\end{lemma}
\leanref{Development/AutomorphicArithmetic/Foundations/GaloisRepresentations/SphericalCoordinates/AlgebraEquivSetup/PTfAdm.lean}{UpperTheory.Gal.W2.Seed.pTfAdm}
\leanref{Development/AutomorphicArithmetic/Foundations/GaloisRepresentations/SphericalCoordinates/AlgebraEquivSetup/SeenShapeHolds.lean}{UpperTheory.Gal.W2.Seed.tfIdem_holds}
\leanref{Development/AutomorphicArithmetic/Foundations/GaloisRepresentations/SphericalCoordinates/AlgebraEquivSetup/SeenShapeHolds.lean}{UpperTheory.Gal.W2.Seed.tfCen_holds}
\leanref{Development/AutomorphicArithmetic/Foundations/GaloisRepresentations/SphericalCoordinates/AlgebraEquivSetup/SeenShapeHolds.lean}{UpperTheory.Gal.W2.Seed.tfAct_holds}

\begin{proof}
(1) Let $k\in K$. Then $gk\in KhK$ if and only if $g\in KhK$. By Lemma~\ref{gal-test-function-type}(1),
$\operatorname{pr}_{w_0}(k)\in J'$, so $\operatorname{pr}_{w_0}(gk)\in J'$ if and only if
$\operatorname{pr}_{w_0}(g)\in J'$, and then
$\psi(\operatorname{pr}_{w_0}(gk))^{-1}=\psi(\operatorname{pr}_{w_0}(g))^{-1}\psi(\operatorname{pr}_{w_0}(k))^{-1}$.
This gives the identity for $gk$, and the identity for $kg$ follows in the same way. For $k\in K_N$ the
factor $\psi(\operatorname{pr}_{w_0}(k))$ is $1$ by Lemma~\ref{gal-test-function-type}(1), so $e_h$ is
bi-$K_N$-invariant. It vanishes off the compact set $KhK$.

(2) Fix $g$. Replacing $x$ by $gx$ permutes $G(\mathbf A_f^p)/K_N$, so $(e_h\ast e)(g)$ is the sum of
$e_h(gx)\,e(x^{-1})$ over $x\in G(\mathbf A_f^p)/K_N$. The factor $e(x^{-1})$ vanishes unless $x\in K$, which
leaves the $[K:K_N]$ cosets contained in $K$. For $x\in K$, (1) gives
$e_h(gx)=\psi(\operatorname{pr}_{w_0}(x))^{-1}e_h(g)$, and $e(x^{-1})=[K:K_N]^{-1}\psi(\operatorname{pr}_{w_0}(x))$.
So each of these terms is $[K:K_N]^{-1}e_h(g)$, and their sum is $e_h(g)$.

(3) By (1) and since $\Xi$ is of type, the function $x\mapsto e_d(x)\,\Xi(gx)$ is right invariant under $K$:
the factors $\psi(\operatorname{pr}_{w_0}(k))^{-1}$ and $\psi(\operatorname{pr}_{w_0}(k))$ cancel. So the sum
over $G(\mathbf A_f^p)/K_N$ defining $(e_d\cdot\Xi)(g)$ is $[K:K_N]$ times the sum of the same terms over
$G(\mathbf A_f^p)/K$. Since $d$ is central, $KdK=dK$, so only the coset $dK$ contributes. Its term is
$e_d(d)\,\Xi(gd)=[K:K_N]^{-1}\Xi(gd)$, as $\operatorname{pr}_{w_0}(d)=1$. Hence $(e_d\cdot\Xi)(g)=\Xi(gd)$.

(4) Put $h'=\operatorname{loc}'(t)d$. Since the two localizations are apart,
$\operatorname{pr}_{w_0}(\operatorname{loc}'(x))=1$ for every $x\in M'$, so
$\operatorname{pr}_{w_0}(h')=1$. As in (3),
\[
 (e_{h'}\cdot\Xi)(g)=[K:K_N]\sum_{y\in G(\mathbf A_f^p)/K}e_{h'}(y)\,\Xi(gy)
 =\sum_{y\in G(\mathbf A_f^p)/K}1_{Kh'K}(y)\,\psi(\operatorname{pr}_{w_0}(y))^{-1}\,\Xi(gy).
\]
Replace $y$ by $dy$, which permutes $G(\mathbf A_f^p)/K$. As $d$ is central with
$\operatorname{pr}_{w_0}(d)=1$, $1_{Kh'K}(dy)=1_{K\operatorname{loc}'(t)K}(y)$ and
$\operatorname{pr}_{w_0}(dy)=\operatorname{pr}_{w_0}(y)$. So the sum is
$\sum_y1_{K\operatorname{loc}'(t)K}(y)\,\Xi'(y)$, where
$\Xi'(y)=\psi(\operatorname{pr}_{w_0}(y))^{-1}\,\Xi(gdy)$ if $\operatorname{pr}_{w_0}(y)\in J'$ and
$\Xi'(y)=0$ otherwise. The function $\Xi'$ is right invariant under $K$, by
Lemma~\ref{gal-test-function-type}(1) and since $\Xi$ is of type.
Lemma~\ref{gal-hecke-double-cosets}, for $M=M'$, $\operatorname{pr}_w=\operatorname{pr}'$,
$\operatorname{loc}_w=\operatorname{loc}'$, $H=K$ and $H_w=K_w$, turns the sum into
$\sum_x1_{K_wtK_w}(x)\,\Xi'(\operatorname{loc}'(x))$, and
$\Xi'(\operatorname{loc}'(x))=\Xi(gd\operatorname{loc}'(x))$ because
$\operatorname{pr}_{w_0}(\operatorname{loc}'(x))=1$.
\end{proof}
% Lean: pTfAct, with pAvgK for the passage from K_N to K, finsum_out_mul for the translation by d, pDosetCosets
% for Lemma gal-hecke-double-cosets, and hAp.1 (pr_{w0} o loc' = 1) for the last step; hAp.2 is not used.

\begin{lemma}\label{gal-hecke-convolution}
In the setting of Definition~\ref{gal-test-function}, let $f$ be a bi-$K_N$-invariant compactly supported
function on $G(\mathbf A_f^p)$, let $\alpha\in\mathbf C$, and let $h\in G(\mathbf A_f^p)$ with
$\operatorname{pr}_{w_0}(h)=1$. Suppose that $f\cdot\Xi=\alpha\,(e_h\cdot\Xi)$ for every function $\Xi$ of type
(Definition~\ref{gal-hecke-type-functions}). Then $f\ast e=\alpha\,e_h$, for the convolution $\ast$ of
Lemma~\ref{gal-test-function-powers}.
\end{lemma}
\leanref{Development/AutomorphicArithmetic/Foundations/GaloisRepresentations/SphericalCoordinates/AlgebraEquivSetup/SeenShapeHolds.lean}{UpperTheory.Gal.W2.Seed.actsConv_holds}

\begin{proof}
Fix $g\in G(\mathbf A_f^p)$ and let $\Xi(x)=e(x^{-1}g)$. For $k\in K$, Lemma~\ref{gal-hecke-test-action}(1)
for $h=1$ gives
$\Xi(xk)=e(k^{-1}x^{-1}g)=\psi(\operatorname{pr}_{w_0}(k^{-1}))^{-1}\,\Xi(x)=\psi(\operatorname{pr}_{w_0}(k))\,\Xi(x)$,
so $\Xi$ is of type, with values in $\mathbf C$. Hence
\[
 (f\ast e)(g)=(f\cdot\Xi)(1)=\alpha\,(e_h\cdot\Xi)(1)=\alpha\,(e_h\ast e)(g)=\alpha\,e_h(g),
\]
the last equality by Lemma~\ref{gal-hecke-test-action}(2).
\end{proof}

\begin{lemma}\label{gal-hecke-seen-shape}
In the setting of Definition~\ref{gal-test-function}, let $g=(K',h,M)$ be a Hecke operator of level $N$ seen
on $U$ with value $\vartheta$ and auxiliary prime $r$ (Definition~\ref{gal-trace-seen-operators}). Then every
$k\in K'$ lies in $K$ and satisfies $\psi(\operatorname{pr}_{w_0}(k))=1$. Moreover:
\begin{enumerate}
\item if $g$ is of kind (2)(a), then $h\in K$ and $\psi(\operatorname{pr}_{w_0}(h))=1$;
\item if $g$ is of kind (2)(b), with the place $w$, the integer $i$, the matrix $t_i$ and the element $z$
  of Definition~\ref{gal-trace-seen-operators}(2)(b), then $w$ is prime to $p$ and lies over a place of $F$
  different from the one under $w_0$ and from those under the places of $W$; $z$ is central in
  $\mathrm{GU}_n(\mathbf A_f)$ with component $1$ at every place over $p$; $j(z)$ is central in
  $G(\mathbf A_f^p)$ with $\operatorname{pr}_{w_0}(j(z))=1$; $h^{-1}=\operatorname{loc}_w(t_i^{-1})\,j(z)^{-i}$,
  and $\operatorname{pr}_u(h)=1$ for $u\in W$; and $\operatorname{pr}_w(K')\subset\operatorname{GL}_n(\mathcal O_w)$
  and $\operatorname{loc}_w(\operatorname{GL}_n(\mathcal O_w))\subset K'$.
\end{enumerate}
\end{lemma}
\leanref{Development/AutomorphicArithmetic/Foundations/GaloisRepresentations/SphericalCoordinates/AlgebraEquivSetup/SeenShapeHolds.lean}{UpperTheory.Gal.W2.Seed.seenShape_holds}
\leanref{Development/AutomorphicArithmetic/Foundations/GaloisRepresentations/SphericalCoordinates/AlgebraEquivSetup/SeenShapeHolds.lean}{UpperTheory.Gal.W2.Seed.stabKer_holds}
\leanref{Development/AutomorphicArithmetic/Foundations/GaloisRepresentations/SphericalCoordinates/AlgebraEquivSetup/SeenShapeHolds.lean}{UpperTheory.Gal.W2.Seed.siteFact_holds}

\begin{proof}
First, every element $k$ that stabilizes $V$ and lies in $K_r$ lies in $K$ and satisfies
$\psi(\operatorname{pr}_{w_0}(k))=1$. Indeed, $k=j(k')$ with $\theta(k')\in K_S$,
$\theta(k')_{w_0}\in J$ and $\chi(\theta(k')_{w_0})=1$. For $u\in W$, $u$ lies over $r$ and $r$ is good, so
$a_u\in\operatorname{GL}_n(\mathcal O_u)$ by Lemma~\ref{gal-identification}(3), and
$\operatorname{pr}_u(k)=a_u^{-1}k_ua_u$ is congruent to $1$ modulo $r$, hence modulo $\varpi_u$, and lies in
$I_u$. So $k\in K$, and $\psi(\operatorname{pr}_{w_0}(k))=\chi(\theta(k')_{w_0})=1$ by the computation at the
start of the proof of Lemma~\ref{gal-test-function-type}.

(1) Let $g$ be of kind (2)(a). Then $K'=K_N$, whose elements stabilize $V$ by the hypotheses of
Definition~\ref{gal-test-function}, and $K_N\subset K_r$ since $r\mid N$. The element $h$ stabilizes $V$ and
lies in $K_r$. So all the claims follow from the first step.

(2) Let $g$ be of kind (2)(b). By the definition of a level at $q$ adapted to $w$, $q\neq p$, and $K'$ is
the group of elements of $K_1$ that act trivially on $\hat\Lambda^p/\mathfrak b\hat\Lambda^p$, where
$\mathfrak b=N'\mathfrak a$, $N'$ is the prime-to-$q$ part of $N$, and $\mathfrak a$ is an ideal of
$\mathcal O_E$ prime to $w$ that contains a power of the ideal $(q,\omega-\rho)$, for an integer $\rho$ with
$q\mid\rho^2-t\rho+m$ and $\omega-\rho\in w$. Since $r\neq q$ and $r\mid N$, $r$ divides $N'$, so
$K'\subset K_r$. The elements of $K'$ stabilize $V$, so the first step gives the first claim of the lemma.

As $q\neq p$ and $q\neq r$, the place $w$ is prime to $p$ and lies over a place of $F$ different from those
under the places of $W$, which lie over $r$. Since $w$ is free for $V$, $w$ and $c(w)$ are not in $S$, while
$w_0\in S$; so $w$ is neither $w_0$ nor $c(w_0)$, and lies over a place of $F$ different from the one under
$w_0$.

The component $z_u$ is the scalar $N(w)$ if $c(u)\sim w$ and $1$ otherwise, so $z$ is central. If
$c(u)\sim w$, then $c(u)$ lies over $q$, and so does $u$. Hence $z_u=1$ at every place $u$ not over $q$, in
particular at the places over $p$ and at the places of $W$. Moreover $z_{w_0}=1$, since $c(w_0)\not\sim w$, $w$
being free for $V$ and $w_0\in S$. As $j$ is surjective, $j(z)$ is central, and
$\operatorname{pr}_{w_0}(j(z))=z_{w_0}=1$.

Since $z$ is central, $h=\operatorname{loc}_w(t_i)\,j(z)^i$ gives
$h^{-1}=\operatorname{loc}_w(t_i^{-1})\,j(z)^{-i}$. For $u\in W$,
$\operatorname{pr}_u(\operatorname{loc}_w(t_i))=1$ because $u$ is neither $w$ nor $c(w)$, and
$\operatorname{pr}_u(j(z))=z_u=1$; so $\operatorname{pr}_u(h)=1$.

An element $k$ of $K'\subset K_1$ preserves $\mathcal O_w^n$, and $a_w\in\operatorname{GL}_n(\mathcal O_w)$ by
Lemma~\ref{gal-identification}(3) since $q$ is good; so
$\operatorname{pr}_w(k)=a_w^{-1}k_wa_w\in\operatorname{GL}_n(\mathcal O_w)$. Conversely, let
$y\in\operatorname{GL}_n(\mathcal O_w)$. The element $\iota_w(y)$ has component $y$ at $w$, the integral
component $c({}^ty^{-1})$ at $c(w)$, and component $1$ elsewhere. As $a_w\in\operatorname{GL}_n(\mathcal O_w)$
and $a_{c(w)}\in\operatorname{GL}_n(\mathcal O_{c(w)})$, $\operatorname{loc}_w(y)$ preserves $\hat\Lambda^p$ and
lies in $K_1$. The ideal $\mathfrak b$ is prime to $w$ and to $c(w)$. Indeed, $N'$ is prime to $q$ and
$\mathfrak a$ is prime to $w$; and $\omega-\rho\notin c(w)$, since otherwise $c(\omega)-\rho\in w$, so that
$t-2\rho\in w$, $q$ divides $t-2\rho$, and then $q\mid\rho^2-t\rho+m$ gives $q\mid t^2-4m$, contrary to
$q$ being split for $\omega$. So at $w$ and $c(w)$ the congruence modulo $\mathfrak b$ only asks for
integrality, and at the other places it holds because $\operatorname{loc}_w(y)$ is $1$ there. Hence
$\operatorname{loc}_w(y)\in K'$.
\end{proof}

\begin{lemma}\label{gal-hecke-seen-action}
In the setting of Definition~\ref{gal-test-function}, let $g=(K',h,M)$ be a Hecke operator of level $N$ seen
on $U$ with value $\vartheta$ and auxiliary prime $r$ (Definition~\ref{gal-trace-seen-operators}), and let
$\Xi$ be a function of type (Definition~\ref{gal-hecke-type-functions}). Then:
\begin{enumerate}
\item $f_g$ is bi-$K_N$-invariant with compact support, and for every $y$ with $f_g(y)\neq0$ there is
  $k\in K'$ with $hky\in K_N$;
\item if $g$ is of kind (2)(a), then $f_g\cdot\Xi=\Xi=e\cdot\Xi$;
\item if $g$ is of kind (2)(b), with the place $w$, the integer $i$, the matrix $t_i$ and the element $z$ of
  Definition~\ref{gal-trace-seen-operators}(2)(b), then $\operatorname{pr}_{w_0}(h)=1$,
  $f_g\cdot\Xi=\operatorname{idx}(g)\,(e_{h^{-1}}\cdot\Xi)$, and for every $x\in G(\mathbf A_f^p)$
  \[
   (e_{h^{-1}}\cdot\Xi)(x)=\sum_{y\in\operatorname{GL}_n(E_w)/\operatorname{GL}_n(\mathcal O_w)}
   1_{\operatorname{GL}_n(\mathcal O_w)t_i^{-1}\operatorname{GL}_n(\mathcal O_w)}(y)\,
   \Xi\bigl(x\,j(z)^{-i}\operatorname{loc}_w(y)\bigr).
  \]
\end{enumerate}
\end{lemma}
\leanref{Development/AutomorphicArithmetic/Foundations/GaloisRepresentations/SphericalCoordinates/AlgebraEquivSetup/SeenShapeHolds.lean}{UpperTheory.Gal.W2.Seed.seenActs_holds}
\leanref{Development/AutomorphicArithmetic/Foundations/GaloisRepresentations/SphericalCoordinates/AlgebraEquivSetup/PMemSum.lean}{UpperTheory.Gal.W2.Seed.pMemSum}
\leanref{Development/AutomorphicArithmetic/Foundations/GaloisRepresentations/SphericalCoordinates/AlgebraEquivSetup/SeenShapeHolds.lean}{UpperTheory.Gal.W2.Seed.memDeck_holds}
\leanref{Development/AutomorphicArithmetic/Foundations/GaloisRepresentations/SphericalCoordinates/AlgebraEquivSetup/SeenShapeHolds.lean}{UpperTheory.Gal.W2.Seed.memAct_holds}
\leanref{Development/AutomorphicArithmetic/Foundations/GaloisRepresentations/SphericalCoordinates/AlgebraEquivSetup/PSumDoset.lean}{UpperTheory.Gal.W2.Seed.pSumDoset}

\begin{proof}
The group $K_1$ preserves $\hat\Lambda^p$, and $K_M$ is the kernel of its action on
$\hat\Lambda^p/M\hat\Lambda^p$, so $K_M$ is a normal subgroup of $K_1$ of finite index. As
$K_M\subset K_N\subset K'\subset K_1$, $K_M$ is a normal subgroup of $K'$ of finite index, and $K'$ is compact
open. By Lemma~\ref{gal-hecke-seen-shape}, $\psi(\operatorname{pr}_{w_0}(k))=1$ for $k\in K'$, so $\Xi$ is right
invariant under $K'$.

(1) By the definition of a Hecke operator, $hkK_M(hk)^{-1}\subset K_N$ for $k\in K'$, so whether $hky\in K_N$
depends only on the class of $k$ in $K'/K_M$, and $f_g(y)$ is the number of classes $k\in K'/K_M$ with
$hky\in K_N$. This number is unchanged when $y$ is multiplied on the right by an element of $K_N$. For
$a\in K_N\subset K'$, $k\mapsto ka$ permutes $K'/K_M$, as $K_M$ is normal in $K'$, so $f_g(ay)=f_g(y)$. If
$f_g(y)\neq0$, then $hky\in K_N$ for some $k\in K'$, so $y$ lies in one of the finitely many cosets
$(hk)^{-1}K_N$, and $f_g$ has compact support. Moreover, each class $k$ contributes to the sum defining
$(f_g\cdot\Xi)(x)$ through the single coset $(hk)^{-1}K_N$, so
\[
 (f_g\cdot\Xi)(x)=\sum_{k\in K'/K_M}\Xi\bigl(x(hk)^{-1}\bigr).\tag{$\ast$}
\]

(2) If $g$ is of kind (2)(a), then $K'=K_M=K_N$, so $(\ast)$ has the single term $\Xi(xh^{-1})$. This is
$\Xi(x)$, because $h\in K$ and $\psi(\operatorname{pr}_{w_0}(h))=1$ by Lemma~\ref{gal-hecke-seen-shape}(1).
And $e\cdot\Xi=\Xi$ by Lemma~\ref{gal-hecke-test-action}(3) for $d=1$.

(3) Let $g$ be of kind (2)(b), and put $t=t_i^{-1}$ and $d=j(z)^{-i}$. By
Lemma~\ref{gal-hecke-seen-shape}(2), $h^{-1}=\operatorname{loc}_w(t)d$ with $d$ central and
$\operatorname{pr}_{w_0}(d)=1$, and $\operatorname{pr}_{w_0}(\operatorname{loc}_w(t))=1$ because $w$ and $w_0$
lie over different places of $F$; so $\operatorname{pr}_{w_0}(h)=1$.

Since $\Xi$ is right invariant under $K'$, the term of $k$ in $(\ast)$ depends only on the coset
$k^{-1}h^{-1}K'$, and these cosets are the cosets of $K'$ contained in $K'h^{-1}K'$. Two classes $k,k'$ give
the same coset if and only if $hk'k^{-1}h^{-1}\in K'$, that is, $k'k^{-1}\in K'\cap h^{-1}K'h$. Since
$hK_Mh^{-1}\subset K_N\subset K'$, the group $K_M$ is contained in $K'\cap h^{-1}K'h$, so each coset comes from
$[K'\cap h^{-1}K'h:K_M]=\operatorname{idx}(g)$ classes. Hence
\[
 (f_g\cdot\Xi)(x)=\operatorname{idx}(g)\sum_{y\in G(\mathbf A_f^p)/K'}1_{K'h^{-1}K'}(y)\,\Xi(xy).
\]
Replacing $y$ by $dy$, which permutes $G(\mathbf A_f^p)/K'$ since $d$ is central, turns $1_{K'h^{-1}K'}(dy)$
into $1_{K'\operatorname{loc}_w(t)K'}(y)$. By Lemma~\ref{gal-hecke-seen-shape}(2), the hypotheses of
Lemma~\ref{gal-hecke-double-cosets} hold for $H=K'$ and $H_w=\operatorname{GL}_n(\mathcal O_w)$. So
$(f_g\cdot\Xi)(x)$ is $\operatorname{idx}(g)$ times the right side of the display in (3).

The place $w$ is split over $F$, prime to $p$, with $w,c(w)\notin S$ as it is free for $V$, and lies over a
place of $F$ different from those under $W$ and from the one under $w_0$ (Lemma~\ref{gal-hecke-seen-shape}(2)).
So Lemma~\ref{gal-test-function-type}(2) gives $\operatorname{pr}_w(K)\subset\operatorname{GL}_n(\mathcal O_w)$
and $\operatorname{loc}_w(\operatorname{GL}_n(\mathcal O_w))\subset K$. By Lemma~\ref{gal-hecke-double-cosets},
$(\operatorname{pr}_w,\operatorname{loc}_w)$ and $(\operatorname{pr}_{w_0},\operatorname{loc}_{w_0})$ are
localizations of $G(\mathbf A_f^p)$ (Definition~\ref{gal-hecke-test-localization}). They are apart: since $w$
and $w_0$ lie over different places of $F$, $\operatorname{loc}_w(y)$ has component $1$ at $w_0$ and
$\operatorname{loc}_{w_0}(y')$ has component $1$ at $w$. So by Lemma~\ref{gal-hecke-test-action}(4), with
$M'=\operatorname{GL}_n(E_w)$ and $K_w=\operatorname{GL}_n(\mathcal O_w)$, the right side of the display in (3) is
$(e_{\operatorname{loc}_w(t)d}\cdot\Xi)(x)=(e_{h^{-1}}\cdot\Xi)(x)$.
\end{proof}
% Lean: the apartness is AtApart, from GlobalTheory.Cong.Apart w w0.

\begin{lemma}\label{gal-hecke-satake-eigenvalue}
In the setting of Lemma~\ref{gal-identification}, let $x\in E$, let $V$ be a vector for an automorphic
representation $U$ of $U_n$ at a place $w_0$ (Definition~\ref{gal-trace-vector}, with the relation $\sim$
defined by $x$), with function $\Phi$ and Hecke character $\kappa$, and let $\theta$ be the automorphism of
Lemma~\ref{gal-identification}(4). Let $w$ be a place free for $V$, let $0\leq i\leq n$, let $\varpi_w$ be a
uniformizer of $E_w$ and
$t_i=\operatorname{diag}(\varpi_w,\ldots,\varpi_w,1,\ldots,1)$ with $i$ entries equal to $\varpi_w$, let
$z\in\mathrm{GU}_n(\mathbf A_f)$ have component the scalar $N(w)$ at the places $u$ with $c(u)\sim w$ and $1$
at the other places, and suppose that $U$ has Satake parameters $B$ at $w$. For
$y\in\operatorname{GL}_n(E_w)$ let $\iota_w(y)\in\mathrm{GU}_n(\mathbf A_f)$ have similitude factor $1$,
component $y$ at $w$, and component $1$ away from $w$ and $c(w)$. Then for every
$g\in\mathrm{GU}_n(\mathbf A_f)$
\[
 \sum_{y\in\operatorname{GL}_n(E_w)/\operatorname{GL}_n(\mathcal O_w)}
 1_{\operatorname{GL}_n(\mathcal O_w)t_i^{-1}\operatorname{GL}_n(\mathcal O_w)}(y)\,
 \Phi\bigl(\theta(gz^{-i}\iota_w(y))\bigr)
 =q_w^{i(n-i)/2}e_i(B)\,\kappa(\varpi_w)^{-i}\,\Phi(\theta(g)).
\]
\end{lemma}
\leanref{Development/AutomorphicArithmetic/Foundations/GaloisRepresentations/SphericalCoordinates/AlgebraEquivSetup/FinsumInvT.lean}{UpperTheory.Gal.W2.Seed.tEig_of_thLoc}

\begin{proof}
Since $w$ and $c(w)$ are not in $S$, Definition~\ref{gal-trace-vector}(4) makes $\Phi$ a spherical vector with
Satake parameters $B$ for right translation through $\iota_w$. So $\Phi$ is invariant under right
translation by $\iota_w(\operatorname{GL}_n(\mathcal O_w))$, and the $i$-th Hecke operator, the characteristic
function of $\operatorname{GL}_n(\mathcal O_w)t_i\operatorname{GL}_n(\mathcal O_w)$, acts on it by
$q_w^{i(n-i)/2}e_i(B)$ (Definition~\ref{gal-definite-unitary}(3)):
\[
 \sum_{y\in\operatorname{GL}_n(E_w)/\operatorname{GL}_n(\mathcal O_w)}
 1_{\operatorname{GL}_n(\mathcal O_w)t_i\operatorname{GL}_n(\mathcal O_w)}(y)\,\Phi(g'\iota_w(y))
 =q_w^{i(n-i)/2}e_i(B)\,\Phi(g')\tag{$\ast\ast$}
\]
for every $g'$. Since $w$ is free for $V$, Definition~\ref{gal-trace-vector}(5) gives
$\Phi(g'z)=\kappa(\varpi_w)^{-1}\Phi(g')$, hence $\Phi(g'z^i)=\kappa(\varpi_w)^{-i}\Phi(g')$.

The automorphism $\theta$ is the inverse transpose on the components and the inverse on the similitude
factor. So $\theta(\iota_w(y))=\iota_w({}^ty^{-1})$, and $\theta(z)=z^{-1}$ because the components of $z$ are
scalars. The term of $y$ on the left is therefore
$1_{\operatorname{GL}_n(\mathcal O_w)t_i^{-1}\operatorname{GL}_n(\mathcal O_w)}(y)\,
\Phi(\theta(g)z^i\iota_w({}^ty^{-1}))$. The map $y\mapsto{}^ty^{-1}$ preserves
$\operatorname{GL}_n(\mathcal O_w)$, so it induces a bijection of
$\operatorname{GL}_n(E_w)/\operatorname{GL}_n(\mathcal O_w)$; and since ${}^tt_i=t_i$, it carries
$\operatorname{GL}_n(\mathcal O_w)t_i^{-1}\operatorname{GL}_n(\mathcal O_w)$ onto
$\operatorname{GL}_n(\mathcal O_w)t_i\operatorname{GL}_n(\mathcal O_w)$. So the left side is the left side of
$(\ast\ast)$ for $g'=\theta(g)z^i$, which is
$q_w^{i(n-i)/2}e_i(B)\,\Phi(\theta(g)z^i)=q_w^{i(n-i)/2}e_i(B)\,\kappa(\varpi_w)^{-i}\Phi(\theta(g))$.
\end{proof}
% Lean: tEig_of_thLoc uses ThLoc (pThLoc_of: th (loc y) = loc (invT y), th z = z^{-1}), IsForm.sat
% (IsSatVecGL with Tj and satEig), IsKap.cen, and finsum_invT, dos_invT for the substitution y -> invT y.

\begin{lemma}\label{gal-test-function-hecke}
In the setting of Definition~\ref{gal-test-function}, suppose that $w_0$ lies over a good prime, that no place of $S$
lies over $p$, and that $\Phi(g_0)\neq0$ for some $g_0$ with $g_{0,w}\in\operatorname{GL}_n(\mathcal O_w)$ for every
$w\mid p$. Let $\phi(y)=\Phi(\theta(s(y)))$ for
$y\in G(\mathbf A_f^p)$, with $s$ the section of Lemma~\ref{gal-identification}(2), and let $g=(K',h,M)$
and $g'$ be Hecke operators of level $N$ seen on $U$ with values $\vartheta$ and $\vartheta'$ and
auxiliary prime $r$ (Definition~\ref{gal-trace-seen-operators}). Then:
\begin{enumerate}
\item $f_g$ is bi-$K_N$-invariant with compact support, and $\operatorname{pr}_w(y)\in I_w$ for $w\in W$
  whenever $f_g(y)\neq0$;
\item either $g$ is of the kind (2)(a) of Definition~\ref{gal-trace-seen-operators} and $f_g\ast e=e$, or
  $\operatorname{pr}_{w_0}(h)=1$, $\operatorname{pr}_w(h)=1$ for $w\in W$, and
  $f_g\ast e=\operatorname{idx}(g)\,e_{h^{-1}}$;
\item $f_g\ast(f_{g'}\ast e)=f_{g'}\ast(f_g\ast e)$;
\item $\sum_{y\in G(\mathbf A_f^p)/K_N}f_g(y)\,\phi(xy)=\vartheta\,\phi(x)$ for every $x$.
\end{enumerate}
The hypothesis on $w_0$ is not used, and (1) to (3) do not use the hypotheses on $S$ and on $\Phi(g_0)$.
\end{lemma}
\leanref{Development/AutomorphicArithmetic/Foundations/GaloisRepresentations/ScalarCommutation/MemFn.lean}{UpperTheory.Gal.Spc.memFn}
\leanref{Development/AutomorphicArithmetic/Foundations/GaloisRepresentations/ScalarCommutation/MemComm.lean}{UpperTheory.Gal.Spc.memComm}
\leanref{Development/AutomorphicArithmetic/Foundations/GaloisRepresentations/ScalarCommutation/PhiEx.lean}{UpperTheory.Gal.Spc.phiEx}

\begin{proof}
(1) By Lemma~\ref{gal-hecke-seen-action}(1), $f_g$ is bi-$K_N$-invariant with compact support. Let
$f_g(y)\neq0$, and take $k\in K'$ with $hky\in K_N$. By Lemma~\ref{gal-hecke-seen-shape} and
Lemma~\ref{gal-test-function-type}(1), $k\in K$ and $hky\in K$. Let $u\in W$. If $g$ is of kind (2)(a), then
$h\in K$ by Lemma~\ref{gal-hecke-seen-shape}(1), so $y\in K$, and $\operatorname{pr}_u(y)\in I_u$ by
Lemma~\ref{gal-test-function-type}(2). If $g$ is of kind (2)(b), then $\operatorname{pr}_u(h)=1$ by
Lemma~\ref{gal-hecke-seen-shape}(2), so $\operatorname{pr}_u(k)\operatorname{pr}_u(y)=\operatorname{pr}_u(hky)$
lies in $I_u$; as $\operatorname{pr}_u(k)\in I_u$ by Lemma~\ref{gal-test-function-type}(2), so does
$\operatorname{pr}_u(y)$.

(2) If $g$ is of kind (2)(a), then $f_g\cdot\Xi=e\cdot\Xi$ for every function $\Xi$ of type
(Definition~\ref{gal-hecke-type-functions}) by Lemma~\ref{gal-hecke-seen-action}(2), and
Lemma~\ref{gal-hecke-convolution}, applied to $\alpha=1$ and to the element $1$, gives $f_g\ast e=e$.
Otherwise $g$ is of kind
(2)(b). Then $\operatorname{pr}_w(h)=1$ for $w\in W$ by Lemma~\ref{gal-hecke-seen-shape}(2), and by
Lemma~\ref{gal-hecke-seen-action}(3), $\operatorname{pr}_{w_0}(h)=1$ and
$f_g\cdot\Xi=\operatorname{idx}(g)\,(e_{h^{-1}}\cdot\Xi)$ for every function $\Xi$ of type.
Lemma~\ref{gal-hecke-convolution}, applied to $\alpha=\operatorname{idx}(g)$ and to $h^{-1}$, gives
$f_g\ast e=\operatorname{idx}(g)\,e_{h^{-1}}$.

For (3), by (2) there are numbers $\alpha,\alpha'$ and elements $h_1,h_1'$, each equal to $1$ or to the
inverse $h^{-1}$ of the element of an operator of kind (2)(b), with $f_g\ast e=\alpha e_{h_1}$ and
$f_{g'}\ast e=\alpha'e_{h'_1}$. Since $e\ast e_{h'_1}=e_{h'_1}$,
\[
 f_g\ast(f_{g'}\ast e)=\alpha'(f_g\ast e)\ast e_{h'_1}=\alpha\alpha'\,e_{h_1}\ast e_{h'_1},
\]
and in the same way $f_{g'}\ast(f_g\ast e)=\alpha\alpha'\,e_{h'_1}\ast e_{h_1}$. An element of kind (2)(b)
is $\operatorname{loc}_w(t)$ times a central element, at a place $w$ free for $V$, over a prime different
from $r$, and the place $w$ satisfies the hypotheses of Lemma~\ref{gal-test-function-type}(2). So
$e_{h_1}$ and $e_{h'_1}$ commute: at places over distinct places of $F$ the factors commute, and at places
over the same place of $F$ the local functions are bi-$\operatorname{GL}_n(\mathcal O_w)$-invariant, and the
spherical Hecke algebra of $\operatorname{GL}_n(E_w)$ is commutative \cite{Car79}.

(4) Suppose that no place of $S$ lies over $p$. First, $\phi$ is of type
(Definition~\ref{gal-hecke-type-functions}). Let $k\in K$, and write $k=j(k')$ with $\theta(k')\in K_S$. By
Lemma~\ref{gal-identification}(2), $s(k)^{-1}k'$ has component $1$ at the places prime to $p$. At a place $u$
over $p$ its component is that of $k'$, which lies in $\operatorname{GL}_n(\mathcal O_u)$ because
$u,c(u)\notin S$ and $\theta(k')\in K_S$. So the similitude factor of $s(k)^{-1}k'$ is a unit, and so is that
of $s(k)$, as that of $k'$ is. The element $\theta(s(k))$ agrees with $\theta(k')$ at the places prime to $p$
and is $1$ at the places over $p$, so it lies in $K_S$. Its component at $w_0$ is
${}^t\operatorname{pr}_{w_0}(k)^{-1}$, which lies in $J$ because $\operatorname{pr}_{w_0}(k)\in J'$
(Lemma~\ref{gal-test-function-type}(1)), and
$\chi({}^t\operatorname{pr}_{w_0}(k)^{-1})=\psi(\operatorname{pr}_{w_0}(k))$. Since $s$ and $\theta$ are
homomorphisms, Definition~\ref{gal-trace-vector}(4) gives
$\phi(xk)=\Phi(\theta(s(x))\,\theta(s(k)))=\psi(\operatorname{pr}_{w_0}(k))\,\phi(x)$.

If $g$ is of kind (2)(a), then $\vartheta=1$, and the sum in (4) is $(f_g\cdot\phi)(x)=\phi(x)$ by
Lemma~\ref{gal-hecke-seen-action}(2). Otherwise $g$ is of kind (2)(b), with the place $w$, the integer $i$,
the matrix $t_i$, the element $z$ and the Satake parameters $B$ of
Definition~\ref{gal-trace-seen-operators}(2)(b). By Lemma~\ref{gal-hecke-seen-action}(3), the sum in (4) is
\[
 (f_g\cdot\phi)(x)=\operatorname{idx}(g)\sum_{y\in\operatorname{GL}_n(E_w)/\operatorname{GL}_n(\mathcal O_w)}
 1_{\operatorname{GL}_n(\mathcal O_w)t_i^{-1}\operatorname{GL}_n(\mathcal O_w)}(y)\,
 \phi\bigl(x\,j(z)^{-i}\operatorname{loc}_w(y)\bigr).
\]
The elements $z^{-i}$ and $\iota_w(y)$ have component $1$ at every place over $p$, by
Lemma~\ref{gal-hecke-seen-shape}(2) and since $w$ is prime to $p$. By Lemma~\ref{gal-identification}(2) they
have the same components away from $p$ as $s(j(z)^{-i})$ and $s(\operatorname{loc}_w(y))$, which also have
component $1$ at the places over $p$. So $s(j(z)^{-i})=z^{-i}$ and $s(\operatorname{loc}_w(y))=\iota_w(y)$,
and $\phi(x\,j(z)^{-i}\operatorname{loc}_w(y))=\Phi(\theta(s(x)z^{-i}\iota_w(y)))$. By
Lemma~\ref{gal-hecke-satake-eigenvalue}, applied to the element $s(x)$ of $\mathrm{GU}_n(\mathbf A_f)$ and to
the uniformizer $\varpi_w$ in $t_i$, the sum over $y$ is $q_w^{i(n-i)/2}e_i(B)\,\kappa(\varpi_w)^{-i}\phi(x)$.
Hence $(f_g\cdot\phi)(x)=\operatorname{idx}(g)\,q_w^{i(n-i)/2}e_i(B)\,\kappa(\varpi_w)^{-i}\phi(x)
=\vartheta\,\phi(x)$.
% Lean: (1) and (2) are memFn = memFn_holds = pMemFn seenShape_holds seenActs_holds actsConv_holds
% tyFact_holds; (4) is the last clause of phiEx = phiEx_holds = pPhiEx phiTy_holds seenShape_holds
% seenActs_holds tyFact_holds tfAct_holds tfCen_holds tEig_of_thLoc atApart_holds pFacEq pJEq pLocComm,
% where the eigenvalue idx g * satEig * kappa^{-i} comes from IsForm.sat and IsKap.cen. (3) is memComm,
% through comm_of_scaled and the commutation lemmas TfComm.
\end{proof}

\begin{lemma}\label{gal-trace-eigenform}
In the setting of Definition~\ref{gal-test-function}, suppose that $w_0$ lies over a good prime, that no place of $S$
lies over $p$ and that $\Phi(g_0)\neq0$ for some $g_0$ with $g_{0,w}\in\operatorname{GL}_n(\mathcal O_w)$ for every
$w\mid p$. Let $\xi=\xi_W\circ\theta$ as a representation of $\Gamma$. Then:
\begin{enumerate}
\item $\xi$ has weight $\lambda_V$ in the sense of Definition~\ref{gal-plan}: for $\gamma\in\Gamma$ whose
  characteristic polynomial is separable, $\operatorname{tr}\xi(\gamma)$ is the product over the
  embeddings $\tau$ of $E$ of the Schur polynomials of $\lambda_V(\tau)$ at the inverses of the eigenvalues
  of $\tau(\gamma)$;
\item there is a nonzero function $\Phi'$ on $G(\mathbf A_f^p)\times G(\mathbf Q_p)$ with values in
  $\mathcal W^\vee$, right invariant under $K_N\times K_p$, with $\Phi'(\gamma x)=\xi(\gamma)\Phi'(x)$ for
  $\gamma\in\Gamma$, embedded by Lemma~\ref{gal-identification}(5), such that $e\cdot\Phi'=\Phi'$ and
  $f_g\cdot\Phi'=\vartheta\Phi'$ for every Hecke operator $g$ of level $N$ seen on $U$ with value $\vartheta$
  and auxiliary prime $r$, where $(f\cdot\Phi')(x)=\sum_{y\in G(\mathbf A_f^p)/K_N}f(y)\Phi'(x(y,1))$.
\end{enumerate}
The hypothesis on $w_0$ is not used, and (1) does not use the hypotheses on $S$ and on $\Phi(g_0)$.
\end{lemma}
\leanref{Development/AutomorphicArithmetic/Foundations/GaloisRepresentations/ScalarCommutation/FrameTh.lean}{UpperTheory.Gal.Spc.eigEx}
\leanref{Development/AutomorphicArithmetic/Foundations/GaloisRepresentations/SphericalCoordinates/AlgebraEquivSetup/PHasEigOf.lean}{UpperTheory.Gal.W2.Seed.pHasEigOf}
\leanref{Development/AutomorphicArithmetic/Foundations/GaloisRepresentations/ProjectiveAutomorphicHecke/CharConjHolds.lean}{UpperTheory.Gal.Int.Geo.Pl.wtOf_holds}

\begin{proof}
(1) Let $\gamma=(g,\nu)\in\Gamma$, and let $\alpha_1,\ldots,\alpha_n$ be the eigenvalues of $\tau_v(g)$. The
relation ${}^tc(g)g=\nu$ gives $c(g)\,{}^tg=\nu$, so the eigenvalues of $(c\circ\tau_v)(g)$ are the
$\nu/\alpha_i$. The dual of $\nabla(W_v)$ composed with $y\mapsto{}^ty^{-1}$ is isomorphic to $\nabla(W_v)$,
so by the Weyl character formula \cite{FH91}, with $s_\mu$ the Schur polynomial of $\mu$,
\[
 \operatorname{tr}\xi(\gamma)=\nu^{-j_z}\prod_vs_{W_v}(\alpha)
 =\prod_vs_{(m_v,\ldots,m_v)}(\alpha^{-1})\,s_{W_v+m_v}\bigl((\nu/\alpha)^{-1}\bigr),
\]
using $\nu^{-\sum_iW_{v,i}}s_{W_v}(\alpha)=s_{W_v}(\alpha/\nu)$, $\nu^{-nm_v}=\prod_i\alpha_i^{-m_v}
\prod_i(\nu/\alpha_i)^{-m_v}$ and $s_{\mu+m}(z)=(z_1\cdots z_n)^ms_\mu(z)$.

(2) Let $\phi$ be the function of Lemma~\ref{gal-test-function-hecke}. For $\gamma\in\Gamma$ whose image at
$p$ lies in $K_p$, the element $s(j(\gamma)y)$ differs from $\gamma s(y)$ by an element that is trivial away
from $p$, lies in $\operatorname{GL}_n(\mathcal O_w)$ at $w\mid p$ by Lemma~\ref{gal-identification}(5), and
commutes with $s(y)$. As $\theta$ of such an element lies in $K_S$, with component $1$ at $w_0$, and no
place of $S$ lies over $p$, the equivariance of $\Phi$ gives $\phi(j(\gamma)y)=\xi(\gamma)\phi(y)$. Hence
the function $\Phi'$ on $G(\mathbf A_f^p)\times G(\mathbf Q_p)$ with $\Phi'(\gamma(y,k))=\xi(\gamma)\phi(y)$
for $\gamma\in\Gamma$ and $k\in K_p$, and $\Phi'=0$ off $\Gamma(G(\mathbf A_f^p)\times K_p)$, is well
defined, right invariant under $K_p$ and satisfies $\Phi'(\gamma x)=\xi(\gamma)\Phi'(x)$ and
$\Phi'(y,1)=\phi(y)$. It is nonzero, because the same argument applied to $\theta(g_0)$ in place of
$\gamma$ gives $\phi(j(\theta(g_0)))=\Phi(g_0)$. For $k\in K$, $\theta(s(k))$ lies in $K_S$ and has the same
component at $w_0$ as $\theta(k')$ for any $k'$ with $j(k')=k$, so
$\phi(yk)=\psi(\operatorname{pr}_{w_0}(k))\phi(y)$ by the equivariance of $\Phi$ and the computation at the
start of the proof of Lemma~\ref{gal-test-function-type}. In particular $\Phi'$ is right invariant under
$K_N$, and summing over $K/K_N$ gives $e\cdot\Phi'=\Phi'$. The eigenvalue equations are
Lemma~\ref{gal-test-function-hecke}(4).
\end{proof}
% Lean: eigEx applies pHasEigOf (pStdForm, pPInt, pPPart, pExtP) to Phi o th o sec; (1) is WtOf =
% wtOf_holds = wtOf_of wtCore_holds, with IsWt stated through schW and endPol = invPol.

\begin{lemma}\label{gal-trace-linear-algebra}
Let $Q$ be a finite set and $J$ a set. For each $q\in Q$ let $V_q$ be a finite-dimensional complex vector
space with endomorphisms $T_{q,j}$ for $j\in J$, $e_q$, $u_q$ and $u'_q$ such that
\[
 e_q^2=e_q,\quad e_qT_{q,j}e_q=T_{q,j}e_q,\quad e_qu_qe_q=u_qe_q,\quad
 T_{q,j}T_{q,j'}e_q=T_{q,j'}T_{q,j}e_q,\quad u_qT_{q,j}e_q=T_{q,j}u_qe_q,\quad u'_qu_qe_q=e_q
\]
for all $j,j'$. Let $\vartheta\in\mathbf C^J$, and suppose that for some $q_0$ there is a nonzero $x\in
V_{q_0}$ with $e_{q_0}x=x$ and $T_{q_0,j}x=\vartheta_jx$ for all $j$. Then there are $m\geq1$ and a
noncommutative polynomial $a$ in variables indexed by $J$ such that for every noncommutative polynomial $s$
with $s(\vartheta)\neq0$,
\[
 \sum_{q\in Q}\operatorname{tr}\bigl(a(T_q)\,s(T_q)\,u_q^m\,e_q\bigr)\neq0.
\]
\end{lemma}
\leanref{Development/AutomorphicArithmetic/Foundations/GaloisRepresentations/ScalarCommutation/SpecAlg.lean}{UpperTheory.Gal.Spc.specAlg}
\leanref{Development/AutomorphicArithmetic/Foundations/GaloisRepresentations/ProjectiveAutomorphicHecke/PProjOne.lean}{UpperTheory.Gal.CertificateRingHomMap.pProjOne_holds}
\leanref{Development/AutomorphicArithmetic/Foundations/GaloisRepresentations/ProjectiveAutomorphicHecke/PTracePow.lean}{UpperTheory.Gal.CertificateRingHomMap.pTracePow_holds}

\begin{proof}
Let $V=\bigoplus_qV_q$ with the diagonal endomorphisms $e$, $u'$, $t_j=T_je$ and $v=ue$. The relations give
that the $t_j$ commute with each other, with $v$ and with $e$, that $v$ commutes with $e$, that $u'v=e$, and
that $e$ is idempotent; for example $t_jv=T_jeue=T_jue=uT_je=ueT_je=vt_j$. Since the $t_j$ commute and $V$
is finite-dimensional, there is a polynomial $a$ in them with $a(\vartheta)=1$ such that $P=a(t)$ is the
projection onto the generalized eigenspace of the $t_j$ for the eigenvalues $\vartheta_j$: $P^2=P$, and
$(t_j-\vartheta_j)^{N_j}P=0$ for some $N_j$. The vector $x$, placed in $V_{q_0}$, satisfies
$t_jx=\vartheta_jx$ and $ex=x$, so $Px=a(\vartheta)x=x$, and $P_1=Pe$ is a nonzero idempotent commuting
with $v$. On the image of $P_1$, $v$ is injective: if $vP_1y=0$ then $P_1y=eP_1y=u'vP_1y=0$. So $v$
restricts to an automorphism of $P_1V\neq0$, its eigenvalues there are nonzero, and by Newton's identities
their power sums $\operatorname{tr}(v^mP_1)$ cannot vanish for all $m=1,\ldots,\dim P_1V$. Fix $m\geq1$
with $\operatorname{tr}(v^mP_1)\neq0$.

Using $eT_je=T_je$ and $eue=ue$ repeatedly, $a(T_q)s(T_q)u_q^me_q=a(t)s(t)v^me$ on $V_q$, so the sum in
the statement is $\operatorname{tr}(s(t)\,v^me\,P)$. On the image of $P$ each $t_j-\vartheta_j$ is
nilpotent and commutes with $v^me$ and $P$, so $\operatorname{tr}(s(t)\,v^meP)=s(\vartheta)
\operatorname{tr}(v^meP)=s(\vartheta)\operatorname{tr}(v^mP_1)\neq0$.
\end{proof}

\begin{lemma}\label{gal-trace-cuspidal-vanishing}
In the setting of Definition~\ref{gal-trace-seen-operators}, let $w_0$ and $w_1$ be places of $E$ split
over $F$ and prime to $p$, let $K\subset G(\mathbf A_f^p)$ be a compact open subgroup that is the product
of its intersections with $\operatorname{loc}_{w_0}(\operatorname{GL}_n(E_{w_0}))$ and with the kernel of
$\operatorname{pr}_{w_0}$, and let $f$ be a bi-$K$-invariant compactly supported function on
$G(\mathbf A_f^p)$ such that:
\begin{enumerate}
\item $f$ is \emph{cuspidal at $w_0$}: $\int_Nf(x\operatorname{loc}_{w_0}(u))\,du=0$ for every $x$ and the
  unipotent radical $N$ of every proper parabolic subgroup of $\operatorname{GL}_n(E_{w_0})$;
\item $\operatorname{pr}_{w_1}(y)$ has separable characteristic polynomial whenever $f(y)\neq0$.
\end{enumerate}
Let $\beta$ be a nondegenerate Hermitian form on $E^n$ that is isomorphic to that of the PEL datum, up to a
scalar, over $\mathbf A_f^p$, and let $\delta$ be a rational similitude of $\beta$ whose minimal polynomial
over $E$ is not separable and irreducible of degree $n$. Then the orbital sum
\[
 \sum_{x\in Z\backslash(G(\mathbf A_f^p)\times G(\mathbf Q_p))/(K\times K_p)}
 \operatorname{vol}\bigl(Z\cap x(K\times K_p)x^{-1}\bigr)^{-1}\,(f\otimes1_{K_p})(x^{-1}\delta x)
\]
of $f\otimes1_{K_p}$ at the image of $\delta$ in $G(\mathbf A_f^p)\times G(\mathbf Q_p)$ vanishes, where
$Z$ is the centralizer of that image and the volumes are relative to $Z\cap(K\times K_p)$. With $G$
presented as in Situation~\ref{gal-count-situation} with $F=E$, the orbital sum $O_P(f,1_{K_p})$
(Definition~\ref{gal-count-pairs}(3)) of every pair $P=(\beta,\gamma,c)$ that has a marking and whose
minimal polynomial $P_\gamma$ is reducible vanishes as well.
\end{lemma}
\leanref{Development/AutomorphicArithmetic/Foundations/GaloisRepresentations/ScalarCommutation/VanForms.lean}{UpperTheory.Gal.Spc.vanForms}
\leanref{Development/AutomorphicArithmetic/Foundations/GaloisRepresentations/ScalarCommutation/VanPairs.lean}{UpperTheory.Gal.Spc.vanPairs}
\leanref{Development/AutomorphicArithmetic/Foundations/GaloisRepresentations/ScalarCommutation/OrbZero.lean}{UpperTheory.Gal.Spc.orbK_eq_zero_of_cuspSlices}

\begin{proof}
If $\operatorname{pr}_{w_1}(\delta)$ has a characteristic polynomial that is not separable, so does every
conjugate, and every term vanishes by (2). Otherwise this characteristic polynomial is the image of the
minimal polynomial of $\delta$, which therefore has degree $n$ and is separable, hence reducible by
hypothesis. Then the centralizer $Z$ is commutative, and $\operatorname{pr}_{w_0}(\delta)$ is a regular
semisimple element of $\operatorname{GL}_n(E_{w_0})$ with reducible characteristic polynomial, so it lies
in a proper Levi subgroup and is not elliptic. By the product decomposition of $K$, the orbital sum is a
sum, over the cosets away from $w_0$, of orbital integrals at $\operatorname{pr}_{w_0}(\delta)$ of the
functions $y\mapsto f(x\operatorname{loc}_{w_0}(y))$, each of which is a cusp form by (1). The orbital
integrals of a cusp form vanish at regular semisimple elements that are not elliptic \cite{HC70}.

For a pair $P=(\beta,\gamma,c)$ with a marking and reducible $P_\gamma$, the element $\delta=\gamma$ is a
rational similitude of $\beta$ whose minimal polynomial $P_\gamma$ is separable of degree $n$ and
reducible, so the orbital sum above at $\delta$ vanishes, and its centralizer $Z$ is commutative. The
groups $Z_P\cap(g,g_p)(K\times K_p)(g,g_p)^{-1}$ in $O_P$ are finite, since $P$ has a marking. Grouping
the terms of $O_P(f,1_{K_p})$ by the image of $(g,g_p)$ in $Z\backslash(G(\mathbf A_f^p)\times
G(\mathbf Q_p))/(K\times K_p)$, on which $f\otimes1_{K_p}$ at the conjugate of $\delta$ is constant, and
using that $Z$ is commutative, writes $O_P(f,1_{K_p})$ as one multiple of that orbital sum, so it
vanishes.
\end{proof}
% Lean: vanForms (definite forms B, through orbK_eq_zero_of_charpoly_not_separable and
% orbK_eq_zero_of_cuspSlices with orbK_prod_eq_zero_of_slices), and vanPairs, the same vanishing for the
% terms P.term of pairs with reducible polynomial. The Lean proves the vanishing for cusp forms itself.

\begin{proposition}\label{gal-trace-nonzero}
In the setting of Lemma~\ref{gal-identification}, let $x\in E$ with $\mathbf Q(x)=\mathbf Q(\omega)$, let
$U$ be an automorphic representation of $U_n$, let $w_0$ be a place of $E$ over a good prime, and let $V$ be
a vector for $U$ at $w_0$ (Definition~\ref{gal-trace-vector}) such that no place of $S$ lies over $p$ and
$\Phi(g_0)\neq0$ for some $g_0$ with $g_{0,w}\in\operatorname{GL}_n(\mathcal O_w)$ for every $w\mid p$.
Then for every finite set $X$ of primes there is a prime $r_0\notin X$ with $r_0\neq p$ such that the
following holds. Let $N\geq3$ be prime to $p$ and divisible by $r_0$, such that every element of $K_N$
stabilizes $V$, and let $(g_j)_{j\in J}$ be Hecke operators of level $N$, each seen on $U$ with a value
$\vartheta_j$ and auxiliary prime $r_0$ (Definition~\ref{gal-trace-seen-operators}). Then there is a
bi-$K_N$-invariant compactly supported function $f$ on $G(\mathbf A_f^p)$, with $y^{-1}$ integral whenever
$f(y)\neq0$, such that the functions $f_{g_j}$ and $f$ satisfy the spectral condition of
Definition~\ref{gal-plan}(1) with the values $\vartheta$ and the weight $\lambda_V$. Moreover, each function
$(ts)(g)\ast f$ for which that condition asserts a nonzero trace can be taken cuspidal at $w_0$ and
supported on elements $y$ whose centralizer in $\operatorname{End}_E(\Lambda\otimes\mathbf Q_{r_0})$ is
commutative and whose components $\operatorname{pr}_w(y)$, at the places $w$ of a set $W$ given by
Lemma~\ref{gal-trace-auxiliary-prime} for $r_0$, have separable characteristic polynomial; in particular its
orbital sums of Lemma~\ref{gal-trace-cuspidal-vanishing} vanish.
\end{proposition}
\leanref{Development/AutomorphicArithmetic/Foundations/GaloisRepresentations/IntegralStructures/GeometricFlagBases/PrimeLevelLocalization/SpecTopHolds.lean}{UpperTheory.Gal.Int.Geo.Pl.specTop_holds}
\leanref{Development/AutomorphicArithmetic/Foundations/GaloisRepresentations/IntegralStructures/GeometricFlagBases/PrimeLevelLocalization/PrimeProj.lean}{UpperTheory.Gal.Int.Geo.Pl.SpecTop}
\leanref{Development/AutomorphicArithmetic/Foundations/GaloisRepresentations/ScalarCommutation/SpecTopG.lean}{UpperTheory.Gal.Spc.specTopG}
\leanref{Development/AutomorphicArithmetic/Foundations/GaloisRepresentations/ScalarCommutation/SpecData.lean}{UpperTheory.Gal.Spc.specData}
\leanref{Development/AutomorphicArithmetic/Foundations/GaloisRepresentations/ScalarCommutation/SpecOfF.lean}{UpperTheory.Gal.Spc.specOfF}

\begin{proof}
Let $r_0=r$ and $W$ be given by Lemma~\ref{gal-trace-auxiliary-prime} for $X$ and the set $S$ of $V$, and
let $N$ and $(g_j)$ be as in the statement. Definition~\ref{gal-test-function} gives the type and the test
function $e$, Lemma~\ref{gal-test-function-powers} gives $u$ and $u'$, and we write $\psi_j=f_{g_j}$.

\emph{Relations.} By Lemma~\ref{gal-test-function-hecke}(2), each $\psi_j\ast e$ is a multiple of some
$e_h$ with $h$ trivial at $w_0$ and at the places of $W$, and $e\ast e_h=e_h$; so
$e\ast(\psi_j\ast e)=\psi_j\ast e$, and $u\ast(\psi_j\ast e)=\psi_j\ast(u\ast e)$ by
Lemma~\ref{gal-test-function-powers}(2). With $e\ast e=e$, Lemma~\ref{gal-test-function-hecke}(3) and
Lemma~\ref{gal-test-function-powers}(1), the functions $\psi_j$, $e$, $u$, $u'$ satisfy the relations of
Lemma~\ref{gal-trace-linear-algebra} for convolution. On functions $Y$ with $e\ast Y=Y$, such as
$p(\psi)\ast u^{\ast m}\ast e$ for a noncommutative polynomial $p$, left convolution by $\psi_j$ agrees
with left convolution by $\psi_j\ast e$, and these commute; so the $\psi_j$ commute on every
$p(\psi)\ast u^{\ast m}\ast e$.

\emph{Supports.} By Lemma~\ref{gal-test-function-type}(5), for each $x$ the function
$F_x(y)=e(x\operatorname{loc}_{w_0}(y))$ on $\operatorname{GL}_n(E_{w_0})$ is compactly supported and
satisfies $F_x(yk)=\psi(k)^{-1}F_x(y)$ for $k\in J'$. So it lies in the representation of
$\operatorname{GL}_n(E_{w_0})$ by left translation compactly induced from $\psi$ on $J'$, whose irreducible
quotients have nonzero vectors on which $J'$ acts by $\psi$ and are cuspidal by
Lemma~\ref{gal-test-function-type}(4). Hence the constant terms of $F_x$ vanish \cite{BZ76}: $e$ is
cuspidal at $w_0$. Left convolution with bi-$K_N$-invariant functions and linear combinations preserve
cuspidality at $w_0$, so every $p(\psi)\ast u^{\ast m}\ast e$ is cuspidal at $w_0$. For $m\geq1$, let
$X_w$ be the sets of Lemma~\ref{gal-test-function-powers}(3). By Lemma~\ref{gal-test-function-hecke}(1) the
set of $y$ with $\operatorname{pr}_w(y)\in X_w$ for all $w\in W$ is stable under left multiplication by the
supports of the $\psi_j$, so it contains the support of every $p(\psi)\ast u^{\ast m}\ast e$. For such $y$,
$\operatorname{pr}_w(y)$ and hence $\operatorname{pr}_{c(w)}(y)$ have separable characteristic polynomial
for $w\in W$, so by Lemma~\ref{gal-trace-auxiliary-prime}(3) the component of $y$ at every place over $r_0$
does, and by Lemma~\ref{gal-identification}(1) the centralizer of $y$ in
$\operatorname{End}_E(\Lambda\otimes\mathbf Q_{r_0})$ is commutative.

\emph{Traces.} The classes of definite forms with the same finite adelic points as $G$ are finite in
number, and for each the forms of weight $\xi$ (Lemma~\ref{gal-trace-eigenform}) and level $K_N\times K_p$
form a finite-dimensional space, on which a bi-$K_N$-invariant function $f'$ acts by an operator $T(f')$
with $T(f'\ast f'')=T(f')T(f'')$; the trace in the spectral condition is the sum of the traces of these
operators. Lemma~\ref{gal-trace-eigenform}(2) gives a nonzero form on one of these spaces fixed by $T(e)$
with eigenvalues $\vartheta_j$ for the $T(\psi_j)$. Lemma~\ref{gal-trace-linear-algebra}, applied to the
operators $T(\psi_j)$, $T(e)$, $T(u)$ and $T(u')$, gives $m\geq1$ and $a$ such that the trace of
$(as)(\psi)\ast u^{\ast m}\ast e$ is nonzero whenever $s(\vartheta)\neq0$.

Put $f=u^{\ast m}\ast e$. It is bi-$K_N$-invariant with compact support, and $y^{-1}$ is integral whenever
$f(y)\neq0$, by Lemma~\ref{gal-test-function-powers}(3). The $\psi_j$ commute on $f$, $\xi$ has weight
$\lambda_V$ by Lemma~\ref{gal-trace-eigenform}(1), and with $t=a$ the function $(ts)(\psi)\ast f$ has
nonzero trace and the cuspidality and supports above. Lemma~\ref{gal-trace-cuspidal-vanishing}, applied
with $w_1\in W$, gives the vanishing of its orbital sums at the elements of groups of definite forms and
at the pairs with a marking and reducible minimal polynomial. Its support consists of regular elements in
the sense of Situation~\ref{gal-count-situation}, with $r_0$ as the prime there, since the centralizer of
each element of the support in $\operatorname{End}_E(\Lambda\otimes\mathbf Q_{r_0})$ is commutative. So
$(ts)(\psi)\ast f$ is elliptic, as Definition~\ref{gal-plan}(1) requires.
\end{proof}
% Lean: specTop_holds = specTopG wtOf_holds. specTopG takes r0 and W from auxiliary_prime_holds, the type
% from tyEx, and Extra (plTr, regTr) from the frame; specData assembles tyTop, auxFn, memFn, memComm,
% eigEx, splitC and cuspIrr_kFinite_rep_eq_centralCuspForm (cuspidality of e), and applies specOfF, which
% uses cuspIdeal, clsMove, opAlg, specAlg, words, vanPairs and vanForms. The trace is dTrace, the sum
% over DFCls G of the traces on HeckeTraceForms.forms; Good is the conjunction of regular support (RegAf),
% vanPairs and vanForms.

\begin{proposition}\label{gal-plans-exist}
In Situation~\ref{gal-seed-situation}:
\begin{enumerate}
\item there is a plan for $U$ with Hodge data at $v'$ (Definition~\ref{gal-plan});
\item for every $d$-split place $w$ not over $\ell$ at which $U$ has Satake parameters, there is a plan
  with Hodge data at $v'$ that is congruent at $w$;
\item for every place $w$ over $\ell$ that lies over the same place of $\mathbf Q(x)$ as $v'$, over a
  place of $F$ different from those under $w_0$ and $v'$, and at which $U$ has Satake parameters, there
  are a plan with Hodge data at $v'$ and at $w$, and a plan with Hodge data at $v'$ that is congruent at
  $w$.
\end{enumerate}
\end{proposition}
\leanref{Development/AutomorphicArithmetic/Foundations/GaloisRepresentations/IntegralStructures/GeometricFlagBases/PackExOf.lean}{UpperTheory.Gal.Int.Geo.PlanEx}
\leanref{Development/AutomorphicArithmetic/Foundations/GaloisRepresentations/IntegralStructures/GeometricFlagBases/PrimeLevelLocalization/SpecTopHolds.lean}{UpperTheory.Gal.Int.Geo.Pl.planEx_holds}
\leanref{Development/AutomorphicArithmetic/Foundations/GaloisRepresentations/IntegralStructures/GeometricFlagBases/PrimeLevelLocalization/PrimeProj.lean}{UpperTheory.Gal.Int.Geo.Pl.Frame}
\leanref{Development/AutomorphicArithmetic/Foundations/GaloisRepresentations/IntegralStructures/GeometricFlagBases/PrimeLevelLocalization/PrimeProj.lean}{UpperTheory.Gal.Int.Geo.Pl.Vec}
\leanref{Development/AutomorphicArithmetic/Foundations/GaloisRepresentations/IntegralStructures/GeometricFlagBases/PrimeLevelLocalization/PrimeProj.lean}{UpperTheory.Gal.Int.Geo.Pl.VecEx}
\leanref{Development/AutomorphicArithmetic/Foundations/GaloisRepresentations/IntegralStructures/GeometricFlagBases/PrimeLevelLocalization/PrimeProj.lean}{UpperTheory.Gal.Int.Geo.Pl.ExpOf}
\leanref{Development/AutomorphicArithmetic/Foundations/GaloisRepresentations/IntegralStructures/GeometricFlagBases/PrimeLevelLocalization/PrimeProj.lean}{UpperTheory.Gal.Int.Geo.Pl.SpecTop}
\leanref{Development/AutomorphicArithmetic/Foundations/GaloisRepresentations/ScalarCommutation/SpecTopG.lean}{UpperTheory.Gal.Spc.specTopG}

\begin{proof}
The construction has four inputs. \emph{The identification.} Since $a$ is odd, the similitude group
of the Hermitian space of the PEL datum and that of the standard definite Hermitian space of rank $a$
are isomorphic at every finite place (Lemma~\ref{gal-identification-isometry}). So there is a surjective homomorphism $j$ from the finite adelic points of the
similitude group of the standard definite space onto those of $G$ away from $p$, whose kernel consists
of the elements trivial at the places prime to $p$, together with matrices $a_w\in\operatorname{GL}_a
(E_w)$, integral at almost all $w$, that conjugate the local components of $j$ into the standard
coordinates on $\Lambda\otimes E_w$ (Lemma~\ref{gal-identification}). \emph{The vector.} From $U$ one chooses its weight $W$ and a twist
$m$ making all weights nonnegative, the subgroup $J$ and character $\chi$ at $w_0$ of hypothesis (2) of
Theorem~\ref{gal-seed-representations} (every representation containing $\chi$ on $J$ is
supercuspidal, and $J$ is torsion-free), a finite set $S$ of places with levels $K_w$, a linear form
$\Phi$ on the weight space nonzero on $U$, and the Hecke character $\kappa$ determined by $W$ and $m$;
these form a vector for $U$ at $w_0$ (Definition~\ref{gal-trace-vector}). \emph{The exponents.} For every embedding $\tau$ there is $\sigma\in\operatorname{Aut}(\mathbf C)$
such that every Hodge--Tate weight $k-\lambda'_k$ of $U$ at $\tau$, shifted by $m_\tau$, is a Hodge
exponent of $\sigma\circ\lambda$, where $\lambda$ is the weight of the plan built from $W$ and $m$.
\emph{The spectral statement.} For an auxiliary prime $r_0$ outside any given finite set and different
from $p$, and every level $N$ divisible by $r_0$ whose elements stabilize the data of the vector, every
family of Hecke operators of level $N$ that act on $U$ through $j$ (deck transformations with value $1$,
and the operators $T_{w,i}$ at free split places with values given by the $i$-th elementary symmetric
function of the Satake parameters of $U$ divided by $\kappa(\varpi_w)^i$) satisfies the spectral
condition with a test function $f$ supported on integral elements (Proposition~\ref{gal-trace-nonzero}).
% Lean: the identification is Pl.Frame (FrameEx); the spectral statement is specTop_holds =
% UpperTheory.Gal.Spc.specTopG wtOf_holds. The Hecke operators entering the congruences are those of
% IsTM, Seen and idx.

Given these, the plan of (1) uses a prime $p$ outside a finite set, the identification, the vector and the
operators at places away from $w_0$ and $v'$; (2) adds the Hecke operators at $w$, and (3) adds those at
$w$ or requires good operators at $w$. The congruence condition at almost all $d$-split places holds
because the operators $T_{w,i}$ act on $U$ by the elementary symmetric functions of the Satake
parameters, which are the coefficients of the Hecke polynomial with the stated roots.
\end{proof}

\begin{proposition}\label{gal-packs-exist}
In Situation~\ref{gal-seed-situation}:
\begin{enumerate}
\item there is a pack for $U$ with Hodge data at $v'$;
\item for every $d$-split place $w$ not over $\ell$ at which $U$ has Satake parameters, there is a pack
  with Hodge data at $v'$ that is congruent at $w$;
\item for every place $w$ over $\ell$ that lies over the same place of $\mathbf Q(x)$ as $v'$, over a
  place of $F$ different from those under $w_0$ and $v'$, and at which $U$ has Satake parameters, there
  are a pack with Hodge data at $v'$ and at $w$, and a pack with Hodge data at $v'$ that is congruent at
  $w$.
\end{enumerate}
\end{proposition}
\leanref{Development/AutomorphicArithmetic/Foundations/GaloisRepresentations/IntegralStructures/GeometricFlagBases/HodgeTrAt.lean}{UpperTheory.Gal.Int.Geo.PackEx}
\leanref{Development/AutomorphicArithmetic/Foundations/GaloisRepresentations/IntegralStructures/GeometricFlagBases/PackExOf.lean}{UpperTheory.Gal.Int.Geo.packEx_of}

\begin{proof}
Each item of Proposition~\ref{gal-plans-exist} gives plans with the corresponding Hodge data and
congruences, and Lemma~\ref{gal-plan-pack} turns each plan into a pack with the same Hodge data and
congruences.
\end{proof}

\section{Construction of the Galois representations}\label{section-gal-construction}

\subsection{Plain components and solvable base change}\label{subsection-gal-plain-components}

\begin{definition}\label{gal-weakly-attached}
Let $E$ be a CM field, $\Omega$ an algebraic closure of $E$, $\ell$ a prime,
$\iota:\bar{\mathbf Q}_\ell\to\mathbf C$ a field isomorphism and $a\geq1$.
\begin{enumerate}
\item Let $w$ be a finite place of $E$, $q_w=|k_w|$ and $R$ a multiset of $a$ complex numbers. A
  homomorphism $r:G_E\to\operatorname{GL}_a(\bar{\mathbf Q}_\ell)$ has \emph{Frobenius eigenvalues $R$
  at $w$} if for every valuation ring $A$ over $w$, $r$ is trivial on $I_A$ and, for every
  $\sigma\in D_A$ acting on the residue field by $z\mapsto z^{q_w}$, the characteristic polynomial of
  $\iota(r(\sigma^{-1}))$ is $\prod_{z\in R}(X-z)$.
\item Let $\pi$ be a cuspidal automorphic representation of $\operatorname{GL}_a(\mathbf A_E)$
  (Definition~\ref{gal-polarized}). It has \emph{Satake parameters} $B$ at $w$ if its local component
  at $w$ is spherical with Satake parameters $B$, normalized as in
  Definition~\ref{gal-definite-unitary}. A homomorphism $r:G_E\to\operatorname{GL}_a(\bar{\mathbf
  Q}_\ell)$ is \emph{weakly attached} to $\pi$ if it is continuous and semisimple and, for all but
  finitely many $w$, there is $B$ such that $\pi$ has Satake parameters $B$ at $w$ and $r$ has
  Frobenius eigenvalues $q_w^{(a-1)/2}B$ at $w$.
\end{enumerate}
\end{definition}
\leanref{Development/AutomorphicArithmetic/AutomorphicVarieties/AutomorphicForms/BlochKatoSelmerVanishes.lean}{UpperTheory.V.Af.FrobEig}
\leanref{Development/AutomorphicArithmetic/AutomorphicVarieties/AutomorphicForms/BlochKatoSelmerVanishes.lean}{UpperTheory.V.Af.AttC}
\leanref{Development/AbelianVarietyTheory/PublishedTheorems/GaloisCohomology₂/StrongMultiplicityOne.lean}{GlobalTheory.CuspRep.CG.HasSat}
\leanref{Development/AbelianVarietyTheory/Definitions/GaloisGeometry/IsWeight.lean}{GlobalTheory.GalCut.Weak}

\begin{definition}\label{gal-spherical-conditions}
In the setting of Definition~\ref{gal-weakly-attached}, let $w$ be a place, $B$ a multiset of $a$
nonzero complex numbers, $Z=(Z_\tau)_{\tau:E\to\mathbf C}$ a family of multisets of $a$ complex
numbers, and $r:G_E\to\operatorname{GL}_a(\bar{\mathbf Q}_\ell)$. For a valuation ring $A$ over $w$,
let $\mathrm{Sat}_A(B)\in\mathrm{WD}(A)$ be the monodromy-free unramified representation on which a
geometric Frobenius element acts with eigenvalues $q_w^{(a-1)/2}b$, $b\in B$. Then $r$ satisfies, at
$w$ with respect to $B$ and $Z$, the condition:
\begin{enumerate}
\item[(unr)] if $w\nmid\ell$, $r$ has Frobenius eigenvalues $q_w^{(a-1)/2}B$ at $w$;
\item[(dR)] for every $A$ over $w$, $r|_{D_A}$ is de Rham and
  $\mathrm{HT}_\tau(r|_{D_A})=\{\frac{a-1}{2}-z:z\in Z_{\iota\circ\tau}\}$ for every
  $\tau:E\to\bar{\mathbf Q}_\ell$ that extends continuously to $E_w$
  (Definition~\ref{gal-p-adic-hodge});
\item[(crys)] for every $A$ over $w$, $r|_{D_A}$ is crystalline;
\item[(phi)] for every $A$ over $w$ and every $\sigma\in W_A$, $\mathrm{WD}_\iota(r|_{D_A})(\sigma)$ and
  $\mathrm{Sat}_A(B)(\sigma)$ have the same characteristic polynomial.
\end{enumerate}
The conditions (dR), (crys) and (phi) say nothing at places not over $\ell$. For a cuspidal $\pi$ with
Satake parameters $B$ at $w$, the conditions \emph{for $\pi$} are those with respect to $B$ and the
infinitesimal characters $Z_\tau(\pi)$, except that (crys) is required only if $\pi$ has a spherical
local component at $w$.
\end{definition}
\leanref{Development/AbelianVarietyTheory/Definitions/GaloisGeometry/CyclicStrongAtCusp.lean}{GlobalTheory.RoadU.satWD}
\leanref{Development/AbelianVarietyTheory/Definitions/GaloisGeometry/CyclicStrongAtCusp.lean}{GlobalTheory.RoadU.SphAt}
\leanref{Development/AbelianVarietyTheory/Definitions/GaloisGeometry/IsWeight.lean}{GlobalTheory.GalCut.ClauseAt}
\leanref{Development/AbelianVarietyTheory/Definitions/GaloisGeometry/IsRegAlgTwistOfHalfIntegral.lean}{GlobalTheory.Row9.SphG}

\begin{definition}\label{gal-scope}
Let $E$ be a CM field with maximal totally real subfield $F$, $\eta$ the quadratic character of
$E/F$, $\ell$ a prime, $a\geq1$, $\pi$ a cuspidal automorphic representation of
$\operatorname{GL}_a(\mathbf A_E)$ and $\chi$ a Hecke character of $F$.
\begin{enumerate}
\item $(\pi,\chi)$ is \emph{in scope} if $\pi$ is regular algebraic and $(\pi,\chi)$ is polarized
  (Definition~\ref{gal-polarized}). It is \emph{in good scope} for $\ell$ if moreover $E$ contains
  $\sqrt{-d}$ for some integer $d>0$, $[F:\mathbf Q]$ is even, and every place $w$ of $E$ not split
  over $F$ is prime to $\ell$, unramified over $F$, and a place where $\pi$ has Satake parameters. It is
  \emph{in self-dual scope} for $\ell$ if it is in good scope and $\chi=\eta^a$.
\item For $s$ one of the three scopes, \emph{existence holds in scope $s$} if for all $E$, $\Omega$,
  $\ell$, $\iota$, $a$ and $(\pi,\chi)$ in scope $s$ there is $r$ weakly attached to $\pi$
  (Definition~\ref{gal-weakly-attached}).
\item For $s$ a scope and $c$ one of (unr), (dR), (crys), (phi), \emph{condition $c$ holds in scope $s$}
  if for all $E$, $\Omega$, $\ell$, $\iota$, $a$, $(\pi,\chi)$ in scope $s$, every $r$ weakly attached
  to $\pi$, and every place $w$ at which $\pi$ has Satake parameters $B$, split over $F$ unless $s$ is
  the first scope, $r$ satisfies condition $c$ for $\pi$ at $w$ (Definition~\ref{gal-spherical-conditions}).
\end{enumerate}
\end{definition}
\leanref{Development/AbelianVarietyTheory/Definitions/GaloisGeometry/IsWeight.lean}{GlobalTheory.GalCut.InScope}
\leanref{Development/AbelianVarietyTheory/Definitions/GaloisGeometry/IsWeight.lean}{GlobalTheory.GalCut.Ex}
\leanref{Development/AbelianVarietyTheory/Definitions/GaloisGeometry/CyclicStrongAtCusp.lean}{GlobalTheory.RoadU.SphericalLevelCondition}
\leanref{Development/AbelianVarietyTheory/Definitions/GaloisGeometry/CyclicStrongAtCusp.lean}{GlobalTheory.RoadU.Core}
\leanref{Development/AbelianVarietyTheory/Definitions/GaloisGeometry/CyclicStrongAtCusp.lean}{GlobalTheory.RoadU.CoreAll}

\begin{definition}\label{gal-plain}
Let $k$ be a finite extension of $\mathbf Q_p$ and $x\in\operatorname{Irr}(k)$. A multiset $s$ of
elements of $\operatorname{Irr}(k)$ is \emph{unlinked} if no two members $y,z$ satisfy $y=z\nu$. Then
$x$ is \emph{plain} if its supercuspidal support consists of characters and is unlinked. For a
cuspidal $\pi$ of $\operatorname{GL}_a(\mathbf A_E)$ and a local component $\sigma$ of $\pi$ at $w$,
write $\sigma^\flat=\sigma\nu^{(1-a)/2}$.
\end{definition}
\leanref{Development/AbelianVarietyTheory/Definitions/GaloisGeometry/CyclicStrongAtCusp.lean}{GlobalTheory.RoadU.Unlinked}
\leanref{Development/AbelianVarietyTheory/Definitions/GaloisGeometry/CyclicStrongAtCusp.lean}{GlobalTheory.RoadU.PlainTor}
\leanref{Development/AbelianVarietyTheory/PublishedTheorems/GaloisCohomology₂/GaloisRepOfPolarizedCuspidal.lean}{GlobalTheory.GalOfCusp.algNorm}

\begin{lemma}\label{gal-generic-monodromy-free}
In Situation~\ref{gal-local-situation}, let $\operatorname{rc}$ satisfy the laws
(Definition~\ref{gal-llc-laws}) and let $x\in\operatorname{Irr}(E_w)$ be generic, with
$\operatorname{rc}(x)$ monodromy-free.
\begin{enumerate}
\item The supercuspidal support $\{y_1,\ldots,y_k\}$ of $x$ is unlinked
  (Definition~\ref{gal-plain}), $x\cong y_1\times\cdots\times y_k$, and
  $\operatorname{rc}(x)=\bigoplus_i\operatorname{rc}(y_i)$.
\item If $\operatorname{rc}(x)$ is abelian, then $x$ is plain.
\item Suppose $E_w\cong\mathbf Q_g$ for a prime $g$, and $\operatorname{rc}(x)$ is tame with
  multiplicity-free tame part (Definition~\ref{gal-wd-classes}). Then $x$ has depth zero, and no two
  members of degree at least $2$ of its supercuspidal support are inertially equivalent, that is,
  differ by a twist by an unramified character.
\end{enumerate}
\end{lemma}
\leanref{Development/AutomorphicArithmetic/Foundations/GaloisRepresentations/IntegralStructures/InertCharpoly/CharPolOfInert.lean}{UpperTheory.Gal.Int.Row.split}
\leanref{Development/AutomorphicArithmetic/Foundations/GaloisRepresentations/IntegralStructures/InertCharpoly/CharPolOfInert.lean}{UpperTheory.Gal.Int.Row.plainTor_of_tor}
\leanref{Development/AutomorphicArithmetic/Foundations/GaloisRepresentations/IntegralStructures/InertCharpoly/CharPolOfInert.lean}{UpperTheory.Gal.Int.Row.dz_of_cusp}
\leanref{Development/AutomorphicArithmetic/Foundations/GaloisRepresentations/IntegralStructures/InertCharpoly/CharPolOfInert.lean}{UpperTheory.Gal.Int.Row.charPol_of_inert}
\leanref{Development/AutomorphicArithmetic/Foundations/GaloisRepresentations/IntegralStructures/InertCharpoly/CharPolOfInert.lean}{UpperTheory.Gal.Int.Row.mixed_of_tame}
\leanref{Development/WeilBound/RepresentationTheoryGL/GaloisPermutationReps/SubquotientEquivalence/IsTypeMap.lean}{PositivityTheory.Rg.DzCount.St3}
\leanref{Development/WeilBound/RepresentationTheoryGL/GaloisPermutationReps/SubquotientEquivalence/IsTypeMap.lean}{PositivityTheory.Rg.DzCount.St8a}

\begin{proof}
(1) Since $x$ is generic, it is the irreducible product of the essentially square-integrable
representations of a multiset of pairwise unlinked segments $(\rho_i,l_i)$, and it is the Langlands
quotient of these segments \cite{Zel80}. By law~(3), $\operatorname{rc}(x)$ is the sum of the
$(\operatorname{rc}(\rho_i)\otimes(|\cdot|^{(l_i-1)/2}\circ a))\otimes\operatorname{Sp}(l_i)$, and
$\operatorname{Sp}(l)$ has nonzero monodromy for $l\geq2$, so every $l_i=1$. Then the $\rho_i$ form the
supercuspidal support, two segments of length one are linked exactly when their supercuspidals differ by
$\nu^{\pm1}$, and $\operatorname{rc}(x)=\bigoplus_i\operatorname{rc}(\rho_i)$.

(2) By law~(4), each $\operatorname{rc}(y_i)$ is irreducible, and by (1) it is a summand of
$\operatorname{rc}(x)$, so the commutators of $W_A$ act unipotently on it. An irreducible
representation of $W_A$ whose commutators act unipotently is a character, since its semisimple image
is then abelian. So $\operatorname{rc}(y_i)$ has dimension $1$, and by law~(1) $y_i$ is a character.
With (1), $x$ is plain.

(3) First let $y$ be a supercuspidal with $\operatorname{rc}(y)$ tame, of degree $n$; we show that $y$
has depth zero. Let $\varphi\in W_A$ have degree $-1$ and $c=\det\operatorname{rc}(y)(\varphi)$. The set
$\mathcal T$ of irreducible tame $V\in\mathrm{WD}(A)$ of dimension $n$ with $\det V(\varphi)=c$ injects
into the finite set $\mathcal L$ of monic irreducible polynomials of degree $n$ over $k_w$ other than
$X$, through the characteristic polynomial of a tame generator reduced through the tame character. There
is an injection $P\mapsto X_P$ from $\mathcal L$ to supercuspidals of degree $n$ and depth zero with
central character taking the value $c$ at the uniformizer $a(\varphi)$, by the theory of depth-zero
types. Each $X_P$ is good (Definition~\ref{gal-good}), so by laws~(4), (11), (1) and (8)
$\operatorname{rc}(X_P)\in\mathcal T$, and $P\mapsto\operatorname{rc}(X_P)$ is injective by law~(10).
Composing, $\mathcal L$ injects into itself through $\mathcal T$, so both maps are bijections. Since
$\operatorname{rc}(y)\in\mathcal T$, $\operatorname{rc}(y)=\operatorname{rc}(X_P)$ for some $P$, and law~(10)
gives $y=X_P$, of depth zero.

Now, by (1) every $\operatorname{rc}(y_i)$ is a summand of the tame $\operatorname{rc}(x)$, so each
$y_i$ has depth zero, and so does $x$, a constituent of a representation induced from depth-zero
supercuspidals \cite{MP96}. Suppose $y_i$ and $y_j=y_i\otimes\mu$, $i\neq j$, have degree $m\geq2$, with
$\mu$ unramified. By law~(7) for the good $y_i$, $\operatorname{rc}(y_j)$ and
$\operatorname{rc}(y_i)$ have the same characteristic polynomial at every element of $I_A$. Let $\gamma$
be a tame generator. The square of the characteristic polynomial of $\operatorname{rc}(y_i)(\gamma)$
divides that of $\operatorname{rc}(x)(\gamma)$, so every eigenvalue of $\operatorname{rc}(y_i)(\gamma)$
has multiplicity at least $2$ in $\operatorname{rc}(x)(\gamma)$ and is a $(q_w-1)$-th root of unity. A
tame irreducible representation on which $\gamma$ acts with eigenvalues of order dividing $q_w-1$ is
abelian, since the Frobenius conjugation $\gamma\mapsto\gamma^{q_w}$ then fixes $\gamma$ on it; as in
(2) it has dimension $1$, contradicting $m\geq2$.
\end{proof}

\begin{definition}\label{gal-solvable-disjoint}
Let $K$ be a number field, $\Omega$ an algebraic closure of $K$, $G_K=\operatorname{Gal}(\Omega/K)$ and
$N$ an open subgroup of $G_K$. A finite extension $M$ of $K$ inside $\Omega$ is \emph{solvable over $K$
and disjoint from $N$} if $M/K$ is Galois with solvable Galois group and every element of $G_K$ is a
product $xn$ with $x\in G_M$ and $n\in N$.
\end{definition}
\leanref{Development/AbelianVarietyTheory/Definitions/GaloisGeometry/RegPolOfRepOf.lean}{GlobalTheory.AboveL.SolvDisj}

\begin{lemma}\label{gal-solvable-cover}
Let $F$ be a number field, $\Omega$ an algebraic closure of $F$, $v$ a finite place of $F$ and $N$ an
open subgroup of $G_F$. There is a finite set $T$ of finite places of $F$, not containing $v$, such that
every $g\in G_F$ is a product $xn$ with $n\in N$ and $x$ in the stabilizer $D_{A'}\subset G_F$ of a
valuation ring $A'$ of $\Omega$ over some place of $T$.
\end{lemma}
\leanref{Development/AbelianVarietyTheory/Proofs/SatakeParameters/FiniteDecompositionGroupsCoverOpenQuotientStatement.lean}{GlobalTheory.T.finite_decompositionGroups_cover_openQuotient}
\leanref{Development/AbelianVarietyTheory/Proofs/SatakeParameters/FiniteDecompositionGroupsCoverOpenQuotient.lean}{GlobalTheory.DecCover.main}

\begin{proof}
Let $U\subset N$ be an open normal subgroup of $G_F$. For $g\in G_F$, the Chebotarev density theorem,
applied to the finite Galois extension of $F$ fixed by $U$, gives a finite place $u_g\neq v$, a
valuation ring $A_g$ over $u_g$ and $\sigma_g\in D_{A_g}$ with $\sigma_gg^{-1}\in U$. Since $U$ is
normal, $\sigma_g^{-1}g=\sigma_g^{-1}(\sigma_gg^{-1})^{-1}\sigma_g$ lies in $U\subset N$. So the open
sets $\{h\in G_F:\sigma_g^{-1}h\in N\}$ cover $G_F$. Since $G_F$ is compact, finitely many of them
cover it, say those for $g_1,\ldots,g_k$; put $T=\{u_{g_1},\ldots,u_{g_k}\}$. If
$\sigma_{g_i}^{-1}h\in N$, then $h=\sigma_{g_i}\cdot\sigma_{g_i}^{-1}h$ has the required form.
\end{proof}

\begin{lemma}\label{gal-solvable-abelian-real}
Let $F$ be a totally real number field, $\Omega$ an algebraic closure of $F$, $v$ a finite place of $F$,
$A$ a valuation ring of $\Omega$ over $v$ and $D_A\subset G_F$ its stabilizer. Let $V$ be an open
subgroup of $D_A$ that contains the commutator subgroup of $D_A$, and $T$ a finite set of finite places
of $F$ not containing $v$. Then there is a finite abelian extension $F_1$ of $F$ inside $\Omega$ that is
totally real, with $D_A\cap G_{F_1}=V$, and with $D_{A'}\subset G_{F_1}$ for every valuation ring $A'$
of $\Omega$ over a place of $T$, where $D_{A'}\subset G_F$ is the stabilizer of $A'$.
\end{lemma}
\leanref{Development/AbelianVarietyTheory/Proofs/SatakeParameters/ExistsAbelianTotallyRealExtPrescribedCompletionStatement.lean}{GlobalTheory.T.exists_abelian_totallyReal_ext_prescribedCompletion}
\leanref{Development/AbelianVarietyTheory/Proofs/SatakeParameters/ExistsAbelianTotallyRealExtPrescribedCompletion.lean}{GlobalTheory.AbStep.exists_abelian_real_field}
\leanref{Development/AbelianVarietyTheory/Proofs/SatakeParameters/ExistsAbelianTotallyRealExtPrescribedCompletion.lean}{GlobalTheory.AbStep.one_char}

\begin{proof}
The subgroup $V$ is normal of finite index in the compact group $D_A$, and $D_A/V$ is abelian. Let
$\psi$ be a character of $D_A/V$. By local class field theory, $\psi$ on the Weil group of $A$ is
$\theta\circ a$ for a character $\theta$ of $F_v^\times$ of finite order, with $a$ the local Artin map.
By the Grunwald--Wang theorem \cite[Chapter~X]{AT09}, there is a Hecke character $\eta$ of $F$ of finite
order whose component at $v$ is $\theta$ and whose components at the places of $T$ and at the infinite
places are trivial. Global class field theory turns $\eta$ into a continuous character
$r_\psi:G_F\to\mathbf C^\times$ of finite order that agrees with $\psi$ on $D_A$, is trivial on $D_{A'}$
for every $A'$ over a place of $T$, and is trivial on every complex conjugation, that is, on every
$c\in G_F$ with $j\circ c=\overline{j}$ for some embedding $j:\Omega\to\mathbf C$. Let $F_1$ be the
fixed field of the intersection of the kernels of the $r_\psi$, a finite abelian extension of $F$. An
element of $D_A$ fixes $F_1$ exactly when every $\psi$ is trivial on it, that is, when it lies in $V$,
since the characters of the finite abelian group $D_A/V$ separate its points. The complex conjugations
fix $F_1$, so $F_1$ is totally real, and $D_{A'}$ fixes $F_1$ for every $A'$ over a place of $T$.
\end{proof}

\begin{lemma}\label{gal-solvable-compositum}
Let $E$ be a CM field with maximal totally real subfield $F$, $\Omega$ an algebraic closure of $E$, $v$ a
finite place of $F$ and $N$ an open subgroup of $G_E$. There is a finite set $T$ of finite places of
$F$, not containing $v$, with the following property. Let $F_1\subset\Omega$ be a finite Galois
extension of $F$ with solvable Galois group that is totally real and such that $D_{A'}\subset G_{F_1}$
for every valuation ring $A'$ of $\Omega$ over a place of $T$, where $D_{A'}\subset G_F$ is the
stabilizer of $A'$. Then the compositum $M=EF_1$ is a CM field, solvable over $E$ and disjoint from $N$
(Definition~\ref{gal-solvable-disjoint}); an element of $G_E$ fixes $M$ if and only if it fixes $F_1$;
and $F_1$ is contained in the maximal totally real subfield of $M$.
\end{lemma}
\leanref{Development/AbelianVarietyTheory/Proofs/SatakeParameters/ExistsCmFieldLiftDisjointOfRealSolvableStatement.lean}{GlobalTheory.T.exists_cmField_lift_disjoint_of_real_solvable}
\leanref{Development/AbelianVarietyTheory/Proofs/SatakeParameters/ExistsCmFieldLiftDisjointOfRealSolvable.lean}{GlobalTheory.KLift.exists_cm_lift}
\leanref{Development/AbelianVarietyTheory/Proofs/SatakeParameters/CmSubfieldsCompositumClosedStatement.lean}{GlobalTheory.T.cm_subfields_compositum_closed}

\begin{proof}
The image of $N$ in $G_F$ is open; let $T$ be the set that Lemma~\ref{gal-solvable-cover} gives for it
and $v$. Let $F_1$ be as in the statement and $M=EF_1$. An element of $G_E$ fixes $M$ if and only if it
fixes $F_1$. Since $F_1/F$ is normal, $M/E$ is a finite Galois extension, and restriction to $F_1$
embeds $\operatorname{Gal}(M/E)$ into $\operatorname{Gal}(F_1/F)$, so $\operatorname{Gal}(M/E)$ is
solvable. The field $M$ is generated over the CM field $E$ by the elements of $F_1$, which are real under
every embedding into $\mathbf C$; so $M$ is a CM field, and $F_1$ lies in its maximal totally real
subfield. Let $g\in G_E$. By the choice of $T$, $g=xn$ in $G_F$ with $n\in N$ and $x\in D_{A'}$ for some
$A'$ over a place of $T$. Then $x=gn^{-1}$ lies in $G_E$ and fixes $F_1$, hence $M$, so $x\in G_M$.
\end{proof}

\begin{lemma}\label{gal-solvable-cm-extensions}
Let $E$ be a CM field with maximal totally real subfield $F$, $\Omega$ an algebraic closure of $E$, $A$ a
valuation ring of $\Omega$ over a finite place $w$ of $E$, and $N$ an open subgroup of $G_E$. For a
subfield $K$ of $\Omega$ of finite degree over $\mathbf Q$ and a valuation ring $A'$ of $\Omega$, let
$D_{A'}(K)\subset G_K$ be the stabilizer of $A'$.
\begin{enumerate}
\item There is a CM field $M_0\subset\Omega$ containing $E$, solvable over $E$ and disjoint from $N$
  (Definition~\ref{gal-solvable-disjoint}), such that $D_A(E)$ and $D_A(M_0^+)$ fix $M_0$, where $M_0^+$
  is the maximal totally real subfield of $M_0$.
\item Suppose that $D_A(F)$ fixes $E$, and let $V$ be an open subgroup of $D_A(E)$ that contains the
  commutator subgroup of $D_A(E)$. There is a CM field $M\subset\Omega$ containing $E$, solvable over $E$
  and disjoint from $N$, with $D_A(E)\cap G_M=V$.
\end{enumerate}
\end{lemma}
\leanref{Development/AbelianVarietyTheory/Proofs/SatakeParameters/ExistsSolvableCmExtensionSplitPlaceStatement.lean}{GlobalTheory.T.exists_solvable_cm_extension_split_place}
\leanref{Development/AbelianVarietyTheory/Proofs/SatakeParameters/ExistsSolvableCmExtensionSplitPlace.lean}{GlobalTheory.AboveL.exists_cm_solvable_split}
\leanref{Development/AbelianVarietyTheory/Proofs/SatakeParameters/ExistsSolvableCmExtensionAbelianCompletionStatement.lean}{GlobalTheory.T.exists_solvable_cm_extension_abelianCompletion}
\leanref{Development/AbelianVarietyTheory/Proofs/SatakeParameters/ExistsSolvableCmExtensionAbelianCompletion.lean}{GlobalTheory.AboveL.exists_cm_solvDisj_dec_eq}

\begin{proof}
Let $v$ be the place of $F$ under $w$, and $T$ the set that Lemma~\ref{gal-solvable-compositum} gives
for $E$, $v$ and $N$.

(1) The subgroup $D_A(F)\cap G_E$ of $D_A(F)$ is open, and it contains the commutator subgroup of
$D_A(F)$ since $\operatorname{Gal}(E/F)$ is abelian. Lemma~\ref{gal-solvable-abelian-real} gives a
finite abelian, hence solvable, totally real extension $F_1$ of $F$ with $D_A(F)\cap
G_{F_1}=D_A(F)\cap G_E$, fixed by $D_{A'}(F)$ for every $A'$ over a place of $T$. By
Lemma~\ref{gal-solvable-compositum}, $M_0=EF_1$ is a CM field, solvable over $E$ and disjoint from $N$.
An element of $D_A(E)$ lies in $D_A(F)\cap G_E$, so it fixes $F_1$, hence $M_0$. An element of
$D_A(M_0^+)$ fixes $F_1\subset M_0^+$, so it lies in $D_A(F)\cap G_{F_1}\subset G_E$; as an element of
$G_E$ that fixes $F_1$, it fixes $M_0$.

(2) Since $D_A(F)$ fixes $E$, $D_A(F)=D_A(E)$, and $V$ is an open subgroup of $D_A(F)$ that contains
its commutator subgroup. Lemma~\ref{gal-solvable-abelian-real} gives a finite abelian totally real
extension $F_1$ of $F$ with $D_A(F)\cap G_{F_1}=V$, fixed by $D_{A'}(F)$ for every $A'$ over a place of
$T$. By Lemma~\ref{gal-solvable-compositum}, $M=EF_1$ is a CM field, solvable over $E$ and disjoint from
$N$, and an element of $D_A(E)$ fixes $M$ exactly when it fixes $F_1$, that is, when it lies in $V$.
\end{proof}

\begin{lemma}\label{gal-solvable-character-kernel}
In Situation~\ref{gal-local-situation}, let $a:W_A\to E_w^\times$ be the local Artin map
(Definition~\ref{gal-local-artin}), $\xi_1,\ldots,\xi_n$ smooth characters of $E_w^\times$ and $g\in W_A$
with $d(g)\neq0$. There is an open subgroup $V$ of $D_A$ that contains $g$ and the commutator subgroup
of $D_A$, such that $\xi_j(a(t))=1$ for every $j$ and every $t\in I_A\cap V$.
\end{lemma}
\leanref{Development/AbelianVarietyTheory/Proofs/SatakeParameters/ExistsOpenSubgroupWeilElementCharacterKernelStatement.lean}{GlobalTheory.T.exists_open_subgroup_weilElement_characterKernel}
\leanref{Development/AbelianVarietyTheory/Proofs/SatakeParameters/ExistsOpenSubgroupWeilElementCharacterKernel.lean}{GlobalTheory.RoadU.exists_open_sub_charKer}
\leanref{Development/AbelianVarietyTheory/Proofs/SatakeParameters/EqOneOfCharpoly.lean}{GlobalTheory.AboveL.AbSub.exists_open_sub}

\begin{proof}
The characters $\chi_j=\xi_j\circ a$ of $W_A$ are trivial on an open subgroup of $I_A$, since $a$ is
continuous on $I_A$ and the $\xi_j$ are smooth. Let $J$ be the subgroup of the $t\in I_A$ with
$\chi_j(t)=1$ for all $j$. There is an open normal subgroup $U_0$ of $D_A$ with $U_0\cap I_A\subset J$,
and since $U_0$ has finite index there is $m\geq1$ with $g^m\in U_0$. Let $\zeta\in\Omega$ be a root of
unity of order $q_w^{m|d(g)|}-1$; an element of $W_A$ fixes $\zeta$ only if its degree is divisible by
$m|d(g)|$. Let $U\subset U_0$ be an open normal subgroup of $D_A$ that fixes $\zeta$.

The subgroup $J$ is normal in $D_A$. Indeed, $W_A$ normalizes $J$, since the $\chi_j$ are characters;
for $u\in U$ and $t\in J$ the commutator $utu^{-1}t^{-1}$ lies in $U_0\cap I_A\subset J$, so $U$
normalizes $J$; and $D_A=W_AU$, since $W_A$ is dense in $D_A$. Put $V=\langle g\rangle JU$. It is a
subgroup since $JU$ is normal, it is open since it contains $U$, and it contains $g$. Modulo the normal
subgroup $JU$, a commutator of two elements of $D_A=W_AU$ is a commutator of two elements of $W_A$,
which has degree $0$ and on which every $\chi_j$ is trivial; so the commutator subgroup of $D_A$ lies in
$JU\subset V$.

Let $t\in I_A\cap V$, and write $t=g^kju$ with $j\in J$ and $u\in U$. Then $u=(g^kj)^{-1}t$ lies in
$W_A$, has degree $-kd(g)$ and fixes $\zeta$, so $m|d(g)|$ divides $kd(g)$ and $m$ divides $k$. Hence
$g^k\in U_0$, and $g^ku=(g^kjg^{-k})^{-1}t$ lies in $U_0\cap I_A\subset J$. So
$t=(g^kjg^{-k})(g^ku)$ lies in $J$, that is, $\xi_j(a(t))=1$ for all $j$.
\end{proof}

\begin{definition}\label{gal-solvable-base-change}
Let $K$ be a totally complex number field, $M$ a totally complex finite extension of $K$, $a\geq1$, and
$\pi$ and $P$ cuspidal automorphic representations of $\operatorname{GL}_a(\mathbf A_K)$ and
$\operatorname{GL}_a(\mathbf A_M)$ (Definition~\ref{gal-polarized}).
\begin{enumerate}
\item $P$ is a \emph{base change} of $\pi$ if:
  \begin{enumerate}
  \item for every finite place $u$ of $M$, with $v$ the place of $K$ under $u$ and $N(u)=N(v)^f$, if
    $\pi$ has Satake parameters $B$ at $v$ then $P$ has Satake parameters $\{b^f:b\in B\}$ at $u$;
  \item $Z_\tau(P)=Z_{\tau|_K}(\pi)$ for every embedding $\tau:M\to\mathbf C$;
  \item for every finite place $v$ of $K$, finite place $u$ of $M$ and continuous field isomorphism
    $i:K_v\to M_u$ extending $K\to M$, if $\sigma$ and $\sigma'$ are local components of $\pi$ at $v$
    and of $P$ at $u$, then $\sigma'\circ i\cong\sigma$.
  \end{enumerate}
\item A base change $P$ of $\pi$ \emph{carries character supports} if for every finite place $v$ of
  $K$, place $u$ of $M$ over $v$, local components $\sigma$ of $\pi$ at $v$ and $\sigma'$ of $P$ at $u$,
  and smooth characters $\xi_1,\ldots,\xi_a$ of $K_v^\times$ that form the supercuspidal support of
  $\sigma$, the characters $\xi_j\circ N_{M_u/K_v}$ form the supercuspidal support of $\sigma'$.
\end{enumerate}
\end{definition}
\leanref{Development/AbelianVarietyTheory/Definitions/GaloisGeometry/IsWeight.lean}{GlobalTheory.GalCut.IsLift}
\leanref{Development/AbelianVarietyTheory/Definitions/GaloisGeometry/CyclicStrongAtCusp.lean}{GlobalTheory.RoadU.CharLift}

\begin{lemma}\label{gal-solvable-local-base-change}
Let $K$ be a totally complex number field, $M$ a totally complex Galois extension of $K$ of prime
degree, $a\geq1$, $\pi$ a cuspidal automorphic representation of $\operatorname{GL}_a(\mathbf A_K)$ and
$P$ a base change of $\pi$ (Definition~\ref{gal-solvable-base-change}). Then $P$ carries character
supports. This follows from the cyclic base change of Arthur and Clozel, whose lift has at every place
the local base change of the local component of $\pi$, from the compatibility of local base change with
parabolic induction \cite{AC89}, and from strong multiplicity one \cite{JS81a}, which identifies $P$
with that lift.
\end{lemma}
\leanref{Development/AutomorphicArithmetic/Foundations/GaloisRepresentations/SphericalGaloisCompatibility/Chl.lean}{UpperTheory.Gal.ToralGalois.Inq}
\leanref{Development/AutomorphicArithmetic/Foundations/GaloisRepresentations/SphericalGaloisCompatibility/GaloisCompatUnlinkedSupportOfSpherical.lean}{UpperTheory.Gal.ToralGalois.inq_holds}
\leanref{Development/AutomorphicArithmetic/Foundations/GaloisRepresentations/MixedUnlinkedExponents/CharLiftPrimeHolds.lean}{UpperTheory.Gal.GamX.charLiftPrime_holds}
\leanref{Development/AutomorphicArithmetic/Foundations/GaloisRepresentations/MixedUnlinkedExponents/CharLiftPrimeOf.lean}{UpperTheory.Gal.GamX.CharSupp}

\begin{lemma}\label{gal-solvable-base-change-facts}
Let $K\subset M\subset L$ be CM fields inside an algebraic closure $\Omega$ of $K$, with maximal totally
real subfields $K^+$, $M^+$ and $L^+$, let $a\geq1$, and let $\pi$, $P$ and $Q$ be cuspidal automorphic
representations of $\operatorname{GL}_a(\mathbf A_K)$, $\operatorname{GL}_a(\mathbf A_M)$ and
$\operatorname{GL}_a(\mathbf A_L)$. Base changes are those of Definition~\ref{gal-solvable-base-change}.
\begin{enumerate}
\item If $P$ is a base change of $\pi$ and $(\pi,\chi)$ is in scope (Definition~\ref{gal-scope}), then
  $(P,\chi\circ N_{M^+/K^+})$ is in scope.
\item If $P$ is a base change of $\pi$, $\ell$ is a prime, $\iota:\bar{\mathbf Q}_\ell\to\mathbf C$ is a
  field isomorphism and $r:G_K\to\operatorname{GL}_a(\bar{\mathbf Q}_\ell)$ is weakly attached to $\pi$
  (Definition~\ref{gal-weakly-attached}), then $r|_{G_M}$ is weakly attached to $P$.
\item If $P$ is a base change of $\pi$ and $Q$ a base change of $P$, both carrying character supports,
  then $Q$ is a base change of $\pi$ that carries character supports.
\end{enumerate}
\end{lemma}
\leanref{Development/AbelianVarietyTheory/Proofs/SatakeParameters/PolarizedPairBaseChangePolarizedSameWeightStatement.lean}{GlobalTheory.T.polarizedPair_baseChange_polarized_same_weight}
\leanref{Development/AbelianVarietyTheory/Proofs/SatakeParameters/PolarizedPairBaseChangePolarizedSameWeight.lean}{GlobalTheory.LiftPolar.main}
\leanref{Development/AbelianVarietyTheory/Proofs/SatakeParameters/WeaklyAttachedGaloisRepresentationBaseChangeRestStatement.lean}{GlobalTheory.T.weaklyAttachedGaloisRepresentation_baseChange_restrict}
\leanref{Development/AbelianVarietyTheory/Proofs/SatakeParameters/WeaklyAttachedGaloisRepresentationBaseChangeRest.lean}{GlobalTheory.GalCut.weak_comp_res}
\leanref{Development/AbelianVarietyTheory/Proofs/SatakeParameters/CuspRepCharacterSupportLiftTransTowerStatement.lean}{GlobalTheory.T.cuspRep_characterSupportLift_trans_tower}
\leanref{Development/AbelianVarietyTheory/Proofs/SatakeParameters/CuspRepCharacterSupportLiftTransTower.lean}{GlobalTheory.RoadU.charLift_trans}

\begin{proof}
(1) By condition~(b) of Definition~\ref{gal-solvable-base-change}, $Z_\tau(P)=Z_{\tau|_K}(\pi)$ consists
of $a$ distinct elements of $\frac{a-1}{2}+\mathbf Z$, so $P$ is regular algebraic. Put
$\chi'=\chi\circ N_{M^+/K^+}$, $\mu=\chi\circ N_{K/K^+}$ and $\mu'=\chi'\circ N_{M/M^+}$. At a real place
of $M^+$ over the real place $v$ of $K^+$, $\chi'(-1)=\chi_v(-1)=(-1)^a$. Let $u$ be a place of $M$,
outside a finite set, with $v$ the place of $K$ under $u$ and $N(u)=N(v)^f$. The polarization of $\pi$
shows that its Satake parameters at $c(v)$ are $\{\mu(\varpi_v)b^{-1}:b\in B\}$, with $B$ those at $v$.
By condition~(a), the Satake parameters of $P$ at $u$ and at $c(u)$ are the $f$-th powers of these, and
$\mu'(\varpi_u)=\mu(\varpi_v)^f$. So $P^c$ and $P^\vee\otimes(\mu'\circ\det)$ have the same Satake
parameters at all but finitely many places, and by strong multiplicity one \cite{JS81a} and the theorem
of Gelfand and Kazhdan on contragredients, $(P,\chi')$ is polarized.

(2) The restriction $r|_{G_M}$ is continuous. It is semisimple: by Clifford's theorem the restriction of
$r$ to the largest normal subgroup of $G_K$ contained in $G_M$ is semisimple, this subgroup has finite
index in $G_M$, and in characteristic $0$ a representation whose restriction to a normal subgroup of
finite index is semisimple is itself semisimple. At all but finitely many places $u$ of $M$, with $v$
and $f$ as in (1), $r$ has Frobenius eigenvalues $q_v^{(a-1)/2}B$ at $v$, with $B$ the Satake
parameters of $\pi$ at $v$. An element of $G_M$ that acts on the residue field of $u$ by $z\mapsto
z^{q_u}$ acts on that of $v$ by $z\mapsto z^{q_v^f}$, so $r|_{G_M}$ has Frobenius eigenvalues
$q_u^{(a-1)/2}\{b^f:b\in B\}$ at $u$, and $\{b^f:b\in B\}$ are the Satake parameters of $P$ at $u$ by
condition~(a).

(3) For places $v$, $u$, $x$ of $K$, $M$, $L$ with $x$ over $u$ over $v$, the residue degree of $x$ over
$v$ is the product of those of $x$ over $u$ and of $u$ over $v$, which gives condition~(a), and
condition~(b) composes directly. A continuous isomorphism $i:K_v\to L_x$ extending $K\to L$ has image
$L_x$, which contains the closure $M_u$ of $M$ in $L_x$, with $u$ the place of $M$ under $x$; so $i$ is
the composite of continuous isomorphisms $K_v\to M_u$ and $M_u\to L_x$ extending $K\to M$ and $M\to L$,
which gives condition~(c). The supports compose since $N_{L_x/K_v}=N_{M_u/K_v}\circ N_{L_x/M_u}$.
\end{proof}

\begin{lemma}\label{gal-solvable-self-twist}
Let $K$ be a CM field, $\Omega$ an algebraic closure of $K$, and $M\subset\Omega$ a CM field that is a
Galois extension of $K$ of prime degree $p$. Let $\ell$ be a prime, $\iota:\bar{\mathbf
Q}_\ell\to\mathbf C$ a field isomorphism, $a\geq1$, $\pi$ a cuspidal automorphic representation of
$\operatorname{GL}_a(\mathbf A_K)$ and $r:G_K\to\operatorname{GL}_a(\bar{\mathbf Q}_\ell)$ weakly
attached to $\pi$ (Definition~\ref{gal-weakly-attached}). Let $\eta$ be a Hecke character of $K$ with
$\eta\neq1$ and $\eta^p=1$, such that $\eta(\varpi_v)^f=1$ for all but finitely many places $u$ of $M$,
with $v$ the place of $K$ under $u$ and $N(u)=N(v)^f$, and such that for all but finitely many places
$v$ of $K$, $\eta$ is unramified at $v$ and $\pi$ has Satake parameters $B$ at $v$ with
$\{\eta(\varpi_v)b:b\in B\}=B$. Then the trace of $r(g)$ is $0$ for every $g\in G_K$ that does not fix
$M$.
\end{lemma}
\leanref{Development/AbelianVarietyTheory/Proofs/SatakeParameters/TraceEqZeroOffSubgroupOfSelfTwistStatement.lean}{GlobalTheory.T.trace_eq_zero_off_subgroup_of_selfTwist}
\leanref{Development/AbelianVarietyTheory/Proofs/SatakeParameters/TraceEqZeroOffSubgroupOfSelfTwist.lean}{GlobalTheory.SelfTwWeak.trace_eq_zero_of_selfTw}

\begin{proof}
By global class field theory there is a continuous character $\chi:G_K\to\mathbf C^\times$ of finite
order, nontrivial since $\eta\neq1$, such that for all but finitely many places $v$ of $K$, every
valuation ring $A$ over $v$ and every $\sigma\in D_A$ acting on the residue field by $z\mapsto z^{q_v}$,
$\chi(\sigma)=\eta(\varpi_v)$. For all but finitely many places $u$ of $M$, with $v$ and $f$ as in the
statement, an element of $G_M$ acting on the residue field of $u$ by $z\mapsto z^{q_u}$ acts on that of
$v$ by $z\mapsto z^{q_v^f}$, so $\chi$, unramified at $v$, takes on it the value $\eta(\varpi_v)^f=1$.
These elements are
dense in $G_M$ by the Chebotarev density theorem, so $\chi$ is trivial on $G_M$. Thus $\chi$ is a
nontrivial character of $\operatorname{Gal}(M/K)$, a group of prime order, so it is injective there, and
$\chi(g^{-1})\neq1$ for every $g\in G_K$ that does not fix $M$.

For all but finitely many places $v$ of $K$ and $\sigma$ as above, $r$ has Frobenius eigenvalues
$q_v^{(a-1)/2}B$ at $v$ with $\{\eta(\varpi_v)b:b\in B\}=B$, so $\eta(\varpi_v)\,\iota(\operatorname{tr}
r(\sigma^{-1}))=\iota(\operatorname{tr}r(\sigma^{-1}))$. The set $X$ of $h\in G_K$ with $\chi(h)\neq1$ is
open, and the elements $\sigma$ at places outside a finite set are dense in $G_K$ by the Chebotarev
density theorem, so they are dense in $X$. At such a $\sigma$ in $X$, $\eta(\varpi_v)=\chi(\sigma)\neq1$,
so $\operatorname{tr}r(\sigma^{-1})=0$. Since $h\mapsto\operatorname{tr}r(h^{-1})$ is continuous, it
vanishes on $X$, which contains $g^{-1}$; so $\operatorname{tr}r(g)=0$.
\end{proof}

\begin{lemma}\label{gal-solvable-cuspidal-lift}
Let $K$ be a CM field, $\Omega$ an algebraic closure of $K$, $\ell$ a prime, $\iota:\bar{\mathbf
Q}_\ell\to\mathbf C$ a field isomorphism, $a\geq1$, $(\pi,\chi)$ in scope (Definition~\ref{gal-scope})
with $\pi$ a cuspidal automorphic representation of $\operatorname{GL}_a(\mathbf A_K)$, and
$r:G_K\to\operatorname{GL}_a(\bar{\mathbf Q}_\ell)$ weakly attached to $\pi$
(Definition~\ref{gal-weakly-attached}). Let $N$ be an open subgroup of $G_K$ such that the trace of
$r(n)$ is nonzero for every $n\in N$, and $M\subset\Omega$ a CM field containing $K$, solvable over $K$
and disjoint from $N$ (Definition~\ref{gal-solvable-disjoint}). Then there are a cuspidal automorphic
representation $P$ of $\operatorname{GL}_a(\mathbf A_M)$ and a Hecke character $\chi'$ of the maximal
totally real subfield of $M$ such that $(P,\chi')$ is in scope and $P$ is a base change of $\pi$ that
carries character supports (Definition~\ref{gal-solvable-base-change}).
\end{lemma}
\leanref{Development/AutomorphicArithmetic/Foundations/GaloisRepresentations/SphericalGaloisCompatibility/ToralDeRham/ExistsCuspidalCharLiftOfSolvableCmExtension.lean}{GalTw.GlobalTheory.T.exists_cuspidal_charLift_of_solvable_cmExtension}
\leanref{Development/AutomorphicArithmetic/Foundations/GaloisRepresentations/SphericalGaloisCompatibility/ToralDeRham/ExistsCuspidalCharLiftOfSolvableCmExtension.lean}{GalTw.GlobalTheory.RoadU.upWeak}
\leanref{Development/AutomorphicArithmetic/Foundations/GaloisRepresentations/SphericalGaloisCompatibility/ToralDeRham/ExistsCuspidalCharLiftOfSolvableCmExtension.lean}{GalTw.GlobalTheory.RoadU.isWeaklyAttached_of_cmBaseChange_of_finrank}
\leanref{Development/AbelianVarietyTheory/Proofs/SatakeParameters/ExistsCharacterSupportBaseChangeLiftDegreeOneStatement.lean}{GlobalTheory.T.exists_characterSupportBaseChangeLift_degreeOne}
\leanref{Development/AbelianVarietyTheory/Proofs/SatakeParameters/SolvableAscentInduction.lean}{GlobalTheory.SolvUp.exists_normal_prime_index}

\begin{proof}
Induct on $d=[M:K]$. If $d=1$, then $M=K$, and $P=\pi$ and $\chi'=\chi$ satisfy the conclusion. Let
$d>1$. The group $\operatorname{Gal}(M/K)$ is finite, solvable and nontrivial, so it has a normal
subgroup $H$ of prime index $p$. Its fixed field $K_1$ is a Galois extension of $K$ of degree $p$, and
a CM field, as a field between the CM fields $K$ and $M$.

Some $n\in N$ does not fix $K_1$. Indeed, let $x\in K_1$ not lie in $K$. If every element of $N$ fixed
$x$, so would every element $yn$ of $G_K$ with $y\in G_M$ and $n\in N$, since $G_M$ fixes $K_1$, and
then $x$ would lie in $K$. The trace of $r(n)$ is nonzero, so by Lemma~\ref{gal-solvable-self-twist} no
Hecke character $\eta$ of $K$ has the properties excluded in Theorem~\ref{gal-cyclic-base-change} for
$K_1/K$ and $\pi$. That theorem gives a cuspidal automorphic representation $P_1$ of
$\operatorname{GL}_a(\mathbf A_{K_1})$ that is a base change of $\pi$, and $P_1$ carries character
supports by Lemma~\ref{gal-solvable-local-base-change}. By
Lemma~\ref{gal-solvable-base-change-facts}(1) and~(2), $(P_1,\chi\circ N_{K_1^+/K^+})$ is in scope and
$r|_{G_{K_1}}$ is weakly attached to $P_1$.

Let $N_1=N\cap G_{K_1}$, an open subgroup of $G_{K_1}$ on which the trace of $r$ does not vanish. The
field $M$ is Galois over $K_1$ with solvable group $H$, and it is disjoint from $N_1$: an element
$g_1\in G_{K_1}$ is a product $yn$ with $y\in G_M$ and $n\in N$, and $n=y^{-1}g_1$ fixes $K_1$, so
$n\in N_1$. Since $[M:K_1]<d$, the induction hypothesis gives $P$ and $\chi'$ over $M$ with $(P,\chi')$
in scope and $P$ a base change of $P_1$ that carries character supports. By
Lemma~\ref{gal-solvable-base-change-facts}(3), $P$ is a base change of $\pi$ that carries character
supports.
\end{proof}

\begin{lemma}\label{gal-solvable-lift}
Let $E$ be a CM field, $\Omega$ an algebraic closure of $E$, $\ell$ a prime, $\iota:\bar{\mathbf
Q}_\ell\to\mathbf C$ a field isomorphism, $n\geq1$, $(\pi,\chi)$ in scope (Definition~\ref{gal-scope})
with $\pi$ a cuspidal automorphic representation of $\operatorname{GL}_n(\mathbf A_E)$, and
$r:G_E\to\operatorname{GL}_n(\bar{\mathbf Q}_\ell)$ weakly attached to $\pi$
(Definition~\ref{gal-weakly-attached}). Let $w$ be a finite place of $E$, $A$ a valuation ring of
$\Omega$ over $w$, $a:W_A\to E_w^\times$ the local Artin map (Definition~\ref{gal-local-artin}),
$\xi_1,\ldots,\xi_n$ smooth characters of $E_w^\times$ and $g\in W_A$ with $d(g)\neq0$. Then there are
a CM field $M\subset\Omega$ containing $E$, a cuspidal automorphic representation $P$ of
$\operatorname{GL}_n(\mathbf A_M)$ and a Hecke character $\chi'$ of the maximal totally real subfield of
$M$ such that $(P,\chi')$ is in scope, $P$ is a base change of $\pi$ that carries character supports
(Definition~\ref{gal-solvable-base-change}), $g$ fixes $M$, and $\xi_j(a(t))=1$ for every $j$ and every
$t\in I_A$ that fixes $M$.
\end{lemma}
\leanref{Development/AutomorphicArithmetic/Foundations/GaloisRepresentations/SphericalGaloisCompatibility/ToralDeRham/WeaklyAttachedToralOfCmLift.lean}{GalTw.GlobalTheory.T.weaklyAttached_toral_of_cm_lift}
\leanref{Development/AutomorphicArithmetic/Foundations/GaloisRepresentations/SphericalGaloisCompatibility/ToralDeRham/WeaklyAttachedToralOfCmLift.lean}{GalTw.GlobalTheory.RoadU.TorLiftC.exists_cm_lift_fix_charKer}
\leanref{Development/AbelianVarietyTheory/Proofs/SatakeParameters/ExistsOpenGaloisSubgroupTraceNeZeroStatement.lean}{GlobalTheory.T.exists_open_galoisSubgroup_trace_ne_zero}
\leanref{Development/AbelianVarietyTheory/Proofs/SatakeParameters/LocalArtinCompNormCompatStatement.lean}{GlobalTheory.T.localArtin_comp_norm_compat}

\begin{proof}
For a subfield $K$ of $\Omega$ finite over $E$, let $D_A(K)\subset G_K$ be the stabilizer of $A$. Since
$r$ is continuous and the trace of $r(1)$ is $n\neq0$, there is an open subgroup $N_0$ of $G_E$ on which
the trace of $r$ does not vanish. Lemma~\ref{gal-solvable-cm-extensions}(1) gives a CM field $M_0$,
solvable over $E$ and disjoint from $N_0$, such that $D_A(E)$ and $D_A(M_0^+)$ fix $M_0$. By
Lemma~\ref{gal-solvable-cuspidal-lift} there are $\pi_0$ and $\chi_0$ over $M_0$ with $(\pi_0,\chi_0)$
in scope and $\pi_0$ a base change of $\pi$ that carries character supports, and $r|_{G_{M_0}}$ is
weakly attached to $\pi_0$ by Lemma~\ref{gal-solvable-base-change-facts}(2).

Let $v_0$ be the place of $M_0$ under $A$ and $a_0$ the local Artin map of $M_0$ at $A$. Since $g$
fixes $M_0$, it lies in the Weil group of $A$ over $M_0$, with nonzero degree there.
Lemma~\ref{gal-solvable-character-kernel}, applied over $M_0$ to the smooth characters
$\xi_j\circ N_{M_{0,v_0}/E_w}$, gives an open subgroup $V$ of $D_A(M_0)$ that contains $g$ and the
commutator subgroup of $D_A(M_0)$, such that $\xi_j(N_{M_{0,v_0}/E_w}(a_0(t)))=1$ for every $t$ in $V$
and in the inertia group of $A$ over $M_0$. Let $N_1$ be an open subgroup of $G_{M_0}$ on which the
trace of $r$ does not vanish. Since $D_A(M_0^+)$ fixes $M_0$, Lemma~\ref{gal-solvable-cm-extensions}(2)
applies to $M_0$, $V$ and $N_1$, and gives a CM field $M$, solvable over $M_0$ and disjoint from $N_1$,
with $D_A(M_0)\cap G_M=V$. By Lemma~\ref{gal-solvable-cuspidal-lift} there are $P$ and $\chi'$ over $M$
with $(P,\chi')$ in scope and $P$ a base change of $\pi_0$ that carries character supports, and by
Lemma~\ref{gal-solvable-base-change-facts}(3), $P$ is a base change of $\pi$ that carries character
supports.

Since $g\in V$, $g$ fixes $M$. Let $t\in I_A$ fix $M$. Then $t$ fixes $M_0$, lies in the inertia group
of $A$ over $M_0$, and lies in $V$ since it fixes $M$; so $\xi_j(N_{M_{0,v_0}/E_w}(a_0(t)))=1$. By the
compatibility of local Artin maps with the norm \cite{Ser67}, $a(t)=N_{M_{0,v_0}/E_w}(a_0(t))$, so
$\xi_j(a(t))=1$.
\end{proof}

\begin{proposition}\label{gal-plain-compatibility}
Suppose existence and the four conditions hold in the first scope (Definition~\ref{gal-scope}). Let
$E$, $\Omega$, $\ell$, $\iota$, $n\geq1$ and $(\pi,\chi)$ be in scope, let $r$ be weakly attached to
$\pi$, and let $\sigma$ be a local component of $\pi$ at a finite place $w$ with $\sigma^\flat$ plain
(Definition~\ref{gal-plain}). Let $A$ be a valuation ring over $w$.
\begin{enumerate}
\item If $w\nmid\ell$, $\mathrm{WD}_\iota(r|_{W_A})$ is monodromy-free and equals
  $\operatorname{rc}(\sigma^\flat)$ for every map $\operatorname{rc}$ that satisfies the laws
  (Definition~\ref{gal-llc-laws}).
\item If $w\mid\ell$, $r$ satisfies (dR) for $\pi$ at $w$, $r|_{D_A}$ is potentially semistable,
  $\mathrm{WD}_\iota(r|_{D_A})$ is monodromy-free, and for every $\operatorname{rc}$ that satisfies the
  laws and every $g\in W_A$, $\mathrm{WD}_\iota(r|_{D_A})(g)$ and $\operatorname{rc}(\sigma^\flat)(g)$
  have the same characteristic polynomial.
\end{enumerate}
\end{proposition}
\leanref{Development/AbelianVarietyTheory/Statements/EllipticCurveLevels/IsQg.lean}{GlobalTheory.DeRhamOfPlainTorus.R1}
\leanref{Development/AutomorphicArithmetic/Foundations/GaloisRepresentations/SphericalGaloisCompatibility/GaloisCompatUnlinkedSupportOfSpherical.lean}{UpperTheory.Gal.ToralGalois.galoisCompat_unlinkedSupport_of_spherical}
\leanref{Development/AutomorphicArithmetic/Foundations/GaloisRepresentations/SphericalGaloisCompatibility/ToralDeRham/GaloisCompatUnlinkedSupportOfSpherical.lean}{GalTw.GlobalTheory.DeRhamOfPlainTorus.galoisCompat_unlinkedSupport_of_spherical}

\begin{proof}
Write the support of $\sigma^\flat$ as characters $\xi_1,\ldots,\xi_n$ of $E_w^\times$. Since they are
unlinked, $\sigma^\flat$ is the full induced representation $\xi_1\times\cdots\times\xi_n$, and by
laws~(2) and~(3), $\operatorname{rc}(\sigma^\flat)=\bigoplus_j\xi_j\circ a$ for every
$\operatorname{rc}$ satisfying the laws, $a$ its local Artin map.

\emph{The lift.} Let $g\in W_A$ have nonzero degree. Lemma~\ref{gal-solvable-lift}, applied to the
characters $\xi_j\nu^{(n-1)/2}$, gives a CM field $M\subset\Omega$ containing $E$ and fixed by $g$, and
$(P,\chi')$ in scope over $M$ with $P$ a base change of $\pi$ that carries character supports
(Definition~\ref{gal-solvable-base-change}), such that $(\xi_j\nu^{(n-1)/2})(a(t))=1$ for every $j$ and
every $t\in I_A$ that fixes $M$. Let $u$ be the place of $M$ under $A$. Every element $t$ of the inertia
group of $A$ over $M$ lies in $I_A$ and fixes $M$, so by the compatibility of local Artin maps with the
norm from $M_u$ to $E_w$ the characters $\xi_j\nu^{(n-1)/2}\circ N_{M_u/E_w}$ are unramified. The
supercuspidal support of $\sigma=\sigma^\flat\nu^{(n-1)/2}$ consists of the $\xi_j\nu^{(n-1)/2}$, and
$P$ carries character supports, so the supercuspidal support of a local component of $P$ at $u$
consists of the $\xi_j\nu^{(n-1)/2}\circ N_{M_u/E_w}$. These characters are pairwise unlinked. Indeed,
since $(\pi,\chi)$ is in scope there is an integer $m$ with $Z_{\tau\circ c}(\pi)=\{n-1-m-z:z\in
Z_\tau(\pi)\}$ for every $\tau$, so the central character of $\pi\otimes|\det|^{-(n-1-m)/2}$ is unitary;
this cuspidal representation is then unitary, and its local component $\sigma^\flat\nu^{m/2}$ at $w$ is
unitarizable. It is generic, with unlinked supercuspidal support the characters $\xi_j\nu^{m/2}$, so by
Tadić's classification of the generic unitary dual \cite{Tad86}, $|\xi_j(x)|=|x|_{E_w}^{e_j-m/2}$ with
$|e_j|<\frac12$. Since $|N_{M_u/E_w}(y)|_{E_w}=|y|_{M_u}$, the quotient of two of the characters
$\xi_j\nu^{(n-1)/2}\circ N_{M_u/E_w}$ has absolute value $|\cdot|_{M_u}^{e}$ with $|e|<1$, so it is not
$\nu^{\pm1}$. A local component of the cuspidal $P$ at $u$ is generic, with supercuspidal support
consisting of unramified, pairwise unlinked characters, so it is the full induced representation, which
is spherical.
% Lean: GalTw.GlobalTheory.RoadU.TorLift.hasSat_charPol_above, through PolarizedWeight.exists_isWeight,
% GalCut7.SizeAway.isUnitarisable, GalTw.GlobalTheory.RoadU.window (the generic unitary dual),
% RoadU.unlinked_chars_comp_normU and GalTw.GlobalTheory.SphUnr.hasSat_of_unlinked_unr.
So $P$ is spherical at $u$ with
Satake parameters $B$ such that, for $g_0\in W_A$ and $g'$ in the Weil group of $M$ at $A$ acting as
$g_0$ on $\Omega$, the characteristic polynomial of $\mathrm{Sat}_A(B)(g')$ is
$\prod_j(X-\xi_j(a(g_0)))$. Since $r$ is weakly attached to $\pi$, its restriction to $G_M$ is weakly
attached to $P$ by Lemma~\ref{gal-solvable-base-change-facts}(2).

(1) Let $w\nmid\ell$. By (unr) for $P$ at $u$ in the first scope, $r|_{G_M}$ has Frobenius eigenvalues
$q_u^{(n-1)/2}B$ at $u$, so its Weil--Deligne representation at $u$ is $\mathrm{Sat}_A(B)$, which is
monodromy-free. The Weil--Deligne representation of $r|_{G_M}$ at $u$ is the restriction of
$\mathrm{WD}_\iota(r|_{W_A})$ to the Weil group of $M$ at $A$, which contains $g$. Hence
$\mathrm{WD}_\iota(r|_{W_A})$ is monodromy-free and its value at $g$ has characteristic polynomial
$\prod_j(X-\xi_j(a(g)))$. This holds for every $g$ of nonzero degree, and a monodromy-free
Weil--Deligne representation is determined by its characteristic polynomials at the elements of nonzero
degree, so $\mathrm{WD}_\iota(r|_{W_A})=\bigoplus_j\xi_j\circ a=\operatorname{rc}(\sigma^\flat)$.

(2) Let $w\mid\ell$. With the same lift, (crys), (dR) and (phi) for $P$ at $u$ show that $r|_{G_M}$ is
crystalline at $u$ with the Hodge--Tate weights given by $Z(P)=Z(\pi)$, and that its Weil--Deligne
representation at $u$ has the characteristic polynomials of $\mathrm{Sat}_A(B)$. So $r|_{D_A}$ becomes
crystalline on an open subgroup, hence is potentially semistable and de Rham with the given
Hodge--Tate weights, which is (dR) for $\pi$ at $w$; and $\mathrm{WD}_\iota(r|_{D_A})$ is
monodromy-free, with characteristic polynomial $\prod_j(X-\xi_j(a(g)))$ at each $g$ of nonzero degree,
hence at every $g\in W_A$, as in (1).
\end{proof}

\subsection{Layered extensions and weak base change}\label{subsection-gal-weak-base-change}

\begin{definition}\label{gal-depth-zero-lift-layered}
Let $E$ be a CM field with maximal totally real subfield $F$, let $\Omega$ be an algebraic closure of
$E$, and let $G_F=\operatorname{Gal}(\Omega/F)$.
\begin{enumerate}
\item A character $\rho:G_F\to\mathbf C^\times$ is \emph{split-ramified} if its kernel is open, if
  $\rho(c)=1$ for every $c\in G_F$ that acts as complex conjugation through some field embedding
  $\Omega\to\mathbf C$, and if $\rho$ is trivial on the inertia group of every valuation ring of
  $\Omega$ over a finite place of $F$ that does not split in $E$. The \emph{field} of $\rho$ is the
  fixed field in $\Omega$ of the elements of $\operatorname{Gal}(\Omega/E)$ in the kernel of $\rho$.
\item Let $S$ be a finite set of finite places of $E$. An \emph{$S$-layered extension} of $E$ is a
  field $X\subset\Omega$ with a chain $E=X_0\subset X_1\subset\cdots\subset X_m=X$ of CM fields, each
  a finite abelian extension of $E$, such that for each $i\geq1$ the field $X_i$ is a Galois extension
  of $X_{i-1}$ of prime degree and there is a finite place $u\notin S$ of $E$ with the following
  properties:
  \begin{enumerate}
  \item for every valuation ring $A$ of $\Omega$ over $u$, every element of the decomposition group of
    $A$ in $G_F$ fixes $E$ and $X_{i-1}$;
  \item for some valuation ring $A$ of $\Omega$ over $u$, some element of the inertia group of $A$ in
    $\operatorname{Gal}(\Omega/X_{i-1})$ does not fix $X_i$.
  \end{enumerate}
  The fields $X_i$ are the \emph{layers} of $X$. Condition~(a) says that the place of $F$ under $u$
  splits in $E$ and splits completely in $X_{i-1}$, and condition~(b) that $X_i$ is ramified over
  $X_{i-1}$ at a place over $u$.
\end{enumerate}
\end{definition}
\leanref{Development/AutomorphicArithmetic/Foundations/GaloisRepresentations/InertiaQuadraticRamification/NumberField.lean}{UpperTheory.Gal.UnramifiedFromSplittingCriterion.Good}
\leanref{Development/AutomorphicArithmetic/Foundations/GaloisRepresentations/InertiaQuadraticRamification/UnrCharAtOf.lean}{UpperTheory.Gal.UnramifiedFromSplittingCriterion.RamSplit}
\leanref{Development/AutomorphicArithmetic/Foundations/GaloisRepresentations/InertiaQuadraticRamification/NumberField.lean}{UpperTheory.Gal.UnramifiedFromSplittingCriterion.Cuts}
\leanref{Development/AutomorphicArithmetic/Foundations/GaloisRepresentations/InertiaQuadraticRamification/NumberField.lean}{UpperTheory.Gal.UnramifiedFromSplittingCriterion.Fld}
\leanref{Development/AutomorphicArithmetic/Foundations/GaloisRepresentations/InertiaQuadraticRamification/NumberField.lean}{UpperTheory.Gal.UnramifiedFromSplittingCriterion.SplitsIn}
\leanref{Development/AutomorphicArithmetic/Foundations/GaloisRepresentations/InertiaQuadraticRamification/NumberField.lean}{UpperTheory.Gal.UnramifiedFromSplittingCriterion.Layer}
\leanref{Development/AutomorphicArithmetic/Foundations/GaloisRepresentations/InertiaQuadraticRamification/NumberField.lean}{UpperTheory.Gal.UnramifiedFromSplittingCriterion.Tower}

\begin{lemma}\label{gal-depth-zero-lift-cyclotomic}
Let $E$ be a CM field with maximal totally real subfield $F$, $\Omega$ an algebraic closure of $E$,
$v$ a finite place of $F$ with residue field of $q_v$ elements, and $f\geq1$. Then there is a
split-ramified character $r$ of $\operatorname{Gal}(\Omega/F)$
(Definition~\ref{gal-depth-zero-lift-layered}) such that, for every valuation ring $A$ of $\Omega$
over $v$, $r$ is trivial on the inertia group of $A$, and $r(\varphi)$ has order $f$ for every
$\varphi$ in the decomposition group of $A$ that acts on the residue field of $A$ by
$z\mapsto z^{q_v}$.
\end{lemma}
\leanref{Development/AutomorphicArithmetic/Foundations/GaloisRepresentations/InertiaQuadraticRamification/UnrCharAtOf.lean}{UpperTheory.Gal.UnramifiedFromSplittingCriterion.UnrCharAt}
\leanref{Development/AutomorphicArithmetic/Foundations/GaloisRepresentations/InertiaQuadraticRamification/UnrCharAtOf.lean}{UpperTheory.Gal.UnramifiedFromSplittingCriterion.unrCharAt_of}
\leanref{Development/AutomorphicArithmetic/Foundations/GaloisRepresentations/InertiaQuadraticRamification/UnrCharPP.lean}{UpperTheory.Gal.UnramifiedFromSplittingCriterion.unrCharPP}
\leanref{Development/AutomorphicArithmetic/Foundations/GaloisRepresentations/InertiaQuadraticRamification/OrdPrime.lean}{UpperTheory.Gal.UnramifiedFromSplittingCriterion.ordPrime}
\leanref{Development/AutomorphicArithmetic/Foundations/GaloisRepresentations/InertiaQuadraticRamification/KumMove.lean}{UpperTheory.Gal.UnramifiedFromSplittingCriterion.kumMove}
\leanref{Development/AutomorphicArithmetic/Foundations/GaloisRepresentations/InertiaQuadraticRamification/CharOrd.lean}{UpperTheory.Gal.UnramifiedFromSplittingCriterion.charOrd}
\leanref{Development/AutomorphicArithmetic/Foundations/GaloisRepresentations/InertiaQuadraticRamification/DecPow.lean}{UpperTheory.Gal.UnramifiedFromSplittingCriterion.decPow}

\begin{proof}
If $f=f_1f_2$ with $f_1$ and $f_2$ coprime, and $r_1$ and $r_2$ have these properties for $f_1$ and
$f_2$, then $r_1r_2$ has them for $f$: its kernel contains $\ker r_1\cap\ker r_2$, and
$r_1(\varphi)r_2(\varphi)$ has order $f_1f_2$ since the orders of the factors are coprime. So we may
assume that $f=\ell^k$ for a prime $\ell$ and some $k\geq0$. For $m\geq0$ let $\zeta_{\ell^m}\in\Omega$ be
a primitive $\ell^m$-th root of unity.

\emph{Kummer theory.} Let $L\subset\Omega$ be the Galois closure of $E$ over $\mathbf Q$ and
$U=\operatorname{Gal}(\Omega/L)$. There is $j\geq1$ such that for every $m\geq j$ no $\ell^j$-th root
of $q_v$ lies in $L(\zeta_{\ell^m})$. Indeed, since $L$ is a finite extension of $\mathbf Q$, there are
$b\geq1$ and $\sigma_0\in U$ with $\sigma_0(\zeta_{\ell^b})\neq\zeta_{\ell^b}$. Let $\mathfrak P$ be a prime of $L$
over a prime $p_0$ dividing $q_v$, and $D=e(\mathfrak P/p_0)\,v_{p_0}(q_v)$; if $z\in L^\times$ and
$z^nq_v^s=1$ with $n,s\geq0$, then $n\,v_{\mathfrak P}(z)=-sD$, so $n$ divides $sD$. Let $\ell^{e_0}$
be the largest power of $\ell$ dividing $D$ and $j=b+e_0$. Suppose that $m\geq j$ and
$y\in L(\zeta_{\ell^m})$ satisfies $y^{\ell^j}=q_v$, and write $\sigma_0(\zeta_{\ell^m})=\zeta_{\ell^m}^a$
with $a\geq1$. For $\tau\in U$,
$\tau(y)=y\chi_\tau$ with $\chi_\tau^{\ell^j}=1$, and $\sigma_0$ and $\tau$ commute on the abelian
extension $L(\zeta_{\ell^m})$ of $L$; since $\sigma_0(\chi_\tau)=\chi_\tau^a$, the element
$z=\sigma_0(y)/y^a$ is fixed by every $\tau\in U$, so $z\in L$. Then $z^{\ell^j}q_v^{a-1}=1$, so
$\ell^j$ divides $(a-1)D$ and $\ell^b$ divides $a-1$. Hence $\sigma_0$ fixes
$\zeta_{\ell^m}^{\ell^{m-b}}$, a primitive $\ell^b$-th root of unity, and so fixes $\zeta_{\ell^b}$, a
contradiction.

\emph{The prime.} Let $m=j+k$ and $y\in\Omega$ with $y^{\ell^j}=q_v$. Since $y\notin
L(\zeta_{\ell^m})$, some $\sigma\in U$ fixes $\zeta_{\ell^m}$ and moves $y$. The field
$L'=L(\zeta_{\ell^m},y)$ is Galois over $\mathbf Q$, as it contains all
roots of $Y^{\ell^j}-q_v$. By the Chebotarev density theorem \cite[Theorem~VII.13.4]{Neu99} there is a
prime $p>q_v$, unramified in $L'$, with a prime of $L'$ over $p$ whose Frobenius element is the
restriction of $\sigma$. This Frobenius element fixes $L$, so $p$ splits completely in $L$ and hence
in $E$. It fixes $\zeta_{\ell^m}$, so $\ell^{j+k}$ divides $p-1$. It moves every root $y\zeta'$, with
$\zeta'$ an $\ell^j$-th root of unity, of $Y^{\ell^j}-q_v$, whose roots are distinct modulo $p$; so
$q_v$ is not an $\ell^j$-th power in $\mathbf F_p^\times$. Write $p-1=\ell^{e_1}h$ with $\ell\nmid h$,
so $e_1\geq j+k$. The order of $q_v$ in the cyclic group $\mathbf F_p^\times$ does not divide
$(p-1)/\ell^j$, so $\ell^{e_1-j+1}$ divides it, and $e_1-j+1\geq k+1$.

\emph{The character.} Let $\alpha\in(\mathbf Z/p)^\times$ be the class of $q_v$. Since $\ell^{k+1}$
divides the order of $\alpha$, $\ell^k$ divides the order of $\alpha^2$; write it as $\ell^ke_2$. Let
$\lambda_0$ be a faithful character of the cyclic group $(\mathbf Z/p)^\times$ and
$\lambda=\lambda_0^{2e_2}$. Then $\lambda(-1)=\lambda_0(-1)^{2e_2}=1$, and
$\lambda(\alpha)=\lambda_0(\alpha^2)^{e_2}$ has order $\ell^k$. Let
$\chi_p:\operatorname{Gal}(\Omega/\mathbf Q)\to(\mathbf Z/p)^\times$ be the cyclotomic character,
$g(\zeta)=\zeta^{\chi_p(g)}$ for every $p$-th root of unity $\zeta$, and let $r$ be the restriction of
$\lambda\circ\chi_p$ to $\operatorname{Gal}(\Omega/F)$. Its kernel contains
$\operatorname{Gal}(\Omega/F(\zeta_p))$ for a primitive $p$-th root of unity $\zeta_p$, so it is
open. A complex conjugation $c$ sends each $p$-th root of unity to its inverse, so
$r(c)=\lambda(-1)=1$. For a finite place $t$ of $F$ not over $p$ and a valuation ring over $t$, the
inertia group fixes the $p$-th roots of unity, and each element of the decomposition group acts on
them by $z\mapsto z^{q_t^i}$ for some $i\geq0$ \cite[Chapter~I, \S10]{Neu99}. Every place of $F$
over $p$ splits in $E$, since $p$ splits completely in $E$; so $r$ is split-ramified. Since $p>q_v$,
the place $v$ is not over $p$, so $r$ is trivial on the inertia groups over $v$. An element $\varphi$
as in the statement sends each $p$-th root of unity $\zeta$ to $\zeta^{q_v}$, since the $p$-th roots of
unity are distinct modulo the maximal ideal of $A$; so $\chi_p(\varphi)=\alpha$ and
$r(\varphi)=\lambda(\alpha)$ has order $\ell^k$.
\end{proof}

\begin{lemma}\label{gal-depth-zero-lift-fresh}
Let $E$ be a CM field with maximal totally real subfield $F$, $\Omega$ an algebraic closure of $E$,
$U$ an open subgroup of $G_F=\operatorname{Gal}(\Omega/F)$, $T$ a finite set of finite places of $F$,
and $n\geq1$. Then there are a split-ramified character $\psi$ of $G_F$
(Definition~\ref{gal-depth-zero-lift-layered}) with $\psi^n=1$ and a finite place $u\notin T$ of $F$
with the following properties.
\begin{enumerate}
\item $\psi$ is trivial on the decomposition group of every valuation ring of $\Omega$ over a place in
  $T$.
\item For every valuation ring $A$ of $\Omega$ over $u$, the decomposition group of $A$ in $G_F$ is
  contained in $U$, and $\psi(\tau)$ has order $n$ for some $\tau$ in the inertia group of $A$ in
  $G_F$.
\end{enumerate}
\end{lemma}
\leanref{Development/AutomorphicArithmetic/Foundations/GaloisRepresentations/InertiaQuadraticRamification/UnrCharAtOf.lean}{UpperTheory.Gal.UnramifiedFromSplittingCriterion.FreshSplit}
\leanref{Development/AutomorphicArithmetic/Foundations/GaloisRepresentations/InertiaQuadraticRamification/FreshSplit.lean}{UpperTheory.Gal.UnramifiedFromSplittingCriterion.freshSplit}
\leanref{Development/AutomorphicArithmetic/Foundations/GaloisRepresentations/InertiaQuadraticRamification/NumberField.lean}{UpperTheory.Gal.UnramifiedFromSplittingCriterion.FreshChar}
\leanref{Development/AutomorphicArithmetic/Foundations/GaloisRepresentations/InertiaQuadraticRamification/FreshChar.lean}{UpperTheory.Gal.UnramifiedFromSplittingCriterion.freshChar}
\leanref{Development/AutomorphicArithmetic/Foundations/GaloisRepresentations/InertiaQuadraticRamification/InfinitelyManyRealQuadraticSteps.lean}{UpperTheory.Gal.UnramifiedFromSplittingCriterion.freshChar_holds}

\begin{proof}
Let $V=U\cap\operatorname{Gal}(\Omega/E)$, an open subgroup of $G_F$ and so of
$\operatorname{Gal}(\Omega/\mathbf Q)$. Let $L$ be the Galois closure over $\mathbf Q$ of the fixed
field of $V$, let $T_0$ be the set of primes under the places in $T$, and let $L'$ be the field
generated over $L$ by the $2n$-th roots of unity and the $n$-th roots of the elements of $T_0$, which
is Galois over $\mathbf Q$. By the Chebotarev density theorem \cite[Theorem~VII.13.4]{Neu99} there is
a prime $p\notin T_0$, unramified in $L'$, that splits completely in $L'$. Then $p\equiv1\pmod{2n}$,
every element of $T_0$ is an $n$-th power modulo $p$, and the decomposition group in
$\operatorname{Gal}(\Omega/\mathbf Q)$ of every valuation ring of $\Omega$ over $p$ is contained in
$\operatorname{Gal}(\Omega/L)\subset V$.

Let $\lambda_0$ be a faithful character of the cyclic group $(\mathbf Z/p)^\times$ and
$\lambda=\lambda_0^{(p-1)/n}$. Then $\lambda^n=1$, $\lambda(-1)=(-1)^{(p-1)/n}=1$ since $(p-1)/n$ is
even, and $\lambda(\alpha)$ has order $n$ for every generator $\alpha$ of $(\mathbf Z/p)^\times$. Let
$\chi_p$ be the cyclotomic character of $\operatorname{Gal}(\Omega/\mathbf Q)$ modulo $p$, as in the
proof of Lemma~\ref{gal-depth-zero-lift-cyclotomic}, and $\psi$ the restriction of
$\lambda\circ\chi_p$ to $G_F$. Then $\psi^n=1$, its kernel is open, and $\psi(c)=\lambda(-1)=1$ for
every complex conjugation $c$. At a finite place $t$ of $F$ not over $p$, $\psi$ is trivial on the
inertia groups, and an element of a decomposition group acts on the $p$-th roots of unity by a power
of $z\mapsto z^{q_t}$ \cite[Chapter~I, \S10]{Neu99}. The decomposition groups over a place of $F$
over $p$ lie in $V\subset\operatorname{Gal}(\Omega/E)$, so that place splits in $E$. Hence $\psi$ is
split-ramified. For $t\in T$, the place $t$ is not over $p$, and $q_t$ is a power of an element of
$T_0$, hence an $n$-th power modulo $p$; so $\psi$ is trivial on the decomposition groups over $t$,
which is (1).

Let $u$ be a place of $F$ over $p$; then $u\notin T$, since $p\notin T_0$. Let $A$ be a valuation
ring of $\Omega$ over $u$. Its decomposition group in $G_F$ lies in that in
$\operatorname{Gal}(\Omega/\mathbf Q)$, hence in $V\subset U$. Since $\mathbf Q_p(\zeta_p)/\mathbf
Q_p$ is totally ramified with group $(\mathbf Z/p)^\times$ \cite[Chapter~I, \S10]{Neu99}, some $\tau$
in the inertia group of $A$ in $\operatorname{Gal}(\Omega/\mathbf Q)$ has $\chi_p(\tau)$ a generator
of $(\mathbf Z/p)^\times$. It lies in the decomposition group, hence in $V\subset G_F$, so in the
inertia group of $A$ in $G_F$, and $\psi(\tau)$ has order $n$. This is (2).
\end{proof}

\begin{lemma}\label{gal-depth-zero-lift-field}
Let $E$ be a CM field with maximal totally real subfield $F$, $\Omega$ an algebraic closure of $E$,
$w$ a finite place of $E$ that is split over $F$, $f\geq1$, and $S$ a finite set of finite places of
$E$. Then there is an $S$-layered extension $X$ of $E$ (Definition~\ref{gal-depth-zero-lift-layered})
with the following properties.
\begin{enumerate}
\item For every valuation ring $A$ of $\Omega$ over $w$, every element of the inertia group $I_A$ of
  $A$ in $\operatorname{Gal}(\Omega/E)$ fixes $X$, and the place of $X$ under $A$ has a residue field
  of $q_w^f$ elements.
\item For every finite place $u$ of $E$ that is not split over $F$ and every valuation ring $A$ of
  $\Omega$ over $u$, every element of the inertia group of $A$ in $\operatorname{Gal}(\Omega/E)$ fixes
  $X$.
\end{enumerate}
\end{lemma}
\leanref{Development/AutomorphicArithmetic/Foundations/GaloisRepresentations/CharpolyFromTrace/UnlinkedOfType.lean}{UpperTheory.Gal.NearKilling.UnrField}
\leanref{Development/AutomorphicArithmetic/Foundations/GaloisRepresentations/InertiaQuadraticRamification/UnrField.lean}{UpperTheory.Gal.UnramifiedFromSplittingCriterion.unrField_holds}
\leanref{Development/AutomorphicArithmetic/Foundations/GaloisRepresentations/InertiaQuadraticRamification/AsmUnr.lean}{UpperTheory.Gal.UnramifiedFromSplittingCriterion.UnrSplit}
\leanref{Development/AutomorphicArithmetic/Foundations/GaloisRepresentations/InertiaQuadraticRamification/AsmUnr.lean}{UpperTheory.Gal.UnramifiedFromSplittingCriterion.unrSplit_of}
\leanref{Development/AutomorphicArithmetic/Foundations/GaloisRepresentations/InertiaQuadraticRamification/MainUnr.lean}{UpperTheory.Gal.UnramifiedFromSplittingCriterion.unrSplit_holds}
\leanref{Development/AutomorphicArithmetic/Foundations/GaloisRepresentations/InertiaQuadraticRamification/InertiaGenAlt.lean}{UpperTheory.Gal.UnramifiedFromSplittingCriterion.GroupStep}
\leanref{Development/AutomorphicArithmetic/Foundations/GaloisRepresentations/InertiaQuadraticRamification/GroupStep.lean}{UpperTheory.Gal.UnramifiedFromSplittingCriterion.groupStep}
\leanref{Development/AutomorphicArithmetic/Foundations/GaloisRepresentations/InertiaQuadraticRamification/NormalOfComm.lean}{UpperTheory.Gal.UnramifiedFromSplittingCriterion.exFld}
\leanref{Development/AutomorphicArithmetic/Foundations/GaloisRepresentations/InertiaQuadraticRamification/InfinitelyManyRealQuadraticSteps.lean}{UpperTheory.Gal.UnramifiedFromSplittingCriterion.layerOf_holds}
\leanref{Development/AutomorphicArithmetic/Foundations/GaloisRepresentations/InertiaQuadraticRamification/ValuationLeOneOfInteger.lean}{UpperTheory.Gal.UnramifiedFromSplittingCriterion.degOf}
\leanref{Development/AutomorphicArithmetic/Foundations/GaloisRepresentations/InertiaQuadraticRamification/SplitIff.lean}{UpperTheory.Gal.UnramifiedFromSplittingCriterion.splitIff}

\begin{proof}
Let $G_F=\operatorname{Gal}(\Omega/F)$, $G_E=\operatorname{Gal}(\Omega/E)$ and $v$ the place of $F$
under $w$. For a split-ramified character $\rho$ of $G_F$ write $X_\rho$ for its field. The fixed
field of $\ker\rho$ is a finite abelian extension of $F$ fixed by every complex conjugation, hence
totally real; so $X_\rho$, which it generates together with $E$, is a CM field and a finite abelian
extension of $E$, and an element $g\in G_E$ fixes $X_\rho$ if and only if $\rho(g)=1$.

By Lemma~\ref{gal-depth-zero-lift-cyclotomic} there is a split-ramified $r$ such that, for every
valuation ring $A$ over $v$, $r$ is trivial on the inertia group of $A$ in $G_F$, and $r(\varphi)$
has order $f$ for every $\varphi$ in the decomposition group of $A$ in $G_F$ acting on the residue
field by $z\mapsto z^{q_v}$. Since $\ker r$ is open, $r^m=1$ for some $m\geq1$. Let $D$ be the union of
the decomposition groups in $G_F$ of the valuation rings over $v$, and $T$ the finite set of places of
$F$ consisting of $v$ and the places under $S$.

We show by induction on $n\geq1$: if $\rho$ is split-ramified and $X_{\rho^n}$ is an $S$-layered
extension of $E$, then there is a split-ramified $\rho'$ that agrees with $\rho$ on $D$ and whose
field $X_{\rho'}$ is an $S$-layered extension of $E$. For $n=1$ take $\rho'=\rho$. Let $n>1$ and
write $n=pn'$ with $p$ prime. Lemma~\ref{gal-depth-zero-lift-fresh}, applied to $U=\ker\rho$, $T$ and
$n$, gives a split-ramified $\psi$ with $\psi^n=1$, a place $u\notin T$ of $F$, a valuation ring $A$
over $u$ and $\tau$ in the inertia group of $A$ in $G_F$ with $\psi(\tau)$ of order $n$. Since
$n>1$ and $\psi$ is split-ramified, $u$ splits in $E$. Put $\rho_1=\rho\psi$. It is split-ramified, it
agrees with $\rho$ on $D$ since $\psi$ is trivial there, and $\rho_1^n=\rho^n$ since $\psi^n=1$. Let
$H=\ker\rho^n$ and $H'=\ker\rho_1^{n'}$. Then $H'\subset H$, and on $H$ the character $\rho_1^{n'}$
takes values in the $p$-th roots of unity. The element $\tau$ lies in the decomposition group of $A$,
which is contained in $\ker\rho\subset H$ and, as $u$ splits in $E$, fixes $E$,
and $\rho_1^{n'}(\tau)=\psi(\tau)^{n'}\neq1$.
Hence $G_E\cap H'$ has index $p$ in $G_E\cap H$, so $X_{\rho_1^{n'}}$ is a Galois extension of
$X_{\rho^n}$ of degree $p$. Let $u_E$ be the place of $E$ under $A$; it is not in $S$, since $u\notin
T$. For every valuation ring over $u_E$, its decomposition group in $G_F$ fixes $E$ and lies in $H$, so
it fixes $X_{\rho^n}$; and $\tau$ lies in the inertia group of $A$ in
$\operatorname{Gal}(\Omega/X_{\rho^n})$ and does not fix $X_{\rho_1^{n'}}$. So $X_{\rho_1^{n'}}$, with
the layers of $X_{\rho^n}$ followed by $X_{\rho_1^{n'}}$, is an $S$-layered extension of $E$. The
induction hypothesis for $\rho_1$ and $n'<n$ gives $\rho'$, which agrees with $\rho_1$, hence with
$\rho$, on $D$.

Apply this to $r$ and $n=m$. Since $r^m=1$, $X_{r^m}=E$, which is an $S$-layered extension of $E$
with no layers beyond $E$. This gives a split-ramified $r'$ that agrees with $r$ on $D$, and
$X=X_{r'}$ is an $S$-layered extension of $E$.

(1) Let $A$ be a valuation ring over $w$, hence over $v$. An element of $I_A$ lies in the inertia
group of $A$ in $G_F$, on which $r'=r$ is trivial; so it fixes $X$. Since $w$ is split over $F$,
$E_w=F_v$, so $q_w=q_v$, and an element $\varphi$ of the decomposition group of $A$ in $G_E$ that
acts on the residue field by $z\mapsto z^{q_w}$ is one as in
Lemma~\ref{gal-depth-zero-lift-cyclotomic}. For $k\geq0$, $\varphi^k$ fixes $X$ if and only
if $r'(\varphi)^k=r(\varphi)^k=1$, that is, if and only if $f$ divides $k$. So the place of $X$ under
$A$ is unramified over $w$ with residue degree $f$, and its residue field has $q_w^f$ elements.

(2) Let $u$ be a finite place of $E$ not split over $F$. The place of $F$ under $u$ does not split in
$E$, so $r'$ is trivial on its inertia groups in $G_F$, which contain those of $u$ in $G_E$. So these
fix $X$.
\end{proof}

\begin{definition}\label{gal-depth-zero-lift-weak}
Let $K\subset M$ be totally complex number fields with $M$ a Galois extension of $K$ of prime degree
$p$, let $a\geq1$, and let $\pi$ and $\Pi$ be cuspidal automorphic representations of
$\operatorname{GL}_a(\mathbf A_K)$ and $\operatorname{GL}_a(\mathbf A_M)$
(Definition~\ref{gal-polarized}).
\begin{enumerate}
\item $\Pi$ is a \emph{weak base change} of $\pi$ if for all but finitely many finite places $u$ of
  $M$, with $v$ the place of $K$ under $u$ and $N(u)=N(v)^f$, there is $B$ such that $\pi$ has Satake
  parameters $B$ at $v$ and $\Pi$ has Satake parameters $\{b^f:b\in B\}$ at $u$
  (Definition~\ref{gal-weakly-attached}).
\item The local components of $\pi$ and $\Pi$ \emph{agree at places of degree one} if for every
  finite place $v$ of $K$, finite place $u$ of $M$ and continuous field isomorphism $i:K_v\to M_u$
  extending $K\to M$, local components $\sigma$ of $\pi$ at $v$ and $\sigma'$ of $\Pi$ at $u$ satisfy
  $\sigma'\circ i\cong\sigma$.
\item A \emph{norm character} for $M$ is a Hecke character $\eta$ of $K$ with $\eta\neq1$,
  $\eta^p=1$, and $\eta(\varpi_v)^f=1$ for all but finitely many finite places $u$ of $M$, with $v$ and
  $f$ as in (1).
\item Let $E\subset X$ be totally complex number fields inside an algebraic closure $\Omega$ of $E$,
  and $\pi_E$ and $P$ cuspidal automorphic representations of $\operatorname{GL}_a(\mathbf A_E)$ and
  $\operatorname{GL}_a(\mathbf A_X)$. A \emph{chain of weak base changes} from $(E,\pi_E)$ to $(X,P)$
  consists of totally complex fields $E=K_0\subset K_1\subset\cdots\subset K_m=X$ inside $\Omega$,
  each $K_i$ a Galois extension of $K_{i-1}$ of prime degree, and cuspidal automorphic
  representations $\pi_i$ of $\operatorname{GL}_a(\mathbf A_{K_i})$ with $\pi_0=\pi_E$, $\pi_m=P$, and
  $\pi_i$ a weak base change of $\pi_{i-1}$ for each $i\geq1$.
\end{enumerate}
\end{definition}
\leanref{Development/AutomorphicArithmetic/AutomorphicVarieties/PublishedTheorems/CyclicBaseChangeC.lean}{UpperTheory.V.Published.WeakBC}
\leanref{PositivityTheory/AutomorphicForms/PublishedInputs/CyclicBaseChangeLifts/SplitLift.lean}{UpperTheory.CBC.SplitLift}
\leanref{Development/AutomorphicArithmetic/AutomorphicVarieties/PublishedTheorems/CyclicBaseChangeC.lean}{UpperTheory.V.Published.IsNormChar}
\leanref{Development/AutomorphicArithmetic/Foundations/GaloisRepresentations/RegularSemisimpleFibres/NumberField.lean}{UpperTheory.Gal.InertialTypesAlt.Pt}
\leanref{Development/AutomorphicArithmetic/Foundations/GaloisRepresentations/RegularSemisimpleFibres/NumberField.lean}{UpperTheory.Gal.InertialTypesAlt.Step}
\leanref{Development/AutomorphicArithmetic/Foundations/GaloisRepresentations/RegularSemisimpleFibres/NumberField.lean}{UpperTheory.Gal.InertialTypesAlt.Lifts}

\begin{theorem}\label{gal-depth-zero-lift-base-change}
Let $K$ be a totally complex number field, $M$ a Galois extension of $K$ of prime degree, and
$a\geq1$. Weak base changes, agreement at places of degree one and norm characters are those of
Definition~\ref{gal-depth-zero-lift-weak}, and $\pi$, $\pi_1$ and $\Pi$ are cuspidal automorphic
representations of $\operatorname{GL}_a(\mathbf A_K)$, $\operatorname{GL}_a(\mathbf A_K)$ and
$\operatorname{GL}_a(\mathbf A_M)$.
\begin{enumerate}
\item For every $\pi$, either $\pi\otimes(\eta\circ\det)\cong\pi$ for some norm character $\eta$ for
  $M$, or there is a $\Pi$ that is a weak base change of $\pi$, with $Z_\tau(\Pi)=Z_{\tau|_K}(\pi)$ for
  every embedding $\tau:M\to\mathbf C$, and whose local components agree with those of $\pi$ at places
  of degree one.
\item For every $\Pi$ such that $\Pi\circ\gamma\cong\Pi$ for every $\gamma\in\operatorname{Gal}(M/K)$,
  acting on $\operatorname{GL}_a(\mathbf A_M)$ entrywise, there is a $\pi$ of which $\Pi$ is a weak
  base change, with $Z_{\tau|_K}(\pi)=Z_\tau(\Pi)$ for every $\tau$, and whose local components agree
  with those of $\Pi$ at places of degree one.
\item If $\Pi$ is as in (2) and is a weak base change of $\pi$, then
  $\pi\otimes(\eta\circ\det)\not\cong\pi$ for every norm character $\eta$ for $M$, and every $\pi_1$
  of which $\Pi$ is a weak base change is isomorphic to $\pi\otimes(\eta\circ\det)$ with $\eta=1$ or
  $\eta$ a norm character for $M$.
\end{enumerate}
This is the cyclic base change of Arthur and Clozel \cite[Chapter~3]{AC89}, together with strong
multiplicity one \cite{JS81a}.
\end{theorem}
\leanref{Development/AutomorphicArithmetic/AutomorphicVarieties/ModularityArgument/BaseChangeUnits/ComplexBaseChangeOfLiterature.lean}{UpperTheory.V.Z.Bc.CBCc}
\leanref{Development/AutomorphicArithmetic/Foundations/Convolution/Cbcc.lean}{UpperTheory.ConvolutionTheory.cbcc}

\begin{lemma}\label{gal-depth-zero-lift-step}
Let $K\subset M$ be CM fields inside an algebraic closure $\Omega$ of $K$, with $M$ a Galois extension
of $K$ of prime degree. Let $\ell$ be a prime, $a\geq1$, and $(\pi,\chi)$ in scope
(Definition~\ref{gal-scope}) with $\pi$ a cuspidal automorphic representation of
$\operatorname{GL}_a(\mathbf A_K)$. Let $A$ be a valuation ring of $\Omega$ over a finite place $v$ of
$K$ at which $\pi$ has Satake parameters, such that some element of the inertia group of $A$ in
$\operatorname{Gal}(\Omega/K)$ does not fix $M$. Then there are a cuspidal automorphic representation
$P$ of $\operatorname{GL}_a(\mathbf A_M)$ and a Hecke character $\chi'$ of the maximal totally real
subfield of $M$ with the following properties.
\begin{enumerate}
\item $P$ is a weak base change of $\pi$, and the local components of $\pi$ and $P$ agree at places
  of degree one (Definition~\ref{gal-depth-zero-lift-weak}).
\item $(P,\chi')$ is in scope.
\item If $\chi=\eta_K^a$, with $\eta_K$ the quadratic character of $K$ over its maximal totally real
  subfield, then $(P,\eta_M^a)$ is polarized, with $\eta_M$ defined in the same way for $M$.
\item For every field isomorphism $\iota:\bar{\mathbf Q}_\ell\to\mathbf C$ and every
  $r:G_K\to\operatorname{GL}_a(\bar{\mathbf Q}_\ell)$ weakly attached to $\pi$
  (Definition~\ref{gal-weakly-attached}), the restriction $r|_{G_M}$ is weakly attached to $P$.
\end{enumerate}
\end{lemma}
\leanref{Development/AutomorphicArithmetic/Foundations/GaloisRepresentations/CharpolyFromTrace/NumberField.lean}{UpperTheory.Gal.NearKilling.StepLift}
\leanref{Development/AutomorphicArithmetic/Foundations/GaloisRepresentations/CharpolyFromTrace/StepLift.lean}{UpperTheory.Gal.NearKilling.stepLift}
\leanref{Development/AutomorphicArithmetic/Foundations/GaloisRepresentations/CharpolyFromTrace/GoodLiftCHolds.lean}{UpperTheory.Gal.NearKilling.stepLift_holds}
\leanref{Development/AutomorphicArithmetic/Foundations/GaloisRepresentations/InertiaQuadraticRamification/InfinitelyManyRealQuadraticSteps.lean}{UpperTheory.Gal.UnramifiedFromSplittingCriterion.noSelfTwist_holds}
\leanref{Development/AutomorphicArithmetic/Foundations/GaloisRepresentations/InertiaQuadraticRamification/NormRam.lean}{UpperTheory.Gal.UnramifiedFromSplittingCriterion.normRam}
\leanref{Development/AutomorphicArithmetic/Foundations/GaloisRepresentations/CharpolyFromTrace/PolNorm.lean}{UpperTheory.Gal.NearKilling.polBC}
\leanref{Development/AutomorphicArithmetic/Foundations/GaloisRepresentations/CharpolyFromTrace/WeakRes.lean}{UpperTheory.Gal.NearKilling.weakRes}

\begin{proof}
\emph{No self-twist.} Let $\eta$ be a norm character for $M$. By global class field theory there is a
continuous character $\chi_\eta$ of $G_K$ of finite order such that, at every finite place $v'$ of
$K$ where $\eta$ is unramified, $\chi_\eta$ is trivial on the inertia groups over $v'$ and takes the
value $\eta(\varpi_{v'})$ on the elements of the decomposition groups over $v'$ that act on the
residue field by $z\mapsto z^{q_{v'}}$. At all but finitely many places $u$ of $M$, with $v'$ and $f$
as in Definition~\ref{gal-depth-zero-lift-weak}, an element of $G_M$ acting on the residue field of
$u$ by $z\mapsto z^{q_u}$ acts on that of $v'$ by $z\mapsto z^{q_{v'}^f}$, so $\chi_\eta$ takes on it
the value $\eta(\varpi_{v'})^f=1$. These elements are dense in $G_M$ by the Chebotarev density
theorem, so $\chi_\eta$ is trivial on $G_M$. It is nontrivial since $\eta\neq1$, so it is a
nontrivial character of $\operatorname{Gal}(M/K)$, a group of prime order, and is injective there.
Some element of the inertia group of $A$ does not fix $M$, so
$\chi_\eta$ is nontrivial on that inertia group, and $\eta$ is ramified at $v$.

Suppose that $\pi\otimes(\eta\circ\det)\cong\pi$, and let $V$ be the space of $\pi$ and
$\mathcal K_v=\operatorname{GL}_a(\mathcal O_v)$. Since $\pi$ has Satake parameters at $v$, there is a
nonzero $x\in V$ fixed by $\mathcal K_v$, and the isomorphism carries it to a nonzero $y\in V$ with
$ky=\eta_v(\det k)y$ for $k\in\mathcal K_v$. We show that $y=0$, a contradiction. Let $P$ be the
projection $z\mapsto\int_{\mathcal K_v}\eta_v(\det k)^{-1}kz\,dk$, for the Haar measure of total mass
$1$. Then $Py=y$; $P(kz)=\eta_v(\det k)Pz$ for $k\in\mathcal K_v$; and $P$ commutes with the action at
the infinite places. So if some $h\in\mathcal K_v$ with $\eta_v(\det h)\neq1$ fixes $z$, then $Pz=0$.
Since $\eta$ is ramified at $v$, there is a unit $u$ of $\mathcal O_v$ with $\eta_v(u)\neq1$. For
$g\in\operatorname{GL}_a(\mathbf A_{K,f})$, write $g_v=k_1dk_2$ with $k_1,k_2\in\mathcal K_v$ and $d$
diagonal (the Cartan decomposition), and put $h=k_1\operatorname{diag}(u,1,\ldots,1)k_1^{-1}$. Then
$h\in\mathcal K_v$, $\det h=u$, and $g_v^{-1}hg_v=k_2^{-1}\operatorname{diag}(u,1,\ldots,1)k_2\in
\mathcal K_v$, since diagonal matrices commute; so $h$ fixes $gx$, and $P(gx)=0$. The subspace $W$ of
$z\in V$ with $P(gz)=0$ for every $g\in\operatorname{GL}_a(\mathbf A_{K,f})$ is stable under the finite
adelic group, and under the action at the infinite places since $P$ commutes with that action. It
contains $x\neq0$, so $W=V$ because $\pi$ is irreducible. Then $y\in W$, so $y=Py=0$. So
$\pi\otimes(\eta\circ\det)\not\cong\pi$.

\emph{The lift.} By Theorem~\ref{gal-depth-zero-lift-base-change}(1) there is a cuspidal $P$ that is a
weak base change of $\pi$, with $Z_\tau(P)=Z_{\tau|_K}(\pi)$ for every $\tau$, and whose local
components agree with those of $\pi$ at places of degree one, which is (1). Since $\pi$ is regular
algebraic, so is $P$.

\emph{Polarization.} Let $K^+$ and $M^+$ be the maximal totally real subfields and
$\chi'=\chi\circ N_{M^+/K^+}$. At a real place of $M^+$ over the real place $v'$ of $K^+$, the norm is
the identity, so $\chi'(-1)=\chi_{v'}(-1)=(-1)^a$. Put $\mu=\chi\circ N_{K/K^+}$ and
$\mu'=\chi'\circ N_{M/M^+}=\mu\circ N_{M/K}$. Let $u$ be a place of $M$ outside a finite set, with $v'$
the place of $K$ under $u$ and $N(u)=N(v')^f$. The polarization of $\pi$ shows that its Satake
parameters at $c(v')$ are $\{\mu(\varpi_{v'})b^{-1}:b\in B\}$, with $B$ those at $v'$. Since $P$ is a
weak base change of $\pi$, its Satake parameters at $u$ and at $c(u)$ are the $f$-th powers of these,
and $\mu'(\varpi_u)=\mu(\varpi_{v'})^f$. So $P^c$ and $P^\vee\otimes(\mu'\circ\det)$ have the same
Satake parameters at all but finitely many places, and by strong multiplicity one \cite{JS81a} and the
description of the contragredient by Gelfand and Kazhdan \cite{GK75}, $(P,\chi')$ is polarized. This
gives (2). If $\chi=\eta_K^a$, then $\chi'=\eta_M^a$, since $\eta_K\circ N_{M^+/K^+}=\eta_M$ because
$M=KM^+$, by class field theory; this is (3).

\emph{Galois representations.} The restriction $r|_{G_M}$ is continuous, and semisimple by Clifford's
theorem, since $G_M$ is a normal subgroup of finite index in $G_K$. At all but finitely many places $u$
of $M$, with $v'$ and $f$ as above, $r$ has Frobenius eigenvalues $q_{v'}^{(a-1)/2}B$ at $v'$, with
$B$ the Satake parameters of $\pi$ at $v'$, and $P$ has Satake parameters $\{b^f:b\in B\}$ at $u$. An
element of $G_M$ acting on the residue field of $u$ by $z\mapsto z^{q_u}$ acts on that of $v'$ by
$z\mapsto z^{q_{v'}^f}$, so $r|_{G_M}$ has Frobenius eigenvalues $q_u^{(a-1)/2}\{b^f:b\in B\}$ at $u$.
This is (4).
\end{proof}

\begin{remark}\label{gal-depth-zero-lift-step-support}
The step \emph{no self-twist} in the proof of Lemma~\ref{gal-depth-zero-lift-step} also follows from the
local component $\sigma$ of $\pi$ at $v$ alone. An isomorphism $\pi\otimes(\eta\circ\det)\cong\pi$ gives
a nonzero vector of $\sigma$ on which $\operatorname{GL}_a(\mathcal O_v)$ acts by $\eta_v\circ\det$, so
$\sigma\otimes(\eta_v^{-1}\circ\det)$ has a nonzero $\operatorname{GL}_a(\mathcal O_v)$-fixed vector and
its supercuspidal support consists of unramified characters \cite{Car79}. That support is the support of
$\sigma$, which consists of unramified characters, twisted by the ramified character $\eta_v^{-1}$; this
contradicts the uniqueness of the supercuspidal support \cite{BZ77}.
\end{remark}

\begin{lemma}\label{gal-depth-zero-lift-tower}
Let $E$ be a CM field, $\Omega$ an algebraic closure of $E$, $\ell$ a prime, $a\geq1$, and
$(\pi,\chi)$ in scope (Definition~\ref{gal-scope}) with $\pi$ a cuspidal automorphic representation
of $\operatorname{GL}_a(\mathbf A_E)$. Let $S$ be a finite set of finite places of $E$ such that $\pi$
has Satake parameters at every finite place not in $S$, and let $X$ be an $S$-layered extension of $E$
(Definition~\ref{gal-depth-zero-lift-layered}). Then there are a cuspidal automorphic representation
$P$ of $\operatorname{GL}_a(\mathbf A_X)$ and a Hecke character $\chi'$ of the maximal totally real
subfield of $X$ such that $(P,\chi')$ is in scope, there is a chain of weak base changes from
$(E,\pi)$ to $(X,P)$ (Definition~\ref{gal-depth-zero-lift-weak}), and for every field isomorphism
$\iota:\bar{\mathbf Q}_\ell\to\mathbf C$ and every $r:G_E\to\operatorname{GL}_a(\bar{\mathbf Q}_\ell)$
weakly attached to $\pi$, the restriction $r|_{G_X}$ is weakly attached to $P$.
\end{lemma}
\leanref{Development/AutomorphicArithmetic/Foundations/GaloisRepresentations/CharpolyFromTrace/GoodLiftAlt.lean}{UpperTheory.Gal.NearKilling.AllLift}
\leanref{Development/AutomorphicArithmetic/Foundations/GaloisRepresentations/CharpolyFromTrace/AllLift.lean}{UpperTheory.Gal.NearKilling.allLift}
\leanref{Development/AutomorphicArithmetic/Foundations/GaloisRepresentations/CharpolyFromTrace/GoodLiftCHolds.lean}{UpperTheory.Gal.NearKilling.allLift_holds}
\leanref{Development/AutomorphicArithmetic/Foundations/GaloisRepresentations/InertiaQuadraticRamification/SatSplit.lean}{UpperTheory.Gal.UnramifiedFromSplittingCriterion.satSplit}

\begin{proof}
Let $F$ be the maximal totally real subfield of $E$ and $E=X_0\subset\cdots\subset X_m=X$ the layers.
By induction on $i$ we construct a cuspidal automorphic representation $P_i$ of
$\operatorname{GL}_a(\mathbf A_{X_i})$ and a Hecke character $\chi_i$ with $(P_i,\chi_i)$ in scope, a
chain of weak base changes from $(E,\pi)$ to $(X_i,P_i)$ through $X_0,\ldots,X_i$, such that
$r|_{G_{X_i}}$ is weakly attached to $P_i$ whenever $r$ is weakly attached to $\pi$, and such that:
\begin{enumerate}
\item[$(\ast)$] for every finite place $u\notin S$ of $E$ and every valuation ring $A$ of $\Omega$
  over $u$ whose decomposition group in $\operatorname{Gal}(\Omega/E)$ fixes $X_i$, $P_i$ has Satake
  parameters at the place of $X_i$ under $A$.
\end{enumerate}
For $i=0$ take $P_0=\pi$ and $\chi_0=\chi$; $(\ast)$ holds since $\pi$ has Satake parameters at every
place not in $S$.

Let $i\geq1$, and let $u\notin S$ and $A$ be as in Definition~\ref{gal-depth-zero-lift-layered}(2)
for the layer $X_i$ over $X_{i-1}$, with an element of the inertia group of $A$ in
$\operatorname{Gal}(\Omega/X_{i-1})$ that does not fix $X_i$. The decomposition group of $A$ in
$\operatorname{Gal}(\Omega/E)$ lies in that in $\operatorname{Gal}(\Omega/F)$, so it fixes $X_{i-1}$,
and by $(\ast)$ the representation $P_{i-1}$ has Satake parameters at the place of $X_{i-1}$ under
$A$. Lemma~\ref{gal-depth-zero-lift-step}, applied to $X_{i-1}\subset X_i$, $(P_{i-1},\chi_{i-1})$
and $A$, gives $(P_i,\chi_i)$ in scope, with $P_i$ a weak base change of $P_{i-1}$ whose local
components agree with those of $P_{i-1}$ at places of degree one, and with $r|_{G_{X_i}}$ weakly
attached to $P_i$ whenever $r|_{G_{X_{i-1}}}$ is weakly attached to $P_{i-1}$.

For $(\ast)$ at $i$, let $u'\notin S$ and $A'$ be a valuation ring over $u'$ whose decomposition group
in $\operatorname{Gal}(\Omega/E)$ fixes $X_i$, and let $w'$ and $v'$ be the places of $X_i$ and
$X_{i-1}$ under $A'$. This decomposition group also fixes $X_{i-1}$, so by $(\ast)$ for $i-1$ the
representation $P_{i-1}$ has Satake parameters $B$ at $v'$. The decomposition group of $A'$ in
$\operatorname{Gal}(\Omega/X_{i-1})$ fixes $X_i$, so there is a continuous field isomorphism from the
completion of $X_{i-1}$ at $v'$ to that of $X_i$ at $w'$ extending $X_{i-1}\to X_i$, and the residue
fields of $v'$ and $w'$ have the same number of elements. The local components of $P_{i-1}$ at $v'$
and of $P_i$ at $w'$ correspond under this isomorphism, so $P_i$ has Satake parameters $B$ at $w'$.
\end{proof}

\begin{definition}\label{gal-depth-zero-lift-stable-spectrum}
Let $K$ be a totally complex number field and $n\geq1$. Let $\operatorname{GL}_n(\mathbf A_K)^1$ be the
group of $g\in\operatorname{GL}_n(\mathbf A_K)$ with $|\det g|=1$ for the idelic norm, and $L^2_n(K)$ the
Hilbert space of square-integrable functions on $\operatorname{GL}_n(K)\backslash\operatorname{GL}_n
(\mathbf A_K)^1$ for a Haar measure. A \emph{positive real scalar} is an element of
$\operatorname{GL}_n(\mathbf A_K)$ that is the same positive real multiple of the identity at every
infinite place and the identity at every finite place. Every $g\in\operatorname{GL}_n(\mathbf A_K)$ is
$zg_1$ with $z$ a positive real scalar and $g_1\in\operatorname{GL}_n(\mathbf A_K)^1$, and it acts on
$L^2_n(K)$ by $(R(g)\varphi)(x)=\varphi(xg_1)$. The \emph{cuspidal part} of $L^2_n(K)$ consists
of the $\varphi$ whose constant terms along the standard maximal parabolic subgroups vanish almost
everywhere.
\begin{enumerate}
\item A \emph{cuspidal summand} for $\operatorname{GL}_n$ over $K$ is a nonzero closed subspace $V$ of the
  cuspidal part of $L^2_n(K)$, stable under $R$, whose only closed $R$-stable subspaces are $0$ and $V$.
\item Let $\pi$ be a cuspidal automorphic representation of $\operatorname{GL}_n(\mathbf A_K)$
  (Definition~\ref{gal-polarized}). A cuspidal summand $V$ \emph{carries} $\pi$ if the forms of $\pi$
  are the cusp forms on $\operatorname{GL}_n(K)\backslash\operatorname{GL}_n(\mathbf A_K)$ that are
  invariant under the positive real scalars and whose restrictions to $\operatorname{GL}_n(\mathbf
  A_K)^1$ lie in $V$.
\item Let $V$ be a cuspidal summand, $\tau:K\to\mathbf C$ an embedding and $m$ a multiset of $n$ complex
  numbers. For $1\leq i\leq n$ let $z_{\tau,i}$ be the element of the centre of the enveloping algebra of
  the $\tau$-factor of the complexified Lie algebra of $\operatorname{GL}_n(K\otimes_{\mathbf Q}\mathbf R)$
  whose Harish-Chandra image is the $i$-th elementary symmetric polynomial $e_i$. Then $V$ has
  \emph{infinitesimal character $m$ at $\tau$} if $R(z_{\tau,i}\beta)=e_i(m)R(\beta)$ on $V$ for every
  $i$ and every smooth compactly supported function $\beta$ on $\operatorname{GL}_n(K\otimes_{\mathbf
  Q}\mathbf R)$, where $R(\beta)=\int\beta(g)R(g)\,dg$ and $z_{\tau,i}$ acts on $\beta$ as a
  differential operator.
\item Let $V$ be a cuspidal summand, $w$ a finite place of $K$ with residue field of $q_w$ elements and
  uniformizer $\varpi_w$, and $B$ a multiset of $n$ complex numbers. For $1\leq i\leq n$ let $a_i$ be the
  diagonal matrix with entries $\varpi_w$ ($i$ times) and $1$ ($n-i$ times), and $T_{w,i}=q_w^{-i(n-i)/2}
  \sum_cR(c)$ on the $\operatorname{GL}_n(\mathcal O_w)$-fixed vectors, the sum over representatives $c$
  of the cosets in $\operatorname{GL}_n(\mathcal O_w)a_i\operatorname{GL}_n(\mathcal O_w)/
  \operatorname{GL}_n(\mathcal O_w)$. Then $V$ has \emph{Satake parameters $B$ at $w$} if $V$ has a
  nonzero $\operatorname{GL}_n(\mathcal O_w)$-fixed vector and every $T_{w,i}$ acts on the
  $\operatorname{GL}_n(\mathcal O_w)$-fixed vectors of $V$ by $e_i(B)$.
\item Let $\Lambda$ assign to every embedding $\tau:K\to\mathbf C$, and $t$ to every finite place of $K$,
  a multiset of complex numbers, and let $S$ be a set of finite places of $K$. A cuspidal summand $V$ has
  \emph{parameters $(\Lambda,t)$ off $S$} if it has infinitesimal character $\Lambda(\tau)$ at every
  $\tau$ and Satake parameters $t(w)$ at every finite place $w\notin S$.
\item Let $\sigma$ be an automorphism of $K$, acting on $\mathbf A_K=K\otimes_{\mathbf Q}\mathbf A_{\mathbf
  Q}$ through the first factor and on $\operatorname{GL}_n(\mathbf A_K)$ entrywise. The \emph{Galois
  twist} $\theta_\sigma$ is the unitary operator of $L^2_n(K)$ given by
  $(\theta_\sigma\varphi)(x)=\varphi(\sigma^{-1}(x))$.
\end{enumerate}
\end{definition}
\leanref{Development/AutomorphicArithmetic/Foundations/Convolution/PowerSumTraces/CuspidalMultiplicityOneTotallyComplex.lean}{UpperTheory.ConvolutionTheory.C.IsSummand}
\leanref{Development/AutomorphicArithmetic/Foundations/Convolution/CentreSplitting/LocalRepresentationL2.lean}{UpperTheory.ConvolutionTheory.Dp.Carries}
\leanref{Development/AutomorphicArithmetic/Foundations/Convolution/TensFun.lean}{UpperTheory.ConvolutionTheory.SInf}
\leanref{Development/AutomorphicArithmetic/Foundations/Convolution/TensFun.lean}{UpperTheory.ConvolutionTheory.SSat}
\leanref{Development/AutomorphicArithmetic/Foundations/Convolution/TensFun.lean}{UpperTheory.ConvolutionTheory.At}
\leanref{Development/AutomorphicArithmetic/Foundations/Convolution/TraceSplitting/NormCharSplit.lean}{UpperTheory.ConvolutionTheory.T.AtAE}
\leanref{Development/WeilBound/RepresentationTheoryGL/TwistedFactors/CuspDecayHeights/HeightLe.lean}{PositivityTheory.Rg.Tf.Core.Twist}
\leanref{Development/WeilBound/RepresentationTheoryGL/TwistedFactors/SelmerRankBounds/IsNormAlong.lean}{PositivityTheory.Rg.Tf.Geo.IsGalTwist}

\begin{definition}\label{gal-depth-zero-lift-stable-parabolic}
Let $K\subset M$ be totally complex number fields and $n\geq1$. Cuspidal summands and their parameters
are those of Definition~\ref{gal-depth-zero-lift-stable-spectrum}.
\begin{enumerate}
\item Let $\Lambda$ assign to every embedding $K\to\mathbf C$, and $t$ to every finite place of $K$, a
  multiset of complex numbers. Their \emph{base changes} to $M$ are $\operatorname{bc}\Lambda(\tau)=
  \Lambda(\tau|_K)$ for an embedding $\tau:M\to\mathbf C$ and $\operatorname{bc}t(w)=\{b^{f_w}:b\in
  t(v)\}$ for a finite place $w$ of $M$, with $v$ the place of $K$ under $w$ and $f_w$ the residue degree
  of $w$ over $v$.
\item Let $L$ be a totally complex number field. A \emph{parabolic datum} of rank $n$ over $L$ consists
  of an integer $r\geq2$, integers $d_1,\ldots,d_r\geq1$ with sum $n$, and for each $i$ a cuspidal summand
  $V_i$ for $\operatorname{GL}_{d_i}$ over $L$ and a complex number $\lambda_i$. Let $\Lambda$ and $t$
  assign multisets of complex numbers to the embeddings and the finite places of $L$, and $S$ be a set of
  finite places of $L$. The parabolic datum has \emph{parameters $(\Lambda,t)$ off $S$} if each $V_i$ has
  parameters $(m_i,B_i)$ off $S$ for some $m_i$ and $B_i$, with $\Lambda(\tau)=\sum_i\{z+\lambda_i:z\in
  m_i(\tau)\}$ for every embedding $\tau:L\to\mathbf C$ and $t(w)=\sum_i\{q_w^{-\lambda_i}b:b\in
  B_i(w)\}$ for every finite place $w\notin S$, as sums of multisets.
\item Let $\Lambda$ assign to every embedding $M\to\mathbf C$, and $t$ to every finite place of $M$, a
  multiset of complex numbers. Then $(\Lambda,t)$ \emph{avoids parabolic data} over $K$ if for every
  finite set $S$ of finite places of $K$, with $S_M$ the set of places of $M$ over $S$:
  \begin{enumerate}
  \item[(i)] no parabolic datum of rank $n$ over $M$ has parameters $(\Lambda,t)$ off $S_M$;
  \item[(ii)] for every parabolic datum of rank $n$ over $K$ with parameters $(\Lambda',t')$ off $S$,
    $\Lambda\neq\operatorname{bc}\Lambda'$ or $t(w)\neq\operatorname{bc}t'(w)$ for some finite place
    $w\notin S_M$.
  \end{enumerate}
\item Let $\pi$ be a cuspidal automorphic representation of $\operatorname{GL}_n(\mathbf A_K)$. A
  \emph{weak base change of $\pi$ to a sum} is a nonempty finite family of cuspidal automorphic
  representations $Q_j$ of $\operatorname{GL}_{m_j}(\mathbf A_M)$, with $m_j\geq1$, such that for all
  but finitely many finite places $w$ of $M$, with $v$ the place of $K$ under $w$ and $N(w)=N(v)^f$, there
  is $B$ such that $\pi$ has Satake parameters $B$ at $v$, each $Q_j$ has Satake parameters $B_j$ at $w$,
  and $\sum_jB_j=\{b^f:b\in B\}$ (Definition~\ref{gal-weakly-attached}).
\end{enumerate}
\end{definition}
\leanref{Development/AutomorphicArithmetic/Foundations/Convolution/TensFun.lean}{UpperTheory.ConvolutionTheory.bcInf}
\leanref{Development/AutomorphicArithmetic/Foundations/Convolution/TensFun.lean}{UpperTheory.ConvolutionTheory.bcSat}
\leanref{Development/AutomorphicArithmetic/Foundations/Convolution/TensFun.lean}{UpperTheory.ConvolutionTheory.Datum}
\leanref{Development/AutomorphicArithmetic/Foundations/Convolution/TensFun.lean}{UpperTheory.ConvolutionTheory.Datum.Pt}
\leanref{Development/AutomorphicArithmetic/Foundations/Convolution/ASmooth.lean}{UpperTheory.ConvolutionTheory.OffFam}
\leanref{Development/AutomorphicArithmetic/Foundations/Convolution/TraceSplitting/NormCharSplit.lean}{UpperTheory.ConvolutionTheory.T.OffAll}
\leanref{Development/AutomorphicArithmetic/Foundations/Convolution/TraceSplitting/NormCharSplit.lean}{UpperTheory.ConvolutionTheory.T.LiftPlus}
\leanref{Development/AutomorphicArithmetic/Foundations/Convolution/NormCharacterPowers/SumSatUnit.lean}{UpperTheory.ConvolutionTheory.R.sumSat}

\begin{lemma}\label{gal-depth-zero-lift-stable-carried}
Let $K$ be a totally complex number field and $n\geq1$, with cuspidal summands, carrying and parameters
as in Definition~\ref{gal-depth-zero-lift-stable-spectrum}.
\begin{enumerate}
\item Every cuspidal summand carries exactly one cuspidal automorphic representation of
  $\operatorname{GL}_n(\mathbf A_K)$.
\item Let $\pi$ be a cuspidal automorphic representation of $\operatorname{GL}_n(\mathbf A_K)$ whose forms
  are invariant under the positive real scalars. Then the closure in
  $L^2_n(K)$ of the restrictions of these forms to $\operatorname{GL}_n(\mathbf A_K)^1$ is a cuspidal
  summand that carries $\pi$, and it is the only cuspidal summand that carries $\pi$.
\item For every cuspidal automorphic representation $\pi$ of $\operatorname{GL}_n(\mathbf A_K)$ there is
  $s\in\mathbf C$ such that $\pi\otimes|\det|^s$, the space of the functions $g\mapsto\varphi(g)|\det
  g|^s$ with $\varphi\in\pi$, is carried by a cuspidal summand.
\item Let $V$ carry $\pi$. Then $V$ has infinitesimal character $m$ at an embedding $\tau$ if and only if
  $Z_\tau(\pi)=m$ (Definition~\ref{gal-polarized}), and $V$ has Satake parameters $B$ at a finite place
  $w$ if and only if $\pi$ has Satake parameters $B$ at $w$ (Definition~\ref{gal-weakly-attached}).
\item Every cuspidal summand has parameters $(\Lambda,t)$ off some finite set of finite places.
\item Two isomorphic cuspidal automorphic representations of $\operatorname{GL}_n(\mathbf A_K)$ are
  equal.
\end{enumerate}
Parts (1) and (2) come from the decomposition of the cuspidal part of $L^2_n(K)$ into cuspidal summands
\cite{GGPS69}, whose smooth vectors are cusp forms \cite{HC68}, and (6) is the multiplicity one theorem
of Shalika \cite{Sha74}.
\end{lemma}
\leanref{Development/AutomorphicArithmetic/Foundations/Convolution/CentreSplitting/LocalRepresentationL2.lean}{UpperTheory.ConvolutionTheory.Dp.ExCusp}
\leanref{Development/AutomorphicArithmetic/Foundations/Convolution/CentreSplitting/ExCusp.lean}{UpperTheory.ConvolutionTheory.Dp.exCusp}
\leanref{Development/AutomorphicArithmetic/Foundations/Convolution/CentreSplitting/LocalRepresentationL2.lean}{UpperTheory.ConvolutionTheory.Dp.ExSummand}
\leanref{Development/AutomorphicArithmetic/Foundations/Convolution/CentreSplitting/ExSummand.lean}{UpperTheory.ConvolutionTheory.Dp.exSummand}
\leanref{Development/AutomorphicArithmetic/Foundations/Convolution/CentreSplitting/LocalRepresentationL2.lean}{UpperTheory.ConvolutionTheory.Dp.ExTwist}
\leanref{Development/AutomorphicArithmetic/Foundations/Convolution/CentreSplitting/ExTwist.lean}{UpperTheory.ConvolutionTheory.Dp.exTwist}
\leanref{Development/AutomorphicArithmetic/Foundations/Convolution/CentreSplitting/DefsObj.lean}{UpperTheory.ConvolutionTheory.Dp.SInfIff}
\leanref{Development/AutomorphicArithmetic/Foundations/Convolution/CentreSplitting/SInfIff.lean}{UpperTheory.ConvolutionTheory.Dp.sInfIff}
\leanref{Development/AutomorphicArithmetic/Foundations/Convolution/CentreSplitting/DefsObj.lean}{UpperTheory.ConvolutionTheory.Dp.SSatIff}
\leanref{Development/AutomorphicArithmetic/Foundations/Convolution/CentreSplitting/SSatIff.lean}{UpperTheory.ConvolutionTheory.Dp.sSatIff}
\leanref{Development/AutomorphicArithmetic/Foundations/Convolution/TraceSplitting/CarryLift.lean}{UpperTheory.ConvolutionTheory.T.carryLift}
\leanref{Development/AutomorphicArithmetic/Foundations/Convolution/CentreSplitting/LocalRepresentationL2.lean}{UpperTheory.ConvolutionTheory.Dp.IsoEq}
\leanref{Development/AutomorphicArithmetic/Foundations/Convolution/CentreSplitting/IsoEq.lean}{UpperTheory.ConvolutionTheory.Dp.isoEq}

\begin{proof}
(3) The positive real scalars are central, so by irreducibility the generator of their Lie algebra
acts on $\pi$ by a scalar, and they act on $\pi$ by $z\mapsto|\det z|^c$ for some $c\in\mathbf C$. For
$s=-c$ the forms of $\pi\otimes|\det|^s$ are invariant under them, and (2) gives the cuspidal summand.

(5) By (1) the cuspidal summand carries some $\pi$. This $\pi$ has an infinitesimal character at every
embedding and Satake parameters at all but finitely many finite places, so (4) gives $(\Lambda,t)$.
\end{proof}

\begin{lemma}\label{gal-depth-zero-lift-stable-twist}
Let $M$ be a totally complex number field, $n\geq1$, $P$ a cuspidal automorphic representation of
$\operatorname{GL}_n(\mathbf A_M)$ and $V$ a cuspidal summand that carries $P$
(Definition~\ref{gal-depth-zero-lift-stable-spectrum}). For an automorphism $\sigma$ of $M$, acting
entrywise on $\operatorname{GL}_n(\mathbf A_M)$, let $P\circ\sigma$ be the space of the functions
$\varphi\circ\sigma$ with $\varphi\in P$.
\begin{enumerate}
\item For every automorphism $\sigma$ of $M$, $P\circ\sigma\cong P$ if and only if $\theta_\sigma V=V$.
\item Let $K\subset M$ be a subfield with $M$ a Galois extension of $K$ of prime degree, and $\sigma$ a
  generator of $\operatorname{Gal}(M/K)$. Then $P\circ\gamma\cong P$ for every
  $\gamma\in\operatorname{Gal}(M/K)$ if and only if $\theta_\sigma V=V$.
\end{enumerate}
\end{lemma}
\leanref{Development/AutomorphicArithmetic/Foundations/Convolution/CentreSplitting/DefsTw.lean}{UpperTheory.ConvolutionTheory.Dp.ConjIff}
\leanref{Development/AutomorphicArithmetic/Foundations/Convolution/CentreSplitting/ConjIff.lean}{UpperTheory.ConvolutionTheory.Dp.conjIff}
\leanref{Development/AutomorphicArithmetic/Foundations/Convolution/CentreSplitting/DefsTw.lean}{UpperTheory.ConvolutionTheory.Dp.ConjCarries}
\leanref{Development/AutomorphicArithmetic/Foundations/Convolution/CentreSplitting/ConjCarries.lean}{UpperTheory.ConvolutionTheory.Dp.conjCarries}
\leanref{Development/AutomorphicArithmetic/Foundations/Convolution/TraceSplitting/NormCharSplit.lean}{UpperTheory.ConvolutionTheory.T.CarryLift}
\leanref{Development/AutomorphicArithmetic/Foundations/Convolution/TraceSplitting/CarryLift.lean}{UpperTheory.ConvolutionTheory.T.carryLift}

\begin{proof}
(1) Since $\sigma$ preserves $\operatorname{GL}_n(M)$, the idelic norm and the positive real scalars,
$P\circ\sigma$ is a cuspidal automorphic representation of
$\operatorname{GL}_n(\mathbf A_M)$, and $\theta_\sigma$ preserves the cuspidal part of $L^2_n(M)$ with
$\theta_\sigma R(g)\theta_\sigma^{-1}=R(\sigma(g))$. For $\varphi\in P$, $\theta_\sigma$ carries the
restriction of $\varphi\circ\sigma$ to that of $\varphi$. So, by
Lemma~\ref{gal-depth-zero-lift-stable-carried}(2), a cuspidal summand $V'$ carries $P\circ\sigma$ if and
only if $\theta_\sigma V'$ carries $P$. If $P\circ\sigma\cong P$, then $P\circ\sigma=P$ by
Lemma~\ref{gal-depth-zero-lift-stable-carried}(6), so $V$ carries $P\circ\sigma$ and $\theta_\sigma V$
carries $P$; since only one cuspidal summand carries $P$, $\theta_\sigma V=V$. Conversely, if
$\theta_\sigma V=V$, then $\theta_\sigma V$ carries $P$, so $V$ carries $P\circ\sigma$, and
$P\circ\sigma=P$ by Lemma~\ref{gal-depth-zero-lift-stable-carried}(1).

(2) If $P\circ\gamma\cong P$ for every $\gamma$, then $\theta_\sigma V=V$ by (1) for $\sigma$.
Conversely, if $\theta_\sigma V=V$, the proof of (1) gives $P\circ\sigma=P$, hence
$P\circ\sigma^k=P$ for every $k\geq0$ by induction on $k$, and every $\gamma\in\operatorname{Gal}(M/K)$
is a power of $\sigma$.
\end{proof}

\begin{proposition}\label{gal-depth-zero-lift-stable-summand}
Let $K\subset M$ be totally complex number fields with $M$ a Galois extension of $K$ of prime degree,
$\sigma$ a generator of $\operatorname{Gal}(M/K)$, and $n\geq1$. Cuspidal summands, their parameters and
the Galois twist $\theta_\sigma$ are those of Definition~\ref{gal-depth-zero-lift-stable-spectrum};
base changes of parameters, avoiding parabolic data and weak base changes to sums are those of
Definition~\ref{gal-depth-zero-lift-stable-parabolic}.
\begin{enumerate}
\item Every cuspidal automorphic representation of $\operatorname{GL}_n(\mathbf A_K)$ has a weak base
  change to a sum.
\item Let $V$ be a cuspidal summand for $\operatorname{GL}_n$ over $K$ with parameters $(\Lambda,t)$ off a
  finite set, such that $(\operatorname{bc}\Lambda,\operatorname{bc}t)$ avoids parabolic data over $K$.
  Then there is a cuspidal summand $V'$ for $\operatorname{GL}_n$ over $M$ with $\theta_\sigma V'=V'$ and
  with parameters $(\operatorname{bc}\Lambda,\operatorname{bc}t)$ off a finite set.
\end{enumerate}
Both are forms of the cyclic base change of Arthur and Clozel \cite[Chapter~3]{AC89}, which they prove by
comparing the trace formula of $\operatorname{GL}_n$ over $K$ with the trace formula of
$\operatorname{GL}_n$ over $M$ twisted by $\sigma$. In (2), condition (i) of avoiding parabolic data
makes the base change of the representation carried by $V$ cuspidal, and by strong multiplicity one
\cite{JS81a} and multiplicity one \cite{Sha74} the cuspidal summand that carries it is stable under
$\theta_\sigma$.
\end{proposition}
\leanref{Development/AutomorphicArithmetic/Foundations/Convolution/TraceSplitting/NormCharSplit.lean}{UpperTheory.ConvolutionTheory.T.LiftAll}
\leanref{Development/AutomorphicArithmetic/Foundations/Convolution/SphericalIdeles/LiftAll.lean}{UpperTheory.ConvolutionTheory.Sph.liftAll}
\leanref{Development/AutomorphicArithmetic/Foundations/Convolution/TraceSplitting/NormCharSplit.lean}{UpperTheory.ConvolutionTheory.T.CoreLift}
\leanref{Development/AutomorphicArithmetic/Foundations/Convolution/TraceSplitting/CoreLift.lean}{UpperTheory.ConvolutionTheory.T.coreLift}

\begin{remark}\label{gal-depth-zero-lift-stable-summand-comparison}
Alpöge proves (2) of Proposition~\ref{gal-depth-zero-lift-stable-summand} with his own comparison of a
trace formula for $\operatorname{GL}_n$ over $K$ with one over $M$ twisted by $\theta_\sigma$, and
deduces (1) from (2) by induction on $n$. The test functions match at every place: matching locally
constant functions at the places of a finite set $S$, the characteristic functions of Iwahori subgroups
at a place $q\notin S$ of $K$ split in $M$ and at the places of $M$ over it, the unit elements of the
spherical Hecke algebras at the other places, matched by the base change fundamental lemma for the unit
element \cite{Kot86}, and matching functions at the infinite places. Spherical Hecke operators that
separate $(\operatorname{bc}\Lambda,\operatorname{bc}t)$ from the parameters of parabolic data make the
contribution of the noncuspidal spectrum negligible, and this is where the hypothesis of (2) enters.
The comparison then identifies, for all large $j$, the trace of the $j$-th power of an Iwahori-Hecke
operator at a place over $q$, composed with a convolution operator and $\theta_\sigma$, over
the $\theta_\sigma$-stable cuspidal summands over $M$ with parameters
$(\operatorname{bc}\Lambda,\operatorname{bc}t)$, with a nonzero multiple of the corresponding trace over
the cuspidal summands over $K$ whose parameters have these base changes. A positivity argument shows
that the second trace is nonzero for infinitely many $j$, since $V$ contributes to it, so some
cuspidal summand contributes to the first; it is the summand $V'$ of (2).
\end{remark}

\begin{lemma}\label{gal-depth-zero-lift-stable-avoid}
Let $K\subset M$ be totally complex number fields with $M$ a Galois extension of $K$ of prime degree,
and $n\geq1$. Let $\pi$ and $P$ be cuspidal automorphic representations of $\operatorname{GL}_n(\mathbf
A_K)$ and $\operatorname{GL}_n(\mathbf A_M)$ with $P$ a weak base change of $\pi$
(Definition~\ref{gal-depth-zero-lift-weak}), and let $V$ be a cuspidal summand that carries $\pi$ and has
parameters $(\Lambda,t)$ off a finite set (Definition~\ref{gal-depth-zero-lift-stable-spectrum}). Then
$(\operatorname{bc}\Lambda,\operatorname{bc}t)$ avoids parabolic data over $K$
(Definition~\ref{gal-depth-zero-lift-stable-parabolic}).
\end{lemma}
\leanref{Development/AutomorphicArithmetic/Foundations/Convolution/SphericalIdeles/SphConj.lean}{UpperTheory.ConvolutionTheory.Sph.Norm}
\leanref{Development/AutomorphicArithmetic/Foundations/Convolution/SphericalIdeles/SameSat.lean}{UpperTheory.ConvolutionTheory.Sph.norm}
\leanref{Development/AutomorphicArithmetic/Foundations/Convolution/TraceSplitting/NormCharSplit.lean}{UpperTheory.ConvolutionTheory.T.FamParts}
\leanref{Development/AutomorphicArithmetic/Foundations/Convolution/TraceSplitting/FamParts.lean}{UpperTheory.ConvolutionTheory.T.famParts}
\leanref{Development/AutomorphicArithmetic/Foundations/Convolution/TraceSplitting/NormCharSplit.lean}{UpperTheory.ConvolutionTheory.T.FamSat}
\leanref{Development/AutomorphicArithmetic/Foundations/Convolution/TraceSplitting/FamSat.lean}{UpperTheory.ConvolutionTheory.T.famSat}
\leanref{Development/AutomorphicArithmetic/Foundations/Convolution/NormCharacterPowers/LiftSum.lean}{UpperTheory.ConvolutionTheory.R.liftSum}
\leanref{Development/AutomorphicArithmetic/Foundations/Convolution/NormCharacterPowers/NormCharPowHolds.lean}{UpperTheory.ConvolutionTheory.R.partA_holds}

\begin{proof}
Suppose not, and let $S$ be a finite set of finite places of $K$ for which (i) or (ii) of
Definition~\ref{gal-depth-zero-lift-stable-parabolic}(3) fails, with $S_M$ the set of places of $M$ over
$S$. We first find cuspidal automorphic representations $Q_1,\ldots,Q_k$ of
$\operatorname{GL}_{m_j}(\mathbf A_M)$, with $k\geq2$ and all $m_j\geq1$, such that at all but finitely
many finite places $w$ of $M$ each $Q_j$ has Satake parameters $B_j$ with
$\sum_jB_j=\operatorname{bc}t(w)$.

\emph{A parabolic datum and its representations.} Let a parabolic datum $(r,d_i,V_i,\lambda_i)$ over a
totally complex field $L$ have parameters $(\Lambda_0,t_0)$ off a set $T$, through $(m_i,B_i)$. By
Lemma~\ref{gal-depth-zero-lift-stable-carried}(1), $V_i$ carries a cuspidal automorphic representation
$\pi_i$ of $\operatorname{GL}_{d_i}(\mathbf A_L)$, and by Lemma~\ref{gal-depth-zero-lift-stable-carried}(4),
$\pi_i$ has Satake parameters $B_i(w)$ at every finite place $w\notin T$. So
$\pi_i\otimes|\det|^{\lambda_i}$ has Satake parameters $\{q_w^{-\lambda_i}b:b\in B_i(w)\}$ at $w$, and
their sum over $i$ is $t_0(w)$.

\emph{If (i) fails.} Apply this to a parabolic datum over $M$ with parameters
$(\operatorname{bc}\Lambda,\operatorname{bc}t)$ off $S_M$, and take $Q_i=\pi_i\otimes|\det|^{\lambda_i}$
for $i\leq r$, so $k=r\geq2$.

\emph{If (ii) fails.} There is a parabolic datum $(r,d_i,V_i,\lambda_i)$ over $K$ with parameters
$(\Lambda',t')$ off $S$ and $\operatorname{bc}t(w)=\operatorname{bc}t'(w)$ for every finite place
$w\notin S_M$. Let $\pi'_i=\pi_i\otimes|\det|^{\lambda_i}$ as above, so that at every $v\notin S$ the sum
over $i$ of the Satake parameters of $\pi'_i$ is $t'(v)$. Since $r\geq2$ and every $d_j\geq1$, each
$d_i<n$. By Proposition~\ref{gal-depth-zero-lift-stable-summand}(1) in rank $d_i$, each $\pi'_i$ has a
weak base change $(Q_{i,j})_j$ to a sum. At all but finitely many finite places $w$ of $M$, with $v\notin
S$ the place of $K$ under $w$ and $N(w)=N(v)^f$, so that $f=f_w$, the Satake parameters of the $Q_{i,j}$
at $w$ add up over $j$ to the $f$-th powers of those of $\pi'_i$ at $v$. Over all pairs $(i,j)$ they
add up to $\{b^f:b\in t'(v)\}=\operatorname{bc}t'(w)=\operatorname{bc}t(w)$, and there are at least
$r\geq2$ pairs.

\emph{The contradiction.} By Lemma~\ref{gal-depth-zero-lift-stable-carried}(4), $\pi$ has Satake
parameters $t(v)$ at all but finitely many finite places $v$ of $K$. Since $P$ is a weak base change of
$\pi$ and the Satake parameters of $\pi$ at a place are unique, at all but finitely many finite places
$w$ of $M$, with $v$ under $w$, $P$ has Satake parameters $\{b^{f_w}:b\in t(v)\}=\operatorname{bc}t(w)$.
So at all but finitely many $w$ the Satake parameters of $P$ are the sum of those of the $Q_j$.
Theorem~\ref{lit-jacquet-shalika}(1), applied to the family consisting of $P$ alone and to the family
$(Q_j)$, gives a bijection from a set with one element to a set with $k\geq2$ elements, which is absurd.
\end{proof}

\begin{lemma}\label{gal-depth-zero-lift-infinity}
Let $E$ be a totally complex number field, $\Omega$ an algebraic closure of $E$, $a\geq1$,
$X\subset\Omega$ a finite extension of $E$, and $\pi$ and $P$ cuspidal automorphic representations of
$\operatorname{GL}_a(\mathbf A_E)$ and $\operatorname{GL}_a(\mathbf A_X)$ with a chain of weak base
changes from $(E,\pi)$ to $(X,P)$ (Definition~\ref{gal-depth-zero-lift-weak}). Then
$Z_\tau(P)=Z_{\tau|_E}(\pi)$ for every embedding $\tau:X\to\mathbf C$.
\end{lemma}
\leanref{Development/AutomorphicArithmetic/Foundations/GaloisRepresentations/SphericalGaloisCompatibility/PAdicSplitting/InfLifts.lean}{UpperTheory.Gal.ToralGalois.R2.InfLifts}
\leanref{Development/AutomorphicArithmetic/Foundations/GaloisRepresentations/SphericalGaloisCompatibility/PAdicSplitting/PInfLifts.lean}{UpperTheory.Gal.ToralGalois.R2.infLifts}
\leanref{Development/AutomorphicArithmetic/Foundations/GaloisRepresentations/SphericalGaloisCompatibility/PAdicSplitting/InfLifts.lean}{UpperTheory.Gal.ToralGalois.R2.InfStep}
\leanref{Development/AutomorphicArithmetic/Foundations/GaloisRepresentations/SphericalGaloisCompatibility/PAdicSplitting/PInfStep.lean}{UpperTheory.Gal.ToralGalois.R2.infStep}
\leanref{Development/AutomorphicArithmetic/Foundations/GaloisRepresentations/RegularSemisimpleFibres/ResEmbCompat.lean}{UpperTheory.Gal.InertialTypesAlt.StableOf}
\leanref{Development/AutomorphicArithmetic/Foundations/GaloisRepresentations/RegularSemisimpleFibres/PStableOf.lean}{UpperTheory.Gal.InertialTypesAlt.stableOf}

\begin{proof}
By induction on the length of the chain it suffices to treat one step: $K\subset M$ totally complex,
$M$ Galois over $K$ of prime degree, and $\Pi$ a weak base change of a cuspidal $\pi_1$ of
$\operatorname{GL}_a(\mathbf A_K)$; we show $Z_\tau(\Pi)=Z_{\tau|_K}(\pi_1)$.

\emph{Stability.} We show that $\Pi\circ\gamma\cong\Pi$ for every $\gamma\in\operatorname{Gal}(M/K)$.
This group has prime order, so it is generated by one element $\sigma$. By
Lemma~\ref{gal-depth-zero-lift-stable-carried}(3) there is $s\in\mathbf C$ such that
$\pi_1\otimes|\det|^s$ is carried by a cuspidal summand $V$ for $\operatorname{GL}_a$ over $K$
(Definition~\ref{gal-depth-zero-lift-stable-spectrum}). Then $\Pi\otimes|\det|^s$ is a weak base change
of $\pi_1\otimes|\det|^s$: if $u$ is a place of $M$ over $v$ with $N(u)=N(v)^f$, and $\pi_1$ and $\Pi$
have Satake parameters $B$ at $v$ and $\{b^f:b\in B\}$ at $u$, then the twists have Satake parameters
$N(v)^{-s}B$ at $v$ and $N(u)^{-s}\{b^f:b\in B\}=\{(N(v)^{-s}b)^f:b\in B\}$ at $u$. By
Lemma~\ref{gal-depth-zero-lift-stable-carried}(5), $V$ has parameters $(\Lambda,t)$ off a finite set. By
Lemma~\ref{gal-depth-zero-lift-stable-avoid}, applied to $\pi_1\otimes|\det|^s$, $\Pi\otimes|\det|^s$
and $V$, the pair $(\operatorname{bc}\Lambda,\operatorname{bc}t)$ avoids parabolic data over $K$
(Definition~\ref{gal-depth-zero-lift-stable-parabolic}). Proposition~\ref{gal-depth-zero-lift-stable-summand}(2)
gives a cuspidal summand $V'$ for $\operatorname{GL}_a$ over $M$ with $\theta_\sigma V'=V'$ and with
parameters $(\operatorname{bc}\Lambda,\operatorname{bc}t)$ off a finite set. By
Lemma~\ref{gal-depth-zero-lift-stable-carried}(1), $V'$ carries a cuspidal automorphic representation
$P'$ of $\operatorname{GL}_a(\mathbf A_M)$. By Lemma~\ref{gal-depth-zero-lift-stable-carried}(4), at all
but finitely many places $u$ of $M$, with $v$ the place of $K$ under $u$ and $N(u)=N(v)^f$, the
representation $\pi_1\otimes|\det|^s$ has Satake parameters $t(v)$ at $v$ and $P'$ has Satake
parameters $\operatorname{bc}t(u)=\{b^f:b\in t(v)\}$ at $u$; so $P'$ is a weak base change of
$\pi_1\otimes|\det|^s$. Two weak base changes of $\pi_1\otimes|\det|^s$ have the same Satake parameters
at all but finitely many places, since the Satake parameters of $\pi_1\otimes|\det|^s$ at a place are
unique. So $P'\cong\Pi\otimes|\det|^s$ by Theorem~\ref{lit-jacquet-shalika}(1), and
$P'=\Pi\otimes|\det|^s$ by Lemma~\ref{gal-depth-zero-lift-stable-carried}(6). Now $V'$ carries
$\Pi\otimes|\det|^s$ and $\theta_\sigma V'=V'$, so $(\Pi\otimes|\det|^s)\circ\gamma\cong
\Pi\otimes|\det|^s$ for every $\gamma$ by Lemma~\ref{gal-depth-zero-lift-stable-twist}(2). Since
$\gamma$ preserves the idelic norm, $(\Pi\otimes|\det|^s)\circ\gamma=(\Pi\circ\gamma)\otimes|\det|^s$,
and twisting by $|\det|^{-s}$ gives $\Pi\circ\gamma\cong\Pi$.
% Lean: stableOf uses UpperTheory.ConvolutionTheory.Sph.weakStable (WeakStableAlt.lean), whose steps
% are those above: Dp.exTwist and T.powTwist for s, Sph.norm for (Lambda, t) and OffAll, T.coreLift,
% Dp.exCusp, the third and last components of T.carryLift, weak_same, and Sph.smoEq.

\emph{Descent.} By Theorem~\ref{gal-depth-zero-lift-base-change}(2) there is a cuspidal $\pi_2$ of which
$\Pi$ is a weak base change, with $Z_\tau(\Pi)=Z_{\tau|_K}(\pi_2)$ for every $\tau$, and by
Theorem~\ref{gal-depth-zero-lift-base-change}(3), $\pi_2\cong\pi_1\otimes(\eta\circ\det)$ with
$\eta=1$ or $\eta$ a norm character.

\emph{Infinity types.} Since $\eta$ has finite order and $K$ is totally complex, $\eta$ is trivial at
the infinite places, so $Z_{\tau|_K}(\pi_2)=Z_{\tau|_K}(\pi_1)$.
\end{proof}

\begin{remark}\label{gal-depth-zero-lift-infinity-stable}
The stability $\Pi\circ\gamma\cong\Pi$ in the proof of Lemma~\ref{gal-depth-zero-lift-infinity} also
follows from strong multiplicity one alone. At all but finitely many places $u$ of $M$, the Satake
parameters of $\Pi\circ\gamma$ at $u$ are those of $\Pi$ at a place conjugate to $u$ by $\gamma$. These two
places lie over the same place of $K$ with the same residue degree $f$, so both sets of parameters are the
$f$-th powers of the Satake parameters of $\pi_1$ there, and $\Pi\circ\gamma\cong\Pi$ by \cite{JS81a}.
\end{remark}

\subsection{Level-zero patterns and the Shintani lift}\label{subsection-gal-shintani-lift}

\begin{definition}\label{gal-support-level-zero}
Let $k$ be a finite field with $q$ elements, $\bar k$ an algebraic closure of $k$, and
$\Theta:\bar k^\times\to\mathbf C^\times$ an injective homomorphism.
\begin{enumerate}
\item Let $H$ be a finite group and $T,\varphi:H\to\mathbf C$. A finite-dimensional complex
  representation $\rho$ of $H$ \emph{meets} $T$ if $\sum_{h\in H}\overline{T(h)}\operatorname{tr}\rho(h)\neq0$.
  The function $\varphi$ is \emph{carried by} $T$ if $\sum_{h\in H}\overline{\varphi(h)}\rho(h)=0$ for
  every finite-dimensional representation $\rho$ of $H$ that does not meet $T$.
\item Let $d=(d_1,\ldots,d_r)$ with $d_i\geq1$ and $n=d_1+\cdots+d_r$. Let $P_d\subset\operatorname{GL}_n(k)$
  be the group of invertible block upper triangular matrices with diagonal blocks of sizes
  $d_1,\ldots,d_r$, and $\pi_i(p)\in\operatorname{GL}_{d_i}(k)$ the $i$-th diagonal block of $p\in P_d$.
  The \emph{Harish-Chandra induction} of functions $Y_i$ on $\operatorname{GL}_{d_i}(k)$ is the
  function on $\operatorname{GL}_n(k)$
  \[
   R_d(Y_1,\ldots,Y_r)(x)=\frac1{|P_d|}\sum_{y\in\operatorname{GL}_n(k),\ y^{-1}xy\in P_d}\ \prod_{i=1}^r
   Y_i\bigl(\pi_i(y^{-1}xy)\bigr).
  \]
\item For $m\geq1$ and $c\geq0$, let $\delta_\Theta(m,c)$ be the function $x\mapsto\Theta(\det x)^c$ on
  $\operatorname{GL}_m(k)$. For an automorphism $h$ of a finite-dimensional $k$-vector space $V$, define
  $\psi_c(h)$ by induction on $\dim V$:
  \[
   \psi_c(h)=\sum_{\alpha}\Theta(\alpha)^c-\sum_{W}\psi_c(h|_W),
  \]
  where $\alpha$ runs over the roots in $\bar k$ of the characteristic polynomial of $h$, with
  multiplicity, and $W$ over the $h$-stable subspaces $W\neq V$; thus $\psi_c$ vanishes on the zero
  space. Let $\kappa_\Theta(m,c)(x)=(-1)^{m-1}\psi_c(x)$ for $x\in\operatorname{GL}_m(k)$ acting on
  $k^m$. The integer $c$ is \emph{regular for} $(q,m)$ if $cq^i\not\equiv c\pmod{q^m-1}$ for $0<i<m$.
\item For finite lists $D=((d_1,c_1),\ldots,(d_r,c_r))$ and $L=((m_1,e_1),\ldots,(m_s,e_s))$ of pairs of
  integers with $d_i,m_j\geq1$ and $c_i,e_j\geq0$, let $n=\sum_id_i+\sum_jm_j$ and
  \[
   J_\Theta(D;L)=R_{(d_1,\ldots,d_r,m_1,\ldots,m_s)}\bigl(\delta_\Theta(d_1,c_1),\ldots,\delta_\Theta(d_r,c_r),
   \kappa_\Theta(m_1,e_1),\ldots,\kappa_\Theta(m_s,e_s)\bigr),
  \]
  a function on $\operatorname{GL}_n(k)$.
\item Let $k'$ be an extension of $k$ of degree $l$ and $s$ the automorphism $z\mapsto z^q$ of $k'$.
  For $\delta\in\operatorname{GL}_n(k')$ put $N_s(\delta)=\delta\,s(\delta)\cdots s^{l-1}(\delta)$, with $s$
  applied to the entries. The \emph{Shintani lift} of a class function $f$ on $\operatorname{GL}_n(k)$ is
  the function $\operatorname{Sh}_s(f)$ on $\operatorname{GL}_n(k')$ with
  $\operatorname{Sh}_s(f)(\delta)=f(\gamma)$ if $N_s(\delta)$ is conjugate in $\operatorname{GL}_n(k')$ to
  an element $\gamma$ of $\operatorname{GL}_n(k)$, and $\operatorname{Sh}_s(f)(\delta)=0$ otherwise. Two
  elements of $\operatorname{GL}_n(k)$ that are conjugate in $\operatorname{GL}_n(k')$ are conjugate in
  $\operatorname{GL}_n(k)$, so this does not depend on $\gamma$.
\item Let $F$ be a nonarchimedean local field with residue field $k$, and $K_1$ the kernel of reduction
  $\operatorname{GL}_n(\mathcal O_F)\to\operatorname{GL}_n(k)$, $g\mapsto\bar g$. The \emph{level-zero part}
  of a smooth representation $\sigma$ of $\operatorname{GL}_n(F)$ on $V$ is the representation $\sigma^0$
  of $\operatorname{GL}_n(k)$ on $V^{K_1}$ with $\sigma^0(\bar g)=\sigma(g)|_{V^{K_1}}$ for
  $g\in\operatorname{GL}_n(\mathcal O_F)$.
\end{enumerate}
\end{definition}
\leanref{Development/AutomorphicArithmetic/Foundations/GaloisRepresentations/RegularSemisimpleFibres/NumberField.lean}{UpperTheory.Gal.InertialTypesAlt.Meets}
\leanref{Development/AbelianVarietyTheory/Definitions/Divisors/OrthSubrepresentation.lean}{GlobalTheory.Shin.Kills}
\leanref{Development/AbelianVarietyTheory/Definitions/Divisors/OrthSubrepresentation.lean}{GlobalTheory.Shin.hc}
\leanref{Development/AutomorphicArithmetic/Foundations/GaloisRepresentations/RegularSemisimpleFibres/NumberField.lean}{UpperTheory.Gal.InertialTypesAlt.detc}
\leanref{Development/AbelianVarietyTheory/Definitions/CommutativeAlgebra/InvariantSubspaceMobiusInverseRecEqOfFinrankLt.lean}{GlobalTheory.CuspGL.psi}
\leanref{Development/AbelianVarietyTheory/Definitions/Divisors/OrthSubrepresentation.lean}{GlobalTheory.Shin.cusp}
\leanref{Development/AbelianVarietyTheory/Definitions/Divisors/OrthSubrepresentation.lean}{GlobalTheory.Shin.Reg}
\leanref{Development/AbelianVarietyTheory/Definitions/Divisors/OrthSubrepresentation.lean}{GlobalTheory.Shin.J}
\leanref{Development/AutomorphicArithmetic/Foundations/GaloisRepresentations/RegularSemisimpleFibres/NumberField.lean}{UpperTheory.Gal.InertialTypesAlt.JD}
\leanref{Development/AbelianVarietyTheory/Definitions/Divisors/OrthSubrepresentation.lean}{GlobalTheory.Shin.shLift}
\leanref{Development/AbelianVarietyTheory/Definitions/NonGaloisDescent/EmbedsLevelZeroDepthZero.lean}{GlobalTheory.LocGL.levelZero}

\begin{lemma}\label{gal-support-finite}
Let $k$, $q$, $\bar k$ and $\Theta$ be as in Definition~\ref{gal-support-level-zero}, and let $m,n\geq1$.
\begin{enumerate}
\item Every cuspidal irreducible representation of $\operatorname{GL}_m(k)$ has character
  $\kappa_\Theta(m,c)$ for some $c$ regular for $(q,m)$. Conversely, if $c$ is regular for $(q,m)$, then
  $\kappa_\Theta(m,c)$ is the character of a cuspidal irreducible representation.
\item If $c'\equiv cq^i\pmod{q^m-1}$ for some $i\geq0$, then $\kappa_\Theta(m,c')=\kappa_\Theta(m,c)$.
  Moreover $\kappa_\Theta(1,c)=\delta_\Theta(1,c)$.
\item A Harish-Chandra induction of characters of representations is the character of a
  representation. In particular $\delta_\Theta(m,c)$ is a character, and $J_\Theta(D;L)$ is a character
  whenever $e_j$ is regular for $(q,m_j)$ for every $j$.
\item For $d=(d_1,\ldots,d_r)$ with sum $n$, all functions $Y_i$ on $\operatorname{GL}_{d_i}(k)$ and every
  class function $T$ on $\operatorname{GL}_n(k)$,
  \[
   \sum_{x\in\operatorname{GL}_n(k)}\overline{R_d(Y_1,\ldots,Y_r)(x)}\,T(x)
   =\frac{|\operatorname{GL}_n(k)|}{|P_d|}\sum_{p\in P_d}\overline{\textstyle\prod_iY_i(\pi_i(p))}\,T(p).
  \]
\item Let $H$ be a finite group and $\rho$, $\lambda$ finite-dimensional representations of $H$. Then
  $\rho$ meets the character of $\lambda$ if and only if there is a nonzero $H$-equivariant map
  $\rho\to\lambda$, if and only if there is a nonzero $H$-equivariant map $\lambda\to\rho$.
\item Let $D$ and $L$ be as in Definition~\ref{gal-support-level-zero}(4), with $e_j$ regular for
  $(q,m_j)$ for every $j$, and let $D^\sharp L$ be the list obtained from $D$ by replacing each $(d,c)$ by
  $d$ copies of $(1,c)$, followed by $L$. Then
  $\sum_x\overline{J_\Theta(\emptyset;D^\sharp L)(x)}\,J_\Theta(D;L)(x)\neq0$.
\end{enumerate}
\end{lemma}
\leanref{Development/AutomorphicArithmetic/Foundations/GaloisRepresentations/RegularSemisimpleFibres/ResEmbCompat.lean}{UpperTheory.Gal.InertialTypesAlt.CuspChar}
\leanref{Development/AutomorphicArithmetic/Foundations/GaloisRepresentations/RegularSemisimpleFibres/PCuspChar.lean}{UpperTheory.Gal.InertialTypesAlt.cuspChar}
\leanref{Development/AutomorphicArithmetic/Foundations/GaloisRepresentations/RegularSemisimpleFibres/SuppBCOfMix.lean}{UpperTheory.Gal.InertialTypesAlt.CuspFrob}
\leanref{Development/AutomorphicArithmetic/Foundations/GaloisRepresentations/RegularSemisimpleFibres/PCuspFrob.lean}{UpperTheory.Gal.InertialTypesAlt.cuspFrob}
\leanref{Development/AutomorphicArithmetic/Foundations/GaloisRepresentations/RegularSemisimpleFibres/ResEmbCompat.lean}{UpperTheory.Gal.InertialTypesAlt.JDEq}
\leanref{Development/AutomorphicArithmetic/Foundations/GaloisRepresentations/RegularSemisimpleFibres/ResEmbCompat.lean}{UpperTheory.Gal.InertialTypesAlt.HcChar}
\leanref{Development/AutomorphicArithmetic/Foundations/GaloisRepresentations/RegularSemisimpleFibres/ResEmbCompat.lean}{UpperTheory.Gal.InertialTypesAlt.JDChar}
\leanref{Development/AutomorphicArithmetic/Foundations/GaloisRepresentations/RegularSemisimpleFibres/PJdChar.lean}{UpperTheory.Gal.InertialTypesAlt.jdChar}
\leanref{Development/AutomorphicArithmetic/Foundations/GaloisRepresentations/RegularSemisimpleFibres/ResEmbCompat.lean}{UpperTheory.Gal.InertialTypesAlt.FrobL}
\leanref{Development/AutomorphicArithmetic/Foundations/GaloisRepresentations/RegularSemisimpleFibres/PFrobL.lean}{UpperTheory.Gal.InertialTypesAlt.frobL}
\leanref{Development/AutomorphicArithmetic/Foundations/GaloisRepresentations/RegularSemisimpleFibres/ResEmbCompat.lean}{UpperTheory.Gal.InertialTypesAlt.MeetsIff}
\leanref{Development/AutomorphicArithmetic/Foundations/GaloisRepresentations/RegularSemisimpleFibres/PMeetsIff.lean}{UpperTheory.Gal.InertialTypesAlt.meetsIff}
\leanref{Development/AutomorphicArithmetic/Foundations/GaloisRepresentations/RegularSemisimpleFibres/SuppBCOfMix.lean}{UpperTheory.Gal.InertialTypesAlt.JDPair}
\leanref{Development/AutomorphicArithmetic/Foundations/GaloisRepresentations/RegularSemisimpleFibres/PJdPair.lean}{UpperTheory.Gal.InertialTypesAlt.jdPair}

\begin{proof}
(1) For $x$ with irreducible characteristic polynomial, the only $x$-stable subspaces of $k^m$ are $0$
and $k^m$, so $\kappa_\Theta(m,c)(x)=(-1)^{m-1}\sum_{i<m}\Theta(z^{q^i})^c$ for a root $z$ of that
polynomial. This is the value of Green's character $\tau_\theta$ for the character $\theta=\Theta^c$ of
the subfield of $\bar k$ with $q^m$ elements, and $\theta$ is regular if and only if $c$ is regular for
$(q,m)$, since $\Theta$ is injective. So (1) is Green's theorem (Theorem~\ref{lit-green-cuspidal}),
together with his formula for the values of cuspidal characters at all elements \cite{Gre55}. Alpöge
proves the first assertion by showing that the character of a cuspidal representation sums to $0$ over
every coset of the unipotent radical of a proper standard parabolic subgroup, and that such a class
function is a combination of the functions $\psi_c$.

(2) The second assertion is the definition for $m=1$. The first is a computation: $\psi_c$ is a class
function whose sums over the cosets of unipotent radicals of proper parabolic subgroups vanish; such a
function on $\operatorname{GL}_m(k)$ is determined by its inner products with the functions
$x\mapsto\sum_\alpha\Theta(\alpha)^e$, $0\leq e<q^m-1$; and these inner products count the $h<m$ with
$cq^h\equiv e\pmod{q^m-1}$, a number that does not change when $c$ is replaced by $c'$.

(3) $R_d(\chi_1,\ldots,\chi_r)$ is the character of the representation induced from $P_d$ by the
inflation of $\rho_1\otimes\cdots\otimes\rho_r$, by the formula for induced characters. The function
$\delta_\Theta(m,c)$ is a one-dimensional character, and (1) gives the rest.

(4) Substitute $p=y^{-1}xy$: the left side is
$|P_d|^{-1}\sum_{y}\sum_{p\in P_d}\overline{\prod_iY_i(\pi_i(p))}\,T(ypy^{-1})$, and $T(ypy^{-1})=T(p)$.

(5) This is Schur orthogonality, since $\rho$ and $\lambda$ are semisimple.

(6) By (4), the representation with character $R_{(1,\ldots,1)}(\delta_\Theta(1,c),\ldots,\delta_\Theta(1,c))$
of $\operatorname{GL}_d(k)$ contains the one with character $\delta_\Theta(d,c)$, since the restriction of
$\delta_\Theta(d,c)$ to the upper triangular matrices is the inflation of
$\delta_\Theta(1,c)\otimes\cdots\otimes\delta_\Theta(1,c)$. By transitivity of induction,
$J_\Theta(\emptyset;D^\sharp L)$ is the character of a representation that contains one with character
$J_\Theta(D;L)$, and $J_\Theta(D;L)(1)\neq0$. Writing both characters as sums of irreducible characters
with multiplicities, the sum is at least $\sum_x|J_\Theta(D;L)(x)|^2>0$.
\end{proof}

\begin{definition}\label{gal-support-pattern}
A \emph{level-zero pattern}, or \emph{pattern}, is a triple $a=(q;D;L)$ of an integer $q\geq2$, its
\emph{base}, and finite lists $D=((d_1,c_1),\ldots,(d_r,c_r))$ and $L=((m_1,e_1),\ldots,(m_s,e_s))$ of
pairs of integers with $d_i,m_j\geq1$, $c_i,e_j\geq0$, and $e_j$ regular for $(q,m_j)$
(Definition~\ref{gal-support-level-zero}(3)) for every $j$. Its \emph{rank} is $\sum_id_i+\sum_jm_j$.
\begin{enumerate}
\item For a prime $l$, the \emph{lift} of $a$ by $l$ is $a^{(l)}=(q^l;D^{(l)};L^{(l)})$. Here
  $D^{(l)}$ replaces each $(d,c)$ by $(d,c(1+q+\cdots+q^{l-1}))$, and $L^{(l)}$ replaces each $(m,e)$ by
  the $l$ pairs $(m/l,eq^i)$, $0\leq i<l$, if $l$ divides $m$, and by the pair
  $(m,e(1+q^m+\cdots+q^{m(l-1)}))$ otherwise. Write $a\leq b$ if $b$ is obtained from $a$ by finitely many
  lifts by primes.
\item $a$ \emph{splits at} $Q$ if $Q=q^f$ for an integer $f\geq1$ divisible by every $m_j$. The
  \emph{blocks} of $a$ at such a $Q$ are the pairs, with multiplicity, consisting of
  $(d_i,c_i(Q-1)/(q-1))$ for each $i$, and of $(1,e_jq^h(Q-1)/(q^{m_j}-1))$ for each $j$ and
  $0\leq h<m_j$, the second entries taken modulo $Q-1$. For $x\in\mathbf C$ with $x^{Q-1}=1$ let
  $\operatorname{pol}_{a,Q}(x)=\prod(Y-x^\beta)^\mu$, the product over the blocks $(\mu,\beta)$.
\item $a$ is \emph{apart} if $c_i<q-1$ for every $i$ and the $c_i$ are pairwise distinct, every
  $m_j\geq2$, and for $j\neq j'$ with $m_j=m_{j'}$ there is no $h\geq0$ with
  $e_{j'}\equiv e_jq^h\pmod{q^{m_j}-1}$.
\item Let $F$ be a nonarchimedean local field with residue field $k$, $\bar k$ an algebraic closure of
  $k$ and $\Theta:\bar k^\times\to\mathbf C^\times$ an injective homomorphism. A smooth representation
  $\sigma$ of $\operatorname{GL}_n(F)$ \emph{has the pattern} $a$ with respect to $\Theta$ if $q=|k|$, $n$ is
  the rank of $a$, the space of $K_1$-fixed vectors of $\sigma$ is finite-dimensional, and $\sigma^0$
  meets $J_\Theta(D;L)$ (Definition~\ref{gal-support-level-zero}).
\end{enumerate}
\end{definition}
\leanref{Development/AutomorphicArithmetic/Foundations/GaloisRepresentations/RegularSemisimpleFibres/NumberField.lean}{UpperTheory.Gal.InertialTypesAlt.Lbl}
\leanref{Development/AutomorphicArithmetic/Foundations/GaloisRepresentations/RegularSemisimpleFibres/NumberField.lean}{UpperTheory.Gal.InertialTypesAlt.Lbl.Ok}
\leanref{Development/AutomorphicArithmetic/Foundations/GaloisRepresentations/RegularSemisimpleFibres/NumberField.lean}{UpperTheory.Gal.InertialTypesAlt.Lbl.lift}
\leanref{Development/AbelianVarietyTheory/Definitions/Divisors/OrthSubrepresentation.lean}{GlobalTheory.Shin.liftAll}
\leanref{Development/AutomorphicArithmetic/Foundations/GaloisRepresentations/RegularSemisimpleFibres/NumberField.lean}{UpperTheory.Gal.InertialTypesAlt.Lbl.Le}
\leanref{Development/AutomorphicArithmetic/Foundations/GaloisRepresentations/RegularSemisimpleFibres/NumberField.lean}{UpperTheory.Gal.InertialTypesAlt.Lbl.Splits}
\leanref{Development/AutomorphicArithmetic/Foundations/GaloisRepresentations/RegularSemisimpleFibres/NumberField.lean}{UpperTheory.Gal.InertialTypesAlt.Lbl.blocks}
\leanref{Development/AutomorphicArithmetic/Foundations/GaloisRepresentations/RegularSemisimpleFibres/NumberField.lean}{UpperTheory.Gal.InertialTypesAlt.Lbl.pol}
\leanref{Development/AutomorphicArithmetic/Foundations/GaloisRepresentations/RegularSemisimpleFibres/SuppBCOfMix.lean}{UpperTheory.Gal.InertialTypesAlt.Lbl.Apart}
\leanref{Development/AutomorphicArithmetic/Foundations/GaloisRepresentations/RegularSemisimpleFibres/NumberField.lean}{UpperTheory.Gal.InertialTypesAlt.HasLbl}

\begin{lemma}\label{gal-support-pattern-arithmetic}
Let $a=(q;D;L)$ be a pattern of rank $n$ (Definition~\ref{gal-support-pattern}).
\begin{enumerate}
\item For every prime $l$, $a^{(l)}$ is a pattern of rank $n$.
\item Let $b=(q';D';L')$ be a pattern with $a\leq b$ such that $a$ splits at $q'$. Then $b$ has rank $n$,
  every pair in $L'$ has first entry $1$, and $a$ and $b$ have the same blocks at $q'$.
\item Let $b=(q';D';L')$ be a pattern whose pairs in $L'$ all have first entry $1$, write $D'$ followed by
  $L'$ as $((d_1,e_1),\ldots,(d_r,e_r))$, and let $x\in\mathbf C$ with $x^{q'-1}=1$. Then the blocks of $b$
  at $q'$ are the pairs $(d_i,e_i)$, and $\operatorname{pol}_{b,q'}(x)=\prod_i(Y-x^{e_i})^{d_i}$.
\item If $D$ is empty, $a$ splits at $Q$ and $x^{Q-1}=1$, then
  $\operatorname{pol}_{a,Q}(x)=\prod_{(m,e)\in L}\prod_{h=0}^{m-1}(Y-x^{eq^h(Q-1)/(q^m-1)})$.
\item If $a$ is apart and splits at $Q$, then the second entries of the blocks of $a$ at $Q$ are pairwise
  distinct modulo $Q-1$.
\end{enumerate}
\end{lemma}
\leanref{Development/AutomorphicArithmetic/Foundations/GaloisRepresentations/RegularSemisimpleFibres/ResEmbCompat.lean}{UpperTheory.Gal.InertialTypesAlt.LiftOk}
\leanref{Development/AutomorphicArithmetic/Foundations/GaloisRepresentations/RegularSemisimpleFibres/PLiftOk.lean}{UpperTheory.Gal.InertialTypesAlt.liftOk}
\leanref{Development/AutomorphicArithmetic/Foundations/GaloisRepresentations/RegularSemisimpleFibres/ResEmbCompat.lean}{UpperTheory.Gal.InertialTypesAlt.LeTop}
\leanref{Development/AutomorphicArithmetic/Foundations/GaloisRepresentations/RegularSemisimpleFibres/PLeTop.lean}{UpperTheory.Gal.InertialTypesAlt.leTop}
\leanref{Development/AutomorphicArithmetic/Foundations/GaloisRepresentations/RegularSemisimpleFibres/ResEmbCompat.lean}{UpperTheory.Gal.InertialTypesAlt.PolTop}
\leanref{Development/AutomorphicArithmetic/Foundations/GaloisRepresentations/RegularSemisimpleFibres/PPolTop.lean}{UpperTheory.Gal.InertialTypesAlt.polTop}
\leanref{Development/AutomorphicArithmetic/Foundations/GaloisRepresentations/RegularSemisimpleFibres/ResEmbCompat.lean}{UpperTheory.Gal.InertialTypesAlt.PolBot}
\leanref{Development/AutomorphicArithmetic/Foundations/GaloisRepresentations/RegularSemisimpleFibres/PPolBot.lean}{UpperTheory.Gal.InertialTypesAlt.polBot}
\leanref{Development/AutomorphicArithmetic/Foundations/GaloisRepresentations/RegularSemisimpleFibres/SuppBCOfMix.lean}{UpperTheory.Gal.InertialTypesAlt.SepOf}
\leanref{Development/AutomorphicArithmetic/Foundations/GaloisRepresentations/RegularSemisimpleFibres/PSepOf.lean}{UpperTheory.Gal.InertialTypesAlt.sepOf}

\begin{proof}
Fix an injective homomorphism $\Theta$ from the multiplicative group of an algebraic closure of
$\mathbf F_q$ to $\mathbf C^\times$, and identify an integer $e$ modulo $q^m-1$ with the character
$\Theta^e$ of $\mathbf F_{q^m}^\times$. Then $e$ is regular for $(q,m)$ if and only if the orbit of
$\Theta^e$ under $z\mapsto z^q$ has $m$ elements, and the residue $e(Q-1)/(q^m-1)$ modulo $Q-1$ corresponds
to $\Theta^e\circ N_{\mathbf F_Q/\mathbf F_{q^m}}$ when $\mathbf F_{q^m}\subset\mathbf F_Q$.

(1) The rank does not change, since the pairs replacing $(m,e)$ have first entries summing to $m$. If
$l\mid m$, the orbit of $\Theta^{eq^i}$ under $z\mapsto z^{q^l}$ has $m/l$ elements. If $l\nmid m$, the
integer $e(1+q^m+\cdots+q^{m(l-1)})$ corresponds to $\Theta^e\circ N_{\mathbf F_{q^{lm}}/\mathbf F_{q^m}}$, whose
orbit under $z\mapsto z^{q^l}$ has $m$ elements, since $l$ is prime to $m$ and the norm is surjective.

(2) By (1) and induction along the lifts, $b$ is a pattern of rank $n$. If a lift $a'^{(l)}$ splits at
$Q$, then $a'$ splits at $Q$ with the same blocks: for $l\mid m$ the blocks of the pairs $(m/l,eq^i)$ have
residues $eq^{i+lh}(Q-1)/(q^m-1)$, $0\leq i<l$, $0\leq h<m/l$; for $l\nmid m$ the blocks of the lifted pair
have residues $eq^{lh}(Q-1)/(q^m-1)$, $0\leq h<m$, and $lh$ runs over the classes modulo $m$; and the pairs
of $D$ give $c(Q-1)/(q-1)$ in both cases. Along $a\leq b$ one also finds $q'=q^c$ with $c\geq1$, and for each
$(m',e')\in L'$ a pair $(m,e)\in L$ with $m'\gcd(m,c)=m$. Since $a$ splits at $q'=q^c$, every $m$ divides
$c$, so $m'=1$.

(3) and (4) are the definition of the blocks: for $Q=q'$ and first entry $1$ the factor
$(Q-1)/(q'-1)$ is $1$.

(5) The block of $(d,c)\in D$ corresponds to $\Theta^c\circ N_{\mathbf F_Q/\mathbf F_q}$, and the blocks of
$(m,e)\in L$ to the $\Theta^{eq^h}\circ N_{\mathbf F_Q/\mathbf F_{q^m}}$, $0\leq h<m$. Since the norms are
surjective, a character of the second kind has an orbit of $m$ elements under $z\mapsto z^q$, by
regularity, and one of the first kind is fixed. So blocks of the first kind are distinct, since the $c$
are distinct and less than $q-1$; a block of the first kind differs from one of the second, since
$m\geq2$; and two blocks of the second kind coincide only if they come from pairs $(m,e)$, $(m,e')$ with
$e'\equiv eq^h\pmod{q^m-1}$, which apartness and regularity exclude.
\end{proof}

\begin{situation}\label{gal-shintani-situation}
Let $k\subset k'$ be finite fields, $l=[k':k]$ and $s$ a generator of $\operatorname{Gal}(k'/k)$, acting on
matrices over $k'$ entry by entry. Let $n\geq0$, $G=\operatorname{GL}_n(k)$ and $G'=\operatorname{GL}_n(k')$, so
that $G$ is the set of elements of $G'$ with entries in $k$. For $\delta\in G'$ put
$N_s(\delta)=\delta\,s(\delta)\cdots s^{l-1}(\delta)$, and for a class function $f$ on $G$ let
$\operatorname{Sh}_s(f)$ be the function on $G'$ given by the formula of
Definition~\ref{gal-support-level-zero}(5) for this $s$. For $\gamma\in G$ let $k'[\gamma]\subset M_n(k')$ be
the $k'$-subalgebra generated by $\gamma$, and $Z_G(\gamma)$ the centralizer of $\gamma$ in $G$. For $a,b\geq0$
with $a+b=n$ let $P'_{(a,b)}\subset G'$ be the group of invertible block upper triangular matrices with
diagonal blocks of sizes $a$ and $b$, let $P_{(a,b)}=P'_{(a,b)}\cap G$, and let $\pi_1(p)$ and $\pi_2(p)$ be
the diagonal blocks of $p\in P'_{(a,b)}$.
\end{situation}
\leanref{Development/AbelianVarietyTheory/Definitions/Divisors/OrthSubrepresentation.lean}{GlobalTheory.Shin.galAct}
\leanref{Development/AbelianVarietyTheory/Definitions/Divisors/OrthSubrepresentation.lean}{GlobalTheory.Shin.incl}
\leanref{Development/AbelianVarietyTheory/Definitions/Divisors/OrthSubrepresentation.lean}{GlobalTheory.Shin.nrm}
\leanref{Development/AbelianVarietyTheory/Definitions/Divisors/OrthSubrepresentation.lean}{GlobalTheory.Shin.IsNorm}
\leanref{Development/AbelianVarietyTheory/Definitions/Divisors/OrthSubrepresentation.lean}{GlobalTheory.Shin.InPar}

\begin{definition}\label{gal-shintani-twisted}
In Situation~\ref{gal-shintani-situation}:
\begin{enumerate}
\item Let $F$ be a finite field and $t$ a ring endomorphism of $F$, acting on matrices over $F$ entry by
  entry. For $a+b=n$ let $P_{(a,b)}(F)\subset\operatorname{GL}_n(F)$ be the group of invertible block upper
  triangular matrices with diagonal blocks of sizes $a$ and $b$, and $\pi_1(p)$, $\pi_2(p)$ the diagonal
  blocks of $p\in P_{(a,b)}(F)$; thus $P_{(a,b)}(k')=P'_{(a,b)}$. For functions $\alpha'$ on
  $\operatorname{GL}_a(F)$ and $\beta'$ on $\operatorname{GL}_b(F)$, the \emph{$t$-twisted Harish-Chandra
  induction} is the function on $\operatorname{GL}_n(F)$
  \[
   R^t_{(a,b)}(\alpha',\beta')(x)=\frac1{|P_{(a,b)}(F)|}\sum_{y\in\operatorname{GL}_n(F),\ y^{-1}xt(y)\in
   P_{(a,b)}(F)}\alpha'\bigl(\pi_1(y^{-1}xt(y))\bigr)\,\beta'\bigl(\pi_2(y^{-1}xt(y))\bigr).
  \]
  For $t=1$ it is the Harish-Chandra induction $R_{(a,b)}$ of Definition~\ref{gal-support-level-zero}(2)
  over $F$. It is used for $F=k'$ and $t\in\operatorname{Gal}(k'/k)$.
\item $\hat G=G'\rtimes\operatorname{Gal}(k'/k)$ is the semidirect product in which $tgt^{-1}=t(g)$ for
  $g\in G'$ and $t\in\operatorname{Gal}(k'/k)$; its elements are written $gt$.
\item Let $F$ be a finite field and $S$ a function from polynomials over $F$ to $\mathbf C$. For an
  automorphism $h$ of a finite-dimensional $F$-vector space $V$, with characteristic polynomial $\chi_h$,
  define $\psi_S(h)$ by induction on $\dim V$:
  \[
   \psi_S(h)=S(\chi_h)-\sum_W\psi_S(h|_W),
  \]
  where $W$ runs over the $h$-stable subspaces $W\neq V$. For $x\in\operatorname{GL}_m(F)$, $\psi_S(x)$ is
  $\psi_S$ of $x$ acting on $F^m$. If $\bar F$ is an algebraic closure of $F$,
  $\Theta:\bar F^\times\to\mathbf C^\times$ an injective homomorphism and $S(P)=\sum_\alpha\Theta(\alpha)^c$,
  with $\alpha$ running over the roots of $P$ in $\bar F$ with multiplicity, then $\psi_S$ is the function
  $\psi_c$ of Definition~\ref{gal-support-level-zero}(3).
\end{enumerate}
\end{definition}
\leanref{Development/AbelianVarietyTheory/Definitions/Divisors/OrthSubrepresentation.lean}{GlobalTheory.Shin.hcTw}
\leanref{Development/AbelianVarietyTheory/Definitions/Divisors/OrthSubrepresentation.lean}{GlobalTheory.Shin.GHat}
\leanref{Development/AbelianVarietyTheory/Definitions/CommutativeAlgebra/InvariantSubspaceMobiusInverseRecEqOfFinrankLt.lean}{GlobalTheory.CuspGL.psi}
\leanref{Development/AbelianVarietyTheory/Definitions/Divisors/OrthSubrepresentation.lean}{GlobalTheory.Shin.psiG}
\leanref{Development/AbelianVarietyTheory/Definitions/Divisors/OrthSubrepresentation.lean}{GlobalTheory.Shin.chG}

\begin{lemma}\label{gal-shintani-norm}
In Situation~\ref{gal-shintani-situation}, let $\delta,\varepsilon,y\in G'$ and $\gamma\in G$.
\begin{enumerate}
\item $N_s(y^{-1}\delta s(y))=y^{-1}N_s(\delta)y$ and $\delta\,s(N_s(\delta))=N_s(\delta)\,\delta$. The
  elements of $G'$ fixed by $s$ are those of $G$. If $s^i(\delta)$ and $s^j(\varepsilon)$ commute for all
  $i,j\geq0$, then $N_s(\delta\varepsilon)=N_s(\delta)N_s(\varepsilon)$.
\item Let $B\subset M_n(k')$ be a $k'$-subalgebra with $s(B)\subset B$. If $\varepsilon\in G'\cap B$ and
  $N_s(\varepsilon)=1$, then $\varepsilon=x^{-1}s(x)$ for some $x\in G'\cap B$. The number of
  $\varepsilon\in G'\cap B$ with $N_s(\varepsilon)=1$ is $|G'\cap B|/|G\cap B|$.
\item There is $\delta_0\in G'\cap k'[\gamma]$ with $N_s(\delta_0)=\gamma$.
\item If $N_s(\delta)$ is conjugate to $\gamma$ in $G'$, there is $y_1\in G'$ such that
  $\delta_1=y_1^{-1}\delta s(y_1)$ lies in $k'[\gamma]$ and $N_s(\delta_1)=\gamma$.
\item $N_s(\delta)$ is conjugate in $G'$ to an element of $G$. The number of $\delta'\in G'$ such that
  $N_s(\delta')$ is conjugate to $\gamma$ in $G'$ is $|G'|/|Z_G(\gamma)|$.
\end{enumerate}
\end{lemma}
\leanref{Development/AbelianVarietyTheory/Proofs/DivisorProofs/GaloisTwistedNormGLBasicPropertiesStatement.lean}{GlobalTheory.T.galoisTwistedNormGL_basic_properties}
\leanref{Development/AbelianVarietyTheory/Proofs/DivisorProofs/ExistsCoboundaryOfNormOneInSubalgebraStatement.lean}{GlobalTheory.T.exists_coboundary_of_norm_one_in_subalgebra}
\leanref{Development/AbelianVarietyTheory/Proofs/DivisorProofs/CardTwistedNormOneMulCardFixedUnitsStatement.lean}{GlobalTheory.T.card_twistedNorm_one_mul_card_fixed_units}
\leanref{Development/AbelianVarietyTheory/Proofs/DivisorProofs/NormMapSurjectiveFiniteField.lean}{GlobalTheory.T.norm_map_surjective_finite_field}
\leanref{Development/AbelianVarietyTheory/Proofs/DivisorProofs/ConjugateIntoAdjoinedAlgebra.lean}{GlobalTheory.T.conjugate_into_adjoined_algebra}
\leanref{Development/AbelianVarietyTheory/Proofs/DivisorProofs/EveryElementIsNorm.lean}{GlobalTheory.T.every_element_is_norm}
\leanref{Development/AbelianVarietyTheory/Proofs/DivisorProofs/CardIsGaloisTwistedNormGLMulCardGLStatement.lean}{GlobalTheory.T.card_isGaloisTwistedNormGL_mul_card_GL}
\leanref{Development/AutomorphicArithmetic/Foundations/GaloisRepresentations/RegularSemisimpleFibres/ResEmbCompat.lean}{UpperTheory.Gal.InertialTypesAlt.NormEx}

\begin{proof}
For $j\geq0$ write $N_j(\delta)=\delta\,s(\delta)\cdots s^{j-1}(\delta)$, so that $N_s=N_l$ and $s^l=1$.

(1) By induction on $j$, $N_j(y^{-1}\delta s(y))=y^{-1}N_j(\delta)s^j(y)$ and
$\delta\,s(N_j(\delta))=N_j(\delta)s^j(\delta)$; take $j=l$. A matrix fixed by $s$ has its entries in the
fixed field $k$ of $s$, and it is invertible over $k$ if it lies in $G'$. Under the commutation hypothesis
$N_j(\varepsilon)$ commutes with $s^j(\delta)$, so by induction on $j$,
$N_{j+1}(\delta\varepsilon)=N_j(\delta)N_j(\varepsilon)s^j(\delta)s^j(\varepsilon)=N_{j+1}(\delta)N_{j+1}(\varepsilon)$.

(2) Since $G'$ is finite, $\varepsilon^{-1}$ is a power of $\varepsilon$, so it lies in $B$. The map
$T(v)=s(v)\varepsilon^{-1}$ of $B$ to itself is additive with $T(cv)=s(c)T(v)$ for $c\in k'$, and
$T^j(v)=s^j(v)N_j(\varepsilon)^{-1}$ by induction on $j$, so $T^l=1$. Let $M=\{v\in B:T(v)=v\}$ and
$B_0=\{v\in B:s(v)=v\}$. Then $B_0$ is a ring and $M$ is a left $B_0$-module. Let $b_1,\ldots,b_l$ be a basis
of $k'$ over $k$. By Galois descent for the semilinear maps $T$ and $s$ \cite[Chapter~X, \S1, Proposition~3]{Ser79}, the
maps $M^l\to B$ and $B_0^l\to B$ sending $(v_i)$ to $\sum_ib_iv_i$ are bijective, and they are
$B_0$-linear. So $M^l\cong B_0^l$ as $B_0$-modules, and $M\cong B_0$ by the Krull--Schmidt theorem
\cite{Lam01}, since both are finite. Let $x\in M$ correspond to $1$, so that $M=B_0x$. Writing
$1=\sum_ib_iv_i$ with $v_i\in M$, and $v_i=r_ix$ with $r_i\in B_0$, gives $1=cx$ with
$c=\sum_ib_ir_i\in B$. So $x\in G'\cap B$, and $s(x)\varepsilon^{-1}=x$ gives $\varepsilon=x^{-1}s(x)$.
For the count, $x\mapsto x^{-1}s(x)$ maps $G'\cap B$ into the set of $\varepsilon\in G'\cap B$ with
$N_s(\varepsilon)=1$, by (1) with $\delta=1$, and onto it by the first assertion. By (1),
$x^{-1}s(x)=y^{-1}s(y)$ if and only if $yx^{-1}$ is fixed by $s$, that is, $y\in(G\cap B)x$. So every
fibre has $|G\cap B|$ elements.

(3) The algebra $B=k'[\gamma]$ is commutative, and $s(B)\subset B$ since $s(\gamma)=\gamma$. So $U=G'\cap B$
is an abelian group with $s(U)\subset U$, and $N_s$ is a homomorphism from $U$ to $U$ by (1). For
$\delta\in U$, (1) and commutativity give $s(N_s(\delta))=N_s(\delta)$, so $N_s(U)\subset G\cap B$. By (2) the
kernel of $N_s$ on $U$ has $|U|/|G\cap B|$ elements, so $N_s(U)=G\cap B$, which contains $\gamma$.

(4) Let $c\in G'$ with $N_s(\delta)=c\gamma c^{-1}$ and put $\delta'=c^{-1}\delta s(c)$. By (1),
$N_s(\delta')=\gamma$, and $\delta'$ commutes with $\gamma$ since $s(\gamma)=\gamma$. Let $\delta_0$ be as in
(3), and let $C\subset M_n(k')$ be the centralizer of $\gamma$, a $k'$-subalgebra with $s(C)\subset C$ that
commutes with $k'[\gamma]$. Then $\varepsilon=\delta_0^{-1}\delta'$ lies in $C$, and since
$s^i(\delta_0)\in k'[\gamma]$ commutes with $s^j(\varepsilon)\in C$, (1) gives
$\gamma=N_s(\delta_0)N_s(\varepsilon)=\gamma N_s(\varepsilon)$, so $N_s(\varepsilon)=1$. By (2),
$\varepsilon=x^{-1}s(x)$ with $x\in G'\cap C$. Since $x$ commutes with $\delta_0$, $y_1=cx^{-1}$ gives
$y_1^{-1}\delta s(y_1)=x\delta's(x)^{-1}=x\delta_0x^{-1}=\delta_0$.

(5) Let $C$ and $\delta_0$ be as in (4). If $N_s(\delta')=\gamma$, then $\delta'\in C$ as in (4), and
$N_s(\delta_0\varepsilon)=\gamma N_s(\varepsilon)$ for $\varepsilon\in G'\cap C$. So
$\delta'\mapsto\delta_0^{-1}\delta'$ is a bijection from the set of $\delta'$ with $N_s(\delta')=\gamma$
onto the set of $\varepsilon\in G'\cap C$ with $N_s(\varepsilon)=1$, which has
$|Z_{G'}(\gamma)|/|Z_G(\gamma)|$ elements by (2), as $G'\cap C$ and $G\cap C$ are the centralizers of
$\gamma$ in $G'$ and in $G$. If $g=y\gamma y^{-1}$ with $y\in G'$, then $\delta'\mapsto y^{-1}\delta's(y)$ is a
bijection from the set of $\delta'$ with $N_s(\delta')=g$ onto the set of $\delta'$ with
$N_s(\delta')=\gamma$, by (1). The class of $\gamma$ in $G'$ has $|G'|/|Z_{G'}(\gamma)|$ elements, which gives
the count. Finally, for $\delta'\in G'$ the elements $\gamma'\in G$ to which $N_s(\delta')$ is conjugate in
$G'$ form one conjugacy class of $G$, or none (Definition~\ref{gal-support-level-zero}(5)). So the sum of
$1/|\gamma'^G|$ over these $\gamma'$, where $\gamma'^G$ is the class of $\gamma'$ in $G$, is $1$ if
$N_s(\delta')$ is conjugate to an element of $G$ and $0$ otherwise. Summed over $\delta'\in G'$, it is
$\sum_{\gamma'\in G}|G'|/(|Z_G(\gamma')|\,|\gamma'^G|)=|G'|$ by the count. So every $\delta'$ contributes
$1$.
\end{proof}

\begin{lemma}\label{gal-shintani-lift}
In Situation~\ref{gal-shintani-situation}, let $f$ and $f'$ be class functions on $G$.
\begin{enumerate}
\item If $\delta\in G'$ and $N_s(\delta)$ is conjugate in $G'$ to $\gamma\in G$, then
  $\operatorname{Sh}_s(f)(\delta)=f(\gamma)$. For $\delta,y\in G'$,
  $\operatorname{Sh}_s(f)(y^{-1}\delta s(y))=\operatorname{Sh}_s(f)(\delta)$.
\item $|G|\sum_{\delta\in G'}\operatorname{Sh}_s(f)(\delta)\overline{\operatorname{Sh}_s(f')(\delta)}
  =|G'|\sum_{\gamma\in G}f(\gamma)\overline{f'(\gamma)}$.
\end{enumerate}
\end{lemma}
\leanref{Development/AbelianVarietyTheory/Proofs/DivisorProofs/BaseChangeLiftClassFunction.lean}{GlobalTheory.T.base_change_lift_class_function}
\leanref{Development/AbelianVarietyTheory/Proofs/DivisorProofs/ShintaniLiftClassFunctionIsometryStatement.lean}{GlobalTheory.T.shintaniLiftClassFunction_isometry}

\begin{proof}
(1) The first assertion is the definition, since the elements of $G$ to which $N_s(\delta)$ is conjugate in
$G'$ form one conjugacy class of $G$ (Definition~\ref{gal-support-level-zero}(5)). The second follows,
since $N_s(y^{-1}\delta s(y))=y^{-1}N_s(\delta)y$ by Lemma~\ref{gal-shintani-norm}(1).

(2) Fix $\delta\in G'$. By Lemma~\ref{gal-shintani-norm}(5), the $\gamma\in G$ to which $N_s(\delta)$ is
conjugate form a conjugacy class $C$ of $G$, and by (1),
$f(\gamma)\overline{f'(\gamma)}=\operatorname{Sh}_s(f)(\delta)\overline{\operatorname{Sh}_s(f')(\delta)}$ for
$\gamma\in C$. So $\operatorname{Sh}_s(f)(\delta)\overline{\operatorname{Sh}_s(f')(\delta)}$ is the sum of
$f(\gamma)\overline{f'(\gamma)}/|C|$ over $\gamma\in C$. Sum over $\delta$ and exchange the sums: by
Lemma~\ref{gal-shintani-norm}(5), each $\gamma\in G$ occurs for $|G'|/|Z_G(\gamma)|$ elements $\delta$, with
weight $1/|C|=|Z_G(\gamma)|/|G|$.
\end{proof}

\begin{lemma}\label{gal-shintani-parabolic}
In Situation~\ref{gal-shintani-situation}, let $a+b=n$, $P'=P'_{(a,b)}$ and $P=P_{(a,b)}$.
\begin{enumerate}
\item Let $F$ be a field and $t$ a ring endomorphism of $F$, acting on matrices over $F$ entry by entry, and
  let $p,x\in\operatorname{GL}_n(F)$ be block upper triangular with diagonal blocks of sizes $a$ and $b$, which
  are written $\pi_1$ and $\pi_2$. Then $\pi_i(p^{-1}xt(p))=\pi_i(p)^{-1}\pi_i(x)t(\pi_i(p))$ for $i=1,2$. If
  $\delta\in P'$, $\gamma\in P$ and $N_s(\delta)=\gamma$, then $N_s(\pi_i(\delta))=\pi_i(\gamma)$ for $i=1,2$.
\item Let $F$ be a finite field and $t$ a ring endomorphism of $F$. For functions $\alpha'$ on
  $\operatorname{GL}_a(F)$ and $\beta'$ on $\operatorname{GL}_b(F)$, and $x,y\in\operatorname{GL}_n(F)$,
  $R^t_{(a,b)}(\alpha',\beta')(y^{-1}xt(y))=R^t_{(a,b)}(\alpha',\beta')(x)$
  (Definition~\ref{gal-shintani-twisted}(1)), and $R_{(a,b)}(\alpha',\beta')$ is a class function on
  $\operatorname{GL}_n(F)$. For every function $S$ from polynomials over $F$ to $\mathbf C$, the function
  $\psi_S$ (Definition~\ref{gal-shintani-twisted}(3)) and the function $x\mapsto S(\chi_x)$, where $\chi_x$ is
  the characteristic polynomial of $x$, are class functions on $\operatorname{GL}_n(F)$.
\item Let $\gamma\in G$ and $\delta\in G'\cap k'[\gamma]$. If $y_0\in G$ and $y_0^{-1}\gamma y_0\in P$, then
  $y_0^{-1}\delta y_0\in P'$. If moreover $N_s(\delta)=\gamma$, then for $y\in G'$ the element
  $y^{-1}\delta s(y)$ lies in $P'$ if and only if $y=y_0p$ with $y_0\in G$, $y_0^{-1}\gamma y_0\in P$ and
  $p\in P'$.
\end{enumerate}
\end{lemma}
\leanref{Development/AbelianVarietyTheory/Proofs/DivisorProofs/DiagonalBlocksTwistedConjBlockUpperTriangularStatement.lean}{GlobalTheory.T.diagonalBlocks_twistedConj_blockUpperTriangular}
\leanref{Development/AbelianVarietyTheory/Proofs/DivisorProofs/GaloisTwistedNormGLDiagonalBlocksStatement.lean}{GlobalTheory.T.galoisTwistedNormGL_diagonal_blocks}
\leanref{Development/AbelianVarietyTheory/Proofs/DivisorProofs/TwistedHarishChandraInductionGLIsTwistedClassFunStatement.lean}{GlobalTheory.T.twistedHarishChandraInductionGL_isTwistedClassFunction}
\leanref{Development/AbelianVarietyTheory/Proofs/DivisorProofs/AdjoinConjInParStatement.lean}{GlobalTheory.T.inPar_conj_of_mem_adjoin_incl_of_inPar_conj}
\leanref{Development/AbelianVarietyTheory/Proofs/DivisorProofs/BlockUpperTwistedConjIffDefinedOverBaseStatement.lean}{GlobalTheory.T.blockUpper_twistedConj_iff_defined_over_base}

\begin{proof}
(1) The invertible block upper triangular matrices over $F$ form a group $H$ that contains $t(p)$, and
$p\mapsto(\pi_1(p),\pi_2(p))$ is a homomorphism on $H$, which commutes with $t$ since $t$ is applied entry by
entry. Applying it to $p^{-1}xt(p)\in H$ gives the first assertion. Over $F=k'$, the product
$N_s(\delta)$ of the $s^j(\delta)\in P'$ therefore has $\pi_i(N_s(\delta))=N_s(\pi_i(\delta))$, and
$N_s(\pi_i(\delta))=\pi_i(\gamma)$.

(2) The substitution $w=yz$ turns the sum defining $R^t_{(a,b)}(\alpha',\beta')(y^{-1}xt(y))$, over
$z\in\operatorname{GL}_n(F)$, into the sum defining $R^t_{(a,b)}(\alpha',\beta')(x)$, since
$z^{-1}(y^{-1}xt(y))t(z)=w^{-1}xt(w)$. The case $t=1$ gives the second assertion. For the third, conjugation
by $y$ is an isomorphism of $F^n$ with $x$ onto $F^n$ with $y^{-1}xy$. It matches their stable subspaces and
preserves characteristic polynomials, so $\psi_S(y^{-1}xy)=\psi_S(x)$ by induction on the dimension, and
$\chi_{y^{-1}xy}=\chi_x$.

(3) Let $B\subset M_n(k')$ be the algebra of block upper triangular matrices with diagonal blocks of sizes $a$
and $b$. Then $s(B)\subset B$ and $P'=G'\cap B$. For $y\in G'$, conjugation by $y$ maps $k'[\gamma]$ onto
$k'[y^{-1}\gamma y]$, so $y^{-1}\delta y\in B$ whenever $y^{-1}\gamma y\in B$. This gives the first assertion.
If $y=y_0p$ as stated, then $y^{-1}\delta s(y)=p^{-1}(y_0^{-1}\delta y_0)s(p)$, since $s(y_0)=y_0$, and this
lies in $P'$. Conversely, suppose that $y^{-1}\delta s(y)\in P'$. By Lemma~\ref{gal-shintani-norm}(1), its norm
is $y^{-1}\gamma y$, which lies in $P'$ since $P'$ is a group with $s(P')=P'$. So $y^{-1}\delta y\in P'$, and
$\varepsilon=y^{-1}s(y)=(y^{-1}\delta y)^{-1}(y^{-1}\delta s(y))$ lies in $P'$. By
Lemma~\ref{gal-shintani-norm}(1) with $\delta=1$, $N_s(\varepsilon)=1$, so $\varepsilon=x^{-1}s(x)$ with
$x\in P'$ by Lemma~\ref{gal-shintani-norm}(2) for $B$. Then $s(yx^{-1})=y\varepsilon s(x)^{-1}=yx^{-1}$, so
$y_0=yx^{-1}$ lies in $G$ by Lemma~\ref{gal-shintani-norm}(1), $y=y_0x$, and
$y_0^{-1}\gamma y_0=x(y^{-1}\gamma y)x^{-1}$ lies in $P'\cap G=P$.
\end{proof}

\begin{lemma}\label{gal-shintani-hc}
In Situation~\ref{gal-shintani-situation}, let $a+b=n$ and let $\alpha$ and $\beta$ be class functions on
$\operatorname{GL}_a(k)$ and $\operatorname{GL}_b(k)$. Then on $G'$
\[
 \operatorname{Sh}_s\bigl(R_{(a,b)}(\alpha,\beta)\bigr)
 =R^s_{(a,b)}\bigl(\operatorname{Sh}_s(\alpha),\operatorname{Sh}_s(\beta)\bigr)
\]
(Definitions~\ref{gal-support-level-zero}(2) and~\ref{gal-shintani-twisted}(1)).
\end{lemma}
\leanref{Development/AutomorphicArithmetic/Foundations/GaloisRepresentations/RegularSemisimpleFibres/ResEmbCompat.lean}{UpperTheory.Gal.InertialTypesAlt.LiftHc}
\leanref{Development/AbelianVarietyTheory/Proofs/DivisorProofs/ShintaniLiftHarishChandraInductionEqTwistedStatement.lean}{GlobalTheory.T.shintaniLift_harishChandraInduction_eq_twisted}
\leanref{Development/AbelianVarietyTheory/Proofs/DivisorProofs/ShintaniLiftHarishChandraInductionEqTwisted.lean}{GlobalTheory.Shin.LiftHc.pointwise}
\leanref{Development/AbelianVarietyTheory/Proofs/DivisorProofs/ShintaniLiftHarishChandraInductionEqTwisted.lean}{GlobalTheory.Shin.LiftHc.term}
\leanref{Development/AbelianVarietyTheory/Proofs/DivisorProofs/CardComapMulFinsumEqFinsumFibresStatement.lean}{GlobalTheory.T.card_comap_mul_finsum_eq_finsum_fibres}

\begin{proof}
Write $P'=P'_{(a,b)}$, $P=P_{(a,b)}$, $\alpha'=\operatorname{Sh}_s(\alpha)$ and $\beta'=\operatorname{Sh}_s(\beta)$.
Let $\delta\in G'$. By Lemma~\ref{gal-shintani-norm}(5) and~(4) there are $\gamma\in G$ and $y_1\in G'$ such
that $\delta_1=y_1^{-1}\delta s(y_1)$ lies in $k'[\gamma]$ and $N_s(\delta_1)=\gamma$. Each side of the
identity takes the same value at $\delta$ and at $\delta_1$: the left side by Lemma~\ref{gal-shintani-lift}(1),
since $R_{(a,b)}(\alpha,\beta)$ is a class function by Lemma~\ref{gal-shintani-parabolic}(2), and the right side
by Lemma~\ref{gal-shintani-parabolic}(2). At $\delta_1$ the left side is $R_{(a,b)}(\alpha,\beta)(\gamma)$, by
Lemma~\ref{gal-shintani-lift}(1). So we may assume that $\delta\in k'[\gamma]$ and $N_s(\delta)=\gamma$, and
we must show that $R^s_{(a,b)}(\alpha',\beta')(\delta)=R_{(a,b)}(\alpha,\beta)(\gamma)$.

For $y\in G'$ let $T(y)=\alpha'(\pi_1(y^{-1}\delta s(y)))\,\beta'(\pi_2(y^{-1}\delta s(y)))$ if
$y^{-1}\delta s(y)\in P'$, and $T(y)=0$ otherwise. For $y_0\in G$ let
$c(y_0)=\alpha(\pi_1(y_0^{-1}\gamma y_0))\,\beta(\pi_2(y_0^{-1}\gamma y_0))$ if $y_0^{-1}\gamma y_0\in P$, and
$c(y_0)=0$ otherwise. Then $R^s_{(a,b)}(\alpha',\beta')(\delta)=|P'|^{-1}\sum_{y\in G'}T(y)$ and
$R_{(a,b)}(\alpha,\beta)(\gamma)=|P|^{-1}\sum_{y_0\in G}c(y_0)$.

\emph{Support.} If $T(y)\neq0$, then $y=y_0p$ with $y_0\in G$ and $p\in P'$, by
Lemma~\ref{gal-shintani-parabolic}(3).

\emph{Counting.} If $y=y_0p$ with $y_0\in G$ and $p\in P'$, the pairs $(y_0',p')\in G\times P'$ with
$y_0'p'=y$ are the $(y_0u,u^{-1}p)$ with $u\in G\cap P'=P$. So
$|P|\sum_{y\in G'}T(y)=\sum_{y_0\in G}\sum_{p\in P'}T(y_0p)$.

\emph{Each term.} We show that $T(y_0p)=c(y_0)$ for $y_0\in G$ and $p\in P'$. Put
$\delta_2=y_0^{-1}\delta y_0$. Since $s(y_0)=y_0$, $(y_0p)^{-1}\delta s(y_0p)=p^{-1}\delta_2s(p)$. Suppose
first that $y_0^{-1}\gamma y_0\in P$. By Lemma~\ref{gal-shintani-parabolic}(3), $\delta_2\in P'$, so
$p^{-1}\delta_2s(p)\in P'$ and $\pi_i(p^{-1}\delta_2s(p))=\pi_i(p)^{-1}\pi_i(\delta_2)s(\pi_i(p))$ by
Lemma~\ref{gal-shintani-parabolic}(1). By Lemma~\ref{gal-shintani-norm}(1), $N_s(\delta_2)=y_0^{-1}\gamma y_0$,
so $N_s(\pi_i(\delta_2))=\pi_i(y_0^{-1}\gamma y_0)$ by Lemma~\ref{gal-shintani-parabolic}(1).
Lemma~\ref{gal-shintani-lift}(1), on $\operatorname{GL}_a(k')$ and $\operatorname{GL}_b(k')$, gives
$\alpha'(\pi_1(p^{-1}\delta_2s(p)))=\alpha'(\pi_1(\delta_2))=\alpha(\pi_1(y_0^{-1}\gamma y_0))$, and the same
for $\beta'$; so $T(y_0p)=c(y_0)$. Suppose now that $y_0^{-1}\gamma y_0\notin P$ and that $T(y_0p)\neq0$. By
Lemma~\ref{gal-shintani-parabolic}(3), $y_0p=y_3p_3$ with $y_3\in G$, $y_3^{-1}\gamma y_3\in P$ and
$p_3\in P'$. Then $y_3^{-1}y_0=p_3p^{-1}$ lies in $P'\cap G=P$, so
$y_0^{-1}\gamma y_0=(y_3^{-1}y_0)^{-1}(y_3^{-1}\gamma y_3)(y_3^{-1}y_0)$ lies in $P$, a contradiction. So
$T(y_0p)=0=c(y_0)$.

Hence $|P|\sum_yT(y)=|P'|\sum_{y_0}c(y_0)$, which is the claim.
\end{proof}

\begin{lemma}\label{gal-shintani-clifford}
Let $H$ be a finite group, $\Gamma$ a group of prime order $l$ acting on $H$ by automorphisms, $s\in\Gamma$ with
$s\neq1$, and $\hat H=H\rtimes\Gamma$, whose elements are written $ht$, with $tht^{-1}=t(h)$. Let $X$ be a
virtual character of $\hat H$ with $\sum_{h\in H}|X(hs)|^2=|H|$, and let $Y$ be a character of $H$ with
$Y(s(h))=Y(h)$ for $h\in H$. Suppose that $\sum_{h\in H}X(h)\overline{Y(h)}=z|H|$ for an integer $z$ not
divisible by $l$. Then the function $h\mapsto X(hs)$ on $H$ is carried by $Y$
(Definition~\ref{gal-support-level-zero}(1)).
\end{lemma}
\leanref{Development/AbelianVarietyTheory/Proofs/DivisorProofs/TwistedVirtualCharacterNormOneSupportStatement.lean}{GlobalTheory.T.twisted_virtualCharacter_norm_one_support}
\leanref{Development/AbelianVarietyTheory/Proofs/DivisorProofs/FourierSumsCosetsSemidirectCyclicStatement.lean}{GlobalTheory.T.fourier_sums_cosets_semidirect_cyclic}
\leanref{Development/AbelianVarietyTheory/Proofs/DivisorProofs/VirtualCharacterCosetNormOneSingleStatement.lean}{GlobalTheory.T.virtualCharacter_coset_normOne_single}
\leanref{Development/AbelianVarietyTheory/Proofs/DivisorProofs/CliffordSemidirectPrimeIndexIrreducibleCharacterStatement.lean}{GlobalTheory.T.clifford_semidirect_primeIndex_irreducibleCharacter}
\leanref{Development/AbelianVarietyTheory/Proofs/DivisorProofs/SemidirectCosetPairingEqZeroOrTwistStatement.lean}{GlobalTheory.T.semidirect_coset_pairing_eq_zero_or_twist}
\leanref{Development/AbelianVarietyTheory/Proofs/DivisorProofs/CyclotomicSumNormSqEqOneSingleNonzeroStatement.lean}{GlobalTheory.T.cyclotomic_sum_normSq_eq_one_single_nonzero}
\leanref{Development/AbelianVarietyTheory/Proofs/DivisorProofs/ExistsSemidirectCharacterExtOfStableStatement.lean}{GlobalTheory.T.exists_semidirect_character_ext_of_stable}
\leanref{Development/AbelianVarietyTheory/Proofs/DivisorProofs/FinsumSemidirectCosetPairingNeZeroStatement.lean}{GlobalTheory.T.finsum_semidirect_coset_pairing_ne_zero}
\leanref{Development/AbelianVarietyTheory/Proofs/DivisorProofs/FinsumSemidirectCosetCharacterOrthZeroStatement.lean}{GlobalTheory.T.finsum_semidirect_coset_character_orth_zero}
\leanref{Development/AbelianVarietyTheory/Proofs/DivisorProofs/TwistedRestrictedCharacterInIsotypicPartStatement.lean}{GlobalTheory.T.twisted_restricted_character_in_isotypic_part}

\begin{proof}
Let $\zeta=e^{2\pi i/l}$, and let $\eta$ be the linear character of $\hat H$ that is trivial on $H$ with
$\eta(s)=\zeta$. For functions $\varphi,\varphi'$ on $\hat H$ and $t\in\Gamma$ put
$\langle\varphi,\varphi'\rangle_t=\sum_{h\in H}\varphi(ht)\overline{\varphi'(ht)}$ and
$[\varphi,\varphi']=|\hat H|^{-1}\sum_{x\in\hat H}\varphi(x)\overline{\varphi'(x)}$. Since $s$ generates
$\Gamma$, Fourier analysis on $\Gamma$ gives
\[
 l\langle\varphi,\varphi'\rangle_1=|\hat H|\sum_{a=0}^{l-1}[\varphi,\varphi'\eta^a],\qquad
 l\langle\varphi,\varphi'\rangle_s=|\hat H|\sum_{a=0}^{l-1}\zeta^a[\varphi,\varphi'\eta^a].
\]
For virtual characters $\varphi,\varphi'$ the numbers $[\varphi,\varphi'\eta^a]$ are integers, and for
characters they are nonnegative integers.

\emph{Step 1.} There are $e\in\mathbf C$ and an irreducible character $\chi$ of $\hat H$ whose restriction to
$H$ is irreducible, with $X(hs)=e\chi(hs)$ for all $h\in H$. By Clifford theory for the normal subgroup $H$ of
prime index \cite[Chapter~6]{Isa76}, an irreducible character $\chi$ of $\hat H$ either is induced from $H$, and
then vanishes on $Hs$, or restricts irreducibly to $H$; in the second case the characters $\chi\eta^a$,
$0\leq a<l$, are distinct, and the second formula gives $\langle\chi,\chi\rangle_s=|H|$. If $\chi$ and $\chi'$
are irreducible and $\langle\chi,\chi'\rangle_s\neq0$, the second formula gives $[\chi,\chi'\eta^a]\neq0$ for
some $a$, so $\chi=\chi'\eta^a$ and $\chi(hs)=\zeta^a\chi'(hs)$. Write $X$ as an integral combination of
irreducible characters, and group those that do not vanish on $Hs$ into classes under $\chi\sim\chi\eta^a$.
This gives $X(hs)=\sum_iC_i\chi_i(hs)$ with $C_i\in\mathbf Z[\zeta]$, $\langle\chi_i,\chi_i\rangle_s=|H|$ and
$\langle\chi_i,\chi_j\rangle_s=0$ for $i\neq j$. The hypothesis on $X$ gives $\sum_i|C_i|^2=1$, and a
computation with the coefficients of the $C_i$ in the powers of $\zeta$ shows that only one $C_i$ is nonzero.

\emph{Step 2.} $Y$ is the restriction to $H$ of a character $Z$ of $\hat H$. Let
$\rho:H\to\operatorname{GL}_d(\mathbf C)$ afford $Y$. Since $Y(s(h))=Y(h)$, the representations $h\mapsto\rho(s(h))$
and $\rho$ are equivalent: $\rho(s(h))=A\rho(h)A^{-1}$ for some $A$. Then $A^l$ commutes with $\rho(H)$, and
there is a polynomial $p$ with $p(A^l)^lA^l=1$. Since $p(A^l)$ commutes with $A$ and with $\rho(H)$, the matrix
$A'=Ap(A^l)$ satisfies $A'^l=1$ and $\rho(s(h))=A'\rho(h)A'^{-1}$. So $\rho$ extends to $\hat H$ by
$s\mapsto A'$.

\emph{Step 3.} $\langle X,Z\rangle_s\neq0$. By the first formula, $lz|H|=l\langle X,Z\rangle_1=|\hat H|\sum_am_a$
with integers $m_a=[X,Z\eta^a]$. If $\langle X,Z\rangle_s=0$, the second formula gives $\sum_am_a\zeta^a=0$, so
all $m_a$ are equal, since the only linear relation over $\mathbf Q$ among $1,\zeta,\ldots,\zeta^{l-1}$ is that
their sum is $0$. Then $lz|H|=l^2m_0|H|$, and $l$ divides $z$, a contradiction.

\emph{Step 4.} $\sum_{h\in H}\chi(h)\overline{Y(h)}\neq0$, that is, $\langle\chi,Z\rangle_1\neq0$. By Steps 1
and 3, $e\neq0$ and $\langle\chi,Z\rangle_s\neq0$. If $\langle\chi,Z\rangle_1=0$, the first formula shows that
the nonnegative integers $[\chi,Z\eta^a]$ have sum $0$, so they all vanish, and the second formula gives
$\langle\chi,Z\rangle_s=0$.

\emph{Step 5.} Let $\hat\rho$ afford $\chi$, let $\lambda$ be its restriction to $H$, which is irreducible, and
let $\sigma$ afford $Y$. Let $\rho$ be a representation of $H$ that does not meet $Y$. By Step 4 and
Lemma~\ref{gal-support-finite}(5), there is a nonzero $H$-map $\sigma\to\lambda$, which is onto since $\lambda$ is
irreducible; and there is no nonzero $H$-map $\sigma\to\rho$. So there is no nonzero $H$-map $\lambda\to\rho$.
Since $\overline{\chi(hs)}=\chi(s^{-1}h^{-1})=\operatorname{tr}(\hat\rho(s^{-1})\lambda(h^{-1}))$, Schur's
orthogonality relations for matrix coefficients give $\sum_h\overline{\chi(hs)}\rho(h)=0$. So
$h\mapsto\chi(hs)$ is carried by $Y$, and so is $h\mapsto X(hs)=e\chi(hs)$.
\end{proof}

\begin{lemma}\label{gal-shintani-virtual}
In Situation~\ref{gal-shintani-situation}, suppose that $l$ is prime. Let $q=|k|$, let $\bar k$ be an algebraic
closure of $k'$, let $\Theta:\bar k^\times\to\mathbf C^\times$ be an injective homomorphism and let $e\geq0$. For a
polynomial $P$ over $k'$ let $S(P)=\prod_{j=0}^{l-1}\sum_\alpha\Theta(\alpha)^{eq^j}$, where $\alpha$ runs over
the roots of $P$ in $\bar k$ with multiplicity. Then there is a virtual character $X$ of $\hat G$
(Definition~\ref{gal-shintani-twisted}(2)) with $X(g)=\psi_S(g)$ and $X(gs)=\operatorname{Sh}_s(\psi_e)(g)$ for
every $g\in G'$. Here $\psi_S$ is taken on $G'$ (Definition~\ref{gal-shintani-twisted}(3)) and $\psi_e$ on $G$
(Definition~\ref{gal-support-level-zero}(3)).
\end{lemma}
\leanref{Development/AbelianVarietyTheory/Proofs/DivisorProofs/ExistsVirtualCharacterShintaniPsiLiftStatement.lean}{GlobalTheory.T.exists_virtualCharacter_shintani_psi_lift}
\leanref{Development/AbelianVarietyTheory/Proofs/DivisorProofs/ProdRootPowerSumMapGaloisEqStatement.lean}{GlobalTheory.T.prod_rootPowerSum_map_galois_eq}
\leanref{Development/AbelianVarietyTheory/Proofs/CharacterTheory/RootPowerSumCharpolyMemVirtualCharactersGLStatement.lean}{GlobalTheory.T.rootPowerSum_charpoly_mem_virtualCharacters_GL}
\leanref{Development/AbelianVarietyTheory/Proofs/DivisorProofs/ExistsTensorInductionVirtualCharacterPrimeStatement.lean}{GlobalTheory.T.exists_tensorInduction_virtualCharacter_prime}
\leanref{Development/AbelianVarietyTheory/Proofs/DivisorProofs/ExistsTensorInductionCharStatement.lean}{GlobalTheory.T.exists_tensorInduction_char'}
\leanref{Development/AbelianVarietyTheory/Proofs/DivisorProofs/ExistsVirtualCharacterSemidirectOrbitCrossTermsStatement.lean}{GlobalTheory.T.exists_virtualCharacter_semidirect_orbit_cross_terms}
\leanref{Development/AbelianVarietyTheory/Proofs/DivisorProofs/ExistsVirtualCharacterTwistedHarishChandraInductStatement.lean}{GlobalTheory.T.exists_virtualCharacter_twistedHarishChandraInductionGL}
\leanref{Development/AbelianVarietyTheory/Proofs/DivisorProofs/PrimitiveFunctionCharpolyGLHarishChandraRecursioStatement.lean}{GlobalTheory.T.primitiveFunctionCharpolyGL_harishChandra_recursion}
\leanref{Development/AbelianVarietyTheory/Proofs/DivisorProofs/ShintaniLiftCharpolyRootPowerSumEqNormStatement.lean}{GlobalTheory.T.shintaniLift_charpoly_rootPowerSum_eq_norm}

\begin{proof}
Each $t\in\operatorname{Gal}(k'/k)$ is $z\mapsto z^{q^j}$ for some $0\leq j<l$, and the roots of the polynomial
$t(P)$ are then the $\alpha^{q^j}$. So $S(P)=\prod_t\chi_0(t(P))$, where $\chi_0(P)=\sum_\alpha\Theta(\alpha)^e$
over the roots $\alpha$ of $P$ in $\bar k$ with multiplicity. We use four facts.
\begin{enumerate}
\item[(i)] The function $x\mapsto\chi_0(\chi_x)$ on $G'$ is a virtual character. It is $x\mapsto\beta(x^e)$,
  where $\beta(x)=\sum_\alpha\Theta(\alpha)$, over the roots of $\chi_x$, is the Brauer lift of the natural
  representation of $G'$, a virtual character by Green \cite{Gre55}; and $x\mapsto\varphi(x^e)$ is a virtual
  character for every virtual character $\varphi$ \cite[\S9.1, Exercise~9.3]{Ser77}. Alpöge proves (i) by Brauer induction.
\item[(ii)] For a virtual character $\chi$ of $G'$ there is a virtual character $A$ of $\hat G$ with
  $A(g)=\prod_t\chi(t(g))$ and $A(gs)=\chi(N_s(g))$ for $g\in G'$. If $\chi$ is the character of a
  representation $\rho$, $A$ is the character of the representation of $\hat G$ tensor induced from $\rho$;
  its value at $gs$ is the trace of a cyclic permutation of $l$ tensor factors composed with a tensor product
  of the matrices $\rho(s^j(g))$, which is $\operatorname{tr}\rho(g\,s(g)\cdots s^{l-1}(g))$. If $\chi=a-b$ with
  characters $a$ and $b$, let $A_a$ and $A_b$ be the characters just described. For a nonempty proper subset
  $T\subset\operatorname{Gal}(k'/k)$ let $f_T(g)=\prod_{t\in T}\chi(t(g))\prod_{t\notin T}b(t(g))$, a virtual
  character of $G'$. Since $l$ is prime, $\operatorname{Gal}(k'/k)$ permutes these subsets freely, so the sum $Z$
  of the characters induced from $G'$ to $\hat G$ by one $f_T$ in each orbit satisfies
  $Z(g)=\sum_Tf_T(g)=\prod_ta(t(g))-\prod_t\chi(t(g))-\prod_tb(t(g))$ and $Z(gs)=0$. Then $A=A_a-A_b-Z$.
\item[(iii)] Let $a<n$ and let $X_a$ be a virtual character of $\hat G_a=\operatorname{GL}_a(k')\rtimes
  \operatorname{Gal}(k'/k)$. Let $\hat P\subset\hat G$ be the subgroup of the $pt$ with $p\in P'_{(a,n-a)}$, and
  $\hat\pi(pt)=\pi_1(p)t$, a homomorphism $\hat P\to\hat G_a$. The virtual character $Z_a$ induced from $\hat P$
  to $\hat G$ by $X_a\circ\hat\pi$ satisfies $Z_a(gt)=R^t_{(a,n-a)}(g'\mapsto X_a(g't),1)(g)$ for $g\in G'$ and
  $t\in\operatorname{Gal}(k'/k)$, with $1$ the constant function on $\operatorname{GL}_{n-a}(k')$. This is the
  formula for induced characters: conjugation by elements of $\operatorname{Gal}(k'/k)$ preserves $\hat P$ and
  $X_a\circ\hat\pi$, and $|\hat P|=l\,|P'_{(a,n-a)}|$.
\item[(iv)] For a finite field $F$, a function $S'$ on polynomials over $F$ and $x\in\operatorname{GL}_n(F)$,
  \[
   \psi_{S'}(x)=S'(\chi_x)-\sum_{a=0}^{n-1}R_{(a,n-a)}(\psi_{S'},1)(x),
  \]
  with $\psi_{S'}$ on $\operatorname{GL}_a(F)$ and $1$ the constant function on $\operatorname{GL}_{n-a}(F)$.
  Indeed, $y^{-1}xy\in P_{(a,n-a)}$ if and only if $x$ stabilizes the span $W$ of the first $a$ columns of $y$,
  and then $\pi_1(y^{-1}xy)$ is the matrix of $x|_W$ in that basis; each $W$ of dimension $a$ arises from
  $|P_{(a,n-a)}|$ matrices $y$. So $R_{(a,n-a)}(\psi_{S'},1)(x)$ is the sum of $\psi_{S'}(x|_W)$ over the
  $x$-stable subspaces $W$ of dimension $a$.
\end{enumerate}
We argue by induction on $n$. By induction, for each $a<n$ there is a virtual character $X_a$ of $\hat G_a$ with
$X_a(g)=\psi_S(g)$ and $X_a(gs)=\operatorname{Sh}_s(\psi_e)(g)$ for $g\in\operatorname{GL}_a(k')$; let $Z_a$ be
as in (iii). Let $A$ be as in (ii) for the virtual character of (i), and put $X=A-\sum_{a<n}Z_a$. For $g\in G'$
and $t\in\operatorname{Gal}(k'/k)$, $\chi_{t(g)}=t(\chi_g)$, so $A(g)=S(\chi_g)$; and
$Z_a(g)=R_{(a,n-a)}(\psi_S,1)(g)$. So $X(g)=\psi_S(g)$ by (iv) over $k'$. Next,
$A(gs)=\chi_0(\chi_{N_s(g)})=\operatorname{Sh}_s(\gamma\mapsto\chi_0(\chi_\gamma))(g)$, since $N_s(g)$ is
conjugate to an element of $G$ (Lemma~\ref{gal-shintani-norm}(5)), which has the same characteristic
polynomial. And $Z_a(gs)=R^s_{(a,n-a)}(\operatorname{Sh}_s(\psi_e),1)(g)=\operatorname{Sh}_s(R_{(a,n-a)}(\psi_e,1))(g)$
by Lemma~\ref{gal-shintani-hc}, since $\psi_e$ is a class function and $\operatorname{Sh}_s(1)=1$ by
Lemma~\ref{gal-shintani-norm}(5). As $\operatorname{Sh}_s$ is linear, (iv) over $k$ gives
$X(gs)=\operatorname{Sh}_s(\psi_e)(g)$.
\end{proof}

\begin{lemma}\label{gal-shintani-pairing}
In Situation~\ref{gal-shintani-situation} with $n=m\geq1$, suppose that $l$ is prime. Let $q=|k|$, let $\bar k$
be an algebraic closure of $k'$, let $\Theta:\bar k^\times\to\mathbf C^\times$ be an injective homomorphism, and let
$e\geq0$ be regular for $(q,m)$ (Definition~\ref{gal-support-level-zero}(3)). For a polynomial $P$ over $k'$ let
$S(P)=\prod_{j=0}^{l-1}\sum_\alpha\Theta(\alpha)^{eq^j}$, where $\alpha$ runs over the roots of $P$ in $\bar k$
with multiplicity, and take $\psi_S$ on $G'$ (Definition~\ref{gal-shintani-twisted}(3)). Take $\kappa_\Theta$ and
$J_\Theta$ over $k'$ (Definition~\ref{gal-support-level-zero}).
\begin{enumerate}
\item If $l$ divides $m$, then $\sum_{x\in G'}\psi_S(x)\overline{J_\Theta(\emptyset;L)(x)}=\pm|G'|$, where $L$
  is the list of the pairs $(m/l,eq^i)$, $0\leq i<l$.
\item If $l$ does not divide $m$, let $e'=e(1+q^m+\cdots+q^{m(l-1)})$, and let $c$ be the number of
  $(a_0,\ldots,a_{l-1})$ with $0\leq a_j<m$ and $e'\equiv\sum_jeq^jq^{la_j}\pmod{q^{lm}-1}$. Then
  $\sum_{x\in G'}\psi_S(x)\overline{\kappa_\Theta(m,e')(x)}=(-1)^{m-1}c\,|G'|$, and $c\equiv1\pmod l$.
\end{enumerate}
\end{lemma}
\leanref{Development/AbelianVarietyTheory/Proofs/DivisorProofs/CuspidalGaloisPairingEqPmOnePrimeDegreeStatement.lean}{GlobalTheory.T.cuspidal_galois_pairing_eq_pm_one_prime_degree}
\leanref{Development/AbelianVarietyTheory/Proofs/DivisorProofs/PsiGPairingEqChGPairingStatement.lean}{GlobalTheory.T.finsum_psiG_mul_star_eq_chG}
\leanref{Development/AbelianVarietyTheory/Proofs/DivisorProofs/PairingCuspidalCharacterGLRootPowerSumCharacterStatement.lean}{GlobalTheory.T.pairing_cuspidalCharacterGL_rootPowerSumCharacter}
\leanref{Development/AbelianVarietyTheory/Proofs/DivisorProofs/PsiSumEqCharSumAgainstCuspForm.lean}{GlobalTheory.T.psi_sum_eq_char_sum_against_cusp_form}
\leanref{Development/AbelianVarietyTheory/Proofs/DivisorProofs/PrimitiveFunctionCharpolyPairProdRootSumStatement.lean}{GlobalTheory.T.primitiveFunctionCharpoly_pair_prod_rootSum}
\leanref{Development/AbelianVarietyTheory/Proofs/DivisorProofs/CardExponentTuplesModEqOneOfPrimeNotDvdStatement.lean}{GlobalTheory.T.card_exponent_tuples_modEq_one_of_prime_not_dvd}

\begin{proof}
Both are computations with characters of general linear groups over $k'$, which has $q^l$ elements; we say what
they rest on.

(1) Write $m'=m/l$ and $J=J_\Theta(\emptyset;L)$, a character by Lemma~\ref{gal-support-finite}(3). Since $e$ is
regular for $(q,m)$, the sum of $J$ over the matrices that fix the first $m-1$ standard basis vectors of
$k'^m$ vanishes, and for a character with this property $\psi_S$ may be replaced by $x\mapsto S(\chi_x)$ in the
pairing. For $0\leq i,j<l$, the pairing of $\kappa_\Theta(m',eq^i)$ with the function
$x\mapsto\sum_\alpha\Theta(\alpha)^{eq^j}$ on $\operatorname{GL}_{m'}(k')$, $\alpha$ over the roots of $\chi_x$, is
$(-1)^{m'-1}|\operatorname{GL}_{m'}(k')|$ if $i=j$ and $0$ otherwise. From these pairings one computes the
pairing of the Harish-Chandra induction $J$ of the $\kappa_\Theta(m',eq^i)$ with $x\mapsto S(\chi_x)$, which is
$(-1)^{(m'-1)l}|G'|$.

(2) Since $S(1)=0$, $\psi_S$ may again be replaced by $x\mapsto S(\chi_x)$ in the pairing with
$\kappa_\Theta(m,e')$. As $\kappa_\Theta(m,e')=(-1)^{m-1}\psi_{e'}$, it remains to see that the pairing of
$x\mapsto S(\chi_x)$ with $\psi_{e'}$ is $c\,|G'|$, and that $c\equiv1\pmod l$; both are computations with
the exponents.
\end{proof}

\begin{lemma}\label{gal-shintani-cuspidal}
In Situation~\ref{gal-shintani-situation} with $n=m\geq1$, suppose that $l$ is prime. Let $q=|k|$, let $\bar k$
be an algebraic closure of $k'$ and let $\Theta:\bar k^\times\to\mathbf C^\times$ be an injective homomorphism. Let
$e\geq0$ be regular for $(q,m)$ (Definition~\ref{gal-support-level-zero}(3)), and let $A$ be the list of the
pairs $(m/l,eq^i)$, $0\leq i<l$, if $l$ divides $m$, and the list of the one pair
$(m,e(1+q^m+\cdots+q^{m(l-1)}))$ otherwise; these are the pairs that replace $(m,e)$ in a lift by $l$
(Definition~\ref{gal-support-pattern}(1)). Then on $G'$ the Shintani lift $\operatorname{Sh}_s(\kappa_\Theta(m,e))$
is carried by $J_\Theta(\emptyset;A)$, taken over $k'$ (Definition~\ref{gal-support-level-zero}).
\end{lemma}
\leanref{Development/AbelianVarietyTheory/Proofs/DivisorProofs/ShintaniLiftKillsCuspDivisibleDegree.lean}{GlobalTheory.T.shintani_lift_kills_cusp_divisible_degree}
\leanref{Development/AbelianVarietyTheory/Proofs/DivisorProofs/ShintaniLiftKillsCuspOfDegreeNotDvd.lean}{GlobalTheory.T.shintani_lift_kills_cusp_of_degree_not_dvd}
\leanref{Development/AbelianVarietyTheory/Proofs/DivisorProofs/CuspidalCharacterGLIsCharacterNormOneOfRegularStatement.lean}{GlobalTheory.T.cuspidalCharacterGL_isCharacter_norm_one_of_regular}
\leanref{Development/AbelianVarietyTheory/Proofs/DivisorProofs/HarishChandraProductShintaniLiftGaloisStatement.lean}{GlobalTheory.T.harishChandraProduct_shintaniLift_galois}
\leanref{Development/AbelianVarietyTheory/Proofs/DivisorProofs/CuspidalCharacterGLFiniteFieldFrobeniusPowStatement.lean}{GlobalTheory.T.cuspidalCharacterGLFiniteField_frobenius_pow}
\leanref{Development/AbelianVarietyTheory/Proofs/DivisorProofs/RegularPairLiftToPrimePowerField.lean}{GlobalTheory.T.regular_pair_lift_to_prime_power_field}

\begin{proof}
Let $Y=J_\Theta(\emptyset;A)$ on $G'$. Let $X$ be the virtual character of $\hat G$ given by
Lemma~\ref{gal-shintani-virtual} for $n=m$ and this $e$, and put $X'=(-1)^{m-1}X$. Since
$\kappa_\Theta(m,e)=(-1)^{m-1}\psi_e$ and $\operatorname{Sh}_s$ is linear,
$X'(gs)=\operatorname{Sh}_s(\kappa_\Theta(m,e))(g)$ for $g\in G'$. We check the hypotheses of
Lemma~\ref{gal-shintani-clifford} for $H=G'$, $\Gamma=\operatorname{Gal}(k'/k)$, $X'$ and $Y$.

\emph{Norm on the coset.} By Lemma~\ref{gal-support-finite}(1), $\kappa_\Theta(m,e)$ is the character of an
irreducible representation of $G$, so $\sum_{\gamma\in G}|\kappa_\Theta(m,e)(\gamma)|^2=|G|$. By
Lemma~\ref{gal-shintani-lift}(2), $\sum_{g\in G'}|X'(gs)|^2=|G'|$.

\emph{The target.} The triple $(q;\emptyset;((m,e)))$ is a pattern, so $(q^l;\emptyset;A)$ is a pattern by
Lemma~\ref{gal-support-pattern-arithmetic}(1), and $Y$ is a character by Lemma~\ref{gal-support-finite}(3). It
satisfies $Y(s(g))=Y(g)$. Indeed, $s$ is $z\mapsto z^{q^i}$ for some $i$, and
$\kappa_\Theta(m'',e'')(s(x))=\kappa_\Theta(m'',e''q^i)(x)$ for $x\in\operatorname{GL}_{m''}(k')$. If $l$ divides
$m$, it follows that $Y(s(g))$ is the Harish-Chandra induction of the $\kappa_\Theta(m/l,eq^{j+i})$,
$0\leq j<l$, at $g$; since $\kappa_\Theta(m/l,eq^{j+l})=\kappa_\Theta(m/l,eq^j)$ by
Lemma~\ref{gal-support-finite}(2) over $k'$, these are the factors of $Y$ in another order, and Alpöge shows
that the order does not change the induction. If $l$ does not divide $m$, then
$Y=\kappa_\Theta(m,e')$ with $e'=e(1+q^m+\cdots+q^{m(l-1)})$, and there is $a$ with
$e'q\equiv e'q^{la}\pmod{q^{lm}-1}$; so $e'q^i\equiv e'q^{lai}$, and
$\kappa_\Theta(m,e'q^i)=\kappa_\Theta(m,e')$ by Lemma~\ref{gal-support-finite}(2) over $k'$.

\emph{Pairing.} By Lemma~\ref{gal-shintani-pairing}, $\sum_{g\in G'}X'(g)\overline{Y(g)}=z|G'|$ with $z=\pm1$ if
$l$ divides $m$, and with $z\equiv1\pmod l$ otherwise. In both cases $l$ does not divide $z$.

Lemma~\ref{gal-shintani-clifford} now shows that $g\mapsto X'(gs)=\operatorname{Sh}_s(\kappa_\Theta(m,e))(g)$ is
carried by $Y$.
\end{proof}

\begin{lemma}\label{gal-support-shintani}
Let $k\subset k'$ be finite fields with $[k':k]=l$ a prime, $q=|k|$, $s$ the automorphism $z\mapsto z^q$ of
$k'$, $\bar k$ an algebraic closure of $k'$ and $\Theta:\bar k^\times\to\mathbf C^\times$ an injective
homomorphism. Let $a=(q;D;L)$ be a pattern of base $q$ and rank $n$ (Definition~\ref{gal-support-pattern}),
and $a^{(l)}=(q^l;D^{(l)};L^{(l)})$ its lift by $l$. Then on $\operatorname{GL}_n(k')$ the Shintani lift
$\operatorname{Sh}_s(J_\Theta(D;L))$ is carried by $J_\Theta(D^{(l)};L^{(l)})$
(Definition~\ref{gal-support-level-zero}).
\end{lemma}
\leanref{Development/AutomorphicArithmetic/Foundations/GaloisRepresentations/RegularSemisimpleFibres/ResEmbCompat.lean}{UpperTheory.Gal.InertialTypesAlt.KillsOn}
\leanref{Development/AutomorphicArithmetic/Foundations/GaloisRepresentations/RegularSemisimpleFibres/PKillsTrue.lean}{UpperTheory.Gal.InertialTypesAlt.killsTrue}
\leanref{Development/AutomorphicArithmetic/Foundations/GaloisRepresentations/RegularSemisimpleFibres/PKillsNil.lean}{UpperTheory.Gal.InertialTypesAlt.killsNil}
\leanref{Development/AutomorphicArithmetic/Foundations/GaloisRepresentations/RegularSemisimpleFibres/NumberField.lean}{UpperTheory.Gal.InertialTypesAlt.KillsAll}
\leanref{Development/AutomorphicArithmetic/Foundations/GaloisRepresentations/RegularSemisimpleFibres/ResEmbCompat.lean}{UpperTheory.Gal.InertialTypesAlt.LiftHc}
\leanref{Development/AutomorphicArithmetic/Foundations/GaloisRepresentations/RegularSemisimpleFibres/ResEmbCompat.lean}{UpperTheory.Gal.InertialTypesAlt.KillsHc}
\leanref{Development/AutomorphicArithmetic/Foundations/GaloisRepresentations/RegularSemisimpleFibres/ResEmbCompat.lean}{UpperTheory.Gal.InertialTypesAlt.LiftDet}
\leanref{Development/AutomorphicArithmetic/Foundations/GaloisRepresentations/RegularSemisimpleFibres/ResEmbCompat.lean}{UpperTheory.Gal.InertialTypesAlt.NormEx}

\begin{proof}
We use the notation of Situation~\ref{gal-shintani-situation} for $k\subset k'$ and the generator $s$, and
$R^s$ is the $s$-twisted Harish-Chandra induction (Definition~\ref{gal-shintani-twisted}(1)). By
Lemma~\ref{gal-support-pattern-arithmetic}(1) and Lemma~\ref{gal-support-finite}(3),
$J_\Theta(D^{(l)};L^{(l)})$ is a character. The functions $J_\Theta(D;L)$ are class functions by
Lemma~\ref{gal-shintani-parabolic}(2). We use four facts, for $d+b=n$.
\begin{enumerate}
\item[(a)] For class functions $\alpha$ on $\operatorname{GL}_d(k)$ and $\beta$ on
  $\operatorname{GL}_b(k)$, $\operatorname{Sh}_s(R_{(d,b)}(\alpha,\beta))=
  R^s_{(d,b)}(\operatorname{Sh}_s(\alpha),\operatorname{Sh}_s(\beta))$. This is
  Lemma~\ref{gal-shintani-hc}.
\item[(b)] For $c\geq0$ and $c'=c(1+q+\cdots+q^{l-1})$,
  $\operatorname{Sh}_s(\delta_\Theta(d,c))=\delta_\Theta(d,c')$. Indeed, by Lemma~\ref{gal-shintani-norm}(5),
  $N_s(\delta)$ is conjugate to some $\gamma\in\operatorname{GL}_d(k)$ for every
  $\delta\in\operatorname{GL}_d(k')$, and
  $\det\gamma=\det N_s(\delta)=\prod_{i<l}s^i(\det\delta)=(\det\delta)^{1+q+\cdots+q^{l-1}}$.
\item[(c)] If $\alpha'$ is carried by a character $A$ and $\beta'$ by a character $B$, then
  $R^s_{(d,b)}(\alpha',\beta')$ is carried by $R_{(d,b)}(A,B)$. Indeed, let $\rho$ be a representation of
  $G'$ that does not meet $R_{(d,b)}(A,B)$, let $U$ be the unipotent radical of $P=P'_{(d,b)}$ and $M$ its
  block diagonal subgroup. Substituting $p=y^{-1}xs(y)$,
  $\sum_x\overline{R^s(\alpha',\beta')(x)}\rho(x)=|P|^{-1}\sum_y\rho(y)S\rho(s(y))^{-1}$ with
  $S=\sum_{p\in P}\overline{\alpha'(\pi_1p)\beta'(\pi_2p)}\rho(p)$, which is
  $\sum_{m\in M}\overline{\alpha'(\pi_1m)\beta'(\pi_2m)}\rho(m)$ composed with $|U|$ times the projection
  onto the $U$-invariants $\rho^U$. By Lemma~\ref{gal-support-finite}(4), the representation $\rho^U$ of
  $M$ does not meet $(m_1,m_2)\mapsto A(m_1)B(m_2)$. Write $\rho^U$ as a sum of irreducible
  representations $\lambda\otimes\mu$; since $A$ and $B$ are characters, each $\lambda$ does not meet $A$ or
  $\mu$ does not meet $B$, so
  $\sum_m\overline{\alpha'(m_1)\beta'(m_2)}(\lambda\otimes\mu)(m)=
  \bigl(\sum\overline{\alpha'}\lambda\bigr)\otimes\bigl(\sum\overline{\beta'}\mu\bigr)=0$. Hence $S=0$.
\item[(d)] A one-dimensional character $\chi$ is carried by itself, since
  $|H|^{-1}\sum_h\overline{\chi(h)}\rho(h)$ is the projection onto the $\chi$-isotypic part of $\rho$.
\end{enumerate}
We argue by induction on the length of $D$.

\emph{$D$ empty.} We argue by induction on the length of $L$. If $L$ is empty, then $n=0$, and
$\operatorname{Sh}_s(J_\Theta(\emptyset;\emptyset))$ and $J_\Theta(\emptyset;\emptyset)$ are the constant
function $1$ on the trivial group, which is carried by itself by (d). Otherwise let $L$ be $(m,e)$ followed by
$L_1$, let $b$ be the rank of $(q;\emptyset;L_1)$, and let $L_0$ be the list of pairs that replace $(m,e)$ in
$L^{(l)}$, so that $L^{(l)}$ is $L_0$ followed by $L_1^{(l)}$. By transitivity of induction,
$J_\Theta(\emptyset;L)=R_{(m,b)}(\kappa_\Theta(m,e),J_\Theta(\emptyset;L_1))$ and
$J_\Theta(\emptyset;L^{(l)})=R_{(m,b)}(J_\Theta(\emptyset;L_0),J_\Theta(\emptyset;L_1^{(l)}))$. By (a),
$\operatorname{Sh}_s(J_\Theta(\emptyset;L))=R^s_{(m,b)}(\operatorname{Sh}_s(\kappa_\Theta(m,e)),
\operatorname{Sh}_s(J_\Theta(\emptyset;L_1)))$. By Lemma~\ref{gal-shintani-cuspidal},
$\operatorname{Sh}_s(\kappa_\Theta(m,e))$ is carried by $J_\Theta(\emptyset;L_0)$, and by induction
$\operatorname{Sh}_s(J_\Theta(\emptyset;L_1))$ is carried by $J_\Theta(\emptyset;L_1^{(l)})$. Both are characters
by Lemma~\ref{gal-support-pattern-arithmetic}(1) and Lemma~\ref{gal-support-finite}(3), so by (c)
$\operatorname{Sh}_s(J_\Theta(\emptyset;L))$ is carried by $J_\Theta(\emptyset;L^{(l)})$.

\emph{Induction step.} Let $D$ be $(d,c)$ followed by $D_1$, let $b$ be the rank of $(q;D_1;L)$ and
$c'=c(1+q+\cdots+q^{l-1})$. By transitivity of induction,
$J_\Theta(D;L)=R_{(d,b)}(\delta_\Theta(d,c),J_\Theta(D_1;L))$ and
$J_\Theta(D^{(l)};L^{(l)})=R_{(d,b)}(\delta_\Theta(d,c'),J_\Theta(D_1^{(l)};L^{(l)}))$. By induction,
$\operatorname{Sh}_s(J_\Theta(D_1;L))$ is carried by the character $J_\Theta(D_1^{(l)};L^{(l)})$. By (a) and (b),
$\operatorname{Sh}_s(J_\Theta(D;L))=R^s_{(d,b)}(\delta_\Theta(d,c'),\operatorname{Sh}_s(J_\Theta(D_1;L)))$, and by
(c) and (d) it is carried by $J_\Theta(D^{(l)};L^{(l)})$.
% Lean: (a) is LiftHc = GlobalTheory.T.shintaniLift_harishChandraInduction_eq_twisted
% (GlobalTheory.Shin.shLift_hc_eq_hcTw); (c) is KillsHc (GlobalTheory.Shin.HcIso.kills_hcTw); the case
% D empty is KillsAll = GlobalTheory.T.shintani_lift_kills_product_of_cusp_forms, proved by
% GlobalTheory.Shin.DescentAll.kills_liftAll by the induction on L above (J_cons, LiftHc, liftAll split
% into A and liftAll L1, twistedHarishChandra_isotypic_part, and the two single-cuspidal lemmas); its
% base case uses isInCharacterIsotypicPart_self; (b) is LiftDet with NormEx
% (GlobalTheory.T.every_element_is_norm).
\end{proof}

\begin{remark}\label{gal-shintani-remark}
The case $D$ empty of Lemma~\ref{gal-support-shintani} also follows from Shintani's theorem \cite{Shn76}. Let
$J=J_\Theta(\emptyset;L)$, $J'=J_\Theta(\emptyset;L^{(l)})$, $G=\operatorname{GL}_n(k)$ and
$G'=\operatorname{GL}_n(k')$. Since $J'$ is a character, a representation of $G'$ that does not meet $J'$ has no
irreducible constituent that meets it, so it suffices to consider an irreducible representation $\rho$ of $G'$
that does not meet $J'$. Put $T=\sum_\delta\overline{\operatorname{Sh}_s(J)(\delta)}\rho(\delta)$. By
Lemma~\ref{gal-shintani-lift}(1), $\rho(y)^{-1}T\rho(s(y))=T$ for $y\in G'$. If $\rho\circ s\not\cong\rho$, then
$T=0$ by Schur's lemma. Otherwise choose $I$ with $I\rho(y)=\rho(s(y))I$; then $TI$ commutes with $\rho$, so
$T=\lambda I^{-1}$ with
$\lambda\dim\rho=\sum_\delta\overline{\operatorname{Sh}_s(J)(\delta)}\operatorname{tr}(\rho(\delta)I)$. By
Shintani's theorem, $I$ can be chosen so that $\operatorname{tr}(\rho(\delta)I)=\chi(\gamma)$ whenever
$N_s(\delta)$ is conjugate to $\gamma\in G$, for an irreducible character $\chi$ of $G$ whose Shintani lift is
$\rho$; and the Shintani lifts of the irreducible constituents of $J$ are constituents of $J'$, since Shintani
lifting replaces the cuspidal parameters $(m,e)$ by the pairs of $L^{(l)}$. By
Lemma~\ref{gal-shintani-norm}(5), every $N_s(\delta)$ is conjugate to an element of $G$, and the number of
$\delta$ with $N_s(\delta)$ conjugate to $\gamma$ is $|G'|/|G|$ times the number of conjugates of $\gamma$ in
$G$. So $\lambda\dim\rho$ is a nonzero multiple of $\sum_{\gamma\in G}\overline{J(\gamma)}\chi(\gamma)$, which
vanishes since $\chi$ is not a constituent of $J$.
\end{remark}

\begin{lemma}\label{gal-support-depth-zero}
Let $F$ be a nonarchimedean local field of characteristic $0$ with residue field $k$, let $n\geq1$, and
take level-zero parts as in Definition~\ref{gal-support-level-zero}(6).
\begin{enumerate}
\item Let $y\in\operatorname{Irr}(F)$ be supercuspidal of degree $m$ and depth zero
  (Definition~\ref{gal-irreducible-classes}). Its space of $K_1$-fixed vectors is finite-dimensional,
  and $y^0$ is an irreducible cuspidal representation of $\operatorname{GL}_m(k)$.
\item Let $\sigma$ be an irreducible smooth representation of $\operatorname{GL}_n(F)$ of depth zero, with
  supercuspidal support $\{x_1,\ldots,x_s\}$, $x_a$ of degree $m_a$. Then each $x_a$ has depth zero,
  $\sigma$ has a finite-dimensional space of $K_1$-fixed vectors, and after renumbering the support,
  with $m=(m_1,\ldots,m_s)$, there is a nonzero linear map $\psi$ from the space of $\sigma^0$ to that of $x_1^0\otimes\cdots\otimes x_s^0$ with
  $\psi\circ\sigma^0(p)=(x_1^0\otimes\cdots\otimes x_s^0)(\pi_1(p),\ldots,\pi_s(p))\circ\psi$ for every
  $p\in P_m$.
\item In (2), suppose the support is unlinked (Definition~\ref{gal-plain}), and number it in any way,
  with $m=(m_1,\ldots,m_s)$. Then $\sigma\cong x_1\times\cdots\times x_s$. If each $x_a$ has depth zero and
  $\chi_a$ is the character of $x_a^0$, then $\sigma^0$ meets every character $T$ of
  $\operatorname{GL}_n(k)$ with $\sum_x\overline{R_m(\chi_1,\ldots,\chi_s)(x)}\,T(x)\neq0$.
\item Suppose there are a prime $g$ and a continuous field isomorphism $\mathbf Q_g\to F$. If
  $y,y'\in\operatorname{Irr}(F)$ are supercuspidal of degree $m$ and depth zero and $y^0$ and $y'^0$ have
  the same character, then $y'=y\otimes\mu$ for an unramified character $\mu$ of $F^\times$.
\item Let $d=(d_1,\ldots,d_r)$ with sum $n$, let $\omega_1,\ldots,\omega_r$ be characters of $k^\times$, and
  let $\sigma$ be an irreducible smooth representation of $\operatorname{GL}_n(F)$ with a nonzero vector $v$
  such that $\sigma(g)v=\prod_i\omega_i(\det\pi_i(\bar g))v$ for every $g\in\operatorname{GL}_n(\mathcal O_F)$
  with $\bar g\in P_d$. Then there are characters $\xi_1,\ldots,\xi_n$ of $F^\times$, trivial on
  $1+\mathfrak m_F$, such that the supercuspidal support of $\sigma$ is $\{\xi_1,\ldots,\xi_n\}$ and
  $\xi_j(u)=\omega_{i(j)}(\bar u)$ for every unit $u$, where $i(j)$ is the index of the block of $d$
  containing $j$.
\item If in (5) the representation $\sigma$ is generic and $\omega_1,\ldots,\omega_r$ are pairwise distinct,
  then the supercuspidal support of $\sigma$ is unlinked.
\end{enumerate}
\end{lemma}
\leanref{Development/AutomorphicArithmetic/Foundations/GaloisRepresentations/RegularSemisimpleFibres/ResEmbCompat.lean}{UpperTheory.Gal.InertialTypesAlt.CuspIrr}
\leanref{Development/AutomorphicArithmetic/Foundations/GaloisRepresentations/RegularSemisimpleFibres/PCuspIrr.lean}{UpperTheory.Gal.InertialTypesAlt.cuspIrr}
\leanref{Development/AutomorphicArithmetic/Foundations/GaloisRepresentations/RegularSemisimpleFibres/ResEmbCompat.lean}{UpperTheory.Gal.InertialTypesAlt.SuppStd}
\leanref{Development/AutomorphicArithmetic/Foundations/GaloisRepresentations/RegularSemisimpleFibres/ResEmbCompat.lean}{UpperTheory.Gal.InertialTypesAlt.TensFix}
\leanref{Development/AutomorphicArithmetic/Foundations/GaloisRepresentations/RegularSemisimpleFibres/ResEmbCompat.lean}{UpperTheory.Gal.InertialTypesAlt.SubqMap}
\leanref{Development/AutomorphicArithmetic/Foundations/GaloisRepresentations/RegularSemisimpleFibres/SuppBCOfMix.lean}{UpperTheory.Gal.InertialTypesAlt.UnlFull}
\leanref{Development/AutomorphicArithmetic/Foundations/GaloisRepresentations/RegularSemisimpleFibres/SuppBCOfMix.lean}{UpperTheory.Gal.InertialTypesAlt.LblFull}
\leanref{Development/AutomorphicArithmetic/Foundations/GaloisRepresentations/RegularSemisimpleFibres/SuppBCOfMix.lean}{UpperTheory.Gal.InertialTypesAlt.InertOf}
\leanref{Development/AutomorphicArithmetic/Foundations/GaloisRepresentations/RegularSemisimpleFibres/PInertOf.lean}{UpperTheory.Gal.InertialTypesAlt.inertOf}
\leanref{Development/AutomorphicArithmetic/Foundations/GaloisRepresentations/IwahoriTypeSupport/HasTypeVecIff.lean}{UpperTheory.Gal.TypesAndIwahori.HasTypeVec}
\leanref{Development/AutomorphicArithmetic/Foundations/GaloisRepresentations/IwahoriTypeSupport/HasTypeVecIff.lean}{UpperTheory.Gal.TypesAndIwahori.SuppOfType}
\leanref{Development/AutomorphicArithmetic/Foundations/GaloisRepresentations/IwahoriTypeSupport/TameSuppUniqueHolds.lean}{UpperTheory.Gal.TypesAndIwahori.suppOfType_holds}
\leanref{Development/AutomorphicArithmetic/Foundations/GaloisRepresentations/RegularSemisimpleFibres/ResEmbCompat.lean}{UpperTheory.Gal.InertialTypesAlt.TameChars}
\leanref{Development/AutomorphicArithmetic/Foundations/GaloisRepresentations/IwahoriTypeSupport/HasTypeVecIff.lean}{UpperTheory.Gal.TypesAndIwahori.TopB}
\leanref{Development/AutomorphicArithmetic/Foundations/GaloisRepresentations/IwahoriTypeSupport/UnlType.lean}{UpperTheory.Gal.TypesAndIwahori.unlinked_of_typeVec}

\begin{proof}
(1) and (2) are the theory of depth-zero types and of Jacquet modules of Moy and Prasad \cite{MP96}: for
a depth-zero supercuspidal $y$, the $K_1$-fixed vectors form an irreducible cuspidal representation of
$\operatorname{GL}_m(k)$; $\sigma$ is a subquotient of $x_1\times\cdots\times x_s$ \cite{BZ77}, the members
of the support of a depth-zero representation have depth zero
(Proposition~\ref{gal-depth-zero-facts}(6)), and taking $K_1$-fixed vectors carries the Jacquet module
along the standard parabolic of type $m$ to the invariants of the unipotent radical of $P_m$.

(3) The product of supercuspidals with unlinked support is irreducible \cite[Theorem~9.7]{Zel80}, so it
is $\sigma$. By the Iwasawa decomposition, the $K_1$-fixed vectors of $x_1\times\cdots\times x_s$ contain
the representation of $\operatorname{GL}_n(k)$ induced from $P_m$ by the inflation of
$x_1^0\otimes\cdots\otimes x_s^0$, whose character is $R_m(\chi_1,\ldots,\chi_s)$
(Lemma~\ref{gal-support-finite}(3)). A character $T$ pairing nonzero with it shares an irreducible
constituent with it, and so $\sigma^0$ meets $T$.

(4) By Proposition~\ref{gal-depth-zero-facts}(1), $y$ and $y'$ contain cuspidal representations $\tau$ and
$\tau'$ of $\operatorname{GL}_m(\mathbf F_g)$. Since $y^0$ is irreducible by (1), the inflation of $\tau$
is $y^0$, and likewise for $\tau'$; so $\tau$ and $\tau'$ have the same character. By
Proposition~\ref{gal-depth-zero-facts}(4), the central characters $\omega_y$ and $\omega_{y'}$ agree on the
units. Let $\mu$ be an unramified character with $\mu^m=\omega_y\omega_{y'}^{-1}$. Then $y'\otimes\mu$ has
central character $\omega_y$ and contains $\tau'$, since $\mu\circ\det$ is trivial on
$\operatorname{GL}_m(\mathbf Z_g)$, so $y'\otimes\mu\cong y$ by Proposition~\ref{gal-depth-zero-facts}(3).

(5) The vector $v$ is fixed up to the character $\prod_j\omega_{i(j)}$ by the Iwahori subgroup, so this is
Roche's theory of types for principal series \cite{Roc98}: the supercuspidal support consists of
characters $\xi_j$ of $F^\times$ with $\xi_j(u)=\omega_{i(j)}(\bar u)$ on units, which are trivial on
$1+\mathfrak m_F$.

(6) Since the $\omega_i$ are distinct, Roche's theory \cite{Roc98} shows that
$\sigma\cong\sigma_1\times\cdots\times\sigma_r$ with $\sigma_i$ an irreducible representation of
$\operatorname{GL}_{d_i}(F)$ such that $\sigma_i\otimes(\eta_i^{-1}\circ\det)$ has a nonzero
$\operatorname{GL}_{d_i}(\mathcal O_F)$-fixed vector, for a character $\eta_i$ trivial on $1+\mathfrak m_F$
that restricts to $\omega_i$ on units. A product is generic only if its factors are \cite{BZ77}, and a generic
irreducible representation with a nonzero spherical vector is an irreducible product of unramified
characters, whose support is then unlinked \cite{Zel80}. So the support of each $\sigma_i$ is unlinked. Two members
from different blocks are not linked, since $\xi=\xi'\nu$ implies that $\xi$ and $\xi'$ agree on units,
while their restrictions $\omega_i\neq\omega_{i'}$ differ.
\end{proof}

\begin{lemma}\label{gal-support-pattern-exists}
Let $F$ be a nonarchimedean local field of characteristic $0$ with residue field $k$ of $q$ elements,
$\bar k$ an algebraic closure of $k$, $\Theta:\bar k^\times\to\mathbf C^\times$ an injective
homomorphism, $n\geq1$, and $\sigma$ an irreducible smooth representation of $\operatorname{GL}_n(F)$ of
depth zero with supercuspidal support $\{x_1,\ldots,x_s\}$, $x_a$ of degree $m_a$.
\begin{enumerate}
\item Each $x_a$ is a supercuspidal of depth zero, and after renumbering the support there are integers
  $c_a$ regular for $(q,m_a)$ such that $x_a^0$ has character $\kappa_\Theta(m_a,c_a)$ and $\sigma$ has the pattern
  $(q;\emptyset;((m_1,c_1),\ldots,(m_s,c_s)))$ with respect to $\Theta$
  (Definition~\ref{gal-support-pattern}).
\item Suppose that there are a prime $g$ and a continuous field isomorphism $\mathbf Q_g\to F$, that the
  support of $\sigma$ is unlinked (Definition~\ref{gal-plain}), and that no two of its members of degree
  at least $2$ are inertially equivalent. Then $\sigma$ has an apart pattern $(q;D;L)$ with respect to
  $\Theta$ such that the first entry of every pair in $L$ is the degree of a member of the support of
  $\sigma$.
\end{enumerate}
\end{lemma}
\leanref{Development/AutomorphicArithmetic/Foundations/GaloisRepresentations/RegularSemisimpleFibres/ResEmbCompat.lean}{UpperTheory.Gal.InertialTypesAlt.LblOf}
\leanref{Development/AutomorphicArithmetic/Foundations/GaloisRepresentations/RegularSemisimpleFibres/PLblOf.lean}{UpperTheory.Gal.InertialTypesAlt.lblOf}
\leanref{Development/AutomorphicArithmetic/Foundations/GaloisRepresentations/RegularSemisimpleFibres/SuppBCOfMix.lean}{UpperTheory.Gal.InertialTypesAlt.LblMixed}
\leanref{Development/AutomorphicArithmetic/Foundations/GaloisRepresentations/RegularSemisimpleFibres/PLblMixed.lean}{UpperTheory.Gal.InertialTypesAlt.lblMixed}

\begin{proof}
(1) By Lemma~\ref{gal-support-depth-zero}(2), the $x_a$ have depth zero, the $K_1$-fixed vectors of
$\sigma$ form a finite-dimensional space, and after renumbering there is a nonzero map $\psi$ from $\sigma^0$ to
$x_1^0\otimes\cdots\otimes x_s^0$, equivariant for $P_m$. By Lemma~\ref{gal-support-depth-zero}(1) each
$x_a^0$ is an irreducible cuspidal representation, so by Lemma~\ref{gal-support-finite}(1) its character
is $\kappa_\Theta(m_a,c_a)$ with $c_a$ regular for $(q,m_a)$. The representation
$p\mapsto(x_1^0\otimes\cdots\otimes x_s^0)(\pi_1(p),\ldots,\pi_s(p))$ of $P_m$ has character
$p\mapsto\prod_a\kappa_\Theta(m_a,c_a)(\pi_a(p))$. By Lemma~\ref{gal-support-finite}(5), $\psi$ shows that
the restriction of $\sigma^0$ to $P_m$ meets this character, and by Lemma~\ref{gal-support-finite}(4),
with $T=\operatorname{tr}\sigma^0$, the representation $\sigma^0$ meets
$J_\Theta(\emptyset;L)=R_m(\kappa_\Theta(m_1,c_1),\ldots,\kappa_\Theta(m_s,c_s))$. The degrees in the support
add up to $n$, so $(q;\emptyset;L)$ is a pattern of rank $n$.

(2) Take $x_a$ and $c_a$ as in (1). For $m_a=1$ replace $c_a$ by its remainder modulo $q-1$; by
Lemma~\ref{gal-support-finite}(2) this does not change $\kappa_\Theta(1,c_a)$, and regularity is empty for
$m=1$. Let $D$ list the pairs $(d,c)$, one for each value $c$ of the $c_a$ with $m_a=1$, with $d$ the
number of such $a$; and let $L$ list the pairs $(m_a,c_a)$ with $m_a\geq2$. Reorder the support so that
the list $((m_a,c_a))_a$ becomes $D^\sharp L$, in the notation of Lemma~\ref{gal-support-finite}(6). Since
the support is unlinked, $\sigma\cong x_1\times\cdots\times x_s$ in this order
(Lemma~\ref{gal-support-depth-zero}(3)). The triple $a=(q;D;L)$ is a pattern of rank $n$, $J_\Theta(D;L)$ is
a character (Lemma~\ref{gal-support-finite}(3)), and $J_\Theta(\emptyset;D^\sharp L)=R_m(\chi_1,\ldots,\chi_s)$
with $\chi_a=\kappa_\Theta(m_a,c_a)$ the character of $x_a^0$. By Lemma~\ref{gal-support-finite}(6) and
Lemma~\ref{gal-support-depth-zero}(3), $\sigma^0$ meets $J_\Theta(D;L)$, so $\sigma$ has the pattern $a$.

The pattern $a$ is apart. The $c$ in $D$ are distinct and less than $q-1$ by construction, and the
pairs in $L$ have first entry at least $2$. Suppose two pairs of $L$, coming from members $x_a$ and $x_b$
of the support with $m_a=m_b=m$, satisfy $c_b\equiv c_aq^h\pmod{q^m-1}$. Then $x_a^0$ and $x_b^0$ have the
same character by Lemma~\ref{gal-support-finite}(2), so $x_a$ and $x_b$ are inertially equivalent by
Lemma~\ref{gal-support-depth-zero}(4), which the hypothesis excludes. Finally, every pair of $L$ is
$(m_a,c_a)$ for a member $x_a$ of degree $m_a$.
\end{proof}

\begin{lemma}\label{gal-support-pattern-type}
Let $F$ be a nonarchimedean local field with residue field $k$, $\bar k$ an algebraic closure of $k$,
$\Theta:\bar k^\times\to\mathbf C^\times$ an injective homomorphism, and $b=(q;D;L)$ a pattern whose pairs
in $L$ all have first entry $1$ (Definition~\ref{gal-support-pattern}). Write $D$ followed by $L$ as
$((d_1,e_1),\ldots,(d_r,e_r))$ and let $\omega_i(z)=\Theta(z)^{e_i}$ for $z\in k^\times$. Let $\sigma$ be a
smooth representation of $\operatorname{GL}_n(F)$ that has the pattern $b$ with respect to $\Theta$. Then
$\sigma$ has a nonzero vector $v$ with $\sigma(g)v=\prod_i\omega_i(\det\pi_i(\bar g))v$ for every
$g\in\operatorname{GL}_n(\mathcal O_F)$ whose reduction $\bar g$ lies in $P_{(d_1,\ldots,d_r)}$.
\end{lemma}
\leanref{Development/AutomorphicArithmetic/Foundations/GaloisRepresentations/RegularSemisimpleFibres/ResEmbCompat.lean}{UpperTheory.Gal.InertialTypesAlt.TypeOfLbl}
\leanref{Development/AutomorphicArithmetic/Foundations/GaloisRepresentations/RegularSemisimpleFibres/PTypeOfLbl.lean}{UpperTheory.Gal.InertialTypesAlt.typeOfLbl}

\begin{proof}
Let $d=(d_1,\ldots,d_r)$. Since $\kappa_\Theta(1,e)=\delta_\Theta(1,e)$ (Lemma~\ref{gal-support-finite}(2)),
$J_\Theta(D;L)=R_d(\delta_\Theta(d_1,e_1),\ldots,\delta_\Theta(d_r,e_r))$. Let $\chi$ be the one-dimensional
character $p\mapsto\prod_i\omega_i(\det\pi_i(p))$ of $P_d$. By Lemma~\ref{gal-support-finite}(4), the sum
$\sum_{p\in P_d}\overline{\chi(p)}\operatorname{tr}\sigma^0(p)$ is a nonzero multiple of
$\sum_x\overline{J_\Theta(D;L)(x)}\operatorname{tr}\sigma^0(x)$, which is nonzero because $\sigma$ has the
pattern $b$. By Lemma~\ref{gal-support-finite}(5) there is a nonzero $P_d$-equivariant map from $\chi$ to
the restriction of $\sigma^0$ to $P_d$; the image $v$ of $1$ is the vector sought.
\end{proof}

\begin{lemma}\label{gal-support-matching}
Let $F$ be a nonarchimedean local field of characteristic $0$, $F'$ an unramified extension of $F$ of
degree $l$, $q=|k_F|$, and $s$ the automorphism of $F'$ over $F$ that induces $z\mapsto z^q$ on $k_{F'}$.
Let $n\geq1$, $G=\operatorname{GL}_n(F')$, $H=\operatorname{GL}_n(F)$, $U=\operatorname{GL}_n(\mathcal O_{F'})$,
$U_0=\operatorname{GL}_n(\mathcal O_F)$, and $N(\delta)=\delta\,s(\delta)\cdots s^{l-1}(\delta)$ for
$\delta\in G$. For a class function $f$ on $\operatorname{GL}_n(k_F)$ let $f^\natural(x)=f(\bar x)$ on $U_0$
and $\varphi(x)=\operatorname{Sh}_s(f)(\bar x)$ on $U$ (Definition~\ref{gal-support-level-zero}), both $0$
elsewhere. For $\delta\in G$ let $G^s_\delta=\{g\in G:g\delta=\delta s(g)\}$, $C$ its subgroup generated by
the elements that generate relatively compact subgroups, and
\[
 O^s_\delta(\varphi)=\sum_x[C:C\cap xUx^{-1}]\int_U\varphi\bigl((xu)^{-1}\delta s(xu)\bigr)\,du,
\]
the sum over representatives $x$ of the double cosets $G^s_\delta xU$ and $du$ the Haar measure of mass
$1$; define $O_\gamma(f^\natural)$ for $\gamma\in H$ in the same way with $H$, $U_0$ and the identity in
place of $G$, $U$ and $s$.
\begin{enumerate}
\item If $\delta\in G$ and $\gamma\in H$ are such that $N(\delta)$ and $\gamma$ have the same characteristic
  polynomial, and it is separable, then $O^s_\delta(\varphi)=O_\gamma(f^\natural)$.
\item If $\gamma\in H$ has separable characteristic polynomial and $O_\gamma(f^\natural)\neq0$, then there
  is $\delta\in G$ such that $N(\delta)$ has the characteristic polynomial of $\gamma$.
\end{enumerate}
\end{lemma}
\leanref{Development/AbelianVarietyTheory/Definitions/Divisors/IntersectionFamilySymmetric.lean}{GlobalTheory.DZA.InflMatching}
\leanref{Development/AutomorphicArithmetic/Foundations/GaloisRepresentations/RegularSemisimpleFibres/FrobeniusLocalPairs/InflMatch.lean}{UpperTheory.Gal.InertialTypesAlt.L.inflMatch}
\leanref{Development/AbelianVarietyTheory/Definitions/Divisors/SimpleTwistedTraceFormulaStable.lean}{GlobalTheory.BCT.twOrb}
\leanref{Development/AutomorphicArithmetic/Foundations/GaloisRepresentations/RegularSemisimpleFibres/FrobeniusLocalPairs/FrobeniusTwistedHilbertNinetyAlt.lean}{UpperTheory.Gal.InertialTypesAlt.L.GoodDelta}
\leanref{Development/AutomorphicArithmetic/Foundations/GaloisRepresentations/RegularSemisimpleFibres/FrobeniusLocalPairs/GoodDelta.lean}{UpperTheory.Gal.InertialTypesAlt.L.goodDelta}
\leanref{Development/AutomorphicArithmetic/Foundations/GaloisRepresentations/RegularSemisimpleFibres/FrobeniusLocalPairs/FrobeniusTwistedHilbertNinetyAlt.lean}{UpperTheory.Gal.InertialTypesAlt.L.LatDesc}
\leanref{Development/AutomorphicArithmetic/Foundations/GaloisRepresentations/RegularSemisimpleFibres/FrobeniusLocalPairs/LatDesc.lean}{UpperTheory.Gal.InertialTypesAlt.L.latDesc}
\leanref{Development/AutomorphicArithmetic/Foundations/GaloisRepresentations/RegularSemisimpleFibres/FrobeniusLocalPairs/FrobeniusTwistedHilbertNinetyAlt.lean}{UpperTheory.Gal.InertialTypesAlt.L.OrbTransfer}
\leanref{Development/AutomorphicArithmetic/Foundations/GaloisRepresentations/RegularSemisimpleFibres/FrobeniusLocalPairs/OrbTransfer.lean}{UpperTheory.Gal.InertialTypesAlt.L.orbTransfer}
\leanref{Development/AutomorphicArithmetic/Foundations/GaloisRepresentations/RegularSemisimpleFibres/FrobeniusLocalPairs/FrobeniusTwistedHilbertNinetyAlt.lean}{UpperTheory.Gal.InertialTypesAlt.L.RedNorm}
\leanref{Development/AutomorphicArithmetic/Foundations/GaloisRepresentations/RegularSemisimpleFibres/FrobeniusLocalPairs/RedNorm.lean}{UpperTheory.Gal.InertialTypesAlt.L.redNorm}
\leanref{Development/AutomorphicArithmetic/Foundations/Convolution/TwistedLocalOrbits/NormInj.lean}{UpperTheory.ConvolutionTheory.Loc.normInj}

\begin{proof}
Call $\gamma\in H$ \emph{integral} if its characteristic polynomial has coefficients in $\mathcal O_F$ and
$\det\gamma$ is a unit. We use five facts.
\begin{enumerate}
\item[(a)] For $\beta\in U$, $N(\beta)\in U$ and $\overline{N(\beta)}=N_s(\bar\beta)$; for $g,\beta\in G$,
  $N(g^{-1}\beta s(g))=g^{-1}N(\beta)g$; and $\varphi(u^{-1}\beta s(u))=\varphi(\beta)$ for $u\in U$. If
  $\beta\in U$ and $N(\beta)=\gamma\in H$, then $\varphi(\beta)=f^\natural(\gamma)$, since
  $N_s(\bar\beta)=\bar\gamma$. An element $h\in H$ lies in $U$ if and only if it lies in $U_0$, and if
  some $g^{-1}\gamma g$ with $g\in G$ lies in $U$, then $\gamma$ is integral.
\item[(b)] If $\gamma\in H$ is integral, there is $\delta_0\in G$ with $N(\delta_0)=\gamma$ that commutes
  with every element of $G$ commuting with $\gamma$, and such that $g^{-1}\delta_0g\in U$ whenever
  $g\in G$ and $g^{-1}\gamma g\in U$. Indeed, let $\chi_\gamma$ be the characteristic polynomial of
  $\gamma$. The class of $Z$ is a unit in $R=\mathcal O_F[Z]/(\chi_\gamma)$, and the norm
  $r\mapsto\prod_{i<l}s^i(r)$ from $(R\otimes_{\mathcal O_F}\mathcal O_{F'})^\times$ to $R^\times$ is surjective,
  because $R$ is a finite product of complete local rings with finite residue fields and $F'/F$ is
  unramified \cite[Chapter~V, \S2]{Ser79}. So there is $P\in\mathcal O_{F'}[Z]$ with
  $\prod_{i<l}s^i(P)\equiv Z$ modulo $\chi_\gamma$; put $\delta_0=P(\gamma)$. Then
  $N(\delta_0)=\prod_is^i(P)(\gamma)=\gamma$ by the Cayley--Hamilton theorem, so $\delta_0$ is invertible
  and $\det\delta_0$ is a unit, and $g^{-1}\delta_0g=P(g^{-1}\gamma g)$.
\item[(c)] If $g\in G$ and $g^{-1}s(g)\in U$, then $g\in HU$. Indeed, $u=g^{-1}s(g)$ is a cocycle of
  $\operatorname{Gal}(F'/F)$ with values in $U$, which is a coboundary $w^{-1}s(w)$, $w\in U$, by Lang's
  theorem and successive approximation \cite{Lan56}; then $gw^{-1}$ is fixed by $s$.
\item[(d)] If $\gamma\in H$ has separable characteristic polynomial and $N(\delta)$ and $N(\delta')$ are
  conjugate to $\gamma$ in $G$, then $\delta'=y^{-1}\delta s(y)$ for some $y\in G$ \cite[Chapter~1,
  \S1]{AC89}. If $\delta_0$ is as in (b), then $G^s_{\delta_0}$ is the centralizer $H_\gamma$ of $\gamma$ in
  $H$: an element $h\in H_\gamma$ commutes with $\delta_0$ and $s(h)=h$; conversely, if
  $g\delta_0=\delta_0s(g)$, then $g\gamma g^{-1}=N(g\delta_0s(g)^{-1})=\gamma$, so $g$ commutes with $\delta_0$,
  and then $s(g)=g$.
\item[(e)] Let $\delta_0$ and $\gamma$ be as in (d). If $\varphi(x^{-1}\delta_0s(x))\neq0$ implies $x\in HU$,
  and $\varphi(h^{-1}\delta_0h)=f^\natural(h^{-1}\gamma h)$ for all $h\in H$, then
  $O^s_{\delta_0}(\varphi)=O_\gamma(f^\natural)$: the double cosets that contribute to $O^s_{\delta_0}(\varphi)$
  are those of elements of $H$, and $h\mapsto h$ identifies $H_\gamma\backslash H/U_0$ with them, with the
  same indices and integrals since $U\cap H=U_0$ and $\varphi$ is invariant under $u\mapsto u^{-1}\beta s(u)$.
\end{enumerate}

(1) The element $\gamma$ is regular semisimple, and $N(\delta)$ is conjugate to $\gamma$ in $G$. Suppose
$\gamma$ is integral, and let $\delta_0$ be as in (b). By (d), $\delta=y^{-1}\delta_0s(y)$, and $O^s$ is
invariant under such twisted conjugation, so $O^s_\delta(\varphi)=O^s_{\delta_0}(\varphi)$. We check the
hypotheses of (e). If $\varphi(x^{-1}\delta_0s(x))\neq0$, then $x^{-1}\delta_0s(x)\in U$, so
$x^{-1}\gamma x=N(x^{-1}\delta_0s(x))\in U$ by (a), hence $x^{-1}\delta_0x\in U$ by (b), and
$x^{-1}s(x)=(x^{-1}\delta_0x)^{-1}(x^{-1}\delta_0s(x))\in U$, so $x\in HU$ by (c). For $h\in H$,
$N(h^{-1}\delta_0h)=h^{-1}\gamma h$. If $h^{-1}\delta_0h\in U$, then
$\varphi(h^{-1}\delta_0h)=f^\natural(h^{-1}\gamma h)$ by (a). Otherwise $\varphi(h^{-1}\delta_0h)=0$, and
$h^{-1}\gamma h\notin U_0$ by (b), so $f^\natural(h^{-1}\gamma h)=0$. So (e) applies. If $\gamma$ is not
integral, no conjugate of $\gamma$ lies in $U$ by (a), so $f^\natural(x^{-1}\gamma x)=0$ for all $x\in H$,
and $\varphi(x^{-1}\delta s(x))=0$ for all $x\in G$, since its norm $x^{-1}N(\delta)x$ is conjugate to
$\gamma$ and would lie in $U$; both sides vanish.

(2) Some $x^{-1}\gamma x$ lies in $U_0$, so $\gamma$ is integral by (a), and $\delta_0$ of (b) has
$N(\delta_0)=\gamma$.
\end{proof}

\begin{lemma}\label{gal-support-stable}
Let $K\subset M$ be totally complex number fields with $M$ a Galois extension of $K$ of prime degree,
$n\geq1$, and $\pi$ and $\Pi$ cuspidal automorphic representations of $\operatorname{GL}_n(\mathbf A_K)$ and
$\operatorname{GL}_n(\mathbf A_M)$ with $\Pi$ a weak base change of $\pi$
(Definition~\ref{gal-depth-zero-lift-weak}). Then $\Pi\circ\gamma\cong\Pi$ for every
$\gamma\in\operatorname{Gal}(M/K)$, and the local components of $\pi$ and $\Pi$ agree at places of degree
one.
\end{lemma}
\leanref{Development/AutomorphicArithmetic/Foundations/GaloisRepresentations/RegularSemisimpleFibres/ResEmbCompat.lean}{UpperTheory.Gal.InertialTypesAlt.StableOf}
\leanref{Development/AutomorphicArithmetic/Foundations/GaloisRepresentations/RegularSemisimpleFibres/PStableOf.lean}{UpperTheory.Gal.InertialTypesAlt.stableOf}

\begin{proof}
The stability is shown in the proof of Lemma~\ref{gal-depth-zero-lift-infinity}. By strong
multiplicity one \cite{JS81a}, $\Pi$ is the base change of $\pi$ of Arthur and Clozel, whose local
component at each place of $M$ is the local base change of the local component of $\pi$ below it
\cite[Chapter~3]{AC89}; at a place of degree one the local base change is the local component itself,
transported along the isomorphism of completions. Alpöge proves this agreement by his own comparison of
trace formulas.
\end{proof}

\begin{lemma}\label{gal-support-inert-trace}
Let $K\subset M$ be totally complex number fields with $M$ a Galois extension of $K$ of prime degree $p$,
$n\geq1$, and $\pi$ and $\Pi$ cuspidal automorphic representations of $\operatorname{GL}_n(\mathbf A_K)$ and
$\operatorname{GL}_n(\mathbf A_M)$ such that $\Pi$ is a weak base change of $\pi$
(Definition~\ref{gal-depth-zero-lift-weak}) and $\Pi\circ\gamma\cong\Pi$ for every
$\gamma\in\operatorname{Gal}(M/K)$. Let $v$ be a finite place of $K$ that is unramified and not split in
$M$, $u$ the place of $M$ over $v$, and $s$ the automorphism $z\mapsto z^{q_v}$ of $k_u$. Let $\sigma$ and
$\sigma'$ be local components of $\pi$ at $v$ and of $\Pi$ at $u$. Then for every class function $f$ on
$\operatorname{GL}_n(k_v)$ with $\sum_\gamma f(\gamma)\operatorname{tr}\sigma^0(\gamma)\neq0$, the operator
$\sum_{\delta\in\operatorname{GL}_n(k_u)}\operatorname{Sh}_s(f)(\delta)\,\sigma'^0(\delta)$ is nonzero
(Definition~\ref{gal-support-level-zero}).
\end{lemma}
\leanref{Development/AutomorphicArithmetic/Foundations/GaloisRepresentations/RegularSemisimpleFibres/NumberField.lean}{UpperTheory.Gal.InertialTypesAlt.FinTrace}
\leanref{Development/AutomorphicArithmetic/Foundations/GaloisRepresentations/RegularSemisimpleFibres/FrobeniusLocalPairs/FinTrace.lean}{UpperTheory.Gal.InertialTypesAlt.L.finTrace}
\leanref{Development/AutomorphicArithmetic/Foundations/GaloisRepresentations/RegularSemisimpleFibres/FrobeniusLocalPairs/InflTraceOfMatching.lean}{UpperTheory.Gal.InertialTypesAlt.L.finTrace_of_matching}
\leanref{Development/AutomorphicArithmetic/Foundations/GaloisRepresentations/RegularSemisimpleFibres/FrobeniusLocalPairs/FrobFactsAlt.lean}{UpperTheory.Gal.InertialTypesAlt.L.InflTrace}
\leanref{Development/AutomorphicArithmetic/Foundations/GaloisRepresentations/RegularSemisimpleFibres/FrobeniusLocalPairs/InflTrace.lean}{UpperTheory.Gal.InertialTypesAlt.L.inflTrace}
\leanref{Development/AutomorphicArithmetic/Foundations/GaloisRepresentations/RegularSemisimpleFibres/FrobeniusLocalPairs/FrobFactsAlt.lean}{UpperTheory.Gal.InertialTypesAlt.L.FinAct}

\begin{proof}
Since $v$ is unramified and not split and $p$ is prime, $u$ is the only place over $v$, $M_u$ is the
unramified extension of $K_v$ of degree $p$, and $s$ is induced by an automorphism of $M_u$ over $K_v$,
also written $s$. Let $f^\natural$ on $\operatorname{GL}_n(K_v)$ and $\varphi$ on $\operatorname{GL}_n(M_u)$
be as in Lemma~\ref{gal-support-matching}; by that lemma their orbital integrals match. By strong
multiplicity one \cite{JS81a}, $\Pi$ is the base change of $\pi$ of Arthur and Clozel, so $\sigma'$ is the
local base change of $\sigma$ \cite[Chapter~3]{AC89}: there are an operator $I$ on the space of $\sigma'$
with $I\sigma'(g)=\sigma'(s(g))I$ for all $g$ and $I^p=1$, and a constant $c\neq0$, such that
$\operatorname{tr}(\sigma'(\varphi)I)=c\operatorname{tr}\sigma(f^\natural)$ for all functions with matching
orbital integrals \cite[Chapter~1, \S6]{AC89}, where $\sigma(f^\natural)=\int f^\natural(g)\sigma(g)\,dg$
for the Haar measure giving $\operatorname{GL}_n(\mathcal O_v)$ volume $1$. Since $f^\natural$ is supported on
$\operatorname{GL}_n(\mathcal O_v)$ and invariant under $K_1$ on both sides,
$\operatorname{tr}\sigma(f^\natural)=|\operatorname{GL}_n(k_v)|^{-1}\sum_\gamma f(\gamma)\operatorname{tr}\sigma^0(\gamma)$.
Since $s(K_1)=K_1$, $I$ preserves the $K_1$-fixed vectors of $\sigma'$, and in the same way
$\operatorname{tr}(\sigma'(\varphi)I)=|\operatorname{GL}_n(k_u)|^{-1}\operatorname{tr}\bigl(\sum_\delta
\operatorname{Sh}_s(f)(\delta)\sigma'^0(\delta)\circ I\bigr)$. The second trace is nonzero, so the operator
is nonzero. Alpöge obtains the identity of traces at $u$ from his own comparison of simple twisted trace
formulas, applied to the matching functions of Lemma~\ref{gal-support-matching}.
\end{proof}

\begin{lemma}\label{gal-support-pattern-step}
Let $K\subset M$ be totally complex number fields inside an algebraic closure $\Omega$ of $K$, with $M$ a
Galois extension of $K$ of prime degree $p$. Let $n\geq1$ and $\pi$ and $\Pi$ cuspidal automorphic
representations of $\operatorname{GL}_n(\mathbf A_K)$ and $\operatorname{GL}_n(\mathbf A_M)$ with $\Pi$ a weak
base change of $\pi$ (Definition~\ref{gal-depth-zero-lift-weak}). Let $A$ be a valuation ring of $\Omega$
with residue field $k_A$, let $v$ and $u$ be the places of $K$ and $M$ under $A$, with $u$ unramified over
$v$, so that $k_v\subset k_u\subset k_A$, and let $\Theta:k_A^\times\to\mathbf C^\times$ be an injective
homomorphism. Let $\sigma$ and $\sigma'$ be local components of $\pi$ at $v$ and of $\Pi$ at $u$, and
suppose $\sigma$ has a pattern $a$ with respect to $\Theta$ (Definition~\ref{gal-support-pattern}).
\begin{enumerate}
\item If $v$ splits in $M$, then $\sigma'$ has the pattern $a$ with respect to $\Theta$.
\item Otherwise $\sigma'$ has the pattern $a^{(p)}$ with respect to $\Theta$.
\end{enumerate}
\end{lemma}
\leanref{Development/AutomorphicArithmetic/Foundations/GaloisRepresentations/RegularSemisimpleFibres/ResEmbCompat.lean}{UpperTheory.Gal.InertialTypesAlt.StepAt}
\leanref{Development/AutomorphicArithmetic/Foundations/GaloisRepresentations/RegularSemisimpleFibres/PStepSplit.lean}{UpperTheory.Gal.InertialTypesAlt.stepSplit}
\leanref{Development/AutomorphicArithmetic/Foundations/GaloisRepresentations/RegularSemisimpleFibres/PStepInert.lean}{UpperTheory.Gal.InertialTypesAlt.stepInert}
\leanref{Development/AutomorphicArithmetic/Foundations/GaloisRepresentations/RegularSemisimpleFibres/ResEmbCompat.lean}{UpperTheory.Gal.InertialTypesAlt.LblIso}
\leanref{Development/AutomorphicArithmetic/Foundations/GaloisRepresentations/RegularSemisimpleFibres/ResEmbCompat.lean}{UpperTheory.Gal.InertialTypesAlt.ResDeg}

\begin{proof}
Write $a=(q;D;L)$, of rank $n$.

(1) The completion map $i:K_v\to M_u$ is a continuous isomorphism extending $K\to M$, so by
Lemma~\ref{gal-support-stable}, $\sigma'\circ i\cong\sigma$. The isomorphism $i$ preserves the rings of
integers and induces the identity of $k_v=k_u$ inside $k_A$. It carries the $K_1$-fixed vectors of
$\sigma$ to those of $\sigma'$ and $\sigma^0$ to $\sigma'^0$, so $\sigma'^0$ meets $J_\Theta(D;L)$.

(2) Since $u$ is unramified over $v$ and $v$ does not split, $k_u$ is an extension of $k_v$ of degree $p$
with generator $s:z\mapsto z^q$, $q=|k_v|$. By Lemma~\ref{gal-support-stable}, $\Pi\circ\gamma\cong\Pi$ for
every $\gamma$. Let $f=\overline{J_\Theta(D;L)}$, a class function on $\operatorname{GL}_n(k_v)$. Since $\sigma$
has the pattern $a$, $\sum_\gamma f(\gamma)\operatorname{tr}\sigma^0(\gamma)\neq0$, so by
Lemma~\ref{gal-support-inert-trace} the operator
$\sum_\delta\overline{\operatorname{Sh}_s(J_\Theta(D;L))(\delta)}\,\sigma'^0(\delta)$ is nonzero. The space of
$K_1$-fixed vectors of $\sigma'$ is finite-dimensional by admissibility. By
Lemma~\ref{gal-support-shintani}, $\operatorname{Sh}_s(J_\Theta(D;L))$ is carried by $J_\Theta(D^{(p)};L^{(p)})$,
so $\sigma'^0$ meets $J_\Theta(D^{(p)};L^{(p)})$. By Lemma~\ref{gal-support-pattern-arithmetic}(1), $a^{(p)}$
is a pattern of rank $n$, and its base $q^p$ is $|k_u|$. So $\sigma'$ has the pattern $a^{(p)}$.
\end{proof}

\begin{lemma}\label{gal-support-pattern-chain}
Let $E$ be a totally complex number field, $\Omega$ an algebraic closure of $E$, $n\geq1$,
$X\subset\Omega$ a finite extension of $E$, and $\pi$ and $P$ cuspidal automorphic representations of
$\operatorname{GL}_n(\mathbf A_E)$ and $\operatorname{GL}_n(\mathbf A_X)$ with a chain of weak base changes
from $(E,\pi)$ to $(X,P)$ (Definition~\ref{gal-depth-zero-lift-weak}). Let $A$ be a valuation ring of
$\Omega$ with residue field $k_A$ such that every element of the inertia group $I_A$ of $A$ in
$\operatorname{Gal}(\Omega/E)$ fixes $X$, let $v$ and $u$ be the places of $E$ and $X$ under $A$, and let
$\Theta:k_A^\times\to\mathbf C^\times$ be an injective homomorphism. Let $\sigma$ and $\sigma'$ be local
components of $\pi$ at $v$ and of $P$ at $u$. If $\sigma$ has a pattern $a$ with respect to $\Theta$
(Definition~\ref{gal-support-pattern}), then $\sigma'$ has a pattern $b\geq a$ with respect to $\Theta$.
\end{lemma}
\leanref{Development/AutomorphicArithmetic/Foundations/GaloisRepresentations/RegularSemisimpleFibres/ResEmbCompat.lean}{UpperTheory.Gal.InertialTypesAlt.LblOn}
\leanref{Development/AutomorphicArithmetic/Foundations/GaloisRepresentations/RegularSemisimpleFibres/PLblOn.lean}{UpperTheory.Gal.InertialTypesAlt.lblOn}
\leanref{Development/AutomorphicArithmetic/Foundations/GaloisRepresentations/RegularSemisimpleFibres/ResEmbCompat.lean}{UpperTheory.Gal.InertialTypesAlt.UnrOfIne}

\begin{proof}
By induction on the length $m$ of the chain $E=K_0\subset\cdots\subset K_m=X$, $\pi_0=\pi,\ldots,\pi_m=P$.
If $m=0$, then $u=v$ and $\sigma$ and $\sigma'$ are local components of $\pi$ at $v$, hence isomorphic,
and $\sigma'$ has the pattern $a$. Let $m\geq1$, let $u_1$ be the place of $K_1$ under $A$ and $\sigma_1$
a local component of $\pi_1$ at $u_1$. Every element of $I_A$ fixes $K_1\subset X$, so the inertia group
of $A$ in $\operatorname{Gal}(\Omega/K_1)$ is $I_A$ and $u_1$ is unramified over $v$. By
Lemma~\ref{gal-support-pattern-step}, $\sigma_1$ has the pattern $a$ or $a^{(p)}$, with $p=[K_1:E]$, and
in both cases a pattern $a_1\geq a$. The inertia group of $A$ in $\operatorname{Gal}(\Omega/K_1)$ fixes
$X$, and $K_1\subset\cdots\subset K_m$ is a chain of length $m-1$ from $(K_1,\pi_1)$ to $(X,P)$, so by
induction $\sigma'$ has a pattern $b\geq a_1\geq a$.
\end{proof}

\begin{lemma}\label{gal-support-inertia}
Let $E$ be a number field, $\Omega$ an algebraic closure of $E$, $A$ a valuation ring of $\Omega$ with
residue field $k_A$, $\Theta:k_A^\times\to\mathbf C^\times$ an injective homomorphism, and $g$ a prime. For
$t$ in an inertia group of $A$ and $\varpi\in\Omega^\times$, $t(\varpi)/\varpi$ lies in $A$; write $\bar z$
for the residue of $z\in A$.
\begin{enumerate}
\item Let $K\subset\Omega$ be a number field and $w$ the place of $K$ under $A$, with $g$ in the prime
  ideal of $w$ but not in its square, let $a'$ be a local Artin map for $K$, $w$ and $A$
  (Definition~\ref{gal-local-artin}), and let $\varpi^{N(w)-1}=g$. For every $t$ in the inertia group of
  $A$ in $\operatorname{Gal}(\Omega/K)$, $a'(t)$ is a unit of $K_w$ whose residue, in $k_w\subset k_A$, is
  $\overline{t(\varpi)/\varpi}$, and $\overline{t(\varpi)/\varpi}^{\,N(w)-1}=1$.
\item Let $v$ be the place of $E$ under $A$, suppose there is a continuous field isomorphism
  $\mathbf Q_g\to E_v$, and let $\operatorname{rc}$ satisfy the laws for $E$, $v$ and $A$
  (Definition~\ref{gal-llc-laws}). Let $y\in\operatorname{Irr}(E_v)$ be supercuspidal of degree $m$ and depth
  zero, such that $y^0$ has character $\kappa_\Theta(m,c)$ with $c$ regular for $(g,m)$
  (Definition~\ref{gal-support-level-zero}). Let $f\geq1$ with $m\mid f$, let $\varpi^{g^f-1}=g$, let $t$
  lie in the inertia group $I_A$ of $A$ in $\operatorname{Gal}(\Omega/E)$, and let $z=t(\varpi)/\varpi$. Then
  the characteristic polynomial of $\operatorname{rc}(y)(t)$ is
  \[
   \prod_{j=0}^{m-1}\Bigl(Y-\Theta(\bar z)^{cg^j(g^f-1)/(g^m-1)}\Bigr).
  \]
\end{enumerate}
\end{lemma}
\leanref{Development/AutomorphicArithmetic/Foundations/GaloisRepresentations/RegularSemisimpleFibres/ResEmbCompat.lean}{UpperTheory.Gal.InertialTypesAlt.ArtRead}
\leanref{Development/AutomorphicArithmetic/Foundations/GaloisRepresentations/RegularSemisimpleFibres/PArtRead.lean}{UpperTheory.Gal.InertialTypesAlt.artRead}
\leanref{Development/AutomorphicArithmetic/Foundations/GaloisRepresentations/RegularSemisimpleFibres/ResEmbCompat.lean}{UpperTheory.Gal.InertialTypesAlt.CuspIne}
\leanref{Development/AutomorphicArithmetic/Foundations/GaloisRepresentations/RegularSemisimpleFibres/PCuspIne.lean}{UpperTheory.Gal.InertialTypesAlt.cuspIne}

\begin{proof}
(1) Since $a'(t)$ has valuation $0$ by Definition~\ref{gal-local-artin}(2), it is a unit. Its residue is
that of $t(\varpi)/\varpi$ by the explicit tame reciprocity law for $K_w$, in which $g$ is a uniformizer
\cite{Ser67}, as in the proof of Lemma~\ref{gal-tame-inertia-character}(4). Finally
$(t(\varpi)/\varpi)^{N(w)-1}=t(g)/g=1$.

(2) Since $E_v\cong\mathbf Q_g$, $k_v=\mathbf F_g$. Take a Green datum of degree $m$, with residue field
$k'\subset k_A$ of $g^m$ elements and local Artin map $a'$ (Lemma~\ref{gal-tame-inertia-character}(1)).
By Proposition~\ref{gal-depth-zero-facts}(1), $y$ contains a cuspidal representation $\tau$ of
$\operatorname{GL}_m(\mathbf F_g)$; as in the proof of Lemma~\ref{gal-support-depth-zero}(4), the
inflation of $\tau$ is $y^0$, so $\tau$ has character $\kappa_\Theta(m,c)$. For $x$ with irreducible
characteristic polynomial and root $z_0\in k'$, $\kappa_\Theta(m,c)(x)=(-1)^{m-1}\sum_{i<m}\Theta(z_0^{g^i})^c$,
so $\tau$ is attached to the character $\theta=\Theta^c$ of $k'^\times$
(Definition~\ref{gal-green-datum}), which is regular since $c$ is regular and $\Theta$ injective. By
law~(12), the characteristic polynomial of $\operatorname{rc}(y)(t)$ is
$\prod_{j<m}(Y-\theta(\bar a'(t))^{g^j})$. Let $N=(g^f-1)/(g^m-1)$, an integer since $m\mid f$. Then
$(\varpi^N)^{g^m-1}=g$, and by Lemma~\ref{gal-tame-inertia-character}(4) applied to $\varpi^N$,
$\bar a'(t)=\overline{t(\varpi^N)/\varpi^N}=\bar z^N$. So
$\theta(\bar a'(t))^{g^j}=\Theta(\bar z)^{cg^jN}$.
\end{proof}

\subsection{Depth-zero compatibility and reduction to the self-dual scope}\label{subsection-gal-depth-zero-compatibility}

\begin{lemma}\label{gal-depth-zero-lift-support}
Let $E$ be a totally complex number field, $\Omega$ an algebraic closure of $E$, $n\geq1$,
$X\subset\Omega$ a finite extension of $E$, and $\pi$ and $P$ cuspidal automorphic representations of
$\operatorname{GL}_n(\mathbf A_E)$ and $\operatorname{GL}_n(\mathbf A_X)$ with a chain of weak base
changes from $(E,\pi)$ to $(X,P)$ (Definition~\ref{gal-depth-zero-lift-weak}). Let $A$ be a valuation
ring of $\Omega$, let $v$ and $u$ be the places of $E$ and $X$ under $A$, and suppose that there are a
prime $g$ and a continuous field isomorphism $\mathbf Q_g\to E_v$, and that every element of the
inertia group $I_A$ of $A$ in $\operatorname{Gal}(\Omega/E)$ fixes $X$. Let $\sigma$ and $\sigma'$ be
local components of $\pi$ at $v$ and of $P$ at $u$, with $\sigma$ of depth zero
(Definition~\ref{gal-irreducible-classes}), and let $f\geq1$ be such that the residue field of $u$ has
$g^f$ elements and the degree of every member of the supercuspidal support of $\sigma$ divides $f$.
Then there are smooth characters $\xi_1,\ldots,\xi_n$ of $X_u^\times$ with the following properties.
\begin{enumerate}
\item Each $\xi_j$ is trivial on the units $x$ of $X_u$ with $|x-1|<1$.
\item The supercuspidal support of $\sigma'$ is $\{\xi_1,\ldots,\xi_n\}$.
\item For every map $\operatorname{rc}$ for $E$, $v$ and $A$ that satisfies the laws
  (Definition~\ref{gal-llc-laws}), every local Artin map $a'$ for $X$, $u$ and $A$
  (Definition~\ref{gal-local-artin}), and every $t$ in the inertia group of $A$ in
  $\operatorname{Gal}(\Omega/X)$, the characteristic polynomial of $\operatorname{rc}(\sigma)(t)$ is
  $\prod_{j=1}^n(Y-\xi_j(a'(t)))$.
\item If the supercuspidal support of $\sigma$ is unlinked (Definition~\ref{gal-plain}) and no two of
  its members of degree at least $2$ are inertially equivalent, then $\{\xi_1,\ldots,\xi_n\}$ is
  unlinked.
\end{enumerate}
\end{lemma}
\leanref{Development/AutomorphicArithmetic/Foundations/GaloisRepresentations/RegularSemisimpleFibres/NumberField.lean}{UpperTheory.Gal.InertialTypesAlt.SuppBC}
\leanref{Development/AutomorphicArithmetic/Foundations/GaloisRepresentations/RegularSemisimpleFibres/NumberField.lean}{UpperTheory.Gal.InertialTypesAlt.ToralOver}
\leanref{Development/AutomorphicArithmetic/Foundations/GaloisRepresentations/RegularSemisimpleFibres/NumberField.lean}{UpperTheory.Gal.InertialTypesAlt.Mixed}
\leanref{Development/AutomorphicArithmetic/Foundations/LocalLanglands/RecTTransport.lean}{UpperTheory.Zero.LLC.TameChar}
\leanref{Development/AutomorphicArithmetic/Foundations/GaloisRepresentations/RegularSemisimpleFibres/Closed.lean}{UpperTheory.Gal.InertialTypesAlt.suppBC_outright}

\begin{proof}
Let $k_A$ be the residue field of $A$, an algebraic closure of $\mathbf F_g$ that contains the residue
fields $k_v\subset k_u$, and fix an injective homomorphism $\Theta:k_A^\times\to\mathbf C^\times$. Patterns
are taken with respect to $\Theta$ (Definition~\ref{gal-support-pattern}). Since $E_v\cong\mathbf Q_g$,
$k_v=\mathbf F_g$ and $g$ is not in the square of the prime ideal of $v$. Since every element of $I_A$
fixes $X$, the inertia group of $A$ in $\operatorname{Gal}(\Omega/X)$ is $I_A$ and $u$ is unramified over
$v$; so $g$ is not in the square of the prime ideal of $u$, and $f$ is the residue degree of $u$.

\emph{Items (1) to (3).} By Lemma~\ref{gal-support-pattern-exists}(1), after renumbering, the support of
$\sigma$ consists of supercuspidals $x_1,\ldots,x_s$ of depth zero and degrees $m_a$, with $x_a^0$ of
character $\kappa_\Theta(m_a,c_a)$, $c_a$ regular for $(g,m_a)$, and $\sigma$ has the pattern
$a=(g;\emptyset;L)$ with $L=((m_1,c_1),\ldots,(m_s,c_s))$. By Lemma~\ref{gal-support-pattern-chain},
$\sigma'$ has a pattern $b=(q';D';L')\geq a$. Then $q'=|k_u|=g^f$, so $f\geq1$. Every $m_a$ divides $f$,
so $a$ splits at $q'$, and by Lemma~\ref{gal-support-pattern-arithmetic}(2) every pair of $L'$ has first
entry $1$, $b$ has rank $n$, and $a$ and $b$ have the same blocks at $q'$. Write $D'$ followed by $L'$ as
$((d_1,e_1),\ldots,(d_r,e_r))$ and $\omega_i=\Theta^{e_i}$ on $k_u^\times$. By
Lemma~\ref{gal-support-pattern-type}, $\sigma'$ has a nonzero vector on which every
$h\in\operatorname{GL}_n(\mathcal O_u)$ with $\bar h\in P_{(d_1,\ldots,d_r)}$ acts by
$\prod_i\omega_i(\det\pi_i(\bar h))$. By Lemma~\ref{gal-support-depth-zero}(5) there are characters
$\xi_1,\ldots,\xi_n$ of $X_u^\times$, trivial on the units $x$ with $|x-1|<1$, such that the supercuspidal
support of $\sigma'$ is $\{\xi_1,\ldots,\xi_n\}$ and $\xi_j(w)=\Theta(\bar w)^{e_{i(j)}}$ for every unit $w$,
with $i(j)$ the block containing $j$. This gives (1) and (2).

For (3), let $\operatorname{rc}$, $a'$ and $t$ be as there, choose $\varpi\in\Omega$ with
$\varpi^{g^f-1}=g$, and let $z=t(\varpi)/\varpi$ and $y=\Theta(\bar z)$. By
Lemma~\ref{gal-support-inertia}(1) for $X$, $u$ and $a'$, $a'(t)$ is a unit with residue $\bar z$, and
$y^{g^f-1}=1$. So $\xi_j(a'(t))=y^{e_{i(j)}}$, and by Lemma~\ref{gal-support-pattern-arithmetic}(3),
\[
 \prod_{j=1}^n\bigl(Y-\xi_j(a'(t))\bigr)=\prod_{i=1}^r(Y-y^{e_i})^{d_i}=\operatorname{pol}_{b,q'}(y).
\]
As an automorphism of $\Omega$, $t$ lies in $I_A$. By Lemma~\ref{gal-support-inertia}(2), the
characteristic polynomial of $\operatorname{rc}(x_a)(t)$ is
$\prod_{h<m_a}(Y-y^{c_ag^h(g^f-1)/(g^{m_a}-1)})$, and by Lemma~\ref{gal-depth-zero-lift-characters}(1) the
characteristic polynomial of $\operatorname{rc}(\sigma)(t)$ is the product of these over $a$. By
Lemma~\ref{gal-support-pattern-arithmetic}(4) this product is $\operatorname{pol}_{a,q'}(y)$, which equals
$\operatorname{pol}_{b,q'}(y)$ since $a$ and $b$ have the same blocks at $q'$. This is (3).

\emph{Item (4).} Suppose the support of $\sigma$ is unlinked and no two of its members of degree at
least $2$ are inertially equivalent. By Lemma~\ref{gal-support-pattern-exists}(2), $\sigma$ has an apart
pattern $a=(g;D;L)$ such that the first entry of every pair in $L$ is the degree of a member of the
support. By Lemma~\ref{gal-support-pattern-chain}, $\sigma'$ has a pattern $b=(q';D';L')\geq a$, with
$q'=g^f$ as before. The first entries of $L$ divide $f$, so $a$ splits at $q'$, and by
Lemma~\ref{gal-support-pattern-arithmetic}(2) the pairs of $L'$ have first entry $1$ and $a$ and $b$ have
the same blocks at $q'$. By Lemma~\ref{gal-support-pattern-arithmetic}(5) their second entries are
pairwise distinct modulo $q'-1$. Write $D'$ followed by $L'$ as $((d_1,e_1),\ldots,(d_r,e_r))$. By
Lemma~\ref{gal-support-pattern-arithmetic}(3), the blocks of $b$ at $q'$ are the $(d_i,e_i)$, so the $e_i$
are pairwise distinct modulo $g^f-1$. Since $\Theta$ is injective and $k_u^\times$ is cyclic of order
$g^f-1$, the characters $\omega_i=\Theta^{e_i}$ of $k_u^\times$ are pairwise distinct. By
Lemma~\ref{gal-support-pattern-type}, $\sigma'$ has a vector as in
Lemma~\ref{gal-support-depth-zero}(5) for $(d_1,\ldots,d_r)$ and these $\omega_i$. As a local component of
a cuspidal automorphic representation of $\operatorname{GL}_n$, $\sigma'$ is generic \cite{Sha74}. So by
Lemma~\ref{gal-support-depth-zero}(6) its supercuspidal support $\{\xi_1,\ldots,\xi_n\}$ is unlinked.
% Lean: suppBC_holds L.inflMatch = suppBC_of_mix (suppAny_holds h) (mixUnlinked (mixType_holds h)).
% Items (1)-(3) follow suppAny with LblOn for patterns with D empty (killsNil); item (4) follows mixType
% (LblOn True via killsTrue, lblMixed, sepOf, leTop, polTop, typeOfLbl) and mixUnlinked
% (TypesAndIwahori.unlinked_of_typeVec with CuspRep.Generic.isGeneric_of_isLocComp).
\end{proof}

\begin{lemma}\label{gal-depth-zero-lift-characters}
In Situation~\ref{gal-local-situation}, let $\operatorname{rc}$ satisfy the laws with respect to the
local Artin map $a$ (Definition~\ref{gal-llc-laws}), let $n\geq1$, and let
$x\in\operatorname{Irr}(E_w)$ be of degree $n$ with supercuspidal support $\{\xi_1,\ldots,\xi_n\}$
for smooth characters $\xi_j$ of $E_w^\times$. Let $x^\flat=x\nu^{(1-n)/2}$.
\begin{enumerate}
\item For every $y\in\operatorname{Irr}(E_w)$ with supercuspidal support $\{\rho_1,\ldots,\rho_k\}$
  and every $t\in I_A$, the characteristic polynomial of $\operatorname{rc}(y)(t)$ is the product of
  those of the $\operatorname{rc}(\rho_i)(t)$.
\item If $\{\xi_1,\ldots,\xi_n\}$ is unlinked, then $x^\flat$ is plain (Definition~\ref{gal-plain}).
\item For every $t\in I_A$, the characteristic polynomial of $\operatorname{rc}(x^\flat)(t)$ is
  $\prod_{j=1}^n(Y-\xi_j(a(t)))$.
\end{enumerate}
\end{lemma}
\leanref{Development/AutomorphicArithmetic/Foundations/GaloisRepresentations/SphericalGaloisCompatibility/PAdicSplitting/InfLifts.lean}{UpperTheory.Gal.ToralGalois.R2.RecChars}
\leanref{Development/AutomorphicArithmetic/Foundations/GaloisRepresentations/SphericalGaloisCompatibility/PAdicSplitting/PRecChars.lean}{UpperTheory.Gal.ToralGalois.R2.recChars}
\leanref{Development/AutomorphicArithmetic/Foundations/GaloisRepresentations/RegularSemisimpleFibres/ResEmbCompat.lean}{UpperTheory.Gal.InertialTypesAlt.SuppIne}
\leanref{Development/AutomorphicArithmetic/Foundations/GaloisRepresentations/RegularSemisimpleFibres/PSuppIne.lean}{UpperTheory.Gal.InertialTypesAlt.suppIne}

\begin{proof}
(1) Let $y$ be the Langlands quotient of the segments $(\rho'_1,l_1),\ldots,(\rho'_s,l_s)$
\cite{Zel80}. Then $y$ is a subquotient of the product of the $\rho'_i\nu^h$, $0\leq h<l_i$, so by the
uniqueness of the supercuspidal support \cite{BZ77} these form the support $\{\rho_1,\ldots,\rho_k\}$.
Let $\sigma\in W_A$. By law~(3), $\operatorname{rc}(y)$ is the sum of the
$(\operatorname{rc}(\rho'_i)\otimes(|\cdot|^{(l_i-1)/2}\circ a))\otimes\operatorname{Sp}(l_i)$. On
$\operatorname{Sp}(l)$ the element $\sigma$ acts on $e_h$ by $q_w^{(h-(l-1)/2)d(\sigma)}$
(Definition~\ref{gal-weil-deligne}), and $|a(\sigma)|=q_w^{d(\sigma)}$ by
Definition~\ref{gal-local-artin}(2). So the eigenvalues of the $i$-th summand at $\sigma$ are those of
$\operatorname{rc}(\rho'_i)(\sigma)$ multiplied by $|a(\sigma)|^h$ for $0\leq h<l_i$, which by law~(5)
are the eigenvalues of the $\operatorname{rc}(\rho'_i\nu^h)(\sigma)$. Hence the characteristic
polynomial of $\operatorname{rc}(y)(\sigma)$ is the product of those of the
$\operatorname{rc}(\rho_i)(\sigma)$, in particular for $\sigma=t\in I_A$.

(2) The supercuspidal support of $x^\flat$ is $\{\xi_j\nu^{(1-n)/2}\}$, which consists of characters.
Since $\xi_i\nu^{(1-n)/2}=\xi_j\nu^{(1-n)/2}\nu$ if and only if $\xi_i=\xi_j\nu$, it is unlinked if
$\{\xi_1,\ldots,\xi_n\}$ is. So $x^\flat$ is plain.

(3) By (1) and law~(2), the characteristic polynomial of $\operatorname{rc}(x)(t)$ is
$\prod_j(Y-\xi_j(a(t)))$. By law~(5),
$\operatorname{rc}(x^\flat)=\operatorname{rc}(x)\otimes(|\cdot|^{(1-n)/2}\circ a)$, and $|a(t)|=1$
since $d(t)=0$. So $\operatorname{rc}(x^\flat)(t)$ has the same characteristic polynomial.
\end{proof}

\begin{lemma}\label{gal-depth-zero-lift}
Let $E$ be a CM field, $\Omega$ an algebraic closure of $E$, $\ell$ a prime,
$\iota:\bar{\mathbf Q}_\ell\to\mathbf C$ a field isomorphism, $n\geq1$, $(\pi,\chi)$ in scope
(Definition~\ref{gal-scope}) with $\pi$ a cuspidal automorphic representation of
$\operatorname{GL}_n(\mathbf A_E)$, and $r:G_E\to\operatorname{GL}_n(\bar{\mathbf Q}_\ell)$ weakly
attached to $\pi$ (Definition~\ref{gal-weakly-attached}). Let $w$ be a finite place of $E$ such that
there are a prime $g$ and a continuous field isomorphism $\mathbf Q_g\to E_w$, and let $\sigma$ be a
local component of $\pi$ at $w$ that has depth zero, whose supercuspidal support is unlinked, and no
two of whose members of degree at least $2$ are inertially equivalent. Let $A$ be a valuation ring of
$\Omega$ over $w$ and $\operatorname{rc}$ a map that satisfies the laws
(Definition~\ref{gal-llc-laws}). Then there are a CM field $X\subset\Omega$ that is a finite abelian
extension of $E$, a cuspidal automorphic representation $P$ of $\operatorname{GL}_n(\mathbf A_X)$, a
Hecke character $\chi'$ of the maximal totally real subfield of $X$, a local component $T$ of $P$ at
the place $u$ of $X$ under $A$, and a map $\operatorname{rc}'$ for $X$, $u$ and $A$ that satisfies the
laws, with the following properties: $(P,\chi')$ is in scope; $r|_{G_X}$ is weakly attached to $P$;
$Z_\tau(P)=Z_{\tau|_E}(\pi)$ for every embedding $\tau:X\to\mathbf C$; $T^\flat$ is plain
(Definition~\ref{gal-plain}); and every $t\in I_A$ lies in the inertia group of $A$ in
$\operatorname{Gal}(\Omega/X)$, where $\operatorname{rc}'(T^\flat)(t)$ and
$\operatorname{rc}(\sigma^\flat)(t)$ have the same characteristic polynomial.
\end{lemma}
\leanref{Development/AutomorphicArithmetic/Foundations/GaloisRepresentations/SphericalGaloisCompatibility/PAdicSplitting/InfLifts.lean}{UpperTheory.Gal.ToralGalois.R2.ToTor}
\leanref{Development/AutomorphicArithmetic/Foundations/GaloisRepresentations/SphericalGaloisCompatibility/PAdicSplitting/PToTor.lean}{UpperTheory.Gal.ToralGalois.R2.toTor}
\leanref{Development/AutomorphicArithmetic/Foundations/GaloisRepresentations/SphericalGaloisCompatibility/PAdicSplitting/LocalGlobalCompatAtIwahoriFixedOfUnlinked.lean}{UpperTheory.Gal.ToralGalois.R2.toTor_holds}
\leanref{Development/AutomorphicArithmetic/Foundations/GaloisRepresentations/CharpolyFromTrace/UnlinkedOfType.lean}{UpperTheory.Gal.NearKilling.IneNorm}
\leanref{Development/AutomorphicArithmetic/Foundations/GaloisRepresentations/CharpolyFromTrace/IneNorm.lean}{UpperTheory.Gal.NearKilling.ineNorm}

\begin{proof}
Let $F$ be the maximal totally real subfield of $E$ and $v$ the place of $F$ under $w$. Since
$E_w\cong\mathbf Q_g$, the residue field of $w$ has $g$ elements and $g$ is not in the square of the
prime ideal of $w$. So $w$ has residue degree and ramification index $1$ over $\mathbf Q$, hence so
does $v$, and $F_v=E_w$; thus $w$ is split over $F$, and $q_w=g$.

The cusp forms in $\pi$ are invariant under an open compact subgroup of
$\operatorname{GL}_n(\mathbf A_{E,f})$, so $\pi$ has a spherical local component, hence Satake
parameters, at all but finitely many finite places \cite{Fla79}. Let $S$ be the finite set of the
others, and $f=n!$. Lemma~\ref{gal-depth-zero-lift-field}
gives an $S$-layered extension $X$ of $E$, a CM field and a finite abelian extension of $E$, such that
every element of $I_A$ fixes $X$ and the place $u$ of $X$ under $A$ has a residue field of $g^{n!}$
elements. Lemma~\ref{gal-depth-zero-lift-tower} gives $P$ and $\chi'$ with $(P,\chi')$ in scope, a
chain of weak base changes from $(E,\pi)$ to $(X,P)$, and $r|_{G_X}$ weakly attached to $P$. By
Lemma~\ref{gal-depth-zero-lift-infinity}, $Z_\tau(P)=Z_{\tau|_E}(\pi)$ for every $\tau$.

Since every element of $I_A$ fixes $X$, $I_A$ is contained in $\operatorname{Gal}(\Omega/X)$. The
inertia group of $A$ in $\operatorname{Gal}(\Omega/X)$ is defined by the same condition on the action
on $A$, so it is $I_A\cap\operatorname{Gal}(\Omega/X)=I_A$: the inertia groups of $A$ over $E$ and over
$X$ are the same group of automorphisms of $\Omega$.

Let $T$ be a local component of $P$ at $u$. Every member of the supercuspidal support of $\sigma$ has
degree at most $n$, so its degree divides $n!$. Lemma~\ref{gal-depth-zero-lift-support}, with
$f=n!$, gives smooth characters $\xi_1,\ldots,\xi_n$ of $X_u^\times$ such that $T$ has supercuspidal
support $\{\xi_1,\ldots,\xi_n\}$, this support is unlinked by the hypotheses on $\sigma$, and
$\operatorname{rc}(\sigma)(t)$ has characteristic polynomial $\prod_j(Y-\xi_j(a'(t)))$ for every
$t\in I_A$, where $a'$ is the local Artin map for $X$, $u$ and $A$
(Proposition~\ref{gal-local-artin-unique}). By Theorem~\ref{gal-llc-exists} there is a map
$\operatorname{rc}'$ for $X$, $u$ and $A$ that satisfies the laws, with respect to $a'$ by the
uniqueness of the local Artin map. By Lemma~\ref{gal-depth-zero-lift-characters}, $T^\flat$ is plain
and $\operatorname{rc}'(T^\flat)(t)$ has characteristic polynomial $\prod_j(Y-\xi_j(a'(t)))$ for every
$t\in I_A$.

Let $t\in I_A$ and let $a$ be the local Artin map for $E$, $w$ and $A$. By law~(5),
$\operatorname{rc}(\sigma^\flat)=\operatorname{rc}(\sigma)\otimes(|\cdot|^{(1-n)/2}\circ a)$, and
$|a(t)|=1$ since $d(t)=0$. So $\operatorname{rc}(\sigma^\flat)(t)$ has the characteristic polynomial
$\prod_j(Y-\xi_j(a'(t)))$ of $\operatorname{rc}(\sigma)(t)$, which is that of
$\operatorname{rc}'(T^\flat)(t)$.
\end{proof}

\begin{proposition}\label{gal-depth-zero-compatibility}
Suppose the hypotheses and conclusions of Proposition~\ref{gal-plain-compatibility}. Let $(\pi,\chi)$
be in scope, $r$ weakly attached to $\pi$, $\sigma$ a local component of $\pi$ at $w$ with nonzero
vectors fixed by $\ker(\operatorname{GL}_n(\mathcal O_w)\to\operatorname{GL}_n(k_w))$, and suppose that
either $\sigma^\flat$ is plain, or $E_w\cong\mathbf Q_g$ for a prime $g$ and no two members of degree at
least $2$ of the supercuspidal support of $\sigma^\flat$ are inertially equivalent. Let $A$ be a
valuation ring over $w$ and $\operatorname{rc}$ a map that satisfies the laws, with
$\operatorname{rc}(\sigma^\flat)$ monodromy-free. Then for every $t\in I_A$:
\begin{enumerate}
\item if $w\nmid\ell$, $\mathrm{WD}_\iota(r|_{W_A})$ is monodromy-free and its value at $t$ has the
  characteristic polynomial of $\operatorname{rc}(\sigma^\flat)(t)$;
\item if $w\mid\ell$, $r$ satisfies (dR) for $\pi$ at $w$, $r|_{D_A}$ is potentially semistable, and
  $\mathrm{WD}_\iota(r|_{D_A})$ is monodromy-free with the characteristic polynomial of
  $\operatorname{rc}(\sigma^\flat)(t)$ at $t$.
\end{enumerate}
\end{proposition}
\leanref{Development/AutomorphicArithmetic/Foundations/GaloisRepresentations/SphericalGaloisCompatibility/PAdicSplitting/LocalGlobalCompatAtIwahoriFixedPlaces.lean}{UpperTheory.Gal.ToralGalois.R2.LocalGlobalCompatAtIwahoriFixedPlaces}
\leanref{Development/AutomorphicArithmetic/Foundations/GaloisRepresentations/SphericalGaloisCompatibility/PAdicSplitting/LocalGlobalCompatAtIwahoriFixedPlaces.lean}{UpperTheory.Gal.ToralGalois.R2.Apart}
\leanref{Development/AutomorphicArithmetic/Foundations/GaloisRepresentations/SphericalGaloisCompatibility/PAdicSplitting/R2g.lean}{UpperTheory.Gal.ToralGalois.R2.r2g}
\leanref{Development/AutomorphicArithmetic/Foundations/GaloisRepresentations/SphericalGaloisCompatibility/PAdicSplitting/LocalGlobalCompatAtIwahoriFixedOfUnlinked.lean}{UpperTheory.Gal.ToralGalois.R2.localGlobalCompatAtIwahoriFixed_of_unlinked}

\begin{proof}
If $\sigma^\flat$ is plain, this is Proposition~\ref{gal-plain-compatibility} restricted to inertia.
Otherwise $E_w\cong\mathbf Q_g$, and since $\sigma$ has nonzero vectors fixed by the first congruence
subgroup and $\operatorname{rc}(\sigma^\flat)$ is monodromy-free, the representation $\sigma$ is
\emph{mixed}: it has depth zero, and its supercuspidal support is unlinked with pairwise inertially
inequivalent members of degree at least $2$. Indeed, $\sigma^\flat$ is generic and
Lemma~\ref{gal-generic-monodromy-free}(1) gives the unlinked support, while the inertial inequivalence
is the hypothesis.

\emph{The lift.} By Lemma~\ref{gal-depth-zero-lift}, applied to $\sigma$, there are an abelian CM
extension $X$ of $E$ inside $\Omega$, a place $u$ of $X$ under
$A$ such that every element of $I_A$ lies in the inertia group of $A$ over $X$, a cuspidal automorphic
representation $P$ of $\operatorname{GL}_n(\mathbf A_X)$ in scope with $Z_\tau(P)=Z_{\tau|_E}(\pi)$ and
with $r|_{G_X}$ weakly attached to $P$, a local component $T$ of $P$ at $u$ with $T^\flat$ plain, and a map
$\operatorname{rc}'$ over $X$ satisfying the laws (Theorem~\ref{gal-llc-exists}) such that
$\operatorname{rc}'(T^\flat)(t)$ and $\operatorname{rc}(\sigma^\flat)(t)$ have the same characteristic
polynomial for every $t\in I_A$. In particular $X/E$ is unramified at the place under $A$, and the
local base change from $E_w$ to $X_u$ turns the depth-zero supercuspidals in the support of
$\sigma^\flat$ into characters.
% Lean: ToTor (Lemma gal-depth-zero-lift), proved by toTor_holds. MixedOf (mixedOf) is the first
% paragraph of the proof and DRDesc (drDesc) the descent of (dR) in (2) below.

(1) If $w\nmid\ell$, Proposition~\ref{gal-plain-compatibility}(1) for $P$ at $u$ gives that the
Weil--Deligne representation of $r|_{G_X}$ at $u$ is $\operatorname{rc}'(T^\flat)$, monodromy-free. It is
the restriction of $\mathrm{WD}_\iota(r|_{W_A})$ to the Weil group of $X$ at $A$, which contains $I_A$.
So $\mathrm{WD}_\iota(r|_{W_A})$ is monodromy-free, and at $t\in I_A$ its characteristic polynomial is
that of $\operatorname{rc}'(T^\flat)(t)$, hence that of $\operatorname{rc}(\sigma^\flat)(t)$.

(2) If $w\mid\ell$, Proposition~\ref{gal-plain-compatibility}(2) for $P$ at $u$ gives (dR) for $P$ at $u$,
potential semistability of $r|_{G_X}$ at $u$, and the characteristic polynomials. Potential
semistability and the inertial characteristic polynomials of $\mathrm{WD}_\iota(r|_{D_A})$ follow by
restriction as in (1). The de Rham property and the Hodge--Tate weights descend from $X$ to $E$: if for
every $A$ over $w$ there are such $X$, $P$ and $u$ with (dR) for $P$ at $u$ and
$Z_\tau(P)=Z_{\tau|_E}(\pi)$, then $r$ satisfies (dR) for $\pi$ at $w$, since $r|_{D_A}$ is de Rham if
its restriction to an open subgroup is, with the same Hodge--Tate weights.
\end{proof}

\begin{proposition}\label{gal-row-from-core}
Suppose existence and the four conditions hold in the first scope. Then for every CM field $E$,
algebraic closure $\Omega$, prime $\ell$, isomorphism $\iota$, $a\geq1$ and $(\pi,\chi)$ in scope,
there is a triple $(r,m,\mathrm{HT})$, with $r$ weakly attached to $\pi$, that satisfies
Definition~\ref{gal-attached} relative to the maps satisfying the laws, with clause~(3) of
Definition~\ref{gal-compatible} required whenever $E_w\cong\mathbf Q_g$ for a prime $g$, without the
bound $g>16n^2$.
\end{proposition}
\leanref{Development/AutomorphicArithmetic/Foundations/GaloisRepresentations/TypeConditions/WeightsConjAlt.lean}{UpperTheory.GalStar.Away}
\leanref{Development/AutomorphicArithmetic/Foundations/GaloisRepresentations/TypeConditions/WeightsConjAlt.lean}{UpperTheory.GalStar.Props}
\leanref{Development/AutomorphicArithmetic/Foundations/GaloisRepresentations/TypeConditions/WeightsConjAlt.lean}{UpperTheory.GalStar.Row0}
\leanref{Development/AutomorphicArithmetic/Foundations/GaloisRepresentations/IntegralStructures/InertCharpoly/GaloisLocalGlobalOfLocalGlobalCompatibility.lean}{UpperTheory.Gal.Int.Row.galoisLocalGlobal_of_placeClassCompatibility}
\leanref{Development/AutomorphicArithmetic/Foundations/GaloisRepresentations/IntegralStructures/InertCharpoly/GaloisLocalGlobalOfLocalGlobalCompatibility.lean}{UpperTheory.Gal.Int.Row.galoisLocalGlobal_of_localGlobalCompatibility}
\leanref{Development/AutomorphicArithmetic/Foundations/GaloisRepresentations/IntegralStructures/InertCharpoly/GaloisLocalGlobalOfLocalGlobalCompatibility.lean}{UpperTheory.Gal.Int.Row.away}
\leanref{Development/AutomorphicArithmetic/Foundations/GaloisRepresentations/IntegralStructures/InertCharpoly/GaloisLocalGlobalOfLocalGlobalCompatibility.lean}{UpperTheory.Gal.Int.Row.above}
\leanref{Development/AutomorphicArithmetic/Foundations/GaloisRepresentations/IntegralStructures/InertCharpoly/GaloisLocalGlobalOfLocalGlobalCompatibility.lean}{UpperTheory.Gal.Int.Row.exists_weights}
\leanref{Development/AutomorphicArithmetic/Foundations/GaloisRepresentations/IntegralStructures/Development.lean}{UpperTheory.Gal.Int.rest}

\begin{proof}
Existence in the first scope gives $r$ weakly attached to $\pi$, continuous and semisimple, which is
clause~(1). Since $(\pi,\chi)$ is polarized and $\pi$ is regular algebraic, there is an integer $m$
with $Z_{\tau\circ c}(\pi)=\{a-1-m-z:z\in Z_\tau(\pi)\}$ for every $\tau$, and the elements of
$\frac{a-1}{2}-Z_\tau(\pi)$ are integers; put $\mathrm{HT}_\tau=\{\frac{a-1}{2}-z:z\in
Z_{\iota\circ\tau}(\pi)\}$, which gives clauses~(3) and~(4).

\emph{Away from $\ell$.} Let $w\nmid\ell$, $A$ over $w$, $\operatorname{rc}$ a map satisfying the laws,
$\sigma$ a local component of $\pi$ at $w$ and $y=\operatorname{rc}(\sigma^\flat)$. The dimension of
$\mathrm{WD}_\iota(r|_{W_A})$ is $a$, which is that of $y$ by law~(1). The representation $\sigma^\flat$
is generic, since local components of cuspidal representations of $\operatorname{GL}_a$ are generic and
twisting preserves genericity. If $y$ is abelian and monodromy-free, Lemma~\ref{gal-generic-monodromy-free}(2)
shows that $\sigma^\flat$ is plain, and Proposition~\ref{gal-plain-compatibility}(1) gives
$\mathrm{WD}_\iota(r|_{W_A})=y$. If $E_w\cong\mathbf Q_g$ and $y$ is tame, monodromy-free and has
multiplicity-free tame part, Lemma~\ref{gal-generic-monodromy-free}(3) shows that $\sigma^\flat$ has
depth zero, hence so does $\sigma=\sigma^\flat\nu^{(a-1)/2}$, since unramified twists preserve depth
zero; so $\sigma$ has nonzero vectors fixed by
$\ker(\operatorname{GL}_a(\mathcal O_w)\to\operatorname{GL}_a(k_w))$, and that the members of degree at
least $2$ of the support of $\sigma^\flat$ are pairwise inertially inequivalent. These are the
hypotheses of Proposition~\ref{gal-depth-zero-compatibility}, whose item (1) gives that
$\mathrm{WD}_\iota(r|_{W_A})$ is monodromy-free with the characteristic polynomials of $y$ at elements
of $I_A$, which are the elements of degree $0$. This is clause~(2) of Definition~\ref{gal-attached}.

\emph{Above $\ell$.} Let $w\mid\ell$. If $\sigma$ is spherical, $\pi$ has Satake parameters at $w$, and
(crys) and (dR) in the first scope give clause~(5)(a). For (5)(b), let $\operatorname{rc}$ satisfy the
laws with $y=\operatorname{rc}(\sigma^\flat)$ monodromy-free and either abelian, or tame with
multiplicity-free tame part where $E_w\cong\mathbf Q_g$. In the abelian case $\sigma^\flat$ is plain as
above, and Proposition~\ref{gal-plain-compatibility}(2) gives (dR), potential semistability, and a
monodromy-free $\mathrm{WD}_\iota(r|_{D_A})$ with the characteristic polynomials of $y$; two
monodromy-free Weil--Deligne representations with the same characteristic polynomials at every element
are equal, so $\mathrm{WD}_\iota(r|_{D_A})=y$, which gives compatibility. In the tame case,
Proposition~\ref{gal-depth-zero-compatibility}(2) gives (dR), potential semistability, and the
characteristic polynomials at the elements of $I_A$; the dimension follows from the characteristic
polynomial at $1$. In both cases the Hodge--Tate multisets are the $\mathrm{HT}_\tau$ by (dR).
\end{proof}

\begin{proposition}\label{gal-core-all}
Suppose existence and the four conditions hold in the self-dual scope (Definition~\ref{gal-scope}).
Then they hold in the first scope.
\end{proposition}
\leanref{Development/AbelianVarietyTheory/Proofs/SatakeParameters/CoreAllOfCore.lean}{GlobalTheory.RoadU.coreAll_of_core}
\leanref{Development/AutomorphicArithmetic/Foundations/GaloisRepresentations/IntegralStructures/CrystallineInputs/CoreAllOfRowStarCore.lean}{UpperTheory.Gal.Int.Inp.coreAll_of_core}
\leanref{Development/AbelianVarietyTheory/Definitions/GaloisGeometry/CyclicStrongAtCusp.lean}{GlobalTheory.RoadU.Inputs}

\begin{proof}
The proof passes through the good scope, and uses the published results listed at the end.

\emph{Existence in the good scope.} Let $(\pi,\chi)$ be in good scope for $\ell$ and $w_0$ a place over
$\ell$, split over $F$. There is an algebraic Hecke character $\kappa$ of $E$ of some infinity type $m$
such that $(\pi\otimes\kappa,\eta^a)$ is in self-dual scope; existence there gives $r'$ weakly attached
to $\pi\otimes\kappa$, and $r'\otimes\mu$, with $\mu$ the $\ell$-adic character attached to
$\kappa^{-1}$ by global class field theory, is continuous and weakly attached to $\pi$.

\emph{(unr) and (crys) in the good scope.} The same twist carries these conditions at split places from
$\pi\otimes\kappa$ to $\pi$, since $\mu$ is unramified and crystalline at the relevant places and its
Frobenius values are those of $\kappa^{-1}$. For (dR) and (phi) the same argument applies, with (crys)
supplying the crystalline twist.

\emph{Existence in the first scope.} Let $(\pi,\chi)$ be in scope. The set $S$ of quadratic CM
extensions $Q$ of $E$ for which the quadratic base change $P_Q$ of $\pi$ to $Q$ is cuspidal and in good
scope is infinite: quadratic base change exists by Theorem~\ref{gal-cyclic-base-change}, it is cuspidal
by the theorem of Jacquet and Shalika, the polarization lifts, and infinitely many $Q$ are split at the
places that the good scope constrains. Existence in the good scope gives $r_Q$ weakly attached to $P_Q$
for each $Q\in S$. For $Q,Q'\in S$, the restrictions of $r_Q$ and $r_{Q'}$ to $G_{QQ'}$ have the same
Frobenius eigenvalues at almost all places, so they are isomorphic by the Chebotarev density theorem and
the Brauer--Nesbitt theorem, and these isomorphisms patch: there is a continuous semisimple
$\rho:G_E\to\operatorname{GL}_a(\bar{\mathbf Q}_\ell)$ with $\rho|_{G_Q}\cong r_Q$ for all but finitely
many $Q\in S$. At almost every place $w$ of $E$, some such $Q$ is split at $w$, so $\rho$ has the
Frobenius eigenvalues at $w$ given by (unr) for $P_Q$; hence $\rho$ is weakly attached to $\pi$.

\emph{The conditions in the first scope.} The four conditions descend in the same way. Let $w$ be a
place where $\pi$ has Satake parameters $B$, $r$ weakly attached to $\pi$, and $A$ a valuation ring over
$w$. Choose a quadratic CM extension $Q$ of $E$ as in the previous step for which moreover
$D_A\subset G_Q$ and the places of $Q$ over $w$ are split over the maximal totally real subfield of
$Q$. Let $P_Q$ be the quadratic base change of $\pi$, polarized and in good scope as before. The
restriction $r|_{G_Q}$ is weakly attached to $P_Q$: by Theorem~\ref{gal-cyclic-base-change}(1), at a
place of $Q$ of relative degree $f$ the Satake parameters of $P_Q$ are the $f$-th powers of those of
$\pi$, and so are the Frobenius eigenvalues of $r|_{G_Q}$. Since $D_A\subset G_Q$, the place $w'$ of
$Q$ under $A$ has the same norm as $w$, and $P_Q$ has Satake parameters $B$ at $w'$. The condition in
the good scope for $P_Q$, $r|_{G_Q}$ and $w'$ is then the condition for $\pi$, $r$ and $w$, because
$D_A$, $I_A$ and $W_A$ are the same subgroups of $G_E$ whether computed over $E$ or over $Q$.

The proof uses the following published results: strong
multiplicity one for $\operatorname{GL}_a$ in the form of Jacquet and Shalika \cite{JS81a}, the theorem
of Gelfand and Kazhdan on contragredients, the classification of irreducible admissible
representations of $\operatorname{GL}_a$ \cite{BZ77}, \cite{Zel80}, the generic unitary dual, cyclic
base change of Arthur and Clozel with its local properties \cite{AC89}, the $\ell$-adic characters
attached to algebraic Hecke characters, and the theorem that de Rham representations are potentially
semistable \cite{Ber02}.
\end{proof}
\leanref{Development/AbelianVarietyTheory/Proofs/SatakeParameters/ImportsAlt.lean}{GlobalTheory.UnrDescent.main}
\leanref{Development/AbelianVarietyTheory/Proofs/SatakeParameters/HoldsSDRAllOfGood.lean}{GlobalTheory.RoadU.holdsS_dR_all_of_good}
\leanref{Development/AbelianVarietyTheory/Proofs/SatakeParameters/CrystallineAllOfGood.lean}{GlobalTheory.RoadU.holdsS_all_crys}

\subsection{Lacunary hulls and division of representations}\label{subsection-gal-lacunary-division}

\begin{situation}\label{gal-definite-situation}
Let $E$ be a CM field with maximal totally real subfield $F$, $\Omega$ an algebraic closure of $E$,
$\ell$ a prime and $\iota:\bar{\mathbf Q}_\ell\to\mathbf C$ a field isomorphism. Let $a\geq1$ be odd and
$d\geq1$, and suppose that $-d$ is a square in $E$ and in $\mathbf Q_\ell$, that $[F:\mathbf Q]$ is
even, and that every place of $E$ not split over $F$ is prime to $\ell$ and unramified over $F$. Let $T$
be a finite set of places of $E$ not split over $F$, and $U$ an automorphic representation of $U_a$
(Definition~\ref{gal-definite-unitary}) unramified at every place not split over $F$ outside $T$. Let
$w_0$ and $v'$ be places of $E$ over $\ell$, split over $F$, that lie over different places of $F$.
\end{situation}
\leanref{Development/AbelianVarietyTheory/Definitions/GaloisGeometry/LatticeDataThicknessBound.lean}{GlobalTheory.CongL.AutOffT}
\leanref{Development/AbelianVarietyTheory/Definitions/GaloisGeometry/ApartAllEmpty.lean}{GlobalTheory.CongL.UsPl}

\begin{definition}\label{gal-generic-parameters}
Let $a\geq1$ and let $B$ be a multiset of $a$ nonzero complex numbers.
\begin{enumerate}
\item $B$ is \emph{multiplicatively generic} if every multiset of $a$ elements of $B$ (with
  repetitions allowed) whose product equals the product of $B$ is $B$ itself.
\item $B$ is \emph{generic} if for some, equivalently every, enumeration $\alpha_1,\ldots,\alpha_a$ of
  $B$: $\alpha_i^c\neq\alpha_j^c$ for $i\neq j$ and $1\leq c\leq a$; $\alpha_i\alpha_k\neq\alpha_j^2$
  for $j\notin\{i,k\}$; and $\prod_i\alpha_i^{m_i}\neq\prod_i\alpha_i$ for every $m\in\{0,\ldots,a\}^a$
  with $\sum_im_i=a$ and $m\neq(1,\ldots,1)$. Equivalently, $D_a(\alpha)\neq0$ for the integer
  polynomial $D_a$ that is the product of the differences in these conditions.
\end{enumerate}
A generic multiset is multiplicatively generic, and so is $cB=\{cb:b\in B\}$ for every $c\neq0$.
\end{definition}
\leanref{Development/AbelianVarietyTheory/Definitions/GaloisGeometry/LimWPhi.lean}{GlobalTheory.CongL.DGen}
\leanref{Development/AbelianVarietyTheory/Definitions/Divisors/GenOf.lean}{GlobalTheory.CongW.GenOf}
\leanref{Development/AutomorphicArithmetic/Foundations/GaloisRepresentations/AnnihilatorModels/GenericMulOfProd.lean}{UpperTheory.Gal.crystallineComparison.GenericMul}
\leanref{Development/AutomorphicArithmetic/Foundations/GaloisRepresentations/IntegralStructures/ZetaCharacterDivisibility/PhiOfDivOf.lean}{UpperTheory.Gal.Int.Div.DGenMul}

\begin{lemma}\label{gal-char-of-ann}
Let $k$ be a field, $a\geq0$, $M$ an $a\times a$ matrix over $k$ and $R$ a multiset of $a$ elements of $k$
such that every multiset of $a$ elements of $R$ (with repetitions allowed) whose product equals the product of $R$ is
$R$ itself; for $a\geq1$ and nonzero elements this says that $R$ is multiplicatively generic
(Definition~\ref{gal-generic-parameters}(1), with $k$ in place of $\mathbf C$). If $\prod_{z\in R}(M-z)=0$ and
$\det M=\prod_{z\in R}z$, then the characteristic polynomial of $M$ is $\prod_{z\in R}(X-z)$.
\end{lemma}
\leanref{Development/AutomorphicArithmetic/Foundations/GaloisRepresentations/AnnihilatorModels/GenericMulOfProd.lean}{UpperTheory.Gal.crystallineComparison.CharOfAnn}

\begin{proof}
Let $L$ be an algebraic closure of $k$. The eigenvalues of $M$ in $L$ are roots of $\prod_{z\in R}(X-z)$, so the
multiset $R'$ of eigenvalues of $M$ in $L$, counted with multiplicity, consists of $a$ elements of $R$. Its product is
$\det M=\prod_{z\in R}z$, so $R'=R$ by the hypothesis on $R$. The characteristic polynomial of $M$ is
$\prod_{z\in R'}(X-z)$ over $L$, which is $\prod_{z\in R}(X-z)$.
% Lean: crystallineComparison.charAnn_of, over L = AlgebraicClosure K: hmem (the roots of the characteristic polynomial
% of M lie in R, through the spectrum), t = the roots pulled back to K, htcard, htsub and htprod, then GenericMul.
\end{proof}
\leanref{Development/AutomorphicArithmetic/Foundations/GaloisRepresentations/AnnihilatorModels/PCharAnn.lean}{UpperTheory.Gal.crystallineComparison.charAnn_of}

\begin{theorem}\label{gal-sen}
Let $E$ be a number field, $\Omega$ an algebraic closure, $\ell$ a prime and
$V:G_E\to\operatorname{GL}_N(\bar{\mathbf Q}_\ell)$ continuous. Suppose that at some place $w$ over
$\ell$, some valuation ring $A$ over $w$ and some embedding $\tau$, $V|_{D_A}$ is Hodge--Tate with
Hodge--Tate multiset at $\tau$ containing a lacunary set of $n$ integers (Definition~\ref{gal-seed}).
Then there is an element $X$ of the Lie algebra of the Zariski closure of $V(G_E)$, over an algebraically
closed extension of $\bar{\mathbf Q}_\ell$, with $n$ distinct eigenvalues forming a lacunary set. Indeed,
by Sen's theorem \cite{Sen73} the Sen operator at $\tau$ lies in the Lie algebra of the Zariski closure
of $V(I_A)$, and its eigenvalues are the Hodge--Tate weights at $\tau$.
\end{theorem}
\leanref{Development/WeilBound/CoreStatements/PadicHodgeComparison/CocharacterLieDivisors/StGalW.lean}{PositivityTheory.AutomorphicGaloisHodgeTate.Div.SenWLac}

\begin{definition}\label{gal-lacunary-hull}
Let $K$ be a field and $N,n\geq1$. Write $K[x]$ for the polynomial ring over $K$ in the entries
$x_{ab}$ of an $N\times N$ matrix $x$, and for $f\in K[x]$ and $x,Y\in M_N(K)$ write
$df_x(Y)=\sum_{a,b}\frac{\partial f}{\partial x_{ab}}(x)\,Y_{ab}$.
\begin{enumerate}
\item Let $\Sigma\subset M_N(K)$. Then $I(\Sigma)\subset K[x]$ is the ideal of polynomials that vanish
  on $\Sigma$, the \emph{Zariski closure} $\bar\Sigma$ is the set of $y\in M_N(K)$ at which every
  $f\in I(\Sigma)$ vanishes, and the \emph{tangent space} $\mathfrak t(\Sigma)$ is the subspace of those
  $Y\in M_N(K)$ with $df_1(Y)=0$ for every $f\in I(\Sigma)$.
\item Let $\Gamma$ be a group and $\rho:\Gamma\to\operatorname{GL}_N(K)$ a homomorphism. The \emph{Lie
  algebra} of $\rho$ is $\mathfrak g_\rho=\mathfrak t(\rho(\Gamma))$, and its \emph{envelope}
  $\mathcal A_\rho$ is the $K$-subalgebra of $M_N(K)$ generated by $\mathfrak g_\rho$. For a field $K'$
  containing $K$, an element of the Lie algebra of $\rho$ \emph{after extension to} $K'$ is an
  $X\in M_N(K')$ with $df_1(X)=0$ for every $f\in I(\rho(\Gamma))$. The Lie algebra of the Zariski
  closure of $\rho(\Gamma)$ in Theorem~\ref{gal-sen} and Proposition~\ref{gal-division} is meant in
  this sense.
\item An element $x\in M_N(K)$ has \emph{degree at most} $n$ if $q(x)=0$ for some nonzero $q\in K[X]$
  of degree at most $n$, that is, if $1,x,\ldots,x^n$ are linearly dependent.
\item A family $e_1,\ldots,e_n$ of elements of $K$ is \emph{lacunary} if
  $|B|\sum_{i\in A}e_i\neq|A|\sum_{i\in B}e_i$ for any two disjoint nonempty subsets $A,B$ of
  $\{1,\ldots,n\}$. A lacunary family is injective (take $A=\{i\}$ and $B=\{j\}$).
\item A \emph{system of idempotents} is a family $p_1,\ldots,p_n$ of nonzero elements of $M_N(K)$ with
  $p_i^2=p_i$, $p_ip_j=0$ for $i\neq j$, and $\sum_ip_i=1$. For $O\subset\{1,\ldots,n\}$, $V_O\subset K^N$
  is the image of $\sum_{k\in O}p_k$, the set of $v$ with $p_kv=0$ for all $k\notin O$.
\item Let $p_1,\ldots,p_n$ be a system of idempotents, $\pi$ a permutation of $\{1,\ldots,n\}$ and
  $y\in M_N(K)$. Then $y$ \emph{separates the cycles} of $\pi$ if for any two different cycles $O\neq O'$
  of $\pi$ no $\lambda\in K$ is an eigenvalue of $y$ with an eigenvector in $V_O$ and an eigenvector in
  $V_{O'}$.
\item Let $\Lambda\subset K$ be finite. A family $\tau_1,\ldots,\tau_n$ of nonzero elements of $K$ is
  \emph{multiplicatively lacunary away from} $\Lambda$ if
  $\bigl(\prod_{i\in O}\tau_i\bigr)^{|O'|}\neq b\bigl(\prod_{i\in O'}\tau_i\bigr)^{|O|}$ for any two
  disjoint nonempty subsets $O,O'$ of $\{1,\ldots,n\}$ and every $b\in\Lambda$.
\item Let $\mathcal A$ be a $K$-subalgebra of $M_N(K)$. A system of idempotents is a \emph{corner
  system} for $\mathcal A$ if every $p_i$ lies in $\mathcal A$ and $p_iap_i\in Kp_i$ for all
  $a\in\mathcal A$ and all $i$. The algebra $\mathcal A$ is \emph{semisimple} if the only $Y\in\mathcal A$
  such that $aY$ is nilpotent for every $a\in\mathcal A$ is $Y=0$; this is the vanishing of its Jacobson
  radical. It is \emph{$\rho$-stable} if $\rho(g)a\rho(g)^{-1}\in\mathcal A$ for all $g\in\Gamma$ and
  $a\in\mathcal A$, and every subspace of $K^N$ stable under $\rho(\Gamma)$ is stable under $\mathcal A$.
\item Let $p_1,\ldots,p_n$ be a corner system for $\mathcal A$. A family of \emph{matrix units} for it
  in $\mathcal A$ is a family of elements $E_{ij}\in\mathcal A$, for $1\leq i,j\leq n$, such that
  $E_{ii}=p_i$, $p_iE_{ij}p_j=E_{ij}$, $E_{ji}\neq0$ whenever $E_{ij}\neq0$, $E_{ij}E_{jk}=E_{ik}\neq0$
  whenever $E_{ij}\neq0$ and $E_{jk}\neq0$, and $p_iap_j\in KE_{ij}$ for all $a\in\mathcal A$ and all
  $i,j$.
\end{enumerate}
\end{definition}
\leanref{Development/EtaleCohomology/DivisorCycles/OpenCoset.lean}{LateralTheory.DivDb.Van}
\leanref{Development/EtaleCohomology/DivisorCycles/OpenCoset.lean}{LateralTheory.DivDb.Zcl}
\leanref{Development/EtaleCohomology/DivisorCycles/OpenCoset.lean}{LateralTheory.DivDb.Tan}
\leanref{Development/EtaleCohomology/DivisorCycles/OpenCoset.lean}{LateralTheory.DivDb.img}
\leanref{Development/EtaleCohomology/DivisorCycles/OpenCoset.lean}{LateralTheory.DivDb.lie}
\leanref{Development/EtaleCohomology/DivisorCycles/OpenCoset.lean}{LateralTheory.DivDb.env}
\leanref{Development/WeilBound/CoreStatements/PadicHodgeComparison/CocharacterLieDivisors/MeetsOpenCosets.lean}{PositivityTheory.AutomorphicGaloisHodgeTate.Div.InLieZar}
\leanref{Development/EtaleCohomology/DivisorCycles/OpenCoset.lean}{LateralTheory.DivDb.Deg}
\leanref{Development/WeilBound/CoreStatements/PadicHodgeComparison/CocharacterLieDivisors/DivWCoreLacUseType.lean}{PositivityTheory.AutomorphicGaloisHodgeTate.Div.IsLac}
\leanref{Development/EtaleCohomology/DivisorCycles/OpenCoset.lean}{LateralTheory.DivDb.isLac_inj}
\leanref{Development/EtaleCohomology/DivisorCycles/OpenCoset.lean}{LateralTheory.DivDb.Sys}
\leanref{Development/EtaleCohomology/DivisorCycles/OpenCoset.lean}{LateralTheory.DivDb.Disj}
\leanref{Development/EtaleCohomology/DivisorCycles/OpenCoset.lean}{LateralTheory.DivDb.MulLac}
\leanref{Development/EtaleCohomology/DivisorCycles/OpenCoset.lean}{LateralTheory.DivDb.Corner}
\leanref{Development/EtaleCohomology/DivisorCycles/OpenCoset.lean}{LateralTheory.DivDb.NoNil}
\leanref{Development/EtaleCohomology/DivisorCycles/OpenCoset.lean}{LateralTheory.DivDb.Normd}
\leanref{Development/EtaleCohomology/DivisorCycles/OpenCoset.lean}{LateralTheory.DivDb.Keeps}
\leanref{Development/EtaleCohomology/DivisorCycles/OpenCoset.lean}{LateralTheory.DivDb.MUnits}

\begin{lemma}\label{gal-lacunary-specialize}
Let $K$ be an algebraically closed field, $R$ an integral domain containing $K$, $m\geq1$, $x\in R^m$,
and $G$ a finite set of polynomials in $K[y_1,\ldots,y_m]$ with $g(x)\neq0$ for every $g\in G$. Then
there is $y\in K^m$ such that $f(y)=0$ for every $f\in K[y_1,\ldots,y_m]$ with $f(x)=0$, and
$g(y)\neq0$ for every $g\in G$.
\end{lemma}
\leanref{Development/EtaleCohomology/DivisorCycles/OpenCoset.lean}{LateralTheory.DivDb.SpecPt}
\leanref{Development/EtaleCohomology/DivisorCycles/SpecPtPf.lean}{LateralTheory.DivDb.SpecPtPf.SpecPt_of}

\begin{proof}
The kernel $\mathfrak p$ of evaluation at $x$ is a prime ideal of $K[y_1,\ldots,y_m]$, and the product
$h$ of the elements of $G$ does not lie in $\mathfrak p$. By Hilbert's Nullstellensatz \cite{Eis95},
$\mathfrak p$ is the ideal of polynomials that vanish on the zero set of $\mathfrak p$ in $K^m$. So some
zero $y$ of $\mathfrak p$ has $h(y)\neq0$.
\end{proof}

\begin{lemma}\label{gal-lacunary-descent}
Let $\Gamma$ be a group, $K$ an algebraically closed field, $K'$ a field containing $K$, $N,n\geq1$, and
$\rho:\Gamma\to\operatorname{GL}_N(K)$ a homomorphism. Let $X$ be an element of the Lie algebra of $\rho$
after extension to $K'$ (Definition~\ref{gal-lacunary-hull}), and let $e_1,\ldots,e_n\in K'$ be a
lacunary family of eigenvalues of $X$. Then there are $X_0\in\mathfrak g_\rho$ and a lacunary family
$e^0_1,\ldots,e^0_n\in K$ of eigenvalues of $X_0$.
\end{lemma}
\leanref{Development/EtaleCohomology/DivisorCycles/OpenCoset.lean}{LateralTheory.DivDb.SpecLie}
\leanref{Development/EtaleCohomology/DivisorCycles/SpecLiePf.lean}{LateralTheory.DivDb.SpecLiePf.SpecLie_of}

\begin{proof}
Apply Lemma~\ref{gal-lacunary-specialize} with $R=K'$ to the point $(X,e)$, in the variables $x_{ab}$ and
$y_1,\ldots,y_n$, with $G$ the set of polynomials $|B|\sum_{i\in A}y_i-|A|\sum_{i\in B}y_i$ for disjoint
nonempty $A,B$. The point $(X,e)$ is a zero of the linear polynomials
$\sum_{a,b}\frac{\partial f}{\partial x_{ab}}(1)x_{ab}$ for $f\in I(\rho(\Gamma))$ and of the polynomials
$\det(y_i-x)$, and no element of $G$ vanishes at it because $e$ is lacunary. The point $(X_0,e^0)$ given
by the lemma is a zero of the same polynomials, so $X_0\in\mathfrak g_\rho$ and each $e^0_i$ is an
eigenvalue of $X_0$, and no element of $G$ vanishes at it, so $e^0$ is lacunary.
\end{proof}

\begin{lemma}\label{gal-lacunary-closure}
Let $K$ be a field, $N,n\geq1$, and $\Sigma\subset M_N(K)$ a submonoid (Definition~\ref{gal-lacunary-hull}).
\begin{enumerate}
\item $\bar\Sigma$ is closed under products.
\item For $Y\in M_N(K)$ let $D_Y$ be the derivation of $K[x]$ with $D_Y(x_{ab})=\sum_kY_{ak}x_{kb}$, so
  that $(D_Yf)(x)=df_x(Yx)$. Then $Y\in\mathfrak t(\Sigma)$ if and only if $D_Yf\in I(\Sigma)$ for every
  $f\in I(\Sigma)$.
\item $\mathfrak t(\Sigma)$ is closed under $(Y,Z)\mapsto YZ-ZY$.
\item There is a set of polynomials in $K[x]$ whose common zeros in $M_N(F)$, for every field $F$
  containing $K$, are the elements of degree at most $n$.
\end{enumerate}
Suppose moreover that $K$ has characteristic $0$, and write $\exp(tY)=\sum_{k\geq0}t^kY^k/k!$ in
$M_N(K[[t]])$.
\begin{enumerate}
\setcounter{enumi}{4}
\item If $Y\in\mathfrak t(\Sigma)$ and $f\in I(\Sigma)$, then $f(\exp(tY))=0$ in $K[[t]]$.
\item If $Y\in\mathfrak t(\Sigma)$ and $Y^2=0$, then $1+Y\in\bar\Sigma$.
\item If every element of $\Sigma$ has degree at most $n$, then so does every element of $\bar\Sigma$
  and of $\mathfrak t(\Sigma)$.
\end{enumerate}
\end{lemma}
\leanref{Development/EtaleCohomology/DivisorCycles/OpenCoset.lean}{LateralTheory.DivDb.Chain}
\leanref{Development/EtaleCohomology/DivisorCycles/ChainPf.lean}{LateralTheory.DivDb.Chain_of}
\leanref{Development/EtaleCohomology/DivisorCycles/OpenCoset.lean}{LateralTheory.DivDb.ZclMul}
\leanref{Development/EtaleCohomology/DivisorCycles/ZclMulPf.lean}{LateralTheory.DivDb.ZclMul_of}
\leanref{Development/EtaleCohomology/DivisorCycles/OpenCoset.lean}{LateralTheory.DivDb.DX}
\leanref{Development/EtaleCohomology/DivisorCycles/OpenCoset.lean}{LateralTheory.DivDb.TanField}
\leanref{Development/EtaleCohomology/DivisorCycles/TangentFieldExtension.lean}{LateralTheory.DivDb.TangentFieldExtension.TanField_of}
\leanref{Development/EtaleCohomology/DivisorCycles/OpenCoset.lean}{LateralTheory.DivDb.TanLie}
\leanref{Development/EtaleCohomology/DivisorCycles/TanLiePf.lean}{LateralTheory.DivDb.TanLie_of}
\leanref{Development/EtaleCohomology/DivisorCycles/OpenCoset.lean}{LateralTheory.DivDb.DegPoly}
\leanref{Development/EtaleCohomology/DivisorCycles/DegPoly2.lean}{LateralTheory.DivDb.DegPoly₂.DegPoly_of}
\leanref{Development/EtaleCohomology/DivisorCycles/OpenCoset.lean}{LateralTheory.DivDb.ExpVan}
\leanref{Development/EtaleCohomology/DivisorCycles/ExpVanPf.lean}{LateralTheory.DivDb.ExpVan_of}
\leanref{Development/EtaleCohomology/DivisorCycles/OpenCoset.lean}{LateralTheory.DivDb.Unip}
\leanref{Development/EtaleCohomology/DivisorCycles/UnipotentPartTrivial.lean}{LateralTheory.DivDb.UnipotentPartTrivial.Unip_of}
\leanref{Development/EtaleCohomology/DivisorCycles/OpenCoset.lean}{LateralTheory.DivDb.TanDeg}
\leanref{Development/EtaleCohomology/DivisorCycles/TangentDegree.lean}{LateralTheory.DivDb.TangentDegree.TanDeg_of}

\begin{proof}
Throughout we use the chain rule: for a linear map $\varphi:M_N(K)\to M_{N'}(K)$ and $f\in K[x']$, the
function $f\circ\varphi$ is a polynomial with $d(f\circ\varphi)_x(Y)=df_{\varphi(x)}(\varphi(Y))$.

(1) Let $s\in\Sigma$ and $f\in I(\Sigma)$. The polynomial $x\mapsto f(sx)$ vanishes on $\Sigma$, so
$f(sy)=0$ for $y\in\bar\Sigma$. Hence for $y\in\bar\Sigma$ the polynomial $x\mapsto f(xy)$ vanishes on
$\Sigma$, and so on $\bar\Sigma$.

(2) Let $f\in I(\Sigma)$ and $s\in\Sigma$. The polynomial $f_s(x)=f(xs)$ lies in $I(\Sigma)$, and by the
chain rule $d(f_s)_1(Y)=df_s(Ys)=(D_Yf)(s)$. So if $Y\in\mathfrak t(\Sigma)$, then $D_Yf$ vanishes on
$\Sigma$. Conversely $(D_Yf)(1)=df_1(Y)$ and $1\in\Sigma$.

(3) On the generators, $D_YD_Z(x_{ab})=(ZYx)_{ab}$, so the commutator of $D_Y$ and $D_Z$ is
$D_{ZY-YZ}$. If $Y,Z\in\mathfrak t(\Sigma)$, it preserves $I(\Sigma)$ by (2), so
$ZY-YZ\in\mathfrak t(\Sigma)$ by (2), and so is its negative.

(4) The elements of degree at most $n$ are those $x$ for which the $N^2\times(n+1)$ matrix with columns
$1,x,\ldots,x^n$ has rank at most $n$. Its minors of size $n+1$ are polynomials in the entries of $x$
with coefficients in $\mathbf Z$.

(5) Since $\frac{d}{dt}\exp(tY)=Y\exp(tY)$, the chain rule gives
$\frac{d}{dt}f(\exp(tY))=(D_Yf)(\exp(tY))$. By induction on $k$, the coefficient of $t^k$ in
$f(\exp(tY))$ is $(D_Y^kf)(1)/k!$. By (2), $D_Y^kf\in I(\Sigma)$, and it vanishes at $1\in\Sigma$.

(6) Here $\exp(tY)=1+tY$, so $f(1+tY)$ is a polynomial in $t$ that vanishes by (5). Setting $t=1$ gives
$f(1+Y)=0$ for every $f\in I(\Sigma)$.

(7) The polynomials of (4) lie in $I(\Sigma)$, so they vanish on $\bar\Sigma$. Let
$Y\in\mathfrak t(\Sigma)$ and let $L$ be the fraction field of $K[[t]]$. By (5) they vanish at
$\exp(tY)=1+tZ$, where $Z=\sum_{k\geq0}t^kY^{k+1}/(k+1)!$. So $1+tZ$ has degree at most $n$ over $L$,
and then so has $Z$: if $q(1+tZ)=0$, then $q(1+tX)$ is a nonzero polynomial of the same degree that
kills $Z$. Hence the polynomials of (4) vanish at $Z\in M_N(K[[t]])$, and taking constant terms, at
$Z(0)=Y$.
\end{proof}

\begin{lemma}\label{gal-lacunary-finite-index}
Let $\Gamma$ be a group, $K$ a field, $\rho:\Gamma\to\operatorname{GL}_N(K)$ a homomorphism, and $U$ a
normal subgroup of finite index in $\Gamma$. Then $\mathfrak g_\rho\subset\mathfrak t(\rho(U))$
(Definition~\ref{gal-lacunary-hull}).
\end{lemma}
\leanref{Development/EtaleCohomology/DivisorCycles/OpenCoset.lean}{LateralTheory.DivDb.TanIndex}
\leanref{Development/EtaleCohomology/DivisorCycles/TangentIndexDvd.lean}{LateralTheory.DivDb.TangentIndexDvd.TanIndex_of}

\begin{proof}
Let $f\in I(\rho(U))$. For each coset $c=gU$ we find $q_c\in K[x]$ with $q_c(1)\neq0$ and $fq_c$
vanishing on $\rho(gU)$. If some polynomial $q$ with $q(1)\neq0$ vanishes on $\rho(gU)$, take $q_c=q$.
Otherwise every polynomial that vanishes on $\rho(gU)$ vanishes at $1$. For $h\in I(\rho(U))$ the
polynomial $x\mapsto h(\rho(g)^{-1}x)$ vanishes on $\rho(gU)$, so $h(\rho(g)^{-1})=0$, and
$\rho(g)^{-1}\in\overline{\rho(U)}$. Some power $g^k$ with $k\geq1$ lies in $U$, so
$\rho(g)=\rho(g^k)\rho(g)^{-(k-1)}$ lies in $\overline{\rho(U)}$ by
Lemma~\ref{gal-lacunary-closure}(1), and so does $\rho(gu)=\rho(g)\rho(u)$ for $u\in U$; take $q_c=1$.
Now $F=f\prod_cq_c$ vanishes on $\rho(\Gamma)$. Since $f(1)=0$, $dF_1=\bigl(\prod_cq_c(1)\bigr)df_1$, so
$df_1(Y)=0$ for $Y\in\mathfrak g_\rho$.
\end{proof}

\begin{lemma}\label{gal-lacunary-envelope}
Let $\Gamma$ be a group, $K$ a field and $\rho:\Gamma\to\operatorname{GL}_N(K)$ a homomorphism. Then the
envelope $\mathcal A_\rho$ is $\rho$-stable (Definition~\ref{gal-lacunary-hull}).
\end{lemma}
\leanref{Development/EtaleCohomology/DivisorCycles/OpenCoset.lean}{LateralTheory.DivDb.EnvStab}
\leanref{Development/EtaleCohomology/DivisorCycles/CoeMulInv.lean}{LateralTheory.DivDb.EnvStab_of}

\begin{proof}
Both properties pass from generators to the algebra they generate, so it suffices to treat
$Y\in\mathfrak g_\rho$. Let $g\in\Gamma$. For $f\in I(\rho(\Gamma))$, the polynomial
$x\mapsto f(\rho(g)x\rho(g)^{-1})$ lies in $I(\rho(\Gamma))$, and by the chain rule its differential at
$1$ is $Y\mapsto df_1(\rho(g)Y\rho(g)^{-1})$; so $\rho(g)Y\rho(g)^{-1}\in\mathfrak g_\rho$. Let
$W\subset K^N$ be stable under $\rho(\Gamma)$, $v\in W$, and $\lambda$ a linear form vanishing on $W$. The
linear polynomial $x\mapsto\lambda(xv)$ vanishes on $\rho(\Gamma)$, so its differential $\lambda(Yv)$ is
$0$. Hence $Yv\in W$.
\end{proof}

\begin{lemma}\label{gal-lacunary-semisimple}
Let $\Gamma$ be a group, $K$ a field, $\rho:\Gamma\to\operatorname{GL}_N(K)$ a semisimple homomorphism,
and $\mathcal A\subset M_N(K)$ a $\rho$-stable subalgebra (Definition~\ref{gal-lacunary-hull}). Then
$\mathcal A$ is semisimple.
\end{lemma}
\leanref{Development/EtaleCohomology/DivisorCycles/CycleCount.lean}{LateralTheory.DivDb.Semis}
\leanref{Development/EtaleCohomology/DivisorCycles/SemisPf.lean}{LateralTheory.DivDb.SemisPf.Semis_of}

\begin{proof}
Let $J$ be the set of $Y\in\mathcal A$ with $aY$ nilpotent for every $a\in\mathcal A$, and $W$ the span
of the images of the elements of $J$. For $g\in\Gamma$ and $Y\in J$, $\rho(g)Y\rho(g)^{-1}$ lies in $J$,
since $a\rho(g)Y\rho(g)^{-1}$ is conjugate to $(\rho(g)^{-1}a\rho(g))Y$. As
$\rho(g)Yv=(\rho(g)Y\rho(g)^{-1})\rho(g)v$, $W$ is stable under $\rho(\Gamma)$, and it has a
$\rho(\Gamma)$-stable complement $W'$. Since $\mathcal A$ preserves $W'$, every $Y\in J$ maps $W'$ into
$W\cap W'=0$. Every $Y\in J$ lies in the Jacobson radical of $\mathcal A$, since $1-aY$ is invertible
for every $a\in\mathcal A$, and this radical is nilpotent because $\mathcal A$ is finite-dimensional
\cite{Lam01}; say its $m$-th power is $0$. For $Y\in J$ and $v=w+w'$ with $w\in W$ and $w'\in W'$, we
have $Yv=Yw$. By induction on $k$, $W$ is therefore spanned by vectors $Y_1\cdots Y_kv$ with all
$Y_i\in J$. For $k=m$ this gives $W=0$, so $J=0$.
\end{proof}

\begin{lemma}\label{gal-lacunary-grading}
Let $K$ be a field and $N,n\geq1$.
\begin{enumerate}
\item Let $X\in M_N(K)$ have degree at most $n$ and $n$ distinct eigenvalues $e_1,\ldots,e_n$. Then there
  is a system of idempotents $p_1,\ldots,p_n$, each a polynomial in $X$, with $X=\sum_ie_ip_i$.
\item Suppose that $K$ has characteristic $0$. Let $p_1,\ldots,p_n$ be a system of idempotents,
  $e_1,\ldots,e_n$ a lacunary family, $X=\sum_ie_ip_i$, and $\mathfrak g\subset M_N(K)$ a subspace with
  $XY-YX\in\mathfrak g$ for every $Y\in\mathfrak g$. Then $p_iYp_j\in\mathfrak g$ for every
  $Y\in\mathfrak g$ and $i\neq j$.
\end{enumerate}
\end{lemma}
\leanref{Development/EtaleCohomology/DivisorCycles/CycleCount.lean}{LateralTheory.DivDb.Proj}
\leanref{Development/EtaleCohomology/DivisorCycles/ProjPf.lean}{LateralTheory.DivDb.Proj_of}
\leanref{Development/EtaleCohomology/DivisorCycles/CycleCount.lean}{LateralTheory.DivDb.Grade}
\leanref{Development/EtaleCohomology/DivisorCycles/GradePf.lean}{LateralTheory.DivDb.Grade_of}

\begin{proof}
(1) The minimal polynomial of $X$ has degree at most $n$ and the $n$ distinct roots $e_i$, so it is
$\prod_i(T-e_i)$. Let $L_i$ be the Lagrange polynomials of degree $n-1$ with $L_i(e_k)=\delta_{ik}$. The
polynomials $L_iL_j-\delta_{ij}L_i$, $\sum_iL_i-1$ and $T-\sum_ie_iL_i$ vanish at every $e_k$, so they kill
$X$. Hence $p_i=L_i(X)$ is a system of orthogonal idempotents with sum $1$ and $X=\sum_ie_ip_i$, and
$p_i\neq0$ because $L_i$ has degree less than $n$.

(2) The map $\operatorname{ad}X:Y\mapsto XY-YX$ multiplies $p_kYp_l$ by $e_k-e_l$. If $i\neq j$ and
$e_i-e_j=e_k-e_l$, then $(k,l)=(i,j)$. Otherwise $k\neq i$, $l\neq j$ and $k\neq l$ by injectivity, and
$e_i+e_l=e_k+e_j$ says that the disjoint sets $\{i,l\}$ and $\{k,j\}$ have the same average (when $i=l$
and $k=j$ it says $e_i=e_j$); this contradicts lacunarity. Let $q(T)=\prod(T-(e_k-e_l))$, over the pairs
$(k,l)\neq(i,j)$. Writing $Y=\sum_{k,l}p_kYp_l$, we get
$q(\operatorname{ad}X)Y=q(e_i-e_j)\,p_iYp_j$, which lies in $\mathfrak g$, and $q(e_i-e_j)\neq0$.
\end{proof}

\begin{lemma}\label{gal-lacunary-cycles}
Let $K$ be an algebraically closed field of characteristic $0$, $p_1,\ldots,p_n$ a system of idempotents
in $M_N(K)$, $\pi$ a permutation of $\{1,\ldots,n\}$, and $y\in\operatorname{GL}_N(K)$ with
$yp_j=p_{\pi(j)}y$ for all $j$. Suppose that $y$ has degree at most $n$ and separates the cycles of $\pi$
(Definition~\ref{gal-lacunary-hull}). Then $y$ has $n$ distinct eigenvalues, and $yp_i\in Kp_i$ for every
$i$ with $\pi(i)=i$.
\end{lemma}
\leanref{Development/EtaleCohomology/DivisorCycles/CycleCount.lean}{LateralTheory.DivDb.CycleCount}
\leanref{Development/EtaleCohomology/DivisorCycles/CycleCountPf.lean}{LateralTheory.DivDb.CycleCountPf.CycleCount_of}
\leanref{Development/EtaleCohomology/DivisorCycles/CycSpecQzkPf.lean}{LateralTheory.DivDb.CycSpec_qzkPf.CycSpec_of}
\leanref{Development/EtaleCohomology/DivisorCycles/FixScalOf.lean}{LateralTheory.DivDb.ScalarFixedPointsUnderAeval.FixScal_of}
\leanref{Development/EtaleCohomology/DivisorCycles/CycleCount2.lean}{LateralTheory.DivDb.CycleCount₂.CycleCount_of}

\begin{proof}
\emph{Eigenvalues.} Let $O$ be a cycle of $\pi$, of length $k$, and $r\in O$. Then $y^k$ commutes with
$p_r$, so it has an eigenvector $v$ in the image of $p_r$, with an eigenvalue $\beta\neq0$ since $y$ is
invertible. Choose $\zeta$ with $\zeta^k=\beta$ and a primitive $k$-th root of unity $\omega$. For
$j=\pi^m(r)$ with $0\leq m<k$ put $\lambda_j=\zeta\omega^m$ and $w_j=\sum_{s<k}\lambda_j^{-s}y^sv$. Since
$\lambda_j^k=\beta$, $yw_j=\lambda_jw_j$. Since $y^sv$ lies in the image of $p_{\pi^s(r)}$, $w_j$ lies in
$V_O$ and $p_rw_j=v\neq0$. The $\lambda_j$ for $j\in O$ are distinct, and $\lambda_j\neq\lambda_{j'}$ for
$j,j'$ in different cycles because $y$ separates the cycles. So $y$ has the $n$ distinct eigenvalues
$\lambda_1,\ldots,\lambda_n$, with eigenvectors $w_j$ in $V_{O(j)}$ for the cycle $O(j)$ of $j$.

\emph{Fixed points.} Let $q\neq0$ of degree at most $n$ kill $y$. Each $\lambda_j$ is a root of $q$, so
$q$ is a nonzero multiple of $\prod_j(T-\lambda_j)$ and $\prod_j(y-\lambda_j)=0$. Let $\pi(i)=i$, so that
$y$ commutes with $p_i$. For $j\neq i$, $y-\lambda_j$ is injective on the image of $p_i$: a nonzero
$w$ there with $yw=\lambda_jw$ would be an eigenvector in $V_{\{i\}}$ for the eigenvalue $\lambda_j$,
which also has the eigenvector $w_j\in V_{O(j)}$, and $O(j)\neq\{i\}$. So $\prod_{j\neq i}(y-\lambda_j)$
is injective on the image of $p_i$, which contains $(y-\lambda_i)p_iu$ for every $u$. Hence
$(y-\lambda_i)p_i=0$.
\end{proof}

\begin{lemma}\label{gal-lacunary-diagonal}
Let $K$ be an algebraically closed field of characteristic $0$, $p_1,\ldots,p_n$ a system of idempotents
in $M_N(K)$, $e_1,\ldots,e_n\in K$ distinct, $X=\sum_ie_ip_i$, and $Y\in M_N(K)$ commuting with every
$p_i$ such that $sY+X$ has degree at most $n$ for every $s\in K$. Then $Yp_i\in Kp_i$ for every $i$.
\end{lemma}
\leanref{Development/EtaleCohomology/DivisorCycles/CycleCount.lean}{LateralTheory.DivDb.DiagScalar}
\leanref{Development/EtaleCohomology/DivisorCycles/DiagScalar2.lean}{LateralTheory.DivDb.DiagScalar_of}

\begin{proof}
Let $\Lambda$ be the finite set of $0$ and the quotients $(e_j-e_i)/(\mu-\mu')$ for eigenvalues
$\mu\neq\mu'$ of $Y$, choose $s\notin\Lambda$, and choose $c$ such that $y=sY+X+c$ is invertible. Then
$y$ commutes with every $p_j$ and has degree at most $n$. It separates the cycles of the identity
permutation: if $\lambda$ is an eigenvalue of $y$ with eigenvectors $v\in V_{\{i\}}$ and $w\in V_{\{j\}}$,
$i\neq j$, then $Xv=e_iv$ and $Xw=e_jw$, so $Yv=\mu v$ and $Yw=\mu'w$ with
$s\mu+e_i=\lambda-c=s\mu'+e_j$; then $\mu\neq\mu'$ and $s=(e_j-e_i)/(\mu-\mu')\in\Lambda$. By
Lemma~\ref{gal-lacunary-cycles}, $yp_i=c'p_i$ for some $c'$, so $Yp_i=s^{-1}(c'-e_i-c)p_i$.
\end{proof}

\begin{lemma}\label{gal-lacunary-algebra}
Let $K$ be a field, $p_1,\ldots,p_n$ a system of idempotents in $M_N(K)$, and $\mathfrak g\subset M_N(K)$
a subspace closed under $(Y,Z)\mapsto YZ-ZY$, such that $p_iYp_j\in\mathfrak g$ for all
$Y\in\mathfrak g$ and $i\neq j$, and $Yp_i\in Kp_i$ for every $i$ whenever $Y\in\mathfrak g$ commutes
with every $p_i$. Then every $a$ in the subalgebra generated by $\mathfrak g$ satisfies
$p_iap_j\in\mathfrak g$ for $i\neq j$ and $p_iap_i\in Kp_i$ for every $i$.
\end{lemma}
\leanref{Development/EtaleCohomology/DivisorCycles/CycleCount.lean}{LateralTheory.DivDb.AlgStruct}
\leanref{Development/EtaleCohomology/DivisorCycles/AlgStruct2.lean}{LateralTheory.DivDb.AlgStruct₂.AlgStruct_of}

\begin{proof}
The set $\mathcal S$ of $a$ with these two properties contains the scalars and is closed under sums. It
contains $\mathfrak g$: for $Y\in\mathfrak g$, $D=\sum_kp_kYp_k=Y-\sum_{k\neq l}p_kYp_l$ lies in
$\mathfrak g$ and commutes with every $p_m$, so $p_iYp_i=Dp_i\in Kp_i$. It is closed under products. Let
$a,b\in\mathcal S$, so that $p_iabp_k=\sum_j(p_iap_j)(p_jbp_k)$. For $i\neq k$, the terms with $j=i$ or
$j=k$ are scalar multiples of $p_ibp_k$ or $p_iap_k$, and for the other $j$,
$(p_iap_j)(p_jbp_k)=[p_iap_j,p_jbp_k]$ lies in $\mathfrak g$, because the product in the other order
contains $p_kp_i=0$. For $i=k$ and $j\neq i$, $Z=[p_iap_j,p_jbp_i]\in\mathfrak g$ is the sum of a
matrix in $p_iM_N(K)p_i$ and one in $p_jM_N(K)p_j$, so it commutes with every $p_m$ and
$(p_iap_j)(p_jbp_i)=Zp_i\in Kp_i$; the term $j=i$ is a product of scalar multiples of $p_i$. So
$\mathcal S$ is a subalgebra containing $\mathfrak g$.
\end{proof}

\begin{proposition}\label{gal-lacunary-structure}
Let $\Gamma$ be a group, $K$ an algebraically closed field of characteristic $0$, $N,n\geq1$, and
$\rho:\Gamma\to\operatorname{GL}_N(K)$ a semisimple homomorphism such that every $\rho(g)$ has degree at
most $n$. Let $X\in\mathfrak g_\rho$ and let $e_1,\ldots,e_n$ be a lacunary family of eigenvalues of $X$
(Definition~\ref{gal-lacunary-hull}). Then there is a corner system $p_1,\ldots,p_n$ for $\mathcal A_\rho$
with $\sum_ie_ip_i\in\mathfrak g_\rho$ and $p_iap_j\in\mathfrak g_\rho$ for all $a\in\mathcal A_\rho$ and
$i\neq j$. Moreover $\mathcal A_\rho$ is semisimple and $\rho$-stable.
\end{proposition}
\leanref{Development/EtaleCohomology/DivisorCycles/CycleCount.lean}{LateralTheory.DivDb.Struct}
\leanref{Development/EtaleCohomology/DivisorCycles/OpenCoset.lean}{LateralTheory.DivDb.Good}
\leanref{Development/EtaleCohomology/DivisorCycles/StructPf.lean}{LateralTheory.DivDb.StructPf.Struct_of}

\begin{proof}
Every element of $\rho(\Gamma)$, hence by Lemma~\ref{gal-lacunary-closure}(7) every element of
$\mathfrak g_\rho$, has degree at most $n$. The family $e$ is injective, so
Lemma~\ref{gal-lacunary-grading}(1) gives a system of idempotents $p_i$, polynomials in $X$ and so in
$\mathcal A_\rho$, with $X=\sum_ie_ip_i$, which lies in $\mathfrak g_\rho$. By
Lemma~\ref{gal-lacunary-closure}(3), $\mathfrak g_\rho$ is closed under commutators, so it is stable under
$Y\mapsto XY-YX$, and Lemma~\ref{gal-lacunary-grading}(2) gives $p_iYp_j\in\mathfrak g_\rho$ for
$Y\in\mathfrak g_\rho$ and $i\neq j$. If $Y\in\mathfrak g_\rho$ commutes with every $p_i$, then
$sY+X\in\mathfrak g_\rho$ has degree at most $n$ for every $s$, and Lemma~\ref{gal-lacunary-diagonal}
gives $Yp_i\in Kp_i$. Lemma~\ref{gal-lacunary-algebra} now gives $p_iap_j\in\mathfrak g_\rho$ for
$i\neq j$ and $p_iap_i\in Kp_i$, for every $a\in\mathcal A_\rho$. Finally $\mathcal A_\rho$ is
$\rho$-stable by Lemma~\ref{gal-lacunary-envelope} and semisimple by
Lemma~\ref{gal-lacunary-semisimple}.
\end{proof}

\begin{lemma}\label{gal-matrix-units}
Let $K$ be a field, $\mathcal A\subset M_N(K)$ a semisimple subalgebra, and $p_1,\ldots,p_n$ a corner
system for $\mathcal A$ (Definition~\ref{gal-lacunary-hull}). Then it has a family of matrix units in
$\mathcal A$.
\end{lemma}
\leanref{Development/EtaleCohomology/DivisorCycles/CycleCount.lean}{LateralTheory.DivDb.MatrixUnits₂}
\leanref{Development/EtaleCohomology/DivisorCycles/MatUnitsOf.lean}{LateralTheory.DivDb.MatrixUnitConstruction.MatUnits_of}

\begin{proof}
Write $i\sim j$ if there are $x\in p_i\mathcal Ap_j$ and $y\in p_j\mathcal Ap_i$ with $xy=p_i$ and
$yx=p_j$. This is an equivalence relation: $(p_i,p_i)$ gives reflexivity, exchanging $x$ and $y$ gives
symmetry, and products of such pairs give transitivity.

\emph{If $p_i\mathcal Ap_j\neq0$, then $i\sim j$.} Let $0\neq Y\in p_i\mathcal Ap_j$. By semisimplicity
some $b\in\mathcal A$ has $bY$ not nilpotent. Put $w=p_jbp_i$. By the corner property $Yw=cp_i$ and
$wY=dp_j$; then $YwY=cY=dY$, so $c=d$. If $c=0$, then $YbY=YwY=0$ and $(bY)^2=0$, which is excluded.
So $Y$ and $c^{-1}w$ give $i\sim j$.

\emph{Coefficients.} If $x,y$ give $i\sim j$ and $a\in\mathcal A$, then $p_iap_jy\in p_i\mathcal Ap_i$
equals $cp_i$, so $p_iap_j=p_iap_jyx=cx$.

\emph{Matrix units.} Choose a representative $\bar\imath$ of the class of each $i$, and $X_i,Y_i$ giving
$i\sim\bar\imath$. Put $E_{ij}=X_iY_j$ if $i\sim j$ and $E_{ij}=0$ otherwise. Then $E_{ii}=p_i$,
$p_iE_{ij}p_j=E_{ij}$, and for $i\sim j\sim k$, $E_{ij}E_{jk}=X_ip_{\bar\jmath}Y_k=E_{ik}$; the pair
$E_{ij},E_{ji}$ gives $i\sim j$, so $E_{ij}\neq0$. Finally $p_iap_j$ is a multiple of $E_{ij}$ by the
coefficient step if $i\sim j$, and is $0$ by the first step otherwise.
\end{proof}

\begin{lemma}\label{gal-matrix-units-conjugate}
Let $K$ be a field of characteristic $0$, $\mathcal A\subset M_N(K)$ a semisimple subalgebra, and
$p_1,\ldots,p_n$ and $q_1,\ldots,q_n$ corner systems for $\mathcal A$
(Definition~\ref{gal-lacunary-hull}). Then there are $v\in\mathcal A\cap\operatorname{GL}_N(K)$ and a
permutation $\pi$ of $\{1,\ldots,n\}$ with $vp_j=q_{\pi(j)}v$ for all $j$, and $\pi(i)=i$ for every $i$
with $q_i\mathcal Ap_i\neq0$.
\end{lemma}
\leanref{Development/EtaleCohomology/DivisorCycles/CycleCount.lean}{LateralTheory.DivDb.IdemConj}
\leanref{Development/EtaleCohomology/DivisorCycles/ConjugateIdempotents.lean}{LateralTheory.DivDb.ConjugateIdempotents.IdemConj_of}

\begin{proof}
Write $q_j\approx p_k$ if there are $\alpha\in q_j\mathcal Ap_k$ and $\beta\in p_k\mathcal Aq_j$ with
$\alpha\beta=q_j$ and $\beta\alpha=p_k$.

\emph{If $q_jap_k\neq0$ for some $a\in\mathcal A$, then $q_j\approx p_k$.} Put $\alpha=q_jap_k$. By
semisimplicity some $a'$ has $p_ka'\alpha\neq0$: otherwise $(b\alpha)^2=bq_ja(p_kb\alpha)=0$ for every
$b$. By the corner property $p_ka'\alpha=cp_k$ with $c\neq0$; put $\beta=c^{-1}p_ka'q_j$, so
$\beta\alpha=p_k$. Then $\alpha\beta=dq_j$ and $(\alpha\beta)^2=\alpha\beta$, so $d\in\{0,1\}$, and
$d\neq0$ because $\alpha=\alpha\beta\alpha$.

\emph{Classes.} Let $Q_k=\{j:q_j\approx p_k\}$. Composing the elements $\alpha,\beta$ shows that
$q_j\approx p_k$, $q_j\approx p_{k'}$ and $q_{j'}\approx p_k$ imply $q_{j'}\approx p_{k'}$, so two sets
$Q_k$, $Q_{k'}$ that meet are equal. If $j\notin Q_k$, then $q_jp_k=0$ by the first step with $a=1$; so
$p_k=\sum_jq_jp_k\neq0$ forces $Q_k\neq\emptyset$. For a class $C=\{k:Q_k=Q\}$, multiplying
$\sum_jq_j=1$ by $\sum_{k\in C}p_k$ gives $\sum_{k\in C}p_k=\sum_{j\in Q}q_j$. Equivalent idempotents
have equal traces ($\operatorname{tr}\alpha\beta=\operatorname{tr}\beta\alpha$), and a nonzero
idempotent has nonzero trace in characteristic $0$, so $|C|=|Q|$.

\emph{The permutation.} Since $|C|=|Q|$, the sets $\{k\in C:k\notin Q\}$ and $\{j\in Q:Q_j\neq Q\}$ have
the same size. Let $\pi(k)=k$ if $k\in Q_k$, and let $\pi$ map the first set bijectively onto the second
for each class. Then $\pi$ is a permutation with $\pi(k)\in Q_k$ for every $k$. Choose
$\alpha_k\in q_{\pi(k)}\mathcal Ap_k$ and $\beta_k$ giving $q_{\pi(k)}\approx p_k$, and put
$v=\sum_k\alpha_k$, $w=\sum_k\beta_k$. By orthogonality $vw=\sum_kq_{\pi(k)}=1$, $wv=\sum_kp_k=1$, and
$vp_j=\alpha_j=q_{\pi(j)}v$. If $q_i\mathcal Ap_i\neq0$, then $i\in Q_i$ by the first step, so
$\pi(i)=i$.
\end{proof}

\begin{lemma}\label{gal-matrix-units-factor}
Let $K$ be a field, $\mathcal A\subset M_N(K)$ a subalgebra, and $p_1,\ldots,p_n$ a corner system for
$\mathcal A$ with a family of matrix units $(E_{ij})$ in $\mathcal A$
(Definition~\ref{gal-lacunary-hull}). Then every $v\in\mathcal A\cap\operatorname{GL}_N(K)$ is
$v=u\sum_jd_jp_j$, with all $d_j\in K^\times$ and $u$ a product of elements $1+p_jap_k$ with
$a\in\mathcal A$ and $j\neq k$.
\end{lemma}
\leanref{Development/EtaleCohomology/DivisorCycles/OpenCoset.lean}{LateralTheory.DivDb.Tv}
\leanref{Development/EtaleCohomology/DivisorCycles/CycleCount.lean}{LateralTheory.DivDb.UnitFactorization}
\leanref{Development/EtaleCohomology/DivisorCycles/UnitFactorizationTheorem.lean}{LateralTheory.DivDb.UnitFactorizationTheorem.UnitFactor_of}

\begin{proof}
Let $\mathcal S$ be the set of products $u\sum_jd_jp_j$ as in the statement. It is closed under
products, because conjugation by $D=\sum_jd_jp_j$ maps $1+p_jap_k$ to $1+(d_j/d_k)p_jap_k$. The relation
$E_{ij}\neq0$ is an equivalence relation; call its classes blocks. For a union $T$ of blocks,
$z_T=\sum_{k\in T}p_k$ commutes with $\mathcal A$, because $p_iap_j\in KE_{ij}$ vanishes for $i,j$ in
different blocks. We show $1-z_T+z_Tv\in\mathcal S$ by induction on $T$; for $T=\{1,\ldots,n\}$ this is
$v\in\mathcal S$. If $T$ is the disjoint union of a block $C$ and a union $T'$ of blocks, then
$1-z_T+z_Tv=(1-z_C+z_Cv)(1-z_{T'}+z_{T'}v)$, so it suffices to treat a block $C$. The map
$\Theta(M)=\sum_{i,j\in C}M_{ij}E_{ij}$ is multiplicative from $M_C(K)$ to $M_N(K)$, with
$\Theta(1)=z_C$, so $\Psi(M)=1-z_C+\Theta(M)$ is multiplicative. Write $p_ivp_j=M_{ij}E_{ij}$ for
$i,j\in C$; then $\Theta(M)=z_Cv$ and $\Psi(M)=1-z_C+z_Cv$. The matrix $M$ is invertible: if $Mw=0$ and
$M''$ is the matrix with all columns $w$, then $v\Theta(M'')=z_Cv\Theta(M'')=\Theta(MM'')=0$, so
$\Theta(M'')=0$ since $v$ is invertible, and $w=0$. By Gaussian elimination $M$ is a product of a
diagonal matrix with nonzero entries and elementary matrices $1+c\,e_{ij}$ with $i\neq j$. Their images
under $\Psi$ are $\sum_jd_jp_j$, with $d_j=1$ off $C$, and $1+p_i(cE_{ij})p_j$, all in $\mathcal S$.
\end{proof}

\begin{lemma}\label{gal-lacunary-generic}
Let $K$ be a field, $p_1,\ldots,p_n$ a system of idempotents in $M_N(K)$, $\pi$ a permutation of
$\{1,\ldots,n\}$, and $x\in\operatorname{GL}_N(K)$ with $xp_j=p_{\pi(j)}x$ for all $j$. Then there is a
finite set $\Lambda\subset K$ such that for every family $\tau_1,\ldots,\tau_n$ multiplicatively
lacunary away from $\Lambda$ (Definition~\ref{gal-lacunary-hull}), $x\sum_j\tau_jp_j$ separates the
cycles of $\pi$.
\end{lemma}
\leanref{Development/EtaleCohomology/DivisorCycles/CycleCount.lean}{LateralTheory.DivDb.GenSpec}
\leanref{Development/EtaleCohomology/DivisorCycles/GenSpecProof.lean}{LateralTheory.DivDb.GenSpecProof.GenSpec_of}

\begin{proof}
Let $\Lambda$ be the set of quotients $\beta'^a/\beta^b$ for $1\leq a,b\leq n$, $\beta$ an eigenvalue of
$x^a$ and $\beta'$ an eigenvalue of $x^b$. Let $y=x\sum_j\tau_jp_j$. By induction on $m$,
$y^mp_k=\bigl(\prod_{s<m}\tau_{\pi^s(k)}\bigr)x^mp_k$; so for a cycle $O$ and $k\in O$,
$y^{|O|}p_k=\tau_Ox^{|O|}p_k$ with $\tau_O=\prod_{i\in O}\tau_i$. If $yv=\lambda v$ with $0\neq v\in V_O$,
it follows that $\lambda^{|O|}v=\tau_Ox^{|O|}v$, so $\lambda^{|O|}=\tau_O\beta$ for an eigenvalue $\beta$
of $x^{|O|}$, and $\lambda\neq0$. If moreover $yw=\lambda w$ with $0\neq w\in V_{O'}$ for another cycle
$O'$, then $\lambda^{|O'|}=\tau_{O'}\beta'$, and raising to the powers $|O'|$ and $|O|$ gives
$\tau_O^{|O'|}=(\beta'^{|O|}/\beta^{|O'|})\tau_{O'}^{|O|}$. This contradicts the choice of $\tau$.
\end{proof}

\begin{lemma}\label{gal-lacunary-torus}
Let $K$ be an algebraically closed field of characteristic $0$, $\Sigma\subset M_N(K)$ a submonoid,
$p_1,\ldots,p_n$ a system of idempotents and $e_1,\ldots,e_n$ a lacunary family with
$\sum_ie_ip_i\in\mathfrak t(\Sigma)$, and $\Lambda\subset K$ finite (Definition~\ref{gal-lacunary-hull}).
Then there is a family $\tau_1,\ldots,\tau_n$ multiplicatively lacunary away from $\Lambda$ with
$\sum_j\tau_jp_j\in\bar\Sigma$.
\end{lemma}
\leanref{Development/EtaleCohomology/DivisorCycles/OpenCoset.lean}{LateralTheory.DivDb.TorusPt}
\leanref{Development/EtaleCohomology/DivisorCycles/TorusPtPf.lean}{LateralTheory.DivDb.TorusPtPf.TorusPt_of}

\begin{proof}
Since the $p_i$ are orthogonal idempotents, $\exp\bigl(t\sum_ie_ip_i\bigr)=\sum_j\exp(e_jt)p_j$. By
Lemma~\ref{gal-lacunary-closure}(5), every polynomial $F_f(X_1,\ldots,X_n)=f\bigl(\sum_jX_jp_j\bigr)$
with $f\in I(\Sigma)$ vanishes at the point $(\exp(e_1t),\ldots,\exp(e_nt))$ of $K[[t]]^n$. The
polynomials $X_j$ and $\bigl(\prod_{i\in O}X_i\bigr)^{|O'|}-b\bigl(\prod_{i\in O'}X_i\bigr)^{|O|}$, for
disjoint nonempty $O,O'$ and $b\in\Lambda$, do not vanish there: the second is
$\exp(|O'|\sum_{O}e_i\,t)-b\exp(|O|\sum_{O'}e_i\,t)$, whose vanishing would force $b=1$ from the
constant term and $|O'|\sum_Oe_i=|O|\sum_{O'}e_i$ from the coefficient of $t$, against lacunarity.
Lemma~\ref{gal-lacunary-specialize}, with $R=K[[t]]$, gives $\tau\in K^n$ at which all $F_f$ vanish and
these finitely many polynomials do not.
\end{proof}

\begin{proposition}\label{gal-lacunary-key}
Let $\Gamma$ be a group, $K$ an algebraically closed field of characteristic $0$, $N,n\geq1$, and
$\rho:\Gamma\to\operatorname{GL}_N(K)$ a homomorphism such that every $\rho(g)$ has degree at most $n$.
Let $e_1,\ldots,e_n$ be a lacunary family and $p_1,\ldots,p_n$ a corner system for $\mathcal A_\rho$
with $\sum_ie_ip_i\in\mathfrak g_\rho$ and $p_iap_j\in\mathfrak g_\rho$ for all $a\in\mathcal A_\rho$
and $i\neq j$, and suppose that $\mathcal A_\rho$ is semisimple and $\rho$-stable
(Definition~\ref{gal-lacunary-hull}). Let $U$ be a normal subgroup of finite index and $g\in\Gamma$.
\begin{enumerate}
\item There is $m\in\overline{\rho(U)}$ such that $\rho(g)m$ has $n$ distinct eigenvalues.
\item For every $i$ such that $p_ia\rho(g)p_i\neq0$ for some $a\in\mathcal A_\rho$, there is
  $w\in\mathcal A_\rho$ with $p_i\rho(g)=p_iw$.
\end{enumerate}
The part of $\rho$-stability that asks every subspace stable under $\rho(\Gamma)$ to be stable under $\mathcal
A_\rho$ is not used.
\end{proposition}
\leanref{Development/EtaleCohomology/DivisorCycles/CycleCount.lean}{LateralTheory.DivDb.Key}
\leanref{Development/EtaleCohomology/DivisorCycles/SumMulSum.lean}{LateralTheory.DivDb.KeyPf.Key_of}

\begin{proof}
Write $\mathcal A=\mathcal A_\rho$ and $R=\rho(g)$. The idempotents $q_j=R^{-1}p_jR$ form a second
corner system for $\mathcal A$, which is $\rho$-stable. Lemma~\ref{gal-matrix-units-conjugate} gives
$v\in\mathcal A\cap\operatorname{GL}_N(K)$ and a permutation $\pi$ with $vp_j=q_{\pi(j)}v$, and
$\pi(i)=i$ whenever $q_i\mathcal Ap_i\neq0$. By Lemma~\ref{gal-matrix-units}, $p$ has matrix units in
$\mathcal A$, and Lemma~\ref{gal-matrix-units-factor} writes $v=uD$ with $D=\sum_jd_jp_j$ and $u$ a
product of elements $1+p_jap_k$, $j\neq k$. Then $up_j=d_j^{-1}vp_j=q_{\pi(j)}u$, so $x=Ru$ satisfies
$xp_j=p_{\pi(j)}x$.

Let $\Lambda$ be the finite set given by Lemma~\ref{gal-lacunary-generic} for $x$ and $\pi$. By
Lemma~\ref{gal-lacunary-finite-index}, $\sum_ie_ip_i\in\mathfrak t(\rho(U))$, and
Lemma~\ref{gal-lacunary-torus} gives $\tau$ multiplicatively lacunary away from $\Lambda$ with
$t=\sum_j\tau_jp_j\in\overline{\rho(U)}$. Each $p_jap_k$ with $j\neq k$ lies in
$\mathfrak g_\rho\subset\mathfrak t(\rho(U))$ and has square $0$, so $1+p_jap_k\in\overline{\rho(U)}$ by
Lemma~\ref{gal-lacunary-closure}(6), and $u\in\overline{\rho(U)}$ by
Lemma~\ref{gal-lacunary-closure}(1). Put $m=ut\in\overline{\rho(U)}$ and $y=Rm=xt$. Then $y$ lies in
$\overline{\rho(\Gamma)}$, so it has degree at most $n$ by Lemma~\ref{gal-lacunary-closure}(7); it is
invertible; $yp_j=p_{\pi(j)}y$ because $t$ commutes with every $p_j$; and $y$ separates the cycles of
$\pi$ by the choice of $\tau$. By Lemma~\ref{gal-lacunary-cycles}, $y=\rho(g)m$ has $n$ distinct
eigenvalues, which is (1), and $yp_i\in Kp_i$ whenever $\pi(i)=i$.

For (2), suppose $p_iaRp_i\neq0$ with $a\in\mathcal A$. Then $q_i(R^{-1}aR)p_i=R^{-1}p_iaRp_i\neq0$
and $R^{-1}aR\in\mathcal A$, so $\pi(i)=i$, and $p_iy=yp_i=cp_i$ for some $c\in K$. The inverse
$m^{-1}=\bigl(\sum_j\tau_j^{-1}p_j\bigr)Dv^{-1}$ lies in $\mathcal A$, since the inverse of an
invertible element of a finite-dimensional algebra is a polynomial in it. Hence
$p_iR=p_iym^{-1}=p_i(cm^{-1})$, and $w=cm^{-1}$ works.
\end{proof}

\begin{lemma}\label{gal-lacunary-big-corner}
Let $\Gamma$ be a group, $K$ a field, $\rho:\Gamma\to\operatorname{GL}_N(K)$ a homomorphism, and
$\mathcal A\subset M_N(K)$ a $\rho$-stable subalgebra with a corner system $p_1,\ldots,p_n$
(Definition~\ref{gal-lacunary-hull}). Suppose that for all $g\in\Gamma$ and $i$ with
$p_ia\rho(g)p_i\neq0$ for some $a\in\mathcal A$, there is $w\in\mathcal A$ with $p_i\rho(g)=p_iw$. Then
the span $\mathcal R$ of the products $a\rho(g)$, for $a\in\mathcal A$ and $g\in\Gamma$, is a
$\rho$-stable subalgebra containing $\rho(\Gamma)$, and $p_1,\ldots,p_n$ is a corner system for
$\mathcal R$.
\end{lemma}
\leanref{Development/EtaleCohomology/DivisorCycles/CycleCount.lean}{LateralTheory.DivDb.BigCorner}
\leanref{Development/EtaleCohomology/DivisorCycles/BigCorner2.lean}{LateralTheory.DivDb.BigCorner₂.BigCorner_of}

\begin{proof}
Since $a\rho(g)b\rho(h)=a\bigl(\rho(g)b\rho(g)^{-1}\bigr)\rho(gh)$, $\mathcal R$ is closed under
products; it contains $1$, every $\rho(g)$ and every $p_i=p_i\rho(1)$. Conjugation by $\rho(h)$ maps
$a\rho(g)$ to $\bigl(\rho(h)a\rho(h)^{-1}\bigr)\rho(hgh^{-1})$, and $a\rho(g)$ preserves every
$\rho(\Gamma)$-stable subspace; so $\mathcal R$ is $\rho$-stable. For the corner property it suffices to
consider $a\rho(g)$ with $p_ia\rho(g)p_i\neq0$. Then $p_i\rho(g)=p_iw$ with $w\in\mathcal A$, and with
$a'=\rho(g)^{-1}a\rho(g)\in\mathcal A$, $p_ia\rho(g)p_i=p_i\rho(g)a'p_i=p_iwa'p_i\in Kp_i$.
\end{proof}

\begin{lemma}\label{gal-matrix-units-representation}
Let $\Gamma$ be a group, $K$ a field, $\rho:\Gamma\to\operatorname{GL}_N(K)$ a homomorphism, and
$\mathcal R\subset M_N(K)$ a subalgebra containing $\rho(\Gamma)$, with a corner system
$p_1,\ldots,p_n$ and a family of matrix units $(E_{ij})$ in $\mathcal R$
(Definition~\ref{gal-lacunary-hull}). For $g\in\Gamma$ let $r(g)\in M_n(K)$ be the matrix with
$p_i\rho(g)p_j=r(g)_{ij}E_{ij}$, and $r(g)_{ij}=0$ when $E_{ij}=0$. Then $r$ is a homomorphism
$\Gamma\to\operatorname{GL}_n(K)$, and there are $A\in M_{N\times n}(K)$ and $B\in M_{n\times N}(K)$
such that:
\begin{enumerate}
\item $BA=1$ and $\rho(g)A=Ar(g)$ for all $g\in\Gamma$;
\item the only $Z\in M_N(K)$ with $ZT=0$ for every $T\in M_{N\times n}(K)$ such that
  $\rho(g)T=Tr(g)$ for all $g$ is $Z=0$.
\end{enumerate}
\end{lemma}
\leanref{Development/EtaleCohomology/DivisorCycles/OpenCoset.lean}{LateralTheory.DivDb.Sub}
\leanref{Development/EtaleCohomology/DivisorCycles/CycleCount.lean}{LateralTheory.DivDb.Rep}
\leanref{Development/EtaleCohomology/DivisorCycles/MinCongr.lean}{LateralTheory.DivDb.RepPf.Rep_of}

\begin{proof}
For $a\in\mathcal R$ let $c(a)\in M_n(K)$ be the matrix with $p_iap_j=c(a)_{ij}E_{ij}$, and $0$ where
$E_{ij}=0$, so that $r(g)=c(\rho(g))$. Since $p_iabp_k=\sum_j(p_iap_j)(p_jbp_k)$, $E_{ij}E_{jk}=E_{ik}$
when both factors are nonzero, and $E_{ik}=0$ forces $E_{ij}=0$ or $E_{jk}=0$ for every $j$, we get
$c(ab)=c(a)c(b)$; also $c(1)=1$. So $r$ is a homomorphism to $\operatorname{GL}_n(K)$.

The relation $E_{ij}\neq0$ is an equivalence relation; let $\bar\jmath$ be the least element of the
class of $j$. For $w\in K^N$ let $T(w)\in M_{N\times n}(K)$ have $j$-th column $E_{j\bar\jmath}w$. For
$a\in\mathcal R$, $aE_{j\bar\jmath}=\sum_ip_iap_jE_{j\bar\jmath}=\sum_ic(a)_{ij}E_{i\bar\imath}$, since
$E_{ij}E_{j\bar\jmath}=E_{i\bar\jmath}=E_{i\bar\imath}$ when $E_{ij}\neq0$. Hence $aT(w)=T(w)c(a)$, and
in particular $\rho(g)T(w)=T(w)r(g)$.

Choose $v_k\neq0$ in the image of $p_k$ and $u_k$ with $u_k^{\mathsf t}v_k=1$, put $w_0=\sum_kv_k$ and
$A=T(w_0)$, and let $B$ have $j$-th row $u_{\bar\jmath}^{\mathsf t}E_{\bar\jmath j}$. Since
$E_{\bar\jmath j}E_{k\bar k}$ is $p_{\bar\jmath}$ for $j=k$ and $0$ otherwise,
$(BA)_{jk}=\delta_{jk}u_{\bar\jmath}^{\mathsf t}p_{\bar\jmath}w_0=\delta_{jk}$. Finally, if $ZT=0$ for
every intertwiner $T$, then $ZE_{j\bar\jmath}w=0$ for all $w$ and $j$, so
$Zp_j=ZE_{j\bar\jmath}E_{\bar\jmath j}=0$ for every $j$, and $Z=Z\sum_jp_j=0$.
\end{proof}

\begin{lemma}\label{gal-matrix-units-subrepresentation}
Let $\Gamma$ be a group, $K$ a field, $\rho:\Gamma\to\operatorname{GL}_N(K)$ a semisimple homomorphism,
$r:\Gamma\to\operatorname{GL}_n(K)$ a homomorphism, and $A\in M_{N\times n}(K)$, $B\in M_{n\times N}(K)$
with properties (1) and (2) of Lemma~\ref{gal-matrix-units-representation}. Then:
\begin{enumerate}
\item $r(g)=B\rho(g)A$ for all $g$, $A$ has rank $n$, and $r$ is semisimple;
\item for $g\in\Gamma$ and $q\in K[X]$, $q(\rho(g))=0$ if and only if $q(r(g))=0$;
\item there are $k\geq0$, $T_1,\ldots,T_k\in M_{n\times N}(K)$ and
  $T'_1,\ldots,T'_k\in M_{N\times n}(K)$ with $T_j\rho(g)=r(g)T_j$ and $\rho(g)T'_j=T'_jr(g)$ for all
  $j$ and $g$, and $\sum_jT'_jT_j=1$.
\end{enumerate}
\end{lemma}
\leanref{Development/EtaleCohomology/DivisorCycles/CycleCount.lean}{LateralTheory.DivDb.FinSub}
\leanref{Development/EtaleCohomology/DivisorCycles/FinSubOf.lean}{LateralTheory.DivDb.FiniteSubrepresentation.FinSub_of}
\leanref{Development/EtaleCohomology/DivisorCycles/CycleCount.lean}{LateralTheory.DivDb.FinEmbed}
\leanref{Development/EtaleCohomology/DivisorCycles/FinEmbedPf.lean}{LateralTheory.DivDb.FinEmbedPf.FinEmbed_of}

\begin{proof}
Call $T$ an intertwiner if $\rho(g)T=Tr(g)$ for all $g$, or $T\rho(g)=r(g)T$ for all $g$ when
$T\in M_{n\times N}(K)$.

(1) $r(g)=BAr(g)=B\rho(g)A$, and $A$ has rank $n$ because $BA=1$. If $W\subset K^n$ is $r$-stable, then
$AW$ is $\rho$-stable and has a $\rho$-stable complement $W''$. Its preimage under $A$ is $r$-stable,
meets $W$ in $0$ because $A$ is injective, and spans $K^n$ together with $W$ because
$AW+W''=K^N$; so it is an $r$-stable complement of $W$.

(2) If $q(\rho(g))=0$, then $q(r(g))=Bq(\rho(g))A=0$. If $q(r(g))=0$, then
$q(\rho(g))T=Tq(r(g))=0$ for every intertwiner $T$, so $q(\rho(g))=0$ by property (2).

(3) Call an idempotent $E\in M_N(K)$ admissible if it commutes with $\rho(\Gamma)$ and is a linear
combination of products $T'T$ of intertwiners. Since $E=0$ is admissible, it suffices to replace an
admissible $E$ with $\ker E\neq0$ by an admissible $E'$ with $\ker E'\subsetneq\ker E$; an admissible
$E$ with $\ker E=0$ is $1$. The idempotent $P=1-E$ is nonzero and commutes with $\rho(\Gamma)$, so by
property (2) some intertwiner $T'$ has $M=PT'\neq0$; $M$ is an intertwiner with $EM=0$. Since $r$ is
semisimple, $\ker M$ has an $r$-stable complement $C$, on which $M$ is injective; $M(C)\neq0$ is
$\rho$-stable and has a $\rho$-stable complement $D$. Let $T$ be the projection onto $M(C)$ along $D$
followed by the inverse of $M|_C$, an intertwiner. Then $H=MTP$, the product of the intertwiners $M$
and $TP$, is an idempotent commuting with $\rho(\Gamma)$ with $EH=HE=0$, so $E'=E+H$ is admissible,
$\ker E'\subset\ker E$, and $E'y=y\neq0$ for $y\in M(C)\subset\ker E$.
\end{proof}

\begin{lemma}\label{gal-lacunary-coset}
Let $\Gamma$ be a compact group, $K$ a field, $r:\Gamma\to\operatorname{GL}_n(K)$ a homomorphism, $f$ a
polynomial in the entries of $n\times n$ matrices, and $O\subset\Gamma$ a nonempty open set with
$f(r(g))=0$ for $g\in O$. Then there are a normal subgroup $U$ of finite index and $g\in\Gamma$ with
$f(r(gu))=0$ for all $u\in U$.
\end{lemma}
\leanref{Development/EtaleCohomology/DivisorCycles/OpenCoset.lean}{LateralTheory.DivDb.OpenCoset}
\leanref{Development/EtaleCohomology/DivisorCycles/OpenCosetPf.lean}{LateralTheory.DivDb.OpenCoset_of}
\leanref{Development/AutomorphicArithmetic/IntegralCasts/CapelliIdentities/CosetOfCover.lean}{UpperTheory.Reformulated.B.Qn18.main}

\begin{proof}
Give $\Gamma$ the topology whose closed sets are the sets $\{g:r(g)\in Z\}$ for Zariski closed
$Z\subset M_n(K)$. It is Noetherian, and left translations are continuous for it, since
$x\mapsto h(r(a)x)$ is a polynomial when $h$ is. Let $Z_f=\{g:f(r(g))=0\}\supset O$. By compactness
finitely many translates $a^{-1}O$, $a\in F$, cover $\Gamma$, so $\Gamma=\bigcup_{a\in F}a^{-1}Z_f$.
The irreducible component $Y$ of $1$ for the new topology lies in one of these closed sets, say
$a^{-1}Z_f$. Left translations permute the finitely many irreducible components, so the stabilizer of
$Y$ has finite index, and so has its normal core $U$. For $u\in U$, $u=u\cdot1\in uY=Y$, so
$f(r(au))=0$.
\end{proof}

\begin{lemma}\label{gal-lacunary-separable}
Let $K$ be a field of characteristic $0$ and $n\geq1$.
\begin{enumerate}
\item There is a polynomial $\Delta$ in the entries of $n\times n$ matrices such that $\Delta(x)\neq0$
  if and only if the characteristic polynomial of $x$ is separable.
\item Let $\Gamma$ be a group, and $\rho$, $r$, $A$ and $B$ as in
  Lemma~\ref{gal-matrix-units-representation} with properties (1) and (2). Let $U$ be a subgroup of
  $\Gamma$, $g\in\Gamma$, and $m\in\overline{\rho(U)}$ such that $\rho(g)m$ has $n$ distinct
  eigenvalues. Then $r(gu)$ has separable characteristic polynomial for some $u\in U$.
\item Let $\Gamma$ be a compact group, $K$ a Hausdorff topological field, and
  $r:\Gamma\to\operatorname{GL}_n(K)$ continuous, such that every coset of every normal subgroup of
  finite index contains an element $g$ for which $r(g)$ has separable characteristic polynomial. Then
  the set of such $g$ is open and dense in $\Gamma$.
\end{enumerate}
\end{lemma}
\leanref{Development/EtaleCohomology/DivisorCycles/OpenCoset.lean}{LateralTheory.DivDb.SepPoly}
\leanref{Development/EtaleCohomology/DivisorCycles/SepPolyProof.lean}{LateralTheory.DivDb.SepPoly_of}
\leanref{Development/EtaleCohomology/DivisorCycles/CycleCount.lean}{LateralTheory.DivDb.DenseAlg}
\leanref{Development/EtaleCohomology/DivisorCycles/DenseAlg2.lean}{LateralTheory.DivDb.DenseAlg₂.DenseAlg_of}
\leanref{Development/EtaleCohomology/DivisorCycles/CycleCount.lean}{LateralTheory.DivDb.FinDense}
\leanref{Development/EtaleCohomology/DivisorCycles/FinDensePf.lean}{LateralTheory.DivDb.FinDense_of}

\begin{proof}
(1) Take for $\Delta(x)$ the resultant of the characteristic polynomial of $x$ and its derivative,
which has degree $n-1$ in characteristic $0$; its coefficients are polynomials in the entries.

(2) Suppose not. Let $\varphi(x)=B\rho(g)xA$, a linear map $M_N(K)\to M_n(K)$. Then
$\Delta(\varphi(x))$ is a polynomial in $x$ that vanishes at $\rho(u)$ for $u\in U$, since
$\varphi(\rho(u))=B\rho(gu)A=r(gu)$; so it vanishes at $m$, and the characteristic polynomial $\chi$ of
$z=\varphi(m)$ is not separable. Let $T$ be an intertwiner as in property (2). The entries of
$(\rho(g)xT-TB\rho(g)xA)B$ are linear polynomials in $x$ that vanish at every $\rho(u)$, so at $m$;
multiplying by $A$ on the right gives $\rho(g)mT=Tz$. Hence $\chi(\rho(g)m)T=T\chi(z)=0$ for every
such $T$, and $\chi(\rho(g)m)=0$. Then the $n$ distinct eigenvalues of $\rho(g)m$ are roots of $\chi$,
which is monic of degree $n$, so $\chi$ is separable, a contradiction.

(3) By (1) the set $V$ of such $g$ is $\{g:\Delta(r(g))\neq0\}$, which is open. If a nonempty open set
$O$ missed $V$, then $\Delta(r(g))=0$ on $O$, and Lemma~\ref{gal-lacunary-coset} would give a coset of
a normal subgroup of finite index on which $\Delta\circ r$ vanishes, against the hypothesis. So $V$ is
dense.
\end{proof}

\begin{proposition}\label{gal-division}
Let $\Gamma$ be a compact group, $K$ an algebraically closed field of characteristic $0$ with a
Hausdorff topology for which it is a topological ring, $N,n\geq1$, and
$\rho:\Gamma\to\operatorname{GL}_N(K)$ continuous and semisimple. Let
$S\subset\Gamma$ be dense, and for $s\in S$ let $P_s\in K[X]$ be monic of degree $n$ with
$P_s(\rho(s))=0$. Suppose that the Lie algebra of the Zariski closure of $\rho(\Gamma)$ contains, after
extension to an algebraically closed field, an element with $n$ distinct eigenvalues forming a lacunary
set, and that $\psi:\Gamma\to K^\times$ is a continuous character with $\psi(s)=(-1)^nP_s(0)$ for
$s\in S$. Then there is a continuous semisimple $r:\Gamma\to\operatorname{GL}_n(K)$ such that:
\begin{enumerate}
\item $r$ is a subrepresentation of $\rho$, and $\rho$ is a quotient of $r^{\oplus k}$ for some $k$;
\item $\det r=\psi$, and the characteristic polynomial of $r(g)$ kills $\rho(g)$ for every $g$;
\item the set of $g$ for which $r(g)$ has separable characteristic polynomial is open and dense, and
  for $s\in S$ the characteristic polynomial of $r(s)$ is $P_s$ if it is separable, and also if $P_s$
  is separable and its roots form a multiplicatively generic multiset
  (Definition~\ref{gal-generic-parameters}, with $K$ in place of $\mathbf C$).
\end{enumerate}
\end{proposition}
\leanref{Development/WeilBound/CoreStatements/PadicHodgeComparison/CocharacterLieDivisors/DivWCoreLacUseType.lean}{PositivityTheory.AutomorphicGaloisHodgeTate.Div.DivWCoreLac}
\leanref{Development/EtaleCohomology/DivisorCycles/GoalHolds.lean}{LateralTheory.DivDb.goal_holds}
\leanref{Development/EtaleCohomology/DivisorCycles/Imports.lean}{LateralTheory.DivDb.goal_of}

\begin{proof}
Write $\mathfrak g=\mathfrak g_\rho$ and $\mathcal A=\mathcal A_\rho$ (Definition~\ref{gal-lacunary-hull}).

\emph{The Lie element over $K$.} By Lemma~\ref{gal-lacunary-descent} there are $X\in\mathfrak g$ and a
lacunary family $e_1,\ldots,e_n\in K$ of eigenvalues of $X$.

\emph{Degrees.} For $s\in S$, $\rho(s)$ is killed by the nonzero polynomial $P_s$ of degree $n$, so it
has degree at most $n$. By Lemma~\ref{gal-lacunary-closure}(4) this is the vanishing of a set of
polynomials in the entries. Composed with $\rho$ they are continuous functions on $\Gamma$ that vanish
on the dense set $S$, and $K$ is Hausdorff, so every $\rho(g)$ has degree at most $n$.

\emph{Corners and matrix units.} Proposition~\ref{gal-lacunary-structure} gives a corner system
$p_1,\ldots,p_n$ for $\mathcal A$ with $\sum_ie_ip_i\in\mathfrak g$ and $p_iap_j\in\mathfrak g$ for
$a\in\mathcal A$ and $i\neq j$, and shows that $\mathcal A$ is semisimple and $\rho$-stable.
Proposition~\ref{gal-lacunary-key}(2), with $U=\Gamma$, gives the hypothesis of
Lemma~\ref{gal-lacunary-big-corner}, which yields a $\rho$-stable subalgebra $\mathcal R$ containing
$\rho(\Gamma)$ for which $p$ is a corner system. It is semisimple by
Lemma~\ref{gal-lacunary-semisimple}, so $p$ has a family of matrix units in $\mathcal R$ by
Lemma~\ref{gal-matrix-units}.

\emph{The representation.} Lemma~\ref{gal-matrix-units-representation} gives a homomorphism
$r:\Gamma\to\operatorname{GL}_n(K)$ and matrices $A$, $B$ with its properties (1) and (2). By
Lemma~\ref{gal-matrix-units-subrepresentation}, $r(g)=B\rho(g)A$, so $r$ is continuous; $r$ is
semisimple; $A$ has rank $n$ and $\rho(g)A=Ar(g)$, so $r$ is a subrepresentation of $\rho$; a polynomial
kills $\rho(g)$ if and only if it kills $r(g)$; and there are intertwiners $T_j$, $T'_j$ with
$\sum_jT'_jT_j=1$. The map $(K^n)^k\to K^N$ given by $(T'_1,\ldots,T'_k)$ is then a surjective
intertwiner from $r^{\oplus k}$ to $\rho$. This is (1).

\emph{Separability.} Let $U$ be a normal subgroup of finite index and $g\in\Gamma$.
Proposition~\ref{gal-lacunary-key}(1) gives $m\in\overline{\rho(U)}$ such that $\rho(g)m$ has $n$
distinct eigenvalues, and Lemma~\ref{gal-lacunary-separable}(2) gives $u\in U$ such that $r(gu)$ has
separable characteristic polynomial. By Lemma~\ref{gal-lacunary-separable}(3), the set $V$ of $g$ for
which $r(g)$ has separable characteristic polynomial is open and dense.

\emph{Characteristic polynomials.} For every $g$, the characteristic polynomial of $r(g)$ kills $r(g)$
by the Cayley--Hamilton theorem, and so kills $\rho(g)$. For $s\in S$, $P_s(r(s))=0$ because
$P_s(\rho(s))=0$. If moreover $s\in V$, the $n$ distinct eigenvalues of $r(s)$ are roots of $P_s$,
which is monic of degree $n$, so the characteristic polynomial of $r(s)$ is $P_s$.

\emph{The determinant.} The set $S\cap V$ is dense, since $S$ is dense and $V$ is open and dense. For
$s\in S\cap V$, $\det r(s)=(-1)^nP_s(0)=\psi(s)$. Both $\det r$ and $\psi$ are continuous, so
$\det r=\psi$. This completes (2).

\emph{Generic roots.} Let $s\in S$ be such that $P_s$ is separable and its roots form a
multiplicatively generic multiset $R$. Their product is $(-1)^nP_s(0)=\psi(s)\neq0$, so they are
nonzero; $\prod_{z\in R}(r(s)-z)=P_s(r(s))=0$; and $\det r(s)=\psi(s)=\prod_{z\in R}z$. By
Lemma~\ref{gal-char-of-ann}, the characteristic polynomial of $r(s)$ is $P_s$. This completes (3). The
separability of $P_s$ is not used in this step.
\end{proof}
\leanref{Development/EtaleCohomology/DivisorCycles/DivPlus.lean}{LateralTheory.DivDb.DivPlus}
\leanref{Development/EtaleCohomology/DivisorCycles/DivPlus.lean}{LateralTheory.DivDb.plus_of}
\leanref{Development/EtaleCohomology/DivisorCycles/CycleCount.lean}{LateralTheory.DivDb.DegAll}
\leanref{Development/EtaleCohomology/DivisorCycles/EntriesCont.lean}{LateralTheory.DivDb.DegAll_of}
\leanref{Development/EtaleCohomology/DivisorCycles/CycleCount.lean}{LateralTheory.DivDb.FinPoly}
\leanref{Development/EtaleCohomology/DivisorCycles/FinPolyPf.lean}{LateralTheory.DivDb.FinPolyPf.FinPoly_of}

\begin{proposition}\label{gal-division-seed}
Let $E$ be a CM field containing $\sqrt{-d}$ for an integer $d>0$, $\Omega$ an algebraic closure,
$\ell$ a prime, $\iota:\bar{\mathbf Q}_\ell\to\mathbf C$ a field isomorphism, $a\geq1$ and $U$ an
automorphic representation of $U_a$ (Definition~\ref{gal-definite-unitary}) that has a seed
representation (Definition~\ref{gal-seed}). Then there is a continuous semisimple
$r:G_E\to\operatorname{GL}_a(\bar{\mathbf Q}_\ell)$, unique up to isomorphism among those with the
property (0) below, such that:
\begin{enumerate}
\item[(0)] there are an open dense $V\subset G_E$ and a finite set $X$ of places such that for every
  $d$-split $w\notin X$ not over $\ell$, every $B$ such that $U$ has Satake parameters $B$ at $w$ and
  every valuation ring $A$ over $w$, $r$ is trivial on $I_A$ and the characteristic polynomial of
  $\iota(r(\sigma^{-1}))$ is $\prod_{b\in B}(X-q_w^{(a-1)/2}b)$ for every arithmetic Frobenius element
  $\sigma\in V$ at $w$;
\end{enumerate}
and, for every seed representation $\rho$ for $U$ and every place $w$:
\begin{enumerate}
\item if $w$ is split over $F$ and not over $\ell$, $U$ has Satake parameters $B$ at $w$ with $B$
  multiplicatively generic, and $\rho$ satisfies the condition of item~(3)(a) of
  Definition~\ref{gal-seed} at $w$, then $r$ has Frobenius eigenvalues $q_w^{(a-1)/2}B$ at $w$;
\item if $\rho|_{D_A}$ is crystalline for every $A$ over $w$ and its Hodge--Tate multisets contain the
  sets $\{k-\lambda_k\}$ for the weights $\lambda$ of $U$, then $r|_{D_A}$ is crystalline and de Rham
  with $\mathrm{HT}_\tau(r|_{D_A})=\{k-\lambda_k:0\leq k<a\}$;
\item if $w$ is split over $F$, $U$ has multiplicatively generic Satake parameters $B$ at $w$ and
  $\rho$ has the property of Lemma~\ref{gal-pack-seed}(3) at $w$, then $r|_{D_A}$ is crystalline and
  satisfies (phi) with respect to $B$ (Definition~\ref{gal-spherical-conditions}).
\end{enumerate}
\end{proposition}
\leanref{Development/AutomorphicArithmetic/Foundations/GaloisRepresentations/IntegralStructures/ZetaCharacterDivisibility/InsideRefl.lean}{UpperTheory.Gal.Int.Div.DivAll}
\leanref{Development/AutomorphicArithmetic/Foundations/GaloisRepresentations/IntegralStructures/ZetaCharacterDivisibility/DivisionOneGeometricRepresentation.lean}{UpperTheory.Gal.Int.Div.divAll_holds}
\leanref{Development/AutomorphicArithmetic/Foundations/GaloisRepresentations/IntegralStructures/ZetaCharacterDivisibility/InsideRefl.lean}{UpperTheory.Gal.Int.Div.IsDivRep}
\leanref{Development/AutomorphicArithmetic/Foundations/GaloisRepresentations/IntegralStructures/ZetaCharacterDivisibility/InsideRefl.lean}{UpperTheory.Gal.Int.Div.IsDetChar}
\leanref{Development/AutomorphicArithmetic/Foundations/GaloisRepresentations/IntegralStructures/ZetaCharacterDivisibility/PhiOfDivOf.lean}{UpperTheory.Gal.Int.Div.DivSet}
\leanref{Development/AutomorphicArithmetic/Foundations/GaloisRepresentations/IntegralStructures/ZetaCharacterDivisibility/PhiOfDivOf.lean}{UpperTheory.Gal.Int.Div.DivOne}
\leanref{Development/AutomorphicArithmetic/Foundations/GaloisRepresentations/IntegralStructures/ZetaCharacterDivisibility/PhiOfDivOf.lean}{UpperTheory.Gal.Int.Div.SubLoc}

\begin{proof}
\emph{The determinant.} There is a continuous homomorphism $\zeta$ from the ideles of $E$ to the
centre of $U_a$, trivial on the ideles fixed by $c$, carrying $E^\times$ into $U_a(F)$, and equal at a
split place $w$ to the scalar matrices under the identification of $U_{a,w}$ with
$\operatorname{GL}_a(E_w)$. By Schur's lemma $\zeta(z)$ acts on $U$ by a scalar $\chi(z)$, and $\chi$ is
a Hecke character of $E$, trivial on the ideles fixed by $c$ and hence algebraic. At a split place $w$
where $U$ has Satake parameters $B$, a spherical vector is fixed by the central units and multiplied by
$\prod_{b\in B}b$ by a central uniformizer, so $\chi$ is unramified at $w$ with
$\chi(\varpi_w)=\prod_{b\in B}b$. Let $\nu=|\cdot|^{-a(a-1)/2}$, so that
$\chi\nu(\varpi_w)=\prod_{b\in B}q_w^{(a-1)/2}b$, and let $\psi:G_E\to\bar{\mathbf Q}_\ell^\times$ be the
$\ell$-adic character attached to the algebraic Hecke character $\chi\nu$ \cite[Chapter~II]{Ser68}.
It is trivial on inertia at the split places $w\nmid\ell$ where $U$ has Satake parameters $B$, with
$\iota\psi(\sigma^{-1})=\prod_{b\in B}q_w^{(a-1)/2}b$ for arithmetic Frobenius elements $\sigma$ at
$w$.

\emph{Division.} Let $\rho$ be a seed representation for $U$, of dimension $N$. The geometric
Frobenius elements at $d$-split places $w\nmid\ell$ outside any finite set, at which $U$ has Satake
parameters $B$, are dense in $G_E$ by the Chebotarev density theorem, and at such an element $s$,
$P_s(X)=\prod_{b\in B}(X-\iota^{-1}(q_w^{(a-1)/2}b))$ is a monic polynomial of degree $a$ with
$P_s(\rho(s))=0$ and $(-1)^aP_s(0)=\psi(s)$. Theorem~\ref{gal-sen} applied to the lacunary Hodge--Tate
weights of $\rho$ gives the Lie element, and Proposition~\ref{gal-division} gives an $a$-dimensional
continuous semisimple subrepresentation $r$ of $\rho$ with $\rho$ a quotient of $r^{\oplus k}$, with
$\det r=\psi$, and with separable characteristic polynomial on an open dense set $V'$; at the
Frobenius elements $s$ above with $s\in V'$, the characteristic polynomial of $r(s)$ is $P_s$. Off a
finite set of places, $\rho$ and hence $r$ is unramified at these places. Taking for $V$ the set of
inverses of elements of $V'$ gives (0).

\emph{Uniqueness.} Two representations with property (0) have the same characteristic polynomials on
a dense set of elements, hence the same traces everywhere by continuity, and they are semisimple, so
they are isomorphic. In particular the $r$ obtained from any seed representation is the same up to
isomorphism, and being a subrepresentation of $\rho$ and having $\rho$ as a quotient of copies of $r$
pass along the isomorphism.

\emph{Local properties.} Since $r$ is a subrepresentation of $\rho$: if $\rho|_{D_A}$ is crystalline,
so is $r|_{D_A}$, and its Hodge--Tate multisets are contained in those of $\rho$, and the characteristic
polynomials of $\mathrm{WD}_\iota(r|_{D_A})$ divide those of $\rho$. For (1), the Frobenius at $w$ kills
$\prod_{b\in B}(X-\iota^{-1}(q_w^{(a-1)/2}b))$ on $\rho$ and hence on $r$, its determinant is given by
$\psi$, and Lemma~\ref{gal-char-of-ann} gives the characteristic polynomial. For (2), the Hodge--Tate
multiset of $r$ has $a$ elements and is contained in that of $\rho$; since $\rho$ is a quotient of
copies of $r$, every Hodge--Tate weight of $\rho$ is one of $r$, and the set $\{k-\lambda_k\}$ has $a$
distinct elements, which forces $\mathrm{HT}_\tau(r|_{D_A})=\{k-\lambda_k\}$. For (3), the
characteristic polynomial of $\mathrm{WD}_\iota(r|_{D_A})$ at a geometric Frobenius element divides a
power of $\prod_{b\in B}(X-q_w^{(a-1)/2}b)$, and its determinant is given by $\psi$, so
Lemma~\ref{gal-char-of-ann} gives the characteristic polynomial at Frobenius elements; at inertia
elements both sides are trivial, since $r|_{D_A}$ is crystalline.
\end{proof}
\leanref{Development/AutomorphicArithmetic/Foundations/GaloisRepresentations/IntegralStructures/ZetaCharacterDivisibility/PSmall.lean}{UpperTheory.Gal.Int.Div.centU_of}
\leanref{Development/AutomorphicArithmetic/Foundations/GaloisRepresentations/IntegralStructures/ZetaCharacterDivisibility/PDetEx.lean}{UpperTheory.Gal.Int.Div.detEx_of}
\leanref{Development/AutomorphicArithmetic/Foundations/GaloisRepresentations/IntegralStructures/ZetaCharacterDivisibility/PJumpSeed.lean}{UpperTheory.Gal.Int.Div.jumpSeed_of}

\subsection{Congruences and trace limits}\label{subsection-gal-congruences}

\begin{lemma}\label{gal-congruences-calculus}
Let $k$ be a field with an ultrametric absolute value. For a set $J$, a function $\theta:J\to k$ and a
set $S$ of functions $J\to k$, \emph{approximated} has the meaning of Lemma~\ref{gal-trace-limit}, with
$k$ in place of $\bar{\mathbf Q}_\ell$.
\begin{enumerate}
\item Let $J$ be a set, $\theta:J\to k$, $S$ and $T$ sets of functions $J\to k$, and $Q$ a polynomial with
  integer coefficients in variables $X_j$, $j\in J$. If $\theta$ is approximated by $S$, and every
  $\eta\in S$ has $Q(\eta)\neq0$ and is approximated by $T$, then $\theta$ is approximated by
  $\{\zeta\in T:Q(\zeta)\neq0\}$.
\item Let $i:J\to J'$ be a map of sets, $\theta:J'\to k$ and $S$ a set of functions $J'\to k$. If
  $\theta$ is approximated by $S$, then $\theta\circ i$ is approximated by $\{\eta\circ i:\eta\in S\}$.
\end{enumerate}
\end{lemma}
\leanref{Development/AutomorphicArithmetic/Foundations/GaloisRepresentations/SeparatingAlgebras/StGlob.lean}{UpperTheory.Gal.Gen.Combine}
\leanref{Development/AutomorphicArithmetic/Foundations/GaloisRepresentations/SeparatingAlgebras/StGlob.lean}{UpperTheory.Gal.Gen.Restrict}
\leanref{Development/AutomorphicArithmetic/Foundations/GaloisRepresentations/SeparatingAlgebras/Combine.lean}{UpperTheory.Gal.Gen.combine_pf}
\leanref{Development/AutomorphicArithmetic/Foundations/GaloisRepresentations/SeparatingAlgebras/Restrict.lean}{UpperTheory.Gal.Gen.restrict_pf}

\begin{proof}
(1) First let $\eta\in S$ and $\varepsilon>0$. Approximate $\eta$ within $\varepsilon|Q(\eta)|$ by a
combination $\sum_\zeta c_\zeta\zeta$ of finitely many members of $T$, and apply the estimate to the
polynomials $PQ$. Dividing by $Q(\eta)$ gives
$|P(\eta)-\sum_\zeta c_\zeta Q(\zeta)Q(\eta)^{-1}P(\zeta)|\leq\varepsilon$ for every $P$, and only the
$\zeta$ with $Q(\zeta)\neq0$ occur with a nonzero coefficient. So $\eta$ is approximated by
$T_Q=\{\zeta\in T:Q(\zeta)\neq0\}$. Now approximate $\theta$ within $\varepsilon/2$ by a combination
$\sum_\eta c_\eta\eta$ of finitely many members of $S$, and each of these $\eta$ within
$\varepsilon/(2(\sum_\eta|c_\eta|+1))$ by a combination of members of $T_Q$. Substituting the second
combinations into the first and using the triangle inequality approximates $\theta$ within $\varepsilon$
by a combination of members of $T_Q$.

(2) For a polynomial $P$ in the variables $X_j$, $j\in J$, let $P'$ be the polynomial in the variables
$X_{j'}$, $j'\in J'$, obtained by replacing each $X_j$ by $X_{i(j)}$. Then $P(\eta\circ i)=P'(\eta)$ for
every $\eta:J'\to k$, so an approximation of $\theta$ that holds for all polynomials $P'$ gives one of
$\theta\circ i$ for all $P$, with the same coefficients.
\end{proof}

\begin{lemma}\label{gal-congruences-eigen}
Let $\ell$ be a prime, $\iota:\bar{\mathbf Q}_\ell\to\mathbf C$ a field isomorphism, and
$|z|_\iota=|\iota^{-1}(z)|$ for $z\in\mathbf C$, where $\bar{\mathbf Q}_\ell$ carries the absolute value
extending that of $\mathbf Q_\ell$. Let $V$ be a complex vector space, $(\mathrm{ev}_x)_{x\in X}$ a family
of linear forms on $V$, and $(T_j)_{j\in J}$ a family of endomorphisms of $V$ such that for every $v\in V$
and $R\in\mathbf R$, if $|\mathrm{ev}_x(v)|_\iota\leq R$ for all $x$, then
$|\mathrm{ev}_x(T_jv)|_\iota\leq R$ for all $j$ and $x$. Let $v_0\in V$ and $x_0\in X$ with
$T_jv_0=\theta_jv_0$ for all $j$ and $|\mathrm{ev}_{x_0}(v_0)|_\iota=1$, and let $v_1,\ldots,v_k\in V$ with
$T_jv_i=\eta_{i,j}v_i$ for all $i$ and $j$. Let $\varepsilon\geq0$ with
$|\mathrm{ev}_x(v_0+v_1+\cdots+v_k)|_\iota\leq\varepsilon$ for all $x$. Then there are
$c_1,\ldots,c_k\in\bar{\mathbf Q}_\ell$ such that
\[
 \Bigl|P\bigl((\iota^{-1}\theta_j)_j\bigr)-\sum_{i=1}^kc_iP\bigl((\iota^{-1}\eta_{i,j})_j\bigr)\Bigr|
 \leq\varepsilon
\]
for every polynomial $P$ with integer coefficients in variables indexed by $J$.
\end{lemma}
\leanref{Development/AbelianVarietyTheory/Proofs/SatakeParameters/EigenvectorSmallSubFiniteEigenApproxInequalityStatement.lean}{GlobalTheory.T.eigenvector_small_sub_finite_eigen_approx_inequality}
\leanref{Development/AbelianVarietyTheory/Proofs/SatakeParameters/EigenvectorSmallSubFiniteEigenApproxInequality.lean}{GlobalTheory.Cong.eigen_small_sub_finite_approx}

\begin{proof}
For a polynomial $P$ with integer coefficients let $u_P=P(\theta)v_0+\sum_iP(\eta_i)v_i$, where
$\theta=(\theta_j)$ and $\eta_i=(\eta_{i,j})_j$. By induction on $P$, $|\mathrm{ev}_x(u_P)|_\iota\leq\varepsilon$
for all $x$: for a constant $n\in\mathbf Z$, $u_n=n(v_0+\sum_iv_i)$ and $|n|_\iota\leq1$; a sum is handled
by the ultrametric inequality; and $u_{PX_j}=T_ju_P$, which the hypothesis on $T_j$ controls. Now
$\mathrm{ev}_{x_0}(u_P)=\mathrm{ev}_{x_0}(v_0)\bigl(P(\theta)+\sum_i\mathrm{ev}_{x_0}(v_i)\mathrm{ev}_{x_0}(v_0)^{-1}P(\eta_i)\bigr)$.
Apply $\iota^{-1}$, which commutes with $P$ since $P$ has integer coefficients, and use
$|\mathrm{ev}_{x_0}(v_0)|_\iota=1$. This gives the estimate with
$c_i=-\iota^{-1}\bigl(\mathrm{ev}_{x_0}(v_i)/\mathrm{ev}_{x_0}(v_0)\bigr)$.
\end{proof}

\begin{definition}\label{gal-congruences-forms}
Let $E$ be a CM field with maximal totally real subfield $F$ and complex conjugation $c$, let $a\geq1$, and
let $U_a$ be as in Definition~\ref{gal-definite-unitary}. An element $h\in U_a\subset\operatorname{GL}_a(\mathbf
A_E)$ has a component $h_u\in\operatorname{GL}_a(E_u)$ at every place $u$ of $E$; for $x\in U_{a,w}$,
$c_w(x)=x_w$. Write $U_a(\mathbf A_{F,f})$ and $U_a(F_\infty)$ for the subgroups of elements equal to $1$
at the infinite places, respectively at the finite places. An \emph{automorphic form} on $U_a$ is a
continuous function on $U_a$, left invariant under $U_a(F)$ and right invariant under an open subgroup of
$U_a(\mathbf A_{F,f})$, whose right translates by $U_a(F_\infty)$ span a finite-dimensional space. Let
$\lambda=(\lambda_\tau)$ be a family of dominant weights indexed by the embeddings $\tau:E\to\mathbf C$,
for each infinite place $v$ of $E$ let $\tau_v$ be the embedding chosen to represent it, and let
$V_\lambda$ be the tensor product over $v$ of the representations $\nabla(\lambda_{\tau_v})$, on which
$U_a(F_\infty)$ acts through its $\tau_v$-factors.
\begin{enumerate}
\item An \emph{algebraic form} of weight $\lambda$ is a linear map $f$ from $V_\lambda$ to the automorphic
  forms on $U_a$ such that $f(xv)$ is the right translate of $f(v)$ by $x$ for every $x\in U_a(F_\infty)$.
  Write $f(v,g)=f(v)(g)$. For a subgroup $C\subset U_a$, $f$ is \emph{fixed by} $C$ if $f(v,gc)=f(v,g)$
  for all $v$, $g$ and $c\in C$.
\item A \emph{level} is a compact open subgroup of $U_a(\mathbf A_{F,f})$. For a finite place $w_0$ of
  $E$, a \emph{level away from} $w_0$ is a compact subgroup $K\subset U_a(\mathbf A_{F,f})$ whose elements
  are $1$ at the places over the place of $F$ under $w_0$ and which contains a neighbourhood of $1$ in the
  group of such elements. For such $K$ and a compact open subgroup $J\subset\operatorname{GL}_a(E_{w_0})$,
  $K[J]$ is the subgroup generated by $K$ and the $x\in U_{a,w_0}$ with $x_{w_0}\in J$.
\item An automorphic representation $U$ of $U_a$ is \emph{of type} $(\lambda,K)$, for a subgroup
  $K\subset U_a$, if it has weight $\lambda_\tau$ at every $\tau$ and a nonzero vector fixed by $K$.
\item For a subgroup $C\subset U_a$ and a property $P$ of automorphic representations,
  $\mathcal A_\lambda(C,P)$ is the span of the algebraic forms of weight $\lambda$ that are fixed by $C$
  and whose values all lie in one automorphic representation with property $P$.
\item A \emph{Hecke operator} is a finitely supported function $t:U_a\to\mathbf C$. It acts on functions
  on $U_a$ by $(t\varphi)(g)=\sum_ht(h)\varphi(gh)$, and on linear maps $f$ from $V_\lambda$ to functions
  by $(tf)(v)=t(f(v))$. A family $(t_j)_{j\in J}$ of Hecke operators \emph{acts on level $K$ by $\eta$},
  where $\eta$ assigns to each automorphic representation $U$ a family $(\eta_j(U))_{j\in J}$ of complex
  numbers, if $t_j\varphi=\eta_j(U)\varphi$ for every $U$, every $\varphi\in U$ fixed by $K$ and every $j$.
\end{enumerate}
\end{definition}
\leanref{Development/AbelianVarietyTheory/Definitions/AutomorphicMembership/ContinuousMapPFin.lean}{GlobalTheory.DefU.algForms}
\leanref{Development/AbelianVarietyTheory/Definitions/AutomorphicMembership/ContinuousMapPFin.lean}{GlobalTheory.DefU.Wt.ofFam}
\leanref{Development/AutomorphicArithmetic/Foundations/GaloisRepresentations/SeparatingAlgebras/StGlob.lean}{UpperTheory.Gal.Gen.IsLevel}
\leanref{Development/AbelianVarietyTheory/Statements/EllipticCurveLevels/ClosedAtLConclusionOfMem.lean}{GlobalTheory.Cong.IsLevelAway}
\leanref{Development/AutomorphicArithmetic/Foundations/GaloisRepresentations/SeparatingAlgebras/GenDataEmpty.lean}{UpperTheory.Gal.Gen.lev}
\leanref{Development/AbelianVarietyTheory/Statements/EllipticCurveLevels/ClosedAtLConclusionOfMem.lean}{GlobalTheory.Cong.OfType}
\leanref{Development/AutomorphicArithmetic/Foundations/GaloisRepresentations/SeparatingAlgebras/GenDataEmpty.lean}{UpperTheory.Gal.Gen.fixedBy}
\leanref{Development/AutomorphicArithmetic/Foundations/GaloisRepresentations/SeparatingAlgebras/GenDataEmpty.lean}{UpperTheory.Gal.Gen.part}
\leanref{Development/AbelianVarietyTheory/Statements/EllipticCurveLevels/ClosedAtLConclusionOfMem.lean}{GlobalTheory.Cong.ActsBy}

\begin{definition}\label{gal-congruences-dense}
Let $E$, $F$, $c$, $a$, $U_a$, $\lambda$ and $V_\lambda$ be as in
Definition~\ref{gal-congruences-forms}, let $\ell$ be a prime and $\iota:\bar{\mathbf Q}_\ell\to\mathbf C$ a
field isomorphism, and write $|z|_\iota=|\iota^{-1}(z)|$ for $z\in\mathbf C$. An embedding
$\tau:E\to\mathbf C$ \emph{lies over} a place $u$ of $E$ over $\ell$ if $\iota^{-1}\circ\tau$ or
$\iota^{-1}\circ\tau\circ c$ induces $u$.
\begin{enumerate}
\item Let $K\subset U_a$ be a subgroup and $P$ a property of automorphic representations. $P$ is
  \emph{$\ell$-adically dense} in weight $\lambda$ and level $K$ if for every algebraic form $f$ of weight
  $\lambda$ fixed by $K$ and every $\varepsilon>0$ there are $G\in\mathcal A_\lambda(K,P)$ and a function
  $e$ on $U_a$ such that $|e(g)|_\iota\leq\varepsilon$ and $f(v,g)-G(v,g)=e(g)f(v,g)$ for all $v$ and $g$.
\item A Hecke operator $t$ is \emph{$\ell$-integral} for $\lambda$ if for every $h$ in its support,
  $h\in U_a(\mathbf A_{F,f})$ and there is $d\geq0$ such that $\ell^dh_u$ and $\ell^dh_u^{-1}$ have entries
  in $\mathcal O_u$ at every place $u$ over $\ell$, and $|t(h)|_\iota\leq\ell^{-ad\,s(h)}$. Here $s(h)$ is
  the sum of $|\lambda_{\tau,0}|+\cdots+|\lambda_{\tau,a-1}|$ over the pairs $(u,\tau)$ of a place $u$ over
  $\ell$ with $h_u\neq1$ and an embedding $\tau$ that lies over $u$.
\item A \emph{gauge} of weight $\lambda$ for a family $(t_j)$ of Hecke operators is a set $X$ with a map
  $b:U_a\times X\to V_\lambda$ such that: for every algebraic form $f$ of weight $\lambda$ the numbers
  $|f(b(g,x),g)|_\iota$ are bounded; $f=0$ if $f(b(g,x),g)=0$ for all $g$ and $x$; and for every linear
  map $F$ from $V_\lambda$ to the functions on $U_a$ and every $R$, if $|F(b(g,x))(g)|_\iota\leq R$ for all
  $g$ and $x$, then $|(t_jF)(b(g,x))(g)|_\iota\leq R$ for all $j$, $g$ and $x$.
\item For a polynomial $\Phi$ in $a$ variables and a finite place $w$ of $E$, an automorphic
  representation $U$ is \emph{$\Phi$-generic} at $w$ if $\Phi(z_1,\ldots,z_a)\neq0$ whenever $U$ has
  Satake parameters $\{z_1,\ldots,z_a\}$ at $w$, in any order.
\end{enumerate}
\end{definition}
\leanref{Development/AutomorphicArithmetic/Foundations/GaloisRepresentations/SeparatingAlgebras/GenDataEmpty.lean}{UpperTheory.Gal.Gen.Dense}
\leanref{Development/AutomorphicArithmetic/Foundations/GaloisRepresentations/SeparatingAlgebras/CAllFin.lean}{UpperTheory.Gal.Gen.Loc.IntL}
\leanref{Development/AutomorphicArithmetic/Foundations/GaloisRepresentations/SeparatingAlgebras/BlockBoundAlt.lean}{UpperTheory.Gal.Gen.FramedX}
\leanref{Development/AutomorphicArithmetic/Foundations/GaloisRepresentations/SeparatingAlgebras/GenDataEmpty.lean}{UpperTheory.Gal.Gen.GenAt}

\begin{lemma}\label{gal-congruences-gauge}
Let $E$ be a CM field, $a\geq1$, $\lambda$ a family of dominant weights indexed by the embeddings
$E\to\mathbf C$, $\ell$ a prime and $\iota:\bar{\mathbf Q}_\ell\to\mathbf C$ a field isomorphism. Every
family of Hecke operators that are $\ell$-integral for $\lambda$ (Definition~\ref{gal-congruences-dense})
has a gauge of weight $\lambda$.
\end{lemma}
\leanref{Development/AutomorphicArithmetic/Foundations/GaloisRepresentations/SeparatingAlgebras/CAllFin.lean}{UpperTheory.Gal.Gen.Loc.frame}
\leanref{Development/AutomorphicArithmetic/Foundations/GaloisRepresentations/SeparatingAlgebras/CAllFin.lean}{UpperTheory.Gal.Gen.Loc.kappa}
\leanref{Development/AutomorphicArithmetic/Foundations/GaloisRepresentations/SeparatingAlgebras/FrameW.lean}{UpperTheory.Gal.Gen.frameW_pf}

\begin{proof}
For each infinite place $v$, $\iota^{-1}\circ\tau_v$ induces a place $u_v$ of $E$ over $\ell$ and extends
to a continuous embedding $\hat\tau_v:E_{u_v}\to\bar{\mathbf Q}_\ell$. Let
$\kappa(g)=(\iota(\hat\tau_v(g_{u_v})))_v\in\prod_v\operatorname{GL}_a(\mathbf C)$, entry by entry. Then
$\kappa$ is a homomorphism, $\kappa(\gamma)=(\tau_v(\gamma))_v$ for $\gamma\in U_a(F)$, $\kappa$ is trivial
on $U_a(F_\infty)$, $\kappa(g)_v=1$ when $g_{u_v}=1$, and if $\ell^dg_u$ and $\ell^dg_u^{-1}$ are integral at
the places $u$ over $\ell$, the entries of $\kappa(g)_v^{\pm1}$ have $|\cdot|_\iota\leq\ell^d$. Let
$g_\infty\in\prod_v\operatorname{GL}_a(\mathbf C)$ be the family of the $\tau_v$-components of $g$ at
infinity and $M(g)=g_\infty^{-1}\kappa(g)$. Realize each $\nabla(\lambda_{\tau_v})$ in the polynomial
functions on matrices, so that $V_\lambda$ is a space of polynomials with a representation $\rho_\lambda$ of
$\prod_v\operatorname{GL}_a(\mathbf C)$; let $X$ be the set of $x\in V_\lambda$ whose coefficients have
$|\cdot|_\iota\leq1$, and $b(g,x)=\rho_\lambda(M(g))x$.

For $\gamma\in U_a(F)$, $s\in U_a$, $y\in U_a(F_\infty)$ and $c'\in U_a(\mathbf A_{F,f})$ one has
$M(\gamma sy c')=y_\infty^{-1}M(s)\kappa(c')$, so for an algebraic form $f$ fixed by $c'$,
$f(b(\gamma syc',x),\gamma syc')=f(\rho_\lambda(M(s))\rho_\lambda(\kappa(c'))x,s)$. Let $C$ be a level fixing
$f$. Finitely many $s$ represent $U_a(F)\backslash U_a/U_a(F_\infty)C$, the elements of $C$ have bounded
denominators at the places over $\ell$, so $\rho_\lambda(\kappa(c'))x$, $c'\in C$, has bounded
coefficients, and the values $f(b(g,x),g)$ are bounded. Every vector of $V_\lambda$ is a nonzero multiple
of some $\rho_\lambda(M(g))x$ with $x\in X$, which gives the second property. For the third, since
$M(g)=M(gh)\kappa(h)^{-1}$ for $h\in U_a(\mathbf A_{F,f})$,
\[
 (t_jF)(b(g,x))(g)=\sum_hF\bigl(\rho_\lambda(M(gh))\,t_j(h)\rho_\lambda(\kappa(h))^{-1}x\bigr)(gh).
\]
The factor $\kappa(h)_v^{-1}$ differs from $1$ only at the $v$ with $h_{u_v}\neq1$, where its entries have
$|\cdot|_\iota\leq\ell^d$; so the coefficients of $\rho_\lambda(\kappa(h))^{-1}x$ have
$|\cdot|_\iota\leq\ell^{ad\,s(h)}$, and $t_j(h)\rho_\lambda(\kappa(h))^{-1}x$ lies in $X$ by
$\ell$-integrality. The ultrametric inequality gives the bound $R$.
\end{proof}

\begin{lemma}\label{gal-congruences-density}
Let $E$ be a CM field with maximal totally real subfield $F$, $a\geq1$, $\lambda$ a family of dominant
weights indexed by the embeddings
$E\to\mathbf C$, $\ell$ a prime and $\iota:\bar{\mathbf Q}_\ell\to\mathbf C$ a field isomorphism. Let
$K\subset U_a(\mathbf A_{F,f})$ be a subgroup and $P$ a property of automorphic representations of $U_a$
that is $\ell$-adically dense in weight $\lambda$ and level $K$ (Definition~\ref{gal-congruences-dense}).
Let $(t_j)_{j\in J}$ be a family of Hecke operators that has a gauge of weight $\lambda$ and acts on level
$K$ by $\eta$ (Definition~\ref{gal-congruences-forms}). Then for every automorphic representation $U$ of
type $(\lambda,K)$, the family $(\iota^{-1}\eta_j(U))_j$ is approximated, in the sense of
Lemma~\ref{gal-trace-limit}, by the families $(\iota^{-1}\eta_j(U'))_j$ for the automorphic
representations $U'$ of type $(\lambda,K)$ with property $P$.
\end{lemma}
\leanref{Development/AutomorphicArithmetic/Foundations/GaloisRepresentations/SeparatingAlgebras/BlockBoundAlt.lean}{UpperTheory.Gal.Gen.ToPolyX}
\leanref{Development/AutomorphicArithmetic/Foundations/GaloisRepresentations/SeparatingAlgebras/ToPolyX.lean}{UpperTheory.Gal.Gen.toPolyX_pf}
\leanref{Development/AutomorphicArithmetic/Foundations/GaloisRepresentations/SeparatingAlgebras/StGlob.lean}{UpperTheory.Gal.Gen.PartFacts}

\begin{proof}
Since $U$ has weight $\lambda$ and a nonzero vector fixed by $K$, there is a nonzero algebraic form $f$
of weight $\lambda$ fixed by $K$ with values in $U$. Let $(X,b)$ be the gauge. Since $f\neq0$,
$f(b(g_0,x_0),g_0)\neq0$ for some $g_0,x_0$; divide $f$ by this number, so that it becomes $1$, and let
$R\geq1$ bound the $|f(b(g,x),g)|_\iota$. Let $\varepsilon>0$. Density gives $G\in\mathcal A_\lambda(K,P)$
and $e$ with $|e(g)|_\iota\leq\varepsilon/R$ and $f-G=ef$. Write $G=\sum_{\pi\in I}F_\pi$, a finite sum of
nonzero algebraic forms fixed by $K$ with values in automorphic representations $\pi$ with property $P$.
Each such $\pi$ is of type $(\lambda,K)$, since it contains the nonzero $K$-fixed values of $F_\pi$ and
$F_\pi$ has weight $\lambda$. Apply Lemma~\ref{gal-congruences-eigen} to the space of linear maps from
$V_\lambda$ to functions on $U_a$, with $\mathrm{ev}_{(g,x)}(F)=F(b(g,x))(g)$ and $T_j$ the action of $t_j$;
the third property of the gauge is the hypothesis on the $T_j$. Take $v_0=f$, with eigenvalues
$\eta_j(U)$, and the $v_\pi=-F_\pi$, with eigenvalues $\eta_j(\pi)$. Then $v_0+\sum_\pi v_\pi=f-G=ef$, whose
evaluations have $|\cdot|_\iota\leq(\varepsilon/R)R=\varepsilon$. So $(\iota^{-1}\eta_j(U))_j$ is
approximated within $\varepsilon$ by a combination of the $(\iota^{-1}\eta_j(\pi))_j$.
\end{proof}

\begin{lemma}\label{gal-congruences-weight-family}
Let $E$ be a CM field with maximal totally real subfield $F$ and complex conjugation $c$, $a\geq1$, $\ell$
a prime and $\iota:\bar{\mathbf Q}_\ell\to\mathbf C$ a field isomorphism. Let $w_0$ and $v'$ be places of
$E$ over $\ell$, split over $F$, that lie over different places of $F$. Let $K$ be a level away from $w_0$
(Definition~\ref{gal-congruences-forms}) that contains every $x\in U_{a,v'}$ with
$x_{v'}\in\operatorname{GL}_a(\mathcal O_{v'})$, let $K_v\subset\operatorname{GL}_a(E_{v'})$ be an open
subgroup, and let $K'$ be the group of the $k\in K$ with $k_{v'}\in K_v$. Let $\lambda$ be a family of
dominant weights, $f$ an algebraic form of weight $\lambda$ fixed by $K'$, $\tau$ an embedding such that
$\iota^{-1}\circ\tau$ induces $v'$, and $b:U_a\times X\to V_\lambda$ a map such that the numbers
$|F(b(g,x),g)|_\iota$ (Definition~\ref{gal-congruences-dense}) are bounded for every algebraic form $F$ of
weight $\lambda$. Then for every $\varepsilon>0$ there are $d\geq0$ and a map
$M:U_a\to\operatorname{GL}_a(\mathbf C)$, with $M(gh)=M(g)$ for $h\in U_a(\mathbf A_{F,f})$ with
$h_{v'}=1$, such that for every dominant weight $\nu$ of $\operatorname{GL}_a$ there are finitely many
automorphic representations $V$ of $U_a$ and, for each, a map $F_V$ from
$\operatorname{GL}_a(\mathbf C)\times\operatorname{GL}_a(\mathbf C)$ to the linear maps $V_\lambda\to V$ with
values fixed by $K$, such that:
\begin{enumerate}
\item $\bigl|f(b(g,x),g)-\sum_VF_V(M(g),M(g))(b(g,x))(g)\bigr|_\iota\leq\varepsilon$ for all $g$ and $x$;
\item every $V$ with $F_V\neq0$ is of type $(\mu,K)$ for a family $\mu$ of dominant weights with
  $\mu_{\tau'}=\lambda_{\tau'}$ for $\tau'\notin\{\tau,\tau\circ c\}$ and $|\mu_{\tau,i}-\nu_i|\leq2d$ for
  $0\leq i<a$.
\end{enumerate}
\end{lemma}
\leanref{Development/AutomorphicArithmetic/Foundations/GaloisRepresentations/WeightCharacters/DeltaNear.lean}{UpperTheory.Gal.Wt.FamEx}
\leanref{Development/AutomorphicArithmetic/Foundations/GaloisRepresentations/WeightCharacters/DeltaNear.lean}{UpperTheory.Gal.Wt.FamProj}
\leanref{Development/AutomorphicArithmetic/Foundations/GaloisRepresentations/WeightCharacters/DeltaNear.lean}{UpperTheory.Gal.Wt.FamWt}
\leanref{Development/AutomorphicArithmetic/Foundations/GaloisRepresentations/WeightCharacters/FamOfApply.lean}{UpperTheory.Gal.Wt.Fam}
\leanref{Development/AutomorphicArithmetic/Foundations/GaloisRepresentations/WeightCharacters/FamExOf.lean}{UpperTheory.Gal.Wt.famEx_of}
\leanref{Development/AutomorphicArithmetic/Foundations/GaloisRepresentations/WeightCharacters/FamProj.lean}{UpperTheory.Gal.Wt.famProj_of}
\leanref{Development/AutomorphicArithmetic/Foundations/GaloisRepresentations/WeightCharacters/FamOffHolds.lean}{UpperTheory.Gal.Wt.famWt_of}

\begin{proof}
The family is Alpöge's twist of $f$ by the $\ell$-adic component at $v'$.

\emph{Construction.} There is a level $C\supset K$ with $C_{v'}\subset\operatorname{GL}_a(\mathcal
O_{v'})$, such that $f$ is fixed by the $c'\in C$ with $c'_{v'}\in K_v$, and that is neat: $\gamma g=gc'y$
with $\gamma\in U_a(F)$, $c'\in C$ and $y\in U_a(F_\infty)$ forces $c'=1$. Neatness gives a map
$\beta:U_a\to C$ with $\beta(\gamma gyc')=\beta(g)c'$. Let $\kappa:U_a\to\operatorname{GL}_a(\mathbf C)$ be
$g\mapsto\iota(\hat\tau(g_{v'}))$, with $\hat\tau:E_{v'}\to\bar{\mathbf Q}_\ell$ the continuous extension of
$\iota^{-1}\circ\tau$; it equals $\tau$ on $U_a(F)$, is trivial on $U_a(F_\infty)$, is $\ell$-adically
continuous, and has $\kappa(g)^{\pm1}$ with entries of $|\cdot|_\iota\leq1$ when
$g_{v'}\in\operatorname{GL}_a(\mathcal O_{v'})$. As $C$ is compact and $f$ is fixed by an open subgroup
of it, there is $R>0$ with $|f(b(g,x),g\beta(g)^{-1}c')|_\iota\leq R$ for all $g$, $x$ and $c'\in C$. As
$C$ is compact and the $c'\in C$ with $c'_{v'}\in K_v$ form an open subgroup containing the kernel of
$\kappa$ on $C$, the separation of disjoint compact subsets of $\bar{\mathbf Q}_\ell^{a\times a}$ by
polynomials gives finitely many $c_q\in C$ and polynomials $P_q$ in the matrix entries such that for every $c'\in C$ some
$q_0$ has: $(c_{q_0}^{-1}c')_{v'}\in K_v$, $\kappa(c_{q_0}^{-1}c')$ within $\varepsilon/R$ of $1$,
$|P_{q_0}(\kappa(c'))-1|_\iota\leq\varepsilon/R$ and $|P_q(\kappa(c'))|_\iota\leq\varepsilon/R$ for $q\neq q_0$.
The twists $F_q(v,g)=f(v,g\beta(g)^{-1}c_q)$ are algebraic forms of weight $\lambda$ fixed by $C$, with
$|F_q(b(g,x),g)|_\iota\leq R$. Let $\tau(g_\infty)$ be the $\tau$-component of $g$ at infinity,
$M(g)=\tau(g_\infty)^{-1}\kappa(g)$ and $n(g)=\kappa(\beta(g))M(g)^{-1}$. For a dominant $\nu$, let
$\delta_\nu(y)$ be the value at the identity matrix of the translate by $y$ of the highest weight vector of
$\nabla(\nu)$, realized in the polynomial functions on matrices. For $y$ whose entries and those of
$y^{-1}$ have $|\cdot|_\iota\leq1$, $|\delta_\nu(y)|_\iota\leq1$, and $|\delta_\nu(y)-1|_\iota\leq\varepsilon/R$
if moreover $y$ is within $\varepsilon/R$ of $1$. Set
\[
 B(y_1,y_2)(v)(g)=\sum_qP_q(n(g)y_2)\,\delta_\nu\bigl(\kappa(c_q)^{-1}n(g)y_1\bigr)\,F_q(v,g).
\]

\emph{Closeness.} Since $n(g)M(g)=\kappa(\beta(g))$, $B(M(g),M(g))(v)(g)$ is
$\sum_qP_q(\kappa(\beta(g)))\delta_\nu(\kappa(c_q^{-1}\beta(g)))F_q(v,g)$. For the index $q_0$ of $\beta(g)$,
$F_{q_0}(v,g)=f(v,g)$, because $\beta(g)^{-1}c_{q_0}$ fixes $f$. So $f-B(M(g),M(g))$ is the sum over $q$ of
coefficients of $|\cdot|_\iota\leq\varepsilon/R$ times $F_q$, which is at most $\varepsilon$ on $b$ by the
ultrametric inequality.

\emph{The family.} Alpöge's computation shows that the $B(y_1,y_2)(v)$ are automorphic forms fixed by $C$,
spanning a finite-dimensional space, with
$B(y_1,y_2)(\rho_\lambda(y)v)$ the right translate of $B(\tau(y)^{-1}y_1,\tau(y)^{-1}y_2)(v)$ for
$y\in U_a(F_\infty)$, with $B(y_1b,y_2)=\prod_ib_{ii}^{\nu_i}B(y_1,y_2)$ for upper triangular $b$, and, for
some $d$ that depends on $f$ and the $P_q$ but not on $\nu$, with the following exponent bounds. Let
$t_k(z)$ be the diagonal matrix with entries $z$ in the first $k$ places and $1$ elsewhere, for $|z|=1$,
and $y(z)\in U_a(F_\infty)$ the element with $\tau$-component $t_k(z)$ and trivial components at the other
infinite places. The Fourier exponents of $z\mapsto B(t_k(z)y_1,t_k(z)y_2)(\rho_\lambda(y(z))v)(g)$ are at
most $\nu_0+\cdots+\nu_{k-1}+d$, and those of $z\mapsto B(1,t_k(z)y_2)(\rho_\lambda(y(z))v)(g)$ are at least
$-d$. Since the span of the values is finite-dimensional
and consists of automorphic forms, it lies in a finite direct sum of automorphic representations $V$, and
projecting gives $F_V$ with $\sum_VF_V=B$ and the same properties.

\emph{Weights.} Let $V$ with $F_V\neq0$. Its values are fixed by $C\supset K$. At an embedding
$\tau'\notin\{\tau,\tau\circ c\}$ the $\tau'$-factor of $U_a(F_\infty)$ acts on the values of $F_V$ only
through $V_\lambda$, so $V$ has weight $\lambda_{\tau'}$ there. At $\tau$ let $\mu$ be the weight of $V$ and
project to $\nabla(\mu)$. The largest exponent of $t_k$ on a $U(a)$-stable subset of $\nabla(\mu)$ is
$\mu_0+\cdots+\mu_{k-1}$, so the upper bound gives $\mu_0+\cdots+\mu_{k-1}\leq\nu_0+\cdots+\nu_{k-1}+d$,
and the equivariance under upper triangular matrices turns the lower bound into
$\mu_0+\cdots+\mu_{k-1}\geq\nu_0+\cdots+\nu_{k-1}-d$. Taking differences of consecutive $k$ gives
$|\mu_i-\nu_i|\leq2d$.
\end{proof}
\leanref{Development/AutomorphicArithmetic/Foundations/GaloisRepresentations/WeightCharacters/DeltaNear.lean}{UpperTheory.Gal.Wt.NeatLev}
\leanref{Development/AutomorphicArithmetic/Foundations/GaloisRepresentations/WeightCharacters/DeltaNear.lean}{UpperTheory.Gal.Wt.EllC}
\leanref{Development/AutomorphicArithmetic/Foundations/GaloisRepresentations/WeightCharacters/DeltaNear.lean}{UpperTheory.Gal.Wt.Kap0}
\leanref{Development/AutomorphicArithmetic/Foundations/GaloisRepresentations/WeightCharacters/DeltaNear.lean}{UpperTheory.Gal.Wt.Pieces}
\leanref{Development/AutomorphicArithmetic/Foundations/GaloisRepresentations/WeightCharacters/DeltaNear.lean}{UpperTheory.Gal.Wt.Twist}
\leanref{Development/AutomorphicArithmetic/Foundations/GaloisRepresentations/WeightCharacters/DeltaNear.lean}{UpperTheory.Gal.Wt.DeltaNear}
\leanref{Development/AutomorphicArithmetic/Foundations/GaloisRepresentations/WeightCharacters/FamOfApply.lean}{UpperTheory.Gal.Wt.famOf}
\leanref{Development/AutomorphicArithmetic/Foundations/GaloisRepresentations/WeightCharacters/DeltaNear.lean}{UpperTheory.Gal.Wt.FamOf}
\leanref{Development/AutomorphicArithmetic/Foundations/GaloisRepresentations/WeightCharacters/DeltaNear.lean}{UpperTheory.Gal.Wt.FamOff}
\leanref{Development/AutomorphicArithmetic/Foundations/GaloisRepresentations/WeightCharacters/DeltaNear.lean}{UpperTheory.Gal.Wt.FamSum}
\leanref{Development/AutomorphicArithmetic/Foundations/GaloisRepresentations/WeightCharacters/FamSum.lean}{UpperTheory.Gal.Wt.famSum_of}

\begin{lemma}\label{gal-congruences-weights}
Let $E$ be a CM field with maximal totally real subfield $F$ and complex conjugation $c$, $a\geq1$, $\ell$ a
prime and
$\iota:\bar{\mathbf Q}_\ell\to\mathbf C$ a field isomorphism. Let $w_0$ and $v'$ be places of $E$ over
$\ell$, split over $F$, that lie over different places of $F$. Let $K$ be a level away from $w_0$
(Definition~\ref{gal-congruences-forms}) that contains every $x\in U_{a,v'}$ with
$x_{v'}\in\operatorname{GL}_a(\mathcal O_{v'})$, let $K_v\subset\operatorname{GL}_a(E_{v'})$ be an open
subgroup, and let $K'$ be the group of the $k\in K$ with $k_{v'}\in K_v$. Let $\lambda$ be a family of
dominant weights, $\tau_0$ an embedding such that $\iota^{-1}\circ\tau_0$ induces $v'$, and $(t_j)_{j\in J}$
a family of Hecke operators that are $\ell$-integral for $\lambda$ (Definition~\ref{gal-congruences-dense}),
supported on elements equal to $1$ at the places over the places of $F$ under $w_0$ and $v'$, and acting on
level $K'$ by $\eta$. Then for every automorphic representation $U$ of type $(\lambda,K')$, the family
$(\iota^{-1}\eta_j(U))_j$ is approximated, in the sense of Lemma~\ref{gal-trace-limit}, by the families
$(\iota^{-1}\eta_j(U'))_j$ for the automorphic representations $U'$ of type $(\mu,K)$, where $\mu$ runs
through the families of dominant weights with $\mu_\tau=\lambda_\tau$ at every $\tau$ that does not lie over
$v'$ and for which the elements $j_0<\cdots<j_{a-1}$ of $\{k-\mu_{\tau_0,k}:0\leq k<a\}$ satisfy
$j_{i+1}-j_0>a(j_i-j_0)$ for $i+1<a$. In particular this set is lacunary (Definition~\ref{gal-seed}).
\end{lemma}
\leanref{Development/AutomorphicArithmetic/Foundations/GaloisRepresentations/IntegralStructures/CongruenceOddLevel/WtDefs.lean}{UpperTheory.Gal.Int.Cong.WtLacEx}
\leanref{Development/AutomorphicArithmetic/Foundations/GaloisRepresentations/IntegralStructures/CongruenceOddLevel/WtDefs.lean}{UpperTheory.Gal.Int.Cong.GapWt}
\leanref{Development/AutomorphicArithmetic/Foundations/GaloisRepresentations/WeightCharacters/TopOf.lean}{UpperTheory.Gal.Wt.top_of}
\leanref{Development/AutomorphicArithmetic/Foundations/GaloisRepresentations/WeightCharacters/LadicWeightEmbeddingOverPlace.lean}{UpperTheory.Gal.Wt.wtLacEx_holds}

\begin{proof}
Let $\varepsilon>0$. As in the proof of Lemma~\ref{gal-congruences-density}, let $f$ be a nonzero algebraic
form of weight $\lambda$ fixed by $K'$ with values in $U$, let $(X,b)$ be a gauge of weight $\lambda$ for
$(t_j)$ (Lemma~\ref{gal-congruences-gauge}), and normalize $f$ so that $f(b(g_0,x_0),g_0)=1$.
Lemma~\ref{gal-congruences-weight-family}, with $\tau=\tau_0$, gives $d$ and $M$. Let $N=(a+1)(4d+a)+1$ and
$\nu_k=-N(a+1)^k$, a dominant weight, and let $V$ and $F_V$ be as in
Lemma~\ref{gal-congruences-weight-family} for $\nu$. Put $F'_V(v)(g)=F_V(M(g),M(g))(v)(g)$. Each $t_j$ is
supported on elements $h\in U_a(\mathbf A_{F,f})$ with $h_{v'}=1$, so $M(gh)=M(g)$, and $t_jF'_V(v)(g)$ is
the value at $g$ of $t_j$ applied to $F_V(M(g),M(g))(v)$, a $K'$-fixed vector of $V$; hence
$t_jF'_V=\eta_j(V)F'_V$. Apply Lemma~\ref{gal-congruences-eigen} as in the proof of
Lemma~\ref{gal-congruences-density}, with $v_0=f$ and the $v_V=-F'_V$: by
Lemma~\ref{gal-congruences-weight-family}(1), $f-\sum_VF'_V$ has evaluations of $|\cdot|_\iota\leq\varepsilon$.
So $(\iota^{-1}\eta_j(U))_j$ is approximated within $\varepsilon$ by the $(\iota^{-1}\eta_j(V))_j$ with
$F_V\neq0$. By Lemma~\ref{gal-congruences-weight-family}(2) each such $V$ is of type $(\mu,K)$ with $\mu$
equal to $\lambda$ away from $\tau_0$ and $\tau_0\circ c$, which lie over $v'$, and with
$|\mu_{\tau_0,k}-\nu_k|\leq2d$. The integers $j_k=k-\mu_{\tau_0,k}$ are increasing in $k$ and differ from
$N(a+1)^k$ by at most $2d+a$; since $N>(a+1)(4d+a)$, a direct computation gives
$j_{i+1}-j_0>a(j_i-j_0)$. This gap condition implies lacunarity: of two disjoint nonempty sets of indices,
the one containing the largest index $m$ has average at least $(j_m-j_0)/a+j_0$, while every element of the
other is at most $j_{m-1}<(j_m-j_0)/a+j_0$.
\end{proof}
\leanref{Development/AutomorphicArithmetic/Foundations/GaloisRepresentations/IntegralStructures/CongruenceOddLevel/LacunaryBaseWeightAntitone.lean}{UpperTheory.Gal.Int.Cong.mu0}
\leanref{Development/AutomorphicArithmetic/Foundations/GaloisRepresentations/IntegralStructures/CongruenceOddLevel/LacunaryBaseWeightAntitone.lean}{UpperTheory.Gal.Int.Cong.gapJ_of_near}

\begin{lemma}\label{gal-congruences-section}
Let $E$ be a CM field with maximal totally real subfield $F$, $a\geq1$, $C$ a level (Definition~\ref{gal-congruences-forms}), $H$ a group and
$\psi:C\to H$ a homomorphism that is trivial on a neighbourhood of $1$ and on the elements of finite order.
Then there is a map $\beta:U_a\to H$ with $\beta(\gamma gyc')=\beta(g)\psi(c')$ for $\gamma\in U_a(F)$,
$g\in U_a$, $y\in U_a(F_\infty)$ and $c'\in C$, and for every function $\phi:H\to\mathbf C$, $\phi\circ\beta$
is an automorphic form on $U_a$.
\end{lemma}
\leanref{Development/AutomorphicArithmetic/Foundations/GaloisRepresentations/SeparatingAlgebras/StGlob.lean}{UpperTheory.Gal.Gen.Ell}
\leanref{Development/AutomorphicArithmetic/Foundations/GaloisRepresentations/SeparatingAlgebras/StGlob.lean}{UpperTheory.Gal.Gen.EllGrp}
\leanref{Development/AutomorphicArithmetic/Foundations/GaloisRepresentations/SeparatingAlgebras/EllPf.lean}{UpperTheory.Gal.Gen.ell_pf}
\leanref{Development/AutomorphicArithmetic/Foundations/GaloisRepresentations/SeparatingAlgebras/EllGrp.lean}{UpperTheory.Gal.Gen.ellGrp_pf}

\begin{proof}
Let $D=CU_a(F_\infty)$, a compact subgroup, and extend $\psi$ to $D$ trivially on $U_a(F_\infty)$. If
$\gamma g=gd$ with $\gamma\in U_a(F)$ and $d\in D$, then $d$ lies in $D\cap g^{-1}U_a(F)g$, a compact
subset of a discrete group, so $d$ has finite order and $\psi(d)=1$. Hence if $\gamma_1rd_1=\gamma_2rd_2$
then $\psi(d_1)=\psi(d_2)$. Choose a representative $r$ in each double coset $U_a(F)gD$, write $g=\gamma rd$
and put $\beta(g)=\psi(d)$; this is well defined and has the stated equivariance. For $\phi:H\to\mathbf C$,
$\phi\circ\beta$ is left invariant under $U_a(F)$, right invariant under $U_a(F_\infty)$, and right invariant
under the open subgroup of $C$ on which $\psi$ is trivial, so it is locally constant and an automorphic form.
\end{proof}

\begin{lemma}\label{gal-congruences-type-character}
Let $E$ be a CM field, $w$ a finite place of $E$ over a prime $\ell$, $a\geq1$, $N$ a neighbourhood of
$1$ in $\operatorname{GL}_a(E_w)$ and $m_0\geq0$. Then there are $m\geq m_0$, a compact open subgroup
$J\subset N$ whose only element of finite order is $1$, and a surjective homomorphism
$\vartheta:J\to\mathbf Z/\ell^m$, such that every irreducible smooth representation of
$\operatorname{GL}_a(E_w)$ with a nonzero vector on which $J$ acts by $\chi\circ\vartheta$, for a nontrivial
character $\chi$ of $\mathbf Z/\ell^m$, is supercuspidal.
\end{lemma}
\leanref{Development/AutomorphicArithmetic/Foundations/GaloisRepresentations/IntegralStructures/CongruenceOddLevel/NotApartOfOverPair.lean}{UpperTheory.Gal.Int.Cong.EllSt}
\leanref{Development/AutomorphicArithmetic/Foundations/GaloisRepresentations/IntegralStructures/CongruenceOddLevel/ExistsEllType.lean}{UpperTheory.Gal.Int.Cong.exists_ellType}
\leanref{Development/AbelianVarietyTheory/Proofs/DivisorProofs/ExistsTorsionFreeSurjZmod.lean}{GlobalTheory.CongW.TypeExists.exists_type_surj}

\begin{proof}
If $a=1$, $\operatorname{GL}_1(E_w)$ has no proper parabolic subgroup, every irreducible representation is
supercuspidal, and any small torsion-free congruence subgroup with a surjection to $\mathbf Z/\ell^{m_0}$
will do. Let $a\geq2$, $\mathcal O=\mathcal O_w$ with residue field $k$, and choose $m\geq\max(m_0,2)$ such
that $J=\ker(\operatorname{GL}_a(\mathcal O)\to\operatorname{GL}_a(\mathcal O/\ell^m\mathcal O))$ lies in
$N$; as $m\geq2$, $J$ has no torsion. Choose $B_0\in M_a(k)$ with no invariant subspace other than $0$ and
$k^a$, a lift $B\in M_a(\mathcal O)$, and an additive map $\varphi:\mathcal O\to\mathbf Z/\ell^m$ with
$\varphi(1)=1$ that vanishes on $\ell^m\mathcal O$. Put $\vartheta(1+\ell^mx)=\varphi(\operatorname{tr}(Bx))$;
since $(1+\ell^mx)(1+\ell^my)=1+\ell^m(x+y+\ell^mxy)$, this is a homomorphism.

Let $W$ be a subspace of $E_w^a$ other than $0$ and $E_w^a$. Since $B_0$ has no invariant subspace, there
are $p\in W\cap\mathcal O^a$ and a linear form $q$ on $\mathcal O^a$ vanishing on $W$ with $q(Bp)$ a unit.
Let $x=q(Bp)^{-1}pq$, of rank one with $x^2=0$ and $\operatorname{tr}(Bx)=1$, and $u=1+\ell^mx\in J$. Then
$\vartheta(u)=1$ generates $\mathbf Z/\ell^m$, so $\vartheta$ is surjective; and $u$ is trivial on $W$ and
on $E_w^a/W$, so it lies in the unipotent radical of every parabolic subgroup stabilizing a flag whose
smallest nonzero member is $W$.

Let $\sigma$ be irreducible smooth with $v\neq0$ on which $J$ acts by $\chi\circ\vartheta$, $\chi\neq1$. Let
$P$ be a proper parabolic subgroup with unipotent radical $N_P$ and $g\in\operatorname{GL}_a(E_w)$. By the
previous paragraph $\chi\circ\vartheta$ is nontrivial on the compact open subgroup $J\cap gN_Pg^{-1}$ of
$gN_Pg^{-1}$, so $\int_{J\cap gN_Pg^{-1}}\sigma(n)v\,dn=0$ and $v$ maps to $0$ in the Jacquet module for
$gN_Pg^{-1}$ \cite{BZ76}. Equivalently $\sigma(g)^{-1}v$ maps to $0$ in the Jacquet module $\sigma_{N_P}$.
These vectors span $\sigma$, so $\sigma_{N_P}=0$ for every proper $P$, and $\sigma$ is supercuspidal.
\end{proof}

\begin{lemma}\label{gal-congruences-type}
Let $E$ be a CM field with maximal totally real subfield $F$, $a\geq1$, $\lambda$ a family of dominant
weights, $\ell$ a prime and $\iota:\bar{\mathbf Q}_\ell\to\mathbf C$ a field isomorphism. Let $w_0$ be a
place of $E$ split over $F$, $K$ a level away from $w_0$ (Definition~\ref{gal-congruences-forms}),
$J\subset\operatorname{GL}_a(E_{w_0})$ a compact open subgroup whose only element of finite order is $1$,
$m\geq0$ and $\vartheta:J\to\mathbf Z/\ell^m$ a homomorphism. Let $P_\vartheta$ be the property of an
automorphic representation $U$ that, for some nontrivial character $\chi$ of $\mathbf Z/\ell^m$, some
nonzero $\varphi\in U$ fixed by $K$ has right translate $\chi(\vartheta(x_{w_0}))\varphi$ by every $x\in
U_{a,w_0}$ with $x_{w_0}\in J$. Then for every algebraic form $f$ of weight $\lambda$ fixed by $K[J]$
there are $G\in\mathcal A_\lambda(K,P_\vartheta)$ and a function $e$ on $U_a$ with $|e(g)|_\iota\leq\ell^{-m}$
(Definition~\ref{gal-congruences-dense}) and $f-G=ef$.
\end{lemma}
\leanref{Development/AutomorphicArithmetic/Foundations/GaloisRepresentations/SeparatingAlgebras/BlockBoundAlt.lean}{UpperTheory.Gal.Gen.TyStep}
\leanref{Development/AutomorphicArithmetic/Foundations/GaloisRepresentations/SeparatingAlgebras/TyStep.lean}{UpperTheory.Gal.Gen.tyStep_pf}
\leanref{Development/AbelianVarietyTheory/Definitions/Divisors/IsIntegralOpW.lean}{GlobalTheory.CongW.HasTypeVec}

\begin{proof}
$K[J]$ is a level, and every $c'\in K[J]$ has $c'_{w_0}\in J$. Let $\psi(c')=\vartheta(c'_{w_0})$. The
$\ell^m$-th powers of elements of $J$ contain a neighbourhood of $1$ and are killed by $\vartheta$, so
$\psi$ is trivial near $1$; it is trivial on elements of finite order since $J$ has none but $1$.
Lemma~\ref{gal-congruences-section} gives $\beta:U_a\to\mathbf Z/\ell^m$ with
$\beta(\gamma gyc')=\beta(g)+\psi(c')$. For a character $\chi$ of $\mathbf Z/\ell^m$,
$F_\chi(v,g)=\chi(\beta(g))f(v,g)$ is an algebraic form of weight $\lambda$ with
$F_\chi(v,gc')=\chi(\psi(c'))F_\chi(v,g)$ for $c'\in K[J]$. Decompose $F_\chi$ into its components in
automorphic representations: each component satisfies the same rule, is fixed by $K$ since $\psi$ is trivial
on $K$, and has $J$ acting through $\chi\circ\vartheta$. So $F_\chi\in\mathcal A_\lambda(K,P_\vartheta)$ for
$\chi\neq1$. Let $G=-\sum_{\chi\neq1}F_\chi$ and $e(g)=\sum_\chi\chi(\beta(g))$, the sum over all characters,
which is $\ell^m$ if $\beta(g)=0$ and $0$ otherwise. Then $f-G=ef$ and $|e(g)|_\iota\leq\ell^{-m}$.
\end{proof}
\leanref{Development/AutomorphicArithmetic/Foundations/GaloisRepresentations/SeparatingAlgebras/BlockBoundAlt.lean}{UpperTheory.Gal.Gen.KerOpen}
\leanref{Development/AutomorphicArithmetic/Foundations/GaloisRepresentations/SeparatingAlgebras/BlockBoundAlt.lean}{UpperTheory.Gal.Gen.CharSum}
\leanref{Development/AutomorphicArithmetic/Foundations/GaloisRepresentations/SeparatingAlgebras/BlockBoundAlt.lean}{UpperTheory.Gal.Gen.TwistPart}

\begin{lemma}\label{gal-congruences-l-group}
Let $\ell$ be a prime, $\iota:\bar{\mathbf Q}_\ell\to\mathbf C$ a field isomorphism, $H$ a finite
$\ell$-group, and $P$ a space of functions $H\to\mathbf C$, stable under right translations, with
$|H|\leq2\dim P$. Then there is $c\in P$ with $|1-c(h)|_\iota^2\leq\ell^{-1}$ for every $h\in H$
(Definition~\ref{gal-congruences-dense}).
\end{lemma}
\leanref{Development/AutomorphicArithmetic/Foundations/GaloisRepresentations/SeparatingAlgebras/StGlob.lean}{UpperTheory.Gal.Gen.AlgC}
\leanref{Development/AutomorphicArithmetic/Foundations/GaloisRepresentations/SeparatingAlgebras/AlgCPf.lean}{UpperTheory.Gal.Gen.algC_pf}
\leanref{Development/AutomorphicArithmetic/Foundations/GaloisRepresentations/SeparatingAlgebras/GroupDetAlt.lean}{UpperTheory.Gal.Gen.Alg}
\leanref{Development/AutomorphicArithmetic/Foundations/GaloisRepresentations/SeparatingAlgebras/AlgPf.lean}{UpperTheory.Gal.Gen.alg_pf}

\begin{proof}
Apply $\iota^{-1}$, so that $P$ becomes a space of functions $H\to\bar{\mathbf Q}_\ell$. By ultrametric
duality on $\bar{\mathbf Q}_\ell^H$ with the maximum norm, it suffices that $|\sum_ha(h)|\leq\ell^{-1/2}$ for
every $a:H\to\bar{\mathbf Q}_\ell$ with $|a(h)|\leq1$ and $\sum_ha(h)c(h)=0$ for all $c\in P$. For such $a$,
the matrix $A=(a(hg^{-1}))_{h,g}$ satisfies $\sum_ha(hg^{-1})c(h)=\sum_ia(i)c(ig)=0$, so its transpose kills
the translation-stable space $P$ and $\operatorname{rank}A\leq|H|-\dim P$. The group determinant bound gives
$|\sum_ha(h)|^{|H|}\leq\ell^{-(|H|-\operatorname{rank}A)}\leq\ell^{-\dim P}$. Indeed, $A=\sum_ka(k)\rho(k)$
for the regular representation $\rho$, and since $H$ is an $\ell$-group there are integral matrices $Q$ and
$Q'$ with $Q\rho(k)Q'$ upper unitriangular modulo $\ell$ for every $k$; so $A'=QAQ'$, of rank at most
$\operatorname{rank}A$, is $(\sum_ha(h))\cdot1+U+\ell B'$ with $U$ strictly upper triangular and $B'$
integral. The determinant of $A'-\ell B'$ is $(\sum_ha(h))^{|H|}$, and expanding it by columns, each nonzero
term takes at least $|H|-\operatorname{rank}A'$ columns from $\ell B'$. Squaring, $|\sum_ha(h)|^{2|H|}\leq\ell^{-2\dim P}\leq
\ell^{-|H|}$, which is the claim.
\end{proof}
\leanref{Development/AutomorphicArithmetic/Foundations/GaloisRepresentations/SeparatingAlgebras/GroupDetAlt.lean}{UpperTheory.Gal.Gen.HB}
\leanref{Development/AutomorphicArithmetic/Foundations/GaloisRepresentations/SeparatingAlgebras/GroupDetAlt.lean}{UpperTheory.Gal.Gen.GroupDet}
\leanref{Development/AutomorphicArithmetic/Foundations/GaloisRepresentations/SeparatingAlgebras/GroupDetAlt.lean}{UpperTheory.Gal.Gen.TriMat}
\leanref{Development/AutomorphicArithmetic/Foundations/GaloisRepresentations/SeparatingAlgebras/GroupDetAlt.lean}{UpperTheory.Gal.Gen.DetLow}

\begin{lemma}\label{gal-congruences-negligible}
Let $E$ be a CM field with maximal totally real subfield $F$ and $a\geq1$. Let $w_0$ and $w$ be places of
$E$, split over $F$, that lie over different places of $F$. Let $K$ be a level away from $w_0$
(Definition~\ref{gal-congruences-forms}) such that an $x\in U_{a,w}$ lies in $K$ if and only if
$x_w\in\operatorname{GL}_a(\mathcal O_w)$, and let $\Phi\neq0$ be a polynomial in $a$ variables. Then for
every $\delta>0$ there is a Hecke operator $h$ supported on $U_{a,w}$, with $g^{-1}g'\notin K$ for distinct
$g,g'$ in its support and $\sum_g|h(g)|^2\leq\delta$, that acts as the identity on the vectors fixed by $K$
of every automorphic representation that is not $\Phi$-generic at $w$
(Definition~\ref{gal-congruences-dense}).
\end{lemma}
\leanref{Development/AutomorphicArithmetic/Foundations/GaloisRepresentations/SeparatingAlgebras/StGlob.lean}{UpperTheory.Gal.Gen.NeglAt}
\leanref{Development/AutomorphicArithmetic/Foundations/GaloisRepresentations/SeparatingAlgebras/GenDataEmpty.lean}{UpperTheory.Gal.Gen.Negl}
\leanref{Development/AutomorphicArithmetic/Foundations/GaloisRepresentations/SeparatingAlgebras/NeglAt.lean}{UpperTheory.Gal.Gen.neglAt_pf}

\begin{proof}
Let $\mathbf T$ be the torus of $(z_1,\ldots,z_a)$ with $|z_i|=1$, with its Haar probability measure.

\emph{A small symmetric polynomial.} For every $\delta'>0$ there is a symmetric Laurent polynomial $p$ with
$p(z)=1$ whenever $\Phi(z)=0$ and $\int_{\mathbf T}|p|^2\leq\delta'$. Indeed, $\Pi(z)=\prod_\sigma\Phi(z_{\sigma(1)},\ldots,z_{\sigma(a)})$,
the product over permutations, is a nonzero symmetric polynomial, and there is a symmetric Laurent
polynomial $\Pi^*$ equal to $\bar\Pi$ on $\mathbf T$. With $\Pi$ bounded by $M_0$ on $\mathbf T$ and
$c_0=1/(M_0^2+1)$, let $p_k=(1-c_0\Pi\Pi^*)^k$. Then $p_k=1$ where $\Phi=0$, and on $\mathbf T$,
$p_k=(1-c_0|\Pi|^2)^k$ lies in $[0,1]$ and tends to $0$ off the zero set of $\Pi$, which has measure $0$. By
dominated convergence $\int_{\mathbf T}|p_k|^2\to0$.

\emph{Plancherel.} By the Satake isomorphism \cite{Car79}, every symmetric Laurent polynomial $p$ is the
Satake transform of a unique bi-$\operatorname{GL}_a(\mathcal O_w)$-invariant compactly supported function
$h_w$ on $\operatorname{GL}_a(E_w)$, which acts on the spherical vector of a representation with Satake
parameters $z$ by $p(z)$. By Macdonald's Plancherel formula for spherical functions \cite{Mac71}, the
$L^2$-norm of $h_w$ is the integral of $|p|^2$ against a density on $\mathbf T$, and this density is
bounded; so there is $C$, depending on $w$ and $a$, with $\sum_g|h_w(g)|^2\leq C\int_{\mathbf T}|p|^2$, the
sum over representatives $g$ of the support modulo $\operatorname{GL}_a(\mathcal O_w)$. Transporting
$h_w$ to $U_{a,w}$ through $x\mapsto x_w$ gives a Hecke operator $h$ with the same support properties that
acts by $p(z)$ on the $K$-fixed vectors of every automorphic representation with Satake parameters $z$ at $w$.

Take $p$ from the first step with $\delta'=\delta/(C+1)$. Then $\sum_g|h(g)|^2\leq\delta$, and on a
representation that is not $\Phi$-generic at $w$, with Satake parameters $z$ such that $\Phi(z)=0$, $h$ acts
by $p(z)=1$.
\end{proof}
\leanref{Development/AutomorphicArithmetic/Foundations/GaloisRepresentations/SeparatingAlgebras/GammaConjAt.lean}{UpperTheory.Gal.Gen.Small}
\leanref{Development/AutomorphicArithmetic/Foundations/GaloisRepresentations/SeparatingAlgebras/SmallPf.lean}{UpperTheory.Gal.Gen.small_pf}
\leanref{Development/AutomorphicArithmetic/Foundations/GaloisRepresentations/SeparatingAlgebras/StGlob.lean}{UpperTheory.Gal.Gen.GlobOp}
\leanref{Development/AutomorphicArithmetic/Foundations/GaloisRepresentations/SeparatingAlgebras/StGlob.lean}{UpperTheory.Gal.Gen.LocOp}
\leanref{Development/AutomorphicArithmetic/Foundations/GaloisRepresentations/SeparatingAlgebras/GammaConjAt.lean}{UpperTheory.Gal.Gen.PlanchAt}

\begin{lemma}\label{gal-congruences-sparse}
Let $E$ be a CM field with maximal totally real subfield $F$, $a\geq1$ and $\lambda$ a family of dominant
weights. Let $w_0$ be a place of $E$ split over $F$, $K$ a level away from $w_0$
(Definition~\ref{gal-congruences-forms}), and $J_0\supset J_1\supset\cdots$ compact open subgroups of
$\operatorname{GL}_a(E_{w_0})$, normal in $J_0$, forming a basis of neighbourhoods of $1$. Let $S$ be a
finite set of places of $E$, split over $F$ and over places of $F$ different from the one under $w_0$, such
that for $w\in S$ an $x\in U_{a,w}$ lies in $K$ if and only if $x_w\in\operatorname{GL}_a(\mathcal O_w)$,
and let $\Phi_w\neq0$, $w\in S$, be polynomials in $a$ variables. Let $\bar P_S$ be the property of not
being $\Phi_w$-generic at some $w\in S$ (Definition~\ref{gal-congruences-dense}). Then for some $n$, the
space $\mathcal A_\lambda(K[J_n],\bar P_S)$ is finite-dimensional and
$2\dim\mathcal A_\lambda(K[J_n],\bar P_S)\leq[J_0:J_n]$.
\end{lemma}
\leanref{Development/AutomorphicArithmetic/Foundations/GaloisRepresentations/SeparatingAlgebras/StGlob.lean}{UpperTheory.Gal.Gen.CountGen}
\leanref{Development/AutomorphicArithmetic/Foundations/GaloisRepresentations/SeparatingAlgebras/CountGen.lean}{UpperTheory.Gal.Gen.countGen_pf}
\leanref{Development/AutomorphicArithmetic/Foundations/GaloisRepresentations/SeparatingAlgebras/StGlob.lean}{UpperTheory.Gal.Gen.Count}
\leanref{Development/AutomorphicArithmetic/Foundations/GaloisRepresentations/SeparatingAlgebras/OpOfFinsetSum.lean}{UpperTheory.Gal.Gen.count_pf}

\begin{proof}
Let $\varepsilon=1/(2(|S|+1))$. Every form in $\mathcal A_\lambda(K[J_n],\bar P_S)$ is a sum of forms in
$\mathcal A_\lambda(K[J_n],\bar P_w)$, $w\in S$, where $\bar P_w$ is not being $\Phi_w$-generic at $w$. It
suffices to show that, for each $w\in S$ and all large $n$, $\dim\mathcal A_\lambda(K[J_n],\bar P_w)\leq
\varepsilon[J_0:J_n]$, since then the dimension of the whole space is at most
$|S|\varepsilon[J_0:J_n]\leq[J_0:J_n]/2$.

\emph{Covering.} By the finiteness of class numbers there is a finite set $S'\subset U_a$ such that, with
$r_q\in U_{a,w_0}$ whose components at $w_0$ represent $J_0/J_n$, every $g\in U_a$ is $\gamma sr_qyc'$ with
$\gamma\in U_a(F)$, $s\in S'$, $y\in U_a(F_\infty)$ and $c'\in K[J_n]$.

\emph{Hilbert--Schmidt bound.} There is $M_\lambda$, depending on $\lambda$, with the following property.
Let $C$ be a level and $(p_i)_{i\in Y}$ finitely many points with $U_a=\bigcup_iU_a(F)p_iU_a(F_\infty)C$, and
let $h$ be a Hecke operator whose support is separated at these points: $p_ig'=\gamma p_igyc'$ with $g,g'$ in
the support, $\gamma\in U_a(F)$, $y\in U_a(F_\infty)$ and $c'\in C$ forces $g=g'$. Then every space of
$C$-fixed algebraic forms of weight $\lambda$ on which $h$ acts as the identity has dimension at most
$M_\lambda|Y|\sum_g|h(g)|^2$. Indeed the values $f(e_\alpha,p_i)$, for a basis $(e_\alpha)$ of $V_\lambda$,
determine such a form $f$; $h$ acts on these values by a matrix whose entries are sums of $h(g)$ times matrix
coefficients of the compact group $U_a(F_\infty)$ on $V_\lambda$, each $h(g)$ entering each row once by the
separation; so its Hilbert--Schmidt norm squared is at most $M_\lambda|Y|\sum_g|h(g)|^2$; and a matrix that
is the identity on a subspace has Hilbert--Schmidt norm squared at least the dimension of the subspace.

\emph{Separation.} For fixed finite $S'$ and a finite set $T'$ of elements of $U_a(\mathbf A_{F,f})$ equal to
$1$ at the places over the place of $F$ under $w_0$, with $g^{-1}g'\notin K$ for distinct $g,g'\in T'$, the
points $sr_q$ are separated for $T'$ and the level $K[J_n]$ once $n$ is large. The elements of $U_a(F)$ that
can occur lie in a fixed compact set, hence form a finite set, and the ones different from $1$ have
components at $w_0$ that avoid $J_n$ for large $n$.

Now let $w\in S$, take $\delta=\varepsilon/(M_\lambda|S'|+1)$ and let $h$ be the operator of
Lemma~\ref{gal-congruences-negligible} at $w$ for $\Phi_w$. It acts as the identity on every form in
$\mathcal A_\lambda(K[J_n],\bar P_w)$, since such a form is a sum of forms with values in representations
not $\Phi_w$-generic at $w$ and fixed by $K\subset K[J_n]$. With the covering by the $|S'|[J_0:J_n]$ points
$sr_q$, separated for $n$ large, the Hilbert--Schmidt bound gives
$\dim\mathcal A_\lambda(K[J_n],\bar P_w)\leq M_\lambda|S'|[J_0:J_n]\delta\leq\varepsilon[J_0:J_n]$.
\end{proof}
\leanref{Development/AutomorphicArithmetic/Foundations/GaloisRepresentations/SeparatingAlgebras/StGlob.lean}{UpperTheory.Gal.Gen.CountCore}
\leanref{Development/AutomorphicArithmetic/Foundations/GaloisRepresentations/SeparatingAlgebras/GroupDetAlt.lean}{UpperTheory.Gal.Gen.HS}
\leanref{Development/AutomorphicArithmetic/Foundations/GaloisRepresentations/SeparatingAlgebras/StGlob.lean}{UpperTheory.Gal.Gen.Cover}
\leanref{Development/AutomorphicArithmetic/Foundations/GaloisRepresentations/SeparatingAlgebras/StGlob.lean}{UpperTheory.Gal.Gen.Sep}

\begin{lemma}\label{gal-congruences-generic}
Let $E$ be a CM field with maximal totally real subfield $F$, $a\geq1$, $\lambda$ a family of dominant
weights, $\ell$ a prime and $\iota:\bar{\mathbf Q}_\ell\to\mathbf C$ a field isomorphism. Let $w_0$ be a
place of $E$ over $\ell$, split over $F$, and $K$ a level away from $w_0$
(Definition~\ref{gal-congruences-forms}). Let $S$ be a finite set of places of $E$, split over $F$ and over
places of $F$ different from the one under $w_0$, such that for $w\in S$ an $x\in U_{a,w}$ lies in $K$ if
and only if $x_w\in\operatorname{GL}_a(\mathcal O_w)$, and let $\Phi_w\neq0$, $w\in S$, be polynomials in
$a$ variables. Then the property of being $\Phi_w$-generic at every $w\in S$ is $\ell$-adically dense in
weight $\lambda$ and level $K$ (Definition~\ref{gal-congruences-dense}).
\end{lemma}
\leanref{Development/AutomorphicArithmetic/Foundations/GaloisRepresentations/SeparatingAlgebras/GenDataEmpty.lean}{UpperTheory.Gal.Gen.DenseGen}
\leanref{Development/AutomorphicArithmetic/Foundations/GaloisRepresentations/SeparatingAlgebras/GenDataEmpty.lean}{UpperTheory.Gal.Gen.GenData}
\leanref{Development/AutomorphicArithmetic/Foundations/GaloisRepresentations/SeparatingAlgebras/DenseGen.lean}{UpperTheory.Gal.Gen.denseGen_pf}
\leanref{Development/AutomorphicArithmetic/Foundations/GaloisRepresentations/SeparatingAlgebras/StGlob.lean}{UpperTheory.Gal.Gen.Step}
\leanref{Development/AutomorphicArithmetic/Foundations/GaloisRepresentations/SeparatingAlgebras/StepPf.lean}{UpperTheory.Gal.Gen.step_pf}

\begin{proof}
Let $P$ be the property of being $\Phi_w$-generic at every $w\in S$.

\emph{One step.} For every algebraic form $f$ of weight $\lambda$ fixed by $K$ there are
$G\in\mathcal A_\lambda(K,P)$ and $e$ with $|e(g)|_\iota^2\leq\ell^{-1}$ and $f-G=ef$. Indeed, $f$ is fixed by
the $x\in U_{a,w_0}$ with $x_{w_0}$ in some neighbourhood $N$ of $1$. Let $J_n$ be the elements of
$\operatorname{GL}_a(\mathcal O_{w_0})$ congruent to $1$ modulo $\ell^{s+n+1}$, with $s$ so large that
$J_0\subset N$ and $J_0$ has no element of finite order but $1$. Then $x^{\ell^n}\in J_n$ for $x\in J_0$, so
$H=J_0/J_n$ is an $\ell$-group. Lemma~\ref{gal-congruences-sparse} gives $n$ with
$2\dim\mathcal A_\lambda(K[J_n],\bar P)\leq|H|$, where $\bar P$ is the negation of $P$. The map
$\psi:K[J_0]\to H$, $c'\mapsto c'_{w_0}J_n$, is surjective with kernel the open subgroup $K[J_n]$, and it is
trivial on elements of finite order since $J_0$ has none but $1$; Lemma~\ref{gal-congruences-section} gives
$\beta:U_a\to H$. For $\phi:H\to\mathbf
C$ let $\Psi(\phi)(v,g)=\phi(\beta(g))f(v,g)$, an algebraic form of weight $\lambda$ fixed by $K[J_n]$, as $f$
is fixed by $K[J_0]$. Every form fixed by $K[J_n]$ lies in
$\mathcal A_\lambda(K[J_n],P)+\mathcal A_\lambda(K[J_n],\bar P)$, so the space $P_H$ of $\phi$ with
$\Psi(\phi)\in\mathcal A_\lambda(K[J_n],P)$ has $\dim P_H\geq|H|-\dim\mathcal A_\lambda(K[J_n],\bar P)\geq|H|/2$.
It is stable under right translations: translating $\phi$ by $\psi(c')$ translates $\Psi(\phi)$ on the right
by $c'\in K[J_0]$, which normalizes $K[J_n]$ and preserves $\mathcal A_\lambda(K[J_n],P)$.
Lemma~\ref{gal-congruences-l-group} gives $c\in P_H$ with $|1-c(h)|_\iota^2\leq\ell^{-1}$. Take
$G=\Psi(c)$, which lies in $\mathcal A_\lambda(K[J_n],P)\subset\mathcal A_\lambda(K,P)$ as $K\subset K[J_n]$,
and $e(g)=1-c(\beta(g))$.

\emph{Iteration.} Applying the step to $f$, then to $f-G_1$, which is again fixed by $K$, and so on, gives
after $m$ steps $f-(G_1+\cdots+G_m)=e_m\cdots e_1f$ with $|e_m\cdots e_1|_\iota\leq\ell^{-m/2}$.
\end{proof}
\leanref{Development/AutomorphicArithmetic/Foundations/GaloisRepresentations/SeparatingAlgebras/StGlob.lean}{UpperTheory.Gal.Gen.StepCore}
\leanref{Development/AutomorphicArithmetic/Foundations/GaloisRepresentations/SeparatingAlgebras/StepCore.lean}{UpperTheory.Gal.Gen.stepCore_pf}
\leanref{Development/AutomorphicArithmetic/Foundations/GaloisRepresentations/SeparatingAlgebras/GroupDetAlt.lean}{UpperTheory.Gal.Gen.Codim}
\leanref{Development/AutomorphicArithmetic/Foundations/GaloisRepresentations/SeparatingAlgebras/StGlob.lean}{UpperTheory.Gal.Gen.Tower}

\begin{lemma}\label{gal-congruences-operators}
Let $E$ be a CM field with maximal totally real subfield $F$, $a\geq1$, and $w$ a place of $E$ split over
$F$, with $q_w$ the order of its residue field.
\begin{enumerate}
\item For every $n\in\mathbf Z$ there is a Hecke operator supported on $U_{a,w}$ that acts on every
  $\varphi$ in an automorphic representation $U$ fixed by the $x\in U_{a,w}$ with
  $x_w\in\operatorname{GL}_a(\mathcal O_w)$ by $\sum_{b\in B}(q_w^{(a-1)/2}b)^n$, where $U$ has Satake
  parameters $B$ at $w$ (Definition~\ref{gal-definite-unitary}).
\item Let $K\subset U_a$ be a subgroup such that an $x\in U_{a,w}$ lies in $K$ if and only if
  $x_w\in\operatorname{GL}_a(\mathcal O_w)$, and $\Phi$ a polynomial in $a$ variables. There is a Hecke
  operator supported on $U_{a,w}$ that acts on level $K$ by some $\eta$ (Definition~\ref{gal-congruences-forms}),
  where for every $U$ with a nonzero vector fixed by $K$, $\eta(U)\neq0$ if and only if $U$ is
  $\Phi$-generic at $w$ (Definition~\ref{gal-congruences-dense}).
\item Let $K_w\subset\operatorname{GL}_a(E_w)$ be a compact open subgroup and $f$ a finitely supported
  function on $\operatorname{GL}_a(E_w)$ that acts by a scalar $z(\sigma)$ on the $K_w$-fixed vectors of
  every irreducible smooth representation $\sigma$. There is a Hecke operator $t$ supported on $U_{a,w}$,
  with $t(g)=f(g_w)$, that acts on the vectors fixed by the $x\in U_{a,w}$ with $x_w\in K_w$ of every
  automorphic representation $U$ by a scalar, equal to $z(\sigma)$ for the local component $\sigma$ of $U$ at
  $w$ when $U$ has such a nonzero vector.
\item Let $\ell$ be a prime, $\iota:\bar{\mathbf Q}_\ell\to\mathbf C$ a field isomorphism and $\lambda$ a
  family of dominant weights. For every Hecke operator $t$ supported on $U_a(\mathbf A_{F,f})$ there is
  $k\geq0$ such that $\ell^kt$ is $\ell$-integral for $\lambda$.
\end{enumerate}
\end{lemma}
\leanref{Development/AutomorphicArithmetic/Foundations/GaloisRepresentations/IntegralStructures/CongruenceOddLevel/ExistsTestOp.lean}{UpperTheory.Gal.Int.Cong.exists_powOp}
\leanref{Development/AutomorphicArithmetic/Foundations/GaloisRepresentations/IntegralStructures/CongruenceOddLevel/ExistsTestOp.lean}{UpperTheory.Gal.Int.Cong.exists_testOp}
\leanref{Development/AutomorphicArithmetic/Foundations/GaloisRepresentations/IntegralStructures/CongruenceOddLevel/ExistsTestOp.lean}{UpperTheory.Gal.Int.Cong.exists_op}
\leanref{Development/AutomorphicArithmetic/Foundations/GaloisRepresentations/IntegralStructures/CongruenceOddLevel/ExistsTestOp.lean}{UpperTheory.Gal.Int.Cong.exists_scale}

\begin{proof}
Since $w$ is split, $x\mapsto x_w$ identifies $U_{a,w}$ with $\operatorname{GL}_a(E_w)$, and $U$ has a unique
local component at $w$; a function $f$ on $\operatorname{GL}_a(E_w)$ transports to the Hecke operator
$g\mapsto f(g_w)$ on $U_{a,w}$, which acts on the relevant fixed vectors of $U$ as $f$ acts on its local
component. This gives (3). For (1), the power sum $\sum_{b\in B}(q_w^{(a-1)/2}b)^n$ is a symmetric Laurent
polynomial in the Satake parameters, hence by the Satake isomorphism \cite{Car79} the eigenvalue of a
spherical Hecke operator on the spherical vectors. For (2), $p(z)=\prod_\sigma\Phi(z_{\sigma(1)},\ldots,z_{\sigma(a)})$
is a symmetric polynomial; the spherical operator with Satake transform $p$, transported to $U_{a,w}$, acts
on the $K$-fixed vectors of $U$ by $p(z)$, where $U$ has Satake parameters $z$ at $w$ when it has a nonzero
$K$-fixed vector, and $p(z)\neq0$ if and only if $\Phi$ does not vanish at any ordering of $z$. For (4), the
support of $t$ is finite, so there is $d$ with $\ell^dh_u^{\pm1}$ integral at the places $u$ over $\ell$ for
all $h$ in the support; take $k$ with $\ell^{-k}|t(h)|_\iota\leq\ell^{-ad\,s(h)}$ for the finitely many $h$
in the support.
\end{proof}

\begin{proposition}\label{gal-congruences}
In Situation~\ref{gal-definite-situation}, there are open compact subgroups $K_w\subset\operatorname{GL}_a(E_w)$ at the split places $w$ away from
$w_0$, equal to $\operatorname{GL}_a(\mathcal O_w)$ where $U$ has Satake parameters away from $v'$, and
fixing vectors of $U$, with the following property. Let $Z$ be a finite set of split places away from
$w_0$, $v'$ and $\ell$ at which $U$ has Satake parameters, and $W$ a finite set of places over $\ell$
away from $w_0$ and $v'$ at which $U$ has Satake parameters. Then there are nonzero constants
$c_{w,n}$ such that for every family of Hecke operators of level $(K_w)$ at split places away from $w_0$
and $\ell$, with values $\theta$ on $U$, the system formed by $\theta$ and the numbers
$\iota^{-1}(c_{w,n}\sum_{b\in B_w}b^n)$, $w\in W$, $n\in\mathbf Z$, is approximated, in the sense
of Lemma~\ref{gal-trace-limit} below, by the same systems for automorphic representations $U'$ of $U_a$
that are unramified at the non-split places outside $T$, have level $(K_w)$, have Satake parameters at
$v'$, have the weights of $U$ at the embeddings that lie over neither $w_0$ nor $v'$ and a lacunary
weight at an embedding over $v'$, have the supercuspidal type of
Theorem~\ref{gal-seed-representations}(2) at $w_0$, and have generic Satake parameters (Definition~\ref{gal-generic-parameters}) at the places of $Z\cup W$. If the
operators have rational values, all such $\theta$ lie in one finite extension of $\mathbf Q_\ell$.
\end{proposition}
\leanref{Development/AbelianVarietyTheory/Definitions/GaloisGeometry/LimWU4OfLimWU3.lean}{GlobalTheory.CongL.CongFinOdd}
\leanref{Development/AbelianVarietyTheory/Definitions/GaloisGeometry/LimWU4OfLimWU3.lean}{GlobalTheory.CongL.CongUW4}
\leanref{Development/AutomorphicArithmetic/Foundations/GaloisRepresentations/IntegralStructures/CongruenceOddLevel/CongFinOddOfLvLac.lean}{UpperTheory.Gal.Int.Cong.congFinOdd_of_ex}
\leanref{Development/AutomorphicArithmetic/Foundations/GaloisRepresentations/IntegralStructures/CongruenceOddLevel/NotApartOfOverPair.lean}{UpperTheory.Gal.Int.Cong.StepSt}
\leanref{Development/AutomorphicArithmetic/Foundations/GaloisRepresentations/IntegralStructures/CongruenceOddLevel/NotApartOfOverPair.lean}{UpperTheory.Gal.Int.Cong.EllSt}
\leanref{Development/AutomorphicArithmetic/Foundations/GaloisRepresentations/IntegralStructures/CongruenceOddLevel/WtDefs.lean}{UpperTheory.Gal.Int.Cong.WtTo}
\leanref{Development/AutomorphicArithmetic/Foundations/GaloisRepresentations/WeightCharacters/LadicWeightEmbeddingOverPlace.lean}{UpperTheory.Gal.Wt.wtLacEx_holds}

\begin{proof}
The proof makes the Satake parameters generic at the places of $Z\cup W$ first, at the weight and level
of $U$; it then changes the weight at $v'$ and the type at $w_0$. Each change approximates systems of
Hecke eigenvalues $\ell$-adically, and Lemma~\ref{gal-congruences-calculus} chains the approximations.
It suffices to treat a finite set $Z$ of split places away from $w_0$ and $v'$ at which $U$ has Satake
parameters, together with $W$, and to obtain genericity at the places of $Z$: the places of $W$ are split
over $F$, as every place over $\ell$ is, so the argument applied to $Z\cup W$ in place of $Z$ gives the
statement. Let $\lambda$ be the family of weights of $U$, and use the notation of
Definitions~\ref{gal-congruences-forms} and~\ref{gal-congruences-dense}.

\emph{Levels.} There are a level $K$ away from $w_0$, an open compact subgroup
$K_v\subset\operatorname{GL}_a(E_{v'})$, and open compact subgroups $K_w\subset\operatorname{GL}_a(E_w)$ at
the split places $w$ away from $w_0$, such that: $U$ has a nonzero vector fixed by the group $K'$ of the
$k\in K$ with $k_{v'}\in K_v$; $K_{v'}=\operatorname{GL}_a(\mathcal O_{v'})$; $K_w=\operatorname{GL}_a(\mathcal
O_w)$ at the split $w$ away from $v'$ where $U$ has Satake parameters; $K$ contains the $x\in U_{a,w}$ with
$x_w\in\operatorname{GL}_a(\mathcal O_w)$ at the places $w$ not split over $F$ and not in $T$; and for split
$w$ away from $w_0$, an $x\in U_{a,w}$ lies in $K$ if and only if $x_w\in K_w$. To build them, fix a nonzero
$\varphi_0\in U$ and take at each place of $F$ other than the one under $w_0$ the integral points if they fix
a nonzero vector of $U$, which is the case at almost all places, and otherwise their intersection with the
stabilizer of $\varphi_0$; a nonzero vector of $U$ is fixed by all of these at once. Then replace the group
at the place under $v'$ by the integral points to obtain $K$, and let $K_v$ be the image of the group chosen
there before. These $K_w$ are the levels of the statement: a representation with a nonzero vector fixed by $K'$ has, at the split $w$
away from $w_0$ and $v'$, a local component with nonzero $K_w$-fixed vectors.

\emph{Operators.} For $w\in W$ and $n\in\mathbf Z$, Lemma~\ref{gal-congruences-operators}(1) gives a Hecke
operator $t_{w,n}$ supported on $U_{a,w}$ that acts on level $K'$ by the power sum
$\sum_{b\in B_w}(q_w^{(a-1)/2}b)^n$, as $K'$ contains the integral points at $w$, and
Lemma~\ref{gal-congruences-operators}(4) gives $k_{w,n}$ with $\ell^{k_{w,n}}t_{w,n}$ $\ell$-integral for
$\lambda$; let $c_{w,n}=\ell^{k_{w,n}}q_w^{n(a-1)/2}$. For $z\in Z$,
Lemma~\ref{gal-congruences-operators}(2) with $\Phi=D_a$ (Definition~\ref{gal-generic-parameters}), scaled by
a power of $\ell$, gives an $\ell$-integral operator $t_z$ acting on level $K'$ by $\eta_z$, with
$\eta_z(U')\neq0$ for $U'$ with a nonzero $K'$-fixed vector if and only if $U'$ has generic Satake parameters
at $z$. Let $(O_j)_{j\in J}$ be Hecke operators of level $(K_w)$ at split places $w_j$ away from $w_0$ and
$\ell$, with values $\theta$ on $U$. Lemma~\ref{gal-congruences-operators}(3) gives operators $t_j$ supported
on $U_{a,w_j}$, acting on level $K'$ by the values $\eta_j(U')$ of $O_j$ on the local components; they are
$\ell$-integral, with $d=0$ and $s=0$, since the coefficients of $O_j$ are integral and $w_j\nmid\ell$. All
these operators are supported on elements equal to $1$ at the places over the places of $F$ under $w_0$ and
$v'$. Let $\mathbf t$ be the family of the $t_j$, the $\ell^{k_{w,n}}t_{w,n}$ and the $t_z$, indexed by
$J\sqcup(W\times\mathbf Z)\sqcup Z$, and $\boldsymbol\eta(U')$ its family of eigenvalues on level $K'$; its
coordinates are the values $\eta_j(U')$, the numbers $c_{w,n}\sum_{b\in B'_w}b^n$ where $U'$ has Satake
parameters $B'_w$ at $w$, and the $\eta_z(U')$.

\emph{Genericity.} By Lemma~\ref{gal-congruences-generic}, with $S=Z$ and $\Phi_z=D_a$, generic Satake
parameters at every $z\in Z$ are $\ell$-adically dense in weight $\lambda$ and level $K'$; the hypothesis at
$z$ holds since $z$ is away from $v'$ and $U$ has Satake parameters there. Lemma~\ref{gal-congruences-gauge}
gives a gauge for $\mathbf t$, and Lemma~\ref{gal-congruences-density} shows that
$\iota^{-1}\boldsymbol\eta(U)$ is approximated by the $\iota^{-1}\boldsymbol\eta(U_1)$ for the $U_1$ of type
$(\lambda,K')$ with generic Satake parameters at every $z\in Z$.

\emph{Weight and type.} Let $U_1$ be such, and let $\tau_0$ be an embedding with $\iota^{-1}\circ\tau_0$
inducing $v'$. By Lemma~\ref{gal-congruences-weights}, $\iota^{-1}\boldsymbol\eta(U_1)$ is approximated by
the $\iota^{-1}\boldsymbol\eta(U_2)$ for the $U_2$ of type $(\mu,K)$, with $\mu$ equal to $\lambda$ at the
embeddings that do not lie over $v'$ and $\{k-\mu_{\tau_0,k}\}$ lacunary. For such $\mu$, the supercuspidal
type of Theorem~\ref{gal-seed-representations}(2) at $w_0$ is $\ell$-adically dense in weight $\mu$ and level
$K$. Indeed, let $f$ be an algebraic form of weight $\mu$ fixed by $K$ and $\varepsilon>0$; $f$ is fixed by
the $x\in U_{a,w_0}$ with $x_{w_0}$ in a neighbourhood $N$ of $1$; choose $m_0$ with $\ell^{-m_0}<\varepsilon$.
Lemma~\ref{gal-congruences-type-character} gives $m\geq m_0$, $J\subset N$ and $\vartheta$; $f$ is fixed by
$K[J]$; and Lemma~\ref{gal-congruences-type} gives $G$ and $e$ with $|e|_\iota\leq\ell^{-m}<\varepsilon$, where
the representations occurring in $G$ have a nonzero $K$-fixed vector on which $J$ acts by $\chi\circ\vartheta$
with $\chi\neq1$, hence have the supercuspidal type at $w_0$ by Lemma~\ref{gal-congruences-type-character}.
The operators $\mathbf t$ are $\ell$-integral for $\mu$, since $\mu=\lambda$ at the embeddings over the places
where elements of their supports are nontrivial, and they act on level $K$ since $K\supset K'$. So Lemmas~\ref{gal-congruences-gauge}
and~\ref{gal-congruences-density} show that $\iota^{-1}\boldsymbol\eta(U_2)$ is approximated by the
$\iota^{-1}\boldsymbol\eta(U')$ for the $U'$ of type $(\mu,K)$ with the supercuspidal type at $w_0$, and
Lemma~\ref{gal-congruences-calculus}(1), with $Q=1$, combines the two approximations for $U_1$.

\emph{Combining.} Let $Q=\prod_{z\in Z}X_z$, the product of the coordinates of the operators $t_z$. Since
$U_1$ has generic Satake parameters at every $z\in Z$, $Q(\iota^{-1}\boldsymbol\eta(U_1))\neq0$.
Lemma~\ref{gal-congruences-calculus}(1) shows that $\iota^{-1}\boldsymbol\eta(U)$ is approximated by the
$\iota^{-1}\boldsymbol\eta(U')$ of the previous step with $Q(\iota^{-1}\boldsymbol\eta(U'))\neq0$, that is,
with generic Satake parameters at every $z\in Z$. Lemma~\ref{gal-congruences-calculus}(2), applied to the
inclusion of $J\sqcup(W\times\mathbf Z)$, removes the coordinates of the $t_z$. What remains is the
statement: $\iota^{-1}\eta_j(U)=\theta_j$, and each $U'$ has level $K$, so it is unramified at the non-split
places outside $T$, has level $(K_w)$, and has Satake parameters at $v'$ since $K$ contains the integral
points at $v'$; it has the weights of $U$ at the embeddings not over $v'$, in particular at those over
neither $w_0$ nor $v'$, and the lacunary weight $\mu_{\tau_0}$ at $\tau_0$, the supercuspidal type at
$w_0$, and generic Satake parameters at the places of $Z$.

\emph{Rationality.} If the $O_j$ have rational values, the $t_j$ have rational coefficients. Let $\varphi\in
U$ be nonzero and fixed by $K'$. For every automorphism $\sigma$ of $\mathbf C$ the conjugate of $\varphi$ is
an automorphic form with $t_j$-eigenvalues $\sigma(\eta_j(U))$; these conjugates span a finite-dimensional
space, and eigenvectors with different systems of eigenvalues are linearly independent, so the family
$(\sigma(\eta_j(U)))_j$ takes finitely many values. Hence the $\eta_j(U)$ lie in one number field, and
their images under $\iota^{-1}$ in one finite extension of $\mathbf Q_\ell$.
\end{proof}
\leanref{Development/AutomorphicArithmetic/Foundations/GaloisRepresentations/IntegralStructures/CongruenceOddLevel/NotApartOfOverPair.lean}{UpperTheory.Gal.Int.Cong.congUWG_of}
\leanref{Development/AutomorphicArithmetic/Foundations/GaloisRepresentations/IntegralStructures/CongruenceOddLevel/NotApartOfOverPair.lean}{UpperTheory.Gal.Int.Cong.stepSt_of_ex}
\leanref{Development/AutomorphicArithmetic/Foundations/GaloisRepresentations/IntegralStructures/CongruenceOddLevel/NotApartOfOverPair.lean}{UpperTheory.Gal.Int.Cong.CongUWG.congUW4}
\leanref{Development/AutomorphicArithmetic/Foundations/GaloisRepresentations/IntegralStructures/CongruenceOddLevel/OpenCompactMap.lean}{UpperTheory.Gal.Int.Cong.exists_level}
\leanref{Development/AbelianVarietyTheory/Proofs/SatakeParameters/LinearIndependentSimulEigen.lean}{GlobalTheory.Cong.finU_holds}

\begin{lemma}\label{gal-trace-limit}
Let $\ell$ be a prime and $\Gamma$ a compact topological group. For a set $J$, a function
$\theta:J\to\bar{\mathbf Q}_\ell$ is \emph{approximated} by a set $S$ of such functions if for every
$\varepsilon>0$ there are finitely many $\eta_1,\ldots,\eta_k\in S$ and $c_1,\ldots,c_k\in\bar{\mathbf
Q}_\ell$ with $|P(\theta)-\sum_ic_iP(\eta_i)|\leq\varepsilon$ for every polynomial $P$ with integer
coefficients in the variables indexed by $J$.
\begin{enumerate}
\item Let $x:J\to\Gamma$, let $\theta:J\to L$ with $L$ a finite extension of $\mathbf Q_\ell$, and let
  $(\rho_i)_{i\in I}$ be continuous representations $\Gamma\to\operatorname{GL}_a(\bar{\mathbf
  Q}_\ell)$ with functions $\eta_i:J\to\bar{\mathbf Q}_\ell$ and sets $D_i\subset J$ such that
  $\eta_i(j)=\operatorname{tr}\rho_i(x_j)$ for $j\in D_i$ and $x(\bigcap_{i\in F}D_i)$ is dense for every
  finite $F\subset I$. If $\theta$ is approximated by $\{\eta_i\}$, then there is a continuous
  $t:\Gamma\to L$ such that every estimate $|P(\theta)-\sum_{i\in F}c_iP(\eta_i)|\leq\varepsilon$ for all
  $P$ implies $|t(g)-\sum_{i\in F}c_i\operatorname{tr}\rho_i(g)|\leq\varepsilon$ for all $g$, and such
  that for every $A\subset I$ by whose $\eta_i$ the function $\theta$ is approximated, $t$ is approximated
  by the $\operatorname{tr}\rho_i$, $i\in A$, and $t(x_j)=\theta(j)$ whenever $j\in D_i$ for all $i\in A$.
\item Let $x:J\to\Gamma$ have dense image and $C$ a set of continuous representations
  $\Gamma\to\operatorname{GL}_a(\bar{\mathbf Q}_\ell)$ such that $\theta$, with values in a finite
  extension of $\mathbf Q_\ell$, is approximated by the functions $j\mapsto\operatorname{tr}\rho(x_j)$,
  $\rho\in C$. Then there is a continuous semisimple $r:\Gamma\to\operatorname{GL}_a(\bar{\mathbf
  Q}_\ell)$ with $\operatorname{tr}r(x_j)=\theta(j)$ for all $j$, whose trace is approximated by the
  traces of members of $C$.
\end{enumerate}
\end{lemma}
\leanref{Development/AbelianVarietyTheory/Definitions/GaloisGeometry/PolyApproxZeroPadic.lean}{GlobalTheory.Cong.PolyApprox}
\leanref{Development/AutomorphicArithmetic/Foundations/GaloisRepresentations/TraceLimits/ExistsTraceLimit.lean}{UpperTheory.Gal.Lim.exists_trace_limit}
\leanref{Development/AbelianVarietyTheory/Proofs/InertiaEigenvalues/ExistsLimitRep.lean}{GlobalTheory.Cong.exists_limit_rep}

\begin{proof}
(1) For each $\varepsilon$ choose $F_\varepsilon$ and $c$ as in the definition. For $g=x_j$ with
$j\in\bigcap_{i\in F_\varepsilon}D_i$, the function $\sum_ic_i\operatorname{tr}\rho_i$ is within
$\varepsilon$ of $\theta(j)$, by the estimate for the polynomial $P=X_j$. For two values of $\varepsilon$
the two combinations differ by at most the larger value on a dense set, hence everywhere by continuity.
So the combinations form a Cauchy family of continuous functions with values in the complete field $L$,
and their uniform limit $t$ has the stated properties.

(2) The approximating traces form a locally Cauchy system on $x(J)$, which extends to a continuous
function $T$ on $\Gamma$ that is a uniform limit of linear combinations of traces of members of $C$. The
identities that define an $a$-dimensional pseudo-character are polynomial in the traces and hold for
each $\operatorname{tr}\rho$; by the uniform estimate for all polynomials they pass to the limit, so
$T$ is a pseudo-character of dimension $a$. By Taylor's theorem \cite{Tay91} it is the trace of a
semisimple representation $r$ over $\bar{\mathbf Q}_\ell$, and $r$ is continuous because its trace is.
\end{proof}

\subsection{Covers and Breuil--Kisin structures}\label{subsection-gal-breuil-kisin}

\begin{definition}\label{gal-cover-definition}
Let $\ell$ be a prime, normalize the absolute value of $\bar{\mathbf Q}_\ell$ by $|\ell|=\ell^{-1}$, let
$\Gamma$ be a topological group and $a\geq0$ an integer, and let
$r:\Gamma\to\operatorname{GL}_a(\bar{\mathbf Q}_\ell)$ be a homomorphism.
\begin{enumerate}
\item A homomorphism $\rho:\Gamma\to\operatorname{GL}_d(\mathbf Z_\ell)$ \emph{reads} $r$ through matrices
  $b_1,\ldots,b_d\in M_a(\bar{\mathbf Q}_\ell)$ and a vector $e\in\mathbf Z_\ell^d$ if
  $r(g)b_j=\sum_i\rho(g)_{ij}b_i$ for all $g\in\Gamma$ and all $j$, and $\sum_ie_ib_i=1$.
\item A \emph{lattice tower} of a homomorphism $\rho:\Gamma\to\operatorname{GL}_d(\mathbf Z_\ell)$
  consists, for every $n\geq0$, of a homomorphism $\rho'_n:\Gamma\to\operatorname{GL}_{d'_n}(\mathbf
  Z_\ell)$ and matrices $\Pi_n\in M_{d\times d'_n}(\mathbf Z_\ell)$ and $f_n\in M_{d'_n\times
  d'_{n+1}}(\mathbf Z_\ell)$, such that $\Pi_n\rho'_n(g)\equiv\rho(g)\Pi_n$ modulo $\ell^n$ for all $g$,
  $\Pi_nX\equiv1$ modulo $\ell^n$ for some $X\in M_{d'_n\times d}(\mathbf Z_\ell)$,
  $f_n\rho'_{n+1}(g)=\rho'_n(g)f_n$ for all $g$, and $\Pi_nf_n\equiv\Pi_{n+1}$ modulo $\ell^n$.
\item Let $C$ be a set of homomorphisms $\Gamma\to\operatorname{GL}_a(\bar{\mathbf Q}_\ell)$ and
  $L_0\subset\bar{\mathbf Q}_\ell$ a finite extension of $\mathbf Q_\ell$. A \emph{layer} for a
  homomorphism $\rho':\Gamma\to\operatorname{GL}_{d'}(\mathbf Z_\ell)$ consists of members
  $\rho_1,\ldots,\rho_m$ of $C$, finite extensions $L_1,\ldots,L_m\subset\bar{\mathbf Q}_\ell$ of $L_0$
  with $\rho_k(\Gamma)\subset\operatorname{GL}_a(L_k)$, elements $b'_j=(b'_{j,k})_k$ of
  $\prod_kM_a(L_k)$ for $1\leq j\leq d'$ that form a $\mathbf Q_\ell$-basis of $\prod_kM_a(L_k)$, a
  vector $e'\in\mathbf Z_\ell^{d'}$ such that $\rho'$ reads every $\rho_k$ through
  $b'_{1,k},\ldots,b'_{d',k}$ and $e'$, and coefficients $c_1,\ldots,c_m\in\bar{\mathbf Q}_\ell$.
\item A \emph{cover} of $r$ relative to $C$ and $L_0$ consists of a continuous homomorphism
  $\rho:\Gamma\to\operatorname{GL}_d(\mathbf Z_\ell)$ that reads $r$ through some $b$ and $e$, a lattice
  tower $(\rho'_n,\Pi_n,f_n)_n$ of $\rho$ with every $\rho'_n$ continuous, and for every $n$ a layer
  $(\rho_{n,k},L_{n,k},b'_{n,j},e'_n,c_{n,k})$ for $\rho'_n$, such that for every $n$ we have
  $\Pi_ne'_n\equiv e$ modulo $\ell^n$ and
  \[
  \Bigl|\operatorname{tr}\Bigl(\sum_i(\Pi_n)_{ij}b_i\Bigr)-\sum_kc_{n,k}\operatorname{tr}b'_{n,j,k}\Bigr|
  \leq\ell^{-n}\qquad(1\leq j\leq d'_n).
  \]
\end{enumerate}
\end{definition}
\leanref{Development/AbelianVarietyTheory/Definitions/NonGaloisDescent/TrivIsCr.lean}{GlobalTheory.BKLim.Reads}
\leanref{Development/AbelianVarietyTheory/Definitions/NonGaloisDescent/TrivIsCr.lean}{GlobalTheory.BKLim.Onto}
\leanref{Development/AbelianVarietyTheory/Definitions/NonGaloisDescent/TrivIsCr.lean}{GlobalTheory.BKLim.Tower}
\leanref{Development/AbelianVarietyTheory/Definitions/NonGaloisDescent/TrivIsCr.lean}{GlobalTheory.BKLim.Layer}
\leanref{Development/AbelianVarietyTheory/Definitions/NonGaloisDescent/TrivIsCr.lean}{GlobalTheory.BKLim.Cover}

\begin{lemma}\label{gal-cover-joint}
Let $\ell$ be a prime, $\Gamma$ a compact topological group, $a\geq0$ an integer,
$L_0\subset\bar{\mathbf Q}_\ell$ a finite extension of $\mathbf Q_\ell$,
$r:\Gamma\to\operatorname{GL}_a(\bar{\mathbf Q}_\ell)$ a continuous semisimple homomorphism, $C$ a set of
continuous homomorphisms $\Gamma\to\operatorname{GL}_a(\bar{\mathbf Q}_\ell)$, $p:C\to\bar{\mathbf Q}_\ell$
a function and $\theta\in\bar{\mathbf Q}_\ell$. Suppose that for every $\varepsilon>0$ there are
$\rho_1,\ldots,\rho_k\in C$ and $c_1,\ldots,c_k\in\bar{\mathbf Q}_\ell$ with
$|\operatorname{tr}r(g)-\sum_ic_i\operatorname{tr}\rho_i(g)|\leq\varepsilon$ for all $g\in\Gamma$ and
$|\theta-\sum_ic_ip(\rho_i)|\leq\varepsilon$. Then $r$ has a cover relative to $C$ and $L_0$
(Definition~\ref{gal-cover-definition}) whose layers satisfy
$|\theta-\sum_kc_{n,k}p(\rho_{n,k})|\leq\ell^{-n}$ for every $n$.
\end{lemma}
\leanref{Development/AutomorphicArithmetic/Foundations/GaloisRepresentations/ClosedPhiFunctoriality/IntBasisAlt.lean}{UpperTheory.Gal.ClosedPhiFromClosedTrace.CoverJ}
\leanref{Development/AutomorphicArithmetic/Foundations/GaloisRepresentations/ClosedPhiFunctoriality/CoverJ.lean}{UpperTheory.Gal.ClosedPhiFromClosedTrace.coverJ_holds}

\begin{proof}
Let $\Lambda\subset M_a(\bar{\mathbf Q}_\ell)$ be the $\mathbf Z_\ell$-span of $r(\Gamma)$. Since $\Gamma$
is compact, the entries of $r(\Gamma)$ are bounded and, by the Baire category theorem, lie in a finite
extension of $\mathbf Q_\ell$, so $\Lambda$ is free of finite rank; choose a basis $b_1,\ldots,b_d$ of
$\mathbf Z_\ell$-combinations of elements of $r(\Gamma)$. Left multiplication by $r(g)$ preserves
$\Lambda$, and its matrix in this basis defines a continuous $\rho:\Gamma\to\operatorname{GL}_d(\mathbf
Z_\ell)$ that reads $r$ through $b$ and the coordinates $e$ of $1\in\Lambda$. Since $r$ is semisimple,
the $\mathbf Q_\ell$-span of $r(\Gamma)$ is a semisimple algebra on which
$(x,y)\mapsto\operatorname{tr}(xy)$ is nondegenerate. Hence there is $c\geq0$ such that every
$x\in\Lambda$ with $|\operatorname{tr}(xr(g))|\leq\ell^{-n-c}$ for all $g$ lies in $\ell^n\Lambda$.
Choose $K\geq0$ with $|\operatorname{tr}x|\leq\ell^K$ for all $x\in\Lambda$. For each $\rho'\in C$
choose a finite extension $L_{\rho'}\supset L_0$ containing the entries of $\rho'(\Gamma)$ and a finite
subset of $M_a(L_{\rho'})$ spanning it over $\mathbf Q_\ell$, whose $\mathbf Z_\ell$-span
$\Lambda_{\rho'}$ contains $1$ and is stable under left multiplication by $\rho'(\Gamma)$; this exists by
compactness.

For every $n$ the hypothesis with $\varepsilon=\ell^{-n-K-c}$ gives a finite set $s_n\subset C$ and
coefficients $c_n(\rho')$, $\rho'\in s_n$. Let $U_n=s_0\cup\cdots\cup s_n$, enumerate it as
$\rho_{n,1},\ldots,\rho_{n,m_n}$, put $L_{n,k}=L_{\rho_{n,k}}$, and let $c_{n,k}$ be
$c_n(\rho_{n,k})$ if $\rho_{n,k}\in s_n$ and $0$ otherwise. Then $\sum_kc_{n,k}p(\rho_{n,k})$ is
$\sum_{\rho'\in s_n}c_n(\rho')p(\rho')$, which is within $\ell^{-n-K-c}\leq\ell^{-n}$ of $\theta$.

For $M\geq0$ let $T_{n,M}\subset\prod_kM_a(L_{n,k})$ be the $\mathbf Z_\ell$-span of the tuples
$(\rho_{n,k}(g))_k$, $g\in\Gamma$, and of $\ell^M\prod_k\Lambda_{\rho_{n,k}}$. It lies between
$\ell^M\prod_k\Lambda_{\rho_{n,k}}$ and $\prod_k\Lambda_{\rho_{n,k}}$, so it is a lattice in
$\prod_kM_a(L_{n,k})$, and it is stable under left multiplication by the tuples $(\rho_{n,k}(g))_k$. In
a $\mathbf Z_\ell$-basis $b'_{n,1},\ldots,b'_{n,d'_n}$ of $T_{n,M}$ this multiplication defines a
continuous $\rho'_n$, which reads every $\rho_{n,k}$ through the $k$-th components of the $b'_{n,j}$ and
the coordinates $e'_n$ of the tuple $(1)_k$.

Send $\sum_gF(g)(\rho_{n,k}(g))_k+y$, with $F:\Gamma\to\mathbf Z_\ell$ finitely supported and
$y\in\ell^M\prod_k\Lambda_{\rho_{n,k}}$, to the class of $\sum_gF(g)r(g)$ in $\Lambda/\ell^{n+K}\Lambda$.
For $M$ large this is well defined. Indeed, let $x=\sum_gF(g)(\rho_{n,k}(g))_k$ lie in
$\ell^M\prod_k\Lambda_{\rho_{n,k}}$, and let $h\in\Gamma$. By the choice of $s_n$ the trace of
$\sum_gF(g)r(gh)$ is within $\ell^{-n-K-c}$ of $\sum_kc_{n,k}\operatorname{tr}(\rho_{n,k}(h)x_k)$, and
this sum has absolute value at most $\ell^{-n-K-c}$ once $M$ is large, since
$\rho_{n,k}(h)x_k\in\ell^M\Lambda_{\rho_{n,k}}$. So $\sum_gF(g)r(g)\in\ell^{n+K}\Lambda$ by the choice of
$c$. Let the $j$-th column of $\Pi_n\in M_{d\times d'_n}(\mathbf Z_\ell)$ be the coordinates of a lift of
the image of $b'_{n,j}$. The map commutes with the actions of $\Gamma$ and is onto, since $\Lambda$ is
spanned by $r(\Gamma)$; so $\Pi_n\rho'_n(g)\equiv\rho(g)\Pi_n$ and $\Pi_nX\equiv1$ modulo $\ell^n$ for
some $X$, and $\Pi_ne'_n\equiv e$ since the tuple $(1)_k$ is sent to $1$. Write
$b'_{n,j}=\sum_gF(g)(\rho_{n,k}(g))_k+y$. Then $\sum_i(\Pi_n)_{ij}b_i$ differs from $\sum_gF(g)r(g)$ by
an element of $\ell^{n+K}\Lambda$, whose trace has absolute value at most $\ell^{-n}$; the trace of
$\sum_gF(g)r(g)$ is within $\ell^{-n-K-c}$ of $\sum_kc_{n,k}\operatorname{tr}(b'_{n,j,k}-y_k)$; and
$|\sum_kc_{n,k}\operatorname{tr}y_k|\leq\ell^{-n}$ once $M$ is large. This is the trace estimate of a
cover.

Choose such an $M=M_n$ for every $n$, increasing in $n$. Since $U_n\subset U_{n+1}$, forgetting the
components of the members of $U_{n+1}$ that are not in $U_n$ maps $T_{n+1,M_{n+1}}$ into $T_{n,M_n}$
and commutes with the actions of $\Gamma$. Its matrix $f_n$ satisfies $f_n\rho'_{n+1}(g)=\rho'_n(g)f_n$,
and $\Pi_nf_n\equiv\Pi_{n+1}$ modulo $\ell^n$, since both send $\sum_gF(g)(\rho_{n+1,k}(g))_k+y$ to the
class of $\sum_gF(g)r(g)$.
\end{proof}

\begin{definition}\label{gal-cover-breuil-kisin}
Let $\ell$ be a prime and $K\subset\bar{\mathbf Q}_\ell$ a finite extension of $\mathbf Q_\ell$, with
$G_K=\operatorname{Gal}(\bar{\mathbf Q}_\ell/K)$. Let $\mathbf C_\ell$ be the completion of
$\bar{\mathbf Q}_\ell$, $\mathcal O^\flat$ the valuation ring of its tilt $\mathbf C_\ell^\flat$,
$A_{\inf}=W(\mathcal O^\flat)\subset W(\mathbf C_\ell^\flat)$ with Fontaine's map
$\theta:A_{\inf}\to\mathcal O_{\mathbf C_\ell}$, and $A_{\mathrm{cris}}\subset B_{\mathrm{cris}}$
Fontaine's rings \cite{Fon94a}. Write $\phi$ for the Frobenius of all of these and $[x]$ for
Teichm\"uller representatives. View the maximal unramified extension $\mathbf Q_\ell^{\mathrm{ur}}$ of
$\mathbf Q_\ell$ in $\bar{\mathbf Q}_\ell$ inside $B_{\mathrm{cris}}$, where $\phi$ acts on it by the
arithmetic Frobenius $\operatorname{Frob}$.
\begin{enumerate}
\item A \emph{Breuil--Kisin datum} for $K$ consists of an integer $h\geq1$, the ring
  $S=W(\mathbf F_{\ell^h})[[u]]$ with the Frobenius $\phi$ that is the Frobenius of $W(\mathbf
  F_{\ell^h})$ on coefficients and sends $u$ to $u^\ell$, a monic polynomial $E\in W(\mathbf
  F_{\ell^h})[u]$ of positive degree whose lower coefficients are divisible by $\ell$, whose constant
  coefficient is $\ell$ times a unit and which is not a polynomial in $u^\ell$, an isomorphism
  $\mathcal O_K\cong S/(E)$, an element $\pi^\flat\in\mathcal O^\flat$, and a homomorphism
  $\iota:S\to A_{\inf}$ with $\iota(u)=[\pi^\flat]$ and $\iota\phi=\phi\iota$, equal on
  $W(\mathbf F_{\ell^h})$ to the map induced by an embedding $\mathbf F_{\ell^h}\to\mathcal O^\flat$
  whose image $G_K$ fixes, such that $\iota(E)$ generates $\ker\theta$ and $\theta\circ\iota$ induces
  the inclusion $\mathcal O_K\subset\mathcal O_{\mathbf C_\ell}$ through $\mathcal O_K\cong S/(E)$.
\item A \emph{cyclotomic element} for such a datum is a $\tau\in A_{\inf}$ with
  $\phi(\tau)=\iota(E)\tau$ and $\phi(\tau)=\eta([\zeta]-1)$ for a unit $\eta$ of $A_{\inf}$ and an
  element $\zeta\in\mathcal O^\flat$ with $\zeta^\sharp=1$ and $(\zeta^{1/\ell})^\sharp\neq1$.
\item Let $\rho:G_K\to M_d(\mathbf Z_\ell)$ be a map and $N\geq0$. A \emph{Breuil--Kisin structure} of
  height $N$ on $\rho$ consists of $T\in\operatorname{GL}_d(W(\mathbf C_\ell^\flat))$, matrices
  $A,A'\in M_d(S)$ with $AA'=E^{2N}$, and a map $\varepsilon:G_K\to M_d(A_{\inf})$, such that
  $\iota(E)^N\phi(T)=T\iota(A)$, $\rho(g)g(T)=T\varepsilon(g)$ for every $g\in G_K$, and
  $\rho(g)g(T)=T$ for every $g\in G_K$ that fixes $\pi^\flat$. It is \emph{crystalline} if $[\pi^\flat]$
  divides every entry of $\varepsilon(g)-1$ in $A_{\inf}$, for every $g\in G_K$.
\item Breuil--Kisin structures with matrices $T$ of size $d$ and $T_1$ of size $d_1$ are \emph{linked}
  at $(n,r_0)$ along $\Pi\in M_{d\times d_1}(\mathbf Z_\ell)$ if there are $R\in M_{d\times d_1}(S)$ and
  $G\in M_{d_1\times d}(S)$ with $\Pi T_1\equiv T\iota(R)$ modulo $\ell^n$ in $W(\mathbf C_\ell^\flat)$
  and $RG\equiv\ell^{r_0}$ modulo $\ell^n$ in $S$.
\item A \emph{period matrix} of a homomorphism $\rho:G_K\to\operatorname{GL}_d(\mathbf Z_\ell)$ is a
  $U\in\operatorname{GL}_d(B_{\mathrm{cris}})$ with $\rho(g)g(U)=U$ for every $g\in G_K$.
\end{enumerate}
\end{definition}
\leanref{Development/AbelianVarietyTheory/Definitions/NonGaloisDescent/IsStratHom.lean}{GlobalTheory.BSP.BK}
\leanref{Development/AbelianVarietyTheory/Definitions/NonGaloisDescent/IsStratHom.lean}{GlobalTheory.BSP.BK.IsPer}
\leanref{Development/AbelianVarietyTheory/Definitions/NonGaloisDescent/TrivIsCr.lean}{GlobalTheory.BKLim.Pre}
\leanref{Development/AbelianVarietyTheory/Definitions/NonGaloisDescent/TrivIsCr.lean}{GlobalTheory.BKLim.Frame}
\leanref{Development/AbelianVarietyTheory/Definitions/NonGaloisDescent/TrivIsCr.lean}{GlobalTheory.BKLim.Frame.IsCr}
\leanref{Development/AbelianVarietyTheory/Definitions/NonGaloisDescent/TrivIsCr.lean}{GlobalTheory.BKLim.Link}
\leanref{Development/AutomorphicArithmetic/Foundations/GaloisRepresentations/ClosedPhiFunctoriality/JointApproxOfMem.lean}{UpperTheory.Gal.ClosedPhiFromClosedTrace.IsPerU}

\begin{lemma}\label{gal-structures-datum}
Let $\ell$ be a prime and $K\subset\bar{\mathbf Q}_\ell$ a finite extension of $\mathbf Q_\ell$. Then $K$ has
a Breuil--Kisin datum (Definition~\ref{gal-cover-breuil-kisin}).
\end{lemma}
\leanref{Development/AbelianVarietyTheory/Proofs/InertiaEigenvalues/NonemptyBreuilKisinDatum.lean}{GlobalTheory.BSP.nonempty_bk_of_finExt}
\leanref{Development/AbelianVarietyTheory/Proofs/InertiaEigenvalues/ExistsWittResidueFieldInRingOfIntegers.lean}{GlobalTheory.BSP.Res.exists_witt_residue}
\leanref{Development/AbelianVarietyTheory/Proofs/InertiaEigenvalues/ExistsUniformiserEisensteinNotInXPow.lean}{GlobalTheory.EisShift.exists_uniformiser}
\leanref{Development/AbelianVarietyTheory/Proofs/InertiaEigenvalues/BreuilKisinQuotientEisensteinEquiv.lean}{GlobalTheory.bkStr_quotient_equiv}
\leanref{Development/AbelianVarietyTheory/Proofs/InertiaEigenvalues/EisensteinDistinguishedPowerSeriesAdicComplete.lean}{GlobalTheory.BKReal.eisPoly_distinguished_complete}
\leanref{Development/AbelianVarietyTheory/Proofs/InertiaEigenvalues/ExistsBreuilKisinEvalTeichmullerTilt.lean}{GlobalTheory.BKReal.exists_ev_of_root}
\leanref{Development/AbelianVarietyTheory/Proofs/InertiaEigenvalues/SpanEqKerThetaOfDistinguishedAinf.lean}{GlobalTheory.BSP.span_eq_ker_theta_of_distinguished}

\begin{proof}
The ring $\mathcal O_K$ is a complete discrete valuation ring, finite over $\mathbf Z_\ell$; let $\ell^h$ be
the number of elements of its residue field $k$. Fix an embedding $\mathbf F_{\ell^h}\to\mathcal O^\flat$, let
$w:W(\mathbf F_{\ell^h})\to A_{\inf}$ be the induced map, and put $w_K=\theta\circ w$. An element of
$\mathcal O^\flat$ in the image of $\mathbf F_{\ell^h}^\times$ is a compatible system of $(\ell^h-1)$-th roots
of unity of $\mathcal O_{\mathbf C_\ell}$, which lie in $K$ by Hensel's lemma. So $G_K$ fixes the image of
$w$, and $w_K$ maps $W(\mathbf F_{\ell^h})$ into $\mathcal O_K$, inducing an isomorphism of residue fields.

\emph{The polynomial $E$.} Let $\pi$ be a uniformizer of $K$. Since $\mathcal O_K$ is totally ramified over
$W(\mathbf F_{\ell^h})$, the minimal polynomial $E$ of $\pi$ over $W(\mathbf F_{\ell^h})$ is an Eisenstein
polynomial, $\mathcal O_K=W(\mathbf F_{\ell^h})[\pi]$, and the polynomials vanishing at $\pi$ are the
multiples of $E$ \cite[Chapter~I, \S6]{Ser79}. Suppose that $E$ is a polynomial in $u^\ell$. Then $\ell$
divides $e=\deg E$, so $e\geq2$ and $\ell=\pi^2t$ with $t\in\mathcal O_K$. The element
$\pi'=\pi+\ell=\pi(1+\pi t)$ is a uniformizer, and its minimal polynomial divides $E(u-\ell)$ and has the
same degree, so it is $E(u-\ell)$. Write $e=n+1$. Then $\ell\nmid n$, so the coefficient of $u^n$ in $E$
vanishes, and the coefficient of $u^n$ in $E(u-\ell)$ is $-(n+1)\ell\neq0$. Replacing $\pi$ by $\pi'$, we may
assume that $E$ is not a polynomial in $u^\ell$. Division by $E$, which is monic with lower coefficients
divisible by $\ell$, gives $S/(E)\cong W(\mathbf F_{\ell^h})[u]/(E)\cong\mathcal O_K$, with $u\mapsto\pi$.

\emph{The map $\iota$.} The element $\pi$ is not a unit of $\mathcal O_{\mathbf C_\ell}$, and a compatible
system of $\ell$-power roots of $\pi$ in $\mathcal O_{\mathbf C_\ell}$ defines $\pi^\flat\in\mathcal O^\flat$
with $\theta([\pi^\flat])=\pi$. Since $A_{\inf}$ is complete and separated for the $[\pi^\flat]$-adic
topology, there is a homomorphism $\iota:S\to A_{\inf}$ equal to $w$ on $W(\mathbf F_{\ell^h})$ with
$\iota(u)=[\pi^\flat]$. The continuous homomorphisms $\iota\circ\phi$ and $\phi\circ\iota$ agree on
$W(\mathbf F_{\ell^h})$, since $w$ commutes with Frobenius, and on $u$, where both give $[\pi^\flat]^\ell$; so
they are equal. For a polynomial $f$ over $W(\mathbf F_{\ell^h})$, $\theta(\iota(f))$ is $f(\pi)$ computed
through $w_K$. Hence $\theta(\iota(E))=0$, and since every element of $S$ is a polynomial plus a multiple of
$E$, the map $\theta\circ\iota$ induces the inclusion $\mathcal O_K\subset\mathcal O_{\mathbf C_\ell}$ through
$S/(E)\cong\mathcal O_K$.

\emph{The kernel of $\theta$.} On $S$ and on $A_{\inf}$, which have no $\ell$-torsion, put
$\delta(x)=(\phi(x)-x^\ell)/\ell$; then $\iota\circ\delta=\delta\circ\iota$. Write $E(0)=\ell c$ with $c$ a unit
of $W(\mathbf F_{\ell^h})$, and $\sigma$ for the Frobenius of $W(\mathbf F_{\ell^h})$. The constant coefficient
of $\delta(E)$ is $\sigma(c)-\ell^{\ell-1}c^\ell$, a unit, so $\delta(E)$ is a unit of $S$. The kernel of
$\theta$ is generated by an element $\xi$ with $\theta(\phi(\xi))=\ell$ \cite{Fon94a}. Write
$\iota(E)=\xi y$ with $y\in A_{\inf}$. Applying $\theta$ to
$\phi(\iota(E))=\iota(E)^\ell+\ell\,\iota(\delta(E))$ gives $\ell\,\theta(\phi(y))=\ell\,\theta(\iota(\delta(E)))$,
so $\theta(\phi(y))$ is a unit of $\mathcal O_{\mathbf C_\ell}$. Hence $\phi(y)$ is a unit modulo $\xi$, and
so a unit, since $A_{\inf}$ is $\xi$-adically complete. As $\phi$ is bijective, $y$ is a unit, and
$\iota(E)$ generates $\ker\theta$.
\end{proof}

\begin{lemma}\label{gal-structures-twist}
Let $\ell$ be a prime, and let $A_{\inf}$ and $\phi$ be as in Definition~\ref{gal-cover-breuil-kisin}. For
every unit $\omega$ of $A_{\inf}$ there is a unit $\eta$ of $A_{\inf}$ with $\phi(\eta)=\omega\eta$.
\end{lemma}
\leanref{Development/AbelianVarietyTheory/Proofs/InertiaEigenvalues/ExistsAinfUnitFrobeniusEqMulStatement.lean}{GlobalTheory.T.exists_ainf_unit_frobenius_eq_mul}
\leanref{Development/AbelianVarietyTheory/Proofs/InertiaEigenvalues/ExistsAinfUnitFrobeniusEqMul.lean}{GlobalTheory.Per.LangUnit.lang_unit}

\begin{proof}
Let $\omega_0\in\mathcal O^\flat$ be the zeroth Witt component of $\omega$, a unit. The ring
$\mathcal O^\flat$ is integrally closed in the algebraically closed field $\mathbf C_\ell^\flat$, so every
monic polynomial of positive degree over $\mathcal O^\flat$ has a root in $\mathcal O^\flat$. The map
$F(x)=\phi(x)-\omega x$ of $A_{\inf}$ is additive, with $F(\ell x)=\ell F(x)$. It is onto modulo $\ell$: for
$y\in A_{\inf}$ with zeroth component $y_0$, a root $x_0$ of $X^\ell-\omega_0X-y_0$ gives
$F([x_0])\equiv y$ modulo $\ell$. Since $A_{\inf}$ is $\ell$-adically complete, successive approximation
shows that $F$ is onto. Choose $t_0\in\mathcal O^\flat$ with $t_0^{\ell-1}=\omega_0$. It is a unit, and
$t_0^\ell=\omega_0t_0$, so $F([t_0])$ has zeroth component $0$ and equals $\ell Y$ for some
$Y\in A_{\inf}$. Choose $Z$ with $F(Z)=Y$ and put $\eta=[t_0]-\ell Z$. Then $F(\eta)=0$, and $\eta$ is a unit,
being congruent to the unit $[t_0]$ modulo $\ell$.
\end{proof}

\begin{lemma}\label{gal-structures-cyclotomic}
Let $\ell$ be a prime and $K\subset\bar{\mathbf Q}_\ell$ a finite extension of $\mathbf Q_\ell$ with a
Breuil--Kisin datum $(S,E,\iota)$ (Definition~\ref{gal-cover-breuil-kisin}). Then the datum has a cyclotomic
element. Every cyclotomic element $\tau$ for it is a unit of $W(\mathbf C_\ell^\flat)$, and so is $\iota(E)$.
\end{lemma}
\leanref{Development/AbelianVarietyTheory/Proofs/InertiaEigenvalues/ExistsBreuilKisinFrobeniusPeriod.lean}{GlobalTheory.BSP.exists_isPer}
\leanref{Development/AbelianVarietyTheory/Proofs/IwasawaDecomposition/FontaineThetaEqZeroIffFontaineXiElementDvd.lean}{GlobalTheory.MemLoc.ker_theta_eq_span_xi}
\leanref{Development/AbelianVarietyTheory/Proofs/InertiaEigenvalues/BreuilKisinPeriodIsUnitWittTiltAndNeZero.lean}{GlobalTheory.BSP.isUnit_inclW_of_isPer}

\begin{proof}
A compatible system of primitive $\ell^n$-th roots of unity in $\mathcal O_{\mathbf C_\ell}$ defines
$\zeta\in\mathcal O^\flat$ with $\zeta^\sharp=1$ and $(\zeta^{1/\ell})^\sharp\neq1$. Put $x=[\zeta^{1/\ell}]$
and $\xi=\sum_{i=0}^{\ell-1}x^i$. Then $\phi(x)=x^\ell=[\zeta]$ and $\xi(x-1)=[\zeta]-1$, and $\xi$ generates
$\ker\theta$ \cite{Fon94a}. Since $\iota(E)$ also generates $\ker\theta$ and $A_{\inf}$ is a domain,
$\iota(E)=\omega\xi$ for a unit $\omega$. Lemma~\ref{gal-structures-twist} gives a unit $\eta$ with
$\phi(\eta)=\omega\eta$. Put $\tau=(x-1)\eta$. Then
$\phi(\tau)=([\zeta]-1)\omega\eta=\omega\xi(x-1)\eta=\iota(E)\tau$, and $\phi(\tau)=\omega\eta([\zeta]-1)$ with
$\omega\eta$ a unit, so $\tau$ is a cyclotomic element.

Let $\tau$ be any cyclotomic element, with $\phi(\tau)=\eta'([\zeta']-1)$ as in
Definition~\ref{gal-cover-breuil-kisin}. Taking zeroth Witt components gives $\tau_0^\ell=\eta'_0(\zeta'-1)$.
Here $\zeta'\neq1$: otherwise $\zeta'^{1/\ell}=1$, as $\mathcal O^\flat$ is perfect, against
$(\zeta'^{1/\ell})^\sharp\neq1$. So $\tau_0\neq0$. An element of $W(\mathbf C_\ell^\flat)$ with nonzero zeroth
component is a unit, as $\mathbf C_\ell^\flat$ is a field. Hence $\tau$ and $\phi(\tau)$ are units of
$W(\mathbf C_\ell^\flat)$, and so is $\iota(E)=\phi(\tau)\tau^{-1}$.
\end{proof}

\begin{lemma}\label{gal-structures-integral}
Let $\ell$ be a prime, with the notation of Definition~\ref{gal-cover-breuil-kisin}, and let $p,q\geq1$. Let
$Z\in M_{p\times q}(W(\mathbf C_\ell^\flat))$, $B\in M_p(A_{\inf})$, $B'\in M_q(A_{\inf})$ and
$C\in M_{p\times q}(A_{\inf})$ satisfy $\phi(Z)=BZB'+C$, with $\phi$ applied entrywise. Then $Z$ has entries in
$A_{\inf}$.
\end{lemma}
\leanref{Development/AbelianVarietyTheory/Proofs/InertiaEigenvalues/FrobeniusIntegralityWittTiltAinf.lean}{GlobalTheory.FrobInt.frobenius_integrality}
\leanref{Development/AbelianVarietyTheory/Proofs/InertiaEigenvalues/FrobRowIntegral.lean}{GlobalTheory.HeightInt.frob_row_integral}

\begin{proof}
Two elements of $W(\mathbf C_\ell^\flat)$ are congruent modulo $\ell^n$ if and only if their first $n$ Witt
components agree. So it suffices to show that $Z$ is congruent modulo $\ell^n$ to a matrix over $A_{\inf}$
for every $n$. Let $\bar Z$, $\bar B$, $\bar B'$ and $\bar C$ be the matrices of zeroth Witt components. Modulo
$\ell$ the equation reads $\bar Z^{(\ell)}=\bar B\bar Z\bar B'+\bar C$, where $\bar Z^{(\ell)}$ has the $\ell$-th
powers of the entries of $\bar Z$. If an entry of $\bar Z$ lies outside $\mathcal O^\flat$, let $-v<0$ be the
least valuation of an entry of $\bar Z$. Then the entries of the right side have valuation at least $-v$,
and an entry of the left side has valuation $-\ell v$, which is impossible. So $\bar Z$ is over
$\mathcal O^\flat$, and
$Z=[\bar Z]+\ell Z_1$, where $[\bar Z]$ is the matrix of Teichm\"uller representatives. The matrix
$\ell(\phi(Z_1)-BZ_1B')=B[\bar Z]B'+C-\phi([\bar Z])$ is over $A_{\inf}$ and divisible by $\ell$ in
$W(\mathbf C_\ell^\flat)$, so it is $\ell C_1$ with $C_1$ over $A_{\inf}$, and $\phi(Z_1)=BZ_1B'+C_1$. By
induction on $n$, $Z_1$ is congruent modulo $\ell^n$ to a matrix over $A_{\inf}$, and then $Z$ is congruent
modulo $\ell^{n+1}$ to one.
\end{proof}

\begin{definition}\label{gal-structures-basis}
Let $\ell$ be a prime and $K\subset\bar{\mathbf Q}_\ell$ a finite extension of $\mathbf Q_\ell$ with a
Breuil--Kisin datum $(S,E,\iota)$ and element $\pi^\flat$ (Definition~\ref{gal-cover-breuil-kisin}). Let
$\rho:G_K\to M_d(\mathbf Z_\ell)$ be a map and $N\geq0$. A \emph{Breuil--Kisin basis} of height $N$ of $\rho$
consists of $T\in\operatorname{GL}_d(W(\mathbf C_\ell^\flat))$ and $A,A'\in M_d(S)$ with $AA'=E^{2N}$,
$\iota(E)^N\phi(T)=T\iota(A)$, and $\rho(g)g(T)=T$ for every $g\in G_K$ that fixes $\pi^\flat$. A
Breuil--Kisin structure $(T,A,A',\varepsilon)$ of height $N$ on $\rho$ has the Breuil--Kisin basis
$(T,A,A')$.
\end{definition}
\leanref{Development/AbelianVarietyTheory/Definitions/NonGaloisDescent/TrivIsCr.lean}{GlobalTheory.BKLim.Pre}

\begin{lemma}\label{gal-structures-poles}
Let $\ell$ be a prime and $K\subset\bar{\mathbf Q}_\ell$ a finite extension of $\mathbf Q_\ell$ with a
Breuil--Kisin datum $(S,E,\iota)$ and a cyclotomic element $\tau$ (Definition~\ref{gal-cover-breuil-kisin}).
Let $\rho:G_K\to M_d(\mathbf Z_\ell)$ be a map, $N\geq0$, and $(T,A,A')$ a Breuil--Kisin basis of height $N$
of $\rho$ (Definition~\ref{gal-structures-basis}). Then there are $Y,Y'\in M_d(A_{\inf})$ with $Y=\tau^NT$ in
$M_d(W(\mathbf C_\ell^\flat))$ and $YY'=\tau^{2N}$.
\end{lemma}
\leanref{Development/AbelianVarietyTheory/Proofs/InertiaEigenvalues/BreuilKisinPreframeBoundedPoles.lean}{GlobalTheory.BKLim.pre_bounded_poles}

\begin{proof}
By Lemma~\ref{gal-structures-cyclotomic}, $\tau$ and $\iota(E)$ are units of $W(\mathbf C_\ell^\flat)$. Since
$AA'=E^{2N}$, $\phi(T)^{-1}=\iota(E)^N\iota(A)^{-1}T^{-1}=\iota(E)^{-N}\iota(A')T^{-1}$, that is,
$\iota(E)^N\phi(T^{-1})=\iota(A')T^{-1}$. Put $Y=\tau^NT$ and $Y'=\tau^NT^{-1}$. As $\phi(\tau)=\iota(E)\tau$,
we get $\phi(Y)=Y\iota(A)$ and $\phi(Y')=\iota(A')Y'$. Lemma~\ref{gal-structures-integral}, with $C=0$, shows
that $Y$ and $Y'$ have entries in $A_{\inf}$, and $YY'=\tau^{2N}$.
\end{proof}

\begin{lemma}\label{gal-structures-saturated}
Let $\ell$ be a prime, with the notation of Definition~\ref{gal-cover-breuil-kisin}. Let $c\in A_{\inf}$ be a
unit of $W(\mathbf C_\ell^\flat)$, let $r\geq0$ and let $x\in W(\mathbf C_\ell^\flat)$. Suppose that for every
$n\geq0$ there is $a_n\in A_{\inf}$ with $\ell^rx\equiv ca_n$ modulo $\ell^nW(\mathbf C_\ell^\flat)$. Then
$x\in cA_{\inf}$.
\end{lemma}
\leanref{Development/AbelianVarietyTheory/Proofs/InertiaEigenvalues/AinfPrincipalIdealSaturatedClosedStatement.lean}{GlobalTheory.T.ainf_principal_ideal_saturated_closed}
\leanref{Development/AbelianVarietyTheory/Proofs/InertiaEigenvalues/AinfPrincipalIdealSaturatedClosed.lean}{GlobalTheory.AinfSat.mem_ideal_of_saturated}

\begin{proof}
Put $y=c^{-1}x$. Then $\ell^ry\equiv a_n$ modulo $\ell^n$ for every $n$. Two elements of
$W(\mathbf C_\ell^\flat)$ are congruent modulo $\ell^n$ if and only if their first $n$ Witt components agree.
Since $\mathbf C_\ell^\flat$ is perfect of characteristic $\ell$, the Witt components of $\ell^ry$ are $r$
zeros followed by $y_0^{\ell^r},y_1^{\ell^r},\ldots$. Hence every $y_i^{\ell^r}$ lies in $\mathcal O^\flat$,
so every $y_i$ does, and $y\in A_{\inf}$.
\end{proof}

\begin{lemma}\label{gal-structures-layer}
Let $\ell$ be a prime and $K\subset\bar{\mathbf Q}_\ell$ a finite extension of $\mathbf Q_\ell$, with
$G_K=\operatorname{Gal}(\bar{\mathbf Q}_\ell/K)$. Let $L_0\subset\bar{\mathbf Q}_\ell$ be a finite extension of
$\mathbf Q_\ell$ containing $\sigma(K)$ for every embedding $\sigma:K\to\bar{\mathbf Q}_\ell$, let $a\geq0$, let
$C$ be a set of homomorphisms $G_K\to\operatorname{GL}_a(\bar{\mathbf Q}_\ell)$, let $J$ be a map from the
embeddings $K\to\bar{\mathbf Q}_\ell$ to finite multisets of integers, and let $N\geq0$ bound the absolute
values of the members of every $J(\sigma)$. Suppose that for every $\rho_1\in C$ and every finite extension
$L\supset L_0$ of $\mathbf Q_\ell$ in $\bar{\mathbf Q}_\ell$ with $\rho_1(G_K)\subset\operatorname{GL}_a(L)$, the
representation $\rho_1$ of $G_K$ on $L^a$ is crystalline with Hodge--Tate multiset $J(\sigma)$ at every
embedding $\sigma:K\to L$, in the sense of Definition~\ref{gal-p-adic-hodge}. Let
$\rho':G_K\to\operatorname{GL}_{d'}(\mathbf Z_\ell)$ be a homomorphism with a layer relative to $C$ and $L_0$
(Definition~\ref{gal-cover-definition}). Then the representation $V'$ of $G_K$ on $\mathbf Q_\ell^{d'}$
through $\rho'$ is crystalline, $\operatorname{Fil}^{-N}D_{\mathrm{dR}}(V')=D_{\mathrm{dR}}(V')$ and
$\operatorname{Fil}^{N+1}D_{\mathrm{dR}}(V')=0$.
\end{lemma}
\leanref{Development/AbelianVarietyTheory/Proofs/InertiaEigenvalues/IsSemistableLatticeInMatrixAlgebraLayerStatement.lean}{GlobalTheory.T.isSemistable_latticeInMatrixAlgebraLayer}
\leanref{Development/AbelianVarietyTheory/Proofs/InertiaEigenvalues/IsSemistableLatticeInMatrixAlgebraLayer.lean}{GlobalTheory.LayerSst.main}
\leanref{Development/AbelianVarietyTheory/Proofs/InertiaEigenvalues/LayerDdRBoundsStatement.lean}{GlobalTheory.T.layer_DdR_bounds_alt}
\leanref{Development/AbelianVarietyTheory/Proofs/InertiaEigenvalues/LayerDdRBounds.lean}{GlobalTheory.LayerJumps.layer_DdR_bounds}

\begin{proof}
Let the layer consist of $\rho_1,\ldots,\rho_m\in C$, fields $L_1,\ldots,L_m$, elements
$b'_j=(b'_{j,k})_k$ and $e'$. The map $\mathbf Q_\ell^{d'}\to\prod_kM_a(L_k)$, $v\mapsto\sum_jv_jb'_j$, is an
isomorphism of $\mathbf Q_\ell$-vector spaces, since the $b'_j$ form a basis. As $\rho'$ reads every $\rho_k$
through the $b'_{j,k}$, it carries the action of $\rho'(g)$ to left multiplication by $(\rho_k(g))_k$. Left
multiplication on $M_a(L_k)$ is the direct sum, over the $a$ columns, of the representation of $G_K$ on
$L_k^a$ through $\rho_k$, viewed over $\mathbf Q_\ell$. So $V'$ is a direct sum of representations that are
crystalline by hypothesis, and it is crystalline \cite{Fon94}. Its module $D_{\mathrm{dR}}(V')$ is the
corresponding direct sum, with the direct sum of the Hodge filtrations. The jumps of the Hodge filtration on
$D_{\mathrm{dR}}$ of $L_k^a$, a module over $K\otimes_{\mathbf Q_\ell}L_k=\prod_{\sigma:K\to L_k}L_k$, are
the members of the Hodge--Tate multisets $J(\sigma)$, which lie in $[-N,N]$. Hence
$\operatorname{Fil}^{-N}D_{\mathrm{dR}}(V')=D_{\mathrm{dR}}(V')$ and
$\operatorname{Fil}^{N+1}D_{\mathrm{dR}}(V')=0$.
\end{proof}

\begin{lemma}\label{gal-structures-kisin}
Let $\ell$ be a prime and $K\subset\bar{\mathbf Q}_\ell$ a finite extension of $\mathbf Q_\ell$ with a
Breuil--Kisin datum (Definition~\ref{gal-cover-breuil-kisin}). Let $N\geq0$ and let
$\rho:G_K\to\operatorname{GL}_d(\mathbf Z_\ell)$ be a continuous homomorphism such that the representation $V$
of $G_K$ on $\mathbf Q_\ell^d$ through $\rho$ is crystalline, with
$\operatorname{Fil}^{-N}D_{\mathrm{dR}}(V)=D_{\mathrm{dR}}(V)$ and
$\operatorname{Fil}^{N+1}D_{\mathrm{dR}}(V)=0$. Then $\rho$ has a crystalline Breuil--Kisin structure of
height $N$.
\end{lemma}
\leanref{Development/AbelianVarietyTheory/Proofs/InertiaEigenvalues/ExistsBreuilKisinFrameSemistableBoundedJumpsStatement.lean}{GlobalTheory.T.exists_breuilKisinFrame_semistable_boundedJumps}
\leanref{Development/AbelianVarietyTheory/Proofs/InertiaEigenvalues/ExistsBreuilKisinFrameSemistableBoundedJumps.lean}{GlobalTheory.BKLim.exists_frame_sst}

\begin{proof}
Let $L$ be $\mathbf Z_\ell^d$ with $G_K$ acting through $\rho$, and let $G_\infty\subset G_K$ be the subgroup
fixing $\pi^\flat$. In terms of modules, a Breuil--Kisin structure of height $N$ on $\rho$ is an $S$-lattice
$M$ in $L\otimes W(\mathbf C_\ell^\flat)$, spanned by the columns of $T$, that is fixed by $G_\infty$, satisfies
$\iota(E)^N\phi(M)\subset M$ with the matrices $A$ and $A'$, and whose $A_{\inf}$-span is stable under $G_K$,
with $\varepsilon(g)$ the matrix of $g$ on it. Twist $L$ by the power of the cyclotomic character that
moves the Hodge--Tate weights from $[-N,N]$ into $[0,2N]$ in Kisin's normalization. Kisin's theorem
\cite[Cor.~(1.3.15), Lem.~(2.1.15)]{Kis06}, applied through his contravariant functor to the dual of the
twist, gives such a lattice for the twist: a Breuil--Kisin module of height $2N$, embedded compatibly with
$\phi$ and $G_\infty$ by Fontaine's theory of \'etale $\phi$-modules \cite{Fon90}. For $\ell$ odd, by Liu's theory
of $(\phi,\hat G)$-modules \cite[Thm.~2.3.1]{Liu10}, the $A_{\inf}$-span of the Frobenius twist of the
Breuil--Kisin module is stable under $G_K$; since $\phi$ is an automorphism of $A_{\inf}$ commuting with
$G_K$, so is the $A_{\inf}$-span of the module itself. For every $\ell$, including $\ell=2$, Gao's theory of
Breuil--Kisin $G_K$-modules \cite[Thm.~7.1.7]{Gao23} shows directly that the $A_{\inf}$-span of the Breuil--Kisin
module is stable under $G_K$. Undoing the twist with a period of the cyclotomic
character, to which the same results apply, gives $T$, $A$, $A'$ and $\varepsilon$. For $\ell$ odd, Gee, Liu
and Savitt \cite[Prop.~5.9]{GLS14} show that, for a lattice in a crystalline representation, a topological
generator of $\operatorname{Gal}(K_\infty(\mu_{\ell^\infty})/K(\mu_{\ell^\infty}))$, where $K_\infty$ is the
field fixed by $G_\infty$, moves every element of the Breuil--Kisin module by an element of
$[\pi^\flat]^\ell\phi(t)$ times the Frobenius twist of its $A_{\inf}$-span, for an element $t\in A_{\inf}$.
Applying $\phi^{-1}$, a lift of this generator to $G_K$ acts on the $A_{\inf}$-span of the module by a matrix
congruent to $1$ modulo $[\pi^\flat]$. For every $\ell$, including $\ell=2$, Du and Liu choose
$\tilde\tau\in G_K$ with $\tilde\tau([\pi^\flat])=[\zeta][\pi^\flat]$, for an element $\zeta\in\mathcal O^\flat$
given by a compatible system of primitive $\ell^n$-th roots of unity, such that $\tilde\tau$ and $G_\infty$
generate a dense subgroup of $G_K$ \cite[\S3.2]{DL26}. As $\theta([\zeta])=1$, the difference
$\tilde\tau([\pi^\flat])-[\pi^\flat]$ lies in $\ker\theta=\iota(E)A_{\inf}$. For a lattice in a crystalline
representation, $\tilde\tau$ then acts on the $A_{\inf}$-span of the Breuil--Kisin module by a matrix with
entries in a subring $A^{(2)}$ of $A_{\inf}$ \cite[Thm.~3.2.1, Cor.~2.4.5]{DL26}. Let $\mathfrak m^\flat$ be the
maximal ideal of $\mathcal O^\flat$ and $W(\mathfrak m^\flat)\subset A_{\inf}$ the ideal of Witt vectors with
components in $\mathfrak m^\flat$. The matrix is congruent to $1$ modulo $W(\mathfrak m^\flat)$, by condition~(5)
of Gao's definition \cite[Def.~7.1.1, Thm.~7.1.7]{Gao23}, and $A^{(2)}\cap W(\mathfrak m^\flat)$ is contained in
$[\pi^\flat]A_{\inf}$ \cite[Lem.~3.3.1(3)]{DL26}. So $\tilde\tau$ acts by a matrix congruent to $1$ modulo
$[\pi^\flat]$, as in the proof of \cite[Cor.~3.3.4]{DL26}. In both cases $G_\infty$ acts trivially, and these
elements generate a dense subgroup of $G_K$. The matrices of the action form a continuous cocycle, and
$[\pi^\flat]$ divides $g([\pi^\flat])$ for every $g$; so every element of $G_K$ acts by a matrix congruent to
$1$ modulo $[\pi^\flat]$. The same holds for the period of the cyclotomic character, so $[\pi^\flat]$ divides
every entry of $\varepsilon(g)-1$, and the structure is crystalline.
% Lean: exists_frame_sst proves the lemma for every l by its own construction: crys_seed
% (SeedExists.exists_seed: period matrix over B_cris, Frobenius matrix over W, periods over A_cris,
% B_dR^+ lattices, modification of the filtration), exists_trivial_cocycle_crystalline_period_data,
% exists_breuilKisinFrame_semistablePeriods, and exists_breuilKisin_frame_height_of_jumps_le with the
% Hodge basis exists_hodge_basis_deRham_lattice to bring the height to N. The crystalline condition IsCr
% comes from the trivial cocycle kappa = 1 of the seed. The text cites the literature instead, as the
% depth floor allows; for l = 2 the condition is DL26 (JEMS 28 (2026)), Thm. 3.2.1 with Lem. 3.3.1(3), as
% in the proof of its Cor. 3.3.4. Gao23, Prop. 7.1.10 alone, (tau - 1)(M) in t W(m_R) (x) M, does not give
% divisibility by [pi^flat].
\end{proof}

\begin{lemma}\label{gal-structures-functorial}
Let $\ell$ be a prime and $K\subset\bar{\mathbf Q}_\ell$ a finite extension of $\mathbf Q_\ell$ with a
Breuil--Kisin datum $(S,E,\iota)$ (Definition~\ref{gal-cover-breuil-kisin}). Let
$\rho:G_K\to M_d(\mathbf Z_\ell)$ and $\rho':G_K\to M_{d'}(\mathbf Z_\ell)$ be maps with Breuil--Kisin bases
$(T,A,A')$ of height $N$ and $(T',B,B')$ of height $N'$ (Definition~\ref{gal-structures-basis}), and let
$f\in M_{d'\times d}(\mathbf Z_\ell)$ satisfy $f\rho(g)=\rho'(g)f$ for every $g\in G_K$. Then there is
$Y\in M_{d'\times d}(S)$ with $T'\iota(Y)=fT$.
\end{lemma}
\leanref{Development/AbelianVarietyTheory/Proofs/InertiaEigenvalues/ExistsEquivariantMapBreuilKisinMatrixStatement.lean}{GlobalTheory.T.exists_equivariantMap_breuilKisinMatrix}
\leanref{Development/AbelianVarietyTheory/Proofs/InertiaEigenvalues/ExistsEquivariantMapBreuilKisinMatrix.lean}{GlobalTheory.BKLim.EquivKisin.main}

\begin{proof}
This is the full faithfulness of Kisin's functor on Breuil--Kisin modules of finite height
\cite[Prop.~(2.1.12)]{Kis06}, which uses the equivariance of $f$ only for the elements of $G_K$ fixing
$\pi^\flat$. In matrices: $X=T'^{-1}fT$ is fixed by those elements and satisfies
$X=\iota(E)^{-N-N'}\iota(B)\phi(X)\iota(A')$, so it has entries in Fontaine's ring of \'etale $\phi$-modules,
the $\ell$-adic closure of $\iota(S)[\iota(u)^{-1}]$ in $W(\mathbf C_\ell^\flat)$; the integrality of $\tau^NT$
and $\tau^{N'}T'^{-1}$ for a cyclotomic element $\tau$ (Lemmas~\ref{gal-structures-cyclotomic}
and~\ref{gal-structures-poles}) then puts $X$ in $\iota(M_{d'\times d}(S))$.
\end{proof}

\begin{definition}\label{gal-structures-etale}
Let $\ell$ be a prime and $K\subset\bar{\mathbf Q}_\ell$ a finite extension of $\mathbf Q_\ell$ with a
Breuil--Kisin datum $(S,E,\iota)$ and element $\pi^\flat$ (Definition~\ref{gal-cover-breuil-kisin}), and write $u$
also for $\iota(u)=[\pi^\flat]$ in $A_{\inf}\subset W(\mathbf C_\ell^\flat)$. \emph{Fontaine's ring}
$\mathcal O_{\mathcal E}$ of the datum is the set of $w\in W(\mathbf C_\ell^\flat)$ such that for every $n\geq0$
there are $k\geq0$ and $s\in S$ with $u^kw\equiv\iota(s)$ modulo $\ell^nW(\mathbf C_\ell^\flat)$. It is a subring of
$W(\mathbf C_\ell^\flat)$ containing $\iota(S)$. The element $\pi^\flat$ is nonzero, as $\theta(\iota(E))=0$ while
$E(0)$ is $\ell$ times a unit, and it is not a unit of $\mathcal O^\flat$, as
$\theta([\pi^\flat])^e\equiv\theta(\iota(E))=0$ modulo $\ell\mathcal O_{\mathbf C_\ell}$ for the degree $e$ of $E$. As
the zeroth Witt component of $u$ is $\pi^\flat$, $u$ is a unit of $W(\mathbf C_\ell^\flat)$ and not a unit of
$A_{\inf}$.
\end{definition}
\leanref{Development/AbelianVarietyTheory/Definitions/NonGaloisDescent/IsStratHom.lean}{GlobalTheory.BSP.BK.OE}
\leanref{Development/AbelianVarietyTheory/Proofs/InertiaEigenvalues/BreuilKisinFlatRadiusIdealFacts.lean}{GlobalTheory.BSP.flat_disc}

\begin{definition}\label{gal-structures-topological}
Let $\ell$ be a prime and $K\subset\bar{\mathbf Q}_\ell$ a finite extension of $\mathbf Q_\ell$ with a
Breuil--Kisin datum $(S,E,\iota)$ (Definition~\ref{gal-cover-breuil-kisin}), and let $\mathfrak j=(\ell,E)\subset S$. A
\emph{topological basis} of $A_{\inf}$ over $S$ consists of a set $I$ with an element $i_0$, elements $b_i\in A_{\inf}$
for $i\in I$ with $b_{i_0}=1$, and additive maps $\gamma_i:A_{\inf}\to S$ for $i\in I$, such that $A_{\inf}$ is
complete and separated for the $\mathfrak jA_{\inf}$-adic topology and:
\begin{enumerate}
\item $\gamma_i(\iota(s)x)=s\gamma_i(x)$ for $s\in S$ and $x\in A_{\inf}$, $\gamma_i(b_i)=1$, and $\gamma_i(b_j)=0$ for
  $j\neq i$;
\item for every $x\in A_{\inf}$ and $n\geq0$, the set of $i$ with $\gamma_i(x)\notin\mathfrak j^n$ is finite, and
  $x-\sum_{i\in F}\iota(\gamma_i(x))b_i\in\mathfrak j^nA_{\inf}$ for every finite set $F\subset I$ containing it.
\end{enumerate}
\end{definition}
\leanref{Development/AbelianVarietyTheory/Definitions/NonGaloisDescent/MemBallOfGaugeLt.lean}{GlobalTheory.BSP.TopBasis}

\begin{lemma}\label{gal-structures-topological-exists}
Let $\ell$ be a prime and $K\subset\bar{\mathbf Q}_\ell$ a finite extension of $\mathbf Q_\ell$ with a
Breuil--Kisin datum $(S,E,\iota)$ (Definition~\ref{gal-cover-breuil-kisin}). Then $A_{\inf}$ has a topological basis
$(b_i,\gamma_i)_{i\in I}$ over $S$ (Definition~\ref{gal-structures-topological}). Every topological basis of $A_{\inf}$
over $S$ has the following properties: an $x\in A_{\inf}$ with $\gamma_i(x)=0$ for every $i$ is $0$;
$\gamma_{i_0}(\iota(s))=s$ for $s\in S$, so $\iota$ is injective; and $\gamma_i(LA_{\inf})\subset L$ for every ideal $L$
of $S$ and every $i$.
\end{lemma}
\leanref{Development/AbelianVarietyTheory/Proofs/InertiaEigenvalues/AinfAdicTopologicalBasisOverBreuilKisinStatement.lean}{GlobalTheory.T.ainf_adicTopologicalBasis_over_breuilKisin}
\leanref{Development/AbelianVarietyTheory/Proofs/InertiaEigenvalues/AinfAdicTopologicalBasisOverBreuilKisin.lean}{GlobalTheory.BSP.AinfTopBasis.main}
\leanref{Development/AbelianVarietyTheory/Proofs/InertiaEigenvalues/NilpotentOfNotIsUnit.lean}{GlobalTheory.Per.AinfComplete.main}
\leanref{Development/AbelianVarietyTheory/Proofs/InertiaEigenvalues/StrictIndepPowOfRegularPair.lean}{GlobalTheory.BSP.strict_indep_pow}
\leanref{Development/AbelianVarietyTheory/Proofs/InertiaEigenvalues/ExistsAdicTopologicalBasisOfSpanOfStrict.lean}{GlobalTheory.BSP.exists_topBasis_of_span_of_strict}
\leanref{Development/AbelianVarietyTheory/Proofs/InertiaEigenvalues/TopologicalBasisBasicFacts.lean}{GlobalTheory.BSP.TopBasis.basic_facts}

\begin{proof}
Write $u$ also for $\iota(u)=[\pi^\flat]$. Let $\mathfrak m=(\ell,u)$, the maximal ideal of $S$, with
$S/\mathfrak m=\mathbf F_{\ell^h}$, let $e$ be the degree of $E$, and $\mathfrak j=(\ell,E)$. Since the lower
coefficients of $E$ are divisible by $\ell$, $\mathfrak j=(\ell,u^e)$, so $\mathfrak m^{1+e}\subset\mathfrak j\subset
\mathfrak m$. The ring $S$ is complete and separated for the $\mathfrak m$-adic topology.

\emph{Completeness of $A_{\inf}$.} The ring $A_{\inf}$ is complete and separated for the $(\ell,y)$-adic topology for
every non-unit $y$ of $A_{\inf}$. Indeed, view $\mathcal O^\flat$ as the ring of sequences $(x_k)_{k\geq0}$ in
$\mathcal O_{\mathbf C_\ell}/\ell$ with $x_{k+1}^\ell=x_k$. The zeroth Witt component $y_0$ of $y$ is not a unit, as an
element of $A_{\inf}$ whose zeroth component is a unit is a unit; so neither is the coordinate $(y_0)_0$, and as a
non-unit of $\mathcal O_{\mathbf C_\ell}$ has a power divisible by $\ell$, $(y_0)_0^m=0$ for some $m$. Then the $k$-th
coordinate of $y_0^N$ vanishes for $N\geq\ell^km$. So $\mathcal O^\flat$ is $y_0$-adically separated, and every series
$\sum_iy_0^ib_i$ converges coordinatewise. As $A_{\inf}/\ell A_{\inf}=\mathcal O^\flat$ and $A_{\inf}$ is $\ell$-adically
complete, lifting gives that every series $\sum_iy^ib_i$ with $b_i\in A_{\inf}$ converges $y$-adically, and that
$A_{\inf}$ is $(\ell,y)$-adically separated; as $(\ell,y)^{2n}\subset(\ell^n,y^n)$, $A_{\inf}$ is
$(\ell,y)$-adically complete. Since $u$ is not a unit (Definition~\ref{gal-structures-etale}), $A_{\inf}$ is complete
and separated for the $\mathfrak mA_{\inf}$-adic topology. As $\mathfrak m^{1+e}\subset\mathfrak j\subset\mathfrak m$, the
$\mathfrak j$-adic and $\mathfrak m$-adic topologies agree on $S$ and on $A_{\inf}$.

\emph{A basis modulo $\mathfrak m$.} The map $\theta\circ\iota$ sends $\ell$ and $E$, hence $u^e$ and $u$, into the
maximal ideal of $\mathcal O_{\mathbf C_\ell}$, so $\mathfrak mA_{\inf}\neq A_{\inf}$, and $A_{\inf}/\mathfrak mA_{\inf}$ is a
nonzero vector space over $S/\mathfrak m$. Choose a basis of it that contains the class of $1$, and lift it to elements
$b_i\in A_{\inf}$, $i\in I$, with $b_{i_0}=1$. Then every $x\in A_{\inf}$ is congruent modulo $\mathfrak mA_{\inf}$ to a
finite sum $\sum_i\iota(a_i)b_i$, and if a finite sum $\sum_i\iota(c_i)b_i$ lies in $\mathfrak mA_{\inf}$, then every
$c_i$ lies in $\mathfrak m$.

\emph{Independence modulo powers of $\mathfrak m$.} The element $\ell$ is not a zero divisor of $A_{\inf}$, and if
$\ell$ divides $ux$ in $A_{\inf}$, then $\ell$ divides $x$, as the zeroth Witt component of $ux$ is $\pi^\flat x_0$.
Hence if $u^n$ divides $\ell z$, then $u^n$ divides $z$: from $\ell z=u^nw$ we get $w=\ell t$ and $z=u^nt$. Let
$(c_i)$ be a finitely supported family in $S$ and $X=\sum_i\iota(c_i)b_i$. First, if $X\in(\ell,u^k)A_{\inf}$, then
every $c_i$ lies in $(\ell,u^k)$. This is clear for $k=0$. If $X=\ell x+u^{k+1}y$, then by induction
$c_i=\ell p_i+u^kq_i$, and with $P=\sum_i\iota(p_i)b_i$ and $Q=\sum_i\iota(q_i)b_i$ we get $u^k(Q-uy)=\ell(x-P)$. So
$\ell$ divides $Q-uy$, $Q$ lies in $\mathfrak mA_{\inf}$, every $q_i$ lies in $\mathfrak m$, and
$c_i\in(\ell,u^{k+1})$. Second, if $X\in(\ell^k,u^n)A_{\inf}$, then every $c_i$ lies in $(\ell^k,u^n)$. This is clear
for $k=0$. If $X=\ell^{k+1}x+u^ny$, then by induction $c_i=\ell^kp_i+u^nq_i$, and $\ell^k(P-\ell x)=u^n(y-Q)$. So
$u^n$ divides $P-\ell x$, every $p_i$ lies in $(\ell,u^n)$ by the first claim, and $c_i\in(\ell^{k+1},u^n)$. Since
$\mathfrak m^{2n}\subset(\ell^n,u^n)\subset\mathfrak m^n$, every $c_i$ lies in $\mathfrak m^n$ if
$X\in\mathfrak m^{2n}A_{\inf}$.

\emph{Expansions.} Call a family $(f_i)_{i\in I}$ in $S$ an expansion of $x\in A_{\inf}$ if for every $n$ there is a
finite set $F$ with $f_i\in\mathfrak m^n$ for $i\notin F$ and $x-\sum_{i\in F}\iota(f_i)b_i\in\mathfrak m^nA_{\inf}$. The
same then holds for every finite set $F'$ with $f_i\in\mathfrak m^n$ for $i\notin F'$, as the two sums differ by terms
$\iota(f_i)b_i$ with $f_i\in\mathfrak m^n$. Every $x$ has an expansion. Indeed, an element of $\mathfrak m^nA_{\inf}$ is
a finite sum $\sum_k\iota(t_k)w_k$ with $t_k\in\mathfrak m^n$, and approximating each $w_k$ modulo $\mathfrak mA_{\inf}$
shows that it is congruent modulo $\mathfrak m^{n+1}A_{\inf}$ to a finite sum $\sum_i\iota(a_i)b_i$ with every
$a_i\in\mathfrak m^n$. Applying this to $x$ and to the successive remainders gives finitely supported families
$(a^{(k)}_i)_i$ with $a^{(k)}_i\in\mathfrak m^k$ and $x-\sum_{k<n}\sum_i\iota(a^{(k)}_i)b_i\in\mathfrak m^nA_{\inf}$ for every
$n$. The sums $f_i=\sum_ka^{(k)}_i$ converge in $S$, and $(f_i)$ is an expansion of $x$. An expansion $(f_i)$ of $0$
vanishes: given $i$ and $n$, take $F$ for $2n$; then $\sum_{j\in F}\iota(f_j)b_j\in\mathfrak m^{2n}A_{\inf}$, so
$f_i\in\mathfrak m^n$ if $i\in F$ by the independence, and $f_i\in\mathfrak m^{2n}$ otherwise; as $S$ is separated,
$f_i=0$. Sums and $S$-multiples of expansions are expansions, so every $x$ has a unique expansion
$(\gamma_i(x))_i$, and the maps $\gamma_i$ are additive with $\gamma_i(\iota(s)x)=s\gamma_i(x)$. The family
$(\delta_{ij})_i$ is the expansion of $b_j$, so $\gamma_i(b_j)=\delta_{ij}$.

\emph{The ideal $\mathfrak j$.} Given $x$ and $n$, the set of $i$ with $\gamma_i(x)\notin\mathfrak j^n$ lies in the finite
set $F_0$ of $i$ with $\gamma_i(x)\notin\mathfrak m^{(1+e)n}$, as $\mathfrak m^{(1+e)n}\subset\mathfrak j^n$. For a finite
set $F$ containing it, $x-\sum_{i\in F\cup F_0}\iota(\gamma_i(x))b_i$ lies in
$\mathfrak m^{(1+e)n}A_{\inf}\subset\mathfrak j^nA_{\inf}$, and so do the terms with $i\in F_0\setminus F$. With the
completeness shown first, $(b_i,\gamma_i)_{i\in I}$ is a topological basis of $A_{\inf}$ over $S$.

Let now $(b_i,\gamma_i)_{i\in I}$ be any topological basis. If every $\gamma_i(x)$ vanishes, then
$x\in\mathfrak j^nA_{\inf}$ for every $n$, by property~(2) with $F=\emptyset$, so $x=0$. Next,
$\gamma_{i_0}(\iota(s))=s\gamma_{i_0}(b_{i_0})=s$. Finally, an element of $LA_{\inf}$ is a finite sum
$\sum_k\iota(t_k)w_k$ with $t_k\in L$, and $\gamma_i$ sends it to $\sum_kt_k\gamma_i(w_k)\in L$.
\end{proof}

\begin{lemma}\label{gal-structures-frobenius-unit}
Let $\ell$ be a prime and $K\subset\bar{\mathbf Q}_\ell$ a finite extension of $\mathbf Q_\ell$ with a
Breuil--Kisin datum $(S,E,\iota)$ (Definition~\ref{gal-cover-breuil-kisin}). Let $g\in S$ be an element not divisible
by $\ell$ such that $\phi(g)$ divides $E^Mg$ in $S$ for some $M\geq0$. Then $g$ is a unit of $S$.
\end{lemma}
\leanref{Development/AbelianVarietyTheory/Proofs/InertiaEigenvalues/IsUnitOfFrobeniusDvdEisensteinPowMulStatement.lean}{GlobalTheory.T.isUnit_of_frobenius_dvd_eisenstein_pow_mul}
\leanref{Development/AbelianVarietyTheory/Proofs/InertiaEigenvalues/IsUnitOfFrobeniusDvdEisensteinPowMul.lean}{GlobalTheory.LemZ.isUnit_of_frob_dvd}

\begin{proof}
Write $W=W(\mathbf F_{\ell^h})$, $\sigma$ for its Frobenius and $K_0$ for its fraction field, and fix an embedding
$K_0\subset\bar{\mathbf Q}_\ell$. Call a monic polynomial over $W$ whose lower coefficients are divisible by $\ell$
\emph{distinguished}. As $\ell$ does not divide $g$, the Weierstrass preparation theorem \cite[\S7.1]{Was97} gives
$g=UP$ with $U$ a unit of $S$ and $P$ a distinguished polynomial. If $P=1$, then $g$ is a unit. Suppose that $P$ has degree $d>0$.

The power series $\phi(P)$ is the polynomial $P^\sigma(u^\ell)$, where $P^\sigma$ is obtained by applying $\sigma$ to
the coefficients of $P$, and it is distinguished. As $U$ and $\phi(U)$ are units, $P^\sigma(u^\ell)$ divides $E^MP$
in $S$. A distinguished polynomial that divides a polynomial in $S$ divides it in $W[u]$, by the uniqueness of division
with remainder \cite[\S7.1]{Was97}; so $P^\sigma(u^\ell)$ divides $E^MP$ in $W[u]$.

First, $P(0)\neq0$. Write $P=u^kP_0$ with $P_0(0)\neq0$. If $k\geq1$, then $u^{k\ell}$ divides $P^\sigma(u^\ell)$ and
$k\ell\geq k+1$, so $u^{k+1}$ divides $E^Mu^kP_0$, and $u$ divides $E^MP_0$. This is impossible, as $E(0)\neq0$.

Next, every root $x\in\bar{\mathbf Q}_\ell$ of $P$ has $|x|<1$: if $|x|\geq1$, then
$|x|^d=|\sum_{i<d}p_ix^i|\leq\ell^{-1}|x|^d$, where the $p_i$ are the lower coefficients of $P$, which is
impossible. Let $x$ be a root of $P$
of largest absolute value $r$. Then $0<r<1$, as $P(0)\neq0$.

Let $\Phi$ be an automorphism of $\bar{\mathbf Q}_\ell$ that extends $\sigma$ on $K_0$. It preserves absolute values,
by the uniqueness of the extension of the absolute value \cite[Chapter~II, \S2]{Ser79}, and for the same reason
conjugates over $K_0$ have the same absolute value. So $z=\Phi(x)$ is a root of $P^\sigma$ with $|z|=r$. Let $m$ be
the minimal polynomial of $z$ over $K_0$; it divides $P^\sigma$. Let $z'$ be a root of $m$ and
$y\in\bar{\mathbf Q}_\ell$ with $y^\ell=z'$. Then $|z'|=r$, and $z'$ is a root of $P^\sigma$, so $y$ is a root of
$P^\sigma(u^\ell)$ and hence of $E^MP$. As $|y|=r^{1/\ell}>r$, $y$ is not a root of $P$, so $y$ is a root of $E$.

So every root of $m(u^\ell)$ is a root of $E$. The polynomial $m(u^\ell)$ is separable, since $m$ is separable,
$m(0)\neq0$ and $\ell$ is invertible in $K_0$. So it has $\ell\deg m$ distinct roots, all of them roots of $E$, and
$\ell\deg m\leq\deg E$. By Eisenstein's criterion $E$ is irreducible over $K_0$ \cite[Chapter~I, \S6]{Ser79}, so it is
the minimal polynomial of each of its roots, and it divides $m(u^\ell)$. Both are monic, so $E=m(u^\ell)$ is a
polynomial in $u^\ell$, against the definition of a Breuil--Kisin datum. Hence $d=0$.
\end{proof}

\begin{lemma}\label{gal-structures-coefficients}
Let $\ell$ be a prime and $K\subset\bar{\mathbf Q}_\ell$ a finite extension of $\mathbf Q_\ell$ with a
Breuil--Kisin datum $(S,E,\iota)$ and a cyclotomic element $\tau$ (Definition~\ref{gal-cover-breuil-kisin}), and let
$(b_i,\gamma_i)_{i\in I}$ be a topological basis of $A_{\inf}$ over $S$ (Definition~\ref{gal-structures-topological}).
Then for every $M\geq0$ there is $r\geq0$ such that $\ell^r$ lies in the ideal of $S$ generated by the
$\gamma_i(\tau^M)$, $i\in I$.
\end{lemma}
\leanref{Development/AbelianVarietyTheory/Proofs/InertiaEigenvalues/ExistsPowMemSpanCoeffPeriodPowStatement.lean}{GlobalTheory.T.exists_pow_mem_span_coeff_period_pow}
\leanref{Development/AbelianVarietyTheory/Proofs/InertiaEigenvalues/ExistsPowMemSpanCoeffPeriodPow.lean}{GlobalTheory.BSP.pow_l_mem_span_coeff}

\begin{proof}
Let $\mathfrak b$ be the ideal generated by the $\gamma_i(\tau^M)$, and $\mathfrak b'$ the ideal of the $x\in S$ with
$\ell^kx\in\mathfrak b$ for some $k$. It contains $\mathfrak b$, and it contains $x$ whenever it contains $\ell x$.

The ideal $\mathfrak b'$ is principal. The ring $S/\ell S=\mathbf F_{\ell^h}[[u]]$ is a principal ideal domain, so the
image of $\mathfrak b'$ is generated by the image of some $g\in\mathfrak b'$. For $x\in\mathfrak b'$ there is $c\in S$ with
$x-cg=\ell y$, and then $y\in\mathfrak b'$. So $\mathfrak b'=gS+\ell\mathfrak b'$, and $\mathfrak b'=gS$ by Nakayama's
lemma, as $S$ is a Noetherian local ring.

Here $g\neq0$. Otherwise every $\gamma_i(\tau^M)$ vanishes, and $\tau^M=0$ by
Lemma~\ref{gal-structures-topological-exists}. But $\tau\neq0$: otherwise $\phi(\tau)=\eta([\zeta]-1)$ in
Definition~\ref{gal-cover-breuil-kisin} gives $[\zeta]=1$, so $\zeta=1$ and $(\zeta^{1/\ell})^\sharp=1$. So $\tau^M\neq0$,
as $A_{\inf}$ is a domain. And $\ell$ does not divide $g$: if $g=\ell y$, then $y\in\mathfrak b'$, so $y=gz$ and
$g=\ell zg$, and $\ell z=1$ in the domain $S$, which is impossible.

Next, let $x\in A_{\inf}$ and let $\mathfrak a\subset S$ be an ideal with $\phi(\gamma_i(x))\in\mathfrak a$ for every
$i$. Then $\gamma_j(\phi(x))\in\mathfrak a$ for every $j$. Indeed, $\phi(\mathfrak j)\subset\mathfrak j$ for
$\mathfrak j=(\ell,E)$, as $\phi(E)\equiv E^\ell$ modulo $\ell$. Given $n$, choose a finite set $F$ as in
Definition~\ref{gal-structures-topological}(2) and apply $\phi$, which commutes with $\iota$:
$\phi(x)-\sum_{i\in F}\iota(\phi(\gamma_i(x)))\phi(b_i)\in\mathfrak j^nA_{\inf}$. Applying $\gamma_j$ and
Lemma~\ref{gal-structures-topological-exists} gives $\gamma_j(\phi(x))\in\mathfrak a+\mathfrak j^n$ for every $n$. By
Krull's intersection theorem for the finite $S$-module $S/\mathfrak a$, $\gamma_j(\phi(x))\in\mathfrak a$.

Apply this to $x=\tau^M$ and $\mathfrak a=\phi(g)S$. As $g$ divides every $\gamma_i(\tau^M)$, $\phi(g)$ divides
$\gamma_j(\phi(\tau^M))=\gamma_j(\iota(E)^M\tau^M)=E^M\gamma_j(\tau^M)$ for every $j$, and so $\phi(g)$ divides $E^Mx$ for
every $x\in\mathfrak b$. As $g\in\mathfrak b'$, $\ell^kg\in\mathfrak b$ for some $k$, and $\phi(g)$ divides $\ell^kE^Mg$.
Since $\phi(g)\equiv g^\ell$ modulo $\ell$ and $\ell$ is a prime element of $S$ that does not divide $g$, $\ell$ does
not divide $\phi(g)$, so $\phi(g)$ divides $E^Mg$. By Lemma~\ref{gal-structures-frobenius-unit}, $g$ is a unit. Hence
$\mathfrak b'=S$, and $\ell^k\in\mathfrak b$ for some $k$.
\end{proof}

\begin{lemma}\label{gal-structures-integral-etale}
Let $\ell$ be a prime and $K\subset\bar{\mathbf Q}_\ell$ a finite extension of $\mathbf Q_\ell$ with a
Breuil--Kisin datum $(S,E,\iota)$, and let $\mathcal O_{\mathcal E}$ be its Fontaine ring
(Definition~\ref{gal-structures-etale}). Then every element of $A_{\inf}\cap\mathcal O_{\mathcal E}$ is $\iota(s)$ for
some $s\in S$.
\end{lemma}
\leanref{Development/AbelianVarietyTheory/Proofs/InertiaEigenvalues/IntegralWittElementOfEvaluationStatement.lean}{GlobalTheory.T.integral_witt_element_of_evaluation}
\leanref{Development/AbelianVarietyTheory/Proofs/InertiaEigenvalues/IntegralWittElementOfEvaluation.lean}{GlobalTheory.BSP.HeightSInt.main}

\begin{proof}
For $s\in S$ let $\bar s\in\mathcal O^\flat$ be the zeroth Witt component of $\iota(s)$, the image of $s$ under
$S\to\mathbf F_{\ell^h}[[u]]\to\mathcal O^\flat$, $u\mapsto\pi^\flat$. If $\ell$ does not divide $s$, let $j$ be least
such that the coefficient of $u^j$ in $s$ is a unit; then $s=u^jt+\ell s'$ with $t$ a unit of $S$, and
$\bar s=(\pi^\flat)^j\bar t$ with $\bar t$ a unit of $\mathcal O^\flat$. In particular $\bar s=0$ only if $\ell$ divides
$s$, as $\pi^\flat\neq0$. Hence, if $x\in\mathcal O^\flat$ and $(\pi^\flat)^kx=\bar s$, then $x=\bar s_0$ for some
$s_0\in S$: if $\ell$ divides $s$, then $x=0$; otherwise $(\pi^\flat)^kx=(\pi^\flat)^j\bar t$, and $k\leq j$ as
$\pi^\flat$ is not a unit (Definition~\ref{gal-structures-etale}), so $x=\overline{u^{j-k}t}$.

If $\ell v\in\mathcal O_{\mathcal E}$ for some $v\in W(\mathbf C_\ell^\flat)$, then $v\in\mathcal O_{\mathcal E}$. Indeed,
given $n$, there are $k$, $s$ and $y$ with $u^k\ell v=\iota(s)+\ell^{n+1}y$; the zeroth components give $\bar s=0$, so
$s=\ell s'$, and $u^kv=\iota(s')+\ell^ny$, as $W(\mathbf C_\ell^\flat)$ has no $\ell$-torsion.

Let $a\in A_{\inf}\cap\mathcal O_{\mathcal E}$. Then $u^ka\equiv\iota(s)$ modulo $\ell$ for some $k$ and $s$, so
$(\pi^\flat)^ka_0=\bar s$ for the zeroth component $a_0$ of $a$, and $a_0=\bar s_0$ for some $s_0\in S$. So
$a-\iota(s_0)$ has zeroth component $0$, and $a=\iota(s_0)+\ell a_1$ with $a_1\in A_{\inf}$. As
$\ell a_1\in\mathcal O_{\mathcal E}$, also $a_1\in\mathcal O_{\mathcal E}$. Repeating this gives $s_0,s_1,\ldots\in S$
with $a\equiv\iota(\sum_{i<n}\ell^is_i)$ modulo $\ell^nA_{\inf}$ for every $n$. The series $s=\sum_i\ell^is_i$
converges in $S$, and $a-\iota(s)\in\ell^nA_{\inf}$ for every $n$, so $a=\iota(s)$.
\end{proof}

\begin{lemma}\label{gal-structures-descent}
Let $\ell$ be a prime and $K\subset\bar{\mathbf Q}_\ell$ a finite extension of $\mathbf Q_\ell$ with a
Breuil--Kisin datum $(S,E,\iota)$ and a cyclotomic element $\tau$ (Definition~\ref{gal-cover-breuil-kisin}), let
$\mathcal O_{\mathcal E}$ be its Fontaine ring (Definition~\ref{gal-structures-etale}), and let $p,q\geq1$.
\begin{enumerate}
\item Let $X$ be a $p\times q$ matrix over $\mathcal O_{\mathcal E}$ such that $\tau^MX$ has entries in $A_{\inf}$ for
  some $M\geq0$. Then $X=\iota(Y)$ for a matrix $Y$ over $S$.
\item Let $R$ be a $p\times q$ matrix over $\mathcal O_{\mathcal E}$, $x\in\mathcal O_{\mathcal E}^p$, $a\in A_{\inf}^q$ and
  $n\geq0$ with $x\equiv Ra$ modulo $\ell^nW(\mathbf C_\ell^\flat)$. Then $x\equiv R\iota(s)$ modulo
  $\ell^nW(\mathbf C_\ell^\flat)$ for some $s\in S^q$.
\end{enumerate}
\end{lemma}
\leanref{Development/AbelianVarietyTheory/Proofs/InertiaEigenvalues/ExistsPowerSeriesLiftOfPeriodPowSmulIntegralMatrStatement.lean}{GlobalTheory.T.exists_powerSeries_lift_of_period_pow_smul_integral_matrix}
\leanref{Development/AbelianVarietyTheory/Proofs/InertiaEigenvalues/ExistsPowerSeriesLiftOfPeriodPowSmulIntegralMatr.lean}{GlobalTheory.BSP.KisinFF.main}
\leanref{Development/AbelianVarietyTheory/Proofs/InertiaEigenvalues/BreuilKisinDescentCoefficientsAinfModPowStatement.lean}{GlobalTheory.T.breuilKisin_descent_coefficients_ainf_mod_pow}
\leanref{Development/AbelianVarietyTheory/Proofs/InertiaEigenvalues/BreuilKisinDescentCoefficientsAinfModPow.lean}{GlobalTheory.CoefDescent.descent}

\begin{proof}
Write $u$ also for $\iota(u)$, a unit of $W(\mathbf C_\ell^\flat)$, and use that
$A_{\inf}\cap\ell^nW(\mathbf C_\ell^\flat)=\ell^nA_{\inf}$. Fix a topological basis $(b_i,\gamma_i)_{i\in I}$ of
$A_{\inf}$ over $S$ (Lemma~\ref{gal-structures-topological-exists}).

(1) It suffices to treat one entry $w\in\mathcal O_{\mathcal E}$, with $a=\tau^Mw\in A_{\inf}$. By
Lemma~\ref{gal-structures-coefficients} there are $r\geq0$ and finitely many nonzero $t_i\in S$ with
$\ell^r=\sum_it_i\gamma_i(\tau^M)$. Put $s_0=\sum_it_i\gamma_i(a)$. Fix $n$. As $w\in\mathcal O_{\mathcal E}$, there
are $k\geq0$, $s\in S$ and $y\in W(\mathbf C_\ell^\flat)$ with $u^kw=\iota(s)+\ell^ny$. Then
$u^ka-\tau^M\iota(s)=\ell^n\tau^My$ lies in $A_{\inf}$, so it is $\ell^nz$ with $z\in A_{\inf}$. Applying $\gamma_i$
gives $u^k\gamma_i(a)-s\gamma_i(\tau^M)=\ell^n\gamma_i(z)$, and multiplying by $t_i$ and summing gives
$u^ks_0-\ell^rs=\ell^nt$ with $t\in S$. Applying $\iota$ and dividing by the unit $u^k$,
\[
\iota(s_0)-\ell^rw=u^{-k}\bigl(\ell^r\iota(s)+\ell^n\iota(t)-\ell^r\iota(s)-\ell^{r+n}y\bigr)
=\ell^nu^{-k}\bigl(\iota(t)-\ell^ry\bigr).
\]
So $\iota(s_0)\equiv\ell^rw$ modulo $\ell^n$ for every $n$, and $\iota(s_0)=\ell^rw$. Hence $\ell^r$ divides
$\iota(s_0)$ in $A_{\inf}$, and $w\in A_{\inf}$. By Lemma~\ref{gal-structures-integral-etale}, $w\in\iota(S)$.

(2) Let $\gamma=\gamma_{i_0}$, so that $\gamma(\iota(s)y)=s\gamma(y)$ and $\gamma(\iota(s))=s$ for $s\in S$ and
$y\in A_{\inf}$. As the entries of $R$ and $x$ lie in $\mathcal O_{\mathcal E}$, there are $M\geq0$ and a matrix
$\varrho$ and a vector $\xi$ over $S$ with $u^MR\equiv\iota(\varrho)$ and $u^Mx\equiv\iota(\xi)$ modulo $\ell^n$; a
common $M$ exists, since $u^kw\equiv\iota(s)$ implies $u^{k+1}w\equiv\iota(us)$. Then
$\iota(\xi)-\iota(\varrho)a\equiv u^M(x-Ra)\equiv0$ modulo $\ell^n$, and the left side is over $A_{\inf}$, so it is
$\ell^nb$ with $b$ over $A_{\inf}$. Applying $\gamma$ entrywise gives $\xi-\varrho\gamma(a)=\ell^n\gamma(b)$. Put
$s=\gamma(a)$. Then $u^M(x-R\iota(s))\equiv\iota(\xi)-\iota(\varrho)\iota(s)=\ell^n\iota(\gamma(b))\equiv0$ modulo
$\ell^n$, and as $u$ is a unit, $x\equiv R\iota(s)$ modulo $\ell^n$.
\end{proof}

\begin{lemma}\label{gal-structures-free}
Let $\ell$ be a prime and $K\subset\bar{\mathbf Q}_\ell$ a finite extension of $\mathbf Q_\ell$ with a
Breuil--Kisin datum $(S,E,\iota)$, and let $\mathcal O_{\mathcal E}$ be its Fontaine ring
(Definition~\ref{gal-structures-etale}). Let $d\geq0$, and let $\mathcal M\subset A_{\inf}^d$ be an $S$-submodule, with
$S$ acting through $\iota$, that contains $y$ whenever it contains $\ell y$ and that contains the limit of every
sequence $(y_n)$ in $\mathcal M$ with $y_{n+1}\equiv y_n$ modulo $\ell^n$. Let $Y_1\in M_d(A_{\inf})$ have nonzero
determinant and satisfy $Y_1\iota(c)\in\mathcal M$ for every $c\in S^d$, and let
$Y'\in\operatorname{GL}_d(W(\mathbf C_\ell^\flat))$ be such that every $y\in\mathcal M$ is $Y'w$ for some
$w\in\mathcal O_{\mathcal E}^d$. Then there is $Y\in M_d(A_{\inf})$, invertible over $W(\mathbf C_\ell^\flat)$, with
$\mathcal M=\{Y\iota(c):c\in S^d\}$.
\end{lemma}
\leanref{Development/AbelianVarietyTheory/Proofs/InertiaEigenvalues/AinfClosedMembershipLimitStatement.lean}{GlobalTheory.T.ainf_closed_membership_limit}
\leanref{Development/AbelianVarietyTheory/Proofs/InertiaEigenvalues/AinfClosedMembershipLimit.lean}{GlobalTheory.BSP.saturated_module_free}

\begin{proof}
For $y$ over $A_{\inf}$ or $W(\mathbf C_\ell^\flat)$ write $y_0$ for the entrywise zeroth Witt components; taking them is
a ring homomorphism, and an element of $A_{\inf}$ with zeroth component $0$ is divisible by $\ell$. Let
$R=S/\ell S=\mathbf F_{\ell^h}[[u]]$, a principal ideal domain, and $e:R\to\mathcal O^\flat$ the homomorphism with
$e(u)=\pi^\flat$, so that $\iota(s)_0=e(s\bmod\ell)$. It is injective, as a nonzero element of $R$ is $u^j$ times a
unit and $\pi^\flat\neq0$.

Let $Q$ be the matrix of zeroth components of $Y'^{-1}$, invertible over $\mathbf C_\ell^\flat$. Every element of
$\mathbf C_\ell^\flat$ becomes integral after multiplication by a power of $\pi^\flat$, as $\mathcal O^\flat$ is a
valuation ring whose only element divisible by every power of $\pi^\flat$ is $0$. Choose $N_1$ with $P=(\pi^\flat)^{N_1}Q$
over $\mathcal O^\flat$.

For every $y\in\mathcal M$ there is $\varrho\in R^d$ with $Py_0=e(\varrho)$. Indeed, $y=Y'w$ with
$w\in\mathcal O_{\mathcal E}^d$, so $Py_0=(\pi^\flat)^{N_1}w_0$. For each entry $w_j$, there are $k$ and $s$ with
$u^kw_j\equiv\iota(s)$ modulo $\ell$, so $(\pi^\flat)^k(Py_0)_j=e(u^{N_1}s\bmod\ell)$, where $(Py_0)_j\in\mathcal O^\flat$.
If $(\pi^\flat)^ko=e(\varrho_1)$ with $o\in\mathcal O^\flat$ and $k\geq1$, then $u$ divides $\varrho_1$, since otherwise
$e(\varrho_1)$ would be a unit divisible by the non-unit $\pi^\flat$; so $(\pi^\flat)^{k-1}o=e(\varrho_1/u)$, and by
induction $o\in e(R)$.

The set $L$ of $\varrho\in R^d$ with $Py_0=e(\varrho)$ for some $y\in\mathcal M$ is an $R$-submodule of $R^d$, as
$(\iota(s)y)_0=e(s\bmod\ell)y_0$. As $R$ is a principal ideal domain, $L$ is free, with a basis
$\varrho_1,\ldots,\varrho_n$ and $n\leq d$. Choose $y_j\in\mathcal M$ with $Py_{j,0}=e(\varrho_j)$, and let $Y$ be the
$d\times n$ matrix with columns $y_1,\ldots,y_n$.

For every $y\in\mathcal M$ there is $c\in S^n$ with $y\equiv Y\iota(c)$ modulo $\ell$: write the element $\varrho$ of $L$
given by $y$ as $\sum_j(c_j\bmod\ell)\varrho_j$ with $c_j\in S$; then $P(y-Y\iota(c))_0=0$, so $(y-Y\iota(c))_0=0$ as
$P$ is invertible, and $\ell$ divides $y-Y\iota(c)$. Hence $\mathcal M=\{Y\iota(c):c\in S^n\}$. Indeed, for
$y\in\mathcal M$ this gives $y=Y\iota(c_0)+\ell y_1$, where $\ell y_1\in\mathcal M$ and so $y_1\in\mathcal M$. Repeating,
$y\equiv Y\iota(\sum_{k<m}\ell^kc_k)$ modulo $\ell^m$ for every $m$; the sum $c=\sum_k\ell^kc_k$ converges in $S^n$, and
$y=Y\iota(c)$, as $A_{\inf}$ is $\ell$-adically separated. Conversely $Y\iota(c)\in\mathcal M$, as the $y_j$ lie in
$\mathcal M$.

The columns of $Y_1$ lie in $\mathcal M$, so $Y_1=Y\iota(C)$ for an $n\times d$ matrix $C$ over $S$; as
$\det Y_1\neq0$, $d\leq n$. So $n=d$. Then $PY_0=e(B)$, where $B$ has columns $\varrho_1,\ldots,\varrho_d$. As these are
linearly independent, $\det B\neq0$, and as $e$ is injective and $P$ invertible, $\det Y_0\neq0$. This is the zeroth
component of $\det Y$, so $\det Y$ is a unit of $W(\mathbf C_\ell^\flat)$, and $Y$ is invertible over
$W(\mathbf C_\ell^\flat)$.
\end{proof}

\begin{definition}\label{gal-structures-points}
Let $\ell$ be a prime, $R$ a commutative ring and $u\in R$. A \emph{perfect point} of $(R,u)$ is a ring
homomorphism $f:R\to O$ to a valuation ring $O$ of characteristic $\ell$ on which $y\mapsto y^\ell$ is
bijective, such that $f(u)$ is nonzero and not a unit and $O$ is complete and separated for the
$f(u)$-adic topology.
\end{definition}
\leanref{Development/AbelianVarietyTheory/Proofs/InertiaEigenvalues/DvdAllPowOfFrobeniusEqKilledByPerfectPointsStatement.lean}{GlobalTheory.T.dvd_all_pow_of_frobenius_eq_killed_by_perfect_points}
\leanref{Development/AbelianVarietyTheory/Proofs/InertiaEigenvalues/DescentOfPointsKilledStatement.lean}{GlobalTheory.T.descent_of_points_killed_alt}

\begin{lemma}\label{gal-structures-points-bounded}
Let $R$ be a commutative ring, $u\in R$ a nonzerodivisor such that $1-ur$ is a unit for every $r\in R$,
$x\in R$ and $M\geq0$. Then either there is $c\geq0$ such that $u^{Mn}$ divides $u^cx^n$ for every
$n\geq0$, or there are an algebraically closed field $L$, a valuation ring $V\subset L$ and a ring
homomorphism $f:R\to V$ such that $f(u)$ is nonzero and not a unit and $f(x)$ divides $f(u)^M$ in $V$.
\end{lemma}
\leanref{Development/AbelianVarietyTheory/Proofs/InertiaEigenvalues/BoundedPowersOrExistsValuationSubringDvdStatement.lean}{GlobalTheory.T.pow_dvd_bound_or_exists_valuationSubring_dvd}
\leanref{Development/AbelianVarietyTheory/Proofs/InertiaEigenvalues/BoundedPowersOrExistsValuationSubringDvd.lean}{GlobalTheory.bounded_powers_or_exists_valuationSubring_dvd}
\leanref{Development/AbelianVarietyTheory/Proofs/InertiaEigenvalues/ExistsValuationSubringAlgClosedNotIsUnitStatement.lean}{GlobalTheory.T.exists_valuationSubring_algClosed_not_isUnit}

\begin{proof}
Let $T=R[1/(ux)]$, in which the images of $u$ and $x$ are units, put $t=u^Mx^{-1}\in T$, so that
$xt=u^M$, and let $R[t]\subset T$ be the $R$-subalgebra generated by $t$. The image of $u$ in $R[t]$ is a
nonzerodivisor, being a unit of $T$.

Suppose first that $u$ is a unit of $R[t]$, say $up(t)=1$ for a polynomial $p$ over $R$. Then
$q=1-up=\sum_{i=0}^nq_iX^i$ satisfies $q(t)=0$, and $q_0=1-up(0)$ is a unit. Put $v=u^M$. Multiplying
$q(t)=0$ by $x^n$ gives $\sum_{i=0}^nq_iv^ix^{n-i}=0$ in $T$, so $(ux)^N\sum_iq_iv^ix^{n-i}=0$ in $R$ for
some $N$, and, as $u$ is a nonzerodivisor, $x^N\sum_{i=0}^nq_iv^ix^{n-i}=0$. By induction on $j$, $v^j$
divides $v^{N+n}x^j$. This is clear for $j\leq N+n$. For $j=N+n+d$ with $d\geq1$, multiplying
the relation by $v^{N+n}x^d$ gives
\[
q_0v^{N+n}x^j=-\sum_{i=1}^nq_iv^i\cdot v^{N+n}x^{j-i}.
\]
By induction $v^{j-i}$ divides $v^{N+n}x^{j-i}$, so $v^j$ divides the right side, and as $q_0$ is a unit,
$v^j$ divides $v^{N+n}x^j$. So $c=M(N+n)$ works.

Otherwise $u$ is not a unit of $R[t]$. Choose a maximal ideal $\mathfrak m$ of $R[t]$ containing $u$ and a
minimal prime $\mathfrak q\subset\mathfrak m$, let $L$ be an algebraic closure of the fraction field of
$R[t]/\mathfrak q$, and let $V\subset L$ be a valuation ring dominating the local ring of
$R[t]/\mathfrak q$ at $\mathfrak m$ \cite[Tag 00IA]{Sta}. Then $R[t]$ maps into $V$, and $u$ into the
maximal ideal of $V$. The elements of the minimal prime $\mathfrak q$ are zero divisors
\cite[Tag 00EU]{Sta}, so the nonzerodivisor $u$ has nonzero image. Let $f$ be the composite
$R\to R[t]\to V$. Then $f(x)f(t)=f(u)^M$, so $f(x)$ divides $f(u)^M$.
\end{proof}

\begin{lemma}\label{gal-structures-points-completion}
\begin{enumerate}
\item Let $V$ be a valuation ring, $x\in V$ nonzero and not a unit, and $\hat V$ the completion of $V$ for
  the $x$-adic topology. Then $\hat V$ is a valuation ring, the image of $x$ in $\hat V$ is nonzero and not
  a unit, and $\hat V$ is complete and separated for its $x$-adic topology.
\item Let $\ell$ be a prime, $W$ a commutative ring with $\ell=0$ in $W$ on which $y\mapsto y^\ell$ is
  bijective, $x\in W$, and $\hat W$ the completion of $W$ for the $x$-adic topology. Then $\ell=0$ in
  $\hat W$ and $y\mapsto y^\ell$ is bijective on $\hat W$.
\end{enumerate}
\end{lemma}
\leanref{Development/AbelianVarietyTheory/Proofs/InertiaEigenvalues/AdicCompletionValuationRingCompleteValuationDomaStatement.lean}{GlobalTheory.T.adicCompletion_valuationRing_complete_valuationDomain}
\leanref{Development/AbelianVarietyTheory/Proofs/InertiaEigenvalues/AdicCompletionValuationRingCompleteValuationDoma.lean}{GlobalTheory.adicCompletion_valuationRing_props}
\leanref{Development/AbelianVarietyTheory/Proofs/InertiaEigenvalues/AdicCompletionValuationRingCompleteValuationDoma.lean}{GlobalTheory.ValCompletion.exists_eq_mul_unit}
\leanref{Development/AbelianVarietyTheory/Proofs/InertiaEigenvalues/AdicCompletionPrincipalPerfectStatement.lean}{GlobalTheory.T.adicCompletion_principal_perfect}
\leanref{Development/AbelianVarietyTheory/Proofs/InertiaEigenvalues/AdicCompletionPrincipalPerfect.lean}{GlobalTheory.PerfectCompletion.perfect}

\begin{proof}
In both parts write $X$ for the image of $x$ in the completion. The completion is complete and separated
for the topology of the extended ideal, as $(x)$ is finitely generated \cite[Tag 05GG]{Sta}, and an element
$y$ is divisible by $X^n$ if and only if its image in $V/x^nV$ (or $W/x^nW$) vanishes; in particular, for $s$
in the ring, $X^n$ divides $s$ in the completion if and only if $x^n$ divides $s$. Every element of the
completion is $s+X^nz$ with $s$ in the ring, for each $n$. An element whose image modulo $X$ is a unit is a
unit.

(1) As $x$ is not a unit, $x^{m+1}$ does not divide $x^m$ in $V$, so $X^m\neq0$ for every $m$; and $X$ is
not a unit, since $Xe=1$ would give $x\mid1$ in $V$. Let $y\in\hat V$ be nonzero. Choose $n$ with
$X^n\nmid y$ and $a\in V$ with $y=a+X^{n+1}z$. Then $x^n\nmid a$, so $a\mid x^n$ in the valuation ring $V$,
say $x^n=ab$, and $y=a(1+bXz)$ with $1+bXz$ a unit. So every nonzero element of $\hat V$ is $au$ with $a\in
V$ dividing a power of $x$ and $u$ a unit. If $y=au$ and $y'=a'u'$ are such, with $a\mid x^n$ and
$a'\mid x^{n'}$, then $yy'\neq0$, since $aa'$ divides $X^{n+n'}\neq0$; and one of $a,a'$ divides the other
in $V$, so one of $y,y'$ divides the other. Hence $\hat V$ is a domain and a valuation ring.

(2) The image of $\ell$ is $0$. For injectivity, let $w^\ell=0$ and $n\geq0$, and write $w=s+X^{\ell n}t$
with $s\in W$. Then $s^\ell=-(X^{\ell n}t)^\ell$ is divisible by $X^{\ell n}$, so $x^{\ell n}$ divides
$s^\ell$ in $W$, and as $y\mapsto y^\ell$ is bijective on $W$, $x^n$ divides $s$. So $X^n$ divides $w$ for
every $n$, and $w=0$. For surjectivity, let $y\in\hat W$, write $y=s_n+X^{\ell n}t_n$ with $s_n\in W$, and
let $b_n\in W$ with $b_n^\ell=s_n$. For $m\leq n$, $(b_m-b_n)^\ell=s_m-s_n$ is divisible by $x^{\ell m}$,
so $x^m$ divides $b_m-b_n$. The $b_n$ therefore converge to some $L\in\hat W$ with $X^n\mid L-b_n$, and
$L^\ell-y=(L-b_n)^\ell-X^{\ell n}t_n$ is divisible by $X^n$ for every $n$, so $L^\ell=y$.
\end{proof}

\begin{lemma}\label{gal-structures-points-killed}
Let $\ell$ be a prime, $R$ a commutative ring in which $\ell=0$, $u\in R$ a nonzerodivisor, $h\geq0$, and
$D_1,\ldots,D_k$ and $c_{ij}$ $(1\leq i,j\leq k)$ elements of $R$ with $u^hD_j=\sum_iD_i^\ell c_{ij}$ for
every $j$. Suppose that $f(D_j)=0$ for every perfect point $f$ of $(R,u)$
(Definition~\ref{gal-structures-points}) and every $j$. Then $u^m$ divides every $D_j$, for every $m\geq0$.
\end{lemma}
\leanref{Development/AbelianVarietyTheory/Proofs/InertiaEigenvalues/DvdAllPowOfFrobeniusEqKilledByPerfectPointsStatement.lean}{GlobalTheory.T.dvd_all_pow_of_frobenius_eq_killed_by_perfect_points}
\leanref{Development/AbelianVarietyTheory/Proofs/InertiaEigenvalues/DvdAllPowOfFrobeniusEqKilledByPerfectPoints.lean}{GlobalTheory.PointsKey.key}
\leanref{Development/AbelianVarietyTheory/Proofs/InertiaEigenvalues/PowDvdOfVanishingAtPerfectValuationPoints.lean}{GlobalTheory.s_points_kpt}
\leanref{Development/AbelianVarietyTheory/Proofs/InertiaEigenvalues/FrobeniusPowEquationDvdAllPow.lean}{GlobalTheory.phiEq_dvd_all_pow}

\begin{proof}
\emph{Localization.} The set $\Sigma=\{1-ur:r\in R\}$ is multiplicative, as
$(1-ur)(1-ur')=1-u(r+r'-urr')$. Let $R'=\Sigma^{-1}R$. In $R'$ we have $\ell=0$ and the relations
$u^hD_j=\sum_iD_i^\ell c_{ij}$. The element $u$ is a nonzerodivisor of $R'$: if $u\cdot a/s=0$, then
$tua=0$ for some $t\in\Sigma$, so $ta=0$ and $a/s=0$. Every $1-ur'$ with $r'=a/s\in R'$ and
$s=1-ub$ is a unit of $R'$, since $(1-ur')s=1-u(a+b)\in\Sigma$. A perfect point of $(R',u)$ composed with
$R\to R'$ is a perfect point of $(R,u)$, so it kills every $D_j$.

\emph{Bounded powers.} For every $j$ and $M\geq0$ there is $c\geq0$ such that $u^{Mn}$ divides $u^cD_j^n$
in $R'$ for every $n$. Otherwise Lemma~\ref{gal-structures-points-bounded}, applied in $R'$ to $x=D_j$, gives
an algebraically closed field $L$, a valuation ring $V\subset L$ and $f:R'\to V$ with $f(u)$ nonzero and not
a unit and $f(D_j)$ dividing $f(u)^M$. The ring $V$ has characteristic $\ell$, and $y\mapsto y^\ell$ is
bijective on it: injective as $V$ is a domain, and surjective since an $\ell$-th root in $L$ of an element of
$V$ lies in $V$. Let $\hat V$ be the $f(u)$-adic completion of $V$. By
Lemma~\ref{gal-structures-points-completion}, $R'\to V\to\hat V$ is a perfect point of $(R',u)$. It kills
$D_j$, so it kills $f(u)^M$; but the image of $f(u)$ in the domain $\hat V$ is nonzero.

\emph{Frobenius iteration.} For every $k\geq0$ there are $a\geq0$ with $a+h\leq h\ell^k$ and $e_{ij}\in R'$
with $u^aD_j=\sum_iD_i^{\ell^k}e_{ij}$ for every $j$. For $k=0$ take $a=0$ and $e_{ij}=\delta_{ij}$. Given
$a$ and $e_{ij}$ for $k$, the additivity of $y\mapsto y^\ell$ in characteristic $\ell$ gives
\[
u^{h+\ell a}D_j=\sum_i(u^aD_i)^\ell c_{ij}=\sum_i\Bigl(\sum_{i'}D_{i'}^{\ell^k}e_{i'i}\Bigr)^\ell c_{ij}
=\sum_{i'}D_{i'}^{\ell^{k+1}}\sum_ie_{i'i}^\ell c_{ij},
\]
and $(h+\ell a)+h\leq\ell h+\ell a\leq h\ell^{k+1}$, as $\ell\geq2$.

\emph{Divisibility in $R'$.} Fix $m$. By the bounded powers with $M=h+1$, choose $C_i$ such that
$u^{(h+1)n}$ divides $u^{C_i}D_i^n$ for every $n$, put $C=\sum_iC_i$, and choose $k$ with $m+C<\ell^k$.
With $a$ and $e_{ij}$ for this $k$, $u^{C+a}D_j=\sum_iu^{C-C_i}\cdot u^{C_i}D_i^{\ell^k}e_{ij}$ is divisible by
$u^{(h+1)\ell^k}$. As $C+a+m\leq(h+1)\ell^k$ and $u$ is a nonzerodivisor of $R'$, $u^m$ divides $D_j$ in $R'$.

\emph{Divisibility in $R$.} If $D_j=u^m\cdot a/s$ in $R'$, then $t(sD_j-u^ma)=0$ in $R$ for some
$t\in\Sigma$, and $ts=1-ur$ for some $r$. So $(1-ur)D_j=u^mb$ with $b=ta$, and
\[
D_j=\Bigl(\sum_{i<m}(ur)^i\Bigr)(1-ur)D_j+(ur)^mD_j=u^m\Bigl(b\sum_{i<m}(ur)^i+r^mD_j\Bigr).
\qedhere
\]
\end{proof}

\begin{lemma}\label{gal-structures-points-descent}
Let $\ell$ be a prime and $A\to B$ a homomorphism of commutative rings that makes $B$ a flat $A$-module, with
$\ell=0$ in $B$. Let $u\in A$ have as image in $B$ a nonzerodivisor, suppose $B$ is separated for the
$u$-adic topology, and let $\operatorname{pr}:B\to A$ be an $A$-linear map with $\operatorname{pr}(1)=1$. Let
$x_1,\ldots,x_k\in B$ satisfy the following.
\begin{enumerate}
\item For every commutative ring $T$ and ring homomorphisms $i_1,i_2:B\to T$ that agree on the image of $A$,
  with $i_1(u)$ a unit, there are $c_{ij}\in T$ with $\delta_j=\sum_i\delta_i^\ell c_{ij}$ for every $j$,
  where $\delta_j=i_1(x_j)-i_2(x_j)$.
\item For every perfect point $\psi:B\to O$ of $(B,u)$ (Definition~\ref{gal-structures-points}) and every
  ring homomorphism $\psi':B\to O$ that agrees with $\psi$ on the image of $A$, $\psi(x_j)=\psi'(x_j)$ for
  every $j$.
\end{enumerate}
Then every $x_j$ is the image of $\operatorname{pr}(x_j)$ in $B$.
\end{lemma}
\leanref{Development/AbelianVarietyTheory/Proofs/InertiaEigenvalues/DescentOfPointsKilledStatement.lean}{GlobalTheory.T.descent_of_points_killed_alt}
\leanref{Development/AbelianVarietyTheory/Proofs/InertiaEigenvalues/DescentOfPointsKilled.lean}{GlobalTheory.descent_of_points_killed}

\begin{proof}
Let $R=B\otimes_AB$, $U=u\otimes1=1\otimes u$ and $D_j=x_j\otimes1-1\otimes x_j$. Then $\ell=0$ in $R$.
Multiplication by $U$ on $R$ is the tensor product of multiplication by $u$ on $B$ with the identity of $B$,
so it is injective, as $B$ is flat over $A$: $U$ is a nonzerodivisor.

Apply~(1) to $T=R[1/U]$ and the maps $i_1(b)=b\otimes1$ and $i_2(b)=1\otimes b$ followed by $R\to T$. They
agree on the image of $A$, and $i_1(u)=U$ is a unit of $T$. The resulting $c_{ij}\in T$ become elements
$c'_{ij}$ of $R$ after multiplication by a power $U^h$, and then $U^hD_j=\sum_iD_i^\ell c'_{ij}$ holds in
$T$, hence in $R$, since $R\to T$ is injective.

Let $f:R\to O$ be a perfect point of $(R,U)$. The maps $\psi(b)=f(b\otimes1)$ and $\psi'(b)=f(1\otimes b)$
agree on the image of $A$, and $\psi(u)=f(U)$, so $\psi$ is a perfect point of $(B,u)$, and~(2) gives
$f(D_j)=\psi(x_j)-\psi'(x_j)=0$. By Lemma~\ref{gal-structures-points-killed}, $U^m$ divides every $D_j$, for
every $m$.

Let $\Phi:R\to B$ be the $A$-linear map with $\Phi(b\otimes b')=\operatorname{pr}(b)b'$, which exists as
$\operatorname{pr}$ is $A$-linear. Then $\Phi(D_j)=\operatorname{pr}(x_j)-\operatorname{pr}(1)x_j=
\operatorname{pr}(x_j)-x_j$ and $\Phi(U^mr)=u^m\Phi(r)$. So $\operatorname{pr}(x_j)-x_j$ lies in $u^mB$ for
every $m$, and it vanishes, as $B$ is $u$-adically separated.
\end{proof}

\begin{lemma}\label{gal-structures-points-untilt}
Let $\ell$ be a prime, $O$ a valuation ring of characteristic $\ell$ on which $y\mapsto y^\ell$ is bijective,
and $x\in O$ nonzero and not a unit, with $O$ complete and separated for the $x$-adic topology. Let $e\geq1$,
let $w$ be a unit of $W(O)$, and put $\xi=[x]^e+\ell w$. Then $\xi W(O)$ is a prime ideal of $W(O)$ that
does not contain $\ell$.
\end{lemma}
\leanref{Development/AbelianVarietyTheory/Proofs/InertiaEigenvalues/TeichmullerDistinguishedSpanIsPrimeAndNotMemStatement.lean}{GlobalTheory.T.teichmullerDistinguished_span_isPrime_and_not_mem}
\leanref{Development/AbelianVarietyTheory/Proofs/InertiaEigenvalues/TeichmullerDistinguishedSpanIsPrimeAndNotMem.lean}{GlobalTheory.teichmuller_distinguished_span_isPrime}
\leanref{Development/AbelianVarietyTheory/Proofs/InertiaEigenvalues/TeichmullerDistinguishedRegularSequenceStatement.lean}{GlobalTheory.T.teichmuller_distinguished_regular_sequence}
\leanref{Development/AbelianVarietyTheory/Proofs/InertiaEigenvalues/WittVectorTeichmullerIsAdicCompleteOfPerfectStatement.lean}{GlobalTheory.T.wittVector_teichmuller_isAdicComplete_of_perfect}

\begin{proof}
Write $W=W(O)$ and $a_0$ for the zeroth Witt component of $a\in W$. The fraction field of the perfect domain
$O$ is perfect, and $W$ embeds into its Witt vectors, a discrete valuation ring with uniformizer $\ell$
\cite[Chapter~II, \S6]{Ser79}; so $W$ is a domain. As $O$ is perfect, $\ell(z_0,z_1,\ldots)=(0,z_0^\ell,
z_1^\ell,\ldots)$, so $a$ is divisible by $\ell$ if and only if $a_0=0$, and $a=[a_0]+\ell a_1$ for some
$a_1\in W$; an element with unit zeroth component is a unit, as $W$ is $\ell$-adically complete. Also
$[z](a_0,a_1,\ldots)=(za_0,z^\ell a_1,\ldots)$. So an element divisible by every power of $[x]$ has every
component divisible by every power of $x$, and vanishes: $W$ is $[x]$-adically separated. It is
$[x]$-adically complete by the lifting argument in the proof of Lemma~\ref{gal-structures-topological-exists},
with $O$ in place of $\mathcal O^\flat$.

\emph{Regularity.} For $k\geq0$ and $a\in W$: if $[x]^k$ divides $\xi a$, then $[x]^k$ divides $a$; and if
$\xi$ divides $[x]^ka$, then $\xi$ divides $a$. For the first, let $k=1$ and $[x]b=\xi a=[x]^ea+\ell wa$.
Then $[x]b'=\ell wa$ for some $b'$, so $xb'_0=0$, $b'_0=0$ and $b'=\ell c$; cancelling $\ell$ gives
$[x]c=wa$, and $[x]$ divides $a$ as $w$ is a unit. For $k+1$, write $a=[x]c$ and cancel $[x]$, which is
nonzero, to get $[x]^k\mid\xi c$; by induction $[x]^k\mid c$. For the second, if $[x]^ka=\xi b$, then
$[x]^k$ divides $\xi b$, so $b=[x]^kc$, and $a=\xi c$.

\emph{Powers of $[x]$.} No power $[x]^N$ lies in $\xi W$: otherwise $\xi$ divides $[x]^N\cdot1$, hence
divides $1$ by regularity; then the zeroth component $x^e$ of $\xi$ is a unit, and so is $x$.

\emph{Normal form.} Let $v=w^{-1}$, so that $\ell\equiv-[x]^ev$ modulo $\xi$. For every $a\in W$, either
$a\equiv[x]^ea'$ modulo $\xi$ for some $a'\in W$, or $a\equiv[y]v'$ modulo $\xi$ for some $y\in O$ dividing
$x^e$ and some unit $v'$. Indeed $a=[a_0]+\ell a_1\equiv[a_0]-[x]^eva_1$. As $O$ is a valuation ring, either
$a_0=x^ec$, and $a\equiv[x]^e([c]-va_1)$; or $a_0c=x^e$ with $c$ a unit, and again $a_0=x^ec^{-1}$; or
$a_0c=x^e$ with $c$ not a unit, and $a\equiv[a_0](1-[c]va_1)$, where $1-[c]va_1$ has unit zeroth component
$1-c(va_1)_0$ and is a unit. Iterating, for every $n$ either $a\equiv[x]^{en}a'$ modulo $\xi$, or
$a\equiv[y]v'$ modulo $\xi$ with $y\mid x^{en}$ and $v'$ a unit: if $a\equiv[x]^{en}a'$ and the step gives
$a'\equiv[y]v'$, then $a\equiv[x^{en}y]v'$ with $x^{en}y\mid x^{e(n+1)}$.

\emph{Closedness.} If for every $k$ there is $c_k$ with $a\equiv[x]^kc_k$ modulo $\xi$, then $a\in\xi W$.
Write $a=\xi b_k+[x]^kc_k$. For $m\leq n$, $\xi(b_m-b_n)=[x]^nc_n-[x]^mc_m$ is divisible by $[x]^m$, so
$[x]^m$ divides $b_m-b_n$ by regularity. So $(b_k)$ converges $[x]$-adically to some $b$, and
$a-\xi b=\xi(b_n-b)+[x]^nc_n$ is divisible by $[x]^n$ for every $n$, so $a=\xi b$.

\emph{Primality.} The ideal $\xi W$ is proper, as $[x]^0=1\notin\xi W$. Let $ab\in\xi W$ with
$a,b\notin\xi W$. By closedness there is $k$ such that $a$ is not congruent to any multiple of $[x]^k$ modulo
$\xi$. In the normal form with $n=k$ the first alternative would give $a\equiv[x]^k[x]^{ek-k}a'$, so
$a\equiv[y]v'$ with $y\mid x^{ek}$ and $v'$ a unit. Likewise $b\equiv[y']v''$ with $y'\mid x^{N'}$ for some
$N'$. Then $[yy']v'v''\in\xi W$, so $[yy']\in\xi W$; writing $x^{ek+N'}=yy's$, the power
$[x]^{ek+N'}=[yy'][s]$ lies in $\xi W$, which is impossible. Finally, if $\ell\in\xi W$, then
$[x]^e=\xi-\ell w\in\xi W$, which is impossible.
\end{proof}

\begin{lemma}\label{gal-structures-residue}
Let $\ell$ be a prime and $K\subset\bar{\mathbf Q}_\ell$ a finite extension of $\mathbf Q_\ell$ with a
Breuil--Kisin datum $(S,E,\iota)$ and element $\pi^\flat$, with $S=W(\mathbf F_{\ell^h})[[u]]$
(Definition~\ref{gal-cover-breuil-kisin}), and let $\mathcal O_{\mathcal E}$ be its Fontaine ring
(Definition~\ref{gal-structures-etale}). Let $\kappa:\mathbf F_{\ell^h}[[u]]\to\mathcal O^\flat$ send the
reduction modulo $\ell$ of each $s\in S$ to the zeroth Witt component of $\iota(s)$, and let
$\overline{\mathcal O}_{\mathcal E}\subset\mathbf C_\ell^\flat$ be the set of zeroth Witt components of the
elements of $\mathcal O_{\mathcal E}$.
\begin{enumerate}
\item The map $\kappa$ is a well-defined injective ring homomorphism with $\kappa(u)=\pi^\flat$, and there is
  an additive map $\operatorname{pr}:\mathcal O^\flat\to\mathbf F_{\ell^h}[[u]]$ with
  $\operatorname{pr}(\kappa(a)y)=a\operatorname{pr}(y)$ for $a\in\mathbf F_{\ell^h}[[u]]$ and
  $y\in\mathcal O^\flat$, and $\operatorname{pr}(1)=1$.
\item An element $y\in\mathbf C_\ell^\flat$ lies in $\overline{\mathcal O}_{\mathcal E}$ if and only if
  $(\pi^\flat)^my\in\kappa(\mathbf F_{\ell^h}[[u]])$ for some $m\geq0$.
\item The ring $\mathcal O_{\mathcal E}$ contains $\iota(s)$ for every $s\in S$; $\iota(u)$ is a unit of
  $W(\mathbf C_\ell^\flat)$; and if $w\in W(\mathbf C_\ell^\flat)$ and $\iota(u)w\in\mathcal O_{\mathcal E}$, then
  $w\in\mathcal O_{\mathcal E}$.
\item The inverse in $\mathbf C_\ell^\flat$ of every nonzero element of $\overline{\mathcal O}_{\mathcal E}$ lies in
  $\overline{\mathcal O}_{\mathcal E}$.
\end{enumerate}
\end{lemma}
\leanref{Development/AbelianVarietyTheory/Proofs/InertiaEigenvalues/ExistsResidueTiltInjectiveRetractStatement.lean}{GlobalTheory.T.exists_residueTilt_injective_retract}
\leanref{Development/AbelianVarietyTheory/Proofs/InertiaEigenvalues/ExistsResidueTiltInjectiveRetract.lean}{GlobalTheory.BSP.exists_resTilt_retract}
\leanref{Development/AbelianVarietyTheory/Proofs/InertiaEigenvalues/BreuilKisinOEModPrimeLaurentFieldStatement.lean}{GlobalTheory.T.breuilKisin_OE_mod_prime_laurent_field}
\leanref{Development/AbelianVarietyTheory/Proofs/InertiaEigenvalues/BreuilKisinOEModPrimeLaurentField.lean}{GlobalTheory.BSP.oe_mod_l}

\begin{proof}
Write $u$ also for $\iota(u)=[\pi^\flat]$ and $\bar s$ for the reduction of $s\in S$ modulo $\ell$.

(1) If $s\in\ell S$, then $\iota(s)\in\ell A_{\inf}$ has zeroth Witt component $0$; so $\kappa$ is well
defined, and it is a ring homomorphism, as taking zeroth Witt components is one. Its value at $u$ is
$\pi^\flat$. Let $(b_i,\gamma_i)_{i\in I}$ be a topological basis of $A_{\inf}$ over $S$
(Lemma~\ref{gal-structures-topological-exists}), let $\gamma=\gamma_{i_0}$, and let $\operatorname{pr}(y)$ be
the reduction modulo $\ell$ of $\gamma([y])$. If $x\in A_{\inf}$ has zeroth Witt component $x_0$, then
$x-[x_0]\in\ell A_{\inf}=\iota(\ell)A_{\inf}$, so $\gamma(x)\equiv\gamma([x_0])$ modulo $\ell$, as $\gamma$ is
$S$-linear. Applied to $[y]+[y']$ and to $\iota(s)[y]$, whose zeroth components are $y+y'$ and $\kappa(\bar
s)y$, this gives $\operatorname{pr}(y+y')=\operatorname{pr}(y)+\operatorname{pr}(y')$ and
$\operatorname{pr}(\kappa(\bar s)y)=\bar s\operatorname{pr}(y)$. Also $\operatorname{pr}(1)$ is the reduction
of $\gamma(b_{i_0})=1$. Hence $\operatorname{pr}(\kappa(a))=a$, and $\kappa$ is injective.

(2) If $y$ is the zeroth component of $w\in\mathcal O_{\mathcal E}$, the definition of $\mathcal O_{\mathcal E}$
with $n=1$ gives $m$, $s$ and $z$ with $u^mw=\iota(s)+\ell z$, and taking zeroth components gives
$(\pi^\flat)^my=\kappa(\bar s)$. Conversely, suppose $(\pi^\flat)^my=\kappa(\bar s)$. As $u$ is a unit of
$W(\mathbf C_\ell^\flat)$ (Definition~\ref{gal-structures-etale}), $w=u^{-m}\iota(s)$ satisfies
$u^mw=\iota(s)$, so it lies in $\mathcal O_{\mathcal E}$. Its zeroth component is
$(\pi^\flat)^{-m}\kappa(\bar s)=y$.

(3) The congruences $u^0\iota(s)\equiv\iota(s)$ show that $\iota(s)\in\mathcal O_{\mathcal E}$. The zeroth Witt
component of $u=[\pi^\flat]$ is $\pi^\flat\neq0$, so $u$ is a unit of $W(\mathbf C_\ell^\flat)$, as
$\mathbf C_\ell^\flat$ is a field. If $uw\in\mathcal O_{\mathcal E}$ and $n\geq0$, there are $k$ and $s$ with
$u^k(uw)\equiv\iota(s)$ modulo $\ell^n$, that is $u^{k+1}w\equiv\iota(s)$; so $w\in\mathcal O_{\mathcal E}$.

(4) Let $y\in\overline{\mathcal O}_{\mathcal E}$ be nonzero. By (2), $(\pi^\flat)^my=\kappa(\bar s)$ for some
$m\geq0$ and $s\in S$, and $\kappa(\bar s)\neq0$, so $s\notin\ell S$. Write $s=\sum_ia_iu^i$ with
$a_i\in W(\mathbf F_{\ell^h})$. A nonunit of $W(\mathbf F_{\ell^h})$ is divisible by $\ell$, so some $a_i$ is a
unit; let $k$ be the least such $i$. Then $s=u^ks_1+\ell s_2$, where $s_1=\sum_ia_{i+k}u^i$ is a unit of $S$,
as its constant term $a_k$ is, and $s_2\in S$. So $\kappa(\bar s)=(\pi^\flat)^k\kappa(\bar s_1)$, and
\[
 (\pi^\flat)^ky^{-1}=(\pi^\flat)^m\kappa(\bar s_1)^{-1}=\kappa\bigl(\overline{u^ms_1^{-1}}\bigr).
\]
By (2), $y^{-1}\in\overline{\mathcal O}_{\mathcal E}$.
\end{proof}

\begin{lemma}\label{gal-structures-points-galois}
Let $\ell$ be a prime and $K\subset\bar{\mathbf Q}_\ell$ a finite extension of $\mathbf Q_\ell$ with a
Breuil--Kisin datum $(S,E,\iota)$ and element $\pi^\flat$ (Definition~\ref{gal-cover-breuil-kisin}), and let
$\kappa:\mathbf F_{\ell^h}[[u]]\to\mathcal O^\flat$ be as in Lemma~\ref{gal-structures-residue}. Let
$\psi:\mathcal O^\flat\to O$ be a perfect point of $(\mathcal O^\flat,\pi^\flat)$
(Definition~\ref{gal-structures-points}) and $\psi':\mathcal O^\flat\to O$ a ring homomorphism with
$\psi\circ\kappa=\psi'\circ\kappa$. Then there is $g\in G_K$ with $g(\pi^\flat)=\pi^\flat$ and
$\psi'=\psi\circ g$.
\end{lemma}
\leanref{Development/AbelianVarietyTheory/Proofs/InertiaEigenvalues/ExistsFixFlatOfTiltMapsAgreeStatement.lean}{GlobalTheory.T.exists_fixFlat_of_tilt_maps_agree}
\leanref{Development/AbelianVarietyTheory/Proofs/InertiaEigenvalues/ExistsFixFlatOfTiltMapsAgree.lean}{GlobalTheory.BSP.exists_galK_fixFlat_of_tilt_maps_agree}
\leanref{Development/AbelianVarietyTheory/Proofs/InertiaEigenvalues/WittVectorMapBreuilKisinEvEqStatement.lean}{GlobalTheory.T.wittVector_map_breuilKisin_ev_eq}
\leanref{Development/AbelianVarietyTheory/Proofs/InertiaEigenvalues/BreuilKisinEisensteinImageTeichmullerPowStatement.lean}{GlobalTheory.T.breuilKisin_eisenstein_image_teichmuller_pow}
\leanref{Development/AbelianVarietyTheory/Proofs/InertiaEigenvalues/ExistsInjectiveMapIntegersClToWittQuotientStatement.lean}{GlobalTheory.T.exists_injective_map_integers_Cl_to_wittQuotient}
\leanref{Development/AbelianVarietyTheory/Proofs/InertiaEigenvalues/ExistsGaloisEqEmbeddingsOCModPrimeStatement.lean}{GlobalTheory.T.exists_galois_eq_embeddings_OC_mod_prime}
\leanref{Development/AbelianVarietyTheory/Proofs/InertiaEigenvalues/UntiltOfTiltRootFlatAlgebraicAndGaloisFixStatement.lean}{GlobalTheory.T.untiltOfTiltRoot_flat_algebraic_and_galois_fix}
\leanref{Development/AbelianVarietyTheory/Proofs/InertiaEigenvalues/RingHomEqOfPerfectOfDvdSubAdicSeparatedStatement.lean}{GlobalTheory.T.ringHom_eq_of_perfect_of_dvd_sub_adicSeparated}

\begin{proof}
Let $x=\psi(\pi^\flat)$, which equals $\psi'(\pi^\flat)$ as $\pi^\flat=\kappa(u)$, let $e$ be the degree of $E$,
and let $W(\psi),W(\psi'):A_{\inf}\to W(O)$ be the maps induced on Witt vectors. For $y\in\mathcal O^\flat$
write $y^\sharp=\theta([y])\in\mathcal O_{\mathbf C_\ell}$.

\emph{The two maps agree on $S$.} $W(\psi)\circ\iota=W(\psi')\circ\iota$. On $W(\mathbf F_{\ell^h})$, $\iota$
is induced by an embedding $\mathbf F_{\ell^h}\to\mathcal O^\flat$, whose composites with $\psi$ and $\psi'$
agree, as $\psi\circ\kappa=\psi'\circ\kappa$; so the two maps agree there, and both send $u$ to $[x]$. For
$s\in S$ and $n\geq0$ write $s=p_n+u^ns_n$ with $p_n$ a polynomial over $W(\mathbf F_{\ell^h})$; then the
difference of the images of $s$ is $[x]^n$ times that of $s_n$. As $[x]^n(z_0,z_1,\ldots)=(x^nz_0,
x^{n\ell}z_1,\ldots)$ and $O$ is $x$-adically separated, an element of $W(O)$ divisible by every $[x]^n$
vanishes.

\emph{The untilt.} Put $\xi=W(\psi)(\iota(E))=W(\psi')(\iota(E))$. Write $E=u^e+\ell(c+uG)$ with $c$ a unit
of $W(\mathbf F_{\ell^h})$ and $G$ a polynomial, as $E$ is monic with lower coefficients divisible by $\ell$
and constant coefficient $\ell$ times a unit. Then $\xi=[x]^e+\ell w$ with $w=W(\psi)(\iota(c+uG))$, whose
zeroth component is a unit of $O$ plus a multiple of $x$, a unit; so $w$ is a unit. By
Lemma~\ref{gal-structures-points-untilt}, $D=W(O)/\xi W(O)$ is a domain in which $\ell\neq0$. Fontaine's map
$\theta:A_{\inf}\to\mathcal O_{\mathbf C_\ell}$ is onto with kernel $\iota(E)A_{\inf}$ \cite{Fon94a}, so there
are ring homomorphisms $\theta_1,\theta_2:\mathcal O_{\mathbf C_\ell}\to D$ with $\theta_1(\theta(a))$ and
$\theta_2(\theta(a))$ the classes of $W(\psi)(a)$ and $W(\psi')(a)$. They are injective: if $\theta_1(y)=0$
with $y\neq0$, then $y$ divides some $\ell^n$ in $\mathcal O_{\mathbf C_\ell}$, so $\ell^n=0$ in the domain
$D$, which is impossible. If $\ell$ divides the class of $b\in W(O)$ in $D$, then $x^e$ divides $b_0$: from
$\ell b'-b=c'\xi$ the zeroth components give $-b_0=c'_0x^e$.

\emph{Agreement on $\mathcal O_K$.} Every $r\in\mathcal O_K$ is $\theta(\iota(s))$ for some $s\in S$, as
$\theta\circ\iota$ induces the inclusion through $S/(E)\cong\mathcal O_K$; so
$\theta_1(r)=\theta_2(r)$ by the first step.

\emph{A Galois element.} There is $g\in G_K$ with $\theta_2(y)=\theta_1(g(y))$ for every
$y\in\mathcal O_{\mathbf C_\ell}$ algebraic over $\mathbf Q_\ell$, and with $\ell$ dividing
$\theta_2(y)-\theta_1(g(y))$ in $D$ for every $y\in\mathcal O_{\mathbf C_\ell}$. Indeed, $\theta_1$ and
$\theta_2$ extend to field embeddings $\tau_1,\tau_2$ of $\mathbf C_\ell$ into the fraction field $F$ of $D$,
which agree on $K$, the fraction field of $\mathcal O_K$. Their restrictions to $\bar{\mathbf Q}_\ell$ are
$K$-embeddings of a normal extension of $K$, so they have the same image, and $\tau_2=\tau_1\circ g$ on
$\bar{\mathbf Q}_\ell$ for the $K$-automorphism $g=(\tau_1|_{\bar{\mathbf Q}_\ell})^{-1}\circ
\tau_2|_{\bar{\mathbf Q}_\ell}$ of $\bar{\mathbf Q}_\ell$, an element of $G_K$. As $\bar{\mathbf Q}_\ell$ is
dense in $\mathbf C_\ell$, every $y\in\mathcal O_{\mathbf C_\ell}$ is $z+\ell w'$ with
$z\in\mathcal O_{\mathbf C_\ell}$ algebraic and $w'\in\mathcal O_{\mathbf C_\ell}$, and
$\theta_2(y)-\theta_1(g(y))=\ell(\theta_2(w')-\theta_1(g(w')))$.

\emph{The element $g$ fixes $\pi^\flat$.} Let $\pi_n=((\pi^\flat)^{1/\ell^n})^\sharp$. Then
$\pi_n^{\ell^n}=(\pi^\flat)^\sharp=\theta(\iota(u))$ lies in $\mathcal O_K$, so $\pi_n$ is algebraic. Also
$\theta_1(\pi_n)$ and $\theta_2(\pi_n)$ are the classes of $[\psi((\pi^\flat)^{1/\ell^n})]$ and
$[\psi'((\pi^\flat)^{1/\ell^n})]$, and these agree, as both elements of $O$ have $\ell^n$-th power $x$ and
$y\mapsto y^\ell$ is injective on $O$. So $\theta_1(g(\pi_n))=\theta_2(\pi_n)=\theta_1(\pi_n)$, and
$g(\pi_n)=\pi_n$ for every $n$. An element $z$ of $\mathcal O^\flat$ is determined by the elements
$(z^{1/\ell^n})^\sharp$, and $g$ sends those of $\pi^\flat$ to those of $g(\pi^\flat)$; so
$g(\pi^\flat)=\pi^\flat$.

\emph{Conclusion.} For $y\in\mathcal O^\flat$, $g(y^\sharp)=g(y)^\sharp$, so $\ell$ divides the class of
$[\psi'(y)]-[\psi(g(y))]$ in $D$, and $x^e$ divides its zeroth component $\psi'(y)-\psi(g(y))$. Given $y$ and
$n$, write $y=z^{\ell^n}$, as $\mathcal O^\flat$ is perfect. Then
$\psi'(y)-\psi(g(y))=(\psi'(z)-\psi(g(z)))^{\ell^n}$ is divisible by $x^{e\ell^n}$, hence by $x^n$. As $O$
is $x$-adically separated, $\psi'(y)=\psi(g(y))$.
\end{proof}

\begin{lemma}\label{gal-structures-residue-descent}
Let $\ell$ be a prime and $K\subset\bar{\mathbf Q}_\ell$ a finite extension of $\mathbf Q_\ell$ with a
Breuil--Kisin datum $(S,E,\iota)$ and element $\pi^\flat$ (Definition~\ref{gal-cover-breuil-kisin}), let
$\overline{\mathcal O}_{\mathcal E}\subset\mathbf C_\ell^\flat$ be as in Lemma~\ref{gal-structures-residue},
and let $G_\infty\subset G_K$ be the subgroup fixing $\pi^\flat$. Let $x_1,\ldots,x_k\in\mathbf C_\ell^\flat$
be fixed by $G_\infty$. Suppose that for every commutative ring $T$ and ring homomorphisms
$i_1,i_2:\mathbf C_\ell^\flat\to T$ that agree on $\overline{\mathcal O}_{\mathcal E}$ there are
$c_{ij}\in T$ with $\delta_j=\sum_i\delta_i^\ell c_{ij}$ for every $j$, where $\delta_j=i_1(x_j)-i_2(x_j)$.
Then every $x_j$ lies in $\overline{\mathcal O}_{\mathcal E}$.
\end{lemma}
\leanref{Development/AbelianVarietyTheory/Proofs/InertiaEigenvalues/TiltInvariantMemReductionOfFrobeniusEquationStatement.lean}{GlobalTheory.T.tilt_invariant_mem_reduction_of_frobenius_equation}
\leanref{Development/AbelianVarietyTheory/Proofs/InertiaEigenvalues/TiltInvariantMemReductionOfFrobeniusEquation.lean}{GlobalTheory.BSP.TwoCopy.mem_reduction_of_invariant}

\begin{proof}
Let $\kappa$ and $\operatorname{pr}$ be as in Lemma~\ref{gal-structures-residue}, and view $\mathcal O^\flat$
as an algebra over the discrete valuation ring $A=\mathbf F_{\ell^h}[[u]]$ through the injective map $\kappa$.
It has no $A$-torsion, so it is flat \cite[Tag 0539]{Sta}; $\ell=0$ in $\mathcal O^\flat$; and
$\kappa(u)=\pi^\flat$ is a nonzerodivisor. As $\pi^\flat$ is nonzero and not a unit
(Definition~\ref{gal-structures-etale}), $\mathcal O^\flat$ is complete and separated for the
$\pi^\flat$-adic topology, every nonzero $b\in\mathcal O^\flat$ divides a power of $\pi^\flat$, and every
element of $\mathbf C_\ell^\flat$ times a power of $\pi^\flat$ lies in $\mathcal O^\flat$. Choose $m$ with
$x'_j=(\pi^\flat)^mx_j\in\mathcal O^\flat$ for every $j$. The hypotheses of
Lemma~\ref{gal-structures-points-descent} hold for $A\to\mathcal O^\flat$, $u$, $\operatorname{pr}$ and the
$x'_j$.

For~(1), let $i_1,i_2:\mathcal O^\flat\to T$ agree on $\kappa(A)$, with $\upsilon=i_1(\pi^\flat)$ a unit; then
$i_2(\pi^\flat)=\upsilon$. Every nonzero $b\in\mathcal O^\flat$ divides a power of $\pi^\flat$, so $i_1(b)$ and
$i_2(b)$ are units, and $i_1,i_2$ extend to $j_1,j_2:\mathbf C_\ell^\flat\to T$. These agree on
$\overline{\mathcal O}_{\mathcal E}$: if $(\pi^\flat)^ny=\kappa(a)$, then $j_1(y)=\upsilon^{-n}i_1(\kappa(a))=
j_2(y)$ (Lemma~\ref{gal-structures-residue}(2)). The hypothesis gives $c_{ij}$ for $j_1,j_2$, and
\[
i_1(x'_j)-i_2(x'_j)=\upsilon^m\bigl(j_1(x_j)-j_2(x_j)\bigr)
=\sum_i\bigl(i_1(x'_i)-i_2(x'_i)\bigr)^\ell\upsilon^{m-m\ell}c_{ij}.
\]

For~(2), let $\psi$ be a perfect point of $(\mathcal O^\flat,\pi^\flat)$ and $\psi'$ agree with it on
$\kappa(A)$. Lemma~\ref{gal-structures-points-galois} gives $g\in G_\infty$ with $\psi'=\psi\circ g$. As $g$
fixes $\pi^\flat$ and the $x_j$, it fixes the $x'_j$, and $\psi'(x'_j)=\psi(x'_j)$.

Lemma~\ref{gal-structures-points-descent} gives $x'_j=\kappa(\operatorname{pr}(x'_j))$, that is,
$(\pi^\flat)^mx_j\in\kappa(A)$, and $x_j\in\overline{\mathcal O}_{\mathcal E}$ by
Lemma~\ref{gal-structures-residue}(2).
\end{proof}

\begin{lemma}\label{gal-structures-echelon}
Let $\ell$ be a prime and $C$ a field of characteristic $\ell$, and write $\phi(Z)$ for the matrix of $\ell$-th
powers of the entries of a matrix $Z$ over $C$. Let $Y,A\in M_d(C)$ with $\det A\neq0$ and $\phi(Y)=YA$. Then
there are a set $I\subset\{1,\ldots,d\}$ and $Q,\Lambda\in M_d(C)$ with the following properties.
\begin{enumerate}
\item $YQ=0$, and $Qv=v$ for every $v\in C^d$ with $Yv=0$.
\item For $i\in I$ the $i$-th row of $Q$ is that of the identity matrix, and for $j\notin I$ the $j$-th column
  of $Q$ is zero.
\item $Q=A\phi(Q)\Lambda$.
\item If $\sigma$ is a ring endomorphism of $C$ and $R\in M_d(C)$ satisfies $R\sigma(Y)=Y$, then
  $\sigma(Q)=Q$.
\item If $T$ is a commutative ring and $i_1,i_2:C\to T$ are ring homomorphisms with $i_1(A)=i_2(A)$, then
  there are $c_{(a',b'),(a,b)}\in T$ with
  $i_1(Q_{ab})-i_2(Q_{ab})=\sum_{a',b'}\bigl(i_1(Q_{a'b'})-i_2(Q_{a'b'})\bigr)^\ell c_{(a',b'),(a,b)}$ for all
  $a,b$.
\end{enumerate}
\end{lemma}
\leanref{Development/AbelianVarietyTheory/Proofs/InertiaEigenvalues/EchelonProjectorKernelFrobeniusStatement.lean}{GlobalTheory.T.echelon_projector_kernel_frobenius}
\leanref{Development/AbelianVarietyTheory/Proofs/InertiaEigenvalues/EchelonProjectorKernelFrobenius.lean}{GlobalTheory.echelon_projector_frobenius}
\leanref{Development/AbelianVarietyTheory/Proofs/InertiaEigenvalues/ExistsKernelVectorMemOEOfFrobeniusStable.lean}{GlobalTheory.BSP.KernelDescent.invariant}
\leanref{Development/AbelianVarietyTheory/Proofs/InertiaEigenvalues/ExistsKernelVectorMemOEOfFrobeniusStable.lean}{GlobalTheory.BSP.KernelDescent.two_copies}

\begin{proof}
If a matrix $P$ has zero columns outside $I$ and a matrix $P'$ has its rows indexed by $I$ equal to those of
the identity, then $PP'=P$: the $(a,b)$ entry of $PP'$ is $\sum_{c\in I}P_{ac}P'_{cb}$, which is $P_{ab}$ if
$b\in I$ and $0=P_{ab}$ if $b\notin I$.

(1) and (2). Let $N=\ker Y$ and $e_1,\ldots,e_d$ the standard basis. Choose a set $J_0$ of indices such that
the images of the $e_j$, $j\in J_0$, form a basis of $C^d/N$, let $I$ be its complement and $W$ the span of
the $e_j$, $j\in J_0$, so that $C^d=N\oplus W$. Let $Q$ be the projection onto $N$ along $W$. Then $YQ=0$ and
$Qv=v$ for $v\in N$. For $j\in J_0$, $Qe_j=0$. For $i\in I$ and every $j$, $e_j-Qe_j\in W$ has $i$-th
coordinate $0$, so $Q_{ij}=\delta_{ij}$.

(3) Let $M=A\phi(Q)$. Then $YM=\phi(Y)\phi(Q)=\phi(YQ)=0$, so the image of $M$ lies in $N$. If $Mv=0$, then
$\phi(Q)v=0$, as $A$ is invertible; the rows of $\phi(Q)$ indexed by $I$ are those of the identity, so $v_i=0$
for $i\in I$, and $v\in W$. Hence the image of $M$ has dimension at least $d-\dim W=\dim N$, and it is $N$. So
every column $Qe_j\in N$ is $M\lambda_j$, and $\Lambda$ with columns $\lambda_j$ satisfies $Q=M\Lambda$.

(4) Here $Y\sigma(Q)=R\sigma(Y)\sigma(Q)=R\sigma(YQ)=0$, so the columns of $\sigma(Q)$ lie in $N$ and
$Q\sigma(Q)=\sigma(Q)$ by~(1). As $\sigma(Q)$ has the shape~(2), the first paragraph gives
$Q\sigma(Q)=Q$. So $\sigma(Q)=Q$.

(5) If $T=0$ there is nothing to show; otherwise $T$ has characteristic $\ell$, as $i_1$ is injective. Write
$Q_k=i_k(Q)$ and $\Phi_k=i_k(\phi(Q))$ for $k=1,2$, and $A'=i_1(A)=i_2(A)$. As $YM=0$, (1) gives $QM=M$. So
$Q_2Q_1=Q_2$ by the first paragraph, $Q_1=A'\Phi_1i_1(\Lambda)$ by~(3), and $Q_2A'\Phi_2=A'\Phi_2$. Hence
\[
(1-Q_2)A'(\Phi_1-\Phi_2)i_1(\Lambda)=(1-Q_2)A'\Phi_1i_1(\Lambda)=(1-Q_2)Q_1=Q_1-Q_2.
\]
The entries of $\Phi_1-\Phi_2$ are $i_1(Q_{ab})^\ell-i_2(Q_{ab})^\ell=(i_1(Q_{ab})-i_2(Q_{ab}))^\ell$.
Expanding the product gives the elements $c_{(a',b'),(a,b)}$.
\end{proof}

\begin{lemma}\label{gal-structures-kernel-vector}
Let $\ell$ be a prime and $K\subset\bar{\mathbf Q}_\ell$ a finite extension of $\mathbf Q_\ell$ with a
Breuil--Kisin datum $(S,E,\iota)$ and element $\pi^\flat$ (Definition~\ref{gal-cover-breuil-kisin}), let
$\overline{\mathcal O}_{\mathcal E}\subset\mathbf C_\ell^\flat$ be as in Lemma~\ref{gal-structures-residue},
let $G_\infty\subset G_K$ be the subgroup fixing $\pi^\flat$, and let $\phi$ act on matrices over
$\mathbf C_\ell^\flat$ by $\ell$-th powers of the entries. Let $Y,A\in M_d(\mathbf C_\ell^\flat)$ with $A$ over
$\overline{\mathcal O}_{\mathcal E}$, $\det A\neq0$ and $\phi(Y)=YA$, and suppose that for every $g\in G_\infty$
there is an invertible $R\in M_d(\mathbf C_\ell^\flat)$ with $Rg(Y)=Y$. If $\det Y=0$, there is a nonzero
$v\in\overline{\mathcal O}_{\mathcal E}^d$ with $Yv=0$.
\end{lemma}
\leanref{Development/AbelianVarietyTheory/Proofs/InertiaEigenvalues/ExistsKernelVectorMemOEOfFrobeniusStableStatement.lean}{GlobalTheory.T.exists_kernel_vector_mem_OE_of_frobenius_stable}
\leanref{Development/AbelianVarietyTheory/Proofs/InertiaEigenvalues/ExistsKernelVectorMemOEOfFrobeniusStable.lean}{GlobalTheory.BSP.exists_ker_vector_mem_OE}

\begin{proof}
Lemma~\ref{gal-structures-echelon} gives $I$, $Q$ and $\Lambda$. By its part~(4), applied to the automorphisms
$g\in G_\infty$ of $\mathbf C_\ell^\flat$, every entry of $Q$ is fixed by $G_\infty$. If two ring
homomorphisms $i_1,i_2:\mathbf C_\ell^\flat\to T$ agree on $\overline{\mathcal O}_{\mathcal E}$, they agree on
the entries of $A$, so part~(5) gives the relations required by Lemma~\ref{gal-structures-residue-descent} for
the entries of $Q$. Hence $Q$ is over $\overline{\mathcal O}_{\mathcal E}$. As $\det Y=0$, there is $v\neq0$
with $Yv=0$, and $Qv=v$; so some column of $Q$ is nonzero, and it lies in $\ker Y$ as $YQ=0$.
\end{proof}

\begin{lemma}\label{gal-structures-unit}
Let $\ell$ be a prime and $K\subset\bar{\mathbf Q}_\ell$ a finite extension of $\mathbf Q_\ell$ with a
Breuil--Kisin datum $(S,E,\iota)$ and element $\pi^\flat$ (Definition~\ref{gal-cover-breuil-kisin}), let
$\mathcal O_{\mathcal E}$ be its Fontaine ring (Definition~\ref{gal-structures-etale}), and let
$G_\infty\subset G_K$ be the subgroup fixing $\pi^\flat$. Let $Y\in M_d(W(\mathbf C_\ell^\flat))$ with
$\det Y\neq0$ and $A\in M_d(\mathcal O_{\mathcal E})$ with $\phi(Y)=YA$. Suppose that for every
$g\in G_\infty$ there is $R\in\operatorname{GL}_d(W(\mathbf C_\ell^\flat))$ with $Rg(Y)=Y$, and that every
$w\in\mathcal O_{\mathcal E}^d$ with $Yw\equiv0$ modulo $\ell$ has all its entries divisible by $\ell$. Then
$Y\in\operatorname{GL}_d(W(\mathbf C_\ell^\flat))$.
\end{lemma}
\leanref{Development/AbelianVarietyTheory/Proofs/InertiaEigenvalues/IsUnitOfSaturated.lean}{GlobalTheory.SatInvUnit.isUnit_of_saturated}

\begin{proof}
Let $c:W(\mathbf C_\ell^\flat)\to\mathbf C_\ell^\flat$ take the zeroth Witt component. It is a ring
homomorphism with kernel $\ell W(\mathbf C_\ell^\flat)$, it commutes with $G_K$, and $c(\phi(w))=c(w)^\ell$.
The ring $W(\mathbf C_\ell^\flat)$ is a discrete valuation ring with uniformizer $\ell$
\cite[Chapter~II, \S6]{Ser79}, so $\det Y=\ell^m\upsilon$ with $\upsilon$ a unit. From
$\phi(\det Y)=\det Y\det A$ we get $\phi(\upsilon)=\upsilon\det A$, so $\det A$ is a unit. Put $\bar Y=c(Y)$
and $\bar A=c(A)$. Then $\det\bar A\neq0$, $\bar A$ is over $\overline{\mathcal O}_{\mathcal E}$
(Lemma~\ref{gal-structures-residue}), $\phi(\bar Y)=\bar Y\bar A$, and $c(R)g(\bar Y)=\bar Y$ with $c(R)$
invertible for $g\in G_\infty$. Suppose $\det\bar Y=0$. Lemma~\ref{gal-structures-kernel-vector} gives a
nonzero $\bar v$ over $\overline{\mathcal O}_{\mathcal E}$ with $\bar Y\bar v=0$. Write $\bar v=c(v)$ with $v$
over $\mathcal O_{\mathcal E}$. Then $c(Yv)=0$, so $Yv\equiv0$ modulo $\ell$, so $\ell$ divides every entry
of $v$, and $\bar v=0$, a contradiction. Hence the zeroth component $\det\bar Y$ of $\det Y$ is nonzero, and
$\det Y$ is a unit, as $\mathbf C_\ell^\flat$ is a field.
\end{proof}

\begin{lemma}\label{gal-structures-saturated-lattice}
Let $\ell$ be a prime and $K\subset\bar{\mathbf Q}_\ell$ a finite extension of $\mathbf Q_\ell$ with a
Breuil--Kisin datum, and let $\mathcal O_{\mathcal E}$ be its Fontaine ring
(Definition~\ref{gal-structures-etale}). Let $Y_1\in M_d(W(\mathbf C_\ell^\flat))$ with $\det Y_1\neq0$. Then
there are $m\geq0$, $P\in M_d(\mathcal O_{\mathcal E})$ with $\det P\neq0$ and
$Y'\in M_d(W(\mathbf C_\ell^\flat))$ with $\ell^mY'=Y_1P$, such that:
\begin{enumerate}
\item for every $v\in\mathcal O_{\mathcal E}^d$ and $n\geq0$ with $Y_1v\equiv0$ modulo $\ell^n$ there is
  $w\in\mathcal O_{\mathcal E}^d$ with $\ell^mv=\ell^nPw$;
\item every $w\in\mathcal O_{\mathcal E}^d$ with $Y'w\equiv0$ modulo $\ell$ has all its entries divisible by
  $\ell$.
\end{enumerate}
\end{lemma}
\leanref{Development/AbelianVarietyTheory/Proofs/InertiaEigenvalues/ExistsIntegralWittLatticeBasis.lean}{GlobalTheory.SatLattice.exists_lattice_basis}

\begin{proof}
By Fontaine's theory \cite{Fon90}, $\mathcal O_{\mathcal E}$ is a discrete valuation ring with uniformizer
$\ell$, and an element of $\mathcal O_{\mathcal E}$ divisible by $\ell$ in $W(\mathbf C_\ell^\flat)$ is
divisible by $\ell$ in $\mathcal O_{\mathcal E}$. Write $\det Y_1=\ell^c\upsilon$ with $\upsilon$ a unit of the
discrete valuation ring $W(\mathbf C_\ell^\flat)$ \cite[Chapter~II, \S6]{Ser79}, and put $m=c$. The
$x\in\mathcal O_{\mathcal E}^d$ with $Y_1x\equiv0$ modulo $\ell^c$ form a submodule containing
$\ell^c\mathcal O_{\mathcal E}^d$, so it is free of rank $d$; let $P$ be the matrix of a basis. Its elements
are the $Pw$ with $w\in\mathcal O_{\mathcal E}^d$, and $Y_1P\equiv0$ modulo $\ell^c$, so $Y_1P=\ell^cY'$.
Writing $\ell^c$ times the standard basis vectors in terms of $P$ gives $PX=\ell^c$, so $\det P\neq0$.

(1) Write $Y_1v=\ell^nz$ and put $x=\upsilon^{-1}\operatorname{adj}(Y_1)z$. Then
$\ell^nx=\upsilon^{-1}\operatorname{adj}(Y_1)Y_1v=\ell^cv$, so $x$ is over $\mathcal O_{\mathcal E}$, and
$Y_1x=\upsilon^{-1}\det(Y_1)z=\ell^cz$. So $x=Pw$ with $w$ over $\mathcal O_{\mathcal E}$, and
$\ell^cv=\ell^nPw$.

(2) If $Y'w=\ell y$, then $Y_1Pw=\ell^{c+1}y$, and~(1) applied to $Pw$ gives $\ell^cPw=\ell^{c+1}Pw'$ with $w'$
over $\mathcal O_{\mathcal E}$. So $P(w-\ell w')=0$, and $w=\ell w'$ as $\det P\neq0$.
\end{proof}

\begin{lemma}\label{gal-structures-isogenous}
Let $\ell$ be a prime and $K\subset\bar{\mathbf Q}_\ell$ a finite extension of $\mathbf Q_\ell$ with a
Breuil--Kisin datum $(S,E,\iota)$ and element $\pi^\flat$ (Definition~\ref{gal-cover-breuil-kisin}), and let
$\mathcal O_{\mathcal E}$ be its Fontaine ring (Definition~\ref{gal-structures-etale}). Let
$\rho:G_K\to M_d(\mathbf Z_\ell)$ be a map with a Breuil--Kisin basis $(T,A,A')$ of height $N$
(Definition~\ref{gal-structures-basis}), let $f\in M_d(\mathbf Z_\ell)$ with $\det f\neq0$, and let
$\rho'':G_K\to\operatorname{GL}_d(\mathbf Z_\ell)$ be a map with $f\rho''(g)=\rho(g)f$ for every $g\in G_K$.
Then there are $T''\in\operatorname{GL}_d(W(\mathbf C_\ell^\flat))$ and $P\in M_d(\mathcal O_{\mathcal E})$
with $fT''=TP$.
\end{lemma}
\leanref{Development/AbelianVarietyTheory/Proofs/InertiaEigenvalues/ExistsEtaleOEBasisOfIsogenousLatticeStatement.lean}{GlobalTheory.T.exists_etale_OE_basis_of_isogenous_lattice}
\leanref{Development/AbelianVarietyTheory/Proofs/InertiaEigenvalues/ExistsEtaleOEBasisOfIsogenousLattice.lean}{GlobalTheory.EtBasis.main}

\begin{proof}
Let $G_\infty\subset G_K$ be the subgroup fixing $\pi^\flat$. By Fontaine's theory \cite{Fon90},
$\mathcal O_{\mathcal E}$ is stable under $\phi$ and fixed by $G_\infty$, $\iota(s)$ is a unit of
$\mathcal O_{\mathcal E}$ for $s\in S$ not divisible by $\ell$, and an element of $\mathcal O_{\mathcal E}$
divisible by $\ell$ in $W(\mathbf C_\ell^\flat)$ is divisible by $\ell$ in $\mathcal O_{\mathcal E}$. Also
$\mathcal O_{\mathcal E}$ contains $\mathbf Z_\ell$, as every element of $\mathbf Z_\ell$ is congruent modulo
$\ell^n$ to an integer.

Let $Y_1=\operatorname{adj}(f)T$, with $\det Y_1=\det(f)^{d-1}\det T\neq0$.
Lemma~\ref{gal-structures-saturated-lattice} gives $m$, $P_0$ over $\mathcal O_{\mathcal E}$ with
$\det P_0\neq0$, and $Y'$ with $\ell^mY'=Y_1P_0$. Then $\det Y'\neq0$.

\emph{Frobenius.} As $E$ is monic, $\ell$ does not divide it, so $\iota(E)$ has an inverse $\epsilon$ in
$\mathcal O_{\mathcal E}$. Put $B=\epsilon^N\iota(A)$, over $\mathcal O_{\mathcal E}$; then $\phi(T)=TB$, and
$\phi(Y_1)=Y_1B$, as $\phi$ fixes the entries of $\operatorname{adj}(f)$. The matrix $V=B\phi(P_0)$ is over
$\mathcal O_{\mathcal E}$, and $\ell^m\phi(Y')=\phi(Y_1)\phi(P_0)=Y_1V$. So each column $v$ of $V$ satisfies
$Y_1v\equiv0$ modulo $\ell^m$, and Lemma~\ref{gal-structures-saturated-lattice}(1) with $n=m$ gives
$v=P_0w$ with $w$ over $\mathcal O_{\mathcal E}$. So $V=P_0A''$ with $A''$ over $\mathcal O_{\mathcal E}$,
$\ell^m\phi(Y')=Y_1P_0A''=\ell^mY'A''$, and $\phi(Y')=Y'A''$.

\emph{Galois.} Let $g\in G_\infty$. Multiplying $f\rho''(g)=\rho(g)f$ on both sides by
$\operatorname{adj}(f)$ and cancelling $\det f$ gives $\rho''(g)\operatorname{adj}(f)=\operatorname{adj}(f)\rho(g)$.
As $g$ fixes the entries of $\operatorname{adj}(f)$ and of $P_0$, and $\rho(g)g(T)=T$,
\[
\ell^m\rho''(g)g(Y')=\rho''(g)\operatorname{adj}(f)g(T)P_0=\operatorname{adj}(f)\rho(g)g(T)P_0=Y_1P_0=\ell^mY',
\]
so $\rho''(g)g(Y')=Y'$.

By Lemma~\ref{gal-structures-unit}, with the condition modulo $\ell$ from
Lemma~\ref{gal-structures-saturated-lattice}(2), $Y'$ is invertible. Put $T''=Y'$ and $P=T^{-1}fY'$. Then
$fT''=TP$, and $\ell^mP=T^{-1}f\operatorname{adj}(f)TP_0=\det(f)P_0$ is over $\mathcal O_{\mathcal E}$, so $P$
is over $\mathcal O_{\mathcal E}$.
\end{proof}

\begin{lemma}\label{gal-structures-kernel-lattice}
Let $\ell$ be a prime, $\Gamma$ a group, $n\geq0$, $\rho:\Gamma\to\operatorname{GL}_d(\mathbf Z_\ell)$ and
$\rho':\Gamma\to\operatorname{GL}_{d'}(\mathbf Z_\ell)$ homomorphisms, and $\Pi\in M_{d\times d'}(\mathbf Z_\ell)$
with $\Pi\rho'(g)\equiv\rho(g)\Pi$ modulo $\ell^n$ for every $g\in\Gamma$ and $\Pi X\equiv1$ modulo $\ell^n$ for
some $X\in M_{d'\times d}(\mathbf Z_\ell)$. Then there are $f\in M_{d'}(\mathbf Z_\ell)$ with $\det f\neq0$ and a
homomorphism $\rho'':\Gamma\to\operatorname{GL}_{d'}(\mathbf Z_\ell)$ with $f\rho''(g)=\rho'(g)f$ for every $g$,
such that the $v\in\mathbf Z_\ell^{d'}$ with $\Pi v\equiv0$ modulo $\ell^n$ are the vectors $fw$,
$w\in\mathbf Z_\ell^{d'}$.
\end{lemma}
\leanref{Development/AbelianVarietyTheory/Proofs/InertiaEigenvalues/ExistsKernelLatticeEquivariantModPowStatement.lean}{GlobalTheory.T.exists_kernel_lattice_equivariant_modPow}
\leanref{Development/AbelianVarietyTheory/Proofs/InertiaEigenvalues/ExistsKernelLatticeEquivariantModPow.lean}{GlobalTheory.KerLat.exists_kernel_lattice}

\begin{proof}
The $v$ with $\Pi v\equiv0$ modulo $\ell^n$ form a submodule of $\mathbf Z_\ell^{d'}$ containing
$\ell^n\mathbf Z_\ell^{d'}$, so it is free of rank $d'$. Let $f$ be the matrix of a basis. Then the elements of
the submodule are the $fw$, and $fw=0$ only for $w=0$, so $\det f\neq0$. The submodule is stable under every
$\rho'(g)$, as $\Pi\rho'(g)v\equiv\rho(g)\Pi v$. So $\rho'(g)f=f\rho''(g)$ for a matrix $\rho''(g)$ over
$\mathbf Z_\ell$, unique as $fX=fX'$ implies $X=X'$. Uniqueness gives $\rho''(1)=1$ and, from
$f\rho''(gh)=\rho'(g)\rho'(h)f=f\rho''(g)\rho''(h)$, $\rho''(gh)=\rho''(g)\rho''(h)$; so $\rho''$ is a
homomorphism to $\operatorname{GL}_{d'}(\mathbf Z_\ell)$.
\end{proof}

\begin{lemma}\label{gal-structures-smith}
Let $\ell$ be a prime and $K\subset\bar{\mathbf Q}_\ell$ a finite extension of $\mathbf Q_\ell$ with a
Breuil--Kisin datum, and let $\mathcal O_{\mathcal E}$ be its Fontaine ring
(Definition~\ref{gal-structures-etale}). Let $n\geq1$, $Q\in M_{d\times d'}(W(\mathbf C_\ell^\flat))$ and
$P\in M_{d'}(\mathcal O_{\mathcal E})$. Suppose that $QS\equiv1$ modulo $\ell^n$ for some
$S\in M_{d'\times d}(W(\mathbf C_\ell^\flat))$, and that the $v\in W(\mathbf C_\ell^\flat)^{d'}$ with
$Qv\equiv0$ modulo $\ell^n$ are the vectors $Pw$, $w\in W(\mathbf C_\ell^\flat)^{d'}$. Then there are
$T_n\in\operatorname{GL}_d(W(\mathbf C_\ell^\flat))$, $R\in M_{d\times d'}(\mathcal O_{\mathcal E})$ and
$R'\in M_{d'\times d}(\mathcal O_{\mathcal E})$ with $Q\equiv T_nR$ and $QR'\equiv T_n$ modulo $\ell^n$.
\end{lemma}
\leanref{Development/AbelianVarietyTheory/Proofs/InertiaEigenvalues/ExistsSmithFactorizationEtaleOEModPrimePowStatement.lean}{GlobalTheory.T.exists_smith_factorization_etale_OE_mod_prime_pow}
\leanref{Development/AbelianVarietyTheory/Proofs/InertiaEigenvalues/ExistsSmithFactorizationEtaleOEModPrimePow.lean}{GlobalTheory.SmithQuot.main}
\leanref{Development/AbelianVarietyTheory/Proofs/InertiaEigenvalues/SmithFormCoefficientRingImageOrPrimeDvdRowStatement.lean}{GlobalTheory.T.smith_form_coefficientRing_image_or_prime_dvd_row}
\leanref{Development/AbelianVarietyTheory/Proofs/InertiaEigenvalues/WittCompletedTiltColumnsOfRightInverseModPowStatement.lean}{GlobalTheory.T.wittCompletedTilt_columns_of_right_inverse_mod_pow}

\begin{proof}
\emph{Elementary divisors.} By Fontaine's theory \cite{Fon90}, $\mathcal O_{\mathcal E}$ is a discrete
valuation ring with uniformizer $\ell$. By the elementary divisor theorem, there are a basis
$b_1,\ldots,b_{d'}$ of $\mathcal O_{\mathcal E}^{d'}$ and $a_1,\ldots,a_r\in\mathcal O_{\mathcal E}$ such that
$a_1b_1,\ldots,a_rb_r$ is a basis of the span of the columns of $P$. Let $U$ be the matrix with columns $b_j$.
For each $j$, either $b_j$ is in that span (if $j\leq r$ and $a_j$ is a unit), or the $j$-th row of
$U^{-1}P$, whose entries are the $j$-th coordinates of the columns of $P$ in the basis $b$, is divisible by
$\ell$.

\emph{The columns that matter.} Let $J$ be the set of $j$ such that $Ue_j$ is not of the form $Pw$ with $w$ over
$W(\mathbf C_\ell^\flat)$. For $j\notin J$, $QUe_j\equiv0$ modulo $\ell^n$. For $j\in J$, $b_j$ is not in the
span of the columns of $P$, so the $j$-th row of $U^{-1}P$ is divisible by $\ell$, and for every $v$ with
$Qv\equiv0$ modulo $\ell^n$, writing $v=Pw$, the $j$-th entry of $U^{-1}v$ is divisible by $\ell$.

\emph{Counting.} Let $M=QU$ and $S'=U^{-1}S$, so $MS'=QS\equiv1$. Let $M_J$ be the columns of $M$ indexed by
$J$ and $S'_J$ the rows of $S'$ indexed by $J$. As the other columns of $M$ vanish modulo $\ell^n$,
$M_JS'_J\equiv1$ modulo $\ell^n$. For $j_0\in J$ let $c\in W(\mathbf C_\ell^\flat)^{d'}$ have entries
$(S'_JM_J-1)_{jj_0}$ at $j\in J$ and $0$ elsewhere. Then $QUc=M_J(S'_JM_J-1)e_{j_0}\equiv0$ modulo $\ell^n$,
so the entries of $U^{-1}Uc=c$ at $j\in J$ are divisible by $\ell$: $S'_JM_J\equiv1$ modulo $\ell$. Over the
residue field $\mathbf C_\ell^\flat$ the reductions of $M_J$ and $S'_J$ are inverse to each other, so $J$ has
$d$ elements.

\emph{Conclusion.} Order $J$ and put $T_n=M_J$. Then $\det(T_n)\det(S'_J)\equiv1$ modulo $\ell^n$ with $n\geq1$,
so $T_n$ is invertible, as $W(\mathbf C_\ell^\flat)$ is local with maximal ideal $\ell W(\mathbf C_\ell^\flat)$.
Let $R$ be the rows of $U^{-1}$ indexed by $J$ and $R'$ the columns of $U$ indexed by $J$, both over
$\mathcal O_{\mathcal E}$. Then $QR'=M_J=T_n$, and $Q=MU^{-1}\equiv M_JR=T_nR$ modulo $\ell^n$, since the
columns of $M$ outside $J$ vanish modulo $\ell^n$.
\end{proof}

\begin{lemma}\label{gal-structures-etale-basis}
Let $\ell$ be a prime and $K\subset\bar{\mathbf Q}_\ell$ a finite extension of $\mathbf Q_\ell$ with a
Breuil--Kisin datum $(S,E,\iota)$ and a cyclotomic element $\tau$ (Definition~\ref{gal-cover-breuil-kisin}), and let
$\mathcal O_{\mathcal E}$ be its Fontaine ring (Definition~\ref{gal-structures-etale}). Let $n\geq1$, let
$\rho:G_K\to\operatorname{GL}_d(\mathbf Z_\ell)$ and $\rho':G_K\to\operatorname{GL}_{d'}(\mathbf Z_\ell)$ be
homomorphisms, let $(T',A,A')$ be a Breuil--Kisin basis of height $N$ of $\rho'$
(Definition~\ref{gal-structures-basis}), and let $\Pi\in M_{d\times d'}(\mathbf Z_\ell)$ satisfy
$\Pi\rho'(g)\equiv\rho(g)\Pi$ modulo $\ell^n$ for every $g\in G_K$ and $\Pi X\equiv1$ modulo $\ell^n$ for some
$X\in M_{d'\times d}(\mathbf Z_\ell)$. Then there are $T_n\in\operatorname{GL}_d(W(\mathbf C_\ell^\flat))$,
$R\in M_{d\times d'}(\mathcal O_{\mathcal E})$ and $R'\in M_{d'\times d}(\mathcal O_{\mathcal E})$ with
$\Pi T'\equiv T_nR$ and $\Pi T'R'\equiv T_n$ modulo $\ell^n$. The cyclotomic element is not used.
\end{lemma}
\leanref{Development/AbelianVarietyTheory/Proofs/InertiaEigenvalues/ExistsEtaleBasisModPowStatement.lean}{GlobalTheory.T.exists_etaleBasis_modPow}
\leanref{Development/AbelianVarietyTheory/Proofs/InertiaEigenvalues/ExistsEtaleBasisModPow.lean}{GlobalTheory.BKLim.etBasis_modPow}

\begin{proof}
Lemma~\ref{gal-structures-kernel-lattice} gives $f$ and $\rho''$, and Lemma~\ref{gal-structures-isogenous},
applied to $\rho'$, $(T',A,A')$, $f$ and $\rho''$, gives $T''$ invertible and $P$ over $\mathcal O_{\mathcal E}$
with $fT''=T'P$. As $W(\mathbf C_\ell^\flat)$ has no $\ell$-torsion, it is flat over $\mathbf Z_\ell$
\cite[Tag 0539]{Sta}; tensoring with it the description of the kernel of $v\mapsto\Pi v$ modulo $\ell^n$ on
$\mathbf Z_\ell^{d'}$ shows that the $u\in W(\mathbf C_\ell^\flat)^{d'}$ with $\Pi u\equiv0$ modulo $\ell^n$ are
the $fw$, $w\in W(\mathbf C_\ell^\flat)^{d'}$. Put $Q=\Pi T'$. Then $Qv\equiv0$ modulo $\ell^n$ if and only if
$T'v=fw$ for some $w$, that is, $v=T'^{-1}fT''T''^{-1}w=PT''^{-1}w$; so these $v$ are the $Pw'$,
$w'\in W(\mathbf C_\ell^\flat)^{d'}$. Also $Q\cdot T'^{-1}X=\Pi X\equiv1$ modulo $\ell^n$.
Lemma~\ref{gal-structures-smith} gives $T_n$, $R$ and $R'$.
\end{proof}

\begin{lemma}\label{gal-structures-limit}
Let $\ell$ be a prime and $K\subset\bar{\mathbf Q}_\ell$ a finite extension of $\mathbf Q_\ell$ with a
Breuil--Kisin datum $(S,E,\iota)$ and a cyclotomic element $\tau$ (Definition~\ref{gal-cover-breuil-kisin}). Let
$N\geq0$, let $\rho:G_K\to\operatorname{GL}_d(\mathbf Z_\ell)$ be a homomorphism with a lattice tower
$(\rho'_n,\Pi_n,f_n)_n$ (Definition~\ref{gal-cover-definition}), and for every $n$ let $(T_n,A_n,A'_n)$ be a
Breuil--Kisin basis of height $N$ of $\rho'_n$ (Definition~\ref{gal-structures-basis}). Then there are a
Breuil--Kisin basis $(T,A,A')$ of height $N$ of $\rho$ and an integer $r\geq0$ such that for every $n\geq0$
there are $m\geq n$, $R\in M_{d\times d'_m}(S)$ and $G\in M_{d'_m\times d}(S)$ with $\Pi_mT_m\equiv T\iota(R)$
modulo $\ell^n$ in $W(\mathbf C_\ell^\flat)$ and $RG\equiv\ell^r$ modulo $\ell^n$ in $S$.
\end{lemma}
\leanref{Development/AbelianVarietyTheory/Proofs/InertiaEigenvalues/ExistsBreuilKisinBasisLimitLatticeStatement.lean}{GlobalTheory.T.exists_breuilKisin_basis_limit_lattice}
\leanref{Development/AbelianVarietyTheory/Proofs/InertiaEigenvalues/ExistsBreuilKisinBasisLimitLattice.lean}{GlobalTheory.BKLim.limit_pre_exists}

\begin{proof}
Congruences modulo $\ell^n$ are entrywise in $W(\mathbf C_\ell^\flat)$. Write $\iota$ also for the map
$S\to W(\mathbf C_\ell^\flat)$, which is injective, $u$ for $\iota(u)=[\pi^\flat]$, and $G_\infty$ for the
subgroup of $G_K$ fixing $\pi^\flat$. Let $\mathcal O_{\mathcal E}\subset W(\mathbf C_\ell^\flat)$ be the Fontaine
ring of the datum (Definition~\ref{gal-structures-etale}), the $\ell$-adic closure of $\iota(S)[u^{-1}]$. By
Fontaine's theory of \'etale $\phi$-modules \cite{Fon90},
$\mathcal O_{\mathcal E}$ is a subring, closed for the $\ell$-adic topology, that contains $w$ whenever it
contains $\ell w$; $\iota(s)$ is a unit of $\mathcal O_{\mathcal E}$ for every $s\in S$ not divisible by $\ell$,
and an element of $\mathcal O_{\mathcal E}$ that is a unit of $W(\mathbf C_\ell^\flat)$ is a unit of
$\mathcal O_{\mathcal E}$; $\mathcal O_{\mathcal E}$ is stable under $\phi$ and fixed by $G_\infty$.
Lemma~\ref{gal-structures-descent} gives two further facts:
\begin{enumerate}
\item[(a)] a matrix $X$ over $\mathcal O_{\mathcal E}$ such that $\tau^MX$ is over $A_{\inf}$ for some $M$ is
  $\iota$ of a matrix over $S$;
\item[(b)] if $x\in\mathcal O_{\mathcal E}^d$, $P$ is a matrix over $\mathcal O_{\mathcal E}$, $a$ is a vector
  over $A_{\inf}$ and $x\equiv Pa$ modulo $\ell^n$, then $x\equiv P\iota(s)$ modulo $\ell^n$ for a vector $s$
  over $S$.
\end{enumerate}
For every $k$, Lemma~\ref{gal-structures-poles} gives $Z_k=\tau^NT_k$ over $A_{\inf}$ and $Z'_k$ over
$A_{\inf}$ with $T_kZ'_k=\tau^N$; Lemma~\ref{gal-structures-functorial} applied to $f_k$ gives $H_k$ over $S$
with $T_k\iota(H_k)=f_kT_{k+1}$; and the tower gives $X_k\in M_{d'_k\times d}(\mathbf Z_\ell)$ with
$\Pi_kX_k\equiv1$ modulo $\ell^k$.

\emph{Step 1: an \'etale basis.} There are $T_0\in\operatorname{GL}_d(W(\mathbf C_\ell^\flat))$ and, for every
$n$, matrices $R_n$ and $R'_n$ over $\mathcal O_{\mathcal E}$ with $\Pi_nT_n\equiv T_0R_n$ and
$R_nR'_n\equiv1$ modulo $\ell^n$. For every $k\geq1$ there are $B_k\in\operatorname{GL}_d(W(\mathbf
C_\ell^\flat))$ and $P_k,P'_k$ over $\mathcal O_{\mathcal E}$ with $\Pi_kT_k\equiv B_kP_k$ and
$\Pi_kT_kP'_k\equiv B_k$ modulo $\ell^k$: as the tower gives $\Pi_k\rho'_k(g)\equiv\rho(g)\Pi_k$ and
$\Pi_kX_k\equiv1$ modulo $\ell^k$, this is Lemma~\ref{gal-structures-etale-basis} with $n=k$, applied to
$\rho'_k$, $\Pi_k$ and the basis $(T_k,A_k,A'_k)$. Then $B_kP_kP'_k\equiv\Pi_kT_kP'_k\equiv B_k$, and
$P_kP'_k\equiv1$ modulo $\ell^k$, as $B_k$ is invertible. Put $W_k=P_k\iota(H_k)P'_{k+1}$. Since
$\Pi_kf_k\equiv\Pi_{k+1}$ modulo $\ell^k$,
\[
B_{k+1}\equiv\Pi_{k+1}T_{k+1}P'_{k+1}\equiv\Pi_kT_k\iota(H_k)P'_{k+1}\equiv B_kW_k\pmod{\ell^k},
\]
so $W_k$ is invertible over $W(\mathbf C_\ell^\flat)$, and hence over $\mathcal O_{\mathcal E}$. Put $C_1=1$ and
$C_{k+1}=W_k^{-1}C_k$. Then
$B_{k+1}C_{k+1}\equiv B_kC_k$ modulo $\ell^k$, so $B_kC_k$ converges to a matrix $T_0$, invertible since
$T_0\equiv B_1$ modulo $\ell$. The matrices $R_k=C_k^{-1}P_k$ and $R'_k=P'_kC_k$ satisfy
$\Pi_kT_k\equiv B_kP_k\equiv T_0R_k$ and $R_kR'_k\equiv1$ modulo $\ell^k$.

\emph{Step 2.} For $g\in G_\infty$, $\rho(g)g(T_0)=T_0$, and $\phi(T_0)=T_0A_0$ with $A_0$ invertible over
$\mathcal O_{\mathcal E}$. Indeed $T_0\equiv\Pi_nT_nR'_n$ modulo $\ell^n$. For $g\in G_\infty$, $g$ fixes
$R'_n$, $\rho(g)\Pi_n\equiv\Pi_n\rho'_n(g)$ modulo $\ell^n$ and $\rho'_n(g)g(T_n)=T_n$, so
$\rho(g)g(T_0)\equiv T_0$ modulo every $\ell^n$. Since $\phi(T_n)=T_n\iota(E)^{-N}\iota(A_n)$ and
$T_n=\phi(T_n)\iota(E)^{-N}\iota(A'_n)$, we get $T_0^{-1}\phi(T_0)\equiv R_n\iota(E)^{-N}\iota(A_n)\phi(R'_n)$ and
$\phi(T_0)^{-1}T_0\equiv\phi(R_n)\iota(E)^{-N}\iota(A'_n)R'_n$ modulo $\ell^n$, both over
$\mathcal O_{\mathcal E}$; as $\mathcal O_{\mathcal E}$ is closed, $A_0$ and $A_0^{-1}$ are over it.

\emph{Step 3.} $R_n\iota(H_n)\equiv R_{n+1}$ modulo $\ell^n$, since
$T_0R_n\iota(H_n)\equiv\Pi_nT_n\iota(H_n)=\Pi_nf_nT_{n+1}\equiv\Pi_{n+1}T_{n+1}\equiv T_0R_{n+1}$.

\emph{Step 4.} For every $n$ there is $c\geq0$ such that for all $m\geq n$ and $s\in S^{d'_n}$ there is
$s'\in S^{d'_m}$ with $R_n\iota(u^cs)\equiv R_m\iota(s')$ modulo $\ell^n$. First, $u^c\equiv\tau^{2N}y$ modulo
$\ell^n$ for some $c\geq0$ and $y\in A_{\inf}$. Indeed, $x=\tau^{2N}$ is a unit of $W(\mathbf C_\ell^\flat)$ by
Lemma~\ref{gal-structures-cyclotomic}, so its zeroth component $x_0\in\mathcal O^\flat$ is nonzero. The element
$\pi^\flat$ is not a unit of $\mathcal O^\flat$ (Definition~\ref{gal-structures-etale}), so $|\pi^\flat|<1$, and
$|\pi^\flat|^{c_0}\leq|x_0|$ for some $c_0$. Then $(\pi^\flat)^{c_0}=x_0b$ with
$b\in\mathcal O^\flat$. Write $x=[x_0]+\ell z$ with $z\in A_{\inf}$. Then $[b]x=u^{c_0}-w$ with $w=-\ell[b]z$.
Since $(u^{c_0}-w)\sum_{i<n}u^{c_0i}w^{n-1-i}=u^{c_0n}-w^n$ and $\ell^n$ divides $w^n$, the element
$y=[b]\sum_{i<n}u^{c_0i}w^{n-1-i}$ satisfies $xy\equiv u^{c_0n}$ modulo $\ell^n$, and $c=c_0n$ works. Then,
modulo $\ell^n$,
\[
T_0R_n\iota(u^cs)\equiv\Pi_nZ_n\tau^Ny\iota(s)\equiv\Pi_mX_m\Pi_nZ_n\tau^Ny\iota(s)
=\Pi_mT_mZ'_mX_m\Pi_nZ_ny\iota(s)\equiv T_0R_ma,
\]
with $a=Z'_mX_m\Pi_nZ_ny\iota(s)$ over $A_{\inf}$. Cancelling $T_0$ gives $R_n\iota(u^cs)\equiv R_ma$, and (b)
replaces $a$ by $\iota(s')$.

\emph{Step 5: limit vectors.} Call $v\in W(\mathbf C_\ell^\flat)^d$ a limit vector if for every $k$ there is
$s_k\in S^{d'_k}$ with $v\equiv R_k\iota(s_k)$ modulo $\ell^k$. For every $n$ there is $m\geq n$ such that for
every $s\in S^{d'_m}$ some limit vector is congruent to $R_m\iota(s)$ modulo $\ell^n$. By Step 3 the images of
$s\mapsto R_k\iota(s)$ modulo $\ell^n$ decrease as $k$ grows, and by Step 4 they all contain the image of
$s\mapsto R_n\iota(u^cs)$. Since $S/(\ell^n,u^c)$ is finite, the images stabilize, and the usual
Mittag-Leffler argument produces the limit vectors.

\emph{Step 6.} Every limit vector $v$ lies in $\mathcal O_{\mathcal E}^d$, and $\tau^NT_0v$ is over $A_{\inf}$.
The first holds as $\mathcal O_{\mathcal E}$ is closed. For the second, modulo $\ell^k$,
$\tau^NT_0v\equiv\tau^NT_0R_k\iota(s_k)\equiv\Pi_kZ_k\iota(s_k)$, which is over $A_{\inf}$; so
Lemma~\ref{gal-structures-saturated}, with $c=1$ and $r=0$, applies to every entry.

\emph{Step 7.} There is a matrix $V$, invertible over $\mathcal O_{\mathcal E}$, whose columns are limit vectors.
Take $m$ for $n=1$ in Step 5. For some $k$, $u^kR'_m\equiv\iota(\Sigma)$ modulo $\ell$ with $\Sigma$ over $S$.
Step 5 gives limit vectors $v_j\equiv R_m\iota(\Sigma e_j)\equiv u^kR_mR'_me_j\equiv u^ke_j$ modulo $\ell$. So the
matrix $V$ with columns $v_j$ is congruent to the unit $u^k$ modulo $\ell$, and it is invertible over the
$\ell$-adically complete ring $\mathcal O_{\mathcal E}$.

\emph{Step 8: the saturated module.} Let $\mathcal M$ be the set of $y\in A_{\inf}^d$ of the form
$\tau^NT_0w$ with $w\in\mathcal O_{\mathcal E}^d$. It is an $S$-module through $\iota$, closed for the
$\ell$-adic topology, containing $y$ whenever it contains $\ell y$, and by Steps 6 and 7 it contains the columns
of $\tau^NT_0V$, a matrix over $A_{\inf}$ with nonzero determinant. Lemma~\ref{gal-structures-free}, applied
with $Y_1=\tau^NT_0V$ and $Y'=\tau^NT_0$, gives a matrix $Y$ over $A_{\inf}$, invertible over
$W(\mathbf C_\ell^\flat)$, with $\mathcal M=Y\iota(S^d)$. Put $W=(\tau^NT_0)^{-1}Y$. Its columns lie in $\mathcal O_{\mathcal E}^d$, and
$V=W\iota(\Gamma)$ with $\Gamma$ over $S$, as the columns of $\tau^NT_0V$ lie in $\mathcal M$; so $W$ is
invertible over $\mathcal O_{\mathcal E}$.

\emph{Step 9: links.} For every $n$ there are $m\geq n$, $R$ over $S$ and $Q$ over $A_{\inf}$ with
$\Pi_mT_m\equiv T_0W\iota(R)$ and $Y\iota(R)Q\equiv\tau^{2N}$ modulo $\ell^n$. Take $m$ from Step 5. For each
basis vector $e_j$ of $S^{d'_m}$, Step 5 gives a limit vector $v_j\equiv R_me_j$ modulo $\ell^n$; by Steps 6 and
8, $\tau^NT_0v_j=Y\iota(c_j)$ with $c_j\in S^d$, that is $v_j=W\iota(c_j)$. Let $R$ have columns $c_j$. Then
$\Pi_mT_m\equiv T_0R_m\equiv T_0W\iota(R)$, and with $Q=Z'_mX_m$,
\[
Y\iota(R)Q=\tau^NT_0W\iota(R)Z'_mX_m\equiv\tau^N\Pi_mT_mZ'_mX_m=\tau^{2N}\Pi_mX_m\equiv\tau^{2N}
\pmod{\ell^n}.
\]

\emph{Step 10.} $Y'=\tau^{2N}Y^{-1}$ is over $A_{\inf}$, so $YY'=Y'Y=\tau^{2N}$. Indeed
$Y'-\iota(R)Q=Y^{-1}(\tau^{2N}-Y\iota(R)Q)\equiv0$ modulo $\ell^n$ by Step 9, and
Lemma~\ref{gal-structures-saturated} applies.

\emph{Step 11.} Put $T=T_0W$. Since $W$ is over $\mathcal O_{\mathcal E}$, $\rho(g)g(T)=T$ for $g\in G_\infty$,
and $\phi(T)=TA_1$ with $A_1=W^{-1}A_0\phi(W)$ invertible over $\mathcal O_{\mathcal E}$.

\emph{Step 12.} Since $Y=\tau^NT$, $Y'=\tau^NT^{-1}$ and $\phi(\tau)=\iota(E)\tau$, the matrix
$\iota(E)^NA_1=\iota(E)^NT^{-1}\phi(T)$ is over $\mathcal O_{\mathcal E}$ and satisfies
$\tau^{2N}\iota(E)^NA_1=Y'\phi(Y)$, which is over $A_{\inf}$. By (a), $\iota(E)^NA_1=\iota(A)$ with $A$ over
$S$. Likewise $\tau^{2N}\iota(E)^NA_1^{-1}=\phi(Y')Y$, so $\iota(E)^NA_1^{-1}=\iota(A')$ with $A'$ over $S$.
Then $\iota(AA')=\iota(E)^{2N}$, so $AA'=E^{2N}$, and $\iota(E)^N\phi(T)=T\iota(A)$. With Step 11,
$(T,A,A')$ is a Breuil--Kisin basis of height $N$ of $\rho$, and by Step 9, $\Pi_mT_m\equiv T\iota(R)$ modulo
$\ell^n$.

\emph{Step 13.} There is $r\geq0$, depending only on $N$, such that if $\iota(R)a\equiv\tau^{2N}$ modulo
$\ell^nA_{\inf}$ for $R$ over $S$ and $a$ over $A_{\inf}$, then $RG\equiv\ell^r$ modulo $\ell^n$ for some $G$ over
$S$: the ring $A_{\inf}$ has a topological basis $(b_k,\gamma_k)_k$ over $S$
(Lemma~\ref{gal-structures-topological-exists}), the ideal of $S$ generated by the $\gamma_k(\tau^{2N})$ contains a
power $\ell^r$ (Lemma~\ref{gal-structures-coefficients}), say $\ell^r=\sum_kt_k\gamma_k(\tau^{2N})$, and
$G=\sum_kt_k\gamma_k(a)$, with $\gamma_k$ applied entrywise, works: as the $\gamma_k$ are $S$-linear,
$RG=\sum_kt_k\gamma_k(\iota(R)a)\equiv\sum_kt_k\gamma_k(\tau^{2N})=\ell^r$ modulo $\ell^n$. Multiplying the congruence of Step 9
by $Y'$ on the left and by $Y$ on the right gives $\tau^{2N}(\iota(R)QY-\tau^{2N})\equiv0$ modulo $\ell^n$;
as $\tau$ is a unit of $W(\mathbf C_\ell^\flat)$ and $A_{\inf}\cap\ell^nW(\mathbf C_\ell^\flat)=\ell^nA_{\inf}$,
$\iota(R)QY\equiv\tau^{2N}$ modulo $\ell^nA_{\inf}$, and $a=QY$ gives $G$.
% Lean: B_k is T.exists_etaleBasis_modPow (BKLim.etBasis_modPow), Lemma gal-structures-etale-basis, which is
% written out through Lemmas gal-structures-points to gal-structures-smith. The properties of O_E cited from
% Fontaine are OEDvr.main, T.breuilKisin_OE_stability_properties and T.mem_OE_of_invariant_frobenius_solution.
\end{proof}

\begin{lemma}\label{gal-structures-limit-frame}
Let $\ell$ be a prime and $K\subset\bar{\mathbf Q}_\ell$ a finite extension of $\mathbf Q_\ell$ with a
Breuil--Kisin datum $(S,E,\iota)$ (Definition~\ref{gal-cover-breuil-kisin}). Let $N,r\geq0$, let
$\rho:G_K\to\operatorname{GL}_d(\mathbf Z_\ell)$ be a homomorphism with a lattice tower $(\rho'_n,\Pi_n,f_n)_n$
(Definition~\ref{gal-cover-definition}), for every $n$ let $(T_n,A_n,A'_n,\varepsilon_n)$ be a Breuil--Kisin
structure of height $N$ on $\rho'_n$, and let $(T,A,A')$ be a Breuil--Kisin basis of height $N$ of $\rho$
(Definition~\ref{gal-structures-basis}). Suppose that for every $n$ there are $m\geq n$,
$R\in M_{d\times d'_m}(S)$ and $G\in M_{d'_m\times d}(S)$ with $\Pi_mT_m\equiv T\iota(R)$ modulo $\ell^n$ in
$W(\mathbf C_\ell^\flat)$ and $RG\equiv\ell^r$ modulo $\ell^n$ in $S$. Then $\varepsilon(g)=T^{-1}\rho(g)g(T)$ has
entries in $A_{\inf}$ for every $g\in G_K$, so $(T,A,A',\varepsilon)$ is a Breuil--Kisin structure of height
$N$ on $\rho$. If every $(T_n,A_n,A'_n,\varepsilon_n)$ is crystalline, so is $(T,A,A',\varepsilon)$.
\end{lemma}
\leanref{Development/AbelianVarietyTheory/Proofs/InertiaEigenvalues/ExistsBreuilKisinLatticeFrameOfLinkedPreframeStatement.lean}{GlobalTheory.T.exists_breuilKisinLatticeFrame_of_linked_preframe}
\leanref{Development/AbelianVarietyTheory/Proofs/InertiaEigenvalues/ExistsBreuilKisinLatticeFrameOfLinkedPreframe.lean}{GlobalTheory.BKLim.frame_of_links}
\leanref{Development/AbelianVarietyTheory/Proofs/InertiaEigenvalues/BreuilKisinFrameIsCrOfLinkedTowerStatement.lean}{GlobalTheory.T.breuilKisinFrame_isCr_of_linked_tower}
\leanref{Development/AbelianVarietyTheory/Proofs/InertiaEigenvalues/BreuilKisinFrameIsCrOfLinkedTower.lean}{GlobalTheory.BKLim.CrLimit.isCr_of_link}
\leanref{Development/AbelianVarietyTheory/Proofs/InertiaEigenvalues/BreuilKisinLatticeFrameCocyclesLinkedStatement.lean}{GlobalTheory.T.breuilKisinLatticeFrame_cocycles_linked}

\begin{proof}
Fix $g\in G_K$ and $n$, and take $m$, $R$ and $G$ as in the hypothesis. The element $g$ fixes the entries of
$\Pi_m$, preserves $\ell^nW(\mathbf C_\ell^\flat)$, and $\Pi_m\rho'_m(g)\equiv\rho(g)\Pi_m$ modulo $\ell^m$.
Modulo $\ell^n$,
\[
T\varepsilon(g)g(\iota(R))=\rho(g)g(T\iota(R))\equiv\rho(g)\Pi_mg(T_m)\equiv\Pi_m\rho'_m(g)g(T_m)
=\Pi_mT_m\varepsilon_m(g)\equiv T\iota(R)\varepsilon_m(g).
\]
As $T$ is invertible, $\varepsilon(g)g(\iota(R))\equiv\iota(R)\varepsilon_m(g)$ modulo $\ell^n$. Multiplying on
the right by $g(\iota(G))$ and using $g(\iota(RG))\equiv\ell^r$ gives
$\ell^r\varepsilon(g)\equiv\iota(R)\varepsilon_m(g)g(\iota(G))$ modulo $\ell^n$, a matrix over $A_{\inf}$. By
Lemma~\ref{gal-structures-saturated} with $c=1$, $\varepsilon(g)$ is over $A_{\inf}$.

Suppose that every $\varepsilon_n(g)-1$ is divisible by $[\pi^\flat]$. Both sides of
$\varepsilon(g)g(\iota(R))\equiv\iota(R)\varepsilon_m(g)$ are now over $A_{\inf}$, and
$A_{\inf}\cap\ell^nW(\mathbf C_\ell^\flat)=\ell^nA_{\inf}$, so the congruence holds modulo $\ell^nA_{\inf}$. The
element $g$ fixes $\iota(W(\mathbf F_{\ell^h}))$, and $g([\pi^\flat])=[g(\pi^\flat)]$ is divisible by
$[\pi^\flat]$, since $g(\pi^\flat)/\pi^\flat$ has absolute value $1$. Writing $s\in S$ as its constant term
plus $u$ times an element of $S$, it follows that $[\pi^\flat]$ divides $g(\iota(s))-\iota(s)$. So
$g(\iota(G))=\iota(G)+[\pi^\flat]Z_2$ and $\varepsilon_m(g)=1+[\pi^\flat]Z_1$ with $Z_1,Z_2$ over $A_{\inf}$.
Modulo $\ell^nA_{\inf}$,
\[
\ell^r\varepsilon(g)\equiv\iota(R)\varepsilon_m(g)g(\iota(G))
=\iota(RG)+[\pi^\flat]\bigl(\iota(R)Z_2+\iota(R)Z_1g(\iota(G))\bigr)
\equiv\ell^r+[\pi^\flat]\bigl(\iota(R)Z_2+\iota(R)Z_1g(\iota(G))\bigr).
\]
The element $[\pi^\flat]$ is a unit of $W(\mathbf C_\ell^\flat)$, as its zeroth component $\pi^\flat$ is
nonzero. Lemma~\ref{gal-structures-saturated} with $c=[\pi^\flat]$ shows that $[\pi^\flat]$ divides every
entry of $\varepsilon(g)-1$.
\end{proof}

\begin{lemma}\label{gal-cover-structures}
Let $\ell$ be a prime and $K\subset\bar{\mathbf Q}_\ell$ a finite extension of $\mathbf Q_\ell$, with the
notation of Definition~\ref{gal-cover-breuil-kisin}.
\begin{enumerate}
\item $K$ has a Breuil--Kisin datum, and every Breuil--Kisin datum for $K$ has a cyclotomic element.
\item Fix a Breuil--Kisin datum for $K$ with a cyclotomic element. Let $L_0\subset\bar{\mathbf Q}_\ell$ be
  a finite extension of $\mathbf Q_\ell$ containing $\sigma(K)$ for every embedding
  $\sigma:K\to\bar{\mathbf Q}_\ell$, let $a\geq0$, let $r:G_K\to\operatorname{GL}_a(\bar{\mathbf Q}_\ell)$
  be a homomorphism, $C$ a set of homomorphisms $G_K\to\operatorname{GL}_a(\bar{\mathbf Q}_\ell)$, $J$ a
  map from the embeddings $K\to\bar{\mathbf Q}_\ell$ to finite multisets of integers, and fix a cover of
  $r$ relative to $C$ and $L_0$ (Definition~\ref{gal-cover-definition}) with representation $\rho$ and
  lattice tower $(\rho'_m,\Pi_m,f_m)_m$. Suppose that for every $\rho_1\in C$ and every finite extension
  $L\supset L_0$ of $\mathbf Q_\ell$ in $\bar{\mathbf Q}_\ell$ with $\rho_1(G_K)\subset\operatorname{GL}_a(L)$,
  the representation $\rho_1$ of $G_K$ on $L^a$ is crystalline with Hodge--Tate multiset
  $J(\sigma)$ at every embedding $\sigma:K\to L$, in the sense of Definition~\ref{gal-p-adic-hodge}.
  Then there are $N,r_0\geq0$ and crystalline Breuil--Kisin structures of height $N$ on $\rho$ and on
  every $\rho'_m$, such that for every $n$ there is $m\geq n$ for which the structure on $\rho$ and the
  structure on $\rho'_m$ are linked at $(n,r_0)$ along $\Pi_m$.
\item Let $\tau$ be a cyclotomic element for a Breuil--Kisin datum for $K$, and $(T,A,A',\varepsilon)$ a
  Breuil--Kisin structure of height $N$ on a map $\rho:G_K\to M_d(\mathbf Z_\ell)$. Then there are
  $Y,Y'\in M_d(A_{\inf})$ with $Y=\tau^NT$ in $M_d(W(\mathbf C_\ell^\flat))$ and $YY'=\tau^{2N}$.
\end{enumerate}
\end{lemma}
\leanref{Development/AbelianVarietyTheory/Proofs/InertiaEigenvalues/NonemptyBreuilKisinDatum.lean}{GlobalTheory.BSP.nonempty_bk_of_finExt}
\leanref{Development/AbelianVarietyTheory/Proofs/InertiaEigenvalues/ExistsBreuilKisinFrobeniusPeriod.lean}{GlobalTheory.BSP.exists_isPer}
\leanref{Development/AutomorphicArithmetic/Foundations/GaloisRepresentations/ClosedPhiFunctoriality/IntBasisAlt.lean}{UpperTheory.Gal.ClosedPhiFromClosedTrace.FrEx}
\leanref{Development/AutomorphicArithmetic/Foundations/GaloisRepresentations/ClosedPhiFunctoriality/FrExHolds.lean}{UpperTheory.Gal.ClosedPhiFromClosedTrace.frEx_holds}
\leanref{Development/AbelianVarietyTheory/Proofs/InertiaEigenvalues/BreuilKisinPreframeBoundedPolesStatement.lean}{GlobalTheory.T.breuilKisinPreframe_bounded_poles}

\begin{proof}
(1) is Lemmas~\ref{gal-structures-datum} and~\ref{gal-structures-cyclotomic}.

For (2), let $N$ bound the absolute values of the members of the multisets $J(\sigma)$. For every $m$, the
cover gives a layer for $\rho'_m$, so by Lemma~\ref{gal-structures-layer} the representation of $G_K$ on
$\mathbf Q_\ell^{d'_m}$ through $\rho'_m$ is crystalline, with $\operatorname{Fil}^{-N}D_{\mathrm{dR}}$ the whole
space and $\operatorname{Fil}^{N+1}D_{\mathrm{dR}}=0$. As $\rho'_m$ is continuous,
Lemma~\ref{gal-structures-kisin} gives a crystalline Breuil--Kisin structure $(T_m,A_m,A'_m,\varepsilon_m)$ of
height $N$ on $\rho'_m$. Lemma~\ref{gal-structures-limit}, applied to the Breuil--Kisin bases
$(T_m,A_m,A'_m)$, gives a Breuil--Kisin basis $(T,A,A')$ of height $N$ of $\rho$ and $r_0\geq0$ such that for
every $n$ there are $m\geq n$, $R$ and $G$ over $S$ with $\Pi_mT_m\equiv T\iota(R)$ and $RG\equiv\ell^{r_0}$
modulo $\ell^n$. By Lemma~\ref{gal-structures-limit-frame}, $\varepsilon(g)=T^{-1}\rho(g)g(T)$ makes
$(T,A,A',\varepsilon)$ a crystalline Breuil--Kisin structure of height $N$ on $\rho$, and these congruences say
that it is linked at $(n,r_0)$ along $\Pi_m$ to the structure on $\rho'_m$.

(3) is Lemma~\ref{gal-structures-poles}, applied to the Breuil--Kisin basis $(T,A,A')$.
% Lean: (1) is nonempty_bk_of_finExt and exists_isPer. (2) is frEx_holds: the layers are crystalline
% and semistable by T.isSemistable_latticeInMatrixAlgebraLayer with jumps in [-N,N] by
% T.layer_DdR_bounds_alt; structures exist by T.exists_breuilKisinFrame_semistable_boundedJumps;
% the limit structure and links are T.exists_deRhamFrame_limitLattice_breuilKisinTower
% (BKLim.limit_frame_exists = limit_pre_exists, frame_of_links, isCr_of_link, plus isDR_of_links',
% whose de Rham condition FrEx does not use). (3) is BKLim.pre_bounded_poles.
\end{proof}

\subsection{Period ring systems and crystalline limits}\label{subsection-gal-crystalline-limits}

\begin{lemma}\label{gal-cover-weil-frobenius}
In Situation~\ref{gal-local-situation}, suppose that $w$ lies over a prime $\ell$, let
$f=[k_w:\mathbf F_\ell]$, and identify $D_A$ with $G_K$, for a finite extension
$K\subset\bar{\mathbf Q}_\ell$ of $\mathbf Q_\ell$ in which $E$ is dense, as in
Definition~\ref{gal-p-adic-hodge}. Let $\operatorname{Frob}$ be the arithmetic Frobenius of the maximal
unramified extension $\mathbf Q_\ell^{\mathrm{ur}}$ of $\mathbf Q_\ell$ in $\bar{\mathbf Q}_\ell$. Then
for every $\sigma\in W_A$, $\operatorname{Frob}^{|fd(\sigma)|}$ fixes every element of
$\mathbf Q_\ell^{\mathrm{ur}}\cap K$.
\end{lemma}
\leanref{Development/AutomorphicArithmetic/Foundations/GaloisRepresentations/ClosedPhiFunctoriality/IsOp2One.lean}{UpperTheory.Gal.ClosedPhiFromClosedTrace.FixOf}
\leanref{Development/AutomorphicArithmetic/Foundations/GaloisRepresentations/ClosedPhiFunctoriality/FixOfHolds.lean}{UpperTheory.Gal.ClosedPhiFromClosedTrace.fixOf_holds}

\begin{proof}
Let $\alpha,\beta\geq0$ with $\sigma(z)^{q_w^\beta}\equiv z^{q_w^\alpha}$ modulo $\mathfrak m_A$ for all
$z\in A$. Then $\sigma$ acts on the residue field of $\bar{\mathbf Q}_\ell$ as the $f(\alpha-\beta)$-th
power of the arithmetic Frobenius, and hence on $\mathbf Q_\ell^{\mathrm{ur}}$ as
$\operatorname{Frob}^{fd(\sigma)}$. Since $\sigma\in G_K$ fixes $K$, $\operatorname{Frob}^{fd(\sigma)}$
fixes $\mathbf Q_\ell^{\mathrm{ur}}\cap K$, and so does its inverse.
\end{proof}

\begin{lemma}\label{gal-cover-zero}
Let $\ell$ be a prime and $K\subset\bar{\mathbf Q}_\ell$ a finite extension of $\mathbf Q_\ell$ with a
Breuil--Kisin datum $(S,E,\iota)$ (Definition~\ref{gal-cover-breuil-kisin}), and write $u$ for the image
of $\iota(u)$ in $A_{\mathrm{cris}}$. For all $c,M\geq0$ there is $n\geq0$ with the following property.
If $D\in M_{d\times d'}(A_{\mathrm{cris}})$ has all entries divisible by $u$, and
$\ell^cD\equiv F\phi(D)F'$ modulo $\ell^n$ for some $F\in M_d(A_{\mathrm{cris}})$ and
$F'\in M_{d'}(A_{\mathrm{cris}})$, then all entries of $D$ are divisible by $\ell^M$. In particular, if
$\ell^cD=F\phi(D)F'$, then $D=0$.
\end{lemma}
\leanref{Development/AutomorphicArithmetic/Foundations/GaloisRepresentations/ClosedPhiFunctoriality/IntBasisAlt.lean}{UpperTheory.Gal.ClosedPhiFromClosedTrace.ZeroAp}
\leanref{Development/AutomorphicArithmetic/Foundations/GaloisRepresentations/ClosedPhiFunctoriality/ZeroAp.lean}{UpperTheory.Gal.ClosedPhiFromClosedTrace.zeroAp_holds}
\leanref{Development/AbelianVarietyTheory/Proofs/InertiaEigenvalues/AcrisFrobeniusMatrixEqZeroOfDvdStatement.lean}{GlobalTheory.T.acris_frobenius_matrix_eq_zero_of_dvd}
\leanref{Development/AbelianVarietyTheory/Proofs/InertiaEigenvalues/PrimePowDvdTiltUniformizerPowAcrisEventuallyStatement.lean}{GlobalTheory.T.primePow_dvd_tiltUniformizer_pow_acris_eventually}

\begin{proof}
Let $e=\deg E$. Since $u^e\equiv\iota(E)$ modulo $\ell$ and $\iota(E)^\ell\in\ell!\,A_{\mathrm{cris}}$,
$\ell$ divides $u^{e\ell}$ in $A_{\mathrm{cris}}$. Hence $\ell^{ct+M}$ divides $u^{\ell^t-1}$ once
$e\ell(ct+M)\leq\ell^t-1$; fix such a $t$ and put $n=ct+M$. Write $\ell^cD=F\phi(D)F'+\ell^nG$. Since
$\phi(u)=u^\ell$, induction on $j$ gives matrices $N_j,H_j$ with $\ell^{cj}D=u^{\ell^j}N_j+\ell^nH_j$:
apply $\phi$ to the identity for $j$ and multiply by $F$ and $F'$. For $j=t$ every term on the right
is divisible by $\ell^{ct+M}$, and $A_{\mathrm{cris}}$ has no $\ell$-torsion, so $\ell^M$ divides $D$. If
$\ell^cD=F\phi(D)F'$, this holds for every $M$, and $A_{\mathrm{cris}}$ is $\ell$-adically separated.
\end{proof}

\begin{lemma}\label{gal-cover-section}
Let $\ell$ be a prime and $K\subset\bar{\mathbf Q}_\ell$ a finite extension of $\mathbf Q_\ell$ with a
Breuil--Kisin datum $(S,E,\iota)$ (Definition~\ref{gal-cover-breuil-kisin}). Write $u$ for the image of
$\iota(u)$ in $A_{\mathrm{cris}}$, and for $x\in S$ put $x^\phi=\phi(\iota(x))$ and
$x^0=\phi(\iota(x(0)))$ in $A_{\mathrm{cris}}$, where $x(0)\in W(\mathbf F_{\ell^h})$ is the constant
term of $x$; apply these to matrices entrywise. For every $N\geq0$ there is $k\geq0$ such that for all
$d$ and all $A,A'\in M_d(S)$ with $AA'=E^{2N}$ there are $Q,Q'\in M_d(A_{\mathrm{cris}})$ with
$QQ'=\ell^{2k}$, $Q\equiv\ell^k$ and $Q'\equiv\ell^k$ modulo $u$, and
\[
(E^\phi)^NQA^0=(E^0)^NA^\phi\phi(Q).
\]
\end{lemma}
\leanref{Development/AutomorphicArithmetic/Foundations/GaloisRepresentations/ClosedPhiFunctoriality/IntBasisAlt.lean}{UpperTheory.Gal.ClosedPhiFromClosedTrace.SecU}
\leanref{Development/AutomorphicArithmetic/Foundations/GaloisRepresentations/ClosedPhiFunctoriality/JointApproxOfMem.lean}{UpperTheory.Gal.ClosedPhiFromClosedTrace.IsSecQ}
\leanref{Development/AutomorphicArithmetic/Foundations/GaloisRepresentations/ClosedPhiFunctoriality/SecUHolds.lean}{UpperTheory.Gal.ClosedPhiFromClosedTrace.secU_holds}
\leanref{Development/AutomorphicArithmetic/Foundations/GaloisRepresentations/ClosedPhiFunctoriality/IntBasisAlt.lean}{UpperTheory.Gal.ClosedPhiFromClosedTrace.CoreU}
\leanref{Development/AutomorphicArithmetic/Foundations/GaloisRepresentations/ClosedPhiFunctoriality/CoreUHolds.lean}{UpperTheory.Gal.ClosedPhiFromClosedTrace.coreU_holds}

\begin{proof}
We have $E^0=\ell w_0$ and $E^\phi=\ell v$ with units $v,w_0$ of $A_{\mathrm{cris}}$: the first since the
constant term of $E$ is $\ell$ times a unit, the second since $\phi(E)=E^\ell+\ell\delta(E)$ with
$\delta(E)$ a unit of $S$ and $\iota(E)^\ell/\ell$ topologically nilpotent in $A_{\mathrm{cris}}$. Put
$\alpha=v^{-N}A^\phi$, $\alpha'=v^{-N}A'^\phi$, $\beta=w_0^{-N}A^0$ and $\beta'=w_0^{-N}A'^0$, so that
$\alpha\alpha'=\beta\beta'=\ell^{2N}$. Since $x^\phi-x^0$ is divisible by $\phi(u)=u^\ell$ for every
$x\in S$, we get $\ell(\alpha-\beta)=u^\ell\Gamma$ for a matrix $\Gamma$ over $A_{\mathrm{cris}}$. As in
the proof of Lemma~\ref{gal-cover-zero}, choose $n_0$ with
$\ell^{(2N+1)n+1}\mid u^{\ell^n-1}$ for all $n>n_0$.

Define $P_0=1$ and $P_{n+1}=\alpha\phi(P_n)\beta'$. By induction, using $\beta\beta'=\ell^{2N}$ and
$\phi(u^{\ell^n})=u^{\ell^{n+1}}$, the matrix $\ell(P_{n+1}-\ell^{2N}P_n)$ is divisible by
$u^{\ell^{n+1}}$, and $\ell^n(P_n-\ell^{2Nn})$ by $u$. Hence $P_{n+1}-\ell^{2N}P_n$ is divisible by
$\ell^{(2N+1)(n+1)}u$ for $n\geq n_0$, so $P_{n_0+m}=\ell^{2Nm}S_m$ with $S_{m+1}-S_m$ divisible by
$\ell^{m+1}u$, and the $S_m$ converge $\ell$-adically to a matrix $S$ with $S\beta=\alpha\phi(S)$ and
$\ell^{n_0}(S-\ell^{2Nn_0})$ divisible by $u$. The same construction with $P'_{n+1}=\beta\phi(P'_n)\alpha'$
gives $S'$, and $P_nP'_n=\ell^{4Nn}$ gives $SS'=\ell^{4Nn_0}$. Put $k=n_0+2Nn_0$, $Q=\ell^{n_0}S$ and
$Q'=\ell^{n_0}S'$. Then $QQ'=\ell^{2k}$, $Q\equiv Q'\equiv\ell^k$ modulo $u$, and multiplying
$Q\beta=\alpha\phi(Q)$ by $(\ell vw_0)^N$ gives the displayed identity.
\end{proof}

\begin{lemma}\label{gal-cover-period}
Let $\ell$ be a prime and $K\subset\bar{\mathbf Q}_\ell$ a finite extension of $\mathbf Q_\ell$ with a
Breuil--Kisin datum $(S,E,\iota)$ and a cyclotomic element $\tau$ (Definition~\ref{gal-cover-breuil-kisin}),
and use the notation of Lemma~\ref{gal-cover-section}. Let $N\geq0$, let $k$ be as in
Lemma~\ref{gal-cover-section} for $N$, let $n_1,n_2\geq0$, let
$\rho:G_K\to\operatorname{GL}_d(\mathbf Z_\ell)$ be a homomorphism, and let $(T,A,A',\varepsilon)$ be a
crystalline Breuil--Kisin structure of height $N$ on $\rho$. Put $\tau_1=\phi(\tau)\in A_{\mathrm{cris}}$.
Then $\tau_1$ is a unit of $B_{\mathrm{cris}}$, and there are $Y,Y'\in M_d(A_{\inf})$ with $Y=\tau^NT$ and
$YY'=\tau^{2N}$, and $Q,Q'$ as in Lemma~\ref{gal-cover-section} for $A$, such that with
$Z=\phi(Y)Q$ and $Z'=Q'\phi(Y')$ in $M_d(A_{\mathrm{cris}})$:
\begin{enumerate}
\item $ZZ'=Z'Z=\tau_1^{2N}\ell^{2k}$;
\item $U=\tau_1^{-N}Z$ is a period matrix of $\rho$;
\item $s=\phi^{n_2}(\tau_1)^{3N}\ell^{4k}\phi^{n_1}(\tau_1)^N$ is a unit of $B_{\mathrm{cris}}$, and
  $V=\phi^{n_2}(\tau_1^{2N}\ell^{2k})\,\phi^{n_1}(Z)\,\phi^{n_2}(Z')$ satisfies
  $V\phi^{n_2}(U)=s\,\phi^{n_1}(U)$.
\end{enumerate}
\end{lemma}
\leanref{Development/AutomorphicArithmetic/Foundations/GaloisRepresentations/ClosedPhiFunctoriality/IntBasisAlt.lean}{UpperTheory.Gal.ClosedPhiFromClosedTrace.InvQ}
\leanref{Development/AutomorphicArithmetic/Foundations/GaloisRepresentations/ClosedPhiFunctoriality/InvQHolds.lean}{UpperTheory.Gal.ClosedPhiFromClosedTrace.invQ_holds}
\leanref{Development/AutomorphicArithmetic/Foundations/GaloisRepresentations/ClosedPhiFunctoriality/IntBasisAlt.lean}{UpperTheory.Gal.ClosedPhiFromClosedTrace.InvU}
\leanref{Development/AutomorphicArithmetic/Foundations/GaloisRepresentations/ClosedPhiFunctoriality/InvUHolds.lean}{UpperTheory.Gal.ClosedPhiFromClosedTrace.invU_holds}
\leanref{Development/AutomorphicArithmetic/Foundations/GaloisRepresentations/ClosedPhiFunctoriality/IsOp2One.lean}{UpperTheory.Gal.ClosedPhiFromClosedTrace.IsOp2}
\leanref{Development/AutomorphicArithmetic/Foundations/GaloisRepresentations/ClosedPhiFunctoriality/IsOp2One.lean}{UpperTheory.Gal.ClosedPhiFromClosedTrace.OpOne2}
\leanref{Development/AutomorphicArithmetic/Foundations/GaloisRepresentations/ClosedPhiFunctoriality/OpOneHolds.lean}{UpperTheory.Gal.ClosedPhiFromClosedTrace.opOne_holds}
\leanref{Development/AbelianVarietyTheory/Proofs/InertiaEigenvalues/FrameAinfMatrixFrobeniusGaloisEqsStatement.lean}{GlobalTheory.T.frame_ainf_matrix_frobenius_galois_eqs}
\leanref{Development/AbelianVarietyTheory/Proofs/InertiaEigenvalues/FrobeniusEquationOfTwistedCocycleDefectStatement.lean}{GlobalTheory.T.frobenius_equation_of_twisted_cocycle_defect}

\begin{proof}
Write $\phi(\tau)=\eta([\zeta]-1)$ as in Definition~\ref{gal-cover-breuil-kisin}. The element
$[\zeta]-1$ divides $t=\log[\zeta]$, a unit of $B_{\mathrm{cris}}$, so $\tau_1$ is a unit. Take $Y,Y'$
from Lemma~\ref{gal-cover-structures}(3) and $Q,Q'$ from Lemma~\ref{gal-cover-section}. Then
$Y'Y=\tau^{2N}$ as well, which gives (1). From $Y=\tau^NT$, $\phi(\tau)=\iota(E)\tau$ and the equations
of the structure one gets in $A_{\inf}$, for $g\in G_K$,
\[
\tau^N\rho(g)g(Y)=g(\tau)^NY\varepsilon(g),\qquad
\iota(E)^N\varepsilon(g)g(\iota(A))=g(\iota(E))^N\iota(A)\phi(\varepsilon(g)).
\]

\emph{Invariance of $Q$.} Let $g\in G_K$ and put $D=Q'\phi(\varepsilon(g))g(Q)-\ell^{2k}$. The entries of
$A^0$ and $E^0$ come from $W(\mathbf F_{\ell^h})$ and are fixed by $G_K$, and $g$ commutes with $\phi$.
Combining the identity of Lemma~\ref{gal-cover-section}, its image under $g$, and the second identity
above mapped to $A_{\mathrm{cris}}$ by $x\mapsto\phi(x)$, a direct computation gives
$(E^0)^{2N}D=A^0\phi(D)A'^0$, that is $\ell^{2N}D=A^0\phi(D)(w_0^{-2N}A'^0)$ with $E^0=\ell w_0$. Modulo
$u$ we have $Q\equiv Q'\equiv\ell^k$, $g(Q)\equiv\ell^k$ since $u$ divides $g(u)$, and
$\phi(\varepsilon(g))\equiv1$ since the structure is crystalline and $u$ divides $\phi(u)=u^\ell$. So
$u$ divides $D$, and $D=0$ by Lemma~\ref{gal-cover-zero}. Multiplying
$Q'\phi(\varepsilon(g))g(Q)=\ell^{2k}$ on the left by $Q$ and cancelling $\ell^{2k}$ gives
$\phi(\varepsilon(g))g(Q)=Q$.

\emph{The period matrix.} By (1), $U$ is invertible in $B_{\mathrm{cris}}$ with inverse
$\tau_1^{-N}\ell^{-2k}Z'$. Applying $\phi$ to the first identity above gives
$\tau_1^N\rho(g)g(\phi(Y))=g(\tau_1)^N\phi(Y)\phi(\varepsilon(g))$. Hence
\[
\rho(g)g(U)=g(\tau_1)^{-N}\rho(g)g(\phi(Y))g(Q)=\tau_1^{-N}\phi(Y)\phi(\varepsilon(g))g(Q)
=\tau_1^{-N}\phi(Y)Q=U,
\]
which is (2).

\emph{The operator.} The element $s$ is a product of units of $B_{\mathrm{cris}}$. Since $\phi$ fixes
$\ell$ and $\phi^{n_2}(Z')\phi^{n_2}(Z)=\phi^{n_2}(\tau_1^{2N}\ell^{2k})$ by (1),
\[
V\phi^{n_2}(U)=\phi^{n_2}(\tau_1)^{-N}\phi^{n_2}(\tau_1^{2N}\ell^{2k})^2\phi^{n_1}(Z)
=\phi^{n_2}(\tau_1)^{3N}\ell^{4k}\phi^{n_1}(\tau_1)^N\phi^{n_1}(U)=s\,\phi^{n_1}(U),
\]
which is (3).
\end{proof}

\begin{lemma}\label{gal-cover-link}
Let $\ell$ be a prime and $K\subset\bar{\mathbf Q}_\ell$ a finite extension of $\mathbf Q_\ell$ with a
Breuil--Kisin datum $(S,E,\iota)$ and a cyclotomic element $\tau$ (Definition~\ref{gal-cover-breuil-kisin}),
and use the notation of Lemma~\ref{gal-cover-section}. For all $N,M\geq0$ there is $n\geq0$ with the
following property. Let $k\geq0$, let $(T,A,A',\varepsilon)$ and $(T_1,A_1,A_1',\varepsilon_1)$ be
Breuil--Kisin structures of height $N$ on maps $\rho:G_K\to M_d(\mathbf Z_\ell)$ and
$\rho_1:G_K\to M_{d_1}(\mathbf Z_\ell)$, let $Y,Y_1$ be matrices over $A_{\inf}$ with $Y=\tau^NT$ and
$Y_1=\tau^NT_1$, let $(Q,Q')$ and $(Q_1,Q_1')$ be as in Lemma~\ref{gal-cover-section} for $A$ and
$A_1$ with the same $k$, and let $\Pi\in M_{d\times d_1}(\mathbf Z_\ell)$ and $R\in M_{d\times d_1}(S)$
with $\Pi T_1\equiv T\iota(R)$ modulo $\ell^n$ in $W(\mathbf C_\ell^\flat)$. Then
\[
\Pi\phi(Y_1)Q_1\equiv\phi(Y)QR^0\quad\text{modulo }\ell^M\text{ in }A_{\mathrm{cris}}.
\]
\end{lemma}
\leanref{Development/AutomorphicArithmetic/Foundations/GaloisRepresentations/ClosedPhiFunctoriality/IsOp2One.lean}{UpperTheory.Gal.ClosedPhiFromClosedTrace.LinkZ}
\leanref{Development/AutomorphicArithmetic/Foundations/GaloisRepresentations/ClosedPhiFunctoriality/LinkZHolds.lean}{UpperTheory.Gal.ClosedPhiFromClosedTrace.linkZ_holds}
\leanref{Development/AutomorphicArithmetic/Foundations/GaloisRepresentations/ClosedPhiFunctoriality/IntBasisAlt.lean}{UpperTheory.Gal.ClosedPhiFromClosedTrace.LinkA}
\leanref{Development/AutomorphicArithmetic/Foundations/GaloisRepresentations/ClosedPhiFunctoriality/PowDvdOfEv.lean}{UpperTheory.Gal.ClosedPhiFromClosedTrace.linkA_holds}
\leanref{Development/AutomorphicArithmetic/Foundations/GaloisRepresentations/ClosedPhiFunctoriality/IntBasisAlt.lean}{UpperTheory.Gal.ClosedPhiFromClosedTrace.LinkQ}
\leanref{Development/AutomorphicArithmetic/Foundations/GaloisRepresentations/ClosedPhiFunctoriality/DvdMulEntriesRight.lean}{UpperTheory.Gal.ClosedPhiFromClosedTrace.linkQ_holds}
\leanref{Development/AutomorphicArithmetic/Foundations/GaloisRepresentations/ClosedPhiFunctoriality/IntBasisAlt.lean}{UpperTheory.Gal.ClosedPhiFromClosedTrace.LinkRec}
\leanref{Development/AutomorphicArithmetic/Foundations/GaloisRepresentations/ClosedPhiFunctoriality/LinkRec.lean}{UpperTheory.Gal.ClosedPhiFromClosedTrace.linkRec_holds}

\begin{proof}
Let $n'$ be given by Lemma~\ref{gal-cover-zero} for $c=3N$ and $M$, and put $n=\max(n',M)$.

\emph{Integral congruences.} Multiplying $\Pi T_1\equiv T\iota(R)$ by $\tau^N$ gives
$\Pi Y_1\equiv Y\iota(R)$ modulo $\ell^n$ in $A_{\inf}$, since $A_{\inf}\cap\ell^nW(\mathbf
C_\ell^\flat)=\ell^nA_{\inf}$. Applying $\phi$ to $\Pi T_1\equiv T\iota(R)$, multiplying by $\iota(E)^N$
and using the Frobenius equations of the two structures gives $\Pi T_1\iota(A_1)\equiv
T\iota(A)\iota(\phi(R))$, and also $\Pi T_1\iota(A_1)\equiv T\iota(R)\iota(A_1)$. Since $T$ is
invertible, $\iota(RA_1-A\phi(R))\equiv0$ modulo $\ell^n$, and so $RA_1\equiv A\phi(R)$ modulo $\ell^n$
in $S$, as $S/\ell^n\to W(\mathbf C_\ell^\flat)/\ell^n$ is injective.

\emph{The sections.} Let $D=R^\phi Q_1-QR^0$. Since $Q\equiv Q_1\equiv\ell^k$ and $R^\phi\equiv R^0$
modulo $u$, the entries of $D$ are divisible by $u$. Since $x\mapsto x^\phi$ and $x\mapsto x^0$
commute with $\phi$, the congruence $RA_1\equiv A\phi(R)$ gives $R^\phi A_1^\phi\equiv A^\phi\phi(R^\phi)$
and $R^0A_1^0\equiv A^0\phi(R^0)$ modulo $\ell^n$. A direct computation from these, from the identities
of Lemma~\ref{gal-cover-section} for $Q$ and $Q_1$, and from $A_1^0A_1'^0=(E^0)^{2N}$ shows that
$\ell^{3N}D-F\phi(D)F'$, for $F=\ell^Nv^{-N}w_0^{-N}A^\phi$ and $F'=A_1'^0$ with $v,w_0$ as in the proof
of Lemma~\ref{gal-cover-section}, is a combination of $R^\phi A_1^\phi-A^\phi\phi(R^\phi)$ and
$R^0A_1^0-A^0\phi(R^0)$. So $\ell^{3N}D\equiv F\phi(D)F'$ modulo $\ell^n$, and $\ell^M$ divides $D$ by
Lemma~\ref{gal-cover-zero}.

\emph{Conclusion.} We have
\[
\Pi\phi(Y_1)Q_1-\phi(Y)QR^0=\bigl(\Pi\phi(Y_1)-\phi(Y)R^\phi\bigr)Q_1+\phi(Y)\bigl(R^\phi Q_1-QR^0\bigr).
\]
The first bracket is the image under $\phi$ of $\Pi Y_1-Y\iota(R)$, since $\phi$ fixes the entries of
$\Pi$, so it is divisible by $\ell^n$; the second is $D$. Both are divisible by $\ell^M$.
\end{proof}

\begin{definition}\label{gal-cover-fontaine-datum}
Let $G$ be a group, $F\to K_0\to K$ homomorphisms of fields, $B$ a commutative $K_0$-algebra and $B_{\mathrm{dR}}$ a
commutative $K$-algebra; through these maps both are $F$-algebras and $B_{\mathrm{dR}}$ is a $K_0$-algebra. A
\emph{Fontaine datum} for $(G,F,K_0,K,B,B_{\mathrm{dR}})$ consists of a prime $p$, a bijective $F$-algebra endomorphism
$\sigma$ of $K_0$, actions of $G$ on $B$ by $K_0$-algebra automorphisms and on $B_{\mathrm{dR}}$ by $K$-algebra
automorphisms, both written $b\mapsto g(b)$, an injective $F$-algebra endomorphism $\phi$ of $B$, an injective
$K_0$-algebra map $\iota:B\to B_{\mathrm{dR}}$, $K$-subspaces $\mathrm{Fil}^a$ of $B_{\mathrm{dR}}$ for $a\in\mathbf Z$, a
homomorphism $\chi:G\to F^\times$ and an element $t\in B$, such that:
% Lean: PeriodData G Qp K₀ K B BdR with F = Qp; the data are p (with hp), σ (with σ_bij), actB, actdR, φ (with φ_inj),
% ι (with ι_inj), Fil, χ and t.
\begin{enumerate}
\item $K_0$ and $K$ are finite-dimensional over $F$; % finK₀, finK
\item every element of $B$ fixed by $G$ is the image of an element of $K_0$, and every element of $B_{\mathrm{dR}}$
  fixed by $G$ is the image of an element of $K$; % invB, invdR
\item $\phi(c)=\sigma(c)$ for $c\in K_0$, and $\phi(g(b))=g(\phi(b))$ and $\iota(g(b))=g(\iota(b))$ for $g\in G$ and
  $b\in B$; % φ_K₀, φ_comm, ι_comm
\item $\mathrm{Fil}^{a+1}\subset\mathrm{Fil}^a$, $1\in\mathrm{Fil}^0$, $\mathrm{Fil}^a\cdot\mathrm{Fil}^b\subset
  \mathrm{Fil}^{a+b}$ and $g(\mathrm{Fil}^a)\subset\mathrm{Fil}^a$ for all $a,b\in\mathbf Z$ and $g\in G$, every element
  of $B_{\mathrm{dR}}$ lies in some $\mathrm{Fil}^a$, and $0$ is the only element that lies in all of them;
  % Fil_anti, one_mem, Fil_mul, Fil_G, Fil_exh, Fil_sep
\item $t$ is a unit of $B$, $g(t)=\chi(g)t$ for $g\in G$, $\phi(t)=pt$, $\iota(t)\in\mathrm{Fil}^1$ and
  $\iota(t)\notin\mathrm{Fil}^2$; % t_unit, t_act, t_φ, mem_filtration_one, not_mem_filtration_two
\item every $b\in B$ with $\phi(b)=b$ and $\iota(b)\in\mathrm{Fil}^0$ is the image of an element of $F$. % fund
\end{enumerate}
\end{definition}
\leanref{Formalism/WeaklyAdmissibleFiltrations/Statement.lean}{Formalism.WeaklyAdmissibleFiltrations.PeriodData}

\begin{definition}\label{gal-cover-period-system}
Let $\ell$ be a prime, $\bar{\mathbf Q}_\ell$ an algebraic closure of $\mathbf Q_\ell$ with the absolute value
normalized by $|\ell|=\ell^{-1}$, $G_{\mathbf Q_\ell}=\operatorname{Gal}(\bar{\mathbf Q}_\ell/\mathbf Q_\ell)$, and
$\mathbf Q_\ell^{\mathrm{ur}}\subset\bar{\mathbf Q}_\ell$ the extension of $\mathbf Q_\ell$ generated by the roots of
unity of order prime to $\ell$, with its arithmetic Frobenius $\operatorname{Frob}$. For a finite extension
$K\subset\bar{\mathbf Q}_\ell$ of $\mathbf Q_\ell$ put $K_0=K\cap\mathbf Q_\ell^{\mathrm{ur}}$ and
$G_K=\operatorname{Gal}(\bar{\mathbf Q}_\ell/K)$, and let $\mathcal O_K$ be the integral closure of $\mathbf Z_\ell$ in
$K$. A \emph{period ring system} $P$ for $\ell$ consists of the following data.
\begin{enumerate}
\item A commutative $\mathbf Q_\ell^{\mathrm{ur}}$-algebra $B_P$ and a commutative $\bar{\mathbf Q}_\ell$-algebra
  $B_{\mathrm{dR},P}$. For every finite extension $K\subset\bar{\mathbf Q}_\ell$ of $\mathbf Q_\ell$ they are a
  $K_0$-algebra and a $K$-algebra through $K_0\subset\mathbf Q_\ell^{\mathrm{ur}}$ and $K\subset\bar{\mathbf Q}_\ell$.
  % LocRings: B, BdR, rB, rdR, aQB, aUB, tUB, aOdR, aQdR, tQdR
\item For every finite extension $K\subset\bar{\mathbf Q}_\ell$ of $\mathbf Q_\ell$, a Fontaine datum $P_K$ for
  $(G_K,\mathbf Q_\ell,K_0,K,B_P,B_{\mathrm{dR},P})$ (Definition~\ref{gal-cover-fontaine-datum}). % PD
\item Actions of $G_{\mathbf Q_\ell}$ on $B_P$ and on $B_{\mathrm{dR},P}$ by $\mathbf Q_\ell$-algebra automorphisms,
  written $b\mapsto g(b)$, $\mathbf Q_\ell$-algebra maps $\phi_P:B_P\to B_P$ and $\iota_P:B_P\to B_{\mathrm{dR},P}$, an
  element $t_P\in B_P$ and a homomorphism $\chi_P:G_{\mathbf Q_\ell}\to\mathbf Q_\ell^\times$. % actQ, actQdR, φB, ιB, tB, χQ
\item Elements $\zeta_m\in\bar{\mathbf Q}_\ell$ for $m\geq0$, and for every finite extension $L\subset\bar{\mathbf
  Q}_\ell$ of $\mathbf Q_\ell$ a valuation $v_L:L\to\mathbf Q\cup\{\infty\}$. % ζQ, v
\end{enumerate}
These are subject to the following conditions.
\begin{enumerate}
\item[(a)] For every $K$, the prime of $P_K$ is $\ell$; the actions of $G_K$ in $P_K$ are the restrictions of the
  actions of (3); the map $\phi$, the map $\iota$ and the element $t$ of $P_K$ are $\phi_P$, $\iota_P$ and $t_P$; the
  homomorphism $\chi$ of $P_K$ is the restriction of $\chi_P$; and the subsets $\mathrm{Fil}^a$ of $B_{\mathrm{dR},P}$
  given by $P_K$ do not depend on $K$. % PD_p, PD_actB, PD_actdR, PD_φ, PD_ι, period_datum_t_eq, PD_χ, PD_Fil
\item[(b)] For $g\in G_{\mathbf Q_\ell}$, $c\in\mathbf Q_\ell^{\mathrm{ur}}$ and $x\in\bar{\mathbf Q}_\ell$, the action
  of $g$ sends the image of $c$ in $B_P$ to the image of $g(c)$, and the image of $x$ in $B_{\mathrm{dR},P}$ to the
  image of $g(x)$. % actQ_Kur, actQdR_Om
\item[(c)] For every $K$ and every $x\in K_0$ with $|x|\leq1$, the map $\sigma$ of $P_K$ satisfies
  $|\sigma(x)-x^\ell|<1$. % PD_σ
\item[(d)] For $g\in G_{\mathbf Q_\ell}$, $m\geq0$ and $\zeta\in\bar{\mathbf Q}_\ell$ with $\zeta^{\ell^m}=1$, there is
  an integer $a\geq0$ with $|\chi_P(g)-a|\leq\ell^{-m}$ and $g(\zeta)=\zeta^a$. % χQ_cyc
\item[(e)] $\zeta_0=1$, $\zeta_{m+1}^\ell=\zeta_m$ for every $m$, and $\zeta_1\neq1$. % ζQ_zero, ζQ_succ, ζQ_one
\item[(f)] For every $L$, $v_L(\ell)=1$, and $v_L(x)\geq0$ exactly when $x\in\mathcal O_L$. % v_l, v_O
\end{enumerate}
\end{definition}
\leanref{Formalism/SharpRestriction/LevelLetterCocycles/DataNLoc.lean}{Formalism.BasicNotions.LocRings}
\leanref{Formalism/SharpRestriction/LevelLetterCocycles/DataNLoc.lean}{Formalism.BasicNotions.LocPrime}
\leanref{Formalism/SharpRestriction/LevelLetterCocycles/DataNLoc.lean}{Formalism.BasicNotions.Loc.MaximalUnramifiedExtension}
\leanref{Formalism/SharpRestriction/LevelLetterCocycles/DataNLoc.lean}{Formalism.BasicNotions.Loc.K0}
\leanref{Development/AbelianVarietyTheory/Definitions/AutomorphicMembership/EqOfPowEqOneOfNormSubLt.lean}{GlobalTheory.Pst.frobUr}

\begin{definition}\label{gal-cover-period-system-ainf}
Let $\ell$, $\bar{\mathbf Q}_\ell$, $\mathbf C_\ell$, $\mathcal O^\flat$, $A_{\inf}=W(\mathcal O^\flat)$ and $\phi$ be as
in Definition~\ref{gal-cover-breuil-kisin}, and $G_{\mathbf Q_\ell}$ and $\mathbf Q_\ell^{\mathrm{ur}}$ as in
Definition~\ref{gal-cover-period-system}. View $\mathcal O^\flat$ as $\varprojlim_{x\mapsto x^\ell}\mathcal O_{\mathbf
C_\ell}/\ell$, and write $x_n\in\mathcal O_{\mathbf C_\ell}/\ell$ for the $n$-th coordinate of $x\in\mathcal O^\flat$.
Let $P$ be a period ring system for $\ell$ (Definition~\ref{gal-cover-period-system}). An \emph{$A_{\inf}$-structure} $A$
on $P$ consists of:
\begin{enumerate}
\item an action of $G_{\mathbf Q_\ell}$ on $\mathbf C_\ell$ by ring automorphisms that preserve the absolute value and
  extend its action on $\bar{\mathbf Q}_\ell$; % actC, actC_ext, actC_norm
\item an action of $G_{\mathbf Q_\ell}$ on $\mathcal O^\flat$ by ring automorphisms such that, for $g\in G_{\mathbf
  Q_\ell}$, $x\in\mathcal O^\flat$, $n\geq0$ and $y\in\mathcal O_{\mathbf C_\ell}$ with $x_n$ the class of $y$, $g(x)_n$
  is the class of $g(y)$; write $W(g)$ for the automorphism of $A_{\inf}$ it induces; % actT, actT_coeff
\item elements $\zeta'_n\in\mathcal O_{\mathbf C_\ell}$ for $n\geq0$ with $\zeta'_0=1$, $\zeta'^{\,\ell}_{n+1}=\zeta'_n$,
  $\zeta'_1\neq1$ and $\zeta'_n=\zeta_n$ (the elements of $P$) for every $n$, and an element $\varepsilon\in\mathcal
  O^\flat$ with $\varepsilon_n$ the class of $\zeta'_n$ for every $n$; % ζ, ζ_zero, ζ_succ, ζ_one, ζ_eq, ε, ε_coeff
\item an injective ring map $\iota_A:A_{\inf}\to B_P$ with $\iota_A(W(g)(a))=g(\iota_A(a))$ and
  $\iota_A(\phi(a))=\phi_P(\iota_A(a))$ for $g\in G_{\mathbf Q_\ell}$ and $a\in A_{\inf}$; % ιA, ιA_inj, ιA_act, ιA_φ
\end{enumerate}
such that, with $\xi=\sum_{i=0}^{\ell-1}[\varepsilon^{1/\ell}]^i\in A_{\inf}$:
\begin{enumerate}
\item[(a)] for every finite extension $K\subset\bar{\mathbf Q}_\ell$ of $\mathbf Q_\ell$ and every $a\in A_{\inf}$,
  $\iota_P(\iota_A(a))$ lies in the subspace $\mathrm{Fil}^0$ of $P_K$, and it lies in $\mathrm{Fil}^1$ exactly when
  $a\in\xi A_{\inf}$; % ainf_image_mem_filtration_zero, ainf_image_mem_filtration_one_iff
\item[(b)] if $c\in\mathbf Q_\ell^{\mathrm{ur}}$ satisfies $c^m=1$ for an integer $m\geq1$ prime to $\ell$, and
  $x\in\mathcal O^\flat$ is such that for every $k\geq0$, $x_k$ is the class of some $y\in\mathcal O_{\mathbf C_\ell}$
  with $y^{\ell^k}=c$ and $y^m=1$, then $\iota_A([x])$ is the image of $c$ in $B_P$. % ιA_teich
\end{enumerate}
\end{definition}
\leanref{Formalism/SharpRestriction/LevelLetterCocycles/DataNLoc.lean}{Formalism.BasicNotions.LocAinf}

\begin{definition}\label{gal-cover-period-system-identification}
Let $\ell$, $\mathbf C_\ell$, $A_{\inf}$, $\theta$, $A_{\mathrm{cris}}\subset B_{\mathrm{cris}}$ and $\phi$ be as in
Definition~\ref{gal-cover-breuil-kisin}, with the actions of $G_{\mathbf Q_\ell}=\operatorname{Gal}(\bar{\mathbf
Q}_\ell/\mathbf Q_\ell)$ on these rings. Let $B_{\mathrm{dR}}^+\supset A_{\mathrm{cris}}$ be Fontaine's ring of de Rham
periods, with $\theta:B_{\mathrm{dR}}^+\to\mathbf C_\ell$, its field of fractions $B_{\mathrm{dR}}$ with the action of
$G_{\mathbf Q_\ell}$, and the filtration $\mathrm{Fil}^aB_{\mathrm{dR}}=(\ker\theta)^aB_{\mathrm{dR}}^+$, $a\in\mathbf Z$
\cite{Fon94a}. Let $P$ be a period ring system for $\ell$ with an $A_{\inf}$-structure $A$
(Definitions~\ref{gal-cover-period-system} and~\ref{gal-cover-period-system-ainf}), and $\varepsilon$ the element of $A$.
A \emph{Fontaine identification} $T$ of $(P,A)$ consists of ring isomorphisms $e_B:B_P\to B_{\mathrm{cris}}$ and
$e_D:B_{\mathrm{dR},P}\to B_{\mathrm{dR}}$ such that:
\begin{enumerate}
\item $e_B(z)=z$ and $e_D(z)=z$ for $z\in\mathbf Z_\ell$; % zB, zD
\item $e_D(x)\in B_{\mathrm{dR}}^+$ and $\theta(e_D(x))=x$ for $x\in\bar{\mathbf Q}_\ell$; % om
\item $e_B(g(b))=g(e_B(b))$, $e_D(g(b'))=g(e_D(b'))$ and $e_B(\phi_P(b))=\phi(e_B(b))$ for $g\in G_{\mathbf Q_\ell}$,
  $b\in B_P$ and $b'\in B_{\mathrm{dR},P}$; % actB, actD, frob
\item $e_D(\iota_P(e_B^{-1}(a)))=a$ for $a\in A_{\mathrm{cris}}$, viewed in $B_{\mathrm{dR}}^+\subset B_{\mathrm{dR}}$; % incl
\item for every finite extension $K\subset\bar{\mathbf Q}_\ell$ of $\mathbf Q_\ell$ and $a\in\mathbf Z$, an element
  $b'\in B_{\mathrm{dR},P}$ lies in the subspace $\mathrm{Fil}^a$ of $P_K$ exactly when $e_D(b')\in\mathrm{Fil}^a
  B_{\mathrm{dR}}$; % fil
\item $e_D(\iota_P(t_P))$ lies in $B_{\mathrm{dR}}^+$ and is congruent, for every $n\geq1$, to
  $\sum_{i=1}^{n-1}(-1)^{i+1}([\varepsilon]-1)^i/i$ modulo $(\ker\theta)^n$; % t (Per.IsLog)
\item $e_B(\iota_A(a))$ is the image of $a$ in $B_{\mathrm{cris}}$ for $a\in A_{\inf}$. % ιA
\end{enumerate}
The \emph{embedding} of $T$ is $u_T:\mathbf Q_\ell^{\mathrm{ur}}\to B_{\mathrm{cris}}$, $u_T(c)=e_B(c)$, where $c$ also
denotes the image of $c$ in $B_P$.
\end{definition}
\leanref{Formalism/LocalRingCardinality/AdicCompletionLocal/LocTied.lean}{Formalism.BasicNotions.LocTied}
\leanref{Development/AutomorphicArithmetic/Foundations/GaloisRepresentations/ClosedPhiFunctoriality/JointApproxOfMem.lean}{UpperTheory.Gal.ClosedPhiFromClosedTrace.uB}
\leanref{Development/WeilBound/CoreStatements/CrystallinePeriodRings/InstCommRingBcris.lean}{Formalism.Per6.Bcris}
\leanref{Development/WeilBound/CoreStatements/CrystallinePeriodRings/InstCommRingBcris.lean}{Formalism.Per6.FilD}
\leanref{Development/WeilBound/CoreStatements/CrystallinePeriodRings/InstCommRingBcris.lean}{Formalism.Per6.IsCrisDR}
\leanref{Development/WeilBound/CoreStatements/CrystallinePeriodRings/InstCommRingBcris.lean}{Formalism.Per6.IsLog}

\begin{lemma}\label{gal-cover-period-system-frobenius}
Let $P$ be a period ring system for a prime $\ell$ (Definition~\ref{gal-cover-period-system}). Then
$\phi_P(c)=\operatorname{Frob}(c)$ in $B_P$ for every $c\in\mathbf Q_\ell^{\mathrm{ur}}$.
\end{lemma}
\leanref{Development/AbelianVarietyTheory/Definitions/AutomorphicMembership/EqOfPowEqOneOfNormSubLt.lean}{GlobalTheory.Pst.phiB_algebraMap}

\begin{proof}
The maps $c\mapsto\phi_P(c)$ and $c\mapsto\operatorname{Frob}(c)$ are $\mathbf Q_\ell$-algebra maps from $\mathbf
Q_\ell^{\mathrm{ur}}$ to $B_P$, and $\mathbf Q_\ell^{\mathrm{ur}}$ is generated over $\mathbf Q_\ell$ by the roots of
unity of order prime to $\ell$. So it suffices to take $c$ with $c^m=1$ for some $m\geq1$ prime to $\ell$; then
$\operatorname{Frob}(c)=c^\ell$. Let $K=\mathbf Q_\ell(c)$, so that $c\in K_0$, and let $\sigma$ be the map of $P_K$. By
Definition~\ref{gal-cover-period-system}(a) and Definition~\ref{gal-cover-fontaine-datum}(3), $\phi_P(c)=\sigma(c)$.
Now $\sigma(c)^m=\sigma(c^m)=1$, and $|\sigma(c)-c^\ell|<1$ by Definition~\ref{gal-cover-period-system}(c), as $|c|=1$.
Two $m$-th roots of unity whose difference has absolute value less than $1$ are equal, since $x^m-1$ has $m$ distinct
roots in the residue field of $\bar{\mathbf Q}_\ell$. So $\sigma(c)=c^\ell$.
\end{proof}
\leanref{Development/AbelianVarietyTheory/Definitions/AutomorphicMembership/EqOfPowEqOneOfNormSubLt.lean}{GlobalTheory.Pst.kur_algHom_ext}
\leanref{Development/AbelianVarietyTheory/Definitions/AutomorphicMembership/EqOfPowEqOneOfNormSubLt.lean}{GlobalTheory.Pst.frobUr_root}
\leanref{Development/AbelianVarietyTheory/Definitions/AutomorphicMembership/EqOfPowEqOneOfNormSubLt.lean}{GlobalTheory.Pst.eq_of_pow_eq_one_of_norm_sub_lt}

\begin{lemma}\label{gal-cover-operators}
Let $\ell$ be a prime and $K\subset\bar{\mathbf Q}_\ell$ a finite extension of $\mathbf Q_\ell$ with a
Breuil--Kisin datum $(S,E,\iota)$ and a cyclotomic element $\tau$ (Definition~\ref{gal-cover-breuil-kisin}).
Let $n_1,n_2\geq0$ be such that $\operatorname{Frob}^{n_1}$ and $\operatorname{Frob}^{n_2}$ fix every
element of $\mathbf Q_\ell^{\mathrm{ur}}\cap K$. Let $\rho:G_K\to\operatorname{GL}_d(\mathbf Z_\ell)$ be a
homomorphism with a lattice tower $(\rho'_m,\Pi_m,f_m)_m$ (Definition~\ref{gal-cover-definition}), let
$N,r_0\geq0$, and let crystalline Breuil--Kisin structures of height $N$ on $\rho$ and on every
$\rho'_m$ be given such that for every $n$ there is $m\geq n$ for which the structure on $\rho$ and the
structure on $\rho'_m$ are linked at $(n,r_0)$ along $\Pi_m$. Then there are $s\in A_{\mathrm{cris}}$
that is a unit of $B_{\mathrm{cris}}$, $V\in M_d(A_{\mathrm{cris}})$ and
$V'_m\in M_{d'_m}(A_{\mathrm{cris}})$ for every $m$, such that:
\begin{enumerate}
\item $V\phi^{n_2}(U)=s\,\phi^{n_1}(U)$ for a period matrix $U$ of $\rho$, and for every $m$,
  $V'_m\phi^{n_2}(U'_m)=s\,\phi^{n_1}(U'_m)$ for a period matrix $U'_m$ of $\rho'_m$;
\item for every $M$ there is $m\geq M$ with $V\Pi_m\equiv\Pi_mV'_m$ modulo $\ell^M$ in
  $A_{\mathrm{cris}}$.
\end{enumerate}
% Lean form: OpEx2 also takes a period ring package (P, AI, LocTied), used only in FixW; this lemma uses Fontaine's rings
% of Definition gal-cover-breuil-kisin. Stating the package is a larger change, not made in the review of 21 November.
\end{lemma}
\leanref{Development/AutomorphicArithmetic/Foundations/GaloisRepresentations/ClosedPhiFunctoriality/IsOp2One.lean}{UpperTheory.Gal.ClosedPhiFromClosedTrace.OpEx2}
\leanref{Development/AutomorphicArithmetic/Foundations/GaloisRepresentations/ClosedPhiFunctoriality/OpExHolds.lean}{UpperTheory.Gal.ClosedPhiFromClosedTrace.opEx_holds}
\leanref{Development/AutomorphicArithmetic/Foundations/GaloisRepresentations/ClosedPhiFunctoriality/IntBasisAlt.lean}{UpperTheory.Gal.ClosedPhiFromClosedTrace.FixW}
\leanref{Development/AutomorphicArithmetic/Foundations/GaloisRepresentations/ClosedPhiFunctoriality/FixWHolds.lean}{UpperTheory.Gal.ClosedPhiFromClosedTrace.fixW_holds}
\leanref{Development/AutomorphicArithmetic/Foundations/GaloisRepresentations/ClosedPhiFunctoriality/IsOp2One.lean}{UpperTheory.Gal.ClosedPhiFromClosedTrace.LinkOp2}
\leanref{Development/AutomorphicArithmetic/Foundations/GaloisRepresentations/ClosedPhiFunctoriality/LinkOp.lean}{UpperTheory.Gal.ClosedPhiFromClosedTrace.linkOp_holds}

\begin{proof}
Let $k$ be given by Lemma~\ref{gal-cover-section} for $N$. Lemma~\ref{gal-cover-period}, applied with
$n_1,n_2$ to the structure on $\rho$ and to the structure on each $\rho'_m$, gives matrices $Z_0,Z_0'$,
$U$, $V$ for $\rho$ and $Z_m,Z_m'$, $U'_m$, $V'_m$ for $\rho'_m$, all with the same
$s=\phi^{n_2}(\tau_1)^{3N}\ell^{4k}\phi^{n_1}(\tau_1)^N$. This gives (1).

For (2), let $M$ be given, let $n$ be given by Lemma~\ref{gal-cover-link} for $N$ and $M$, and choose
$m\geq\max(n,M)$ for which the structures on $\rho$ and $\rho'_m$ are linked at $(\max(n,M),r_0)$ along
$\Pi_m$, with a matrix $R$. Lemma~\ref{gal-cover-link} gives $\Pi_mZ_m\equiv Z_0R^0$ modulo $\ell^M$. The
entries of $R^0$ are fixed by $G_K$, so in $B_{\mathrm{cris}}$ they lie in
$B_{\mathrm{cris}}^{G_K}=\mathbf Q_\ell^{\mathrm{ur}}\cap K$ \cite{Fon94a}, on which $\phi$ acts as
$\operatorname{Frob}$; hence $\phi^{n_1}$ and $\phi^{n_2}$ fix them. The entries of $\Pi_m$ are fixed by
$\phi$. Applying $\phi^{n_i}$ therefore gives $\Pi_m\phi^{n_i}(Z_m)\equiv\phi^{n_i}(Z_0)R^0$ modulo
$\ell^M$ for $i=1,2$. Put $c=\phi^{n_2}(\tau_1^{2N}\ell^{2k})$, so that
$\phi^{n_2}(Z_0')\phi^{n_2}(Z_0)=c$ and $\phi^{n_2}(Z_m)\phi^{n_2}(Z_m')=c$ by
Lemma~\ref{gal-cover-period}(1). Then, modulo $\ell^M$,
\[
V\Pi_m=\phi^{n_1}(Z_0)\phi^{n_2}(Z_0')\Pi_m\,\phi^{n_2}(Z_m)\phi^{n_2}(Z_m')
\equiv\phi^{n_1}(Z_0)\phi^{n_2}(Z_0')\phi^{n_2}(Z_0)R^0\phi^{n_2}(Z_m')
=c\,\phi^{n_1}(Z_0)R^0\phi^{n_2}(Z_m'),
\]
and $c\,\phi^{n_1}(Z_0)R^0\phi^{n_2}(Z_m')\equiv c\,\Pi_m\phi^{n_1}(Z_m)\phi^{n_2}(Z_m')=\Pi_mV'_m$.
\end{proof}

\begin{lemma}\label{gal-cover-trace-identity}
In Situation~\ref{gal-local-situation}, suppose that $w$ lies over a prime $\ell$, let
$f=[k_w:\mathbf F_\ell]$, let $\iota:\bar{\mathbf Q}_\ell\to\mathbf C$ be a field isomorphism, and
identify $D_A$ with $G_K$, for a finite extension $K\subset\bar{\mathbf Q}_\ell$ of $\mathbf Q_\ell$ in
which $E$ is dense, as in Definition~\ref{gal-p-adic-hodge}. Use the period rings of
Definition~\ref{gal-cover-breuil-kisin}. Let $\sigma\in W_A$ and $n_1,n_2\geq0$ with
$fd(\sigma)=n_2-n_1$. Let $r:D_A\to\operatorname{GL}_a(\bar{\mathbf Q}_\ell)$ be a homomorphism with
entries in a finite extension $L$ of $\mathbf Q_\ell$ whose representation on $L^a$ is crystalline, and
let $\mathrm{WD}_\iota(r)$ be its Weil--Deligne representation as in Definition~\ref{gal-p-adic-hodge}.
Let $\rho:G_K\to\operatorname{GL}_d(\mathbf Z_\ell)$ be a homomorphism that reads $r$ through $b$ and $e$
(Definition~\ref{gal-cover-definition}), and let $s\in B_{\mathrm{cris}}$ and
$V\in M_d(B_{\mathrm{cris}})$ satisfy $V\phi^{n_2}(U)=s\,\phi^{n_1}(U)$ for a period matrix $U$ of
$\rho$. Then for every $\mathbf Q_\ell^{\mathrm{ur}}$-linear map
$\lambda:\bar{\mathbf Q}_\ell\to\mathbf Q_\ell^{\mathrm{ur}}$,
\[
\sum_i(Ve)_i\,\lambda(\operatorname{tr}b_i)=s\,\lambda\bigl(\iota^{-1}\operatorname{tr}\mathrm{WD}_\iota(r)(\sigma)\bigr)
\quad\text{in }B_{\mathrm{cris}}.
\]
\end{lemma}
\leanref{Development/AutomorphicArithmetic/Foundations/GaloisRepresentations/ClosedPhiFunctoriality/IsOp2One.lean}{UpperTheory.Gal.ClosedPhiFromClosedTrace.Tie2}
\leanref{Development/AutomorphicArithmetic/Foundations/GaloisRepresentations/ClosedPhiFunctoriality/TraceEqNextCoeff.lean}{UpperTheory.Gal.ClosedPhiFromClosedTrace.tie_holds}
\leanref{Development/AutomorphicArithmetic/Foundations/GaloisRepresentations/ClosedPhiFunctoriality/IsOp2One.lean}{UpperTheory.Gal.ClosedPhiFromClosedTrace.TieRep2}
\leanref{Development/AutomorphicArithmetic/Foundations/GaloisRepresentations/ClosedPhiFunctoriality/TieRep.lean}{UpperTheory.Gal.ClosedPhiFromClosedTrace.tieRep_holds}
\leanref{Development/AutomorphicArithmetic/Foundations/GaloisRepresentations/ClosedPhiFunctoriality/IntBasisAlt.lean}{UpperTheory.Gal.ClosedPhiFromClosedTrace.TrLin}
\leanref{Development/AutomorphicArithmetic/Foundations/GaloisRepresentations/ClosedPhiFunctoriality/TrLinHolds.lean}{UpperTheory.Gal.ClosedPhiFromClosedTrace.trLin_holds}
\leanref{Development/AutomorphicArithmetic/Foundations/GaloisRepresentations/ClosedPhiFunctoriality/JointApproxOfMem.lean}{UpperTheory.Gal.ClosedPhiFromClosedTrace.IsTr}

\begin{proof}
Enlarge $L$ so that it contains the entries of the $b_i$, and let $V_L=L^a$ with the action of
$G_K$ through $r$. Since $\phi^{n_2}(U)$ is invertible, there is $y\in B_{\mathrm{cris}}^d$ with
$\phi^{n_2}(U)y=e$, and then $Ve=s\,\phi^{n_1}(U)y$.

\emph{Invariant vectors.} For $1\leq j\leq d$ and $1\leq k\leq a$ let
$m_{jk}=\sum_ib_i^{(k)}\otimes U_{ij}\in V_L\otimes_{\mathbf Q_\ell}B_{\mathrm{cris}}$, where
$b_i^{(k)}$ is the $k$-th column of $b_i$. Since $r(g)b_j=\sum_i\rho(g)_{ij}b_i$ and
$\rho(g)g(U)=U$, every $m_{jk}$ is fixed by $G_K$, so it lies in $D_{\mathrm{st}}(V_L)$, and
$\phi^n(m_{jk})=\sum_ib_i^{(k)}\otimes\phi^n(U_{ij})$. In Fontaine's construction of
$\mathrm{WD}_\iota(r)$ \cite{Fon94}, $\sigma$ acts on $D_{\mathrm{pst}}(V_L)$ as
$\sigma\circ\phi^{-fd(\sigma)}$ with $\sigma$ acting on both factors. As $\phi$ commutes with $G_K$ and
$fd(\sigma)=n_2-n_1$, it sends $\phi^{n_2}(m)$ to $\phi^{n_1}(m)$ for every $m\in D_{\mathrm{st}}(V_L)$.

\emph{Coordinates.} The map $V_L\otimes_{\mathbf Q_\ell}B_{\mathrm{st}}\to(\bar{\mathbf
Q}_\ell\otimes_{\mathbf Q_\ell^{\mathrm{ur}}}B_{\mathrm{st}})^a$, $v\otimes x\mapsto(v_i\otimes x)_i$,
induces an additive map $\psi$ on the space $D_{\mathrm{pst}}(V_L)\otimes_{L\otimes\mathbf
Q_\ell^{\mathrm{ur}}}\mathbf C$ of $\mathrm{WD}_\iota(r)$ that is semilinear for
$\kappa:\mathbf C\to\bar{\mathbf Q}_\ell\otimes_{\mathbf Q_\ell^{\mathrm{ur}}}B_{\mathrm{st}}$,
$z\mapsto\iota^{-1}(z)\otimes1$.

\emph{A trace formula.} Let $W$ be a vector space of dimension $a$ over a field $k_0$,
$\kappa:k_0\to B'$ a ring homomorphism to a commutative ring, $\psi:W\to B'^a$ additive and
$\kappa$-semilinear, $x_{j,i}\in W$ and $\omega_j\in B'$ ($j$ in a finite set, $1\leq i\leq a$) with
$\sum_j\omega_j\psi(x_{j,i'})_i=\delta_{ii'}$. Then
$\sum_i\sum_j\omega_j\psi(Tx_{j,i})_i=\kappa(\operatorname{tr}T)$ for every endomorphism $T$ of $W$.
Indeed, in a basis of $W$ let $\Psi\in M_a(B')$ have the images under $\psi$ of the basis vectors as
columns, and let $M\in M_a(B')$ have as $i'$-th column $\sum_j\omega_j\kappa(\xi_{j,i'})$, where
$\xi_{j,i'}$ is the coordinate vector of $x_{j,i'}$. The hypothesis says $\Psi M=1$, so $M\Psi=1$, and
the sum is $\operatorname{tr}(\Psi\kappa(T)M)=\operatorname{tr}(\kappa(T)M\Psi)=\kappa(\operatorname{tr}T)$.

\emph{Conclusion.} Apply the trace formula to the space of $\mathrm{WD}_\iota(r)$,
$T=\mathrm{WD}_\iota(r)(\sigma)$, the images $x_{j,i}$ of the $\phi^{n_2}(m_{ji})$ and the weights
$\omega_j=1\otimes y_j$. The hypothesis holds:
$\sum_j\omega_j\psi(x_{j,i'})_i=\sum_{i_2}(b_{i_2})_{ii'}\otimes(\phi^{n_2}(U)y)_{i_2}
=\bigl(\sum_{i_2}e_{i_2}b_{i_2}\bigr)_{ii'}\otimes1=\delta_{ii'}$. Since $Tx_{j,i}$ is the image of
$\phi^{n_1}(m_{ji})$, the formula gives
\[
\iota^{-1}\bigl(\operatorname{tr}\mathrm{WD}_\iota(r)(\sigma)\bigr)\otimes1
=\sum_{i_2}\operatorname{tr}(b_{i_2})\otimes(\phi^{n_1}(U)y)_{i_2}
\quad\text{in }\bar{\mathbf Q}_\ell\otimes_{\mathbf Q_\ell^{\mathrm{ur}}}B_{\mathrm{st}}.
\]
Apply $\lambda\otimes1$ and multiply by $s$. Since $s\,\phi^{n_1}(U)y=Ve$, this is the asserted identity
in $B_{\mathrm{st}}$, between elements of $B_{\mathrm{cris}}$.
\end{proof}

\begin{lemma}\label{gal-cover-functional}
Let $\ell$ be a prime, $F$ a field with $\mathbf Q_\ell\subset F\subset\bar{\mathbf Q}_\ell$, and
$x_1,\ldots,x_n,z\in\bar{\mathbf Q}_\ell$ with $|x_i|\leq1$ and $|z|\leq1$. Then there is an $F$-linear
map $\lambda:\bar{\mathbf Q}_\ell\to F$ with $|\lambda(x_i)|\leq1$ for all $i$ and
$|z|\leq|\lambda(z)|\leq1$.
\end{lemma}
\leanref{Development/AutomorphicArithmetic/Foundations/GaloisRepresentations/ClosedPhiFunctoriality/IntBasisAlt.lean}{UpperTheory.Gal.ClosedPhiFromClosedTrace.Func}
\leanref{Development/AutomorphicArithmetic/Foundations/GaloisRepresentations/ClosedPhiFunctoriality/FuncHolds.lean}{UpperTheory.Gal.ClosedPhiFromClosedTrace.func_holds}
\leanref{Development/AutomorphicArithmetic/Foundations/GaloisRepresentations/ClosedPhiFunctoriality/IntBasisAlt.lean}{UpperTheory.Gal.ClosedPhiFromClosedTrace.IntBasis}
\leanref{Development/AutomorphicArithmetic/Foundations/GaloisRepresentations/ClosedPhiFunctoriality/IntBasis.lean}{UpperTheory.Gal.ClosedPhiFromClosedTrace.intBasis_holds}

\begin{proof}
Let $\mathcal O_F$ be the valuation ring of $F$. If $z=0$, take $\lambda=0$; so let $z\neq0$.

\emph{A family.} For every finite family $u_1,\ldots,u_k\in\bar{\mathbf Q}_\ell$ with $|u_i|\leq1$ there
are $e_1,\ldots,e_m\in\bar{\mathbf Q}_\ell$, linearly independent over $F$ with $|e_j|\leq1$, such that
each $u_i$ is a combination of them with coefficients in $\mathcal O_F$. We argue by induction on $k$,
the case $k=0$ being the empty family. Given $e_1,\ldots,e_m$ for $u_1,\ldots,u_{k-1}$, let $u=u_k$. If
$u$ is not in the $F$-span of the $e_j$, add it to the family. If $u=\sum_ja_je_j$ with every
$|a_j|\leq1$, keep the family. Otherwise choose $j_0$ with $|a_{j_0}|$ maximal, so $|a_{j_0}|>1$, and
replace $e_{j_0}$ by $u$. The new family is linearly independent over $F$ since $a_{j_0}\neq0$, and $u$
is one of its members. For $i<k$, if $u_i=\sum_jc_je_j$ with $c_j\in\mathcal O_F$, then
$e_{j_0}=a_{j_0}^{-1}(u-\sum_{j\neq j_0}a_je_j)$ gives
$u_i=(c_{j_0}/a_{j_0})u+\sum_{j\neq j_0}(c_j-c_{j_0}a_j/a_{j_0})e_j$, and these coefficients lie in
$\mathcal O_F$ by the ultrametric inequality, since $|a_j|\leq|a_{j_0}|$ and $|a_{j_0}|>1$.

\emph{The map.} Apply this to $z$ and the $x_i$, and write $z=\sum_jc_je_j$; then $m\geq1$ since
$z\neq0$. Choose $j_0$ with $|c_{j_0}|$ maximal. Extend the coordinate
$\sum_jy_je_j\mapsto y_{j_0}$ from $Fe_1\oplus\cdots\oplus Fe_m$ to an $F$-linear map $\lambda$ on
$\bar{\mathbf Q}_\ell$. Then $\lambda(x_i)\in\mathcal O_F$, $\lambda(z)=c_{j_0}\in\mathcal O_F$, and
$|z|\leq\max_j|c_j|=|\lambda(z)|$ by the ultrametric inequality.
\end{proof}

\begin{remark}\label{gal-cover-functional-free}
The family in the proof of Lemma~\ref{gal-cover-functional} can also be taken to be a basis of the
$\mathcal O_F$-submodule of $\bar{\mathbf Q}_\ell$ generated by $z$ and the $x_i$. This module is finitely
generated and torsion-free over the valuation ring $\mathcal O_F$, hence free, and a basis consists of
elements of absolute value at most $1$ that are linearly independent over $F$.
\end{remark}

\begin{lemma}\label{gal-cover-trace-limit}
Let $\ell$ be a prime, $\Gamma$ a topological group, $r:\Gamma\to\operatorname{GL}_a(\bar{\mathbf Q}_\ell)$
a homomorphism and $C$ a set of homomorphisms $\Gamma\to\operatorname{GL}_a(\bar{\mathbf Q}_\ell)$, and fix
a cover of $r$ relative to $C$ and a finite extension $L_0$ (Definition~\ref{gal-cover-definition}), with
$\rho$, $b$, $e$, lattice tower $(\rho'_m,\Pi_m,f_m)_m$ and layers
$(\rho_{m,k},L_{m,k},b'_{m,j},e'_m,c_{m,k})$. Use the period rings of
Definition~\ref{gal-cover-breuil-kisin}. Let $s\in A_{\mathrm{cris}}$ be a unit of $B_{\mathrm{cris}}$,
$V\in M_d(A_{\mathrm{cris}})$, $V'_m\in M_{d'_m}(A_{\mathrm{cris}})$ for every $m$, and
$p,\theta,p_{m,k}\in\bar{\mathbf Q}_\ell$. Suppose that for every $\mathbf Q_\ell^{\mathrm{ur}}$-linear
$\lambda:\bar{\mathbf Q}_\ell\to\mathbf Q_\ell^{\mathrm{ur}}$
\[
\sum_i(Ve)_i\,\lambda(\operatorname{tr}b_i)=s\lambda(p),\qquad
\sum_j(V'_me'_m)_j\,\lambda(\operatorname{tr}b'_{m,j,k})=s\lambda(p_{m,k})\ \text{for all }m,k,
\]
that for every $M$ there is $m\geq M$ with $V\Pi_m\equiv\Pi_mV'_m$ modulo $\ell^M$ in
$A_{\mathrm{cris}}$, and that $|\theta-\sum_kc_{m,k}p_{m,k}|\leq\ell^{-m}$ for every $m$. Then $p=\theta$.
% Lean form: TrLim also carries a period ring package and a Breuil-Kisin datum for some K, used only for the facts that
% A_cris has no l-torsion, is l-adically separated and injects into B_cris; this lemma uses Fontaine's rings and leaves
% the datum out. Stating the package is a larger change, not made in the review of 21 November.
\end{lemma}
\leanref{Development/AutomorphicArithmetic/Foundations/GaloisRepresentations/ClosedPhiFunctoriality/IntBasisAlt.lean}{UpperTheory.Gal.ClosedPhiFromClosedTrace.TrLim}
\leanref{Development/AutomorphicArithmetic/Foundations/GaloisRepresentations/ClosedPhiFunctoriality/TrLimHolds.lean}{UpperTheory.Gal.ClosedPhiFromClosedTrace.trLim_holds}
\leanref{Development/AutomorphicArithmetic/Foundations/GaloisRepresentations/ClosedPhiFunctoriality/IntBasisAlt.lean}{UpperTheory.Gal.ClosedPhiFromClosedTrace.TrStep}
\leanref{Development/AutomorphicArithmetic/Foundations/GaloisRepresentations/ClosedPhiFunctoriality/TrStep.lean}{UpperTheory.Gal.ClosedPhiFromClosedTrace.trStep_holds}
\leanref{Development/AutomorphicArithmetic/Foundations/GaloisRepresentations/ClosedPhiFunctoriality/IntBasisAlt.lean}{UpperTheory.Gal.ClosedPhiFromClosedTrace.TrId}
\leanref{Development/AutomorphicArithmetic/Foundations/GaloisRepresentations/ClosedPhiFunctoriality/TrIdHolds.lean}{UpperTheory.Gal.ClosedPhiFromClosedTrace.trId_holds}

\begin{proof}
Suppose $p\neq\theta$, and write $t_i=\operatorname{tr}b_i$. Choose $\nu_1$ with
$|\ell^{\nu_1}(p-\theta)|\leq1$ and $|\ell^{\nu_1}t_i|\leq1$ for all $i$, and $\nu_2$ with
$|\ell^{\nu_1}(p-\theta)|\geq\ell^{-\nu_2}$. We show that $\ell^{\nu_2}s\in\ell^MA_{\mathrm{cris}}$ for
every $M$. Since $A_{\mathrm{cris}}$ is $\ell$-adically separated and has no $\ell$-torsion, this forces
$s=0$, which is not a unit of $B_{\mathrm{cris}}$.

Let $M$ be given and choose $m\geq M$ with $V\Pi_m\equiv\Pi_mV'_m$ modulo $\ell^M$; write $\Pi=\Pi_m$,
$e'=e'_m$, $V'=V'_m$, $c_k=c_{m,k}$ and $t'_{j,k}=\operatorname{tr}b'_{m,j,k}$. Put
\[
\delta_j=\sum_i\Pi_{ij}t_i-\sum_kc_kt'_{j,k},\qquad\eta=\theta-\sum_kc_kp_{m,k}.
\]
By the trace estimate of the cover and the last hypothesis, $|\delta_j|\leq\ell^{-M}$ and
$|\eta|\leq\ell^{-M}$. Lemma~\ref{gal-cover-functional}, with $F=\mathbf Q_\ell^{\mathrm{ur}}$, gives a
$\mathbf Q_\ell^{\mathrm{ur}}$-linear $\lambda$ with $\lambda(\ell^{\nu_1}t_i)$,
$\lambda(\ell^{-M}\delta_j)$ and $\lambda(\ell^{-M}\eta)$ integral and $\ell^{-\nu_2}\leq|q|\leq1$ for
$q=\lambda(\ell^{\nu_1}(p-\theta))$.

Since $\lambda$ is $\mathbf Z_\ell$-linear, the hypotheses for $p$ and for the $p_{m,k}$ (the latter
applied to $x\mapsto\lambda(c_kx)$) give the identity
\[
s\lambda(p-\theta)=\sum_i\bigl((V\Pi-\Pi V')e'-V(\Pi e'-e)\bigr)_i\lambda(t_i)
+\sum_j(V'e')_j\lambda(\delta_j)-s\lambda(\eta).
\]
Indeed, $s\lambda(p)=\sum_i(Ve)_i\lambda(t_i)$, the first sum on the right is this minus
$\sum_i(\Pi V'e')_i\lambda(t_i)=\sum_j(V'e')_j\lambda(\sum_i\Pi_{ij}t_i)$, and
$\sum_j(V'e')_j\lambda(\sum_kc_kt'_{j,k})=s\sum_k\lambda(c_kp_{m,k})=s\lambda(\theta)-s\lambda(\eta)$.
Multiply by $\ell^{\nu_1}$. Integral elements of $\mathbf Q_\ell^{\mathrm{ur}}$ lie in $A_{\mathrm{cris}}$,
$V\Pi-\Pi V'$ and $\Pi e'-e$ are divisible by $\ell^M$, and $\lambda(\delta_j)$, $\lambda(\eta)$ are
$\ell^M$ times integral elements. So $sq\in\ell^MA_{\mathrm{cris}}$, and since $\ell^{\nu_2}/q$ is
integral, $\ell^{\nu_2}s\in\ell^MA_{\mathrm{cris}}$, using that $A_{\mathrm{cris}}\to B_{\mathrm{cris}}$
is injective.
\end{proof}

\begin{lemma}\label{gal-closed-crystalline}
Let $E$ be a number field, $\Omega$ an algebraic closure, $\ell$ a prime, $w$ a place over $\ell$ and
$A$ a valuation ring over $w$. Let $r:G_E\to\operatorname{GL}_a(\bar{\mathbf Q}_\ell)$ be continuous and
semisimple, and $C$ a set of continuous representations of $G_E$ of dimension $a$ whose traces
approximate that of $r$ uniformly by linear combinations.
\begin{enumerate}
\item If every $\rho|_{D_A}$, $\rho\in C$, is crystalline with Hodge--Tate multiset $J_\tau$ at every
  $\tau$, then $r|_{D_A}$ is crystalline with Hodge--Tate multiset $J_\tau$ at every $\tau$.
\item In the setting of (1), let $\sigma\in W_A$ and $S_n\in\bar{\mathbf Q}_\ell$ be such that for every
  $n_0$ and $\varepsilon$ a single linear combination approximates both the trace of $r$ within
  $\varepsilon$ and each $S_n$, $n\leq n_0$, by the same combination of the power sums
  $\operatorname{tr}\mathrm{WD}_\iota(\rho|_{D_A})(\sigma)^n$. Then
  $\iota^{-1}\operatorname{tr}\mathrm{WD}_\iota(r|_{D_A})(\sigma)^n=S_n$ for all $n$.
\end{enumerate}
\end{lemma}
\leanref{Development/AbelianVarietyTheory/Proofs/SatakeParameters/ClosedAtL.lean}{GlobalTheory.LimitTop.closedAtL}
\leanref{Development/AbelianVarietyTheory/Proofs/InertiaEigenvalues/CoverExists.lean}{GlobalTheory.TraceConjugationCover.cover_exists}
\leanref{Development/AbelianVarietyTheory/Definitions/NonGaloisDescent/JointApproxOfMem.lean}{GlobalTheory.CongPhi.ClosedPhi}
\leanref{Development/AbelianVarietyTheory/Definitions/NonGaloisDescent/JointApproxOfMem.lean}{GlobalTheory.CongPhi.JointApprox}
\leanref{Development/AutomorphicArithmetic/Foundations/GaloisRepresentations/ClosedPhiFunctoriality/ClosedPhiHolds.lean}{UpperTheory.Gal.ClosedPhiFromClosedTrace.congPhi_holds}

\begin{proof}
For (1), the approximation of traces gives a \emph{cover} of $r$: a continuous representation
$\rho:G_E\to\operatorname{GL}_d(\mathbf Z_\ell)$ whose image spans a lattice in the algebra spanned by
the image of $r$, from which $r$ is read off by fixed matrices, and a tower of lattices in it whose
$n$-th layer is a continuous representation in the algebra spanned by finitely many members of $C$, over
a finite extension of $\mathbf Q_\ell$, with traces within $\ell^{-n}$ of those of $r$. The cover
restricts to $D_A$. Each layer is a lattice in the matrix algebra spanned by the images of finitely
many members of $C$, on which $D_A$ acts through them, so it is semistable at $A$, crystalline if the
members of $C$ are, with Hodge--Tate weights in the finite set of the $J_\tau$. A
representation of $D_A$ covered by such layers is de Rham with Hodge--Tate multisets $J_\tau$, and
crystalline if the layers are. This is the closedness of crystalline representations with bounded
Hodge--Tate weights under $\ell$-adic limits, proved by Berger \cite{Ber04}; Kisin proves it in the
form of crystalline deformation rings \cite{Kis08}. Realizing $r$ over a finite extension of
$\mathbf Q_\ell$ that contains the entries of its image gives (1).

For (2), fix $n$ and put $g=\sigma^n$. For a representation $x$ of $D_A$ with a Weil--Deligne
representation, $\mathrm{WD}_\iota(x)(g)=\mathrm{WD}_\iota(x)(\sigma)^n$, so the $n$-th power sum of the
eigenvalues of $\mathrm{WD}_\iota(x)(\sigma)$ is $\operatorname{tr}\mathrm{WD}_\iota(x)(g)$. Put
$p(\rho)=\iota^{-1}\operatorname{tr}\mathrm{WD}_\iota(\rho|_{D_A})(g)$ for $\rho\in C$ and
$p_r=\iota^{-1}\operatorname{tr}\mathrm{WD}_\iota(r|_{D_A})(g)$. The hypothesis with $n_0=n$ says that
for every $\varepsilon>0$ there are $\rho_1,\ldots,\rho_k\in C$ and $c_1,\ldots,c_k$ with
$|\operatorname{tr}r(h)-\sum_ic_i\operatorname{tr}\rho_i(h)|\leq\varepsilon$ for all $h\in G_E$ and
$|S_n-\sum_ic_ip(\rho_i)|\leq\varepsilon$, and it suffices to show $p_r=S_n$.

By (1), $r|_{D_A}$ is crystalline. Identify $D_A$ with $G_K$ as in Definition~\ref{gal-p-adic-hodge},
let $f=[k_w:\mathbf F_\ell]$, and let $L_0\subset\bar{\mathbf Q}_\ell$ be a finite extension of
$\mathbf Q_\ell$ containing the images of all embeddings of $K$ into $\bar{\mathbf Q}_\ell$.
Lemma~\ref{gal-cover-joint}, with $\Gamma=G_E$ and $\theta=S_n$, gives a cover of $r$ relative to $C$
and $L_0$ (Definition~\ref{gal-cover-definition}), with representation $\rho_0$ reading $r$ through $b$
and $e$, lattice tower $(\rho'_m,\Pi_m,f_m)_m$ and layers $(\rho_{m,k},L_{m,k},b'_{m,j},e'_m,c_{m,k})$,
such that $|S_n-\sum_kc_{m,k}p(\rho_{m,k})|\leq\ell^{-m}$ for every $m$. Restricting its
representations to $G_K$ gives a cover of $r|_{G_K}$ relative to the restrictions of the members of
$C$. Lemma~\ref{gal-cover-structures}(1) gives a Breuil--Kisin datum for $K$ with a cyclotomic element
(Definition~\ref{gal-cover-breuil-kisin}), and Lemma~\ref{gal-cover-structures}(2), with $J$ sending an
embedding of $K$ to the multiset $J_\tau$ of its restriction $\tau$ to $E$, gives $N$, $r_0$ and
crystalline Breuil--Kisin structures of height $N$ on $\rho_0|_{G_K}$ and on every $\rho'_m|_{G_K}$,
linked along the $\Pi_m$ as there. Put
$n_2=\max(fd(g),0)$ and $n_1=\max(-fd(g),0)$, so that $fd(g)=n_2-n_1$; by
Lemma~\ref{gal-cover-weil-frobenius}, $\operatorname{Frob}^{n_1}$ and $\operatorname{Frob}^{n_2}$ fix
$\mathbf Q_\ell^{\mathrm{ur}}\cap K$. Lemma~\ref{gal-cover-operators} gives $s$, $V$ and the $V'_m$,
with period matrices $U$ of $\rho_0|_{G_K}$ and $U'_m$ of $\rho'_m|_{G_K}$.
Lemma~\ref{gal-cover-trace-identity}, applied with $U$ to $r|_{D_A}$, which is crystalline and read by
$\rho_0$ through $b$ and $e$, and with $U'_m$ to every $\rho_{m,k}|_{D_A}$, which is crystalline and read
by $\rho'_m$ through the $b'_{m,j,k}$ and $e'_m$, gives the two families of identities of
Lemma~\ref{gal-cover-trace-limit} with $p=p_r$ and $p_{m,k}=p(\rho_{m,k})$. The congruences of
Lemma~\ref{gal-cover-operators}(2) and the bound for $S_n$ are its other hypotheses, so $p_r=S_n$ by
Lemma~\ref{gal-cover-trace-limit}.
% Lean: (2) is ClosedPhi = congPhi_of_closedTr ClosedTr_holds powTr_holds, with
% ClosedTr_holds = closedTr_of coverJ_holds frEx_holds opEx_holds fixOf_holds tie_holds trLim_holds.
% closedTr_of takes the datum from BSP.nonempty_bk_of_finExt and the cyclotomic element from
% BSP.exists_isPer, and pulls the cover back to D_A by Cover.pull.
\end{proof}
\leanref{Development/AbelianVarietyTheory/Proofs/InertiaEigenvalues/ClosedOfCover.lean}{GlobalTheory.CovClose.closed_of_cover}
\leanref{Development/AbelianVarietyTheory/Definitions/NonGaloisDescent/TrivIsCr.lean}{GlobalTheory.BKLim.Cover}
\leanref{Development/AutomorphicArithmetic/Foundations/GaloisRepresentations/ClosedPhiFunctoriality/ClosedPhiHolds.lean}{UpperTheory.Gal.ClosedPhiFromClosedTrace.ClosedTr_holds}
\leanref{Development/AutomorphicArithmetic/Foundations/GaloisRepresentations/ClosedPhiFunctoriality/IsOp2One.lean}{UpperTheory.Gal.ClosedPhiFromClosedTrace.ClosedTr}
\leanref{Development/AutomorphicArithmetic/Foundations/GaloisRepresentations/ClosedPhiFunctoriality/ClosedPhiOfClosedTr.lean}{UpperTheory.Gal.ClosedPhiFromClosedTrace.closedTr_of}
\leanref{Development/AutomorphicArithmetic/Foundations/GaloisRepresentations/ClosedPhiFunctoriality/ClosedPhiOfClosedTr.lean}{UpperTheory.Gal.ClosedPhiFromClosedTrace.congPhi_of_closedTr}
\leanref{Development/AutomorphicArithmetic/Foundations/GaloisRepresentations/ClosedPhiFunctoriality/IsOp2One.lean}{UpperTheory.Gal.ClosedPhiFromClosedTrace.PowTr}
\leanref{Development/AutomorphicArithmetic/Foundations/GaloisRepresentations/ClosedPhiFunctoriality/PowTrHolds.lean}{UpperTheory.Gal.ClosedPhiFromClosedTrace.powTr_holds}

\begin{proposition}\label{gal-limits}
In Situation~\ref{gal-definite-situation}, there is a continuous semisimple
$r:G_E\to\operatorname{GL}_a(\bar{\mathbf Q}_\ell)$ such that:
\begin{enumerate}
\item for every $d$-split place $w$ not over $\ell$ (Definition~\ref{gal-seed}) and every $B$ such
  that $U$ has Satake parameters $B$ at $w$, $r$ has Frobenius eigenvalues $q_w^{(a-1)/2}B$ at $w$;
\item for every place $w$ over $\ell$, over a place of $F$ different from those under $w_0$ and $v'$,
  at which $U$ has Satake parameters, and every $A$ over $w$, $r|_{D_A}$ is crystalline and de Rham,
  and for every $\tau:E\to\bar{\mathbf Q}_\ell$ extending continuously to $E_w$ and every weight
  $\lambda$ of $U$ at $\iota\circ\tau$, $\mathrm{HT}_\tau(r|_{D_A})=\{k-\lambda_k:0\leq k<a\}$;
\item at every such $w$ with Satake parameters $B$, $r$ satisfies (phi) with respect to $B$
  (Definition~\ref{gal-spherical-conditions}).
\end{enumerate}
\end{proposition}
\leanref{Development/AbelianVarietyTheory/Definitions/GaloisGeometry/LimWU4OfLimWU3.lean}{GlobalTheory.CongL.LimFinOdd}
\leanref{Development/AutomorphicArithmetic/Foundations/GaloisRepresentations/TraceLimits/LimFinOddOf.lean}{UpperTheory.Gal.Lim.limFinOdd_of}
\leanref{Development/AutomorphicArithmetic/Foundations/GaloisRepresentations/TraceLimits/LimFinOddOf.lean}{UpperTheory.Gal.Lim.limWU4_of}
\leanref{Development/AutomorphicArithmetic/Foundations/GaloisRepresentations/IntegralStructures/CrystallineDescent/LimOddOfLimWU3.lean}{UpperTheory.Gal.Int.Core.LimU}
\leanref{Development/AutomorphicArithmetic/Foundations/GaloisRepresentations/IntegralStructures/CrystallineDescent/LimOddOfLimWU3.lean}{UpperTheory.Gal.Int.Core.limU_of_limOdd}
\leanref{Development/AutomorphicArithmetic/Foundations/GaloisRepresentations/IntegralStructures/ZetaCharacterDivisibility/MainFin.lean}{UpperTheory.Gal.Int.Div.seedFinOdd_of_geoOne}
\leanref{Development/AbelianVarietyTheory/Definitions/GaloisGeometry/LimWU4OfLimWU3.lean}{GlobalTheory.CongL.SeedFinOdd}

\begin{proof}
\emph{Seeds of rank $a$.} Let $U'$ satisfy the hypotheses of Theorem~\ref{gal-seed-representations},
with lacunary weight at $v'$, the supercuspidal type at $w_0$, and generic Satake parameters at the
places of a finite set $Y$. Theorem~\ref{gal-seed-representations}(a)--(c) gives seed representations,
and Proposition~\ref{gal-division-seed} gives one $a$-dimensional $r'$, independent of the seed up to
isomorphism. By Proposition~\ref{gal-division-seed}(0) and (1), with
Theorem~\ref{gal-seed-representations}(b), $r'$ has Frobenius eigenvalues $q_w^{(a-1)/2}B$ at the
$d$-split places $w\in Y$, $w\nmid\ell$, and its characteristic polynomials have this form at Frobenius
elements in an open dense set at almost all $d$-split places. By
Proposition~\ref{gal-division-seed}(2) and (3), with Theorem~\ref{gal-seed-representations}(c), $r'$ is
crystalline and de Rham with the Hodge--Tate weights of (2) and satisfies (phi) at the places of $Y$
over $\ell$ away from $w_0$ and $v'$. Here $w_0$ and $v'$ may be replaced within their classes over
$\mathbf Q(x)$ by complex conjugation, as the hypotheses of Theorem~\ref{gal-seed-representations}
require.

\emph{Congruences.} Proposition~\ref{gal-congruences} gives the levels $K_w$ and, for every finite $Y$,
approximations of the eigenvalue system $\theta$ of $U$ by those of forms $U'$ as in the first step,
generic at the places of $Y$.

\emph{Limit.} Choose Hecke operators $(O_j)_{j\in J}$ of level $(K_w)$ at the $d$-split places
$w\nmid\ell$ where $U$ has Satake parameters, and elements $x_j\in G_E$, such that $x_j$ is a Frobenius
element at the place of $O_j$, the $x_j$ at places outside any finite set are dense, and the value of
$O_j$ on a form with Satake parameters $B$ is the trace of $x_j$ on any representation with Frobenius
eigenvalues $q_w^{(a-1)/2}B$ there. For a form $U'$ of the first step let $\eta_{U'}=\theta'$ and
$D_{U'}$ the set of $j$ with $x_j$ in the open dense set of Proposition~\ref{gal-division-seed}(0) for
$U'$ at a place outside its exceptional set; then $\eta_{U'}(j)=\operatorname{tr}r'(x_j)$ for $j\in
D_{U'}$, and finite intersections of the $x(D_{U'})$ are dense. Lemma~\ref{gal-trace-limit}(1) gives a
continuous $t$ with values in a finite extension of $\mathbf Q_\ell$, and Lemma~\ref{gal-trace-limit}(2)
gives a continuous semisimple $r$ with trace $t$, whose trace is approximated by the traces of the
$r'$ attached to forms generic at any given finite set $Y$.

\emph{Properties.} For (1), take $Y=\{w\}$: the representations $r'$ are trivial on $I_A$ and have
Frobenius eigenvalues $q_w^{(a-1)/2}B'$, the approximation of traces shows that $r$ is trivial on $I_A$,
and the value of $t$ at the Frobenius elements $x_j$ at $w$ gives the Hecke eigenvalues of $U$, which
determine the characteristic polynomial. For (2), take $Y=\{w\}$ and apply
Lemma~\ref{gal-closed-crystalline}(1): the $r'$ are crystalline at $w$ with the Hodge--Tate multisets of
(2), since the weights of $U'$ agree with those of $U$ at the embeddings that lie over neither $w_0$ nor
$v'$, and $w$ lies over a place of $F$ different from those under $w_0$ and $v'$. For (3), with
$Y=\{w\}$, Proposition~\ref{gal-congruences} approximates
the power sums of the Satake parameters of $U$ at $w$, scaled by the constants $c_{w,n}$, by the same
combinations of those of the $U'$, which by (phi) for $r'$ are the power sums of the Frobenius
eigenvalues of $\mathrm{WD}_\iota(r'|_{D_A})$. Lemma~\ref{gal-closed-crystalline}(2) shows that the power
sums of $\mathrm{WD}_\iota(r|_{D_A})(\sigma)$ are those of $\mathrm{Sat}_A(B)(\sigma)$ for every
$\sigma\in W_A$, and power sums of all orders determine the characteristic polynomial.
\end{proof}

\subsection{Enlargements and constituents}\label{subsection-gal-enlargements}

\begin{definition}\label{gal-larger-enlargement}
Let $E$ be a CM field with maximal totally real subfield $F$. A Hecke character $\psi$ of $E$ is
\emph{algebraic of infinity type} $p=(p_\tau)_{\tau:E\to\mathbf C}$ if
$\psi_\infty(z)=\prod_\tau\tau(z)^{p_\tau}$ for $z\in E^\times$, where $\psi_\infty$ is the restriction of
$\psi$ to $(E\otimes\mathbf R)^\times$. Let $n\geq1$ and let $U$ be an automorphic representation of $U_n$
with weight $\lambda_\tau$ at every $\tau:E\to\mathbf C$ (Definition~\ref{gal-definite-unitary}).
\begin{enumerate}
\item Let $t\geq0$ and let $\kappa$ be a Hecke character of $E$ algebraic of infinity type $m$. An
  \emph{enlargement} of $(U,\lambda)$ along $\kappa$ in rank $n+t$ consists of an automorphic
  representation $U'$ of $U_{n+t}$ with weight $\lambda'_\tau$ at every $\tau$ and Hecke characters
  $\psi_1,\ldots,\psi_t$ of $E$ algebraic of infinity types $p_1,\ldots,p_t$, such that:
  \begin{enumerate}
  \item for every $\tau$, as multisets,
    \[
     \{k-\lambda'_{\tau,k}:0\leq k<n+t\}=\{k-\lambda_{\tau,k}+m_\tau:0\leq k<n\}
     +\{-p_{1,\tau},\ldots,-p_{t,\tau}\};
    \]
  \item for every place $w$ of $E$ split over $F$ at which $\kappa,\psi_1,\ldots,\psi_t$ are unramified and
    every $B$ such that $U$ has Satake parameters $B$ at $w$, there is $B'$ such that $U'$ has Satake
    parameters $B'$ at $w$ and
    \[
     q_w^{(n+t-1)/2}B'=\kappa(\varpi_w)^{-1}q_w^{(n-1)/2}B+\{\psi_1(\varpi_w),\ldots,\psi_t(\varpi_w)\},
    \]
    where $xB=\{xb:b\in B\}$ for a complex number $x$.
  \end{enumerate}
\item Let $w_0$ be a place of $E$ split over $F$. A \emph{pair of enlargements} of $(U,\lambda)$ at $w_0$
  consists of an integer $t\geq0$ with $n+t$ odd, a Hecke character $\kappa$ of $E$ algebraic of some
  infinity type and unramified at $w_0$, and for $i=1,2$ an enlargement
  $(U_i,\lambda_i,\psi_{i,1},\ldots,\psi_{i,t})$ of $(U,\lambda)$ along $\kappa$ in rank $n+t$, such that
  $\psi_{1,j}\neq\psi_{2,j'}$ for all $j,j'$ and $\psi_{1,1},\ldots,\psi_{1,t}$ are unramified at $w_0$.
\end{enumerate}
\end{definition}
\leanref{Development/AbelianVarietyTheory/PublishedTheorems/AutomorphicForms₂/ExistsUniqueHcharOfFiniteOrder.lean}{GlobalTheory.HeckeGal.IsAlgPos}
\leanref{Development/AutomorphicArithmetic/Foundations/GaloisRepresentations/ThetaParseval/SplitLineSubtraction/RowDefOfOdd.lean}{UpperTheory.Gal.TP.Sub.IsUp}
\leanref{Development/AutomorphicArithmetic/Foundations/GaloisRepresentations/ThetaParseval/SplitLineSubtraction/RowDefOfOdd.lean}{UpperTheory.Gal.TP.Sub.wtJumps}
\leanref{Development/AutomorphicArithmetic/Foundations/GaloisRepresentations/ThetaParseval/SplitLineSubtraction/RowDefOfOdd.lean}{UpperTheory.Gal.TP.Sub.TwoUp}

\begin{lemma}\label{gal-larger-good-base}
Let $E$ be a CM field with maximal totally real subfield $F$ and complex conjugation $c$. Suppose that $E$
contains $\sqrt{-D}$ for some integer $D>0$, that $E/F$ is unramified at every finite place, and that
every finite place $v$ of $F$ inert in $E$ satisfies $v\nmid2$, $N(v)>3$ and $v\nmid\mathfrak d_F$, where
$\mathfrak d_F$ is the different of $F$. Then there are a good base for $E$, with embeddings $\tau_v$ at the
infinite places and $P=\emptyset$ (Definition~\ref{gal-good-base}), an element $\delta\in E^\times$ with
$c(\delta)=-\delta$ and $i\tau_v(\delta)>0$ for every infinite place $v$ of $F$, and a fractional ideal
$\mathfrak a$ of $E$, such that for all $n,d\geq1$, $E^n$ with the lattice $\mathcal O_E^n$ and the matrices
$\theta_v=1$ and $\beta_v=1$ is a standard Hermitian space of rank $n$, and $E^d$ with the form
$\delta\sum_ix_ic(y_i)$, the lattice $\mathfrak a^d$ and the matrices $\theta_v=1$ and $\beta_v=1$ is a scalar
skew-Hermitian space of rank $d$.
\end{lemma}
\leanref{Development/AutomorphicArithmetic/Foundations/GaloisRepresentations/ThetaParseval/SplitLineSubtraction/TwoUpOKOf.lean}{UpperTheory.Gal.TP.Sub.BaseNG}
\leanref{Development/AutomorphicArithmetic/Foundations/GaloisRepresentations/ThetaParseval/SplitLineSubtraction/PBaseNG.lean}{UpperTheory.Gal.TP.Sub.baseNG_holds}
\leanref{Development/AutomorphicArithmetic/Foundations/GaloisRepresentations/ThetaParseval/CartanIdeles/SpacesExOf.lean}{UpperTheory.Gal.TP.Skew.SkewLine}
\leanref{Development/AutomorphicArithmetic/Foundations/GaloisRepresentations/ThetaParseval/CartanIdeles/SpacesExOf.lean}{UpperTheory.Gal.TP.Skew.IsLine}
\leanref{Development/AutomorphicArithmetic/Foundations/GaloisRepresentations/ThetaParseval/CartanIdeles/LineEx.lean}{UpperTheory.Gal.TP.Skew.lineEx}
\leanref{Development/AutomorphicArithmetic/Foundations/GaloisRepresentations/ThetaParseval/CartanIdeles/OrdLine.lean}{UpperTheory.Gal.TP.Skew.ordLine}
\leanref{Development/AutomorphicArithmetic/Foundations/GaloisRepresentations/ThetaParseval/CartanIdeles/IdeleStep.lean}{UpperTheory.Gal.TP.Skew.ideleStep}

\begin{proof}
\emph{The additive character.} Let $\psi_{\mathbf Q}$ be the character of $\mathbf Q\backslash\mathbf A_{\mathbf Q}$
with $\psi_{\mathbf Q}(x)=\exp(-2\pi ix)$ on $\mathbf R$ and trivial on $\hat{\mathbf Z}$, and put
$\psi=\psi_{\mathbf Q}\circ\operatorname{Tr}_{F/\mathbf Q}$. Then $\psi(x)=\exp(-2\pi i\operatorname{Tr}(x))$ on
$F\otimes\mathbf R$. At a finite place $v$, the largest fractional ideal of $F_v$ on which $\psi_v$ is
trivial is the local inverse different, so $\psi_v$ has conductor $\mathcal O_v$ when $v\nmid\mathfrak d_F$,
in particular at every place inert in $E$.

\emph{The element $\delta$ and the ideal $\mathfrak a$.} Let
$\mathfrak d_E^{-1}=\{z\in E:\operatorname{Tr}_{E/\mathbf Q}(z\mathcal O_E)\subset\mathbf Z\}$ be the inverse
different of $E$ and $\delta_0=\sqrt{-D}$, so that $c(\delta_0)=-\delta_0$. For a finite place $w$ of $E$ put
$e(w)=\operatorname{ord}_w(\mathfrak d_E^{-1})-\operatorname{ord}_w(\delta_0)$. Since $c$ preserves
$\mathfrak d_E$ and $\delta_0$ up to sign, $e(cw)=e(w)$; since $E/F$ is unramified at the finite places,
there is an idele $a$ of $F$ with $\operatorname{ord}_v(a)=e(w)$ for every finite place $w$ over every
finite place $v$. By global class field theory, the ideles of $F$ whose Artin symbol for $E/F$ is trivial
form the group $F^\times N_{E/F}(\mathbf A_E^\times)$ \cite[Chapter~VI]{Neu99}. A real place of $F$ ramifies
in $E$, and there the local symbol of $-1$ is nontrivial; multiplying $a$ by $-1$ at one real place if
necessary, which changes no $\operatorname{ord}_v(a)$, gives $a=xN_{E/F}(b)$ with $x\in F^\times$ and
$b\in\mathbf A_E^\times$. Put $\delta=x\delta_0$ and let $\mathfrak a$ be the fractional ideal with
$\operatorname{ord}_w(\mathfrak a)=\operatorname{ord}_w(b_w)$. Then $c(\delta)=-\delta$, and for every finite
place $w$ over $v$,
\[
 \operatorname{ord}_w(\delta)+\operatorname{ord}_w(\mathfrak a)+\operatorname{ord}_{cw}(\mathfrak a)
 =\operatorname{ord}_v(x)+\operatorname{ord}_w(\delta_0)+\operatorname{ord}_v(N_{E/F}b)
 =\operatorname{ord}_w(\mathfrak d_E^{-1}),
\]
because $\operatorname{ord}_v(N_{E/F}b)$ is $\operatorname{ord}_w(b_w)+\operatorname{ord}_{cw}(b_{cw})$ if $v$
splits and $2\operatorname{ord}_w(b_w)$ if $v$ is inert.

\emph{Self-duality.} For $y\in E^\times$, the condition that
$\operatorname{Tr}_{E/\mathbf Q}(oy\delta c(x))\in\mathbf Z$ for all $x\in\mathfrak a$ and $o\in\mathcal O_E$
says that $y\delta c(\mathfrak a)\subset\mathfrak d_E^{-1}$, that is,
$\operatorname{ord}_w(y)+\operatorname{ord}_w(\delta)+\operatorname{ord}_{cw}(\mathfrak a)\geq
\operatorname{ord}_w(\mathfrak d_E^{-1})$ for every $w$. By the identity above this is
$\operatorname{ord}_w(y)\geq\operatorname{ord}_w(\mathfrak a)$ for every $w$, that is, $y\in\mathfrak a$.
Applied in each coordinate, this shows that $\mathfrak a^d$ is self-dual for $\delta\sum_ix_ic(y_i)$ in the
sense of Definition~\ref{gal-good-base}(4). The lattice $\mathcal O_E^n$ is self-dual for
$h(x,y)=\sum_ix_ic(y_i)$, since $h(y,e_i)=y_i$ for the standard basis $(e_i)$.

\emph{The embeddings.} Since $c(\delta)=-\delta\neq0$, every complex embedding takes $\delta$ to a nonzero
purely imaginary number. The two embeddings over an infinite place $v$ of $F$ are complex conjugate, so
one of them, $\tau_v$, satisfies $i\tau_v(\delta)>0$. Let $m_\alpha(v)=m_\beta(v)=0$, and let $v_1$ be a finite
place of $F$ split in $E$, which exists by the Chebotarev density theorem \cite[Chapter~VII, \S13]{Neu99},
with an $F$-embedding $\tau_1:E\to F_{v_1}$. With these data, the character $\psi$ and $P=\emptyset$, the
conditions of Definition~\ref{gal-good-base}(1) hold by the hypotheses on $E$ and the conductors of the
$\psi_v$.

\emph{The Cartan involutions.} The identity matrix is a Cartan involution of the identity matrix, since
$\operatorname{Re}({}^tx\bar x)=\sum_i|x_i|^2>0$ for $x\neq0$, and of the scalar matrix $i\tau_v(\delta)$, since
$i\tau_v(\delta)\sum_i|x_i|^2>0$ for $x\neq0$. As $P=\emptyset$, the matrices $\beta_v=1$ are subject to no
condition.
\end{proof}
% Lean: SkewLine (skewLine_of from SpacesEx and BaseEx with m_alpha = m_beta = 0; BaseEx takes v1 from
% GlobalTheory.Orb.exists_split_place_notMem, which the Lean proves its own way), with stdSpc and skewSpc
% (theta = 1, bas = 1) and the Cartan clauses of stdHerm and LineSkew.

\begin{lemma}\label{gal-larger-conjugate-character}
Let $E$ be a CM field with maximal totally real subfield $F$ such that $E/F$ is unramified at every finite
place, and let $\eta$ be the quadratic character of $\mathbf A_F^\times$ attached to $E/F$. Let
$\tau_v:E\to\mathbf C$ be an embedding over each infinite place $v$ of $F$, let $m(v)$ be odd integers, and
let $w_0$ be a finite place of $E$. Then there is a unitary Hecke character $\chi$ of $E$, unramified at
$w_0$, with $\chi|_{\mathbf A_F^\times}=\eta$ and $\chi(z)=\prod_v(\tau_v(z_v)/|\tau_v(z_v)|)^{m(v)}$ on
$(E\otimes\mathbf R)^\times$.
\end{lemma}
\leanref{Development/AutomorphicArithmetic/Foundations/GaloisRepresentations/ThetaParseval/IdeleNormCharacters/AbsFactsAlt.lean}{UpperTheory.Gal.TP.Chr.MuEmb}
\leanref{Development/AutomorphicArithmetic/Foundations/GaloisRepresentations/ThetaParseval/IdeleNormCharacters/ConjDivNorm.lean}{UpperTheory.Gal.TP.Chr.muEmb}
\leanref{Development/AutomorphicArithmetic/Foundations/Convolution/HermitianDescent/OffUnitsLeUnitIdl.lean}{UpperTheory.ConvolutionTheory.D.muEx}
\leanref{Development/AbelianVarietyTheory/Proofs/IwasawaDecomposition/ExistsUnitaryCharOfOpenSubgroupAndHom.lean}{GlobalTheory.exists_unitary_char_of_open_subgroup_and_hom}

\begin{proof}
Write $c$ also for the action of complex conjugation on the ideles $\mathbf A_E^\times$, and put
$\psi_\infty(x)=\prod_v(\tau_v(x_v)/|\tau_v(x_v)|)^{m(v)}$, a continuous unitary character of
$\mathbf A_E^\times$ through its infinite component. Choose a finite place $w_1$ of $E$, split over $F$,
that lies over a place of $F$ different from the one under $w_0$, and an open subgroup
$V\subset\mathcal O_{w_1}^\times$ that contains no root of unity other than $1$. Let $U'$ be the subgroup
of $\mathbf A_E^\times$ generated by the units of $E_w$ at the finite places $w$ not over the place of $F$
under $w_1$, and $O$ the open subgroup of the $x\in\mathbf A_E^\times$ such that $x$ and $c(x)$ are units
at every finite place and the component of $x/c(x)$ at $w_1$ lies in $V$.

\emph{Consistency.} Let $k\in E^\times$, $b\in\mathbf A_F^\times$ and $u\in U'$ with $x=kbu\in O$; we show
that $\eta(b)=\psi_\infty(x)$. Put $t=k/c(k)$. Since $c(b)=b$, $t=(x/c(x))(u/c(u))^{-1}$, so $t$ is a
unit at every finite place; and $|\sigma(t)|=1$ for every embedding $\sigma:E\to\mathbf C$, since
$\sigma(c(k))$ is the complex conjugate of $\sigma(k)$. So $t$ is a root of unity, by Kronecker's
theorem. The components of $u$ and $c(u)$ at $w_1$ are $1$, so the component of $t$ at $w_1$ is that of
$x/c(x)$, which lies in $V$;
hence $t=1$, and $k\in F^\times$. Put $b'=kb\in\mathbf A_F^\times$. Then $\eta(b')=\eta(b)$, and
$b'=xu^{-1}$ is a unit at every finite place. Since $E/F$ is unramified at every finite place, $\eta$ is
trivial on these units, so $\eta(b')$ is the product of the signs of the components $b'_v$ at the real
places $v$ of $F$. As $\psi_\infty(u)=1$ and each $m(v)$ is odd,
$\psi_\infty(x)=\psi_\infty(b')=\prod_v\operatorname{sgn}(b'_v)^{m(v)}=\eta(b')=\eta(b)$.

\emph{Extension.} Let $A=E^\times\times\mathbf A_F^\times\times U'$, with the homomorphism
$\alpha(k,b,u)=kbu$ to $\mathbf A_E^\times$ and the unitary character $a(k,b,u)=\eta(b)$. By the
consistency, $a(y)=\psi_\infty(\alpha(y))$ whenever $\alpha(y)\in O$. Hence the continuous unitary
character $(g,y)\mapsto\psi_\infty(g)a(y)^{-1}$ of the open subgroup $O\times A$ of
$\mathbf A_E^\times\times A$, with $A$ discrete, is trivial on its intersection with the subgroup of the
$(\alpha(y),y)$. By Pontryagin duality it extends to a continuous unitary character of
$\mathbf A_E^\times\times A$ trivial on that subgroup. Its restriction $\chi$ to $\mathbf A_E^\times$ is
continuous and unitary, equal to $\psi_\infty$ on $O$, and satisfies $\chi(\alpha(y))=a(y)$ for $y\in A$.
So $\chi$ is trivial on $E^\times$, a Hecke character; $\chi|_{\mathbf A_F^\times}=\eta$; $\chi$ is trivial
on $U'$, so unramified at every finite place not over the place of $F$ under $w_1$, in particular at
$w_0$; and $\chi=\psi_\infty$ on $(E\otimes\mathbf R)^\times\subset O$.
\end{proof}

\begin{definition}\label{gal-larger-attuned}
Let $E$ be a CM field with maximal totally real subfield $F$ and complex conjugation $c$, let
$\tau_v:E\to\mathbf C$ be an embedding over each infinite place $v$ of $F$, and let $n,t\geq0$. Let
$\chi_X,\chi_W$ be unitary Hecke characters of $E$ and $m_X(v),m_W(v)$ integers. Hecke characters
$\kappa,\psi_0,\ldots,\psi_{t-1}$ of $E$, algebraic of infinity types $m,p_0,\ldots,p_{t-1}$
(Definition~\ref{gal-larger-enlargement}), are \emph{attuned} to $(\chi_X,\chi_W)$ in rank $(n,t)$ if:
\begin{enumerate}
\item $2m_{\tau_v}=m_W(v)-m_X(v)+t$ and $2m_{c\tau_v}=m_X(v)-m_W(v)+t$ for every $v$;
\item $2p_{i,\tau_v}=m_X(v)+2i-(n+2t-2)$ and $2p_{i,c\tau_v}=-m_X(v)+2i-(n+2t-2)$ for every $v$ and
  every $i$;
\item at every finite place $w$ at which $\kappa,\psi_0,\ldots,\psi_{t-1}$ are unramified, $\chi_X$ and
  $\chi_W$ are unramified, $\kappa(\varpi_w)=\chi_W(\varpi_w)\chi_X(\varpi_w)^{-1}q_w^{-t/2}$ and
  $\psi_i(\varpi_w)=\chi_X(\varpi_w)q_w^{(n+2t-2)/2-i}$ for every $i$.
\end{enumerate}
\end{definition}
\leanref{Development/AutomorphicArithmetic/Foundations/GaloisRepresentations/ThetaParseval/SplitLineSubtraction/TwoUpOKOf.lean}{UpperTheory.Gal.TP.Sub.Dict}

\begin{lemma}\label{gal-larger-characters}
Let $E$ be a CM field with maximal totally real subfield $F$ such that $E/F$ is unramified at every finite
place, and let $\eta$ be the quadratic character of $\mathbf A_F^\times$ attached to $E/F$. Let
$\tau_v:E\to\mathbf C$ be an embedding over each infinite place $v$ of $F$, let $n\geq0$ be even and
$t\geq1$ odd, let $b$ be an integer, and let $w_0$ be a place of $E$ split over $F$. Then there are a Hecke
character $\kappa$ of $E$ algebraic of some infinity type (Definition~\ref{gal-larger-enlargement}), and
for $j=0,1$ unitary Hecke characters $\chi_{X,j},\chi_{W,j}$ of $E$, integers $m_{X,j}(v),m_{W,j}(v)$ and
Hecke characters $\psi_{j,0},\ldots,\psi_{j,t-1}$ of $E$ algebraic of some infinity types, such that:
\begin{enumerate}
\item $\chi_{X,j}|_{\mathbf A_F^\times}=\eta^n$ and $\chi_{W,j}|_{\mathbf A_F^\times}=\eta^{n+t}$; on
  $(E\otimes\mathbf R)^\times$, $\chi_{X,j}(z)=\prod_v(\tau_v(z_v)/|\tau_v(z_v)|)^{m_{X,j}(v)}$ and
  $\chi_{W,j}(z)=\prod_v(\tau_v(z_v)/|\tau_v(z_v)|)^{m_{W,j}(v)}$; and $m_{W,j}(v)\leq b$ for every $v$;
\item $\kappa,\psi_{j,0},\ldots,\psi_{j,t-1}$ are attuned to $(\chi_{X,j},\chi_{W,j})$ in rank $(n,t)$
  (Definition~\ref{gal-larger-attuned});
\item $\psi_{0,i}\neq\psi_{1,i'}$ for all $i,i'$, and $\kappa,\psi_{0,0},\ldots,\psi_{0,t-1}$ are unramified
  at $w_0$.
\end{enumerate}
\end{lemma}
\leanref{Development/AutomorphicArithmetic/Foundations/GaloisRepresentations/ThetaParseval/SplitLineSubtraction/TwoUpOKOf.lean}{UpperTheory.Gal.TP.Sub.ChrPairs}
\leanref{Development/AutomorphicArithmetic/Foundations/GaloisRepresentations/ThetaParseval/IdeleNormCharacters/ChrPairsPf.lean}{UpperTheory.Gal.TP.Chr.chrPairs}
\leanref{Development/AutomorphicArithmetic/Foundations/GaloisRepresentations/ThetaParseval/IdeleNormCharacters/ChrPairsHolds.lean}{UpperTheory.Gal.TP.Chr.chrPairs_holds}
\leanref{Development/AutomorphicArithmetic/Foundations/GaloisRepresentations/ThetaParseval/IdeleNormCharacters/AbsFactsAlt.lean}{UpperTheory.Gal.TP.Chr.ChrEx}
\leanref{Development/AutomorphicArithmetic/Foundations/GaloisRepresentations/ThetaParseval/IdeleNormCharacters/ChrPairsHolds.lean}{UpperTheory.Gal.TP.Chr.chrEx_holds}

\begin{proof}
Lemma~\ref{gal-larger-conjugate-character} with $m(v)=-1$ gives a unitary Hecke character $\chi$ of $E$,
unramified at $w_0$, with $\chi|_{\mathbf A_F^\times}=\eta$ and
$\chi(z)=\prod_v(\tau_v(z_v)/|\tau_v(z_v)|)^{-1}$ on $(E\otimes\mathbf R)^\times$. Let $\|\cdot\|$ be the
idelic norm of $E$, so that $\|\varpi_w\|=q_w^{-1}$ and $\|z\|=\prod_v|\tau_v(z_v)|^2$ on
$(E\otimes\mathbf R)^\times$. For integers $k\geq0$ and $s$ put $\psi(k,s)=\chi^k\|\cdot\|^{s/2}$. Then:
\begin{enumerate}
\item[(a)] if $k+s$ is even, $\psi(k,s)$ is algebraic of infinity type $p$ with $p_{\tau_v}=(s-k)/2$ and
  $p_{c\tau_v}=(s+k)/2$, since for $z\in E^\times$ its value at infinity is
  $\prod_v\tau_v(z)^{-k}|\tau_v(z)|^{k+s}=\prod_v\tau_v(z)^{(s-k)/2}\,c\tau_v(z)^{(s+k)/2}$;
\item[(b)] $\psi(k,s)$ is unramified at a finite place $w$ if and only if $\chi^k$ is, and its value at
  $\varpi_w$ is $\chi_w(\varpi_w)^kq_w^{-s/2}$;
\item[(c)] if $\psi(k,s)=\psi(k',s')$ with $k\geq k'$, then $\chi^{k-k'}=\|\cdot\|^{(s'-s)/2}$; as $\chi$ is
  unitary, $s=s'$ and $\chi^{k-k'}=1$, and since $\chi^{k-k'}(z)=\prod_v(\tau_v(z_v)/|\tau_v(z_v)|)^{k'-k}$ on
  $(E\otimes\mathbf R)^\times$, $k=k'$.
\end{enumerate}
Choose an odd $K\geq1$ with $-K\leq b$, and put
\[
 \kappa=\psi(K,t),\quad\chi_{X,j}=\chi^{2j},\quad\chi_{W,j}=\chi^{K+2j},\quad m_{X,j}(v)=-2j,\quad
 m_{W,j}(v)=-K-2j,\quad\psi_{j,i}=\psi\bigl(2j,2i-(n+2t-2)\bigr).
\]
The integers $K+t$ and $2j+2i-(n+2t-2)$ are even, since $K$ and $t$ are odd and $n$ is even, so these
characters are algebraic by (a). Item (1) holds because $\eta^{2j}=1=\eta^n$, $\eta^{K+2j}=\eta=\eta^{n+t}$,
by the infinity type of $\chi$, and because $-K-2j\leq-K\leq b$. For (2): the infinity type $m$ of $\kappa$
has $2m_{\tau_v}=t-K$ and $2m_{c\tau_v}=t+K$, while $m_{W,j}(v)-m_{X,j}(v)=-K$, which is condition (1) of
Definition~\ref{gal-larger-attuned}; the infinity type of $\psi_{j,i}$ has
$2p_{\tau_v}=2i-(n+2t-2)-2j$ and $2p_{c\tau_v}=2i-(n+2t-2)+2j$, which is condition (2). For condition
(3), if $\kappa$ and $\psi_{j,0}$ are unramified at $w$, then so are $\chi^K$ and $\chi^{2j}$ by (b), hence
$\chi_{X,j}$ and $\chi_{W,j}$, and
\[
 \kappa(\varpi_w)=\chi_w(\varpi_w)^Kq_w^{-t/2}=\chi_{W,j}(\varpi_w)\chi_{X,j}(\varpi_w)^{-1}q_w^{-t/2},\qquad
 \psi_{j,i}(\varpi_w)=\chi_{X,j}(\varpi_w)q_w^{(n+2t-2)/2-i}.
\]
Item (3) follows from (c), since $0\neq2$, and from (b), since $\chi$ is unramified at $w_0$.
\end{proof}

\begin{lemma}\label{gal-satake-support-gap}
Let $L$ be a finite set and $g:L\to\mathbf Z$ a function with $g(a)\leq0$ for some $a\in L$. Then there is
a nonempty subset $S\subset L$ such that $|g(a)-g(b)|\geq2$ for all $a\in S$ and $b\in L\setminus S$, and
$2\sum_{a\in S}g(a)\leq|S|(|S|-1)$.
\end{lemma}
\leanref{Development/AutomorphicArithmetic/ThetaCorrespondence/GeneralLinearGroups/Imports.lean}{UpperTheory.Lit.Mix.comb}

\begin{proof}
Let $g_0\leq0$ be the minimum of $g$. For an integer $T$ let $S(T)$ be the set of $a\in L$ with
$g(a)\leq T$, and let $(\ast)_T$ be the condition that $g_0\leq T$, $T+1\leq|S(T)|$ and
$2\sum_{a\in S(T)}g(a)\leq|S(T)|(|S(T)|-1)$. The condition $(\ast)_{g_0}$ holds, since $S(g_0)$ is nonempty
and $g=g_0\leq0$ on it. Suppose that $(\ast)_T$ holds and that $T+1$ is a value of $g$. Let $e\geq1$ be the
number of $a$ with $g(a)=T+1$, and $s=|S(T)|$. Then $|S(T+1)|=s+e\geq T+2$, and since $T+1\leq s$,
\[
 2\sum_{a\in S(T+1)}g(a)\leq s(s-1)+2e(T+1)\leq s(s-1)+2es+e(e-1)=(s+e)(s+e-1),
\]
so $(\ast)_{T+1}$ holds. As $|S(T)|\leq|L|$, the condition $(\ast)_T$ fails for $T\geq|L|$. Hence there is
$T$ with $(\ast)_T$ such that $T+1$ is not a value of $g$. The set $S=S(T)$ is nonempty and satisfies the
inequality on sums; the values of $g$ on $S$ are at most $T$, and those on $L\setminus S$ are at least
$T+2$.
\end{proof}

\begin{lemma}\label{gal-satake-support-front}
Let $k$ be a non-archimedean local field, let $\nu=|\det|$ for its absolute value, and let $m,n\geq0$. Let
$\sigma$ be an irreducible smooth representation of $\operatorname{GL}_{m+n}(k)$ whose supercuspidal support
(Definition~\ref{gal-irreducible-classes}, with $k$ in place of $E_w$) is $A_0+A_1$, where $A_0$ and $A_1$
are multisets of supercuspidal representations and the degrees of the members of $A_0$ add up to $m$.
Suppose that no $\rho\in A_0$ and $\rho'\in A_1$ satisfy $\rho\cong\rho'\nu^j$ with $j\in\{-1,0,1\}$. Then
there are an ordering $\rho_1,\ldots,\rho_h$ of $A_0$, irreducible smooth representations
$\rho'_1,\ldots,\rho'_{h'}$ whose degrees add up to $n$, and an injective intertwining map from $\sigma$ to
$\pi_1\times\pi_2$, where $\pi_1=\rho_1\times\cdots\times\rho_h$ and $\pi_2=\rho'_1\times\cdots\times\rho'_{h'}$
are normalized parabolic inductions.
\end{lemma}
\leanref{Development/AutomorphicArithmetic/ThetaCorrespondence/GeneralLinearGroups/MFront.lean}{UpperTheory.Lit.Mix.front}

\begin{proof}
By Jacquet's subrepresentation theorem \cite{BZ77}, $\sigma$ embeds into a product
$\rho''_1\times\cdots\times\rho''_s$ of supercuspidal representations, and by the uniqueness of the
supercuspidal support \cite{BZ77} the multiset of the $\rho''_i$ is $A_0+A_1$. By the case $j=0$ of the
hypothesis, no member of $A_0$ is isomorphic to a member of $A_1$, so each $\rho''_i$ lies in one of the
two. If $\rho''_i$ lies in $A_1$ and $\rho''_{i+1}$ in $A_0$, then $\rho''_i\times\rho''_{i+1}$ is
irreducible, since $\rho''_{i+1}\not\cong\rho''_i\nu^{\pm1}$; it has the same Jordan--Hölder constituents as
$\rho''_{i+1}\times\rho''_i$, so the two are isomorphic \cite{BZ77,Zel80}. By transitivity of parabolic
induction, exchanging these two neighbours does not change the product up to isomorphism. Repeated
exchanges move the members of $A_0$ to the front, keeping the order among the members of $A_0$ and among
those of $A_1$. By transitivity of induction once more, the result is $\pi_1\times\pi_2$, with
$\rho_1,\ldots,\rho_h$ the members of $A_0$ and $\rho'_1,\ldots,\rho'_{h'}$ those of $A_1$ in the order in
which they occur; the degrees of the latter add up to $n$.
\end{proof}

\begin{lemma}\label{gal-satake-unitarizable-exponent}
Let $k$ be a non-archimedean local field with absolute value $|\cdot|$, let $m,n\geq0$, and let
$P\subset\operatorname{GL}_{m+n}(k)$ be the group of block upper triangular matrices
$p=\left(\begin{smallmatrix}a&b\\0&d\end{smallmatrix}\right)$ with $a\in\operatorname{GL}_m(k)$ and
$d\in\operatorname{GL}_n(k)$, with $\delta_P(p)=|\det a|^n|\det d|^{-m}$. Let $\pi_1$ and $\pi_2$ be
representations of $\operatorname{GL}_m(k)$ and $\operatorname{GL}_n(k)$, and let $\pi_1\times\pi_2$ be the
space of smooth functions $f$ on $\operatorname{GL}_{m+n}(k)$ with values in $\pi_1\otimes\pi_2$ such that
$f(pg)=\delta_P(p)^{1/2}(\pi_1(a)\otimes\pi_2(d))f(g)$, with $\operatorname{GL}_{m+n}(k)$ acting by right
translation. Let $\sigma$ be a nonzero admissible unitarizable representation of
$\operatorname{GL}_{m+n}(k)$ with an injective intertwining map to $\pi_1\times\pi_2$. Let $z\in k^\times$
with $|z|<1$, and suppose that the scalar matrix $z\cdot1_m$ acts on $\pi_1$ as multiplication by $c$. Then
$|c|\,|z|^{mn/2}\leq1$.
\end{lemma}
\leanref{Development/AutomorphicArithmetic/ThetaCorrespondence/GeneralLinearGroups/MBound.lean}{UpperTheory.Lit.Mix.bound}

\begin{proof}
Let $f$ be the map, and put $\varphi(v)=f(v)(1)\in\pi_1\otimes\pi_2$, so that $f(v)(g)=\varphi(\sigma(g)v)$.
For $p\in P$, $\varphi(\sigma(p)v)=f(v)(p)=\delta_P(p)^{1/2}(\pi_1(a)\otimes\pi_2(d))\varphi(v)$. Hence
$\varphi\circ\sigma(u)=\varphi$ for $u$ in the unipotent radical $N$ of $P$, and for
$t_0=\operatorname{diag}(z\cdot1_m,1_n)$ and $i\geq0$, $\varphi(\sigma(t_0^i)v)=\gamma^i\varphi(v)$ with
$\gamma=c\,|z|^{mn/2}$. As $f$ is injective and $\sigma\neq0$, there is $v$ with $\varphi(v)\neq0$. Choose a
linear form $l$ on $\pi_1\otimes\pi_2$ with $l(\varphi(v))\neq0$, and put $\lambda=l\circ\varphi$.

The stabilizer $H$ of $v$ is open, as $\sigma$ is smooth. For
$x=\left(\begin{smallmatrix}A&B\\C&D\end{smallmatrix}\right)$ near $1$ the block $D$ is invertible; with
$Z=BD^{-1}$, $u=\left(\begin{smallmatrix}1&Z\\0&1\end{smallmatrix}\right)\in N$ and
$y_i=\left(\begin{smallmatrix}A-ZC&0\\z^iC&D\end{smallmatrix}\right)$ one has $xt_0^i=ut_0^iy_i$. The
$z^i$ with $i\geq0$, together with $0$, form a compact set, since $z^i\to0$. So there is a compact open
subgroup $K$ such that $y_i\in H$ for all $x\in K$ and $i\geq0$, and then
$\lambda(\sigma(x)\sigma(t_0^i)v)=\lambda(\sigma(u)\sigma(t_0^i)\sigma(y_i)v)=\lambda(\sigma(t_0^i)v)$ for
$x\in K$.

Let $\langle\cdot,\cdot\rangle$ be a positive definite invariant Hermitian form on $\sigma$, with norm
$\|\cdot\|$. Let $g\in\operatorname{GL}_{m+n}(k)$ satisfy $\lambda(\sigma(x)\sigma(g)v)=\lambda(\sigma(g)v)$
for all $x\in K$, and let $w$ be the average over $K$ of the vectors $\sigma(x)\sigma(g)v$, a finite average
since $\sigma(g)v$ is fixed by an open subgroup of $K$. Then $w$ is fixed by $K$,
$\lambda(w)=\lambda(\sigma(g)v)$ and $\langle w,w\rangle=\langle\sigma(g)v,w\rangle$, so
$\|w\|\leq\|\sigma(g)v\|=\|v\|$ by the Cauchy--Schwarz inequality. As $\sigma$ is admissible, its
$K$-fixed vectors form a finite-dimensional space, on which $|\lambda|\leq C_0\|\cdot\|$ for some $C_0$.
Hence $|\lambda(\sigma(g)v)|\leq C_0\|v\|$. For $g=t_0^i$ this gives $|\gamma|^i|\lambda(v)|\leq C_0\|v\|$
for all $i\geq0$, and since $\lambda(v)\neq0$, $|\gamma|\leq1$.
\end{proof}

\begin{lemma}\label{gal-satake-unitarizable-support}
Let $k$ be a non-archimedean local field with absolute value $|\cdot|$, let $0\leq r<t$, let $\sigma$ be an
irreducible smooth unitarizable representation of $\operatorname{GL}_r(k)$, and let $\chi$ be a unitary
smooth character of $k^\times$. Then the character $\chi|\cdot|^{-(t-1)/2}$ of $\operatorname{GL}_1(k)$ does
not lie in the supercuspidal support of $\sigma$ (Definition~\ref{gal-irreducible-classes}, with $k$ in
place of $E_w$).
\end{lemma}
\leanref{Development/AutomorphicArithmetic/ThetaCorrespondence/GeneralLinearGroups/MNolow.lean}{UpperTheory.Lit.Mix.nolow}

\begin{proof}
Suppose that it does, and let $(\rho_a)_{a\in A}$ be the supercuspidal support of $\sigma$, indexed by a
finite set $A$. Put $c_0=-(t-1)/2$ and $\nu=|\det|$. Let $L\subset A$ be the set of $a$ with
$\rho_a\cong\chi|\cdot|^{c_0+j}$ for an integer $j$, and for $a\in L$ let $g(a)$ be this $j$, which is
unique since $|\cdot|^s=1$ only for $s=0$. By hypothesis $g(a)=0$ for some $a\in L$. Let $S\subset L$ be the
subset given by Lemma~\ref{gal-satake-support-gap}. The $\rho_a$ with $a\in S$ are characters, so
$m=|S|$ satisfies $1\leq m\leq r$; put $n=r-m$, and let $A_0$ and $A_1$ be the multisets of the $\rho_a$
with $a\in S$ and with $a\notin S$.

Let $a\in S$ and $b\in A\setminus S$, and suppose that $\rho_a\cong\rho_b\nu^j$ with $|j|\leq1$. Then
$\rho_b\cong\chi|\cdot|^{c_0+g(a)-j}$, so $b\in L$ with $g(b)=g(a)-j$ and $|g(a)-g(b)|\leq1$, which
contradicts the choice of $S$. So Lemma~\ref{gal-satake-support-front} applies: $\sigma$ embeds into
$\pi_1\times\pi_2$, where $\pi_1=\chi|\cdot|^{e_1}\times\cdots\times\chi|\cdot|^{e_m}$ and
$e_1,\ldots,e_m$ are the numbers $c_0+g(a)$ for $a\in S$, in some order.

Choose $z\in k^\times$ with $|z|<1$. The scalar matrix $z\cdot1_m$ acts on $\pi_1$ as multiplication by
$c=\prod_i\chi(z)|z|^{e_i}$, the product of the central characters of the factors. The representation
$\sigma$ is admissible \cite{BZ76} and nonzero, so Lemma~\ref{gal-satake-unitarizable-exponent} gives
$|c|\,|z|^{mn/2}\leq1$. As $\chi$ is unitary, $|c|\,|z|^{mn/2}=|z|^{\sum_ie_i+mn/2}$, and as $|z|<1$,
$\sum_ie_i+mn/2\geq0$. By the choice of $S$,
\[
 \sum_ie_i=mc_0+\sum_{a\in S}g(a)\leq-\frac{m(t-1)}2+\frac{m(m-1)}2.
\]
Hence $0\leq m(m+n-t)/2=m(r-t)/2$, which is negative since $m\geq1$ and $r<t$.
\end{proof}

\begin{definition}\label{gal-satake-weil-equivariant}
Let $k$ be a non-archimedean local field with absolute value $|\cdot|$, let $N,r\geq0$, and let $\mu$ be a
smooth character of $k^\times$. Let $\mathcal S_{N,r}$ be the space of Schwartz--Bruhat functions on
$M_{N\times r}(k)$, the locally constant functions with compact support. For
$g\in\operatorname{GL}_r(k)$, $h\in\operatorname{GL}_N(k)$ and $\Phi\in\mathcal S_{N,r}$ let
\[
 \omega_\mu(g,h)\Phi(x)=\mu(\det g)\,\mu(\det h)^{-1}|\det g|^{N/2}|\det h|^{-r/2}\,\Phi(h^{-1}xg).
\]
Let $\sigma$ and $\pi$ be representations of $\operatorname{GL}_r(k)$ and $\operatorname{GL}_N(k)$ on complex
vector spaces $V$ and $W$.
\begin{enumerate}
\item A bilinear map $B:\mathcal S_{N,r}\times W\to V$ is \emph{$\omega_\mu$-equivariant} if
  $B(\omega_\mu(g,h)\Phi,\pi(h)w)=\sigma(g)B(\Phi,w)$ for all $g$, $h$, $\Phi$ and $w$.
\item For $j\geq0$, such a map $B$ \emph{vanishes from rank $j$} if $B(\Phi,w)=0$ for every $w\in W$ and
  every $\Phi\in\mathcal S_{N,r}$ that vanishes at all matrices of rank less than $j$.
\item For $c\in k^\times$ and $x_0\in M_{N\times r}(k)$, $D_c(x_0)$ is the set of $x\in M_{N\times r}(k)$
  with $|x_{il}-(x_0)_{il}|\leq|c|$ for all $i,l$, and $\mathbf 1_{D_c(x_0)}$ is its indicator function. For
  $j\leq\min(N,r)$, $e_j\in M_{N\times r}(k)$ is the matrix whose entries at $(i,i)$ for $i\leq j$ are $1$
  and whose other entries are $0$.
\end{enumerate}
\end{definition}
\leanref{Development/AbelianVarietyTheory/Definitions/Foundations/MatrixMeasurableSpace.lean}{GlobalTheory.SplitWeil.om}
\leanref{Development/AbelianVarietyTheory/Definitions/Foundations/MatrixMeasurableSpace.lean}{GlobalTheory.SplitWeil.SB}
\leanref{Development/WeilBound/CoreStatements/TemperedRepresentationTheory/GermIntegrals/IsGerm.lean}{PositivityTheory.Tls2.Ming3.Equi}
\leanref{Development/WeilBound/CoreStatements/TemperedRepresentationTheory/GermIntegrals/IsGerm.lean}{PositivityTheory.Tls2.Ming3.VanRk}
\leanref{Development/WeilBound/CoreStatements/TemperedRepresentationTheory/GermIntegrals/IsGerm.lean}{PositivityTheory.Tls2.Ming3.ballM}
\leanref{Development/WeilBound/CoreStatements/TemperedRepresentationTheory/GermIntegrals/IsGerm.lean}{PositivityTheory.Tls2.Ming3.idem}

\begin{lemma}\label{gal-satake-weil-rank-ball}
Let $k$ be a non-archimedean local field with ring of integers $\mathcal O$ and residue field of $q$
elements, let $N,r\geq0$, and let $0\leq j<r$ with $2j<N$. Let $\mu$ be a unitary smooth character of
$k^\times$, let $\sigma$ be a unitarizable representation of $\operatorname{GL}_r(k)$ on $V$, let $\pi$ be a
representation of $\operatorname{GL}_N(k)$ on $W$, and let $B$ be an $\omega_\mu$-equivariant bilinear map
that vanishes from rank $j+1$ (Definition~\ref{gal-satake-weil-equivariant}). Let $c\in k^\times$ with
$|c|<1$. Then $B(\mathbf 1_{D_c(e_j)},w)=0$ for every $w\in W$.
\end{lemma}
\leanref{Development/WeilBound/RepresentationTheoryGL/PointAveraging/MinkowskiBalls/ValLeOneOfLc.lean}{PositivityTheory.Rg.Ming.core}
\leanref{Development/WeilBound/CoreStatements/TemperedRepresentationTheory/GermIntegrals/StepSt.lean}{PositivityTheory.Tls2.Ming3.CoreSt}

\begin{proof}
Let $\varpi$ be a uniformizer, put $r'=r-j\geq1$, and fix a set $R\subset\mathcal O$ of representatives of
the residue field. For a $j\times r'$ matrix $\rho$ with entries in $R$ let
\[
 M_\rho=\begin{pmatrix}1_j&\varpi^{-1}c\rho\\0&\varpi^{-1}1_{r'}\end{pmatrix}\in\operatorname{GL}_r(k).
\]
There are $q^{jr'}$ such $\rho$, and $\det M_\rho=\varpi^{-r'}$, so
$\omega_\mu(M_\rho,1)\Phi(x)=s\,\Phi(xM_\rho)$ with $s=\mu(\varpi^{-r'})|\varpi|^{-r'N/2}$; as $\mu$ is
unitary, $|s|=q^{r'N/2}$.

We claim that $\mathbf 1_{D_c(e_j)}(x)=\sum_\rho\mathbf 1_{D_c(e_j)}(xM_\rho)$ for every $x$ of rank at
most $j$. Write $x=(x'\;x'')$ with $x'$ its first $j$ columns, so that
$xM_\rho=(x'\;\varpi^{-1}(cx'\rho+x''))$. If $x'$ is not within $|c|$ of the first $j$ columns of $e_j$,
both sides vanish. Otherwise the top $j\times j$ block of $x'$ is congruent to $1_j$ modulo the maximal
ideal and the entries of $x'$ lie in $\mathcal O$. So $x'$ has rank $j$, and as $x$ has rank at most $j$,
$x''=x'T$ for a $j\times r'$ matrix $T$. For $y\in k^j$ and $\gamma\geq0$, the entries of $x'y$ all have
absolute value at most $\gamma$ if and only if those of $y$ do: if $|y_a|$ is the largest of the $|y_b|$,
the $a$-th entry of $x'y$ is $y_a$ plus a sum of terms of smaller absolute value. Hence $x\in D_c(e_j)$ if
and only if $|T_{al}|\leq|c|$ for all $a,l$, and $xM_\rho\in D_c(e_j)$ if and only if
$|c\rho_{al}+T_{al}|\leq|c\varpi|$ for all $a,l$. If some $|T_{al}|>|c|$, no $\rho$ satisfies the second
condition. If all $|T_{al}|\leq|c|$, the second condition says that $\rho_{al}\equiv-T_{al}/c$ modulo
$\varpi$ for all $a,l$, and exactly one $\rho$ satisfies it. This proves the claim.

By the claim, the function
$\mathbf 1_{D_c(e_j)}-s^{-1}\sum_\rho\omega_\mu(M_\rho,1)\mathbf 1_{D_c(e_j)}$ vanishes at all matrices of
rank at most $j$, so $B$ kills it. By equivariance with $h=1$, the vector $v=B(\mathbf 1_{D_c(e_j)},w)$
satisfies $v=s^{-1}\sum_\rho\sigma(M_\rho)v$. Let $\langle\cdot,\cdot\rangle$ be a positive definite
invariant Hermitian form on $\sigma$. By the Cauchy--Schwarz inequality
$|\langle\sigma(M_\rho)v,v\rangle|\leq\langle v,v\rangle$, so
\[
 \langle v,v\rangle\leq q^{jr'}|s|^{-1}\langle v,v\rangle=q^{r'(j-N/2)}\langle v,v\rangle.
\]
As $2j<N$ and $r'\geq1$, $q^{r'(j-N/2)}<1$, so $\langle v,v\rangle=0$ and $v=0$.
\end{proof}

\begin{lemma}\label{gal-satake-weil-rank-step}
Let $k$ be a non-archimedean local field, let $N,r\geq0$, and let $0\leq j<r$ with $2j<N$. Let $\mu$ be a
unitary smooth character of $k^\times$, let $\sigma$ be a unitarizable representation of
$\operatorname{GL}_r(k)$ on $V$, let $\pi$ be a representation of $\operatorname{GL}_N(k)$ on $W$, and let
$B$ be an $\omega_\mu$-equivariant bilinear map that vanishes from rank $j+1$
(Definition~\ref{gal-satake-weil-equivariant}). Then $B$ vanishes from rank $j$.
\end{lemma}
\leanref{Development/WeilBound/RepresentationTheoryGL/PointAveraging/MinkowskiBalls/MemBallMSelf.lean}{PositivityTheory.Rg.Ming.step}
\leanref{Development/WeilBound/CoreStatements/TemperedRepresentationTheory/GermIntegrals/StepSt.lean}{PositivityTheory.Tls2.Ming3.StepSt}
\leanref{Development/WeilBound/RepresentationTheoryGL/PointAveraging/MinkowskiBalls/IsCompactSupport.lean}{PositivityTheory.Rg.Ming.ball}
\leanref{Development/WeilBound/RepresentationTheoryGL/PointAveraging/MinkowskiBalls/IsCompactSupportSB.lean}{PositivityTheory.Rg.Ming.loc}
\leanref{Development/WeilBound/RepresentationTheoryGL/PointAveraging/MinkowskiBalls/ExtBasis.lean}{PositivityTheory.Rg.Ming.nf}
\leanref{Development/WeilBound/CoreStatements/TemperedRepresentationTheory/GermIntegrals/DetNeZeroOfNearOne.lean}{PositivityTheory.Tls2.Ming3.rk}

\begin{proof}
Let $\mathcal O$ be the ring of integers of $k$, $K_N=\operatorname{GL}_N(\mathcal O)$ and
$K_r=\operatorname{GL}_r(\mathcal O)$. The sets $D_c(x_0)$ are compact and open; two of them with the same
$c$ are equal or disjoint; $hD_c(x_0)g=D_c(hx_0g)$ for $h\in K_N$ and $g\in K_r$; they form a basis of
neighbourhoods of $x_0$ as $|c|\to0$; and every $\Phi\in\mathcal S_{N,r}$ is, for $|c|$ small enough, a
finite sum $\sum_{x\in X}\Phi(x)\mathbf 1_{D_c(x)}$ with $\Phi$ constant on each $D_c(x)$. Every matrix of
rank $m$ is $he_mg$ with $h\in\operatorname{GL}_N(k)$ and $g\in\operatorname{GL}_r(k)$. If $|c|<1$, then:
\begin{enumerate}
\item[(a)] every matrix in $D_c(e_m)$ has rank at least $m$, as its top left $m\times m$ block is congruent
  to $1_m$ modulo the maximal ideal;
\item[(b)] every $x\in D_c(e_m)$ of rank $m$ is $he_mg$ with $h\in K_N$ and $g\in K_r$: writing
  $x=\left(\begin{smallmatrix}A&B'\\C&D'\end{smallmatrix}\right)$ with $A$ of size $m\times m$, the block $A$
  lies in $\operatorname{GL}_m(\mathcal O)$, the rank forces $D'=CA^{-1}B'$, and one can take
  $h=\left(\begin{smallmatrix}A&0\\C&1\end{smallmatrix}\right)$ and
  $g=\left(\begin{smallmatrix}1&A^{-1}B'\\0&1\end{smallmatrix}\right)$.
\end{enumerate}
Translation preserves the functions that $B$ kills: if $B(\Phi,\cdot)=0$ and $\Psi(y)=\Phi(hyg)$, then
$\omega_\mu(g^{-1},h)\Psi$ is a nonzero multiple of $\Phi$, and
$B(\Psi,w)=\sigma(g)B(\omega_\mu(g^{-1},h)\Psi,\pi(h)w)=0$.

Let $\Phi$ vanish at all matrices of rank less than $j$, and let $w\in W$. By the compactness of the support
of $\Phi$ and the decomposition into indicator functions of small sets $D_c(x)$, it suffices to find, for
each $x$ with $\Phi(x)\neq0$, an open neighbourhood $U$ of $x$ such that $B(\Psi,w)=0$ for every
$\Psi\in\mathcal S_{N,r}$ that vanishes outside $U$ and at the zeros of $\Phi$. Let $m\geq j$ be the rank of
such an $x$, write $x=he_mg$, choose $c$ with $|c|<1$, and let $U=hD_c(e_m)g$. If $m>j$, every matrix in $U$
has rank at least $j+1$ by (a), so $B(\Psi,\cdot)=0$ by hypothesis. If $m=j$, the function
$\Psi'(y)=\Psi(hyg)$ vanishes outside $D_c(e_j)$ and at all matrices of rank less than $j$, and it suffices
that $B(\Psi',\cdot)=0$. Write $\Psi'=\sum_{x'\in X}\Psi'(x')\mathbf 1_{D_{c'}(x')}$ with $|c'|\leq|c|$ and
$\Psi'$ constant and nonzero on each $D_{c'}(x')$. If $D_{c'}(x')$ contains a matrix $y$ of rank $j$, then
$y\in D_c(e_j)$, so $y=h'e_jg'$ with $h'\in K_N$ and $g'\in K_r$ by (b), and
$D_{c'}(x')=D_{c'}(y)=h'D_{c'}(e_j)g'$; then $B(\mathbf 1_{D_{c'}(x')},\cdot)=0$ by
Lemma~\ref{gal-satake-weil-rank-ball} and translation. Otherwise every matrix in $D_{c'}(x')$ has rank at
least $j$, as $\Psi'$ does not vanish there, and not $j$, so at least $j+1$; then
$B(\mathbf 1_{D_{c'}(x')},\cdot)=0$ by hypothesis.
\end{proof}

\begin{lemma}\label{gal-satake-weil-germ}
Let $k$ be a non-archimedean local field with ring of integers $\mathcal O$ and absolute value $|\cdot|$, let
$r,t\geq0$ and $N=r+t$, and let $dx$ be the Haar measure on $M_{N\times r}(k)$ for which
$M_{N\times r}(\mathcal O)$ has volume $1$. Let $\mu$ be a smooth character of $k^\times$, let $\sigma$ be a
representation of $\operatorname{GL}_r(k)$ on $V$, let $\pi$ be a smooth representation of
$\operatorname{GL}_N(k)$ on $W$, and let $B$ be an $\omega_\mu$-equivariant bilinear map
(Definition~\ref{gal-satake-weil-equivariant}). Let $e=e_r$, and let $P$ be the group of block upper
triangular matrices $p=\left(\begin{smallmatrix}a&b\\0&d\end{smallmatrix}\right)$ with
$a\in\operatorname{GL}_r(k)$ and $d\in\operatorname{GL}_t(k)$, with $\delta_P(p)=|\det a|^t|\det d|^{-r}$.
Then:
\begin{enumerate}
\item there is a unique linear map $\mathrm{ev}:W\to V$ such that for every $w\in W$ there is
  $c\in k^\times$ with $B(\Phi,w)=\bigl(\int\Phi\,dx\bigr)\,\mathrm{ev}(w)$ for every $\Phi\in\mathcal S_{N,r}$
  that vanishes outside $D_c(e)$;
\item if $\mathrm{ev}=0$, then $B$ vanishes from rank $r$;
\item $\mathrm{ev}(\pi(p)w)=\delta_P(p)^{1/2}\mu(\det d)\,\sigma(a)\,\mathrm{ev}(w)$ for $p\in P$ and
  $w\in W$.
\end{enumerate}
\end{lemma}
\leanref{Development/WeilBound/RepresentationTheoryGL/PointAveraging/MinkowskiBalls/MeasurableSetBallM.lean}{PositivityTheory.Rg.Ming.germ}
\leanref{Development/WeilBound/CoreStatements/TemperedRepresentationTheory/GermIntegrals/StepSt.lean}{PositivityTheory.Tls2.Ming3.GermSt}
\leanref{Development/WeilBound/CoreStatements/TemperedRepresentationTheory/GermIntegrals/IsGerm.lean}{PositivityTheory.Tls2.Ming3.IsGerm}
\leanref{Development/WeilBound/RepresentationTheoryGL/PointAveraging/MinkowskiBalls/GermUnique.lean}{PositivityTheory.Rg.Ming.uniq}
\leanref{Development/WeilBound/CoreStatements/TemperedRepresentationTheory/GermIntegrals/StepSt.lean}{PositivityTheory.Tls2.Ming3.UniqSt}
\leanref{Development/WeilBound/RepresentationTheoryGL/PointAveraging/MinkowskiBalls/RepInvApply.lean}{PositivityTheory.Rg.Ming.nz}
\leanref{Development/WeilBound/CoreStatements/TemperedRepresentationTheory/GermIntegrals/StepSt.lean}{PositivityTheory.Tls2.Ming3.NzSt}
\leanref{Development/WeilBound/CoreStatements/TemperedRepresentationTheory/GermIntegrals/RealPowerProductEqOne.lean}{PositivityTheory.Tls2.Ming3.law}
\leanref{Development/WeilBound/CoreStatements/TemperedRepresentationTheory/GermIntegrals/StepSt.lean}{PositivityTheory.Tls2.Ming3.LawSt}
\leanref{Development/WeilBound/CoreStatements/TemperedRepresentationTheory/GermIntegrals/IsGerm.lean}{PositivityTheory.Tls2.Ming3.ParLaw}
\leanref{Development/WeilBound/CoreStatements/TemperedRepresentationTheory/GermIntegrals/Imports.lean}{PositivityTheory.Tls2.Ming3.cov}

\begin{proof}
The facts about the sets $D_c(x_0)$ in the proof of Lemma~\ref{gal-satake-weil-rank-step} apply; in
particular $\int\mathbf 1_{D_c(x_0)}dx$ is nonzero and does not depend on $x_0$.

(1) Let $w\in W$. As $\pi$ is smooth and $\mu$ is locally constant, there is $c$ with $|c|<1$ such that every
$h\in\operatorname{GL}_N(\mathcal O)$ with entries within $|c|$ of those of $1_N$ fixes $w$ and has
$\mu(\det h)=1$. Let $x\in D_c(e)$, and let $h$ be the matrix whose first $r$ columns are those of $x$ and
whose other columns are those of $1_N$. Then $h$ is such a matrix, $he=x$ and $|\det h|=1$, so
$\mathbf 1_{D_{c'}(x)}=\omega_\mu(1,h)\mathbf 1_{D_{c'}(e)}$ for every $c'$, and
$B(\mathbf 1_{D_{c'}(x)},w)=B(\mathbf 1_{D_{c'}(e)},w)$ by equivariance. Let $\Phi$ vanish outside $D_c(e)$.
For $|c'|$ small, $\Phi=\sum_{x\in X}\Phi(x)\mathbf 1_{D_{c'}(x)}$ with $X\subset D_c(e)$ finite, and so
\[
 B(\Phi,w)=\Bigl(\sum_{x\in X}\Phi(x)\Bigr)B(\mathbf 1_{D_{c'}(e)},w)
 =\Bigl(\int\Phi\,dx\Bigr)\operatorname{vol}(D_{c'}(e))^{-1}B(\mathbf 1_{D_{c'}(e)},w).
\]
Choosing $c'$ small enough for $\Phi$ and for $\mathbf 1_{D_c(e)}$, and applying the display to
$\mathbf 1_{D_c(e)}$, shows that the last factor is
$v=\operatorname{vol}(D_c(e))^{-1}B(\mathbf 1_{D_c(e)},w)$. So $v$ has the property required of
$\mathrm{ev}(w)$. If $v_1$ and $v_2$ both have it, with $c_1$ and $c_2$, applying it to
$\mathbf 1_{D_c(e)}$ with $|c|\leq|c_1|,|c_2|$ gives $v_1=v_2$, as the integral is nonzero. Linearity
follows from uniqueness: if $v_1$ and $v_2$ serve for $w_1$ and $w_2$, then $av_1+v_2$ serves for
$aw_1+w_2$, with the smaller $c$.

(2) Suppose that $\mathrm{ev}=0$, and let $\Phi$ vanish at all matrices of rank less than $r$, and $w\in W$.
As in the proof of Lemma~\ref{gal-satake-weil-rank-step}, it suffices to find, for each $x$ with
$\Phi(x)\neq0$, an open neighbourhood $U$ of $x$ with $B(\Psi,w)=0$ for every $\Psi\in\mathcal S_{N,r}$ that
vanishes outside $U$. Such an $x$ has rank $r$, so $x=heg$ with $h\in\operatorname{GL}_N(k)$ and
$g\in\operatorname{GL}_r(k)$. Let $c$ be given by (1) for the vector $\pi(h^{-1})w$, and let $U=hD_c(e)g$.
If $\Psi$ vanishes outside $U$, then $\Psi'=\omega_\mu(g,h^{-1})\Psi$ vanishes outside $D_c(e)$ and
$\Psi=\omega_\mu(g^{-1},h)\Psi'$, so
\[
 B(\Psi,w)=B\bigl(\omega_\mu(g^{-1},h)\Psi',\pi(h)\pi(h^{-1})w\bigr)=\sigma(g^{-1})B(\Psi',\pi(h^{-1})w)
 =\Bigl(\int\Psi'\,dx\Bigr)\sigma(g^{-1})\,\mathrm{ev}(\pi(h^{-1})w)=0.
\]

(3) Let $p\in P$, so that $pe=ea$. Let $c_0$ be given by (1) for $w$, and choose $c$ with
$p^{-1}D_c(e)a\subset D_{c_0}(e)$, which exists as $p^{-1}ea=e$. Let $\Phi'$ vanish outside $D_c(e)$, and let
$\Phi=\omega_\mu(a^{-1},p^{-1})\Phi'$, which vanishes outside $p^{-1}D_c(e)a$; then
$\Phi'=\omega_\mu(a,p)\Phi$, and
\[
 B(\Phi',\pi(p)w)=\sigma(a)B(\Phi,w)=\Bigl(\int\Phi\,dx\Bigr)\sigma(a)\,\mathrm{ev}(w).
\]
The substitution $x\mapsto hxg$ multiplies $dx$ by $|\det h|^r|\det g|^N$. With $\det p=\det a\det d$ and
$N=r+t$ this gives
\[
 \int\Phi'\,dx=\mu(\det a)\mu(\det p)^{-1}|\det a|^{N/2}|\det p|^{-r/2}\,|\det p|^r|\det a|^{-N}\int\Phi\,dx
 =\mu(\det d)^{-1}\delta_P(p)^{-1/2}\int\Phi\,dx.
\]
Hence $B(\Phi',\pi(p)w)=\bigl(\int\Phi'\,dx\bigr)\delta_P(p)^{1/2}\mu(\det d)\sigma(a)\mathrm{ev}(w)$ for all
such $\Phi'$, and (3) follows from the uniqueness in (1).
\end{proof}

\begin{lemma}\label{gal-satake-weil-lift}
Let $k$ be a non-archimedean local field, let $0\leq r<t$ and $N=r+t$, let $\mu$ be a unitary smooth
character of $k^\times$, let $\sigma$ be a unitarizable representation of $\operatorname{GL}_r(k)$, and let
$\pi$ be an irreducible smooth representation of $\operatorname{GL}_N(k)$. Suppose that the normalized
parabolic induction $\sigma\times(\mu\circ\det)$ of $\sigma\otimes(\mu\circ\det)$ from the standard
parabolic subgroup with Levi factor $\operatorname{GL}_r\times\operatorname{GL}_t$ is irreducible, and that
there is a nonzero $\omega_\mu$-equivariant bilinear map $B$ (Definition~\ref{gal-satake-weil-equivariant}).
Then $\pi\cong\sigma\times(\mu\circ\det)$.
\end{lemma}
\leanref{Development/AutomorphicArithmetic/ThetaCorrespondence/FockModels/MingZel.lean}{UpperTheory.Lit.HoweDuality.lift_of_irr}
\leanref{Development/WeilBound/CoreStatements/TemperedRepresentationTheory/GermIntegrals/LFrob.lean}{PositivityTheory.Tls2.Ming3.frob}
\leanref{Development/WeilBound/CoreStatements/TemperedRepresentationTheory/GermIntegrals/StepSt.lean}{PositivityTheory.Tls2.Ming3.FrobSt}
\leanref{Development/AbelianVarietyTheory/Definitions/GaloisGeometry/IsParahoricVec.lean}{GlobalTheory.ThetaGL.ThetaPair}

\begin{proof}
Let $\mathrm{ev}$ be the linear map of Lemma~\ref{gal-satake-weil-germ}(1). If $\mathrm{ev}=0$, then $B$
vanishes from rank $r$ by Lemma~\ref{gal-satake-weil-germ}(2), and Lemma~\ref{gal-satake-weil-rank-step},
applied for $j=r-1,r-2,\ldots,0$, shows that $B$ vanishes from rank $0$, that is, $B=0$; each such $j$
satisfies $j<r$ and $2j<r+t=N$. So $\mathrm{ev}\neq0$. For $w$ in the space of $\pi$ let
$f_w(g)=\mathrm{ev}(\pi(g)w)$. By Lemma~\ref{gal-satake-weil-germ}(3),
$f_w(pg)=\delta_P(p)^{1/2}\mu(\det d)\sigma(a)f_w(g)$ for
$p=\left(\begin{smallmatrix}a&b\\0&d\end{smallmatrix}\right)$ in the group $P$ of that lemma, and $f_w$ is
invariant under right translation by the stabilizer of $w$, which is open. So $w\mapsto f_w$ is an
intertwining map from $\pi$ to $\sigma\times(\mu\circ\det)$, and it is nonzero as $f_w(1)=\mathrm{ev}(w)$.
Its kernel and image are subrepresentations, and both representations are irreducible, so it is an
isomorphism.
\end{proof}

\begin{lemma}\label{gal-larger-split-satake}
Let $E$ be a CM field with maximal totally real subfield $F$. Let $n\geq1$ and $d>2n$, let
$\chi_X,\chi_W$ be unitary Hecke characters of $E$, and let $U$ and $U'$ be automorphic representations of
$U_n$ and $U_d$ (Definition~\ref{gal-definite-unitary}) that satisfy condition (2) of
Theorem~\ref{gal-theta-lift-large-rank} for $\chi_X$ and $\chi_W$. Let $w$ be a place of $E$ split over $F$
at which $\chi_X$ and $\chi_W$ are unramified, and suppose that $U$ has Satake parameters $B$ at $w$. Then
$U'$ has Satake parameters
\[
 \chi_X(\varpi_w)\chi_W(\varpi_w)^{-1}B+\{\chi_X(\varpi_w)q_w^{(d-n-1)/2-i}:0\leq i<d-n\}
\]
at $w$.
\end{lemma}
\leanref{Development/AutomorphicArithmetic/Foundations/GaloisRepresentations/ThetaParseval/SatakeLifts/CompUnit.lean}{UpperTheory.Gal.TP.Loc.LiftSatU}
\leanref{Development/AutomorphicArithmetic/Foundations/GaloisRepresentations/ThetaParseval/SatakeLifts/LiftSatUPf.lean}{UpperTheory.Gal.TP.Loc.liftSatU}
\leanref{Development/AutomorphicArithmetic/Foundations/GaloisRepresentations/ThetaParseval/SatakeLifts/LiftSatUHolds.lean}{UpperTheory.Gal.TP.Loc.liftSatU_holds}
\leanref{Development/AutomorphicArithmetic/Foundations/GaloisRepresentations/ThetaParseval/SatakeLifts/CompUnit.lean}{UpperTheory.Gal.TP.Loc.PairSat}
\leanref{Development/AutomorphicArithmetic/Foundations/GaloisRepresentations/ThetaParseval/SatakeLifts/SrcApply.lean}{UpperTheory.Gal.TP.Loc.LiftIso}
\leanref{Development/AutomorphicArithmetic/Foundations/GaloisRepresentations/ThetaParseval/SatakeLifts/SrcApply.lean}{UpperTheory.Gal.TP.Loc.LiftSat}
\leanref{Development/AutomorphicArithmetic/ThetaCorrespondence/ThetaPairLifts/LocalTypeDictionary/ThetaLiftGeneralLinear.lean}{GlobalTheory.T.theta_lift_general_linear}
\leanref{Development/AutomorphicArithmetic/Foundations/GaloisRepresentations/ThetaParseval/SatakeLifts/CompUnitPf.lean}{UpperTheory.Gal.TP.Loc.compUnit}

\begin{proof}
Let $v$ be the place of $F$ under $w$ and $t:E\to F_v$ the $F$-embedding that defines $w$; it identifies
$E_w$ with $F_v$ and $U_{a,w}$ with $\operatorname{GL}_a(F_v)$. An automorphic representation of $U_a$ has
a local component at every split place, unique up to isomorphism, by the tensor product theorem
\cite{Fla79}; let $\tau'$ be the local component of $U'$ at $w$, and $\tau$ a local component of $U$ at $w$
with Satake parameters $B$. The representation $\tau$ is unitarizable: $U$ consists of functions on the
compact space $U_n(F)\backslash U_n$, the $L^2$ inner product on $U$ is positive definite and invariant
under right translation, and it pulls back along the map $\tau\to U$, which is injective since $\tau$ is
irreducible.

By condition (2) of Theorem~\ref{gal-theta-lift-large-rank} there is a nonzero map $\beta$ from pairs
$(\Phi,x)$ to $\tau'$, linear in $\Phi$, conjugate-linear in $x$ and equivariant for the twisted Weil
representation $\omega$ on Schwartz--Bruhat functions on $M_{n\times d}(F_v)$. Composing with the
conjugate-linear isomorphism of $\tau$ onto its smooth contragredient $\tau^\vee$ given by the invariant inner
product turns $\beta$ into a nonzero bilinear map $\beta'$ on $\omega\times\tau^\vee$ with values in
$\tau'$ and equivariant for $\operatorname{GL}_n(F_v)\times\operatorname{GL}_d(F_v)$:
$\beta'(\omega(g,h)\Phi,\tau^\vee(g)x)=\tau'(h)\beta'(\Phi,x)$. So $\tau'$ is a theta lift of
$\tau^\vee$ for the dual pair $(\operatorname{GL}_n,\operatorname{GL}_d)$ of type II, with $d-n>n$. Let
$P$ be the standard parabolic subgroup of $\operatorname{GL}_d$ with Levi factor
$\operatorname{GL}_n\times\operatorname{GL}_{d-n}$. Alpöge identifies the smooth contragredient
$\tau'^\vee$ of $\tau'$ with an induced representation and then inverts Satake parameters, in three steps.

\emph{Irreducibility.} Let $\rho$ be an irreducible unitarizable representation of
$\operatorname{GL}_n(F_v)$, $\mu$ a unitary character of $F_v^\times$, and $I_\rho$ the normalized parabolic
induction of $\rho\otimes(\mu\circ\det)$ from $P$. Then $I_\rho$ is irreducible. It is induced from a
unitarizable representation, so it is unitarizable and a direct sum of irreducible representations, and it
suffices that its endomorphisms are scalars. By Frobenius reciprocity they correspond to the maps from the
Jacquet module of $I_\rho$ along $P$ to $\rho\otimes(\mu\circ\det)$. By the geometric lemma of Bernstein
and Zelevinsky \cite{BZ77}, that Jacquet module has a filtration whose cell for the identity contributes
$\rho\otimes(\mu\circ\det)$ once, and it suffices that no other cell contributes it. Another cell would
make $\rho$ a subquotient of a product $\rho_1\times\rho_2$ with $\rho_2$ a constituent of a proper Jacquet
module of the character $\mu\circ\det$ of $\operatorname{GL}_{d-n}(F_v)$, whose supercuspidal support
contains $\mu|\cdot|^{-(d-n-1)/2}$. This character would then lie in the supercuspidal support of $\rho$,
which is the union of those of $\rho_1$ and $\rho_2$ \cite{BZ77}. That contradicts
Lemma~\ref{gal-satake-unitarizable-support}, with $r=n$ and $t=d-n$, since $\rho$ is unitarizable and $\mu$
is unitary.
% Lean: UpperTheory.Lit.Mix.top: ss (unitarizable and Jacquet multiplicity at most 1), cellone (geometric
% lemma, jacquet_multiplicity_single_cell), detjac (Jacquet modules of a character) and nolow.

\emph{The contragredient of $\tau'$.} Let $\sigma'$ be the representation
$g\mapsto(\chi_{X,w}\chi_{W,w}^{-1})(\det g)\,\tau^\vee({}^tg^{-1})$ of $\operatorname{GL}_n(F_v)$ on the
space of $\tau^\vee$. It is irreducible and unitarizable, as $\tau^\vee$ is, and so is its smooth
contragredient $\sigma'^\vee$. Let $\omega'=\omega_{\chi_{X,w}^{-1}}$ be the representation of
$\operatorname{GL}_n(F_v)\times\operatorname{GL}_d(F_v)$ on Schwartz--Bruhat functions on
$M_{d\times n}(F_v)$ of Definition~\ref{gal-satake-weil-equivariant}, with $N=d$ and $r=n$, and for such a
function $\Psi$ let $\Psi^t(Y)=\Psi({}^tY)$ for $Y\in M_{n\times d}(F_v)$. Comparing the formulas gives
\[
 \omega({}^tg^{-1},h)\Psi^t=(\chi_{X,w}\chi_{W,w}^{-1})(\det g)\,(\omega'(g,h)\Psi)^t.
\]
For a smooth linear form $\lambda$ on $\tau'$ let $\beta''(\Psi,\lambda)$ be the linear form
$x\mapsto\lambda(\beta'(\Psi^t,x))$ on the space of $\sigma'$. The identity and the equivariance of $\beta'$
give $\beta''(\omega'(g,h)\Psi,\tau'^\vee(h)\lambda)=\sigma'^\vee(g)\beta''(\Psi,\lambda)$. In particular
$\beta''(\Psi,\lambda)$ is fixed by $\sigma'^\vee$ of the stabilizer of $\Psi$, which is open, so it is
smooth. Thus $\beta''$ is an $\omega'$-equivariant bilinear map, for the representations $\sigma'^\vee$ of
$\operatorname{GL}_n(F_v)$ and $\tau'^\vee$ of $\operatorname{GL}_d(F_v)$. It is nonzero, since
$\beta'\neq0$ and the smooth linear forms on $\tau'$ separate its vectors. The first step, with
$\sigma'^\vee$ and $\chi_{X,w}^{-1}$ in place of $\rho$ and $\mu$, shows that
$\sigma'^\vee\times(\chi_{X,w}^{-1}\circ\det)$ is irreducible, and $\tau'^\vee$ is irreducible as $\tau'$
is. Lemma~\ref{gal-satake-weil-lift}, with $r=n$, $t=d-n$ and $\mu=\chi_{X,w}^{-1}$, gives
$\tau'^\vee\cong\sigma'^\vee\times(\chi_{X,w}^{-1}\circ\det)$.
% Lean: B is UpperTheory.Gal.TP.Loc.pairOfLift (LeftSmooth, KeepsSB, ModelDict, Separates); LiftIso
% applies theta_lift_general_linear = HoweDuality.thetaLift_gl_of_zelevinsky_classification (Mix.top
% and lift_of_irr) to contra (src (chi_X chi_W^-1) tau^vee) and chi_X^-1.

\emph{Satake parameters.} Write $\alpha=\chi_X(\varpi_w)\chi_W(\varpi_w)^{-1}$ and, for $x\in\mathbf
C^\times$, $\Sigma(x)=\{x\,q_w^{(d-n-1)/2-i}:0\leq i<d-n\}$, a multiset. Alpöge uses five facts about Satake
parameters: those of the smooth contragredient of an irreducible unitarizable representation are the
inverses of its own; so are those of $g\mapsto\rho({}^tg^{-1})$; those of a twist by $\nu\circ\det$ with
$\nu$ unramified are multiplied by $\nu(\varpi_w)$; those of the character $\nu\circ\det$ of
$\operatorname{GL}_{d-n}(F_v)$, for $\nu$ unramified, are $\Sigma(\nu(\varpi_w))$; and those of the
normalized parabolic induction from $P$ of $\rho_1\otimes\rho_2$, with $\rho_1$ and $\rho_2$ irreducible,
are the union of those of $\rho_1$ and $\rho_2$. So $\tau^\vee$ has Satake parameters $B^{-1}$, $\sigma'$
has $\alpha B$, $\sigma'^\vee$ has $\alpha^{-1}B^{-1}$, and $\tau'^\vee$ has
$\alpha^{-1}B^{-1}+\Sigma(\chi_X(\varpi_w)^{-1})$. The representation $\tau'^\vee$ is unitarizable, being
isomorphic to a representation induced from a unitarizable one, and $\tau'$ is isomorphic to the smooth
contragredient of $\tau'^\vee$, since $\tau'$ is irreducible and smooth. So the Satake parameters of
$\tau'$ are the inverses of those of $\tau'^\vee$, which are $\alpha B$ together with
$\Sigma(\chi_X(\varpi_w)^{-1})^{-1}=\Sigma(\chi_X(\varpi_w))$, the last equality by the substitution
$i\mapsto d-n-1-i$.
% Lean: LiftSat (liftSat_of) with SatContra, SatInvT, SatTwist, SatDet, SatProd and SegInv; the
% contragredient of tau is taken in pairSat, and the double contragredient is ContraIrr.
\end{proof}

\begin{remark}\label{gal-larger-split-satake-minguez}
The proof of Lemma~\ref{gal-larger-split-satake} can also identify $\tau'$ itself. Let
$\sigma=\tau\otimes(\chi_{X,w}\chi_{W,w}^{-1}\circ\det)$, and let $I$ be the normalized parabolic induction
of $\sigma\otimes(\chi_{X,w}\circ\det)$ from $P$. Smooth contragredients commute with parabolic induction
\cite{BZ77}, so $\tau'\cong\sigma'\times(\chi_{X,w}\circ\det)$, and $\sigma'\cong\sigma$, since the
representation $g\mapsto\tau^\vee({}^tg^{-1})$ is isomorphic to $\tau$ by the theorem of Gelfand and
Kazhdan \cite{GK75}. So $\tau'\cong I$, and the Satake parameters of $\tau'$ are those of $\sigma$ together
with those of $\chi_{X,w}\circ\det$. The identification of $\tau'$ with $I$ also follows from the
literature. The explicit theta correspondence for dual pairs of type II \cite{Min08} describes the theta
lift of $\tau^\vee$, and Bernstein's theorem that parabolic induction from irreducible unitarizable
representations of general linear groups is irreducible \cite{Ber84} shows that $I$ is irreducible.
\end{remark}

\begin{lemma}\label{gal-larger-lift-enlargement}
Let $E$ be a CM field with maximal totally real subfield $F$, complex conjugation $c$ and a good base with
embeddings $\tau_v$ (Definition~\ref{gal-good-base}). Let $1\leq n<t$, let $\chi_X,\chi_W$ be unitary Hecke
characters of $E$ and $m_X(v),m_W(v)$ integers, let $U$ be an automorphic representation of $U_n$ with
weight $\lambda_\tau$ at every $\tau:E\to\mathbf C$, and let $U'$ be an automorphic representation of
$U_{n+t}$ that satisfies conditions (1) and (2) of Theorem~\ref{gal-theta-lift-large-rank} for $d=n+t$,
these characters and integers, and $\lambda(v)=\lambda_{\tau_v}$. Let $\kappa,\psi_0,\ldots,\psi_{t-1}$ be
Hecke characters of $E$, algebraic of infinity types $m,p_0,\ldots,p_{t-1}$ and attuned to
$(\chi_X,\chi_W)$ in rank $(n,t)$ (Definition~\ref{gal-larger-attuned}). Then there are weights
$\lambda'_\tau$ such that $(U',\lambda',\psi_0,\ldots,\psi_{t-1})$ is an enlargement of $(U,\lambda)$ along
$\kappa$ in rank $n+t$ (Definition~\ref{gal-larger-enlargement}).
\end{lemma}
\leanref{Development/AutomorphicArithmetic/Foundations/GaloisRepresentations/ThetaParseval/SplitLineSubtraction/TwoUpOKOf.lean}{UpperTheory.Gal.TP.Sub.UpOf}
\leanref{Development/AutomorphicArithmetic/Foundations/GaloisRepresentations/ThetaParseval/SplitLineSubtraction/UpOfHolds.lean}{UpperTheory.Gal.TP.Sub.upOf_holds}
\leanref{Development/AutomorphicArithmetic/Foundations/GaloisRepresentations/ThetaParseval/SplitLineSubtraction/UpMain.lean}{UpperTheory.Gal.TP.Sub.upOf}
\leanref{Development/AutomorphicArithmetic/Foundations/GaloisRepresentations/ThetaParseval/SplitLineSubtraction/TwoUpOKOf.lean}{UpperTheory.Gal.TP.Sub.WtId}
\leanref{Development/AutomorphicArithmetic/Foundations/GaloisRepresentations/ThetaParseval/SplitLineSubtraction/MapUnivAdd.lean}{UpperTheory.Gal.TP.Sub.wtId_holds}
\leanref{Development/AutomorphicArithmetic/Foundations/GaloisRepresentations/ThetaParseval/SplitLineSubtraction/TwoUpOKOf.lean}{UpperTheory.Gal.TP.Sub.SatId}
\leanref{Development/AutomorphicArithmetic/Foundations/GaloisRepresentations/ThetaParseval/SplitLineSubtraction/PSatId.lean}{UpperTheory.Gal.TP.Sub.satId_holds}

\begin{proof}
\emph{Weights.} Every embedding $E\to\mathbf C$ is $\tau_v$ or $c\tau_v$ for one infinite place $v$ of $F$,
and these $2[F:\mathbf Q]$ embeddings are distinct. By condition (1) of
Theorem~\ref{gal-theta-lift-large-rank}, $U'$ has weight $\lambda'(v)$ at $\tau_v$, with
$2\lambda'(v)_k=2\lambda_{\tau_v,k}-(n+t)-m_W(v)+n+m_X(v)$ for $k<n$ and $2\lambda'(v)_k=n+m_X(v)$ for
$n\leq k<n+t$. An automorphic representation of $U_a$ with weight $\mu$ at $\tau$ has weight $\mu^\vee$ at
$c\tau$, where $\mu^\vee_k=-\mu_{a-1-k}$, since the $c\tau$-factor at infinity acts through the complex
conjugate of the $\tau$-factor; and its weight at each embedding is unique. Hence
$\lambda_{c\tau_v}=\lambda_{\tau_v}^\vee$, and $U'$ has weight $\lambda'_\tau$ at every $\tau$ for
$\lambda'_{\tau_v}=\lambda'(v)$ and $\lambda'_{c\tau_v}=\lambda'(v)^\vee$.

At $\tau=\tau_v$: for $k<n$, condition (1) of Definition~\ref{gal-larger-attuned} gives
$k-\lambda'_{\tau,k}=k-\lambda_{\tau,k}+(t+m_W(v)-m_X(v))/2=k-\lambda_{\tau,k}+m_\tau$; for $k=n+i$ with
$0\leq i<t$, $k-\lambda'_{\tau,k}=(n+2i-m_X(v))/2$, which is $-p_{t-1-i,\tau}$ by condition (2). This is
(1)(a) of Definition~\ref{gal-larger-enlargement} at $\tau_v$. At $\tau=c\tau_v$ the entries
$k-\lambda'_{\tau,k}$ are $n+t-1-k'+\lambda'(v)_{k'}$ for $0\leq k'<n+t$. For $k'<n$ this is
$n-1-k'+\lambda_{\tau_v,k'}+(t+m_X(v)-m_W(v))/2$, which is the entry of index $n-1-k'$ of
$\{k-\lambda_{c\tau_v,k}\}$ shifted by $m_{c\tau_v}$; for $k'=n+i$ it is $t-1-i+(n+m_X(v))/2=-p_{i,c\tau_v}$.
This is (1)(a) at $c\tau_v$.

\emph{Satake parameters.} Let $w$ be a place split over $F$ at which $\kappa,\psi_0,\ldots,\psi_{t-1}$ are
unramified, and let $U$ have Satake parameters $B$ at $w$. By condition (3) of
Definition~\ref{gal-larger-attuned}, $\chi_X$ and $\chi_W$ are unramified at $w$, and since $n+t>2n$,
Lemma~\ref{gal-larger-split-satake} shows that $U'$ has Satake parameters
$B'=\chi_X(\varpi_w)\chi_W(\varpi_w)^{-1}B+\{\chi_X(\varpi_w)q_w^{(t-1)/2-i}:0\leq i<t\}$ at $w$. By
condition (3), $q_w^{(n+t-1)/2}\chi_X(\varpi_w)\chi_W(\varpi_w)^{-1}=\kappa(\varpi_w)^{-1}q_w^{(n-1)/2}$ and
$q_w^{(n+t-1)/2}\chi_X(\varpi_w)q_w^{(t-1)/2-i}=\chi_X(\varpi_w)q_w^{(n+2t-2)/2-i}=\psi_i(\varpi_w)$. This is
(1)(b) of Definition~\ref{gal-larger-enlargement}.
\end{proof}

\begin{proposition}\label{gal-larger-pairs}
Let $E$ be a CM field with maximal totally real subfield $F$. Suppose that $E$ contains $\sqrt{-D}$ for some
integer $D>0$, that $[F:\mathbf Q]$ is even, that $E/F$ is unramified at every finite place, and that every
finite place $v$ of $F$ inert
in $E$ satisfies $v\nmid2$, $N(v)>3$ and $v\nmid\mathfrak d_F$, where $\mathfrak d_F$ is the different of
$F$. Let $n\geq1$ be even, let $U$ be an automorphic representation of $U_n$ with weight $\lambda_\tau$ at
every $\tau:E\to\mathbf C$, and let $w_0$ be a place of $E$ split over $F$. Then there is a pair of
enlargements of $(U,\lambda)$ at $w_0$ (Definition~\ref{gal-larger-enlargement}).
\end{proposition}
\leanref{Development/AutomorphicArithmetic/Foundations/GaloisRepresentations/ThetaParseval/AllCoreOf.lean}{UpperTheory.Gal.TP.TwoUpOK}
\leanref{Development/AutomorphicArithmetic/Foundations/GaloisRepresentations/ThetaParseval/AllCoreOf.lean}{UpperTheory.Gal.TP.InertOK}
\leanref{Development/AutomorphicArithmetic/Foundations/GaloisRepresentations/ThetaParseval/TwoUpOKOfSteps.lean}{UpperTheory.Gal.TP.twoUpOK_of_steps}
\leanref{Development/AutomorphicArithmetic/Foundations/GaloisRepresentations/ThetaParseval/TwoUpOKOfSteps.lean}{UpperTheory.Gal.TP.twoUpOK_of_k}

\begin{proof}
Lemma~\ref{gal-larger-good-base} gives a good base for $E$ with embeddings $\tau_v$, a standard Hermitian
space $X$ of rank $n$, and for every $d\geq1$ a scalar skew-Hermitian space $W_d$ of rank $d$.
Theorem~\ref{gal-theta-lift-large-rank} for $X$ gives an integer $d_0$. Put $t=2(d_0+n)+1$; then $t$ is
odd, $n<t$, $n+t\geq d_0$, and $n+t$ is odd since $n$ is even. Let $b$ be an integer with
$b\leq\lambda_{\tau_v,k}$ for all $v$ and $k$. Lemma~\ref{gal-larger-characters}, with the $\tau_v$, $n$,
$t$, the integer $2b-(n+t)$ and $w_0$, gives $\kappa$ and, for $j=0,1$, characters
$\chi_{X,j},\chi_{W,j}$, integers $m_{X,j}(v),m_{W,j}(v)$, and Hecke characters
$\psi_{j,0},\ldots,\psi_{j,t-1}$ attuned to $(\chi_{X,j},\chi_{W,j})$ in rank $(n,t)$.

For each $j$, Theorem~\ref{gal-theta-lift-large-rank} applies with $d=n+t$, $W=W_{n+t}$,
$\chi_X=\chi_{X,j}$, $\chi_W=\chi_{W,j}$ and $\lambda(v)=\lambda_{\tau_v}$: the restrictions to
$\mathbf A_F^\times$ and the infinity types are those of Lemma~\ref{gal-larger-characters}(1), and
$(n+t)+m_{W,j}(v)\leq(n+t)+2b-(n+t)=2b\leq2\lambda_{\tau_v,k}$ for all $v$ and $k$. It gives an
automorphic representation $U'_j$ of $U_{n+t}$, and Lemma~\ref{gal-larger-lift-enlargement} gives weights
$\lambda'_j$ such that $(U'_j,\lambda'_j,\psi_{j,0},\ldots,\psi_{j,t-1})$ is an enlargement of $(U,\lambda)$
along $\kappa$ in rank $n+t$. By Lemma~\ref{gal-larger-characters}(3), $t$, $\kappa$ and the enlargements
for $j=0$ and $j=1$, numbered $i=1$ and $i=2$, form a pair of enlargements at $w_0$. The hypothesis that
$[F:\mathbf Q]$ is even is not used.
\end{proof}

\begin{definition}\label{gal-constituent-data}
Let $E$ be a number field, $\Omega$ an algebraic closure of $E$, $\ell$ a prime,
$\iota:\bar{\mathbf Q}_\ell\to\mathbf C$ a field isomorphism, $a\geq1$ and $w$ a finite place of $E$. Let
$S$ be a multiset of $a$ nonzero complex numbers and $J=(J_\tau)_{\tau:E\to\bar{\mathbf Q}_\ell}$ a family of
multisets of $a$ integers. For a valuation ring $A$ of $\Omega$ over $w$, let
$\mathrm{Un}_A(S)\in\mathrm{WD}(A)$ be the monodromy-free unramified representation on which a geometric
Frobenius element acts with eigenvalues $S$. A continuous homomorphism
$r:G_E\to\operatorname{GL}_a(\bar{\mathbf Q}_\ell)$ has \emph{spherical data} $(S,J)$ at $w$ if:
\begin{enumerate}
\item when $w\nmid\ell$, $r$ has Frobenius eigenvalues $S$ at $w$ (Definition~\ref{gal-weakly-attached});
\item when $w\mid\ell$, for every valuation ring $A$ over $w$, $r|_{D_A}$ is crystalline and de Rham,
  $\mathrm{HT}_\tau(r|_{D_A})=J_\tau$ for every $\tau$ that extends continuously to $E_w$, and
  $\mathrm{WD}_\iota(r|_{D_A})(\sigma)$ and $\mathrm{Un}_A(S)(\sigma)$ have the same characteristic
  polynomial for every $\sigma\in W_A$ (Definition~\ref{gal-p-adic-hodge}).
\end{enumerate}
If $E$ is a CM field, $B$ a multiset of $a$ nonzero complex numbers and $Z=(Z_\tau)_{\tau:E\to\mathbf C}$ a
family of multisets of $a$ complex numbers, then $\mathrm{Un}_A(q_w^{(a-1)/2}B)=\mathrm{Sat}_A(B)$, so $r$
satisfies (unr), (dR), (crys) and (phi) at $w$ with respect to $B$ and $Z$
(Definition~\ref{gal-spherical-conditions}) exactly when it has spherical data $(q_w^{(a-1)/2}B,J)$ at
$w$ with $J_\tau=\{\frac{a-1}{2}-z:z\in Z_{\iota\circ\tau}\}$.
\end{definition}
\leanref{Development/AutomorphicArithmetic/Foundations/GaloisRepresentations/ThetaParseval/SplitLineSubtraction/RowDefOfOdd.lean}{UpperTheory.Gal.TP.Sub.Cl}
\leanref{Development/AutomorphicArithmetic/Foundations/GaloisRepresentations/ThetaParseval/SplitLineSubtraction/RowDefOfOdd.lean}{UpperTheory.Gal.TP.Sub.Above}
\leanref{Development/AutomorphicArithmetic/Foundations/GaloisRepresentations/ThetaParseval/CoreAtIffRowAt.lean}{UpperTheory.Gal.TP.sphG_iff_cl}

\begin{lemma}\label{gal-constituent-rules}
Let $E$ be a number field, $\Omega$ an algebraic closure of $E$, $\ell$ a prime,
$\iota:\bar{\mathbf Q}_\ell\to\mathbf C$ a field isomorphism and $w$ a finite place of $E$. Let $a,b\geq1$,
let $r,r':G_E\to\operatorname{GL}_a(\bar{\mathbf Q}_\ell)$, $s:G_E\to\operatorname{GL}_b(\bar{\mathbf Q}_\ell)$
and $\mu:G_E\to\bar{\mathbf Q}_\ell^\times$ be continuous homomorphisms, let $S$ be a multiset of $b$ nonzero
complex numbers, $J=(J_\tau)$ a family of multisets of $b$ integers, $c$ a nonzero complex number and
$(m_\tau)_{\tau:E\to\bar{\mathbf Q}_\ell}$ a family of integers, and suppose that $\mu$ has spherical data
$(\{c\},(\{m_\tau\})_\tau)$ at $w$ (Definition~\ref{gal-constituent-data}).
\begin{enumerate}
\item If $r\cong r'$, then $r$ and $r'$ have the same spherical data at $w$.
\item If $s\oplus\mu$ has spherical data $(S+\{c\},(J_\tau+\{m_\tau\})_\tau)$ at $w$, then $s$ has spherical
  data $(S,J)$ at $w$.
\item If $s$ has spherical data $(S,J)$ at $w$, then $s\otimes\mu$ has spherical data
  $(cS,(\{x+m_\tau:x\in J_\tau\})_\tau)$ at $w$.
\end{enumerate}
\end{lemma}
\leanref{Development/AutomorphicArithmetic/Foundations/GaloisRepresentations/ThetaParseval/SplitLineSubtraction/SubtractOf.lean}{UpperTheory.Gal.TP.Sub.ClConj}
\leanref{Development/AutomorphicArithmetic/Foundations/GaloisRepresentations/ThetaParseval/SplitLineSubtraction/SubtractOf.lean}{UpperTheory.Gal.TP.Sub.ClSub}
\leanref{Development/AutomorphicArithmetic/Foundations/GaloisRepresentations/ThetaParseval/SplitLineSubtraction/SubtractOf.lean}{UpperTheory.Gal.TP.Sub.ClTw}
\leanref{Development/AutomorphicArithmetic/Foundations/GaloisRepresentations/ThetaParseval/SplitLineSubtraction/SubtractHolds.lean}{UpperTheory.Gal.TP.Sub.clConj}
\leanref{Development/AutomorphicArithmetic/Foundations/GaloisRepresentations/ThetaParseval/SplitLineSubtraction/SubtractHolds.lean}{UpperTheory.Gal.TP.Sub.clSub}
\leanref{Development/AutomorphicArithmetic/Foundations/GaloisRepresentations/ThetaParseval/SplitLineSubtraction/SubtractHolds.lean}{UpperTheory.Gal.TP.Sub.clTw}

\begin{proof}
Away from $\ell$ the data are the triviality on $I_A$ and the characteristic polynomials of
$\iota(r(\sigma^{-1}))$ for $\sigma\in D_A$ acting on the residue field by $z\mapsto z^{q_w}$. Both are
invariant under isomorphism. For $s\oplus\mu$ the characteristic polynomial is the product of those of $s$
and $\mu$, and the factor $X-c$ cancels; the eigenvalues of $(s\otimes\mu)(\sigma^{-1})$ are those of
$s(\sigma^{-1})$ multiplied by $\mu(\sigma^{-1})$, whose image under $\iota$ is $c$.

Over $\ell$, being crystalline or de Rham is invariant under isomorphism, passes to direct summands and is
preserved by twisting with a crystalline character \cite{Fon94}; Hodge--Tate multisets add over direct sums
and are shifted by $m_\tau$ under the twist by $\mu$. De Rham representations are potentially semistable
\cite{Ber02}, so $\mathrm{WD}_\iota$ is computed from $D_{\mathrm{pst}}$, which is compatible with
isomorphisms, direct sums and tensor products. So the characteristic polynomials of
$\mathrm{WD}_\iota(\cdot)(\sigma)$ multiply over direct sums, and for $s\otimes\mu$ the eigenvalues are
multiplied by the scalar $\mathrm{WD}_\iota(\mu|_{D_A})(\sigma)$, which is $\mathrm{Un}_A(\{c\})(\sigma)$.
Since $\mathrm{Un}_A(S+\{c\})=\mathrm{Un}_A(S)\oplus\mathrm{Un}_A(\{c\})$ and
$\mathrm{Un}_A(cS)=\mathrm{Un}_A(S)\otimes\mathrm{Un}_A(\{c\})$, this gives (2) and (3).
\end{proof}

\begin{lemma}\label{gal-constituent-characters}
Let $E$ be a CM field, $\Omega$ an algebraic closure of $E$, $\ell$ a prime and
$\iota:\bar{\mathbf Q}_\ell\to\mathbf C$ a field isomorphism. Let $\psi$ be a Hecke character of $E$
algebraic of infinity type $p$ (Definition~\ref{gal-larger-enlargement}). Then there is a continuous
character $\chi_\psi:G_E\to\bar{\mathbf Q}_\ell^\times$ that has spherical data
$(\{\psi(\varpi_w)\},(\{-p_{\iota\circ\tau}\})_\tau)$ at every finite place $w$ at which $\psi$ is unramified
(Definition~\ref{gal-constituent-data}). If $\psi'$ is a second Hecke character of $E$ algebraic of some
infinity type and $\chi_\psi=\chi_{\psi'}$, then $\psi=\psi'$.
\end{lemma}
\leanref{Development/AutomorphicArithmetic/Foundations/GaloisRepresentations/ThetaParseval/SplitLineSubtraction/SubtractOf.lean}{UpperTheory.Gal.TP.Sub.CharEx}
\leanref{Development/AutomorphicArithmetic/Foundations/GaloisRepresentations/ThetaParseval/SplitLineSubtraction/PCharEx.lean}{UpperTheory.Gal.TP.Sub.charEx_holds}
\leanref{Development/AutomorphicArithmetic/Foundations/GaloisRepresentations/ThetaParseval/SplitLineSubtraction/SubtractHolds.lean}{UpperTheory.Gal.TP.Sub.charEx}
\leanref{Development/AbelianVarietyTheory/PublishedTheorems/AutomorphicForms₂/ExistsUniqueHcharOfFiniteOrder.lean}{GlobalTheory.Published.HeckeLAdicCharacter}

\begin{proof}
The character $\chi_\psi$ is the $\ell$-adic character attached to $\psi$ by Weil and Serre
\cite[Chapter~II]{Ser68}, transported by $\iota$. At a finite place $w\nmid\ell$ at which $\psi$ is
unramified it is unramified, and $\iota$ of its value at a geometric Frobenius element is $\psi(\varpi_w)$.
At a place $w\mid\ell$ at which $\psi$ is unramified, its restriction to $D_A$ is, through local class field
theory, the product of an unramified character and of a character of the inertia group given by an
algebraic character of $\mathcal O_w^\times$ determined by $p$; so it is crystalline, with Hodge--Tate weight
$-p_{\iota\circ\tau}$ at $\tau$, and its Weil--Deligne representation is $\mathrm{Un}_A(\{\psi(\varpi_w)\})$
\cite[Chapter~III]{Ser68}, \cite{Fon94}. If $\chi_\psi=\chi_{\psi'}$, then $\psi(\varpi_w)=\psi'(\varpi_w)$ at
all but finitely many $w$, so $\psi'\psi^{-1}$ is trivial on $E_w^\times$ for all places $w$ outside a finite
set $Y$ that contains the infinite places. As $E^\times$ is dense in $\prod_{w\in Y}E_w^\times$ by weak
approximation, $\psi'\psi^{-1}$ is trivial.
\end{proof}

\begin{lemma}\label{gal-constituent-peel}
Let $\Gamma$ be a group and $K$ an algebraically closed field of characteristic $0$. Let $n,t,t',m'\geq0$,
let $r:\Gamma\to\operatorname{GL}_{n+t}(K)$ and $r':\Gamma\to\operatorname{GL}_{m'}(K)$ be semisimple
representations, and let $a_1,\ldots,a_t,a'_1,\ldots,a'_{t'}:\Gamma\to K^\times$ be characters with
$a_i\neq a'_j$ for all $i,j$. Suppose that for every $g\in\Gamma$
\[
 \det(X-r(g))\prod_{j=1}^{t'}(X-a'_j(g))=\det(X-r'(g))\prod_{i=1}^{t}(X-a_i(g)).
\]
\begin{enumerate}
\item There is a semisimple $s:\Gamma\to\operatorname{GL}_n(K)$ with
  $r\cong s\oplus a_1\oplus\cdots\oplus a_t$. If $\Gamma$ and $K$ are topological and $r$ is continuous, $s$
  is continuous.
\item Let $E$ be a number field, $\Omega$ an algebraic closure of $E$, $\ell$ a prime,
  $\iota:\bar{\mathbf Q}_\ell\to\mathbf C$ a field isomorphism and $w$ a finite place of $E$, and suppose
  that $\Gamma=G_E$, $K=\bar{\mathbf Q}_\ell$, and $r$ and the $a_i$ are continuous. Let $S$ be a multiset of
  $n$ nonzero complex numbers, $J=(J_\tau)$ a family of multisets of $n$ integers, $c_1,\ldots,c_t$ nonzero
  complex numbers and $m_{i,\tau}$ integers, and suppose that each $a_i$ has spherical data
  $(\{c_i\},(\{m_{i,\tau}\})_\tau)$ at $w$ (Definition~\ref{gal-constituent-data}) and that $r$ has spherical data
  $(S+\{c_1,\ldots,c_t\},(J_\tau+\{m_{1,\tau},\ldots,m_{t,\tau}\})_\tau)$ at $w$. Then $s$ has spherical
  data $(S,J)$ at $w$.
\end{enumerate}
\end{lemma}
\leanref{Development/AutomorphicArithmetic/Foundations/GaloisRepresentations/ThetaParseval/SplitLineSubtraction/SubtractOf.lean}{UpperTheory.Gal.TP.Sub.PeelOne}
\leanref{Development/AutomorphicArithmetic/Foundations/GaloisRepresentations/ThetaParseval/SplitLineSubtraction/SubtractOf.lean}{UpperTheory.Gal.TP.Sub.PeelAll}
\leanref{Development/AutomorphicArithmetic/Foundations/GaloisRepresentations/ThetaParseval/SplitLineSubtraction/PPeelOne.lean}{UpperTheory.Gal.TP.Sub.peelOne_holds}
\leanref{Development/AutomorphicArithmetic/Foundations/GaloisRepresentations/ThetaParseval/SplitLineSubtraction/PPeelAll.lean}{UpperTheory.Gal.TP.Sub.peelAll_holds}
\leanref{Development/AutomorphicArithmetic/Foundations/GaloisRepresentations/ThetaParseval/SplitLineSubtraction/SubtractHolds.lean}{UpperTheory.Gal.TP.Sub.peelAll}

\begin{proof}
By induction on $t$; for $t=0$ take $s=r$. Let $t\geq1$ and $c=a_t$. The representations
$r\oplus a'_1\oplus\cdots\oplus a'_{t'}$ and $r'\oplus a_1\oplus\cdots\oplus a_t$ are semisimple and, by the
hypothesis, have the same characteristic polynomial at every $g$, so they are isomorphic by the
Brauer--Nesbitt theorem. The second has a line on which $\Gamma$ acts by $c$, hence so does the first. Write
a generator of this line as $(v_0,v_1,\ldots,v_{t'})$ with $v_0$ in the space of $r$; then
$a'_j(g)v_j=c(g)v_j$ for all $g$, and $a'_j\neq c$, so $v_j=0$ for $j\geq1$. Hence $r$ has a line on which
$\Gamma$ acts by $c$, and since $r$ is semisimple the line has an invariant complement:
$r\cong s_1\oplus c$ with $s_1$ semisimple. Then $\det(X-r(g))=\det(X-s_1(g))(X-c(g))$, and cancelling the
factor $X-c(g)$ in $K[X]$ gives the hypothesis for $s_1$, $r'$, $a_1,\ldots,a_{t-1}$ and the $a'_j$. The
induction hypothesis gives $s$ with $s_1\cong s\oplus a_1\oplus\cdots\oplus a_{t-1}$, so
$r\cong s\oplus a_1\oplus\cdots\oplus a_t$. If $r$ is continuous, so is
$s\oplus a_1\oplus\cdots\oplus a_t$, which is conjugate to $r$, and so is its block $s$.

For (2), in the induction step Lemma~\ref{gal-constituent-rules}(1) gives the spherical data of $r$ to
$s_1\oplus a_t$, Lemma~\ref{gal-constituent-rules}(2) shows that $s_1$ has spherical data
$(S+\{c_1,\ldots,c_{t-1}\},(J_\tau+\{m_{1,\tau},\ldots,m_{t-1,\tau}\})_\tau)$, and the induction hypothesis
applies to $s_1$.
\end{proof}

\begin{proposition}\label{gal-constituent}
Let $E$ be a CM field with maximal totally real subfield $F$, $\Omega$ an algebraic closure of $E$, $\ell$ a
prime, $\iota:\bar{\mathbf Q}_\ell\to\mathbf C$ a field isomorphism, and $d>0$ an integer with
$\sqrt{-d}\in E$. Suppose that for every odd $a$ and every automorphic representation $U'$ of $U_a$ with
weight $\mu_\tau$ at every $\tau:E\to\mathbf C$ (Definition~\ref{gal-definite-unitary}) there is a
continuous semisimple $r':G_E\to\operatorname{GL}_a(\bar{\mathbf Q}_\ell)$ that, at every $d$-split place
$w$ (Definition~\ref{gal-seed}) and for every $B'$ such that $U'$ has Satake parameters $B'$ at $w$,
satisfies (unr), (dR), (crys) and (phi) with respect to $B'$ and the multisets
$Z_\tau=\{\mu_{\tau,k}+\frac{a-1}{2}-k:0\leq k<a\}$ (Definition~\ref{gal-spherical-conditions}). Let
$n\geq1$, let $U$ be an automorphic representation of $U_n$ with weight $\lambda_\tau$ at every $\tau$, and
suppose that $(U,\lambda)$ has a pair of enlargements at every $d$-split place
(Definition~\ref{gal-larger-enlargement}). Then there is a continuous semisimple
$r:G_E\to\operatorname{GL}_n(\bar{\mathbf Q}_\ell)$ that, at every $d$-split place $w$ and for every $B$ such
that $U$ has Satake parameters $B$ at $w$, satisfies (unr), (dR), (crys) and (phi) with respect to $B$ and
the multisets $\{\lambda_{\tau,k}+\frac{n-1}{2}-k:0\leq k<n\}$.
\end{proposition}
\leanref{Development/AutomorphicArithmetic/Foundations/GaloisRepresentations/ThetaParseval/CoreAtQIffRowQ.lean}{UpperTheory.Gal.TP.SubtractQ}
\leanref{Development/AutomorphicArithmetic/Foundations/GaloisRepresentations/ThetaParseval/CoreAtQIffRowQ.lean}{UpperTheory.Gal.TP.subtractQ_of}
\leanref{Development/AutomorphicArithmetic/Foundations/GaloisRepresentations/ThetaParseval/CoreAtQIffRowQ.lean}{UpperTheory.Gal.TP.atPlaceQ}
\leanref{Development/AutomorphicArithmetic/Foundations/GaloisRepresentations/ThetaParseval/CoreAtQIffRowQ.lean}{UpperTheory.Gal.TP.frobIdQ}
\leanref{Development/AutomorphicArithmetic/Foundations/GaloisRepresentations/ThetaParseval/CoreAtQIffRowQ.lean}{UpperTheory.Gal.TP.weakIsoQ}
\leanref{Development/AutomorphicArithmetic/Foundations/GaloisRepresentations/ThetaParseval/CoreAtQIffRowQ.lean}{UpperTheory.Gal.TP.coreAtQ_iff_rowQ}
\leanref{Development/AutomorphicArithmetic/Foundations/GaloisRepresentations/ThetaParseval/CoreAtQIffRowQ.lean}{UpperTheory.Gal.TP.GoodGuard}
\leanref{Development/AutomorphicArithmetic/Foundations/GaloisRepresentations/IntegralStructures/CrystallineDescent/JoinAll.lean}{UpperTheory.Gal.Int.Core.goodGuard_usPl}
\leanref{Development/AutomorphicArithmetic/Foundations/GaloisRepresentations/ThetaParseval/SplitLineSubtraction/PSatAE.lean}{UpperTheory.Gal.TP.Sub.satAE_holds}
\leanref{Development/AutomorphicArithmetic/Foundations/GaloisRepresentations/IntegralStructures/CrystallineDescent/LimOddOfLimFinOdd.lean}{UpperTheory.Gal.Int.Core.subtractQ_holds}

\begin{proof}
For weights $\mu$ of rank $a$ let $J(\mu)$ be the family $J(\mu)_\tau=\{k-\mu_{\iota\circ\tau,k}:0\leq k<a\}$
indexed by $\tau:E\to\bar{\mathbf Q}_\ell$. By Definition~\ref{gal-constituent-data}, the four conditions
at $w$ with respect to $B$ and the multisets $\{\mu_{\tau,k}+\frac{a-1}{2}-k\}$ say that $r$ has spherical
data $(q_w^{(a-1)/2}B,J(\mu))$ at $w$.

\emph{Density.} For every finite set $Y$ of places of $E$, the geometric Frobenius elements at the $d$-split
places $w\notin Y$ with $w\nmid\ell$ are dense in $G_E$. Indeed, let $w$ be a place of $E$ unramified over
$\mathbf Q$ whose norm is a prime $p$. Then $E_w=\mathbf Q_p$ contains $\sqrt{-d}$, so $-d$ is a square in
$\mathbf Q_p$; and $w$ has residue degree one over the place of $F$ under it, so this place splits in $E$
unless it ramifies. Hence all but finitely many places of prime norm are $d$-split, and the Frobenius
elements at the places of prime norm outside a finite set are dense by the Chebotarev density theorem. In
particular $d$-split places exist.

\emph{A representation for each place.} Let $w_0$ be a $d$-split place, and take a pair of enlargements of
$(U,\lambda)$ at $w_0$: an integer $t$ with $n+t$ odd, a Hecke character $\kappa$ algebraic of infinity
type $m$ and unramified at $w_0$, and enlargements $(U_i,\lambda_i,\psi_{i,1},\ldots,\psi_{i,t})$ for
$i=1,2$, with $\psi_{i,j}$ of infinity type $p_{i,j}$. Let $\chi_\kappa$ and $a_{i,j}=\chi_{\psi_{i,j}}$ be the
characters of Lemma~\ref{gal-constituent-characters}; then $a_{1,j}\neq a_{2,j'}$ for all $j,j'$, since
$\psi_{1,j}\neq\psi_{2,j'}$. As $n+t$ is odd, the hypothesis gives continuous semisimple
$r_i:G_E\to\operatorname{GL}_{n+t}(\bar{\mathbf Q}_\ell)$ with spherical data
$(q_w^{(n+t-1)/2}B_i,J(\lambda_i))$ at every $d$-split place $w$ at which $U_i$ has Satake parameters $B_i$.

Let $Y$ be the set of finite places that lie over $\ell$, at which $\kappa$ or some $\psi_{i,j}$ is
ramified, or that are split over $F$ and at which $U$ has no Satake parameters. It is finite: a nonzero
vector of $U$ is fixed by $c_w^{-1}(\operatorname{GL}_n(\mathcal O_w))$ for all but finitely many split
$w$, and then $U$ has a spherical local component at $w$, which has Satake parameters. Let $w\notin Y$ be
$d$-split, let $U$ have Satake parameters $B$ at $w$, and put $S=\kappa(\varpi_w)^{-1}q_w^{(n-1)/2}B$. By
(1)(b) of Definition~\ref{gal-larger-enlargement}, $U_i$ has Satake parameters $B_i$ at $w$ with
$q_w^{(n+t-1)/2}B_i=S+\{\psi_{i,1}(\varpi_w),\ldots,\psi_{i,t}(\varpi_w)\}$; so $r_i$ has these Frobenius
eigenvalues at $w$, and $a_{i,j}$ has Frobenius eigenvalue $\psi_{i,j}(\varpi_w)$. Hence, for every
geometric Frobenius element $g$ at $w$,
\[
 \det(X-r_1(g))\prod_{j=1}^t(X-a_{2,j}(g))=\det(X-r_2(g))\prod_{j=1}^t(X-a_{1,j}(g)),
\]
since both sides have image
$\prod_{z\in S}(X-z)\prod_j(X-\psi_{1,j}(\varpi_w))\prod_j(X-\psi_{2,j}(\varpi_w))$ under $\iota$. The
coefficients of both sides are continuous functions of $g\in G_E$, and these Frobenius elements are dense,
so the identity holds for every $g\in G_E$.

Lemma~\ref{gal-constituent-peel}(1), with $r=r_1$, $r'=r_2$, $a_j=a_{1,j}$ and $a'_j=a_{2,j}$, gives a
continuous semisimple $s:G_E\to\operatorname{GL}_n(\bar{\mathbf Q}_\ell)$ with
$r_1\cong s\oplus a_{1,1}\oplus\cdots\oplus a_{1,t}$. Put $r_{w_0}=s\otimes\chi_\kappa$, which is continuous
and semisimple. Let $w$ be a $d$-split place at which $\kappa,\psi_{1,1},\ldots,\psi_{1,t}$ are unramified
and $U$ has Satake parameters $B$, and let $S$ be as before. By (1)(b) of
Definition~\ref{gal-larger-enlargement}, $U_1$ has Satake parameters $B_1$ at $w$ with
$q_w^{(n+t-1)/2}B_1=S+\{\psi_{1,j}(\varpi_w)\}_j$, and by (1)(a),
$J(\lambda_1)_\tau=\{x+m_{\iota\circ\tau}:x\in J(\lambda)_\tau\}+\{-p_{1,j,\iota\circ\tau}\}_j$. So $r_1$ has
spherical data $(S+\{\psi_{1,j}(\varpi_w)\}_j,J(\lambda_1))$ at $w$, and each $a_{1,j}$ has spherical data
$(\{\psi_{1,j}(\varpi_w)\},(\{-p_{1,j,\iota\circ\tau}\})_\tau)$ by Lemma~\ref{gal-constituent-characters}.
Lemma~\ref{gal-constituent-peel}(2) shows that $s$ has spherical data
$(S,(\{x+m_{\iota\circ\tau}:x\in J(\lambda)_\tau\})_\tau)$ at $w$. As $\chi_\kappa$ has spherical data
$(\{\kappa(\varpi_w)\},(\{-m_{\iota\circ\tau}\})_\tau)$, Lemma~\ref{gal-constituent-rules}(3) shows that
$r_{w_0}$ has spherical data $(q_w^{(n-1)/2}B,J(\lambda))$ at $w$: it satisfies the four conditions for $U$
and $\lambda$ at $w$. This applies at $w=w_0$ and at every $d$-split $w\notin Y$; in particular $r_{w_0}$
has Frobenius eigenvalues $q_w^{(n-1)/2}B$ at every $d$-split $w\notin Y$.

\emph{One representation for all places.} For two $d$-split places $w_0$ and $w_1$, the representations
$r_{w_0}$ and $r_{w_1}$ have the same Frobenius eigenvalues at all but finitely many $d$-split places not
over $\ell$. So their characteristic polynomials agree on a dense set of Frobenius elements, hence on $G_E$
by continuity, and $r_{w_0}\cong r_{w_1}$ by the Brauer--Nesbitt theorem, both being semisimple. Fix a
$d$-split place $w_1$ and put $r=r_{w_1}$. For every $d$-split place $w$ at which $U$ has Satake parameters
$B$, $r\cong r_w$, and $r_w$ has spherical data $(q_w^{(n-1)/2}B,J(\lambda))$ at $w$, so $r$ has them by
Lemma~\ref{gal-constituent-rules}(1).
\end{proof}

\subsection{Patching and the self-dual scope}\label{subsection-gal-patching}

\begin{lemma}\label{gal-patch-index-two}
Let $\Gamma$ be a topological group, $k$ an algebraically closed field of characteristic $0$ with a
topology for which it is a topological ring, and $n\geq0$. Let $I$ be an infinite set and
$(\eta_i)_{i\in I}$ pairwise distinct homomorphisms $\Gamma\to\{\pm1\}$ whose kernels $H_i$ are open.
For each $i\in I$ let $r_i:H_i\to\operatorname{GL}_n(k)$ be a continuous semisimple homomorphism, and
suppose:
\begin{enumerate}
\item $\operatorname{tr}r_i(ghg^{-1})=\operatorname{tr}r_i(h)$ for all $i\in I$, $g\in\Gamma$ and
  $h\in H_i$;
\item $\operatorname{tr}r_i(h)=\operatorname{tr}r_j(h)$ for all $i,j\in I$ and $h\in H_i\cap H_j$.
\end{enumerate}
Then there is a continuous semisimple homomorphism $\rho:\Gamma\to\operatorname{GL}_n(k)$ such that,
for all but finitely many $j\in I$, $\operatorname{tr}\rho(h)=\operatorname{tr}r_j(h)$ for every
$h\in H_j$. For the absolute Galois group of a number field and $k=\bar{\mathbf Q}_\ell$, this is the
form of Sorensen's patching lemma that Steps 1 to 6 of its proof give for an infinite family of
quadratic extensions \cite[Section~4.3, Lemma~1]{Sor10}.
\end{lemma}
\leanref{PositivityTheory/AutomorphicForms/PublishedInputs/RealGaloisRepresentations/IwahoriSphericalTransfer/EqOneOrEqOfKer.lean}{GlobalTheory.IndexTwoPatch.main}

\begin{lemma}\label{gal-patch-unique}
Let $K$ be a number field, $\Omega$ an algebraic closure of $K$, $\ell$ a prime,
$\iota:\bar{\mathbf Q}_\ell\to\mathbf C$ a field isomorphism, $a\geq0$, and
$r,r':G_K\to\operatorname{GL}_a(\bar{\mathbf Q}_\ell)$ continuous semisimple homomorphisms. Suppose
that for all but finitely many places $w$ of $K$ of prime norm there is a multiset $R$ such that $r$
and $r'$ both have Frobenius eigenvalues $R$ at $w$ (Definition~\ref{gal-weakly-attached}(1), with
$K$ in place of $E$). Then $r\cong r'$.
\end{lemma}
\leanref{Development/AutomorphicArithmetic/Foundations/GaloisRepresentations/IntegralStructures/CrystallineDescent/PatchOneAlt.lean}{UpperTheory.Gal.Int.Core.IsoOne}
\leanref{Development/AutomorphicArithmetic/Foundations/GaloisRepresentations/IntegralStructures/CrystallineDescent/IsoOne.lean}{UpperTheory.Gal.Int.Core.isoOne_holds}

\begin{proof}
A place of $K$ has prime norm if and only if it has absolute residue degree one. Every nonempty open
subset of $G_K$ contains a coset of $G_L$ for a finite Galois extension $L$ of $K$ in $\Omega$, so by
Theorem~\ref{lit-chebotarev-degree-one} it contains a Frobenius element at a place of prime norm
outside any given finite set. Hence the inverses $\varphi$ of the Frobenius elements at the places $w$
of the hypothesis are dense in $G_K$. At such $\varphi$, $\iota(r(\varphi))$ and $\iota(r'(\varphi))$
both have characteristic polynomial $\prod_{z\in R}(X-z)$. The characteristic polynomials of $r$ and
$r'$ depend continuously on the element of $G_K$ and agree on a dense set, so they agree everywhere,
and $r\cong r'$ by the Brauer--Nesbitt theorem.
\end{proof}

\begin{definition}\label{gal-patch-system}
Let $F$ be a number field, $\Omega$ an algebraic closure of $F$, $\ell$ a prime,
$\iota:\bar{\mathbf Q}_\ell\to\mathbf C$ a field isomorphism and $a\geq0$. A \emph{Satake system} for
$F$ is a family $\mathfrak B=(\mathfrak B_v)_v$, indexed by the finite places $v$ of $F$, of sets of
multisets of nonzero complex numbers, each with at most one element. For a finite extension $K$ of $F$
in $\Omega$, a homomorphism $r:G_K\to\operatorname{GL}_a(\bar{\mathbf Q}_\ell)$ \emph{realizes}
$\mathfrak B$ if for all but finitely many places $w$ of $K$ of prime norm there is $B\in\mathfrak B_v$,
with $v$ the place of $F$ under $w$, such that $r$ has Frobenius eigenvalues $q_w^{(a-1)/2}B$ at $w$
(Definition~\ref{gal-weakly-attached}(1), with $K$ in place of $E$).
\end{definition}
\leanref{Development/AutomorphicArithmetic/Foundations/GaloisRepresentations/IntegralStructures/CrystallineDescent/PatchOneAlt.lean}{UpperTheory.Gal.Int.Core.AttOne}

\begin{lemma}\label{gal-patch-quadratic}
Let $F$ be a number field, $\Omega$ an algebraic closure of $F$, $\ell$ a prime,
$\iota:\bar{\mathbf Q}_\ell\to\mathbf C$ a field isomorphism, $a\geq0$ and $\mathfrak B$ a Satake
system for $F$ (Definition~\ref{gal-patch-system}). Let $S$ be an infinite set of quadratic extensions
of $F$ in $\Omega$ that are CM fields, and for each $Q\in S$ let
$r_Q:G_Q\to\operatorname{GL}_a(\bar{\mathbf Q}_\ell)$ be a continuous semisimple homomorphism that
realizes $\mathfrak B$. Then there is a continuous semisimple
$\rho:G_F\to\operatorname{GL}_a(\bar{\mathbf Q}_\ell)$ such that $\rho|_{G_Q}\cong r_Q$ for all but
finitely many $Q\in S$.
\end{lemma}
\leanref{Development/AutomorphicArithmetic/Foundations/GaloisRepresentations/IntegralStructures/CrystallineDescent/PatchOneAlt.lean}{UpperTheory.Gal.Int.Core.PatchOne}
\leanref{Development/AutomorphicArithmetic/Foundations/GaloisRepresentations/IntegralStructures/CrystallineDescent/PatchOne.lean}{UpperTheory.Gal.Int.Core.patchOne_holds}
\leanref{Development/AutomorphicArithmetic/Foundations/GaloisRepresentations/IntegralStructures/CrystallineDescent/PatchOneAlt.lean}{UpperTheory.Gal.Int.Core.OneConj}
\leanref{Development/AutomorphicArithmetic/Foundations/GaloisRepresentations/IntegralStructures/CrystallineDescent/OneConj.lean}{UpperTheory.Gal.Int.Core.oneConj_holds}
\leanref{Development/AutomorphicArithmetic/Foundations/GaloisRepresentations/IntegralStructures/CrystallineDescent/PatchOneAlt.lean}{UpperTheory.Gal.Int.Core.OneTower}
\leanref{Development/AutomorphicArithmetic/Foundations/GaloisRepresentations/IntegralStructures/CrystallineDescent/OneTower.lean}{UpperTheory.Gal.Int.Core.oneTower_holds}
\leanref{PositivityTheory/AutomorphicForms/PublishedInputs/RealGaloisRepresentations/IwahoriSphericalTransfer/RangeRes.lean}{GlobalTheory.RealGal.SignChar.main}
\leanref{PositivityTheory/AutomorphicForms/PublishedInputs/RealGaloisRepresentations/IwahoriSphericalTransfer/ImportsTer.lean}{GlobalTheory.TraceConj.main}

\begin{proof}
For a finite extension $K$ of $F$ in $\Omega$, two continuous semisimple representations of $G_K$ that
realize $\mathfrak B$ are isomorphic. Indeed each $\mathfrak B_v$ has at most one element, so the two
have the same Frobenius eigenvalues at all but finitely many places of $K$ of prime norm, and
Lemma~\ref{gal-patch-unique} applies.

The $r_Q$ satisfy condition (1) of Lemma~\ref{gal-patch-index-two}. Let $Q\in S$ and $g\in G_F$. Since
$Q$ is Galois over $F$, conjugation by $g$ preserves $G_Q$, and $r_Q^g(s)=r_Q(gsg^{-1})$ defines a
continuous semisimple representation of $G_Q$. It realizes $\mathfrak B$: for a valuation ring $A$ of
$\Omega$, conjugation by $g$ carries $D_A$, $I_A$ and the Frobenius elements of $D_A$ to $D_{gA}$,
$I_{gA}$ and the Frobenius elements of $D_{gA}$, and the places of $Q$ under $A$ and under $gA$ have the
same norm and lie over the same place of $F$. So $r_Q^g\cong r_Q$ by the first paragraph, and
$\operatorname{tr}r_Q(gsg^{-1})=\operatorname{tr}r_Q(s)$ for every $s\in G_Q$.

The $r_Q$ satisfy condition (2). Let $Q,Q'\in S$ and $M=QQ'$, a finite Galois extension of $F$ with
$G_M=G_Q\cap G_{Q'}$. The restriction $r_Q|_{G_M}$ is continuous, semisimple by Clifford's theorem
since $G_M$ is normal in $G_Q$, and realizes $\mathfrak B$: a place $u$ of $M$ of prime norm has the
same norm as the place $w$ of $Q$ under it, so for a valuation ring $A$ over $u$ the Frobenius elements
of $D_A\cap G_M$ are Frobenius elements at $w$ and $I_A\cap G_M\subset I_A$, while only finitely many
places of $M$ lie over the finitely many places of $Q$ excluded for $r_Q$. The same holds for
$r_{Q'}|_{G_M}$, so the two restrictions are isomorphic by the first paragraph and have the same trace
on $G_M$.

Finally, $G_Q$ is the kernel of the quadratic character $\eta_Q:G_F\to\{\pm1\}$ of $Q/F$, this kernel is
open, and distinct $Q$ have distinct $\eta_Q$ by Galois theory. Lemma~\ref{gal-patch-index-two}, with
$\Gamma=G_F$, $k=\bar{\mathbf Q}_\ell$ and the family $(\eta_Q,r_Q)_{Q\in S}$, gives a continuous
semisimple $\rho$ with $\operatorname{tr}\rho=\operatorname{tr}r_Q$ on $G_Q$ for all but finitely many
$Q\in S$. For such $Q$, $\rho|_{G_Q}$ is semisimple by Clifford's theorem, so $\rho|_{G_Q}\cong r_Q$ by
the Brauer--Nesbitt theorem.
\end{proof}

\begin{definition}\label{gal-patch-serves}
Let $E$ be a CM field, $\Omega$ an algebraic closure of $E$, $\ell$ a prime,
$\iota:\bar{\mathbf Q}_\ell\to\mathbf C$ a field isomorphism, $a\geq1$, $\pi$ a cuspidal automorphic
representation of $\operatorname{GL}_a(\mathbf A_E)$ (Definition~\ref{gal-polarized}), and
$Q\subset\Omega$ a quadratic extension of $E$ that is a CM field.
\begin{enumerate}
\item A \emph{base change} of $\pi$ to $Q$ is a cuspidal automorphic representation $P$ of
  $\operatorname{GL}_a(\mathbf A_Q)$ with properties (1) to (3) of
  Theorem~\ref{gal-cyclic-base-change} for $K=E$ and $M=Q$.
\item Let $w_0$ be a finite place of $E$. The field $Q$ \emph{serves} $\pi$ at $w_0$ if:
  \begin{enumerate}
  \item $D_A\subset G_Q$ for every valuation ring $A$ of $\Omega$ over $w_0$;
  \item some place of $Q$ ramified over $E$ lies over a place where $\pi$ has Satake parameters;
  \item for every base change $P$ of $\pi$ to $Q$ there are a continuous semisimple
    $r_P:G_Q\to\operatorname{GL}_a(\bar{\mathbf Q}_\ell)$ and a set $W$ of places of $Q$, containing
    every place over $w_0$ and all but finitely many places of prime norm, such that $r_P$ satisfies
    the four conditions for $P$ (Definition~\ref{gal-spherical-conditions}) at every $w\in W$ where $P$
    has Satake parameters.
  \end{enumerate}
\end{enumerate}
\end{definition}
\leanref{Development/AutomorphicArithmetic/Foundations/GaloisRepresentations/IntegralStructures/CrystallineDescent/PatchOneAlt.lean}{UpperTheory.Gal.Int.Core.Serves}
\leanref{Development/AutomorphicArithmetic/Foundations/GaloisRepresentations/IntegralStructures/CrystallineDescent/LimOddOfLimWU3.lean}{UpperTheory.Gal.Int.Core.RowP}

\begin{lemma}\label{gal-patch-descent}
Let $E$ be a CM field, $\Omega$ an algebraic closure of $E$, $\ell$ a prime,
$\iota:\bar{\mathbf Q}_\ell\to\mathbf C$ a field isomorphism, $a\geq1$ and $\pi$ a cuspidal automorphic
representation of $\operatorname{GL}_a(\mathbf A_E)$. Suppose that for every finite place $w_0$ of $E$
infinitely many quadratic extensions $Q\subset\Omega$ of $E$ that are CM fields serve $\pi$ at $w_0$
(Definition~\ref{gal-patch-serves}). Then there is a continuous semisimple
$r:G_E\to\operatorname{GL}_a(\bar{\mathbf Q}_\ell)$ that satisfies the four conditions for $\pi$
(Definition~\ref{gal-spherical-conditions}) at every place where $\pi$ has Satake parameters.
\end{lemma}
\leanref{Development/AutomorphicArithmetic/Foundations/GaloisRepresentations/IntegralStructures/CrystallineDescent/PatchOneAlt.lean}{UpperTheory.Gal.Int.Core.StepDown}
\leanref{Development/AutomorphicArithmetic/Foundations/GaloisRepresentations/IntegralStructures/CrystallineDescent/JoinEndHolds.lean}{UpperTheory.Gal.Int.Core.stepDown_of}
\leanref{Development/AutomorphicArithmetic/Foundations/GaloisRepresentations/IntegralStructures/CrystallineDescent/LimOddOfLimFinOdd.lean}{UpperTheory.Gal.Int.Core.stepDown_holds}
\leanref{Development/AutomorphicArithmetic/Foundations/GaloisRepresentations/IntegralStructures/CrystallineDescent/PatchOneAlt.lean}{UpperTheory.Gal.Int.Core.SphDown}
\leanref{Development/AutomorphicArithmetic/Foundations/GaloisRepresentations/IntegralStructures/CrystallineDescent/LimOddOfLimFinOdd.lean}{UpperTheory.Gal.Int.Core.sphDown_holds}
\leanref{Development/AbelianVarietyTheory/Definitions/GaloisGeometry/IsWeight.lean}{GlobalTheory.GalCut.QuadLift}
\leanref{Development/AbelianVarietyTheory/Proofs/SatakeParameters/QuadLift.lean}{GlobalTheory.QuadLiftBridge.quadLift}

\begin{proof}
Let $\mathfrak B$ be the Satake system for $E$ (Definition~\ref{gal-patch-system}) with $\mathfrak B_v$
the set of Satake parameters of $\pi$ at $v$, which has at most one element.

Let $Q$ serve $\pi$ at some place. Then $\pi$ has a base change $P$ to $Q$. Indeed, by condition (2)
$Q$ is ramified over $E$ at a place $v$ where $\pi$ has Satake parameters. Let $\eta$ be a Hecke
character as in the hypothesis of Theorem~\ref{gal-cyclic-base-change} for $K=E$ and $M=Q$. Then $\eta$
is the quadratic character of $Q/E$, so it is ramified at $v$, and the local component of
$\pi\otimes\eta$ at $v$ is ramified. But $\pi$ and $\pi\otimes\eta$ have the same Satake parameters at
almost all places, so $\pi\cong\pi\otimes\eta$ by strong multiplicity one \cite{JS81a}, a contradiction.
So no such $\eta$ exists, and Theorem~\ref{gal-cyclic-base-change} gives $P$.
Every $r_P$ as in condition (3) realizes $\mathfrak B$. For all but finitely many places $w$ of $Q$ of
prime norm, $w\in W$, $w\nmid\ell$, and $\pi$ has Satake parameters $B$ at the place $v$ under $w$.
Since $N(w)$ is prime, $w$ has residue degree one over $v$, so $P$ has Satake parameters $B$ at $w$ by
property (1) of base change, and condition (unr) for $P$ at $w$ says that $r_P$ has Frobenius
eigenvalues $q_w^{(a-1)/2}B$ at $w$.

Let $S$ be the set of quadratic extensions $Q\subset\Omega$ of $E$ that are CM fields and carry a
continuous semisimple $r_Q:G_Q\to\operatorname{GL}_a(\bar{\mathbf Q}_\ell)$ realizing $\mathfrak B$, and
choose such an $r_Q$ for each $Q\in S$. By the previous paragraph $S$ contains every $Q$ that serves
$\pi$ at some place, so $S$ is infinite. Lemma~\ref{gal-patch-quadratic} gives a continuous semisimple
$\rho:G_E\to\operatorname{GL}_a(\bar{\mathbf Q}_\ell)$ and a finite subset $S_0\subset S$ with
$\rho|_{G_Q}\cong r_Q$ for every $Q\in S\setminus S_0$.

Let $w$ be a place where $\pi$ has Satake parameters $B$. Infinitely many $Q$ serve $\pi$ at $w$, so
one of them lies outside $S_0$; call it $Q$, and let $P$, $r_P$ and $W$ be as above. Both $r_P$ and
$r_Q$ realize $\mathfrak B$, and each $\mathfrak B_v$ has at most one element, so
$r_P\cong r_Q\cong\rho|_{G_Q}$ by Lemma~\ref{gal-patch-unique}. The four
conditions are unchanged under isomorphism, so $\rho|_{G_Q}$ satisfies them for $P$ at every place of
$Q$ over $w$ where $P$ has Satake parameters, since these places lie in $W$. They descend to $\rho$ at
$w$. Let $A$ be a valuation ring over $w$ and $w'$ the place of $Q$ under $A$. Since $D_A\subset G_Q$
by condition (1), $w'$ has the same completion and the same norm as $w$. So $P$ has Satake parameters
$B$ at $w'$ by property (1) of base change, its local component at $w'$ is that of $\pi$ at $w$ by
property (3), its infinitesimal characters are those of $\pi$ by property (2), and $D_A$, $I_A$ and
$W_A$ are the same subgroups whether computed over $E$ or over $Q$. Hence each of the four conditions
for $P$, $\rho|_{G_Q}$ and $w'$ is the same condition for $\pi$, $\rho$ and $w$.
\end{proof}

\begin{definition}\label{gal-patch-steps}
Let $Y$ be a CM field, $\Omega$ an algebraic closure of $Y$ and $T$ a finite set of finite places of
$Y$. For a CM field $X$ write $X^+$ for its maximal totally real subfield.
\begin{enumerate}
\item A \emph{real step} of $Y$ for $T$ is a quadratic extension $Q=Y(\delta)\subset\Omega$ of $Y$ that
  is a CM field, where $\delta^2$ is a positive integer and $\delta$ does not lie in the Galois closure
  of $Y$ over $\mathbf Q$, such that the places of $Y^+$ under $T$ split in $Y^+(\delta)$, the places
  in $T$ split in $Q$, and $Q$ is ramified over $Y$ at some place outside $T$.
\item A \emph{CM step} of $Y$ for $T$ is a quadratic extension $Q\subset\Omega$ of $Y$ that is a CM
  field, such that: $[Q^+:\mathbf Q]$ is even; the places in $T$ split in $Q$; $Q$ is ramified over $Y$
  at some place outside $T$; every place of $Q$ not split over $Q^+$ is unramified over $Q^+$, lies
  over a place of $Y$ outside $T$ that splits in $Q$, and lies over a place $v$ of $Q^+$ with
  $v\nmid2$, $N(v)>3$ and $v$ prime to the different of $Q^+$; and $Q$ contains $\sqrt{-d}$ for an
  integer $d>0$ such that every place of $Q$ over $T$ is $d$-split (Definition~\ref{gal-seed}).
\end{enumerate}
\end{definition}
\leanref{Development/AutomorphicArithmetic/Foundations/GaloisRepresentations/InertiaQuadraticRamification/InertSplitAlt.lean}{UpperTheory.Gal.UnramifiedFromSplittingCriterion.RealStep}
\leanref{Development/AutomorphicArithmetic/Foundations/GaloisRepresentations/InertiaQuadraticRamification/InertSplitAlt.lean}{UpperTheory.Gal.UnramifiedFromSplittingCriterion.Step0}
\leanref{Development/AutomorphicArithmetic/Foundations/GaloisRepresentations/InertiaQuadraticRamification/InertSplitAlt.lean}{UpperTheory.Gal.UnramifiedFromSplittingCriterion.Step1}
\leanref{Development/AutomorphicArithmetic/Foundations/GaloisRepresentations/IntegralStructures/CrystallineDescent/PatchOneAlt.lean}{UpperTheory.Gal.Int.Core.Step2}

\begin{lemma}\label{gal-patch-steps-kummer}
Let $E$ be a number field, $\Omega$ an algebraic closure of $E$, and $A\neq\Omega$ a valuation ring of
$\Omega$, with maximal ideal $\mathfrak m_A$ and valuation $v_A$; let $u$ be the place of $E$ with prime
ideal $\mathfrak m_A\cap\mathcal O_E$, and $I$ the group of $\sigma\in G_E$ with $\sigma(A)=A$ and
$v_A(\sigma(z)-z)>0$ for all $z\in A$. Let $g$ be a prime number with $v_u(g)=1$, let $m\geq1$ be an
integer prime to $g$, and let $Y,\mu\in\Omega$ with $Y^m=g$ and $\mu^m=1$. Then $\sigma(Y)=\mu Y$ for some
$\sigma\in I$.
\end{lemma}
\leanref{Development/AbelianVarietyTheory/Proofs/IwasawaDecomposition/TameKummerInertiaSurjective.lean}{GlobalTheory.TameKummerInertia.tame_kummer_surj}
\leanref{Development/AbelianVarietyTheory/Proofs/IwasawaDecomposition/TameKummerInertiaSurjective.lean}{GlobalTheory.TameKummerInertia.index_one_root}
\leanref{Development/AbelianVarietyTheory/Proofs/IwasawaDecomposition/TameKummerInertiaSurjective.lean}{GlobalTheory.TameKummerInertia.unram_of_fixed}
\leanref{Development/AbelianVarietyTheory/Proofs/IwasawaDecomposition/ValSubOneLtOfRed.lean}{GlobalTheory.TameCoupledEx.fix_rou}
\leanref{Development/AbelianVarietyTheory/Proofs/IwasawaDecomposition/ValSubOneLtOfRed.lean}{GlobalTheory.TameCoupledEx.ratio_mul}

\begin{proof}
For $\sigma\in G_E$, $(\sigma(Y)/Y)^m=\sigma(g)/g=1$. The elements of $I$ fix the $m$-th roots of unity.
Indeed, $m$ is a unit of $A$, as it is prime to $g\in\mathfrak m_A$. If $\zeta^m=1$ and $\sigma\in I$, then
$\xi=\sigma(\zeta)/\zeta$ is an $m$-th root of unity with $v_A(\xi-1)>0$, as $\zeta$ is a unit of $A$. Each
$\xi^i-1$ is a multiple of $\xi-1$ in $A$, so lies in $\mathfrak m_A$. If $\xi\neq1$, then
$\sum_{i<m}\xi^i=0$, and $m=-\sum_{i<m}(\xi^i-1)$ lies in $\mathfrak m_A$, which is false. So
$\sigma(\zeta)=\zeta$. Hence $\chi(\sigma)=\sigma(Y)/Y$ is a homomorphism from $I$ to the group of $m$-th
roots of unity: for $\sigma\in I$ and $\tau\in G_E$, $\sigma(\tau(Y))=\sigma(\chi(\tau)Y)=\chi(\tau)\sigma(Y)$.

The image of $\chi$ is cyclic; let $\sigma_0\in I$ be such that $\chi(\sigma_0)$ generates it, and let $d$
be its order, a divisor of $m$. Then $Z=Y^d$ satisfies $\sigma(Z)=\chi(\sigma)^dZ=Z$ for every $\sigma\in
I$, and $Z^{m/d}=g$. The field $K=E(Z)$ is fixed by $I$, so the prime of $K$ under $A$ is unramified over
$u$ \cite[Chapter~I, \S9]{Neu99}. Its normalized valuation $w$ therefore has $w(g)=v_u(g)=1$, and
$(m/d)\,w(Z)=1$ gives $m/d=1$. So $\chi(\sigma_0)$ is a primitive $m$-th root of unity, $\mu=\chi(\sigma_0)^i$
for some $i\geq0$, and $\sigma=\sigma_0^i\in I$ satisfies $\sigma(Y)=\chi(\sigma_0)^iY=\mu Y$.
\end{proof}

\begin{lemma}\label{gal-patch-steps-squares}
Let $\Omega$ be an algebraic closure of $\mathbf Q$ and $A\neq\Omega$ a valuation ring of $\Omega$, with
maximal ideal $\mathfrak m_A$ and valuation $v_A$. Let $D_A$ be the group of automorphisms $\sigma$ of
$\Omega$ with $\sigma(A)=A$, and $I_A\subset D_A$ the subgroup of the $\sigma$ with
$v_A(\sigma(z)-z)>0$ for all $z\in A$. Let $a$ be a nonzero integer and $x\in\Omega$ with $x^2=a$.
\begin{enumerate}
\item If $t\in\mathfrak m_A$ for a prime $t$ with $8t\mid a-1$, then $D_A$ fixes $x$.
\item If $2a$ is a unit of $A$, then $I_A$ fixes $x$.
\item Let $F\subset\Omega$ be a number field and $u$ the place of $F$ with prime ideal
  $\mathfrak m_A\cap\mathcal O_F$. If $a=pb$ for an odd prime $p\in\mathfrak m_A$ with $v_u(p)=1$ and an
  integer $b$ prime to $p$, then some $\sigma\in I_A\cap G_F$ has $\sigma(x)=-x$.
\item If $t$ is a prime and $d$ a positive integer with $8t\mid d+1$, then $-d$ is a square in
  $\mathbf Q_t$.
\end{enumerate}
\end{lemma}
\leanref{Development/AutomorphicArithmetic/Foundations/GaloisRepresentations/InertiaQuadraticRamification/InertSplitAlt.lean}{UpperTheory.Gal.UnramifiedFromSplittingCriterion.SqrtFix}
\leanref{Development/AutomorphicArithmetic/Foundations/GaloisRepresentations/InertiaQuadraticRamification/SqrtFix.lean}{UpperTheory.Gal.UnramifiedFromSplittingCriterion.sqrtFix}
\leanref{Development/AbelianVarietyTheory/Proofs/IwasawaDecomposition/TameKummerInertiaSurjective.lean}{GlobalTheory.TameKummerInertia.tame_kummer_surj}
\leanref{PositivityTheory/AutomorphicForms/PublishedInputs/RealGaloisRepresentations/IwahoriSphericalTransfer/InfinitePrimesNegIsSquare.lean}{GlobalTheory.RealGal.PrimesNegSquare.isSquare_neg_of_dvd}

\begin{proof}
Every $\sigma\in D_A$ has $\sigma(x)^2=a$, so $\sigma(x)=\pm x$.

For (1), let $\sigma\in D_A$ with $\sigma(x)=-x$, and put $y=(x-1)/2$. Then $\sigma(y)=(-x-1)/2$, so
$y+\sigma(y)=-1$ and $y\sigma(y)=(1-a)/4$, which lies in $2t\mathbf Z\subset\mathfrak m_A$. Since
$\sigma$ maps $\mathfrak m_A$ onto itself, either $y$ and $\sigma(y)$ both lie in $\mathfrak m_A$, and
then so does $-1$, or neither does, and then $v_A(y\sigma(y))\leq0$. Both are absurd, so
$\sigma(x)=x$.

For (2), $x$ is a unit of $A$ since $x^2=a$ is. If some $\sigma\in I_A$ had $\sigma(x)=-x$, then
$\sigma(x)-x=-2x$ would be a unit, against $v_A(\sigma(x)-x)>0$.

For (3), let $y\in\Omega$ with $y^2=p$. By Lemma~\ref{gal-patch-steps-kummer}, with $E=F$, $g=p$, $m=2$,
$Y=y$ and $\mu=-1$, some $\sigma\in I_A\cap G_F$ has $\sigma(y)=-y$. Put $z=x/y$, so that $z^2=b$. Since $\mathfrak m_A\cap\mathbf Z=p\mathbf Z$ and $p$ is odd and prime to
$b$, $2b$ is a unit of $A$, and $\sigma(z)=z$ by (2). Hence $\sigma(x)=\sigma(z)\sigma(y)=-x$.

For (4), $-d\equiv1\pmod{8t}$. If $t$ is odd, $-d$ is a nonzero square modulo $t$; if $t=2$, then
$-d\equiv1\pmod{16}$. In both cases $-d$ is a square in $\mathbf Q_t$ by Hensel's lemma
\cite[Chapter~II, \S4]{Neu99}.
\end{proof}

\begin{remark}\label{gal-patch-steps-squares-inertia}
The element $\sigma$ in the proof of Lemma~\ref{gal-patch-steps-squares}(3) can also be found directly. Let
$L\subset\Omega$ be a finite Galois extension of $F$ containing $y$. The image of $I_A\cap G_F$ in
$\operatorname{Gal}(L/F)$ is the inertia group of the place of $L$ under $A$, whose fixed field is
unramified over $u$ \cite[Chapter~I, \S9]{Neu99}. The values of $v_A$ on that fixed field are therefore the
multiples of $v_A(p)$, since $v_u(p)=1$, while $v_A(y)=v_A(p)/2$. So $y$ is not in it, and some
$\sigma\in I_A\cap G_F$ has $\sigma(y)=-y$.
\end{remark}
% Lean: (1) and (2) are sqrtFix, whose proof is the one given. (3) is exists_moves in
% ExistsPlaceOfPrime.lean for b = 1 and WeylLatticeInDetCoeffSpan.moved_at_ram
% (NonsplitOfSqrtConductor.lean) for b = -p_2, both from tame_kummer_surj; (4) is isSquare_neg_of_dvd.

\begin{lemma}\label{gal-patch-steps-frobenius}
Let $F$ be a number field, $\Omega$ an algebraic closure of $F$, $v$ a finite place of $F$ with norm
$q_v$, and $A$ a valuation ring of $\Omega$ over $v$, with maximal ideal $\mathfrak m_A$ and
decomposition and inertia groups $D_A\supset I_A$ in $G_F$ (Situation~\ref{gal-local-situation}). Let
$\varphi\in D_A$ with $\varphi(z)\equiv z^{q_v}\pmod{\mathfrak m_A}$ for all $z\in A$, and let
$X\subset\Omega$ be a finite set fixed pointwise by $I_A$. Then for every $g\in D_A$ there is an
integer $k$ with $g(x)=\varphi^k(x)$ for all $x\in X$.
\end{lemma}
\leanref{Development/AutomorphicArithmetic/Foundations/GaloisRepresentations/InertiaQuadraticRamification/InertSplitAlt.lean}{UpperTheory.Gal.UnramifiedFromSplittingCriterion.DecGen}
\leanref{Development/AutomorphicArithmetic/Foundations/GaloisRepresentations/InertiaQuadraticRamification/DecGen.lean}{UpperTheory.Gal.UnramifiedFromSplittingCriterion.decGen}
\leanref{Development/AbelianVarietyTheory/PublishedTheorems/GaloisCohomology₂/ExistsCoprimePowOfMemIneKer.lean}{GlobalTheory.TameQuot.closure_frob_sup_Ine}

\begin{proof}
The subgroup $H$ of $D_A$ generated by $\varphi$ and $I_A$ is dense in $D_A$, since $D_A/I_A$ is the
Galois group of the residue field extension of $A$ over $v$, which is topologically generated by the
image of $\varphi$ \cite[Chapter~I, \S9]{Neu99}. The $h\in D_A$ that agree with $g$ on $X$ form the
set $g(D_A\cap G_{F(X)})$, which is open in $D_A$ since $F(X)$ is finite over $F$. So some $h\in H$
agrees with $g$ on $X$. Let $\Omega^{I_A}$ be the set of elements of $\Omega$
fixed by $I_A$. Since $I_A$ is normal in $D_A$, every power of $\varphi$ maps $\Omega^{I_A}$ to itself,
and $I_A$ acts trivially on it. Writing $h$ as a product of powers of $\varphi$ and elements of $I_A$
shows that $h$ acts on $\Omega^{I_A}$ as some $\varphi^k$. As $X\subset\Omega^{I_A}$, $g=h=\varphi^k$
on $X$.
\end{proof}

\begin{lemma}\label{gal-patch-steps-adjoin}
Let $Y$ be a CM field, $\Omega$ an algebraic closure of $Y$, $a$ an integer and $x\in\Omega$ with
$x^2=a$ and $x\notin Y$. Then $Y(x)=Y+Yx$ is a quadratic extension of $Y$ and a CM field.
\end{lemma}
\leanref{Development/AutomorphicArithmetic/Foundations/GaloisRepresentations/InertiaQuadraticRamification/InertSplitAlt.lean}{UpperTheory.Gal.UnramifiedFromSplittingCriterion.MkQ}
\leanref{Development/AutomorphicArithmetic/Foundations/GaloisRepresentations/InertiaQuadraticRamification/ConjOfSqNonneg.lean}{UpperTheory.Gal.UnramifiedFromSplittingCriterion.mkQ}

\begin{proof}
A root $b\in Y$ of $X^2-a$ would give $x=\pm b\in Y$, so $X^2-a$ is irreducible over $Y$,
$[Y(x):Y]=2$ and $Y(x)=Y+Yx$. Since $x\notin Y$, $a\neq0$; let $\varepsilon=\pm1$ be the sign of $a$.
For every field embedding $\iota:Y(x)\to\mathbf C$, $\iota(x)^2=a$, so
$\overline{\iota(x)}=\varepsilon\iota(x)$. Let $c$ be the automorphism of $Y(x)$ that is complex
conjugation on $Y$ and sends $x$ to $\varepsilon x$; it exists since $(\varepsilon x)^2=a$. Then
$\overline{\iota(z)}=\iota(c(z))$ for every $z\in Y(x)$ and every $\iota$, since this holds on $Y$ and
at $x$. So the fixed field of $c$ is totally real, and $Y(x)$, which is totally imaginary since $Y$ is,
is a quadratic extension of it. Hence $Y(x)$ is a CM field.
\end{proof}

\begin{lemma}\label{gal-patch-steps-cm-places}
Let $L$ be a CM field with maximal totally real subfield $L^+$, $\Omega$ an algebraic closure of $L$,
$w$ a finite place of $L$ and $A$ a valuation ring of $\Omega$ over $w$, with maximal ideal
$\mathfrak m_A$ and decomposition and inertia groups $D_A\supset I_A$ in $G_{L^+}$
(Situation~\ref{gal-local-situation}).
\begin{enumerate}
\item The place $w$ is split over $L^+$ if and only if $D_A\subset G_L$.
\item Let $x\in\Omega$ with $x^2\in\mathbf Z$. If $I_A$ fixes $L$ and $x$ pointwise and
  $D_A\not\subset G_L$, then $D_A\cap G_L$ fixes $x$.
\end{enumerate}
\end{lemma}
\leanref{Development/AutomorphicArithmetic/Foundations/GaloisRepresentations/InertiaQuadraticRamification/InertSplitAlt.lean}{UpperTheory.Gal.UnramifiedFromSplittingCriterion.SplitIff}
\leanref{Development/AutomorphicArithmetic/Foundations/GaloisRepresentations/InertiaQuadraticRamification/SplitIff.lean}{UpperTheory.Gal.UnramifiedFromSplittingCriterion.splitIff}
\leanref{Development/AutomorphicArithmetic/Foundations/GaloisRepresentations/InertiaQuadraticRamification/InertSplitAlt.lean}{UpperTheory.Gal.UnramifiedFromSplittingCriterion.InertSplit}
\leanref{Development/AutomorphicArithmetic/Foundations/GaloisRepresentations/InertiaQuadraticRamification/InertSplit.lean}{UpperTheory.Gal.UnramifiedFromSplittingCriterion.inertSplit}

\begin{proof}
Let $s$ be the nontrivial automorphism of $L$ over $L^+$ and $v$ the place of $L^+$ under $w$.

For (1), the places of $L$ over $v$ are $w$ and $s(w)$, and $w$ is split over $L^+$ if and only if
$s(w)\neq w$. An element of $G_{L^+}$ outside $G_L$ restricts to $s$ on $L$, so if it lies in $D_A$,
then $s(w)=w$. Conversely, if $s(w)=w$ and $\sigma\in G_{L^+}$ is outside $G_L$, then $A$ and
$\sigma(A)$ both lie over $w$, so $\sigma(A)=\tau(A)$ for some $\tau\in G_L$
\cite[Chapter~I, \S9]{Neu99}, and $\tau^{-1}\sigma$ is an element of $D_A$ outside $G_L$.

For (2), let $\varphi\in D_A$ with $\varphi(z)\equiv z^{q_v}\pmod{\mathfrak m_A}$ for all $z\in A$,
which exists \cite[Chapter~I, \S9]{Neu99}. Choose $e\in L$ with $e\neq0$ and $s(e)=-e$. Then
$L=L^+(e)$ and $e^2\in L^+$, so $\varphi(e)=\pm e$, and $\varphi(x)=\pm x$ since $x^2\in\mathbf Z$.
The set $X=\{x,e\}$ is fixed by $I_A$. Let $g\in D_A$ be outside $G_L$. Since $g$ fixes $L^+$ and
$L=L^+(e)$, $g(e)\neq e$. By Lemma~\ref{gal-patch-steps-frobenius}, $g=\varphi^k$ on $X$ for some $k$,
so $\varphi(e)\neq e$, that is, $\varphi(e)=-e$. Let $h\in D_A\cap G_L$. By
Lemma~\ref{gal-patch-steps-frobenius}, $h=\varphi^m$ on $X$ for some $m$. Then $\varphi^m(e)=h(e)=e$,
so $m$ is even, and $h(x)=\varphi^m(x)=x$ since $\varphi^2(x)=x$.
\end{proof}

\begin{lemma}\label{gal-patch-steps-chebotarev}
Let $Y$ be a number field, $\Omega$ an algebraic closure of $Y$, $q$ a positive integer and
$\delta\in\Omega$ with $\delta^2=q$. Let $g$ be an automorphism of $\Omega$ with $g(\delta)=-\delta$
such that $g^2$ fixes $j(Y)$ pointwise for every field embedding $j:Y\to\Omega$. Let $M\geq1$,
$\zeta\in\Omega$ a primitive $M$-th root of unity, and $B$ an integer. Then there is a prime $p>B$ with
$g(\zeta)=\zeta^p$ such that for every valuation ring $A$ of $\Omega$ whose maximal ideal contains $p$,
every automorphism $h$ of $\Omega$ with $h(A)=A$ and $h(\delta)=\delta$ fixes $Y$ pointwise.
\end{lemma}
\leanref{Development/AutomorphicArithmetic/Foundations/GaloisRepresentations/InertiaQuadraticRamification/InertSplitAlt.lean}{UpperTheory.Gal.UnramifiedFromSplittingCriterion.PrimeWith}
\leanref{Development/AutomorphicArithmetic/Foundations/GaloisRepresentations/InertiaQuadraticRamification/PrimeWith.lean}{UpperTheory.Gal.UnramifiedFromSplittingCriterion.primeWith}

\begin{proof}
Let $\alpha$ be a primitive element of $Y$ over $\mathbf Q$, let $X\subset\Omega$ be the finite set
made of $\delta$, $\zeta$ and the $j(\alpha)$ for the field embeddings $j:Y\to\Omega$, and let
$L\subset\Omega$ be the Galois closure of $\mathbf Q(X)$ over $\mathbf Q$.
Theorem~\ref{lit-chebotarev-degree-one}, for $L/\mathbf Q$, the restriction of $g$ to $L$ and the
primes at most $B+M$, gives a prime $p>B+M$ unramified in $L$, a valuation ring $A_0$ of $\Omega$ whose
maximal ideal $\mathfrak m_0$ contains $p$, and $\varphi\in\operatorname{Gal}(\Omega/\mathbf Q)$ with
$\varphi(A_0)=A_0$, $\varphi(z)\equiv z^p\pmod{\mathfrak m_0}$ for $z\in A_0$, and $\varphi=g$ on $L$.
Since $L$ is unramified at $p$, the inertia group of $A_0$ in $\operatorname{Gal}(\Omega/\mathbf Q)$
fixes $L$, and so fixes $X$.

The elements $\varphi(\zeta)$ and $\zeta^p$ are $M$-th roots of unity, congruent modulo
$\mathfrak m_0$. Two distinct $M$-th roots of unity are not congruent modulo $\mathfrak m_0$, since
$\prod_{\xi\neq1}(1-\xi)=M$, over the $M$-th roots of unity $\xi\neq1$, and $p\nmid M$. Hence
$g(\zeta)=\varphi(\zeta)=\zeta^p$.

Let $A$ be a valuation ring of $\Omega$ whose maximal ideal contains $p$, and $h$ an automorphism of
$\Omega$ with $h(A)=A$ and $h(\delta)=\delta$. There is $\sigma\in\operatorname{Gal}(\Omega/\mathbf Q)$
with $\sigma(A_0)=A$ \cite[Chapter~I, \S9]{Neu99}. Then $h'=\sigma^{-1}h\sigma$ stabilizes $A_0$, and
$h'(\delta)=\delta$ since $\sigma(\delta)=\pm\delta$. By Lemma~\ref{gal-patch-steps-frobenius} for
$F=\mathbf Q$, $h'=\varphi^k$ on $X$ for some integer $k$. On $X\subset L$, $\varphi^k=g^k$, and
$g^k(\delta)=(-1)^k\delta$, so $k$ is even. Let $j$ be the restriction of $\sigma^{-1}$ to $Y$. Since
$g^2$ fixes $j(\alpha)$, so does $g^k$, so $h'(j(\alpha))=j(\alpha)$, that is, $h(\alpha)=\alpha$. As
$Y=\mathbf Q(\alpha)$, $h$ fixes $Y$ pointwise.
\end{proof}

\begin{lemma}\label{gal-patch-steps-two-primes}
Let $Y_0\subset Y$ be number fields, $\Omega$ an algebraic closure of $Y$, and $\delta\in Y$ with
$Y=Y_0(\delta)$, such that $\delta^2$ is a positive integer and $\delta$ does not lie in the Galois
closure of $Y_0$ over $\mathbf Q$ in $\Omega$. Let $M\geq1$ and $B$ be integers. Then there are primes
$p_1<p_2$ with $B<p_1$ and $M\mid p_1p_2+1$ such that for every valuation ring $A$ of $\Omega$ whose
maximal ideal contains $p_1p_2$, every automorphism $h$ of $\Omega$ with $h(A)=A$ and $h(\delta)=\delta$
fixes $Y$ pointwise.
\end{lemma}
\leanref{Development/AutomorphicArithmetic/Foundations/GaloisRepresentations/InertiaQuadraticRamification/InertSplitAlt.lean}{UpperTheory.Gal.UnramifiedFromSplittingCriterion.TwoPrimes}
\leanref{Development/AutomorphicArithmetic/Foundations/GaloisRepresentations/InertiaQuadraticRamification/TwoPrimes.lean}{UpperTheory.Gal.UnramifiedFromSplittingCriterion.twoPrimes}

\begin{proof}
Let $q=\delta^2$ and let $Y_0'$ be the Galois closure of $Y_0$ over $\mathbf Q$ in $\Omega$. Since
$\delta\notin Y_0'$ and $\delta^2\in\mathbf Q$, some automorphism $s$ of $\Omega$ fixes $Y_0'$ pointwise
and has $s(\delta)=-\delta$. Fix a field embedding $\iota:\Omega\to\mathbf C$. Its image is stable
under complex conjugation, so there is an automorphism $c$ of $\Omega$ with
$\iota(c(z))=\overline{\iota(z)}$ for all $z$. Then $c(\delta)=\delta$, since $\iota(\delta)^2=q>0$, and
$c(\xi)=\xi^{-1}$ for every root of unity $\xi$.

Put $g_1=c\circ s$ and $g_2=s^{-1}$. Both send $\delta$ to $-\delta$. Let $j:Y\to\Omega$ be a field
embedding. Then $s$ fixes $j(Y_0)$ and $c(j(Y_0))$ pointwise, as both lie in $Y_0'$, and $c^2=1$, so
$g_1^2=c\,s\,c\,s$ and $g_2^2=s^{-2}$ fix $j(Y_0)$ pointwise. They also fix $j(\delta)=\pm\delta$, and $j(Y)=j(Y_0)(j(\delta))$, so they fix $j(Y)$
pointwise. Let $\zeta\in\Omega$ be a primitive $M$-th root of unity.
Lemma~\ref{gal-patch-steps-chebotarev} for $g_1$ and $B$ gives a prime $p_1>B$ with
$s(\zeta)^{-1}=c(s(\zeta))=\zeta^{p_1}$. The same lemma for $g_2$ and $p_1$ gives a prime $p_2>p_1$ with
$s^{-1}(\zeta)=\zeta^{p_2}$, that is, $\zeta=s(\zeta)^{p_2}$. Then
$\zeta^{p_1p_2}=s(\zeta)^{-p_2}=\zeta^{-1}$, so $M\mid p_1p_2+1$. A valuation ring whose maximal ideal
contains $p_1p_2$ has $p_1$ or $p_2$ in its maximal ideal, and the last property is that of
Lemma~\ref{gal-patch-steps-chebotarev} for $p_1$ or for $p_2$.
\end{proof}

\begin{lemma}\label{gal-patch-steps-exist}
Let $Y$ be a CM field, $\Omega$ an algebraic closure of $Y$ and $T$ a finite set of finite places of
$Y$.
\begin{enumerate}
\item There are infinitely many real steps of $Y$ for $T$ (Definition~\ref{gal-patch-steps}).
\item Suppose that $Y=Y_0(\delta)$ for a subfield $Y_0$ with $[Y:Y_0]=2$ and an element $\delta$ such
  that $\delta^2$ is a positive integer and $\delta$ does not lie in the Galois closure of $Y_0$ over
  $\mathbf Q$. Then there are infinitely many CM steps of $Y$ for $T$.
\end{enumerate}
\end{lemma}
\leanref{Development/AutomorphicArithmetic/Foundations/GaloisRepresentations/InertiaQuadraticRamification/InertSplitAlt.lean}{UpperTheory.Gal.UnramifiedFromSplittingCriterion.Quad0}
\leanref{Development/AutomorphicArithmetic/Foundations/GaloisRepresentations/InertiaQuadraticRamification/InfinitelyManyRealQuadraticSteps.lean}{UpperTheory.Gal.UnramifiedFromSplittingCriterion.infinitely_many_real_quadratic_steps}
\leanref{Development/AutomorphicArithmetic/Foundations/GaloisRepresentations/IntegralStructures/CrystallineDescent/PatchOneAlt.lean}{UpperTheory.Gal.Int.Core.Quad2Ex}
\leanref{Development/AutomorphicArithmetic/Foundations/GaloisRepresentations/IntegralStructures/CrystallineDescent/ExistsQuadraticCMExtSplitAt.lean}{UpperTheory.Gal.Int.Core.exists_quadratic_CM_ext_split_at}
\leanref{Development/AutomorphicArithmetic/Foundations/GaloisRepresentations/InertiaQuadraticRamification/InfOf.lean}{UpperTheory.Gal.UnramifiedFromSplittingCriterion.infOf}
\leanref{Development/AutomorphicArithmetic/Foundations/GaloisRepresentations/InertiaQuadraticRamification/InertSplitAlt.lean}{UpperTheory.Gal.UnramifiedFromSplittingCriterion.Quad0Ex}
\leanref{Development/AutomorphicArithmetic/Foundations/GaloisRepresentations/InertiaQuadraticRamification/ExistsPlaceOfPrime.lean}{UpperTheory.Gal.UnramifiedFromSplittingCriterion.quad0Ex}
\leanref{Development/AutomorphicArithmetic/Foundations/GaloisRepresentations/InertiaQuadraticRamification/QuadraticCharTwistOfExistsOfInfinity.lean}{UpperTheory.Gal.UnramifiedFromSplittingCriterion.quadraticCharTwist_of_exists_of_infinity}

\begin{proof}
For a finite place $u$ of a number field write $N(u)$ for its norm. For a valuation ring $A$ of
$\Omega$ with maximal ideal $\mathfrak m_A$, let $D_A\supset I_A$ be as in
Lemma~\ref{gal-patch-steps-squares}; for a subfield $K$ of $\Omega$, $D_A\cap G_K$ and $I_A\cap G_K$
are the decomposition and inertia groups of $A$ in $G_K$. Let $A$ lie over a place $u$ of a number
field $K$ and let $K'\subset\Omega$ be a finite extension of $K$. The place of $K'$ under $A$ is
unramified over $u$ if and only if $I_A\cap G_K\subset G_{K'}$, and if $K'$ is quadratic over $K$, then
$u$ splits in $K'$ if and only if $D_A\cap G_K\subset G_{K'}$ \cite[Chapter~I, \S9]{Neu99}. By Dedekind's discriminant theorem, a number
field is ramified at finitely many primes, those dividing its discriminant
\cite[Chapter~III, \S2]{Neu99}.

In both parts it suffices to construct one step of $Y$ for $T\cup S$ for every finite set $S$ of
finite places of $Y$. Such a step is a step for $T$ that is ramified over $Y$ at a place outside $S$.
Finitely many quadratic extensions of $Y$ are ramified at finitely many places in all, so taking for
$S$ these places gives a step outside any finite set of steps, and there are infinitely many. So it
suffices to construct, for every finite set $T$, one real step of $Y$ for $T$, and in (2) one CM step.
The constructions are Alpöge's.

\emph{Real steps.} Let $c=\prod_{t\in T}N(t)$ and let $Y'\subset\Omega$ be the Galois closure of $Y$
over $\mathbf Q$. Choose a prime $q\equiv1\pmod{8c}$ with $q>\max(c,2)$ that is unramified in $Y'$,
and so in $Y$. It exists by Theorem~\ref{lit-chebotarev-degree-one} for the field of $8c$-th roots of
unity, the identity, and the set of primes that are at most $\max(c,2)$ or ramified in $Y'$, since a
prime unramified in that field with a Frobenius element fixing its roots of unity is
$\equiv1\pmod{8c}$. Let $\delta\in\Omega$ with $\delta^2=q$, and $Q=Y(\delta)$.

Let $A'$ be a valuation ring of $\Omega$ whose maximal ideal contains $q$, and $u'$ the place of $Y'$
under it. Since $q$ is odd and $v_{u'}(q)=1$, Lemma~\ref{gal-patch-steps-squares}(3) for $F=Y'$ and
$b=1$ gives an element of $G_{Y'}$ that sends $\delta$ to $-\delta$. So $\delta\notin Y'$.

Let $u$ be a place of $Y$ over $q$ and $A$ a valuation ring of $\Omega$ over $u$. In the same way,
Lemma~\ref{gal-patch-steps-squares}(3) for $F=Y$ gives $\tau\in I_A\cap G_Y$ with
$\tau(\delta)=-\delta$. So $\delta\notin Y$, and $Q$ is a quadratic extension of $Y$ and a CM field by
Lemma~\ref{gal-patch-steps-adjoin}. Since $\tau$ does not fix $Q$, $Q$ is ramified over $Y$ at $u$.
Also $u\notin T$, since $N(u)$ is a power of $q$ and $q>c\geq N(t)$ for $t\in T$.

Let $t\in T$, $A$ a valuation ring of $\Omega$ over $t$, and $r$ the prime under $t$. Then
$r\mid N(t)\mid c$, so $8r\mid q-1$, and $D_A$ fixes $\delta$ by
Lemma~\ref{gal-patch-steps-squares}(1). So the place of $Y^+$ under $t$ splits in $Y^+(\delta)$, and
$t$ splits in $Q$. Hence $Q$ is a real step of $Y$ for $T$.

\emph{CM steps.} Let $Y=Y_0(\delta)$ be as in (2), $q_0=\delta^2$ and $F=Y^+$. Let $c$ be the product
of $6$, the norms $N(t)$ for $t\in T$, the norms of the places of $F$ ramified in $Y$, and the primes
ramified in $Y$. Let $B_0$ be such that every prime $p>B_0$ is unramified in $Y$.
Lemma~\ref{gal-patch-steps-two-primes}, for $M=8c$ and the bound $\max(B_0,c)$, gives primes
$p_1<p_2$ with $p_1>\max(B_0,c)$ and $8c\mid p_1p_2+1$ such that for every valuation ring $A$ of
$\Omega$ with $p_1p_2\in\mathfrak m_A$, every element of $D_A$ that fixes $\delta$ fixes $Y$. Let
$p=p_1p_2$, so that $8c\mid -p-1$, let $x\in\Omega$ with $x^2=-p$, and let $Q=Y(x)$.

The field $Q$ is a quadratic extension of $Y$ and a CM field, ramified over $Y$ at a place outside
$T$. Indeed, let $u$ be a place of $Y$ over $p_1$ and $A_0$ a valuation ring of $\Omega$ over $u$.
Since $p_1>B_0$, $v_u(p_1)=1$, and since $p_1>c\geq6$, $p_1$ is odd and prime to $p_2$.
Lemma~\ref{gal-patch-steps-squares}(3) for $F=Y$ and $b=-p_2$ gives $\tau_0\in I_{A_0}\cap G_Y$ with
$\tau_0(x)=-x$. So $x\notin Y$, $Q$ is a quadratic extension of $Y$ and a CM field by
Lemma~\ref{gal-patch-steps-adjoin}, and $Q$ is ramified over $Y$ at $u$. Also $u\notin T$, since $N(u)$
is a power of $p_1$ and $p_1>c$.

The places in $T$ split in $Q$. Indeed, let $t\in T$ and $A$ a valuation ring of $\Omega$ over $t$.
The prime $r$ under $t$ divides $c$, so $8r\mid -p-1$, and $D_A$ fixes $x$ by
Lemma~\ref{gal-patch-steps-squares}(1). So $D_A\cap G_Y$ fixes $Y$ and $x$, hence $Q$. The degree
$[Q^+:\mathbf Q]=[Y:\mathbf Q]=2[F:\mathbf Q]$ is even.

Let $w$ be a place of $Q$ not split over $Q^+$, and $A$ a valuation ring of $\Omega$ over $w$. By
Lemma~\ref{gal-patch-steps-cm-places}(1) for $L=Q$, some $g\in D_A\cap G_{Q^+}$ does not fix $Q$. As
$[Q:Q^+]=2$, $Q$ is generated over $Q^+$ by any of its elements outside $Q^+$. The field $Y$ is totally
imaginary, so it is not contained in $Q^+$, and $x\notin Q^+$ since $x^2<0$. Hence $g$ fixes neither
$Y$ nor $x$. So $c\notin\mathfrak m_A$: otherwise some prime $r\mid c$ lies in $\mathfrak m_A$,
$8r\mid -p-1$, and $g$ fixes $x$ by Lemma~\ref{gal-patch-steps-squares}(1). Next,
$p\notin\mathfrak m_A$: since $\delta^2=q_0>0$, $\delta$ is totally real, so $\delta\in Q^+$ and
$g(\delta)=\delta$, and if $p\in\mathfrak m_A$ the choice of $p_1$ and $p_2$ would make $g$ fix $Y$.
As $2\mid c$, $2p$ is a unit of $A$, and $I_A$ fixes $x$ by Lemma~\ref{gal-patch-steps-squares}(2).
And $I_A\cap G_F$ fixes $Y$: otherwise the place of $F$ under $w$ is ramified in $Y$, its norm
divides $c$ and lies in $\mathfrak m_A$, and $c\in\mathfrak m_A$.

Since $F\subset Q^+$, $I_A\cap G_{Q^+}\subset I_A\cap G_F$ fixes $Y$ and $x$, hence $Q$, so $w$ is
unramified over $Q^+$. The place $u_w$ of $Y$ under $w$ is outside $T$: for $t\in T$, $N(t)$ divides
$c$, so $N(t)\notin\mathfrak m_A$, while $N(u_w)\in\mathfrak m_A$. Since $g\in
D_A\cap G_F$ does not fix $Y$, Lemma~\ref{gal-patch-steps-cm-places}(2) for $L=Y$ shows that
$D_A\cap G_Y$ fixes $x$, hence $Q$, so $u_w$ splits in $Q$. Let $v$ be the place of $Q^+$ under $w$
and $r$ the prime under $v$. Since $6\mid c$ and $c\notin\mathfrak m_A$, $r\geq5$, so $v\nmid2$ and
$N(v)\geq5$. Also $r\nmid c$, so $r$ is unramified in $Y$ and $I_A\cap G_{\mathbf Q}$ fixes $Y$; it
fixes $x$ as well. So $Q$, and with it $Q^+$, is unramified over $\mathbf Q$ at $w$, and $v$ is prime
to the different of $Q^+$ by Dedekind's different theorem \cite[Chapter~III, \S2]{Neu99}.

Finally, $x\in Q$ with $x^2=-p$, and every place $w$ of $Q$ over $T$ is $p$-split
(Definition~\ref{gal-seed}). Indeed $w$ is split over $Q^+$, since the places not split over $Q^+$ lie
over places of $Y$ outside $T$. The prime $r$ under $w$ divides $N(t)$ for the place $t\in T$ under
$w$, so $r\mid c$ and $8r\mid p+1$, and $-p$ is a square in $\mathbf Q_r$ by
Lemma~\ref{gal-patch-steps-squares}(4). Hence $Q$ is a CM step of $Y$ for $T$.
\end{proof}
% Lean: Quad0 is infinitely_many_real_quadratic_steps = quadraticCharTwist_of_exists_of_infinity
% (quad0Ex sqrtFix mkQ unrLarge) infOf. In quad0Ex, B1 and B2 are unrLarge for E and its normal
% closure E', q comes from Nat.exists_prime_gt_modEq_one, and sigma' (the element of G_{E'}) and tau
% come from exists_moves. Quad2Ex is exists_quadratic_CM_ext_split_at: R and R' are the places of F
% and of Q ramified in Y (finite_places_inertia_not_le_openSubgroup), p1 and p2 come from
% twoPrimes (primeWith decGen), the paragraph on w is WeylLatticeInDetCoeffSpan.nonsplit_of_sqrt_conductor
% (with realSpan and fix_of_real) and hall, and InertOK uses ramificationIdx_eq_one and
% not_dvd_differentIdeal_iff.

\begin{proposition}\label{gal-core}
Existence and the four conditions hold in the self-dual scope (Definition~\ref{gal-scope}).
\end{proposition}
\leanref{Development/AutomorphicArithmetic/Foundations/GaloisRepresentations/IntegralStructures/CrystallineDescent/ToCoreK.lean}{UpperTheory.Gal.Int.Core.toCoreK}
\leanref{Development/AutomorphicArithmetic/Foundations/GaloisRepresentations/IntegralStructures/CrystallineDescent/LimOddOfLimFinOdd.lean}{UpperTheory.Gal.Int.Core.toCore}
\leanref{Development/AutomorphicArithmetic/Foundations/GaloisRepresentations/IntegralStructures/CrystallineDescent/PatchOneAlt.lean}{UpperTheory.Gal.Int.Core.JoinU}
\leanref{Development/AutomorphicArithmetic/Foundations/GaloisRepresentations/IntegralStructures/CrystallineDescent/PatchOneAlt.lean}{UpperTheory.Gal.Int.Core.JoinAll}
\leanref{Development/AutomorphicArithmetic/Foundations/GaloisRepresentations/IntegralStructures/CrystallineDescent/PatchOneAlt.lean}{UpperTheory.Gal.Int.Core.JoinGL}
\leanref{Development/AutomorphicArithmetic/Foundations/GaloisRepresentations/IntegralStructures/CrystallineDescent/PatchOneAlt.lean}{UpperTheory.Gal.Int.Core.JoinUp}
\leanref{Development/AutomorphicArithmetic/Foundations/GaloisRepresentations/IntegralStructures/CrystallineDescent/LimOddOfLimWU3.lean}{UpperTheory.Gal.Int.Core.RowCsd}
\leanref{Development/AbelianVarietyTheory/Statements/EllipticCurveLevels/X0OfStarCondition.lean}{GlobalTheory.RoadW.RowStarCore}
\leanref{Development/AutomorphicArithmetic/Foundations/GaloisRepresentations/ThetaParseval/SplitLineSubtraction/RowDefOfOdd.lean}{UpperTheory.Gal.TP.Sub.TwoUp}
\leanref{Development/AutomorphicArithmetic/Foundations/GaloisRepresentations/ThetaParseval/TwoUpOKOfSteps.lean}{UpperTheory.Gal.TP.twoUpOK_of_k}
\leanref{Development/AutomorphicArithmetic/Foundations/GaloisRepresentations/ThetaParseval/CoreAtQIffRowQ.lean}{UpperTheory.Gal.TP.SubtractQ}
\leanref{Development/AbelianVarietyTheory/Definitions/GaloisGeometry/MatchOfMatchOffTEmpty.lean}{GlobalTheory.Cong.DescOffT}
\leanref{Development/AutomorphicArithmetic/Foundations/GaloisRepresentations/BaseChangeDescent/BcBackOffT.lean}{UpperTheory.Gal.Desc.descOffT}
\leanref{Development/AutomorphicArithmetic/Foundations/GaloisRepresentations/IntegralStructures/CrystallineDescent/PatchOneAlt.lean}{UpperTheory.Gal.Int.Core.PatchOne}

\begin{proof}
Call a CM field $E$ with maximal totally real subfield $F$ \emph{admissible} for $\ell$ if $E$ contains
$\sqrt{-d}$ for some integer $d>0$ that is a square in $\mathbf Q_\ell$, $[F:\mathbf Q]$ is even, $E/F$
is unramified at every finite place, every place over $\ell$ is split over $F$, and there are three
places over $\ell$ lying over distinct places of $F$; it is \emph{inert-admissible} if moreover every
place $v$ of $F$ inert in $E$ has $v\nmid2$, $N(v)>3$, and does not divide the different of $F$. The
proof goes through five stages, each giving a continuous semisimple $r$ with the four conditions of
Definition~\ref{gal-spherical-conditions} at the places it names.

\emph{Stage 1: definite groups of odd rank.} Let $m$ be odd, $E$ admissible, and $U$ an automorphic
representation of $U_m$ with weights $\lambda$. Then there is $r$ satisfying, at every $d$-split place
where $U$ has Satake parameters $B$, the four conditions with respect to $B$ and the infinitesimal
characters of $\lambda$. Indeed $U$ is unramified at the non-split places outside a finite set $T$.
For each pair $u,v$ of places over $\ell$ over distinct places of $F$, Proposition~\ref{gal-limits}
gives $r_{u,v}$ with (unr) at the $d$-split places not over $\ell$, and (crys), (dR) and (phi) at the
places over $\ell$ away from $u$ and $v$; for (phi) the power sums of Frobenius determine the
characteristic polynomials. Any two $r_{u,v}$ have the same Frobenius eigenvalues at the $d$-split places
of prime norm, which are all but finitely many places of prime norm, so they are isomorphic by the
Chebotarev density theorem and the Brauer--Nesbitt theorem. Given a place $w$ over $\ell$, one of the
three pairs among the three places of the definition avoids $w$, which gives the conditions at $w$.

\emph{Stage 2: definite groups of all ranks.} Let $E$ be inert-admissible, $n\geq1$, and $U$ an
automorphic representation of $U_n$ with weights $\lambda$; there is $r$ as in Stage 1. For odd $n$ this
is Stage 1. For even $n$, $E$ satisfies the hypotheses of Proposition~\ref{gal-larger-pairs}, which
gives a pair of enlargements of $(U,\lambda)$ at every place $w_0$ split over $F$
(Definition~\ref{gal-larger-enlargement}): an integer $t$ with $n+t$ odd, an algebraic Hecke character
$\kappa$ unramified at $w_0$, and two automorphic representations $U_1$, $U_2$ of $U_{n+t}$, obtained by
theta lifts, with algebraic Hecke characters $\psi_{i,1},\ldots,\psi_{i,t}$ such that
$\psi_{1,j}\neq\psi_{2,j'}$ for all $j,j'$. Stage 1, for all odd ranks, is the hypothesis of
Proposition~\ref{gal-constituent}, which removes the characters of the $\psi_{1,j}$ from the representation
attached to $U_1$, twists by the character of $\kappa$, and gives $r$.
% Lean: joinAll_of (with goodGuard_usPl for the d-split places) is allCoreQ_of with subtractQ_holds
% (Proposition gal-constituent) and twoUpOK_of_k (Proposition gal-larger-pairs).

\emph{Stage 3: conjugate self-dual cuspidal representations.} Let $E$ be inert-admissible and
$(\pi,\eta^a)$ in self-dual scope over $E$. Then there is $r$ weakly attached to $\pi$ satisfying the four
conditions for $\pi$ at the $d$-split places where $\pi$ has Satake parameters. Indeed, by descent
from $\operatorname{GL}_a$ to the definite unitary group there are a finite set $T$ of non-split places,
disjoint from any given finite set, and an automorphic $U$ of $U_a$, unramified at the non-split places
outside $T$, whose local components at the split places are those of $\pi$ and whose weights match the
infinitesimal characters of $\pi$. This descent is due to Labesse \cite{Lab11}; Alpöge proves it with the
base change of Theorem~\ref{gal-base-change-definite}, the bound of Jacquet and Shalika and
multiplicity one for cuspidal representations. Stage 2 for $U$ gives $r$.

\emph{Stage 4: all places.} Let $E$ be any CM field and $(\pi,\eta^a)$ in scope. Then there is a
continuous semisimple $r$ satisfying the four conditions for $\pi$ at every place where $\pi$ has Satake
parameters. The proof applies Lemma~\ref{gal-patch-descent} three times, along two real steps and then a
CM step (Definition~\ref{gal-patch-steps}). For a CM field $Y$ write $Y^+$ for its maximal totally real
subfield and $\eta_Y$ for the quadratic character of $Y/Y^+$. For a cuspidal automorphic representation
$\pi'$ of $\operatorname{GL}_a(\mathbf A_Y)$ and a place $w_0$ of $Y$, let $T(\pi',w_0)$ be the finite
set made of $w_0$, the places over $\ell$ and the places where $\pi'$ has no Satake parameters. Let
$(\pi',\eta_Y^a)$ be in scope and $Q$ a real step or a CM step of $Y$ for $T(\pi',w_0)$. Then $Q$
satisfies conditions (1) and (2) of Definition~\ref{gal-patch-serves} for $\pi'$ at $w_0$: $w_0$ splits
in $Q$, and $Q$ is ramified over $Y$ at a place outside $T(\pi',w_0)$, where $\pi'$ has Satake
parameters. For every base change $P$ of $\pi'$ to $Q$, the pair $(P,\eta_Q^a)$ is in scope by the
theorems of Jacquet--Shalika and Gelfand--Kazhdan.

First, let $Y$ be as in Lemma~\ref{gal-patch-steps-exist}(2), and suppose that every CM extension $M$ of
$Y$ in $\Omega$ has three places over $\ell$ lying over distinct places of $M^+$. Then the conclusion of
Stage 4 holds for every $(\pi',\eta_Y^a)$ in scope over $Y$. Indeed, let $w_0$ be a place of $Y$, $Q$ a
CM step of $Y$ for $T=T(\pi',w_0)$, with $d$ as in Definition~\ref{gal-patch-steps}(2), and $P$ a base
change of $\pi'$ to $Q$. The places of $Q$ over $\ell$ lie over $T$, so they are $d$-split: they are
split over $Q^+$, and $-d$ is a square in $\mathbf Q_\ell$. The field $Q$ is unramified over $Q^+$ at
every finite place, since the places not split over $Q^+$ are unramified over it. With the hypothesis on
$Y$ and the remaining properties of a CM step, $Q$ is inert-admissible. The pair $(P,\eta_Q^a)$ is in
self-dual scope: a place $u$ of $Q$ not split over $Q^+$ lies over a place of $Y$ outside $T$, so
$u\nmid\ell$, $u$ is unramified over $Q^+$, and $P$ has Satake parameters at $u$ by property (1) of base
change, since $\pi'$ has them at the place of $Y$ under $u$. Stage 3 for $P$ gives $r_P$ satisfying the four conditions for $P$ at the $d$-split places where
$P$ has Satake parameters. The $d$-split places include the places over $w_0$, which lie over $T$, and
all but finitely many places of prime norm: such a place, if unramified over $\mathbf Q$, has completion
$\mathbf Q_p$, so it is split over $Q^+$ and $-d$ is a square in $\mathbf Q_p$. So $Q$ serves $\pi'$ at
$w_0$. There are infinitely many CM steps for $T$ by Lemma~\ref{gal-patch-steps-exist}(2), and
Lemma~\ref{gal-patch-descent} gives the conclusion for $\pi'$.

Next, let $w_0$ be a place of $E$, $Q_1=E(\delta_1)$ a real step of $E$ for $T(\pi,w_0)$, $P_1$ a base
change of $\pi$ to $Q_1$, $w_1$ a place of $Q_1$, and $Q_2=Q_1(\delta_2)$ a real step of $Q_1$ for
$T(P_1,w_1)$. Let $t$ be a place of $E$ over $\ell$. Then $t\in T(\pi,w_0)$, so the place of $E^+$
under $t$ splits in $Q_1^+=E^+(\delta_1)$, and $t$ splits in $Q_1$. The two places of $Q_1^+$ over it
lie under the places of $Q_1$ over $t$, since the nontrivial automorphism of $Q_1/E$ restricts to that
of $Q_1^+/E^+$. These places of $Q_1$ lie over $\ell$, hence in $T(P_1,w_1)$, so the two places of
$Q_1^+$ split in $Q_2^+=Q_1^+(\delta_2)$. Hence $Q_2^+$ has four places over the place of $E^+$ under $t$.
For every CM extension $M$ of $Q_2$, $M^+$ contains $Q_2^+$ and so has at least four places over that
place of $E^+$, and places of $M$ over three of them are three places over $\ell$ lying over distinct places of $M^+$. Since
$(\pi,\eta^a)$ is in scope, so are $(P_1,\eta_{Q_1}^a)$ and $(P_2,\eta_{Q_2}^a)$ for every base change
$P_2$ of $P_1$ to $Q_2$. By the first step for $Y=Q_2$ and $Y_0=Q_1$, the conclusion of Stage 4 holds
for every such $P_2$. So $Q_2$ serves $P_1$ at $w_1$, with $W$ the set of all places of $Q_2$. There are
infinitely many real steps by Lemma~\ref{gal-patch-steps-exist}(1), so Lemma~\ref{gal-patch-descent}
gives the conclusion for $P_1$. Hence $Q_1$ serves $\pi$ at $w_0$, with $W$ the set of all places of
$Q_1$, and Lemma~\ref{gal-patch-descent} gives the conclusion for $\pi$.
% Lean: joinUp_of' (BadFinite.lean). T(pi', w0) is bad l pi' w0. The first step is the local statement
% top: Quad2Ex with infOf; serves with RamUnr (ramUnr_holds) for conditions (1) and (2); FieldOK from
% Step2, allUnr_holds and the three places; Good; CoreGLu (Stage 3) and OneUs. The next step applies
% StepDown twice with infinitely_many_real_quadratic_steps and threeApart_holds. Scope of base changes
% is CsdLift (csdLift_holds); existence of base changes is QuadLift (quadLift), used inside StepDown.

\emph{Stage 5: conclusion.} For $(\pi,\eta^a)$ in self-dual scope, Stage 4 gives $r$ satisfying the four
conditions at every place where $\pi$ has Satake parameters; it is weakly attached to $\pi$ since
(unr) holds at almost all places and $\pi$ is spherical at almost all places. Any weakly attached
representation is isomorphic to $r$ by the Chebotarev density theorem and the Brauer--Nesbitt theorem,
so the conditions hold for all of them.
\end{proof}
\leanref{Development/AutomorphicArithmetic/Foundations/GaloisRepresentations/BaseChangeDescent/BcBackOffT.lean}{UpperTheory.Gal.Desc.desc}
